\ifdefined\GinibreFullProof
  \documentclass[11pt,reqno]{amsart}
\else
  \documentclass[10pt,reqno]{amsart}
\fi

\newif\ifginibrefullproof
\ifdefined\GinibreFullProof
  \ginibrefullprooftrue
\else
  \ginibrefullprooffalse
\fi

\usepackage[utf8]{inputenc}
\usepackage[T1]{fontenc}
\usepackage{amsmath,amssymb,amsthm,mathtools}
\usepackage{graphicx}
\usepackage[protrusion=true,expansion=false]{microtype}
\ifginibrefullproof
  \usepackage[margin=1in]{geometry}
  \setlength{\emergencystretch}{3em}
\else
  \usepackage[margin=0.85in]{geometry}
\fi
\usepackage{xcolor}
\usepackage{hyperref}
\usepackage{cleveref}
\usepackage{array,booktabs}
\usepackage{longtable}

\hypersetup{
  colorlinks=true,
  linkcolor=blue!50!black,
  citecolor=blue!50!black,
  urlcolor=blue!70!black,
  pdfauthor={Leslie P. Polzer},
  pdfsubject={Two-minus adjoint-character Ginibre Q3 inequality},
  pdfkeywords={Ginibre Q3, compact Lie group, adjoint character, exact certificates},
}

\theoremstyle{plain}
\newtheorem{theorem}{Theorem}[section]
\newtheorem{proposition}[theorem]{Proposition}
\newtheorem{lemma}[theorem]{Lemma}
\newtheorem{corollary}[theorem]{Corollary}

\theoremstyle{definition}
\newtheorem{definition}[theorem]{Definition}
\newtheorem{remark}[theorem]{Remark}

\DeclareMathOperator{\Haar}{Haar}
\DeclareMathOperator{\mult}{mult}
\DeclareMathOperator{\rank}{rank}
\newcommand{\Qthree}{Q_3}
\newcommand{\chiad}{\chi_{\mathrm{ad}}}
\newcommand{\triv}{\mathbf{1}}
\newcommand{\ad}{\mathrm{ad}}
\newcommand{\SU}{\mathrm{SU}}
\newcommand{\SO}{\mathrm{SO}}
\newcommand{\Sp}{\mathrm{Sp}}
\newcommand{\Spin}{\mathrm{Spin}}
\newcommand{\Gtwo}{G_2}
\newcommand{\Ffour}{F_4}
\newcommand{\Esix}{E_6}
\newcommand{\Eseven}{E_7}
\newcommand{\Eeight}{E_8}
\newcommand{\Z}{\mathbb{Z}}
\newcommand{\N}{\mathbb{N}}
\newcommand{\DEnv}{\mathsf{DEnv}}
\newcolumntype{L}[1]{>{\raggedright\arraybackslash}p{#1}}
\newcommand{\hashsplit}[2]{%
  \begin{tabular}{@{}l@{}}%
    \texttt{\scriptsize #1}\\[-1pt]
    \texttt{\scriptsize #2}%
  \end{tabular}}
\ifginibrefullproof
  \newcommand{\contractref}[1]{\cref{#1}}
\else
  % The compact build prints the exact supplement label without creating an
  % unresolved cross-reference to a result that is typeset only in the
  % detailed build.
  \newcommand{\contractref}[1]{\path{#1}}
\fi

% Part II contains exact integer ledgers whose widest rows are intentionally
% much longer than prose width.  Ordinary displays retain their natural size;
% only a display whose measured box exceeds \linewidth is reduced to fit.
\newsavebox{\ginibremathbox}
\newcommand{\ginibrefitmath}[1]{%
  \sbox{\ginibremathbox}{$\displaystyle #1$}%
  \ifdim\wd\ginibremathbox>\linewidth
    \resizebox{\linewidth}{!}{\usebox{\ginibremathbox}}%
  \else
    \usebox{\ginibremathbox}%
  \fi}
\ifginibrefullproof
  \let\ginibreopendisplay\[
  \let\ginibreclosedisplay\]
  \long\def\[#1\]{%
    \ginibreopendisplay\ginibrefitmath{#1}\ginibreclosedisplay}
\fi

\ifginibrefullproof
\title[Adjoint two-minus Ginibre integral]{The two-minus adjoint-character case of Ginibre's $Q_3$ condition\\\large Detailed computational supplement to Parts I--II}
\hypersetup{pdftitle={The two-minus adjoint-character case of Ginibre's Q3 condition, detailed computational supplement}}
\else
\title[Adjoint two-minus Ginibre integral]{The two-minus adjoint-character case of Ginibre's $Q_3$ condition\\\large Parts I--II: Classification and exact certificates}
\hypersetup{pdftitle={The two-minus adjoint-character case of Ginibre's Q3 condition, Parts I-II}}
\fi

\author{Leslie P. Polzer}
\address{}
\email{polzer@fastmail.com}

\date{July 19, 2026}

\subjclass[2020]{Primary 22E46, 43A75; Secondary 40E05, 44A60, 60B15}
\keywords{compact Lie group, Haar measure, adjoint character, Ginibre inequality, character rings, exact certificates}

\begin{document}

\begin{abstract}
Let $G$ be a compact connected Lie group with simple Lie algebra, let
$\chiad$ be its adjoint character, and define
\[
  \Qthree^G(n)=
  \int_{G\times G}(\chiad(g)-\chiad(h))^2
  (\chiad(g)+\chiad(h))^n\,d\Haar(g)\,d\Haar(h).
\]
We prove $\Qthree^G(n)\ge0$ for every such $G$ and every integer $n\ge0$.
This is the specialization of Ginibre's condition~(Q3) in which every
function is $\chiad$, exactly two factors carry a minus sign, and the other
$n$ factors carry plus signs; it is not a proof of the full cone-valued Q3
condition.  The moment formula turns the odd-exponent problem into inequalities
for invariant multiplicities in tensor powers of the adjoint representation.
The proof follows the classification: type~$A$ uses transition ratios;
types~$B,C,D$ use stable moments, finite correction prefixes, and uniform
trace tails; and the exceptional types use exact character-ring prefixes and
rectangular analytic tails.  All finite decisions are made with exact
integer or rational arithmetic or with outward-rounded intervals.  The
computer-assisted steps are stated as finite mathematical contracts, with
their active domains and proof dependencies separated from certificate
generation and transcript authentication.
\end{abstract}

\maketitle

\ifginibrefullproof
\tableofcontents
\fi

\begin{remark}[Companion architecture]
\ifginibrefullproof
This detailed computational supplement expands the $B/C$ Pieri-correction
development and prints the detailed certificate, transcript, and row
ledgers.  It is a load-bearing component of the proof of the residual
$B/C$ closure stated compactly in \texttt{paper.pdf}.  The journal-facing
Parts~I--II remain consolidated in \texttt{paper.pdf}; the present supplement
supplies the numbered proofs of the correction-prefix contracts consumed
there and the half-stable bridge reduction that converts those contracts into
Chain inequalities.  The compact contract table identifies every active result by label
and proves that offsets beyond $27$ are not consumed.  Conditional and
superseded development routes are excluded from the formal supplement; they
can be typeset separately from the source archive by defining
\texttt{\string\GinibreDevelopmentArchive}.  They are not part of that
indexed proof.  Part~III is the separate full signed-product extension
\texttt{full\_q3\_extension.pdf}.
\else
This journal manuscript consolidates Parts~I--II of the proof: the
classification reduction, family closures, compact correction-prefix
contract, and accepted-certificate boundary.  The detailed computational supplement
\texttt{paper\_full.pdf} expands the correction-prefix derivations and
certificate ledgers from the same source and proves the exact active results
listed in the compact contract table.  It also proves the explicitly indexed
half-stable bridge reduction that converts those correction intervals into
Chain inequalities.  Only that indexed subgraph is load-bearing for the
residual $B/C$ closure.
Part~III, \texttt{full\_q3\_extension.pdf}, imports the theorem proved here
and establishes the full signed-product hierarchy on the adjoint-generated
cone.
\fi
\end{remark}

\begin{remark}[Publication proof spine and author self-audits]
The load-bearing dependency spine for Parts~I--III and their detailed
supplement is listed in
\path{PUBLICATION_PROOF_SPINE.md} and, for the present theorem, in the
final-facing dependency map below.  The supplement's numbered
correction-prefix proofs and half-stable bridge reduction are load-bearing;
explicitly marked diagnostic and superseded routes remain provenance and are
not consumed.  The archived
historical report and the programs under
\path{referee_audits/} are author-controlled self-audits.  They provide
redundant checks and source mappings; they are not independent peer review and
are not premises of any theorem.
\end{remark}

%======================================================================
\section{Introduction}\label{sec:introduction}
%======================================================================

Ginibre's condition~(Q3) asks that, for a probability measure $\sigma$ and
every finite family $f_1,\ldots,f_r$ in a specified multiplicative convex
cone, each signed product integral
\[
 \int\!\!\int \prod_{i=1}^r\bigl(f_i(x)\pm f_i(y)\bigr)
 \,d\sigma(x)\,d\sigma(y)
\]
be nonnegative; see
\cite[\S1, condition~(Q3), equation~(1.2)]{Ginibre1970}.  The present theorem
addresses one precise noncommutative specialization.  We take normalized Haar
measure on a compact connected group with simple Lie algebra, set every
$f_i$ equal to the adjoint character, choose exactly two minus signs, and
choose $n$ plus signs.  The resulting integral is $\Qthree^G(n)$.  We do not
assert Q3 for a cone of functions, for unequal characters, or for four or
more minus factors.

Even $n$ is immediate from pointwise nonnegativity.  Odd $n$ is not: the
adjoint character takes negative values, and expansion produces differences
of large invariant-tensor multiplicities.  The moment formula in
\cref{lem:moment-formula} converts the problem into exact inequalities for
$m_k^G=\mult(\triv,\ad^{\otimes k})$.  The proof then separates into the
classification families.  Type~$A$ is governed by stable derangement moments
and finite-rank transition ratios.  The classical types use the exact stable
trace law of Diaconis--Shahshahani, orthogonal and symplectic identities of
Rains, finite correction bridges, and uniform trace tails.  The exceptional
types combine finite character-ring intervals with analytic rectangular
tails.

Computer assistance enters only after the mathematical reduction has
specified a finite integer, rational, or interval inequality.  A source table
is not accepted merely because its hash matches.  For every moment ledger
passed to an active source-replay stage, the clean replay regenerates the
claimed exceptional or low classical moments from the Cartan datum by
\cref{lem:racah-speiser-step,lem:character-pairing}; in
particular, the $E_8$ values through $m_{100}$ require the decomposition only
through $\ad^{\otimes50}$.  The ancillary $E_8$ values of
Bourjaily--Plesser--Vergu
\cite[ancillary file \texttt{adjoint\_tensor\_ranks.tex}]{BPV2024} are
therefore comparison data, not an unverified premise.  Likewise, every stable
classical moment used in a finite correction is regenerated from the exact
\ifginibrefullproof
recurrence of \cref{lem:bc-stable-moment-recurrence};
\else
three-term recurrence reconstructed by the exact source replay;
\fi
the archived OEIS row is checked against that regeneration and is not used as
a mathematical premise.
Exact discrete calculations use GMP integers or rationals, and analytic
comparisons use outward-rounded MPFR intervals.

Parts~I--II are consolidated in this manuscript.  The final-facing reductions
and family closures consume a compact finite-correction contract and the
accepted certificate ledger below.  The detailed computational supplement prints
the line-by-line correction derivations and historical ledgers; its numbered
correction-prefix proofs and half-stable bridge reduction are incoming nodes
for the compact residual closure, while explicitly diagnostic and superseded
material adds no hypothesis or incoming edge.  The replay
contract states which inputs are regenerated, which legacy certificates are
audited without a full generator rerun, and which archived files are retained
only as diagnostic history.

The choice of the adjoint character is not merely a source of numerical
examples.  Its powers are the characters of tensor powers of the adjoint
module, so every Haar moment is an invariant-tensor multiplicity and can be
attacked uniformly through the classification.  Part~III takes the smallest
uniformly closed multiplicative convex cone generated by this character.
Ginibre's Proposition~1 promotes the single-generator signed-product
inequalities to exactly this generated cone; any claim on a larger cone
requires additional input not supplied by that promotion.  The resulting
theorem is thus a
complete statement about the adjoint tensor category, while remaining
strictly narrower than a theorem for all central positive-definite functions.

\subsection*{Contribution and relation to prior results}
Ginibre's Proposition~1 supplies a promotion mechanism once the required
signed products are known on a self-conjugate generating set; it does not
establish the repeated-adjoint inequality proved here.  The exact stable
trace laws of Diaconis--Shahshahani, the power-map identities of Rains, and
the character bounds of Garibaldi--Guralnick--Rains provide inputs to
individual family arguments, but none of those results compares the two
moment diagonals in \cref{lem:moment-formula} for every exponent and every
simple type.  Likewise, the adjoint tensor-rank calculations of
Bourjaily--Plesser--Vergu provide comparison data for exceptional moments,
not the sign theorem.

The contribution of Parts~I--II is the uniform classification theorem and
the mechanism joining three regimes that do not overlap automatically:
stable invariant-tensor moments, finite-rank corrections, and analytic
tails.  The repeated-adjoint and exactly-two-minus restrictions are essential
parts of the statement, not abbreviations for a result about arbitrary
characters.  In particular, this paper makes no claim for unequal
irreducible characters or for all positive-definite class functions.  The
separate Part~III uses Ginibre's promotion only after importing the theorem
proved here; it is not used to justify the present theorem or its novelty.

The computer calculations are proof devices rather than the mathematical
claim.  Each active calculation has a numbered finite contract, and the
compact residual $B/C$ contract table below records the exact theorem labels,
rank ranges, and moment offsets supplied by the detailed proof.

%======================================================================
\section{Statement and General Reductions}
%======================================================================

\begin{definition}[Adjoint $Q_3$ functional]
For a compact connected Lie group $G$ with simple Lie algebra and normalized
Haar probability measure, set
\[
  \Qthree^G(n)=
  \int_{G\times G}(\chiad(g)-\chiad(h))^2
  (\chiad(g)+\chiad(h))^n\,d\Haar(g)\,d\Haar(h)
\]
for $n\in\Z_{\ge0}$, where $\chiad$ is the adjoint character.
\end{definition}

\begin{theorem}[Two-minus adjoint-character Q3 nonnegativity]\label{thm:main}
For every compact connected Lie group $G$ with simple Lie algebra and every
integer $n\ge0$,
\[
  \Qthree^G(n)\ge0.
\]
\end{theorem}

\begin{lemma}[Moment formula]\label{lem:moment-formula}
Let
\[
  m_k^G=\int_G \chiad(g)^k\,d\Haar(g)
       =\mult(\triv,\ad^{\otimes k}).
\]
Then
\[
  \Qthree^G(n)
  =
  2\sum_{k=0}^n \binom{n}{k}
  \bigl(m_{k+2}^Gm_{n-k}^G-m_{k+1}^Gm_{n-k+1}^G\bigr).
\]
\end{lemma}

\begin{proof}
Expand $(\chiad(g)+\chiad(h))^n$ by the binomial theorem, multiply by
$(\chiad(g)-\chiad(h))^2$, and integrate term by term over $G\times G$.
Schur orthogonality identifies $\int_G\chiad(g)^k\,d\Haar(g)$ with the
invariant multiplicity $\mult(\triv,\ad^{\otimes k})$.
\end{proof}

\begin{lemma}[Racah--Speiser adjoint tensor step]
\label{lem:racah-speiser-step}
Let $V_\lambda$ be an irreducible representation of a compact connected
semisimple group, let $\rho$ be the half-sum of the positive roots, and let
$a(\mu)$ be the multiplicity of the weight $\mu$ in the adjoint
representation.  For every adjoint weight $\mu$, proceed as follows.
If $\lambda+\mu+\rho$ lies on a Weyl wall, assign zero.  Otherwise let $w$ be
the unique Weyl element for which $w(\lambda+\mu+\rho)$ is strictly dominant,
put
\[
 \nu=w(\lambda+\mu+\rho)-\rho,
\]
and assign $\operatorname{sgn}(w)a(\mu)$ to $V_\nu$.  After equal dominant
weights are collected, the resulting signed sum is the irreducible
decomposition of $V_\lambda\otimes\ad$.
\end{lemma}

\begin{proof}
For a weight $\eta$ write
$A_\eta=\sum_{w\in W}\operatorname{sgn}(w)e^{w\eta}$.
The Weyl character formula gives
$\chi_\lambda=A_{\lambda+\rho}/A_\rho$.  Since the adjoint weight multiset is
Weyl invariant,
\[
 A_{\lambda+\rho}\chiad
   =\sum_\mu a(\mu)A_{\lambda+\mu+\rho}.
\]
An alternant vanishes when its index is singular.  If the index is regular,
moving it to the strict dominant chamber changes the alternant by
$\operatorname{sgn}(w)$, and division by $A_\rho$ gives the character
specified in the statement.  Collecting equal dominant indices proves the
formula.  This is the Racah--Speiser procedure; the same shifted-wall
convention is described in
\cite[Appendix~B, especially (B.1)--(B.6)]{BPV2024}.
\end{proof}

\begin{lemma}[Character-pairing reconstruction of moments]
\label{lem:character-pairing}
Write the exact irreducible decomposition
\[
 \chiad^j=\sum_{\lambda\in P_+}c_j(\lambda)\chi_\lambda .
\]
Then, for every $j\ge0$,
\[
 m_{2j}^G=\sum_{\lambda\in P_+}c_j(\lambda)^2,
 \qquad
 m_{2j+1}^G=
 \sum_{\lambda\in P_+}c_j(\lambda)c_{j+1}(\lambda).
\]
Consequently all moments through degree $M$ can be reconstructed after
iterating \cref{lem:racah-speiser-step} only through tensor power
$\lceil M/2\rceil$.
\end{lemma}

\begin{proof}
The adjoint character is real, so
\[
 m_{a+b}^G
 =\int_G\chiad(g)^a\overline{\chiad(g)^b}\,d\Haar(g)
 =\left\langle\chiad^a,\chiad^b\right\rangle_{L^2(G)}.
\]
Irreducible-character orthonormality makes the last expression
$\sum_\lambda c_a(\lambda)c_b(\lambda)$.  Taking $(a,b)=(j,j)$ and
$(j,j+1)$ proves the two identities.
\end{proof}

\begin{lemma}[Positivity of adjoint moments]\label{lem:adjoint-moment-positive}
The adjoint character is real-valued, and
\[
  m_0^G=1,\qquad m_1^G=0,\qquad m_2^G=1,
\]
while
\[
  m_k^G=\dim\bigl((\mathfrak g_{\mathbf C}^{\otimes k})^G\bigr)
  \quad\text{is a positive integer for every }k\ge2.
\]
In particular, every transition ratio used below has a nonzero positive
denominator.
\end{lemma}

\begin{proof}
The adjoint action preserves a positive-definite real inner product $B$ on
the compact real form $\mathfrak g$, so its complexification is self-dual
and its character is real.  The compact-semisimple/complex-semisimple
correspondence and decomposition into simple root systems in
\cite[Chapter~V, \S5, pp.~209--215, especially (5.5) and
Theorem~(5.8)]{BrockertomDieck} show that
$\mathfrak g_{\mathbf C}$ is a complex simple Lie algebra.  If a complex
subspace of $\mathfrak g_{\mathbf C}$ is invariant under $G$, connectedness
of $G$ and differentiation make it invariant under
$\operatorname{ad}(\mathfrak g_{\mathbf C})$; it is therefore an ideal.
Thus the complexified adjoint representation is irreducible.

Character orthogonality now gives the displayed invariant-space formula and,
because the character is real,
\[
  m_2^G=\int_G|\chiad(g)|^2\,d\Haar(g)=1.
\]
Degree zero consists of the scalars.  A degree-one
invariant is a vector fixed by the adjoint action, hence lies in the center of
$\mathfrak g_{\mathbf C}$; simplicity gives $m_1^G=0$.  The tensor $B$ is a
nonzero invariant tensor of
degree two.  Simplicity and non-abelianness imply that the Cartan tensor
\[
  \omega(X,Y,Z)=B(X,[Y,Z])
\]
is a nonzero invariant tensor of degree three.  Every integer $k\ge2$ has
the form $2a+3b$ with $a,b\ge0$; tensoring $a$ copies of $B$ and $b$ copies
of $\omega$, and using $B$ to identify $\mathfrak g$ with its dual, gives a
nonzero invariant tensor in degree $k$.  Hence $m_k^G\ge1$.
\end{proof}

\begin{lemma}[Central quotient invariance]\label{lem:central-quotient}
If $\pi:\widetilde G\to G$ is a finite central cover of compact connected
Lie groups with simple Lie algebra, then
\[
  \Qthree^{\widetilde G}(n)=\Qthree^G(n)
\]
for every $n\ge0$.
\end{lemma}

\begin{proof}
The adjoint characters satisfy
$\chiad^{\widetilde G}=\chiad^G\circ\pi$.  The push-forward of normalized
Haar measure under $\pi$ is normalized Haar measure.  Apply this in the two
variables of the defining integral.
\end{proof}

\begin{lemma}[Reduction to the simple Lie algebra]\label{lem:global-form}
Let $G$ be a compact connected Lie group with simple Lie algebra
$\mathfrak g$.  There is a compact simply connected Lie group $G_{\rm sc}$
with Lie algebra $\mathfrak g$ and a finite central covering
$G_{\rm sc}\to G$.  Consequently, if $G_1$ and $G_2$ are compact connected
Lie groups with isomorphic simple Lie algebras, then
\[
  \Qthree^{G_1}(n)=\Qthree^{G_2}(n)
\]
for every $n\ge0$.
\end{lemma}

\begin{proof}
A simple Lie algebra is semisimple.  For a compact connected semisimple
group, the universal covering group is compact and its kernel is finite and
central; moreover, the simply connected compact group is determined by its
Lie algebra.  These statements are recorded in
\cite[Chapter~V, Remark~(7.13) and \S8, p.~232]{BrockertomDieck}.
Thus $G_{\rm sc}\to G$ has the asserted properties.  Isomorphic simple Lie
algebras give isomorphic simply connected compact covering groups, so two
applications of \cref{lem:central-quotient} prove the displayed equality.
\end{proof}

\begin{lemma}[Even exponents]\label{lem:even}
For every $G$ and every even $n\ge0$, $\Qthree^G(n)\ge0$.
\end{lemma}

\begin{proof}
For even $n$, both $(\chiad(g)-\chiad(h))^2$ and
$(\chiad(g)+\chiad(h))^n$ are pointwise nonnegative.
\end{proof}

\begin{definition}[Chain difference]\label{def:chain-diff}
For odd exponents define
\[
  D_G(n)=\Qthree^G(n+2)-4\Qthree^G(n).
\]
Equivalently, with $n=2m+1$,
\[
  D_G(2m+1)=\Qthree^G(2m+3)-4\Qthree^G(2m+1).
\]
\end{definition}

\begin{lemma}[Chain propagation]\label{lem:chain-propagation}
Let $0\le m_0\le m_1$.  Assume $D_G(2m+1)\ge0$ for
$m=m_0,m_0+1,\ldots,m_1$ and $\Qthree^G(2m_0+1)\ge0$.  Then
\[
  \Qthree^G(2m+1)\ge0
\]
for every $m_0\le m\le m_1+1$.
\end{lemma}

\begin{proof}
The recurrence
\[
  \Qthree^G(2m+3)=4\Qthree^G(2m+1)+D_G(2m+1)
\]
propagates nonnegativity one odd exponent at a time.
\end{proof}

%======================================================================
\section{Type A: The Adjoint Representation of \texorpdfstring{$\SU(N)$}{SU(N)}}
%======================================================================

This section proves the type $A$ branch.  Since finite central quotients do
not change $\Qthree$ by \cref{lem:global-form}, it suffices to use
$\SU(N)$ as the representative of type $A_{N-1}$.

\begin{theorem}[Stable-rank type $A$ branch]\label{thm:su-stable-rank}
For every $N\ge2$ and every $n\ge0$ satisfying $N\ge n+2$,
\[
  \Qthree^{\SU(N),\ad}(n)\ge0.
\]
\end{theorem}

\begin{proof}
Write $X=|\operatorname{tr}U|^2$ for $U$ Haar-distributed in $\SU(N)$.
For $0\le k\le N$, scalar invariance identifies the $\SU(N)$ integral of
$X^k$ with the $\mathrm{U}(N)$ integral.  Schur expansion and character
orthogonality give
\[
  \mathbf E X^k
  =
  \sum_{\lambda\vdash k,\ \ell(\lambda)\le N}
  (\dim \operatorname{Sp}_{\lambda})^2.
\]
When $k\le N$ the length condition is void, and the regular
representation of the symmetric group gives $\mathbf E X^k=k!$.
Since $\chiad(U)=X-1$, this gives
\[
  m_r^{\SU(N)}=\sum_{j=0}^r(-1)^{r-j}\binom rj j!=!r
  \qquad(0\le r\le N),
\]
where $!r$ is the derangement number.  If $N\ge n+2$, all moments
appearing in \cref{lem:moment-formula} are in this stable range.

Let $Y$ have the exponential distribution of mean $1$, and set
$A=Y-1$.  The exponential generating function
$M(z)=\mathbf E e^{zA}=e^{-z}(1-z)^{-1}$ satisfies
\[
  M''(z)M(z)-M'(z)^2
  =M(z)^2(\log M)''(z)
  =\frac{e^{-2z}}{(1-z)^4}.
\]
Thus $\Qthree^{\SU(N)}(n)/2$ is the $n$th exponential coefficient of
$e^{-2z}(1-z)^{-4}$.  If these coefficients are denoted $q_n$, then
$q_0=1$, $q_1=2$, and differentiating gives
\[
  q_{n+1}=(n+2)q_n+2nq_{n-1}\qquad(n\ge1).
\]
Hence $q_n>0$ for all $n\ge0$.
\end{proof}

\begin{definition}[Type $A$ transition ratios]\label{def:su-ratio}
For $N\ge2$, put
\[
  m_k^{(N)}=\int_{\SU(N)}\chiad(U)^k\,dU,\qquad
  \rho_k^{(N)}=\frac{m_{k+1}^{(N)}}{m_k^{(N)}}\quad(k\ge3).
\]
The transition-ratio inequality at $(N,k)$ is
\[
  R_{N,k}:=m_{k+2}^{(N)}m_k^{(N)}-(m_{k+1}^{(N)})^2\ge0,
\]
equivalently $\rho_k^{(N)}\le \rho_{k+1}^{(N)}$.
\end{definition}

\begin{lemma}[Boundary reduction for type $A$]\label{lem:su-boundary}
Fix $N\ge4$.  If
\[
  \rho_3^{(N)}=\frac92,\qquad
  \frac92\le \rho_k^{(N)}\le \rho_{k+1}^{(N)}
  \quad(k\ge3),
\]
then $\Qthree^{\SU(N)}(n)>0$ for every odd $n\ge5$.
\end{lemma}

\begin{proof}
Use the symmetric form
\[
  \frac12\Qthree(n)=
  \frac12\sum_{j=0}^n\binom nj
  \bigl(m_{j+2}m_{n-j}+m_jm_{n+2-j}
       -2m_{j+1}m_{n+1-j}\bigr).
\]
For $3\le j\le n-3$, division by
$m_{j+1}m_{n+1-j}$ gives
\[
  \frac{\rho_{j+1}}{\rho_{n-j}}
  +\frac{\rho_{n+1-j}}{\rho_j}-2\ge0
\]
by monotonicity and AM-GM.  The remaining depth-two boundary
contribution
\[
  L_2=B_0+nB_1+\binom n2 B_2
\]
uses only $m_0=1,m_1=0,m_2=1,m_3=2,m_4=9$ and the four ratios
$R=\rho_{n-2}$, $S=\rho_{n-1}$, $T=\rho_n$, $U=\rho_{n+1}$.  Substitution
and $U\ge T$ give
\[
  \frac{L_2}{m_{n-2}}\ge
  RS\left(T^2+\frac{n^2-5n+2}{2}\right)-2n(n-2)R
  +\frac92 n(n-1).
\]
Since $R,S,T\ge9/2$, the right hand side is at least
\[
  \frac{9(10n^2-66n+765)}{16}>0.
\]
Thus the boundary contribution is positive and all interior contributions
are nonnegative.
\end{proof}

\begin{lemma}[Exact formal derivation of the low-rank type \(A\) recurrences]
\label{lem:su-low-rank-recurrence-derivation}
Put
\[
 E_k^{(N)}=\mathbf E_{\mathrm U(N)}|\operatorname{tr}U|^{2k},
 \qquad
 m_k^{(N)}=\sum_{j=0}^k(-1)^{k-j}\binom kjE_j^{(N)}.
\]
Then the following recurrences hold:
\[
\begin{array}{c|l}
N&\text{recurrence}\\
\hline
3&(k+3)^2m_{k+1}=k(7k+11)m_k+8k(k+1)m_{k-1},
   \quad k\ge1,\\[2pt]
4&\begin{aligned}
0={}&(k^3+22k^2+161k+392)m_{k+4}\\
&-(16k^3+230k^2+1054k+1524)m_{k+3}\\
&+(10k^3-340k-750)m_{k+2}\\
&+(72k^3+522k^2+1242k+972)m_{k+1}\\
&+(45k^3+270k^2+495k+270)m_k,
\qquad k\ge0,
\end{aligned}\\[12pt]
5&\begin{aligned}
0={}&192(k+1)(k+2)(k+3)(k+6)m_k\\
&+16(k+2)(k+3)(22k^2+184k+309)m_{k+1}\\
&+(k+3)(129k^3+1245k^2+2570k-1112)m_{k+2}\\
&-(k+3)(30k^3+642k^2+4487k+10196)m_{k+3}\\
&+(k+8)^2(k+10)^2m_{k+4},
\qquad k\ge0.
\end{aligned}
\end{array}
\]
\end{lemma}

\begin{proof}
Gessel's determinant
\cite[\S7, Theorem~16 and the formula following (26)]{Gessel1990} gives
\[
 F_N(t):=\sum_{k\ge0}E_k^{(N)}\frac{t^k}{(k!)^2}
 =\det\bigl(I_{|i-j|}(2\sqrt t)\bigr)_{1\le i,j\le N}.
\]
Put \(y=I_0(2\sqrt t)\) and \(z=y'\).  The Bessel recurrence and
\(ty''+y'-y=0\) give
\[
 I_1(2\sqrt t)=\sqrt t\,z,\qquad
 I_{r+1}(2\sqrt t)=I_{r-1}(2\sqrt t)
 -\frac r{\sqrt t}I_r(2\sqrt t),
\qquad
 z'=\frac{y-z}{t}.
\]
Expanding the three determinants, with no series truncation, gives
\[
\begin{aligned}
F_3={}&z(2y^2-yz-2tz^2),\\
F_4={}&-\frac1t\bigl(
4t^2z^4-8ty^2z^2+8tyz^3-tz^4
+4y^4-8y^3z+4y^2z^2\bigr),\\
F_5={}&-\frac4{t^2}
\bigl(2y^2-yz-(2t+1)z^2\bigr)\\
&\hspace{28mm}\cdot
\bigl(4y^3-5y^2z+yz^2-4tyz^2+3tz^3\bigr).
\end{aligned}
\tag{\ref{lem:su-low-rank-recurrence-derivation}.1}
\]

Here is a finite differential certificate for the coefficient recurrences.
On \(\mathbb Q[t,t^{-1},y,z]\), let \(\partial\) be the derivation
\[
 \partial t=1,\qquad \partial y=z,\qquad
 \partial z=(y-z)/t,
\]
and put
\[
 \vartheta=t\partial,\qquad
 \mathcal S=(\vartheta+1)\partial.
\]
If \(F=\sum_{k\ge0}a_kt^k/(k!)^2\), then
\[
 \mathcal SF=\sum_{k\ge0}a_{k+1}\frac{t^k}{(k!)^2},
 \qquad
 q(\vartheta)F=\sum_{k\ge0}q(k)a_k\frac{t^k}{(k!)^2}.
\tag{\ref{lem:su-low-rank-recurrence-derivation}.2}
\]
For the following explicitly displayed polynomials, direct calculation in
this Laurent-polynomial differential ring gives
\[
 \sum_jq_{N,j}(\vartheta)\mathcal S^jF_N=0
 \qquad(N=3,4,5).
\tag{\ref{lem:su-low-rank-recurrence-derivation}.3}
\]
They are
\[
\begin{array}{c|l}
N& (q_{N,0}(x),q_{N,1}(x),\ldots)\\
\hline
3&
\begin{aligned}
(&9(x+1)^2,\,
 -(19x^2+78x+77),\\
 &11x^2+70x+109,\,-(x+5)^2);
\end{aligned}\\[4pt]
4&
\begin{aligned}
(&-64(x+1)^2(x+2),\\
 &2(106x^3+763x^2+1791x+1366),\\
 &-(253x^3+2662x^2+9223x+10482),\\
 &127x^3+1795x^2+8378x+12884,\\
 &-(23x^3+428x^2+2625x+5312),\\
 &(x+8)^2(x+9));
\end{aligned}\\[4pt]
5&
\begin{aligned}
(&225(x+1)^2(x+2)^2,\\
 &-(x+2)^2(934x^2+6226x+10317),\\
 &1487x^4+22936x^3+130422x^2+324256x+297420,\\
 &-2(554x^4+11290x^3+85109x^2+281496x+344526),\\
 &367x^4+9494x^3+91097x^2+384340x+601200,\\
 &-(38x^4+1282x^3+15993x^2+87412x+176688),\\
 &(x+10)^2(x+12)^2).
\end{aligned}
\end{array}
\tag{\ref{lem:su-low-rank-recurrence-derivation}.4}
\]
Thus \eqref{lem:su-low-rank-recurrence-derivation}.2--.3 give the
\(E_k^{(N)}\)-recurrences with coefficients
\(q_{N,j}(k)\).

For completeness, we verify that these are exactly the binomial transforms
of the three recurrences in the statement.  Let
\[
 A_N(w)=\sum_{k\ge0}E_k^{(N)}\frac{w^k}{k!},
 \qquad
 M_N(w)=\sum_{k\ge0}m_k^{(N)}\frac{w^k}{k!}.
\]
The binomial-transform identity is \(M_N(w)=e^{-w}A_N(w)\).
If a proposed recurrence is written
\(\sum_rp_{N,r}(k)m_{k+r}=0\), its exponential-generating-function
operator is
\[
 \sum_rp_{N,r}(wD)D^r,\qquad D=\frac d{dw}.
\]
Conjugating by \(e^{-w}\) replaces \(D\) by \(D-1\) and \(wD\) by
\(wD-w\).  Normal ordering with the single relation
\[
 Dw=wD+1
\]
and then taking coefficients gives the polynomials in
\eqref{lem:su-low-rank-recurrence-derivation}.4, after index shifts
\(1,1,2\) for \(N=3,4,5\), respectively.  The finitely many coefficients
below those shifts vanish by the hook-length values for
\(E_k^{(N)}\).

Every algebraic operation just described is replayed without floating point
or symbolic-library assumptions by
\[
\texttt{character\_ring\_iter/verify\_sun\_recurrence\_derivation\_exact.py}.
\]
It represents a monomial as \(t^{a/2}y^bz^c\), expands the determinants over
all permutations, normal-orders the Weyl algebra by the displayed
commutation relation, applies \(\partial z=(y-z)/t\), and requires the final
Laurent polynomial to have zero terms.  It also regenerates the initial
\(E_k^{(N)}\) from the hook-length formula and checks the omitted
coefficients.  The status-zero transcript reports zero reduced-identity
terms and zero finite residuals for all three ranks.  The source and
transcript SHA-256 values are, respectively,
\[
\hashsplit{e08536805bd04122b6cf87ac1ba083bb}{d5076cb2c1620baad6002ae3f42fb100},
\qquad
\hashsplit{71423c65e6c02d1862e14d3008120816}{e324d15a9642dd7af5143bf530c4a7ec}.
\]
This proves the three formal recurrences for every stated index.
\end{proof}

\begin{theorem}[Low-rank type $A$ ratio theorem]\label{thm:su-low-rank-ratios}
For $N=3$ one has
\[
  \rho_3^{(3)}=\rho_4^{(3)}=4,
  \qquad 4\le\rho_k^{(3)}\le\rho_{k+1}^{(3)}\quad(k\ge3).
\]
For $N=4,5$ one has
\[
  \rho_3^{(N)}=\frac92,
  \qquad \frac92\le\rho_k^{(N)}\le\rho_{k+1}^{(N)}\quad(k\ge3).
\]
\end{theorem}

\begin{proof}
Use the exact formal recurrences of
\cref{lem:su-low-rank-recurrence-derivation}.

For $N=3$, set $U_k=8-\rho_k$ and
\[
 G_k=\frac{32k+100}{(k+4)^2},\qquad
 H_k=\frac{32(k+11/5)}{(k+3)^2}.
\]
The first recurrence gives $U_{k+1}=\Phi_{k+1}(U_k)$, where
\[
 \Phi_k(u)=\frac{k^2+37k+72-8k(k+1)/(8-u)}{(k+3)^2}.
\]
This map is decreasing for $u<8$.  The exact base values are
$U_5=111/32$, $G_5=260/81$, and $H_5=18/5$.  The two induction margins are
\[
\begin{aligned}
 \Phi_{k+1}(H_k)-G_{k+1}
 &=\frac{4(34k^4+288k^3+683k^2+278k-403)}
 {(k+4)^2(k+5)^2(5k^2+10k+1)},\\
 H_{k+1}-\Phi_{k+1}(G_k)
 &=\frac{29k^2+91k+54}{5(k+4)^2(2k^2+8k+7)}.
\end{aligned}
\]
They are positive for $k\ge5$, so $G_k\le U_k\le H_k$.  If $r_k^*$ is
the positive root of
\[
 (k+4)^2r^2-(k+1)(7k+18)r-8(k+1)(k+2)=0,
\]
then, on writing $\epsilon_k^*=8-r_k^*$, the discriminant identity
\[
 (9k^2+103k+238)^2-4(k+4)^2(288k+864)
 =(9k^2+39k+38)^2+16(k-3)(k+2)
\]
gives $\epsilon_k^*\le G_k\le U_k$.  Hence
$\rho_{k+1}\ge\rho_k$ for $k\ge5$; the exact initial ratios are
$\rho_3=\rho_4=4$ and $\rho_5=145/32$.

For $N=4$, set $U_k=15-\rho_k$ and
\[
 G_k=\frac{1125}{10k+71},\qquad H_k=\frac{1125}{10k+66}.
\]
Solving the second recurrence for $U_{k+3}$ gives a rational map increasing
in $U_k,U_{k+1},U_{k+2}$ on the box $0\le U_j\le H_j$ for $j\ge7$.
The exact base margins are
\[
\begin{array}{c|cc}
j&U_j-G_j&H_j-U_j\\
\hline
7&4847/73179&48089/211752\\
8&120402/1635179&286983/1581034\\
9&331053/4344746&103677/701636.
\end{array}
\]
Clearing the
positive denominators in the lower and upper induction margins gives,
respectively,
\[
\begin{aligned}
P_-(s)={}&2432000s^5+100233600s^4+1609345435s^3\\
 &+12643460157s^2+48855354258s+74691137880,\\
P_+(s)={}&128000s^5+24662400s^4+750722465s^3\\
 &+8917574583s^2+46522820852s+89208809220,
\end{aligned}
\]
at $s=k-7$.  The denominators are
\[
\begin{aligned}
40(k+7)^2(k+8)(5k-2)(5k+3)(5k+8)(10k+101),\\
10(k+7)^2(k+8)(5k+48)(10k-9)(10k+1)(10k+11),
\end{aligned}
\]
so the induction proves $G_k\le U_k\le H_k$ for $k\ge7$.  Since
$H_{k+1}<G_k$, the ratios increase.  Direct substitution of
$m_3=2,m_4=9,m_5=43,m_6=245,m_7=1557,m_8=10829$ checks the preceding
indices and the lower bound $\rho_k\ge9/2$.

For $N=5$, set $U_k=24-\rho_k$ and
\[
 G_k=\frac{2880}{10k+109},\qquad H_k=\frac{288}{k+10}.
\]
The third recurrence similarly gives an increasing three-variable map for
$k\ge6$.  Its exact base margins are
\[
\begin{array}{c|cc}
j&U_j-G_j&H_j-U_j\\
\hline
6&92/1859&10/11\\
7&19597/163248&11345/15504\\
8&3424/21315&4272/7105.
\end{array}
\]
After clearing the positive denominators, the two induction numerators at
$s=k-6$ are
\[
\begin{aligned}
R_-(s)={}&3750000s^6+174475000s^5+3214281000s^4\\
 &+30095015750s^3+153613752915s^2\\
 &+416461589985s+482743428300,\\
R_+(s)={}&1900s^5+76662s^4+1135088s^3+7550562s^2\\
 &+22473804s+23833584,
\end{aligned}
\]
with denominators
\[
\begin{aligned}
2(k+8)^2(k+10)^2(10k-11)(10k-1)(10k+9)(10k+139),\\
k(k-2)(k-1)(k+8)^2(k+10)^2(k+13).
\end{aligned}
\]
Thus $G_k\le U_k\le H_k$ for $k\ge6$, and again $H_{k+1}<G_k$.
The exact initial ratios
$9/2,44/9,6,76/11$ complete the $N=5$ proof.
\end{proof}

\begin{lemma}[The $SU(3)$ boundary]\label{lem:su3-boundary}
The ratio inequalities in \cref{thm:su-low-rank-ratios} imply
$\Qthree^{\SU(3)}(n)>0$ for every odd $n\ge11$.
\end{lemma}

\begin{proof}
The interior terms in the symmetric moment formula are nonnegative by the
same AM--GM argument as in \cref{lem:su-boundary}.  For the depth-two
boundary put
\[
 \rho_{n-2}=4+s,\qquad \rho_{n-1}=4+s+t,\qquad s,t\ge0.
\]
Substitution of the $N=3$ recurrence from
\cref{thm:su-low-rank-ratios} gives
\[
 L_2=\frac{m_{n-2}}{2(n+3)^2(n+4)^2}P(n,s,t),
\]
where
\[
\begin{aligned}
P={}&n^6(s^2+st+4s+4t+8)
 +n^5(9s^2+9st+24s+36t+56)\\
&+n^4(119s^2+119st+884s+476t+2104)\\
&+n^3(479s^2+479st+4256s+1916t+10120)\\
&+n^2(716s^2+716st+7184s+2864t+17088)\\
&+n(636s^2+636st+6528s+2544t+14784)\\
&+576s^2+576st+4608s+2304t+9216.
\end{aligned}
\]
Every coefficient is nonnegative and the constant term is positive.
Thus $L_2>0$, while every interior term is nonnegative.
\end{proof}

\begin{theorem}[The $SU(2)$ row]\label{thm:su2-row}
For every $n\ge0$, $\Qthree^{\SU(2)}(n)\ge0$.
\end{theorem}

\begin{proof}
Parametrize $SU(2)$ conjugacy classes by $\alpha\in[0,\pi]$.  The class
density is $(2/\pi)\sin^2\alpha\,d\alpha$, and
$\chiad(\alpha)=1+2\cos(2\alpha)$.  For independent angles $\alpha,\beta$,
put $\theta=\alpha+\beta$ and $\phi=\alpha-\beta$.  Folding the resulting
diamond uses the identities
\[
\begin{gathered}
 \chiad(\alpha)+\chiad(\beta)=2+4\cos\theta\cos\phi,\\
 (\chiad(\alpha)-\chiad(\beta))^2
   =16\sin^2\theta\sin^2\phi,\\
 4\sin^2\alpha\sin^2\beta=(\cos\theta-\cos\phi)^2,
 \qquad d\alpha\,d\beta=\tfrac12d\theta\,d\phi.
\end{gathered}
\]
Thus folding by the $C_2$ Weyl symmetries sends
\[
 \frac12(\chiad(\alpha)-\chiad(\beta))^2\,d\mu(\alpha)d\mu(\beta)
\]
to the normalized $\operatorname{USp}(4)$ Weyl density
\[
 \frac8{\pi^2}(\cos\theta-\cos\phi)^2
 \sin^2\theta\sin^2\phi\,d\theta d\phi
 \quad(0\le\theta,\phi\le\pi)
\]
from \cite[Chapter~IV, (1.11)]{BrockertomDieck}.  The normalization follows
at $n=0$ from Schur orthogonality.  If $V$ is
the defining representation of $\operatorname{USp}(4)$, then
\[
 \chiad(\alpha)+\chiad(\beta)
 =2+4\cos\theta\cos\phi=\chi_{\Lambda^2V}(\theta,\phi).
\]
Consequently
\[
 \frac12\Qthree^{\SU(2)}(n)
 =\dim\bigl(((\Lambda^2V)^{\otimes n})^{\operatorname{USp}(4)}\bigr)
 \ge0.
\]
\end{proof}

\begin{lemma}[Derangement-ratio gap]\label{lem:derangement-ratio-gap}
If $D_r=!r$, then for every $r\ge3$,
\[
 \frac{D_{r+2}}{D_{r+1}}-\frac{D_{r+1}}{D_r}\ge\frac7{18}.
\]
\end{lemma}

\begin{proof}
The recurrence $D_{r+1}=(r+1)D_r+(-1)^{r+1}$ gives
\[
 \frac{D_{r+2}}{D_{r+1}}-\frac{D_{r+1}}{D_r}
 =1+\frac{(-1)^{r+2}}{D_{r+1}}-\frac{(-1)^{r+1}}{D_r}.
\]
For even $r$ this is greater than $1$.  For odd $r\ge3$, use
$D_r\ge D_3=2$ and $D_{r+1}\ge D_4=9$ to obtain
$1-1/9-1/2=7/18$.
\end{proof}

\begin{lemma}[Uniform type $A$ transition strip]\label{lem:su-transition-strip}
For every $N\ge6$ and every integer $k$ with
\[
  N-1\le k\le N+33,
\]
one has
\[
  R_{N,k}=m_{k+2}^{(N)}m_k^{(N)}-(m_{k+1}^{(N)})^2>0.
\]
\end{lemma}

\begin{proof}
Let $D_r=!r$.  By Schur--Weyl duality and RSK,
\[
 a!-E_a^{(N)}
 =\#\{\pi\in S_a:\operatorname{lds}(\pi)>N\}.
\]
If the longest decreasing subsequence has length greater than $N$, some
fixed set of $N+1$ positions appears in decreasing relative order.  The
union bound therefore gives
\[
 0\le a!-E_a^{(N)}
 \le a!\frac{\binom a{N+1}}{(N+1)!}.
\]
Fix $k$ in the displayed strip and put $K=k+2$.  Comparing the binomial
transforms for $m_r^{(N)}$ and $D_r$, for $r\le K$ one obtains
\[
 \left|m_r^{(N)}-D_r\right|
 \le3r!\frac{\binom K{N+1}}{(N+1)!}
 \le \tau D_r,
 \qquad
 \tau:=9\frac{\binom K{N+1}}{(N+1)!},
\]
where $\sum_{j\ge0}1/j!<3$ and $D_r\ge r!/3$ for $r\ge2$.
The latter bound follows directly from
$D_r/r!=\sum_{j=0}^r(-1)^j/j!$: it is an alternating sum, with the
minimum $1/3$ attained at $r=3$.

By \cref{lem:derangement-ratio-gap},
\[
 D_{k+2}D_k-D_{k+1}^2\ge\frac7{18}D_kD_{k+1}.
\]
Writing $m_r^{(N)}=D_r+e_r$, expanding $R_{N,k}$, and using
$|e_r|\le\tau D_r$ together with
\[
 D_{k+1}\le(k+2)D_k,
 \qquad D_{k+2}\le(k+3)D_{k+1},
\]
gives
\[
 R_{N,k}\ge D_kD_{k+1}
 \left[\frac7{18}-\tau(4k+10)-\tau^2(2k+5)\right].
 \tag{\(*\)}
\]

For $6\le N\le18$, the exact C++/GMP hook-length replay evaluates all
$35$ values $N-1\le k\le N+33$ and finds every integer $R_{N,k}$ strictly
positive.  Now assume $N\ge19$ and write $k=N+d$, $-1\le d\le33$.
The bracket in \((*)\) is smallest at $d=33$.  In that case
\[
 \tau_N=9\frac{\binom{N+35}{N+1}}{(N+1)!},
 \qquad
 \frac{\tau_{N+1}}{\tau_N}=\frac{N+36}{(N+2)^2}.
\]
Both error terms decrease for $N\ge19$.  Indeed, after putting $s=N-19$,
the cleared numerator for the first ratio comparison is
\[
 4s^3+382s^2+10478s+83928,
\]
and for the second it is
\[
 2s^5+277s^4+14446s^3+362171s^2
 +4408498s+20862654.
\]
All coefficients are positive.  At the initial value $N=19,d=33$, exact
rational arithmetic gives
\[
 \tau_{19}=\frac{3737279683}{3143467008000},
\]
and the bracket in \((*)\) equals
\[
 \frac{1280170982560916731574699}
 {9881384830384472064000000}>\frac1{10}.
\]
Thus \((*)\) is strictly positive for every $N\ge19$ and every allowed
$d$.  The GMP checker independently reconstructs the finite moments, the
base rational margin, and the two shifted coefficient lists.
\end{proof}

\begin{lemma}[Fixed-band determinant criterion]\label{lem:su-fixed-band}
Let $\mu$ be a probability measure on $[-1,A]$, where
$A>c>u>b>1$, and put $m_j=\int x^j\,d\mu(x)$.  If
\[
 p_B=\mu([b,u]),\qquad p_C=\mu([c,A])
\]
satisfy, for an odd integer $k$,
\[
 p_Bp_Cb^k(c-u)^2>2(c+1)^2,
 \tag{\ref{lem:su-fixed-band}.1}
\]
then $m_{k+2}m_k-m_{k+1}^2>0$.
\end{lemma}

\begin{proof}
For independent $X,Y$ with law $\mu$,
\[
 m_{k+2}m_k-m_{k+1}^2
 =\frac12\mathbf E\bigl[X^kY^k(X-Y)^2\bigr].
\]
The two symmetric rectangles $[b,u]\times[c,A]$ contribute at least
$P=p_Bp_Cb^kc^k(c-u)^2$.  The negative contribution for which the positive
variable $x$ lies in $[c,A]$ has absolute value at most
$H=\int_{[c,A]}x^k(x+1)^2\,d\mu(x)$.  For $x\ge c$, the quotient
$(x+1)/(x-u)$ is decreasing.  Since $p_C\le1$,
\eqref{lem:su-fixed-band}.1 gives
\[
 p_Bb^k(x-u)^2
 \ge p_Bb^k(c-u)^2\frac{(x+1)^2}{(c+1)^2}>2(x+1)^2.
\]
Integrating first over $[b,u]$ and then over $[c,A]$ shows $P>2H$.
The remaining negative contribution, with positive variable in $(0,c)$,
is at most $L=(c+1)^2c^k$, while \eqref{lem:su-fixed-band}.1 gives
$P/2>L$.  All same-sign contributions are nonnegative, proving the claim.
\end{proof}

\begin{lemma}[Quadratic type-$A$ transition window]
\label{lem:su-quadratic-window}
For $N\ge25$,
\[
 R_{N,k}>0
 \qquad\left(N+16\le k\le
 \left\lfloor\frac{N^2}{11}\right\rfloor-2\right).
\]
\end{lemma}

\begin{proof}
Put $K=\lfloor N^2/11\rfloor$, $\theta=121/176$, and $D_r=!r$.
The RSK formula in the proof of \cref{lem:su-transition-strip} gives, for
$r\le K$,
\[
 |m_r^{(N)}-D_r|
 \le3r!\frac{\binom K{N+1}}{(N+1)!}.
\]
Since $K<(N+1)^2/11$, $(N+1)!\ge((N+1)/e)^{N+1}$, and $e<11/4$,
\[
 \frac{\binom K{N+1}}{(N+1)!}\le\theta^{N+1}.
\]
Also $D_r\ge r!/3$ for $r\ge2$, so
$|m_r^{(N)}-D_r|\le\tau D_r$ with $\tau=9\theta^{N+1}$.
Expanding $R_{N,k}$ about the derangement determinant as in
\cref{lem:su-transition-strip} yields
\[
 R_{N,k}\ge D_kD_{k+1}
 \left[\frac7{18}-\tau(4k+10)-\tau^2(2k+5)\right].
 \tag{\ref{lem:su-quadratic-window}.1}
\]
For $k\le K$ the error terms are bounded by
\[
 9\theta^{N+1}\left(\frac{4N^2}{11}+10\right),\qquad
 81\theta^{2N+2}\left(\frac{2N^2}{11}+5\right).
\]
They decrease for $N\ge25$: after clearing denominators, the two successive
ratio inequalities reduce respectively to
\[
 220N^2-968N+5566>0,
 \qquad 32670N^2-58564N+869143>0.
\]
At $N=25$ their exact rational sum is less than $129/1000<7/18$.
Thus the bracket in \eqref{lem:su-quadratic-window}.1 is positive.  The
C++/GMP replay independently checks the two cleared polynomials and the
exact $N=25$ comparison.
\end{proof}

\begin{proposition}[Uniform type-$A$ endpoint overlaps]
\label{prop:su-endpoint-overlaps}
The transition determinant satisfies $R_{N,k}>0$ in both ranges
\[
 \begin{array}{c|c}
 N&\text{odd }k\\ \hline
 N\ge24&k\ge71,\\
 6\le N\le23&k\ge53.
 \end{array}
\]
\end{proposition}

\begin{proof}
Let $U$ be Haar-distributed in $\SU(N)$, put
$Y=|\operatorname{tr}U|^2$, $X=Y-1$, and let $\mu_N$ be the law of $X$.
Then $X\in[-1,N^2-1]$.  The $\SU(N)$ and $\mathrm U(N)$ moments of $Y$
agree, and
\[
 \mathbf E Y^r=
 \sum_{\lambda\vdash r,\ \ell(\lambda)\le N}(f^\lambda)^2\le r!,
 \tag{\ref{prop:su-endpoint-overlaps}.1}
\]
with equality for $r\le N$.

\ifginibrefullproof
First suppose $N\ge24$.  Set $z=y-7/2$ and
\[
\begin{aligned}
P(z)={}&925503648+348262482z-142832255z^2-40966535z^3\\
 &+8723312z^4+880405z^5-242416z^6+14686z^7-274z^8,
\end{aligned}
\]
and $T=1467426013$.  On $y\in[5/2,9/2]$, $|z|\le1$,
$0\le(y-5/2)(9/2-y)\le1$, and $|P(z)|\le T$; outside that interval the
quadratic factor is nonpositive.  Consequently
\[
 \mathbf1_{[5/2,9/2]}(y)\ge
 \frac{(y-5/2)(9/2-y)P(y-7/2)^2}{T^2}.
\]
Using the factorial moments through degree $18$ gives
\[
 \mu_N([3/2,7/2])\ge
 \eta_B:=
 \frac{122923229137548665407}{65536\cdot1467426013^2}.
\]
Paley--Zygmund applied to $Y^{12}$, using the moments $12!$ and $24!$,
gives
\[
 \mu_N([4,N^2-1])\ge
 \eta_C:=\frac{88255484124721}{992717442773183102976}.
\]
Exact rational arithmetic gives $\eta_B\eta_C(3/2)^{71}>200$.
Thus \cref{lem:su-fixed-band} applies with
$(b,u,c)=(3/2,7/2,4)$ for every odd $k\ge71$.

Now let $6\le N\le23$.  Put $z=y-43/10$ and
\[
 Q(z)=959324176457+276017837985z-58083597219z^2
       -11664458272z^3+1230335979z^4,
\]
and $S=935204207737834/625$.  On $[29/10,57/10]$ one has
$|z|\le7/5$, $0\le(y-29/10)(57/10-y)\le49/25$, and $|Q(z)|\le S$.
The same signed-polynomial minorant gives
\[
 \mu_N([19/10,47/10])\ge
 \delta_B:=
 \frac{44525786465344542780431067563007}
 {43884276324717505323728973379833856}.
\]
For $N\ge10$ this uses the stable moments through degree $10$; for
$6\le N\le9$ the exact hook calculation has nonnegative excess, whose
minimum is $1513726621221888441$ at $N=9$.

Paley--Zygmund applied to $Y^{15}$ gives
\[
 \mu_N([5,N^2-1])\ge
 \delta_C:=\frac{471756427}{178410035553750000}.
\]
For $N\ge15$, this follows from $\mathbf EY^{15}=15!$ and
$\mathbf EY^{30}\le30!$.  For $6\le N\le14$, the exact hook calculation
checks both $\mathbf EY^{15}>6^{15}$ and
\[
 \left(1-\frac{6^{15}}{\mathbf EY^{15}}\right)^2
 \frac{(\mathbf EY^{15})^2}{\mathbf EY^{30}}\ge\delta_C;
\]
the smallest excess occurs at $N=14$ and is strictly positive.
Finally, exact rational arithmetic gives
\[
 \delta_B\delta_C(19/10)^{52}(3/10)^2>72.
\]
For every odd $k\ge53$, \cref{lem:su-fixed-band} applies with
$(b,u,c)=(19/10,47/10,5)$.  The C++/GMP replay reconstructs all finite
hook moments, polynomial expectations, Paley--Zygmund domain checks, and
the two terminal rational inequalities used above.
\else
For $N\ge24$, the degree-$18$ signed-polynomial minorant reconstructed from
the manifested verifier source by the C++/GMP replay gives
$\mu_N([3/2,7/2])\ge
\eta_B=122923229137548665407/(65536\cdot1467426013^2)$.
Paley--Zygmund applied to $Y^{12}$ gives
$\mu_N([4,N^2-1])\ge
\eta_C=88255484124721/992717442773183102976$, and exact rational arithmetic
gives $\eta_B\eta_C(3/2)^{71}>200$.  Thus
\cref{lem:su-fixed-band} applies with $(b,u,c)=(3/2,7/2,4)$.

For $6\le N\le23$, the analogous degree-$10$ minorant and exact hook
moments give
$\mu_N([19/10,47/10])\ge
\delta_B=44525786465344542780431067563007/
43884276324717505323728973379833856$.
Paley--Zygmund applied to $Y^{15}$ gives
$\mu_N([5,N^2-1])\ge
\delta_C=471756427/178410035553750000$.  The replay checks the nonstable
ranks, including $\mathbf EY^{15}>6^{15}$, and verifies
$\delta_B\delta_C(19/10)^{52}(3/10)^2>72$.  Hence
\cref{lem:su-fixed-band} applies with $(b,u,c)=(19/10,47/10,5)$ for every
odd $k\ge53$.
\fi
\end{proof}

\begin{proposition}[Finite type-$A$ overlap rows]
\label{prop:su-finite-overlaps}
One has $R_{N,k}>0$ for the following odd transition indices:
\[
\begin{array}{c|l@{\qquad}c|l}
N&k&N&k\\ \hline
24&57,59,\ldots,69&6&39,41,\ldots,51\\
25&59,61,\ldots,69&7,8&41,43,\ldots,51\\
26&61,63,\ldots,69&9,10&43,45,\ldots,51\\
27&65,67,69&11,12&45,47,49,51\\
&&13,14&47,49,51\\
&&15,16&49,51\\
&&17,18&51.
\end{array}
\]
\end{proposition}

\begin{proof}
For $24\le N\le27$, put $K=k+2$ and
$\tau=9\binom K{N+1}/(N+1)!$.  The RSK union bound and determinant
expansion in \cref{lem:su-transition-strip} give
\[
 R_{N,k}\ge D_kD_{k+1}
 \left[\frac7{18}-\tau(4k+10)-\tau^2(2k+5)\right].
\]
Exact rational evaluation over the displayed finite rows gives a bracket
greater than $3/10$ in every case.

For $6\le N\le18$, the exact hook formula in
\eqref{prop:su-endpoint-overlaps}.1 supplies $\mathbf EY^r$ through $r=53$;
binomial transformation gives the adjoint moments, and direct integer
evaluation gives $R_{N,k}>0$ at every displayed pair.  The C++/GMP checker
recomputes these partitions and hook dimensions, requires exact division at
every step, matches a hard-coded positive row minimum for each $N$, and
rejects any missing transition index.
\end{proof}

\begin{theorem}[Finite type $A$ ratio theorem]\label{thm:su-ratio}
For every $N\ge6$,
\[
  \rho_3^{(N)}=\frac92,
  \qquad
  \frac92\le \rho_k^{(N)}\le\rho_{k+1}^{(N)}
  \quad(k\ge3).
\]
\end{theorem}

\begin{proof}
The stable part $3\le k\le N-2$ follows from the derangement moments
above.  Indeed $m_r^{(N)}=!r$ for $r\le N$, and
\cref{lem:derangement-ratio-gap} gives both monotonicity and the lower
bound $\rho_k^{(N)}\ge\rho_3^{(N)}=9/2$.

For even $k=2q$, Cauchy--Schwarz applied to the real functions
$\chiad^q$ and $\chiad^{q+1}$ gives
\[
  (m_{k+1}^{(N)})^2\le m_k^{(N)}m_{k+2}^{(N)}.
\]
The denominators are positive by \cref{lem:adjoint-moment-positive}, so
this proves the even transition inequalities.

The uniform strip \cref{lem:su-transition-strip} covers
$N-1\le k\le N+33$.  The remaining odd transition range is covered by the
following proved overlaps:
\[
\begin{array}{c|c}
\text{range} & \text{certificate}\\
\hline
N-1\le k\le N+33
  & \cref{lem:su-transition-strip}\\
k\text{ even}
  & \text{Cauchy--Schwarz, as above}\\
6\le N\le18,\ k\le51
  & \cref{prop:su-finite-overlaps}\\
6\le N\le23,\ k\ge53
  & \cref{prop:su-endpoint-overlaps}\\
24\le N\le27,\ k<71
  & \cref{lem:su-quadratic-window,prop:su-finite-overlaps}\\
N\ge24,\ k\ge71
  & \cref{prop:su-endpoint-overlaps}\\
N\ge28,\ k\le69
  & \cref{lem:su-quadratic-window}.
\end{array}
\]
These alternatives cover all odd $k\ge N-1$: the first line covers through
$N+33$.  For $6\le N\le18$, the finite rows meet the endpoint range
$k\ge53$; for $19\le N\le23$, that endpoint already begins before the first
remaining odd index.  For $24\le N\le27$, the finite and quadratic rows meet
$k\ge71$.  If $N\ge28$ and a remaining odd index is at most $69$, then it is
at least $N+34>N+16$, while
$\lfloor N^2/11\rfloor-2\ge69$, so
\cref{lem:su-quadratic-window} applies.  Every later odd index is covered by
\cref{prop:su-endpoint-overlaps}.

\ifginibrefullproof
The certificate output for the finite-overlap replay is:
\[
\begin{array}{ll}
\text{uniform strip }N-1\le k\le N+33 & \text{ok},\\
\text{endpoint constants} & \text{ok},\\
N=6,\ldots,18              & \text{ok},\\
\text{final line}          & \texttt{All SU(N) repair certificates verified.}
\end{array}
\]
\else
The replay reports the uniform strip, endpoint constants, and every finite
row as exact successes.
\fi
All arithmetic in the finite replay is integer or rational arithmetic; no
floating-point comparison is used in a sign decision.  The infinite parts
are the displayed inequalities in
\cref{lem:su-transition-strip,lem:su-quadratic-window,prop:su-endpoint-overlaps}.
\end{proof}

\begin{corollary}[Finite-rank type $A$ branch]\label{cor:su-finite}
For every $N\ge2$ and every $n\ge0$ satisfying $N<n+2$,
\[
  \Qthree^{\SU(N),\ad}(n)\ge0.
\]
\end{corollary}

\begin{proof}
The $N=2$ row is \cref{thm:su2-row}.  For $N=3$,
\cref{thm:su-low-rank-ratios,lem:su3-boundary} cover every odd $n\ge11$;
the hook-length formula gives the exact positive values
\[
\begin{array}{c|rrrrr}
n&1&3&5&7&9\\
\hline
\Qthree(n)/2&2&28&654&21312&902860.
\end{array}
\]
For $N=4,5$, \cref{thm:su-low-rank-ratios,lem:su-boundary} give
positivity for odd $n\ge5$.  For $N\ge6$,
\cref{thm:su-ratio,lem:su-boundary} do the same.  The remaining odd cases
$n=1,3$ follow from the first adjoint moments:
\[
\begin{array}{c|ccccc|cc}
N&m_0&m_1&m_2&m_3&m_5&\Qthree(1)&\Qthree(3)\\
\hline
2&1&0&1&1&6&2&8\\
3&1&0&1&2&32&4&56\\
4&1&0&1&2&43&4&78\\
\ge5&1&0&1&2&44&4&80.
\end{array}
\]
Here $\Qthree(1)=2m_3$ and $\Qthree(3)=2(m_5-2m_3)$.
Even $n$ is \cref{lem:even}.
\end{proof}

\begin{theorem}[Type $A$ closure]\label{thm:type-a}
For every $N\ge2$ and every $n\ge0$,
\[
  \Qthree^{\SU(N),\ad}(n)\ge0.
\]
\end{theorem}

\begin{proof}
If $N\ge n+2$, use \cref{thm:su-stable-rank}.  If $N<n+2$, use
\cref{cor:su-finite}.
\end{proof}

%======================================================================
\section{The Non-A Classical Families}
%======================================================================

The non-$A$ classical Lie algebras are represented by the standard
$B_b$, $C_b$, and $D_b$ families.  The proof uses exact adjoint moments,
Rains stabilization, and a final finite certificate.
We use the irreducible-root-system rank conventions

\[
 B_b\ (b\ge2),\qquad C_b\ (b\ge2),\qquad D_b\ (b\ge4).
\]
The low-rank coincidences $B_1=C_1=A_1$ and $D_3=A_3$ are already covered
by \cref{thm:type-a}; $D_2=A_1\times A_1$ is not simple.  The coincidence
$B_2=C_2$ causes harmless overlap between the two finite classical rows.

\begin{definition}[Classical constants and first hits]\label{def:classical-constants}
Let
\[
  C_G=\max_{g\in G}(-\chiad(g)).
\]
Define
\[
  s(G)=
  \begin{cases}
    2b-1, & G=B_b,C_b,\\
    \text{largest odd integer }\le b-3, & G=D_b,
  \end{cases}
\]
the largest odd exponent covered by the stable moment comparison.  After
the row-gated bridge through $m=29$ the first possible unbridged odd exponent is
\[
  N_{\mathrm{hit}}^{(29)}(G)=\max(63,s(G)+2).
\]
Equivalently,
\[
\begin{array}{ll}
B_b,C_b: & N_{\mathrm{hit}}^{(29)}=\max(63,2b+1),\\
D_b:     & N_{\mathrm{hit}}^{(29)}=
  \max(63,b-1)\text{ for even }b,\quad
  \max(63,b-2)\text{ for odd }b.
\end{array}
\]
\end{definition}

\begin{lemma}[Classical negative endpoint]\label{lem:classical-negative-endpoint}
For the standard $B_b,C_b,D_b$ representatives,
\[
  C_{B_b}=b,\qquad C_{C_b}=b,\qquad
  C_{D_b}=
  \begin{cases}
    b, & b\text{ even},\\
    b-2, & b\text{ odd}.
  \end{cases}
\]
\end{lemma}

\begin{proof}
Write the defining eigenvalue pairs as $z_j,z_j^{-1}$ and put
$x_j=(z_j+z_j^{-1})/2\in[-1,1]$ and $S=\sum_j x_j$, $Q=\sum_jx_j^2$.
For $B_b$ the defining trace is $T_1=1+2S$ and
$T_2=1+4Q-2b$.  Since the adjoint representation is
$\Lambda^2$ of the defining representation,
\[
  \chiad=\frac{T_1^2-T_2}{2}
        =b+2S+2(S^2-Q).
\]
This expression is affine in each $x_j$ separately, so its minimum occurs at a
vertex.  At a vertex $Q=b$ and $S=b-2k$, hence
$\chiad=2S^2+2S-b$; the minimum is $-b$.

For $D_b$, $T_1=2S$ and $T_2=4Q-2b$, again with
$\chiad=(T_1^2-T_2)/2$.  Thus
\[
  \chiad=b+2(S^2-Q).
\]
At a vertex this is $2S^2-b$.  If $b$ is even, $S=0$ is allowed and the
minimum is $-b$; if $b$ is odd, the nearest allowed values are $S=\pm1$ and
the minimum is $2-b$.

For $C_b$, the adjoint representation is $\operatorname{Sym}^2$ of the
defining representation, so
\[
  \chiad=\frac{T_1^2+T_2}{2}=2S^2+2Q-b.
\]
This is bounded below by $-b$, and equality is attained at
$x_1=\cdots=x_b=0$.
\end{proof}

\begin{lemma}[Stable classical edge]\label{lem:classical-stable-edge}
For the classical families the stable comparison proves the following
Chain differences:
\[
  D_G(2r+1)=\Qthree^G(2r+3)-4\Qthree^G(2r+1)>0
\]
whenever
\[
\begin{array}{c|c}
G & \text{stable range}\\
\hline
B_b,C_b & b\ge r+2\\
D_b & b\ge 2r+6 .
\end{array}
\]
Equivalently, the stable edge covers all odd exponents at most $s(G)$,
where $s(G)$ is defined in \cref{def:classical-constants}.
\end{lemma}

\begin{proof}
Write $T_j(U)=\operatorname{Tr}(U^j)$ in the defining representation.  The
adjoint characters are
\[
 \chiad^{\SO}=(T_1^2-T_2)/2,
 \qquad
 \chiad^{\Sp}=(T_1^2+T_2)/2.
\]
Diaconis--Shahshahani
\cite[Theorem~4, p.~56, and Theorem~6, p.~59]{DiaconisShahshahani1994}
give the exact stable joint trace moments
\[
 (T_1,T_2)_{\SO}\longmapsto(Z_1,1+\sqrt2Z_2),
 \qquad
 (T_1,T_2)_{\Sp}\longmapsto(Z_1,-1+\sqrt2Z_2),
\]
where $Z_1,Z_2$ are independent standard normals.  Hence the stable
orthogonal and symplectic adjoint variables are
\[
  X_{\mathrm{SO}}=\frac{Z_1^2-1-\sqrt2 Z_2}{2},\qquad
  X_{\mathrm{Sp}}=\frac{Z_1^2-1+\sqrt2 Z_2}{2},
\]
which have the same distribution.  Their common moment generating function
is
\[
  Z(t)=\frac{\exp(t^2/4-t/2)}{\sqrt{1-t}}.
\]

It remains to justify the stated finite-rank ranges.  For $B_b$, the
adjoint character is the Schur function $s_{(1,1)}$.  Expand
$s_{(1,1)}^s=\sum_\lambda c_\lambda s_\lambda$.  Rains' orthogonal Schur
integral criterion
\cite[\S3, paragraph preceding Lemma~3.1]{Rains1997} says that only
partitions with even row lengths contribute to the $O(2b+1)$ integral.
Such a partition has size $2s$ and at most $s$ rows; thus $2b+1\ge s$
removes the length cutoff.  The integrand is unchanged by $U\mapsto-U$, so
the $SO(2b+1)$ and $O(2b+1)$ integrals agree.

For $C_b$, the adjoint character is $s_{(2)}$.  Every partition in the
Pieri expansion of $s_{(2)}^s$ has at most $s$ rows.  Rains' symplectic
criterion in the same source paragraph retains precisely the partitions in
which every row length has even multiplicity.  Their number of rows is even
and at most $2\lfloor s/2\rfloor$, so $b\ge\lfloor s/2\rfloor$ removes the
symplectic length cutoff.

For $D_b$, expand $\chiad^s$ into monomials
$T_1^{2j}T_2^{s-j}$, each of weighted trace degree $2s$.  If $b\ge s+1$,
then $2s<2b$, so Diaconis--Shahshahani's exact orthogonal trace-moment
identity applies.  In this strict degree range the $SO(2b)$ integral equals
the $O(2b)$ integral: the difference is the determinant-isotypic integral,
and the determinant representation first occurs in matrix-entry degree
$2b$.  This proves the finite-rank stabilization for every adjoint moment
$m_s$ in the three stated ranges.

For independent copies of the common stable variable, the exponential
generating function of $Q_3$ is
\[
  F(t)=\mathbf E[(X-Y)^2e^{t(X+Y)}]
      =\bigl((1-t)^{-3}+(1-t)^{-1}\bigr)e^{-t+t^2/2}.
\]
Hence
\[
  F''(t)-4F(t)=P(t)G(t),
\]
where
\[
  G(t)=e^{-t+t^2/2},\qquad
  P(t)=12(1-t)^{-5}-7(1-t)^{-3}-4(1-t)^{-1}+(1-t).
\]
Write $P(t)=\sum p_k t^k$ and
$G(t)=\sum (-1)^k h_k t^k$.  The coefficients satisfy
\[
  p_{k+1}>p_k,\qquad h_k\ge h_{k+1}>0 \qquad(k\ge0),
\]
because
\[
  p_k=12\binom{k+4}{4}-7\binom{k+2}{2}-4+[k=0]-[k=1],
\]
and
\[
  (k+1)h_{k+1}=h_k+h_{k-1},\qquad h_{-1}=0,\quad h_0=1.
\]
Indeed, $p_0=2$, while $p_1-p_0=32$, $p_2-p_1=100$, and for $k\ge2$,
\[
 p_{k+1}-p_k
 =12\binom{k+4}{3}-7(k+2)
 =(k+2)\bigl(2(k+4)(k+3)-7\bigr)>0.
\]
The recurrence gives $h_1=1$ and positivity of every $h_k$ by induction.
For $k\ge1$, its shifted form
\[
  k h_k=h_{k-1}+h_{k-2}\ge h_{k-1}
\]
then implies
\[
  h_{k+1}=\frac{h_k+h_{k-1}}{k+1}\le h_k.
\]
For $N=2r+1$,
\[
  D_{\mathrm{st}}(N)=
  N!\sum_{j=0}^{(N-1)/2}
  \bigl(p_{2j+1}h_{N-2j-1}-p_{2j}h_{N-2j}\bigr)>0.
\]
Here every summand is strictly positive: with
$a=N-2j-1$, the coefficient inequalities give
\[
 p_{2j+1}h_a-p_{2j}h_{a+1}
 =(p_{2j+1}-p_{2j})h_a+p_{2j}(h_a-h_{a+1})>0.
\]
The Chain difference at $2r+1$ uses only
$m_0^G,\ldots,m_{2r+5}^G$.  For $B_b,C_b$, the condition $b\ge r+2$
implies $b\ge\lfloor(2r+5)/2\rfloor$; for $D_b$, the condition
$b\ge2r+6$ is exactly $b\ge(2r+5)+1$.  The preceding stabilization
therefore identifies every required finite-rank moment with the stable
one.  Substitution gives the displayed inequality.
\end{proof}

\begin{definition}[Classical row gate]\label{def:row-gate}
For an odd Chain target $2m+3$ and a finite row $G$ with $C_G<296$, the
row is active in the exact middle bridge if
\[
  0\le m\le29,\qquad s(G)<2m+3\le61.
\]
The stable edge \cref{lem:classical-stable-edge} covers odd exponents
$\le s(G)$, while the post-bridge hit begins at
$N_{\mathrm{hit}}^{(29)}(G)=\max(63,s(G)+2)$.  Thus the row gate is exactly
the finite interval between stable coverage and the post-$m=29$ tail.
\end{definition}

\begin{lemma}[Pushforward tail criterion]\label{lem:pushforward-tail}
Let $X=\chiad(g)$ and $Y=\chiad(h)$ for independent Haar variables, and
let $\psi_G$ be the positive finite measure on $\mathbb R$ defined by
\[
  \int f(u)\,d\psi_G(u)
  =
  \mathbf E\bigl[(X-Y)^2f(X+Y)\bigr].
\]
Then $\psi_G(\mathbb R)=2$ and
\[
  \Qthree^G(n)=\int u^n\,d\psi_G(u).
\]
For $c>2C_G$ and $P_G(c)=\psi_G([c,2\dim G])$,
\[
  \Qthree^G(n)\ge c^nP_G(c)-2(2C_G)^n
\]
for every odd $n$.  More generally, for $0<a<2C_G<c$ and
$N_G(a)=\psi_G([-2C_G,-a])$,
\[
  \Qthree^G(n)\ge
  c^nP_G(c)-(2C_G)^nN_G(a)-2a^n .
\]
Therefore the inequality
\[
  P_G(c)\ge N_G(a)\left(\frac{2C_G}{c}\right)^n
        +2\left(\frac{a}{c}\right)^n
\]
implies $\Qthree^G(n)\ge0$, and if it holds at an odd $n_0$ it holds for
all later odd exponents.
\end{lemma}

\begin{proof}
The first identity is the definition of $\psi_G$ with
$f(u)=u^n$.  By \cref{lem:adjoint-moment-positive},
$\mathbf E X=0$ and $\mathbf E X^2=1$, hence
$\psi_G(\mathbb R)=\mathbf E(X-Y)^2=2$.
For odd $n$, split the integral into $u\ge c$, $0\le u<c$, and $u<0$.
The middle piece is nonnegative.  Since $\chiad\ge -C_G$, the negative
support is contained in $[-2C_G,0]$.  On that interval
$u^n\ge-(2C_G)^n$, and its $\psi_G$-mass is at most
$\psi_G(\mathbb R)=2$, giving the first lower bound.  For the second, on
$[-2C_G,-a]$ one has $u^n\ge-(2C_G)^n$, while on $[-a,0]$ one has
$u^n\ge-a^n$.  Therefore
\[
 \int_{-2C_G}^{0}u^n\,d\psi_G(u)
 \ge -(2C_G)^nN_G(a)-a^n\psi_G([-a,0])
 \ge -(2C_G)^nN_G(a)-2a^n.
\]
After division by $c^n$, the two terms in the sufficient condition are
nonnegative multiples of $(2C_G/c)^n$ and $(a/c)^n$.  Both bases lie in
$(0,1)$, proving the monotonicity assertion.
\end{proof}

\begin{lemma}[Weyl caps and elementary pushforward conversion]
\label{lem:weyl-cap-pushforward}
Let $G=D_b=\SO(2b)$, let $r=b$, $d=b(2b-1)$, and let
$d_1,\ldots,d_r$ be the degrees of its Weyl group $W$.  Use the
type-$D_b$ roots $e_i\pm e_j$, all of squared length $2$, and write
\[
 \sum_{\alpha>0}\alpha(\theta)^2=\kappa|\theta|^2.
\]
Suppose $R>0$ is small enough that the ball
$\kappa|\theta|^2\le R$ lies in one torus eigenangle chart, and suppose
\[
 \prod_{\alpha>0}
 \left(\frac{\sin(\alpha(\theta)/2)}{\alpha(\theta)/2}\right)^2
 \ge\eta_R>0
 \tag{\ref{lem:weyl-cap-pushforward}.1}
\]
throughout that ball.  If $X=\chi_{\rm ad}(g)$ for Haar $g$, then
\[
 \Pr\{X\ge d-R\}\ge
 \frac{\eta_R}{|W|}
 \frac{\prod_{i=1}^r d_i!}
 {(2\pi)^{r/2}\Gamma(d/2+1)}
 \left(\frac{R}{2\kappa}\right)^{d/2}.
 \tag{\ref{lem:weyl-cap-pushforward}.2}
\]

Write $C=C_G$, put $x_0=d-R$, and let $P_G$ be the pushforward mass of
\cref{lem:pushforward-tail}.  If $x_0-C>2C$ and $x_0\ge1$, then
\[
 P_G(x_0-C)\ge
 \Pr\{X\ge x_0\}\frac{(x_0-1)^2}{C+1}.
 \tag{\ref{lem:weyl-cap-pushforward}.3}
\]
If $t>1$ and $x_0-t>2C$, then
\[
 P_G(x_0-t)\ge
 \Pr\{X\ge x_0\}(x_0-t)^2(1-t^{-2}).
 \tag{\ref{lem:weyl-cap-pushforward}.4}
\]
\end{lemma}

\begin{proof}
For $g=\exp(\theta)$ in the chosen torus chart,
\[
 d-\chi_{\rm ad}(g)
 =2\sum_{\alpha>0}(1-\cos\alpha(\theta))
 \le\sum_{\alpha>0}\alpha(\theta)^2
 =\kappa|\theta|^2.
\]
Weyl integration, \eqref{lem:weyl-cap-pushforward}.1, and the
Macdonald--Mehta integral therefore reduce the cap mass to
\[
 \frac{\eta_R}{|W|(2\pi)^r}
 \int_{\kappa|\theta|^2\le R}
 \prod_{\alpha>0}\alpha(\theta)^2\,d\theta.
\]
The root product is homogeneous of degree $d-r$.  Taking the radial part of
the Macdonald--Mehta Gaussian integral gives
\[
 \int_{\kappa|\theta|^2\le R}
 \prod_{\alpha>0}\alpha(\theta)^2\,d\theta
 =\frac{(2\pi)^{r/2}\prod_i d_i!}{\Gamma(d/2+1)}
  \left(\frac{R}{2\kappa}\right)^{d/2},
\]
which proves \eqref{lem:weyl-cap-pushforward}.2.

Let $Y$ be an independent copy of $X$.  Since $Y\ge-C$ and
$\mathbf EY=0$, if $p=\Pr\{Y>1\}$ then
$0=\mathbf EY\ge p-C(1-p)$; hence
$\Pr\{Y\le1\}\ge1/(C+1)$.  On
$\{X\ge x_0,Y\le1\}$ one has
$X+Y\ge x_0-C$ and $(X-Y)^2\ge(x_0-1)^2$.  Independence proves
\eqref{lem:weyl-cap-pushforward}.3.  Also $\mathbf EY^2=1$ by
\cref{lem:adjoint-moment-positive}, so Chebyshev gives
$\Pr\{|Y|\le t\}\ge1-t^{-2}$.  On
$\{X\ge x_0,|Y|\le t\}$, both $X+Y$ and $X-Y$ are at least $x_0-t$.
This proves \eqref{lem:weyl-cap-pushforward}.4.
\end{proof}

\begin{lemma}[Root-normalized type-$B$ Weyl caps]
\label{lem:b-weyl-cap-pushforward}
Let $G=B_b=\SO(2b+1)$, let $r=b$, $d=b(2b+1)$, and let
$d_i=2i$ for $1\le i\le b$.  Use the actual type-$B_b$ roots
\[
 e_i\pm e_j\quad(i<j),\qquad e_i\quad(1\le i\le b),
\]
and put $\kappa=2b-1$.  Suppose the ball
$\kappa|\theta|^2\le R$ lies in one torus eigenangle chart and
\[
 \prod_{\alpha>0}
 \left(\frac{\sin(\alpha(\theta)/2)}{\alpha(\theta)/2}\right)^2
 \ge\eta_R>0
 \qquad(\kappa|\theta|^2\le R).
 \tag{\ref{lem:b-weyl-cap-pushforward}.1}
\]
If $X=\chi_{\rm ad}(g)$ for Haar $g\in G$, then
\[
 \Pr\{X\ge d-R\}\ge
 \frac{2^{-b}\eta_R}{|W|}
 \frac{\prod_{i=1}^b d_i!}
 {(2\pi)^{b/2}\Gamma(d/2+1)}
 \left(\frac{R}{2\kappa}\right)^{d/2}.
 \tag{\ref{lem:b-weyl-cap-pushforward}.2}
\]
With $C=C_G$, $x_0=d-R$, and $P_G$ as in
\cref{lem:pushforward-tail}, the conversion bounds
\eqref{lem:weyl-cap-pushforward}.3 and
\eqref{lem:weyl-cap-pushforward}.4 also hold.
\end{lemma}

\begin{proof}
For the displayed positive roots,
\[
 \sum_{\alpha>0}\alpha(\theta)^2
 =\sum_{i<j}\bigl((\theta_i+\theta_j)^2
                  +(\theta_i-\theta_j)^2\bigr)
  +\sum_i\theta_i^2
 =(2b-1)|\theta|^2.
\]
The Macdonald--Mehta normalization used with the degrees $2,4,\ldots,2b$
replaces each short root $e_i$ by the squared-length-$2$ normal
$\sqrt2e_i$, while leaving the roots $e_i\pm e_j$ unchanged.  Hence the
square of the actual Weyl discriminant is $2^{-b}$ times the
Macdonald--Mehta discriminant square.  Repeating the radial calculation in
the proof of \cref{lem:weyl-cap-pushforward} with this factor and with
$\kappa=2b-1$ gives \eqref{lem:b-weyl-cap-pushforward}.2.  The proofs of
\eqref{lem:weyl-cap-pushforward}.3 and
\eqref{lem:weyl-cap-pushforward}.4 use only the lower support bound for the
adjoint character, $\mathbf EX=0$, $\mathbf EX^2=1$, and independence, so
the same calculation applies to $B_b$.
\end{proof}

\begin{lemma}[One-sided moment pushforward criterion]
\label{lem:one-sided-moment-pushforward}
In the setting of \cref{lem:pushforward-tail}, put
\[
 K_j^G=\int_{\mathbb R}u^j\,d\psi_G(u)=\Qthree^G(j).
\]
Let $k,m$ be nonnegative integers, let $0<a<2C_G<c$, and suppose
$2k+m\le n$.  Define
\[
 H_{k,m}^G(a)=
 \sum_{j=0}^{2m}(-1)^j\binom{2m}{j}a^{-j}K_{2k+j}^G.
\]
Then
\[
 \int_{-a}^{0}|u|^n\,d\psi_G(u)
 \le \frac{a^{n-2k}}{4^m}H_{k,m}^G(a).
\]
Consequently,
\[
 P_G(c)\ge
 N_G(a)\left(\frac{2C_G}{c}\right)^n
 +\frac{a^{n-2k}}{4^mc^n}H_{k,m}^G(a)
 \tag{\ref{lem:one-sided-moment-pushforward}.1}
\]
implies $\Qthree^G(n)\ge0$.  If the quantities on the right are fixed and
$a,2C_G<c$, validity at one odd exponent implies validity at every later odd
exponent.
\end{lemma}

\begin{proof}
For $-a\le u\le0$, write $x=-u/a\in[0,1]$.  The arithmetic--geometric
mean inequality gives
\[
 \frac{(1+x)^2}{4}\ge x.
\]
Since $n-2k\ge m$, it follows that
\[
 |u|^n=a^nx^n
 \le \frac{a^{n-2k}}{4^m}u^{2k}(1-u/a)^{2m}.
\]
The polynomial on the right is nonnegative on all of $\mathbb R$.
Integration and binomial expansion give the first assertion.  Splitting the
negative part of the integral defining $\Qthree^G(n)$ at $-a$ gives
\[
 \Qthree^G(n)\ge
 c^nP_G(c)-(2C_G)^nN_G(a)
 -\frac{a^{n-2k}}{4^m}H_{k,m}^G(a),
\]
which proves the criterion.  The two ratios in
\eqref{lem:one-sided-moment-pushforward}.1 are strictly below one, proving
the last assertion.
\end{proof}

\begin{lemma}[Exact push moments from adjoint moments]
\label{lem:exact-push-from-adjoint}
Let $X,Y$ be independent copies of the adjoint character of a compact group,
and write $m_j=\mathbf E[X^j]$.  Then the pushforward moments in
\cref{lem:one-sided-moment-pushforward} satisfy, for every $q\ge0$,
\[
 K_q^G
 =2\sum_{i=0}^{q}\binom qi
 \bigl(m_{i+2}m_{q-i}-m_{i+1}m_{q-i+1}\bigr).
 \tag{\ref{lem:exact-push-from-adjoint}.1}
\]
Consequently $K_0^G,\ldots,K_q^G$ are determined exactly by
$m_0,\ldots,m_{q+2}$.
\end{lemma}

\begin{proof}
Expand $(X-Y)^2(X+Y)^q$ and use independence.  The two pure-square sums are
equal after replacing $i$ by $q-i$; collecting them with the mixed term gives
\eqref{lem:exact-push-from-adjoint}.1.
\end{proof}

\begin{lemma}[Stable finite-rank push moments]
\label{lem:stable-finite-rank-push-moments}
Define the integers
\[
 K_j=j![t^j]\left(
   \bigl((1-t)^{-3}+(1-t)^{-1}\bigr)e^{-t+t^2/2}
 \right).
\]
For $G=B_b$ or $C_b$ and every $0\le j\le2b-2$, one has
\[
 K_j^G=K_j.
\]
For $G=D_b$ the same identity holds for every $0\le j\le b-3$.
In particular, if $2k+2m\le2b-2$ in types $B,C$, or if
$2k+2m\le b-3$ in type $D$, then
\[
 H_{k,m}^G(a)=
 \sum_{j=0}^{2m}(-1)^j\binom{2m}{j}a^{-j}K_{2k+j}.
\]
\end{lemma}

\begin{proof}
The proof of \cref{lem:classical-stable-edge} computes the exponential
generating function
\[
 \mathbf E[(X-Y)^2e^{t(X+Y)}]
 =\bigl((1-t)^{-3}+(1-t)^{-1}\bigr)e^{-t+t^2/2}
\]
for the stable orthogonal/symplectic adjoint variable.  The same proof invokes
Rains' Schur-integral stabilization to identify the finite $B_b,C_b$ adjoint
moments through degree $2r+5$ whenever $b\ge r+2$.  Take $r=b-2$.  Since
$K_j^G$ uses adjoint moments only through degree $j+2\le2b$, the finite and
stable values agree in types $B,C$.  In type $D_b$, the exact decomposition
of \cref{lem:d-finite-moment-decomposition} has $A_q(b)=B_q(b)=0$ for
$q<b$.  Thus the finite adjoint moments used by $K_j^G$ agree with their
stable values whenever $j+2<b$, equivalently $j\le b-3$.  The final formula
is the binomial expansion in \cref{lem:one-sided-moment-pushforward}.
\end{proof}

\begin{lemma}[Trace-tail reduction]\label{lem:trace-tail}
Let $U$ be Haar-distributed in the defining representation of
$B_b,C_b,D_b$, and put
\[
  T_1(U)=\operatorname{Tr}(U),\qquad
  \tau_G(U)=
  \begin{cases}
    \operatorname{Tr}(U^2), & G=B_b,D_b,\\
    -\operatorname{Tr}(U^2),& G=C_b.
  \end{cases}
\]
For $1<a<2C_G<c$, define
\[
  p_1(G;c,a)=
  \Pr\{|T_1(U)|\ge\sqrt{2c+3a}\},\qquad
  p_2(G;a)=\Pr\{\tau_G(U)\ge a\}.
\]
Then
\[
  P_G(c)\ge c^2(1-a^{-2})(p_1(G;c,a)-p_2(G;a)),
  \qquad
  N_G(a)\le 2(C_G^2+1)p_2(G;a).
\]
\end{lemma}

\begin{proof}
The defining-character identities are
\[
  \chiad(U)=\frac{T_1(U)^2-T_2(U)}2\quad(B_b,D_b),\qquad
  \chiad(U)=\frac{T_1(U)^2+T_2(U)}2\quad(C_b).
\]
On the event $|T_1(U)|\ge\sqrt{2c+3a}$ and $\tau_G(U)\le a$ one has
$\chiad(U)\ge c+a$.  If $V$ is an independent Haar element with
$|\chiad(V)|\le a$, then $\chiad(U)+\chiad(V)\ge c$ and
$|\chiad(U)-\chiad(V)|\ge c$.  Chebyshev and Schur orthogonality give
$\Pr\{|\chiad(V)|\le a\}\ge1-a^{-2}$, proving the lower bound for $P_G(c)$.
For the negative tail, $\chiad(U)+\chiad(V)\le -a$ implies one of the two
characters is at most $-a/2$.  Put $X=\chiad(U)$, $Y=\chiad(V)$, and
$A=\{X\le-a/2\}$.  Symmetry, followed by independence and
\cref{lem:adjoint-moment-positive}, gives
\begin{align*}
 N_G(a)
 &\le 2\,\mathbf E\bigl[(X-Y)^2\mathbf 1_A\bigr]\\
 &=2\left(\mathbf E[X^2\mathbf 1_A]
          -2\mathbf E[X\mathbf 1_A]\mathbf E[Y]
          +\Pr(A)\mathbf E[Y^2]\right)\\
 &\le2(C_G^2+1)\Pr(A).
\end{align*}
The displayed adjoint
formulas imply $\{\chiad(U)\le-a/2\}\subseteq\{\tau_G(U)\ge a\}$.
\end{proof}

\begin{lemma}[Decoupled trace-tail reduction and excess-square negative tail]
\label{lem:decoupled-trace-tail}
Use the notation of \cref{lem:trace-tail}.  Let
\[
 1<a<2C_G<c,\qquad d>1,\qquad q>1,
\]
and put
\[
 p_1(G;c,d,q)=
 \Pr\{|T_1(U)|\ge\sqrt{2c+2d+q}\},
 \qquad p_2(G;q)=\Pr\{\tau_G(U)\ge q\}.
\]
Then
\[
 P_G(c)\ge
 c^2(1-d^{-2})\bigl(p_1(G;c,d,q)-p_2(G;q)\bigr).
 \tag{\ref{lem:decoupled-trace-tail}.1}
\]
Moreover,
\[
 N_G(a)\le2\mathbf E\!\left[(\tau_G(U)-a)_+^2\right].
 \tag{\ref{lem:decoupled-trace-tail}.2}
\]
For every $v>0$, this implies
\[
 N_G(a)\le
 \frac{8e^{-2}}{v^2}
 e^{-va}\mathbf E e^{v\tau_G(U)}.
 \tag{\ref{lem:decoupled-trace-tail}.3}
\]
\end{lemma}

\begin{proof}
On the event
$|T_1(U)|\ge\sqrt{2c+2d+q}$ and $\tau_G(U)\le q$, the two defining-character
identities in \cref{lem:trace-tail} give $\chiad(U)\ge c+d$.  For an
independent Haar element $V$, the event $|\chiad(V)|\le d$ therefore gives
both $\chiad(U)+\chiad(V)\ge c$ and
$|\chiad(U)-\chiad(V)|\ge c$.  Schur orthogonality and Chebyshev's inequality
give probability at least $1-d^{-2}$ for the latter event, proving
\eqref{lem:decoupled-trace-tail}.1.

Write $X=\chiad(U)$ and $Y=\chiad(V)$.  Partition the event
$X+Y\le-a$ into $X\le Y$ and $Y<X$.  The two contributions are equal, and
the diagonal has zero $(X-Y)^2$ weight.  On the first part,
\[
 0\le Y-X\le-a-2X.
\]
For all three classical families the defining-character identity is
$X=(T_1(U)^2-\tau_G(U))/2$, so $-2X\le\tau_G(U)$.  Hence
\[
 (X-Y)^2\le(\tau_G(U)-a)^2,
 \qquad \tau_G(U)\ge a,
\]
which proves \eqref{lem:decoupled-trace-tail}.2.  Finally, for $z\ge a$,
the maximum of $(z-a)^2e^{-vz}$ is
$4e^{-2}v^{-2}e^{-va}$, attained at $z=a+2/v$.  Applying this pointwise and
then \eqref{lem:decoupled-trace-tail}.2 proves
\eqref{lem:decoupled-trace-tail}.3.
\end{proof}

\begin{lemma}[Rains-square high-rank trace certificate]\label{lem:rains-square}
Set
\[
  r=\frac{2000001}{1000000},\qquad
  \eta=\frac{359}{2000},\qquad
  A=\frac{4571}{2000},\qquad
  L=2r+3\eta.
\]
For every $B_b,C_b,D_b$ row with $C_G\ge296$, the hypotheses of
\cref{lem:trace-tail} hold at $a=\eta C_G$ and $c=rC_G$, and the
comparison in \cref{lem:pushforward-tail} holds at the stable edge.
\end{lemma}

\begin{proof}
The positive $T_1$ tail is supplied by the one-trace density estimate of
Courteaut--Johansson--Lambert, Theorem 1.4, combined with the Mills lower
bound, giving
\[
  \Pr\{|T_1|\ge\sqrt{LC_G}\}\ge \exp(-AC_G)
\]
	throughout the checked range.  The determinant and parity shifts in the
	orthogonal cases change $T_1$ by at most a deterministic unit, and the
	one-trace $L^1$ error is absorbed by the replayed inequality
	$E_1(C_G)\le \exp(-AC_G)$.

		For the $\tau_G$ tail, use the power-map decomposition at power $2$.
		Rains' Theorem 1.5, equations (22)--(23), and Theorem 1.7 give the
		following centered case split.  Here
		$X^U_N:=2\operatorname{Re}\operatorname{Tr}V$ for Haar $V\in U(N)$;
		thus it is the trace of Rains' real-projection eigenvalue multiset and is
		centered at mean $0$,
		and $X^{\pm}_N$ is the centered trace in the nonfixed $N$-pair part of
		the orthogonal determinant component $O^\pm$.  The deterministic
		fixed-trace contribution deleted by Rains' convention is
		\[
		  \kappa(O^+(2N))=\kappa(O^-(2N))=0,\qquad
		  \kappa(O^+(2N+1))=1,\qquad
		  \kappa(O^-(2N+1))=-1.
		\]
		The full square-trace means are
		$\mathbf E\operatorname{Tr}(U^2)=1$ for $B_b,D_b$ and
		$\mathbf E\operatorname{Tr}(U^2)=-1$ for $C_b$, by the adjoint character
		identities in \cref{lem:trace-tail} and Schur orthogonality.  Thus
		centering $\tau_G$ at $\mathbf E\tau_G=1$ gives:
	\[
		\begin{array}{c|c|c}
		G & \tau_G-1 & \text{component pair-count lower bound}\\
		\hline
		B_b & X^U_b & b=C_G,\\
	D_b & X^+_{\lceil b/2\rceil}+X^-_{\lceil(b-1)/2\rceil}
	    & \lfloor b/2\rfloor\ge\lfloor C_G/2\rfloor,\\
	C_b & -X^-_{\lceil(b-1)/2\rceil}-X^+_{\lceil b/2\rceil}
		    & \lfloor b/2\rfloor=\lfloor C_G/2\rfloor .
		\end{array}
		\]
			The deterministic shifts in this table are componentwise as follows.
			For $B_b$, equation (23) with $p=2$ gives
			$O^+(2b+1)^2\sim \operatorname{Re}U(b)$ after the source-side fixed
			$+1$ eigenvalue is omitted.  Restoring that source contribution gives
			\[
			  \operatorname{Tr}(U^2)=1+X^U_b.
			\]
			For $D_b$, equation (22) with $p=2$ gives, with
			$n_0=\lceil b/2\rceil$ and $n_1=\lceil(b-1)/2\rceil$,
			\[
			  O^+(2b)^2\sim
			  O^+(2n_0)\oplus O^-(2n_1+1),
			\]
			where the source component has fixed contribution $0$, the
			$O^+(2n_0)$ component has fixed contribution $0$, and the
			$O^-(2n_1+1)$ component has a deleted fixed contribution $-1$.  Since
			$X^-_{n_1}$ denotes only the centered nonfixed trace, the centered full
			source trace is
			\[
			  \operatorname{Tr}(U^2)-1
			  \sim X^+_{\lceil b/2\rceil}
			      +X^-_{\lceil(b-1)/2\rceil}.
			\]
			For $C_b$, Rains' equation (36) identifies
			$\operatorname{Sp}(2b)$ with $O^-(2b+2)$ after omitting fixed
			$\pm1$ eigenvalues, whose trace contribution is $+1-1=0$.  Applying
			equation (22) with $n=b+1$ gives
			\[
			  O^-(2b+2)^2\sim
			  O^-(2\lceil(b+1)/2\rceil)\oplus O^+(2\lceil b/2\rceil+1).
			\]
			The first component again has deleted fixed contribution $0$, while the
			second has deleted fixed contribution $+1$.  With the centered
			nonfixed traces denoted by $X^-_{\lceil(b-1)/2\rceil}$ and
			$X^+_{\lceil b/2\rceil}$, and with
			$\mathbf E\operatorname{Tr}(U^2)=-1$, the convention
			$\tau_G=-\operatorname{Tr}(U^2)$ gives
			\[
			  -\operatorname{Tr}(U^2)-1
			  \sim -X^-_{\lceil(b-1)/2\rceil}
		       -X^+_{\lceil b/2\rceil}.
		\]
		Thus every non-unitary component has at least
		$N(C_G)=\lfloor C_G/2\rfloor$ nontrivial conjugate pairs.  In the $B_b$
	row, Courteaut--Johansson--Lambert, Theorem 1.1, applies to the unitary
	trace and projection cannot increase total variation; in the $C_b,D_b$
	rows, Courteaut--Johansson--Lambert, Theorem 1.4, applies componentwise.
	The replay checks that the $B_b$ unitary-projection error is bounded by
	the same two-component envelope used below.  For
	$N(C)=\lfloor C/2\rfloor$, Theorem 1.4 gives the component bound
	\[
	\begin{gathered}
	  E_1(N)=\frac{5(\log(2N))^{1/4}}{\sqrt{\Gamma(2N+1)}},
	  \qquad N(C)=\lfloor C/2\rfloor .
\end{gathered}
\]
The $L^1$ distance of the convolution from the corresponding Gaussian
convolution is at most the sum of the two component errors.  The Gaussian
convolution has variance $2$, so the one-sided Gaussian tail gives
\[
\begin{gathered}
	\Pr\{\tau_G\ge \eta C_G\}
	  \le
	  \exp\!\left(-\frac{(\eta C_G-1)^2}{4}\right)+2E_1(N(C_G)).
\end{gathered}
\]
For completeness, we print the scalar certificate rather than leaving it as
an unnamed replay.  Put
\[
 \widehat E(C)=5(\log(2C))^{1/4}
   \exp\{-C(\log(2C)-1)\},
 \qquad
 E_U(C)=\frac{19C^{1/4}(\log C)^{1/2}}{\Gamma(C+2)},
\]
\[
 T(C)=\exp\!\left(-\frac{(\eta C-1)^2}{4}\right)
       +2\widehat E(\lfloor C/2\rfloor),qquad
 d(C)=1-\frac1{\eta^2C^2}.
\]
The elementary bound $(2C)!\ge(2C/e)^{2C}$ gives
$E_1(C)\le\widehat E(C)$.  With the rational constants $r,\eta,A,L$
displayed in the statement, define
\[
\begin{aligned}
 G(C)={}&(A-L/2)C+\tfrac12\log(LC)-\log(LC+1)
          -\tfrac32\log2,\\
 R_0(C)={}&e^{AC}T(C),\\
 R_1(C)={}&\frac{r(1+C^{-2})}{2d(C)}e^{AC}T(C)
             \left(\frac2r\right)^C,\\
 R_2(C)={}&\frac{4r}{\eta^3C^2d(C)}
       \exp\!\left(-\bigl(\log(r/\eta)-A\bigr)C\right).
\end{aligned}
\tag{\ref{lem:rains-square}.1}
\]
The finite scalar certificate consists exactly of
\[
 G(C)\ge0,\quad
 \widehat E(C)e^{AC}\le1,\quad
 E_U(C)\le\widehat E(\lfloor C/2\rfloor),\quad
 R_i(C)\le\frac12\ (i=0,1,2).
 \tag{\ref{lem:rains-square}.2}
\]
The first two comparisons control the Gaussian lower tail and its total
variation error, the third absorbs the type-$B$ unitary projection error, and
the last three are precisely the positive-tail/negative-tail ratios left by
\cref{lem:trace-tail,lem:pushforward-tail}.  The directed-MPFR replay evaluates
\eqref{lem:rains-square}.2 for every integer $296\le C\le10000$.  At
$C=296$ it also checks, with outward rounding, the sign conditions
\[
 A-L/2-(2C)^{-1}>0,quad
 A-\log(2C)+(4C\log(2C))^{-1}<0,
\]
\[
 A-\eta(\eta C-1)/2<0,quad
 A-\eta(\eta C-1)/2+\log(2/r)<0,quad
 A-\log(r/\eta)<0.
 \tag{\ref{lem:rains-square}.3}
\]
These are the derivative upper bounds used for the five monotone factors in
\eqref{lem:rains-square}.2; their left sides only improve with $C$.
Consequently the endpoint check propagates beyond $10000$.  This lists every
scalar predicate decided by the high-rank replay.
\end{proof}

\begin{proposition}[Stable and high-rank classical tail]\label{prop:classical-high}
For the classical families $B_b,C_b,D_b$, all rows with the classical
size parameter $C_G\ge296$ satisfy
\[
  \Qthree^G(n)\ge0
\]
for every odd exponent $n\ge s(G)+2$.
\end{proposition}

\begin{proof}
Rains moment stabilization identifies the adjoint moment prefix with the
stable orthogonal/symplectic moment sequence through the required degree,
and the stable Chain differences are exact positive integers.  The
finite-rank correction is controlled by \cref{lem:rains-square}: after
\cref{lem:trace-tail} reduces the correction to the two defining-trace
probabilities, the exact one-variable replay verifies the pushforward
comparison in \cref{lem:pushforward-tail} at the stable edge for every
$C_G\ge296$.  The monotonicity clause of \cref{lem:pushforward-tail} then
gives all later odd exponents.
\end{proof}

\begin{proposition}[Source-generated $B/C$ row-gated bridge]
\label{prop:bc-source-row-gated-bridge}
For $G=B_b$ or $C_b$, $2\le b\le30$, and every row-gated index
$b-1\le m\le29$, one has
\[
  D_G(2m+1)>0.
\]
All $870$ inequalities are regenerated from the Pieri rule and checked with
exact integer arithmetic.
\end{proposition}

\begin{proof}
Write
\[
 e_2^j=\sum_{\lambda\vdash2j}a_j(\lambda)s_\lambda,
 \qquad a_j(\lambda)\in\mathbb Z_{\ge0}.
 \tag{\ref{prop:bc-source-row-gated-bridge}.1}
\]
The orthogonal and symplectic Schur-integral criteria used in
\cref{prop:bc-pieri-correction-sign} give
\[
 m_j^{B_b}=\sum_{\substack{\lambda:\ \lambda_i\ {\rm even}\\
                         \ell(\lambda)\le2b+1}}a_j(\lambda),
 \qquad
 m_j^{C_b}=\sum_{\substack{\lambda:\ \lambda_i\ {\rm even}\\
                         \lambda_1\le2b}}a_j(\lambda).
 \tag{\ref{prop:bc-source-row-gated-bridge}.2}
\]
Indeed, the first identity is the finite odd-orthogonal criterion.  For the
second, apply the standard involution $\omega$: it sends $h_2^j$ to $e_2^j$
and $s_\alpha$ to $s_{\alpha'}$.  A partition $\alpha$ has its row lengths
in equal pairs precisely when every row length of $\alpha'$ is even, while
$\ell(\alpha)\le2b$ is equivalent to $(\alpha')_1\le2b$.  This proves
\eqref{prop:bc-source-row-gated-bridge}.2.

Pieri's rule gives the exact recurrence
\[
 a_{j+1}(\mu)=
 \sum_{\substack{\lambda:\ \mu/\lambda\text{ is a vertical two-strip}}}
 a_j(\lambda),
 \qquad a_0(\varnothing)=1.
 \tag{\ref{prop:bc-source-row-gated-bridge}.3}
\]
For a band with largest active rank $r$, the verifier retains exactly the
union
\[
 \ell(\lambda)\le2r+1\quad\hbox{or}\quad\lambda_1\le2r.
 \tag{\ref{prop:bc-source-row-gated-bridge}.4}
\]
A discarded partition has both height and width beyond these bounds.  Adding
boxes cannot decrease either quantity, so none of its descendants can enter
either sum in \eqref{prop:bc-source-row-gated-bridge}.2.  Thus this truncation
is exact.  A retained partition is stored in ordinary coordinates when its
height is bounded and in conjugate coordinates otherwise; in the latter
coordinates \eqref{prop:bc-source-row-gated-bridge}.3 becomes horizontal
two-strip Pieri.  This is a bijective change of representation, not a further
truncation.

The three simultaneous bands and their exact moment windows are
\[
\begin{array}{c|c|c}
\text{band}&b&\text{finite moments generated through}\\ \hline
\text{low}&2,3,4&61\\
\text{middle}&5,6,7,8&47,45,43,43\\
\text{high}&9,10,\ldots,19&41.
\end{array}
\tag{\ref{prop:bc-source-row-gated-bridge}.5}
\]
The recurrence is evaluated modulo the first five, four, or three primes,
respectively, from
\[
\begin{aligned}
 &(2305843009213694963,2305843009214695031,
   2305843009215694967,\\
 &\hspace{35mm}2305843009216695041,2305843009217694971).
\end{aligned}
\tag{\ref{prop:bc-source-row-gated-bridge}.6}
\]
The exact inversions in each CRT reconstruction verify that the moduli are
pairwise coprime; primality is not required.  The relevant products are
\begingroup\scriptsize
\[
\begin{aligned}
P_3&=12259964326943078111827907482335947634252293072155482851>s_{41},\\
P_4&=28269553036527760477584768395351886838659733871027006344612486973772241891>s_{47},\\
P_5&=65185151242986397660135212167169058667750503816445442673844289788983521680612515978066230161>s_{61}.
\end{aligned}
\tag{\ref{prop:bc-source-row-gated-bridge}.7}
\]
\endgroup
Each sum in \eqref{prop:bc-source-row-gated-bridge}.2 lies in $[0,s_j]$ by
nonnegativity and the stable Schur sum.  Hence its modular residues have a
unique representative in that interval, and the GMP Chinese-remainder
reconstruction is the exact finite moment.  The code substitutes the result
back into every prime and also verifies stabilization before $j=2b+2$.

Put $m_j^G=s_j+\delta_j^G$.  Outside the generated windows,
\cref{prop:bc-pieri-correction-sign} and nonnegativity give the proved box
\[
 -s_j\le\delta_j^G\le0.
 \tag{\ref{prop:bc-source-row-gated-bridge}.8}
\]
For each row the verifier expands, over $\mathbb Z$,
\[
 D_G(2m+1)=D_{\rm st}(2m+1)
 +\sum_jL_{m,j}\delta_j^G
 +\sum_{i\le j}Q_{m,i,j}\delta_i^G\delta_j^G.
 \tag{\ref{prop:bc-source-row-gated-bridge}.9}
\]
An exact correction is substituted.  For an interval correction, each linear
term is minimized at the appropriate endpoint and each quadratic term at all
four endpoint products; for a diagonal term the value at zero is included.
These are the extrema because every interval in
\eqref{prop:bc-source-row-gated-bridge}.8 is one-sided.  Thus the resulting
integer is a rigorous lower bound for \eqref{prop:bc-source-row-gated-bridge}.9.

There are
\[
 2\sum_{b=2}^{30}(31-b)=870
\]
row-gated checks.  The replay generates $832$ correction values, evaluates
$416$ steps exactly and $454$ by the stated interval substitution, and reports
no failure.  Its least lower bound is the exact value $1046$ at $B_2,m=1$.
The source and the deterministic machine-C output have SHA-256 hashes
\[
\begin{aligned}
&\hashsplit{fbd5565215a7b976b67a4529ad21daa2}{f235f2988da463c502370e3dab3ad797},\\
&\hashsplit{d019c183f7dedff1049199b31a50d681}{3e0c17a0c565df1e845119e2a9fba054}.
\end{aligned}
\tag{\ref{prop:bc-source-row-gated-bridge}.10}
\]
This proves every asserted strict inequality.
\end{proof}

\begin{proposition}[Correction-aware classical row-gated bridge]
\label{prop:classical-bridge}
The classical rows in the finite bridge through $m=29$ satisfy
\[
  D_G(2m+1)\ge0
\]
for every Chain step $m=0,\ldots,29$ whose target exponent $2m+3$ is
covered by the row gate.
\end{proposition}

\begin{proof}
The $B_b,C_b$ rows are exactly \cref{prop:bc-source-row-gated-bridge}.
For type $D$, the ranks $4\le b\le24$ are
\cref{prop:d-low-exact-bridge}, the ranks $25\le b\le52$ lie in the
correction-aware interval range proved in
\cref{prop:post29-d25-52-exact}, and the ranks $b\ge53$ are
\cref{prop:d-a-free-row-gated-bridge}.  These ranges exhaust every classical
row admitted by \cref{def:row-gate}; all inequalities are strict, hence in
particular nonnegative.
\end{proof}

\begin{proposition}[First-hit trace rows]\label{prop:first-hit}
The 57 top finite rows at the first post-bridge hit satisfy
\[
  \Qthree^G(n)\ge0
\]
for every later odd exponent.
\end{proposition}

\begin{proof}
At the first hit
\[
  N_{\mathrm{hit}}^{(29)}(G)=\max(63,s(G)+2),
\]
the $384$-bit directed-rounding C++/MPFR certificate checks the one-trace
lower bound against the negative correction terms for exactly 57 rows:
\[
  B:276,278,\ldots,295,\quad
  C:276,278,\ldots,295,\quad
  D:276,278,280,\ldots,295,297.
\]
Every interval lower endpoint is positive.  The replay output is:
\[
\begin{array}{ll}
\text{rows interval-certified} & 57,\\
\text{minimum row} & D_{281},\\
\text{minimum log-margin interval}
 & [0.6644444976107690011484081126,\\
 & \phantom{[}0.6644444976107690011484081126].
\end{array}
\]
The checker uses outward rounding for every rational operation,
logarithm, exponential, and logarithmic subtraction.  Its archived output
has SHA-256
\[
\hashsplit{ca3a938c9d45d65eb892f6ba637ed8c3}{a7b1d32a571a5be52e2f5efb80477176}.
\]
The monotonicity clause of \cref{lem:pushforward-tail} then carries the
inequality to all later odd exponents.
\end{proof}

\begin{proposition}[Direct post-$m=29$ low-rank tails]
\label{prop:post29-direct}
The following low-rank post-$m=29$ rows satisfy
\[
  \Qthree^G(n)\ge0
\]
for every odd $n\ge N_{\mathrm{hit}}^{(29)}(G)$:
\[
\begin{array}{c|c}
\text{family} & \text{ranks}\\
\hline
B & 2\le b\le13\\
C & 2\le b\le19\\
D & 4\le b\le24 .
\end{array}
\]
\end{proposition}

\begin{proof}
The exact box integrals for $B_2,C_2,B_3,C_3$ are replayed by the flags
\[
\begin{array}{c}
\texttt{--b2-c2-post-m29-tail-certificate},\\
\texttt{--b3-c3-post-m29-tail-certificate}.
\end{array}
\]
For $B_4,B_5,B_6,B_8$, the
fixed rational cap replays apply
\cref{lem:b-weyl-cap-pushforward}; in particular they include the factor
$2^{-b}$ absent from an unnormalized non-simply-laced Macdonald--Mehta
product.  Their cap data are
\[
\begin{array}{c|c|c|c|c}
G&R&x_0&\text{independent-variable event}&c\\ \hline
B_4&9&27&Y\le1&23\\
B_5&9&46&Y\le1&41\\
B_6&1053/50&2847/50&|Y|\le3/2&1386/25\\
B_8&1609/25&1791/25&|Y|\le3/2&3507/50.
\end{array}
\]
The corrected exact-rational margins are positive in all four rows.

The rows
\[
B_7,B_9,B_{10},B_{11},B_{12},B_{13},\qquad C_4,\ldots,C_{19}
\]
use the exact determinant tails of
\cref{prop:post29-low-bc-exact-mgf}.  Since $b\le19$ here,
\cref{def:classical-constants} gives
$N_{\mathrm{hit}}^{(29)}(G)=63$.  For $B_{11},B_{12},B_{13}$ and
$C_{13},C_{14},C_{17}$, the first-hit certificate
\cref{prop:post29} supplies $\Qthree^G(63)\ge0$ and
\cref{prop:post29-low-bc-exact-bridges} supplies every odd exponent before
the exact-MGF onset.  For $C_{16}$, the same first-hit certificate supplies
the exponent $63$ and the exact-MGF tail starts at $65$.  Every other row in
the display has exact-MGF onset $63$.  Thus these certificates cover every
odd exponent at least $63$ in all $B/C$ rows listed in the proposition.

For $D_4,\ldots,D_9$, the fixed rational ball-cap and product-sine replays
\[
\begin{gathered}
\texttt{--d4-post-m29-tail-certificate},\qquad
\texttt{--d5-post-m29-tail-certificate},\\
\texttt{--d6-b6-post-m29-tail-certificate},\qquad
\texttt{--d7-b7-c7-post-m29-tail-certificate},\\
\texttt{--d8-post-m29-tail-certificate},\qquad
\texttt{--b8-c9-d9-post-m29-tail-certificate}
\end{gathered}
\]
instantiate \cref{lem:weyl-cap-pushforward}.  Their type-$D$ cap and
central-variable data are
\[
\begin{array}{c|c|c|c}
b&x_0&\text{independent-variable event}&c\\ \hline
4&26&Y\le1&22\\
5&41&Y\le1&38\\
6&949/20&|Y|\le7/5&921/20\\
7&73&|Y|\le8/5&357/5\\
8&84&|Y|\le8/5&412/5\\
9&767/10&|Y|\le8/5&751/10.
\end{array}
\]
For $D_4,D_5$, the replay uses the per-root bound
$\bigl(\sin(\alpha/2)/(\alpha/2)\bigr)^2\ge1/4$.  For $D_6,D_7$ it uses
\[
 \prod_{\alpha>0}
 \left(\frac{\sin(\alpha/2)}{\alpha/2}\right)^2
 \ge\left(1-\frac{R}{24}\right)^2,
\]
which follows from $\sin z/z\ge1-z^2/6$ and
$\prod_i(1-u_i)\ge1-\sum_i u_i$.  For $D_8$, the Euler product and
$\sum_{\alpha>0}(\alpha/2)^2\le9$ give the lower bound
$(\sin3/3)^2$; the replay inserts the rational inequality
$\sin3\ge15/106-(15/106)^3/6$.  For $D_9$, put
$L=|\alpha|^2=2$.  Concavity of $\log(1-u)$ and
$\alpha(\theta)^2\le LR/\kappa$ give
\[
 \prod_{\alpha>0}
 \left(1-\frac{\alpha(\theta)^2}{24}\right)^2
 \ge\left(1-\frac{LR}{24\kappa}\right)^{2\kappa/L}.
\]
The exact rational replays substitute these bounds in
\eqref{lem:weyl-cap-pushforward}.2--\eqref{lem:weyl-cap-pushforward}.4 and
prove the pushforward-tail inequality from odd exponent $63$.  For $D_{10}$
and $D_{12},\ldots,D_{24}$, the exact finite-rank certificates in
\cref{prop:post29-d10-24-exact-tail} also start at odd exponent $63$.
The row $D_{11}$ instead has an exact correction-aware Chain bridge through
odd exponent $75$ and an exact-MGF tail from exponent $75$ onward.  Finally,
\cref{prop:d-low-exact-bridge} gives every row-gated Chain difference through
$m=29$, so the stable edge and \cref{lem:chain-propagation} reach odd exponent
$61$ in every row.  These three cases cover every odd exponent at least $63$.
\end{proof}

\begin{proposition}[Post-$m=29$ high-edge Chernoff trace tails]\label{prop:post29-chernoff}
For $G=B_b$ or $G=C_b$ with $218\le b\le275$, or $b=277$, one has
\[
  \Qthree^G(n)\ge0
\]
for every odd $n\ge N_{\mathrm{hit}}^{(29)}(G)$.
\end{proposition}

\begin{proof}
This is Proposition 24.17E in the proof ledger.  The active replay uses
$384$-bit directed rounding in
\path{character_ring_iter/trace_square_cutoff_mpfr.cpp}.  Its
source-supported CJL/Rains Chernoff estimate uses
\[
(r,\eta)\in
\left\{
(2001/1000,9/20),
(10001/5000,56/125),
(99/50,56/125)
\right\}.
\]
It checks the one-trace source range, the
Chernoff admissibility window, and the trace-pushforward inequality of
\cref{lem:pushforward-tail} at $N_{\mathrm{hit}}^{(29)}(G)$, with outward
rounding for every rational operation, logarithm, exponential, and logarithmic
subtraction.  The local log
\[
\texttt{certificates/classical\_trace/first\_hit\_trace\_replay.log}
\]
exits with status $0$ and has SHA-256 hash
\[
\hashsplit{ca3a938c9d45d65eb892f6ba637ed8c3}{a7b1d32a571a5be52e2f5efb80477176}.
\]
All $118$ interval lower endpoints are positive.  The smallest is attained
at $B_{219}$ and $C_{219}$, and is enclosed by
\[
[0.3074299695438132583903429588,
  0.3074299695438132583903429588].
\]
\end{proof}

\begin{lemma}[Type $B$ unitary square-trace tail]\label{lem:b-unitary-square-tail}
Let $O$ be Haar-distributed in the defining representation of $B_b$, and put
\(\tau_B(O)=\operatorname{Tr}(O^2)\).  For every real $a>1$,
\[
  \Pr\{\tau_B(O)\ge a\}
  \le 4\exp\!\left(-\frac{(a-1)^2}{4}\right).
\]
\end{lemma}

\begin{proof}
Use the standard representative $O\in\operatorname{SO}(2b+1)$.  Rains
\cite[Theorem~1.5, equation~(23)]{RainsPowerMaps}, specialized to power
$p=2$, identifies the nonfixed eigenvalues of $O^2$ with those of
$\operatorname{Re}U(b)$.  Thus, if $V$ is Haar-distributed in $U(b)$,
restoring the fixed
$+1$ eigenvalue gives the distributional identity
\[
  \tau_B(O)\overset{d}=1+2\operatorname{Re}\operatorname{Tr}V.
\]
Courteaut--Johansson--Lambert
\cite[Lemma~2.2, with $m=1$]{CJL2022} gives, for every $L>0$,
\[
 \Pr\!\left\{
   (\operatorname{Re}\operatorname{Tr}V,
    \operatorname{Im}\operatorname{Tr}V)
   \notin[-L/2,L/2]^2
 \right\}
 <4e^{-L^2/4}.
\]
For every $0<L<a-1$, the event
$\{2\operatorname{Re}\operatorname{Tr}V\ge a-1\}$ is contained in the
displayed complement of the square.  Hence its probability is at most
$4e^{-L^2/4}$.  Letting $L\uparrow a-1$ proves the claim.
\end{proof}

\begin{proposition}[Post-$m=29$ middle $B$ tails from the unitary square
trace]\label{prop:post29-b-unitary-square}
For every $124\le b\le217$ and every odd $n\ge2b+1$,
\[
  \Qthree^{B_b}(n)\ge0.
\]
\end{proposition}

\begin{proof}
Set
\[
 r=\frac{2001}{1000},\qquad
 \eta=\frac7{20},\qquad
 c=rb,\qquad a=\eta b,
 \qquad L=2r+3\eta=\frac{1263}{250}.
\]
Then $1<a<2b<c$.  In the notation of \cref{lem:trace-tail}, put
\[
 p_1=\Pr\{|T_1(O)|\ge\sqrt{Lb}\},\qquad
 p_2=\Pr\{\tau_B(O)\ge\eta b\}.
\]
The defining trace has mean zero.  Courteaut--Johansson--Lambert
\cite[Theorem~1.4]{CJL2022}, which applies because $b\ge124$, and the
two-sided Mills lower bound with $\pi<4$ give
\[
 p_1\ge G_b-E_b,
\]
where
\[
 G_b=\frac{e^{-Lb/2}\sqrt{Lb}}{\sqrt2(Lb+1)},
 \qquad
 E_b=\frac{5(\log(2b))^{1/4}}{\sqrt{(2b)!}}.
\]
By \cref{lem:b-unitary-square-tail},
\[
 p_2\le U_b:=4\exp\!\left(-\frac{(\eta b-1)^2}{4}\right).
\]

The directed-rounding MPFR replay
\[
\texttt{certificates/post\_m29/
post\_m29\_b\_unitary\_square\_trace\_mpfr\_certificate.log}
\]
checks, for every integer $124\le b\le217$, the four inequalities
\[
 E_b\le\frac14G_b,\qquad U_b\le\frac14G_b,
\]
\[
 (rb)^2\bigl(1-(\eta b)^{-2}\bigr)\frac{G_b}{2}
 \ge
 2\left(2(b^2+1)U_b(2/r)^{2b+1}\right),
\]
and
\[
 (rb)^2\bigl(1-(\eta b)^{-2}\bigr)\frac{G_b}{2}
 \ge
 2\left(2(\eta/r)^{2b+1}\right).
\]
It uses $384$-bit MPFR intervals, rounds every lower bound downward and every
upper bound upward, exits with status $0$, and reports
\texttt{rows\_checked=94} and \texttt{failures=0}.  The four smallest
natural-log margins are all attained at $B_{124}$ and are respectively
\[
\begin{aligned}
 &241.28842853377649\ldots,\qquad
  129.87520977348063\ldots,\\
 &130.69322950465274\ldots,\qquad
  126.27950899088545\ldots.
\end{aligned}
\]
The log has SHA-256
\[
\hashsplit{7780bfe3253bf9e030a9a544dcd7327f}{cd905aed8a2890ded191dc38efae8230},
\]
and the C++ source has SHA-256
\[
\hashsplit{d0d503ffc4bd6cb0485a97a4cef2178a}{930f3f00c38ffd4486cb66aa01652aa7}.
\]

The first two inequalities imply $p_1-p_2\ge G_b/2$.  Hence
\cref{lem:trace-tail} gives
\[
 P_{B_b}(c)
 \ge (rb)^2\bigl(1-(\eta b)^{-2}\bigr)\frac{G_b}{2},
 \qquad
 N_{B_b}(a)\le2(b^2+1)U_b.
\]
The final two checked inequalities imply
\[
 P_{B_b}(c)
 \ge
 N_{B_b}(a)(2/r)^{2b+1}+2(\eta/r)^{2b+1}.
\]
Thus \cref{lem:pushforward-tail} applies at the odd exponent $2b+1$ and,
by its monotonicity clause, at every later odd exponent.
\end{proof}

\begin{definition}[CJL Fredholm remainder]\label{def:cjl-fredholm-remainder}
For an integer $d\ge1$ and $0\le v<d+1$, set
\[
 \mathfrak b_d(v)=
 \frac{\exp\!\left(v^2/(d+1)\right)}
      {1-v^2/(d+1)^2}\,
 \frac{v^d}{d!}.
\]
\end{definition}

\begin{definition}[CJL trace-norm envelope]
\label{def:cjl-trace-norm-envelope}
For an integer $d\ge1$ and $0\le v<d+1$, set
\[
 \mathfrak a_d(v)=
 \frac{\exp(v^2/(d+1))v^d}{d!}
 \frac1{1-v/(d+1)}
 \frac1{\sqrt{1-v^2/(d+1)^2}}.
\]
\end{definition}

\begin{lemma}[Extended small-Fredholm trace MGF bounds]
\label{lem:extended-cjl-trace-mgf}
Let $Z$ be a centered full defining trace in one of the four orthogonal
Jacobi components of Courteaut--Johansson--Lambert, or in the corresponding
symplectic component, and let $d$ be the total matrix dimension attached to
that component.  If
\[
 0\le v<d+1,\qquad \mathfrak a_d(v)<1,
\]
then $\mathfrak b_d(v)<1$ and
\[
 e^{v^2/2}(1-\mathfrak b_d(v))
 \le \mathbf E e^{\pm vZ}
 \le \frac{e^{v^2/2}}{1-\mathfrak b_d(v)}.
 \tag{\ref{lem:extended-cjl-trace-mgf}.1}
\]
\end{lemma}

\begin{proof}
Courteaut--Johansson--Lambert
\cite[Proposition~3.3]{CJL2022} writes the centered transform as
\[
 \mathbf E e^{zZ}=e^{z^2/2}\det(I+QK_zQ).
\]
The proof of their Lemma~3.6 starts from the trace-norm inequality in their
equation~(3.4), inserts their Bessel bound (2.16), and sums two geometric
series.  If the least Bessel index is denoted by the total dimension $d$,
the calculation before their equation~(3.5) gives, without using the later
sufficient restriction $|z|<2e^{-5/4}n$,
\[
 \|QK_zQ\|_{\mathcal J_1}
 \le
 \frac{\exp(|z|^2/(d+1))|z|^d}{d!}
 \frac1{1-|z|/(d+1)}
 \frac1{\sqrt{1-|z|^2/(d+1)^2}}
 =\mathfrak a_d(|z|)
\]
whenever $|z|<d+1$.  The four kernels listed in Proposition~3.3 have least
indices $2n,2n+1,2n+1,2n+2$, respectively, exactly their corresponding
total dimensions, so the same displayed calculation applies to every case.

The proof of their Lemma~3.7, in particular equation~(3.6), also gives
directly
\[
 \sum_{j\ge1}\frac1j
 \left|\operatorname{Tr}(QK_zQ)^j\right|
 \le-\log(1-\mathfrak b_d(|z|)).
\]
Here $\mathfrak b_d(v)<\mathfrak a_d(v)<1$.  Thus the trace-norm hypothesis
justifies Plemelj's series, and the last display bounds the absolute value of
the logarithm of the determinant.  For real $z=\pm v$, the determinant is
positive because it equals $e^{-v^2/2}\mathbf E e^{\pm vZ}$.  Exponentiation
proves \eqref{lem:extended-cjl-trace-mgf}.1.
\end{proof}

\begin{lemma}[Layered exponential-tilt lower tail]
\label{lem:layered-trace-tilt}
Let $Z$ satisfy \eqref{lem:extended-cjl-trace-mgf}.1 at every argument used
below, with the same dimension $d$.  Fix $s,l>0$, put $t=s-l$, and choose
$\lambda_->0$ and pairs $(u_i,\lambda_i)$, $i=0,\ldots,L-1$, such that
\[
 0<u_0<\cdots<u_{L-1}.
\]
Write $D(v)=1-\mathfrak b_d(v)$ and define
\[
 A_-=
 \frac{\exp(\lambda_-^2/2-\lambda_-l)}
      {D(s)D(s-\lambda_-)},
 \qquad
 A_i^+=
 \frac{\exp(\lambda_i^2/2-\lambda_i u_i)}
      {D(s)D(s+\lambda_i)},
\]
\[
 g_i=1-A_--A_i^+.
\]
Suppose
\[
 0<g_0\le g_1\le\cdots\le g_{L-1}.
\]
Then, with $w_i=e^{-su_i}$,
\[
 \Pr\{Z\ge t\}\ge V_d(s,l):=
 D(s)e^{-s^2/2}
 \left(
   w_{L-1}g_{L-1}
   +\sum_{i=0}^{L-2}(w_i-w_{i+1})g_i
 \right).
 \tag{\ref{lem:layered-trace-tilt}.1}
\]
If the same MGF bounds hold after replacing $Z$ by $-Z$, then
\[
 \Pr\{|Z|\ge t\}\ge2V_d(s,l).
 \tag{\ref{lem:layered-trace-tilt}.2}
\]
\end{lemma}

\begin{proof}
Tilt the law by $d\mathbf P_s=M(s)^{-1}e^{sZ}\,d\mathbf P$, where
$M(s)=\mathbf E e^{sZ}$.  Exponential Markov at the independently chosen
parameters $\lambda_-$ and $\lambda_i$, followed by
\cref{lem:extended-cjl-trace-mgf}, gives
\[
 \mathbf P_s\{Z-s\le-l\}\le A_-,
 \qquad
 \mathbf P_s\{Z-s\ge u_i\}\le A_i^+.
\]
Consequently
\[
 \mathbf P_s\{t<Z\le s+u_i\}\ge g_i.
\]
For a decreasing function, the least integral compatible with these nested
cumulative-mass lower bounds places each required mass increment at the
right endpoint.  Hence
\[
 \mathbf E_s[e^{-s(Z-s)}\mathbf1_{\{Z>t\}}]
 \ge
 w_{L-1}g_{L-1}
 +\sum_{i=0}^{L-2}(w_i-w_{i+1})g_i.
\]
Multiplying by $M(s)e^{-s^2}$ and using
$M(s)\ge e^{s^2/2}D(s)$ proves
\eqref{lem:layered-trace-tilt}.1.  Applying the same argument to $-Z$ and
adding the two disjoint tails proves \eqref{lem:layered-trace-tilt}.2.
\end{proof}

\begin{lemma}[Arbitrary-parameter $B/C/D$ square-trace MGF bounds]
\label{lem:bc-square-arbitrary-mgf}
For $B_b$, let $\tau_B=\operatorname{Tr}(O^2)$ and, for $x>1$, $v>0$, put
\[
 U_B(x,v)=\exp\bigl(-v(x-1)+v^2\bigr).
\]
Then
\[
 e^{-vx}\mathbf E e^{v\tau_B}\le U_B(x,v).
 \tag{\ref{lem:bc-square-arbitrary-mgf}.1}
\]
For $C_b$, use the integers $d_0,d_1$ of
\cref{lem:c-square-mgf-tail} and suppose
$\mathfrak a_{d_i}(v)<1$ for $i=0,1$.  Put
\[
 U_C(x,v)=
 \frac{\exp(-v(x-1)+v^2)}
 {(1-\mathfrak b_{d_0}(v))(1-\mathfrak b_{d_1}(v))}.
\]
Then
\[
 e^{-vx}\mathbf E e^{v\tau_C}\le U_C(x,v).
 \tag{\ref{lem:bc-square-arbitrary-mgf}.2}
\]
For $D_b$, put
\[
 m_0=\left\lceil\frac b2\right\rceil,
 \qquad m_1=\left\lceil\frac{b-1}{2}\right\rceil,
 \qquad d_0=2m_0,
 \qquad d_1=2m_1+1,
\]
and suppose $\mathfrak a_{d_i}(v)<1$ for $i=0,1$.  With
$\tau_D=\operatorname{Tr}(O^2)$, put
\[
 U_D(x,v)=
 \frac{\exp(-v(x-1)+v^2)}
 {(1-\mathfrak b_{d_0}(v))(1-\mathfrak b_{d_1}(v))}.
\]
Then
\[
 e^{-vx}\mathbf E e^{v\tau_D}\le U_D(x,v).
 \tag{\ref{lem:bc-square-arbitrary-mgf}.3}
\]
\end{lemma}

\begin{proof}
For $B_b$, the Rains identity used in
\cref{lem:b-unitary-square-tail} gives
$\tau_B\overset d=1+2\operatorname{Re}\operatorname{Tr}V$ with
$V$ Haar in $U(b)$.  Courteaut--Johansson--Lambert
\cite[Lemma~2.12]{CJL2022}, applied to the real function
$f(\theta)=2v\cos\theta$, gives
\[
 \mathbf E e^{2v\operatorname{Re}\operatorname{Tr}V}\le e^{v^2}.
\]
This proves \eqref{lem:bc-square-arbitrary-mgf}.1.  For $C_b$, the independent
Rains components $Y_0,Y_1$ in the proof of
\cref{lem:c-square-mgf-tail} satisfy
$\tau_C-1\overset d=-(Y_0+Y_1)$.  Apply the upper bound in
\cref{lem:extended-cjl-trace-mgf} to each centered component at $-v$, multiply
the two bounds, and include the deterministic $+1$.  This gives
\eqref{lem:bc-square-arbitrary-mgf}.2.

For $D_b$, Rains' power-two decomposition
\cite[Theorem~1.5, equation~(22)]{RainsPowerMaps}, with the fixed
eigenvalues restored as in the proof of \cref{lem:rains-square}, gives
independent centered full orthogonal traces $Y_0,Y_1$ of total dimensions
$d_0,d_1$ such that
\[
 \tau_D-1\overset d=Y_0+Y_1.
\]
Apply the upper bound of \cref{lem:extended-cjl-trace-mgf} at $v$ to both
components, multiply, and include the deterministic $+1$.  This proves
\eqref{lem:bc-square-arbitrary-mgf}.3.
\end{proof}

\begin{lemma}[Type $B$ defining-trace MGF sandwich]\label{lem:b-trace-mgf}
Let $O$ be Haar-distributed in $\operatorname{SO}(2b+1)$, put
$T_b=\operatorname{Tr}O$, and, for $v\ge0$, define
\[
 \beta_b(v)=\mathfrak b_{2b+1}(v).
\]
If $0\le v<2e^{-5/4}b$, then $\beta_b(v)<1$ and
\[
 e^{v^2/2}\bigl(1-\beta_b(v)\bigr)
 \le \mathbf E e^{vT_b}
 \le \frac{e^{v^2/2}}{1-\beta_b(v)}.
\]
\end{lemma}

\begin{proof}
In the notation of Courteaut--Johansson--Lambert, the $b$ nonfixed
eigenvalue pairs of $O^+(2b+1)$ have Jacobi parameters
$(1/2,-1/2)$.  Their trace variable in Section~3 omits the deterministic
$+1$ eigenvalue and has mean $-1$.  Thus the linear factor in
\cite[Proposition~3.3]{CJL2022} cancels when that eigenvalue is restored.
Specializing the proposition at $\xi=-iv$ gives
\[
 \mathbf E e^{vT_b}
 =e^{v^2/2}\det\!\left(I+Q_bK^{+-}_{v}Q_b\right).
\]
For $v<2e^{-5/4}b$, Lemmas~3.6--3.7 of the same source, with total matrix
size $d=2b+1$, give
\[
 \left|\log\det\!\left(I+Q_bK^{+-}_{v}Q_b\right)\right|
 \le -\log\bigl(1-\beta_b(v)\bigr).
\]
The determinant is positive because it is
$e^{-v^2/2}\mathbf E e^{vT_b}$.  Exponentiating the last display proves both
bounds.
\end{proof}

\begin{lemma}[Exponentially tilted type $B$ defining-trace tail]
\label{lem:b-trace-tilt}
In the setting of \cref{lem:b-trace-mgf}, let $t,h>0$, put $s=t+h$, and
suppose $t+2h<2e^{-5/4}b$.  Define
\[
 \rho_b(t,h)=1-
 \frac{e^{-h^2/2}}{1-\beta_b(s)}
 \left(
   \frac1{1-\beta_b(t)}+
   \frac1{1-\beta_b(t+2h)}
 \right).
\]
If $\rho_b(t,h)>0$, then
\[
 \Pr\{T_b\ge t\}
 \ge
 V_b(t,h):=
 \bigl(1-\beta_b(s)\bigr)\rho_b(t,h)
 \exp\!\left(-\frac{s^2}{2}-sh\right).
\]
\end{lemma}

\begin{proof}
Write $M(v)=\mathbf E e^{vT_b}$ and tilt Haar probability by
\[
 d\mathbf P_s=M(s)^{-1}e^{sT_b}\,d\mathbf P.
\]
Exponential Markov and \cref{lem:b-trace-mgf} give
\[
 \mathbf P_s\{T_b-s\ge h\}
 \le
 \frac{e^{-h^2/2}}
 {(1-\beta_b(s))(1-\beta_b(s+h))}
\]
and
\[
 \mathbf P_s\{T_b-s\le-h\}
 \le
 \frac{e^{-h^2/2}}
 {(1-\beta_b(s))(1-\beta_b(s-h))}.
\]
Consequently
$\mathbf P_s\{s-h<T_b<s+h\}\ge\rho_b(t,h)$.  Since $s-h=t$,
\begin{align*}
 \Pr\{T_b\ge t\}
 &\ge M(s)\,
   \mathbf E_s\!\left[e^{-sT_b}
   \mathbf 1_{\{s-h<T_b<s+h\}}\right]\\
 &\ge M(s)e^{-s(s+h)}\rho_b(t,h)\\
 &\ge
 (1-\beta_b(s))\rho_b(t,h)e^{-s^2/2-sh},
\end{align*}
as claimed.
\end{proof}

\begin{proposition}[Post-$m=29$ middle $B$ tails from an MGF tilt]
\label{prop:post29-b-tilted-mgf}
For every $62\le b\le123$ and every odd $n\ge2b+29$,
\[
 \Qthree^{B_b}(n)\ge0.
\]
\end{proposition}

\begin{proof}
Set
\[
 r=\frac{20001}{10000},\qquad
 \alpha=\frac{979}{250},\qquad
 h=\frac{601}{500},\qquad
 c_b=rb,\qquad a_b=\alpha\sqrt b,
\]
and put
\[
 t_b=\sqrt{2c_b+3a_b},\qquad s_b=t_b+h,
 \qquad V_b=V_b(t_b,h),
\]
where the final quantity is defined in \cref{lem:b-trace-tilt}.  One has
$1<a_b<2b<c_b$.  In the notation of \cref{lem:trace-tail},
\[
 p_1\ge\Pr\{T_b\ge t_b\}\ge V_b,
 \qquad
 p_2\le U_b:=4\exp\!\left(-\frac{(a_b-1)^2}{4}\right),
\]
where the second inequality is \cref{lem:b-unitary-square-tail}.

The $384$-bit directed-rounding MPFR replay
\[
\texttt{certificates/post\_m29/
post\_m29\_bc\_tilted\_mgf\_mpfr\_certificate.log}
\]
checks, for every integer $62\le b\le123$, that
\[
 t_b+2h<2e^{-5/4}b,\qquad \rho_b(t_b,h)>0,
 \qquad V_b\ge2U_b,
\]
and that
\[
 c_b^2(1-a_b^{-2})\frac{V_b}{2}
 \ge
 2\left(2(b^2+1)U_b(2/r)^{2b+29}\right),
\]
\[
 c_b^2(1-a_b^{-2})\frac{V_b}{2}
 \ge
 2\left(2(a_b/c_b)^{2b+29}\right).
\]
Every lower quantity is rounded downward and every upper quantity upward.
The combined $B/C$ log exits with status $0$, reports
\texttt{B\_rows\_checked=62} and \texttt{failures=0}, and places all five
$B$-side minima at $B_{62}$.  The minimum
domain and tilted-mass lower bounds are respectively
\[
 14.669512561242725\ldots,\qquad
 0.028831039797142247\ldots,
\]
and the three minimum natural-log comparison margins are
\[
 0.11377835996790686\ldots,\qquad
 0.12021572562760810\ldots,\qquad
 0.15645659337912115\ldots.
\]
The log has SHA-256
\[
\hashsplit{ca6860204466e4a6b38d9c5597511cd1}{1cfebb4b1b65cab88db6ee45f9477111},
\]
and the C++ source has SHA-256
\[
\hashsplit{0f5ad1037bcf738ab4ddd9c5b12f151a}{adc1bff2b5c4726c60622f608ef9e0e7}.
\]

The first comparison gives $p_1-p_2\ge V_b/2$.  Hence
\cref{lem:trace-tail} gives
\[
 P_{B_b}(c_b)\ge c_b^2(1-a_b^{-2})\frac{V_b}{2},
 \qquad N_{B_b}(a_b)\le2(b^2+1)U_b.
\]
The final two checked comparisons imply the hypothesis of
\cref{lem:pushforward-tail} at the odd exponent $2b+29$.  Its monotonicity
clause proves every later odd exponent.
\end{proof}

\begin{lemma}[Type $C$ defining-trace MGF sandwich]
\label{lem:c-trace-mgf}
Let $U$ be Haar-distributed in $\operatorname{Sp}(2b)$, put
$S_b=\operatorname{Tr}U$, and define
\[
 \gamma_b(v)=\mathfrak b_{2b}(v).
\]
If $0\le v<2e^{-5/4}b$, then $\gamma_b(v)<1$ and
\[
 e^{v^2/2}\bigl(1-\gamma_b(v)\bigr)
 \le \mathbf E e^{vS_b}
 \le \frac{e^{v^2/2}}{1-\gamma_b(v)}.
\]
\end{lemma}

\begin{proof}
For $\operatorname{Sp}(2b)$, the Jacobi parameters in
\cite[Section~3, equation~(3.1)]{CJL2022} are
$(1/2,1/2)$, $b$ nontrivial conjugate pairs, and total matrix dimension
$d=2b$.  There are no deterministic eigenvalues and the linear factor in
\cite[Proposition~3.3]{CJL2022} vanishes.  Substitution of $\xi=-iv$ in
that proposition therefore gives
\[
 \mathbf E e^{vS_b}
 =e^{v^2/2}\det\!\left(I+Q_bK^{++}_{v}Q_b\right).
\]
For $v<2e^{-5/4}b$, Lemmas~3.6--3.7 of the same source, with $d=2b$,
give
\[
 \left|\log\det\!\left(I+Q_bK^{++}_{v}Q_b\right)\right|
 \le-\log\bigl(1-\mathfrak b_{2b}(v)\bigr).
\]
The determinant is positive because it equals
$e^{-v^2/2}\mathbf E e^{vS_b}$.  Exponentiation proves both bounds.
\end{proof}

\begin{lemma}[Exponentially tilted type $C$ defining-trace tail]
\label{lem:c-trace-tilt}
In the setting of \cref{lem:c-trace-mgf}, let $t,h>0$, put $s=t+h$, and
suppose $t+2h<2e^{-5/4}b$.  Define
\[
 \rho_b^C(t,h)=1-
 \frac{e^{-h^2/2}}{1-\gamma_b(s)}
 \left(
   \frac1{1-\gamma_b(t)}+
   \frac1{1-\gamma_b(t+2h)}
 \right).
\]
If $\rho_b^C(t,h)>0$, then
\[
 \Pr\{S_b\ge t\}
 \ge W_b(t,h):=
 \bigl(1-\gamma_b(s)\bigr)\rho_b^C(t,h)
 \exp\!\left(-\frac{s^2}{2}-sh\right).
\]
\end{lemma}

\begin{proof}
The proof of \cref{lem:b-trace-tilt} uses only the two displayed MGF
bounds in \cref{lem:b-trace-mgf}.  Apply its exponential tilt at $s=t+h$
with $S_b$ in place of $T_b$ and with $\gamma_b$ in place of $\beta_b$.
The two exponential-Markov estimates become
\[
 \mathbf P_s\{S_b-s\ge h\}
 \le
 \frac{e^{-h^2/2}}
 {(1-\gamma_b(s))(1-\gamma_b(s+h))},
\]
\[
 \mathbf P_s\{S_b-s\le-h\}
 \le
 \frac{e^{-h^2/2}}
 {(1-\gamma_b(s))(1-\gamma_b(s-h))}.
\]
Their sum leaves interval mass at least $\rho_b^C(t,h)$.  On
$s-h<S_b<s+h$, one has $e^{-sS_b}\ge e^{-s(s+h)}$; the lower MGF bound
of \cref{lem:c-trace-mgf} then gives exactly the displayed $W_b(t,h)$.
\end{proof}

\begin{lemma}[Type $C$ square-trace MGF tail]
\label{lem:c-square-mgf-tail}
Let $U$ be Haar-distributed in $\operatorname{Sp}(2b)$, where $b\ge2$,
and put $\tau_C(U)=-\operatorname{Tr}(U^2)$.  Define
\[
 m_0=\left\lceil\frac{b+1}{2}\right\rceil,
 \quad m_1=\left\lceil\frac b2\right\rceil,
 \quad n_0=m_0-1,
 \quad n_1=m_1,
\]
\[
 d_0=2m_0,\qquad d_1=2m_1+1.
\]
For $a>1$, put $v=(a-1)/2$.  If
\[
 v<2e^{-5/4}\min(n_0,n_1),
\]
then $\mathfrak b_{d_i}(v)<1$ for $i=0,1$ and
\[
 \Pr\{\tau_C(U)\ge a\}
 \le
 \frac{\exp\!\left(-(a-1)^2/4\right)}
 {(1-\mathfrak b_{d_0}(v))(1-\mathfrak b_{d_1}(v))}.
\]
\end{lemma}

\begin{proof}
Rains \cite[Theorem~1.5 and the subsequent remark, equation~(36)]{RainsPowerMaps} identifies the
nonfixed eigenvalues of $\operatorname{Sp}(2b)$ with those of
$O^-(2b+2)$.  The omitted eigenvalues are $+1,-1$, so after squaring their
trace contribution is $2$.  Equation~(22) of the same theorem at $p=2$
decomposes the remaining eigenvalue distribution into independent
components
\[
 O^-(2m_0)\oplus O^+(2m_1+1).
\]
Let $Y_0,Y_1$ be the full defining traces of these two components.  The
first component has fixed eigenvalues $+1,-1$, of total trace $0$; the
second has a fixed $+1$.  Restoring these fixed eigenvalues in Rains'
identity therefore gives
\[
 \operatorname{Tr}(U^2)\overset d=Y_0+Y_1-1,
 \qquad
 \tau_C(U)-1\overset d=-(Y_0+Y_1).
\]

By Courteaut--Johansson--Lambert
\cite[Section~3, equation~(3.1)]{CJL2022}, the first component has Jacobi
parameters $(1/2,1/2)$ and $n_0$ nontrivial pairs, while the second has
parameters $(1/2,-1/2)$ and $n_1$ nontrivial pairs.  Proposition~3.3 shows
that the nonfixed trace means are respectively $0$ and $-1$; after the
fixed traces $0$ and $+1$ are restored, both $Y_0,Y_1$ are centered.
Lemmas~3.6--3.7, at the argument $-v$, now give
\[
 \mathbf E e^{-vY_i}
 \le\frac{e^{v^2/2}}{1-\mathfrak b_{d_i}(v)}
 \qquad(i=0,1).
\]
Independence and exponential Markov yield
\[
 \Pr\{-(Y_0+Y_1)\ge a-1\}
 \le
 \frac{e^{-v(a-1)+v^2}}
 {(1-\mathfrak b_{d_0}(v))(1-\mathfrak b_{d_1}(v))}.
\]
Since $v=(a-1)/2$, the exponent is $-(a-1)^2/4$.
\end{proof}

\begin{proposition}[Post-$m=29$ middle $C$ tails from MGF tilts]
\label{prop:post29-c-tilted-mgf}
For every $62\le b\le217$ and every odd $n\ge2b+29$,
\[
 \Qthree^{C_b}(n)\ge0.
\]
\end{proposition}

\begin{proof}
Set
\[
 r=\frac{20001}{10000},\qquad
 \alpha=\frac{979}{250},\qquad
 h=\frac{601}{500},\qquad
 c_b=rb,\qquad a_b=\alpha\sqrt b,
\]
and put
\[
 t_b=\sqrt{2c_b+3a_b},\qquad s_b=t_b+h,
 \qquad W_b=W_b(t_b,h),
\]
where the final quantity is defined in \cref{lem:c-trace-tilt}.  For the
square trace, use the integers $m_i,n_i,d_i$ of
\cref{lem:c-square-mgf-tail}, put $v_b=(a_b-1)/2$, and define
\[
 U_b^C=
 \frac{\exp\!\left(-(a_b-1)^2/4\right)}
 {(1-\mathfrak b_{d_0}(v_b))(1-\mathfrak b_{d_1}(v_b))}.
\]
One has $1<a_b<2b<c_b$.  In the notation of \cref{lem:trace-tail},
\[
 p_1\ge\Pr\{S_b\ge t_b\}\ge W_b,
 \qquad p_2\le U_b^C,
\]
by \cref{lem:c-trace-tilt,lem:c-square-mgf-tail}.

The $384$-bit directed-rounding MPFR replay
\[
\texttt{certificates/post\_m29/
post\_m29\_bc\_tilted\_mgf\_mpfr\_certificate.log}
\]
checks, for every integer $62\le b\le217$, that
\[
 t_b+2h<2e^{-5/4}b,\qquad
 v_b<2e^{-5/4}\min(n_0,n_1),\qquad
 \rho_b^C(t_b,h)>0,\qquad W_b\ge2U_b^C,
\]
and that
\[
 c_b^2(1-a_b^{-2})\frac{W_b}{2}
 \ge
 2\left(2(b^2+1)U_b^C(2/r)^{2b+29}\right),
\]
\[
 c_b^2(1-a_b^{-2})\frac{W_b}{2}
 \ge
 2\left(2(a_b/c_b)^{2b+29}\right).
\]
Every lower quantity is rounded downward and every upper quantity upward.
The replay exits with status $0$, reports
\texttt{C\_rows\_checked=156} and \texttt{failures=0}, and gives the
minimum defining-trace domain, square-trace domain, and tilted-mass lower
bounds
\[
 14.669512561242725\ldots,\qquad
 2.7221542041383810\ldots,\qquad
 0.028831039797142247\ldots.
\]
The three minimum natural-log comparison margins are
\[
 1.5000727210859520\ldots,\qquad
 1.5065100867456532\ldots,\qquad
 0.15645659337912115\ldots.
\]
The log has SHA-256
\[
\hashsplit{ca6860204466e4a6b38d9c5597511cd1}{1cfebb4b1b65cab88db6ee45f9477111},
\]
and the C++ source has SHA-256
\[
\hashsplit{0f5ad1037bcf738ab4ddd9c5b12f151a}{adc1bff2b5c4726c60622f608ef9e0e7}.
\]

The comparison $W_b\ge2U_b^C$ gives $p_1-p_2\ge W_b/2$.  Hence
\cref{lem:trace-tail} gives
\[
 P_{C_b}(c_b)\ge c_b^2(1-a_b^{-2})\frac{W_b}{2},
 \qquad N_{C_b}(a_b)\le2(b^2+1)U_b^C.
\]
The final two checked comparisons imply the hypothesis of
\cref{lem:pushforward-tail} at the odd exponent $2b+29$.  Its monotonicity
clause proves every later odd exponent.
\end{proof}

\begin{lemma}[Exact odd-orthogonal defining-trace MGFs]
\label{lem:b-exact-toeplitz-hankel-mgf}
Let $O$ be Haar-distributed in $\operatorname{SO}(2b+1)$, put
$Z_b=\operatorname{Tr}O$, and, for $v>0$ and $\epsilon\in\{+1,-1\}$, set
$M_{b,\epsilon}(v)=\mathbf E e^{\epsilon vZ_b}$.  Then
\[
 M_{b,+}(v)=e^v\det\left(
 I_{|j-k|}(2v)-I_{j+k+1}(2v)
 \right)_{0\le j,k<b},
 \tag{\ref{lem:b-exact-toeplitz-hankel-mgf}.1}
\]
\[
 M_{b,-}(v)=e^{-v}\det\left(
 I_{|j-k|}(2v)+I_{j+k+1}(2v)
 \right)_{0\le j,k<b},
 \tag{\ref{lem:b-exact-toeplitz-hankel-mgf}.2}
\]
where $I_m$ is the modified Bessel function.  Both displayed matrices are
positive definite.
\end{lemma}

\begin{proof}
For $O^+(2b+1)$, Courteaut--Johansson--Lambert
\cite[Section~3, equation~(3.1)]{CJL2022} identify the $b$ nontrivial
eigenvalue pairs with Jacobi parameters $(1/2,-1/2)$.  Their exact
Toeplitz--Hankel identity
\cite[Lemma~3.1, the $E_n^{+-}$ formula]{CJL2022}, applied to
$\psi(x)=e^{2vx}$, has Fourier coefficients
$\widehat\phi_m=I_m(2v)$.  Restoring the deterministic $+1$ eigenvalue
gives \eqref{lem:b-exact-toeplitz-hankel-mgf}.1.

For the argument $-v$, use $I_m(-2v)=(-1)^mI_m(2v)$.  Conjugating the
resulting matrix by $\operatorname{diag}(1,-1,1,-1,\ldots)$ changes it to
the plus-Hankel matrix in
\eqref{lem:b-exact-toeplitz-hankel-mgf}.2 and does not change its
determinant.  The two matrices are moment Gram matrices for the strictly
positive Jacobi weights occurring in Lemma~3.1 of the cited source, so they
are positive definite.
\end{proof}

\begin{lemma}[Exact even-orthogonal defining-trace MGFs]
\label{lem:d-exact-even-orthogonal-mgf}
For $b\ge2$, let $Z_b^+=\operatorname{Tr}O$ with $O$ Haar-distributed in
$\SO(2b)=O^+(2b)$.  For $r\ge1$, let $Z_r^-=\operatorname{Tr}O$ with $O$
Haar-distributed in the determinant-minus component $O^-(2r+2)$.  For $v>0$,
\[
 \mathbf E e^{vZ_b^+}
 =\frac12\det\left(
 I_{|j-k|}(2v)+I_{j+k}(2v)
 \right)_{0\le j,k<b},
 \tag{\ref{lem:d-exact-even-orthogonal-mgf}.1}
\]
and
\[
 \mathbf E e^{vZ_r^-}
 =\det\left(
 I_{|j-k|}(2v)-I_{j+k+2}(2v)
 \right)_{0\le j,k<r}.
 \tag{\ref{lem:d-exact-even-orthogonal-mgf}.2}
\]
Both displayed matrices are positive definite.
\end{lemma}

\begin{proof}
Courteaut--Johansson--Lambert
\cite[Section~3, equation~(3.1)]{CJL2022} identify $O^+(2b)$ with the
Jacobi parameters $(-1/2,-1/2)$ and $O^-(2r+2)$ with the parameters
$(1/2,1/2)$.  Apply their exact Toeplitz--Hankel identities
\cite[Lemma~3.1, the $E_b^{--}$ and $E_r^{++}$ formulas]{CJL2022} to
$\psi(x)=e^{2vx}$, whose Fourier coefficients are $I_j(2v)$.

For completeness, the factor $1/2$ in
\eqref{lem:d-exact-even-orthogonal-mgf}.1 is forced by the normalized
$(-1/2,-1/2)$ Jacobi measure.  In the orthonormal cosine basis
$1,\sqrt2T_1,\sqrt2T_2,\ldots$, Andreief's identity gives a Gram determinant.
The unscaled Toeplitz--Hankel matrix in the display is obtained by multiplying
the first Gram row and column by $\sqrt2$, so its determinant is twice the
normalized expectation.  In particular it equals $2$ at $v=0$.  The
$(1/2,1/2)$ formula is already normalized at $v=0$.  Finally, both matrices
are Gram matrices, up to this positive diagonal rescaling, for their strictly
positive Jacobi weights.  Hence every principal pivot is positive.
\end{proof}

\begin{lemma}[Exact low-rank trace supports]
\label{lem:exact-classical-trace-supports}
For $O\in\SO(2b+1)$,
\[
 -(2b-1)\le\operatorname{Tr}O\le2b+1.
 \tag{\ref{lem:exact-classical-trace-supports}.1}
\]
For $O\in O^-(2r+2)$,
\[
 |\operatorname{Tr}O|\le2r.
 \tag{\ref{lem:exact-classical-trace-supports}.2}
\]
For $U\in\Sp(2b)$ and $\tau_C=-\operatorname{Tr}(U^2)$,
\[
 |\operatorname{Tr}U|\le2b,\qquad |\tau_C|\le2b.
 \tag{\ref{lem:exact-classical-trace-supports}.3}
\]
\end{lemma}

\begin{proof}
The nonreal eigenvalues occur in conjugate unit-circle pairs and each pair
contributes a number in $[-2,2]$ to the trace.  An element of
$\SO(2b+1)$ has a forced $+1$ eigenvalue, giving the upper endpoint
$2b+1$ and the lower endpoint $1-2b$.  An element of $O^-(2r+2)$ has
forced eigenvalues $+1,-1$, whose trace contributions cancel, leaving $r$
conjugate pairs.  A symplectic matrix has $b$ such pairs and no forced real
eigenvalue; squaring preserves the unit-circle pair structure.  These
observations give all three displays.
\end{proof}

\begin{lemma}[Exact type-$C$ defining-trace MGF]
\label{lem:c-exact-defining-mgf}
Let $U$ be Haar-distributed in $\Sp(2b)$, $b\ge1$, and put
$S_b=\operatorname{Tr}U$.  For every $v>0$,
\[
 \mathbf E e^{vS_b}=\mathbf E e^{-vS_b}
 =\det\left(
 I_{|j-k|}(2v)-I_{j+k+2}(2v)
 \right)_{0\le j,k<b}.
 \tag{\ref{lem:c-exact-defining-mgf}.1}
\]
The displayed matrix is positive definite.
\end{lemma}

\begin{proof}
The remark following Rains' Theorem~1.5
\cite[equation~(36)]{RainsPowerMaps} identifies the nonfixed eigenvalues of
$\Sp(2b)$ with those of $O^-(2b+2)$.  The latter component has fixed
eigenvalues $+1,-1$, of total trace zero, so the full traces have the same
distribution.  Equation~\eqref{lem:d-exact-even-orthogonal-mgf}.2 with
$r=b$ gives the determinant in \eqref{lem:c-exact-defining-mgf}.1 and its
positive definiteness.  Multiplication of $U$ by the central element $-I$
preserves Haar measure and changes $S_b$ to $-S_b$, proving the equality of
the two MGFs.
\end{proof}

\begin{lemma}[Exact type-$C$ square-trace MGF]
\label{lem:c-exact-square-mgf}
Let $U$ be Haar-distributed in $\Sp(2b)$, $b\ge2$, and put
$\tau_C=-\operatorname{Tr}(U^2)$.  Define
\[
 m_0=\left\lceil\frac{b+1}{2}\right\rceil,
 \qquad m_1=\left\lceil\frac b2\right\rceil,
\]
and
\[
 H_m(v)=\det\left(
 I_{|j-k|}(2v)-I_{j+k+2}(2v)
 \right)_{0\le j,k<m-1}.
\]
With $M_{m_1,-}$ as in
\eqref{lem:b-exact-toeplitz-hankel-mgf}.2, one has, for every $v>0$,
\[
 \mathbf E e^{v\tau_C}=e^vH_{m_0}(v)M_{m_1,-}(v).
 \tag{\ref{lem:c-exact-square-mgf}.1}
\]
Both determinant matrices on the right are positive definite.
\end{lemma}

\begin{proof}
Use Rains' identification
$\Sp(2b)\sim O^-(2b+2)$ after omitting the fixed eigenvalues $+1,-1$
\cite[remark following Theorem~1.5, equation~(36)]{RainsPowerMaps}.
Theorem~1.5 at $p=2$, specifically equation~(22), decomposes the nonfixed
eigenvalues after squaring into independent components
\[
 O^-(2m_0)\oplus O^+(2m_1+1).
\]
If $Y_0,Y_1$ are their full traces, then the first component's fixed
eigenvalues have total trace zero and the second component has a fixed
$+1$.  Restoring these eigenvalues gives
\[
 \tau_C-1\overset d=-(Y_0+Y_1).
\]
The law of $Y_0$ is symmetric because multiplication by $-I$ preserves the
determinant-minus component in even dimension.  Thus
$\mathbf E e^{-vY_0}=H_{m_0}(v)$ by
\eqref{lem:d-exact-even-orthogonal-mgf}.2.  By definition,
$\mathbf E e^{-vY_1}=M_{m_1,-}(v)$.  Independence proves
\eqref{lem:c-exact-square-mgf}.1.  Positive definiteness is supplied by
\cref{lem:b-exact-toeplitz-hankel-mgf,lem:d-exact-even-orthogonal-mgf}.
\end{proof}

\begin{lemma}[Exact type-$D$ square-trace MGF]
\label{lem:d-exact-square-mgf}
Let $O$ be Haar-distributed in $\SO(2b)$ and put
$\tau_D=\operatorname{Tr}(O^2)$.  Define
\[
 m_0=\left\lceil\frac b2\right\rceil,
 \qquad m_1=\left\lceil\frac{b-1}{2}\right\rceil.
\]
If $G_{m,+}(v)=\mathbf E e^{vZ_m^+}$ denotes the even-orthogonal MGF in
\eqref{lem:d-exact-even-orthogonal-mgf}.1 and $M_{m,-}(v)$ denotes the
odd-orthogonal MGF in
\eqref{lem:b-exact-toeplitz-hankel-mgf}.2, then, for $v>0$,
\[
 \mathbf E e^{v\tau_D}=e^vG_{m_0,+}(v)M_{m_1,-}(v).
 \tag{\ref{lem:d-exact-square-mgf}.1}
\]
\end{lemma}

\begin{proof}
Rains' power-two decomposition
\cite[Theorem~1.5, equation~(22)]{RainsPowerMaps}, with the fixed
eigenvalues restored, gives independent full traces $Y_0,Y_1$ on
$O^+(2m_0)$ and $O^-(2m_1+1)$.  The latter component has a forced
eigenvalue $-1$, which is omitted in Rains' convention, and therefore
$\tau_D-1\overset d=Y_0+Y_1$; this is the type-$D$ identity used in the proof
of \cref{lem:bc-square-arbitrary-mgf}.  Independence and
\cref{lem:b-exact-toeplitz-hankel-mgf,lem:d-exact-even-orthogonal-mgf}
give \eqref{lem:d-exact-square-mgf}.1.
\end{proof}

\begin{lemma}[Layered tilt with exact two-sided MGFs]
\label{lem:exact-mgf-layered-trace-tilt}
Let $Z$ be a real random variable and write
$M_\epsilon(v)=\mathbf E e^{\epsilon vZ}$ for
$\epsilon\in\{+1,-1\}$.  Fix $s>0$ and $t>0$, and choose
$\lambda_->0$ and pairs $(x_i,\lambda_i)$, $i=0,\ldots,L-1$, with
$t<x_0<\cdots<x_{L-1}$ and $\lambda_i>0$.  Assume the displayed MGFs are
finite.  For each sign define
\[
 A_{\epsilon,-}=
 e^{\lambda_-t}
 \frac{M_\epsilon(s-\lambda_-)}{M_\epsilon(s)},
 \qquad
 A_{\epsilon,i}^+=
 e^{-\lambda_i x_i}
 \frac{M_\epsilon(s+\lambda_i)}{M_\epsilon(s)},
\]
\[
 g_{\epsilon,i}=1-A_{\epsilon,-}-A_{\epsilon,i}^+,
 \qquad w_i=e^{-sx_i}.
\]
If $\epsilon Z\le x_i$ almost surely, one may instead set
$A_{\epsilon,i}^+=0$.
If $0<g_{\epsilon,0}\le\cdots\le g_{\epsilon,L-1}$, then
\[
 \Pr\{\epsilon Z\ge t\}\ge
 V_\epsilon:=M_\epsilon(s)
 \left(
  w_{L-1}g_{\epsilon,L-1}
  +\sum_{i=0}^{L-2}(w_i-w_{i+1})g_{\epsilon,i}
 \right).
 \tag{\ref{lem:exact-mgf-layered-trace-tilt}.1}
\]
Consequently $\Pr\{|Z|\ge t\}\ge V_++V_-$.
\end{lemma}

\begin{proof}
For fixed $\epsilon$, tilt the law by
$d\mathbf P_{\epsilon,s}=M_\epsilon(s)^{-1}e^{\epsilon sZ}\,d\mathbf P$.
Exponential Markov gives
\[
 \mathbf P_{\epsilon,s}\{\epsilon Z<t\}\le A_{\epsilon,-},
 \qquad
 \mathbf P_{\epsilon,s}\{\epsilon Z>x_i\}\le A_{\epsilon,i}^+.
\]
The second inequality is replaced by the deterministic zero bound when
$\epsilon Z\le x_i$ almost surely.  Thus the tilted mass of
$t\le\epsilon Z\le x_i$ is at least
$g_{\epsilon,i}$.  The same summation-by-parts argument as in
\cref{lem:layered-trace-tilt} bounds the integral of the decreasing weight
$e^{-s\epsilon Z}$ by the bracket in
\eqref{lem:exact-mgf-layered-trace-tilt}.1.  Since
\[
 \Pr\{\epsilon Z\ge t\}
 =M_\epsilon(s)\mathbf E_{\epsilon,s}
   [e^{-s\epsilon Z}\mathbf1_{\{\epsilon Z\ge t\}}],
\]
multiplication by $M_\epsilon(s)$ proves the displayed inequality.  The
positive and negative tail events are disjoint for $t>0$, so their lower
bounds add.
\end{proof}

\begin{proposition}[Low-rank $B/C$ tails from exact determinant MGFs]
\label{prop:post29-low-bc-exact-mgf}
For the rows and onsets in the following table,
\[
 \Qthree^G(n)\ge0\qquad
 \text{for every odd }n\ge N_G.
\]
\[
\begin{array}{c|cccccc}
G&B_7&B_9&B_{10}&B_{11}&B_{12}&B_{13}\\ \hline
N_G&63&63&63&75&81&83
\end{array}
\]
\[
\begin{array}{c|cccccccc}
G&C_4&C_5&C_6&C_7&C_8&C_9&C_{10}&C_{11}\\ \hline
N_G&63&63&63&63&63&63&63&63
\end{array}
\qquad
\begin{array}{c|cccccccc}
G&C_{12}&C_{13}&C_{14}&C_{15}&C_{16}&C_{17}&C_{18}&C_{19}\\ \hline
N_G&63&71&75&63&65&73&63&63
\end{array}
\]
\end{proposition}

\begin{proof}
\ifginibrefullproof
For each row put
\[
 c=rb,\qquad q=\alpha_q\sqrt b,\qquad
 a=\alpha_a\sqrt b,\qquad
 t=\sqrt{2c+2d+q},\qquad s=t+l,
\]
\[
 u_i=u_0+i\Delta\quad(0\le i\le7),\qquad
 x_i=t+u_i,\qquad
 \lambda_-=f_-l,\qquad \lambda_i=f_+u_i.
\]
The fixed type-$B$ parameters are
\[
\begin{array}{c|ccccccccc|c}
G&r&\alpha_q&\alpha_a&d&l&u_0&f_-&f_+&(k,m)&N_G\\ \hline
B_7&45/7&7&35/10&14/10&10&25/10&1/10&33/100&(3,3)&63\\
B_9&60/9&9&4333/1000&14/10&5&25/10&1/10&1/2&(2,6)&63\\
B_{10}&75/10&7&4250/1000&14/10&4&25/10&1/10&1/2&(2,7)&63\\
B_{11}&80/11&11&4523/1000&14/10&3&3&5/100&1/2&(3,7)&75\\
B_{12}&90/12&12&4619/1000&14/10&3&3&5/100&1/2&(3,8)&81\\
B_{13}&100/13&13&4715/1000&14/10&2&3&5/100&75/100&(4,8)&83
\end{array}
\tag{\ref{prop:post29-low-bc-exact-mgf}.1}
\]
Here $\Delta=1/10$, $f_q=1$, and $f_a=509/500$ in every type-$B$ row.
The fixed type-$C$ parameters are
{\scriptsize
\[
\begin{array}{c|ccccccccc|c}
G&r&\alpha_q&\alpha_a&d&l&u_0&f_-&f_+&(k,m)&N_G\\ \hline
C_4&32/10&401/100&35/10&14/10&35/10&10&5/100&15/10&(1,2)&63\\
C_5&4&5&35/10&14/10&3&100&5/100&15/10&(2,2)&63\\
C_6&32/6&6&35/10&14/10&5&15/10&1/5&15/10&(2,3)&63\\
C_7&6&7&35/10&14/10&5&125/100&1/5&2&(3,3)&63\\
C_8&60/8&8&46/10&14/10&7&15/10&5/100&2&(3,4)&63\\
C_9&60/9&9&42/10&14/10&5&175/100&1/10&2&(4,4)&63\\
C_{10}&6&7&4&14/10&3&25/10&5/100&2&(2,7)&63\\
C_{11}&70/11&7&45/10&14/10&1&25/10&5/100&15/10&(3,7)&63\\
C_{12}&80/12&7&45/10&14/10&2&2&1/2&1/2&(3,8)&63\\
C_{13}&80/13&8&45/10&14/10&1/2&3&1/10&15/10&(3,9)&71\\
C_{14}&20/3&8&45/10&14/10&2&3&1/2&1&(3,10)&75\\
C_{15}&76/15&8&44/10&14/10&1&25/10&3/4&3/4&(5,9)&63\\
C_{16}&75/16&7&45/10&14/10&75/100&250/100&75/100&75/100&(4,11)&65\\
C_{17}&5&9&45/10&14/10&1&25/10&3/4&1/2&(3,13)&73\\
C_{18}&226172/10^5&460576/10^5&438971/10^5&138552/10^5&
105713/10^5&240817/10^5&1&56/100&(5,8)&63\\
C_{19}&226172/10^5&460576/10^5&438971/10^5&138552/10^5&
105713/10^5&240817/10^5&1&56/100&(5,8)&63
\end{array}
\tag{\ref{prop:post29-low-bc-exact-mgf}.2}
\]
}
Here $f_q=75841/10^5$ and $f_a=79538/10^5$.  Also
$\Delta=1/10$ except in the last two rows, where
$\Delta=30177/10^6$.

For $B_b$, apply \cref{lem:exact-mgf-layered-trace-tilt} separately to
$\epsilon=+1,-1$, with the exact MGFs of
\cref{lem:b-exact-toeplitz-hankel-mgf}, and call the sum of the resulting
lower bounds $L_G$.  For $C_b$, use the symmetry and exact determinant in
\cref{lem:c-exact-defining-mgf}; twice the positive-sign lower bound is
$L_G$.  Whenever an endpoint reaches the deterministic support, use the
zero upper-tail clause justified by
\cref{lem:exact-classical-trace-supports}.

Put $v_q=f_q(q-1)/2$ and $v_a=f_a(a-1)/2$.  For type $B$, set
\[
 U_B(z,v)=\exp\{-v(z-1)+v^2\}.
\]
For type $C$, let
\[
 U_C(z,v)=e^{-v(z-1)}H_{m_0}(v)M_{m_1,-}(v),
\]
with $m_0,m_1,H$ as in \cref{lem:c-exact-square-mgf}.  Chernoff's
inequality and
\cref{lem:bc-square-arbitrary-mgf,lem:c-exact-square-mgf} give the required
square-trace upper bounds; if $z>2b$, the bound is zero by
\eqref{lem:exact-classical-trace-supports}.3.  Define
\[
 P_G^*=c^2(1-d^{-2})\bigl(L_G-U_G(q,v_q)\bigr),
 \qquad
 N_G^*=\frac{8e^{-2}}{v_a^2}U_G(a,v_a),
\]
and
\[
 H_{k,m}(a)=
 \sum_{j=0}^{2m}(-1)^j\binom{2m}{j}a^{-j}K_{2k+j}.
\]

The fixed $384$-bit directed-rounding replay
\[
\texttt{certificates/post\_m29/
post\_m29\_low\_bc\_exact\_mgf\_mpfr\_certificate.log}
\]
checks every geometry condition, determinant pivot, tilted-layer mass,
layer monotonicity, positive difference $L_G-U_G(q,v_q)$, and polynomial
expectation.  At the displayed onset it checks
\[
 P_G^*\ge
 N_G^*\left(\frac{2b}{c}\right)^{N_G}
 +\frac{a^{N_G-2k}}{4^mc^{N_G}}H_{k,m}(a).
 \tag{\ref{prop:post29-low-bc-exact-mgf}.3}
\]
The replay constructs every $K_j$ as an exact GMP integer, encloses every
Bessel series remainder, and evaluates each determinant by interval
Cholesky decomposition.  It checks $22$ rows, reports no failure, and has
minimum final natural-log margin
\[
 0.07437283746004294208266355122>0
\]
at $C_{15}$.  Its log and source SHA-256 hashes are respectively
\[
\hashsplit{5bf958c5ce217dfc64efc08c1ef72bd6}{98aab30f1259eb3d93c9bcb0c8e0f3d2},
\]
\[
\hashsplit{71cc2293c4ad5ba50d216c8f939cd1ec}{604ce7b86d6a19e9be1fcb27adc421fe}.
\]
The degree check $2k+2m\le2b-2$ holds in every row, so the stable moments
used in \eqref{prop:post29-low-bc-exact-mgf}.3 are supplied by
\cref{lem:stable-finite-rank-push-moments}.  Thus
\cref{lem:one-sided-moment-pushforward} proves the claim at $N_G$.  Since
$2b/c<1$ and $a/c<1$ in every checked row, the monotonicity clause of that
lemma proves every later odd exponent.
\else
The manifested verifier source contains the exact rational parameter row for
every displayed onset.  The $384$-bit directed-rounding C++/MPFR replay evaluates
the exact Toeplitz--Hankel defining-trace MGFs, the square-trace Chernoff
bounds, and the stable push moments, and checks the sufficient inequality of
\cref{lem:one-sided-moment-pushforward}.  It checks all $22$ rows with no
failure; the minimum final natural-log margin is
$0.07437283746004294208266355122>0$ at $C_{15}$.  It also checks
$2k+2m\le2b-2$, $2b/c<1$, and $a/c<1$ in every row.  Thus the criterion
holds at the listed onset and its monotonicity clause gives every later odd
exponent.  The accepted log and current C++ source are authenticated by the
replay ledger.
\fi
\end{proof}

\begin{proposition}[Finite Chain bridges to the low-rank exact-MGF onsets]
\label{prop:post29-low-bc-exact-bridges}
The Chain differences are nonnegative in the following ranges:
\[
\begin{array}{c|c}
G&\{m:D_G(2m+1)\ge0\}\\ \hline
B_{11}&31\le m\le35\\
B_{12}&31\le m\le38\\
B_{13}&31\le m\le39\\
C_{13}&31\le m\le33\\
C_{14}&31\le m\le35\\
C_{17}&31\le m\le34.
\end{array}
\tag{\ref{prop:post29-low-bc-exact-bridges}.1}
\]
Consequently, if $\Qthree^G(63)\ge0$ in one of these six rows, then
$\Qthree^G(n)\ge0$ for every odd $63\le n<N_G$, where $N_G$ is the onset
in \cref{prop:post29-low-bc-exact-mgf}.
\end{proposition}

\begin{proof}
\ifginibrefullproof
Write the finite-row moment as
$m_j^G=m_j^{\mathrm{st}}+\delta_j^G$.  Expanding the Chain difference gives
the exact identity
\[
 D_G(2m+1)=D_{\mathrm{st}}(2m+1)
 +\sum_jL_{m,j}\delta_j^G
 +\sum_{i\le j}Q_{m,i,j}\delta_i^G\delta_j^G.
 \tag{\ref{prop:post29-low-bc-exact-bridges}.2}
\]
The authenticated source logs supply consecutive exact corrections through
the following degrees:
\[
\begin{array}{c|cccccc}
G&B_{11}&B_{12}&B_{13}&C_{13}&C_{14}&C_{17}\\ \hline
\text{last degree}&51&53&57&43&51&45.
\end{array}
\]
Every supplied correction is
nonpositive.  By \cref{prop:bc-pieri-correction-sign}, every missing higher
correction is also nonpositive; nonnegativity of the actual moment therefore
gives $-m_j^{\mathrm{st}}\le\delta_j^G\le0$.  The exact-integer replay
retains every known linear and quadratic term in
\eqref{prop:post29-low-bc-exact-bridges}.2, bounds a missing linear term
below by $-L_{m,j}m_j^{\mathrm{st}}$ when $L_{m,j}>0$, bounds a missing
negative quadratic term using
$|\delta_i^G\delta_j^G|\le m_i^{\mathrm{st}}m_j^{\mathrm{st}}$, and drops
only nonnegative terms.

The replay
\[
\texttt{certificates/post\_m29/
post\_m29\_low\_bc\_exact\_onset\_bridge\_certificate.log}
\]
checks all $34$ differences in
\eqref{prop:post29-low-bc-exact-bridges}.1.  Its smallest exact integer
lower bound is
\[
64688975233529723032612107788184688805229314688529030019703101781365210144991608409847821836>0
\]
at $C_{13},m=31$.  It records status zero and has SHA-256
\[
\hashsplit{c156d1296b8f74afc24f9f48ee615f3a}{ba57e5e9d87f671258139226f8f57d28}.
\]
This proves \eqref{prop:post29-low-bc-exact-bridges}.1.  Starting from odd
exponent $63$, repeated use of \cref{lem:chain-propagation} gives the final
assertion.
\else
Expanding each Chain difference in
$m_j^G=m_j^{\mathrm{st}}+\delta_j^G$ gives an exact linear--quadratic
polynomial in the corrections.  The exact-integer replay consumes the
authenticated consecutive exact-correction logs and uses
$-m_j^{\mathrm{st}}\le\delta_j^G\le0$ only beyond the supplied prefix, as
given by the Pieri omitted-summand formula and moment nonnegativity.  It
checks all $34$ displayed differences.  Its smallest exact integer lower
bound is positive and occurs at $C_{13},m=31$.  Hence
\cref{lem:chain-propagation} carries the value at odd exponent $63$ to each
exact-MGF onset.
\fi
\end{proof}

\begin{proposition}[Residual $B/C$ tails from layered MGFs and one-sided
moments]
\label{prop:post29-bc-layered-mgf}
For every $B_b$ with $14\le b\le61$ and every $C_b$ with
$20\le b\le61$, one has
\[
 \Qthree^G(n)\ge0
\]
for every odd $n\ge2b+29$.
\end{proposition}

\begin{proof}
The following table gives four exact rational parameter sets.  In each row
put
\[
 c=rb,\quad q=\alpha_q\sqrt b,\quad a=\alpha_a\sqrt b,
 \quad t=\sqrt{2c+2d+q},\quad s=t+l,
\]
\[
 u_i=u_0+i\Delta\quad(0\le i\le7),\qquad
 \lambda_-=l,\qquad \lambda_i=f_+u_i,
\]
\[
 v_q=f_q(q-1)/2,\qquad v_a=f_a(a-1)/2.
\]
\[
\begin{array}{c|c|c|c|c|c|c|c|c|c|c|c}
G&r&\alpha_q&\alpha_a&d&l&u_0&\Delta&f_+&f_q&f_a&(k,m)\\ \hline
B_{14}\text{--}B_{15}&\frac{16}{5}&\frac{529}{100}&\frac{213}{50}&
\frac{143}{100}&\frac{109}{100}&\frac{13}{10}&\frac{31}{1000}&
1&1&\frac{509}{500}&(5,b-6)\\
B_{16}&\frac{200001}{10^5}&\frac{45546}{10^4}&\frac{4284}{10^3}&
\frac{1277}{10^3}&\frac{856}{10^3}&\frac{1847}{10^3}&\frac{34}{10^3}&
\frac{567}{10^3}&1&\frac{101331}{10^5}&(5,10)\\
B_{17}&\frac{20446}{10^4}&\frac{44739}{10^4}&\frac{42902}{10^4}&
\frac{1341}{10^3}&\frac{972}{10^3}&\frac{15366}{10^4}&\frac{332}{10^4}&
\frac{7291}{10^4}&1&\frac{101331}{10^5}&(5,10)\\
B_{18}\text{--}B_{61}&\frac{20674}{10^4}&\frac{441721}{10^5}&
\frac{429465}{10^5}&\frac{13501}{10^4}&\frac{103311}{10^5}&
\frac{139822}{10^5}&\frac{31205}{10^6}&1&1&\frac{101331}{10^5}&(5,8)\\
C_{20}\text{--}C_{61}&\frac{226172}{10^5}&\frac{460576}{10^5}&
\frac{438971}{10^5}&\frac{138552}{10^5}&\frac{105713}{10^5}&
\frac{135104}{10^5}&\frac{30177}{10^6}&1&\frac{75841}{10^5}&
\frac{79538}{10^5}&(5,8)
\end{array}
\]
For $B_{14},B_{15}$, let $V_{b,+},V_{b,-}$ be the exact-MGF lower bounds
in \cref{lem:exact-mgf-layered-trace-tilt}, using
\cref{lem:b-exact-toeplitz-hankel-mgf}, $L=8$, $\lambda_-=l$, and
$x_i=s+u_i$, $\lambda_i=u_i$, and put
$L_b=V_{b,+}+V_{b,-}$.  In the other four rows,
let $V_b$ be the lower bound in \cref{lem:layered-trace-tilt} with the
displayed parameters and put $L_b=2V_b$.  For $G=B_b,C_b$, respectively,
write $U_G$ for the corresponding function in
\cref{lem:bc-square-arbitrary-mgf}.  Define
\[
 P_b^*=c^2(1-d^{-2})\bigl(L_b-U_G(q,v_q)\bigr),
 \qquad
 N_b^*=\frac{8e^{-2}}{v_a^2}U_G(a,v_a),
\]
and, using the stable push moments of
\cref{lem:stable-finite-rank-push-moments},
\[
 H_{k,m}(a)=
 \sum_{j=0}^{2m}(-1)^j\binom{2m}{j}a^{-j}K_{2k+j}.
\]

The $384$-bit directed-rounding MPFR replay
\[
\texttt{certificates/post\_m29/
post\_m29\_bc\_layered\_mgf\_mpfr\_certificate.log}
\]
checks every intermediate domain and sign condition and, at
$n_b=2b+29$, the final comparison
\[
 P_b^*\ge
 N_b^*\left(\frac{2b}{c}\right)^{n_b}
 +\frac{a^{n_b-2k}}{4^mc^{n_b}}H_{k,m}(a).
 \tag{\ref{prop:post29-bc-layered-mgf}.1}
\]
For $B_{14},B_{15}$ it also evaluates
\[
 I_\nu(2v)=\sum_{r=0}^{180}
 \frac{v^{2r+\nu}}{r!(r+\nu)!}+R_{\nu,v},
 \qquad
 0<R_{\nu,v}\le
 \frac{v^{362+\nu}}{181!(181+\nu)!}
 \frac1{1-v^2/(181(181+\nu))},
\]
after checking that the final ratio is below one for every argument and
$0\le\nu\le29$.  It evaluates both exact Toeplitz--Hankel determinants by
interval Cholesky decomposition and checks every pivot strictly positive.
It constructs the integers $K_j$
exactly in GMP from their generating function, converts them to
outward-rounded MPFR intervals, and uses outward rounding for every
subsequent arithmetic operation.  The combined replay checks $333$ rows:
$48$ type-$B$, $42$ type-$C$, and the $243$ type-$D$ rows used later in
\cref{prop:post29-d-a-free}.  It exits with status $0$ and
reports \texttt{failures=0}.  Its $B/C$ lower margins include
\[
\begin{array}{c|c}
\text{condition}&\text{minimum directed lower bound}\\ \hline
c-2b,\ 2b-a,\ a-1,\ d-1&0.00016\\
1-\mathfrak a_d&0.24176531795234819\ldots\\
\text{exact Cholesky pivot}&0.85461901487040600\ldots\\
g_0&0.0028157129657487152\ldots\\
\min_i(g_{i+1}-g_i)&0.008012263371171611\ldots\\
\log P_b^*-\log(\text{right side of
\eqref{prop:post29-bc-layered-mgf}.1})&
1.358823017167476129\ldots
\end{array}
\]
The two exact-determinant rows have final natural-log margins
\[
6.045003937030598629\ldots
\quad\text{and}\quad
7.590922674135796189\ldots .
\]
Among the $B/C$ rows the final minimum is attained at $B_{16}$.  The combined-log
minimum is the positive $D_{220}$ margin quoted in
\cref{prop:post29-d-a-free}.  The log has SHA-256
\[
\hashsplit{3e4900e5e2de2ec1a7fb04830ac59dea}{7369f8eebb1932ec08ac9ef73ba11654},
\]
and the C++ source has SHA-256
\[
\hashsplit{71cc2293c4ad5ba50d216c8f939cd1ec}{604ce7b86d6a19e9be1fcb27adc421fe}.
\]

By \cref{lem:exact-mgf-layered-trace-tilt,lem:layered-trace-tilt}, as
applicable, the defining-trace event in
\cref{lem:decoupled-trace-tail} has probability at least $L_b$.  Equations
\eqref{lem:bc-square-arbitrary-mgf}.1--\eqref{lem:bc-square-arbitrary-mgf}.2
bound the positive square-trace loss by $U_G(q,v_q)$.  Therefore
\eqref{lem:decoupled-trace-tail}.1 gives $P_G(c)\ge P_b^*$.  Equations
\eqref{lem:decoupled-trace-tail}.2--\eqref{lem:decoupled-trace-tail}.3 give
$N_G(a)\le N_b^*$.  The degree conditions for
\cref{lem:stable-finite-rank-push-moments} are
$2k+2m=2b-2$ for $B_{14},B_{15}$,
$2k+2m=30=2(16)-2$ for $B_{16},B_{17}$, and
$2k+2m=26\le2b-2$ in the uniform rows.  Thus
\eqref{prop:post29-bc-layered-mgf}.1 and
\cref{lem:one-sided-moment-pushforward} prove the claim at $n_b$.  Both ratios
$2b/c$ and $a/c$ are strictly below one, so the monotonicity clause of that
lemma proves every later odd exponent.
\end{proof}

\begin{proposition}[Post-$m=29$ first-hit lower bound]
\label{prop:post29}
For the remaining post-$m=29$ finite residue rows
\[
  B_{11},\ldots,B_{275},B_{277},\qquad
  C_{13},\ldots,C_{275},C_{277},\qquad
  D_{25},\ldots,D_{275},D_{277},D_{279},
\]
one has
\[
  D_G\bigl(N_{\mathrm{hit}}^{(29)}(G)-2\bigr)\ge0
\]
where
\[
  N_{\mathrm{hit}}^{(29)}(G)=
  \max(63,s(G)+2).
\]
Consequently $\Qthree^G(N_{\mathrm{hit}}^{(29)}(G))\ge0$ for each of these
rows.
\end{proposition}

\begin{proof}
The 783 rows split into five arithmetic buckets:
\[
\begin{array}{c|c|c}
\text{bucket} & \text{rows} & \text{minimum positive margin row}\\
\hline
\text{BC exact }m=30\text{ window} & 14 & B_{11}\\
\text{BC stable-tail }m=30 & 26 & B_{19}\\
\text{BC stable-tail sloped} & 490 & B_{32}\\
\text{D exact two-sided interval} & 28 & D_{25}\\
\text{D }A\text{-free frontier} & 225 & D_{53}
\end{array}
\]
For a row whose first hit is $n=2m+3$, let
\[
  \Delta_{\mathrm{st}}(m)
  =D_{\mathrm{st}}(2m+1)
\]
be the stable Chain difference, and write the finite-row moment correction
as $m_j^G=m_j^{\mathrm{st}}+\delta_j^G$.  Expanding
\cref{def:chain-diff} in the variables $\delta_j^G$ gives
\[
  D_G(2m+1)
  =
  \Delta_{\mathrm{st}}(m)
  +\sum_j L_{m,j}\delta_j^G
  +\sum_{i\le j}Q_{m,i,j}\delta_i^G\delta_j^G
\]
with integer coefficients $L_{m,j}$ and $Q_{m,i,j}$ determined from
\cref{lem:moment-formula}.  The $B/C$ checker computes these coefficients
exactly and proves a positive integer lower bound in each of its $530$ rows.
Those rows use the proved nonpositive boundary correction, so the positive
linear tail is compared with the negative quadratic terms.  The accepted
replay is restricted to $B/C$ and does not parse a determinant-only type-$D$
row.  It reports
\[
\begin{array}{ll}
\text{rows checked} & 530,\\
\text{stable moments generated} & m_0,\ldots,m_{557},\\
\text{scope} & \text{one Chain step at }N_{\mathrm{hit}}^{(29)}(G),\\
\text{exit status} & 0.
\end{array}
\]
Its three minimum margins have the following decimal sizes:
\[
\begin{array}{c|c|c}
\text{bucket} & \text{minimum row} & \log_{10}(\text{minimum margin})\\
\hline
\text{BC exact }m=30\text{ window} & B_{11} &
85.79\\
\text{BC stable-tail }m=30 & B_{19} &
85.79\\
\text{BC stable-tail sloped} & B_{32} &
93.85
\end{array}
\]
The replay log has SHA-256
\[
\hashsplit{21642c87b90affa8e618bb9f97bf6ad3}{46289380832690546488c6dcc7a7300e}.
\]

For $D_{25},\ldots,D_{52}$, the exact two-sided interval replay of
\cref{prop:post29-d25-52-exact} checks every Chain step from the stable edge
through its direct-tail onset, hence in particular the displayed first-hit
step.  Its $707$ exact checks have minimum lower endpoint
$4695986280539928492451072>0$.  For every displayed $D_b$ with $b\ge53$,
the short-frontier replay of \cref{prop:post29-d-a-free} checks the required
moment ceiling $N_D(b)+2\le2b$ and every Chain step from $m=30$ to the
direct-tail onset using \cref{lem:d-determinant-component}.3.  Its worst
positive logarithmic lower bound is $167.515421966088354\ldots$ at
$D_{53},m=50$.  The two status-$0$ logs have SHA-256 hashes
\[
\begin{array}{c}
\hashsplit{e47196c8097fb2a087b4673d47a5e34c}{aae1cd1a316058759e17b13708e2d888},\\[3pt]
\hashsplit{876d3f752d03b8f35183f0b28af7c15e}{9e28d104396fc65c215e5d854d8ae31f}.
\end{array}
\]

Every one of the five buckets therefore has a strictly positive lower
margin.  Thus
$D_G(N_{\mathrm{hit}}^{(29)}(G)-2)\ge0$ in every displayed row.  The preceding
odd exponent is either covered by the stable edge or by the row-gated bridge,
so \cref{lem:chain-propagation} gives
$\Qthree^G(N_{\mathrm{hit}}^{(29)}(G))\ge0$ in every displayed row.  The $B/C$ first-hit replay supplies only this one Chain step and is not used as
an all-later tail. The all-later claims are proved separately below by the
interval, layered-MGF, tilted-MGF, and correction-bridge propositions.
\end{proof}

\begin{lemma}[Exact finite-$D$ moment decomposition]
\label{lem:d-finite-moment-decomposition}
For $D_b=\SO(2b)$ write
\[
  e_2^j=\sum_{\lambda\vdash 2j}c_j(\lambda)s_\lambda,
\]
where the nonnegative integers $c_j(\lambda)$ are given by the iterated
vertical-two-strip Pieri rule.  Put
\[
 A_j(b)=
 \sum_{\substack{\nu\vdash j\\\ell(\nu)>2b}}
 c_j(2\nu_1,2\nu_2,\ldots)
\]
and
\[
 B_j(b)=
 \begin{cases}
 0,&j<b,\\[2pt]
 \displaystyle
 \sum_{\substack{\nu\vdash j-b\\\ell(\nu)\le2b}}
 c_j(1+2\nu_1,\ldots,1+2\nu_{\ell(\nu)},
              1^{,2b-\ell(\nu)}),&j\ge b.
 \end{cases}
\]
If $s_j=m_j^{\mathrm{st}}$, then the exact correction is
\[
  m_j^{D_b}-s_j=B_j(b)-A_j(b).
\]
In particular, the determinant-isotypic sum $B_j(b)$ alone is not the
finite-rank correction once $A_j(b)$ is nonzero.
\end{lemma}

\begin{proof}
The adjoint character of $\SO(2b)$ is $e_2$ in the defining eigenvalues.
Normalized Haar measure on $\SO(2b)$ is obtained from normalized Haar measure
on $O(2b)$ by multiplication by $1+\det$.  Expand $e_2^j$ by Pieri's rule.
Rains' orthogonal Schur-integral criterion \cite[\S3]{Rains1997} says that
the untwisted $O(2b)$ integral retains exactly the Schur summands with even
row lengths and at most $2b$ rows.  The stable sum retains all even-row
summands, so the omitted part is exactly $A_j(b)$.

For the determinant-twisted integral, use
$\det\,s_\lambda=s_{\lambda+(1^{2b})}$ on $O(2b)$.  The same criterion is
nonzero precisely when $\lambda$ has exactly $2b$ rows, all odd.  Such a
partition has the unique displayed form with $\nu\vdash j-b$; none exists
for $j<b$.  Its total contribution is $B_j(b)$.  Adding the untwisted and
determinant-twisted integrals gives
$m_j^{D_b}=s_j-A_j(b)+B_j(b)$.
\end{proof}

\begin{lemma}[Type-$D$ determinant component and the $A$-free envelope]
\label{lem:d-determinant-component}
Let $b\ge4$, and use the notation of
\cref{lem:d-finite-moment-decomposition}.  Put
\[
 p_j(b)=s_j-A_j(b).
\]
Then
\[
 p_j(b)-B_j(b)
 =\int_{\operatorname{Sp}(2b-2)}\chi_{\omega_2}(V)^j\,dV,
 \tag{\ref{lem:d-determinant-component}.1}
\]
where $\omega_2$ is the second fundamental representation of
$\operatorname{Sp}(2b-2)$.  Consequently, for every $j\ge0$,
\[
 0\le B_j(b)\le p_j(b)\le s_j.
 \tag{\ref{lem:d-determinant-component}.2}
\]
Moreover $A_j(b)=0$ for $0\le j\le2b$.  Hence throughout this window the
actual type-$D$ correction satisfies
\[
 0\le m_j^{D_b}-s_j=B_j(b)\le s_j.
 \tag{\ref{lem:d-determinant-component}.3}
\]
\end{lemma}

\begin{proof}
Normalized Haar measure on $O(2b)$ is the half-sum of normalized Haar
measure on its determinant components.  Since
\[
 p_j(b)=\int_{O(2b)}e_2(U)^j\,dU,
 \qquad
 B_j(b)=\int_{O(2b)}\det(U)e_2(U)^j\,dU,
\]
the normalized $O^-(2b)$ component moment is $p_j(b)-B_j(b)$.

Rains \cite[remark following Theorem~1.5, equation~(36)]{RainsPowerMaps}
identifies the nonfixed eigenvalue distribution of $O^-(2b)$ with that of
$\operatorname{Sp}(2b-2)$.  Restoring the fixed eigenvalues $+1,-1$, let
$V$ denote the defining $2b-2$ dimensional symplectic representation.  On
eigenvalue multisets,
\[
 e_2(V\oplus\mathbf1\oplus(-\mathbf1))
 =e_2(V)-1.
\]
The symplectic decomposition
$\Lambda^2V\cong\mathbf1\oplus V_{\omega_2}$ therefore turns the last
character into $\chi_{\omega_2}$.  This proves
\eqref{lem:d-determinant-component}.1.  Its right side is the multiplicity
of the trivial representation in $V_{\omega_2}^{\otimes j}$, so it is a
nonnegative integer.  Since $A_j(b),B_j(b)\ge0$, this proves
\eqref{lem:d-determinant-component}.2.

Finally, a partition $\nu$ with more than $2b$ nonzero parts has
$|\nu|\ge2b+1$.  The defining sum for $A_j(b)$ is therefore empty when
$j\le2b$.  Substitution in
\cref{lem:d-finite-moment-decomposition} proves
\eqref{lem:d-determinant-component}.3.
\end{proof}

\begin{lemma}[Two-sided type-$D$ correction box]
\label{lem:d-two-sided-correction-box}
Use the notation of \cref{lem:d-finite-moment-decomposition} and put
$\delta_j(b)=m_j^{D_b}-s_j$ and $U_j=s_j$.
Then
\[
 \delta_j(b)=0\quad(0\le j<b),
 \tag{\ref{lem:d-two-sided-correction-box}.1}
\]
\[
 0\le\delta_j(b)\le U_j\quad(b\le j\le2b),
 \tag{\ref{lem:d-two-sided-correction-box}.2}
\]
and
\[
 -s_j\le\delta_j(b)\le U_j\quad(j>2b).
 \tag{\ref{lem:d-two-sided-correction-box}.3}
\]
\end{lemma}

\begin{proof}
For $j<b$, both sums $A_j(b),B_j(b)$ vanish.  For $j\le2b$, the sum
$A_j(b)$ vanishes, so \cref{lem:d-determinant-component} gives
$\delta_j(b)=B_j(b)\ge0$.  For all $j$, the same lemma gives
$0\le B_j(b)\le s_j$ and the definition gives $0\le A_j(b)\le s_j$;
hence $\delta_j(b)\ge-s_j$.  The same inequalities give
$\delta_j(b)=B_j(b)-A_j(b)\le B_j(b)\le s_j=U_j$ for every parity.
Together with the preceding vanishing and nonnegativity statements, this
proves all three displays.
\end{proof}

\begin{proposition}[Unconditional $A$-free type-$D$ row-gated bridge]
\label{prop:d-a-free-row-gated-bridge}
For every $b\ge53$, the type-$D_b$ rows in the finite bridge through
$m=29$ satisfy
\[
 D_{D_b}(2m+1)\ge0
\]
at every row-gated Chain step.
\end{proposition}

\begin{proof}
At a row-gated step $m\le29$, the two diagonals in the exact moment formula
use moment indices only through $2m+5\le63$.  Since $63\le2b$ for
$b\ge53$, \cref{lem:d-two-sided-correction-box} supplies
\[
 \delta_j^{D_b}=0\quad(j<b),\qquad
 0\le\delta_j^{D_b}\le s_j\quad(b\le j\le63).
\]
The exact GMP/OpenMP verifier expands each Chain difference as in
\eqref{prop:bc-source-row-gated-bridge}.9 and minimizes every linear and
quadratic monomial over these intervals.  For $53\le b\le63$ it checks the
$36$ row-gated pairs
\[
 \left\lfloor\frac{b-4}{2}\right\rfloor\le m\le29.
\]
Every lower bound is positive; the least is
\[
1295951275047754385311107869606245254765288973412864475150697991241728
>0
\]
at $D_{53},m=24$.  For $b\ge64$ the row-gated range through $m=29$ is empty.
The verifier source and deterministic machine-C output have SHA-256 hashes
\[
\begin{aligned}
&\hashsplit{1e574e0c53e624df30e9c69bc24280b}{8e6c56a099923428f0a27036b74a7454b},\\
&\hashsplit{a84df851316bf376f00933ff958b7ad0d}{ac444acfa8f45d05d4221f17ddadd12}.
\end{aligned}
\]
Thus every nonvacuous row in the proposition has a source-generated exact
integer certificate.
\end{proof}

\begin{proposition}[Unconditional short-onset type-$D$ tails]
\label{prop:post29-d-a-free}
For every $53\le b\le295$ and every odd integer
\[
 n\ge N_{\mathrm{hit}}^{(29)}(D_b),
\]
one has
\[
 \Qthree^{D_b}(n)\ge0.
\]
\end{proposition}

\begin{proof}
For each rank, let $N_D(b)$ be the odd onset in
\[
\texttt{certificates/post\_m29/
post\_m29\_bc\_layered\_mgf\_mpfr\_certificate.log}.
\]
The directed-rounding verifier uses
\[
\begin{gathered}
 c_b=\frac{20674}{10000}b,
 \quad q_b=\frac{441721}{100000}\sqrt b,
 \quad a_b=\frac{429465}{100000}\sqrt b,
 \quad d=\frac{13501}{10000},\\
 l=\frac{103311}{100000},
 \quad u_i=\frac{139822}{100000}
       +i\frac{31205}{1000000}\quad(0\le i<8),
 \quad k=5,
 \quad m=8.
\end{gathered}
\]
The lower and upper defining-trace tilt fractions and the positive
square-trace tilt fraction are $1$; the negative square-trace fraction is
$101331/100000$.  The full defining trace of $\SO(2b)$ is centered, and
\cref{lem:extended-cjl-trace-mgf,lem:layered-trace-tilt} give the checked
two-sided defining-trace lower bound.  The type-$D$ clause
\eqref{lem:bc-square-arbitrary-mgf}.3 gives both checked square-trace upper
bounds.  Since $2k+2m=26\le b-3$, the type-$D$ clause of
\cref{lem:stable-finite-rank-push-moments} supplies the exact polynomial
$H_{5,8}^{D_b}(a_b)$.

For each $b$, the verifier starts at
$N_{\mathrm{hit}}^{(29)}(D_b)$ and increases the odd exponent until the
strict comparison
\[
 P_{D_b}(c_b)>
 N_{D_b}(a_b)\left(\frac{2b}{c_b}\right)^{N_D(b)}
 +\frac{a_b^{N_D(b)-10}}{4^8c_b^{N_D(b)}}
   H_{5,8}^{D_b}(a_b)
 \tag{\ref{prop:post29-d-a-free}.1}
\]
holds.  For odd $b$, $C_{D_b}=b-2$, so the first ratio in
\cref{lem:one-sided-moment-pushforward} is no larger than the conservative
$2b/c_b$ used in \eqref{prop:post29-d-a-free}.1.  All geometry,
trace-norm, determinant, nested-mass, polynomial, and final comparisons are
outward-rounded at $384$-bit MPFR precision.  The machine-C run checks all
$243$ ranks, reports
\[
\texttt{SUMMARY rows\_checked=333 B\_rows\_checked=48
C\_rows\_checked=42 D\_rows\_checked=243 failures=0},
\]
and has minimum final natural-log margin
\[
 0.007880461551348326759916693134>0
\]
at $D_{220}$.  Its log and C++ source have respective SHA-256 hashes
\[
\begin{aligned}
&\hashsplit{3e4900e5e2de2ec1a7fb04830ac59dea}{7369f8eebb1932ec08ac9ef73ba11654},\\
&\hashsplit{71cc2293c4ad5ba50d216c8f939cd1ec}{604ce7b86d6a19e9be1fcb27adc421fe}.
\end{aligned}
\]
The verifier also checks $N_D(b)+2\le2b$ in every row, with minimum slack
$1$.  Thus \cref{lem:one-sided-moment-pushforward} proves the assertion for
every odd $n\ge N_D(b)$.

It remains to reach these onsets.  The exact GMP/OpenMP replay
\[
\texttt{certificates/post\_m29/
post\_m29\_d53\_295\_short\_frontier\_gmp\_machine\_c\_certificate.log}
\]
reads each onset exactly once, rejects an even or missing onset, and
independently checks $N_D(b)+2\le2b$.  It expands every Chain difference in
the stable moments and the correction variables and uses
\cref{lem:d-frontier-stable-envelope,cor:d-frontier-bridge} together with
\cref{lem:d-determinant-component}.3.  It verifies every active step
$m=30,\ldots,163$ for all $D_{53},\ldots,D_{295}$, reports
\texttt{failures=0}, and records minimum $A$-free slack $1$.  The worst
positive frontier is $m=50$, odd target $103$, frontier $D_{53}$, with
logarithmic lower bound
\[
 167.515421966088354.
\]
The log and source hashes are
\[
\begin{aligned}
&\hashsplit{876d3f752d03b8f35183f0b28af7c15e}{9e28d104396fc65c215e5d854d8ae31f},\\
&\hashsplit{444cf4dbe59d5000a3315f31054ef07e}{c1cf710edf91a7c685f973644d65ad74}.
\end{aligned}
\]

The stable edge and \cref{prop:d-a-free-row-gated-bridge} give
$\Qthree^{D_b}(61)>0$.  The replayed Chain differences and
\cref{lem:chain-propagation} reach $N_D(b)$, after which the monotone direct
tail just proved supplies every later odd exponent.  This proves the
proposition.
\end{proof}

\begin{proposition}[Unconditional $D_{19}$ closure]
\label{prop:post29-d19-exact}
For every odd integer $n\ge s(D_{19})+2=17$, one has
\[
 \Qthree^{D_{19}}(n)\ge0.
\]
\end{proposition}

\begin{proof}
\ifginibrefullproof
\begingroup\footnotesize
For the monotone direct tail, use the rational geometry and layered tilts in
the proof of \cref{prop:post29-d-a-free}, but take polynomial degrees
$(k,m)=(6,10)$, square-trace parameter $\alpha=9/2$, and square-tilt fraction
$1$.  The required push moments have degrees at most $2k+2m=32$.
The exact C++/GMP source replay computes the finite-$D_{19}$ adjoint moments
$m_{19},\ldots,m_{34}$ from the Pieri sums in
\cref{lem:d-finite-moment-decomposition}; the lower moments are stable by
the same lemma.  Formula \eqref{lem:exact-push-from-adjoint}.1 then gives
every required $K_q^{D_{19}}$ exactly.  The sixteen status-$0$ machine-C
source logs and the two final certificates are authenticated by
\[
 \texttt{certificates/post\_m29/d19\_exact\_mgf\_machine\_c.sha256},
 \qquad
 \operatorname{SHA256}=
 \hashsplit{ead8966c1769cbaaaf6d7a6928df831d}{65887f5286d76c0fa18a19fcceb05d5b}.
 \tag{\ref{prop:post29-d19-exact}.1}
\]
The exact-moment generator source has SHA-256
\[
 \hashsplit{bc22efcafafe2d40b5f5430cb8f9cb04}{6bc6b9b188da684e61dc6fdfaa27077e}.
\]

The directed $384$-bit C++/MPFR verifier applies
\cref{lem:d-exact-even-orthogonal-mgf,lem:d-exact-square-mgf,%
lem:exact-mgf-layered-trace-tilt,lem:one-sided-moment-pushforward}.
At odd exponent $63$ it reports exact defining- and square-determinant pivot
lower bounds $0.8516811782\ldots$ and $0.8449575175\ldots$ and the strict
natural-log margin
\[
 6.7370668175795237151\ldots>0.
\]
Its log and source hashes are
\[
\begin{aligned}
&\hashsplit{42e987f476bacc0e7422adb382556222}{7aa0511a58a6984f61adcf4c08b3e4b5},\\
&\hashsplit{1f296ce3a32756fbee27bdc44eb9749d}{03dc88e74107eb7f15d91f8fda3a15e5}.
\end{aligned}
\tag{\ref{prop:post29-d19-exact}.2}
\]
Thus the monotonicity clause of
\cref{lem:one-sided-moment-pushforward} supplies every odd $n\ge63$.

It remains to reach $63$.  For $19\le j\le38$, the accepted determinant
rows are the full corrections because $A_j(19)=0$ by
\cref{lem:d-determinant-component}.  At $j=39$ the machine-C source replay
instead records the full identity
\[
 B_{39}(19)-A_{39}(19)
 =2849263974779776701152303067614034147.
 \tag{\ref{prop:post29-d19-exact}.3}
\]
The exact-GMP verifier parses legacy determinant rows only through $j=38$,
parses \eqref{prop:post29-d19-exact}.3 as a full correction, and uses the
proved two-sided box of \cref{lem:d-two-sided-correction-box} at every other
index.  It checks all $24$ Chain steps $m=7,\ldots,30$, reports no failure,
and has minimum exact lower bound
\[
 36778114931645730>0.
\]
The bridge log and source hashes are
\[
\begin{aligned}
&\hashsplit{a0a46e0274a2fa4ded548edfbec2f2d8}{3cc7d5cecc13be155aa37d77fe9e11a0},\\
&\hashsplit{edec7e651f112a342909230068651c9a}{59dc44a58bffcb5fa9e0420bf9f7ad8b}.
\end{aligned}
\tag{\ref{prop:post29-d19-exact}.4}
\]
The stable edge supplies the odd exponents through $s(D_{19})=15$;
these positive Chain differences and \cref{lem:chain-propagation} reach
$63$, where the monotone direct tail begins.
\endgroup
\else
\begingroup\scriptsize
By \cref{lem:d-finite-moment-decomposition}, the status-$0$ machine-C
C++/GMP logs give the adjoint moments through degree $34$; expanding
$(X-Y)^2(X+Y)^q$ determines exactly the push moments through degree $32$.
With $(k,m)=(6,10)$, the directed
$384$-bit MPFR verifier applies
\cref{lem:d-exact-even-orthogonal-mgf,lem:d-exact-square-mgf,%
lem:exact-mgf-layered-trace-tilt,lem:one-sided-moment-pushforward} and proves
the monotone tail from odd $n=63$, with strict log margin
$6.7370668175795237151\ldots$.  For the bridge,
\cref{lem:d-determinant-component} makes the determinant rows full
corrections through $j=38$; the source replay supplies the full
$B_{39}-A_{39}$ value.  With the remaining boxes of
\cref{lem:d-two-sided-correction-box}, the exact-GMP verifier checks all $24$
steps $m=7,\ldots,30$, with minimum lower bound
$36778114931645730>0$.  The identity ledger and
\texttt{d19\_exact\_mgf\_machine\_c.sha256} authenticate the two status-$0$
outputs and every exact input.  Finally,
\cref{lem:classical-stable-edge,lem:chain-propagation} supplies the prefix
through $15$ and propagates these steps to $63$.
\endgroup
\fi
\end{proof}

\begin{proposition}[Unconditional $D_{25}$--$D_{52}$ closure]
\label{prop:post29-d25-52-exact}
For every $25\le b\le52$ and every odd integer $n\ge s(D_b)+2$, one has
\[
 \Qthree^{D_b}(n)\ge0.
\]
\end{proposition}

\begin{proof}
\begingroup\footnotesize
First consider the monotone direct tails.  Use the rational parameters and
tilt fractions in the proof of \cref{prop:post29-d-a-free}; only the MGF
supplier and the following polynomial pairs $(k_b,m_b)$ change:
\[
 (5,6),\ (4,7),\ (5,7),\ (5,7)
 \quad\text{at }b=25,26,27,28,
 \qquad (k_b,m_b)=(5,8)\quad(29\le b\le52).
 \tag{\ref{prop:post29-d25-52-exact}.1}
\]
Thus $2k_b+2m_b\le b-3$ in every row, and
\cref{lem:stable-finite-rank-push-moments} supplies the exact polynomial
$H_{k_b,m_b}^{D_b}(a_b)$.

For $25\le b\le28$, the defining-trace lower bound is obtained from
\cref{lem:d-exact-even-orthogonal-mgf,lem:exact-mgf-layered-trace-tilt}; the
trace law is symmetric because $-I\in\SO(2b)$.  For $29\le b\le52$, it is
obtained from
\cref{lem:extended-cjl-trace-mgf,lem:layered-trace-tilt}.  In every row,
\cref{lem:d-exact-square-mgf} supplies both square-trace upper bounds.  The
directed $384$-bit MPFR verifier evaluates each Bessel series through term
$180$, adds the positive geometric remainder bound displayed in the proof of
\cref{prop:post29-bc-layered-mgf}, and uses interval Cholesky decomposition
for every exact determinant.  It verifies all geometry, pivot, Fredholm,
nested-mass, polynomial, and final comparison inequalities.

Let $T_b$ be the resulting odd onset.  The exact list is
\ifginibrefullproof
\[
\begin{aligned}
(T_{25},\ldots,T_{38})
 &=(71,73,71,73,73,73,75,77,79,79,81,83,83,85),\\
(T_{39},\ldots,T_{52})
 &=(85,87,89,89,91,93,93,95,97,97,99,99,101,103).
\end{aligned}
\tag{\ref{prop:post29-d25-52-exact}.2}
\]
\else
recorded row by row in the machine-C log cited below.
\fi
At $T_b$ the verifier proves the strict inequality
\[
 P_{D_b}(c_b)>
 N_{D_b}(a_b)\left(\frac{2b}{c_b}\right)^{T_b}
 +\frac{a_b^{T_b-2k_b}}{4^{m_b}c_b^{T_b}}
 H_{k_b,m_b}^{D_b}(a_b).
 \tag{\ref{prop:post29-d25-52-exact}.3}
\]
The conservative factor $2b/c_b$ dominates the actual endpoint factor also
when $b$ is odd.  Hence \cref{lem:one-sided-moment-pushforward} proves every
odd exponent at least $T_b$.  The machine-C log
\[
 \texttt{certificates/post\_m29/
 post\_m29\_d25\_52\_exact\_mgf\_mpfr\_machine\_c\_certificate.log}
\]
checks $28$ rows with no failures.  Its minimum final natural-log margin is
\[
 0.01375109087492612411266244772>0
\]
at $D_{39}$.  The exact defining- and square-determinant pivot minima are,
respectively, $0.8522815995\ldots$ and $0.8481405784\ldots$.  The log and source hashes
are respectively
\[
\begin{aligned}
&\hashsplit{f87ab4bce0b189878ef4c3528ce34687}{1721bff05cb00ccb553a6c4e84e1f049},\\
&\hashsplit{71cc2293c4ad5ba50d216c8f939cd1ec}{604ce7b86d6a19e9be1fcb27adc421fe}.
\end{aligned}
\tag{\ref{prop:post29-d25-52-exact}.4}
\]

It remains to connect the stable edge to $T_b$.  Expand each Chain difference
in the correction variables as in \cref{lem:d-frontier-stable-envelope}.
The exact C++/GMP/OpenMP verifier substitutes the boxes of
\cref{lem:d-two-sided-correction-box}; for every bilinear term it takes the
minimum over the four interval endpoints, and for a diagonal term it also
accounts for an interval containing zero.  It checks every Chain index
\[
 \left\lfloor\frac{b-4}{2}\right\rfloor
 \le m\le\frac{T_b-3}{2}
 \tag{\ref{prop:post29-d25-52-exact}.5}
\]
for every $25\le b\le52$.  These are exactly the steps whose target exponents
run from $s(D_b)+2$ through $T_b$.  The replay performs $707$ exact integer
checks and reports no failure.  Its smallest exact lower bound is
\[
 4695986280539928492451072>0
\]
at $D_{25},m=10$.  The machine-C log
\[
 \texttt{certificates/post\_m29/
 post\_m29\_d25\_52\_correction\_interval\_gmp\_machine\_c\_certificate.log}
\]
and its C++ source have respective SHA-256 hashes
\[
\begin{aligned}
&\hashsplit{e47196c8097fb2a087b4673d47a5e34c}{aae1cd1a316058759e17b13708e2d888},\\
&\hashsplit{9a6cf4da5b866d3b8fc2ce6d1d1a442f}{d3b419702fda2296da83d2fb788379e4}.
\end{aligned}
\tag{\ref{prop:post29-d25-52-exact}.6}
\]
The stable edge, these positive Chain differences, and
\cref{lem:chain-propagation} reach $T_b$; the direct tail then supplies every
later odd exponent.
\endgroup
\end{proof}

\begin{definition}[Legacy finite type-$D$ envelope]
\label{def:d-interval-envelope}
For $4\le b\le52$, let $T_b$ be the first odd exponent of the monotone
pushforward tail in the accepted row certificate cited in
\cref{prop:post29-direct}.  Let $\DEnv$ denote the following finite
statement: for $4\le b\le52$ and every $b\le j\le T_b+2$, the quantities of
\cref{lem:d-finite-moment-decomposition} satisfy
\[
  0\le B_j(b)-A_j(b)\le s_j.
\]
It contains only finitely many integer inequalities and no asymptotic
assertion.
\end{definition}

\begin{proposition}[The legacy envelope is false]
\label{prop:denv-counterexample}
The statement $\DEnv$ of \cref{def:d-interval-envelope} is false.
Specifically, at $b=4$ and $j=18$,
\[
\begin{gathered}
 s_{18}=658906772228668,
 \qquad A_{18}(4)=369382858189470,\\
 B_{18}(4)=288601773705246,
 \qquad B_{18}(4)-A_{18}(4)=-80781084484224<0.
\end{gathered}
\]
\end{proposition}

\begin{proof}
The exact C++/GMP Pieri replay on machine C reports
\[
 m_{18}^{D_4}=578125687744444
 =s_{18}-369382858189470+288601773705246.
\]
The archived log
\[
\texttt{certificates/post\_m29/
d4\_j18\_denv\_counterexample\_gmp\_machine\_c.log}
\]
has SHA-256
\[
\hashsplit{463e3c9632bb92936620fa65c03d3ff5}{abdb86940624f1621ffa976156d7d02f},
\]
and the C++ source has SHA-256
\[
\hashsplit{bc22efcafafe2d40b5f5430cb8f9cb04}{6bc6b9b188da684e61dc6fdfaa27077e}.
\]
All displayed identities are exact integer identities.  The negative
correction contradicts the lower inequality in $\DEnv$.
\end{proof}

\begin{proposition}[Exact $D_4$--$D_{24}$ row-gated bridge]
\label{prop:d-low-exact-bridge}
For every $4\le b\le24$, the type-$D_b$ row satisfies
\[
 D_{D_b}(2m+1)\ge0
\]
at every row-gated Chain step $0\le m\le29$.
\end{proposition}

\begin{proof}
The exact signed adjoint-moment sources are generated from the two Pieri sums
in \cref{lem:d-finite-moment-decomposition}.  They supply the following full
moment prefixes:
\[
\begin{array}{c|ccccccc}
b&4&5&6,7&8,\ldots,17&18&19&20,\ldots,24\\
\hline
\text{exact through }m_j&59&44&40&38&38&34&30.
\end{array}
\tag{\ref{prop:d-low-exact-bridge}.1}
\]
At $D_4$, five distinct deterministic-primality-checked $63$-bit moduli have
product greater than $28^{59}$.  The Chinese remainder reconstruction is
therefore the unique integer in the character bound
$0\le m_j^{D_4}\le28^j$ for every $0\le j\le59$.  It agrees with two
independent archived reconstruction summaries through $m_{59}$ and with the
full Pieri source at all $37$ overlapping indices $m_4,\ldots,m_{40}$.

For every index after the prefix in
\eqref{prop:d-low-exact-bridge}.1, the verifier uses the proved intervals of
\cref{lem:d-two-sided-correction-box}.  It expands each Chain difference
exactly as an integer polynomial in the corrections; a bilinear term is
minimized over all four interval endpoints, and a diagonal term also tests
zero when its interval contains zero.  The fixed-mode GMP/OpenMP replay
\[
\texttt{certificates/post\_m29/
post\_m29\_d4\_24\_exact\_prefix\_interval\_gmp\_certificate.log}
\]
checks all $530$ row-gated steps and reports \texttt{failures=0}.  Its
smallest exact lower endpoint is $34>0$, at $D_5,m=0$.  The separate D4 CRT
replay reports five primes, $60$ reconstructed moments, $37$ full-source
cross-checks, and no failure.  Every signed source log records exactly one
zero status, machine-C host \texttt{nb1cb2f}, and source SHA-256
\[
\hashsplit{bc22efcafafe2d40b5f5430cb8f9cb04}{6bc6b9b188da684e61dc6fdfaa27077e}.
\]
Before a correction is used, the C++ loader verifies both
$O_j^{\rm even}+O_j^{\det}=\delta_j$ and $s_j+\delta_j=m_j$.
The parent accepted-artifact manifest authenticates both replay outputs and
all source logs.  Thus every exact or interval-substituted integer lower
endpoint used in the stated rows is positive.
\end{proof}

\begin{proposition}[Exact $D_{10}$--$D_{24}$ terminal tails]
\label{prop:post29-d10-24-exact-tail}
For every $10\le b\le24$ and every odd integer $n\ge63$,
\[
 \Qthree^{D_b}(n)\ge0.
\]
\end{proposition}

\begin{proof}
First consider $D_{10}$.  Put $d_{10}=\dim D_{10}$ and use the following
rational cap and squared-Chebyshev data:
\[
\begin{array}{c|rrrrrrrrr}
b&d_b&C_{D_b}&R_b&\kappa_b&M_b&\ell_b&S_b&a_b&J_b\\
\hline
10&190&10&107&18&12400&9&4650&14&38.
\end{array}
\tag{\ref{prop:post29-d10-24-exact-tail}.1}
\]
Here $\sum_{\alpha>0}\alpha(\theta)^2=(2b-2)|\theta|^2$, which gives the
displayed values of $\kappa_b$.  More explicitly, if
$p_{2j}=\sum_i\theta_i^{2j}$, expansion over the positive roots
$e_i\pm e_j$ gives
\[
 \sum_{\alpha>0}\alpha(\theta)^4
 =6p_2^2+(2b-8)p_4\le\kappa_b p_2^2
\]
and
\[
 \sum_{\alpha>0}\alpha(\theta)^6
 =30p_2p_4+(2b-32)p_6\le\kappa_b p_2^3.
\]
For the second inequality, after putting $y_i=\theta_i^2/p_2$, the
difference between the right and left sides, divided by $p_2^3$, is
\[
 \sum_i y_i(1-y_i)\bigl((2b-2)+(2b-32)y_i\bigr)\ge0
\]
for $b=10$.  On the cap
$Q(\theta)=\kappa_b|\theta|^2\le R_b$, these formulas also give
$|\alpha(\theta)|^2\le2R_b/\kappa_b$.

The cap width is $w=6/5$ and the positive threshold is
$c_{10}=d_{10}-R_{10}-w=409/5$.  Weyl integration in exponential
coordinates \cite[Chapter~IV, (1.11)]{BrockertomDieck}, the
Macdonald--Mehta normalization \cite[Theorem~3.1(i)]{Etingof2009}, and
\[
 \prod_{\alpha>0}
 \left(\frac{\sin(\alpha(\theta)/2)}{\alpha(\theta)/2}\right)^2
 \ge M_b^{-1}
\]
give an exact rational lower bound for $P_{D_b}(c_b)$.  The checker proves
the sine-product inequality by the rational exponential majorant
\[
 \frac{R_b}{12}+\frac{R_b^2}{1440\kappa_b}
 +\frac{R_b^3}{90720\kappa_b^2(1-2R_b/(39\kappa_b))},
\]
using $90$ Taylor terms.  All remaining cap,
quartic-tail, monotonicity, Weyl-order, and factorial comparisons are exact
integer or rational inequalities.

For completeness, the displayed sine exponent follows directly from the
Euler product for $\sin z/z$.  The quadratic, quartic, and sextic terms are
bounded by
\[
 \frac{R_b}{12},\qquad
 \frac{R_b^2}{1440\kappa_b},\qquad
 \frac{R_b^3}{90720\kappa_b^2},
\]
respectively.  Since $\pi^2>39/4$, each remaining Euler-product term is at
most $2R_b/(39\kappa_b)$ times its predecessor; the geometric sum is the
third term in the stated majorant.  Also
$R_b<\kappa_b\pi^2$, so the cap is contained in one eigenangle chart.
Writing $(d_i)=(2,4,\ldots,2b-2,b)$, its resulting probability lower bound
is
\[
 \frac1{|W(D_b)|M_b}
 \frac{\prod_i d_i!}{(2\pi)^{b/2}\Gamma(d_b/2+1)}
 \left(\frac{R_b}{2\kappa_b}\right)^{d_b/2}.
 \tag{\ref{prop:post29-d10-24-exact-tail}.2}
\]
The verifier replaces the powers of $\pi$ by an explicit rational lower
bound.  Namely, $\pi<22/7$ gives
\[
 \frac1{(2\pi)^5\Gamma(96)}>
 \frac1{(44/7)^5\,95!}
 \tag{\ref{prop:post29-d10-24-exact-tail}.2a}
\]
for $D_{10}$.  The checker uses this bound only in the direction decreasing
the cap probability.

Let $Y$ be an independent adjoint character.  By
\cref{lem:d-finite-moment-decomposition}, its first four moments are stable;
since
$\mathbf EY=0$, $\mathbf EY^2=\mathbf EY^3=1$, and
$\mathbf EY^4=6$, for $r=1/2$, $s=15/4$ one has
\[
 \mathbf E(Y-r)^2(Y-s)^2
 =6-2(r+s)+r^2+s^2+4rs+r^2s^2=:V_4.
\]
If $y<-w$ and $x\ge d_b-R_b\ge r$, then
\[
 \frac{x-y}{x+w}\le
 \frac{(r-y)(s-y)}{(r+w)(s+w)}.
\]
Thus the lost $(X-Y)^2$-weight is at most
$V_4(x+w)^2/((r+w)^2(s+w)^2)$.  The verifier checks that the resulting
quadratic function of $x$ is positive and increasing from
$x=d_b-R_b$ onward.  This proves that the cap contribution is a lower bound
for $P_{D_b}(c_b)$, rather than merely for unweighted Haar mass.

For a real variable $u$, define
\[
 R_b^{\rm Ch}(u)=
 T_{\ell_b}\!\left(1-\frac{4d_b}{S_b}u+\frac2{S_b}u^2\right)^2.
\]
For $u\le-a_b$, the quantity $u(u-2d_b)$ is at least
$a_b(a_b+2d_b)$; since $T_{\ell_b}(z)^2$ increases for $z\ge1$, this
polynomial is at least $R_b^{\rm Ch}(-a_b)>0$.  Hence
\[
 N_{D_b}(a_b)\le
 \frac{\int R_b^{\rm Ch}(u)\,d\psi_{D_b}(u)}
      {R_b^{\rm Ch}(-a_b)}.
 \tag{\ref{prop:post29-d10-24-exact-tail}.3}
\]
The numerator has degree $4\ell_b$ and is an exact linear combination of
$K_0^{D_b},\ldots,K_{4\ell_b}^{D_b}$.  By
\cref{lem:exact-push-from-adjoint}, these are reconstructed from the full
adjoint moments through $m_{J_b}^{D_b}$.  The machine-C C++/GMP sources
provide the gap-free range $m_{10},\ldots,m_{38}$ and verify, row by row,
\[
 m_j^{D_b}=s_j+O^{\rm even}_j(D_b)+O^{\det}_j(D_b).
 \tag{\ref{prop:post29-d10-24-exact-tail}.4}
\]
The exact-rational C++/GMP replay then proves
\[
 P_{D_b}(c_b)>
 N_{D_b}(a_b)\left(\frac{2C_{D_b}}{c_b}\right)^{63}
 +2\left(\frac{a_b}{c_b}\right)^{63}.
 \tag{\ref{prop:post29-d10-24-exact-tail}.5}
\]
Its output is
\[
 \texttt{certificates/post\_m29/
 post\_m29\_d10\_exact\_tail\_gmp\_certificate.log}.
\]

For the exact-MGF rows $11\le b\le24$, use the common parameters
\[
\begin{gathered}
 c_b=\frac{20674}{10000}b,
 \quad a_b=\frac{429465}{100000}\sqrt b,
 \quad d_0=\frac{13501}{10000},
 \quad l=\frac{103311}{100000},\\
 u_i=\frac{139822}{100000}
      +i\frac{31205}{1000000}\quad(0\le i<8),
 \qquad f_-=f_i=f_{\rm sq}=1,
 \qquad f_{\mathrm{neg}}=\frac{101331}{100000},
\end{gathered}
\tag{\ref{prop:post29-d10-24-exact-tail}.6}
\]
Here the four \(f\)'s are respectively the lower and upper defining-trace
tilt fractions and the positive and negative square-trace tilt fractions.
Use the following rank-dependent polynomial and square-trace parameters:
\[
\begin{array}{c|c|c}
b&(k_b,m_b)&q_b/\sqrt b\\
\hline
11&(3,7)&24/5\\
12,13,14&(4,6)&19/4\\
15,16&(4,6)&9/2\\
17&(5,5)&441721/100000\\
18&(4,7)&441721/100000\\
19&(5,6)&441721/100000\\
20,21&(5,7)&441721/100000\\
22&(5,8)&441721/100000\\
23&(5,8)&17/4\\
24&(5,9)&441721/100000.
\end{array}
\tag{\ref{prop:post29-d10-24-exact-tail}.7}
\]
The source loader first requires one status-zero machine-C provenance line,
the pinned C++/GMP generator hash, and both exact identities in
\eqref{prop:post29-d10-24-exact-tail}.4 for every consumed moment.  It then
uses \cref{lem:exact-push-from-adjoint} to reconstruct the finite-rank
$K_j^{D_b}$ needed by $H_{k_b,m_b}^{D_b}(a_b)$.

Put $n_{11}=75$ and $n_b=63$ for $12\le b\le24$.  The exact
even-orthogonal defining-trace determinant of
\cref{lem:d-exact-even-orthogonal-mgf}, the exact square-trace determinant of
\cref{lem:d-exact-square-mgf}, and the layered tilt of
\cref{lem:exact-mgf-layered-trace-tilt} supply the positive trace tail.
The same square-trace determinant and
\eqref{lem:decoupled-trace-tail}.3 supply $N_{D_b}(a_b)$.  With directed
$384$-bit MPFR intervals, the verifier proves at $n=n_b$
\[
 P_{D_b}(c_b)>
 N_{D_b}(a_b)\left(\frac{2b}{c_b}\right)^{n_b}
 +\frac{a_b^{n_b-2k_b}}{4^{m_b}c_b^{n_b}}
  H_{k_b,m_b}^{D_b}(a_b).
 \tag{\ref{prop:post29-d10-24-exact-tail}.8}
\]
The factor $2b/c_b$ conservatively dominates $2C_{D_b}/c_b$.
For $D_{11}$, the machine-C replay reports defining- and square-MGF pivot
lower bounds $0.8426844977700063107\ldots$ and
$0.9814906225170877019\ldots$, and final natural-log margin
$0.091132792490649002559\ldots>0$.  Its output is
\[
 \texttt{certificates/post\_m29/
 post\_m29\_d11\_exact\_mgf\_mpfr\_machine\_c\_certificate.log}.
\]
The separate exact C++/GMP/OpenMP bridge consumes the gap-free machine-C
corrections through $m_{46}^{D_{11}}$, checks all $34$ Chain steps from the
stable edge through odd exponent $75$, and reports no failure.  Its minimum
exact lower endpoint is $11976538>0$.  Its output is
\[
 \texttt{certificates/post\_m29/
 post\_m29\_d11\_onset75\_bridge\_gmp\_machine\_c\_certificate.log}.
\]
By the stable edge and \cref{lem:chain-propagation}, the bridge proves the
required odd exponents through $75$; the exact-MGF comparison and
\cref{lem:one-sided-moment-pushforward} prove every odd exponent from $75$
onward.

For $D_{12},\ldots,D_{24}$, the replay checks all $13$ rows with no failure.
Its minimum defining- and square-MGF Cholesky pivots are respectively
$0.8380062150\ldots$ and $0.8377469953\ldots$, and its minimum final
natural-log margin is
\[
 0.0920085575002693989685\ldots>0
\]
at $D_{24}$.  The output is
\[
 \texttt{certificates/post\_m29/
 post\_m29\_d12\_24\_exact\_mgf\_mpfr\_certificate.log}.
\]
Finally, \cref{lem:pushforward-tail} propagates
\eqref{prop:post29-d10-24-exact-tail}.5 for $D_{10}$, while
\cref{lem:one-sided-moment-pushforward} propagates
\eqref{prop:post29-d10-24-exact-tail}.8 for $D_{11},\ldots,D_{24}$.
\end{proof}

\begin{lemma}[D stable-envelope frontier monotonicity]\label{lem:d-frontier-stable-envelope}
Fix $m\ge31$.  Let $s_j=m_j^{\mathrm{st}}$ be the stable $D$-row
moments, let $\Delta_{\mathrm{st}}(m)$ be the stable Chain difference, and
let $L_{m,j}$ and $Q_{m,i,j}$ be the integer coefficients in the expansion
\[
  D_G(2m+1)
  =
  \Delta_{\mathrm{st}}(m)
  +\sum_j L_{m,j}\delta_j^G
  +\sum_{i\le j}Q_{m,i,j}\delta_i^G\delta_j^G,
\]
where $m_j^G=s_j+\delta_j^G$.  For an integer rank boundary $r$, define
\[
  F_m(r)=
  \Delta_{\mathrm{st}}(m)
  +\sum_{\substack{j\ge r\\ L_{m,j}<0}}L_{m,j}s_j
  +\sum_{\substack{i\le j,\ i,j\ge r\\ Q_{m,i,j}<0}}
       Q_{m,i,j}s_i s_j .
\]
Then $F_m(r')\ge F_m(r)$ whenever $r'\ge r$.  Moreover, if a type $D_b$
row has determinant corrections satisfying
\[
  \delta_j^{D_b}=0\quad(j<b),\qquad
  0\le \delta_j^{D_b}\le s_j\quad(j\ge b),
\]
then
\[
  D_{D_b}(2m+1)\ge F_m(b).
\]
\end{lemma}

\begin{proof}
The difference $F_m(r')-F_m(r)$ is the negative of the sum of those displayed
linear and quadratic summands whose indices lie in the set removed when the
boundary is raised from $r$ to $r'$.  Every removed displayed summand is
nonpositive because $s_j\ge0$ and only negative coefficients are displayed;
hence the difference is nonnegative.

For the row bound, expand $D_{D_b}(2m+1)$ in the corrections
$\delta_j^{D_b}$.  Terms with positive coefficient are nonnegative and may be
discarded in a lower bound.  If a negative linear coefficient occurs, then
$0\le\delta_j^{D_b}\le s_j$ gives
$L_{m,j}\delta_j^{D_b}\ge L_{m,j}s_j$.  If a negative quadratic coefficient
occurs, then the product vanishes unless both indices are at least $b$, and
$0\le\delta_i^{D_b}\delta_j^{D_b}\le s_i s_j$ gives
$Q_{m,i,j}\delta_i^{D_b}\delta_j^{D_b}\ge Q_{m,i,j}s_i s_j$.
Adding these lower bounds gives $D_{D_b}(2m+1)\ge F_m(b)$.
\end{proof}

\begin{corollary}[D frontier bridge reduction]\label{cor:d-frontier-bridge}
Let $\mathcal I$ be a finite interval of type $D$ ranks for which the
stable-envelope hypotheses of \cref{lem:d-frontier-stable-envelope} hold.
For a fixed $m$, let $r_m$ be the smallest rank in $\mathcal I$ whose
finite bridge still includes the Chain step $2m+1$.  If $F_m(r_m)>0$, then
\[
  D_{D_b}(2m+1)>0
\]
for every active rank $b\in\mathcal I$ at that Chain step.
\end{corollary}

\begin{proof}
Every active rank satisfies $b\ge r_m$.  By
\cref{lem:d-frontier-stable-envelope},
$D_{D_b}(2m+1)\ge F_m(b)\ge F_m(r_m)>0$.
\end{proof}

\begin{proposition}[Post-$m=29$ $D$-interval finite tails]
\label{prop:post29-d-interval}
For $G=D_b$ with $98\le b\le275$ or $b=277,279$, one has
\[
  \Qthree^G(n)\ge0
\]
for every odd $n\ge N_{\mathrm{hit}}^{(29)}(G)$.
\end{proposition}

\begin{proof}
Every displayed rank lies in $53\le b\le295$, so this is the corresponding
subrange of \cref{prop:post29-d-a-free}.  The former $R=38b$ onset and long
frontier artifacts are retained only as superseded provenance; no theorem
now consumes their false stable-envelope hypothesis.
\end{proof}

\begin{proposition}[Post-$m=29$ $B/C$ interval direct tails]\label{prop:post29-bc-interval-direct}
For $G=B_b$ with $14\le b\le217$ or $G=C_b$ with
$20\le b\le217$, let $N_{\mathrm{dir}}^{BC}(G)$ be the odd onset recorded
for the row in
\[
\texttt{certificates/post\_m29/bc\_interval\_formula\_mpfr\_certificate.log}.
\]
Then
\[
  \Qthree^G(n)\ge0
\]
for every odd $n\ge N_{\mathrm{dir}}^{BC}(G)$.
\end{proposition}

\begin{proof}
The accepted C++/Optimus verifier
\[
\texttt{certificates/post\_m29/bc\_cheb80\_negative\_tail\_gmp\_certificate.log}
\]
first supplies the common negative-tail input used by the interval direct-tail
template.  It uses exact GMP rational arithmetic with the degree-$8$
Chebyshev majorant, scale $12000$, and cutoff $80$.  The log exits with
status $0$, reports \texttt{failures=0}, and has SHA-256
\[
\hashsplit{80fce36378ec2192dca2c0b616c083b4}{6168cf5134b66e73775156b995ea303b}.
\]
Its worst row is $B_{217}$, with exact upper mass log
$-15.1757329227092583<-15$.

The directed-rounding MPFR verifier
\[
\texttt{certificates/post\_m29/bc\_interval\_formula\_mpfr\_certificate.log}
\]
then checks the hypothesis of \cref{lem:pushforward-tail} at the logged row
onset for each of the $402$ rows.  It exits with status $0$, reports
\texttt{failures=0}, and has SHA-256
\[
\hashsplit{53530cfe3d4b1947a77bf45c5c5aaa56}{a052ec1cde93c0704b68e004f22f2cb5}.
\]
The largest logged onset is $N_{\mathrm{dir}}^{BC}(B_{217})=30897$.  The
smallest directed lower log margin at an onset, attained at $B_{20}$, is
\[
0.0061301728071582916875791>0.
\]
By the monotonicity clause of \cref{lem:pushforward-tail}, the same
inequality holds for every later odd exponent in the row.
\end{proof}

\begin{corollary}[Uniform residual $B/C$ direct-tail onset bound]\label{cor:post29-bc-direct-onset-bound}
For every row $G=B_b$ with $14\le b\le217$ or $G=C_b$ with
$20\le b\le217$,
\[
  N_{\mathrm{dir}}^{BC}(G)\le30897 .
\]
\end{corollary}

\begin{proof}
This is the largest-onset clause in the directed-rounding MPFR verifier used
in \cref{prop:post29-bc-interval-direct}: the largest logged onset is
$N_{\mathrm{dir}}^{BC}(B_{217})=30897$.
\end{proof}

\begin{corollary}[Initial residual $B$ direct-tail onset values]\label{cor:post29-bc-initial-b-onsets}
The directed-rounding MPFR verifier used in
\cref{prop:post29-bc-interval-direct} records
\[
  N_{\mathrm{dir}}^{BC}(B_{14})=925,\qquad
  N_{\mathrm{dir}}^{BC}(B_{15})=415,\qquad
  N_{\mathrm{dir}}^{BC}(B_{16})=329,\qquad
  N_{\mathrm{dir}}^{BC}(B_{17})=303 .
\]
\end{corollary}

\begin{proof}
These are the row entries for $B_{14}$, $B_{15}$, $B_{16}$, and $B_{17}$ in the
status-$0$ directed-rounding MPFR verifier log
\[
\texttt{certificates/post\_m29/bc\_interval\_formula\_mpfr\_certificate.log},
\]
whose SHA-256 hash and failure-free summary are cited in
\cref{prop:post29-bc-interval-direct}.
\end{proof}

\begin{proposition}[$B/C$ Pieri correction sign]\label{prop:bc-pieri-correction-sign}
For $G=B_b$ or $G=C_b$, put $s_j=m_j^{\mathrm{st}}$ and
$m_j^G=s_j+\delta_j^G$ in the first nonstable moment range $j\ge2b+2$.
Then
\[
  \delta_j^G\le0.
\]
More precisely, writing
\[
  e_2^j=\sum_\lambda c_j(\lambda)s_\lambda,\qquad
  h_2^j=\sum_\lambda d_j(\lambda)s_\lambda,
\]
the coefficients $c_j(\lambda)$ and $d_j(\lambda)$ are nonnegative integers,
and
\[
  \delta_j^{B_b}
  =
  -\sum_{\substack{\nu\vdash j\\ \ell(\nu)>2b+1}}
    c_j(2\nu_1,2\nu_2,\ldots),
\]
while
\[
  \delta_j^{C_b}
  =
  -\sum_{\substack{\mu\vdash j\\ \ell(\mu)>b}}
    d_j(\mu_1,\mu_1,\ldots,\mu_{\ell(\mu)},\mu_{\ell(\mu)}).
\]
\end{proposition}

\begin{proof}
For $B_b$, the adjoint character is $e_2=s_{(1,1)}$ in the defining
eigenvalues.  Pieri's rule expands $e_2^j$ by adding vertical two-strips, so
each $c_j(\lambda)$ is a nonnegative integer.  Rains' orthogonal
Schur-integral criterion \cite[\S3]{Rains1997} says that the stable
orthogonal integral keeps exactly the Schur summands whose row lengths are
even, while the finite $\SO(2b+1)$ integral keeps exactly those among them
with at most $2b+1$ nonzero rows.  The omitted stable summands are precisely
the partitions $(2\nu_1,2\nu_2,\ldots)$ with $\ell(\nu)>2b+1$, giving the
displayed negative sum.

For $C_b$, the adjoint character is $h_2=s_{(2)}$.  Pieri's rule expands
$h_2^j$ by adding horizontal two-strips, so each $d_j(\lambda)$ is
nonnegative.  Rains' symplectic Schur-integral criterion
\cite[\S3]{Rains1997} says that the stable symplectic integral keeps exactly
the Schur summands whose row lengths occur with even multiplicity, while the
finite $\Sp(2b)$ integral keeps exactly those among them with at most $2b$
nonzero rows.  The omitted summands are exactly
$(\mu_1,\mu_1,\ldots,\mu_{\ell(\mu)},\mu_{\ell(\mu)})$ with
$\ell(\mu)>b$, giving the second negative sum.
\end{proof}

\ifginibrefullproof

\begin{corollary}[Residual window-size bounds from the direct-tail table]\label{cor:post29-bc-window-size-strata}
In the residual $B/C$ windows of \cref{rem:former-post29-tail}, with $q=b+1$, the
quantity $j-q$ satisfies the following bounds:
\[
\begin{array}{c|c}
\text{range of }q & \text{upper bound for } j-q\\ \hline
q=15 & 912\\
16\le q\le17 & 401\\
q=18 & 287\\
19\le q\le20 & 283\\
21\le q\le22 & 303\\
23\le q\le24 & 337\\
25\le q\le28 & 423\\
q\ge29 & 30870 .
\end{array}
\]
\end{corollary}

\begin{proof}
The residual window ends at $j=N_{\mathrm{dir}}^{BC}(G)+2$.  For the finite
small-$q$ ranges, the displayed maxima of $N_{\mathrm{dir}}^{BC}(G)+2-q$ are
read directly from the status-$0$ directed-rounding MPFR verifier table cited
in \cref{prop:post29-bc-interval-direct}; the attaining rows are respectively
$B_{14},B_{15},B_{17},B_{19},B_{21},B_{23}$, and $B_{27}$, with the $C$ rows in
the same $q$-ranges no larger.  For $q\ge29$,
\cref{cor:post29-bc-direct-onset-bound} gives
\[
  j-q\le30899-q\le30870 .
\]
\end{proof}

\begin{lemma}[Rightmost even-column wall removal]\label{lem:bc-rightmost-even-column-wall-removal}
Let $q\ge2$, and let $\lambda$ be a partition such that $\lambda'$ is even
and $\lambda_1\ge q$.  Put $L=\lambda_1$, and define a new partition $\mu$ by
its column lengths:
\[
  \mu'_c=
  \begin{cases}
    \lambda'_c,&1\le c\le L-q,\\
    \lambda'_c-2,&L-q<c\le L,
  \end{cases}
\]
omitting zero columns at the right.  Then $\mu'$ is even and
$|\mu|=|\lambda|-2q$.  Moreover, the skew shape $\lambda/\mu$ admits an
iterated Pieri filling with content $(2^q)$ whose labels each occupy a
horizontal two-strip.  Consequently $s_\lambda$ occurs with positive
coefficient in $s_\mu h_2^q$; in particular the width-overflow Schur sum over
column-even $\lambda$ is Schur-supported inside
\[
  h_2^q\sum_{\nu':\,\mathrm{even}} s_\nu
\]
after this rightmost-column deletion.  Applying the involution $\omega$ gives
the corresponding row-even vertical-strip statement with $e_2^q$.
\end{lemma}

\begin{proof}
Since $\lambda'$ is even, every nonzero column length $\lambda'_c$ is a
positive even integer.  Thus the rightmost $q$ nonzero columns all have length
at least $2$, so the displayed definition is nonnegative.  The sequence
$(\mu'_c)$ is weakly decreasing: the unchanged part is weakly decreasing, the
rightmost part remains weakly decreasing after subtracting $2$ from each of
its entries, and at the boundary
\[
  \mu'_{L-q}=\lambda'_{L-q}\ge\lambda'_{L-q+1}\ge\lambda'_{L-q+1}-2
  =\mu'_{L-q+1}
\]
when $L>q$.  Hence $\mu$ is a partition.  All nonzero $\mu'_c$ are even, and
the total removed size is two boxes in each of $q$ columns, giving
$|\mu|=|\lambda|-2q$.

List the removed columns from left to right as
\[
  c_d=L-q+d\qquad(1\le d\le q).
\]
The two boxes removed from column $c_d$ lie in rows
$\lambda'_{c_d}-1$ and $\lambda'_{c_d}$.  Let
\[
  a_i=\left\lceil\frac{i}{2}\right\rceil
  \qquad(1\le i\le2q).
\]
Fill the upper removed box of column $c_d$ by $a_d$ and the lower removed box
by $a_{q+d}$.  The word
\[
  a_1,\ldots,a_{2q}=1,1,2,2,\ldots,q,q
\]
has content $(2^q)$.  Since $q\ge2$, one has $a_d<a_{q+d}$ for every
$d$, so columns are strictly increasing.

It remains to check rows.  If two upper removed boxes lie in the same row,
their columns have the same even height and their labels are
$a_d$ in increasing column order, hence weakly increase.  The same argument
applies to lower removed boxes, using the weakly increasing labels
$a_{q+d}$.  An upper removed row has odd index, while a lower removed row has
even index, because all column heights are even; hence an upper removed box
and a lower removed box cannot lie in the same row.  Thus the filling is
semistandard.

For each label $r$, the two boxes carrying $r$ lie in distinct columns,
because columns are strictly increasing.  Therefore the boxes of each label
form a horizontal two-strip.  Reading labels $1,\ldots,q$ gives an iterated
Pieri chain from $\mu$ to $\lambda$, so Pieri's rule gives a positive
coefficient of $s_\lambda$ in $s_\mu h_2^q$.  Summing over the possible
column-even $\mu$ gives the stated Schur-support containment.  Applying
$\omega(s_\alpha)=s_{\alpha'}$ exchanges horizontal and vertical two-strips
and gives the final row-even statement.
\end{proof}

\begin{lemma}[Column-even Littlewood graph model]\label{lem:bc-littlewood-graph-model}
Let
\[
  K(x)=\prod_{1\le a<b}(1-x_ax_b)^{-1}.
\]
For every $j\ge0$, the stable $B/C$ adjoint moment $s_j=m_j^{\mathrm{st}}$
appearing in \cref{prop:bc-pieri-correction-sign} is
\[
  s_j=[x_1^2\cdots x_j^2]\,K(x_1,\ldots,x_j).
\]
Equivalently, $s_j$ is the number of loopless labelled multigraphs on
$\{1,\ldots,j\}$ in which every vertex has degree $2$, with parallel edges
allowed.
\end{lemma}

\begin{proof}
Albion's classical Littlewood identity
\cite[Introduction, (1.1c)]{Albion2024} gives
\[
  \sum_{\lambda':\,\mathrm{even}} s_\lambda(x)
  =
  \prod_{1\le a<b}(1-x_ax_b)^{-1},
\]
where the paragraph following (1.1c) defines $\lambda'$ as the conjugate
partition and ``even'' as having only even parts.  For the stable $C$-side,
the Rains criterion used in \cref{prop:bc-pieri-correction-sign} gives
\[
  s_j=
  \left\langle h_2^j,\sum_{\lambda':\,\mathrm{even}}s_\lambda\right\rangle.
\]
For the stable $B$-side, the same criterion gives the Hall pairing of
$e_2^j$ with the row-even Schur sum; applying the standard involution
$\omega(s_\lambda)=s_{\lambda'}$ sends this to the same displayed pairing
with $h_2^j$ and the column-even Schur sum.  Since
$h_2^j=h_{(2^j)}$ is Hall-dual to the monomial symmetric function
$m_{(2^j)}$, this pairing is the coefficient of $m_{(2^j)}$ in $K$, hence
the coefficient of the monomial $x_1^2\cdots x_j^2$.

Finally, expanding
\[
  \prod_{1\le a<b\le j}(1-x_ax_b)^{-1}
  =
  \sum_{(m_{ab})_{a<b}}\prod_{a<b}(x_ax_b)^{m_{ab}}
\]
chooses a nonnegative edge multiplicity $m_{ab}$ for each unordered pair
$a<b$.  The coefficient of $x_1^2\cdots x_j^2$ is therefore exactly the
number of such multiplicity arrays with
$\sum_{b\ne a}m_{\min(a,b),\max(a,b)}=2$ for every $a$, i.e. the stated
loopless labelled $2$-regular multigraphs.
\end{proof}

\begin{corollary}[Direct full-wall overflow deletion has a parity-zero obstruction]\label{cor:bc-full-wall-overflow-obstruction}
For every $q\ge1$, the set of column-even semistandard tableaux of content
$(2^{2q+1})$ whose terminal first row has length at least $2q$ is nonempty,
whereas the stable one-label moment satisfies
\[
  s_1=0 .
\]
Consequently no direct full-wall estimate which bounds every such
width-overflow tableau at $j=2q+1$ by a finite multiple of $s_{j-2q}$ can be
true.
\end{corollary}

\begin{proof}
\Cref{lem:bc-row-packing-fails} constructs, for each $q\ge1$, a column-even
semistandard tableau of content $(2^{2q+1})$ and shape
$(2q+1,2q+1)$.  Its first row has length $2q+1$, hence at least $2q$, so the
displayed width-overflow family is nonempty.

By \cref{lem:bc-littlewood-graph-model}, $s_1$ counts loopless labelled
$2$-regular multigraphs on the one-label set.  There is no unordered pair
$a<b$ on one label, so the only loopless multigraph has degree $0$ at its
unique vertex.  Hence there is no loopless $2$-regular multigraph and
$s_1=0$.

At $j=2q+1$, deleting a full wall of $2q$ labels leaves the stable index
$j-2q=1$.  Any finite-multiple bound of the stated form therefore has right
side $0$, while the nonempty family above gives a positive left side.
\end{proof}

\begin{lemma}[Stable Littlewood moments satisfy a three-term recurrence]\label{lem:bc-stable-moment-recurrence}
Let $s_n$ be the stable moment sequence of
\cref{lem:bc-littlewood-graph-model}.  Put $s_{-1}=0$, $s_0=1$, and
$s_1=0$.  For every $n\ge1$,
\[
  s_{n+1}=n s_n+n s_{n-1}-\binom n2 s_{n-2}.
\]
\end{lemma}

\begin{proof}
By \cref{lem:bc-littlewood-graph-model}, $s_n$ counts loopless labelled
$2$-regular multigraphs on an $n$-element label set.  Such a graph is a set of
connected components.  A component of size $2$ is the unique doubled edge on
two labels.  A component of size $r\ge3$ is a simple unoriented cycle on its
$r$ labels, because any parallel edge would already use both degrees at its
two endpoints.

Thus the exponential generating function
\[
  A(z)=\sum_{n\ge0}s_n\frac{z^n}{n!}
\]
is
\[
  A(z)=
  \exp\left(
    \frac{z^2}{2}+\sum_{r\ge3}\frac{z^r}{2r}
  \right)
  =
  \exp\left(-\frac z2+\frac{z^2}{4}\right)(1-z)^{-1/2}.
\]
Therefore
\[
  \frac{A'(z)}{A(z)}
  =
  -\frac12+\frac z2+\frac{1}{2(1-z)}
  =
  \frac{z-\frac12 z^2}{1-z},
\]
or equivalently
\[
  (1-z)A'(z)=\left(z-\frac12 z^2\right)A(z).
\]
Comparing the coefficient of $z^n/n!$ gives
\[
  s_{n+1}-n s_n=n s_{n-1}-\binom n2 s_{n-2},
\]
which is the displayed recurrence.
\end{proof}

\begin{lemma}[Cumulative deleted-only graph counts have a factorial bound]\label{lem:bc-cumulative-deleted-only-factorial-bound}
For the stable moment sequence $s_n$ of
\cref{lem:bc-littlewood-graph-model} and every $K\ge0$,
\[
  \sum_{d=0}^K (d+1)s_d \le (K+2)! .
\]
\end{lemma}

\begin{proof}
By \cref{lem:bc-littlewood-graph-model}, $s_d$ counts loopless labelled
$2$-regular multigraphs on a fixed $d$-element label set.  Each connected
component is either the doubled edge on two labels or a simple unoriented
cycle on at least three labels.  Encode such a graph as a permutation by
writing each doubled edge as a transposition and each simple cycle as the
cycle whose least label is written first and whose second label is smaller
than its last label.  The components are disjoint, so this gives a
well-defined permutation, and the graph is recovered from the cycle
decomposition of that permutation by forgetting the chosen orientation of
cycles of length at least three and replacing every transposition by a doubled
edge.  Hence $s_d\le d!$.

Therefore
\[
  \sum_{d=0}^K(d+1)s_d
  \le
  \sum_{d=0}^K(d+1)d!
  =
  \sum_{r=1}^{K+1}r! .
\]
The last sum is at most $(K+2)!$: there are $K+1$ summands, each at most
$(K+1)!$, and $K+1<K+2$.
\end{proof}

\begin{lemma}[Loopless graph deficit repair]\label{lem:bc-graph-deficit-repair}
Let $\Gamma$ be a loopless labelled multigraph on a finite vertex set $V$ in
which every vertex has degree $2$, with parallel edges allowed.  Let
$S\subset V$, put $V_0=V\setminus S$, and let $\Gamma_0$ be the graph obtained
from $\Gamma$ by deleting $S$ and all incident edges.  For $v\in V_0$ put
\[
  c_v=2-\deg_{\Gamma_0}(v).
\]
If the nonzero deficits are not concentrated as $c_v=2$ at a single vertex,
then there exists a loopless multigraph $R$ on $V_0$ such that
\[
  \deg_R(v)=c_v\qquad(v\in V_0).
\]
Consequently $\Gamma_0\cup R$ is a loopless labelled $2$-regular multigraph
on $V_0$.
\end{lemma}

\begin{proof}
For each $v\in V_0$, the number $c_v$ is the number of edges of $\Gamma$
joining $v$ to a vertex of $S$, counted with multiplicity.  Thus
$c_v\in\{0,1,2\}$.  The sum
\[
  C=\sum_{v\in V_0}c_v
\]
is the number of edges from $S$ to $V_0$.  Since the total degree of the
vertices in $S$ is $2|S|$ and each edge internal to $S$ contributes $2$ to
that degree, $C$ is even.

Let $A=\{v:c_v=2\}$ and $B=\{v:c_v=1\}$.  Then $|B|$ is even.  If
$A=B=\varnothing$, take $R$ empty.  If $A=\varnothing$, pair the vertices of
$B$ arbitrarily and put one edge on each pair.

Assume next that $A$ is nonempty.  If $|A|$ is even, pair the vertices of
$A$ arbitrarily and put two parallel edges on each pair; also pair the
vertices of $B$ arbitrarily with one edge on each pair.  If $|A|$ is odd,
then the excluded concentration case is exactly $|A|=1$ and $B=\varnothing$.
When $|A|\ge3$, choose three vertices of $A$ and connect them by a triangle,
then pair the remaining vertices of $A$ with two parallel edges on each pair
and pair $B$ arbitrarily.  When $|A|=1$, the exclusion forces $|B|\ge2$;
connect the unique vertex of $A$ to two vertices of $B$, and pair the
remaining vertices of $B$ arbitrarily.

In every case the constructed graph $R$ has no loops and has degree $c_v$ at
each $v$.  Adding it to $\Gamma_0$ preserves looplessness and gives total
degree $\deg_{\Gamma_0}(v)+c_v=2$ at every vertex of $V_0$.
\end{proof}

\begin{lemma}[Canonical graph deficit repair away from the obstruction]\label{lem:bc-canonical-graph-deficit-repair}
In the setting of \cref{lem:bc-graph-deficit-repair}, assume the label set
$V_0$ is linearly ordered and that the nonzero deficits are not concentrated
as $c_v=2$ at a single vertex.  Then there is a canonically chosen loopless
multigraph
\[
  R_{\mathrm{can}}(\Gamma,S)
\]
on $V_0$, determined only by the ordered deficit word
$(c_v)_{v\in V_0}$, such that
\[
  \deg_{R_{\mathrm{can}}(\Gamma,S)}(v)=c_v\qquad(v\in V_0).
\]
Consequently
\[
  \Gamma_0\cup R_{\mathrm{can}}(\Gamma,S)
\]
is a loopless labelled $2$-regular multigraph on $V_0$, and the repair graph
adds exactly
\[
  \frac12\sum_{v\in V_0}c_v
\]
edges.
\end{lemma}

\begin{proof}
Let
\[
  A=\{v\in V_0:c_v=2\},\qquad B=\{v\in V_0:c_v=1\},
\]
each listed in the fixed increasing order.  As in the proof of
\cref{lem:bc-graph-deficit-repair}, $|B|$ is even.

If $A=B=\varnothing$, set $R_{\mathrm{can}}=\varnothing$.  If
$A=\varnothing$, pair the ordered list $B$ consecutively and put one edge on
each pair.

Suppose $A\ne\varnothing$.  If $|A|$ is even, pair the ordered list $A$
consecutively and put two parallel edges on each pair; also pair the ordered
list $B$ consecutively and put one edge on each pair.  If $|A|$ is odd and
$|A|\ge3$, take the first three vertices of $A$ and put a triangle on them,
then pair the remaining vertices of $A$ consecutively with two parallel edges
on each pair, and pair $B$ consecutively with one edge on each pair.  Finally,
if $|A|=1$, the excluded obstruction implies $|B|\ge2$; connect the unique
vertex of $A$ to the first two vertices of $B$, and pair the remaining
vertices of $B$ consecutively.

All choices are made by increasing order, so the construction is determined
by the ordered deficit word.  No loop is introduced: every edge joins two
distinct vertices in a pair, in a triangle, or between the unique vertex of
$A$ and a vertex of $B$.  The degree check is case-by-case from the
construction.  Vertices of $A$ receive degree $2$, vertices of $B$ receive
degree $1$, and all other vertices receive degree $0$.  This proves
$\deg_{R_{\mathrm{can}}}(v)=c_v$.

Since the sum of degrees of $R_{\mathrm{can}}$ is $\sum_v c_v$, the number of
added edges is half that sum.  Adding $R_{\mathrm{can}}$ to $\Gamma_0$ gives
degree $\deg_{\Gamma_0}(v)+c_v=2$ at every retained vertex and preserves
looplessness.
\end{proof}

\begin{lemma}[Deleted graph attachments are two-port data]\label{lem:bc-graph-attachment-port-code}
Let $\Gamma$ be a loopless labelled $2$-regular multigraph on a finite
linearly ordered label set $V$, and let $S\subset V$.  Put
$V_0=V\setminus S$.  For each $s\in S$, introduce two ordered ports
\[
  (s,1),\qquad (s,2).
\]
The multiset of all edges of $\Gamma$ with at least one endpoint in $S$ is
equivalently encoded by the following port datum:
\begin{enumerate}
\item a matching of some ports into unordered pairs
\[
  \{(s,a),(t,b)\}\qquad(s,t\in S,\ s\ne t),
\]
representing the edges internal to $S$; and
\item a map from every unmatched port to $V_0$, representing the cross-edge
endpoint of that port.
\end{enumerate}
For each $v\in V_0$, the number of unmatched ports mapped to $v$ is the
deficit
\[
  c_v=2-\deg_{\Gamma_0}(v)
\]
of \cref{lem:bc-graph-deficit-repair}.  Moreover this port datum is
canonically determined by the edge multiset using the ambient label order, and
it determines the edge multiset.
\end{lemma}

\begin{proof}
For each unordered pair of distinct labels $x<y$, let $m_{xy}$ be the edge
multiplicity and label the formal parallel copies by
$1,\ldots,m_{xy}$.  For every edge copy incident to $S$, record at each
endpoint $s\in S$ the opposite endpoint and this copy index.  Order the two
records at a fixed $s$ lexicographically by the opposite endpoint and then by
the copy index, and assign them to the ports $(s,1),(s,2)$ in that order.  If
an edge copy has both endpoints in $S$, pair the two ports assigned to its two
endpoint records.
Looplessness gives distinct labels $s\ne t$ for such a pair.  If an edge joins
$s\in S$ to $v\in V_0$, leave the assigned port unmatched and map it to $v$.
Since $\Gamma$ has degree $2$ at every $s\in S$, every port is used exactly
once.  The construction is canonical from the ordered labels and the
edge-copy ordering.

Conversely, from such a port datum, each matched pair
$\{(s,a),(t,b)\}$ gives one internal edge between $s$ and $t$, and each
unmatched port $(s,a)$ mapped to $v$ gives one cross edge between $s$ and
$v$.  This reconstructs the multiset of all edges incident to $S$.  Finally,
the cross edges incident to a retained vertex $v$ are exactly the edges lost
from $v$ when $S$ is deleted, so their number is
$2-\deg_{\Gamma_0}(v)=c_v$.
\end{proof}

\begin{proposition}[Canonical graph deficit splices are recovered from attachment data]\label{prop:bc-canonical-graph-splice-recovery}
Work in the setting of \cref{lem:bc-canonical-graph-deficit-repair}.  Let
\[
  \Gamma^\ast=\Gamma_0\cup R_{\mathrm{can}}(\Gamma,S)
\]
be the repaired loopless labelled $2$-regular multigraph on
$V_0=V\setminus S$.  Suppose a datum $A$ determines the multiset of all
edges of $\Gamma$ with at least one endpoint in $S$, split as
\[
  E_S(\Gamma)=\{\text{edges with both endpoints in }S\},
  \qquad
  E_\times(\Gamma)=\{\text{edges with one endpoint in }S
        \text{ and one endpoint in }V_0\}.
\]
Then $\Gamma$ is recovered from
\[
  S,\quad \Gamma^\ast,\quad A .
\]
In particular, in the non-obstructed graph-splice route the remaining
branch-data problem is exactly the recovery of the deleted internal edges and
cross-attachments incident to $S$, equivalently the two-port datum of
\cref{lem:bc-graph-attachment-port-code}; the canonical degree repair itself
carries no additional choice.
\end{proposition}

\begin{proof}
From $E_\times(\Gamma)$ we compute, for every $v\in V_0$,
\[
  c_v=\#\{\text{edges in }E_\times(\Gamma)\text{ incident to }v\},
\]
with multiplicity.  Therefore
\cref{lem:bc-canonical-graph-deficit-repair} reconstructs the same canonical
repair multigraph $R_{\mathrm{can}}(\Gamma,S)$ from the ordered deficit word
$(c_v)_{v\in V_0}$.

Since $\Gamma^\ast$ was formed as the multigraph union
$\Gamma_0\cup R_{\mathrm{can}}(\Gamma,S)$, subtracting these known
edge-multiplicities recovers $\Gamma_0$.  The original graph $\Gamma$ is the
disjoint union of the retained part $\Gamma_0$, the deleted internal edge
multiset $E_S(\Gamma)$, and the cross-edge multiset $E_\times(\Gamma)$.
All three pieces are now known, so $\Gamma$ is recovered.
\end{proof}

\begin{lemma}[Canonical graph repair is not a witness-free supplier]\label{lem:bc-canonical-graph-repair-large-fiber}
Let $q\ge9$.  Let
\[
  S=\{s_1,\ldots,s_q\},\qquad
  V_0=\{a_1,\ldots,a_q\}
\]
be disjoint ordered label sets.  For each permutation
$\pi\in\mathfrak S_q$, let $\Gamma_\pi$ be the loopless labelled
$2$-regular multigraph on $S\cup V_0$ obtained by putting two parallel edges
between $s_i$ and $a_{\pi(i)}$ for every $i$.

Deleting $S$ from every $\Gamma_\pi$ gives the same retained graph
$\Gamma_0=\varnothing$ on $V_0$, the same deficit word
$c_{a_r}=2$ for $1\le r\le q$, and hence the same canonical repaired graph
of \cref{lem:bc-canonical-graph-deficit-repair}.  The family has cardinality
$q!$, and
\[
  q!>2^{2q}.
\]
Consequently the canonical graph deficit repair, together with only
$2q$ unrestricted branch bits, cannot be an injective recovery map on the
unrestricted graph family.  In particular, for the residual fixed-witness
range $q\ge15$, the target-size graph repair must use additional structure
from the fixed witness crossings; the witness-free graph repair is not the
final supplier.
\end{lemma}

\begin{proof}
Each $\Gamma_\pi$ is loopless and every vertex has degree $2$: each $s_i$ is
incident to the two parallel edges joining it to $a_{\pi(i)}$, and each
$a_r$ is hit by the unique $s_i$ with $\pi(i)=r$, again with multiplicity
two.  After deleting $S$, no edge remains on $V_0$, so $\Gamma_0$ is empty and
$c_{a_r}=2$ for all $r$.  Since $q\ge9$, the deficits are not concentrated at
a single vertex, and \cref{lem:bc-canonical-graph-deficit-repair} applies.
Its input ordered deficit word is the same for every $\pi$, so its repaired
graph is also the same for every $\pi$.

The graphs $\Gamma_\pi$ are distinct because the neighbour of $s_1$ in
$V_0$ determines $\pi(1)$, and similarly for the other labels.  Hence the
fiber over that common repaired graph has size $q!$.

It remains only to compare with the branch budget.  One has
\[
  9! = 362880 > 262144 = 2^{18}.
\]
If $q!\!>2^{2q}$ and $q\ge9$, then
\[
  (q+1)!=(q+1)q!>(q+1)2^{2q}>4\cdot 2^{2q}=2^{2(q+1)}.
\]
Thus induction gives $q!>2^{2q}$ for every $q\ge9$.  Therefore
$2q$ bits, which distinguish at most $2^{2q}$ cases, cannot separate this
fiber.  The final sentence is the specialization to residual $q\ge15$.
\end{proof}

\begin{lemma}[The single-deficit obstruction is component-local]\label{lem:bc-single-deficit-component}
In the setting of \cref{lem:bc-graph-deficit-repair}, suppose that the only
nonzero deficit is $c_v=2$ at a single vertex $v\in V_0$.  Put
$U=V_0\setminus\{v\}$.  Then no edge of $\Gamma$ joins a vertex of
$S\cup\{v\}$ to a vertex of $U$.  Consequently the induced graph on
$S\cup\{v\}$ is a union of connected components of $\Gamma$, the induced graph
on $U$ is a loopless labelled $2$-regular multigraph, and the component
containing $v$ has all two edges incident to $v$ joining $v$ to vertices of
$S$, counted with multiplicity.

Conversely, if for some $v\in V_0$ no edge joins $S\cup\{v\}$ to
$V_0\setminus\{v\}$ and $v$ has no neighbour in $V_0$, then the deletion of
$S$ has exactly the single nonzero deficit $c_v=2$.
\end{lemma}

\begin{proof}
For any $u\in U$, the equality $c_u=0$ means that no edge of $\Gamma$ joins
$u$ to $S$.  The equality $c_v=2$, together with $\deg_\Gamma(v)=2$, means that
both edges incident to $v$ join $v$ to $S$, counted with multiplicity; in
particular no edge joins $v$ to $U$.  Thus no edge joins $S\cup\{v\}$ to $U$.
All degrees inside the induced graph on $U$ are therefore still $2$, and all
degrees inside the induced graph on $S\cup\{v\}$ are also still $2$.  This
proves the component decomposition and the asserted description of the
component containing $v$.

For the converse, the absence of edges from $S\cup\{v\}$ to
$V_0\setminus\{v\}$ gives $c_u=0$ for every $u\in V_0\setminus\{v\}$, while
the absence of a neighbour of $v$ in $V_0$ gives $\deg_{\Gamma_0}(v)=0$ and
hence $c_v=2$.
\end{proof}

\begin{lemma}[Single-deficit attachments are forced port matchings]\label{lem:bc-single-deficit-port-code}
Work in the setting of \cref{lem:bc-graph-attachment-port-code}.  Suppose that
the deletion of $S$ has exactly one nonzero deficit, namely
\[
  c_v=2
\]
for a retained vertex $v\in V_0$.  Then the port datum of
\cref{lem:bc-graph-attachment-port-code} has the following form:
\begin{enumerate}
\item exactly two ports are unmatched and mapped to $v$;
\item no port is mapped to any vertex of $V_0\setminus\{v\}$; and
\item every remaining port is matched to a port of a distinct deleted label.
\end{enumerate}
This restricted port datum determines the induced graph on $S\cup\{v\}$, and
the full graph is recovered from that induced graph together with the
loopless $2$-regular graph induced on $V_0\setminus\{v\}$.
\end{lemma}

\begin{proof}
By \cref{lem:bc-graph-attachment-port-code}, the number of ports mapped to a
retained vertex $u\in V_0$ is exactly the deficit $c_u$.  Since the only
nonzero deficit is $c_v=2$, exactly two ports are mapped to $v$, and no port
is mapped to any vertex of $V_0\setminus\{v\}$.  Every port is used exactly
once in the port datum, so all remaining ports are matched internally.  The
definition of the port datum permits an internal matched pair only between
distinct deleted labels, proving the third item.

The two ports mapped to $v$ give the two edges from $v$ to $S$, and each
matched pair gives one internal edge among the deleted labels.  Hence the
restricted port datum determines precisely the induced graph on $S\cup\{v\}$.
By \cref{lem:bc-single-deficit-component}, no edge joins $S\cup\{v\}$ to
$V_0\setminus\{v\}$, and the induced graph on $V_0\setminus\{v\}$ is already
loopless and $2$-regular.  The original graph is therefore the disjoint union
of these two induced graphs.
\end{proof}

\begin{lemma}[Single-deficit port codes split by connected component]\label{lem:bc-single-deficit-port-components}
Work in the setting of \cref{lem:bc-single-deficit-port-code}, and let $H$
be the loopless labelled multigraph on $S\cup\{v\}$ determined by the
restricted port datum.  Let $H_v$ be the connected component of $H$ containing
$v$, and put
\[
  S_v=V(H_v)\cap S.
\]
Then $H_v$ is a loopless connected labelled $2$-regular multigraph on
$S_v\cup\{v\}$, and its two edges incident to $v$ are exactly the two ports
mapped to $v$.  Every other connected component of $H$ is a loopless connected
labelled $2$-regular multigraph on a subset of $S$ and is encoded only by
internally matched deleted-label ports.  The restricted port datum decomposes
uniquely as the $v$-component subdatum together with these deleted-only
component subdata.
\end{lemma}

\begin{proof}
By \cref{lem:bc-single-deficit-port-code}, the restricted port datum has
exactly two unmatched ports, both mapped to $v$, no port mapped to any other
retained label, and all remaining ports internally matched between distinct
deleted labels.  Hence the graph $H$ determined by the datum is loopless.  The
two unmatched ports give the two edges incident to $v$, and the two ports at
each deleted label are each used once, so every vertex of $H$ has degree $2$.

Connected components of a loopless labelled $2$-regular multigraph are again
loopless labelled connected $2$-regular multigraphs.  The component containing
$v$ therefore has vertex set $S_v\cup\{v\}$ and contains the two edges coming
from the two ports mapped to $v$.  Since no port is mapped to a retained label
other than $v$, every component not containing $v$ has all vertices in $S$ and
all of its edges come from internally matched deleted-label ports.

Finally, connected components partition the edge multiset of $H$.  Restricting
the port datum to the ports whose endpoints lie in one component gives the
corresponding component subdatum, and the union of these subdata reconstructs
the original restricted port datum.  The partition into connected components
is unique, proving the claimed decomposition.
\end{proof}

\begin{lemma}[The single-deficit component is an anchored cycle]\label{lem:bc-single-deficit-anchored-cycle}
Work in the setting of \cref{lem:bc-single-deficit-port-components}.  If
$|S_v|=1$, then $H_v$ consists of the unique deleted label in $S_v$, the
retained vertex $v$, and two parallel edges between them.  If $|S_v|\ge2$,
then $H_v$ has no parallel edges and its underlying simple graph is a cycle on
$S_v\cup\{v\}$.  Equivalently, deleting $v$ and its two incident edges leaves a
simple path through all labels of $S_v$, whose endpoints are the two deleted
neighbours of $v$.

The same dichotomy holds for every deleted-only component in
\cref{lem:bc-single-deficit-port-components}: it is either two deleted labels
joined by two parallel edges, or a simple cycle on at least three deleted
labels.
\end{lemma}

\begin{proof}
Let $C$ be any connected component of the loopless labelled $2$-regular
multigraph $H$.  If $C$ has a pair of parallel edges between two labels $a$ and
$b$, those two edge copies already contribute degree $2$ at both $a$ and $b$.
Since every vertex of $C$ has degree $2$ and $C$ is connected, no further
vertex or edge can lie in $C$.  Thus a component with parallel edges is exactly
the two-vertex double-edge component.

If $C$ has no parallel edges, it is a finite connected simple graph in which
every vertex has degree $2$.  Start at any vertex and follow an edge without
immediately backtracking.  Finiteness gives a first repeated vertex, and the
degree-two condition forces the resulting cycle to contain every vertex of
$C$; otherwise a vertex on the cycle would need an additional edge to the
remaining connected part, giving degree at least $3$.  Hence $C$ is a simple
cycle.

Apply this classification to the component $H_v$.  Since $v$ has degree $2$ in
$H_v$, the two-vertex case with $v$ is exactly the double edge from $v$ to the
unique deleted label in $S_v$.  In the simple-cycle case, removing $v$ and its
two incident edges breaks the cycle into the asserted simple path through all
deleted labels of $S_v$.  Applying the same component classification to the
components not containing $v$ gives the deleted-only statement.
\end{proof}

\begin{lemma}[Counting single-deficit anchored components]\label{lem:bc-single-deficit-anchored-count}
Let $S$ be a finite label set with $|S|=q$, and let $v\notin S$.  For a
loopless labelled $2$-regular multigraph $H$ on $S\cup\{v\}$, let $S_v$ be
the set of labels of $S$ lying in the connected component of $v$.  Fix
$T\subset S$ with $|T|=h$.  The number of such graphs with $S_v=T$ is
\[
  a_h\,s_{q-h},
  \qquad
  a_1=1,\qquad a_h=\frac{h!}{2}\quad(h\ge2),
\]
and there are no graphs with $h=0$.  Consequently the number of such graphs
whose $v$-component has exactly $h$ deleted labels is
\[
  \binom qh a_h\,s_{q-h}\qquad(1\le h\le q).
\]
\end{lemma}

\begin{proof}
The case $h=0$ is impossible because $v$ would have degree $0$ in a graph in
which every vertex has degree $2$ and loops are forbidden.  Fix
$T\subset S$ with $|T|=h\ge1$.

The component classification proved in
\cref{lem:bc-single-deficit-anchored-cycle} applies to every connected
component of a loopless labelled $2$-regular multigraph.  Hence the component
containing $v$ is determined as follows.  If $h=1$, it is forced to be the two
parallel edges between $v$ and the unique label of $T$.  If $h\ge2$, deleting
$v$ leaves a simple path through all labels of $T$, and adjoining $v$ to the
two endpoints recovers the simple cycle.  The number of unoriented simple
paths through the $h$ labelled vertices of $T$ is $h!/2$, since a linear order
and its reversal give the same path and these are the only repetitions.

The components not containing $v$ form an arbitrary loopless labelled
$2$-regular multigraph on $S\setminus T$, independently of the anchored
component.  By \cref{lem:bc-littlewood-graph-model}, the number of such
deleted-only graphs is $s_{q-h}$.  This gives $a_hs_{q-h}$ for fixed $T$.
There are $\binom qh$ choices for $T$, giving the displayed component-size
count.
\end{proof}

\begin{lemma}[Anchored single-deficit components have a support-path code]\label{lem:bc-single-deficit-support-path-code}
Let $S$ be a finite label set and let $v\notin S$.  A loopless labelled
$2$-regular multigraph $H$ on $S\cup\{v\}$ is recovered from the following
data:
\begin{enumerate}
\item the support $T\subset S$ of the connected component containing $v$;
\item if $|T|\ge2$, an ordering $(t_1,\ldots,t_{|T|})$ of the elements of $T$,
taken modulo reversal; and
\item the loopless labelled $2$-regular multigraph induced on $S\setminus T$.
\end{enumerate}
The recovery rule is canonical: if $|T|=1$, put two parallel edges between
$v$ and the unique label of $T$; if $|T|\ge2$, put the cycle
\[
  v-t_1-t_2-\cdots-t_{|T|}-v
\]
and then take the disjoint union with the graph on $S\setminus T$.  Conversely,
the graph $H$ determines exactly these data.
\end{lemma}

\begin{proof}
For a given graph $H$, let $T$ be the set of labels of $S$ in the connected
component of $v$.  If this component contains a pair of parallel edges between
two labels, then those two edge copies already give degree $2$ at both
endpoints; connectedness and the degree-two condition force the component to
be exactly that double-edge component.  Thus, when $|T|=1$, the component of
$v$ is the forced double edge from $v$ to the unique label of $T$.

If there are no parallel edges in the component of $v$, it is a finite
connected simple graph in which every vertex has degree $2$, hence a simple
cycle.  For $|T|\ge2$, deleting $v$ from this cycle leaves a simple path
through all labels of $T$.  Reading that path from one endpoint to the other
gives an ordering of $T$, and reversing the reading gives the only other
ordering describing the same path.  The remaining connected components are
precisely the induced loopless labelled $2$-regular multigraph on
$S\setminus T$.

Conversely, the displayed rule constructs a loopless labelled $2$-regular
multigraph on $T\cup\{v\}$ whose component containing $v$ has support $T$.
The reversal of $(t_1,\ldots,t_{|T|})$ gives the same unoriented cycle, so the
rule is well defined on reversal classes.  Taking the disjoint union with the
given graph on $S\setminus T$ gives a loopless labelled $2$-regular multigraph
on $S\cup\{v\}$.  The two constructions are inverse by the component
classification in the preceding paragraphs.
\end{proof}

\begin{lemma}[A single retained deficit is repaired by inserting into one retained edge]\label{lem:bc-single-deficit-retained-edge-insertion}
Let $U$ be a finite linearly ordered label set with $|U|\ge2$, let
$v\notin U$, and let $\Gamma_U$ be a loopless labelled $2$-regular multigraph
on $U$.  Choose the lexicographically least edge $ab$ of $\Gamma_U$, counting
parallel copies by their multiplicity order.  Let
\[
  \mathcal I_v(\Gamma_U)
\]
be the multigraph on $U\cup\{v\}$ obtained by deleting that one copy of $ab$
and adding the two edges $va$ and $vb$.

Then $\mathcal I_v(\Gamma_U)$ is a loopless labelled $2$-regular multigraph
on $U\cup\{v\}$.  Moreover $\Gamma_U$ is recovered from
$\mathcal I_v(\Gamma_U)$ and $v$ by deleting the two edges incident to $v$ and
restoring one edge between their other endpoints.
\end{lemma}

\begin{proof}
Since $\Gamma_U$ is loopless and every vertex has degree $2$, it has at least
one edge; the chosen edge joins two distinct labels $a,b\in U$.  Deleting one
copy of $ab$ lowers the degrees of $a$ and $b$ by one.  Adding $va$ and $vb$
restores the degrees of $a$ and $b$ and gives degree $2$ at $v$.  Every other
vertex has unchanged degree.  No loop is introduced because
$v\notin U$ and $a\ne b$.

In $\mathcal I_v(\Gamma_U)$ the vertex $v$ is incident exactly to the two
edges $va$ and $vb$.  Deleting those two edges lowers the degrees of $a$ and
$b$ by one and removes $v$; adding one copy of $ab$ restores exactly the edge
deleted in the construction.  Thus the stated inverse recovers $\Gamma_U$.
\end{proof}

\begin{proposition}[Single-deficit graphs have a stable retained repair]\label{prop:bc-single-deficit-stable-retained-repair}
Work in the setting of \cref{lem:bc-single-deficit-component}.  Suppose that
the deletion of $S$ has exactly one nonzero deficit $c_v=2$ and that
$U=V_0\setminus\{v\}$ has at least two vertices.  Let $\Gamma_U$ be the graph
induced by $\Gamma$ on $U$, and set
\[
  \Gamma^\ast=\mathcal I_v(\Gamma_U),
\]
using \cref{lem:bc-single-deficit-retained-edge-insertion}.  Then
$\Gamma^\ast$ is a loopless labelled $2$-regular multigraph on $V_0$.
Furthermore, the original graph $\Gamma$ is recovered from
\[
  S,\quad v,\quad \Gamma^\ast,
  \quad\text{and the support-path code of }
  \cref{lem:bc-single-deficit-support-path-code}.
\]
Thus, in the single-deficit case with $|V_0\setminus\{v\}|\ge2$, the stable
retained output carries no additional choice; the remaining branch-data
problem is the recovery of the retained vertex $v$ and the support-path code,
including the deleted-only graph on the complement of the component containing
$v$.
\end{proposition}

\begin{proof}
By \cref{lem:bc-single-deficit-component}, no edge joins $S\cup\{v\}$ to $U$,
and the graph induced on $U$ is loopless and $2$-regular.  Since $|U|\ge2$,
\cref{lem:bc-single-deficit-retained-edge-insertion} applies and gives a
loopless labelled $2$-regular graph $\Gamma^\ast$ on $U\cup\{v\}=V_0$.

The same lemma recovers $\Gamma_U$ from $\Gamma^\ast$ and $v$.  The
support-path code of \cref{lem:bc-single-deficit-support-path-code} recovers
the induced graph on $S\cup\{v\}$.  Again by
\cref{lem:bc-single-deficit-component}, the original graph is the disjoint
union of these two induced graphs.  Hence the displayed data recover
$\Gamma$.
\end{proof}

\begin{lemma}[Deleted-only single-deficit components are independent data]\label{lem:bc-deleted-only-single-deficit-data}
Let $S$ contain five distinct labels
\[
  s_0,d_1,d_2,d_3,d_4,
\]
let $v\notin S$, and let $U=\{u_1,u_2\}$ be disjoint from
$S\cup\{v\}$.  There exist two distinct loopless labelled $2$-regular
multigraphs on $S\cup U\cup\{v\}$ such that
\begin{enumerate}
\item deleting $S$ gives the same single retained deficit $c_v=2$;
\item the stable retained repair of
      \cref{prop:bc-single-deficit-stable-retained-repair} is the same
      graph $\Gamma^\ast$;
\item the component containing $v$ has the same support $\{s_0\}$ and the
      same two parallel edges from $v$ to $s_0$; and
\item the two graphs differ only in the deleted-only component on
      $\{d_1,d_2,d_3,d_4\}$.
\end{enumerate}
Consequently a recovery statement for the whole original graph cannot omit
the induced deleted-only graph on $S\setminus T$, where $T$ is the support of
the component containing $v$, unless an additional argument proves
$T=S$ or supplies that deleted-only graph by some other target datum.
\end{lemma}

\begin{proof}
Let $\Gamma_U$ be the graph with two parallel edges between $u_1$ and $u_2$.
Let $H_0$ be the graph with two parallel edges between $v$ and $s_0$.
On $D=\{d_1,d_2,d_3,d_4\}$, take the two simple cycles
\[
  d_1-d_2-d_3-d_4-d_1
  \qquad\text{and}\qquad
  d_1-d_3-d_2-d_4-d_1 .
\]
Let $\Gamma$ and $\Gamma'$ be the disjoint unions of $\Gamma_U$, $H_0$, and
these two respective cycles on $D$.  Both graphs are loopless and every
vertex has degree $2$.

After deleting $S$, both graphs leave $\Gamma_U$ on $U$ and the isolated
vertex $v$, so the only nonzero retained deficit is $c_v=2$.  Since the
retained complement graph $\Gamma_U$ is the same in both cases,
\cref{prop:bc-single-deficit-stable-retained-repair} produces the same
repaired retained graph $\Gamma^\ast=\mathcal I_v(\Gamma_U)$.  The component
containing $v$ is $H_0$ in both cases, hence has support $\{s_0\}$ and the
same anchored double edge.  The two graphs are nevertheless distinct because
the cycle on $D$ is different.  This proves the four displayed properties.

The final assertion follows because the displayed data
$S,v,\Gamma^\ast$ and the anchored support/path data are identical for
$\Gamma$ and $\Gamma'$, while the original graphs are not identical.  The
missing datum is exactly the induced deleted-only graph on $S\setminus T$.
\end{proof}

\begin{lemma}[Deleted-only complements exceed the retained-repair branch budget]\label{lem:bc-deleted-only-complement-budget-guard}
Let $q\ge15$.  Let $S$ be a linearly ordered $q$-element set, choose
$s_0\in S$, and put $D=S\setminus\{s_0\}$.  Let $v\notin S$, and let $U$ be
a linearly ordered set disjoint from $S\cup\{v\}$ which carries a fixed
loopless labelled $2$-regular multigraph $\Gamma_U$.  Assume $|U|\ge2$, and
put
\[
  \Gamma^\ast=\mathcal I_v(\Gamma_U)
\]
as in \cref{lem:bc-single-deficit-retained-edge-insertion}.

There are more than $2^{2q}$ loopless labelled $2$-regular multigraphs
$\Gamma$ on $S\cup U\cup\{v\}$ such that
\begin{enumerate}
\item deleting $S$ gives exactly one nonzero deficit, namely $c_v=2$;
\item the stable retained repair of
      \cref{prop:bc-single-deficit-stable-retained-repair} is the fixed graph
      $\Gamma^\ast$;
\item the component containing $v$ has support $\{s_0\}$, so the anchored
      support is a fixed one-point witness interval and the adjacent-rank
      internal-edge condition is vacuous; and
\item the only varying datum is the deleted-only component on $D$.
\end{enumerate}
Consequently the retained-repair target $\Gamma^\ast$, even together with
$v$, the anchored interval support, and an unrestricted $2q$-bit branch word,
cannot recover this family.  A single-deficit fixed-witness supplier must
therefore either prove the deleted-only complement is absent, recover it from
the actual target-size path, or encode it in a target object which is not just
the stable retained repair.
\end{lemma}

\begin{proof}
For every unoriented simple cycle $C$ through all labels of $D$, let
$\Gamma_C$ be the disjoint union of $\Gamma_U$, the two parallel edges between
$v$ and $s_0$, and the cycle $C$ on $D$.  This graph is loopless and every
vertex has degree $2$.  Deleting $S$ leaves $\Gamma_U$ on $U$ and isolates
$v$, so the only nonzero retained deficit is $c_v=2$.  Since $\Gamma_U$ is
fixed, the stable retained repair is the same graph
$\Gamma^\ast=\mathcal I_v(\Gamma_U)$ for every $C$.

The component containing $v$ is always the double edge between $v$ and $s_0$.
Thus its support is the fixed one-point interval $\{s_0\}$, and there are no
anchored internal edges.  Different cycles $C$ give different graphs
$\Gamma_C$, and the only varying part is the deleted-only component on $D$.

The set $D$ has $q-1$ labels.  The number of unoriented simple cycles on a
fixed $(q-1)$-element labelled set is
\[
  \frac{(q-2)!}{2}.
\]
At $q=15$ this is
\[
  \frac{13!}{2}=3113510400>1073741824=2^{30}=2^{2q}.
\]
If $(q-2)!/2>2^{2q}$ and $q\ge15$, then
\[
  \frac{(q-1)!}{2}
  =(q-1)\frac{(q-2)!}{2}
  >4\cdot2^{2q}=2^{2(q+1)}.
\]
Induction gives $(q-2)!/2>2^{2q}$ for every $q\ge15$.

For the displayed family, $S$, $v$, $\Gamma^\ast$, and the anchored interval
support are fixed.  A $2q$-bit word has only $2^{2q}$ values, fewer than the
number of graphs in the family.  Hence these data cannot be injective on the
family.  The final sentence is exactly this obstruction applied to any
single-deficit supplier based only on the retained repair target.
\end{proof}

\begin{lemma}[Unrestricted anchored path orders exceed the branch budget]\label{lem:bc-anchored-path-order-budget-guard}
Let $q\ge10$.  Let $S$ be a $q$-element linearly ordered label set, let
$v\notin S$, and let $U$ be a linearly ordered label set disjoint from
$S\cup\{v\}$ which carries a fixed loopless labelled $2$-regular multigraph
$\Gamma_U$.  Assume $|U|\ge2$, and put
\[
  \Gamma^\ast=\mathcal I_v(\Gamma_U)
\]
as in \cref{lem:bc-single-deficit-retained-edge-insertion}.

There are $q!/2$ loopless labelled $2$-regular multigraphs $\Gamma$ on
$S\cup U\cup\{v\}$ such that
\begin{enumerate}
\item deleting $S$ gives exactly one nonzero deficit, namely $c_v=2$;
\item the stable retained repair of
\cref{prop:bc-single-deficit-stable-retained-repair} is the same graph
$\Gamma^\ast$ for every member of the family; and
\item the support of the $v$-component is all of $S$, while the anchored path
order varies over all orderings of $S$ modulo reversal.
\end{enumerate}
Moreover
\[
  \frac{q!}{2}>2^{2q}.
\]
Consequently the repaired retained graph, the vertex $v$, and an unrestricted
$2q$-bit branch word cannot by themselves recover the unrestricted anchored
path order in the residual range $q\ge15$.
\end{lemma}

\begin{proof}
For each ordering $(s_1,\ldots,s_q)$ of $S$, taken modulo reversal, let
$H(s_1,\ldots,s_q)$ be the cycle
\[
  v-s_1-s_2-\cdots-s_q-v
\]
on $S\cup\{v\}$.  Let $\Gamma$ be the disjoint union of this cycle with the
fixed graph $\Gamma_U$ on $U$.  It is loopless and $2$-regular on
$S\cup U\cup\{v\}$.  After deleting $S$, the vertex $v$ is isolated, while the
graph on $U$ is still $\Gamma_U$; hence the only nonzero deficit is
$c_v=2$.  Since the retained complement graph is the same fixed $\Gamma_U$ in
every case, \cref{prop:bc-single-deficit-stable-retained-repair} produces the
same repaired retained graph $\Gamma^\ast=\mathcal I_v(\Gamma_U)$ for every
member of the family.  Distinct orderings modulo reversal give distinct
unoriented cycles through the labelled set $S\cup\{v\}$, so the family has
cardinality $q!/2$.

For the inequality, the base case is
\[
  \frac{10!}{2}=1814400>1048576=2^{20}.
\]
If $q!/2>2^{2q}$ and $q\ge10$, then
\[
  \frac{(q+1)!}{2}
  =(q+1)\frac{q!}{2}
  >(q+1)2^{2q}
  >4\cdot2^{2q}
  =2^{2(q+1)}.
\]
Thus induction proves $q!/2>2^{2q}$ for every $q\ge10$.  A $2q$-bit word has
only $2^{2q}$ possible values, so it cannot distinguish all members of this
common repaired-output fiber.  Since the residual B/C range has $q\ge15$, the
final sentence follows.
\end{proof}

\begin{proposition}[Monotone anchored order leaves only support bits for the anchored component]\label{prop:bc-monotone-anchored-order-support-bits}
Work in the setting of
\cref{prop:bc-single-deficit-stable-retained-repair}.  Assume that, in the
support-path code of \cref{lem:bc-single-deficit-support-path-code}, the
anchored path order on the support $T\subset S$ is the increasing order of
$T$ in the ambient order on $S$, taken modulo reversal.  Then the original
graph $\Gamma$ is recovered from
\[
  S,\quad v,\quad \Gamma^\ast,\quad
  \Gamma_{S\setminus T},\quad \mathbf 1_T\in\{0,1\}^{S},
\]
where $\Gamma_{S\setminus T}$ is the loopless labelled $2$-regular graph
induced by $\Gamma$ on $S\setminus T$.
Consequently, under any fixed-witness crossing theorem which determines the
retained deficit vertex $v$ and forces this monotone anchored order, the
single-deficit graph splice needs no separate path-order branch data for the
anchored component.  A full supplier must still either supply or absorb the
deleted-only graph $\Gamma_{S\setminus T}$, or prove $T=S$.
\end{proposition}

\begin{proof}
The characteristic word $\mathbf 1_T$ and the ordered set $S$ recover the
support $T$.  If $|T|=1$, \cref{lem:bc-single-deficit-support-path-code}
prescribes the unique double edge from $v$ to the single label of $T$.  If
$|T|\ge2$, the monotonicity hypothesis says that the ordering of $T$ modulo
reversal is the increasing order, so the same lemma reconstructs the anchored
cycle
\[
  v-t_1-\cdots-t_{|T|}-v
\]
where $t_1<\cdots<t_{|T|}$ are the elements of $T$ in the ambient order.

The graph $\Gamma^\ast$ and the vertex $v$ recover the retained complement
graph $\Gamma_U$ by
\cref{lem:bc-single-deficit-retained-edge-insertion}, as used in
\cref{prop:bc-single-deficit-stable-retained-repair}.  Taking the disjoint
union of this recovered retained complement, the reconstructed anchored
component on $T\cup\{v\}$, and the supplied deleted-only graph
$\Gamma_{S\setminus T}$ gives $\Gamma$, by
\cref{lem:bc-single-deficit-component,lem:bc-single-deficit-port-components}.
The final sentence is the same recovery statement with the independent
deleted-only datum separated, as required by
\cref{lem:bc-deleted-only-single-deficit-data}.
\end{proof}

\begin{lemma}[Adjacent support edges force monotone anchored order]\label{lem:bc-adjacent-support-edges-force-monotone}
Let $S$ be a finite linearly ordered label set, let $v\notin S$, and let
$T\subset S$ with $|T|=h\ge2$.  Let $P$ be a simple path through all labels of
$T$.  Suppose every edge of $P$ joins two labels which are adjacent in the
induced order on $T$.  Then $P$ is the monotone path through $T$, uniquely up
to reversal.
\end{lemma}

\begin{proof}
Write the elements of $T$ in increasing order as
\[
  t_1<t_2<\cdots<t_h .
\]
The unordered pairs of adjacent labels in the induced order are exactly
\[
  \{t_i,t_{i+1}\}\qquad(1\le i<h).
\]
There are $h-1$ such pairs.  A simple path through all $h$ labels of $T$ also
has exactly $h-1$ edges.  Since every edge of $P$ is one of the displayed
adjacent pairs and $P$ is connected on all vertices of $T$, it must use all of
them.  Hence $P$ is the path
\[
  t_1-t_2-\cdots-t_h,
\]
read in one direction or the other.
\end{proof}

\begin{proposition}[Adjacent support-edge recovery criterion]\label{prop:bc-adjacent-support-edge-recovery}
Work in the setting of
\cref{prop:bc-single-deficit-stable-retained-repair}.  Assume that, in the
support-path code of \cref{lem:bc-single-deficit-support-path-code}, every
internal edge of the anchored path on $T$ joins adjacent labels in the induced
order on $T$.  Then $\Gamma$ is recovered from
\[
  S,\quad v,\quad \Gamma^\ast,\quad
  \Gamma_{S\setminus T},\quad \mathbf 1_T\in\{0,1\}^{S}.
\]
Consequently this leaf removes the anchored path-order choice once
$v$, the support, and the deleted-only component have been supplied; it does
not by itself encode the independent deleted-only graph on $S\setminus T$.
\end{proposition}

\begin{proof}
If $|T|=1$, the support-path code has the forced double edge from $v$ to the
unique label of $T$, so there is no internal path-order choice.  If
$|T|\ge2$, \cref{lem:bc-adjacent-support-edges-force-monotone} shows that the
anchored path order on $T$ is the increasing order, modulo reversal.  Therefore
\cref{prop:bc-monotone-anchored-order-support-bits} applies and gives recovery
from the displayed data.  The final sentence is precisely this recovery
criterion interpreted for the fixed-witness supplier.
\end{proof}

\begin{lemma}[Consecutive frontier events give adjacent support labels]\label{lem:bc-frontier-consecutive-gives-support-adjacent}
Let $S$ be a finite linearly ordered label set and let $T\subset S$.  Suppose
each $t\in T$ is assigned a frontier event $e(t)$ in a linearly ordered set
and that this assignment is order-preserving in the sense that, for
$t,t'\in T$,
\[
  t<t' \quad\Longleftrightarrow\quad e(t)<e(t').
\]
If $a,b\in T$ and no event $e(t)$ with $t\in T$ lies strictly between
$e(a)$ and $e(b)$, then $a$ and $b$ are adjacent in the induced order on $T$.
\end{lemma}

\begin{proof}
Assume, after interchanging $a$ and $b$ if necessary, that $a<b$.  The
order-preserving hypothesis gives $e(a)<e(b)$.  If a label
$t\in T$ satisfied $a<t<b$, then the same hypothesis would give
\[
  e(a)<e(t)<e(b),
\]
contradicting the assumption that no support event lies strictly between
$e(a)$ and $e(b)$.  Hence no element of $T$ lies strictly between $a$ and $b$,
which is exactly adjacency in the induced order on $T$.
\end{proof}

\begin{proposition}[Consecutive frontier-event edges give single-deficit recovery]\label{prop:bc-frontier-consecutive-edge-recovery}
Work in the setting of
\cref{prop:bc-single-deficit-stable-retained-repair}.  Let $T\subset S$ be the
support of the anchored component in
\cref{lem:bc-single-deficit-support-path-code}.  Suppose the labels of $T$ are
assigned order-preservingly to their fixed-witness frontier events, and every
internal edge of the anchored path on $T$ has endpoints whose assigned
frontier events are consecutive among the events assigned to $T$.  Then
$\Gamma$ is recovered from
\[
  S,\quad v,\quad \Gamma^\ast,\quad
  \Gamma_{S\setminus T},\quad \mathbf 1_T\in\{0,1\}^{S}.
\]
Consequently the frontier-event condition settles only the anchored
path-order part of the single-deficit recovery; the deleted-only graph remains
an explicit supplier input unless $T=S$.
\end{proposition}

\begin{proof}
By \cref{lem:bc-frontier-consecutive-gives-support-adjacent}, each internal
edge of the anchored path joins adjacent labels in the induced order on $T$.
Therefore \cref{prop:bc-adjacent-support-edge-recovery} applies and gives the
displayed recovery.  The final sentence is this criterion in the language of
the fixed-witness crossing construction.
\end{proof}

\begin{lemma}[Full frontier adjacency restricts to support adjacency]\label{lem:bc-full-frontier-adjacency-restricts-support}
Let $S$ be a finite linearly ordered label set, let $T\subset S$, and let
$e:S\to E$ be a map into a linearly ordered set such that, for
$s,s'\in S$,
\[
  s<s'\quad\Longleftrightarrow\quad e(s)<e(s').
\]
If $a,b\in T$ and no event $e(s)$ with $s\in S$ lies strictly between
$e(a)$ and $e(b)$, then $a$ and $b$ are adjacent in the induced order on $T$.
\end{lemma}

\begin{proof}
The restriction of $e$ to $T$ satisfies the same order equivalence.  Since
$T\subset S$,
the absence of an event $e(s)$ with $s\in S$ strictly between $e(a)$ and
$e(b)$ implies the same absence for events $e(t)$ with $t\in T$.  Hence
\cref{lem:bc-frontier-consecutive-gives-support-adjacent} applies to the
restricted map $T\to E$ and gives the asserted adjacency.
\end{proof}

\begin{proposition}[Full frontier-adjacent edges give single-deficit recovery]\label{prop:bc-full-frontier-adjacent-edge-recovery}
Work in the setting of
\cref{prop:bc-single-deficit-stable-retained-repair}.  Let $T\subset S$ be the
support of the anchored component in
\cref{lem:bc-single-deficit-support-path-code}.  Suppose all labels of $S$ are
assigned to their fixed-witness frontier events by a map $e:S\to E$ satisfying
\[
  s<s'\quad\Longleftrightarrow\quad e(s)<e(s')
  \qquad(s,s'\in S).
\]
If every internal edge of the anchored path on $T$ has endpoints whose assigned
frontier events are consecutive among the full event set assigned to $S$, then
$\Gamma$ is recovered from
\[
  S,\quad v,\quad \Gamma^\ast,\quad
  \Gamma_{S\setminus T},\quad \mathbf 1_T\in\{0,1\}^{S}.
\]
Consequently full-frontier adjacency removes the path-order choice without
pre-filtering by the unknown support.  It still leaves the independent
deleted-only component on $S\setminus T$ to be supplied, absorbed into the
target, or ruled out by proving $T=S$.
\end{proposition}

\begin{proof}
Let $ab$ be an internal edge of the anchored path on $T$.  By hypothesis, no
frontier event assigned to a label of $S$ lies strictly between the events
assigned to $a$ and $b$.  By
\cref{lem:bc-full-frontier-adjacency-restricts-support}, the labels $a$ and
$b$ are adjacent in the induced order on $T$.  Thus every internal edge of the
anchored path joins adjacent support labels.  Applying
\cref{prop:bc-adjacent-support-edge-recovery} gives the displayed recovery.
The final sentence is the same criterion expressed in the full frontier-event
list visible in the fixed-witness crossing construction.
\end{proof}

\begin{lemma}[Single-deficit components factor from the retained graph]\label{lem:bc-single-deficit-factorization}
Let $S$ be a finite label set, let $v\notin S$, and let $U$ be a finite label
set disjoint from $S\cup\{v\}$.  Among loopless labelled $2$-regular
multigraphs on $S\cup\{v\}\cup U$, the graphs whose deletion of $S$ has exactly
one nonzero deficit, namely $c_v=2$, are in bijection with pairs
\[
  (\Gamma_1,\Gamma_2),
\]
where $\Gamma_1$ is a loopless labelled $2$-regular multigraph on
$S\cup\{v\}$ and $\Gamma_2$ is a loopless labelled $2$-regular multigraph on
$U$.  Consequently their number is
\[
  s_{|S|+1}s_{|U|}.
\]
\end{lemma}

\begin{proof}
By \cref{lem:bc-single-deficit-component}, a graph with exactly the deficit
$c_v=2$ has no edge joining $S\cup\{v\}$ to $U$, and the induced graphs on
$S\cup\{v\}$ and on $U$ are both loopless labelled $2$-regular multigraphs.
Thus restriction to the two induced subgraphs gives the displayed pair.

Conversely, the disjoint union of any such pair is loopless and has degree $2$
at every vertex.  After deleting $S$, every vertex of $U$ still has degree $2$,
while $v$ has no neighbour in $U$ and hence has degree $0$ in the retained
graph.  Therefore the unique nonzero deficit is $c_v=2$.  The two constructions
are inverse to one another.  The product count follows from
\cref{lem:bc-littlewood-graph-model}, which identifies $s_n$ with the number
of loopless labelled $2$-regular multigraphs on any fixed $n$-element label
set.
\end{proof}

\begin{corollary}[The residual-floor all-graph single-deficit family exceeds the fixed-fiber budget]\label{cor:bc-residual-floor-single-deficit-budget-obstruction}
Let $|S|=15$, fix $v\notin S$, and let $U$ be a $14$-element label set
disjoint from $S\cup\{v\}$.  The all-graph single-deficit family of
\cref{lem:bc-single-deficit-factorization} has cardinality
\[
  s_{16}s_{14}
  =
  23184038368042864094640,
\]
whereas a target consisting of an arbitrary stable retained graph on
$15$ labels together with a $30$-bit branch word has size
\[
  2^{30}s_{15}
  =
  158126931013997166592.
\]
In particular,
\[
  s_{16}s_{14}>2^{30}s_{15}.
\]
Consequently a supplier which only compresses the witness-free all-graph
single-deficit family to a stable retained graph/path on $j-q=15$ labels and
$2q$ branch bits cannot be injective at the first residual boundary
$q=15$, $j=30$.
\end{corollary}

\begin{proof}
By \cref{lem:bc-single-deficit-factorization}, the stated family has
cardinality $s_{|S|+1}s_{|U|}=s_{16}s_{14}$.  Iterating
\cref{lem:bc-stable-moment-recurrence} gives
\[
  s_{14}=10157945044,\qquad
  s_{15}=147267180508,\qquad
  s_{16}=2282355168060.
\]
Hence
\[
  s_{16}s_{14}
  =
  23184038368042864094640
\]
and
\[
  2^{30}s_{15}
  =
  158126931013997166592.
\]
The displayed strict inequality follows.  The final sentence is then the
pigeonhole principle applied to a domain larger than its proposed target set.
\end{proof}

\begin{lemma}[$B/C$ corrections as Littlewood bad-shape sums]\label{lem:bc-correction-bad-shapes}
Let $a_j(\lambda)$ be the Schur coefficient
\[
  h_2^j=\sum_\lambda a_j(\lambda)s_\lambda .
\]
For $b\ge1$ and $j\ge0$,
\[
  |\delta_j^{B_b}|
  =
  \sum_{\substack{\lambda':\,\mathrm{even}\\ \lambda_1\ge 2b+2}}
    a_j(\lambda),
\]
and
\[
  |\delta_j^{C_b}|
  =
  \sum_{\substack{\lambda':\,\mathrm{even}\\ \ell(\lambda)\ge 2b+2}}
    a_j(\lambda).
\]
\end{lemma}

\begin{proof}
For $B_b$, \cref{prop:bc-pieri-correction-sign} gives
\[
  |\delta_j^{B_b}|
  =
  \sum_{\substack{\nu\vdash j\\ \ell(\nu)>2b+1}}
    c_j(2\nu_1,2\nu_2,\ldots),
\]
where $e_2^j=\sum_\lambda c_j(\lambda)s_\lambda$.  Applying the Hall
involution $\omega$ gives
\[
  c_j(2\nu_1,2\nu_2,\ldots)
  =
  a_j\bigl((2\nu_1,2\nu_2,\ldots)'\bigr),
\]
because $\omega(e_2)=h_2$ and $\omega(s_\lambda)=s_{\lambda'}$.  The map
\[
  \nu\longmapsto \lambda=(2\nu_1,2\nu_2,\ldots)'
\]
is a bijection from partitions with $\ell(\nu)\ge2b+2$ to partitions
$\lambda$ whose conjugate has only even parts and whose width satisfies
$\lambda_1\ge2b+2$.  This proves the $B_b$ identity.

For $C_b$, \cref{prop:bc-pieri-correction-sign} gives
\[
  |\delta_j^{C_b}|
  =
  \sum_{\substack{\mu\vdash j\\ \ell(\mu)>b}}
    a_j(\mu_1,\mu_1,\ldots,\mu_{\ell(\mu)},\mu_{\ell(\mu)}).
\]
The partitions of the displayed paired-row form are exactly the partitions
$\lambda$ for which $\lambda'$ has only even parts.  The condition
$\ell(\mu)>b$ is equivalent to
$\ell(\lambda)=2\ell(\mu)\ge2b+2$.  This proves the $C_b$ identity.
\end{proof}

\begin{lemma}[Canonical witnesses for residual bad shapes]\label{lem:bad-shape-witnesses}
Let $a_j(\lambda)$ be as in \cref{lem:bc-correction-bad-shapes}.  Equivalently,
$a_j(\lambda)$ counts semistandard tableaux of shape $\lambda$ and content
$(2^j)$.  Fix $q\ge1$.

If $\lambda'$ is even and $\lambda_1\ge2q$, then every such tableau $T$ has a
canonical $q$-element witness
\[
  W_q(T)=\{T(1,1),T(1,3),\ldots,T(1,2q-1)\}\subset\{1,\ldots,j\}.
\]
If $\lambda'$ is even and $\ell(\lambda)\ge2q$, then every such tableau $T$
has a canonical $q$-element witness
\[
  H_q(T)=\{T(1,1),T(3,1),\ldots,T(2q-1,1)\}\subset\{1,\ldots,j\}.
\]
\end{lemma}

\begin{proof}
The coefficient statement is the standard complete-symmetric expansion:
$h_2^j=h_{(2^j)}$ is the generating function for semistandard tableaux with
content $(2^j)$.

For the width witness, the first row of $T$ is weakly increasing and each
label occurs exactly twice in the whole tableau.  If two entries among
$T(1,1),T(1,3),\ldots,T(1,2q-1)$ were equal, then weak increase would force at
least three copies of that label in the first row between the two odd
positions.  This contradicts content $(2^j)$.  Hence $W_q(T)$ has cardinality
$q$.

For the height witness, the first column of $T$ is strictly increasing.
Therefore the entries $T(1,1),T(3,1),\ldots,T(2q-1,1)$ are distinct, and
$H_q(T)$ has cardinality $q$.
\end{proof}

\begin{definition}[Pieri two-strip path]\label{def:bc-pieri-path}
A Pieri two-strip path of length $j$ is a sequence of partitions
\[
  \varnothing=\lambda^{(0)}\subset\lambda^{(1)}\subset\cdots
  \subset\lambda^{(j)}
\]
such that each skew difference
$\lambda^{(t)}/\lambda^{(t-1)}$ is a horizontal strip of size $2$.  The path
is column-even if $(\lambda^{(j)})'$ has only even parts.  For $q\ge1$ define
\[
  C^w_q(\lambda^\bullet)
  =
  \{t:\min(q,\lceil\lambda^{(t)}_1/2\rceil)
      -\min(q,\lceil\lambda^{(t-1)}_1/2\rceil)=1\}
\]
and
\[
  C^h_q(\lambda^\bullet)
  =
  \{t:\min(q,\lceil\ell(\lambda^{(t)})/2\rceil)
      -\min(q,\lceil\ell(\lambda^{(t-1)})/2\rceil)=1\}.
\]
\end{definition}

\begin{lemma}[Pieri paths encode the coefficient tableaux]\label{lem:bc-pieri-path-bijection}
For every $j\ge0$, semistandard tableaux of content $(2^j)$ are in bijection
with Pieri two-strip paths of length $j$.  Under this bijection, the tableau
shape is $\lambda^{(j)}$.
\end{lemma}

\begin{proof}
Given a tableau $T$, let $\lambda^{(t)}$ be the shape of the subtableau
formed by entries at most $t$.  The two boxes labelled $t$ are not in the same
column, because columns are strictly increasing, so
$\lambda^{(t)}/\lambda^{(t-1)}$ is a horizontal strip of size $2$.

Conversely, given such a path, fill every box of
$\lambda^{(t)}/\lambda^{(t-1)}$ with the label $t$.  Along each row, boxes are
added only at the right end as $t$ increases, so rows are weakly increasing.
If a box in row $r>1$ and column $c$ is added at time $t$, then the box
directly above it belongs to $\lambda^{(t)}$ by the partition condition; the
horizontal-strip condition says that this upper box was not added at the same
time, so it was added earlier and has smaller label.  Thus columns are
strictly increasing.  The two constructions are inverse to each other.
\end{proof}

\begin{lemma}[Witnesses as odd-frontier crossings]\label{lem:bc-witness-frontier-crossings}
Let $T$ be a semistandard tableau of content $(2^j)$, and let
\[
  \varnothing=\lambda^{(0)}\subset\lambda^{(1)}\subset\cdots
  \subset\lambda^{(j)}=\lambda(T)
\]
be the Pieri shape path in which $\lambda^{(t)}$ is the shape of the
subtableau of entries at most $t$.  Put
\[
  \omega_q(t)=\min\{q,\lceil \lambda^{(t)}_1/2\rceil\},\qquad
  \eta_q(t)=\min\{q,\lceil \ell(\lambda^{(t)})/2\rceil\}.
\]
If $\lambda_1(T)\ge2q$, then
\[
  W_q(T)=\{t:\omega_q(t)-\omega_q(t-1)=1\}.
\]
If $\ell(\lambda(T))\ge2q$, then
\[
  H_q(T)=\{t:\eta_q(t)-\eta_q(t-1)=1\}.
\]
\end{lemma}

\begin{proof}
The boxes labelled $t$ form the skew shape
$\lambda^{(t)}/\lambda^{(t-1)}$.  Since columns in $T$ are strictly
increasing and the label $t$ occurs exactly twice, this skew shape is a
horizontal two-strip of size $2$.  Hence the first-row length can increase by
at most $2$, and the height can increase by at most $1$.

For the width statement, the box $(1,c)$ has label $t$ exactly when
\[
  \lambda^{(t-1)}_1<c\le\lambda^{(t)}_1 .
\]
Taking $c=2r-1$ with $1\le r\le q$, this is equivalent, because the odd
frontier levels are spaced by $2$ and the first row grows by at most $2$ at
one step, to
\[
  \omega_q(t)-\omega_q(t-1)=1
  \quad\text{and}\quad
  \omega_q(t)=r.
\]
Thus the labels at $(1,1),(1,3),\ldots,(1,2q-1)$ are precisely the labels
whose steps increase $\omega_q$.

For the height statement, the box $(r,1)$ has label $t$ exactly when
\[
  \ell(\lambda^{(t-1)})<r\le\ell(\lambda^{(t)}).
\]
Taking $r=2a-1$ with $1\le a\le q$, and using that a horizontal strip can
create at most one new row, gives the analogous equivalence with
$\eta_q(t)-\eta_q(t-1)=1$.  Hence the first-column witnesses are exactly the
steps increasing $\eta_q$.
\end{proof}

\begin{lemma}[Terminal threshold after all frontier crossings]\label{lem:bc-crossing-terminal-threshold}
Let $\lambda^\bullet$ be a column-even Pieri two-strip path of length $j$,
and fix $q\ge1$.
If $|C^w_q(\lambda^\bullet)|=q$, then $\lambda^{(j)}_1\ge2q-1$, and the
width bad-shape condition is exactly the exclusion
$\lambda^{(j)}_1\ne2q-1$.
If $|C^h_q(\lambda^\bullet)|=q$, then automatically
$\ell(\lambda^{(j)})\ge2q$.
\end{lemma}

\begin{proof}
For the width statistic, the quantity
\[
  \min(q,\lceil\lambda^{(t)}_1/2\rceil)
\]
starts at $0$, is nondecreasing, and can increase by at most $1$ at each
Pieri step, because a horizontal two-strip adds at most two boxes to the first
row.  Therefore $|C^w_q(\lambda^\bullet)|=q$ implies
\[
  \min(q,\lceil\lambda^{(j)}_1/2\rceil)=q,
\]
hence $\lambda^{(j)}_1\ge2q-1$.  Under this conclusion, the inequality
$\lambda^{(j)}_1\ge2q$ is equivalent to $\lambda^{(j)}_1\ne2q-1$.

The height statistic is analogous.  A horizontal strip creates at most one
new row, so
\[
  \min(q,\lceil\ell(\lambda^{(t)})/2\rceil)
\]
also starts at $0$, is nondecreasing, and increases by at most $1$ at each
step.  Thus $|C^h_q(\lambda^\bullet)|=q$ implies
$\ell(\lambda^{(j)})\ge2q-1$.  Since $\lambda^\bullet$ is column-even, the
first column length $\ell(\lambda^{(j)})$ is even whenever it is nonzero.
Here it is nonzero because $q\ge1$, so $\ell(\lambda^{(j)})\ge2q$.
\end{proof}

\begin{lemma}[Local frontier states at crossing steps]\label{lem:bc-local-crossing-states}
Let $\lambda^\bullet$ be a Pieri two-strip path, and list
$C^w_q(\lambda^\bullet)=\{s_1<\cdots<s_m\}$.  For each $1\le r\le m$,
\[
  \lambda^{(s_r)}_1\in\{2r-1,2r\},
\]
and
\[
  \lambda^{(s_r-1)}_1=
  \begin{cases}
  0, & r=1,\\
  2r-3\text{ or }2r-2, & r>1.
  \end{cases}
\]
Similarly, if $C^h_q(\lambda^\bullet)=\{u_1<\cdots<u_m\}$, then for each
$1\le r\le m$,
\[
  \ell(\lambda^{(u_r-1)})=2r-2,\qquad
  \ell(\lambda^{(u_r)})=2r-1.
\]
\end{lemma}

\begin{proof}
For the width statistic, put
\[
  \omega_q(t)=\min(q,\lceil \lambda^{(t)}_1/2\rceil).
\]
The sequence $\omega_q(t)$ starts at $0$, is nondecreasing, and increases by
at most $1$ at each step, since a horizontal two-strip adds at most two boxes
to the first row.  Therefore, at the $r$-th time $s_r$ where it increases,
\[
  \omega_q(s_r-1)=r-1,\qquad \omega_q(s_r)=r.
\]
For $r=1$, the first equality gives $\lambda^{(s_1-1)}_1=0$.  For $r>1$ it
gives $\lambda^{(s_r-1)}_1\in\{2r-3,2r-2\}$.  The second equality gives
$\lambda^{(s_r)}_1\ge2r-1$, while the horizontal two-strip bound gives
$\lambda^{(s_r)}_1\le\lambda^{(s_r-1)}_1+2\le2r$, proving the width claims.

For the height statistic, put
\[
  \eta_q(t)=\min(q,\lceil \ell(\lambda^{(t)})/2\rceil).
\]
At the $r$-th time $u_r$ where it increases, $\eta_q(u_r-1)=r-1$ and
$\eta_q(u_r)=r$.  A horizontal strip creates at most one new row.  If
$\ell(\lambda^{(u_r-1)})\le2r-3$, then
$\ell(\lambda^{(u_r)})\le2r-2$, so $\eta_q$ would not increase to $r$.
Thus $\ell(\lambda^{(u_r-1)})=2r-2$, and the one-row bound forces
$\ell(\lambda^{(u_r)})=2r-1$.
\end{proof}

\begin{lemma}[Local crossing strips]\label{lem:bc-local-crossing-strips}
Let $\lambda^\bullet$ be a Pieri two-strip path.
If $s_r$ is the $r$-th element of $C^w_q(\lambda^\bullet)$, then the
horizontal strip
\[
  \lambda^{(s_r)}/\lambda^{(s_r-1)}
\]
contains the frontier box $(1,2r-1)$, and exactly one of the following three
alternatives holds:
\[
\begin{array}{ll}
\mathrm{(i)}&
\lambda^{(s_r-1)}_1=2r-3,\quad
\lambda^{(s_r)}_1=2r-1,\\
&\text{and the two new first-row boxes are }(1,2r-2),(1,2r-1);\\[2mm]
\mathrm{(ii)}&
\lambda^{(s_r-1)}_1=2r-2,\quad
\lambda^{(s_r)}_1=2r-1,\\
&\text{and the second new box is not in the first row and not in column }2r-1;\\[2mm]
\mathrm{(iii)}&
\lambda^{(s_r-1)}_1=2r-2,\quad
\lambda^{(s_r)}_1=2r,\\
&\text{and the two new first-row boxes are }(1,2r-1),(1,2r).
\end{array}
\]
For $r=1$, only alternative $\mathrm{(iii)}$ can occur.

If $u_r$ is the $r$-th element of $C^h_q(\lambda^\bullet)$, then
$\lambda^{(u_r)}/\lambda^{(u_r-1)}$ contains the frontier box
$(2r-1,1)$, and its other box is not in column $1$.
\end{lemma}

\begin{proof}
By \cref{lem:bc-local-crossing-states}, a width crossing has
\[
  \lambda^{(s_r)}_1\in\{2r-1,2r\}.
\]
Since $\lambda^{(s_r-1)}_1<2r-1\le\lambda^{(s_r)}_1$, the strip contains
the box $(1,2r-1)$.  The same lemma gives
$\lambda^{(s_r-1)}_1=0$ for $r=1$ and
$\lambda^{(s_r-1)}_1\in\{2r-3,2r-2\}$ for $r>1$.  Combining these
possibilities with the fact that the strip has size $2$ gives exactly the
three listed alternatives.  In alternative $\mathrm{(ii)}$, the second new
box cannot lie in the first row because the first-row length increases by
only one, and it cannot lie in column $2r-1$ because the strip is horizontal.
For $r=1$, the path starts with first-row length $0=2r-2$, and a horizontal
strip of size $2$ from the empty partition gives alternative $\mathrm{(iii)}$.

For height, \cref{lem:bc-local-crossing-states} gives
\[
  \ell(\lambda^{(u_r-1)})=2r-2,\qquad
  \ell(\lambda^{(u_r)})=2r-1.
\]
Thus the strip contains the new first-column box $(2r-1,1)$.  Since the strip
is horizontal, its other box cannot also lie in column $1$.
\end{proof}

\begin{lemma}[Odd alternating height crossings create every row]\label{lem:bc-odd-alternating-height-every-row}
Let $q\ge1$ and let
\[
  \lambda^\bullet:\quad
  \varnothing=\lambda^{(0)}\subset\lambda^{(1)}\subset\cdots
  \subset\lambda^{(2q)}
\]
be a column-even Pieri two-strip path of length $2q$.
If
\[
  C^h_q(\lambda^\bullet)=\{1,3,\ldots,2q-1\},
\]
then
\[
  \ell(\lambda^{(t)})=t\qquad(0\le t\le2q).
\]
Consequently every two-strip step creates exactly one new row; at time $t$
the strip $\lambda^{(t)}/\lambda^{(t-1)}$ contains the new first-column box
$(t,1)$, and its other box is not in column $1$.

If instead
\[
  C^h_q(\lambda^\bullet)=\{2,4,\ldots,2q\},
\]
then no such path exists.
\end{lemma}

\begin{proof}
Assume first that $C^h_q(\lambda^\bullet)=\{1,3,\ldots,2q-1\}$.  For
$1\le r\le q$, \cref{lem:bc-local-crossing-states} applied to the $r$-th
height crossing $u_r=2r-1$ gives
\[
  \ell(\lambda^{(2r-2)})=2r-2,
  \qquad
  \ell(\lambda^{(2r-1)})=2r-1 .
\]
This proves the asserted length formula at all odd times and at all even
times up to $2q-2$.  The terminal even time follows from
\cref{lem:bc-crossing-terminal-threshold}, which gives
$\ell(\lambda^{(2q)})\ge2q$, while a horizontal two-strip can create at most
one new row from $\lambda^{(2q-1)}$, whose length is $2q-1$.  Hence
$\ell(\lambda^{(2q)})=2q$.

Thus every step raises the length by exactly one.  The new row at time $t$ is
row $t$, so the strip contains the box $(t,1)$.  Since the strip is horizontal,
its other box cannot lie in the same column.

For the even alternating set, the first Pieri step starts from the empty
partition and adds a horizontal strip of size $2$.  Such a strip must create
the first row, so
\[
  \ell(\lambda^{(0)})=0,\qquad \ell(\lambda^{(1)})=1 .
\]
Therefore $\min(q,\lceil\ell(\lambda^{(1)})/2\rceil)-
\min(q,\lceil\ell(\lambda^{(0)})/2\rceil)=1$, so $1\in C^h_q(\lambda^\bullet)$.
This contradicts $C^h_q(\lambda^\bullet)=\{2,4,\ldots,2q\}$.
\end{proof}

\begin{lemma}[Off-column boxes encode odd alternating height paths]\label{lem:bc-odd-alternating-height-second-box-encoding}
Let $\lambda^\bullet$ be as in
\cref{lem:bc-odd-alternating-height-every-row} with
\[
  C^h_q(\lambda^\bullet)=\{1,3,\ldots,2q-1\}.
\]
For each $1\le t\le2q$, let
\[
  b_t=(r_t,c_t)
\]
be the unique box of
$\lambda^{(t)}/\lambda^{(t-1)}$ different from $(t,1)$.  Then $c_t\ge2$,
and the sequence
\[
  (b_1,\ldots,b_{2q})
\]
uniquely determines the whole Pieri path $\lambda^\bullet$.

Moreover, the terminal shape $\lambda^{(2q)}$ is column-even if and only if,
for every $c\ge2$, the number of indices $t$ with $c_t=c$ is even.  In that
case the off-column boxes have a canonical same-column pairing, obtained by
listing the times with fixed column $c$ increasingly and pairing consecutive
times.
\end{lemma}

\begin{proof}
By \cref{lem:bc-odd-alternating-height-every-row}, the $t$-th two-strip
contains $(t,1)$ and exactly one further box outside column $1$; hence
$c_t\ge2$.  Starting from $\lambda^{(0)}=\varnothing$, the recursion
\[
  \lambda^{(t)}
  =
  \lambda^{(t-1)}\cup\{(t,1),b_t\}
  \qquad(1\le t\le2q)
\]
recovers every intermediate partition from the displayed sequence of boxes.
Thus the sequence determines the path.

At the terminal time, column $1$ has exactly the boxes
$(1,1),\ldots,(2q,1)$, so its length is $2q$.  Every terminal box outside
column $1$ is one of the off-column boxes $b_t$, because each step adds only
the first-column box and $b_t$, and boxes are never removed.  Therefore, for
each $c\ge2$, the length of column $c$ in $\lambda^{(2q)}$ is precisely
\[
  \#\{t:c_t=c\}.
\]
The terminal shape is column-even exactly when all these numbers are even.
When they are even, increasing order of the finitely many times in each fixed
column gives a canonical pairing of those times, and hence of the
corresponding off-column boxes.
\end{proof}

\begin{lemma}[Row choices encode the off-column boxes]\label{lem:bc-odd-alternating-height-row-choice-encoding}
In the setting of
\cref{lem:bc-odd-alternating-height-second-box-encoding}, the row indices
\[
  r_1,\ldots,r_{2q}
\]
determine the off-column boxes $b_t=(r_t,c_t)$ and hence determine the whole
Pieri path.  More explicitly, for each $t$ one has $1\le r_t\le t$, and
\[
  c_t=
  \begin{cases}
    \lambda^{(t-1)}_{r_t}+1, & r_t<t,\\
    2, & r_t=t.
  \end{cases}
\]
Thus, recursively from $\lambda^{(0)}=\varnothing$, the row-choice sequence
determines
\[
  \lambda^{(t)}
  =
  \lambda^{(t-1)}
  \cup\{(t,1),(r_t,c_t)\}.
\]
\end{lemma}

\begin{proof}
By \cref{lem:bc-odd-alternating-height-every-row},
$\ell(\lambda^{(t)})=t$ and $\ell(\lambda^{(t-1)})=t-1$.  Hence the
off-column box added at time $t$ lies in one of the rows
$1,\ldots,t$, proving $1\le r_t\le t$.

If $r_t<t$, then row $r_t$ already exists in $\lambda^{(t-1)}$.  Since
$\lambda^{(t)}/\lambda^{(t-1)}$ is obtained by adding boxes to a Young
diagram, the new box in row $r_t$ must be the first box after the previous row
end, namely at column $\lambda^{(t-1)}_{r_t}+1$.

If $r_t=t$, then the row did not exist in $\lambda^{(t-1)}$.  The step already
adds $(t,1)$, and the off-column box is in the same new row.  A Young diagram
cannot contain $(t,c)$ with $c>2$ unless it also contains $(t,2)$, but only
two boxes are added at this step.  Since the off-column box is not $(t,1)$,
it must be $(t,2)$.

Thus $c_t$ is determined by the previous partition and $r_t$ in the displayed
way.  Starting from $\lambda^{(0)}=\varnothing$, this recursion determines all
$b_t$ and all $\lambda^{(t)}$.  The final claim also follows from
\cref{lem:bc-odd-alternating-height-second-box-encoding}.
\end{proof}

\begin{lemma}[Odd alternating height paths are even-column standard tableaux]\label{lem:bc-odd-alternating-height-standard-tableau}
Let $q\ge1$.  The map that sends a column-even Pieri two-strip path
\[
  \lambda^\bullet:\quad
  \varnothing=\lambda^{(0)}\subset\cdots\subset\lambda^{(2q)}
\]
with
\[
  C^h_q(\lambda^\bullet)=\{1,3,\ldots,2q-1\}
\]
to its shifted off-column recording tableau is a bijection onto standard Young
tableaux of size $2q$ whose columns all have even length.

Explicitly, if $b_t=(r_t,c_t)$ is the off-column box of the $t$-th strip, put
the entry $t$ in the shifted box $(r_t,c_t-1)$.  The resulting tableau is
standard and has only even columns.  Conversely, given a standard Young
tableau $U$ of size $2q$ with even columns, let $\nu^{(t)}$ be the shape of
the boxes of $U$ with entries at most $t$, and define
\[
  \lambda^{(t)}
  =
  \{(i,1):1\le i\le t\}
  \cup
  \{(i,c+1):(i,c)\in\nu^{(t)}\}.
\]
Then $\lambda^\bullet$ is a column-even Pieri two-strip path with the displayed
odd alternating height-crossing set, and the two constructions are inverse.
\end{lemma}

\begin{proof}
Start with an odd alternating height path.  By
\cref{lem:bc-odd-alternating-height-every-row}, each step adds the first-column
box $(t,1)$ and one off-column box $b_t=(r_t,c_t)$.  Shift all off-column
boxes one column left.  At time $t$, the shifted shape is obtained from the
previous shifted shape by adding the single box $(r_t,c_t-1)$, because the
unshifted off-column boxes are exactly the boxes outside column $1$ of
$\lambda^{(t)}$.  These shifted shapes are therefore a saturated Young-lattice
chain from $\varnothing$ to a partition of size $2q$.  Filling the box added
at time $t$ with $t$ gives a standard Young tableau.  The terminal shifted
shape has column lengths equal to the terminal column lengths of
$\lambda^{(2q)}$ in columns $2,3,\ldots$, so its columns are even by
column-evenness of $\lambda^{(2q)}$.

Conversely, let $U$ be a standard Young tableau of size $2q$ with even-column
shape, and let $\nu^{(t)}$ be its standard growth chain.  The displayed formula
for $\lambda^{(t)}$ adds a first column of height $t$ and shifts
$\nu^{(t)}$ one column to the right.  Since $\nu^{(t)}$ is a partition, the
result is a partition.  From $t-1$ to $t$, the first-column box $(t,1)$ is
added, and the unique box of $\nu^{(t)}/\nu^{(t-1)}$ is added one column to
the right.  Thus $\lambda^{(t)}/\lambda^{(t-1)}$ is a horizontal strip of
size $2$: the two new boxes lie in columns $1$ and at least $2$.

The terminal first column has length $2q$, and every other terminal column of
$\lambda^{(2q)}$ is a column of the even-column shape of $U$.  Hence
$\lambda^\bullet$ is column-even.  Also $\ell(\lambda^{(t)})=t$ for every
$t$, so
\[
  \min(q,\lceil t/2\rceil)-\min(q,\lceil (t-1)/2\rceil)
\]
is $1$ exactly when $t$ is odd and $1\le t\le2q-1$.  Therefore
$C^h_q(\lambda^\bullet)=\{1,3,\ldots,2q-1\}$.

The two constructions are inverse because one deletes the forced first column
and shifts the remaining boxes left, while the other restores that first
column and shifts the standard growth chain right.
\end{proof}

\begin{lemma}[Schensted contraction of the odd alternating height model]\label{lem:bc-odd-alternating-height-schensted-contraction}
Let $q\ge1$ and let $\lambda^\bullet$ be a column-even Pieri two-strip path
of length $2q$ with
\[
  C^h_q(\lambda^\bullet)=\{1,3,\ldots,2q-1\}.
\]
Let $U$ be the even-column standard tableau associated to $\lambda^\bullet$
by \cref{lem:bc-odd-alternating-height-standard-tableau}, and let $\pi$ be the
fixed-point-free involution on $\{1,\ldots,2q\}$ whose Schensted insertion
tableau is $U$.  Contract the adjacent pairs
\[
  \{1,2\},\{3,4\},\ldots,\{2q-1,2q\}
\]
to vertices $1,\ldots,q$.  Each transposition of $\pi$ then gives an edge
between the corresponding vertices, with a loop exactly when the transposition
is one of the adjacent pairs $\{2v-1,2v\}$.

The resulting contracted object is a labelled degree-two multigraph
$\Gamma^\circ(\lambda^\bullet)$ on $\{1,\ldots,q\}$, with parallel edges and
loops allowed.  Moreover there is a word
\[
  \epsilon(\lambda^\bullet)\in\{0,1\}^{2q}
\]
such that the pair $(\Gamma^\circ(\lambda^\bullet),\epsilon(\lambda^\bullet))$
recovers $\lambda^\bullet$.

If $\Gamma^\circ(\lambda^\bullet)$ has no loops, then it is a loopless
labelled $2$-regular multigraph counted by the stable moment $s_q$ in
\cref{lem:bc-littlewood-graph-model}.  Thus the loopless-contraction subfamily
of the odd alternating height fiber satisfies the required
$2^{2q}s_q$ bound.  The only obstruction to this direct stable-graph target is
the presence of adjacent transpositions $\{2v-1,2v\}$ in $\pi$.
\end{lemma}

\begin{proof}
For a standard tableau $U$, the inverse Robinson--Schensted correspondence
with equal insertion and recording tableaux gives a unique involution $\pi$
with insertion tableau $U$.  Rains \cite[Lemma 3.3]{Rains1997} records the
classical Schensted fact that the number of fixed points of this involution is
the number of odd columns of $U$.  Since $U$ has only even columns, $\pi$ has
no fixed points; hence it is a fixed-point-free involution, i.e. a perfect
matching of $\{1,\ldots,2q\}$.

After contracting the adjacent pairs to vertices $1,\ldots,q$, every vertex
has two incident ports, namely $2v-1$ and $2v$, and each port is matched
exactly once by $\pi$.  Therefore the contracted multigraph has degree two at
every vertex, counting a loop as degree two.  A transposition contributes a
loop precisely when both of its endpoints lie in the same adjacent pair,
which is exactly the displayed adjacent-pair case.

It remains to record why $2q$ bits are enough to recover the original matching
from the contracted multigraph.  For each vertex $v$, order its two incident
edge-slots canonically from the contracted multigraph: first by the other
endpoint label, then by the multiplicity number of the edge, and for a loop by
its two loop-slots.  Let the two bits attached to the ports $2v-1$ and $2v$
record which of these two ordered slots receives each port.  The word
$\epsilon(\lambda^\bullet)$ is the collection of these $2q$ slot choices.  The
contracted multigraph together with this slot assignment reconstructs every
matched pair of ports, hence reconstructs $\pi$.

The Schensted map recovers $U$ from $\pi$, and
\cref{lem:bc-odd-alternating-height-standard-tableau} then recovers
$\lambda^\bullet$ from $U$.  Hence the contracted graph and the word
$\epsilon(\lambda^\bullet)$ recover the original path.

If there are no loops, then $\Gamma^\circ(\lambda^\bullet)$ is loopless and
degree two at every vertex, with parallel edges allowed.  By
\cref{lem:bc-littlewood-graph-model}, the number of possible loopless
contracted graphs is $s_q$, and the number of possible bit words is at most
$2^{2q}$.  Hence the loopless-contraction subfamily has size at most
$2^{2q}s_q$.  The preceding loop criterion gives the final sentence.
\end{proof}

\begin{lemma}[Loop repair for contracted degree-two graphs]\label{lem:bc-looped-degree-two-graph-repair}
Let $q\ge2$, and let $\Gamma^\circ$ be a labelled degree-two multigraph on
$\{1,\ldots,q\}$, with parallel edges and loops allowed, where each loop
contributes degree $2$.  Let
\[
  L=\{v:\Gamma^\circ\text{ has a loop at }v\}.
\]
Then there is a canonically chosen loopless labelled degree-two multigraph
\[
  \mathcal R(\Gamma^\circ)
\]
on $\{1,\ldots,q\}$, determined by $\Gamma^\circ$ and the usual order on
$\{1,\ldots,q\}$, such that $\Gamma^\circ$ is recoverable from
$\mathcal R(\Gamma^\circ)$ and $L$.
\end{lemma}

\begin{proof}
Because a loop contributes degree $2$, a vertex in $L$ has exactly one loop
and no other incident edge.  Hence the subgraph on the complement
$L^c$ is already loopless and degree two at every vertex of $L^c$.

If $L=\varnothing$, set $\mathcal R(\Gamma^\circ)=\Gamma^\circ$.  Suppose
next that $|L|$ is even.  Delete the loops on $L$, list the vertices of $L$
in increasing order, pair them consecutively, and put two parallel edges on
each pair.  This gives a loopless degree-two graph on $L$, disjoint from the
unchanged graph on $L^c$.

If $|L|$ is odd and at least $3$, delete the loops on $L$, put a triangle on
the first three vertices of $L$, and put two parallel edges on each
consecutive pair of the remaining vertices of $L$.  Again this gives a
loopless degree-two graph on $L$, disjoint from the unchanged graph on $L^c$.

It remains to handle $|L|=1$, say $L=\{v\}$.  If $L^c$ had only one vertex,
then that vertex would have degree two in the loopless graph on $L^c$, which
is impossible.  Thus $L^c$ has at least two vertices.  Its loopless degree-two
graph therefore has at least one non-loop edge.  Choose the lexicographically
least edge $ab$ in $L^c$, counting parallel edges by their multiplicity order.
Delete one copy of $ab$, delete the loop at $v$, and add the two edges
$va$ and $vb$.  The resulting graph is loopless and degree two at every
vertex.

The recovery from $\mathcal R(\Gamma^\circ)$ and $L$ is inverse to these
constructions.  If $L=\varnothing$, do nothing.  If $|L|$ is even or at least
$3$, the repair edges are exactly the edges whose two endpoints lie in $L$;
delete them and restore one loop at every vertex of $L$.  If $L=\{v\}$, the
two edges incident to $v$ in $\mathcal R(\Gamma^\circ)$ have other endpoints
$a$ and $b$ with $a\ne b$; delete those two edges, restore the loop at $v$,
and add one edge $ab$.  This recovers the original graph in every case.
\end{proof}

\begin{proposition}[Odd alternating height full-zone decoder]\label{prop:bc-odd-alternating-height-full-zone-decoder}
For every $q\ge2$, the height fixed-witness fiber with
\[
  j=2q,\qquad S=\{1,3,\ldots,2q-1\}
\]
satisfies
\[
  \#\mathcal F^h(S)\le 2^{2q}s_q .
\]
Equivalently, the odd alternating height full-zone case of
\cref{prop:bc-alternating-full-zone-decoder-criterion} has the required
decoder.
\end{proposition}

\begin{proof}
By \cref{cor:bc-fixed-fibers-pieri-paths}, this fixed-witness fiber is
equinumerous with the column-even Pieri paths of length $2q$ with
$C^h_q=\{1,3,\ldots,2q-1\}$.  For such a path $\lambda^\bullet$, apply
\cref{lem:bc-odd-alternating-height-standard-tableau} and then
\cref{lem:bc-odd-alternating-height-schensted-contraction}.  This gives a
contracted degree-two graph $\Gamma^\circ(\lambda^\bullet)$ with loops
allowed.  Apply \cref{lem:bc-looped-degree-two-graph-repair} to get the
loopless graph
\[
  H(\lambda^\bullet)=\mathcal R(\Gamma^\circ(\lambda^\bullet)).
\]
Fix once and for all the lexicographic order on loopless labelled degree-two
multigraphs and the lexicographic order on the finite set
$\mathcal P_q^{\mathrm{even}}$ of two-step growth paths.  Since
\cref{lem:bc-littlewood-graph-model,lem:bc-growth-path-model} identify both
sets as having cardinality $s_q$, these orders give a fixed bijection between
loopless graphs and paths in $\mathcal P_q^{\mathrm{even}}$.  Let
$P^\ast(\lambda^\bullet)$ be the path corresponding to $H(\lambda^\bullet)$.

The branch word has two parts.  The first $q$ bits are the characteristic
function of the loop set
\[
  L(\lambda^\bullet)=\{v:\Gamma^\circ(\lambda^\bullet)\text{ has a loop at }v\}.
\]
Given $P^\ast(\lambda^\bullet)$, the fixed inverse bijection recovers
$H(\lambda^\bullet)$.  After $H(\lambda^\bullet)$ and this loop set are known,
\cref{lem:bc-looped-degree-two-graph-repair} recovers
$\Gamma^\circ(\lambda^\bullet)$.  The remaining $q$ bits record, for each
vertex $v$, which of the two canonically ordered incident edge-slots of
$\Gamma^\circ(\lambda^\bullet)$ receives the port $2v-1$; the port $2v$ then
receives the other slot.  This recovers the fixed-point-free involution
$\pi$, hence its Schensted tableau $U$, and finally recovers
$\lambda^\bullet$ by
\cref{lem:bc-odd-alternating-height-standard-tableau}.

Thus the assignment
\[
  \lambda^\bullet\longmapsto
  \bigl(P^\ast(\lambda^\bullet),\text{the }2q\text{-bit branch word}\bigr)
\]
is injective, with a decoder depending only on the alternating set $S$.  Since
$|\mathcal P_q^{\mathrm{even}}|=s_q$ and there are $2^{2q}$ possible branch
words, this gives the claimed bound.
\end{proof}

\begin{proposition}[Even alternating height full-zone fiber is empty]\label{prop:bc-even-alternating-height-full-zone-empty}
For every $q\ge1$, the height fixed-witness fiber with
\[
  j=2q,\qquad S=\{2,4,\ldots,2q\}
\]
is empty.  In particular it satisfies
\[
  \#\mathcal F^h(S)\le 2^{2q}s_q .
\]
\end{proposition}

\begin{proof}
By \cref{cor:bc-fixed-fibers-pieri-paths}, this fixed-witness fiber is
equinumerous with the column-even Pieri paths of length $2q$ with
$C^h_q=\{2,4,\ldots,2q\}$.  The second half of
\cref{lem:bc-odd-alternating-height-every-row} proves that no such path
exists.  The cardinality is therefore zero, giving the displayed bound.
\end{proof}

\begin{corollary}[Height alternating full-zone decoder]\label{cor:bc-height-alternating-full-zone-decoder}
For every $q\ge2$ and
\[
  S\in
  \bigl\{\{1,3,\ldots,2q-1\},\{2,4,\ldots,2q\}\bigr\},
\]
the height fixed-witness fiber $\mathcal F^h(S)$ in the alternating full-zone
case of \cref{prop:bc-alternating-full-zone-decoder-criterion} satisfies
\[
  \#\mathcal F^h(S)\le 2^{2q}s_q .
\]
\end{corollary}

\begin{proof}
For $S=\{1,3,\ldots,2q-1\}$ this is
\cref{prop:bc-odd-alternating-height-full-zone-decoder}.  For
$S=\{2,4,\ldots,2q\}$ it is
\cref{prop:bc-even-alternating-height-full-zone-empty}.
\end{proof}

\begin{lemma}[Odd alternating width full-zone path is forced]\label{lem:bc-odd-alternating-width-forced-path}
Let $q\ge1$, and let $\lambda^\bullet$ be a column-even Pieri two-strip path
of length $2q$ such that
\[
  C^w_q(\lambda^\bullet)=\{1,3,\ldots,2q-1\},
  \qquad
  \lambda^{(2q)}_1\ge2q .
\]
Then, for $0\le r\le q$,
\[
  \lambda^{(2r)}=(2r,2r),
\]
with the convention $\lambda^{(0)}=\varnothing$, and for $1\le r\le q$,
\[
  \lambda^{(2r-1)}=(2r,2r-2).
\]
In particular there is exactly one such path.
\end{lemma}

\begin{proof}
The terminal partition has size $4q$.  Since it is column-even, every column
has even positive height, hence height at least $2$.  The hypothesis
$\lambda^{(2q)}_1\ge2q$ therefore gives
\[
  4q=|\lambda^{(2q)}|\ge 2\lambda^{(2q)}_1\ge4q.
\]
Thus $\lambda^{(2q)}_1=2q$, every column has height exactly $2$, and
\[
  \lambda^{(2q)}=(2q,2q).
\]
All intermediate partitions are contained in this rectangle, so they have at
most two rows.

Fix $1\le r\le q$.  The time $2r-1$ is the $r$-th width crossing.  By
\cref{lem:bc-local-crossing-states},
\[
  \lambda^{(2r-1)}_1\in\{2r-1,2r\}.
\]
If $\lambda^{(2r-1)}_1=2r-1$, then the size
$|\lambda^{(2r-1)}|=4r-2$ and the two-row bound force
\[
  \lambda^{(2r-1)}=(2r-1,2r-1).
\]
For $r<q$, the next time $2r$ is not a width crossing, so
$\lambda^{(2r)}_1\le2r$.  From $(2r-1,2r-1)$, a horizontal two-strip inside
the terminal two-row rectangle and with first-row length at most $2r$ cannot
be added: the only possible first-row addition is $(1,2r)$, while the
matching second-row addition $(2,2r)$ lies in the same column.  For $r=q$,
the same obstruction prevents the final step from reaching
$(2q,2q)$.  Hence this case is impossible, and
\[
  \lambda^{(2r-1)}_1=2r.
\]
The size $4r-2$ and the two-row bound then force
\[
  \lambda^{(2r-1)}=(2r,2r-2).
\]

For $r<q$, the time $2r$ is not a width crossing, so
$\lambda^{(2r)}_1\le2r$.  Since $|\lambda^{(2r)}|=4r$ and
$\lambda^{(2r)}$ has at most two rows, this forces
$\lambda^{(2r)}=(2r,2r)$.  The case $r=q$ is the terminal rectangle already
proved.  This gives the displayed path at every time, and uniqueness follows.
\end{proof}

\begin{proposition}[Width alternating full-zone decoder]\label{prop:bc-width-alternating-full-zone-decoder}
For every $q\ge2$ and
\[
  S\in
  \bigl\{\{1,3,\ldots,2q-1\},\{2,4,\ldots,2q\}\bigr\},
\]
the width fixed-witness fiber $\mathcal F^w(S)$ in the alternating full-zone
case of \cref{prop:bc-alternating-full-zone-decoder-criterion} satisfies
\[
  \#\mathcal F^w(S)\le 2^{2q}s_q .
\]
\end{proposition}

\begin{proof}
If $S=\{2,4,\ldots,2q\}$, the fiber is empty.  Indeed, the first Pieri step
from $\varnothing$ is necessarily the horizontal strip giving the partition
$(2)$, so $1\in C^w_q(\lambda^\bullet)$ for every path.  By
\cref{cor:bc-fixed-fibers-pieri-paths}, no width fixed-witness path has the
even alternating crossing set.

It remains to consider $S=\{1,3,\ldots,2q-1\}$.  By
\cref{cor:bc-fixed-fibers-pieri-paths}, this fiber is equinumerous with
column-even Pieri paths of length $2q$ with this crossing set and terminal
first-row length at least $2q$.  By
\cref{lem:bc-odd-alternating-width-forced-path}, there is exactly one such
path.

For $q\ge2$ the stable graph model is nonempty: for $q=2$ take the graph with
two parallel edges between the two vertices, and for $q\ge3$ take the
labelled cycle on $\{1,\ldots,q\}$.  Hence $s_q\ge1$ by
\cref{lem:bc-littlewood-graph-model}.  Therefore the singleton odd fiber, and
the empty even fiber, both have cardinality at most $2^{2q}s_q$.
\end{proof}

\begin{corollary}[Alternating full-zone decoder is closed]\label{cor:bc-alternating-full-zone-decoder-closed}
For every $q\ge2$, for both alternating sets
\[
  S\in
  \bigl\{\{1,3,\ldots,2q-1\},\{2,4,\ldots,2q\}\bigr\},
\]
and for either fixed-witness fiber $\mathcal F^w(S)$ or $\mathcal F^h(S)$ in
\cref{prop:bc-alternating-full-zone-decoder-criterion}, the fixed fiber has
cardinality at most $2^{2q}s_q$.
\end{corollary}

\begin{proof}
The height fibers are handled by
\cref{cor:bc-height-alternating-full-zone-decoder}; the width fibers are
handled by \cref{prop:bc-width-alternating-full-zone-decoder}.
\end{proof}

\begin{corollary}[Fixed fibers as Pieri crossing paths]\label{cor:bc-fixed-fibers-pieri-paths}
Fix $q\ge1$, $j\ge q$, and $S\subset\{1,\ldots,j\}$ with $|S|=q$.  The width
fiber
\[
 \{T:\lambda(T)' \text{ even},\ \lambda_1(T)\ge2q,\ W_q(T)=S\}
\]
is equinumerous with the column-even Pieri two-strip paths
$\lambda^\bullet$ of length $j$ such that $C^w_q(\lambda^\bullet)=S$ and
$\lambda^{(j)}_1\ge2q$.  The height fiber
\[
 \{T:\lambda(T)' \text{ even},\ \ell(\lambda(T))\ge2q,\ H_q(T)=S\}
\]
is equinumerous with the column-even Pieri two-strip paths
$\lambda^\bullet$ of length $j$ such that $C^h_q(\lambda^\bullet)=S$.
\end{corollary}

\begin{proof}
Use the bijection of \cref{lem:bc-pieri-path-bijection}.  The terminal
column-even condition is exactly $\lambda(T)'$ even.  By
\cref{lem:bc-witness-frontier-crossings}, the width witness is the crossing
set $C^w_q(\lambda^\bullet)$ whenever $\lambda_1(T)\ge2q$, and the height
witness is $C^h_q(\lambda^\bullet)$ whenever $\ell(\lambda(T))\ge2q$.  In
the height case the latter terminal condition is automatic from
\cref{lem:bc-crossing-terminal-threshold}, because $|S|=q$.
\end{proof}

\begin{lemma}[Bounded column defect after label deletion]\label{lem:bc-deletion-defect}
Let $T$ be a semistandard tableau of shape $\lambda$ and content $(2^j)$,
with $\lambda'$ even.  For a subset $S\subset\{1,\ldots,j\}$, let
$T\setminus S$ be the labelled retained subdiagram obtained by deleting from
$T$ the two boxes carrying each label in $S$, while keeping the original row
and column coordinates of all remaining boxes.  Then $T\setminus S$ has
content $(2^{j-|S|})$ on the complementary labels, its remaining entries are
weakly increasing along each retained row and strictly increasing along each
retained column, and at most $2|S|$ columns contain an odd number of remaining
boxes.
\end{lemma}

\begin{proof}
The row weak-increase and column strict-increase conditions are inherited by
subsets of boxes in each row and column.  The content statement is immediate
from the definition of content $(2^j)$.
For a column $c$, let $d_c$ be the number of deleted boxes in that column.
Since $\lambda'$ is even, column $c$ of $\lambda$ contains an even number of
boxes.  Hence the parity of the number of boxes remaining in column $c$ of
$T\setminus S$ is the parity of $d_c$.  Therefore a column of $T\setminus S$
can be odd only if it contains at least one deleted box.  The total number of
deleted boxes is $2|S|$, so at most $2|S|$ columns can be odd.
\end{proof}

\begin{lemma}[Label deletion need not give a skew shape]\label{lem:bc-deletion-not-skew}
The retained subdiagram in \cref{lem:bc-deletion-defect} is not, in general,
a skew Young diagram.  This remains true for column-even semistandard tableaux
of content $(2^j)$.
\end{lemma}

\begin{proof}
Take the semistandard tableau of shape $(4,4)$
\[
  \begin{array}{cccc}
  1&1&2&2\\
  3&3&4&4
  \end{array}
\]
and delete the two boxes labelled $3$.  The original shape has four columns,
each of height $2$, hence $\lambda'$ is even, and the content is $(2^4)$.
After deletion, the second row has boxes only in columns $3$ and $4$, while
the first row has boxes in columns $1,2,3,4$.  If this retained set were
$\alpha/\beta$ for partitions $\beta\subset\alpha$, then the number of boxes
deleted from the left of the first two rows would be
$\beta_1=0$ and $\beta_2=2$, contradicting $\beta_1\ge\beta_2$.  Thus the
retained set is not a skew Young diagram.
\end{proof}

\begin{lemma}[Row packing does not repair witness deletion]\label{lem:bc-row-packing-fails}
For every $q\ge1$ there is a column-even semistandard tableau $T$ of content
$(2^{2q+1})$ with $\lambda_1(T)\ge2q$ such that, for its width witness
$S=W_q(T)$, deleting all labels in $S$ and then left-packing each retained row
does not produce a Young diagram.
\end{lemma}

\begin{proof}
Let $T$ have shape $(2q+1,2q+1)$, with first row
\[
  1,1,2,2,\ldots,q,q,q+1
\]
and second row
\[
  q+1,q+2,q+2,\ldots,2q+1,2q+1.
\]
Rows are weakly increasing, columns are strictly increasing, and each label
$1,\ldots,2q+1$ occurs exactly twice.  The shape has two boxes in every
column, so it is column-even.  The width witness is
$W_q(T)=\{1,\ldots,q\}$, read from the first-row odd positions
$1,3,\ldots,2q-1$.  Deleting the two copies of every label in $W_q(T)$ leaves
one box in the first row, labelled $q+1$, and all $2q+1$ boxes in the second
row.  Left-packing rows therefore gives row lengths $(1,2q+1)$, which are not
weakly decreasing.  Hence the packed rows do not form a Young diagram.
\end{proof}

\begin{lemma}[Column-even tableaux as two-step growth paths]\label{lem:bc-growth-path-model}
Let $T$ be a semistandard tableau of content $(2^j)$ whose shape has only even
column lengths, and let $k$ be any integer such that every column of $T$ has
length at most $2k$.  Krattenthaler's semistandard growth-diagram bijection
\cite[Theorem 4]{Krattenthaler2016} encodes $T$ by a sequence of partitions
\[
  \varnothing=\rho_0,\rho_1,\ldots,\rho_{2j}=\varnothing
\]
such that every $\rho_t$ has at most $k$ columns, each adjacent pair differs
by a vertical strip, and for every $i=1,\ldots,j$ the two-step block
\[
  \rho_{2i-2}\supseteq \rho_{2i-1}\subseteq\rho_{2i}
\]
corresponding, under Krattenthaler's index convention, to tableau label
$j-i+1$ satisfies
\[
  |\rho_{2i-2}|-2|\rho_{2i-1}|+|\rho_{2i}|=2 .
\]
Conversely, such sequences are counted by the same tableaux.
\end{lemma}

\begin{proof}
Krattenthaler's Theorem 4 gives a bijective equality between semistandard
tableaux with prescribed content $j_1,\ldots,j_n$, $m$ odd columns, and all
columns of length at most $2k$, and sequences
\[
  \varnothing=\rho_0,\rho_1,\ldots,\rho_{2n}=(1^m)
\]
with at most $k$ columns, vertical-strip adjacent differences, and
\[
  |\rho_{2i-2}|-2|\rho_{2i-1}|+|\rho_{2i}|=j_{n-i+1}.
\]
Apply this with $n=j$, $m=0$, and $j_1=\cdots=j_j=2$.  Since all prescribed
contents are equal to $2$, the reversal of the index in the displayed source
formula does not change the two-step condition; it only identifies source
block $i$ with tableau label $j-i+1$, or equivalently label $\ell$ with
block $j-\ell+1$.
\end{proof}

\begin{lemma}[Knuth growth-diagram cells are locally reversible]\label{lem:bc-knuth-growth-local-reversibility}
In the first arbitrary-filling variant of Krattenthaler's growth diagrams,
\cite[\S4.1, Theorem 7]{Krattenthaler2006}, a cell filled by a nonnegative
integer $m$ with corner labels
\[
  \begin{matrix}
  \nu & \lambda\\
  \rho & \mu
  \end{matrix}
\]
is governed by mutually inverse forward and backward algorithms.  The forward
algorithm determines $\lambda$ from $(m,\rho,\mu,\nu)$ when
$\rho\subseteq\mu$ and $\rho\subseteq\nu$ differ by horizontal strips.  The
backward algorithm determines $(m,\rho)$ from $(\lambda,\mu,\nu)$ when
$\mu\subseteq\lambda$ and $\nu\subseteq\lambda$ differ by horizontal strips.
Consequently, on any finite Ferrers cell arrangement in this variant, the
left and bottom boundary labels together with the cell entries determine the
top and right boundary labels, and the top and right boundary labels determine
the cell entries and the remaining corner labels.
\end{lemma}

\begin{proof}
Krattenthaler \cite[\S4.1]{Krattenthaler2006} constructs arbitrary-filling
growth diagrams by separating each cell entry into a $0$--$1$ blow-up, applying
the ordinary Fomin local rules, and shrinking back to the original diagram.
Immediately after Figure 9, the source records the equivalent direct cell
algorithms: the forward algorithm computes the northeast corner $\lambda$ from
the cell entry and the other three corners, while the backward algorithm
computes the southwest corner $\rho$ and the cell entry from the other three
corners.  The theorem following those algorithms states the corresponding
boundary bijection for Ferrers shapes.  Iterating the forward cell rule across
the cells gives the first global determination statement, and iterating the
backward cell rule from the top-right boundary gives the second.
\end{proof}

\begin{corollary}[Parsed boundaries recover Knuth subfillings]\label{cor:bc-parsed-boundary-recovers-knuth-subfilling}
Let $F$ be a finite Ferrers cell arrangement equipped with the first
arbitrary-filling growth diagram of
\cref{lem:bc-knuth-growth-local-reversibility}, and let $E\subseteq F$ be a
Ferrers subarrangement.  Suppose that a parser has recovered the partition
labels on the top/right boundary of $E$.  Then the shape of $E$ and those
boundary labels determine, by a deterministic backward sweep, every cell entry
of $E$ and every corner label on or inside $E$.

Consequently, in the right-boundary auxiliary setting of
\cref{prop:bc-right-boundary-auxiliary-global-recovery}, any auxiliary datum
which is a fixed function of such a recovered subfilling is already a fixed
function of the parsed outside growth-boundary data and requires no additional
branch word.
\end{corollary}

\begin{proof}
The subarrangement $E$ is itself a finite Ferrers cell arrangement in the
orientation used by Krattenthaler's local rules.  By
\cref{lem:bc-knuth-growth-local-reversibility}, the backward cell algorithm
determines the southwest corner and the cell entry of a cell from its
northeast, southeast, and northwest corner labels.  Starting from the known
top/right boundary of $E$ and sweeping cells leftward and downward therefore
recovers every cell entry and every remaining corner label in $E$.

If the right-boundary auxiliary datum is defined from one of these recovered
subfillings, then the preceding deterministic recovery expresses that datum as
a fixed function of the already parsed outside boundary labels.  Thus it is an
auxiliary input of the type allowed in
\cref{prop:bc-right-boundary-auxiliary-global-recovery}, not an extra encoded
choice.
\end{proof}

\begin{lemma}[Dual RSK' growth-diagram cells are locally reversible]\label{lem:bc-dual-rsk-prime-growth-local-reversibility}
In the fourth arbitrary-filling variant of Krattenthaler's growth diagrams,
\cite[\S4.4, Theorem 11]{Krattenthaler2006}, a cell filled by a
nonnegative integer $m$ with corner labels
\[
  \begin{matrix}
  \nu & \lambda\\
  \rho & \mu
  \end{matrix}
\]
is governed by mutually inverse forward and backward algorithms with vertical
strip adjacent differences.  The forward algorithm determines $\lambda$ from
$(m,\rho,\mu,\nu)$ when $\rho\subseteq\mu$ and $\rho\subseteq\nu$ differ by
vertical strips.  The backward algorithm determines $(m,\rho)$ from
$(\lambda,\mu,\nu)$ when $\mu\subseteq\lambda$ and
$\nu\subseteq\lambda$ differ by vertical strips.

Consequently, on any finite Ferrers cell arrangement in this variant, the
left and bottom boundary labels together with the cell entries determine the
top and right boundary labels, and the top and right boundary labels determine
the cell entries and the remaining corner labels.
\end{lemma}

\begin{proof}
Krattenthaler \cite[\S4.4]{Krattenthaler2006} recalls the fourth
arbitrary-filling growth-diagram variant, denoted dual RSK'.  Immediately
before Theorem 11 the source gives the direct cell algorithms (F4) and (B4):
(F4) computes the northeast corner $\lambda$ from the cell entry and the
other three corners, while (B4) computes the southwest corner $\rho$ and the
cell entry from the other three corners.  The hypotheses in those algorithms
are exactly the vertical-strip containment conditions displayed above.
Theorem 11 then states the corresponding boundary bijection for Ferrers
shapes.  Iterating (F4) across the cells gives the forward determination, and
iterating (B4) from the top/right boundary gives the backward determination.
\end{proof}

\begin{lemma}[Dual-RSK' cell size conservation]\label{lem:bc-dual-rsk-prime-cell-size-conservation}
In the one-cell fourth arbitrary-filling growth diagram of
\cref{lem:bc-dual-rsk-prime-growth-local-reversibility}, let the cell entry be
$m$ and write the corner labels as
\[
  \begin{matrix}
  \nu & \lambda\\
  \rho & \mu .
  \end{matrix}
\]
Then
\[
  |\lambda|-|\mu|-|\nu|+|\rho|=m .
\]
\end{lemma}

\begin{proof}
Use Krattenthaler's forward algorithm (F4) from
\cite[\S4.4]{Krattenthaler2006}.  Extend all partitions by zero parts.  Let
$c_i$ be the carry entering row $i$, with $c_1=m$.  Put
\[
  a_i=\chi(\rho_i=\mu_i=\nu_i),\qquad
  d_i=\min(\mu_i,\nu_i)-\rho_i .
\]
Because $\mu/\rho$ and $\nu/\rho$ are vertical strips, $d_i$ is either $0$ or
$1$.  The forward rule sets
\[
  b_i=\min(a_i,c_i),\qquad
  \lambda_i=\max(\mu_i,\nu_i)+b_i,
\]
and updates
\[
  c_{i+1}=c_i-b_i+d_i .
\]
For all sufficiently large $i$, the partition parts are zero and the carry has
fallen to zero, so the following finite sum is legitimate.  Since
\[
  \max(\mu_i,\nu_i)=\mu_i+\nu_i-\min(\mu_i,\nu_i)
  =\mu_i+\nu_i-\rho_i-d_i,
\]
we have
\[
  \lambda_i-\mu_i-\nu_i+\rho_i=b_i-d_i=c_i-c_{i+1}.
\]
Summing over $i$ telescopes to
\[
  |\lambda|-|\mu|-|\nu|+|\rho|
  =
  c_1-\lim_i c_i
  =
  m .
\]
\end{proof}

\begin{corollary}[Diagonal-two cells have padded carrier valleys]\label{cor:bc-diagonal-two-cell-has-padded-valley}
In the setting of
\cref{lem:bc-dual-rsk-prime-cell-size-conservation}, suppose that
\[
  |\lambda|-|\rho|=2 .
\]
Then the left/bottom valley satisfies
\[
  |\nu|-2|\rho|+|\mu|=2-m .
\]
\end{corollary}

\begin{proof}
By \cref{lem:bc-dual-rsk-prime-cell-size-conservation},
\[
  m=|\lambda|-|\mu|-|\nu|+|\rho|.
\]
Using $|\lambda|-|\rho|=2$ gives
\[
  |\nu|-2|\rho|+|\mu|
  =
  |\mu|+|\nu|-|\rho|-|\lambda|+2
  =
  2-m .
\]
\end{proof}

\begin{corollary}[Parsed dual-RSK' boundaries recover subfillings]\label{cor:bc-parsed-dual-rsk-prime-boundary-recovers-subfilling}
Let $F$ be a finite Ferrers cell arrangement equipped with the fourth
arbitrary-filling growth diagram of
\cref{lem:bc-dual-rsk-prime-growth-local-reversibility}, and let
$E\subseteq F$ be a Ferrers subarrangement.  Suppose that a parser has
recovered the partition labels on the top/right boundary of $E$.  Then the
shape of $E$ and those boundary labels determine, by a deterministic backward
sweep, every cell entry of $E$ and every corner label on or inside $E$.
\end{corollary}

\begin{proof}
The proof is the vertical-strip analogue of
\cref{cor:bc-parsed-boundary-recovers-knuth-subfilling}.  By
\cref{lem:bc-dual-rsk-prime-growth-local-reversibility}, the backward cell
algorithm of the fourth variant determines the southwest corner and the cell
entry from the northeast, southeast, and northwest corner labels.  Sweeping
the cells of $E$ leftward and downward from the recovered top/right boundary
therefore determines every entry and every remaining corner label.
\end{proof}

\begin{lemma}[Krattenthaler's semistandard path is a dual-RSK' boundary]\label{lem:bc-krattenthaler-semistandard-path-boundary}
Use the concrete growth-diagram construction in the proof sketch of
\cite[Theorem 4, Steps 2--3]{Krattenthaler2016}.  For a semistandard tableau
$T$ with content $(j_1,\ldots,j_n)$ and even-column endpoint parameter
$m=0$, the partition sequence counted in \cref{lem:bc-growth-path-model} is
read from the top/right boundary of the Step~3 half-square growth diagram.
That Step~3 diagram is the fourth arbitrary-filling, dual-RSK' variant of
\cref{lem:bc-dual-rsk-prime-growth-local-reversibility}.

Consequently, if a later decoder identifies a Ferrers subarrangement of this
Step~3 diagram whose top/right boundary lies in a recovered segment of
\[
  \rho_0,\rho_1,\ldots,\rho_{2j},
\]
then \cref{cor:bc-parsed-dual-rsk-prime-boundary-recovers-subfilling}
recovers the corresponding subfilling and all its internal corner labels
without additional branch choices.
\end{lemma}

\begin{proof}
In the proof sketch of \cite[Theorem 4]{Krattenthaler2016}, Step~2 applies a
Knuth-type inverse growth-diagram construction to obtain a symmetric matrix
with nonnegative integer entries.  Step~3 then applies the Knuth-type
extension cited there from Krattenthaler's fourth variant, and the output
sequence is obtained by reading the partitions along the top/right boundary
of the resulting half-square cell arrangement.  Krattenthaler's
\cite[\S4.4, Theorem 11]{Krattenthaler2006} is precisely that fourth
arbitrary-filling, dual-RSK' boundary bijection.  Applying this description
in the $m=0$ and $j_1=\cdots=j_j=2$ specialization used in
\cref{lem:bc-growth-path-model} gives the first assertion.  The final
assertion is then exactly
\cref{cor:bc-parsed-dual-rsk-prime-boundary-recovers-subfilling} applied to
the identified subarrangement.
\end{proof}

\begin{lemma}[Full Step-3 half-square entries determine the semistandard path]\label{lem:bc-step3-full-half-square-determines-rho}
Use the Step~3 dual-RSK' half-square in the $m=0$, content $(2^j)$
specialization of \cref{lem:bc-growth-path-model}.  Any datum which determines
every cell entry of this full Step~3 half-square determines the whole
semistandard boundary sequence
\[
  \rho_0,\rho_1,\ldots,\rho_{2j}.
\]
\end{lemma}

\begin{proof}
By \cref{lem:bc-krattenthaler-semistandard-path-boundary}, the sequence
$\rho_0,\rho_1,\ldots,\rho_{2j}$ is the top/right boundary of the Step~3
half-square growth diagram.  In Krattenthaler's Step~3 construction the
left and bottom boundary labels of the full half-square are the fixed empty
boundary labels of the growth diagram.  Hence the known cell entries together
with this fixed left/bottom boundary determine the top/right boundary by the
forward part of
\cref{lem:bc-dual-rsk-prime-growth-local-reversibility}.  This top/right
boundary is exactly the displayed $\rho$-sequence.
\end{proof}

\begin{lemma}[Step-3 cells are Step-2 half-matrix entries]\label{lem:bc-step3-cells-are-step2-half-matrix}
Use the construction in \cite[Theorem 4, Steps 2--3]{Krattenthaler2016} in the
$m=0$ specialization of \cref{lem:bc-growth-path-model}.  Let $M$ be the
symmetric nonnegative integer matrix produced by the Step~2 first-variant
inverse growth diagram.  The Step~3 dual-RSK' half-square is filled by one
off-diagonal half of $M$: each Step~3 cell entry is the corresponding entry of
$M$, and the transposed Step~2 entry is the same integer by symmetry.
Consequently, a recovered Step~3 subfilling determines the entries of every
off-diagonal Step~2 cell whose cell or transposed cell lies in that recovered
half-square; the Step~2 diagonal entries are known to be zero.
\end{lemma}

\begin{proof}
In \cite[Theorem 4]{Krattenthaler2016}, Step~2 is explicitly the Knuth-type
inverse growth-diagram construction of \cite[\S4.1]{Krattenthaler2006}; in the
semistandard extension it outputs a symmetric matrix with nonnegative integer
entries and zeros on the diagonal.  Step~3 then replaces the ordinary final
growth-diagram pass by the Knuth-type extension of
\cite[\S4.4]{Krattenthaler2006}: one half of the symmetric Step~2 matrix is
used as the filling of the half-square, and the partitions are read along the
top/right boundary after applying the dual-RSK' local rules.  In the $m=0$
case there is no extra-letter region to delete; the only discarded cells are
the diagonal and the opposite half of the symmetric matrix.  Thus each Step~3
cell entry is literally the corresponding off-diagonal entry of $M$, and the
opposite off-diagonal entry is equal to it because $M$ is symmetric.
\end{proof}

\begin{lemma}[Step-2 matrix degrees are the doubled content]\label{lem:bc-step2-matrix-doubled-content-degrees}
Use the construction in \cite[Theorem 4, Step 2]{Krattenthaler2016} and
\cite[\S4.1, Theorem 7]{Krattenthaler2006} for a semistandard tableau of
content $(2^j)$.  Let
\[
  M=(M_{ab})_{1\le a,b\le j}
\]
be the Step~2 first-variant Knuth matrix.  Then $M$ is symmetric, has zero
diagonal, and satisfies
\[
  \sum_{b=1}^j M_{ab}=2=\sum_{b=1}^j M_{ba}
  \qquad(1\le a\le j).
\]
\end{lemma}

\begin{proof}
The cited Step~2 construction applies the arbitrary-filling Knuth
correspondence of \cite[\S4.1, Theorem 7]{Krattenthaler2006}.  In the
rectangular form of that correspondence, an entry $M_{ab}$ contributes
$M_{ab}$ copies of the biletter with row label $a$ and column label $b$; the
row sums give the weight of one semistandard tableau and the column sums give
the weight of the other.  In Krattenthaler's Theorem~4 construction these two
tableaux are the equal insertion/recording pair obtained from the
jeu-de-taquin-preprocessed tableau, and the source records that the resulting
matrix is symmetric with zero diagonal.  Specializing to content $(2^j)$
means that each ordinary label $a$ occurs exactly twice.  Hence both the row
and column sums of the Step~2 matrix are equal to $2$ for every label.
\end{proof}

\begin{lemma}[HKKO even-bound up-down boundary events]\label{lem:bc-hkko-updown-boundary-events}
In the notation of Huh--Kim--Krattenthaler--Okada, let
$\mathrm{MUD}^*_n(w;\mu\to\nu)$ be the marked up-down class used for the
even-bound Littlewood identities.  Its marked set is a subset of
\[
  E_w(T)=
  \{j:\ell(T^{2j-2})<\ell(T^{2j-1})
        =\ell(T^{2j})=w\}.
 \]
After the coefficient extraction from Theorem 7.15 to Theorem 8.7 of
\cite{HKKO2025}, an index $j\in E_w(T)$ becomes an eligible boundary-up step
in the corresponding marked vacillating class
$\mathrm{MVT}^*_n(w;\mu\to\nu)$; that is, the corresponding marked
vacillating index is a $+\epsilon_w$ step whose starting point has
$w$-th coordinate $0$.
\end{lemma}

\begin{proof}
Definition 7.12 of \cite{HKKO2025} defines
$E_w(T)$ by the displayed formula and defines
$\mathrm{MUD}^*_n(w;\mu\to\nu)$ by requiring the marking set to be a subset of
$E_w(T)$.  Thus a marked up-down event first raises the length from below $w$
to $w$ at time $2j-1$ and keeps length $w$ at time $2j$.

Theorem 8.7 is obtained from the even-bound up-down identities of
Theorem 7.15 by coefficient extraction into marked vacillating tableaux.
Definition 8.6 then states the admissibility condition for the resulting
marks: a marked index in $\mathrm{MVT}^*$ must be a $+\epsilon_w$ step whose
starting point lies on $x_w=0$.  Hence the marked events coming from
$E_w(T)$ are exactly eligible boundary-up steps in the vacillating model, as
recorded in \cref{lem:bc-hkko-marked-boundary}.
\end{proof}

\begin{lemma}[HKKO boundary-event splice]\label{lem:bc-hkko-boundary-splice}
In the even-bound proof of \cite[Theorem 7.15]{HKKO2025}, let $Z_0$ be the
fixed-point class appearing in the proof of equation (7.32).  Thus an element
$(T,S)\in Z_0$ satisfies
\[
  S\ne\varnothing,\qquad
  j\in S\Rightarrow \ell(T^{2j-1})=w-1,\qquad
  j\notin S\Rightarrow 2j-1\text{ is not a length peak of }T .
\]
For every $j\in S$, the local operation $\psi$ of Figure 7 changes the
configuration at the $j$-th up-down pair, removes the mark, and produces an
unmarked up-down tableau $T'$.  Applied for all $j\in S$, this operation is a
bijection from $Z_0$ onto the class $Y$ in the proof of (7.32).  It preserves
the global endpoints $\mu,\nu$ and the length $2n$, and it is reversible from
$T'$ and the selected set $S$.
\end{lemma}

\begin{proof}
In the proof of \cite[Theorem 7.15]{HKKO2025}, after equation (7.36), the
authors define the fixed-point set $Z_0$ by exactly the three displayed
conditions.  They then define $\psi(T,S)=(T',\varnothing)$ by applying the
operation shown in Figure 7 for each $j\in S$.  The paragraph following the
definition states that $\psi$ is a bijection from $Z_0$ to the class $Y$.
Both $Z_0$ and $Y$ are classes of $w$-up-down tableaux from the same $\mu$ to
the same $\nu$ and of the same length $2n$, so the operation preserves these
global data.  Bijection gives the asserted reversibility from the output
together with the selected set.
\end{proof}

\begin{lemma}[Explicit $Z_0$ landing criterion]\label{lem:bc-z0-landing-criterion}
Let
\[
  U=(U^0,U^1,\ldots,U^{2N})
\]
be a $w$-up-down tableau from $\mu$ to $\nu$, and let $M\subseteq[N]$.
Suppose first that $M=\varnothing$.  Then the identity operation gives an
unmarked up-down tableau with the same endpoints and length.

Suppose instead that $M\ne\varnothing$ and the following three conditions
hold:
\[
  \ell(U^r)<w\qquad(0\le r\le2N),
\]
\[
  m\in M\quad\Longrightarrow\quad \ell(U^{2m-1})=w-1,
\]
and
\[
  m\notin M\quad\Longrightarrow\quad
  \neg\bigl(\ell(U^{2m-2})<\ell(U^{2m-1})>\ell(U^{2m})\bigr).
\]
Then $(U,M)$ lies in the $Z_0$ domain of
\cref{lem:bc-hkko-boundary-splice}.  Consequently the HKKO boundary-event
splice produces an unmarked $w$-up-down tableau with endpoints $\mu,\nu$ and
length $2N$, and the operation is reversible from the output and $M$.
\end{lemma}

\begin{proof}
If $M=\varnothing$, there are no selected boundary events to splice, and the
identity operation has the asserted properties.

Assume $M\ne\varnothing$.  The first displayed condition says precisely that
$(U,M)$ lies in the class $\mathrm{MUD}^{<}_N(w;\mu\to\nu)$ of Definition 7.12
of \cite{HKKO2025}.  The second and third displayed conditions are the two
fixed-point conditions defining $Z_0$ in the proof of
\cite[Theorem 7.15, equation (7.32)]{HKKO2025}: selected indices have midpoint
length $w-1$, and unselected odd midpoints are not length peaks.  Hence
$(U,M)\in Z_0$.  Applying \cref{lem:bc-hkko-boundary-splice} gives the
unmarked output with the same endpoints and length and gives reversibility
from that output together with $M$.
\end{proof}

\begin{lemma}[Penultimate-row creations give the selected $Z_0$ midpoint]\label{lem:bc-z0-selected-midpoint-creation}
Let
\[
  U=(U^0,\ldots,U^{2N})
\]
be a $w$-up-down tableau, and let $M\subseteq[N]$.  Suppose that, for every
$m\in M$, the partition $U^{2m-2}$ has at most $w-2$ nonzero parts and
$U^{2m-1}$ is obtained from $U^{2m-2}$ by adding one box in row $w-1$.  Then
\[
  m\in M\quad\Longrightarrow\quad \ell(U^{2m-1})=w-1.
\]
Consequently, if the same $U$ also satisfies the below-wall and unselected
non-peak conditions of \cref{lem:bc-z0-landing-criterion}, then the selected
midpoint condition in that criterion is automatic.
\end{lemma}

\begin{proof}
Fix $m\in M$.  Since $U^{2m-2}$ has at most $w-2$ nonzero parts, its
$(w-1)$-st part is zero.  Adding one box in row $w-1$ makes the
$(w-1)$-st part of $U^{2m-1}$ positive.  As $U^{2m-1}$ is a partition, all
parts below row $w-1$ are zero.  Hence its number of nonzero parts is exactly
$w-1$.  The final sentence is immediate after substituting this implication
into the hypotheses of \cref{lem:bc-z0-landing-criterion}.
\end{proof}

\begin{lemma}[$Z_0$ landing is stable under even-time concatenation]\label{lem:bc-z0-concatenation}
Let
\[
  U=(U^0,\ldots,U^{2N})
  \quad\text{and}\quad
  V=(V^0,\ldots,V^{2M})
\]
be $w$-up-down tableaux with $U^{2N}=V^0$.  Let
$A\subseteq[N]$ and $B\subseteq[M]$.  Suppose that $U$ satisfies
\[
  \ell(U^r)<w\quad(0\le r\le2N),
  \qquad
  a\in A\Rightarrow \ell(U^{2a-1})=w-1,
\]
and
\[
  a\notin A\Rightarrow
  \neg\bigl(\ell(U^{2a-2})<\ell(U^{2a-1})>\ell(U^{2a})\bigr),
\]
and that $V$ satisfies the analogous three conditions with $B$ in place of
$A$.  Let $W$ be the even-time concatenation
\[
  W=(U^0,\ldots,U^{2N},V^1,\ldots,V^{2M})
\]
and put
\[
  C=A\cup\{N+b:b\in B\}\subseteq[N+M].
\]
Then $W$ satisfies the same three $Z_0$ landing conditions with selected set
$C$.
\end{lemma}

\begin{proof}
Every state of $W$ is either a state of $U$ or a state of $V$, with the common
endpoint counted once.  Hence $\ell(W^r)<w$ for all $0\le r\le2(N+M)$.

If $c\in C$, then either $c=a\in A$, in which case
$W^{2c-1}=U^{2a-1}$, or $c=N+b$ with $b\in B$, in which case
$W^{2c-1}=V^{2b-1}$.  The selected midpoint condition follows from the
corresponding condition for $U$ or $V$.

Finally let $c\notin C$.  If $c\le N$, then $c\notin A$ and the three states
$W^{2c-2},W^{2c-1},W^{2c}$ are the corresponding three states of $U$.  If
$c>N$, write $c=N+b$ with $1\le b\le M$; then $b\notin B$ and the same three
states are the shifted corresponding states of $V$.  In both cases the
non-peak condition follows from the appropriate segment.
\end{proof}

\begin{lemma}[HKKO even-bound marked boundary model]\label{lem:bc-hkko-marked-boundary}
Let $\mathrm{SYT}_{n,2w}[r]$ be the set of standard Young tableaux of size
$n$, width at most $2w$, and exactly $r$ odd columns.  In the notation of
Huh--Kim--Krattenthaler--Okada, for $0\le k\le w-1$,
\[
 |\mathrm{SYT}_{n,2w}[k]|
 =
 |\mathrm{MVT}^0_n(w;(1^k,0^{w-k})\to(0^w))|
\]
and
\[
 |\mathrm{SYT}_{n,2w}[2w-k]|
 =
 |\mathrm{MVT}^1_n(w;(1^k,0^{w-k})\to(0^w))|.
\]
Moreover
\[
 |\mathrm{SYT}_{n,2w}[k]|+|\mathrm{SYT}_{n,2w}[2w-k]|
 =
 |\mathrm{MVT}^*_n(w;(1^k,0^{w-k})\to(0^w))|,
\]
and for $k=w$,
\[
 |\mathrm{SYT}_{n,2w}[w]|
 =
 |\mathrm{MVT}^*_n(w;(1^w)\to(0^w))|.
\]
In $\mathrm{MVT}^*_n$, every marked index is a nonzero vacillating step
$+\epsilon_w$ whose starting point lies on the boundary hyperplane $x_w=0$,
and the underlying vacillating path has no zero steps.  For
$0\le k\le w-1$, the two refinements $\mathrm{MVT}^0_n$ and
$\mathrm{MVT}^1_n$ are separated by whether the first $+\epsilon_w$ step is
unmarked or marked.
\end{lemma}

\begin{proof}
Definitions 8.1, 8.5, and 8.6 of
\cite{HKKO2025} define $\mathrm{SYT}_{n,w}[r]$, marked vacillating tableaux,
the subset $\mathrm{MVT}^*_n$, and the refinements
$\mathrm{MVT}^0_n,\mathrm{MVT}^1_n$.  In Definition 8.6, a mark is allowed
only on a step $\epsilon_w$ whose starting point is on $x_w=0$, and the same
definition requires the underlying path to have no zero step.  The $0/1$
split is exactly the unmarked/marked state of the first such step.
The four displayed cardinality identities are precisely Theorem 8.7 of
\cite{HKKO2025}.
\end{proof}

\begin{lemma}[Marked boundary branch count]\label{lem:bc-marked-boundary-branch-count}
Let $T=(T^0,\ldots,T^n)$ be a $w$-vacillating tableau with no zero steps, and
write it as a walk in
\[
  \{x_1\ge\cdots\ge x_w\ge0\}.
\]
Let
\[
  B(T)=\{r: T^r-T^{r-1}=+\epsilon_w
       \text{ and }T^{r-1}_w=0\}
\]
be the set of eligible boundary-up steps.  For this fixed underlying
vacillating path $T$, the elements of
$\mathrm{MVT}^*_n(w;\mu\to\nu)$ with underlying path $T$ are exactly the
choices of a subset of $B(T)$.  In particular, if a repair construction
designates a set $B_0\subseteq B(T)$ of size at most $2h$ and leaves all
eligible steps outside $B_0$ unmarked, then the number of possible marking
choices is at most $2^{2h}$.
\end{lemma}

\begin{proof}
By Definition 8.5 of \cite{HKKO2025}, a marked vacillating tableau is a pair
$(T,S)$ with $S\subseteq[n]$.  Definition 8.6 says that membership in
$\mathrm{MVT}^*_n$ requires the underlying path to have no zero step and
requires every marked index to be a $+\epsilon_w$ step whose starting point
lies on $x_w=0$.  For fixed $T$, this condition is exactly $S\subseteq B(T)$.
Thus the possible markings are precisely the subsets of $B(T)$.  The final
claim is the same statement restricted to the designated set $B_0$, whose
number of subsets is at most $2^{2h}$.
\end{proof}

\begin{lemma}[Eligible boundary-up steps are last-row creations]\label{lem:bc-hkko-boundary-up-characterization}
Let
\[
  V=(V^0,\ldots,V^N)
\]
be a $w$-vacillating tableau, viewed as a walk in the chamber
\[
  x_1\ge x_2\ge\cdots\ge x_w\ge0 .
\]
For a time $t$, the step $t$ belongs to the eligible boundary-up set
\[
  B(V)=\{t:V^t-V^{t-1}=+\epsilon_w\text{ and }V^{t-1}_w=0\}
\]
if and only if $V^{t-1}$ has at most $w-1$ nonzero parts and $V^t$ is obtained
from $V^{t-1}$ by adding one box in row $w$.
\end{lemma}

\begin{proof}
In the chamber notation of Definition 8.2 of \cite{HKKO2025}, recalled in
\cref{lem:bc-hkko-marked-boundary}, the vector $\epsilon_w$ adds one to the
last coordinate and leaves all other coordinates unchanged.  The condition
$V^{t-1}_w=0$ is exactly that the starting partition has no nonzero $w$-th
part.  Thus the step is $+\epsilon_w$ from the boundary hyperplane precisely
when it creates a single box in row $w$.
\end{proof}

\begin{lemma}[Odd-column removal for standard tableaux]\label{lem:bc-odd-column-removal}
Let $U$ be a standard Young tableau of size $n$ with $m$ odd-length columns
and with all column lengths at most $2k+1$.  Krattenthaler's reversible
odd-column removal in \cite[\S4, item (2)]{Krattenthaler2016} gives a bijection
between such tableaux $U$ and pairs $(U^{\mathrm{ev}},A)$, where
$U^{\mathrm{ev}}$ is a standard Young tableau of size $n-m$ with only even
columns and all column lengths at most $2k$, and
$A\subset\{1,\ldots,n\}$ has cardinality $m$.
\end{lemma}

\begin{proof}
Krattenthaler \cite[\S4, item (2)]{Krattenthaler2016} treats standard Young
tableaux whose columns have length at most $2k+1$.  Starting with the last
entry in the rightmost odd column and proceeding right-to-left through the
currently rightmost odd column, he applies the inverse Robinson--Schensted
insertion until all odd columns have disappeared.  The cited paragraph states
that the output is a standard Young tableau of size $n-m$ with only even
columns, all of length at most $2k$, together with a subset of
$\{1,\ldots,n\}$ of cardinality $m$, and that all steps are reversible.  This
is exactly the displayed bijection.
\end{proof}

\begin{lemma}[Odd-column removal is not the final fixed-witness map]\label{lem:bc-odd-column-removal-size-guard}
Let $U$ be a standard Young tableau of size $2(j-q)$ with $m$ odd-length
columns.  If $m>0$, the output tableau $U^{\mathrm{ev}}$ of
\cref{lem:bc-odd-column-removal} has size strictly smaller than $2(j-q)$.
Consequently, applying odd-column removal as the final operation after
witness-label deletion cannot by itself produce an element of
$\mathcal P_{j-q}^{\mathrm{even}}$, unless the retained shape is already
column-even.
\end{lemma}

\begin{proof}
\Cref{lem:bc-odd-column-removal} sends a standard tableau of size $n$ with
$m$ odd columns to a column-even standard tableau of size $n-m$.  Substituting
$n=2(j-q)$ gives output size $2(j-q)-m$, which is strictly smaller than
$2(j-q)$ when $m>0$.  But the paths in
$\mathcal P_{j-q}^{\mathrm{even}}$ correspond to semistandard tableaux of
content $(2^{j-q})$ and therefore have total size $2(j-q)$.  Hence the
odd-column-removal operation can only serve as an internal reversible splice
or boundary-mark mechanism; it is already target-size preserving only in the
case $m=0$.
\end{proof}

\begin{lemma}[Standardization of doubled-content tableaux]\label{lem:bc-doubled-standardization}
Let $T$ be a semistandard tableau of shape $\lambda$ and content $(2^j)$.
For each label $a\in\{1,\ldots,j\}$, replace the leftmost box carrying $a$
by $2a-1$ and the rightmost box carrying $a$ by $2a$.  The resulting filling
$\operatorname{std}(T)$ is a standard Young tableau of shape $\lambda$.
The map $T\mapsto\operatorname{std}(T)$ is injective.  It preserves the shape,
and hence preserves the column-even condition, the width and height
thresholds, and the canonical witness boxes of
\cref{lem:bad-shape-witnesses}.  Deleting the two boxes carrying a
semistandard label $a$ in $T$ is the same box deletion as deleting the
standard adjacent pair $\{2a-1,2a\}$ in $\operatorname{std}(T)$.
\end{lemma}

\begin{proof}
The two boxes carrying a fixed label $a$ cannot lie in the same column, because
columns in $T$ are strictly increasing.  Thus the leftmost and rightmost boxes
are well-defined.

Consider two boxes in the same row, with the left box carrying label $a$ and
the right box carrying label $b$.  Since rows of $T$ are weakly increasing,
$a\le b$.  If $a<b$, then every standardized value assigned to the left box is
at most $2a<2b-1$, which is at most the standardized value assigned to the
right box.  If $a=b$, then the two boxes are the leftmost and rightmost
occurrences of $a$ in that row and receive $2a-1$ and $2a$.  Hence rows of
$\operatorname{std}(T)$ are strictly increasing.

Now consider two boxes in the same column, with the upper box carrying $a$ and
the lower box carrying $b$.  Since columns of $T$ are strictly increasing,
$a<b$.  The standardized value in the upper box is at most $2a$, while the
standardized value in the lower box is at least $2b-1$, and
$2a<2b-1$.  Hence columns of $\operatorname{std}(T)$ are strictly increasing.
The entries of $\operatorname{std}(T)$ are exactly
$1,2,\ldots,2j$, so it is standard.

Injectivity follows by replacing each pair $\{2a-1,2a\}$ in
$\operatorname{std}(T)$ by $a$.  The remaining assertions are immediate
because the operation changes only labels, not boxes.
\end{proof}

\begin{lemma}[Standardized witness-pair local states]\label{lem:bc-standardized-witness-pairs}
Let $T$ be a semistandard tableau of content $(2^j)$, and let
$\operatorname{std}(T)$ be the standard tableau of
\cref{lem:bc-doubled-standardization}.

If $\lambda_1(T)\ge2q$ and $s_r=T(1,2r-1)$ for $1\le r\le q$, then the two
boxes deleted by the width witness label $s_r$ are exactly the boxes carrying
the adjacent standard pair $\{2s_r-1,2s_r\}$ in $\operatorname{std}(T)$.
The frontier box $(1,2r-1)$ carries one of these two entries.  Thus the bit
\[
  \epsilon_r^w(T)=
  \begin{cases}
  0,&\operatorname{std}(T)(1,2r-1)=2s_r-1,\\
  1,&\operatorname{std}(T)(1,2r-1)=2s_r
  \end{cases}
\]
records which member of the adjacent pair lies at the width frontier.  If the
other copy of $s_r$ lies in the first row, then it is in column $2r$ when
$\epsilon_r^w(T)=0$, and in column $2r-2$ when $\epsilon_r^w(T)=1$.

If $\ell(\lambda(T))\ge2q$ and $t_r=T(2r-1,1)$ for $1\le r\le q$, then the two
boxes deleted by the height witness label $t_r$ are exactly the boxes carrying
the adjacent standard pair $\{2t_r-1,2t_r\}$ in $\operatorname{std}(T)$, and
the frontier box $(2r-1,1)$ carries $2t_r-1$.
\end{lemma}

\begin{proof}
For every label $a$, \cref{lem:bc-doubled-standardization} replaces the two
boxes carrying $a$ by the adjacent standard entries $2a-1$ and $2a$.  This
proves the adjacent-pair statement for both width and height witnesses.

For the width statement, the frontier box $(1,2r-1)$ is one of the two boxes
carrying $s_r$, hence in the standardized tableau it carries either
$2s_r-1$ or $2s_r$, giving the displayed bit.  If the other copy lies in the
first row, \cref{lem:bc-width-witness-second-copy} places it in column
$2r-2$ or $2r$.  The standardization assigns $2s_r-1$ to the leftmost copy and
$2s_r$ to the rightmost copy, so the other first-row copy is in column $2r$
in the first case and in column $2r-2$ in the second.

For the height statement, \cref{lem:bc-height-witness-second-copy} says that
the second copy of $t_r$ is not in the first column.  Since the frontier copy
is in column $1$, it is the leftmost copy and therefore receives the
standardized label $2t_r-1$.
\end{proof}

\begin{lemma}[Frontier crossing time inside a deleted standard pair]\label{lem:bc-standard-frontier-time}
Let $T$ be a semistandard tableau of content $(2^j)$, let
$U=\operatorname{std}(T)$, and let
\[
  \varnothing=\sigma^{(0)}\subset\sigma^{(1)}\subset\cdots\subset
  \sigma^{(2j)}=\lambda(T)
\]
be the standard shape chain of $U$.

If $\lambda_1(T)\ge2q$ and $s_r=T(1,2r-1)$ for $1\le r\le q$, then the
standard time
\[
  \kappa^w_r=2s_r-1+\epsilon_r^w(T)
\]
is the unique time in the adjacent pair $\{2s_r-1,2s_r\}$ at which the box
$(1,2r-1)$ is added.  Equivalently,
\[
  \sigma^{(\kappa^w_r-1)}_1<2r-1\le\sigma^{(\kappa^w_r)}_1 .
\]

If $\ell(\lambda(T))\ge2q$ and $t_r=T(2r-1,1)$ for $1\le r\le q$, then the
standard time
\[
  \kappa^h_r=2t_r-1
\]
is the unique time in the adjacent pair $\{2t_r-1,2t_r\}$ at which the box
$(2r-1,1)$ is added.  Equivalently,
\[
  \ell(\sigma^{(\kappa^h_r-1)})<2r-1\le\ell(\sigma^{(\kappa^h_r)}) .
\]
\end{lemma}

\begin{proof}
In a standard tableau, the box carrying entry $m$ is added exactly when the
standard shape chain passes from $\sigma^{(m-1)}$ to $\sigma^{(m)}$.

For the width statement, \cref{lem:bc-standardized-witness-pairs} says that
the frontier box $(1,2r-1)$ carries $2s_r-1$ when
$\epsilon_r^w(T)=0$ and carries $2s_r$ when $\epsilon_r^w(T)=1$.  Thus it is
added at $\kappa^w_r$.  Before that time the first row does not contain column
$2r-1$; after that time it does.  This is exactly the displayed inequality.
The time is unique because one standard entry labels one box.

For the height statement, \cref{lem:bc-standardized-witness-pairs} says that
the frontier box $(2r-1,1)$ carries $2t_r-1$.  Therefore it is added at
$\kappa^h_r=2t_r-1$.  Before that time the first column has length below
$2r-1$; after that time it has length at least $2r-1$.  Again uniqueness is
the uniqueness of the standard entry in that box.
\end{proof}

\begin{lemma}[Frontier crossing times are ordered deleted times]\label{lem:bc-standard-frontier-order}
Let $T$ be a semistandard tableau of content $(2^j)$, let
$U=\operatorname{std}(T)$, and let
\[
  \varnothing=\sigma^{(0)}\subset\sigma^{(1)}\subset\cdots\subset
  \sigma^{(2j)}=\lambda(T)
\]
be the standard shape chain of $U$.

Assume $\lambda_1(T)\ge2q$.  Put $s_r=T(1,2r-1)$,
$S=W_q(T)$, and
\[
  D(S)=\bigcup_{s\in S}\{2s-1,2s\},\qquad
  \kappa^w_r=2s_r-1+\epsilon_r^w(T)
  \quad(1\le r\le q).
\]
Then
\[
  2s_r-1\le\kappa^w_r\le2s_r,\qquad
  \kappa^w_r\in D(S),\qquad
  \kappa^w_1<\kappa^w_2<\cdots<\kappa^w_q .
\]
Moreover,
\[
 \{m:\min(q,\lceil\sigma^{(m)}_1/2\rceil)
       -\min(q,\lceil\sigma^{(m-1)}_1/2\rceil)=1\}
 =
 \{\kappa^w_r:1\le r\le q\}.
\]

Assume instead that $\ell(\lambda(T))\ge2q$.  Put
$t_r=T(2r-1,1)$, $S=H_q(T)$, and
\[
  D(S)=\bigcup_{t\in S}\{2t-1,2t\},\qquad
  \kappa^h_r=2t_r-1
  \quad(1\le r\le q).
\]
Then
\[
  2t_r-1\le\kappa^h_r\le2t_r,\qquad
  \kappa^h_r\in D(S),\qquad
  \kappa^h_1<\kappa^h_2<\cdots<\kappa^h_q .
\]
Moreover,
\[
 \{m:\min(q,\lceil\ell(\sigma^{(m)})/2\rceil)
       -\min(q,\lceil\ell(\sigma^{(m-1)})/2\rceil)=1\}
 =
 \{\kappa^h_r:1\le r\le q\}.
\]
\end{lemma}

\begin{proof}
For the width case, $\epsilon_r^w(T)\in\{0,1\}$ by definition, so
$2s_r-1\le\kappa^w_r\le2s_r$ and $\kappa^w_r\in D(S)$.  The entries
$s_1,\ldots,s_q$ occur in weakly increasing order along the first row, and
they are distinct by \cref{lem:bad-shape-witnesses}; hence
$s_1<\cdots<s_q$.  Therefore
\[
  \kappa^w_r\le2s_r<2s_{r+1}-1\le\kappa^w_{r+1}
\]
for $1\le r<q$.

By \cref{lem:bc-standard-frontier-time}, $\kappa^w_r$ is exactly the standard
time at which the box $(1,2r-1)$ is added.  At that time the first row length
passes from $2r-2$ to $2r-1$, so the displayed minimum increases by one.
Conversely, if the displayed minimum increases at time $m$, then the single
box added from $\sigma^{(m-1)}$ to $\sigma^{(m)}$ must be the first-row box
$(1,2r-1)$ for a unique $1\le r\le q$.  The same cited lemma identifies its
standard time with $\kappa^w_r$.  This proves the width equality of crossing
sets.

For the height case, the inequalities
$2t_r-1\le\kappa^h_r\le2t_r$ and $\kappa^h_r\in D(S)$ are immediate from
$\kappa^h_r=2t_r-1$.  The first-column entries are strictly increasing, so
$t_1<\cdots<t_q$, and hence
$\kappa^h_1<\cdots<\kappa^h_q$.

By \cref{lem:bc-standard-frontier-time}, $\kappa^h_r$ is exactly the standard
time at which the box $(2r-1,1)$ is added.  At that time the first-column
length passes from $2r-2$ to $2r-1$, so the displayed height minimum increases
by one.  Conversely, if that minimum increases at time $m$, then the added
box is the first-column box $(2r-1,1)$ for a unique $1\le r\le q$, and
\cref{lem:bc-standard-frontier-time} again gives $m=\kappa^h_r$.
\end{proof}

\begin{lemma}[Fixed-witness frontier assignments are strictly ordered]\label{lem:bc-fixed-witness-frontier-assignment-ordered}
In the width setting of \cref{lem:bc-standard-frontier-order}, the assignment
\[
  e^w:S=W_q(T)\longrightarrow D(S),\qquad e^w(s_r)=\kappa^w_r
\]
satisfies
\[
  s<s'\quad\Longleftrightarrow\quad e^w(s)<e^w(s')
  \qquad(s,s'\in S).
\]
In the height setting, the assignment
\[
  e^h:S=H_q(T)\longrightarrow D(S),\qquad e^h(t_r)=\kappa^h_r
\]
satisfies the same equivalence.
\end{lemma}

\begin{proof}
We prove the width case; the height case is identical.  By
\cref{lem:bad-shape-witnesses}, the labels
\[
  s_r=T(1,2r-1)\qquad(1\le r\le q)
\]
are distinct, and they occur weakly increasingly along the first row.  Hence
\[
  s_1<s_2<\cdots<s_q .
\]
\Cref{lem:bc-standard-frontier-order} gives
\[
  \kappa^w_1<\kappa^w_2<\cdots<\kappa^w_q .
\]
Thus, for $s=s_r$ and $s'=s_{r'}$, one has $s<s'$ exactly when
$r<r'$, and this is exactly when
$e^w(s)=\kappa^w_r<\kappa^w_{r'}=e^w(s')$.
\end{proof}

\begin{corollary}[Actual frontier-consecutive edges give single-deficit recovery]\label{cor:bc-actual-frontier-consecutive-edge-recovery}
Work in the setting of
\cref{prop:bc-single-deficit-stable-retained-repair}, and let
$A\subset S$ be the support of the anchored component in
\cref{lem:bc-single-deficit-support-path-code}.

In the width setting of \cref{lem:bc-standard-frontier-order}, suppose
$S=W_q(T)$, write $S=\{s_1<\cdots<s_q\}$ as there, and assign
$s_r\mapsto\kappa^w_r$.  If every internal edge of the anchored path on $A$
has endpoints whose assigned times are consecutive in the full ordered list
\[
  \kappa^w_1<\kappa^w_2<\cdots<\kappa^w_q,
\]
then $\Gamma$ is recovered from
\[
  S,\quad v,\quad \Gamma^\ast,\quad
  \Gamma_{S\setminus A},\quad \mathbf 1_A\in\{0,1\}^{S},
\]
where $\Gamma_{S\setminus A}$ is the graph induced by $\Gamma$ on
$S\setminus A$.

In the height setting, the same assertion holds with
$S=H_q(T)$, $S=\{t_1<\cdots<t_q\}$, and the full ordered list
\[
  \kappa^h_1<\kappa^h_2<\cdots<\kappa^h_q.
\]
Thus, in the actual fixed-witness fibers, the single-deficit crossing leaf is
the identification of $v$ together with this full-frontier consecutiveness of
the anchored internal edges, plus a separate supplier for the deleted-only
graph unless $A=S$.
\end{corollary}

\begin{proof}
In the width case, \cref{lem:bc-fixed-witness-frontier-assignment-ordered}
shows that the assignment $s_r\mapsto\kappa^w_r$ satisfies the order
equivalence required in \cref{prop:bc-full-frontier-adjacent-edge-recovery}.
The displayed consecutiveness hypothesis is exactly consecutiveness among the
full event set assigned to $S$.  That proposition therefore gives the stated
recovery.  The height case is identical, using the assignment
$t_r\mapsto\kappa^h_r$.  The final sentence just records the two cases in the
language of the fixed-witness fibers.
\end{proof}

\begin{lemma}[Residual windows have a $q$-bit retained-vertex code]\label{lem:bc-residual-qbit-retained-vertex-code}
Let $G=B_b$ with $14\le b\le217$ or $G=C_b$ with $20\le b\le217$, and put
$q=b+1$.  Let $j$ be in the residual window of \cref{rem:former-post29-tail}, and let
$S\subset\{1,\ldots,j\}$ have cardinality $q$.  Then the ordered retained set
\[
  V_0=\{1,\ldots,j\}\setminus S
\]
has cardinality at most $2^q$.  Consequently any retained vertex $v\in V_0$
is encoded by a word in $\{0,1\}^q$, using its rank in the increasing order of
$V_0$.
\end{lemma}

\begin{proof}
For all residual rows, $q\ge15$.  By
\cref{cor:post29-bc-direct-onset-bound}, and since $j$ lies in the residual
window of \cref{rem:former-post29-tail},
\[
  j\le N_{\mathrm{dir}}^{BC}(G)+2\le30899.
\]
Therefore
\[
  |V_0|=j-q\le30899-15=30884<32768=2^{15}\le2^q .
\]
Listing $V_0$ increasingly identifies its elements with
$1,\ldots,|V_0|$, and the displayed inequality embeds this rank set into the
$2^q$ binary words of length $q$.
\end{proof}

\begin{proposition}[$q$ spare bits remove $v$-identification from the single-deficit leaf]\label{prop:bc-qbit-v-code-single-deficit-recovery}
Work in one residual row and moment of \cref{rem:former-post29-tail}.  Put
$q=b+1$, let $S\subset\{1,\ldots,j\}$ have cardinality $q$, and put
$V_0=\{1,\ldots,j\}\setminus S$.  In the setting of
\cref{prop:bc-single-deficit-stable-retained-repair}, let
$A\subset S$ be the support of the anchored component.

Assume that the actual width fixed-witness fiber has
$S=W_q(T)$ and every internal edge of the anchored path on $A$ has endpoints
whose assigned times are consecutive in the full ordered list
$\kappa^w_1<\cdots<\kappa^w_q$.  Then $\Gamma$ is recovered from
\[
  S,\quad \Gamma^\ast,\quad
  \Gamma_{S\setminus A},\quad
  \mathbf 1_A\in\{0,1\}^{S},\quad
  \rho(v)\in\{0,1\}^q,
\]
where $\rho(v)$ is the $q$-bit rank code of
\cref{lem:bc-residual-qbit-retained-vertex-code}.  The same assertion holds in
the height fixed-witness fiber with $S=H_q(T)$ and
$\kappa^h_1<\cdots<\kappa^h_q$.

Consequently, on the single-deficit leaf, the remaining fixed-witness crossing
theorem no longer has to identify $v$: it only has to prove the displayed
full-$\kappa$ consecutiveness of anchored internal edges.  The support word and
the retained-vertex code together use at most $2q$ branch bits, while the
deleted-only graph $\Gamma_{S\setminus A}$ remains an explicit supplier
obligation unless it is absorbed into the target or absent.
\end{proposition}

\begin{proof}
\Cref{lem:bc-residual-qbit-retained-vertex-code} recovers $v$ from
$S$ and the word $\rho(v)$.  In the width case, after $v$ is recovered,
\cref{cor:bc-actual-frontier-consecutive-edge-recovery} applies to the
displayed full-$\kappa^w$ consecutiveness hypothesis and recovers
$\Gamma$ from $S$, $v$, $\Gamma^\ast$, $\Gamma_{S\setminus A}$, and
$\mathbf 1_A$.  The height case is identical with $\kappa^h$ in place of
$\kappa^w$.  The last sentence counts the $q$ support bits and the $q$
retained-vertex bits, and records the separate deleted-only input.
\end{proof}

\begin{lemma}[Consecutive frontier times are adjacent witness ranks]\label{lem:bc-consecutive-kappa-adjacent-ranks}
Let
\[
  \kappa_1<\kappa_2<\cdots<\kappa_q
\]
be a strictly ordered list.  If $1\le r,r'\le q$ and no $\kappa_i$ lies
strictly between $\kappa_r$ and $\kappa_{r'}$, then $|r-r'|=1$ unless
$r=r'$.  Conversely, if $|r-r'|=1$, then no $\kappa_i$ lies strictly between
$\kappa_r$ and $\kappa_{r'}$.
\end{lemma}

\begin{proof}
Assume $r<r'$.  If $r'-r\ge2$, then $\kappa_{r+1}$ lies strictly between
$\kappa_r$ and $\kappa_{r'}$, contradiction.  Hence $r'=r+1$.  The case
$r'>r$ after interchanging $r,r'$ is the same.  Conversely, if $r'=r+1$, the
strict ordering leaves no index strictly between $r$ and $r'$, and hence no
listed value lies strictly between $\kappa_r$ and $\kappa_{r'}$.  The reversed
adjacent case is identical.
\end{proof}

\begin{proposition}[Witness-rank adjacent edges give residual single-deficit recovery]\label{prop:bc-rank-adjacent-single-deficit-recovery}
Work in one residual row and moment of \cref{rem:former-post29-tail}.  Put
$q=b+1$, let $S\subset\{1,\ldots,j\}$ have cardinality $q$, and use the
notation of \cref{prop:bc-qbit-v-code-single-deficit-recovery}.

In the width fixed-witness fiber, write
\[
  S=W_q(T)=\{s_1<\cdots<s_q\}.
\]
If every internal edge of the anchored path on $A$ joins two labels
$s_r,s_{r'}$ with $|r-r'|=1$, then $\Gamma$ is recovered from
\[
  S,\quad \Gamma^\ast,\quad
  \Gamma_{S\setminus A},\quad
  \mathbf 1_A\in\{0,1\}^{S},\quad
  \rho(v)\in\{0,1\}^q .
\]
The same assertion holds in the height fixed-witness fiber with
$S=H_q(T)=\{t_1<\cdots<t_q\}$ and the condition that every internal anchored
edge joins $t_r,t_{r'}$ with $|r-r'|=1$.

Consequently the remaining single-deficit crossing theorem may be stated
without frontier-time notation: after the repaired retained graph is fixed,
the geometric condition for the anchored component is that anchored internal
edges join adjacent witness ranks; the deleted-only graph remains separate
unless it is otherwise supplied.
\end{proposition}

\begin{proof}
In the width case, \cref{lem:bc-standard-frontier-order} gives the strictly
ordered frontier list
\[
  \kappa^w_1<\kappa^w_2<\cdots<\kappa^w_q .
\]
If an internal anchored edge joins $s_r$ to $s_{r'}$ with $|r-r'|=1$, then
\cref{lem:bc-consecutive-kappa-adjacent-ranks} shows that the corresponding
frontier times $\kappa^w_r$ and $\kappa^w_{r'}$ are consecutive in the full
ordered list.  Thus the hypotheses of
\cref{prop:bc-qbit-v-code-single-deficit-recovery} hold, and that proposition
gives the displayed recovery.  The height case is identical, using
\[
  \kappa^h_1<\kappa^h_2<\cdots<\kappa^h_q .
\]
\end{proof}

\begin{lemma}[Full-rank adjacent anchored paths have interval support]\label{lem:bc-rank-adjacent-path-interval-support}
Let
\[
  S=\{s_1<\cdots<s_q\}
\]
be a linearly ordered witness set, let $T\subset S$, and let $P$ be a simple
path through all labels of $T$.  Suppose every edge of $P$ joins two labels
$s_r,s_{r'}$ with $|r-r'|=1$ in the full ordered set $S$.  Then $T$ is a
consecutive interval of $S$: either $T=\varnothing$, or
\[
  T=\{s_a,s_{a+1},\ldots,s_b\}
\]
for uniquely determined $1\le a\le b\le q$.  If $|T|\ge2$, then $P$ is the
monotone path through this interval, uniquely up to reversal.
\end{lemma}

\begin{proof}
If $T=\varnothing$ there is nothing to prove.  Put
\[
  I=\{r:s_r\in T\}\subset\{1,\ldots,q\}.
\]
The path $P$ gives a connected graph on the vertex set $I$.  Every edge joins
indices differing by one.  If $I$ had a gap, say $r<r'$ in $I$ with no
element of $I$ strictly between them, then no path using only unit-difference
edges could connect the part of $I$ at or below $r$ to the part at or above
$r'$.  This contradicts connectedness.  Hence $I$ is an interval
$\{a,\ldots,b\}$, proving the support statement and uniqueness of $a,b$.

When $|T|\ge2$, the path $P$ has $|T|-1=b-a$ edges.  The only possible edges
inside the interval are
\[
  \{s_i,s_{i+1}\}\qquad(a\le i<b),
\]
also $b-a$ edges.  Since $P$ is connected on every vertex of the interval and
uses only these possible edges, it uses all of them.  Thus it is the monotone
path through $s_a,s_{a+1},\ldots,s_b$, read in one direction or the other.
\end{proof}

\begin{corollary}[Adjacent-rank single-deficit recovery has interval support code]\label{cor:bc-rank-adjacent-interval-code-recovery}
Work in one residual row and moment of \cref{rem:former-post29-tail}.  Put
$q=b+1$, let $S\subset\{1,\ldots,j\}$ have cardinality $q$, and use the
notation of \cref{prop:bc-rank-adjacent-single-deficit-recovery}.  Suppose
the width fixed-witness fiber has
\[
  S=W_q(T)=\{s_1<\cdots<s_q\}
\]
and every internal edge of the anchored path joins two labels
$s_r,s_{r'}$ with $|r-r'|=1$.  Then the anchored support is a consecutive
witness interval, and $\Gamma$ is recovered from
\[
  S,\quad \Gamma^\ast,\quad
  \Gamma_{S\setminus A},\quad
  \rho(v)\in\{0,1\}^q,\quad
  \iota(A)\in\{0,1\}^q,
\]
where $\rho(v)$ is the retained-vertex rank code of
\cref{lem:bc-residual-qbit-retained-vertex-code} and $\iota(A)$ is the
lexicographic rank code of the support interval
$\{s_a,\ldots,s_b\}$ among all nonempty intervals of $S$.  The same assertion
holds in the height fixed-witness fiber with
$S=H_q(T)=\{t_1<\cdots<t_q\}$.

Thus, under the adjacent-rank crossing condition, the support data in the
single-deficit leaf are interval endpoints rather than an arbitrary subset of
the witness set; the retained vertex and interval support still use at most
the allowed $2q$ branch bits in the residual range.  This does not encode the
deleted-only graph on the complement of the interval.
\end{corollary}

\begin{proof}
Apply \cref{lem:bc-rank-adjacent-path-interval-support} to the anchored path
on its support.  It shows that the support is a consecutive interval of the
full witness list and that the anchored path is the monotone path through that
interval up to reversal.  The number of nonempty intervals in a $q$-element
ordered set is
\[
  \frac{q(q+1)}2 .
\]
In the residual range $q\ge15$, this number is at most $2^q$, so the
lexicographic rank of the interval is encoded by a $q$-bit word: at $q=15$
one has $15\cdot16/2=120<2^{15}$, and the implication from $q$ to $q+1$
follows because
\[
  \frac{(q+1)(q+2)}{q(q+1)}=\frac{q+2}{q}\le2 .
\]

The interval code recovers the support characteristic vector
$\mathbf 1_A\in\{0,1\}^S$.  Therefore
\cref{prop:bc-rank-adjacent-single-deficit-recovery} applies and recovers
$\Gamma$ from $S$, $\Gamma^\ast$, $\Gamma_{S\setminus A}$, the recovered
support vector, and the retained-vertex code $\rho(v)$.  The height case is
identical with the ordered list $t_1<\cdots<t_q$.
\end{proof}

\begin{corollary}[Full-support adjacent-rank single-deficit recovery]\label{cor:bc-rank-adjacent-full-support-recovery}
Work in one residual row and moment of \cref{rem:former-post29-tail}.  Put
$q=b+1$, let $S\subset\{1,\ldots,j\}$ have cardinality $q$, and use the
notation of \cref{prop:bc-rank-adjacent-single-deficit-recovery}.  Suppose
the width fixed-witness fiber has
\[
  S=W_q(T)=\{s_1<\cdots<s_q\},
\]
the anchored support is all of $S$, and every internal edge of the anchored
path joins two labels $s_r,s_{r'}$ with $|r-r'|=1$.  Then $\Gamma$ is
recovered from
\[
  S,\quad \Gamma^\ast,\quad \rho(v)\in\{0,1\}^q,
\]
where $\rho(v)$ is the retained-vertex rank code of
\cref{lem:bc-residual-qbit-retained-vertex-code}.  The same assertion holds
in the height fixed-witness fiber with
$S=H_q(T)=\{t_1<\cdots<t_q\}$.

Consequently, in the adjacent-rank single-deficit leaf, the deleted-only
component obstruction is absent once the anchored support is the whole witness
set.  The remaining proper-support case is exactly the case in which the
deleted-only graph on $S\setminus A$ has to be supplied, absorbed into the
target, or ruled out.
\end{corollary}

\begin{proof}
If $A=S$, then the support interval in
\cref{cor:bc-rank-adjacent-interval-code-recovery} is the whole ordered
witness set, so its interval code is determined by $S$.  The graph
$\Gamma_{S\setminus A}$ is the unique empty graph on the empty label set.
Thus the data required by
\cref{cor:bc-rank-adjacent-interval-code-recovery} are fixed functions of
\[
  S,\quad \Gamma^\ast,\quad \rho(v).
\]
That corollary recovers $\Gamma$ in the width case, and the height case is
identical.  The final sentence is the contrapositive localization of the only
datum removed in this full-support case.
\end{proof}

\begin{corollary}[Proper-support adjacent-rank leaves reduce to deleted-only readout]\label{cor:bc-rank-adjacent-proper-support-deleted-only-readout}
Work in one residual row and moment of \cref{rem:former-post29-tail}.  Put
$q=b+1$, let $S\subset\{1,\ldots,j\}$ have cardinality $q$, and use the
notation of \cref{cor:bc-rank-adjacent-interval-code-recovery}.  Suppose the
width fixed-witness fiber has
\[
  S=W_q(T)=\{s_1<\cdots<s_q\},
\]
the anchored support $A$ is a proper nonempty interval of $S$, and every
internal edge of the anchored path joins two labels $s_r,s_{r'}$ with
$|r-r'|=1$.  Let $\Delta_A$ be any datum, determined without using an
additional branch word, which determines the deleted-only graph
\[
  \Gamma_{S\setminus A}.
\]
Then $\Gamma$ is recovered from
\[
  S,\quad \Gamma^\ast,\quad \rho(v)\in\{0,1\}^q,
  \quad \iota(A)\in\{0,1\}^q,\quad \Delta_A .
\]
The same assertion holds in the height fixed-witness fiber with
$S=H_q(T)=\{t_1<\cdots<t_q\}$.

Consequently, after the adjacent-rank crossing condition has been proved, the
proper-support single-deficit leaf has no remaining path-order, support-set,
or retained-vertex ambiguity.  Its remaining supplier obligation is exactly a
branch-free target-path readout of $\Gamma_{S\setminus A}$, or a theorem
excluding proper support.
\end{corollary}

\begin{proof}
The interval code $\iota(A)$ recovers the proper support interval $A$.
Together with the retained-vertex code $\rho(v)$, the datum $\Delta_A$
therefore supplies every input listed in
\cref{cor:bc-rank-adjacent-interval-code-recovery}: the deleted-only graph is
the graph recovered from $\Delta_A$.  Applying that corollary recovers
$\Gamma$ in the width case, and the height case is identical.

The final sentence is just the same recovery statement read against the branch
accounting in \cref{cor:bc-rank-adjacent-interval-code-recovery}: the allowed
$2q$ branch bits are already the $q$-bit retained-vertex code and the $q$-bit
interval code, so the remaining deleted-only graph must be supplied by the
target-path/parser data or eliminated by proving $A=S$.
\end{proof}

\begin{lemma}[Line-adjacent deleted-only complements are forced]\label{lem:bc-line-adjacent-deleted-only-complement-forced}
Let
\[
  S=\{s_1<\cdots<s_q\}
\]
be a finite linearly ordered label set, and let
\[
  A=\{s_a,s_{a+1},\ldots,s_b\}
  \qquad(1\le a\le b\le q)
\]
be a nonempty interval.  Put $D=S\setminus A$.  Let $\Gamma_D$ be a loopless
labelled $2$-regular multigraph on $D$, with parallel edges allowed.  Suppose
every edge of $\Gamma_D$ joins two labels $s_r,s_{r'}$ with
\[
  |r-r'|=1 .
\]
Then the following dichotomy holds.
\begin{enumerate}
\item If either $a-1$ or $q-b$ is odd, no such graph $\Gamma_D$ exists.
\item If both $a-1$ and $q-b$ are even, then $\Gamma_D$ is uniquely determined
by $S$ and $A$: it consists of two parallel edges between
$s_{2u-1}$ and $s_{2u}$ for $1\le u\le(a-1)/2$, and two parallel edges between
$s_{b+2u-1}$ and $s_{b+2u}$ for $1\le u\le(q-b)/2$.
\end{enumerate}
\end{lemma}

\begin{proof}
The deleted set $D$ is the disjoint union of the left interval
$\{s_1,\ldots,s_{a-1}\}$ and the right interval
$\{s_{b+1},\ldots,s_q\}$.  Since $A$ is nonempty, no label in the left interval
is adjacent in full witness rank to a label in the right interval.  Hence the
adjacency hypothesis makes $\Gamma_D$ the disjoint union of its restrictions to
these two intervals.

It remains to analyze one contiguous interval
\[
  I=\{s_\ell,s_{\ell+1},\ldots,s_r\}
\]
of deleted labels.  Write $d=r-\ell+1$.  If $d=0$, there is nothing to prove.
If $d\ge1$, let $m_i$ be the multiplicity of the edge between $s_i$ and
$s_{i+1}$ for $\ell\le i<r$.  The adjacency hypothesis and looplessness allow
no other edges inside $I$.  The endpoint $s_\ell$ has degree $m_\ell$, so
$m_\ell=2$ when $d\ge2$, while for $d=1$ degree $2$ is impossible.
For an interior vertex $s_i$ one has
\[
  m_{i-1}+m_i=2 .
\]
Thus the multiplicities are forced alternately as
\[
  m_\ell=2,\quad m_{\ell+1}=0,\quad m_{\ell+2}=2,\quad\ldots .
\]
At the right endpoint the degree condition is $m_{r-1}=2$.  This holds exactly
when $d$ is even.  Therefore an odd-length interval is impossible, while an
even-length interval has the unique doubled-adjacent-pair graph
\[
  (s_\ell,s_{\ell+1}),\ (s_{\ell+2},s_{\ell+3}),\ldots,\ (s_{r-1},s_r),
\]
with two parallel copies of each listed edge.

Applying this interval analysis to the left and right components of $D$ gives
the stated parity obstruction and the displayed unique graph.
\end{proof}

\begin{corollary}[Line-adjacent proper-support leaves need no deleted-only readout]\label{cor:bc-rank-adjacent-line-deleted-only-recovery}
Work in one residual row and moment of \cref{rem:former-post29-tail}.  Put
$q=b+1$, let $S\subset\{1,\ldots,j\}$ have cardinality $q$, and use the
notation of \cref{cor:bc-rank-adjacent-interval-code-recovery}.  Suppose the
width fixed-witness fiber has
\[
  S=W_q(T)=\{s_1<\cdots<s_q\},
\]
the anchored support is a proper nonempty interval
$A=\{s_a,\ldots,s_b\}$, every internal edge of the anchored path joins two
labels $s_r,s_{r'}$ with $|r-r'|=1$, and every edge of the deleted-only graph
$\Gamma_{S\setminus A}$ also joins two labels $s_r,s_{r'}$ with
$|r-r'|=1$ in the full witness list.  If either $a-1$ or $q-b$ is odd, these
hypotheses are inconsistent.  If both are even, then $\Gamma$ is recovered from
\[
  S,\quad \Gamma^\ast,\quad \rho(v)\in\{0,1\}^q,
  \quad \iota(A)\in\{0,1\}^q .
\]
The same assertion holds in the height fixed-witness fiber with
$S=H_q(T)=\{t_1<\cdots<t_q\}$.

Consequently, after the adjacent-rank crossing condition has been proved for
both the anchored component and the deleted-only complement, the proper-support
single-deficit leaf has no deleted-only readout obligation.
\end{corollary}

\begin{proof}
The interval code $\iota(A)$ recovers $A$, hence the integers $a$ and $b$.  If
one of $a-1$ and $q-b$ is odd, then
\cref{lem:bc-line-adjacent-deleted-only-complement-forced} says that no
deleted-only graph satisfying the full-rank adjacency hypothesis exists.  This
gives the inconsistency statement.

If both integers are even, the same lemma recovers the deleted-only graph
$\Gamma_{S\setminus A}$ from $S$ and $A$ alone.  Use this forced graph as the
branch-free datum $\Delta_A$ in
\cref{cor:bc-rank-adjacent-proper-support-deleted-only-readout}.  That
corollary recovers $\Gamma$ in the width case.  The height case is identical,
with the ordered witness list $t_1<\cdots<t_q$ replacing
$s_1<\cdots<s_q$.
\end{proof}

\begin{corollary}[Actual frontier-consecutive deleted-only complements need no readout]\label{cor:bc-actual-frontier-consecutive-deleted-only-recovery}
Work in one residual row and moment of \cref{rem:former-post29-tail}.  Put
$q=b+1$, let $S\subset\{1,\ldots,j\}$ have cardinality $q$, and use the
notation of \cref{cor:bc-rank-adjacent-interval-code-recovery}.  Suppose the
width fixed-witness fiber has
\[
  S=W_q(T)=\{s_1<\cdots<s_q\},
\]
and that the anchored support is a proper nonempty interval
$A=\{s_a,\ldots,s_b\}$.  Assume:
\begin{enumerate}
\item every internal edge of the anchored path on $A$ has endpoints whose
assigned times are consecutive in the full ordered list
\[
  \kappa^w_1<\cdots<\kappa^w_q;
\]
\item every edge of the deleted-only graph $\Gamma_{S\setminus A}$ has endpoints
whose assigned times are consecutive in the same full ordered list.
\end{enumerate}
If either $a-1$ or $q-b$ is odd, these hypotheses are inconsistent.  If both
are even, then $\Gamma$ is recovered from
\[
  S,\quad \Gamma^\ast,\quad \rho(v)\in\{0,1\}^q,
  \quad \iota(A)\in\{0,1\}^q .
\]
The same assertion holds in the height fixed-witness fiber with
$S=H_q(T)=\{t_1<\cdots<t_q\}$ and
$\kappa^h_1<\cdots<\kappa^h_q$.

Consequently the proper-support single-deficit leaf is closed once the actual
fixed-witness parser proves full-frontier consecutiveness for the anchored
internal edges and the deleted-only complement edges.
\end{corollary}

\begin{proof}
For the width case, \cref{lem:bc-consecutive-kappa-adjacent-ranks} turns item
(1) into the adjacent-rank condition for every anchored internal edge.  The
same lemma turns item (2) into the condition that every deleted-only edge joins
labels $s_r,s_{r'}$ with $|r-r'|=1$ in the full witness list.  Therefore the
hypotheses of \cref{cor:bc-rank-adjacent-line-deleted-only-recovery} are
satisfied, and that corollary gives the inconsistency/recovery dichotomy.

The height case is identical, using the ordered list
$\kappa^h_1<\cdots<\kappa^h_q$ and the height witness labels
$t_1<\cdots<t_q$.
\end{proof}

\begin{proposition}[Actual frontier-consecutive single-deficit graphs use only the interval code]\label{prop:bc-actual-frontier-consecutive-single-deficit-interval-code}
Work in one residual row and moment of \cref{rem:former-post29-tail}.  Put
$q=b+1$, let $S\subset\{1,\ldots,j\}$ have cardinality $q$, and work in the
setting of \cref{prop:bc-single-deficit-stable-retained-repair}.  Let
$A\subset S$ be the support of the anchored component in
\cref{lem:bc-single-deficit-support-path-code}.

In the width fixed-witness fiber suppose
\[
  S=W_q(T)=\{s_1<\cdots<s_q\},
\]
with full frontier list
\[
  \kappa^w_1<\cdots<\kappa^w_q .
\]
Assume that every internal edge of the anchored path on $A$ has endpoints whose
assigned times are consecutive in this full list, and that every edge of the
deleted-only graph $\Gamma_{S\setminus A}$ also has endpoints whose assigned
times are consecutive in this full list.  Then $A$ is a nonempty interval in
the ordered witness set, and $\Gamma$ is recovered from
\[
  S,\quad \Gamma^\ast,\quad \rho(v)\in\{0,1\}^q,
  \quad \iota(A)\in\{0,1\}^q .
\]
The same assertion holds in the height fixed-witness fiber with
$S=H_q(T)=\{t_1<\cdots<t_q\}$ and
$\kappa^h_1<\cdots<\kappa^h_q$.

Consequently the single-deficit graph leaf, under actual full-frontier
consecutiveness for both the anchored component and the deleted-only
complement, consumes only the already budgeted retained-vertex rank code and
support-interval code.
\end{proposition}

\begin{proof}
Consider the width case.  The component containing $v$ has nonempty support
$A$, because the single retained deficit has the two ports at $v$ attached to
deleted labels by \cref{lem:bc-single-deficit-port-code}.  By
\cref{lem:bc-consecutive-kappa-adjacent-ranks}, the full-list consecutiveness of
the anchored internal edges implies that every such edge joins adjacent witness
ranks.  Then \cref{lem:bc-rank-adjacent-path-interval-support} shows that $A$
is a nonempty interval.

If $A=S$, then \cref{cor:bc-rank-adjacent-full-support-recovery} recovers
$\Gamma$ from $S$, $\Gamma^\ast$, and $\rho(v)$, so the displayed data recover
$\Gamma$ as well.  If $A$ is a proper interval, then
\cref{cor:bc-actual-frontier-consecutive-deleted-only-recovery} applies to the
same full-list consecutiveness hypotheses.  Its inconsistent parity alternative
cannot occur for the actual graph under discussion, and its recovery
alternative gives $\Gamma$ from the displayed data.

The height case is identical with the ordered height witness list and
$\kappa^h_1<\cdots<\kappa^h_q$ replacing the width data.
\end{proof}

\begin{lemma}[Small deleted-only complements are forced]\label{lem:bc-small-deleted-only-complements}
Let $D$ be a finite label set, and let $\Gamma_D$ be a loopless labelled
$2$-regular multigraph on $D$, with parallel edges allowed.  If $|D|=0$, then
$\Gamma_D$ is the unique empty graph.  If $|D|=1$, no such graph exists.  If
$|D|=2$, then $\Gamma_D$ is the unique doubled edge on the two labels.  If
$|D|=3$, then $\Gamma_D$ is the unique simple triangle on the three labels.
\end{lemma}

\begin{proof}
For $|D|=0$, the empty edge multiset is the only possible graph and the degree
condition is vacuous.  For $|D|=1$, looplessness forbids an edge, so the unique
vertex would have degree $0$, not degree $2$.

For $|D|=2$, each edge must join the two distinct labels.  Degree $2$ at either
label therefore forces exactly two parallel copies of that edge.

For $|D|=3$, a pair of parallel edges between two labels would already use
degree $2$ at those two labels and would leave the third label with degree
$0$, so no parallel edge can occur.  Thus the graph is simple.  The only
simple graph on three labels in which every vertex has degree $2$ is the
triangle.
\end{proof}

\begin{corollary}[Near-full proper-support leaves need no deleted-only readout]\label{cor:bc-rank-adjacent-near-full-proper-support}
Work in one residual row and moment of \cref{rem:former-post29-tail}.  Put
$q=b+1$, let $S\subset\{1,\ldots,j\}$ have cardinality $q$, and use the
notation of
\cref{cor:bc-rank-adjacent-proper-support-deleted-only-readout}.  Suppose the
width fixed-witness fiber has
\[
  S=W_q(T)=\{s_1<\cdots<s_q\},
\]
the anchored support $A$ is a proper nonempty interval of $S$, every internal
edge of the anchored path joins two labels $s_r,s_{r'}$ with $|r-r'|=1$, and
\[
  |S\setminus A|\le3 .
\]
If $|S\setminus A|=1$, then these hypotheses are inconsistent.  If
$|S\setminus A|=2$ or $|S\setminus A|=3$, then $\Gamma$ is recovered from
\[
  S,\quad \Gamma^\ast,\quad \rho(v)\in\{0,1\}^q,
  \quad \iota(A)\in\{0,1\}^q .
\]
The same assertion holds in the height fixed-witness fiber with
$S=H_q(T)=\{t_1<\cdots<t_q\}$.

Consequently, after the adjacent-rank crossing condition has been proved, the
only proper-support deleted-only readout still not covered by forced small
complements alone is the case
\[
  |S\setminus A|\ge4 .
\]
\end{corollary}

\begin{proof}
Put $D=S\setminus A$.  By
\cref{lem:bc-single-deficit-port-components}, the deleted-only part induced on
$D$ is a loopless labelled $2$-regular multigraph.  If $|D|=1$, this
contradicts \cref{lem:bc-small-deleted-only-complements}, so the hypotheses are
inconsistent.

If $|D|=2$ or $|D|=3$, then
\cref{lem:bc-small-deleted-only-complements} says that the deleted-only graph
on $D$ is uniquely determined by the label set $D$, hence by $S$ and
$\iota(A)$.  Use this forced graph as the branch-free datum $\Delta_A$ in
\cref{cor:bc-rank-adjacent-proper-support-deleted-only-readout}.  That
corollary recovers $\Gamma$ in the width case, and the height case is
identical.

The final sentence combines these two cases with
\cref{cor:bc-rank-adjacent-full-support-recovery}, which already covers
$D=\varnothing$.
\end{proof}

\begin{corollary}[Small deleted-only complements fit in the support code]\label{cor:bc-rank-adjacent-small-complement-support-code}
Work in one residual row and moment of \cref{rem:former-post29-tail}.  Put
$q=b+1$, let $S\subset\{1,\ldots,j\}$ have cardinality $q$, and use the
notation of \cref{cor:bc-rank-adjacent-interval-code-recovery}.  Suppose the
width fixed-witness fiber has
\[
  S=W_q(T)=\{s_1<\cdots<s_q\},
\]
the anchored support $A$ is a nonempty interval of $S$, every internal edge of
the anchored path joins two labels $s_r,s_{r'}$ with $|r-r'|=1$, and
\[
  |S\setminus A|\le7 .
\]
Then there is a $q$-bit code
\[
  \zeta(A,\Gamma_{S\setminus A})\in\{0,1\}^q,
\]
depending only on the interval $A$ and on the deleted-only graph induced on
$S\setminus A$, such that $\Gamma$ is recovered from
\[
  S,\quad \Gamma^\ast,\quad \rho(v)\in\{0,1\}^q,
  \quad \zeta(A,\Gamma_{S\setminus A})\in\{0,1\}^q .
\]
The same assertion holds in the height fixed-witness fiber with
$S=H_q(T)=\{t_1<\cdots<t_q\}$.

Consequently, after the adjacent-rank crossing condition has been proved, the
proper-support deleted-only readout that is not already absorbed by the
existing support-code slot begins only in the case
\[
  |S\setminus A|\ge8 .
\]
\end{corollary}

\begin{proof}
For a fixed ordered $q$-element witness set $S$, the number of nonempty
intervals $A\subseteq S$ with $|S\setminus A|=d$ is $d+1$.  For such an
interval, \cref{lem:bc-littlewood-graph-model} gives exactly $s_d$ possible
loopless labelled $2$-regular graphs on $S\setminus A$, where $s_d$ is the
stable moment sequence of \cref{lem:bc-stable-moment-recurrence}.  Iterating
that recurrence gives
\[
  s_0=1,\quad s_1=0,\quad s_2=1,\quad s_3=1,\quad
  s_4=6,\quad s_5=22,\quad s_6=130,\quad s_7=822 .
\]
Therefore the number of pairs $(A,\Gamma_{S\setminus A})$ with
$|S\setminus A|\le7$ is
\[
  \sum_{d=0}^{7}(d+1)s_d
  =
  7656
  <
  32768
  =
  2^{15}
  \le 2^q ,
\]
because every residual row has $q\ge15$.  Listing these pairs in
lexicographic order gives the advertised $q$-bit code
$\zeta(A,\Gamma_{S\setminus A})$.

The code recovers both $A$ and the deleted-only graph
$\Gamma_{S\setminus A}$.  Applying
\cref{cor:bc-rank-adjacent-interval-code-recovery} with these recovered data
recovers $\Gamma$ in the width case, and the height case is identical.
The final sentence is the same branch accounting: the $q$ retained-vertex bits
and the $q$ packed support/deleted-only bits are the whole allowed $2q$ branch
budget.
\end{proof}

\begin{corollary}[Residual joint branch code absorbs complements of size nine]\label{cor:bc-rank-adjacent-nine-complement-joint-code}
Work in one residual row and moment of \cref{rem:former-post29-tail}.  Put
$q=b+1$, let $S\subset\{1,\ldots,j\}$ have cardinality $q$, and use the
notation of \cref{cor:bc-rank-adjacent-interval-code-recovery}.  Suppose the
width fixed-witness fiber has
\[
  S=W_q(T)=\{s_1<\cdots<s_q\},
\]
the anchored support $A$ is a nonempty interval of $S$, every internal edge of
the anchored path joins two labels $s_r,s_{r'}$ with $|r-r'|=1$, and
\[
  |S\setminus A|\le9 .
\]
Then there is a $2q$-bit code
\[
  \xi(v,A,\Gamma_{S\setminus A})\in\{0,1\}^{2q},
\]
depending only on the retained deficit vertex $v$, the interval $A$, and the
deleted-only graph induced on $S\setminus A$, such that $\Gamma$ is recovered
from
\[
  S,\quad \Gamma^\ast,\quad
  \xi(v,A,\Gamma_{S\setminus A})\in\{0,1\}^{2q}.
\]
The same assertion holds in the height fixed-witness fiber with
$S=H_q(T)=\{t_1<\cdots<t_q\}$.

Consequently, after the adjacent-rank crossing condition has been proved, the
proper-support deleted-only readout that is not already absorbed by the total
$2q$ branch budget begins only in the case
\[
  |S\setminus A|\ge10 .
\]
\end{corollary}

\begin{proof}
For fixed $S$ and a fixed complement size $d$, there are $d+1$ nonempty
intervals $A\subseteq S$ with $|S\setminus A|=d$.  For each such interval,
\cref{lem:bc-littlewood-graph-model} gives exactly $s_d$ possible deleted-only
graphs on $S\setminus A$.  Iterating
\cref{lem:bc-stable-moment-recurrence} one more step beyond
\cref{cor:bc-rank-adjacent-small-complement-support-code} gives
\[
  s_8=6202,\qquad s_9=52552,
\]
and hence
\[
  M_9:=\sum_{d=0}^{9}(d+1)s_d
  =
  588994 .
\]
Thus, for fixed $S$, the number of triples
\[
  (v,A,\Gamma_{S\setminus A})
  \qquad\text{with}\qquad |S\setminus A|\le9
\]
is at most
\[
  (j-q)M_9 .
\]
Since $j$ lies in the residual window of \cref{rem:former-post29-tail}, one has
$j\le N_{\mathrm{dir}}^{BC}(G)+2$.

For the first three residual $B$ rows, \cref{cor:post29-bc-initial-b-onsets}
gives the exact onset values.  Therefore
\[
\begin{array}{c|c|c|c}
q & j-q\le & (j-q)M_9 & 2^{2q}\\ \hline
15 & 912 & 537162528 & 1073741824\\
16 & 401 & 236186594 & 4294967296\\
17 & 314 & 184944116 & 17179869184 .
\end{array}
\]
In each of these rows, $(j-q)M_9<2^{2q}$.  For all remaining residual rows
one has $q\ge18$, and \cref{cor:post29-bc-direct-onset-bound} gives
\[
  j-q\le30899-q\le30881 .
\]
Consequently
\[
  (j-q)M_9
  \le
  30881\cdot588994
  =
  18188723714
  <
  68719476736
  =
  2^{36}
  \le 2^{2q}.
\]

Listing all triples $(v,A,\Gamma_{S\setminus A})$ in lexicographic order gives
the advertised $2q$-bit code $\xi$.  The code recovers $v$, $A$, and the
deleted-only graph.  Applying
\cref{cor:bc-rank-adjacent-interval-code-recovery} with these recovered inputs
recovers $\Gamma$ in the width case, and the height case is identical.  The
final sentence is exactly the displayed packing statement applied to proper
support.
\end{proof}

\begin{corollary}[Stratified residual joint branch code]\label{cor:bc-rank-adjacent-stratified-complement-joint-code}
Work in one residual row and moment of \cref{rem:former-post29-tail}.  Put
$q=b+1$, let $S\subset\{1,\ldots,j\}$ have cardinality $q$, and use the
notation of \cref{cor:bc-rank-adjacent-interval-code-recovery}.  Define
\[
K(q)=
\begin{cases}
9, & q=15,\\
10, & 16\le q\le17,\\
11, & q=18,\\
12, & 19\le q\le20,\\
13, & 21\le q\le22,\\
14, & 23\le q\le24,\\
15, & q\ge25 .
\end{cases}
\]
Suppose the width fixed-witness fiber has
\[
  S=W_q(T)=\{s_1<\cdots<s_q\},
\]
the anchored support $A$ is a nonempty interval of $S$, every internal edge of
the anchored path joins two labels $s_r,s_{r'}$ with $|r-r'|=1$, and
\[
  |S\setminus A|\le K(q).
\]
Then there is a $2q$-bit code
\[
  \Xi(v,A,\Gamma_{S\setminus A})\in\{0,1\}^{2q},
\]
depending only on the retained deficit vertex $v$, the interval $A$, and the
deleted-only graph induced on $S\setminus A$, such that $\Gamma$ is recovered
from
\[
  S,\quad \Gamma^\ast,\quad
  \Xi(v,A,\Gamma_{S\setminus A})\in\{0,1\}^{2q}.
\]
The same assertion holds in the height fixed-witness fiber with
$S=H_q(T)=\{t_1<\cdots<t_q\}$.

Consequently, after the adjacent-rank crossing condition has been proved, the
unabsorbed proper-support deleted-only readout is restricted to the cases
\[
  |S\setminus A|\ge K(q)+1 .
\]
\end{corollary}

\begin{proof}
Let
\[
  M_K=\sum_{d=0}^{K}(d+1)s_d .
\]
As in the proof of \cref{cor:bc-rank-adjacent-nine-complement-joint-code},
for fixed $S$ the number of triples
$(v,A,\Gamma_{S\setminus A})$ with $|S\setminus A|\le K$ is at most
$(j-q)M_K$.  Iterating \cref{lem:bc-stable-moment-recurrence} gives
\[
\begin{array}{c|c}
K & M_K\\ \hline
9 & 588994\\
10 & 6080128\\
11 & 68940568\\
12 & 851837540\\
13 & 11388066236\\
14 & 163757241896\\
15 & 2520032130024 .
\end{array}
\]
Combining these values with the residual window bounds of
\cref{cor:post29-bc-window-size-strata} gives
\[
\begin{array}{c|c|c|c}
\text{range of }q & K(q) & (j-q)M_{K(q)}\le & 2^{2q}\text{ lower bound}\\
\hline
q=15 & 9 & 537162528 & 1073741824\\
16\le q\le17 & 10 & 2438131328 & 4294967296\\
q=18 & 11 & 19785943016 & 68719476736\\
19\le q\le20 & 12 & 241070023820 & 274877906944\\
21\le q\le22 & 13 & 3450584069508 & 4398046511104\\
23\le q\le24 & 14 & 55186190518952 & 70368744177664\\
25\le q\le28 & 15 & 1065973591000152 & 1125899906842624\\
q\ge29 & 15 & 77793391853840880 & 288230376151711744 .
\end{array}
\]
In every row of the table the third entry is strictly smaller than the fourth.
Thus the triples $(v,A,\Gamma_{S\setminus A})$ with
$|S\setminus A|\le K(q)$ inject into the $2^{2q}$ possible branch words.
Listing them in lexicographic order gives the code $\Xi$.

The code recovers $v$, $A$, and the deleted-only graph.  Applying
\cref{cor:bc-rank-adjacent-interval-code-recovery} with these recovered inputs
recovers $\Gamma$ in the width case; the height case is identical.  The final
sentence is the complement of the displayed packing range.
\end{proof}

\begin{corollary}[High-rank linear joint branch code]\label{cor:bc-rank-adjacent-linear-complement-joint-code}
Work in one residual row and moment of \cref{rem:former-post29-tail}.  Put
$q=b+1$, assume $q\ge45$, let $S\subset\{1,\ldots,j\}$ have cardinality $q$,
and use the notation of \cref{cor:bc-rank-adjacent-interval-code-recovery}.
Set
\[
  L(q)=\left\lfloor\frac q4\right\rfloor+10 .
\]
Suppose the width fixed-witness fiber has
\[
  S=W_q(T)=\{s_1<\cdots<s_q\},
\]
the anchored support $A$ is a nonempty interval of $S$, every internal edge of
the anchored path joins two labels $s_r,s_{r'}$ with $|r-r'|=1$, and
\[
  |S\setminus A|\le L(q).
\]
Then there is a $2q$-bit code
\[
  \Theta(v,A,\Gamma_{S\setminus A})\in\{0,1\}^{2q},
\]
depending only on the retained deficit vertex $v$, the interval $A$, and the
deleted-only graph induced on $S\setminus A$, such that $\Gamma$ is recovered
from
\[
  S,\quad \Gamma^\ast,\quad
  \Theta(v,A,\Gamma_{S\setminus A})\in\{0,1\}^{2q}.
\]
The same assertion holds in the height fixed-witness fiber with
$S=H_q(T)=\{t_1<\cdots<t_q\}$.

Consequently, for residual rows with $q\ge45$, after the adjacent-rank
crossing condition has been proved, the unabsorbed proper-support deleted-only
readout is restricted to
\[
  |S\setminus A|\ge \left\lfloor\frac q4\right\rfloor+11 .
\]
\end{corollary}

\begin{proof}
By \cref{cor:post29-bc-direct-onset-bound}, every residual row has
$j\le30899$.  Hence the retained deficit vertex has at most $j-q\le30899$
choices.  Let $L=L(q)$.  As in the proof of
\cref{cor:bc-rank-adjacent-stratified-complement-joint-code}, for fixed $S$
the number of triples
\[
  (v,A,\Gamma_{S\setminus A})
  \qquad\text{with}\qquad |S\setminus A|\le L
\]
is at most
\[
  30899\sum_{d=0}^L(d+1)s_d .
\]
By \cref{lem:bc-cumulative-deleted-only-factorial-bound}, this is at most
\[
  30899\,(L+2)! .
\]

We claim that
\[
  30899\,(L(q)+2)!<2^{2q}
  \qquad(45\le q\le218).
\]
At $q=45$, $L(q)=21$ and the exact integer inequality is
\[
  30899\cdot23!
  =
  798801465214806893199360000
  <
  1237940039285380274899124224
  =
  2^{90}.
\]
When $q$ first reaches $48$, $L(q)$ increases from $21$ to $22$; the left side
is multiplied by $24$, while the right side from $q=45$ to $q=48$ is
multiplied by $2^6=64$.  After that, each further increase of $L(q)$ by one
requires four increases of $q$, so the right side is multiplied by $2^8=256$.
Throughout the residual range $q\le218$, one has $L(q)\le64$, so the left
side is multiplied by at most $66$ at each such step.  Hence the displayed
inequality holds for every $45\le q\le218$.

Thus the triples $(v,A,\Gamma_{S\setminus A})$ with
$|S\setminus A|\le L(q)$ inject into the $2^{2q}$ possible branch words.
Listing them in lexicographic order gives the code $\Theta$.  The code
recovers $v$, $A$, and the deleted-only graph.  Applying
\cref{cor:bc-rank-adjacent-interval-code-recovery} with these recovered inputs
recovers $\Gamma$ in the width case; the height case is identical.  The final
sentence is the complement of the displayed packing range.
\end{proof}

\begin{corollary}[Residual one-third joint branch code]\label{cor:bc-rank-adjacent-one-third-complement-joint-code}
Work in one residual row and moment of \cref{rem:former-post29-tail}.  Put
$q=b+1$, let $S\subset\{1,\ldots,j\}$ have cardinality $q$, and use the
notation of \cref{cor:bc-rank-adjacent-interval-code-recovery}.  Set
\[
  H(q)=\left\lceil\frac q3\right\rceil
      =\left\lfloor\frac{q+2}{3}\right\rfloor .
\]
Suppose the width fixed-witness fiber has
\[
  S=W_q(T)=\{s_1<\cdots<s_q\},
\]
the anchored support $A$ is a nonempty interval of $S$, every internal edge of
the anchored path joins two labels $s_r,s_{r'}$ with $|r-r'|=1$, and
\[
  |S\setminus A|\le H(q).
\]
Then there is a $2q$-bit code
\[
  \Upsilon(v,A,\Gamma_{S\setminus A})\in\{0,1\}^{2q},
\]
depending only on the retained deficit vertex $v$, the interval $A$, and the
deleted-only graph induced on $S\setminus A$, such that $\Gamma$ is recovered
from
\[
  S,\quad \Gamma^\ast,\quad
  \Upsilon(v,A,\Gamma_{S\setminus A})\in\{0,1\}^{2q}.
\]
The same assertion holds in the height fixed-witness fiber with
$S=H_q(T)=\{t_1<\cdots<t_q\}$.

Consequently, after the adjacent-rank crossing condition has been proved, the
unabsorbed proper-support deleted-only readout is restricted to
\[
  |S\setminus A|\ge \left\lceil\frac q3\right\rceil+1 .
\]
\end{corollary}

\begin{proof}
As in \cref{cor:bc-rank-adjacent-linear-complement-joint-code}, it is enough
to prove
\[
  30899\,(H(q)+2)!<2^{2q}
  \qquad(15\le q\le218).
\]
For $H(q)=5$, the only residual case is $q=15$, and
\[
  30899\cdot7!
  =
  155730960
  <
  1073741824
  =
  2^{30}.
\]
For $H(q)=6$, the worst residual case is $q=16$, and
\[
  30899\cdot8!
  =
  1245847680
  <
  4294967296
  =
  2^{32}.
\]

It remains to handle $7\le h:=H(q)\le73$.  For fixed $h$, the smallest
residual $q$ with $H(q)=h$ is $q=3h-2$, so it suffices to prove
\[
  30899\,(h+2)!<2^{6h-4}.
\]
At $h=7$ one has the stronger margin
\[
  8\cdot30899\cdot9!
  =
  89701032960
  <
  274877906944
  =
  2^{38}.
\]
For $7\le h\le61$, increasing $h$ by one multiplies the left side by
$h+3\le64$ and the right side by $2^6=64$, so this factor-$8$ margin is
preserved through $h=61$.  For $62\le h\le73$, the total possible loss of
margin from $h=61$ to $h=73$ is at most
\[
  \left(\frac{75}{64}\right)^{12}<8 .
\]
Thus the strict inequality still holds for every $62\le h\le73$.  Since
$q\le218$ gives $H(q)\le73$, the displayed bound holds in the whole residual
range.

By \cref{lem:bc-cumulative-deleted-only-factorial-bound}, the number of triples
$(v,A,\Gamma_{S\setminus A})$ with $|S\setminus A|\le H(q)$ is then strictly
less than $2^{2q}$.  Listing those triples lexicographically gives the code
$\Upsilon$, and the recovery follows from
\cref{cor:bc-rank-adjacent-interval-code-recovery} exactly as in
\cref{cor:bc-rank-adjacent-linear-complement-joint-code}.  The height case is
identical.
\end{proof}

\begin{corollary}[High-rank one-third-plus-eight joint branch code]\label{cor:bc-rank-adjacent-onethird-plus-eight-complement-joint-code}
Work in one residual row and moment of \cref{rem:former-post29-tail}.  Put
$q=b+1$, assume $q\ge80$, let $S\subset\{1,\ldots,j\}$ have cardinality $q$,
and use the notation of \cref{cor:bc-rank-adjacent-interval-code-recovery}.
Set
\[
  P(q)=\left\lceil\frac q3\right\rceil+8 .
\]
Suppose the width fixed-witness fiber has
\[
  S=W_q(T)=\{s_1<\cdots<s_q\},
\]
the anchored support $A$ is a nonempty interval of $S$, every internal edge of
the anchored path joins two labels $s_r,s_{r'}$ with $|r-r'|=1$, and
\[
  |S\setminus A|\le P(q).
\]
Then there is a $2q$-bit code
\[
  \Pi(v,A,\Gamma_{S\setminus A})\in\{0,1\}^{2q},
\]
depending only on the retained deficit vertex $v$, the interval $A$, and the
deleted-only graph induced on $S\setminus A$, such that $\Gamma$ is recovered
from
\[
  S,\quad \Gamma^\ast,\quad
  \Pi(v,A,\Gamma_{S\setminus A})\in\{0,1\}^{2q}.
\]
The same assertion holds in the height fixed-witness fiber with
$S=H_q(T)=\{t_1<\cdots<t_q\}$.

Consequently, for residual rows with $q\ge80$, after the adjacent-rank
crossing condition has been proved, the unabsorbed proper-support deleted-only
readout is restricted to
\[
  |S\setminus A|\ge \left\lceil\frac q3\right\rceil+9 .
\]
\end{corollary}

\begin{proof}
As in \cref{cor:bc-rank-adjacent-linear-complement-joint-code}, it is enough
to prove
\[
  30899\,(P(q)+2)!<2^{2q}
  \qquad(80\le q\le218).
\]
Put $h=P(q)$.  Then $35\le h\le81$.  For $h=35$, the worst residual case is
$q=80$, and the exact integer inequality is
\[
  30899\cdot37!<2^{160}.
\]
For $h=36$, the worst residual case is $q=82$, and
\[
  30899\cdot38!<2^{164}.
\]
For $36\le h\le61$, increasing $h$ by one increases the worst-case $q$ by
three, so the right side is multiplied by $2^6=64$, while the left side is
multiplied by $h+3\le64$.  Thus the displayed inequality propagates from
$h=36$ through $h=61$.

At $h=61$, we have the stronger margin
\[
  256\cdot30899\cdot63!<2^{314}.
\]
For $62\le h\le81$, each further step multiplies the left side by at most
$84$ and the right side by $64$.  The total possible loss of margin is at most
\[
  \left(\frac{84}{64}\right)^{20}
  =
  \left(\frac{21}{16}\right)^{20}
  <
  256,
\]
because $(21/16)^4<3$.  Hence the strict inequality remains valid through
$h=81$, which covers the residual endpoint $q=218$.

By \cref{lem:bc-cumulative-deleted-only-factorial-bound}, the number of triples
$(v,A,\Gamma_{S\setminus A})$ with $|S\setminus A|\le P(q)$ is then strictly
less than $2^{2q}$.  Listing those triples lexicographically gives the code
$\Pi$, and the recovery follows from
\cref{cor:bc-rank-adjacent-interval-code-recovery} exactly as in
\cref{cor:bc-rank-adjacent-linear-complement-joint-code}.  The height case is
identical.
\end{proof}

\begin{corollary}[Exact high-rank one-third-plus-ten joint branch code]\label{cor:bc-rank-adjacent-exact-onethird-plus-ten-complement-joint-code}
Work in one residual row and moment of \cref{rem:former-post29-tail}.  Put
$q=b+1$, assume $q\ge86$, let $S\subset\{1,\ldots,j\}$ have cardinality $q$,
and use the notation of \cref{cor:bc-rank-adjacent-interval-code-recovery}.
Set
\[
  Q(q)=\left\lceil\frac q3\right\rceil+10 .
\]
Suppose the width fixed-witness fiber has
\[
  S=W_q(T)=\{s_1<\cdots<s_q\},
\]
the anchored support $A$ is a nonempty interval of $S$, every internal edge of
the anchored path joins two labels $s_r,s_{r'}$ with $|r-r'|=1$, and
\[
  |S\setminus A|\le Q(q).
\]
Then there is a $2q$-bit code
\[
  \Psi(v,A,\Gamma_{S\setminus A})\in\{0,1\}^{2q},
\]
depending only on the retained deficit vertex $v$, the interval $A$, and the
deleted-only graph induced on $S\setminus A$, such that $\Gamma$ is recovered
from
\[
  S,\quad \Gamma^\ast,\quad
  \Psi(v,A,\Gamma_{S\setminus A})\in\{0,1\}^{2q}.
\]
The same assertion holds in the height fixed-witness fiber with
$S=H_q(T)=\{t_1<\cdots<t_q\}$.

Consequently, for residual rows with $q\ge86$, after the adjacent-rank
crossing condition has been proved, the unabsorbed proper-support deleted-only
readout is restricted to
\[
  |S\setminus A|\ge \left\lceil\frac q3\right\rceil+11 .
\]
\end{corollary}

\begin{proof}
Let
\[
  M_K=\sum_{d=0}^K(d+1)s_d
\]
for the stable moment sequence of \cref{lem:bc-stable-moment-recurrence}.  As
in \cref{cor:bc-rank-adjacent-linear-complement-joint-code}, the number of
triples $(v,A,\Gamma_{S\setminus A})$ with $|S\setminus A|\le Q(q)$ is at
most
\[
  30899\,M_{Q(q)}
\]
because $j-q\le30899$ in the residual window.  It remains to check
\[
  30899\,M_{\lceil q/3\rceil+10}<2^{2q}
  \qquad(86\le q\le218).
\]
The exact C++/GMP/OpenMP replay
\[
\texttt{certificates/post\_m29/bc\_complement\_packing\_gmp\_certificate.log}
\]
computes the $s_d$ from \cref{lem:bc-stable-moment-recurrence}, forms the
cumulative sums $M_K$, and checks the displayed integer inequality for all
$133$ integers $q$ in the range $86\le q\le218$.  The replay was run on
\texttt{optimus}; it exits with status $0$, reports \texttt{failures=0},
\texttt{q\_checked=133}, \texttt{k\_min=39}, \texttt{k\_max=83}, and has
SHA-256
\[
\hashsplit{bbbfd5c129c4f8d70e10d38dd1dbc5cf}{24ce4a4cd3491c5e579aeaeb90253125}.
\]
The log reports its closest approximate ratio at $q=88$, $K=40$, with
\[
\log_{10}\left(30899\,M_{40}/2^{176}\right)
\approx -0.115519519802795401<0.
\]
Therefore the triples with $|S\setminus A|\le Q(q)$ inject into the
$2^{2q}$ possible branch words.  Listing them lexicographically gives the
code $\Psi$, and the recovery follows from
\cref{cor:bc-rank-adjacent-interval-code-recovery} exactly as in
\cref{cor:bc-rank-adjacent-linear-complement-joint-code}.  The height case is
identical.
\end{proof}

\begin{definition}[Exact residual complement frontier]\label{def:bc-exact-complement-frontier}
For $15\le q\le218$, let
\[
  M_K=\sum_{d=0}^K(d+1)s_d
\]
for the stable moment sequence of \cref{lem:bc-stable-moment-recurrence}, and
let $W(q)$ be the residual window bound for $j-q$:
\[
W(q)=
\begin{cases}
912, & q=15,\\
401, & 16\le q\le17,\\
287, & q=18,\\
283, & 19\le q\le20,\\
303, & 21\le q\le22,\\
337, & 23\le q\le24,\\
423, & 25\le q\le28,\\
30899-q, & q\ge29.
\end{cases}
\]
Define $F(q)$ to be the largest integer $K$ with $0\le K\le q-2$ such that
\[
  W(q)M_K<2^{2q}.
\]
\end{definition}

\begin{corollary}[Exact residual complement frontier joint branch code]\label{cor:bc-rank-adjacent-exact-frontier-complement-joint-code}
Work in one residual row and moment of \cref{rem:former-post29-tail}.  Put
$q=b+1$, let $S\subset\{1,\ldots,j\}$ have cardinality $q$, and use the
notation of \cref{cor:bc-rank-adjacent-interval-code-recovery}.  Let
$F(q)$ and $W(q)$ be as in \cref{def:bc-exact-complement-frontier}.
The exact C++/GMP/OpenMP replay
\[
\texttt{certificates/post\_m29/bc\_complement\_frontier\_gmp\_certificate.log}
\]
computes these frontiers for every $15\le q\le218$ from the recurrence of
\cref{lem:bc-stable-moment-recurrence}; it exits with status $0$, reports
\texttt{failures=0}, \texttt{q\_checked=204}, \texttt{min\_frontier=9},
\texttt{max\_frontier=83}, and \texttt{max\_frontier\_minus\_q=-6}.  Its
SHA-256 is
\[
\hashsplit{09ae87166af7f93ae6ef4df0322fcacd}{ec86fc6ae5701ce0117851a655643b70}.
\]
Its closest accepted ratio is at $q=62$, $F(q)=30$, with
\[
  \log_{10}\!\left(W(62)M_{30}/2^{124}\right)
  =
  -0.00654409170066116985<0.
\]

Suppose the width fixed-witness fiber has
\[
  S=W_q(T)=\{s_1<\cdots<s_q\},
\]
the anchored support $A$ is a nonempty interval of $S$, every internal edge of
the anchored path joins two labels $s_r,s_{r'}$ with $|r-r'|=1$, and
\[
  |S\setminus A|\le F(q).
\]
Then there is a $2q$-bit code
\[
  \Phi(v,A,\Gamma_{S\setminus A})\in\{0,1\}^{2q},
\]
depending only on the retained deficit vertex $v$, the interval $A$, and the
deleted-only graph induced on $S\setminus A$, such that $\Gamma$ is recovered
from
\[
  S,\quad \Gamma^\ast,\quad
  \Phi(v,A,\Gamma_{S\setminus A})\in\{0,1\}^{2q}.
\]
The same assertion holds in the height fixed-witness fiber with
$S=H_q(T)=\{t_1<\cdots<t_q\}$.

Consequently, after the adjacent-rank crossing condition has been proved, the
unabsorbed proper-support deleted-only readout is restricted to
\[
  |S\setminus A|\ge F(q)+1 .
\]
\end{corollary}

\begin{proof}
The residual window bound is \cref{cor:post29-bc-window-size-strata}; for
$q\ge29$ the final line is the sharper bound
$j-q\le30899-q$ from \cref{cor:post29-bc-direct-onset-bound}.  Hence the
retained deficit vertex has at most $W(q)$ possible ranks.

For fixed $S$, the number of triples
\[
  (v,A,\Gamma_{S\setminus A})
  \qquad\text{with}\qquad |S\setminus A|\le F(q)
\]
is at most $W(q)M_{F(q)}$, since for each complement size $d$ there are
$d+1$ interval supports $A$ and $s_d$ loopless $2$-regular deleted-only graphs
on $S\setminus A$.  By the defining inequality for $F(q)$, this number is
strictly less than $2^{2q}$.  Listing those triples lexicographically gives
the code $\Phi$.

The code recovers $v$, $A$, and the deleted-only graph.  Applying
\cref{cor:bc-rank-adjacent-interval-code-recovery} with these recovered inputs
recovers $\Gamma$ in the width case; the height case is identical.  The final
sentence is the complement of the exact packed range.
\end{proof}

\begin{corollary}[Combined residual complement threshold]\label{cor:bc-rank-adjacent-combined-complement-joint-code}
Work in one residual row and moment of \cref{rem:former-post29-tail}.  Put
$q=b+1$, let $S\subset\{1,\ldots,j\}$ have cardinality $q$, and use the
notation of \cref{cor:bc-rank-adjacent-interval-code-recovery}.  Define
\[
K_0(q)=
\begin{cases}
9, & q=15,\\
10, & 16\le q\le17,\\
11, & q=18,\\
12, & 19\le q\le20,\\
13, & 21\le q\le22,\\
14, & 23\le q\le24,\\
15, & q\ge25,
\end{cases}
\]
\[
L_0(q)=
\begin{cases}
\left\lfloor q/4\right\rfloor+10, & q\ge45,\\
0, & 15\le q\le44,
\end{cases}
\qquad
H_0(q)=\left\lceil q/3\right\rceil,
\]
\[
P_0(q)=
\begin{cases}
\left\lceil q/3\right\rceil+8, & q\ge80,\\
0, & 15\le q\le79,
\end{cases}
\]
\[
Q_0(q)=
\begin{cases}
\left\lceil q/3\right\rceil+10, & q\ge86,\\
0, & 15\le q\le85,
\end{cases}
\]
let $F_0(q)=F(q)$ be the exact complement frontier of
\cref{def:bc-exact-complement-frontier}, and
\[
  R(q)=\max\{K_0(q),L_0(q),H_0(q),P_0(q),Q_0(q),F_0(q)\}.
\]
Suppose the width fixed-witness fiber has
\[
  S=W_q(T)=\{s_1<\cdots<s_q\},
\]
the anchored support $A$ is a nonempty interval of $S$, every internal edge of
the anchored path joins two labels $s_r,s_{r'}$ with $|r-r'|=1$, and
\[
  |S\setminus A|\le R(q).
\]
Then there is a $2q$-bit code
\[
  \Omega(v,A,\Gamma_{S\setminus A})\in\{0,1\}^{2q},
\]
depending only on the retained deficit vertex $v$, the interval $A$, and the
deleted-only graph induced on $S\setminus A$, such that $\Gamma$ is recovered
from
\[
  S,\quad \Gamma^\ast,\quad
  \Omega(v,A,\Gamma_{S\setminus A})\in\{0,1\}^{2q}.
\]
The same assertion holds in the height fixed-witness fiber with
$S=H_q(T)=\{t_1<\cdots<t_q\}$.

Consequently, after the adjacent-rank crossing condition has been proved, the
unabsorbed proper-support deleted-only readout is restricted to
\[
  |S\setminus A|\ge R(q)+1 .
\]
\end{corollary}

\begin{proof}
Choose one of $K_0(q)$, $L_0(q)$, $H_0(q)$, $P_0(q)$, $Q_0(q)$, or
$F_0(q)$ attaining the maximum $R(q)$.
If $K_0(q)$ attains the maximum, the assertion is exactly
\cref{cor:bc-rank-adjacent-stratified-complement-joint-code}.  If $H_0(q)$
attains the maximum, it is exactly
\cref{cor:bc-rank-adjacent-one-third-complement-joint-code}.  If $L_0(q)$
attains the maximum, then $q\ge45$ and the assertion is exactly
\cref{cor:bc-rank-adjacent-linear-complement-joint-code}.  If $P_0(q)$
attains the maximum, then $q\ge80$ and the assertion is exactly
\cref{cor:bc-rank-adjacent-onethird-plus-eight-complement-joint-code}.  If
$Q_0(q)$ attains the maximum, then $q\ge86$ and the assertion is exactly
\cref{cor:bc-rank-adjacent-exact-onethird-plus-ten-complement-joint-code}.
If $F_0(q)$ attains the maximum, the assertion is exactly
\cref{cor:bc-rank-adjacent-exact-frontier-complement-joint-code}.
The final sentence is the complement of the combined packing range.
\end{proof}

\begin{lemma}[Witness deletion is consecutive standard-time deletion]\label{lem:bc-standard-time-pair-deletion}
Let $T$ be a semistandard tableau of content $(2^j)$, let
$U=\operatorname{std}(T)$, and let
\[
  \varnothing=\sigma^{(0)}\subset\sigma^{(1)}\subset\cdots\subset
  \sigma^{(2j)}=\lambda(T)
\]
be the standard shape chain of $U$, where $\sigma^{(r)}$ is the shape formed by
the entries at most $r$.  Let
\[
  \varnothing=\lambda^{(0)}\subset\lambda^{(1)}\subset\cdots\subset
  \lambda^{(j)}=\lambda(T)
\]
be the Pieri two-strip path of \cref{lem:bc-pieri-path-bijection}.  Then
\[
  \sigma^{(2a)}=\lambda^{(a)}
  \quad\text{for }0\le a\le j,
\]
and each pair
\[
  \sigma^{(2a-2)}\subset\sigma^{(2a-1)}\subset\sigma^{(2a)}
\]
adds exactly the two boxes of semistandard label $a$.

Consequently, for $S\subset\{1,\ldots,j\}$, deleting all semistandard labels in
$S$ is the same box deletion as deleting the consecutive standard-time pairs
\[
  D(S)=\bigcup_{s\in S}\{2s-1,2s\}
\]
from $U$.  The set $D(S)$ is a union of at most $|S|$ maximal consecutive
intervals of standard times, and these intervals are exactly the doubled
standard-time form of the deletion joints in \cref{def:bc-raw-block-deletion}.
\end{lemma}

\begin{proof}
By \cref{lem:bc-doubled-standardization}, the boxes carrying semistandard label
$a$ in $T$ are precisely the boxes carrying standard entries $2a-1$ and $2a$
in $U$.  Therefore the shape formed by entries at most $2a$ in $U$ is exactly
the shape formed by semistandard labels at most $a$ in $T$, which is
$\lambda^{(a)}$ by \cref{lem:bc-pieri-path-bijection}.  This proves
$\sigma^{(2a)}=\lambda^{(a)}$ and identifies the two single-box additions
between times $2a-2$ and $2a$ with the two boxes of label $a$.

The deletion statement follows from the same identification of boxes.  Since
each deleted semistandard label contributes the adjacent standard pair
$\{2s-1,2s\}$, $D(S)$ is a union of at most $|S|$ maximal consecutive intervals.
Under the label-to-block conversion used in \cref{def:bc-raw-block-deletion},
maximal consecutive deleted labels are the same deleted objects as maximal
consecutive deleted two-step blocks, with the order reversed by
\cref{lem:bc-growth-path-model}.  Reversing order preserves the number and
endpoint nature of the intervals, so these are exactly the same deletion
joints in standard-time coordinates.
\end{proof}

\begin{lemma}[Low-label boxes determine the standard even readout]\label{lem:bc-low-label-boxes-standard-readout}
Let $T$ be a semistandard tableau of content $(2^j)$, let
$U=\operatorname{std}(T)$, and let
\[
  \varnothing=\sigma^{(0)}\subset\sigma^{(1)}\subset\cdots\subset
  \sigma^{(2j)}=\lambda(T)
\]
be the standard shape chain of $U$.  Fix $M\le j$.  Any datum which determines,
for each $1\le a\le M$, the two boxes of $T$ carrying semistandard label $a$,
determines the even standard-time endpoints
\[
  \sigma^{(0)},\sigma^{(2)},\sigma^{(4)},\ldots,\sigma^{(2M)}.
\]
The readout is
\[
  \sigma^{(2a)}
  =
  \bigcup_{1\le u\le a}
    \{\text{boxes of }T\text{ carrying label }u\}
  \qquad(0\le a\le M),
\]
with the empty union for $a=0$.
\end{lemma}

\begin{proof}
By \cref{lem:bc-doubled-standardization}, the boxes carrying semistandard label
$u$ in $T$ are exactly the boxes carrying the adjacent standard pair
$\{2u-1,2u\}$ in $U$.  Therefore the boxes of $U$ with entries at most $2a$
are precisely the union of the boxes of $T$ carrying labels $1,\ldots,a$.
Since $\sigma^{(2a)}$ is the shape formed by entries at most $2a$ in $U$, the
displayed formula follows for every $a\le M$.
\end{proof}

\begin{lemma}[Step-2 Knuth prefixes determine low-label boxes]\label{lem:bc-step2-knuth-prefix-low-label-boxes}
Use the Step~2 first arbitrary-filling Knuth growth diagram in the proof of
\cite[Theorem 4]{Krattenthaler2016} for a semistandard tableau $T$ of content
$(2^j)$.  Let
\[
  \varnothing=\lambda^{(0)}\subset\lambda^{(1)}\subset\cdots\subset
  \lambda^{(j)}
\]
be the Pieri two-strip path of $T$.  Fix $M\le j$.  Any datum which determines
the Step~2 boundary prefix
\[
  \lambda^{(0)},\lambda^{(1)},\ldots,\lambda^{(M)}
\]
determines, for each $1\le a\le M$, the two boxes of $T$ carrying
semistandard label $a$.

Moreover, the same conclusion holds if the datum determines a Ferrers
subarrangement $C$ of the Step~2 first-variant growth diagram, the cell
entries of $C$, and the left/bottom boundary labels of $C$, and the forward
boundary sweep on $C$ contains the displayed prefix.
\end{lemma}

\begin{proof}
By \cref{lem:bc-pieri-path-bijection}, the two boxes carrying label $a$ are
exactly the skew difference
\[
  \lambda^{(a)}\setminus\lambda^{(a-1)}.
\]
Thus the displayed boundary prefix determines those two-box sets for every
$a\le M$.

For the final assertion, apply the forward part of
\cref{lem:bc-knuth-growth-local-reversibility} to the known cell entries and
known left/bottom boundary labels of $C$.  Sweeping across $C$ determines its
top/right boundary labels.  If that boundary contains the displayed prefix,
the first paragraph recovers the low-label boxes.
\end{proof}

\begin{lemma}[Step-2 Knuth intervals determine label boxes]\label{lem:bc-step2-knuth-interval-label-boxes}
Use the Step~2 first arbitrary-filling Knuth growth diagram in the proof of
\cite[Theorem 4]{Krattenthaler2016} for a semistandard tableau $T$ of content
$(2^j)$.  Let
\[
  \varnothing=\lambda^{(0)}\subset\lambda^{(1)}\subset\cdots\subset
  \lambda^{(j)}
\]
be the Pieri two-strip path of $T$.  Fix $1\le a\le b\le j$.  Any datum which
determines the Step~2 boundary interval
\[
  \lambda^{(a-1)},\lambda^{(a)},\ldots,\lambda^{(b)}
\]
determines, for every $a\le s\le b$, the two boxes of $T$ carrying
semistandard label $s$.

Moreover, the same conclusion holds if the datum determines a Ferrers
subarrangement $C$ of the Step~2 first-variant growth diagram, the cell entries
of $C$, and the left/bottom boundary labels of $C$, and the forward boundary
sweep on $C$ contains the displayed boundary interval.
\end{lemma}

\begin{proof}
By \cref{lem:bc-pieri-path-bijection}, the two boxes carrying label $s$ are
exactly the skew difference
\[
  \lambda^{(s)}\setminus\lambda^{(s-1)} .
\]
The displayed boundary interval contains both endpoints of this difference for
every $a\le s\le b$, so it determines all listed two-box label sets.

For the final assertion, apply the forward part of
\cref{lem:bc-knuth-growth-local-reversibility} to the known cell entries and
known left/bottom boundary labels of $C$.  Sweeping across $C$ determines its
top/right boundary labels.  If that boundary contains the displayed interval,
the first paragraph recovers the label boxes.
\end{proof}

\begin{lemma}[Step-2 Knuth continuations determine continuation boxes]\label{lem:bc-step2-knuth-continuation-boxes}
Use the Step~2 first arbitrary-filling Knuth growth diagram in the proof of
\cite[Theorem 4]{Krattenthaler2016} for a semistandard tableau $T$ of content
$(2^j)$.  Let
\[
  \varnothing=\lambda^{(0)}\subset\lambda^{(1)}\subset\cdots\subset
  \lambda^{(j)}
\]
be the Pieri two-strip path of $T$.  Fix $M\le j$.  Any datum which determines
the Step~2 boundary continuation
\[
  \lambda^{(M)},\lambda^{(M+1)},\ldots,\lambda^{(j)}
\]
determines, for each $M<t\le j$, the two boxes of $T$ carrying semistandard
label $t$.

Moreover, the same conclusion holds if the datum determines a Ferrers
subarrangement $C$ of the Step~2 first-variant growth diagram, the cell
entries of $C$, and the left/bottom boundary labels of $C$, and the forward
boundary sweep on $C$ contains the displayed continuation segment.
\end{lemma}

\begin{proof}
By \cref{lem:bc-pieri-path-bijection}, the two boxes carrying label $t$ are
exactly the skew difference
\[
  \lambda^{(t)}\setminus\lambda^{(t-1)}.
\]
Thus the displayed boundary continuation determines those two-box sets for
every $M<t\le j$.

For the final assertion, apply the forward part of
\cref{lem:bc-knuth-growth-local-reversibility} to the known cell entries and
known left/bottom boundary labels of $C$.  Sweeping across $C$ determines its
top/right boundary labels.  If that boundary contains the displayed
continuation segment, the first paragraph recovers the labelled continuation
boxes.
\end{proof}

\begin{lemma}[Step-2 Knuth matrix is unique from the labelled tableau]\label{lem:bc-step2-matrix-unique-from-tableau}
Use the Step~2 first arbitrary-filling Knuth growth diagram in the proof of
\cite[Theorem 4]{Krattenthaler2016}.  For a fixed semistandard tableau $T$ in
that construction, the Step~2 symmetric nonnegative integer matrix is uniquely
determined by $T$.
\end{lemma}

\begin{proof}
In \cite[Theorem 4, Step 2]{Krattenthaler2016}, the first arbitrary-filling
Knuth construction of \cite[\S4.1, Theorem 7]{Krattenthaler2006} is applied in
the inverse direction to a pair of equal semistandard boundary tableaux.  The
equality of the two tableaux is exactly the symmetry statement recorded in the
source for the resulting matrix.  Hence the labelled tableau $T$ determines
both the $P$- and $Q$-side top/right boundary paths of the full Step~2
rectangular growth diagram.

By \cref{lem:bc-knuth-growth-local-reversibility}, the backward local
algorithm for this first arbitrary-filling variant determines the southwest
corner and the cell entry of each cell from the other three corner labels.
Starting from the full top/right boundary determined by $T$ and sweeping
backward through the whole rectangle therefore recovers every Step~2 cell
entry.  This recovered filling is the unique Step~2 matrix associated to
$T$.
\end{proof}

\begin{lemma}[Step-2 Q-prefix strips determine boundary prefixes]\label{lem:bc-step2-q-prefix-strip-determines-prefix}
Use the Step~2 first arbitrary-filling Knuth growth diagram in the proof of
\cite[Theorem 4]{Krattenthaler2016}, and orient it as the rectangular
first-variant RSK diagram of \cite[\S4.1, Theorem 7]{Krattenthaler2006}.
Let
\[
  \varnothing=\lambda^{(0)}\subset\lambda^{(1)}\subset\cdots\subset
  \lambda^{(j)}
\]
be the $Q$-side semistandard boundary path of this diagram; equivalently,
$\lambda^{(a)}/\lambda^{(a-1)}$ is the set of boxes of the recording tableau
carrying label $a$.  For $M\le j$, let $C_M$ be the Step~2 rectangular
subarrangement consisting of the first $M$ cells in the $Q$-coordinate direction
and all cells in the opposite coordinate direction.  If the entries of $C_M$
are known, then
\[
  \lambda^{(0)},\lambda^{(1)},\ldots,\lambda^{(M)}
\]
are determined.
\end{lemma}

\begin{proof}
In the rectangular case, \cite[\S4.1, Theorem 7]{Krattenthaler2006} identifies
the first arbitrary-filling growth diagram with RSK.  The first half of the
top/right boundary sequence is the recording tableau $Q$: the boxes added in
the difference $\lambda^{(a)}/\lambda^{(a-1)}$ are precisely the boxes carrying
the label $a$ in $Q$.  The subarrangement $C_M$ is anchored on the global empty
left/bottom boundary, so those boundary labels are all $\varnothing$.  Applying
the forward local rules of \cref{lem:bc-knuth-growth-local-reversibility} to
the known entries of $C_M$ therefore determines the corresponding top/right
boundary labels, and these include exactly
$\lambda^{(0)},\ldots,\lambda^{(M)}$.
\end{proof}

\begin{lemma}[Step-2 Q-suffix strips determine boundary continuations]\label{lem:bc-step2-q-suffix-strip-determines-continuation}
Use the Step~2 first arbitrary-filling Knuth growth diagram and notation of
\cref{lem:bc-step2-q-prefix-strip-determines-prefix}.  For $M\le j$, let
$C_M^+$ be the rectangular subarrangement consisting of the cells whose
$Q$-coordinate label lies in
\[
  \{M+1,\ldots,j\}
\]
and whose opposite coordinate is arbitrary.  Suppose a datum determines the cell
entries of $C_M^+$ and the partition labels on the left/bottom boundary of
$C_M^+$.  Then the same datum determines the $Q$-side continuation boundary
\[
  \lambda^{(M)},\lambda^{(M+1)},\ldots,\lambda^{(j)} .
\]
\end{lemma}

\begin{proof}
The subarrangement $C_M^+$ is a finite rectangular Ferrers cell arrangement in
the first arbitrary-filling Knuth growth diagram.  Applying the forward local
rules of \cref{lem:bc-knuth-growth-local-reversibility} to the known entries and
known left/bottom boundary labels determines the top/right boundary labels of
$C_M^+$.  In the orientation of
\cref{lem:bc-step2-q-prefix-strip-determines-prefix}, the part of this boundary
in the $Q$-coordinate direction is exactly the segment of the recording
boundary after the first $M$ labels, namely
\[
  \lambda^{(M)},\lambda^{(M+1)},\ldots,\lambda^{(j)} .
\]
\end{proof}

\begin{lemma}[Only the internal Q-cut of a suffix strip is non-global]\label{lem:bc-q-suffix-boundary-internal-cut}
Use the Step~2 first arbitrary-filling Knuth growth diagram and notation of
\cref{lem:bc-step2-q-suffix-strip-determines-continuation}.  The left/bottom
boundary of the suffix strip $C_M^+$ is the union of:
\begin{enumerate}
\item the part lying on the global empty left/bottom boundary of the full
Step~2 rectangle; and
\item the internal separating boundary between the prefix strip $C_M$ and the
suffix strip $C_M^+$.
\end{enumerate}
The labels on the first part are all $\varnothing$.  Consequently any datum
which determines the ordered partition labels on the internal separating
boundary determines the full left/bottom boundary of $C_M^+$.
\end{lemma}

\begin{proof}
The strips $C_M$ and $C_M^+$ are obtained by cutting the full Step~2 rectangle
between the $Q$-coordinate labels $M$ and $M+1$.  The suffix strip still
contains every cell in the opposite coordinate direction.  Hence the incoming
boundary for a forward sweep on $C_M^+$ has only two kinds of edges: those
which already lie on the global incoming boundary of the full rectangle, and
those created by the cut from the removed prefix strip.  The first arbitrary
filling RSK rectangle is anchored on the global empty left/bottom boundary, so
the labels of the first kind are all $\varnothing$.  The labels of the second
kind are exactly the ordered internal cut labels.  Knowing those ordered cut
labels therefore supplies every non-fixed label on the left/bottom boundary of
$C_M^+$.
\end{proof}

\begin{lemma}[Q-prefix entries determine the internal Q-cut]\label{lem:bc-q-prefix-entries-determine-internal-q-cut}
Use the Step~2 first arbitrary-filling Knuth growth diagram and notation of
\cref{lem:bc-step2-q-prefix-strip-determines-prefix,lem:bc-q-suffix-boundary-internal-cut}.
If a datum determines all cell entries of the $Q$-prefix strip $C_M$, then it
determines the ordered partition labels on the internal separating boundary
between $C_M$ and $C_M^+$.
\end{lemma}

\begin{proof}
The prefix strip $C_M$ is anchored on the global empty left/bottom boundary of
the full Step~2 rectangle.  Applying the forward local rules of
\cref{lem:bc-knuth-growth-local-reversibility} to the known entries of $C_M$
therefore determines every boundary label produced by the forward sweep on
$C_M$.  The boundary component of $C_M$ created by cutting the full rectangle
between the $Q$-coordinate labels $M$ and $M+1$ is exactly the internal
separating boundary between $C_M$ and $C_M^+$, so its ordered labels are
determined.
\end{proof}

\begin{lemma}[Low-label incidence cells cover the Step-2 Q-prefix strip]\label{lem:bc-step3-low-label-incidence-covers-q-prefix}
Use the notation of \cref{lem:bc-step2-q-prefix-strip-determines-prefix} and
\cref{lem:bc-step3-cells-are-step2-half-matrix}.  For $M\le j$, let
$H_M$ be the Step~3 half-square cell set consisting of the cells corresponding
to unordered off-diagonal Step~2 label pairs
\[
  \{a,b\}\subset\{1,\ldots,j\},
  \qquad a\ne b,
  \qquad \min(a,b)\le M .
\]
Then every non-diagonal cell of the Step~2 $Q$-prefix strip $C_M$ has its
cell or its transpose represented in $H_M$.  Consequently, recovering the
Step~3 subfilling on any carrier containing $H_M$ determines all non-diagonal
entries of $C_M$.
\end{lemma}

\begin{proof}
The strip $C_M$ consists of the Step~2 cells whose $Q$-coordinate label lies in
$\{1,\ldots,M\}$.  A non-diagonal cell of this strip therefore has ordered
label coordinates $(a,b)$ with $a\le M$ and $a\ne b$.  The Step~2 matrix is
symmetric, and Step~3 keeps exactly one off-diagonal representative of the
unordered pair $\{a,b\}$.  This representative belongs to $H_M$ because
$\min(a,b)\le a\le M$.  The last sentence is then
\cref{lem:bc-step3-cells-are-step2-half-matrix}.
\end{proof}

\begin{corollary}[Low-label incidence entries determine the internal Q-cut]\label{cor:bc-low-label-incidence-determines-internal-q-cut}
Use the Step~2 and Step~3 notation of
\cref{lem:bc-step3-low-label-incidence-covers-q-prefix}.  If a datum determines
the Step~3 cell entries on every cell of $H_M$, then it determines the ordered
partition labels on the internal separating boundary between the Step~2
$Q$-prefix strip $C_M$ and the suffix strip $C_M^+$.
\end{corollary}

\begin{proof}
By \cref{lem:bc-step3-low-label-incidence-covers-q-prefix}, the recovered
entries on $H_M$ determine every non-diagonal entry of the Step~2 $Q$-prefix
strip $C_M$.  The diagonal entries of the Step~2 matrix are zero by
\cref{lem:bc-step3-cells-are-step2-half-matrix}.  Hence all cell entries of
$C_M$ are determined.  Applying
\cref{lem:bc-q-prefix-entries-determine-internal-q-cut} gives the ordered
internal cut labels.
\end{proof}

\begin{lemma}[Continuation incidence cells cover the Step-2 Q-suffix strip]\label{lem:bc-step3-continuation-incidence-covers-q-suffix}
Use the notation of
\cref{lem:bc-step2-q-suffix-strip-determines-continuation} and
\cref{lem:bc-step3-cells-are-step2-half-matrix}.  For $M\le j$, let $J_M$ be
the Step~3 half-square cell set corresponding to unordered off-diagonal Step~2
label pairs
\[
  \{a,b\}\subset\{1,\ldots,j\},
  \qquad a\ne b,\qquad \max(a,b)>M .
\]
Then every non-diagonal cell of the Step~2 $Q$-suffix strip $C_M^+$ has its cell
or its transpose represented in $J_M$.  Consequently, recovering the Step~3
subfilling on any carrier containing $J_M$ determines all non-diagonal entries
of $C_M^+$; the diagonal entries of $C_M^+$ are zero.

Moreover, in the Step~3 label order $j,j-1,\ldots,1$, the set $J_M$ is exactly
the first $j-M$ strict-left columns, namely the columns incident to the
continuation labels $M+1,\ldots,j$.
\end{lemma}

\begin{proof}
A non-diagonal cell of the Step~2 $Q$-suffix strip has ordered label coordinates
$(a,b)$ with $a>M$ and $a\ne b$.  Since the Step~2 matrix is symmetric and
Step~3 retains one representative of each unordered off-diagonal pair, that
cell or its transpose is represented by the Step~3 cell for the unordered pair
\[
  \{a,b\}.
\]
This pair belongs to $J_M$ because $\max(a,b)\ge a>M$.  The entries agree by
\cref{lem:bc-step3-cells-are-step2-half-matrix}, and the diagonal entries of the
Step~2 matrix are zero by the same construction.

In the Step~3 half-square the labels are ordered as $j,j-1,\ldots,1$.  A cell
whose unordered pair has maximum label greater than $M$ is incident to one of
the labels $M+1,\ldots,j$, and in the retained strict half-square it lies in the
column of such a continuation label.  These are precisely the first $j-M$
strict-left columns.
\end{proof}

\begin{lemma}[Low-label incidence splits into cross and low-low cells]\label{lem:bc-low-label-incidence-split}
Use the Step~3 half-square labelling of
\cref{lem:bc-step3-low-label-incidence-covers-q-prefix}.  For $M\le j$, let
$X_M$ be the cell set corresponding to unordered pairs
\[
  \{a,b\},\qquad \min(a,b)\le M<\max(a,b),
\]
and let $L_M$ be the cell set corresponding to unordered pairs
\[
  \{a,b\},\qquad 1\le a<b\le M .
\]
Then
\[
  H_M=X_M\sqcup L_M .
\]
Moreover, if the Step~3 half-square is ordered by labels
$j,j-1,\ldots,1$ as in \cite[Theorem 4]{Krattenthaler2016}, then $X_M$ lies in
the first $j-M$ strict-left columns, namely the columns incident to the
continuation labels $M+1,\ldots,j$.
\end{lemma}

\begin{proof}
The definition of $H_M$ says that an unordered off-diagonal pair belongs to it
exactly when at least one of its two labels is at most $M$.  Such a pair has
either both labels at most $M$, which is precisely $L_M$, or exactly one label
at most $M$, which is precisely $X_M$.  These alternatives are disjoint and
exhaust $H_M$.

In Krattenthaler's Step~3 half-square the labels are ordered
$j,j-1,\ldots,1$.  A cell of $X_M$ has one incident label in
$M+1,\ldots,j$; the half-square representative of that unordered pair appears
in the strict-left column indexed by that continuation label.  Hence every cell
of $X_M$ lies in the first $j-M$ strict-left columns.
\end{proof}

\begin{lemma}[Low-low cells are a terminal dual-RSK' corner]\label{lem:bc-low-low-terminal-corner-boundary}
Use the Step~3 dual-RSK' half-square of
\cref{lem:bc-krattenthaler-semistandard-path-boundary} and the notation of
\cref{lem:bc-low-label-incidence-split}.  For $M\le j$, the cell set $L_M$ is
the terminal triangular Ferrers subarrangement on the labels
$M,M-1,\ldots,1$ in the Step~3 label order $j,j-1,\ldots,1$.  Its top/right
boundary in the full Step~3 diagram is the terminal segment
\[
  \rho_{2(j-M)},\rho_{2(j-M)+1},\ldots,\rho_{2j}
\]
of the semistandard growth path.  Consequently, any data determining this
boundary segment determine every Step~3 entry on $L_M$ and every corner label
on or inside $L_M$.
\end{lemma}

\begin{proof}
The Step~3 half-square is indexed by the ordered labels
$j,j-1,\ldots,1$; its cells are the off-diagonal unordered pairs of distinct
labels, represented in one triangular half.  The subset with both incident
labels at most $M$ is therefore the principal triangular corner cut out by the
terminal ordered label list $M,M-1,\ldots,1$.  This is exactly the cell set
$L_M$ from \cref{lem:bc-low-label-incidence-split}.

By \cref{lem:bc-krattenthaler-semistandard-path-boundary}, the sequence
$\rho_0,\rho_1,\ldots,\rho_{2j}$ is read along the top/right boundary of this
Step~3 diagram.  The labels $j,j-1,\ldots,M+1$ occupy the first $j-M$ label
positions, each contributing two boundary steps in the content $(2^j)$
specialization of \cref{lem:bc-growth-path-model}.  Hence the boundary of the
remaining terminal $M$-label corner begins at $\rho_{2(j-M)}$ and ends at
$\rho_{2j}$, giving the displayed segment.  Applying
\cref{cor:bc-parsed-dual-rsk-prime-boundary-recovers-subfilling} to this
terminal Ferrers subarrangement recovers its cell entries and internal corner
labels from that top/right boundary.
\end{proof}

\begin{corollary}[The terminal low-low carrier is exactly the terminal-boundary readout]\label{cor:bc-terminal-low-low-carrier-equivalent-boundary}
Use the notation of \cref{lem:bc-low-low-terminal-corner-boundary}.  Let
$\mathfrak A$ be any datum.  The datum $\mathfrak A$ determines the terminal
boundary segment
\[
  \rho_{2(j-M)},\rho_{2(j-M)+1},\ldots,\rho_{2j}
\]
if and only if $\mathfrak A$ identifies the terminal Ferrers subarrangement
$L_M$ and determines the partition labels on its top/right boundary.  Thus the
direct low-low carrier route with $E=L_M$ is precisely the terminal-boundary
readout.  A larger parsed carrier can be a source only by recovering this
terminal boundary after the backward sweep.
\end{corollary}

\begin{proof}
\Cref{lem:bc-low-low-terminal-corner-boundary} identifies $L_M$ as the terminal
triangular Ferrers subarrangement and identifies its top/right boundary with the
displayed boundary segment.  Therefore a datum determines that boundary segment
exactly when it determines the top/right boundary labels of this particular
Ferrers subarrangement.  The final sentence is this equivalence applied to the
carrier choice $E=L_M$.
\end{proof}

\begin{corollary}[Any parsed carrier containing the low-low corner recovers the terminal boundary]\label{cor:bc-low-low-carrier-forces-terminal-boundary}
Use the notation of \cref{lem:bc-low-low-terminal-corner-boundary}.  Let
\(E\) be a Ferrers subarrangement of the Step~3 dual-RSK' half-square which
contains the terminal low-low subarrangement \(L_M\).  If a datum determines
\(E\) and the partition labels on the top/right boundary of \(E\), then it
determines the terminal boundary segment
\[
  \rho_{2(j-M)},\rho_{2(j-M)+1},\ldots,\rho_{2j}.
\]
Consequently every direct carrier source for \(L_M\) is, after the dual-RSK'
backward sweep, a terminal-boundary source.
\end{corollary}

\begin{proof}
Applying
\cref{cor:bc-parsed-dual-rsk-prime-boundary-recovers-subfilling} to the
determined carrier \(E\) recovers every corner label on or inside \(E\).  Since
\(L_M\subseteq E\), the top/right boundary labels of the subarrangement
\(L_M\) are among these recovered corner labels.  By
\cref{lem:bc-low-low-terminal-corner-boundary}, that top/right boundary is
exactly the displayed terminal segment of the Step~3 semistandard path.
\end{proof}

\begin{lemma}[Consecutive Step--3 label intervals have boundary segments]\label{lem:bc-step3-consecutive-label-interval-boundary}
Use the Step--3 dual-RSK' half-square of
\cref{lem:bc-krattenthaler-semistandard-path-boundary}.  Let
\[
  1\le a\le b\le j
\]
and let $E_{a,b}$ be the Ferrers subarrangement on the consecutive Step--3
labels
\[
  b,b-1,\ldots,a
\]
inside the global label order $j,j-1,\ldots,1$.  Then the top/right boundary
of $E_{a,b}$ in the full Step--3 diagram is the boundary segment
\[
  \rho_{2(j-b)},\rho_{2(j-b)+1},\ldots,\rho_{2(j-a+1)}.
\]
Consequently, if a parser determines the partition labels on this boundary,
then it determines every Step--3 cell entry and every corner label on or
inside $E_{a,b}$.
\end{lemma}

\begin{proof}
In the Step--3 half-square, the boundary labels are ordered by
$j,j-1,\ldots,1$, and in the content $(2^j)$ specialization each label
contributes two boundary steps to the path
\[
  \rho_0,\rho_1,\ldots,\rho_{2j}.
\]
The labels preceding the interval $b,b-1,\ldots,a$ are
$j,j-1,\ldots,b+1$, of which there are $j-b$.  Hence the interval begins at
the boundary label $\rho_{2(j-b)}$.  The interval itself has
$b-a+1$ labels, so its boundary segment ends after
$2(b-a+1)$ additional steps, namely at
\[
  \rho_{2(j-b)+2(b-a+1)}
  =
  \rho_{2(j-a+1)}.
\]
This proves the displayed segment.

The final statement is exactly
\cref{cor:bc-parsed-dual-rsk-prime-boundary-recovers-subfilling} applied to
the consecutive-label Ferrers subarrangement $E_{a,b}$.
\end{proof}

\begin{corollary}[Consecutive Step--3 carrier data recover interval boundaries]\label{cor:bc-step3-consecutive-carrier-data-recovers-boundary}
In the setting of
\cref{lem:bc-step3-consecutive-label-interval-boundary}, suppose a parser
determines the cell entries of $E_{a,b}$ and the partition labels on the
left/bottom boundary of $E_{a,b}$.  Then the same data determine the boundary
segment
\[
  \rho_{2(j-b)},\rho_{2(j-b)+1},\ldots,\rho_{2(j-a+1)}.
\]
\end{corollary}

\begin{proof}
Apply the forward determination in
\cref{lem:bc-dual-rsk-prime-growth-local-reversibility} to the finite
subarrangement $E_{a,b}$.  Its cell entries together with its left/bottom
boundary labels determine its top/right boundary labels.  By
\cref{lem:bc-step3-consecutive-label-interval-boundary}, that top/right
boundary is exactly the displayed segment.
\end{proof}

\begin{lemma}[Deleted intervals carry ordered frontier events]\label{lem:bc-frontier-events-per-joint}
Let $T$ be a semistandard tableau of content $(2^j)$, let
$U=\operatorname{std}(T)$, and let
\[
  \varnothing=\sigma^{(0)}\subset\sigma^{(1)}\subset\cdots\subset
  \sigma^{(2j)}=\lambda(T)
\]
be the standard shape chain of $U$.

In the width case assume $\lambda_1(T)\ge2q$, put
$s_r=T(1,2r-1)$ and $S=W_q(T)$.  Write the maximal consecutive label
intervals in $S$ as
\[
  J_\alpha=\{a_\alpha,a_\alpha+1,\ldots,b_\alpha\},
  \qquad 1\le\alpha\le L,
\]
with $b_\alpha+1<a_{\alpha+1}$ for $\alpha<L$.  Then the corresponding
maximal standard-time components of $D(S)$ are
\[
  \widetilde J_\alpha
  =\{2a_\alpha-1,2a_\alpha,\ldots,2b_\alpha\}.
\]
For
\[
  E^w_\alpha=\{\kappa^w_r:s_r\in J_\alpha\},
\]
one has
\[
  E^w_\alpha\subset \widetilde J_\alpha,\qquad
  |E^w_\alpha|=|J_\alpha|,
\]
and the elements of $E^w_\alpha$ occur in the same order as the labels in
$J_\alpha$.  If $\alpha<\beta$, then every element of $E^w_\alpha$ is smaller
than every element of $E^w_\beta$.  Under the raw block deletion
\cref{def:bc-raw-block-deletion}, the same interval $J_\alpha$ becomes the
deleted block interval
\[
  \{j-b_\alpha+1,\ldots,j-a_\alpha+1\}.
\]

The same conclusions hold in the height case, with
$t_r=T(2r-1,1)$, $S=H_q(T)$, and
$E^h_\alpha=\{\kappa^h_r:t_r\in J_\alpha\}$.
\end{lemma}

\begin{proof}
The maximal consecutive intervals in $S$ are disjoint, ordered by label, and
their union is $S$.  Since
\[
  D(S)=\bigcup_{s\in S}\{2s-1,2s\},
\]
the labels in $J_\alpha=\{a_\alpha,\ldots,b_\alpha\}$ contribute exactly
$\{2a_\alpha-1,2a_\alpha,\ldots,2b_\alpha\}$.  For $\alpha<L$, the gap
$b_\alpha+1<a_{\alpha+1}$ leaves at least the missing pair
$\{2b_\alpha+1,2b_\alpha+2\}$ before the next component.  These are therefore
precisely the maximal standard-time components of $D(S)$.

For the width case, \cref{lem:bc-standard-frontier-order} gives
$\kappa^w_r\in\{2s_r-1,2s_r\}\subset D(S)$ and
$\kappa^w_1<\cdots<\kappa^w_q$.  Hence the events whose labels lie in
$J_\alpha$ lie in $\widetilde J_\alpha$, one event for each label, and their
order is the label order.  The separation of the $\widetilde J_\alpha$ gives
the stated ordering between different intervals.  The height case is
identical, using the height half of
\cref{lem:bc-standard-frontier-order}.

Finally, the raw block deletion uses the label-to-block conversion
$s\mapsto j-s+1$ from \cref{def:bc-raw-block-deletion}.  Thus a consecutive
label interval $\{a_\alpha,\ldots,b_\alpha\}$ maps to the consecutive block
interval $\{j-b_\alpha+1,\ldots,j-a_\alpha+1\}$, with order reversed.
\end{proof}

\begin{lemma}[Consecutive witness labels have consecutive frontier ranks]\label{lem:bc-consecutive-witness-ranks}
In the width case of \cref{lem:bc-frontier-events-per-joint}, write
\[
  W_q(T)=\{s_1<s_2<\cdots<s_q\}.
\]
Let $J_\alpha=\{a_\alpha,\ldots,b_\alpha\}$ be a maximal consecutive label
interval in $W_q(T)$, put $h_\alpha=|J_\alpha|$, and let $p_\alpha$ be the
unique index with $s_{p_\alpha}=a_\alpha$.  Then
\[
  s_{p_\alpha+i-1}=a_\alpha+i-1
  \qquad(1\le i\le h_\alpha),
\]
and the ordered frontier-event list at $J_\alpha$ is
\[
  E^w_\alpha=\{\kappa^w_{p_\alpha},
        \kappa^w_{p_\alpha+1},\ldots,
        \kappa^w_{p_\alpha+h_\alpha-1}\}.
\]
Moreover, for $1\le i\le h_\alpha$,
\[
  \sigma^{(\kappa^w_{p_\alpha+i-1}-1)}_1
  <2(p_\alpha+i-1)-1
  \le
  \sigma^{(\kappa^w_{p_\alpha+i-1})}_1 .
\]

In the height case, with
\[
  H_q(T)=\{t_1<t_2<\cdots<t_q\},
\]
the analogous assertion holds: if $t_{p_\alpha}=a_\alpha$, then
$t_{p_\alpha+i-1}=a_\alpha+i-1$, the ordered event list is
\[
  E^h_\alpha=\{\kappa^h_{p_\alpha},
        \kappa^h_{p_\alpha+1},\ldots,
        \kappa^h_{p_\alpha+h_\alpha-1}\},
\]
and
\[
  \ell(\sigma^{(\kappa^h_{p_\alpha+i-1}-1)})
  <2(p_\alpha+i-1)-1
  \le
  \ell(\sigma^{(\kappa^h_{p_\alpha+i-1})})
\]
for $1\le i\le h_\alpha$.
\end{lemma}

\begin{proof}
We prove the width statement; the height statement is identical with
$t_r,\kappa^h_r$ in place of $s_r,\kappa^w_r$.  The labels
$s_1<\cdots<s_q$ list $W_q(T)$ in increasing order.  Since
$J_\alpha=\{a_\alpha,\ldots,b_\alpha\}$ is an interval contained in
$W_q(T)$, the labels of $W_q(T)$ that belong to $J_\alpha$ are exactly
\[
  a_\alpha,a_\alpha+1,\ldots,b_\alpha.
\]
They occur in the same increasing order in the list $s_1,\ldots,s_q$.
Therefore, starting from the unique index $p_\alpha$ with
$s_{p_\alpha}=a_\alpha$, one has
$s_{p_\alpha+i-1}=a_\alpha+i-1$ for $1\le i\le h_\alpha$.

By definition in \cref{lem:bc-frontier-events-per-joint},
$E^w_\alpha=\{\kappa^w_r:s_r\in J_\alpha\}$, so the preceding identification
of the ranks gives the displayed event list.  Finally,
\cref{lem:bc-standard-frontier-time} gives the displayed first-row crossing
inequality for each rank $p_\alpha+i-1$.
\end{proof}

\begin{corollary}[Joint frontier-slot concatenation is the full frontier list]\label{cor:bc-joint-frontier-slots-concatenate-full-kappa}
In the width case of \cref{lem:bc-frontier-events-per-joint}, write the maximal
consecutive label intervals in $S=W_q(T)$ as
\[
  J_\alpha=\{a_\alpha,\ldots,b_\alpha\}\qquad(1\le\alpha\le L),
\]
and write
\[
  E^w_\alpha=\{e^w_{\alpha,1}<\cdots<e^w_{\alpha,h_\alpha}\}
\]
for the ordered frontier-event set at $J_\alpha$.  Then the concatenated list
\[
  e^w_{1,1},\ldots,e^w_{1,h_1},
  e^w_{2,1},\ldots,e^w_{2,h_2},
  \ldots,
  e^w_{L,1},\ldots,e^w_{L,h_L}
\]
is exactly
\[
  \kappa^w_1<\kappa^w_2<\cdots<\kappa^w_q .
\]
Consequently two entries adjacent in this concatenated joint-slot list are
consecutive entries in the full width frontier list.

The same assertion holds in the height case, with $S=H_q(T)$ and the ordered
frontier list $\kappa^h_1<\cdots<\kappa^h_q$.
\end{corollary}

\begin{proof}
The intervals $J_1,\ldots,J_L$ are the maximal consecutive components of $S$,
so they are disjoint and their union is $S$.  By
\cref{lem:bc-frontier-events-per-joint}, each $E^w_\alpha$ consists exactly of
the frontier events attached to labels in $J_\alpha$, has cardinality
$h_\alpha=|J_\alpha|$, and every event in $E^w_\alpha$ is smaller than every
event in $E^w_\beta$ when $\alpha<\beta$.  Therefore the displayed
concatenation is strictly increasing and contains precisely one frontier event
for each label in $S$.

\Cref{lem:bc-standard-frontier-order} says that the full width frontier events
attached to the labels of $S=W_q(T)$ are precisely
$\kappa^w_1<\cdots<\kappa^w_q$ in increasing order.  The concatenated list and
the full frontier list are increasing lists with the same elements, hence they
are equal term by term.  Adjacent entries of the concatenated list are therefore
adjacent entries of the full frontier list.  The height case uses the height
part of \cref{lem:bc-frontier-events-per-joint} and
\cref{lem:bc-standard-frontier-order} in the same way.
\end{proof}

\begin{proposition}[Joint-slot adjacent graph edges close the single-deficit graph leaf]\label{prop:bc-joint-slot-adjacent-single-deficit-graph-leaf}
Work in one residual row and moment of \cref{rem:former-post29-tail}.  Put $q=b+1$
and let $S\subset\{1,\ldots,j\}$ have cardinality $q$.  Work in the setting of
\cref{prop:bc-single-deficit-stable-retained-repair}, and let $A\subset S$ be
the support of the anchored component in
\cref{lem:bc-single-deficit-support-path-code}.

In the width fixed-witness fiber, suppose $S=W_q(T)$ and form the concatenated
joint-slot list
\[
  e^w_{1,1},\ldots,e^w_{1,h_1},
  e^w_{2,1},\ldots,e^w_{2,h_2},
  \ldots,
  e^w_{L,1},\ldots,e^w_{L,h_L}
\]
from the ordered frontier-event sets of the maximal witness-label intervals, as
in \cref{cor:bc-joint-frontier-slots-concatenate-full-kappa}.  Assume that every
internal edge of the anchored path on $A$ and every edge of the deleted-only
graph $\Gamma_{S\setminus A}$ has endpoints whose assigned frontier events are
adjacent entries of this concatenated joint-slot list.  Then $\Gamma$ is
recovered from
\[
  S,\quad \Gamma^\ast,\quad \rho(v)\in\{0,1\}^q,
  \quad \iota(A)\in\{0,1\}^q .
\]

The same assertion holds in the height fixed-witness fiber, with
$S=H_q(T)$ and the height joint-slot list.  Consequently the graph part of the
repaired fixed-witness supplier is reduced to proving this joint-slot adjacency
condition for the actual target-size parser.
\end{proposition}

\begin{proof}
In the width case,
\cref{cor:bc-joint-frontier-slots-concatenate-full-kappa} identifies the
concatenated joint-slot list with the full ordered frontier list
\[
  \kappa^w_1<\cdots<\kappa^w_q .
\]
Thus the displayed joint-slot adjacency hypothesis is exactly the
full-frontier consecutiveness hypothesis for the anchored internal edges and the
deleted-only complement edges in
\cref{prop:bc-actual-frontier-consecutive-single-deficit-interval-code}.  That
proposition recovers $\Gamma$ from the displayed data.

The height case is identical, using the height half of
\cref{cor:bc-joint-frontier-slots-concatenate-full-kappa} and the list
$\kappa^h_1<\cdots<\kappa^h_q$.
\end{proof}

\begin{proposition}[Combined complement packing leaves only large joint-slot adjacency]\label{prop:bc-combined-complement-large-joint-slot-graph-dichotomy}
Work in one residual row and moment of \cref{rem:former-post29-tail}.  Put $q=b+1$,
let $S\subset\{1,\ldots,j\}$ have cardinality $q$, and let
$R(q)$ be the complement threshold of
\cref{cor:bc-rank-adjacent-combined-complement-joint-code}.  Work in the
single-deficit setting of
\cref{prop:bc-single-deficit-stable-retained-repair}, and let
$A\subset S$ be the support of the anchored component in
\cref{lem:bc-single-deficit-support-path-code}.

In the width fixed-witness fiber with
\[
  S=W_q(T)=\{s_1<\cdots<s_q\},
\]
the graph recovery has the following two supplier alternatives.
\begin{enumerate}
\item If $A$ is a nonempty interval of $S$, every internal edge of the
anchored path joins two labels $s_r,s_{r'}$ with $|r-r'|=1$, and
\[
  |S\setminus A|\le R(q),
\]
then $\Gamma$ is recovered from
\[
  S,\quad \Gamma^\ast,\quad
  \Omega(v,A,\Gamma_{S\setminus A})\in\{0,1\}^{2q}
\]
for the combined complement code of
\cref{cor:bc-rank-adjacent-combined-complement-joint-code}.
\item If $|S\setminus A|\ge R(q)+1$ and every internal edge of the anchored
path on $A$ and every edge of the deleted-only graph
$\Gamma_{S\setminus A}$ has endpoints whose assigned frontier events are
adjacent entries of the concatenated joint-slot list of
\cref{cor:bc-joint-frontier-slots-concatenate-full-kappa}, then $\Gamma$ is
recovered from
\[
  S,\quad \Gamma^\ast,\quad \rho(v)\in\{0,1\}^q,
  \quad \iota(A)\in\{0,1\}^q .
\]
\end{enumerate}
The same two alternatives hold in the height fixed-witness fiber with
$S=H_q(T)=\{t_1<\cdots<t_q\}$.

Consequently, after the adjacent-rank crossing condition has been proved and
the combined complement code has been used, the graph-side source-extraction
leaf for the repaired fixed-witness supplier is precisely the actual
target-size parser theorem asserting the joint-slot adjacency in item~(2) for
the complementary range $|S\setminus A|\ge R(q)+1$.
\end{proposition}

\begin{proof}
The first width alternative is exactly
\cref{cor:bc-rank-adjacent-combined-complement-joint-code}.  The second width
alternative is exactly
\cref{prop:bc-joint-slot-adjacent-single-deficit-graph-leaf}.  The height
alternatives are the height halves of the same two results.  The final sentence
is the partition of the proper-support deleted-only complement into
$|S\setminus A|\le R(q)$ and $|S\setminus A|\ge R(q)+1$.
\end{proof}

\begin{lemma}[The all-but-one complement is beyond the combined threshold]\label{lem:bc-all-but-one-complement-large}
For every residual rank parameter $15\le q\le218$, the combined threshold
$R(q)$ of \cref{cor:bc-rank-adjacent-combined-complement-joint-code} satisfies
\[
  R(q)\le q-2 .
\]
Consequently, if $A$ is a one-point subset of a $q$-element witness set $S$,
then
\[
  |S\setminus A|=q-1\ge R(q)+1 .
\]
\end{lemma}

\begin{proof}
Use the definitions of $K_0,L_0,H_0,P_0,Q_0,F_0$ in
\cref{cor:bc-rank-adjacent-combined-complement-joint-code}.  The exact
frontier verifier cited in
\cref{cor:bc-rank-adjacent-exact-frontier-complement-joint-code} reports
\[
  F_0(q)-q\le -6
\]
throughout the residual range, so $F_0(q)\le q-2$.
For
$15\le q\le24$, the displayed values of $K_0$ give
\[
\begin{array}{c|cccccc}
q&15&16,17&18&19,20&21,22&23,24\\
\hline
K_0(q)&9&10&11&12&13&14,
\end{array}
\]
and the functions $L_0,H_0,P_0,Q_0$ are no larger than these values or zero
in this range, while each listed value is at most $q-2$.
For $25\le q\le44$, one has $K_0(q)=15$, $L_0(q)=P_0(q)=Q_0(q)=0$, and
$H_0(q)=\lceil q/3\rceil\le15$, so $R(q)\le15\le q-2$.
For $45\le q\le79$,
\[
  L_0(q)=\left\lfloor q/4\right\rfloor+10\le q-2,
  \qquad
  H_0(q)=\left\lceil q/3\right\rceil\le q-2,
\]
and $K_0(q)=15\le q-2$, while $P_0(q)=Q_0(q)=0$.
For $80\le q\le85$, the only new threshold is
$P_0(q)=\lceil q/3\rceil+8\le q-2$, and the earlier bounds remain valid.
For $86\le q\le218$, the largest possible new threshold is
$Q_0(q)=\lceil q/3\rceil+10\le q-2$, and the earlier bounds remain valid.
Together with the bound for $F_0$, every entry in the maximum defining $R(q)$
is at most $q-2$.
\end{proof}

\begin{proposition}[Large-complement adjacency is not graph-formal]\label{prop:bc-large-complement-adjacency-not-graph-formal}
Let $15\le q\le218$ and let
\[
  S=\{s_1<s_2<\cdots<s_q\}
\]
be a linearly ordered witness set.  There is a loopless labelled
$2$-regular multigraph $\Gamma$ in the single-deficit retained-repair setting
of \cref{prop:bc-single-deficit-stable-retained-repair} with anchored support
\[
  A=\{s_1\},
\]
such that
\[
  |S\setminus A|\ge R(q)+1,
\]
but the deleted-only graph $\Gamma_{S\setminus A}$ has an edge whose endpoints
are not adjacent entries of the full ordered frontier list, equivalently not
adjacent entries of the concatenated joint-slot list of
\cref{cor:bc-joint-frontier-slots-concatenate-full-kappa}.

Consequently the large-complement joint-slot adjacency hypothesis in
\cref{prop:bc-combined-complement-large-joint-slot-graph-dichotomy} is not a
formal consequence of the retained-repair graph data.  A source-extraction
proof of that hypothesis must use the actual target-size parser to exclude
these deleted-only graph choices, or must replace adjacency by an explicit
target-window datum that recovers the deleted-only graph.
\end{proposition}

\begin{proof}
Let $U=\{u_1,u_2\}$ be disjoint from $S\cup\{v\}$ and put a double edge between
$u_1$ and $u_2$.  Put
\[
  D=S\setminus\{s_1\}=\{d_1<d_2<\cdots<d_{q-1}\}.
\]
Since $q\ge15$, the set $D$ has at least four elements.  On $D$, take the
simple cycle
\[
  d_1-d_3-d_2-d_4-d_5-\cdots-d_{q-1}-d_1 .
\]
Let $\Gamma$ be the disjoint union of this cycle, the double edge between
$u_1$ and $u_2$, and the double edge between $v$ and $s_1$.  This graph is
loopless and every vertex has degree $2$.  Deleting $S$ leaves the fixed
retained double edge on $U$ and an isolated retained vertex $v$, so the only
nonzero retained deficit is $c_v=2$; hence this is the single-deficit setting
of \cref{prop:bc-single-deficit-stable-retained-repair}.  The component
containing $v$ has support $A=\{s_1\}$.

By \cref{lem:bc-all-but-one-complement-large},
\[
  |S\setminus A|=q-1\ge R(q)+1 .
\]
The deleted-only graph contains the edge $d_1d_3$.  In the full ordered list
of $S$, the label $d_2$ lies strictly between $d_1$ and $d_3$, so these
endpoints are not adjacent witness ranks.  By
\cref{cor:bc-joint-frontier-slots-concatenate-full-kappa}, the concatenated
joint-slot list is the same full ordered frontier list, so the same edge is not
joint-slot adjacent.  This proves the example.

The final assertion follows because all retained-repair graph data in the
construction are valid before any target-size parser restriction is imposed,
while the displayed deleted-only edge violates joint-slot adjacency.
\end{proof}

\begin{lemma}[Step--2 submatrices recover induced Littlewood graphs]\label{lem:bc-step2-submatrix-recovers-induced-graph}
Let
\[
  M=(M_{ab})_{1\le a,b\le j}
\]
be a symmetric nonnegative integer matrix with zero diagonal and row sums equal
to $2$.  Let $\Gamma_M$ be the loopless labelled multigraph on
$\{1,\ldots,j\}$ with $M_{ab}$ parallel edges between $a$ and $b$ for
$a<b$.  Then $\Gamma_M$ is a loopless labelled $2$-regular multigraph.
For every subset $D\subset\{1,\ldots,j\}$, the entries
\[
  M_{ab}\qquad(a,b\in D,\ a<b)
\]
determine the induced deleted-only graph $(\Gamma_M)_D$.

In the content $(2^j)$ Step--2/Step--3 construction of
\cref{lem:bc-step2-matrix-doubled-content-degrees,lem:bc-step3-cells-are-step2-half-matrix},
any recovered Step--3 subfilling containing the cell for each unordered pair
$\{a,b\}\subset D$ determines this same induced graph.
\end{lemma}

\begin{proof}
The zero diagonal forbids loops, and the symmetry lets $M_{ab}$ be interpreted
as the multiplicity of the unordered edge $\{a,b\}$.  The degree of vertex
$a$ in $\Gamma_M$ is
\[
  \sum_{b<a}M_{ba}+\sum_{a<b}M_{ab}
  =
  \sum_{b=1}^jM_{ab}
  =
  2,
\]
so $\Gamma_M$ is loopless and $2$-regular.  The induced graph on $D$ keeps
exactly the edges whose two endpoints lie in $D$, hence exactly the displayed
entries.

For a semistandard tableau of content $(2^j)$, the Step--2 matrix has the
displayed symmetry, zero diagonal, and degree condition by
\cref{lem:bc-step2-matrix-doubled-content-degrees}.  By
\cref{lem:bc-step3-cells-are-step2-half-matrix}, each recovered Step--3 cell
entry is the corresponding off-diagonal entry of this Step--2 matrix, up to
the symmetric transpose.  Therefore recovering all Step--3 cells indexed by
unordered pairs inside $D$ recovers all entries $M_{ab}$ with
$a,b\in D$, $a<b$, and hence recovers $(\Gamma_M)_D$.
\end{proof}

\begin{proposition}[Target-window submatrices recover the large deleted-only graph]\label{prop:bc-target-window-submatrix-recovers-large-deleted-only-graph}
Work in one residual row and moment of \cref{rem:former-post29-tail}.  Put $q=b+1$,
let $S\subset\{1,\ldots,j\}$ have cardinality $q$, and work in the
single-deficit setting of
\cref{prop:bc-single-deficit-stable-retained-repair}.  Let
$A\subset S$ be the support of the anchored component and put
\[
  D=S\setminus A.
\]
Assume that $A$ is a nonempty proper interval in the ordered witness set, that
every internal edge of the anchored path joins adjacent witness ranks, and
that the Littlewood graph $\Gamma$ is the graph $\Gamma_M$ associated by
\cref{lem:bc-step2-submatrix-recovers-induced-graph} to the Step--2 matrix of
the original tableau.

Suppose that a datum recovered from the target-size parser, without using an
additional branch word, determines the Step--3 cell entry for every unordered
pair $\{a,b\}\subset D$.  Then $\Gamma$ is recovered from
\[
  S,\quad \Gamma^\ast,\quad \rho(v)\in\{0,1\}^q,
  \quad \iota(A)\in\{0,1\}^q
\]
and that target-window datum.  The same assertion holds in the height
fixed-witness fiber.

Consequently, in the large-complement branch
$|S\setminus A|\ge R(q)+1$, one sufficient replacement for the false
graph-formal adjacency theorem is a target-size parser/carrier theorem
recovering the Step--3 subfilling on all unordered pairs inside
$S\setminus A$.
\end{proposition}

\begin{proof}
The target-window datum recovers the Step--3 entries on all unordered pairs in
$D$.  By \cref{lem:bc-step2-submatrix-recovers-induced-graph}, these entries
determine the induced deleted-only graph $\Gamma_D=\Gamma_{S\setminus A}$.
Use this recovered graph as the datum $\Delta_A$ in
\cref{cor:bc-rank-adjacent-proper-support-deleted-only-readout}.  The
retained-vertex code $\rho(v)$ and the interval code $\iota(A)$ supply the
remaining inputs of that corollary, while the adjacent-rank hypothesis on the
anchored internal edges is one of its hypotheses.  Therefore $\Gamma$ is
recovered.  The height case is the height half of the same corollary.

The final sentence is this recovery statement specialized to the
complementary range left by
\cref{prop:bc-combined-complement-large-joint-slot-graph-dichotomy}.
\end{proof}

\begin{proposition}[Ferrers carrier covers recover the large deleted-only graph]\label{prop:bc-ferrers-carrier-cover-recovers-large-deleted-only-graph}
Work in one residual row and moment of \cref{rem:former-post29-tail}.  Put $q=b+1$,
let $S\subset\{1,\ldots,j\}$ have cardinality $q$, and work in the
single-deficit setting of
\cref{prop:bc-single-deficit-stable-retained-repair}.  Let
$A\subset S$ be the support of the anchored component and put
\[
  D=S\setminus A.
\]
Assume that $A$ is a nonempty proper interval in the ordered witness set, that
every internal edge of the anchored path joins adjacent witness ranks, and
that the Littlewood graph $\Gamma$ is the graph $\Gamma_M$ associated by
\cref{lem:bc-step2-submatrix-recovers-induced-graph} to the Step--2 matrix of
the original tableau.

Suppose the Step--3 half-square admits a finite family of Ferrers
subarrangements
\[
  E_1,\ldots,E_N
\]
such that the target-size parser determines, without an additional branch
word, the partition labels on the top/right boundary of each $E_r$, and such
that every Step--3 cell corresponding to an unordered pair
\[
  \{a,b\}\subset D,\qquad a<b,
\]
belongs to at least one $E_r$.  Then $\Gamma$ is recovered from
\[
  S,\quad \Gamma^\ast,\quad \rho(v)\in\{0,1\}^q,
  \quad \iota(A)\in\{0,1\}^q
\]
and those carrier boundaries.  The same assertion holds in the height
fixed-witness fiber.

Consequently, in the large-complement branch
$|S\setminus A|\ge R(q)+1$, it is enough to prove a target-size carrier-cover
theorem: the parsed Step--3 Ferrers carriers cover every unordered pair inside
$S\setminus A$.
\end{proposition}

\begin{proof}
Apply \cref{cor:bc-parsed-dual-rsk-prime-boundary-recovers-subfilling} to each
Ferrers subarrangement $E_r$.  The recovered top/right boundary of $E_r$
determines every Step--3 cell entry in $E_r$.  Since the family covers every
cell indexed by an unordered pair inside $D$, the same parser datum determines
the Step--3 cell entry for every unordered pair $\{a,b\}\subset D$.

The hypotheses of
\cref{prop:bc-target-window-submatrix-recovers-large-deleted-only-graph} are
therefore satisfied, with the recovered carrier boundaries supplying the
target-window datum.  That proposition recovers $\Gamma$ from the displayed
data.  The height case is identical, using the height half of the same
proposition.
\end{proof}

\begin{corollary}[Interval-hull carrier boundaries recover the large deleted-only graph]\label{cor:bc-interval-hull-carriers-recover-large-deleted-only-graph}
Work in the setting of
\cref{prop:bc-ferrers-carrier-cover-recovers-large-deleted-only-graph}.  Write
the maximal consecutive label intervals of
\[
  D=S\setminus A
\]
as
\[
  D=\bigsqcup_{\alpha=1}^L
  \{a_\alpha,a_\alpha+1,\ldots,b_\alpha\},
  \qquad b_\alpha+1<a_{\alpha+1}.
\]
Suppose that for every $1\le\alpha\le\beta\le L$ the target-size parser
determines, without an additional branch word, the boundary segment
\[
  \rho_{2(j-b_\beta)},\rho_{2(j-b_\beta)+1},
  \ldots,\rho_{2(j-a_\alpha+1)}
\]
of the Step--3 semistandard path.  Then $\Gamma$ is recovered from
\[
  S,\quad \Gamma^\ast,\quad \rho(v)\in\{0,1\}^q,
  \quad \iota(A)\in\{0,1\}^q
\]
and those boundary segments.  The same assertion holds in the height
fixed-witness fiber.
\end{corollary}

\begin{proof}
For each pair $\alpha\le\beta$, let $E_{\alpha,\beta}$ be the consecutive
Step--3 label-interval subarrangement
\[
  E_{a_\alpha,b_\beta}
\]
of \cref{lem:bc-step3-consecutive-label-interval-boundary}.  That lemma
identifies the top/right boundary of $E_{\alpha,\beta}$ with the displayed
segment.  Hence the assumed parser data recover the top/right boundary of
each $E_{\alpha,\beta}$.

Let $\{x,y\}\subset D$ with $x<y$.  Then $x$ lies in a unique interval
$\{a_\alpha,\ldots,b_\alpha\}$ and $y$ lies in a unique interval
$\{a_\beta,\ldots,b_\beta\}$ with $\alpha\le\beta$.  Therefore both labels
belong to the consecutive interval
\[
  \{a_\alpha,a_\alpha+1,\ldots,b_\beta\},
\]
so the Step--3 cell for $\{x,y\}$ lies in $E_{\alpha,\beta}$.  Thus the
family $\{E_{\alpha,\beta}\}_{1\le\alpha\le\beta\le L}$ covers every unordered
pair inside $D$.  Applying
\cref{prop:bc-ferrers-carrier-cover-recovers-large-deleted-only-graph}
recovers $\Gamma$.  The height case is identical.
\end{proof}

\begin{corollary}[Terminal boundaries above the deleted maximum recover the large deleted-only graph]\label{cor:bc-terminal-max-boundary-recovers-large-deleted-only-graph}
Work in the setting of
\cref{prop:bc-target-window-submatrix-recovers-large-deleted-only-graph}, and
assume $D=S\setminus A$ is nonempty.  Let
\[
  M_D=\max D .
\]
Let $M$ be any integer with
\[
  M_D\le M\le j.
\]
Suppose the target-size parser determines, without an additional branch word,
the Step--3 terminal boundary segment
\[
  \rho_{2(j-M)},\rho_{2(j-M)+1},\ldots,\rho_{2j}.
\]
Then $\Gamma$ is recovered from
\[
  S,\quad \Gamma^\ast,\quad \rho(v)\in\{0,1\}^q,
  \quad \iota(A)\in\{0,1\}^q
\]
and that terminal boundary segment.  The same assertion holds in the height
fixed-witness fiber.
\end{corollary}

\begin{proof}
Every unordered pair $\{a,b\}\subset D$ has $a,b\le M_D\le M$.  Hence its Step--3
cell lies in the low-low terminal cell set $L_{M_D}$ of
\cref{lem:bc-low-label-incidence-split} and also in $L_M$.  By
\cref{lem:bc-low-low-terminal-corner-boundary}, the displayed terminal
boundary segment determines every Step--3 entry on $L_M$.  In particular
it determines every Step--3 cell entry indexed by an unordered pair inside
$D$.  Applying
\cref{prop:bc-target-window-submatrix-recovers-large-deleted-only-graph}
recovers $\Gamma$.  The height case is the height half of that proposition.
\end{proof}

\begin{proposition}[Continuation columns plus terminal boundaries recover the large deleted-only graph]\label{prop:bc-continuation-terminal-boundary-recovers-large-deleted-only-graph}
Work in the setting of
\cref{prop:bc-target-window-submatrix-recovers-large-deleted-only-graph}, and
assume $D=S\setminus A$ is nonempty.  Let $M$ be an integer with
\[
  1\le M\le j.
\]
Suppose the target-size parser determines, without an additional branch word,
the following two pieces of Step--3 data:
\begin{enumerate}
\item every Step--3 cell entry corresponding to an unordered pair
  $\{a,b\}\subset\{1,\ldots,j\}$ with $\max(a,b)>M$;
\item the terminal boundary segment
\[
  \rho_{2(j-M)},\rho_{2(j-M)+1},\ldots,\rho_{2j}.
\]
\end{enumerate}
Then $\Gamma$ is recovered from
\[
  S,\quad \Gamma^\ast,\quad \rho(v)\in\{0,1\}^q,
  \quad \iota(A)\in\{0,1\}^q
\]
and those parser data.  The same assertion holds in the height fixed-witness
fiber.
\end{proposition}

\begin{proof}
Let $\{a,b\}\subset D$ with $a<b$.  If $b>M$, then the corresponding Step--3
cell entry is determined by item~(1).  If $b\le M$, then both labels are at
most $M$, so the cell lies in the low-low terminal cell set $L_M$ of
\cref{lem:bc-low-label-incidence-split}.  By
\cref{lem:bc-low-low-terminal-corner-boundary}, the boundary segment in
item~(2) determines every Step--3 entry on $L_M$, and hence determines this
cell entry.

Thus the parser data determine the Step--3 entry for every unordered pair
inside $D$.  Applying
\cref{prop:bc-target-window-submatrix-recovers-large-deleted-only-graph}
recovers $\Gamma$.  The height case is the height half of that proposition.
\end{proof}

\begin{corollary}[Strict-left data plus a terminal boundary recover the large deleted-only graph]\label{cor:bc-strict-left-terminal-boundary-recovers-large-deleted-only-graph}
Work in the setting of
\cref{prop:bc-target-window-submatrix-recovers-large-deleted-only-graph}, and
also in the right-boundary auxiliary setting of
\cref{prop:bc-right-boundary-auxiliary-global-recovery}.  Assume the final
zone is in the non-full tight-height setting of
\cref{lem:bc-tight-outside-recovers-left-growth-segment}.  Put
\[
  M=2m,\qquad Z_L=\{j-2m+1,\ldots,j\},
\]
and assume $D=S\setminus A$ is nonempty.  Suppose the outside datum
$\mathcal D_{\mathrm{out}}(P)$ has been recovered as in
\cref{prop:bc-right-boundary-auxiliary-global-recovery}, from the retained
blocks and the earlier-zone inverses using only the branch words already
present in that recovery.  Suppose further that the remaining target-window
data determine, without an additional branch word, the terminal Step--3
boundary segment
\[
  \rho_{2(j-2m)},\rho_{2(j-2m)+1},\ldots,\rho_{2j}.
\]
Then $\Gamma$ is recovered from
\[
  S,\quad \Gamma^\ast,\quad \rho(v)\in\{0,1\}^q,
  \quad \iota(A)\in\{0,1\}^q
\]
together with the recovered outside datum and this terminal boundary segment.
\end{corollary}

\begin{proof}
By \cref{lem:bc-strict-left-recovers-continuation-columns}, the outside datum
$\mathcal D_{\mathrm{out}}(P)$ determines every Step--3 entry in the
continuation column subarrangement $J_M$, whose columns are incident to the
labels $M+1,\ldots,j$.  Hence it determines every Step--3 cell entry indexed by
an unordered pair $\{a,b\}$ with $\max(a,b)>M$.

The displayed terminal boundary segment is the segment required in
\cref{prop:bc-continuation-terminal-boundary-recovers-large-deleted-only-graph}
for this value of $M$.  Applying that proposition gives the claimed recovery
of $\Gamma$.
\end{proof}

\begin{corollary}[Full-budget terminal boundary lists recover the large deleted-only graph]\label{cor:bc-full-budget-terminal-boundary-list-recovers-large-deleted-only-graph}
Work in the setting of
\cref{prop:bc-target-window-submatrix-recovers-large-deleted-only-graph} and
in the right-boundary auxiliary setting of
\cref{prop:bc-right-boundary-auxiliary-global-recovery}.  Assume the final
zone is in the non-full tight-height setting of
\cref{lem:bc-tight-outside-recovers-left-growth-segment}.  Put
\[
  M=2m,\qquad Z_L=\{j-2m+1,\ldots,j\},
\]
and assume $D=S\setminus A$ is nonempty and $m\ge1$.  Suppose the outside
datum $\mathcal D_{\mathrm{out}}(P)$ has been recovered as in
\cref{prop:bc-right-boundary-auxiliary-global-recovery}, using only the branch
words already present in that recovery.  Suppose further that the remaining
target-window data determine a canonically ordered finite set
\[
  \mathcal B_{\mathrm{tw}}
\]
of possible terminal Step--3 boundary segments
\[
  \tau=(\tau_{2(j-2m)},\tau_{2(j-2m)+1},\ldots,\tau_{2j}),
\]
that the true segment
\[
  (\rho_{2(j-2m)},\rho_{2(j-2m)+1},\ldots,\rho_{2j})
\]
belongs to $\mathcal B_{\mathrm{tw}}$, and that
\[
  |\mathcal B_{\mathrm{tw}}|\le2^{2m}.
\]
Then $\Gamma$ is recovered from
\[
  S,\quad \Gamma^\ast,\quad \rho(v)\in\{0,1\}^q,\quad
  \iota(A)\in\{0,1\}^q,
\]
the recovered outside datum, $\mathcal B_{\mathrm{tw}}$, and the already
budgeted final-zone word $\beta_L\in\{0,1\}^{2m}$.
\end{corollary}

\begin{proof}
Encode the rank of the true terminal boundary segment in the canonical order
on $\mathcal B_{\mathrm{tw}}$ by the word $\beta_L$.  Since
$|\mathcal B_{\mathrm{tw}}|\le2^{2m}$, this uses no branch data beyond the
already budgeted final-zone word and selects the true segment from
$\mathcal B_{\mathrm{tw}}$.

With that selected segment in hand, the hypotheses of
\cref{cor:bc-strict-left-terminal-boundary-recovers-large-deleted-only-graph}
are satisfied.  That corollary recovers $\Gamma$ from the displayed data.
\end{proof}

\begin{corollary}[Full-budget low-low candidate lists recover the large deleted-only graph]\label{cor:bc-full-budget-low-low-candidate-list-recovers-large-deleted-only-graph}
Work in the setting of
\cref{prop:bc-target-window-submatrix-recovers-large-deleted-only-graph} and
in the right-boundary auxiliary setting of
\cref{prop:bc-right-boundary-auxiliary-global-recovery}.  Assume the final
zone is in the non-full tight-height setting of
\cref{lem:bc-tight-outside-recovers-left-growth-segment}.  Put
\[
  M=2m,\qquad Z_L=\{j-2m+1,\ldots,j\},
\]
and let $L_{2m}$ be the low-low cell set of
\cref{lem:bc-low-label-incidence-split}.  Assume $D=S\setminus A$ is nonempty
and $m\ge1$.  Suppose the outside datum $\mathcal D_{\mathrm{out}}(P)$ has
been recovered as in \cref{prop:bc-right-boundary-auxiliary-global-recovery},
using only the branch words already present in that recovery.  Suppose further
that the remaining target-window data determine a canonically ordered finite
set
\[
  \mathcal L_{\mathrm{tw}}
\]
of possible Step--3 entry assignments on $L_{2m}$, that the true assignment
belongs to $\mathcal L_{\mathrm{tw}}$, and that
\[
  |\mathcal L_{\mathrm{tw}}|\le2^{2m}.
\]
Then $\Gamma$ is recovered from
\[
  S,\quad \Gamma^\ast,\quad \rho(v)\in\{0,1\}^q,\quad
  \iota(A)\in\{0,1\}^q,
\]
the recovered outside datum, $\mathcal L_{\mathrm{tw}}$, and the already
budgeted final-zone word $\beta_L\in\{0,1\}^{2m}$.
\end{corollary}

\begin{proof}
Use the $2m$ bits of $\beta_L$ to encode the rank of the true low-low
assignment in the canonical order on $\mathcal L_{\mathrm{tw}}$.  Since
$|\mathcal L_{\mathrm{tw}}|\le2^{2m}$, this selects the true Step--3 entries
on $L_{2m}$ without any branch data beyond the already budgeted final-zone
word.

By \cref{lem:bc-strict-left-recovers-continuation-columns}, the recovered
outside datum determines every Step--3 entry in the continuation column
subarrangement $J_M$, whose columns are incident to the labels
$M+1,\ldots,j$.  Let $\{a,b\}\subset D$ with $a<b$.  If $b>M$, then the
corresponding Step--3 entry is determined by the recovered continuation
columns.  If $b\le M$, then the cell lies in the low-low cell set $L_M=L_{2m}$,
and the selected low-low assignment determines it.

Thus every Step--3 cell indexed by an unordered pair inside $D$ is recovered.
By \cref{lem:bc-step2-submatrix-recovers-induced-graph}, these entries
determine the induced deleted-only graph $\Gamma_{S\setminus A}$.  Using this
graph as the datum $\Delta_A$ in
\cref{cor:bc-rank-adjacent-proper-support-deleted-only-readout} recovers
$\Gamma$ from the displayed data.
\end{proof}

\begin{corollary}[Full-budget odd-packet classes recover the large deleted-only graph]\label{cor:bc-full-budget-omega-recovers-large-deleted-only-graph}
Work in the setting of
\cref{cor:bc-full-budget-low-low-candidate-list-recovers-large-deleted-only-graph},
and assume $m\ge9$.  Suppose the remaining target-window data determine
\[
  \lambda^{(2m)},\qquad E_1,\ldots,E_m,
\]
the labelled continuation boxes, and a canonically ordered set
\[
  \Omega_{\mathrm{tw}}\subseteq
  \Omega(\lambda^{(2m)};E_1,\ldots,E_m)
\]
which contains the true deleted odd-packet assignment and satisfies
\[
  |\Omega_{\mathrm{tw}}|\le2^{2m}.
\]
Then $\Gamma$ is recovered from
\[
  S,\quad \Gamma^\ast,\quad \rho(v)\in\{0,1\}^q,\quad
  \iota(A)\in\{0,1\}^q,
\]
the recovered outside datum, $\Omega_{\mathrm{tw}}$, and the already budgeted
final-zone word $\beta_L\in\{0,1\}^{2m}$.
\end{corollary}

\begin{proof}
Let $\mathcal L_{\mathrm{tw}}$ be the low-low candidate set constructed in
\cref{cor:bc-refined-omega-induces-low-low-candidates}, with the displayed
set $\Omega_{\mathrm{tw}}$ in place of the half-budget class there.  That
construction is a fixed function of the target-window data.  Its membership
and injectivity proof does not use the half-budget cardinality assumption:
the true low-low assignment belongs to $\mathcal L_{\mathrm{tw}}$, and the
map from $\mathcal L_{\mathrm{tw}}$ to the corresponding deleted odd-packet
assignment is injective into $\Omega_{\mathrm{tw}}$.  Hence
\[
  |\mathcal L_{\mathrm{tw}}|
  \le |\Omega_{\mathrm{tw}}|
  \le 2^{2m}.
\]
The hypotheses of
\cref{cor:bc-full-budget-low-low-candidate-list-recovers-large-deleted-only-graph}
are therefore satisfied with this $\mathcal L_{\mathrm{tw}}$, and that
corollary recovers $\Gamma$.
\end{proof}

\begin{corollary}[Bounded signature lists recover the large deleted-only graph]\label{cor:bc-bounded-signature-list-recovers-large-deleted-only-graph}
Work in the setting of
\cref{cor:bc-full-budget-omega-recovers-large-deleted-only-graph}, and assume
\(m\ge9\).  Suppose the remaining target-window data determine
\[
  \lambda^{(2m)},\qquad E_1,\ldots,E_m,
\]
the labelled continuation boxes, and a canonically ordered finite set
\[
  \Sigma_{\mathrm{tw}}
\]
of commuting normal-form signature words such that the true fixed-even standard
tableau has signature in \(\Sigma_{\mathrm{tw}}\) and
\[
  |\Sigma_{\mathrm{tw}}|\le2^m .
\]
Then \(\Gamma\) is recovered from
\[
  S,\quad \Gamma^\ast,\quad \rho(v)\in\{0,1\}^q,\quad
  \iota(A)\in\{0,1\}^q,
\]
the recovered outside datum, these target-window data, and the already
budgeted final-zone word \(\beta_L\in\{0,1\}^{2m}\).
\end{corollary}

\begin{proof}
\Cref{cor:bc-bounded-signature-list-full-budget-omega} constructs from the
displayed target-window data a canonically ordered set
\[
  \Omega_{\mathrm{tw}}(\Sigma_{\mathrm{tw}})
  \subseteq \Omega(\lambda^{(2m)};E_1,\ldots,E_m)
\]
which contains the true deleted odd-packet assignment and has cardinality at
most \(2^{2m}\).  Applying
\cref{cor:bc-full-budget-omega-recovers-large-deleted-only-graph} to this
full-budget odd-packet class recovers \(\Gamma\).
\end{proof}

\begin{corollary}[Nine-to-twelve broad odd-packet classes recover the large deleted-only graph]\label{cor:bc-nine-to-twelve-broad-omega-recovers-large-deleted-only-graph}
Work in the setting of
\cref{cor:bc-full-budget-low-low-candidate-list-recovers-large-deleted-only-graph}.
Suppose
\[
  9\le m\le12
\]
and suppose the remaining target-window data determine
\[
  \lambda^{(2m)},\qquad E_1,\ldots,E_m,
\]
and the labelled continuation boxes.  Then $\Gamma$ is recovered from the
data
\[
  S,\quad \Gamma^\ast,\quad \rho(v)\in\{0,1\}^q,\quad
  \iota(A)\in\{0,1\}^q,
\]
the recovered outside datum, these target-window data, and the already
budgeted final-zone word
\[
  \beta_L\in\{0,1\}^{2m}.
\]
\end{corollary}

\begin{proof}
Let
\[
  \Omega_{\mathrm{tw}}
  =
  \Omega(\lambda^{(2m)};E_1,\ldots,E_m)
\]
with the lexicographic order of \cref{def:bc-tight-odd-packet-omega}.  The
displayed target-window data determine this set.  By
\cref{lem:bc-true-tight-packets-in-omega}, it contains the true deleted
odd-packet assignment.  By \cref{lem:bc-binary-prefix-factorial-bound},
\[
  |\Omega_{\mathrm{tw}}|\le2(m-2)! .
\]
For $9\le m\le12$,
\[
  2(m-2)!\le2^{2m};
\]
the four cases give
\[
  10080\le262144,\quad
  80640\le1048576,\quad
  725760\le4194304,\quad
  7257600\le16777216 .
\]
Thus the hypotheses of
\cref{cor:bc-full-budget-omega-recovers-large-deleted-only-graph} hold, and
that corollary recovers $\Gamma$.
\end{proof}

\begin{corollary}[Thirteen-to-twenty-eight broad odd-packet classes recover the large deleted-only graph]\label{cor:bc-thirteen-to-twenty-eight-envelope-omega-recovers-large-deleted-only-graph}
Work in the setting of
\cref{cor:bc-full-budget-low-low-candidate-list-recovers-large-deleted-only-graph}.
Suppose
\[
  13\le m\le28
\]
and suppose the remaining target-window data determine
\[
  \lambda^{(2m)},\qquad E_1,\ldots,E_m,
\]
and the labelled continuation boxes.  Then $\Gamma$ is recovered from the
data
\[
  S,\quad \Gamma^\ast,\quad \rho(v)\in\{0,1\}^q,\quad
  \iota(A)\in\{0,1\}^q,
\]
the recovered outside datum, these target-window data, and the already
budgeted final-zone word
\[
  \beta_L\in\{0,1\}^{2m}.
\]
\end{corollary}

\begin{proof}
Let
\[
  \Omega_{\mathrm{tw}}
  =
  \Omega(\lambda^{(2m)};E_1,\ldots,E_m)
\]
with the lexicographic order of \cref{def:bc-tight-odd-packet-omega}.  The
displayed target-window data determine this set, and
\cref{lem:bc-true-tight-packets-in-omega} puts the true deleted odd-packet
assignment in it.

\Cref{lem:bc-omega-fixed-cell-completions,prop:bc-completions-equal-reduced-growth-chains,lem:bc-reduced-growth-reverse-corner-tree,lem:bc-reverse-tree-fixed-even-standard-tableaux}
identify its cardinality with the cardinality of the associated fixed-even
standard-tableau fiber
\[
  \mathcal S(\gamma;e_1,\ldots,e_{m-1}).
\]
By \cref{lem:bc-fixed-even-recursive-envelope-bounds-fibers},
\[
  |\Omega_{\mathrm{tw}}|
  =
  |\mathcal S(\gamma;e_1,\ldots,e_{m-1})|
  \le W(\gamma).
\]
Since $13\le m\le28$,
\cref{prop:bc-fixed-even-envelope-m28-certificate} gives
\[
  W(\gamma)\le2^{2m}.
\]
Thus the hypotheses of
\cref{cor:bc-full-budget-omega-recovers-large-deleted-only-graph} hold with
the unrefined compatibility set $\Omega_{\mathrm{tw}}$, and that corollary
recovers $\Gamma$.
\end{proof}

\begin{corollary}[Top-fixed envelope non-frontier broad odd-packet classes recover the large deleted-only graph]\label{cor:bc-envelope-nonfrontier-omega-recovers-large-deleted-only-graph}
Work in the setting of
\cref{cor:bc-full-budget-low-low-candidate-list-recovers-large-deleted-only-graph}.
Suppose \(m\ge13\), and suppose the remaining target-window data determine
\[
  \lambda^{(2m)},\qquad E_1,\ldots,E_m,
\]
and the labelled continuation boxes.  Let
\[
  \Omega_{\mathrm{tw}}
  =
  \Omega(\lambda^{(2m)};E_1,\ldots,E_m)
\]
with the lexicographic order of \cref{def:bc-tight-odd-packet-omega}.  If
\(\Omega_{\mathrm{tw}}\ne\varnothing\), let
\((\gamma;e_1,\ldots,e_{m-1})\) be the associated fixed-even datum of
\cref{def:bc-odd-fixed-even-standard-tableaux}, and assume
\[
  (\gamma,e_{m-1})\notin\mathcal P^{\mathrm{env}}_m.
\]
Then \(\Gamma\) is recovered from the data
\[
  S,\quad \Gamma^\ast,\quad \rho(v)\in\{0,1\}^q,\quad
  \iota(A)\in\{0,1\}^q,
\]
the recovered outside datum, these target-window data, and the already
budgeted final-zone word
\[
  \beta_L\in\{0,1\}^{2m}.
\]
\end{corollary}

\begin{proof}
The displayed target-window data determine
\(\Omega_{\mathrm{tw}}\), and
\cref{lem:bc-true-tight-packets-in-omega} puts the true deleted odd-packet
assignment in it.  If \(\Omega_{\mathrm{tw}}=\varnothing\), there is no true
assignment in the present fiber, so the recovery assertion is vacuous.

Assume \(\Omega_{\mathrm{tw}}\ne\varnothing\).  As in the proof of
\cref{cor:bc-thirteen-to-twenty-eight-envelope-omega-recovers-large-deleted-only-graph},
\cref{lem:bc-omega-fixed-cell-completions,prop:bc-completions-equal-reduced-growth-chains,lem:bc-reduced-growth-reverse-corner-tree,lem:bc-reverse-tree-fixed-even-standard-tableaux}
identify its cardinality with
\[
  |\mathcal S(\gamma;e_1,\ldots,e_{m-1})|.
\]
By \cref{lem:bc-fixed-even-top-envelope-bounds-fibers},
\[
  |\Omega_{\mathrm{tw}}|\le W^{\mathrm{top}}(\gamma;e_{m-1}).
\]
Since \((\gamma,e_{m-1})\notin\mathcal P^{\mathrm{env}}_m\),
\cref{def:bc-fixed-even-envelope-frontier} gives
\[
  W^{\mathrm{top}}(\gamma;e_{m-1})\le2^{2m}.
\]
Thus the hypotheses of
\cref{cor:bc-full-budget-omega-recovers-large-deleted-only-graph} hold with
the unrefined compatibility set \(\Omega_{\mathrm{tw}}\), and that corollary
recovers \(\Gamma\).
\end{proof}

\begin{corollary}[Top-two envelope non-frontier broad odd-packet classes recover the large deleted-only graph]\label{cor:bc-envelope-top-two-nonfrontier-omega-recovers-large-deleted-only-graph}
Work in the setting of
\cref{cor:bc-full-budget-low-low-candidate-list-recovers-large-deleted-only-graph}.
Suppose \(m\ge13\), and suppose the remaining target-window data determine
\[
  \lambda^{(2m)},\qquad E_1,\ldots,E_m,
\]
and the labelled continuation boxes.  Let
\[
  \Omega_{\mathrm{tw}}
  =
  \Omega(\lambda^{(2m)};E_1,\ldots,E_m)
\]
with the lexicographic order of \cref{def:bc-tight-odd-packet-omega}.  If
\(\Omega_{\mathrm{tw}}\ne\varnothing\), let
\((\gamma;e_1,\ldots,e_{m-1})\) be the associated fixed-even datum of
\cref{def:bc-odd-fixed-even-standard-tableaux}, and assume
\[
  (\gamma,e_{m-1},e_{m-2})\notin\mathcal P^{\mathrm{env},2}_m.
\]
Then \(\Gamma\) is recovered from the data
\[
  S,\quad \Gamma^\ast,\quad \rho(v)\in\{0,1\}^q,\quad
  \iota(A)\in\{0,1\}^q,
\]
the recovered outside datum, these target-window data, and the already
budgeted final-zone word
\[
  \beta_L\in\{0,1\}^{2m}.
\]
\end{corollary}

\begin{proof}
The displayed target-window data determine \(\Omega_{\mathrm{tw}}\), and
\cref{lem:bc-true-tight-packets-in-omega} puts the true deleted odd-packet
assignment in it.  If \(\Omega_{\mathrm{tw}}=\varnothing\), there is no true
assignment in the present fiber, so the recovery assertion is vacuous.

Assume \(\Omega_{\mathrm{tw}}\ne\varnothing\).  As in the proof of
\cref{cor:bc-envelope-nonfrontier-omega-recovers-large-deleted-only-graph},
\cref{lem:bc-omega-fixed-cell-completions,prop:bc-completions-equal-reduced-growth-chains,lem:bc-reduced-growth-reverse-corner-tree,lem:bc-reverse-tree-fixed-even-standard-tableaux}
identify its cardinality with
\[
  |\mathcal S(\gamma;e_1,\ldots,e_{m-1})|.
\]
By \cref{lem:bc-fixed-even-top-two-envelope-bounds-fibers},
\[
  |\Omega_{\mathrm{tw}}|
  \le W^{\mathrm{top},2}(\gamma;e_{m-1},e_{m-2}).
\]
Since \((\gamma,e_{m-1},e_{m-2})\notin\mathcal P^{\mathrm{env},2}_m\),
\cref{def:bc-fixed-even-envelope-frontier} gives
\[
  W^{\mathrm{top},2}(\gamma;e_{m-1},e_{m-2})\le2^{2m}.
\]
Thus the hypotheses of
\cref{cor:bc-full-budget-omega-recovers-large-deleted-only-graph} hold with
the unrefined compatibility set \(\Omega_{\mathrm{tw}}\), and that corollary
recovers \(\Gamma\).
\end{proof}

\begin{corollary}[Fixed-depth envelope non-frontier broad odd-packet classes recover the large deleted-only graph]\label{cor:bc-fixed-depth-nonfrontier-omega-recovers-large-deleted-only-graph}
Work in the setting of
\cref{cor:bc-full-budget-low-low-candidate-list-recovers-large-deleted-only-graph}.
Let \(d\ge3\), and suppose \(m\ge d+1\).  Suppose the remaining target-window
data determine
\[
  \lambda^{(2m)},\qquad E_1,\ldots,E_m,
\]
and the labelled continuation boxes.  Let
\[
  \Omega_{\mathrm{tw}}
  =
  \Omega(\lambda^{(2m)};E_1,\ldots,E_m)
\]
with the lexicographic order of \cref{def:bc-tight-odd-packet-omega}.  If
\(\Omega_{\mathrm{tw}}\ne\varnothing\), let
\((\gamma;e_1,\ldots,e_{m-1})\) be the associated fixed-even datum of
\cref{def:bc-odd-fixed-even-standard-tableaux}, and assume
\[
  (\gamma,e_{m-1},e_{m-2},\ldots,e_{m-d})
  \notin\mathcal P^{\mathrm{env},d}_m.
\]
Then \(\Gamma\) is recovered from the data
\[
  S,\quad \Gamma^\ast,\quad \rho(v)\in\{0,1\}^q,\quad
  \iota(A)\in\{0,1\}^q,
\]
the recovered outside datum, these target-window data, and the already
budgeted final-zone word
\[
  \beta_L\in\{0,1\}^{2m}.
\]
\end{corollary}

\begin{proof}
The displayed target-window data determine \(\Omega_{\mathrm{tw}}\), and
\cref{lem:bc-true-tight-packets-in-omega} puts the true deleted odd-packet
assignment in it.  If \(\Omega_{\mathrm{tw}}=\varnothing\), there is no true
assignment in the present fiber, so the recovery assertion is vacuous.

Assume \(\Omega_{\mathrm{tw}}\ne\varnothing\).  As in the proof of
\cref{cor:bc-envelope-top-two-nonfrontier-omega-recovers-large-deleted-only-graph},
\cref{lem:bc-omega-fixed-cell-completions,prop:bc-completions-equal-reduced-growth-chains,lem:bc-reduced-growth-reverse-corner-tree,lem:bc-reverse-tree-fixed-even-standard-tableaux}
identify its cardinality with
\[
  |\mathcal S(\gamma;e_1,\ldots,e_{m-1})|.
\]
If \(d=3\), then \cref{lem:bc-fixed-even-top-three-envelope-bounds-fibers}
gives
\[
  |\Omega_{\mathrm{tw}}|
  \le W^{\mathrm{top},3}(\gamma;e_{m-1},e_{m-2},e_{m-3}).
\]
If \(d\ge4\), then
\cref{lem:bc-fixed-even-fixed-depth-envelope-bounds-fibers} gives
\[
  |\Omega_{\mathrm{tw}}|
  \le
  W^{\mathrm{top},d}(\gamma;e_{m-1},e_{m-2},\ldots,e_{m-d}).
\]
In both cases, the displayed non-frontier hypothesis and
\cref{def:bc-fixed-even-envelope-frontier} give
\[
  |\Omega_{\mathrm{tw}}|\le2^{2m}.
\]
Thus the hypotheses of
\cref{cor:bc-full-budget-omega-recovers-large-deleted-only-graph} hold with
the unrefined compatibility set \(\Omega_{\mathrm{tw}}\), and that corollary
recovers \(\Gamma\).
\end{proof}

\begin{corollary}[Twenty-nine top-three envelope broad odd-packet classes recover the large deleted-only graph]\label{cor:bc-twenty-nine-top-three-envelope-omega-recovers-large-deleted-only-graph}
Work in the setting of
\cref{cor:bc-full-budget-low-low-candidate-list-recovers-large-deleted-only-graph}.
Suppose \(m=29\), and suppose the remaining target-window data determine
\[
  \lambda^{(58)},\qquad E_1,\ldots,E_{29},
\]
and the labelled continuation boxes.  Then \(\Gamma\) is recovered from the
data
\[
  S,\quad \Gamma^\ast,\quad \rho(v)\in\{0,1\}^q,\quad
  \iota(A)\in\{0,1\}^q,
\]
the recovered outside datum, these target-window data, and the already
budgeted final-zone word
\[
  \beta_L\in\{0,1\}^{58}.
\]
\end{corollary}

\begin{proof}
Let
\[
  \Omega_{\mathrm{tw}}
  =
  \Omega(\lambda^{(58)};E_1,\ldots,E_{29})
\]
with the lexicographic order of \cref{def:bc-tight-odd-packet-omega}.  The
displayed target-window data determine this set, and
\cref{lem:bc-true-tight-packets-in-omega} puts the true deleted odd-packet
assignment in it.  If \(\Omega_{\mathrm{tw}}=\varnothing\), the assertion is
vacuous.

Assume \(\Omega_{\mathrm{tw}}\ne\varnothing\), and let
\((\gamma;e_1,\ldots,e_{28})\) be the associated fixed-even datum.  As in the
proof of
\cref{cor:bc-thirteen-to-twenty-eight-envelope-omega-recovers-large-deleted-only-graph},
\cref{lem:bc-omega-fixed-cell-completions,prop:bc-completions-equal-reduced-growth-chains,lem:bc-reduced-growth-reverse-corner-tree,lem:bc-reverse-tree-fixed-even-standard-tableaux}
identify the cardinality of \(\Omega_{\mathrm{tw}}\) with
\[
  |\mathcal S(\gamma;e_1,\ldots,e_{28})|.
\]
By \cref{lem:bc-fixed-even-top-three-envelope-bounds-fibers},
\[
  |\Omega_{\mathrm{tw}}|
  \le W^{\mathrm{top},3}(\gamma;e_{28},e_{27},e_{26}).
\]
\Cref{prop:bc-fixed-even-envelope-m29-obstruction} gives
\[
  \mathcal P^{\mathrm{env},3}_{29}=\varnothing,
\]
so the right side is at most \(2^{58}\).  Thus the hypotheses of
\cref{cor:bc-full-budget-omega-recovers-large-deleted-only-graph} hold with
the unrefined compatibility set \(\Omega_{\mathrm{tw}}\), and that corollary
recovers \(\Gamma\).
\end{proof}

\begin{corollary}[Thirty fixed-depth envelope broad odd-packet classes recover the large deleted-only graph]\label{cor:bc-thirty-fixed-depth-envelope-omega-recovers-large-deleted-only-graph}
Work in the setting of
\cref{cor:bc-full-budget-low-low-candidate-list-recovers-large-deleted-only-graph}.
Suppose \(m=30\), and suppose the remaining target-window data determine
\[
  \lambda^{(60)},\qquad E_1,\ldots,E_{30},
\]
and the labelled continuation boxes.  Then \(\Gamma\) is recovered from the
data
\[
  S,\quad \Gamma^\ast,\quad \rho(v)\in\{0,1\}^q,\quad
  \iota(A)\in\{0,1\}^q,
\]
the recovered outside datum, these target-window data, and the already
budgeted final-zone word
\[
  \beta_L\in\{0,1\}^{60}.
\]
\end{corollary}

\begin{proof}
Let
\[
  \Omega_{\mathrm{tw}}
  =
  \Omega(\lambda^{(60)};E_1,\ldots,E_{30})
\]
with the lexicographic order of \cref{def:bc-tight-odd-packet-omega}.  The
displayed target-window data determine this set, and
\cref{lem:bc-true-tight-packets-in-omega} puts the true deleted odd-packet
assignment in it.  If \(\Omega_{\mathrm{tw}}=\varnothing\), the assertion is
vacuous.

Assume \(\Omega_{\mathrm{tw}}\ne\varnothing\), and let
\((\gamma;e_1,\ldots,e_{29})\) be the associated fixed-even datum.  As in the
proof of
\cref{cor:bc-thirteen-to-twenty-eight-envelope-omega-recovers-large-deleted-only-graph},
\cref{lem:bc-omega-fixed-cell-completions,prop:bc-completions-equal-reduced-growth-chains,lem:bc-reduced-growth-reverse-corner-tree,lem:bc-reverse-tree-fixed-even-standard-tableaux}
identify the cardinality of \(\Omega_{\mathrm{tw}}\) with
\[
  |\mathcal S(\gamma;e_1,\ldots,e_{29})|.
\]
By \cref{lem:bc-fixed-even-fixed-depth-envelope-bounds-fibers},
\[
  |\Omega_{\mathrm{tw}}|
  \le W^{\mathrm{top},7}
  (\gamma;e_{29},e_{28},e_{27},e_{26},e_{25},e_{24},e_{23}).
\]
\Cref{prop:bc-fixed-even-envelope-m30-frontier-certificate} gives
\[
  \mathcal P^{\mathrm{env},7}_{30}=\varnothing,
\]
so the right side is at most \(2^{60}\).  Thus the hypotheses of
\cref{cor:bc-full-budget-omega-recovers-large-deleted-only-graph} hold with
the unrefined compatibility set \(\Omega_{\mathrm{tw}}\), and that corollary
recovers \(\Gamma\).
\end{proof}

\begin{corollary}[Thirty-one fixed-depth envelope broad odd-packet classes recover the large deleted-only graph]\label{cor:bc-thirty-one-fixed-depth-envelope-omega-recovers-large-deleted-only-graph}
Work in the setting of
\cref{cor:bc-full-budget-low-low-candidate-list-recovers-large-deleted-only-graph}.
Suppose \(m=31\), and suppose the remaining target-window data determine
\[
  \lambda^{(62)},\qquad E_1,\ldots,E_{31},
\]
and the labelled continuation boxes.  Then \(\Gamma\) is recovered from the
data
\[
  S,\quad \Gamma^\ast,\quad \rho(v)\in\{0,1\}^q,\quad
  \iota(A)\in\{0,1\}^q,
\]
the recovered outside datum, these target-window data, and the already
budgeted final-zone word
\[
  \beta_L\in\{0,1\}^{62}.
\]
\end{corollary}

\begin{proof}
Let
\[
  \Omega_{\mathrm{tw}}
  =
  \Omega(\lambda^{(62)};E_1,\ldots,E_{31})
\]
with the lexicographic order of \cref{def:bc-tight-odd-packet-omega}.  The
displayed target-window data determine this set, and
\cref{lem:bc-true-tight-packets-in-omega} puts the true deleted odd-packet
assignment in it.  If \(\Omega_{\mathrm{tw}}=\varnothing\), the assertion is
vacuous.

Assume \(\Omega_{\mathrm{tw}}\ne\varnothing\), and let
\((\gamma;e_1,\ldots,e_{30})\) be the associated fixed-even datum.  As in the
proof of
\cref{cor:bc-thirteen-to-twenty-eight-envelope-omega-recovers-large-deleted-only-graph},
\cref{lem:bc-omega-fixed-cell-completions,prop:bc-completions-equal-reduced-growth-chains,lem:bc-reduced-growth-reverse-corner-tree,lem:bc-reverse-tree-fixed-even-standard-tableaux}
identify the cardinality of \(\Omega_{\mathrm{tw}}\) with
\[
  |\mathcal S(\gamma;e_1,\ldots,e_{30})|.
\]
By \cref{lem:bc-fixed-even-fixed-depth-envelope-bounds-fibers},
\[
  |\Omega_{\mathrm{tw}}|
  \le W^{\mathrm{top},10}
  (\gamma;e_{30},e_{29},e_{28},e_{27},e_{26},
  e_{25},e_{24},e_{23},e_{22},e_{21}).
\]
\Cref{prop:bc-fixed-even-envelope-m31-frontier-certificate} gives
\[
  \mathcal P^{\mathrm{env},10}_{31}=\varnothing,
\]
so the right side is at most \(2^{62}\).  Thus the hypotheses of
\cref{cor:bc-full-budget-omega-recovers-large-deleted-only-graph} hold with
the unrefined compatibility set \(\Omega_{\mathrm{tw}}\), and that corollary
recovers \(\Gamma\).
\end{proof}

\begin{corollary}[Column-even tight prefixes recover the large deleted-only graph]\label{cor:bc-column-even-prefix-recovers-large-deleted-only-graph}
Work in the setting of
\cref{cor:bc-full-budget-low-low-candidate-list-recovers-large-deleted-only-graph}
and in the setting of \cref{cor:bc-tight-height-prefix-column-even-gate}.
Let
\[
  \lambda^{(0)}\subset\lambda^{(1)}\subset\cdots\subset\lambda^{(2m)}
\]
be the true tight-prefix Pieri path.  Assume $m\ge2$ and that
$\lambda^{(2m)}$ is column-even.  Suppose the final-zone replacement window
$V^\ast_L$ and branch word
\[
  \beta_L\in\{0,1\}^{2m}
\]
are chosen by the fixed odd-alternating height full-zone decoder of
\cref{prop:bc-odd-alternating-height-full-zone-decoder} applied to this tight
prefix.  Then $\Gamma$ is recovered from
\[
  S,\quad \Gamma^\ast,\quad \rho(v)\in\{0,1\}^q,\quad
  \iota(A)\in\{0,1\}^q,
\]
the recovered outside datum, $V^\ast_L$, and $\beta_L$.
\end{corollary}

\begin{proof}
By \cref{cor:bc-tight-height-prefix-column-even-gate}, the displayed prefix is
an odd alternating height full-zone path in the domain of
\cref{prop:bc-odd-alternating-height-full-zone-decoder}.  The decoder therefore
recovers the whole prefix path
\[
  \lambda^{(0)},\lambda^{(1)},\ldots,\lambda^{(2m)}
\]
from $V^\ast_L$ and $\beta_L$.

The recovered prefix path determines the boxes carrying the semistandard
labels $1,\ldots,2m$ by \cref{lem:bc-pieri-path-bijection}.  Hence it
determines the prefix labelled tableau.  By
\cref{lem:bc-step2-matrix-unique-from-tableau}, that labelled tableau
determines the Step--2 symmetric matrix entries whose two labels are at most
$2m$; by \cref{lem:bc-step3-cells-are-step2-half-matrix}, these are exactly
the Step--3 entries on the low-low cell set $L_{2m}$.

The recovered outside datum determines the continuation-column entries
$J_{2m}$ by \cref{lem:bc-strict-left-recovers-continuation-columns}.  Let
$\{a,b\}\subset D$ with $a<b$.  If $b>2m$, the corresponding Step--3 entry is
recovered from $J_{2m}$; if $b\le2m$, it is recovered from $L_{2m}$.  Thus
every Step--3 cell indexed by an unordered pair inside $D$ is recovered.  As
in \cref{prop:bc-target-window-submatrix-recovers-large-deleted-only-graph},
\cref{lem:bc-step2-submatrix-recovers-induced-graph} and
\cref{cor:bc-rank-adjacent-proper-support-deleted-only-readout} recover
$\Gamma$ from these entries and the displayed graph data.
\end{proof}

\begin{corollary}[Strict-left low-low carriers recover the large deleted-only graph]\label{cor:bc-strict-left-low-low-carrier-recovers-large-deleted-only-graph}
Work in the setting of
\cref{cor:bc-full-budget-low-low-candidate-list-recovers-large-deleted-only-graph}.
Suppose that, from
\[
  S,\quad Z_L,\quad V^\ast_L,\quad
  \rho_{2a_L-2},\rho_{2j},\quad
  \rho_0,\rho_1,\ldots,\rho_{2(j-2m)},
\]
a deterministic rule identifies a Ferrers subarrangement $E$ of the Step--3
dual-RSK' half-square of
\cref{lem:bc-krattenthaler-semistandard-path-boundary}, identifies the
partition labels on the top/right boundary of $E$, and verifies
\[
  L_{2m}\subseteq E .
\]
Then $\Gamma$ is recovered from
\[
  S,\quad \Gamma^\ast,\quad \rho(v)\in\{0,1\}^q,\quad
  \iota(A)\in\{0,1\}^q,
\]
the recovered outside datum, the carrier $E$, its recovered top/right
boundary labels, and the verified inclusion $L_{2m}\subseteq E$.
\end{corollary}

\begin{proof}
By \cref{lem:bc-tight-outside-recovers-left-growth-segment}, the recovered
outside datum determines
\[
  \rho_0,\rho_1,\ldots,\rho_{2(j-2m)} .
\]
Thus the displayed deterministic rule is applied to data already recovered
from the target-window and outside data.  It identifies the carrier $E$, its
top/right boundary labels, and the inclusion $L_{2m}\subseteq E$.

The recovered top/right boundary labels determine every Step--3 entry on $E$
by \cref{cor:bc-parsed-dual-rsk-prime-boundary-recovers-subfilling}.  Hence
they determine the true Step--3 assignment on $L_{2m}$.

By \cref{lem:bc-strict-left-recovers-continuation-columns}, the recovered
outside datum determines every Step--3 entry in the continuation column
subarrangement $J_{2m}$.  Let $\{a,b\}\subset D$ with $a<b$.  If $b>2m$,
then the corresponding Step--3 entry is determined by $J_{2m}$; if
$b\le2m$, then the corresponding cell lies in $L_{2m}$ and is determined by
the recovered carrier subfilling.  Hence every Step--3 cell indexed by an
unordered pair inside $D$ is recovered.  As in
\cref{prop:bc-target-window-submatrix-recovers-large-deleted-only-graph},
\cref{lem:bc-step2-submatrix-recovers-induced-graph} and
\cref{cor:bc-rank-adjacent-proper-support-deleted-only-readout} recover
$\Gamma$ from these entries and the displayed graph data.
\end{proof}

\begin{corollary}[Strict-left incidence-cover carriers recover the large deleted-only graph]\label{cor:bc-strict-left-incidence-cover-carrier-recovers-large-deleted-only-graph}
Work in the setting of
\cref{cor:bc-full-budget-low-low-candidate-list-recovers-large-deleted-only-graph}
and in the setting of \cref{def:bc-degree-low-low-candidates}, and also in the
non-full tight-height setting of
\cref{lem:bc-tight-outside-recovers-left-growth-segment}.  Put
\[
  M=2m,\qquad Z_L=\{j-2m+1,\ldots,j\},
\]
and assume \(m\ge29\).  Let \(A(m)\) and \(B(m)\) be the frontier functions of
\cref{prop:bc-active-twenty-low-low-profile-full-budget}.  Let \(d_a\) be the
residual low-low degrees of \cref{def:bc-degree-low-low-candidates}.  Suppose
that, from
\[
  S,\quad Z_L,\quad V^\ast_L,\quad
  \rho_{2a_L-2},\rho_{2j},\quad
  \rho_0,\rho_1,\ldots,\rho_{2(j-2m)},
\]
a deterministic rule identifies a target-window-determined set
\[
  U_{\mathrm{tw}}\subseteq\{1\le a\le2m:d_a>0\},
\]
identifies a Ferrers subarrangement \(E\) of the Step--3 dual-RSK' half-square
of \cref{lem:bc-krattenthaler-semistandard-path-boundary}, identifies the
partition labels on the top/right boundary of \(E\), and verifies
\[
  \{\{a,b\}:1\le a<b\le2m,\ \{a,b\}\cap U_{\mathrm{tw}}\ne\varnothing\}
  \subseteq E .
\]
Assume moreover that one positive-residual incidence-cover threshold holds.
Namely, put
\[
  r_i=\#\{1\le a\le2m:d_a=i\}\quad(i=1,2),\qquad
  D=\sum_{a=1}^{2m}d_a,\qquad
  D(U_{\mathrm{tw}})=\sum_{a\in U_{\mathrm{tw}}}d_a.
\]
Assume one of
\[
  |U_{\mathrm{tw}}|>r_1+r_2-A(m)
  \qquad\text{or}\qquad
  D(U_{\mathrm{tw}})>D-B(m).
\]
Then \(\Gamma\) is recovered from
\[
  S,\quad \Gamma^\ast,\quad \rho(v)\in\{0,1\}^q,\quad
  \iota(A)\in\{0,1\}^q,
\]
the recovered outside datum, the carrier \(E\), its recovered top/right
boundary labels, the incidence-cover data, and the already budgeted final-zone
word \(\beta_L\in\{0,1\}^{2m}\).
\end{corollary}

\begin{proof}
By \cref{lem:bc-tight-outside-recovers-left-growth-segment}, the recovered
outside datum determines
\[
  \rho_0,\rho_1,\ldots,\rho_{2(j-2m)} .
\]
Thus the displayed deterministic rule is applied to data already recovered
from the target-window and outside data.  It identifies
\(U_{\mathrm{tw}}\), the carrier \(E\), its top/right boundary labels, and the
displayed incidence inclusion.

The recovered top/right boundary labels determine every Step--3 entry on
\(E\) by
\cref{cor:bc-parsed-dual-rsk-prime-boundary-recovers-subfilling}.  Hence the
target-window data determine the true Step--3 entries on every low-low cell
incident to \(U_{\mathrm{tw}}\).  If
\[
  |U_{\mathrm{tw}}|>r_1+r_2-A(m),
\]
then the positivity of \(U_{\mathrm{tw}}\) gives
\[
  \#\{a\notin U_{\mathrm{tw}}:d_a>0\}
  =
  r_1+r_2-|U_{\mathrm{tw}}|
  <A(m),
\]
so \cref{cor:bc-incidence-covered-low-labels-frontier-full-budget} gives a
canonical full-budget low-low candidate set containing the true assignment.  If
instead \(D(U_{\mathrm{tw}})>D-B(m)\), then
\cref{cor:bc-large-incidence-cover-frontier-full-budget} gives the same
candidate set.
The hypotheses of
\cref{cor:bc-full-budget-low-low-candidate-list-recovers-large-deleted-only-graph}
are therefore satisfied with this candidate set, and that corollary recovers
\(\Gamma\).
\end{proof}

\begin{corollary}[Hall-augmented non-column-even carriers recover the large deleted-only graph]\label{cor:bc-hall-augmented-low-low-carrier-recovers-large-deleted-only-graph}
Work in the non-full tight-height setting of
\cref{cor:bc-tight-height-zone-endpoints-visible,lem:bc-tight-outside-recovers-left-growth-segment}
and in the setting of
\cref{cor:bc-strict-left-low-low-carrier-recovers-large-deleted-only-graph}.
Let
\[
  \lambda^{(0)}\subset\lambda^{(1)}\subset\cdots\subset\lambda^{(2m)}
\]
be the true tight-prefix Pieri path, and assume that $\lambda^{(2m)}$ is not
column-even.  Suppose that the target-window data
\[
  S,\quad Z_L,\quad V^\ast_L,\quad \mathcal D_{\mathrm{out}}(P),
  \quad \rho_{2a_L-2},\rho_{2j}
\]
determine $\lambda^{(2m)}$, the paired endpoint-defect set of
\cref{lem:bc-tight-height-continuation-parity}, the continuation-incidence
sets
\[
  N(a)=\{t\in\{2m+1,\ldots,j\}: C_t\cap P_a\ne\varnothing\},
\]
where $P_a$ are the adjacent endpoint-defect pairs and $C_t$ is the
two-column set hit by continuation label $t$, as in
\cref{lem:bc-tight-continuation-label-cover,lem:bc-continuation-incidence-semistandard-local},
and the lexicographically least canonical Hall charge of
\cref{lem:bc-canonical-continuation-hall-charge}.  Suppose that, from
\[
  S,\quad Z_L,\quad V^\ast_L,\quad
  \rho_{2a_L-2},\rho_{2j},\quad
  \rho_0,\rho_1,\ldots,\rho_{2(j-2m)},
\]
together with $\lambda^{(2m)}$, the endpoint-defect pairs, the continuation
incidences, the canonical continuation Hall charge, and the Hall-charged
continuation blocks recovered by
\cref{cor:bc-direct-hall-charge-selected-blocks-recovered}, a deterministic rule
identifies a Ferrers subarrangement $E$ of the Step--3 dual-RSK' half-square
of \cref{lem:bc-krattenthaler-semistandard-path-boundary}, identifies the
partition labels on the top/right boundary of $E$, and verifies
\[
  L_{2m}\subseteq E .
\]
Then $\Gamma$ is recovered from
\[
  S,\quad \Gamma^\ast,\quad \rho(v)\in\{0,1\}^q,\quad
  \iota(A)\in\{0,1\}^q,
\]
the recovered outside datum and these Hall-augmented carrier data.
\end{corollary}

\begin{proof}
The first recovery hypothesis makes $\lambda^{(2m)}$, the endpoint-defect
pairs, the continuation-incidence sets, and the canonical Hall charge fixed
functions of the displayed target-window data.  By
\cref{lem:bc-tight-outside-recovers-left-growth-segment}, the recovered
outside datum determines
\[
  \rho_0,\rho_1,\ldots,\rho_{2(j-2m)}.
\]
Therefore the displayed data determine the Hall-charged continuation blocks by
\cref{cor:bc-direct-hall-charge-selected-blocks-recovered}, and hence every
input to the deterministic Hall-augmented carrier rule.  Thus they determine
$E$, its top/right boundary labels, and the inclusion $L_{2m}\subseteq E$.
Applying
\cref{cor:bc-strict-left-low-low-carrier-recovers-large-deleted-only-graph}
recovers $\Gamma$.
\end{proof}

\begin{corollary}[Hall-augmented incidence-cover carriers recover the large deleted-only graph]\label{cor:bc-hall-augmented-incidence-cover-carrier-recovers-large-deleted-only-graph}
Work in the non-full tight-height setting of
\cref{cor:bc-tight-height-zone-endpoints-visible,lem:bc-tight-outside-recovers-left-growth-segment}
and in the setting of
\cref{cor:bc-strict-left-incidence-cover-carrier-recovers-large-deleted-only-graph}
and \cref{def:bc-degree-low-low-candidates}.
Let
\[
  \lambda^{(0)}\subset\lambda^{(1)}\subset\cdots\subset\lambda^{(2m)}
\]
be the true tight-prefix Pieri path, assume \(m\ge29\), and assume that
\(\lambda^{(2m)}\) is not column-even.  Suppose that the target-window data
\[
  S,\quad Z_L,\quad V^\ast_L,\quad \mathcal D_{\mathrm{out}}(P),
  \quad \rho_{2a_L-2},\rho_{2j}
\]
determine \(\lambda^{(2m)}\), the paired endpoint-defect set of
\cref{lem:bc-tight-height-continuation-parity}, the continuation-incidence
sets
\[
  N(a)=\{t\in\{2m+1,\ldots,j\}: C_t\cap P_a\ne\varnothing\},
\]
where \(P_a\) are the adjacent endpoint-defect pairs and \(C_t\) is the
two-column set hit by continuation label \(t\), as in
\cref{lem:bc-tight-continuation-label-cover,lem:bc-continuation-incidence-semistandard-local},
and the lexicographically least canonical Hall charge of
\cref{lem:bc-canonical-continuation-hall-charge}.  Let \(A(m)\) and \(B(m)\)
be the frontier functions of
\cref{prop:bc-active-twenty-low-low-profile-full-budget}.  Let \(d_a\) be the
residual low-low degrees of \cref{def:bc-degree-low-low-candidates}.  Suppose
that, from
\[
  S,\quad Z_L,\quad V^\ast_L,\quad
  \rho_{2a_L-2},\rho_{2j},\quad
  \rho_0,\rho_1,\ldots,\rho_{2(j-2m)},
\]
together with \(\lambda^{(2m)}\), the endpoint-defect pairs, the continuation
incidences, the canonical continuation Hall charge, and the Hall-charged
continuation blocks recovered by
\cref{cor:bc-direct-hall-charge-selected-blocks-recovered}, a deterministic
rule identifies a target-window-determined set
\[
  U_{\mathrm{tw}}\subseteq\{1\le a\le2m:d_a>0\},
\]
identifies a Ferrers subarrangement \(E\) of the Step--3 dual-RSK' half-square
of \cref{lem:bc-krattenthaler-semistandard-path-boundary}, identifies the
partition labels on the top/right boundary of \(E\), and verifies
\[
  \{\{a,b\}:1\le a<b\le2m,\ \{a,b\}\cap U_{\mathrm{tw}}\ne\varnothing\}
  \subseteq E .
\]
Assume moreover that one positive-residual incidence-cover threshold holds.
Namely, put
\[
  r_i=\#\{1\le a\le2m:d_a=i\}\quad(i=1,2),\qquad
  D=\sum_{a=1}^{2m}d_a,\qquad
  D(U_{\mathrm{tw}})=\sum_{a\in U_{\mathrm{tw}}}d_a.
\]
Assume one of
\[
  |U_{\mathrm{tw}}|>r_1+r_2-A(m)
  \qquad\text{or}\qquad
  D(U_{\mathrm{tw}})>D-B(m).
\]
Then \(\Gamma\) is recovered from
\[
  S,\quad \Gamma^\ast,\quad \rho(v)\in\{0,1\}^q,\quad
  \iota(A)\in\{0,1\}^q,
\]
the recovered outside datum and these Hall-augmented incidence-cover carrier
data.
\end{corollary}

\begin{proof}
The first recovery hypothesis makes \(\lambda^{(2m)}\), the endpoint-defect
pairs, the continuation-incidence sets, and the canonical Hall charge fixed
functions of the displayed target-window data.  By
\cref{lem:bc-tight-outside-recovers-left-growth-segment}, the recovered
outside datum determines
\[
  \rho_0,\rho_1,\ldots,\rho_{2(j-2m)}.
\]
Therefore the displayed data determine the Hall-charged continuation blocks by
\cref{cor:bc-direct-hall-charge-selected-blocks-recovered}, and hence every
input to the deterministic Hall-augmented incidence-cover carrier rule.  Thus
they determine \(U_{\mathrm{tw}}\), \(E\), its top/right boundary labels, the
incidence inclusion, and the displayed positive-residual threshold.  Applying
\cref{cor:bc-strict-left-incidence-cover-carrier-recovers-large-deleted-only-graph}
recovers \(\Gamma\).
\end{proof}

\begin{corollary}[Continuation segments supply the Hall-augmented graph carrier]\label{cor:bc-continuation-segment-hall-carrier-recovers-large-deleted-only-graph}
Work in the setting of
\cref{cor:bc-hall-augmented-low-low-carrier-recovers-large-deleted-only-graph}.
Suppose that the target-window data
\[
  S,\quad Z_L,\quad V^\ast_L,\quad \mathcal D_{\mathrm{out}}(P),
  \quad \rho_{2a_L-2},\rho_{2j}
\]
determine the Pieri continuation segment
\[
  \lambda^{(2m)}\subset\lambda^{(2m+1)}\subset\cdots\subset\lambda^{(j)} .
\]
Using this recovered segment, form the endpoint-defect pairs, the
continuation incidences, and the canonical continuation Hall charge of
\cref{lem:bc-pieri-continuation-determines-hall-charge}, and recover the
Hall-charged continuation blocks of
\cref{cor:bc-hall-charge-selected-blocks-recovered}.  Suppose that, from
\[
  S,\quad Z_L,\quad V^\ast_L,\quad
  \rho_{2a_L-2},\rho_{2j},\quad
  \rho_0,\rho_1,\ldots,\rho_{2(j-2m)},
\]
together with the recovered continuation segment, its canonical Hall data, and
the Hall-charged continuation blocks,
a deterministic rule identifies a Ferrers subarrangement $E$ of the Step--3
dual-RSK' half-square of
\cref{lem:bc-krattenthaler-semistandard-path-boundary}, identifies the
partition labels on the top/right boundary of $E$, and verifies
\[
  L_{2m}\subseteq E .
\]
Then $\Gamma$ is recovered from
\[
  S,\quad \Gamma^\ast,\quad \rho(v)\in\{0,1\}^q,\quad
  \iota(A)\in\{0,1\}^q,
\]
the recovered outside datum, the recovered continuation segment, and this
carrier rule.
\end{corollary}

\begin{proof}
The first partition in the recovered continuation segment is
$\lambda^{(2m)}$.  For each continuation label $t=2m+1,\ldots,j$, the
difference
\[
  \lambda^{(t)}\setminus\lambda^{(t-1)}
\]
is exactly the pair of boxes carrying label $t$ by
\cref{lem:bc-pieri-path-bijection}.  Thus the recovered continuation segment
supplies the prefix endpoint and the continuation label boxes.

By \cref{lem:bc-pieri-continuation-determines-hall-charge}, the same segment
determines the endpoint-defect pairs, the continuation incidences, and the
canonical Hall charge.  Moreover,
\cref{cor:bc-hall-charge-selected-blocks-recovered,cor:bc-direct-hall-charge-selected-blocks-recovered}
identify the Hall-charged continuation blocks as fixed functions of the
canonical Hall charge and the recovered outside datum.  Therefore the
deterministic carrier rule in the present statement supplies the
Hall-augmented carrier rule required by
\cref{cor:bc-hall-augmented-low-low-carrier-recovers-large-deleted-only-graph}.
Applying that corollary recovers $\Gamma$.
\end{proof}

\begin{corollary}[Internal Q-cuts supply the Hall-augmented graph carrier]\label{cor:bc-q-suffix-hall-carrier-recovers-large-deleted-only-graph}
Work in the setting of
\cref{cor:bc-continuation-segment-hall-carrier-recovers-large-deleted-only-graph}.
Let $C_{2m}^+$ be the Step~2 $Q$-suffix strip of
\cref{lem:bc-step2-q-suffix-strip-determines-continuation}, and let $C_{2m}$
be the corresponding Step~2 $Q$-prefix strip.  Suppose that the
target-window data
\[
  S,\quad Z_L,\quad V^\ast_L,\quad \mathcal D_{\mathrm{out}}(P),
  \quad \rho_{2a_L-2},\rho_{2j}
\]
determine the ordered partition labels on the internal separating boundary
between the Step~2 $Q$-prefix strip $C_{2m}$ and the suffix strip
$C_{2m}^+$.  Use
\cref{cor:bc-strict-left-q-suffix-boundary-recovers-continuation-segment} to
recover the Pieri continuation segment and then form its canonical Hall data
and Hall-charged continuation blocks.  Suppose that, from
\[
  S,\quad Z_L,\quad V^\ast_L,\quad
  \rho_{2a_L-2},\rho_{2j},\quad
  \rho_0,\rho_1,\ldots,\rho_{2(j-2m)},
\]
together with the recovered continuation segment, its canonical Hall data, and
the Hall-charged continuation blocks, a deterministic rule identifies a Ferrers
subarrangement $E$ of the Step~3 dual-RSK' half-square of
\cref{lem:bc-krattenthaler-semistandard-path-boundary}, identifies the
partition labels on the top/right boundary of $E$, and verifies
\[
  L_{2m}\subseteq E .
\]
Then $\Gamma$ is recovered from
\[
  S,\quad \Gamma^\ast,\quad \rho(v)\in\{0,1\}^q,\quad
  \iota(A)\in\{0,1\}^q,
\]
the recovered outside datum, the recovered internal $Q$-cut, and this
Hall-augmented carrier rule.
\end{corollary}

\begin{proof}
\Cref{cor:bc-strict-left-q-suffix-boundary-recovers-continuation-segment}
recovers the Pieri continuation segment from the displayed target-window data
and the internal cut.  The remaining hypotheses are then exactly the
continuation-segment Hall-carrier hypotheses of
\cref{cor:bc-continuation-segment-hall-carrier-recovers-large-deleted-only-graph}.
Applying that corollary recovers $\Gamma$.
\end{proof}

\begin{lemma}[Column-parity defects are local to deletion intervals]\label{lem:bc-column-defects-per-joint}
Let $T$ be a semistandard tableau of content $(2^j)$ whose shape has only even
column lengths.  Let $S\subset\{1,\ldots,j\}$, and write the maximal
consecutive intervals in $S$ as
\[
  J_\alpha=\{a_\alpha,a_\alpha+1,\ldots,b_\alpha\},
  \qquad 1\le\alpha\le L .
\]
For each interval put
\[
  O_\alpha(T)=
  \{c:\#\{\text{boxes of }T\text{ in column }c
       \text{ whose label lies in }J_\alpha\}\text{ is odd}\}.
\]
Then
\[
  |O_\alpha(T)|\le 2|J_\alpha|
\]
for every $\alpha$.  If
\[
  O_{\mathrm{ret}}(T,S)=
  \{c:\#\{\text{boxes of }T\setminus S\text{ in column }c\}
       \text{ is odd}\},
\]
then
\[
  O_{\mathrm{ret}}(T,S)=
  O_1(T)\triangle O_2(T)\triangle\cdots\triangle O_L(T),
\]
where $\triangle$ denotes symmetric difference.  Consequently, every odd
column of the retained diagram $T\setminus S$ lies in at least one
$O_\alpha(T)$, and
the final column-parity defect is controlled by at most
\[
  \sum_{\alpha=1}^L2|J_\alpha|=2|S|
\]
joint-local column candidates.
\end{lemma}

\begin{proof}
For a fixed $\alpha$, the interval $J_\alpha$ contains $|J_\alpha|$ labels,
and each label occurs exactly twice in $T$.  Hence the total number of boxes
whose label lies in $J_\alpha$ is $2|J_\alpha|$.  A column can belong to
$O_\alpha(T)$ only if it contains at least one of these boxes, so
$|O_\alpha(T)|\le2|J_\alpha|$.

Let $d_{\alpha,c}$ be the number of boxes in column $c$ whose label lies in
$J_\alpha$.  Since the original shape of $T$ has even column lengths, the
parity of column $c$ after deleting all labels in $S$ is
\[
  \sum_{\alpha=1}^L d_{\alpha,c}\pmod 2.
\]
If this parity is odd, then at least one summand $d_{\alpha,c}$ is odd, and
therefore $c\in O_\alpha(T)$ for that $\alpha$.  More precisely, the parity is
odd exactly when $c$ belongs to an odd number of the sets $O_\alpha(T)$, which
is the stated symmetric-difference identity.  Summing the preceding size bounds
over $\alpha$ gives the final estimate.
\end{proof}

\begin{lemma}[Joint endpoints determine the parity defects]\label{lem:bc-joint-endpoint-defects}
Let $T$ be a semistandard tableau of content $(2^j)$, and let
\[
  \varnothing=\lambda^{(0)}\subset\lambda^{(1)}\subset\cdots\subset
  \lambda^{(j)}=\lambda(T)
\]
be its Pieri two-strip path.  Let $U=\operatorname{std}(T)$, and let
\[
  \varnothing=\sigma^{(0)}\subset\sigma^{(1)}\subset\cdots\subset
  \sigma^{(2j)}=\lambda(T)
\]
be the standard shape chain of $U$.  Let
$J_\alpha=\{a_\alpha,a_\alpha+1,\ldots,b_\alpha\}$ be one of the maximal
consecutive intervals in a subset $S\subset\{1,\ldots,j\}$, and define
$O_\alpha(T)$ as in \cref{lem:bc-column-defects-per-joint}.  Then
\[
  O_\alpha(T)=
  \{c\ge1:\,
    (\lambda^{(b_\alpha)})'_c
     -(\lambda^{(a_\alpha-1)})'_c
     \equiv 1 \pmod 2\}.
\]
Equivalently,
\[
  O_\alpha(T)=
  \{c\ge1:\,
    (\sigma^{(2b_\alpha)})'_c
     -(\sigma^{(2a_\alpha-2)})'_c
     \equiv 1 \pmod 2\}.
\]
Here missing column lengths are interpreted as $0$.
In particular, the local parity-defect candidates at the joint are determined
by the two endpoint partitions of the deleted standard-time component
\[
  \{2a_\alpha-1,2a_\alpha,\ldots,2b_\alpha\}.
\]
Moreover $|O_\alpha(T)|$ is even.
\end{lemma}

\begin{proof}
By \cref{lem:bc-pieri-path-bijection}, the boxes whose labels lie in
$J_\alpha$ are exactly the skew difference
$\lambda^{(b_\alpha)}/\lambda^{(a_\alpha-1)}$.  Hence, for each column $c$,
the number of boxes in column $c$ whose labels lie in $J_\alpha$ is
\[
  (\lambda^{(b_\alpha)})'_c-(\lambda^{(a_\alpha-1)})'_c .
\]
Taking this equality modulo $2$ gives the displayed description of
$O_\alpha(T)$.
By \cref{lem:bc-standard-time-pair-deletion},
\[
  \sigma^{(2b_\alpha)}=\lambda^{(b_\alpha)}
  \quad\text{and}\quad
  \sigma^{(2a_\alpha-2)}=\lambda^{(a_\alpha-1)}.
\]
Substitution gives the equivalent standard-time formula.

The total number of boxes in the deleted interval is
$2|J_\alpha|$, which is even.  This total is the sum, over all columns, of the
displayed column differences.  Therefore the number of columns for which the
difference is odd is even.
\end{proof}

\begin{corollary}[Joint parity defects have canonical pairs]\label{cor:bc-joint-defect-pairs}
In the setting of \cref{lem:bc-joint-endpoint-defects}, write
\[
  O_\alpha(T)=\{o_1<o_2<\cdots<o_{2r_\alpha}\}.
\]
Then $r_\alpha\le |J_\alpha|$, and the pairs
\[
  \{o_1,o_2\},\ \{o_3,o_4\},\ \ldots,\ \{o_{2r_\alpha-1},o_{2r_\alpha}\}
\]
are determined by the two endpoint partitions
$\sigma^{(2a_\alpha-2)}$ and $\sigma^{(2b_\alpha)}$ of the deleted
standard-time component.
\end{corollary}

\begin{proof}
\Cref{lem:bc-joint-endpoint-defects} proves that $|O_\alpha(T)|$ is even, so
the displayed increasing enumeration and adjacent pairing are defined.  The
same lemma identifies $O_\alpha(T)$ from the two standard-time endpoint
partitions, hence the increasing order and the adjacent pairing are also
determined by those endpoints.  Finally,
\cref{lem:bc-column-defects-per-joint} gives
$|O_\alpha(T)|\le2|J_\alpha|$, so $r_\alpha\le|J_\alpha|$.
\end{proof}

\begin{corollary}[Frontier events cover paired joint defects]\label{cor:bc-joint-events-cover-defects}
In either the width or height setting of
\cref{lem:bc-frontier-events-per-joint}, let
\[
  E_\alpha(T)=\{e_1<e_2<\cdots<e_{h_\alpha}\}
\]
be the ordered frontier-event set at $J_\alpha$, and let
\[
  \mathcal O_\alpha(T)=\{\{o_1,o_2\},\ldots,
      \{o_{2r_\alpha-1},o_{2r_\alpha}\}\}
\]
be the paired defect set of \cref{cor:bc-joint-defect-pairs}.  Then
$h_\alpha=|J_\alpha|$ and $r_\alpha\le h_\alpha$.  Hence the rule
\[
  \{o_{2u-1},o_{2u}\}\longmapsto e_u\qquad(1\le u\le r_\alpha)
\]
is an order-preserving injection from the paired defect slots at the joint
into the frontier-event slots at the same joint.
\end{corollary}

\begin{proof}
\Cref{lem:bc-frontier-events-per-joint} gives
$|E_\alpha(T)|=|J_\alpha|$, so $h_\alpha=|J_\alpha|$.
\Cref{cor:bc-joint-defect-pairs} gives $r_\alpha\le |J_\alpha|$.
Therefore $r_\alpha\le h_\alpha$, and the displayed assignment is well
defined.  It is order-preserving by construction from the increasing
enumerations.
\end{proof}

\begin{lemma}[Local branch words encode orientation and selected defect slots]\label{lem:bc-local-branch-word}
In either the width or height setting of
\cref{lem:bc-frontier-events-per-joint}, fix one deletion interval
$J_\alpha$ and put $h_\alpha=|J_\alpha|$.  Let
\[
  E_\alpha(T)=\{e_1<e_2<\cdots<e_{h_\alpha}\}
\]
be its ordered frontier-event set, and let
$\mathcal O_\alpha(T)$ be the paired defect set of
\cref{cor:bc-joint-defect-pairs}, matched to the first
$r_\alpha$ event slots by \cref{cor:bc-joint-events-cover-defects}.

In the width case, let $\theta_i$ be the bit $\epsilon^w_r(T)$ attached by
\cref{lem:bc-standardized-witness-pairs} to the unique witness label whose
frontier event is $e_i$.  In the height case, put $\theta_i=0$ for every
$i$.  Then, for every subset
$M_\alpha\subseteq\{1,\ldots,r_\alpha\}$ of paired defect slots, the data
\[
  \bigl((\theta_1,\ldots,\theta_{h_\alpha}),M_\alpha\bigr)
\]
inject into $\{0,1\}^{2h_\alpha}$ by using the first $h_\alpha$ bits for
the word $\theta$ and the last $h_\alpha$ bits for the characteristic function
of $M_\alpha$ inside the matched event slots.
\end{lemma}

\begin{proof}
The ordered event list has length $h_\alpha$ by
\cref{lem:bc-frontier-events-per-joint}.  In the width case each event belongs
to a unique witness label in $J_\alpha$, so
\cref{lem:bc-standardized-witness-pairs} supplies one bit $\theta_i$ for that
event.  In the height case the same lemma forces the frontier copy to be the
lower member of the adjacent standard pair, so the all-zero word records the
orientation data.

By \cref{cor:bc-joint-events-cover-defects}, the paired defect slots are
matched injectively to the first $r_\alpha\le h_\alpha$ ordered event slots.
Thus the characteristic function of any subset $M_\alpha$ extends uniquely by
zeros to a word of length $h_\alpha$.  Concatenating this selection word with
$\theta$ gives an element of $\{0,1\}^{2h_\alpha}$.  The two halves of the
concatenated word recover $\theta$ and $M_\alpha$, so the encoding is
injective.
\end{proof}

\begin{lemma}[Branch words recover the local repair data]\label{lem:bc-local-branch-recovery}
In the setting of \cref{lem:bc-local-branch-word}, write
$J_\alpha=\{a_\alpha,a_\alpha+1,\ldots,b_\alpha\}$ and
$h_\alpha=|J_\alpha|$.  Let
\[
  \beta_\alpha=(\beta_1,\ldots,\beta_{2h_\alpha})
  \in\{0,1\}^{2h_\alpha}
\]
be the branch word produced there from the orientation word and a subset
$M_\alpha$ of paired defect slots.  Put
\[
  \sigma^-=\sigma^{(2a_\alpha-2)},\qquad
  \sigma^+=\sigma^{(2b_\alpha)} .
\]
Then $J_\alpha$, $\sigma^-$, $\sigma^+$, and $\beta_\alpha$ determine the
following local data:
\begin{enumerate}
\item the ordered frontier-event list
\[
  E_\alpha(T)=\{e_1<\cdots<e_{h_\alpha}\},
\]
with
\[
  e_i=2(a_\alpha+i-1)-1+\beta_i
  \quad\text{in the width case,}
\]
and
\[
  e_i=2(a_\alpha+i-1)-1
  \quad\text{in the height case;}
\]
\item the paired local defect set, obtained by listing increasingly the
columns
\[
  O_\alpha(T)=\{c\ge1:\,(\sigma^+)'_c-(\sigma^-)'_c\equiv1\pmod2\}
\]
and pairing adjacent entries;
\item the selected subset $M_\alpha$, by reading the last $h_\alpha$ bits of
$\beta_\alpha$ on the matched defect slots of
\cref{cor:bc-joint-events-cover-defects}.
\end{enumerate}
Consequently, after the branch word is fixed, the only missing input for
\cref{lem:bc-local-branch-hkko-marks} is an actual zero-free vacillating
segment and an injection of the selected paired defect slots into its eligible
boundary-up steps.
\end{lemma}

\begin{proof}
The labels in the interval $J_\alpha$ occur, in increasing order, as
$a_\alpha,a_\alpha+1,\ldots,b_\alpha$.  In the width case the first
$h_\alpha$ bits of $\beta_\alpha$ are the orientation bits $\theta_i$ of
\cref{lem:bc-local-branch-word}.  The event attached to the label
$a_\alpha+i-1$ is therefore
$2(a_\alpha+i-1)-1+\theta_i$ by
\cref{lem:bc-standard-frontier-time}.  Since $\theta_i=\beta_i$, this gives
the displayed width formula.  In the height case
\cref{lem:bc-standard-frontier-time} gives the lower time
$2(a_\alpha+i-1)-1$, and the first half of the branch word is the all-zero
word by \cref{lem:bc-local-branch-word}.  The ordering and the fact that these
are exactly the events of $E_\alpha(T)$ are supplied by
\cref{lem:bc-frontier-events-per-joint}.

\Cref{lem:bc-joint-endpoint-defects} identifies $O_\alpha(T)$ from the two
endpoint partitions $\sigma^-$ and $\sigma^+$.  Then
\cref{cor:bc-joint-defect-pairs} determines the adjacent pairing.  Finally,
\cref{lem:bc-local-branch-word} encodes $M_\alpha$ in the second half of
$\beta_\alpha$ after the matched event-slot injection of
\cref{cor:bc-joint-events-cover-defects}; reading those bits recovers exactly
$M_\alpha$.  The last sentence is just the remaining hypothesis in
\cref{lem:bc-local-branch-hkko-marks} after the recovered data have been
fixed.
\end{proof}

\begin{lemma}[Local branch words give admissible HKKO markings]\label{lem:bc-local-branch-hkko-marks}
In the setting of \cref{lem:bc-local-branch-word}, suppose a joint repair
construction also supplies a zero-free $w$-vacillating path
\[
  V=(V^0,\ldots,V^N)
\]
and an injection $\iota_\alpha$ from the paired defect slots
$\{1,\ldots,r_\alpha\}$ into the eligible boundary-up set
\[
  B(V)=\{t:V^t-V^{t-1}=+\epsilon_w\text{ and }V^{t-1}_w=0\}.
\]
Then, for every subset $M_\alpha\subseteq\{1,\ldots,r_\alpha\}$, the pair
\[
  \bigl(V,\iota_\alpha(M_\alpha)\bigr)
\]
is an admissible element of the HKKO marked class
$\mathrm{MVT}^*_N(w;V^0\to V^N)$.  For fixed $V$ and $\iota_\alpha$, different
subsets $M_\alpha$ give different markings, and the second half of the local
branch word of \cref{lem:bc-local-branch-word} recovers $M_\alpha$.
\end{lemma}

\begin{proof}
By hypothesis $V$ has no zero steps, and every element in the image of
$\iota_\alpha$ is a step $+\epsilon_w$ whose starting point has $w$-th
coordinate $0$.  Therefore Definition 8.6 of \cite{HKKO2025}, recorded in
\cref{lem:bc-hkko-marked-boundary}, says that
$(V,\iota_\alpha(M_\alpha))$ belongs to
$\mathrm{MVT}^*_N(w;V^0\to V^N)$.

Since $\iota_\alpha$ is injective, distinct subsets of
$\{1,\ldots,r_\alpha\}$ have distinct images.  Finally,
\cref{lem:bc-local-branch-word} encodes $M_\alpha$ as the last
$h_\alpha$ bits of the branch word, extended by zeros outside the matched
defect slots.  Hence that half of the branch word recovers $M_\alpha$.
\end{proof}

\begin{lemma}[Frontier-event eligibility supplies the local marks]\label{lem:bc-frontier-events-supply-marks}
In the setting of \cref{lem:bc-local-branch-recovery}, suppose a joint repair
construction supplies a zero-free $w$-vacillating segment
\[
  V=(V^0,\ldots,V^N)
\]
and an order-preserving assignment of the recovered frontier-event slots
$e_1<\cdots<e_{h_\alpha}$ to distinct local times
\[
  \tau_1<\cdots<\tau_{h_\alpha}
  \quad\text{in }\{1,\ldots,N\}.
\]
Let $r_\alpha$ be the number of paired defect slots.  If
\[
  \tau_u\in B(V)\qquad(1\le u\le r_\alpha),
\]
then the rule $u\mapsto\tau_u$ is the required injection of paired defect
slots into eligible boundary-up steps in
\cref{lem:bc-local-branch-hkko-marks}.  Consequently every selected subset
$M_\alpha$ encoded in the second half of the branch word gives an admissible
HKKO marking, and $M_\alpha$ is recovered from that half of the branch word.
\end{lemma}

\begin{proof}
The times $\tau_1,\ldots,\tau_{h_\alpha}$ are distinct by hypothesis.  Since
$r_\alpha\le h_\alpha$ by \cref{cor:bc-joint-events-cover-defects}, the rule
$u\mapsto\tau_u$ is an injection on $\{1,\ldots,r_\alpha\}$.  The displayed
eligibility condition says that its image lies in $B(V)$.  Therefore the
hypotheses of \cref{lem:bc-local-branch-hkko-marks} hold with
$\iota_\alpha(u)=\tau_u$.  That lemma gives the admissible marking statement
and recovery of the selected subset from the second half of the branch word.
\end{proof}

\begin{corollary}[Fixed-witness defect slots use exactly the local branch budget]\label{cor:bc-fixed-witness-defect-slots-budgeted}
Let $S$ be a fixed witness set and decompose it into maximal consecutive
label intervals
\[
  J_1,\ldots,J_L .
\]
For each $J_\alpha$, let $h_\alpha=|J_\alpha|$ and let
\[
  \beta_\alpha\in\{0,1\}^{2h_\alpha}
\]
be the local branch word of \cref{lem:bc-local-branch-word}.  Then the
endpoint partitions at the deletion joint and the word $\beta_\alpha$
determine the paired local odd-column defect slots and the selected subset of
those slots.  The number of paired slots at $J_\alpha$ is at most
$h_\alpha$, and the selected subset uses only the last $h_\alpha$ bits of
$\beta_\alpha$.

Moreover, if the local repair segment realizes the recovered ordered
frontier-event list by distinct local times whose first paired-defect slots
are eligible boundary-up steps as in
\cref{lem:bc-frontier-events-supply-marks}, then the selected subset gives
admissible HKKO marks with no branch data beyond $\beta_\alpha$.  Across all
deletion joints the total branch budget is
\[
  \prod_{\alpha=1}^L 2^{2h_\alpha}=2^{2|S|}.
\]
\end{corollary}

\begin{proof}
For a fixed joint, \cref{lem:bc-local-branch-recovery} recovers the ordered
frontier-event list, the paired local defect set, and the selected subset from
the joint endpoints and $\beta_\alpha$.  The same lemma obtains the paired
defect set by applying \cref{lem:bc-joint-endpoint-defects} and then
\cref{cor:bc-joint-defect-pairs}; \cref{cor:bc-joint-events-cover-defects}
gives at most $h_\alpha$ paired slots and matches them to the first
$r_\alpha\le h_\alpha$ ordered frontier events.  The construction of
\cref{lem:bc-local-branch-word} stores the selected subset in the last
$h_\alpha$ bits, so no other branch source is used for the defect selection.

Under the stated realization hypothesis,
\cref{lem:bc-frontier-events-supply-marks} turns those selected slots into
admissible HKKO markings and recovers the selected subset from the same last
half of $\beta_\alpha$.  Finally, the maximal intervals $J_\alpha$ partition
$S$, so
\[
  \sum_{\alpha=1}^L h_\alpha=|S|.
\]
Multiplying the local branch counts gives the displayed total budget.
\end{proof}

\begin{lemma}[Decoded branch data suffice for the local $Z_0$ splice]\label{lem:bc-local-z0-splice-package}
In the setting of \cref{lem:bc-local-branch-recovery}, let
\[
  \beta_\alpha\in\{0,1\}^{2h_\alpha}
\]
be a local branch word, and let $M_\alpha\subseteq\{1,\ldots,r_\alpha\}$ be
the selected paired-defect subset recovered from its second half.  Suppose a
joint repair construction supplies a $w$-up-down tableau
\[
  U=(U^0,U^1,\ldots,U^{2N})
\]
from $\mu$ to $\nu$, together with distinct indices
\[
  \tau_1,\ldots,\tau_{h_\alpha}\in[N]
\]
assigned to the ordered frontier-event slots.  Put
\[
  M^\tau_\alpha=\{\tau_u:u\in M_\alpha\}.
\]
If $M^\tau_\alpha=\varnothing$, set $U^\ast=U$.  If
$M^\tau_\alpha\ne\varnothing$, assume
\[
  \ell(U^r)<w\qquad(0\le r\le2N),
\]
\[
  u\in M_\alpha\quad\Longrightarrow\quad
  \ell(U^{2\tau_u-1})=w-1,
\]
and
\[
  m\notin M^\tau_\alpha\quad\Longrightarrow\quad
  \neg\bigl(\ell(U^{2m-2})<\ell(U^{2m-1})>\ell(U^{2m})\bigr).
\]
Then the HKKO boundary-event splice produces an unmarked $w$-up-down tableau
$U^\ast$ with endpoints $\mu,\nu$ and length $2N$.  Moreover the operation is
reversible from $U^\ast$, the branch word $\beta_\alpha$, and the assigned
indices $\tau_1,\ldots,\tau_{h_\alpha}$.
\end{lemma}

\begin{proof}
By \cref{lem:bc-local-branch-recovery}, the branch word
$\beta_\alpha$ recovers the selected subset $M_\alpha$.  Since the assigned
indices $\tau_1,\ldots,\tau_{h_\alpha}$ are distinct, the set
$M^\tau_\alpha=\{\tau_u:u\in M_\alpha\}$ is then also recovered from
$\beta_\alpha$ and the assignment $\tau$.

If $M^\tau_\alpha=\varnothing$, the identity operation has the asserted
endpoint, length, and reversibility properties.  If
$M^\tau_\alpha\ne\varnothing$, the three displayed hypotheses are exactly the
conditions of \cref{lem:bc-z0-landing-criterion} with
$M=M^\tau_\alpha$.  That lemma applies
\cref{lem:bc-hkko-boundary-splice} and gives an unmarked output $U^\ast$ with
the same endpoints and length, reversible from $U^\ast$ and
$M^\tau_\alpha$.  Since $M^\tau_\alpha$ is recovered from
$\beta_\alpha$ and the assigned indices, the claimed reversibility follows.
\end{proof}

\begin{lemma}[Deleted intervals have canonical HKKO event indices]\label{lem:bc-canonical-hkko-event-indices}
In the setting of \cref{lem:bc-local-branch-recovery}, write
$J_\alpha=\{a_\alpha,\ldots,b_\alpha\}$ and
$h_\alpha=|J_\alpha|$.  Normalize the deleted standard-time component
\[
  \{2a_\alpha-1,2a_\alpha,\ldots,2b_\alpha\}
\]
to the local interval $\{1,\ldots,2h_\alpha\}$ by subtracting
$2a_\alpha-2$.  Then the adjacent standard pair attached to the label
$a_\alpha+i-1$ is the local pair
\[
  \{2i-1,2i\}\qquad(1\le i\le h_\alpha).
\]
Thus, for a local HKKO segment whose $h_\alpha$ up-down pairs are indexed by
these deleted adjacent pairs in label order, the canonical event-index
assignment is
\[
  \tau_i=i\qquad(1\le i\le h_\alpha).
\]
This assignment is determined by $J_\alpha$ alone.  Consequently, in
\cref{lem:bc-local-z0-splice-package} for such a local segment one has
$M^\tau_\alpha=M_\alpha$, and the HKKO splice is reversible from the output,
the fixed interval $J_\alpha$, and the branch word $\beta_\alpha$.
\end{lemma}

\begin{proof}
The standard-time component associated with $J_\alpha$ is
$\{2a_\alpha-1,\ldots,2b_\alpha\}$ by
\cref{lem:bc-frontier-events-per-joint}.  For the $i$-th label in the
interval, namely $a_\alpha+i-1$, the adjacent standardized pair is
\[
  \{2(a_\alpha+i-1)-1,\ 2(a_\alpha+i-1)\}.
\]
Subtracting $2a_\alpha-2$ gives $\{2i-1,2i\}$.  Hence the natural local
up-down pair index of that event slot is $i$, giving $\tau_i=i$.  Since
$J_\alpha$ gives $a_\alpha$ and $h_\alpha$, no additional data are needed to
recover this assignment.

With $\tau_i=i$, the selected HKKO index set
$M^\tau_\alpha=\{\tau_u:u\in M_\alpha\}$ equals $M_\alpha$.  The final
reversibility assertion is therefore the reversibility assertion of
\cref{lem:bc-local-z0-splice-package}, with the assigned indices recovered
from the fixed interval $J_\alpha$.
\end{proof}

\begin{lemma}[Top-wall peak-loop segment]\label{lem:bc-top-wall-peak-loop}
Let $w\ge3$, let $N\ge0$, and let $\tau$ be a partition with exactly
$w-2$ nonzero parts.  For a subset $M\subseteq[N]$, define a sequence
\[
  U=(U^0,U^1,\ldots,U^{2N})
\]
by setting $U^{2m}=\tau$ for every $0\le m\le N$ and, for
$1\le m\le N$, setting $U^{2m-1}$ to be obtained from $\tau$ by adding one
box in row $w-1$ if $m\in M$, and by adding one box in row $1$ if
$m\notin M$.  Then $U$ is a $w$-up-down tableau from $\tau$ to $\tau$,
\[
  \ell(U^r)<w\qquad(0\le r\le2N),
\]
and its odd length-peak set is exactly $M$:
\[
  \{m\in[N]:\ell(U^{2m-2})<\ell(U^{2m-1})>\ell(U^{2m})\}=M.
\]
Moreover, every selected midpoint has length $w-1$ and is obtained by adding
one box in row $w-1$ from a partition with exactly $w-2$ nonzero parts.
\end{lemma}

\begin{proof}
Since $\tau$ has exactly $w-2$ nonzero parts, adding one box in row $w-1$
creates a partition of length $w-1$, and adding one box in row $1$ creates a
partition still having length $w-2$.  In each two-step block the first step
adds one box to $\tau$ and the second removes that same box, so the sequence is
a $w$-up-down tableau from $\tau$ back to $\tau$.

All even states have length $w-2$.  If $m\in M$, then the odd state
$U^{2m-1}$ has length $w-1$, so
\[
  \ell(U^{2m-2})<\ell(U^{2m-1})>\ell(U^{2m}).
\]
If $m\notin M$, then the odd state has length $w-2$, equal to the adjacent
even-state lengths, so there is no length peak at $2m-1$.  This proves the
displayed equality for the odd length-peak set.  The below-wall bound follows
from the lengths $w-2$ and $w-1$, both strictly less than $w$.  The final
assertion is the construction in the selected case.
\end{proof}

\begin{lemma}[Below-wall peakless connectors]\label{lem:bc-below-wall-peakless-connector}
Let $w\ge1$, and let $\mu,\nu$ be partitions with
\[
  \ell(\mu)<w,\qquad \ell(\nu)<w.
\]
Then there is a $w$-up-down tableau
\[
  C=(C^0,\ldots,C^{2R})
\]
from $\mu$ to $\nu$ such that
\[
  \ell(C^r)<w\qquad(0\le r\le2R)
\]
and $C$ has no odd length peaks:
\[
  \neg\bigl(\ell(C^{2a-2})<\ell(C^{2a-1})>\ell(C^{2a})\bigr)
  \qquad(1\le a\le R).
\]
Moreover $C$ can be chosen from $\mu$ and $\nu$ alone.
\end{lemma}

\begin{proof}
Put $\zeta=\mu\vee\nu$, the componentwise maximum of $\mu$ and $\nu$.
Then $\zeta$ is a partition and $\ell(\zeta)<w$.  Starting from $\mu$, add
the boxes of $\zeta/\mu$ one at a time, row by row from top to bottom, to
obtain a chain of partitions from $\mu$ to $\zeta$.  Starting from $\zeta$,
remove the boxes of $\zeta/\nu$ one at a time, row by row from bottom to top,
to obtain a chain from $\zeta$ to $\nu$.  Concatenating these two chains gives
partitions
\[
  \gamma_0=\mu,\gamma_1,\ldots,\gamma_R=\nu
\]
all below the wall, such that each adjacent pair is comparable and differs by
one box.

If $\gamma_{a-1}\subseteq\gamma_a$, set
\[
  (C^{2a-2},C^{2a-1},C^{2a})
  =
  (\gamma_{a-1},\gamma_a,\gamma_a).
\]
If $\gamma_a\subseteq\gamma_{a-1}$, set
\[
  (C^{2a-2},C^{2a-1},C^{2a})
  =
  (\gamma_{a-1},\gamma_{a-1},\gamma_a).
\]
The overlaps agree at the even times, so this defines a $w$-up-down tableau
from $\mu$ to $\nu$.  Each nontrivial adjacent difference is a single box and
hence a vertical strip; the other adjacent difference in the same two-step
block is empty.  Since every displayed state is one of the $\gamma_a$, all
states remain below the wall.

In an increasing two-step block the midpoint equals the right endpoint, so
its length cannot be strictly larger than both endpoint lengths.  In a
decreasing two-step block the midpoint equals the left endpoint, with the same
conclusion.  Thus no odd midpoint is a length peak.  The construction used
only $\mu$ and $\nu$.
\end{proof}

\begin{lemma}[Peakless connectors attach the top-wall loop]\label{lem:bc-peakless-connector-attachment}
Let $w\ge3$.  Let
\[
  A=(A^0,\ldots,A^{2R}),\qquad
  L=(L^0,\ldots,L^{2N}),\qquad
  B=(B^0,\ldots,B^{2S})
\]
be $w$-up-down tableaux with
\[
  A^{2R}=L^0,\qquad L^{2N}=B^0.
\]
Assume all three segments stay below the wall:
\[
  \ell(A^r)<w,\qquad \ell(L^r)<w,\qquad \ell(B^r)<w
\]
for all admissible $r$.  Assume also that $A$ and $B$ have no odd length
peaks:
\[
  \neg\bigl(\ell(A^{2a-2})<\ell(A^{2a-1})>\ell(A^{2a})\bigr)
  \quad(1\le a\le R),
\]
and
\[
  \neg\bigl(\ell(B^{2b-2})<\ell(B^{2b-1})>\ell(B^{2b})\bigr)
  \quad(1\le b\le S).
\]
Let
\[
  M_L=\{\ell_0\in[N]:
        \ell(L^{2\ell_0-2})<\ell(L^{2\ell_0-1})>\ell(L^{2\ell_0})\}
\]
be the odd peak set of $L$.  Form the even-time concatenation
\[
  W=(A^0,\ldots,A^{2R},L^1,\ldots,L^{2N},B^1,\ldots,B^{2S}).
\]
Then $W$ is a $w$-up-down tableau from $A^0$ to $B^{2S}$, it stays below the
wall, and its odd length-peak set is exactly
\[
  \{R+\ell_0:\ell_0\in M_L\}.
\]
\end{lemma}

\begin{proof}
Even-time concatenation of up-down tableaux with matching endpoints is again
an up-down tableau, with start $A^0$ and end $B^{2S}$.  Every state of $W$ is
a state of one of the three input segments, so the below-wall condition is
inherited.

For $1\le m\le R$, the three states
$W^{2m-2},W^{2m-1},W^{2m}$ are the corresponding three states of $A$, so there
is no odd peak by hypothesis on $A$.  For $R<m\le R+N$, write
$m=R+\ell_0$; the three states are the shifted corresponding states of $L$,
so $m$ is an odd peak of $W$ exactly when $\ell_0\in M_L$.  For
$R+N<m\le R+N+S$, write $m=R+N+b$; the three states are the shifted
corresponding states of $B$, so there is no odd peak by hypothesis on $B$.
These three cases prove the displayed peak set.
\end{proof}

\begin{corollary}[Top-wall loops attach to arbitrary below-wall endpoints]\label{cor:bc-top-wall-loop-endpoints}
Let $w\ge3$, let $\mu,\nu$ be partitions with
\[
  \ell(\mu)<w,\qquad \ell(\nu)<w,
\]
and let $\tau$ be a partition with exactly $w-2$ nonzero parts.  For every
$N\ge0$ and every $M\subseteq[N]$, there are integers $R,S\ge0$ and a
$w$-up-down tableau
\[
  W=(W^0,\ldots,W^{2(R+N+S)})
\]
from $\mu$ to $\nu$, staying below the wall, whose odd length-peak set is
\[
  \{R+m:m\in M\}.
\]
The connector parts of $W$ may be chosen from $(\mu,\tau)$ and $(\tau,\nu)$
alone.
\end{corollary}

\begin{proof}
Apply \cref{lem:bc-below-wall-peakless-connector} to obtain a peakless
below-wall connector $A$ from $\mu$ to $\tau$ and a peakless below-wall
connector $B$ from $\tau$ to $\nu$.  Apply
\cref{lem:bc-top-wall-peak-loop} to $\tau$ and $M$ to obtain a below-wall loop
$L$ from $\tau$ to $\tau$ whose odd length peaks are exactly $M$.  The
concatenation statement is then exactly
\cref{lem:bc-peakless-connector-attachment}.  The final assertion follows
from the canonical connector construction in
\cref{lem:bc-below-wall-peakless-connector}.
\end{proof}

\begin{lemma}[Peak-supported branch selections land in $Z_0$]\label{lem:bc-peak-supported-z0-landing}
In the setting of \cref{lem:bc-local-branch-recovery}, suppose a local repair
construction supplies a $w$-up-down tableau
\[
  U=(U^0,U^1,\ldots,U^{2N})
\]
and distinct assigned event indices
\[
  \tau_1,\ldots,\tau_{h_\alpha}\in[N].
\]
Let
\[
  P(U)=
  \{m\in[N]:\ell(U^{2m-2})<\ell(U^{2m-1})>\ell(U^{2m})\}
\]
be the set of odd length peaks.  Assume
\[
  \ell(U^r)<w\qquad(0\le r\le2N),
\]
\[
  P(U)\subseteq\{\tau_u:1\le u\le r_\alpha\},
\]
and
\[
  m\in P(U)\quad\Longrightarrow\quad \ell(U^{2m-1})=w-1.
\]
Define
\[
  M_\alpha=\{u\in\{1,\ldots,r_\alpha\}:\tau_u\in P(U)\}.
\]
If $\beta_\alpha$ is the local branch word whose second half encodes this
$M_\alpha$, then the HKKO boundary-event splice of
\cref{lem:bc-local-z0-splice-package} applies to $U$ and is reversible from
the unmarked output, $\beta_\alpha$, and the assigned indices.  In particular,
the unselected non-peak condition in \cref{lem:bc-z0-landing-criterion} is
automatic from the peak-support hypothesis.
\end{lemma}

\begin{proof}
The set $M^\tau_\alpha=\{\tau_u:u\in M_\alpha\}$ of
\cref{lem:bc-local-z0-splice-package} is exactly $P(U)$ by the definition of
$M_\alpha$ and the inclusion
$P(U)\subseteq\{\tau_u:1\le u\le r_\alpha\}$.  Therefore the displayed
top-wall condition on $P(U)$ gives
\[
  u\in M_\alpha\quad\Longrightarrow\quad
  \ell(U^{2\tau_u-1})=w-1 .
\]
If $m\notin M^\tau_\alpha$, then $m\notin P(U)$, which is precisely the
negation of the odd length-peak condition at $m$.  The below-wall condition is
one of the hypotheses.  Thus all hypotheses of
\cref{lem:bc-local-z0-splice-package} hold.  That lemma gives the splice and
the stated reversibility.  The last sentence is the same implication
$m\notin M^\tau_\alpha\Rightarrow m\notin P(U)$.
\end{proof}

\begin{lemma}[Canonical loop segments align the local event indices]\label{lem:bc-canonical-loop-event-indices}
Let $w\ge3$, let $h\ge0$, and let $\mu,\nu$ be partitions with
\[
  \ell(\mu)<w,\qquad \ell(\nu)<w .
\]
Let $\eta_w=(1^{w-2})$, the partition with $w-2$ parts all equal to $1$.
For $M\subseteq[h]$, choose the connector $A$ from $\mu$ to $\eta_w$ by the
canonical construction in \cref{lem:bc-below-wall-peakless-connector}, choose
the loop $L$ from $\eta_w$ to itself by
\cref{lem:bc-top-wall-peak-loop}, and choose the connector $B$ from $\eta_w$
to $\nu$ by the same canonical connector construction.  Let $R$ be the
half-length of $A$, and concatenate $A,L,B$ as in
\cref{lem:bc-peakless-connector-attachment}.  If
\[
  \vartheta_i=R+i\qquad(1\le i\le h),
\]
then the resulting segment $W$ is a below-wall $w$-up-down tableau from
$\mu$ to $\nu$ with odd length-peak set
\[
  P(W)=\{\vartheta_i:i\in M\}.
\]
Every midpoint indexed by $P(W)$ has length $w-1$.  The indices
$\vartheta_1,\ldots,\vartheta_h$ are distinct and are determined by $w$,
$h$, and the left endpoint $\mu$; in particular they carry no branch data.

Consequently, in the setting of \cref{lem:bc-local-branch-recovery}, if
$h=h_\alpha$, $M=M_\alpha\subseteq\{1,\ldots,r_\alpha\}$, and the local
repair segment is chosen as above, then the hypotheses of
\cref{lem:bc-peak-supported-z0-landing} hold with assigned indices
$\tau_i=\vartheta_i$.  Thus the HKKO splice applies with no additional
event-index storage.
\end{lemma}

\begin{proof}
The partition $\eta_w$ has exactly $w-2$ nonzero parts, so
\cref{lem:bc-top-wall-peak-loop} applies to $\eta_w$ and $M$.  The two
connector segments are below-wall and have no odd length peaks by
\cref{lem:bc-below-wall-peakless-connector}.  Therefore
\cref{lem:bc-peakless-connector-attachment} gives a below-wall segment $W$
from $\mu$ to $\nu$ whose odd peak set is the loop peak set shifted by the
half-length $R$ of the left connector.  Since the loop peak set is $M$, this
is exactly
\[
  P(W)=\{R+i:i\in M\}=\{\vartheta_i:i\in M\}.
\]
The selected midpoints are the selected midpoints of the loop, so
\cref{lem:bc-top-wall-peak-loop} gives length $w-1$ at each of them.

The connector construction of \cref{lem:bc-below-wall-peakless-connector} is
canonical once its two endpoints are fixed: it passes through the
componentwise maximum and adds or removes boxes in the prescribed row order.
Hence the half-length $R$ of the connector from $\mu$ to $\eta_w$ is
determined by $w$ and $\mu$, and the displayed indices are determined by
$w,h,\mu$.  Their distinctness is immediate from $\vartheta_i=R+i$.

For the final assertion, the displayed peak identity gives
\[
  P(W)\subseteq\{\vartheta_i:1\le i\le r_\alpha\}
\]
because $M_\alpha\subseteq\{1,\ldots,r_\alpha\}$.  The below-wall condition
and the length-$w-1$ condition at peaks have just been proved.  Thus
\cref{lem:bc-peak-supported-z0-landing} applies with
$\tau_i=\vartheta_i$.  Since the assigned indices are determined by the local
endpoint data and $h_\alpha$, they require no additional storage.
\end{proof}

\begin{corollary}[Canonical loop splice is locally invertible]\label{cor:bc-canonical-loop-splice-invertible}
In the setting of \cref{lem:bc-local-branch-recovery}, suppose
$h=h_\alpha$, the local endpoints $\mu,\nu$ satisfy
\[
  \ell(\mu)<w,\qquad \ell(\nu)<w,
\]
and the selected subset
$M_\alpha\subseteq\{1,\ldots,r_\alpha\}$ is the subset recovered from the
second half of the local branch word $\beta_\alpha$.  Construct the local
segment $W$ from $\mu$ to $\nu$ by
\cref{lem:bc-canonical-loop-event-indices}, and let $W^\ast$ be its HKKO
boundary-event splice.  Then $W^\ast$ is an unmarked $w$-up-down tableau from
$\mu$ to $\nu$, and the marked preimage segment $W$ is recoverable from
$W^\ast$, $\beta_\alpha$, $w$, $h_\alpha$, and the local endpoint $\mu$.
Thus the canonical loop splice introduces no hidden local data beyond the
branch word.
\end{corollary}

\begin{proof}
\Cref{lem:bc-canonical-loop-event-indices} supplies assigned indices
\[
  \tau_i=\vartheta_i=R+i\qquad(1\le i\le h_\alpha)
\]
and proves that they are determined by $w,h_\alpha,\mu$.  The same lemma
checks the hypotheses of \cref{lem:bc-peak-supported-z0-landing}; hence
\cref{lem:bc-local-z0-splice-package} applies to $W$ and produces an unmarked
segment $W^\ast$ with the same endpoints $\mu,\nu$.  Its reversibility
statement recovers $W$ from $W^\ast$, $\beta_\alpha$, and the assigned
indices.  Since the assigned indices are themselves recovered from
$w,h_\alpha,\mu$, the displayed data recover $W$.  No further local datum is
used.
\end{proof}

\begin{lemma}[Local endpoints canonically choose a below-wall parameter]\label{lem:bc-local-endpoints-canonical-wall}
Let $\mu$ and $\nu$ be partitions.  Define
\[
  w(\mu,\nu)=\max\{3,\ell(\mu)+1,\ell(\nu)+1\}.
\]
Then $w(\mu,\nu)\ge3$,
\[
  \ell(\mu)<w(\mu,\nu),\qquad \ell(\nu)<w(\mu,\nu),
\]
and the integer $w(\mu,\nu)$ is a fixed function of the endpoint pair
$(\mu,\nu)$.
\end{lemma}

\begin{proof}
The inequalities $w(\mu,\nu)\ge \ell(\mu)+1$ and
$w(\mu,\nu)\ge \ell(\nu)+1$ are immediate from the definition of the maximum,
and they imply the two strict below-wall inequalities.  The same definition
also gives $w(\mu,\nu)\ge3$ and shows that no additional choice is involved.
\end{proof}

\begin{proposition}[Below-wall endpoints supply the local HKKO splice package]\label{prop:bc-below-wall-local-hkko-splice-package}
Work in the setting of \cref{lem:bc-local-branch-recovery}.  Put
$h=h_\alpha$, and let $\mu,\nu$ be the local endpoint partitions for a
deletion joint.  Suppose that
\[
  \ell(\mu)<w,\qquad \ell(\nu)<w .
\]
Then, for every local branch word
\[
  \beta_\alpha\in\{0,1\}^{2h_\alpha},
\]
there is a canonical assigned index list
\[
  \tau_1<\cdots<\tau_{h_\alpha}
\]
and a canonical $w$-up-down segment $W$ from $\mu$ to $\nu$ whose selected
indices are exactly the subset $M_\alpha$ recovered from the second half of
$\beta_\alpha$.  The HKKO boundary-event splice produces an unmarked
$w$-up-down segment $W^\ast$ from $\mu$ to $\nu$, and the marked preimage
$W$ is recoverable from
\[
  W^\ast,\quad \beta_\alpha,\quad w,\quad h_\alpha,\quad \mu .
\]
No event-index choice or marking choice beyond the branch word is introduced.
\end{proposition}

\begin{proof}
By \cref{lem:bc-local-branch-recovery}, the branch word recovers the selected
subset $M_\alpha\subseteq\{1,\ldots,r_\alpha\}$ from its second half.  Apply
\cref{lem:bc-canonical-loop-event-indices} to the endpoints $\mu,\nu$, the
integer $h=h_\alpha$, and this subset.  It gives a canonical below-wall
$w$-up-down segment $W$ from $\mu$ to $\nu$ and a canonical assigned index
list $\tau_i=\vartheta_i$ determined by $w,h_\alpha,\mu$; its peak set is
exactly the image of $M_\alpha$ under that list.

The same lemma checks the hypotheses of
\cref{lem:bc-peak-supported-z0-landing}.  Hence
\cref{lem:bc-local-z0-splice-package} applies and the HKKO boundary-event
splice produces an unmarked segment $W^\ast$ with the same endpoints.  Its
reversibility statement recovers $W$ from $W^\ast$, $\beta_\alpha$, and the
assigned indices.  Since the assigned indices are themselves fixed functions
of $w,h_\alpha,\mu$, the displayed data recover $W$.
\end{proof}

\begin{lemma}[Canonical below-wall packet outputs are image-compatible]\label{lem:bc-canonical-below-wall-packet-output}
Work in the setting of \cref{lem:bc-local-branch-recovery}.  Put
$h=h_\alpha$, let $\mu,\nu$ be endpoint partitions for a local packet, and
put
\[
  w=w(\mu,\nu)
\]
as in \cref{lem:bc-local-endpoints-canonical-wall}.  For every local branch
word
\[
  \beta_\alpha\in\{0,1\}^{2h},
\]
there are canonically chosen segments
\[
  W(\mu,\nu,\beta_\alpha),\qquad
  W^\ast(\mu,\nu,\beta_\alpha),
\]
such that $W^\ast(\mu,\nu,\beta_\alpha)$ is an unmarked HKKO packet output
from $\mu$ to $\nu$ in the image of the canonical below-wall splice package
of \cref{prop:bc-below-wall-local-hkko-splice-package}, and the marked
preimage $W(\mu,\nu,\beta_\alpha)$ is recoverable from
\[
  W^\ast(\mu,\nu,\beta_\alpha),\quad
  \beta_\alpha,\quad w,\quad h,\quad \mu .
\]
No packet-output image choice is introduced beyond the endpoint pair and the
branch word.
\end{lemma}

\begin{proof}
\Cref{lem:bc-local-endpoints-canonical-wall} gives
\[
  \ell(\mu)<w,\qquad \ell(\nu)<w .
\]
Therefore \cref{prop:bc-below-wall-local-hkko-splice-package} applies to the
endpoint pair $\mu,\nu$, the canonical wall $w$, and the branch word
$\beta_\alpha$.  In the proof of that proposition the endpoint pair and
branch word canonically choose the assigned index list, the below-wall marked
segment $W$, and its unmarked HKKO splice output $W^\ast$.  The same
proposition states that $W$ is recovered from
\[
  W^\ast,\quad \beta_\alpha,\quad w,\quad h,\quad \mu .
\]
Since $W^\ast$ is produced by that splice, it lies in the image of the
canonical below-wall splice package.  All choices used in its construction are
fixed functions of $\mu,\nu$ and $\beta_\alpha$.
\end{proof}

\begin{lemma}[Connector insertion is not target-size preserving]\label{lem:bc-connector-target-size-guard}
Let $P$ be a two-step path with $j$ labelled blocks, and let
$S\subset\{1,\ldots,j\}$ have cardinality $q$.  Suppose one forms a new
windowed object by keeping every block whose label is not in $S$ and, for the
maximal deleted intervals, inserting local two-step windows of half-lengths
$N_1,\ldots,N_L$.  Then the resulting windowed object has
\[
  (j-q)+\sum_{\alpha=1}^L N_\alpha
\]
two-step windows.  Hence it belongs to
$\mathcal P_{j-q}^{\mathrm{even}}$ by block count only if every inserted
window has half-length $0$.

Consequently the connector-loop segments of
\cref{cor:bc-top-wall-loop-endpoints,lem:bc-canonical-loop-event-indices,cor:bc-canonical-loop-splice-invertible}
cannot be the final fixed-witness map by raw insertion when any connector-loop
window has positive half-length.  They may only be used inside a
target-size-preserving local splice or deterministic parser whose output still
has exactly $j-q$ two-step blocks.
\end{lemma}

\begin{proof}
Deleting the $q$ labelled blocks leaves exactly $j-q$ retained blocks.  Raw
insertion of local windows adds $N_\alpha$ two-step windows at the
$\alpha$-th deleted interval and does not remove any retained block.  Summing
over the maximal deleted intervals gives the displayed count.  The paths in
$\mathcal P_{j-q}^{\mathrm{even}}$ have, by definition, exactly $j-q$
two-step blocks, so equality of block counts forces
$\sum_\alpha N_\alpha=0$.  Since each $N_\alpha$ is nonnegative, this is
equivalent to $N_\alpha=0$ for all $\alpha$.

The connector-loop segments have positive half-length whenever they contain a
nonempty connector or loop.  Therefore raw insertion of such a segment cannot
be the final target-size map.  A valid use must occur inside a construction
whose output block count is restored to $j-q$, which is precisely the
target-size-preserving splice or parser alternative stated in the last
sentence.
\end{proof}

\begin{lemma}[Unchanged retained windows leave no target-size slots]\label{lem:bc-target-size-parser-no-spare-windows}
Let $P$ be a two-step path with $j$ labelled blocks, let
$S\subset\{1,\ldots,j\}$ have cardinality $q$, and put
$I(S)=\{j-s+1:s\in S\}$.  Let $P^\ast$ be a two-step path with exactly
$j-q$ two-step windows.  Suppose a deterministic parser of $P^\ast$ returns
pairwise disjoint output windows consisting of:
\begin{enumerate}
\item one retained window for each $i\notin I(S)$, equal to the original
block $B_i(P)$; and
\item local output windows $U^\ast_\alpha$ attached to the maximal deleted
intervals of $I(S)$,
\end{enumerate}
and suppose these returned windows exhaust the $j-q$ two-step windows of
$P^\ast$.  Then every local output window $U^\ast_\alpha$ has half-length
$0$.

Consequently, a target-size fixed-witness repair cannot both keep every
nondeleted block as a full unchanged output window and also insert any
positive-length local HKKO or connector segment.  Any positive-length local
repair must replace or merge some neighbouring retained windows, or the local
inverse must use a zero-window slot together with the endpoint and branch
data.
\end{lemma}

\begin{proof}
The set $I(S)$ has $q$ elements, so there are exactly $j-q$ indices
$i\notin I(S)$.  The parser therefore already returns $j-q$ retained
two-step windows in item (1).  Since these windows are pairwise disjoint and
the whole target path $P^\ast$ has exactly $j-q$ two-step windows, no
additional returned window can have positive half-length.  Hence each
$U^\ast_\alpha$ has half-length $0$.

The consequence is the same counting statement applied to any proposed
repair whose output lies in the target set
$\mathcal P_{j-q}^{\mathrm{even}}$: a positive-length local segment would add
at least one extra two-step window unless it replaces or merges windows that
were otherwise counted among the unchanged retained blocks.  If no such
replacement or merger occurs, the local parser has only a zero-window slot at
that joint, so inversion at the joint can use only the endpoint data and the
branch word.
\end{proof}

\begin{lemma}[Zero-window unchanged repairs force endpoint matching]\label{lem:bc-zero-window-endpoint-matching}
Let
\[
  P=(\rho_0,\rho_1,\ldots,\rho_{2j})
\]
be a two-step path with $j$ labelled blocks, let
$S\subset\{1,\ldots,j\}$, and put $I(S)=\{j-s+1:s\in S\}$.  Write the maximal
consecutive intervals of $I(S)$ as
\[
  K_\alpha=\{u_\alpha,u_\alpha+1,\ldots,v_\alpha\}.
\]
Suppose a target two-step path $P^\ast$ is parsed as the unchanged retained
blocks $B_i(P)$, $i\notin I(S)$, in their original order, with no
positive-length local output window at any deleted interval.  Then each
deleted interval satisfies the corresponding endpoint matching condition:
\[
  \rho_{2u_\alpha-2}=\rho_{2v_\alpha}
  \qquad\text{if }1<u_\alpha\le v_\alpha<j,
\]
\[
  \rho_{2v_\alpha}=\varnothing
  \qquad\text{if }u_\alpha=1\le v_\alpha<j,
\]
and
\[
  \rho_{2u_\alpha-2}=\varnothing
  \qquad\text{if }1<u_\alpha\le v_\alpha=j.
\]
There is no endpoint condition when $K_\alpha=\{1,\ldots,j\}$.

Consequently, an unchanged-retained zero-window repair can apply only to
deletion joints whose boundary partitions already match the start/end
requirements above.  At any nonmatching joint, a target-size repair must
change or merge neighbouring retained windows, or else use a local mechanism
that is not parsed as an unchanged zero-window deletion.
\end{lemma}

\begin{proof}
Consider first an internal interval $1<u_\alpha\le v_\alpha<j$.  Since
$K_\alpha$ is maximal, the neighbouring block indices $u_\alpha-1$ and
$v_\alpha+1$ are retained.  With no positive local output window between
them, these two retained output blocks are adjacent in the parsed target path.
The right endpoint of $B_{u_\alpha-1}(P)$ is $\rho_{2u_\alpha-2}$, and the
left endpoint of $B_{v_\alpha+1}(P)$ is $\rho_{2v_\alpha}$.  Adjacent
two-step windows in a single path share their common endpoint, so
$\rho_{2u_\alpha-2}=\rho_{2v_\alpha}$.

If $u_\alpha=1$ and $v_\alpha<j$, then the first retained block is
$B_{v_\alpha+1}(P)$ and, with no preceding local window, it must start at the
initial partition of $P^\ast$, namely $\varnothing$.  Its left endpoint is
$\rho_{2v_\alpha}$, proving the second displayed equality.  The terminal case
$1<u_\alpha\le v_\alpha=j$ is the same argument at the right end: the last
retained block before the deleted interval has right endpoint
$\rho_{2u_\alpha-2}$, and the target path ends at $\varnothing$.  If all
blocks are deleted, there are no retained neighbouring windows and hence no
endpoint equality is forced.

The consequence is the contrapositive of these necessary conditions.
\end{proof}

\begin{lemma}[Endpoint matching is sufficient for raw retained concatenation]\label{lem:bc-endpoint-matching-sufficient}
Let
\[
  P=(\rho_0,\rho_1,\ldots,\rho_{2j})
\]
be a two-step path with $j$ labelled blocks, let
$S\subset\{1,\ldots,j\}$, and put $I(S)=\{j-s+1:s\in S\}$.  Suppose every
maximal deleted interval $K_\alpha=\{u_\alpha,\ldots,v_\alpha\}$ of $I(S)$
satisfies the matching condition displayed in
\cref{lem:bc-zero-window-endpoint-matching}.  Then concatenating the retained
blocks
\[
  B_i(P)=(\rho_{2i-2},\rho_{2i-1},\rho_{2i}),\qquad i\notin I(S),
\]
in their original order gives a two-step path with exactly $j-|S|$ blocks.
Each retained block keeps its original vertical-strip adjacent differences and
its original weight.
\end{lemma}

\begin{proof}
Order the retained indices increasingly.  If two consecutive retained indices
$i<i'$ satisfy $i'=i+1$, then the right endpoint of $B_i(P)$ is
$\rho_{2i}$ and the left endpoint of $B_{i'}(P)$ is
$\rho_{2i'-2}=\rho_{2i}$, so the blocks concatenate.

If $i'>i+1$, then the missing indices between them form a maximal deleted
interval $K_\alpha=\{i+1,\ldots,i'-1\}$.  This is an internal interval, so the
matching condition gives
\[
  \rho_{2i}=\rho_{2(i+1)-2}=\rho_{2(i'-1)}=\rho_{2i'-2},
\]
which is again exactly the endpoint equality needed to concatenate
$B_i(P)$ and $B_{i'}(P)$.

If the first retained index is $i>1$, then the deleted prefix
$\{1,\ldots,i-1\}$ is a maximal deleted interval and the boundary matching
condition gives $\rho_{2i-2}=\varnothing$, so the first retained block starts
at the required initial partition.  If there is no deleted prefix, then
$i=1$ and $\rho_0=\varnothing$ already.  The terminal suffix is identical:
if the last retained index is $i<j$, the matching condition gives
$\rho_{2i}=\varnothing$, and otherwise $i=j$ gives the original terminal
partition $\rho_{2j}=\varnothing$.

Thus the retained blocks concatenate to a single two-step path.  The number
of retained indices is $j-|S|$, so this path has exactly $j-|S|$ blocks.  The
vertical-strip and weight assertions are inherited from the original blocks.
\end{proof}

\begin{lemma}[Endpoint-matched zero-window codes are not universal]\label{lem:bc-zero-window-loop-fiber-large}
Let $h\ge1$ and let
\[
  \mu=(5,4,3,2,1).
\]
There are at least $5^h$ local two-step segments
\[
  U=(U^0,U^1,\ldots,U^{2h})
\]
from $\mu$ to $\mu$ such that every adjacent difference is a vertical strip
and every two-step block has weight $2$:
\[
  |U^{2a-2}|-2|U^{2a-1}|+|U^{2a}|=2
  \qquad(1\le a\le h).
\]
Consequently no zero-window local inverse that is given only the endpoint
$\mu$, the length $h$, and a branch word in $\{0,1\}^{2h}$ can recover all
endpoint-matched local segments of this type.

Thus a matched zero-window step in the fixed-witness repair cannot use a
universal endpoint-only inverse for arbitrary local growth segments; it must
use the fixed witness-crossing restrictions, or else the target path must
carry additional replacement or merger data.
\end{lemma}

\begin{proof}
The partition $\mu=(5,4,3,2,1)$ has five removable corner boxes, one at the
end of each row.  For each word
\[
  c=(c_1,\ldots,c_h)\in\{1,2,3,4,5\}^h,
\]
define a local segment by
\[
  U^{2a-2}=\mu,\qquad U^{2a}=\mu
  \quad(1\le a\le h),
\]
and let $U^{2a-1}$ be the partition obtained from $\mu$ by deleting the
corner box in row $c_a$.  Removing one corner box and adding it back are both
vertical-strip moves.  Moreover
\[
  |U^{2a-2}|-2|U^{2a-1}|+|U^{2a}|
  =
  |\mu|-2(|\mu|-1)+|\mu|
  =
  2.
\]
Different words $c$ give different midpoint sequences, hence different local
segments.  Therefore there are at least $5^h$ such segments.

The number of branch words of length $2h$ is $2^{2h}=4^h$, and
$5^h>4^h$.  If the local output window has length zero, all these segments
have the same output and the same endpoint data.  A map using only those data
and a $2h$-bit branch word therefore has fewer possible codes than local
segments, so it cannot be injective on this family.  The final sentence is
this counting obstruction applied to any proposed zero-window local inverse.
\end{proof}

\begin{corollary}[Singleton zero-window packet inverses are not universal]\label{cor:bc-singleton-zero-window-packet-obstruction}
There is no fixed inverse which recovers every one-block local segment
\[
  U=(U^0,U^1,U^2)
\]
from only the common endpoint pair
\[
  U^0=U^2=(5,4,3,2,1),
\]
the known singleton interval, and a branch word in $\{0,1\}^2$, when the
local output window has half-length zero and the segment is required only to
have vertical-strip adjacent differences and weight
\[
  |U^0|-2|U^1|+|U^2|=2 .
\]
Consequently, a height singleton packet cannot be discharged by a universal
endpoint-plus-branch zero-window inverse in the ambient growth-path category.
Any singleton-packet supplier must use the fixed witness-crossing restrictions
or carry replacement/merged target data that recover the midpoint.
\end{corollary}

\begin{proof}
This is the case $h=1$ of
\cref{lem:bc-zero-window-loop-fiber-large}.  There are at least five such
segments with the displayed common endpoints, obtained by deleting and
restoring one of the five removable corner boxes of $(5,4,3,2,1)$.  A
zero-window code with the displayed endpoint pair and a two-bit branch word
has only four possible values.  Therefore it cannot be injective on this
five-element family, so no fixed inverse using only those data can recover all
segments.  The final sentence is the same obstruction applied to a proposed
singleton-packet zero-window decoder before any fixed-witness or merged-target
information has been used.
\end{proof}

\begin{corollary}[Path formulation of the repaired deletion problem]\label{cor:bc-path-repair-problem}
To prove the fixed-witness estimates in
\cref{prop:bc-bad-shape-fiber-reduction}, it is enough to construct, for the
growth paths of \cref{lem:bc-growth-path-model} attached to a fixed witness
set $S$, an injection into
\[
  \{0,1\}^{2q}\times \mathcal P_{j-q}^{\mathrm{even}},
\]
where $\mathcal P_{j-q}^{\mathrm{even}}$ denotes the two-step growth paths
corresponding to column-even semistandard tableaux of content $(2^{j-q})$.
The path operation must remove the $q$ two-step blocks whose tableau labels
lie in $S$--equivalently, blocks $j-s+1$ for $s\in S$--repair the resulting
joints, and recover the original path from the repaired path, $S$, and the
$2q$ branch bits.  The source-backed wall-event model available for this
repair is the up-down boundary-event model of
\cref{lem:bc-hkko-updown-boundary-events}, the reversible boundary-event
splice of \cref{lem:bc-hkko-boundary-splice}, and the marked boundary-step
specialization in \cref{lem:bc-hkko-marked-boundary};
the standard-tableau odd-column removal available after a joint has been
converted into a genuine standard tableau is \cref{lem:bc-odd-column-removal},
but \cref{lem:bc-odd-column-removal-size-guard} shows that this removal cannot
be the final fixed-witness map unless no odd columns remain;
the semistandard-to-standard bridge is the doubled-label standardization of
\cref{lem:bc-doubled-standardization}, with the local adjacent-pair frontier
states of \cref{lem:bc-standardized-witness-pairs} and the crossing times of
\cref{lem:bc-standard-frontier-time}, ordered inside the deleted pairs by
\cref{lem:bc-standard-frontier-order}, distributed across deletion joints by
\cref{lem:bc-frontier-events-per-joint}, and assigned their consecutive local
frontier ranks by \cref{lem:bc-consecutive-witness-ranks}; the matching
column-parity defects are
localized joint-by-joint by \cref{lem:bc-column-defects-per-joint} and are
determined by the adjacent endpoint partitions by
\cref{lem:bc-joint-endpoint-defects}, with canonical paired defect slots from
\cref{cor:bc-joint-defect-pairs} assigned to ordered frontier-event slots by
\cref{cor:bc-joint-events-cover-defects}; the allowed local branch word is
made explicit in \cref{lem:bc-local-branch-word}, recovers the local
event/defect data by \cref{lem:bc-local-branch-recovery}, and becomes an
admissible HKKO marking once the selected slots are realized as eligible
boundary-up steps by
\cref{lem:bc-local-branch-hkko-marks,lem:bc-frontier-events-supply-marks}.
By
\cref{lem:bc-standard-time-pair-deletion}, it remains to connect the deleted
consecutive standard-time pairs, equivalently the deleted two-step blocks, and
their ordered per-joint frontier times to the penultimate-row creations of
\cref{lem:bc-z0-selected-midpoint-creation} in the up-down splice segment and
to the marked boundary-up steps of
\cref{lem:bc-hkko-boundary-up-characterization} after coefficient extraction,
to place the resulting local segment in the domain of the HKKO splice
\cref{lem:bc-hkko-boundary-splice} using the explicit criterion
\cref{lem:bc-z0-landing-criterion} and the decoded-splice package
\cref{lem:bc-local-z0-splice-package}, with the canonical local event indices
of \cref{lem:bc-canonical-hkko-event-indices}, the explicit top-wall
peak-loop model of \cref{lem:bc-top-wall-peak-loop}, the peakless connector
attachment of \cref{lem:bc-peakless-connector-attachment}, the
arbitrary-endpoint connector and index-alignment leaves
\cref{cor:bc-top-wall-loop-endpoints,lem:bc-canonical-loop-event-indices}, the
local splice inversion of
\cref{cor:bc-canonical-loop-splice-invertible}, the below-wall package
\cref{prop:bc-below-wall-local-hkko-splice-package}, and
the peak-supported landing criterion of \cref{lem:bc-peak-supported-z0-landing},
while respecting the target-size guard of
\cref{lem:bc-connector-target-size-guard} and the no-spare-window constraint
of \cref{lem:bc-target-size-parser-no-spare-windows}, as well as the
endpoint-matching obstruction of \cref{lem:bc-zero-window-endpoint-matching}.
Conversely, \cref{lem:bc-endpoint-matching-sufficient} says that matched
zero-window joints are exactly the cases where the retained windows already
concatenate as a target-size path, while
\cref{lem:bc-zero-window-loop-fiber-large} rules out a universal
endpoint-only inverse for all matched local loops with only the $2h$ branch
bits.
Thus the remaining construction must output an object of
$\mathcal P_{j-q}^{\mathrm{even}}$ through a target-size-preserving splice or
parser, cannot simply append positive-length local windows while keeping all
retained windows unchanged, cannot simply skip nonmatching deleted intervals,
cannot recover arbitrary matched local loops from only endpoint data and the
branch word, and must recover the original path from that repaired path and
the decoded branch words.  The replacement-zone recovery criterion of
\cref{prop:bc-replacement-zones-global-recovery} is the remaining positive
parser shape: a local repair may absorb neighbouring retained blocks only if
the parsed replacement zone outputs exactly the number of retained blocks in
that zone and has a local inverse.  The canonical halo zones of
\cref{lem:bc-canonical-halo-replacement-zones} provide a deterministic
one-block-neighbour replacement-zone choice from $S$ with the exact global
block-count identity needed by that parser.  The branch-accounting form
\cref{prop:bc-replacement-zone-fiber-criterion} shows that two branch bits per
deleted witness block inside each replacement zone are exactly sufficient for
the required $2^{2q}s_{j-q}$ fixed-fiber bound.
\cref{lem:bc-canonical-halo-short-gaps} shows that the retained material
inside one zone consists only of endpoint halo blocks and one- or two-block
gaps between deleted packets.
\cref{lem:bc-canonical-halo-retained-budget} bounds the replacement output in
each zone by twice the number of localized deleted blocks.
\cref{lem:bc-canonical-halo-zone-balance} gives the exact retained-minus-deleted
balance of a zone from its endpoint halo bits, internal retained gap
lengths, and deleted packet lengths.
\cref{cor:bc-height-halo-singleton-nonnegative} sharpens this on the height
side in the residual range: height witness labels are separated, so the
canonical halo zones contain only singleton deleted packets and have no
negative local balance.
\cref{cor:bc-height-zero-surplus-zones} further classifies the tight height
zones: they are exactly right-boundary alternating singleton strings; every
other height zone has positive surplus.
\cref{cor:bc-height-zero-surplus-label-prefix} translates such a tight zone
back to labels: the deleted labels are precisely
$\{1,3,\ldots,2m-1\}$ and the retained labels are
$\{2,4,\ldots,2m\}$.
\cref{cor:bc-tight-height-prefix-column-even-gate} gives the exact remaining
gate: if $m\ge2$ and the prefix endpoint at label $2m$ is already
column-even, the closed odd alternating height decoder applies; otherwise the
missing datum is the singleton local case or the endpoint/continuation parity
repair.
\cref{cor:bc-singleton-tight-height-forced-strip} removes the only
combinatorial freedom in that singleton case: the deleted label-$1$ strip is
the forced first-row domino, so the unresolved singleton work is only the
target-size merge with the retained label-$2$ block and the surrounding
endpoint/continuation recovery.
\cref{lem:bc-tight-height-continuation-parity} makes the continuation datum
explicit: the columns in which the later labels must contribute odd parity
are exactly the odd columns of the tight prefix endpoint, with canonical
adjacent pairing.
\cref{lem:bc-tight-prefix-defects-exceed-local-slots} also shows that this
full endpoint-defect set cannot simply be routed through the tight zone's own
$m$ frontier slots: a concrete tight $m=2$ prefix has three paired endpoint
defects.
\cref{lem:bc-tight-continuation-label-cover} gives the replacement structure:
those paired endpoint defects are nevertheless chargeable injectively to
distinct continuation labels whose strips hit the corresponding defect pair.
\cref{cor:bc-tight-height-visible-separator} gives the parser-facing geometry
of every non-full tight height zone: the block immediately to its left is
outside the one-block halo and is therefore returned unchanged, so the left
endpoint of the tight zone is visible without branch data.
\cref{cor:bc-tight-height-zone-endpoints-visible} packages this for the
right-boundary auxiliary criterion: the parsed separator block and the
terminal endpoint $\varnothing$ determine the endpoint pair of the tight zone.
\cref{lem:bc-tight-continuation-labels-external} then locates the
Hall-matched continuation labels on the parser side: their corresponding
blocks are never inside the tight zone, but lie either in the unchanged
complement or in strictly earlier halo-zone windows.
\cref{lem:bc-canonical-continuation-hall-charge} removes the remaining hidden
choice in this continuation charge: the matching is the lexicographically
least Hall injection determined by the recovered incidence relation.
\cref{lem:bc-continuation-incidence-semistandard-local} identifies that
incidence relation in semistandard coordinates: a continuation label is
incident to a defect pair exactly when one of its two tableau boxes lies in
one of the paired columns.
\cref{lem:bc-continuation-boxes-determine-pieri-segment} records the exact
conversion from labelled boxes to the path segment: the prefix endpoint plus
the two boxes of every continuation label determine
$\lambda^{(2m)}\subset\cdots\subset\lambda^{(j)}$.
\cref{lem:bc-pieri-continuation-determines-hall-charge} packages the same
information at the path level: recovering the Pieri continuation segment from
label $2m$ onward determines the paired endpoint defects, all continuation
incidences, and the canonical Hall charge.
\cref{prop:bc-tight-height-auxiliary-from-continuation-boxes} combines these
pieces into the tight-zone auxiliary criterion: outside recovery of the prefix
endpoint and labelled continuation boxes supplies the tight-zone endpoint
pair, the incidence graph, and the canonical Hall charge with no new branch
data.
\cref{prop:bc-target-window-hall-augmented-low-low-carrier-supplies-large-right-boundary-auxiliary}
uses the smaller direct Hall-data packaging for the current carrier route: it is
enough for the target-window data to recover the non-column-even prefix
endpoint, paired endpoint defects, continuation incidences, canonical Hall
charge, and a Hall-augmented Step~3 carrier containing $L_{2m}$.
\Cref{cor:bc-direct-hall-charge-selected-blocks-recovered} proves that the
Hall-charged continuation blocks selected by this direct Hall charge are already
part of the recovered outside datum, without first recovering the full Pieri
continuation segment.
\cref{prop:bc-target-window-continuation-segment-hall-carrier-supplies-large-right-boundary-auxiliary}
remains one source for this direct interface: a recovered Pieri continuation
segment determines the Hall data, and
\cref{cor:bc-hall-charge-selected-blocks-recovered} recovers the selected
outside blocks from that segment.
\cref{lem:bc-step2-q-suffix-strip-determines-continuation,lem:bc-step3-continuation-incidence-covers-q-suffix,cor:bc-strict-left-q-suffix-boundary-recovers-continuation-segment}
make the continuation-segment half smaller: the outside strict-left subfilling
already supplies the Step~2 $Q$-suffix entries, so the target window only has to
recover the internal separating boundary between the Step~2 $Q$-prefix strip
and that suffix strip; the remaining incoming suffix boundary is the fixed
global empty boundary by
\cref{lem:bc-q-suffix-boundary-internal-cut}.
\Cref{lem:bc-q-prefix-entries-determine-internal-q-cut,cor:bc-low-label-incidence-determines-internal-q-cut,cor:bc-strict-left-low-low-determines-internal-q-cut}
identify this cut from the opposite side: the $Q$-prefix entries determine the
cut, the low-label incidence entries determine those prefix entries, and in the
non-full tight setting the recovered strict-left columns reduce that readout to
the low-low cell set $L_{2m}$.
\cref{lem:bc-tight-outside-recovers-left-growth-segment} identifies the
growth-path data already supplied by the right-boundary outside parser: it
recovers the strict-left initial segment of the Krattenthaler growth path, the
segment corresponding to continuation labels $2m+1,\ldots,j$.
\cref{prop:bc-strict-left-growth-bridge-supplies-tight-auxiliary} names the
remaining bridge in a closure-ready form: if this strict-left segment decodes
the Pieri prefix endpoint and continuation boxes, then all tight-zone
auxiliary data are parser-supplied and branch-free.
\cref{lem:bc-strict-left-alone-prefix-obstruction} shows that this decoding is
not a consequence of the strict-left segment alone: a tight-zone inverse must
also use the constraints or branch data native to the final right-boundary
zone.
\cref{lem:bc-strict-left-obstruction-packet-defects} localizes the missing
information to the endpoint-defect data of a deleted singleton packet, so the
next inverse target is recovery of that packet datum from the final-zone
constraints or right-boundary auxiliary data before applying the packet-local
branch decoder.
\cref{lem:bc-tight-retained-closure-obstruction} rules out the next tempting
shortcut: retained boxes together with their smallest column-even closure can
miss an all-deleted even-column pair.  Thus the final tight-zone inverse must
recover the full tight-prefix packet data, not only the endpoint-defect
parity.
\cref{lem:bc-tight-even-retained-prefix-obstruction} sharpens this once more:
even the prefix endpoint, the retained even-label boxes, and the continuation
boxes can agree while two odd deleted packets are interchanged.  Thus the
residual odd-packet matching itself must be recovered from the repaired target
window or from the already budgeted branch word.
\cref{lem:bc-height-branch-first-half-slack} identifies the available bit
source: in height packets the first half of the local branch word is unused by
the local recovery formulas.  \Cref{prop:bc-tight-odd-packet-rank-criterion}
therefore reduces the residual odd-packet matching to a finite ambiguity
bound $|\Omega|\le2^m$ for the parser-visible tight-prefix datum.
\Cref{def:bc-tight-odd-packet-omega,lem:bc-true-tight-packets-in-omega,cor:bc-tight-omega-bound-closes-rank-leaf}
make this $\Omega$ explicit from the tight-prefix endpoint and retained even
strips, and give a first sufficient rank leaf: a uniform cardinality bound for
this endpoint/even-strip compatibility set.
\Cref{def:bc-tight-fixed-cell-completions,lem:bc-omega-fixed-cell-completions,cor:bc-fixed-cell-bound-closes-omega}
replace the recursive Pieri-path formulation of $\Omega$ by an equivalent
fixed-cell semistandard completion problem with ordinary row and column
inequalities.
\Cref{lem:bc-fixed-cell-first-column-forced,def:bc-reduced-fixed-even-growth-chains,prop:bc-completions-equal-reduced-growth-chains,cor:bc-reduced-growth-chain-bound-closes-fixed-cell}
strip off the now forced first column and reduce the remaining cardinality
leaf to a fixed-even growth-chain count in the non-first-column diagram.
\Cref{lem:bc-reduced-even-strip-sizes} records that every nonterminal reduced
even strip is a singleton; only the terminal even strip can have two boxes.
\Cref{def:bc-reduced-reverse-corner-tree,lem:bc-reduced-growth-reverse-corner-tree,cor:bc-reverse-corner-tree-bound-closes-reduced-leaf}
then reverse that count into an exact corner-removal tree; the unresolved
estimate is the uniform leaf bound for this tree.
\Cref{def:bc-odd-fixed-even-standard-tableaux,lem:bc-reverse-tree-fixed-even-standard-tableaux,cor:bc-odd-fixed-even-standard-bound-closes-reverse-tree}
remove the terminal fixed strip and identify this tree with an odd-size
standard-tableau fiber in which all even entries are fixed.
\Cref{def:bc-odd-cell-interval-extension-poset,lem:bc-fixed-even-standard-interval-extensions,cor:bc-interval-extension-bound-closes-fixed-even-leaf}
then identify the same fiber with interval-constrained linear extensions of
the odd-cell Young subposet.
\Cref{def:bc-last-pair-deletion-candidates,lem:bc-last-pair-deletion-recurrence,cor:bc-last-pair-blocked-binary}
give the exact terminal recurrence and show that a non-removable final fixed
even cell has only two possible final odd blockers.
\Cref{lem:bc-broad-fixed-even-bound-false} shows that the full fixed-even
$2^m$ bound is false in this breadth: a compatible $m=26$ datum has
$254531860>2^{26}$ fixed-even tableaux.  Thus the residual supplier cannot be
the uniform fixed-even estimate alone; it must use the extra parser-visible
outside data of the canonical halo/right-boundary construction.
\Cref{lem:bc-singleton-tight-omega-forced,cor:bc-singleton-tight-rank-leaf-closed}
and
\cref{lem:bc-two-packet-tight-omega-binary,lem:bc-binary-prefix-factorial-bound}
and
\cref{cor:bc-small-tight-rank-leaves-closed,prop:bc-seven-packet-tight-rank-leaf-certificate,prop:bc-eight-packet-tight-rank-leaf-certificate}
remove every case $m\le8$ from that live ambiguity: when $m=1$ the odd packet
is forced, and when $2\le m\le6$ the binary first-two-coordinate prefix and a
residual permutation bound already fit inside the $2^m$ branch budget; the
$m=7,8$ cases are closed by exact finite enumeration of the tight data.
\Cref{cor:bc-refined-tight-packet-criterion} records the corrected live leaf:
construct, from those parser-visible outside data, a canonically ordered
refined subset of the broad compatibility set that contains the true
odd-packet assignment and has size at most $2^m$.
\Cref{def:bc-commuting-normal-form-signature,prop:bc-commuting-normal-form-fibers-binary,cor:bc-commuting-signature-closes-refined-suffix}
then isolate the exact remaining overflow datum: after the branch-free
parser-visible right-boundary data fix the commuting normal-form signature,
every remaining signature class is binary at each non-commuting stage and
hence has size at most $2^m$.
\Cref{lem:bc-one-step-commuting-signature-symbol} gives the first local
decoder rule for that signature: the symbol is $\ast$ exactly when the fixed
even cell is blocked, and otherwise it is recovered from the lower shape
after deleting the top odd/even pair.
\Cref{def:bc-odd-sublevel-shape-chain,prop:bc-odd-sublevel-chain-determines-signature,cor:bc-target-window-odd-sublevel-chain-closes-signature}
iterate this rule: the whole commuting signature is determined by the
associated fixed-even datum and its odd sublevel shape chain.
\Cref{prop:bc-odd-sublevels-are-reduced-halfsteps,cor:bc-target-window-halfstep-chain-closes-signature}
translate this shape chain back to the reduced half-step diagrams already
present in the fixed-cell growth-chain model.
\Cref{lem:bc-retained-even-strips-determine-reduced-datum,cor:bc-target-window-retained-even-halfstep-closes-signature}
show that the reduced fixed-even datum itself is not a separate mystery once
the tight-prefix endpoint and retained even strips are visible; the live
remaining decoder is the true reduced half-step chain.
\Cref{lem:bc-pieri-prefix-shapes-are-reduced-halfsteps,cor:bc-target-window-prefix-path-closes-signature}
then identify those half-steps as first-column deletions of the true odd-time
tight-prefix Pieri shapes, so a decoder for the true tight-prefix shape chain
also supplies the signature input.
\Cref{cor:bc-target-window-standard-even-chain-closes-signature} translates
the same target into doubled-standardization coordinates: it is enough to
recover the even standard-time endpoints through time $4m$.
\Cref{cor:bc-target-window-dual-rsk-carrier-closes-signature} then gives the
growth-diagram form of the same task: a parsed dual-RSK' carrier
subarrangement, together with a deterministic standardization readout from
the recovered subfilling and internal corner labels, supplies those endpoints
by a backward sweep.
\Cref{cor:bc-target-window-low-label-carrier-closes-signature} makes that
readout concrete: it is enough for the recovered carrier data to determine
the two boxes carrying each semistandard label $1,\ldots,2m$.
\Cref{lem:bc-low-label-boxes-determine-tight-packet-datum} also identifies
the same readout with the tight-prefix packet datum
$(\lambda^{(2m)};E_1,\ldots,E_m;x_1,\ldots,x_m)$ used in the finite
odd-packet rank criterion.
\Cref{lem:bc-step2-knuth-prefix-low-label-boxes,cor:bc-target-window-step2-prefix-carrier-closes-signature}
then restate the readout in the actual two-stage Krattenthaler construction:
it is enough for the recovered Step~3 carrier cells to determine a Step~2
Knuth subarrangement whose forward boundary contains the prefix
$\lambda^{(0)},\ldots,\lambda^{(2m)}$.
\Cref{lem:bc-step3-cells-are-step2-half-matrix,cor:bc-target-window-step3-half-matrix-prefix-closes-signature}
remove the cell-entry part of that readout: since Step~3 is run on one
off-diagonal half of the symmetric Step~2 matrix, a recovered Step~3 carrier
already supplies the entries of any Step~2 cells it contains, up to transpose.
\Cref{lem:bc-step2-q-prefix-strip-determines-prefix,cor:bc-target-window-step2-q-prefix-strip-closes-signature}
then remove the left/bottom-boundary part for the anchored $Q$-prefix strip:
the strip starts on the global empty boundary, so its entries alone recover
$\lambda^{(0)},\ldots,\lambda^{(2m)}$ by the first-variant RSK boundary
description.
\Cref{lem:bc-step3-low-label-incidence-covers-q-prefix,cor:bc-target-window-low-label-incidence-carrier-closes-signature}
make the required Step~3 carrier explicit: it is enough to recover the
half-square cells representing unordered label pairs with at least one label in
$\{1,\ldots,2m\}$.
\Cref{lem:bc-low-label-incidence-split,lem:bc-strict-left-recovers-continuation-columns}
then split this target: the cross cells with one continuation label are already
recovered from the strict-left boundary segment, so the live target is the
low-low triangle $L_{2m}$.
\Cref{cor:bc-target-window-low-label-incidence-entries-close-signature,cor:bc-strict-left-reduces-incidence-to-low-low}
record the corresponding entry-level criteria.
\Cref{lem:bc-low-low-terminal-corner-boundary,cor:bc-target-window-low-low-boundary-closes-signature}
sharpen this once more: $L_{2m}$ is the terminal low-label dual-RSK' corner,
so it is enough to recover its boundary segment
$\rho_{2(j-2m)},\ldots,\rho_{2j}$ from the target-window data.
\Cref{cor:bc-target-window-terminal-zone-decoder-closes-signature} translates
that boundary target back into replacement-zone language: a branch-free
decoder for the original terminal right-boundary zone subpath is sufficient,
but stronger than the rank-budget route requires.
\Cref{cor:bc-target-window-refined-omega-closes-suffix} records the smaller
target compatible with the existing branch accounting: the target-window data
may instead determine a canonical at-most-$2^m$ refined ambiguity class
$\Omega_{\mathrm{tw}}$ for the odd packets.
\Cref{cor:bc-target-window-signature-fiber-omega} identifies one such class:
the fiber of the commuting normal-form signature inside the broad
compatibility set.
\Cref{cor:bc-target-window-retained-boxes-signature-fiber} rewrites the base
datum in parser-facing terms: the retained strips are just the two-box
readouts for labels $2,4,\ldots,2m$.
\Cref{cor:bc-target-window-even-chain-signature-fiber} rewrites this once more
in growth-chain terms: the even standard-time endpoint chain, together with
the continuation boxes, supplies the base datum and the signature fiber.
\Cref{cor:bc-target-window-low-label-boxes-signature-fiber} reduces the same
criterion to the low-label box readout: recovering the two boxes carrying each
label $1,\ldots,2m$, together with the labelled continuation boxes, already
recovers that even chain and therefore the signature-fiber ambiguity class.
\Cref{cor:bc-target-window-low-label-boxes-singleton-rank} records the sharper
conclusion used by the current route: the same low-label readout determines
the actual deleted odd-packet assignment, so the refined ambiguity class may be
chosen to be a singleton.
\Cref{cor:bc-target-window-continuation-sweep-singleton-rank} replaces the
abstract continuation-box hypothesis by a Step~2 continuation sweep: once a
recovered Knuth subarrangement contains the boundary segment
$\lambda^{(2m)},\ldots,\lambda^{(j)}$, the labelled continuation boxes are
read off from consecutive differences.
\Cref{cor:bc-strict-left-low-low-singleton-rank} then pulls this back to the
existing low-low triangle target: the strict-left segment recovers every
Step~3 cell incident to a continuation label, so the remaining low-low entries
recover the whole Step~2 boundary path and close the singleton rank route.
\Cref{cor:bc-bounded-low-low-candidates-fit-branch-budget} is the
branch-budget version of the same target: exact recovery of $L_{2m}$ is not
necessary if the target-window data determine a canonical at-most-$2^m$ list
of possible low-low subfillings containing the true one.
\Cref{def:bc-degree-low-low-candidates,prop:bc-strict-left-degree-low-low-pruning}
make this row-free: because the Step~2 matrix has doubled-content row sums and
the strict-left segment already recovers every continuation-column entry, the
true low-low subfilling must lie in an explicitly determined finite
degree-constrained ambient set.  The remaining supplier is therefore a uniform
target-window refinement of that ambient set, not a row-by-row residual
computation.
\Cref{prop:bc-sparse-active-low-low-degree-full-budget} closes the sparse
residual-degree part of this ambient set under the relaxed final-zone budget:
if the total residual low-low degree \(D=\sum_a d_a\) satisfies
\((D-1)!!\le2^{2m}\), then the whole degree-constrained ambient set already
fits in the full branch word.  In particular it also closes the active-label
case \((2r-1)!!\le2^{2m}\), where \(r=\#\{a:d_a>0\}\).  Thus the live low-low
source extraction starts only after this double-factorial bound fails.
\Cref{prop:bc-degree-two-symmetry-low-low-full-budget} shrinks that remaining
branch further: if \(r_2=\#\{a:d_a=2\}\), then each low-low assignment has at
least \(2^{\lceil r_2/2\rceil}\) labelled-port realizations, so the whole
degree-pruned ambient set still fits the full final-zone budget whenever
\[
  (D-1)!!\le 2^{2m+\lceil r_2/2\rceil}.
\]
\Cref{def:bc-exact-low-low-degree-profile-count,prop:bc-exact-degree-profile-low-low-full-budget}
replace these sufficient double-factorial tests by the exact labelled
max-degree-two profile count
\(\mathsf N(r_1,r_2)\), where \(r_i=\#\{a:d_a=i\}\).  Only the exact
degree-profile dense cases
\(\mathsf N(r_1,r_2)>2^{2m}\) still need an additional target-window low-low
candidate source.  The certificate-backed
\cref{prop:bc-active-twenty-low-low-profile-full-budget} further shows that
every such dense case has at least \(A(m)\) positive residual labels, where
\(A(m)\) is the exact active-support dense-profile frontier defined there, and
total residual low-low degree at least \(B(m)\), where \(B(m)\) is the
corresponding exact total-degree frontier.  Equivalently,
\cref{cor:bc-continuation-heavy-low-low-profile-full-budget} closes the
continuation-heavy complement: if \(r_0=\#\{a:d_a=0\}\) and
\(C=\sum_a c_a\), the live dense branch has
\[
  r_0\le2m-A(m),\qquad C\le4m-B(m).
\]
\Cref{prop:bc-degree-low-low-ambient-too-large} records the matching
obstruction to a degree-only refinement: residual degrees $d_a=1$ already give
$(2m-1)!!$ low-low assignments, exceeding the $2^m$ branch budget for
$m\ge3$.  Thus the refinement has to use the target-window signature or
equivalent semistandard data, not only the recovered residual degrees.
\Cref{lem:bc-strict-left-alone-low-low-obstruction} gives the concrete
four-label version of the same obstruction inside the Step~3 construction:
two examples have the same strict-left continuation subfilling and the same
residual low-label degrees, but different $L_4$ fillings.
\Cref{cor:bc-refined-omega-induces-low-low-candidates} gives the current
refinement mechanism: any parser-visible refined odd-packet class
$\Omega_{\mathrm{tw}}$ of size at most $2^m$ induces, by inverse Step~2
uniqueness, a low-low candidate class of the same size bound.  Thus the live
target can be phrased either as a bounded refined odd-packet class or as the
induced bounded low-low class.
\Cref{cor:bc-target-window-signature-fiber-low-low-candidates} eliminates the
separate refined-class construction from this live target: once the
target-window data recover the base tight-prefix endpoint, retained even
strips, labelled continuation boxes, and the true commuting signature
$\kappa$, the bounded low-low class is the induced image of the signature
fiber $\Omega_{\mathrm{tw}}(\kappa)$.
\Cref{cor:bc-bounded-signature-list-full-budget-omega} records the relaxed
full-branch version: exact recovery of \(\kappa\) is stronger than necessary if
the target-window data instead recover an at-most-\(2^m\) list of possible
commuting signatures containing the true one.
\Cref{lem:bc-step3-full-half-square-determines-rho,cor:bc-bounded-low-low-candidates-recover-right-boundary}
then plug this branch-budgeted low-low recovery into the actual local inverse:
after the first $m$ branch bits select the true low-low filling, the
strict-left continuation columns complete the Step~3 half-square, whose
forward boundary is the original terminal $\rho$-subpath.
\Cref{cor:bc-target-window-signature-fiber-right-boundary-inverse} packages
the resulting final-facing sufficient condition: base data plus $\kappa$
from the target window give the tight right-boundary inverse.
\Cref{prop:bc-target-window-signature-supplies-right-boundary-auxiliary}
then inserts this inverse into the global right-boundary auxiliary criterion:
with $A(P)=\mathcal D_{\mathrm{out}}(P)$, the final tight-zone inverse is a
fixed function of the allowed parser data and the already budgeted branch word.
\Cref{cor:bc-small-omega-right-boundary-inverse,prop:bc-small-omega-supplies-right-boundary-auxiliary}
supplies the matching right-boundary inverse for the already closed small
leaves $m\le8$ and inserts it into the global auxiliary criterion: in those
cases the full tight compatibility set is itself small enough, so no
target-window signature is needed.
\Cref{prop:bc-split-target-window-supplies-right-boundary-auxiliary} packages
the split condition actually needed by the global supplier: base target-window
data suffice for $m\le8$ and, by the full-branch broad-$\Omega$ bound, also
for $9\le m\le28$; only the range $m\ge29$ still requires the commuting
normal-form signature $\kappa$.
\Cref{prop:bc-even-chain-supplies-large-right-boundary-auxiliary} removes
$\kappa$ as an independent large-zone readout target: in the remaining
supplier range $m\ge29$ (the statement is available already for $m\ge9$), it
is enough that the same target-window tuple recover the even standard-time
endpoint chain through time $4m$ and the labelled continuation boxes.
\Cref{prop:bc-low-label-boxes-supply-large-right-boundary-auxiliary} pushes
this to the concrete box readout: recovering the two boxes carrying each
label $1,\ldots,2m$, together with the labelled continuation boxes, supplies
the large final-zone inverse.
\Cref{prop:bc-low-low-entries-supply-large-right-boundary-auxiliary} pushes
the same large-zone criterion to the Step~3 data already split by the
strict-left machinery: it is enough to recover the low-low triangle entries
$L_{2m}$, since the strict-left outside datum supplies the continuation
columns.
\Cref{prop:bc-terminal-boundary-supplies-large-right-boundary-auxiliary}
records the branch-free boundary version of this criterion: recovering the
terminal low-label boundary segment
$\rho_{2(j-2m)},\ldots,\rho_{2j}$ recovers $L_{2m}$ and hence supplies the
large final-zone inverse.
\Cref{prop:bc-full-budget-terminal-boundaries-supply-large-right-boundary-auxiliary}
uses the whole final-zone branch word directly: a canonical at-most-$2^{2m}$
list of possible terminal boundary segments containing the true one is already
enough to recover the original final-zone subpath.
\Cref{prop:bc-bounded-terminal-boundaries-supply-large-right-boundary-auxiliary}
is the half-budget boundary version used by the refined-rank route: a
canonical at-most-$2^m$ list of possible terminal boundary segments is enough,
because each candidate has a deterministic low-low subfilling.
\Cref{prop:bc-bounded-low-low-candidates-supply-large-right-boundary-auxiliary}
records the corresponding half-budget low-low form.
\Cref{prop:bc-full-budget-low-low-candidates-supply-large-right-boundary-auxiliary}
is the direct final-zone version: if the target-window data determine a
canonical at-most-$2^{2m}$ list of possible $L_{2m}$ fillings containing the
true one, all $2m$ branch bits may select that filling and the full Step~3
half-square then recovers the final-zone subpath.
\Cref{prop:bc-refined-omega-supplies-large-right-boundary-auxiliary} gives the
current supplier-side form: it is enough for the target-window data to recover
a bounded refined odd-packet class $\Omega_{\mathrm{tw}}$, since its induced
low-low candidates inject into it.
\Cref{prop:bc-full-budget-omega-supplies-large-right-boundary-auxiliary}
records the relaxed full-branch version: an at-most-$2^{2m}$ odd-packet class
also suffices when all final-zone branch bits are used directly for the
right-boundary inverse.
\Cref{cor:bc-full-budget-omega-induces-low-low-candidates} records the same
source in the low-low language: the induced low-low candidates still inject
into the full-budget odd-packet class.
\Cref{prop:bc-bounded-signature-list-supplies-large-right-boundary-auxiliary}
is the corresponding parser-facing signature-list supplier: an at-most-\(2^m\)
list of commuting signatures gives an at-most-\(2^{2m}\) odd-packet class by
the binary signature-fiber theorem, and
\cref{cor:bc-bounded-signature-list-low-low-candidates} converts it to the
same full-budget low-low candidate source.
\Cref{prop:bc-nine-to-twelve-broad-omega-supplies-large-right-boundary-auxiliary}
uses the existing binary-prefix factorial bound to close the unrefined broad
compatibility set for $9\le m\le12$ under this relaxed budget; those cases no
longer require the commuting-signature refinement once the base target-window
data are recovered.
\Cref{prop:bc-thirteen-to-twenty-eight-envelope-omega-supplies-large-right-boundary-auxiliary}
uses the exact recursive fixed-even envelope certificate to close the
unrefined broad compatibility set for $13\le m\le28$ under the same relaxed
budget.
\Cref{prop:bc-fixed-even-envelope-m29-obstruction} verifies exactly that the
same recursive envelope already exceeds the full branch budget at $144$
non-column-even frontier shapes for $m=29$, and that all other $m=29$ shapes
remain below budget.  It also verifies exactly $346$ over-budget top-fixed
pairs \((\gamma,e)\), exactly $62$ over-budget top-two triples, and no
over-budget top-three quadruple at \(m=29\).  Thus
\cref{prop:bc-envelope-nonfrontier-omega-supplies-large-right-boundary-auxiliary}
closes every broad-\(\Omega\) case off the top-fixed envelope frontier
\(\mathcal P^{\mathrm{env}}_m\), while
\cref{prop:bc-twenty-nine-top-three-envelope-omega-supplies-large-right-boundary-auxiliary}
closes the remaining \(m=29\) frontier by using the next two retained even
cells.
\Cref{prop:bc-fixed-even-envelope-m30-frontier-certificate,prop:bc-thirty-fixed-depth-envelope-omega-supplies-large-right-boundary-auxiliary}
closes the \(m=30\) frontier at fixed depth seven, and
\cref{prop:bc-fixed-even-envelope-m31-frontier-certificate,prop:bc-thirty-one-fixed-depth-envelope-omega-supplies-large-right-boundary-auxiliary}
closes the \(m=31\) frontier at fixed depth ten.
\Cref{prop:bc-envelope-top-two-nonfrontier-omega-supplies-large-right-boundary-auxiliary,prop:bc-fixed-depth-nonfrontier-omega-supplies-large-right-boundary-auxiliary}
then close every \(m\ge32\) case outside some target-window-determined
fixed-depth envelope frontier.  The live broad-\(\Omega\) branch therefore
starts at \(m\ge32\) and must remain over budget at every fixed depth.
\Cref{lem:bc-binary-prefix-full-budget-stops-at-twelve} records the matching
threshold check: the same factorial estimate exceeds the full branch budget
from $m=13$ onward; the certified recursive envelope supplies the replacement
bound through $m=28$.
\Cref{prop:bc-strict-left-incidence-carrier-supplies-large-right-boundary-auxiliary}
gives the concrete carrier form: a target-window/strict-left rule that
identifies a Step~3 carrier containing $H_{2m}$ supplies the needed low-low
entries and hence the large final-zone inverse.
\Cref{prop:bc-strict-left-low-low-carrier-supplies-large-right-boundary-auxiliary}
sharpens the live carrier target: because the strict-left outside datum already
supplies the cross-incidence part, the carrier for the remaining large range
only has to contain the terminal low-low cell set $L_{2m}$.
\Cref{cor:bc-target-window-strict-left-carrier-closes-signature} inserts the
already recovered strict-left $\rho$-segment into this criterion, leaving only
the carrier-identification, Step~2 subarrangement, and left/bottom boundary
readout step from that segment together with the target window and endpoint
pair.
\Cref{cor:bc-target-window-strict-left-q-prefix-strip-closes-signature} records
the corresponding anchored-strip version: it is enough to identify a Step~3
carrier whose recovered half-matrix contains the Step~2 $Q$-prefix strip
$C_{2m}$ up to transpose.
\Cref{cor:bc-target-window-strict-left-low-label-incidence-closes-signature}
is the same statement in incidence form: the carrier must contain the cell set
$H_{2m}$.
\Cref{prop:bc-target-window-signature-closes-refined-suffix} makes the data
source explicit: it is enough for the right-boundary replacement window
$V^\ast_L$, together with the already recovered outside datum, to determine
that signature.
\Cref{def:bc-viable-active-odd-cells,prop:bc-binary-viable-branching-closes-interval-leaf}
give a sufficient binary-branching criterion for restricted instances of this
count, while \cref{lem:bc-binary-viable-branching-obstruction} shows that
this criterion is too strong even before the broader counterexample above.
\Cref{lem:bc-reverse-corner-commuting-or-blocked} gives the local
branch-control dichotomy: each odd corner removal either commutes past the
fixed even removal, or is immediately blocked by a fixed even box to its east
or south.
\cref{prop:bc-right-boundary-auxiliary-global-recovery,prop:bc-right-boundary-auxiliary-fiber-criterion}
record the resulting recovery criterion: a right-boundary tight zone may use
only such already recovered outside data as an auxiliary input, with no new
branch words and with the same $2^{2q}s_{j-q}$ fixed-fiber count.
The development archive's split target-window supplier inserts the split
inverse into this route: base target-window data handle
the right-boundary $C$-side alternative through $m=28$, through the certified
\(m=29,30,31\) envelope frontiers, and through every fixed-depth non-frontier
case.  Only the \(m\ge32\) branch that remains in every fixed-depth envelope
frontier still needs the commuting-signature or bounded signature-list readout
in this split form.
Its terminal-source variant gives the last readouts: for $m\ge29$ in the
non-column-even final-zone branch, the degree-pruned ambient set already closes
the sparse residual-degree cases of
\cref{prop:bc-sparse-active-low-low-degree-full-budget} and the degree-profile
cases of \cref{prop:bc-degree-two-symmetry-low-low-full-budget}, and then the
exact degree-profile cases of
\cref{prop:bc-exact-degree-profile-low-low-full-budget}; only the cases with
\(\mathsf N(r_1,r_2)>2^{2m}\), hence at least \(A(m)\) positive residual
labels and total residual degree at least \(B(m)\) by
\cref{prop:bc-active-twenty-low-low-profile-full-budget}, still need a
source beyond the strict-left degree equations.  The supplier now accepts
either a partial low-low readout whose residual profile falls back under the
exact full-budget frontier, by
\cref{prop:bc-partial-low-low-profile-full-budget}, a full-budget odd-packet
class, a bounded commuting-signature list, or a canonical full-budget list of
possible \(L_{2m}\) low-low fillings containing the true one; the odd-packet
and signature-list alternatives now both feed this last artifact through
\cref{cor:bc-full-budget-omega-induces-low-low-candidates,cor:bc-bounded-signature-list-low-low-candidates}.
\Cref{cor:bc-partial-low-low-edge-mass-frontier-full-budget} gives a
checkable sufficient form of the partial-readout alternative: enough recovered
low-low edge mass forces the residual total degree below \(B(m)\).
\Cref{cor:bc-incidence-covered-low-labels-frontier-full-budget} gives the
parser-facing sufficient form of this partial-readout alternative: an
incidence readout on enough low labels closes the source as soon as the
unexposed complement is below the certified active-support or degree frontier.
In the same branch,
\cref{cor:bc-large-incidence-cover-frontier-full-budget} records the
corresponding cardinality and covered-degree tests for this complement
condition, and
\cref{cor:bc-incidence-cover-positive-projection} normalizes those tests to the
positive-residual part of the cover.  Finally,
\cref{cor:bc-continuation-heavy-low-low-profile-full-budget} gives the
parser-side equivalent bounds \(r_0\le2m-A(m)\) and \(C\le4m-B(m)\) for the
fully continued labels and total low-continuation incidence.  A full-budget
list of possible terminal Step~3 boundary segments
\(\rho_{2(j-2m)},\ldots,\rho_{2j}\), and in particular exact recovery of that
segment, induces such a low-low candidate list by
\cref{cor:bc-full-budget-terminal-boundaries-induce-low-low-candidates}.  The
earlier carrier and Hall-data formulations have all been reduced to this
single terminal artifact.
\Cref{prop:bc-column-even-prefix-supplies-large-right-boundary-auxiliary}
removes one large subcase from that terminal target: if the tight prefix endpoint
is already column-even, the closed odd-alternating full-zone decoder supplies
the final right-boundary inverse.  The live $m\ge29$ carrier problem is
therefore the non-column-even tight-prefix endpoint case.  In that remaining
case the useful auxiliary object is not a row-indexed table: it is the direct
endpoint-defect/continuation incidence structure, the canonical Hall charge, and
the Hall-selected continuation blocks recovered from that Hall charge and the
outside datum, as isolated in
\cref{prop:bc-target-window-hall-augmented-low-low-carrier-supplies-large-right-boundary-auxiliary,prop:bc-target-window-hall-augmented-incidence-cover-carrier-supplies-large-right-boundary-auxiliary,cor:bc-direct-hall-charge-selected-blocks-recovered}.
\Cref{cor:bc-shared-direct-hall-incidence-cover-carrier-terminal-source}
records the corresponding shared terminal-source form: the same
Hall-augmented incidence-cover carrier also supplies the large-complement
deleted-only graph branch.
\Cref{cor:bc-target-window-q-suffix-hall-carrier-supplies-large-right-boundary-auxiliary}
narrows the continuation-segment part further to recovery of the internal
Step~2 $Q$-cut between the prefix and suffix strips; the entries of that suffix
are already supplied by the strict-left Step~3 columns, and the rest of the
incoming suffix boundary is globally empty.
\Cref{cor:bc-strict-left-low-low-determines-internal-q-cut} ties this boundary
target back to the carrier target: once the target-window mechanism supplies the
low-low entries on $L_{2m}$, the internal $Q$-cut and the continuation segment
are also recovered.  Thus a direct carrier containing $L_{2m}$ closes the
non-column-even branch without first recovering the continuation segment, and
\cref{cor:bc-shared-low-low-carrier-terminal-source} records that the same
carrier also supplies the target-window signature input.
\Cref{prop:bc-strict-left-incidence-cover-carrier-supplies-large-right-boundary-auxiliary}
is the corresponding partial-carrier version: it is enough for the carrier to
contain all low-low cells incident to a determined set of low labels, provided
the unexposed complement is below the certified \(A(m),B(m)\) frontier.  The
weaker full-budget candidate-list route closes it after the final-zone branch
word selects the true low-low filling.  Its graph counterpart,
\cref{cor:bc-strict-left-incidence-cover-carrier-recovers-large-deleted-only-graph},
shows that this same partial carrier also recovers the large deleted-only graph;
the Hall route remains available when the internal $Q$-cut or direct Hall data
are recovered first.
\Cref{cor:bc-hall-data-alone-packet-obstruction} shows the limitation of this
auxiliary object: Hall data alone do not determine the residual odd-packet
matching, so the remaining carrier rule must also use the actual final
target-window or Step~3 boundary information.
\Cref{cor:bc-direct-hall-blocks-alone-packet-obstruction} shows that adding the
Hall-selected continuation pieces does not remove this limitation; those pieces
are an interface for a carrier or bounded-candidate source, not the source
itself.
\cref{lem:bc-krattenthaler-semistandard-path-boundary} identifies the
semistandard growth path used here with the top/right boundary of the Step~3
dual-RSK' half-square in Krattenthaler's Theorem~4 construction.
\cref{cor:bc-parsed-dual-rsk-prime-boundary-recovers-subfilling} is the
matching parser-to-filling bridge for that boundary: once the relevant
top/right boundary of an outside dual-RSK' subdiagram is parsed,
Krattenthaler's backward local rules recover the subfilling without further
choices.
\cref{lem:bc-halo-zones-frontier-packets} assigns the ordered frontier-event
packets to those same canonical zones, so the remaining local construction can
be made zone-by-zone.  The packet branch data are still available inside a
zone by the deterministic splitting of
\cref{lem:bc-zone-branch-word-splitting}.
If the zone replacement window exposes the retained halo blocks, then
\cref{lem:bc-zone-retained-blocks-expose-packet-endpoints} recovers the
endpoint pairs needed by the packet-local HKKO inverses.
\cref{lem:bc-packet-endpoints-adjacent-halo} sharpens this endpoint recovery:
each packet endpoint needs only the adjacent retained halo block on that side,
or the zone endpoint if the packet touches the zone boundary.
\cref{lem:bc-all-retained-halo-packet-adjacent} records the limitation of
this saving: if separate packet inverses are called for every deleted
component, the adjacent endpoint blocks range over all retained halo indices.
Thus a true target-size saving must come from a merged endpoint rule, a
packet-compression rule, or a zero-window packet mechanism, not merely from
packetwise endpoint lookup.
\cref{prop:bc-zone-packet-inverses-assemble} then assembles those packet
inverses and retained halo blocks into the full zone inverse.
\cref{lem:bc-zone-no-spare-packet-windows} shows that positive packet outputs
cannot be exposed as additional disjoint windows after all retained halo blocks
have already been exposed unchanged.
\cref{lem:bc-halo-zone-parser-budget} gives the sharper budget: disjoint
positive packet output windows must be paid for by the same number of retained
halo windows that are not exposed unchanged, and raw length-$|K|$ packet
outputs cannot fit in zones with fewer retained than deleted blocks.
\cref{cor:bc-packet-compression-budget} refines this into an exact compression
accounting rule: each unit by which packet outputs are shortened buys one more
visible retained halo block, after the local surplus has been spent.
\cref{cor:bc-raw-packet-zone-surplus} makes this exact: raw packet outputs can
expose at most the local surplus
$|Z_\gamma\setminus I|-|Z_\gamma\cap I|$ retained blocks.
\cref{cor:bc-tight-zone-raw-packets-hide-retained-blocks} records the tight
case: when the raw packet lengths exactly fill the target-size window, no
retained halo block is exposed unchanged, so packet endpoints must come from a
merged decoder rather than from visible retained blocks.
\cref{lem:bc-alternating-witness-full-halo-zone} shows that the alternating
witness sets at $j=2q$ are exactly such single-zone tight cases.
\cref{prop:bc-alternating-full-zone-decoder-criterion} isolates the resulting
first hard positive target: a full-zone alternating repair is precisely an
injection into one target path of length $q$ and one $2q$-bit branch word.
\cref{lem:bc-odd-alternating-height-every-row} proves the first structural
fact for that target: on the odd alternating height side, every step creates a
new row, while the even alternating height side is empty.
\cref{lem:bc-odd-alternating-height-second-box-encoding} then shows that the
off-column boxes determine the path and are canonically paired by column in
the terminal column-even case.
\cref{lem:bc-odd-alternating-height-row-choice-encoding} compresses those
off-column boxes further to their row choices.
\cref{lem:bc-odd-alternating-height-standard-tableau} identifies the resulting
odd alternating height fiber with even-column standard Young tableaux of size
$2q$.
\cref{lem:bc-odd-alternating-height-schensted-contraction} contracts the
Schensted matching of such a tableau to the stable graph target whenever no
adjacent-pair loop is present.
\cref{lem:bc-looped-degree-two-graph-repair} repairs those adjacent-pair loops
by a canonical loopless graph plus the loop-set bits, and
\cref{prop:bc-odd-alternating-height-full-zone-decoder} closes the odd
alternating height full-zone decoder within the same $2q$-bit branch budget.
\cref{prop:bc-even-alternating-height-full-zone-empty,cor:bc-height-alternating-full-zone-decoder}
then close the height side of the alternating full-zone obstruction; the
\cref{lem:bc-odd-alternating-width-forced-path,prop:bc-width-alternating-full-zone-decoder}
close the width side: the odd width case has a single forced rectangular path,
and the even width case is empty.  Thus
\cref{cor:bc-alternating-full-zone-decoder-closed} closes the alternating
full-zone obstruction.  The remaining canonical-halo work is the
non-alternating zone-local decoder.
\cref{prop:bc-count-compatible-halo-zone-package} isolates the remaining
count-compatible form: the target-size zone window must decode the retained
halo blocks and the packet outputs from the same replacement blocks, either by
merging them or by using zero-window packet data.
\cref{prop:bc-below-wall-packet-splices-give-halo-zone-inverses} then plugs
the below-wall HKKO splice into this count-compatible form: once the
target-size window recovers the retained halo blocks, the packet endpoints
determined from those retained blocks and the existing packet branch words
canonically construct the image-compatible packet outputs.  Thus the parser no
longer has to output those packet objects separately; it must supply the
retained blocks and the fixed readouts from the resulting marked preimages.
\cref{lem:bc-local-branch-boundary-readout} closes the bad-boundary part of
those fixed readouts: width boundary second-copy positions are selected by the
orientation half of the local branch word, and height boundary second copies
never remain in the bad first column.
\cref{lem:bc-canonical-halo-complement-exposes-zone-endpoints} removes the
remaining internal zone-endpoint datum: unchanged complement blocks determine
all canonical-zone endpoints, using the fixed initial and terminal endpoints
of the growth-path model at boundary zones.
\cref{prop:bc-below-wall-packet-splice-halo-local-repair-criterion} promotes
this zone-local inverse to the fixed-witness fiber estimate for any canonical
halo parser with those below-wall packet-splice readouts.
\end{corollary}

\begin{proof}
By \cref{lem:bc-growth-path-model}, column-even semistandard tableaux of
content $(2^j)$ are equivalently represented by two-step growth paths with one
two-step block of total weight $2$ for each label.  Thus a fixed witness set
$S$ identifies $q$ labelled two-step blocks, using the label-to-block reversal
in \cref{lem:bc-growth-path-model}.  If the displayed injection exists
for either the width-witness or height-witness family, then the target set has
cardinality $2^{2q}s_{j-q}$, because \cref{lem:bc-littlewood-graph-model} and
\cref{lem:bc-growth-path-model} count the same column-even tableaux of content
$(2^{j-q})$.  This is exactly the corresponding fiber estimate in
\cref{prop:bc-bad-shape-fiber-reduction}.
\end{proof}

\begin{definition}[Raw block deletion datum]\label{def:bc-raw-block-deletion}
Let
\[
  P=(\rho_0,\rho_1,\ldots,\rho_{2j})
\]
be a two-step growth path as in \cref{lem:bc-growth-path-model}, and let
$S\subset\{1,\ldots,j\}$.  Put
\[
  I(S)=\{j-s+1:s\in S\}.
\]
The raw block deletion datum $R(P,S)$ is the ordered list of triples
\[
  B_i(P)=(\rho_{2i-2},\rho_{2i-1},\rho_{2i}),\qquad i\notin I(S),
\]
together with deletion-joint markers at the following places: before the first
retained index $i$ when $i>1$, with boundaries $\varnothing$ and
$\rho_{2i-2}$; between two consecutive retained indices $i<i'$ when
$i'>i+1$, with boundaries $\rho_{2i}$ and $\rho_{2i'-2}$; and after the last
retained index $i$ when $i<j$, with boundaries $\rho_{2i}$ and $\varnothing$.
\end{definition}

\begin{lemma}[Raw deletion localizes the repair]\label{lem:bc-raw-deletion-joints}
Let $P$ and $S$ be as in \cref{def:bc-raw-block-deletion}.  Every retained
triple in $R(P,S)$ is a two-step block with vertical-strip adjacent differences
and weight $2$:
\[
  |\rho_{2i-2}|-2|\rho_{2i-1}|+|\rho_{2i}|=2 .
\]
The only failures of $R(P,S)$ to be a single two-step growth path with
$j-|S|$ blocks can occur at deletion joints, and the number of deletion joints
is at most $|S|$.
\end{lemma}

\begin{proof}
For $i\notin I(S)$, the retained triple is one of the original two-step
blocks of $P$, so the vertical-strip and weight-$2$ conditions are inherited
from \cref{lem:bc-growth-path-model}.  If two consecutive retained indices
$i<i'$ satisfy $i'=i+1$, then the right endpoint of $B_i(P)$ is
$\rho_{2i}=\rho_{2i'-2}$, so the two blocks concatenate without a new
condition.  If the first retained index is $1$, the raw datum begins at
$\rho_0=\varnothing$; otherwise the initial endpoint condition is exactly the
first marked deletion joint.  Similarly, if the last retained index is $j$,
the raw datum ends at $\rho_{2j}=\varnothing$; otherwise the terminal endpoint
condition is exactly the last marked deletion joint.  Finally, if two
consecutive retained indices satisfy $i'>i+1$, the blocks were separated by at
least one deleted block, and the only missing internal path condition is the
compatibility of the two boundary partitions $\rho_{2i}$ and $\rho_{2i'-2}$ at
that splice.  Thus all repair data are concentrated at the marked deletion
joints.  These joints are in bijection with maximal consecutive intervals of
deleted indices, so there are at most $|I(S)|=|S|$ of them.
\end{proof}

\begin{proposition}[Local inverses give global fixed-witness recovery]\label{prop:bc-local-inverses-global-recovery}
Let
\[
  P=(\rho_0,\rho_1,\ldots,\rho_{2j})
\]
be a two-step growth path as in \cref{lem:bc-growth-path-model}, let
$S\subset\{1,\ldots,j\}$ have cardinality $q$, and put
$I(S)=\{j-s+1:s\in S\}$.  Write the maximal consecutive intervals of $I(S)$
as
\[
  K_\alpha=\{u_\alpha,u_\alpha+1,\ldots,v_\alpha\}
  \qquad(1\le\alpha\le L)
\]
in increasing block order.  Suppose a repair assignment sends $P$ to
\[
  \Bigl(P^\ast,(\beta_\alpha)_{\alpha=1}^L\Bigr)
\]
and satisfies the following three properties.

\begin{enumerate}
\item There is a deterministic parser, depending only on $S$, which applied to
$P^\ast$ returns the retained windows outside the intervals $K_\alpha$ and
local output windows $U^\ast_\alpha$ for the intervals $K_\alpha$.
\item The retained windows returned by the parser are exactly the original
two-step blocks
\[
  (\rho_{2i-2},\rho_{2i-1},\rho_{2i})\qquad(i\notin I(S))
\]
in their original order.
\item For each $\alpha$, the parser returns the endpoint pair
\[
  \mu_\alpha=\rho_{2u_\alpha-2},\qquad
  \nu_\alpha=\rho_{2v_\alpha},
\]
and a fixed local inverse recovers the deleted subpath
\[
  (\rho_{2u_\alpha-2},\rho_{2u_\alpha-1},\ldots,\rho_{2v_\alpha})
\]
from $U^\ast_\alpha$, $\beta_\alpha$, $\mu_\alpha$, $\nu_\alpha$, and
$K_\alpha$.
\end{enumerate}

Then the original path $P$ is uniquely recoverable from
$S$, $P^\ast$, and the words $\beta_1,\ldots,\beta_L$.
Consequently, after composing with the bijection
\cref{lem:bc-pieri-path-bijection}, the original semistandard tableau is
uniquely recoverable from the same data.
\end{proposition}

\begin{proof}
The set $S$ determines $I(S)$ and hence the ordered list of maximal intervals
$K_\alpha$.  Applying the deterministic parser to $P^\ast$ gives all retained
windows and all local output windows $U^\ast_\alpha$ in this same order.  By
the second hypothesis, every original two-step block with index outside
$I(S)$ is therefore recovered.

For each maximal deleted interval $K_\alpha$, the third hypothesis recovers
the whole deleted subpath from the corresponding local output window, branch
word, endpoints, and the known interval $K_\alpha$.  The intervals $K_\alpha$
and the retained block indices are disjoint and together cover
$\{1,\ldots,j\}$.  Interleaving the recovered deleted subpaths with the
recovered retained two-step blocks in block order reconstructs the full
sequence
\[
  \rho_0,\rho_1,\ldots,\rho_{2j}.
\]
The endpoint overlaps agree because the recovered deleted subpaths and
retained windows are all subpaths of the same original path at their common
boundary partitions.  Hence $P$ is uniquely recovered.

Finally, \cref{lem:bc-pieri-path-bijection} is a bijection between the
semistandard tableaux of content $(2^j)$ and their Pieri two-strip paths, so
recovering $P$ recovers the original tableau.
\end{proof}

\begin{proposition}[Replacement zones give global fixed-witness recovery]\label{prop:bc-replacement-zones-global-recovery}
Let
\[
  P=(\rho_0,\rho_1,\ldots,\rho_{2j})
\]
be a two-step growth path as in \cref{lem:bc-growth-path-model}, let
$S\subset\{1,\ldots,j\}$ have cardinality $q$, and put
$I(S)=\{j-s+1:s\in S\}$.  Let
\[
  Z_\alpha=\{a_\alpha,a_\alpha+1,\ldots,b_\alpha\}
  \qquad(1\le\alpha\le L)
\]
be disjoint intervals, in increasing order, such that every index in $I(S)$
belongs to exactly one $Z_\alpha$.  Put
\[
  R=\{1,\ldots,j\}\setminus\bigcup_{\alpha=1}^L Z_\alpha .
\]
Suppose a repair assignment sends $P$ to
\[
  \Bigl(P^\ast,(\beta_\alpha)_{\alpha=1}^L\Bigr)
\]
and satisfies the following properties.

\begin{enumerate}
\item The intervals $Z_\alpha$ are determined by $S$.
\item $P^\ast$ is a two-step path with exactly $j-q$ two-step windows, and a
deterministic parser depending only on $S$ decomposes those windows, in order,
into unchanged retained blocks for indices $i\in R$ and local replacement
windows $V^\ast_\alpha$ for the intervals $Z_\alpha$.
\item For each $\alpha$, the local replacement window $V^\ast_\alpha$ has
exactly
\[
  |Z_\alpha\setminus I(S)|
\]
two-step blocks.
\item The unchanged retained blocks returned by the parser are exactly
\[
  (\rho_{2i-2},\rho_{2i-1},\rho_{2i})\qquad(i\in R)
\]
in their original order.
\item For each $\alpha$, the parser returns the endpoint pair
\[
  \mu_\alpha=\rho_{2a_\alpha-2},\qquad
  \nu_\alpha=\rho_{2b_\alpha},
\]
and a fixed local inverse recovers the original subpath
\[
  (\rho_{2a_\alpha-2},\rho_{2a_\alpha-1},\ldots,\rho_{2b_\alpha})
\]
from $V^\ast_\alpha$, $\beta_\alpha$, $\mu_\alpha$, $\nu_\alpha$, and
$Z_\alpha$.
\end{enumerate}

Then the original path $P$ is uniquely recoverable from
$S$, $P^\ast$, and the words $\beta_1,\ldots,\beta_L$.
\end{proposition}

\begin{proof}
The set $S$ determines $I(S)$ and, by hypothesis, the intervals
$Z_\alpha$ and the complement $R$.  Applying the deterministic parser to
$P^\ast$ gives the unchanged retained windows for $R$ and the replacement
windows $V^\ast_\alpha$ in their original order.  By the fourth hypothesis,
all original blocks with indices in $R$ are recovered.

For each replacement interval $Z_\alpha$, the fifth hypothesis recovers the
whole original subpath from the corresponding parsed replacement window,
branch word, endpoints, and known interval.  The intervals $Z_\alpha$ and the
singleton indices in $R$ partition $\{1,\ldots,j\}$.  Interleaving the
recovered replacement-zone subpaths with the recovered retained blocks in
block order reconstructs the full sequence
\[
  \rho_0,\rho_1,\ldots,\rho_{2j}.
\]
The endpoint overlaps agree because every recovered piece is recovered with
the endpoint pair it had in the original path, and every unchanged retained
block is the original block.  Thus $P$ is uniquely recovered.
\end{proof}

\begin{proposition}[Right-boundary auxiliary replacement-zone recovery]\label{prop:bc-right-boundary-auxiliary-global-recovery}
Work in the setting and notation of
\cref{prop:bc-replacement-zones-global-recovery}.  Suppose that the zones are
ordered increasingly and that the last one is a right-boundary zone,
\[
  Z_L=\{a_L,\ldots,j\}.
\]
Assume properties (1)--(4) of
\cref{prop:bc-replacement-zones-global-recovery}.  Assume also that the
parser returns the endpoint pair
\[
  \rho_{2a_\alpha-2},\qquad \rho_{2b_\alpha}
\]
for every zone.  For every $\alpha<L$, suppose a fixed local inverse recovers
the original subpath on $Z_\alpha$ from
$V^\ast_\alpha$, $\beta_\alpha$, that endpoint pair, and $Z_\alpha$.

Let the recovered outside datum be
\[
  \mathcal D_{\mathrm{out}}(P)
  =
  \left(
    (B_i(P))_{i\in R},
    \bigl((\rho_{2a_\alpha-2},\ldots,\rho_{2b_\alpha})\bigr)_{\alpha<L}
  \right).
\]
Suppose that an auxiliary datum $A(P)$ for the right-boundary zone is a fixed
function of $S$ and $\mathcal D_{\mathrm{out}}(P)$, and that a fixed inverse
recovers the original subpath on $Z_L$ from
\[
  V^\ast_L,\quad \beta_L,\quad
  \rho_{2a_L-2},\rho_{2j},\quad Z_L,\quad A(P).
\]
Then the original path $P$ is uniquely recoverable from
$S$, $P^\ast$, and the words $\beta_1,\ldots,\beta_L$.
\end{proposition}

\begin{proof}
Applying the deterministic parser to $P^\ast$ gives the unchanged retained
blocks for $R$, the replacement windows $V^\ast_\alpha$, and the endpoint
pair of every zone.  By property (4) of
\cref{prop:bc-replacement-zones-global-recovery}, the retained blocks in $R$
are the original blocks.

For each $\alpha<L$, the assumed local inverse recovers the original subpath
on $Z_\alpha$.  Therefore the datum
$\mathcal D_{\mathrm{out}}(P)$ is determined from
$S$, $P^\ast$, and the branch words: its first component is supplied by the
parsed retained blocks in $R$, and its second component is supplied by the
subpaths just recovered for the zones strictly left of $Z_L$.

Since $A(P)$ is a fixed function of $S$ and this recovered outside datum,
$A(P)$ is also recovered.  The right-boundary inverse then recovers the
original subpath on $Z_L$.  The singleton indices in $R$ and the intervals
$Z_1,\ldots,Z_L$ partition $\{1,\ldots,j\}$.  Interleaving the recovered
subpaths and retained blocks in block order reconstructs the full sequence
\[
  \rho_0,\rho_1,\ldots,\rho_{2j}.
\]
The endpoint overlaps agree for the same reason as in
\cref{prop:bc-replacement-zones-global-recovery}: each recovered piece carries
the endpoint pair it had in the original path.  Hence $P$ is uniquely
recoverable.
\end{proof}

\begin{lemma}[Canonical halo replacement zones]\label{lem:bc-canonical-halo-replacement-zones}
Let $j\ge1$ and let $S\subset\{1,\ldots,j\}$.  Put
\[
  I(S)=\{j-s+1:s\in S\}
\]
and define the one-block halo of $I(S)$ by
\[
  H(S)=\{r\in\{1,\ldots,j\}:\ |r-i|\le1
          \text{ for some }i\in I(S)\}.
\]
Write the maximal consecutive intervals of $H(S)$ as
\[
  Z_\alpha=\{a_\alpha,a_\alpha+1,\ldots,b_\alpha\}.
\]
Then the intervals $Z_\alpha$ are determined by $S$, are pairwise disjoint,
and cover every deleted block index in $I(S)$.  If
\[
  R=\{1,\ldots,j\}\setminus\bigcup_\alpha Z_\alpha,
\]
then
\[
  |R|+\sum_\alpha |Z_\alpha\setminus I(S)|=j-|S|.
\]
Moreover every internal boundary of a zone is separated from every other zone
by at least one index not in $H(S)$.
\end{lemma}

\begin{proof}
The set $I(S)$ is determined by $S$, and $H(S)$ is determined by $I(S)$.
The maximal consecutive intervals of a finite ordered set are uniquely
defined, so the intervals $Z_\alpha$ are determined by $S$, pairwise disjoint,
and their union is $H(S)$.  Since every $i\in I(S)$ satisfies $|i-i|\le1$,
one has $I(S)\subseteq H(S)$, so the zones cover every deleted block index.

The sets $R$ and the intervals $Z_\alpha$ partition $\{1,\ldots,j\}$.  Hence
\[
  j=|R|+\sum_\alpha |Z_\alpha|.
\]
Because $I(S)\subseteq\bigcup_\alpha Z_\alpha$ and the zones are disjoint,
\[
  \sum_\alpha |Z_\alpha|
  =
  |I(S)|+\sum_\alpha |Z_\alpha\setminus I(S)|.
\]
Substituting $|I(S)|=|S|$ gives the displayed count.

Finally, if two maximal intervals of $H(S)$ are distinct and ordered, then the
index immediately after the first interval, if it exists before the second,
does not belong to $H(S)$; otherwise the two intervals would be part of the
same maximal consecutive interval.  This is the stated separation.
\end{proof}

\begin{lemma}[Complement blocks expose canonical halo zone endpoints]\label{lem:bc-canonical-halo-complement-exposes-zone-endpoints}
Let
\[
  P=(\rho_0,\rho_1,\ldots,\rho_{2j})
\]
be a two-step growth path as in \cref{lem:bc-growth-path-model}, let
$S\subset\{1,\ldots,j\}$, and let $Z_\gamma=\{a_\gamma,\ldots,b_\gamma\}$ be
the canonical halo replacement zones of
\cref{lem:bc-canonical-halo-replacement-zones}.  Let
\[
  R=\{1,\ldots,j\}\setminus\bigcup_\gamma Z_\gamma .
\]
Suppose one is given $S$ and all complement blocks
\[
  B_r(P)=(\rho_{2r-2},\rho_{2r-1},\rho_{2r})\qquad(r\in R),
\]
and note that \cref{lem:bc-growth-path-model} gives
$\rho_0=\rho_{2j}=\varnothing$.
Then the endpoint pair
\[
  \rho_{2a_\gamma-2},\qquad \rho_{2b_\gamma}
\]
is determined for every canonical halo zone.
\end{lemma}

\begin{proof}
The set $S$ determines the canonical halo zones and the complement $R$ by
\cref{lem:bc-canonical-halo-replacement-zones}.  Fix one zone
$Z_\gamma=\{a_\gamma,\ldots,b_\gamma\}$.

For the left endpoint, if $a_\gamma=1$, then
$\rho_{2a_\gamma-2}=\rho_0=\varnothing$ is fixed.  If $a_\gamma>1$, then
$a_\gamma-1$ is not in the one-block halo; otherwise $Z_\gamma$ would not be
the maximal halo component beginning at $a_\gamma$.  Hence
$a_\gamma-1\in R$, and the recovered complement block $B_{a_\gamma-1}(P)$ has
right endpoint
\[
  \rho_{2(a_\gamma-1)}=\rho_{2a_\gamma-2}.
\]

For the right endpoint, if $b_\gamma=j$, then
$\rho_{2b_\gamma}=\rho_{2j}=\varnothing$ is fixed.  If $b_\gamma<j$, then
$b_\gamma+1$ is not in the one-block halo, again by maximality of the halo
component.  Thus
$b_\gamma+1\in R$, and the recovered complement block $B_{b_\gamma+1}(P)$ has
left endpoint
\[
  \rho_{2(b_\gamma+1)-2}=\rho_{2b_\gamma}.
\]
These formulas determine both endpoints of every zone.
\end{proof}

\begin{lemma}[Alternating witnesses form one tight halo zone]\label{lem:bc-alternating-witness-full-halo-zone}
Let $q\ge1$, put $j=2q$, and let
\[
  S_{\mathrm{odd}}=\{1,3,\ldots,2q-1\},
  \qquad
  S_{\mathrm{even}}=\{2,4,\ldots,2q\}.
\]
For either choice $S\in\{S_{\mathrm{odd}},S_{\mathrm{even}}\}$, the canonical
halo construction of \cref{lem:bc-canonical-halo-replacement-zones} has a
single replacement zone
\[
  Z=\{1,2,\ldots,2q\}.
\]
Moreover
\[
  |Z\cap I(S)|=|Z\setminus I(S)|=q,
\]
and the components of $Z\cap I(S)$ are the $q$ singleton deleted blocks of
one parity.
\end{lemma}

\begin{proof}
Since $j=2q$, the label-to-block map is $s\mapsto2q-s+1$.  It sends
$S_{\mathrm{odd}}$ to the even block indices
\[
  I(S_{\mathrm{odd}})=\{2,4,\ldots,2q\},
\]
and sends $S_{\mathrm{even}}$ to the odd block indices
\[
  I(S_{\mathrm{even}})=\{1,3,\ldots,2q-1\}.
\]
Thus, in either case, $I(S)$ consists of all block indices of one parity and
has cardinality $q$.

The one-block halo of either parity class is all of $\{1,\ldots,2q\}$: every
index of the opposite parity is adjacent to an element of $I(S)$, with the
endpoints covered by the endpoint parity element itself or by its adjacent
neighbour.  Hence the maximal consecutive intervals of the halo consist of the
single interval $Z=\{1,\ldots,2q\}$.

Inside this zone, $Z\cap I(S)=I(S)$ has cardinality $q$, and
$Z\setminus I(S)$ is the complementary parity class, also of cardinality
$q$.  Since no two elements of one parity class are consecutive, the maximal
components of $Z\cap I(S)$ are exactly the $q$ singleton deleted blocks.
\end{proof}

\begin{lemma}[Canonical halo zones have short retained gaps]\label{lem:bc-canonical-halo-short-gaps}
Let $j\ge1$, let $S\subset\{1,\ldots,j\}$, put $I=I(S)$, and let
$Z_\alpha=\{a_\alpha,\ldots,b_\alpha\}$ be one canonical halo zone of
\cref{lem:bc-canonical-halo-replacement-zones}.  Write the maximal consecutive
components of $Z_\alpha\cap I$ as
\[
  K_1=\{u_1,\ldots,v_1\},\ldots,K_m=\{u_m,\ldots,v_m\},
  \qquad u_1\le v_1<u_2\le v_2<\cdots<u_m\le v_m .
\]
Then
\[
  a_\alpha=\max(1,u_1-1),\qquad
  b_\alpha=\min(j,v_m+1),
\]
and, for $1\le t<m$,
\[
  u_{t+1}-v_t\in\{2,3\}.
\]
Equivalently, inside a canonical halo zone the retained gaps between
successive deleted components have length one or two, and the retained prefix
and suffix have length at most one.
\end{lemma}

\begin{proof}
The full halo $H(S)$ is the union of the one-block halos of the maximal
consecutive components of $I$.  The zone $Z_\alpha$ is one maximal consecutive
component of this union, and the intervals $K_1,\ldots,K_m$ are exactly the
components of $I$ whose one-block halos lie in this component of $H(S)$.

The one-block halo of $K_t=\{u_t,\ldots,v_t\}$ inside $\{1,\ldots,j\}$ is
\[
  \{\max(1,u_t-1),\ldots,\min(j,v_t+1)\}.
\]
Since the halos of $K_1,\ldots,K_m$ form one consecutive component, the left
endpoint of their union is the left endpoint of the first halo and the right
endpoint is the right endpoint of the last halo.  This proves the displayed
formulas for $a_\alpha$ and $b_\alpha$.

For consecutive components of $I$, maximality gives $u_{t+1}\ge v_t+2$.
If instead $u_{t+1}\ge v_t+4$, then the index $v_t+2$ lies strictly between
the two components and has distance at least two from every element of both
$K_t$ and $K_{t+1}$.  It also has distance at least two from every other
component of $I$, because the other components lie farther to the left or
right.  Hence $v_t+2\notin H(S)$, which would split the two halos into
different components of $H(S)$, contradicting that both lie in $Z_\alpha$.
Thus $u_{t+1}\le v_t+3$, and
$u_{t+1}-v_t\in\{2,3\}$.

The final sentence is just a restatement: the retained gap between $K_t$ and
$K_{t+1}$ is $\{v_t+1,\ldots,u_{t+1}-1\}$, whose size is
$u_{t+1}-v_t-1\in\{1,2\}$, and the endpoint formulas show that at most one
retained halo index can occur before $K_1$ or after $K_m$.
\end{proof}

\begin{lemma}[Canonical halo zones have local retained budget]\label{lem:bc-canonical-halo-retained-budget}
Let $j\ge1$, let $S\subset\{1,\ldots,j\}$, and let $Z_\alpha$ be the canonical
halo replacement zones of \cref{lem:bc-canonical-halo-replacement-zones}.  Put
$I=I(S)$.  Then every nonempty zone satisfies
\[
  |Z_\alpha\setminus I|\le 2\,|Z_\alpha\cap I|.
\]
Consequently
\[
  \sum_\alpha |Z_\alpha\setminus I|\le 2|S|.
\]
\end{lemma}

\begin{proof}
Fix a zone $Z_\alpha$.  Since $Z_\alpha$ is a maximal component of the
one-block halo of $I$, it contains at least one element of $I$.  Every retained
index $r\in Z_\alpha\setminus I$ belongs to the halo, so there is an
$i\in Z_\alpha\cap I$ with $|r-i|=1$.  Assign $r$ to one such adjacent deleted
index, choosing the smaller one if two are available.  A fixed deleted index
$i$ can receive only the retained neighbours $i-1$ and $i+1$, so it receives at
most two retained indices.  This gives an injection from
$Z_\alpha\setminus I$ into two labelled slots over $Z_\alpha\cap I$, proving
the first displayed inequality.

The zones are disjoint and cover $I$, so summing the first inequality over
$\alpha$ gives
\[
  \sum_\alpha |Z_\alpha\setminus I|
  \le
  2\sum_\alpha |Z_\alpha\cap I|
  =
  2|I|=2|S|.
\]
\end{proof}

\begin{lemma}[Canonical halo zone balance]\label{lem:bc-canonical-halo-zone-balance}
Let $j\ge1$, let $S\subset\{1,\ldots,j\}$, put $I=I(S)$, and let
$Z_\alpha=\{a_\alpha,\ldots,b_\alpha\}$ be one canonical halo zone of
\cref{lem:bc-canonical-halo-replacement-zones}.  Write the maximal consecutive
components of $Z_\alpha\cap I$ as
\[
  K_t=\{u_t,\ldots,v_t\}\qquad(1\le t\le m),
\]
ordered increasingly.  Put $h_t=|K_t|$, and for $1\le t<m$ put
\[
  g_t=u_{t+1}-v_t-1.
\]
Finally set
\[
  \epsilon_L=
  \begin{cases}
  1,&u_1>1,\\
  0,&u_1=1,
  \end{cases}
  \qquad
  \epsilon_R=
  \begin{cases}
  1,&v_m<j,\\
  0,&v_m=j.
  \end{cases}
\]
Then
\[
  |Z_\alpha\cap I|=\sum_{t=1}^m h_t,
  \qquad
  |Z_\alpha\setminus I|
  =
  \epsilon_L+\epsilon_R+\sum_{t=1}^{m-1}g_t .
\]
Equivalently, the local retained-minus-deleted balance is
\[
  |Z_\alpha\setminus I|-|Z_\alpha\cap I|
  =
  \epsilon_L+\epsilon_R+\sum_{t=1}^{m-1}g_t-\sum_{t=1}^m h_t .
\]
\end{lemma}

\begin{proof}
The first identity holds because the intervals $K_t$ are precisely the
connected components of $Z_\alpha\cap I$.

By \cref{lem:bc-canonical-halo-short-gaps}, the left endpoint of the zone is
$a_\alpha=\max(1,u_1-1)$.  Hence there is one retained prefix index before
$K_1$ exactly when $u_1>1$, and none when $u_1=1$; this contributes
$\epsilon_L$.  Similarly, the right endpoint
$b_\alpha=\min(j,v_m+1)$ gives one retained suffix index exactly when
$v_m<j$, contributing $\epsilon_R$.

Between $K_t$ and $K_{t+1}$ the retained indices are exactly
\[
  \{v_t+1,\ldots,u_{t+1}-1\},
\]
whose cardinality is $g_t$.  The retained prefix, the internal retained gaps,
and the retained suffix are disjoint and exhaust $Z_\alpha\setminus I$.
Summing their sizes gives the second identity, and subtracting the first
identity gives the balance formula.
\end{proof}

\begin{lemma}[Height witness labels are separated]\label{lem:bc-height-witness-labels-separated}
Let $T$ be a semistandard tableau of content $(2^j)$ with
$\ell(\lambda(T))\ge2q$, and put
\[
  t_r=T(2r-1,1)\qquad(1\le r\le q).
\]
Then
\[
  t_r+1<t_{r+1}\qquad(1\le r<q).
\]
Consequently the height witness set $H_q(T)=\{t_1,\ldots,t_q\}$ contains no
two consecutive labels.
\end{lemma}

\begin{proof}
The first column of a semistandard tableau is strictly increasing.  For
$1\le r<q$ the entries in rows $2r-1$, $2r$, and $2r+1$ therefore satisfy
\[
  t_r=T(2r-1,1)<T(2r,1)<T(2r+1,1)=t_{r+1}.
\]
All entries are integers, so $t_r+1<t_{r+1}$.  This is exactly the assertion
that no two labels in $H_q(T)$ are consecutive.
\end{proof}

\begin{lemma}[Height witness blocks have retained collars]\label{lem:bc-height-witness-retained-collars}
Let $T$ be a semistandard tableau of content $(2^j)$ with
$\ell(\lambda(T))\ge2q$, put
\[
  t_r=T(2r-1,1)\qquad(1\le r\le q),
  \qquad S=H_q(T),
\]
and fix $s=t_r\in S$.  Let
\[
  u=j-s+1
\]
be the corresponding block index in the growth path of
\cref{lem:bc-growth-path-model}.  Then $u>1$ and $u-1\notin I(S)$.
Moreover, if $u<j$, then $u+1\notin I(S)$.  Equivalently, the successor
label $s+1$ is retained, and the predecessor label $s-1$ is retained whenever
it exists.
\end{lemma}

\begin{proof}
Since $\ell(\lambda(T))\ge2q$, the box $(2r,1)$ belongs to $T$.  Column
strictness gives
\[
  s=T(2r-1,1)<T(2r,1)\le j,
\]
so $s<j$.  Also $s+1\notin S$: for $k\le r$ one has $t_k\le s$, while for
$k>r$ the separation lemma gives $s+1<t_{r+1}\le t_k$.  Thus no $t_k$ equals
$s+1$.  The inequality $s<j$ gives $u=j-s+1>1$, and $s+1\notin S$ is
equivalent, under $a\mapsto j-a+1$, to $u-1\notin I(S)$.

It remains to prove the predecessor assertion.  Assume $s>1$.  If
$s-1\in S$, then $s-1=t_k$ for some $k$.  Since $t_1<\cdots<t_q$ and
$s-1<s=t_r$, one has $k<r$.  If $k<r$, then
\cref{lem:bc-height-witness-labels-separated}, applied along the increasing
list from $k$ to $r$, gives in particular
\[
  t_k+1<t_{k+1}\le t_r=s,
\]
so $t_k\le s-2$, contradicting $t_k=s-1$.  Hence $s-1\notin S$.
The condition $u<j$ is equivalent to $s>1$, and $s-1\notin S$ is equivalent
to $u+1\notin I(S)$.
\end{proof}

\begin{lemma}[Height successor pairs have forced frontier form]\label{lem:bc-height-successor-pair-frontier-form}
Let $T$ be a semistandard tableau of content $(2^j)$ with
$\ell(\lambda(T))\ge2q$, and let
\[
  \varnothing=\lambda^{(0)}\subset\lambda^{(1)}\subset\cdots
  \subset\lambda^{(j)}=\lambda(T)
\]
be its Pieri two-strip path.  Put
\[
  t_r=T(2r-1,1)\qquad(1\le r\le q),
  \qquad S=H_q(T).
\]
For a fixed $r$, set $s=t_r$.  Then $s<j$, $s+1\notin S$, and
\[
  \ell(\lambda^{(s-1)})=2r-2,\qquad
  \ell(\lambda^{(s)})=2r-1 .
\]
Moreover
\[
  \lambda^{(s)}\setminus\lambda^{(s-1)}
  =
  \{(2r-1,1),x_s\}
\]
for a box $x_s$ not in the first column, while
\[
  \ell(\lambda^{(s+1)})\in\{2r-1,2r\}.
\]
If $\ell(\lambda^{(s+1)})=2r$, then the retained successor strip
$\lambda^{(s+1)}\setminus\lambda^{(s)}$ contains the new first-column box
$(2r,1)$; otherwise it contains no first-column box below row $2r-1$.
\end{lemma}

\begin{proof}
Since $\ell(\lambda(T))\ge2q$, the first-column box $(2r,1)$ belongs to
$T$.  Column strictness gives
\[
  s=T(2r-1,1)<T(2r,1)\le j,
\]
so $s<j$.  The labels $t_1<\cdots<t_q$ are strictly increasing, and
\cref{lem:bc-height-witness-labels-separated} gives
$t_r+1<t_{r+1}$ for $r<q$.  Hence $s+1$ is not one of the labels
$t_1,\ldots,t_q$, proving $s+1\notin S$.

The box $(2r-1,1)$ carries label $s$, so it is absent from
$\lambda^{(s-1)}$ and present in $\lambda^{(s)}$.  Thus
\[
  \ell(\lambda^{(s-1)})\le2r-2,\qquad
  \ell(\lambda^{(s)})\ge2r-1 .
\]
For $r=1$ the first inequality forces
$\ell(\lambda^{(s-1)})=0$.  For $r>1$, the box $(2r-2,1)$ lies above
$(2r-1,1)$ in the first column, so its label is smaller than $s$ and it is
already present in $\lambda^{(s-1)}$.  Hence
$\ell(\lambda^{(s-1)})=2r-2$ in every case.

The strip $\lambda^{(s)}\setminus\lambda^{(s-1)}$ is horizontal of size two by
\cref{lem:bc-pieri-path-bijection}.  It contains $(2r-1,1)$ and therefore
cannot also contain $(2r,1)$, because those two boxes are in the same column.
Thus $\ell(\lambda^{(s)})=2r-1$.  The second box $x_s$ carrying label $s$ is
therefore not in the first column.

The next strip $\lambda^{(s+1)}\setminus\lambda^{(s)}$ is again horizontal of
size two, so it can create at most one new row.  Since
$\ell(\lambda^{(s)})=2r-1$, this gives
$\ell(\lambda^{(s+1)})\in\{2r-1,2r\}$.  If the length becomes $2r$, the new
row is created precisely by the first-column box $(2r,1)$, which must belong
to the successor strip.  If the length stays $2r-1$, no first-column box below
row $2r-1$ is added at that step.
\end{proof}

\begin{corollary}[Height anchor pairs reduce to off-frontier data]\label{cor:bc-height-anchor-pair-off-frontier-reduction}
Work in the height setting of
\cref{cor:bc-height-surplus-canonical-right-anchors}.  Let
$u\in Z_\gamma\cap I(S)$ be a deleted singleton block, and put
\[
  s=j-u+1 .
\]
Let $r$ be determined by $s=t_r$ in the ordered height witness set
$S=H_q(T)$.  Then the retained left-neighbour block $u-1$ corresponds to the
successor label $s+1$, and the two-label Pieri segment
\[
  \lambda^{(s-1)}\subset\lambda^{(s)}\subset\lambda^{(s+1)}
\]
has the forced frontier form of
\cref{lem:bc-height-successor-pair-frontier-form}.  Consequently, a
target-window readout which determines $\lambda^{(s-1)}$, the off-first-column
box $x_s$, and the retained successor strip
\[
  \lambda^{(s+1)}\setminus\lambda^{(s)}
\]
determines this two-label Pieri segment.  Equivalently, the height
anchor-pair inverse need not encode the odd first-column witness box; the
remaining Pieri-side input is off-frontier data together with the retained
successor strip.
\end{corollary}

\begin{proof}
The label-to-block convention in \cref{lem:bc-growth-path-model} sends a label
$a$ to the block $j-a+1$.  Thus the deleted singleton block $u$ corresponds to
the deleted height witness label $s=j-u+1$, while the block $u-1$ corresponds
to $s+1$.  Since $u-1\in Z_\gamma\setminus I(S)$ by
\cref{cor:bc-height-anchor-reduced-left-pairs}, this is the retained
left-neighbour block.

Applying \cref{lem:bc-height-successor-pair-frontier-form} to the witness
label $s=t_r$ gives
\[
  \lambda^{(s)}=
  \lambda^{(s-1)}\cup\{(2r-1,1),x_s\}
\]
with $x_s$ off the first column, and gives the stated restriction on the
successor strip.  Therefore the data
$\lambda^{(s-1)}$, $x_s$, and
$\lambda^{(s+1)}\setminus\lambda^{(s)}$ reconstruct first
$\lambda^{(s)}$ and then $\lambda^{(s+1)}$.  The final sentence is just this
reconstruction read as a restriction on what the merged height anchor-pair
inverse has to supply.
\end{proof}

\begin{corollary}[Height anchor pairs have two-sided retained collar data]\label{cor:bc-height-anchor-pair-retained-collars}
Work in the height setting of
\cref{cor:bc-height-surplus-canonical-right-anchors}.  Let
$u_t\in Z_\gamma\cap I(S)$ be a deleted singleton block, and put
$s_t=j-u_t+1$.  Then the left-neighbour block $u_t-1$ exists and is retained.
If $u_t<j$, then the right-neighbour block $u_t+1$ also exists and is
retained.  Moreover, whenever these neighbours exist, they lie in the same
canonical halo zone $Z_\gamma$.

Consequently, a positive-height merged-window construction for the pair
$(u_t-1,u_t)$ has deterministic access, from the canonical halo parser, to
the retained left collar and, except when $u_t=j$, to the retained right
collar.  The exceptional case $u_t=j$ is precisely the terminal right boundary
side, where the zone endpoint $\rho_{2j}$ is part of the endpoint data.
\end{corollary}

\begin{proof}
Apply \cref{lem:bc-height-witness-retained-collars} to the height witness
label $s_t$.  It gives $u_t>1$ and $u_t-1\notin I(S)$, so the left-neighbour
block exists and is retained.  If $u_t<j$, the same lemma gives
$u_t+1\notin I(S)$, so the right-neighbour block exists and is retained.

The canonical halo zone $Z_\gamma$ is a connected component of the one-block
halo of $I(S)$.  Hence every existing index at distance one from
$u_t\in Z_\gamma\cap I(S)$ lies in that same component.  Thus the retained
neighbours just identified lie in $Z_\gamma$.  If $u_t=j$, the right
neighbour does not exist and the right side of the singleton block is the
global endpoint $\rho_{2j}$, which is included as the right endpoint of any
terminal zone containing $u_t$.
\end{proof}

\begin{corollary}[Step--3 pair boundary is the anchor-pair segment]\label{cor:bc-height-anchor-pair-step3-boundary-segment}
Work in the setting of
\cref{cor:bc-height-anchor-pair-off-frontier-reduction}.  Let
$P=(\rho_0,\rho_1,\ldots,\rho_{2j})$ be the Step--3 boundary growth path of
\cref{lem:bc-growth-path-model}.  Then the retained-successor/deleted-witness
pair of labels $(s+1,s)$ is represented on the Step--3 top/right boundary by
the consecutive two-block segment
\[
  \rho_{2u-4},\rho_{2u-3},\rho_{2u-2},
  \rho_{2u-1},\rho_{2u}.
\]
Consequently, any target-window readout which determines this Step--3 boundary
segment determines the pair segment required by the local inverse hypothesis of
\cref{prop:bc-height-anchor-pair-zone-recovery}.
\end{corollary}

\begin{proof}
By \cref{lem:bc-growth-path-model}, tableau label $\ell$ corresponds to block
$j-\ell+1$ of the Step--3 boundary growth path.  The deleted witness label is
$s=j-u+1$, hence it corresponds to block $u$.  The retained successor label
$s+1$ corresponds to block
\[
  j-(s+1)+1=u-1 .
\]
Thus the adjacent label pair $(s+1,s)$ is exactly the adjacent block pair
$(u-1,u)$.  The Step--3 boundary segment covered by these two blocks is
\[
  \rho_{2(u-1)-2},\rho_{2(u-1)-1},\rho_{2(u-1)},
  \rho_{2u-1},\rho_{2u},
\]
which is the displayed segment.  This is precisely the segment appearing in
the pair-inverse hypothesis of
\cref{prop:bc-height-anchor-pair-zone-recovery}.
\end{proof}

\begin{proposition}[Step--3 pair-boundary carriers supply anchor-pair inverses]\label{prop:bc-step3-pair-boundary-carriers-supply-anchor-pair}
Work in the height setting of
\cref{prop:bc-height-anchor-pair-zone-recovery}.  Suppose that the parser
there decomposes $V^\ast_\gamma$ into the unchanged anchor blocks and one
merged window $M_t$ for each deleted singleton $u_t$, and suppose that for
each $t$ the data
\[
  M_t,\quad \beta_t,\quad S,\quad Z_\gamma,\quad t,\quad
  \rho_{2a_\gamma-2},\rho_{2b_\gamma}
\]
determine the Step--3 boundary segment
\[
  \rho_{2u_t-4},\rho_{2u_t-3},\rho_{2u_t-2},
  \rho_{2u_t-1},\rho_{2u_t}.
\]
Then the pair-inverse hypothesis of
\cref{prop:bc-height-anchor-pair-zone-recovery} holds for every $t$, and hence
the original height halo-zone subpath is recovered from
$V^\ast_\gamma$, $\beta_\gamma$, the zone endpoint pair, and $Z_\gamma$.
\end{proposition}

\begin{proof}
For each deleted singleton $u_t$, put $s_t=j-u_t+1$.  The pair
$(u_t-1,u_t)$ is the retained-successor/deleted-witness pair for labels
$(s_t+1,s_t)$ by
\cref{cor:bc-height-anchor-pair-off-frontier-reduction}.  Therefore
\cref{cor:bc-height-anchor-pair-step3-boundary-segment} identifies the
displayed boundary segment with the consecutive original two-block segment
required in the local pair-inverse hypothesis.  The assumed deterministic
readout from $M_t$, $\beta_t$, the zone data, and the endpoint pair is
therefore a valid fixed pair inverse for this $t$.  Applying this to every
$t$ supplies all pair inverses required by
\cref{prop:bc-height-anchor-pair-zone-recovery}, and that proposition recovers
the full original zone subpath.
\end{proof}

\begin{corollary}[Two-label Step--3 carrier data supply anchor-pair inverses]\label{cor:bc-two-label-step3-carriers-supply-anchor-pair}
Work in the setting of
\cref{prop:bc-step3-pair-boundary-carriers-supply-anchor-pair}.  For the
deleted singleton $u_t$, put $s_t=j-u_t+1$, and let
$E_{s_t,s_t+1}$ be the two-label Step--3 subarrangement of
\cref{lem:bc-step3-consecutive-label-interval-boundary}.  Suppose that, for
each $t$, the data
\[
  M_t,\quad \beta_t,\quad S,\quad Z_\gamma,\quad t,\quad
  \rho_{2a_\gamma-2},\rho_{2b_\gamma}
\]
determine the cell entries of $E_{s_t,s_t+1}$ and the partition labels on its
left/bottom boundary.  Then the hypotheses of
\cref{prop:bc-step3-pair-boundary-carriers-supply-anchor-pair} hold, and hence
the original height halo-zone subpath is recovered from
$V^\ast_\gamma$, $\beta_\gamma$, the zone endpoint pair, and $Z_\gamma$.
\end{corollary}

\begin{proof}
Apply
\cref{cor:bc-step3-consecutive-carrier-data-recovers-boundary} with
$a=s_t$ and $b=s_t+1$.  Since $u_t=j-s_t+1$, the boundary segment recovered
there is
\[
  \rho_{2(j-(s_t+1))},\ldots,\rho_{2(j-s_t+1)}
  =
  \rho_{2u_t-4},\ldots,\rho_{2u_t}.
\]
Thus each merged-window datum determines the Step--3 pair-boundary segment
required in
\cref{prop:bc-step3-pair-boundary-carriers-supply-anchor-pair}.  That
proposition then gives the anchor-pair inverses and recovers the zone subpath.
\end{proof}

\begin{lemma}[Two-label Step--3 carrier entries are three-valued]\label{lem:bc-two-label-step3-cell-entry-three-valued}
Use the Step--3 dual-RSK' half-square in the content $(2^j)$ specialization of
\cref{lem:bc-growth-path-model}.  For $1\le s<j$, the two-label
subarrangement $E_{s,s+1}$ has a single cell, and the entry of that cell
belongs to $\{0,1,2\}$.
\end{lemma}

\begin{proof}
The Step--3 half-square has one off-diagonal cell for each unordered pair of
distinct labels in the chosen half of the symmetric Step--2 matrix.  Therefore
the subarrangement on the two labels $s+1,s$ has exactly one cell.

Let $N=(N_{ab})$ be the Step--2 matrix of
\cref{lem:bc-step2-matrix-doubled-content-degrees}.  By
\cref{lem:bc-step3-cells-are-step2-half-matrix}, the unique Step--3 cell entry
of $E_{s,s+1}$ is either $N_{s,s+1}$ or $N_{s+1,s}$, which are equal by
symmetry.  The entries of $N$ are nonnegative integers and every row sum is
$2$ by \cref{lem:bc-step2-matrix-doubled-content-degrees}; hence
$0\le N_{s,s+1}\le2$.  Thus the unique cell entry lies in $\{0,1,2\}$.
\end{proof}

\begin{lemma}[Two-label Step--3 carriers have diagonal jump two]\label{lem:bc-two-label-step3-diagonal-jump-two}
Use the Step--3 dual-RSK' half-square in the content $(2^j)$ specialization of
\cref{lem:bc-growth-path-model}.  For $1\le s<j$, let
\[
  \begin{matrix}
  \nu & \lambda\\
  \rho & \mu
  \end{matrix}
\]
be the four corner labels of the one-cell carrier $E_{s,s+1}$, with
$\rho$ the southwest corner and $\lambda$ the northeast corner in
\cref{lem:bc-dual-rsk-prime-growth-local-reversibility}.  Then
\[
  |\lambda|-|\rho|=2 .
\]
\end{lemma}

\begin{proof}
In Krattenthaler's Step--3 half-square the labels are ordered
$j,j-1,\ldots,1$, and the cell for the unordered adjacent pair
$\{s,s+1\}$ lies next to the deleted diagonal.  Consider the southwest hook of
this cell in the retained strict half-square: the cell itself, the cells in
the same retained row between it and the left boundary, and the cells in the
same retained column between it and the bottom boundary.  Because $s$ and
$s+1$ are adjacent in the label order, this hook is exactly the set of
off-diagonal representatives incident to one of the two labels $s,s+1$; which
one occurs depends only on which of the two symmetric halves was retained in
Step~3.

Sum the one-cell conservation identity of
\cref{lem:bc-dual-rsk-prime-cell-size-conservation} over this hook.  All
internal edge labels cancel.  The two exposed hook ends lie on the fixed empty
left and bottom boundaries of the Step--3 construction, while the remaining
outer diagonal corners are precisely $\rho$ and $\lambda$.  Hence
\[
  |\lambda|-|\rho|
\]
is the sum of the Step--3 entries in the hook.

By \cref{lem:bc-step3-cells-are-step2-half-matrix}, those entries are exactly
one off-diagonal half of the corresponding Step--2 symmetric matrix entries
incident to the selected label.  The diagonal Step--2 entry is zero, and the
row and column sums of the Step--2 matrix are both $2$ by
\cref{lem:bc-step2-matrix-doubled-content-degrees}.  Therefore the hook sum is
$2$, giving the displayed diagonal jump.
\end{proof}

\begin{corollary}[Raw one-cell top/right corners are not target blocks]\label{cor:bc-one-cell-top-right-corners-not-target-blocks}
Use the Step--3 dual-RSK' half-square in the content $(2^j)$ specialization
of \cref{lem:bc-growth-path-model}.  Let
\[
  \begin{matrix}
  \nu & \lambda\\
  \rho & \mu
  \end{matrix}
\]
be the four corner labels of a two-label one-cell carrier, and let $e$ be its
cell entry.  If the three non-southwest corners are read as a two-step growth
block with middle label $\lambda$, then
\[
  |\nu|-2|\lambda|+|\mu|=-2-e .
\]
Consequently this triple is not a content-$2$ target block.  Thus the
one-cell corner readout in
\cref{prop:bc-one-cell-step3-corner-readouts-supply-grouped-carrier-decoders}
cannot be supplied by placing the raw top/right carrier corners as a single
target block; it needs a content-normalized or genuinely collective grouped
encoding.
\end{corollary}

\begin{proof}
By \cref{lem:bc-two-label-step3-diagonal-jump-two},
\[
  |\lambda|-|\rho|=2 .
\]
The one-cell conservation identity of
\cref{lem:bc-dual-rsk-prime-cell-size-conservation} gives
\[
  e=|\lambda|-|\mu|-|\nu|+|\rho|,
\]
so
\[
  |\nu|+|\mu|=|\lambda|+|\rho|-e=2|\lambda|-2-e .
\]
Subtracting $2|\lambda|$ gives the displayed identity.  Finally
\cref{lem:bc-two-label-step3-cell-entry-three-valued} gives
$e\in\{0,1,2\}$, so the displayed value is negative and is not $2$.
\end{proof}

\begin{proposition}[Two-bit cell-entry carriers reduce anchor pairs to boundary recovery]\label{prop:bc-two-bit-cell-entry-carriers-supply-anchor-pair}
Work in the setting of
\cref{prop:bc-step3-pair-boundary-carriers-supply-anchor-pair}.  Fix once and
for all an injection
\[
  \kappa:\{0,1,2\}\hookrightarrow\{0,1\}^2 .
\]
For the deleted singleton $u_t$, put $s_t=j-u_t+1$ and let $e_t$ be the
unique Step--3 cell entry of $E_{s_t,s_t+1}$.  Suppose that, for each $t$, the
data
\[
  M_t,\quad S,\quad Z_\gamma,\quad t,\quad
  \rho_{2a_\gamma-2},\rho_{2b_\gamma}
\]
determine the partition labels on the left/bottom boundary of
$E_{s_t,s_t+1}$, and suppose the two-bit branch word satisfies
$\beta_t=\kappa(e_t)$.  Then the hypotheses of
\cref{cor:bc-two-label-step3-carriers-supply-anchor-pair} hold.  Consequently
the original height halo-zone subpath is recovered from
$V^\ast_\gamma$, $\beta_\gamma$, the zone endpoint pair, and $Z_\gamma$.
\end{proposition}

\begin{proof}
By \cref{lem:bc-two-label-step3-cell-entry-three-valued}, $e_t\in\{0,1,2\}$,
so the fixed injection $\kappa$ lets the inverse read $e_t$ from the two-bit
word $\beta_t$.  Since $E_{s_t,s_t+1}$ has a single cell, this determines all
cell entries of that two-label subarrangement.  The left/bottom boundary
labels are determined by hypothesis from $M_t$, the zone data, the endpoint
pair, and $t$.  Thus the data
\[
  M_t,\quad \beta_t,\quad S,\quad Z_\gamma,\quad t,\quad
  \rho_{2a_\gamma-2},\rho_{2b_\gamma}
\]
determine both the cell entries and the left/bottom boundary labels of
$E_{s_t,s_t+1}$ for every $t$.
\Cref{cor:bc-two-label-step3-carriers-supply-anchor-pair} applies and gives
the stated recovery of the original height halo-zone subpath.
\end{proof}

\begin{proposition}[Two-label left/bottom carriers recover positive height zones]\label{prop:bc-two-label-left-bottom-carrier-zone-recovery}
Work in the height setting of
\cref{prop:bc-height-anchor-pair-zone-recovery}.  Suppose that the parser
there decomposes $V^\ast_\gamma$ into the unchanged anchor blocks and one
merged window $M_t$ for each deleted singleton $u_t$, listed in increasing
$t$.  Fix an injection
\[
  \kappa:\{0,1,2\}\hookrightarrow\{0,1\}^2 .
\]
For each $t$, put $s_t=j-u_t+1$, let $E_{s_t,s_t+1}$ be the corresponding
two-label Step--3 carrier, and let $e_t$ be its unique cell entry.  Assume
that the data
\[
  M_t,\quad S,\quad Z_\gamma,\quad t,\quad
  \rho_{2a_\gamma-2},\rho_{2b_\gamma}
\]
determine the ordered partition labels on the left/bottom boundary of
$E_{s_t,s_t+1}$, and that the deterministic two-bit split of
$\beta_\gamma$ satisfies $\beta_t=\kappa(e_t)$ for every $t$.  Then the
local inverse hypothesis of
\cref{prop:bc-height-anchor-pair-zone-recovery} holds for every deleted
singleton in $Z_\gamma$, and the original height halo-zone subpath is
recovered from $V^\ast_\gamma$, $\beta_\gamma$, the zone endpoint pair, and
$Z_\gamma$.
\end{proposition}

\begin{proof}
The hypotheses for a fixed $t$ are exactly the hypotheses of
\cref{prop:bc-two-bit-cell-entry-carriers-supply-anchor-pair}: the merged
window and endpoint/parser data determine the left/bottom boundary of
$E_{s_t,s_t+1}$, while the two-bit word records the unique cell entry through
the fixed injection $\kappa$.  Therefore
\cref{prop:bc-two-bit-cell-entry-carriers-supply-anchor-pair} supplies the
anchor-pair inverse for this $t$.  Applying this argument to every deleted
singleton supplies all local pair inverses required in
\cref{prop:bc-height-anchor-pair-zone-recovery}.  That proposition then
recovers the full original height halo-zone subpath.
\end{proof}

\begin{corollary}[Raw carrier boundaries are not automatically target windows]\label{cor:bc-raw-carrier-boundary-not-target-window}
The construction required in
\cref{prop:bc-two-label-left-bottom-carrier-zone-recovery} cannot be supplied,
in the ambient one-cell dual-RSK' category, by simply declaring the
left/bottom boundary of the two-label carrier $E_{s,s+1}$ to be a
content-$2$ target two-step window.  A target merged window must use the
doubled-content restrictions, neighbouring parser data, or an explicit
content-normalized encoding of that carrier boundary.
\end{corollary}

\begin{proof}
Consider a single cell in the fourth arbitrary-filling dual-RSK' local rule
of \cref{lem:bc-dual-rsk-prime-growth-local-reversibility}.  Put
\[
  m=0,\qquad \rho=\varnothing,\qquad \mu=(1),\qquad \nu=\varnothing .
\]
The hypotheses of the forward local rule hold: $\rho\subseteq\mu$ and
$\rho\subseteq\nu$, and the two differences are vertical strips.  Applying
the forward rule gives $\lambda=(1)$, so this is a valid one-cell
dual-RSK' local datum.

The natural left/bottom valley of this cell is
\[
  \nu\supseteq\rho\subseteq\mu
  \qquad\text{or, equivalently,}\qquad
  \mu\supseteq\rho\subseteq\nu .
\]
In either order its two-step content weight is
\[
  |\nu|-2|\rho|+|\mu|=1 .
\]
But every target two-step block in the content $(2^j)$ growth-path model has
weight $2$ by \cref{lem:bc-growth-path-model}.  Thus the raw left/bottom
boundary of a dual-RSK' carrier cell is not, by the local rules alone, a
valid target two-step window.  Any valid merged-window construction must
therefore use additional structure beyond this raw boundary identification.
\end{proof}

\begin{definition}[Bottom padding]\label{def:bc-bottom-padding}
For a partition $\alpha=(\alpha_1,\ldots,\alpha_\ell)$ and
$e\in\{0,1,2\}$, define
\[
  \operatorname{pad}_e(\alpha)
  =
  (\alpha_1,\ldots,\alpha_\ell,\underbrace{1,\ldots,1}_{e\text{ times}}),
\]
with the convention that $\operatorname{pad}_0(\alpha)=\alpha$.
For a valley of partitions $\nu\supseteq\rho\subseteq\mu$, put
\[
  \operatorname{Pad}_e(\nu,\rho,\mu)
  =
  \bigl(\operatorname{pad}_e(\nu),\rho,\mu\bigr).
\]
\end{definition}

\begin{lemma}[Bottom padding normalizes a deficient carrier valley]\label{lem:bc-bottom-padding-normalizes-carrier-valley}
Let
\[
  \nu\supseteq\rho\subseteq\mu
\]
be partitions such that $\nu/\rho$ and $\mu/\rho$ are vertical strips, and let
$e\in\{0,1,2\}$.  If
\[
  |\nu|-2|\rho|+|\mu|=2-e,
\]
then
 $\operatorname{Pad}_e(\nu,\rho,\mu)$ is a valid content-$2$ target two-step
block.  Moreover
$\operatorname{Pad}_e(\nu,\rho,\mu)$ together with $e$ determines
$(\nu,\rho,\mu)$.
\end{lemma}

\begin{proof}
Because $\nu/\rho$ is a vertical strip, for every row index $i$ the difference
$\nu_i-\rho_i$ is either $0$ or $1$.  The partition
$\operatorname{pad}_e(\nu)$ is obtained by appending $e$ new bottom rows of
length $1$.  Against the omitted zero rows of $\rho$, each appended row
therefore contributes one new box, and the old rows still contribute the
vertical strip $\nu/\rho$.  Hence
$\operatorname{pad}_e(\nu)/\rho$ is a vertical strip.  The strip
$\mu/\rho$ is vertical by hypothesis.

The content weight of the padded block is
\[
  |\operatorname{pad}_e(\nu)|-2|\rho|+|\mu|
  =
  |\nu|+e-2|\rho|+|\mu|
  =
  2.
\]
Thus the displayed triple is a content-$2$ target two-step block in the sense
of \cref{lem:bc-growth-path-model}.  Finally, since $e$ is known and the
padding consists of the last $e$ rows of
$\operatorname{pad}_e(\nu)$, deleting those rows recovers $\nu$; the middle
and right entries of the triple are already $\rho$ and $\mu$.
\end{proof}

\begin{proposition}[Padded carrier valleys recover positive height zones]\label{prop:bc-padded-carrier-valleys-recover-positive-height-zones}
Work in the height setting of
\cref{prop:bc-height-anchor-pair-zone-recovery}.  Suppose that the parser
there decomposes $V^\ast_\gamma$ into the unchanged anchor blocks and one
merged window $M_t$ for each deleted singleton $u_t$, listed in increasing
$t$.  Fix an injection
\[
  \kappa:\{0,1,2\}\hookrightarrow\{0,1\}^2 .
\]
For each $t$, put $s_t=j-u_t+1$.  Let
\[
  \begin{matrix}
  \nu_t & \lambda_t\\
  \rho_t & \mu_t
  \end{matrix}
\]
be the corner labels of the one-cell Step--3 carrier $E_{s_t,s_t+1}$, and
let $e_t$ be that cell entry.  Assume that
the merged window is the padded carrier valley
\[
  M_t=\operatorname{Pad}_{e_t}(\nu_t,\rho_t,\mu_t),
\]
and that the deterministic two-bit split of $\beta_\gamma$ satisfies
$\beta_t=\kappa(e_t)$ for every $t$.  Then the original height halo-zone
subpath is recovered from $V^\ast_\gamma$, $\beta_\gamma$, the zone endpoint
pair, and $Z_\gamma$.
\end{proposition}

\begin{proof}
For each $t$, the word $\beta_t$ determines $e_t$ through the fixed injection
$\kappa$.  By \cref{lem:bc-two-label-step3-diagonal-jump-two}, the carrier
satisfies $|\lambda_t|-|\rho_t|=2$.  Therefore
\cref{cor:bc-diagonal-two-cell-has-padded-valley} gives
\[
  |\nu_t|-2|\rho_t|+|\mu_t|=2-e_t.
\]
Therefore \cref{lem:bc-bottom-padding-normalizes-carrier-valley} shows that
the target block
$M_t$ and this recovered value of $e_t$ determine the left/bottom valley
$(\nu_t,\rho_t,\mu_t)$ of $E_{s_t,s_t+1}$.  Since the carrier has one cell by
\cref{lem:bc-two-label-step3-cell-entry-three-valued}, this is exactly the
left/bottom boundary and all cell entries required in
\cref{cor:bc-two-label-step3-carriers-supply-anchor-pair}.  Applying that
corollary for every deleted singleton supplies the local pair inverses in
\cref{prop:bc-height-anchor-pair-zone-recovery}, and that proposition recovers
the whole original zone subpath.
\end{proof}

\begin{lemma}[Endpoint-compatible blocks form a replacement window]\label{lem:bc-endpoint-compatible-blocks-form-replacement-window}
Let $S\subset\{1,\ldots,j\}$, put $I=I(S)$, let
$Z_\gamma=\{a_\gamma,\ldots,b_\gamma\}$ be a replacement zone, and let
\[
  h_1<h_2<\cdots<h_N
\]
be the retained slot indices in $Z_\gamma\setminus I$.  For each $r$ let
\[
  C_r=(L_r,X_r,R_r)
\]
be a two-step block with vertical-strip adjacent differences and content
weight $2$.  Suppose that
\[
  L_{h_1}=\rho_{2a_\gamma-2},\qquad
  R_{h_N}=\rho_{2b_\gamma},
\]
and
\[
  R_{h_i}=L_{h_{i+1}}\qquad(1\le i<N).
\]
Then the concatenation
\[
  L_{h_1},X_{h_1},R_{h_1}=L_{h_2},X_{h_2},\ldots,
  R_{h_{N-1}}=L_{h_N},X_{h_N},R_{h_N}
\]
is a two-step replacement window from $\rho_{2a_\gamma-2}$ to
$\rho_{2b_\gamma}$ with exactly $|Z_\gamma\setminus I|$ two-step blocks.
The deterministic parser which cuts this concatenation after each block
recovers the ordered blocks $C_{h_1},\ldots,C_{h_N}$.
\end{lemma}

\begin{proof}
The displayed endpoint equalities are exactly the compatibility conditions
needed to concatenate the triples $C_{h_i}$ into a single partition sequence.
Each triple contributes one two-step block, and its vertical-strip and
weight-$2$ conditions are part of the hypotheses.  Since
$N=|Z_\gamma\setminus I|$, the concatenated sequence has the required
target size.  Its first and last endpoints are the prescribed zone endpoints.
The parser is fixed by the known ordered slot list $h_1<\cdots<h_N$, so
cutting the sequence into consecutive triples recovers each displayed block.
\end{proof}

\begin{proposition}[Endpoint-compatible padded carriers recover positive height zones]\label{prop:bc-endpoint-compatible-padded-carriers-recover-positive-height-zones}
Work in the height setting of
\cref{prop:bc-height-anchor-pair-zone-recovery}.  Let
\[
  I=I(S),
\]
fix an injection
\[
  \kappa:\{0,1,2\}\hookrightarrow\{0,1\}^2 ,
\]
and let
\[
  Z_\gamma\cap I=\{u_1<\cdots<u_m\},
\]
and let $A_\gamma$ be the canonical right-anchor set.  For each retained slot
$r\in Z_\gamma\setminus I$, define a candidate target block $C_r$ as follows.
If $r\in A_\gamma$, let
\[
  C_r=B_r(P).
\]
If $r=u_t-1$ for some $t$, put $s_t=j-u_t+1$, let
\[
  \begin{matrix}
  \nu_t & \lambda_t\\
  \rho_t & \mu_t
  \end{matrix}
\]
be the corner labels of $E_{s_t,s_t+1}$, let $e_t$ be its cell entry, and set
\[
  C_r=\operatorname{Pad}_{e_t}(\nu_t,\rho_t,\mu_t).
\]
Then each $C_r$ is a valid content-$2$ two-step block.  Assume also that these
ordered blocks satisfy the endpoint equalities in
\cref{lem:bc-endpoint-compatible-blocks-form-replacement-window}.  If the
two-bit split of the branch word satisfies $\beta_t=\kappa(e_t)$ for every
$t$, then the original height halo-zone subpath is recovered from the
resulting replacement window, $\beta_\gamma$, the zone endpoint pair, and
$Z_\gamma$.
\end{proposition}

\begin{proof}
For $r\in A_\gamma$, the block $C_r=B_r(P)$ is an original block of a path in
\cref{lem:bc-growth-path-model}, hence has vertical-strip adjacent differences
and content weight $2$.  For $r=u_t-1$, the diagonal-jump lemma
\cref{lem:bc-two-label-step3-diagonal-jump-two} and
\cref{cor:bc-diagonal-two-cell-has-padded-valley} give
\[
  |\nu_t|-2|\rho_t|+|\mu_t|=2-e_t.
\]
Thus \cref{lem:bc-bottom-padding-normalizes-carrier-valley} makes
$C_r=\operatorname{Pad}_{e_t}(\nu_t,\rho_t,\mu_t)$ a valid content-$2$
two-step block.

By \cref{lem:bc-endpoint-compatible-blocks-form-replacement-window}, the
ordered list of candidate blocks concatenates to a target-size replacement
window $V^\ast_\gamma$, and the deterministic parser recovers the unchanged
anchor blocks $C_r=B_r(P)$ and the merged padded-carrier windows
$C_{u_t-1}$ in increasing order.  Therefore the hypotheses of
\cref{prop:bc-padded-carrier-valleys-recover-positive-height-zones} are
satisfied, and that proposition recovers the original height halo-zone
subpath.
\end{proof}

\begin{proposition}[Tagged collective padded splices recover positive height zones]\label{prop:bc-tagged-collective-padded-splice-recovers-positive-height-zones}
Work in the height setting of
\cref{prop:bc-height-anchor-pair-zone-recovery}.  Let
\[
  Z_\gamma\cap I=\{u_1<\cdots<u_m\},
\]
and let $A_\gamma$ be the canonical right-anchor set.  Fix an injection
\[
  \kappa:\{0,1,2\}\hookrightarrow\{0,1\}^2 .
\]
For each deleted singleton $u_t$, put $s_t=j-u_t+1$, let
\[
  \begin{matrix}
  \nu_t & \lambda_t\\
  \rho_t & \mu_t
  \end{matrix}
\]
be the corner labels of $E_{s_t,s_t+1}$, and let $e_t$ be its cell entry.

Suppose that the replacement window is parsed, by a deterministic parser
depending only on $S$, $Z_\gamma$, and the displayed replacement window, as an
ordered tagged list
\[
  (\tau_1,C_1),\ldots,(\tau_N,C_N),
  \qquad C_i=(L_i,X_i,R_i),
\]
whose tag set is exactly the disjoint union
\[
  \{\operatorname{anc}(r):r\in A_\gamma\}
  \sqcup
  \{\operatorname{car}(t):1\le t\le m\}.
\]
Assume that the blocks are defined by their tags as follows:
\[
  C_i=B_r(P)\quad\text{if }\tau_i=\operatorname{anc}(r),
\]
and
\[
  C_i=\operatorname{Pad}_{e_t}(\nu_t,\rho_t,\mu_t)
  \quad\text{if }\tau_i=\operatorname{car}(t).
\]
Assume also that the ordered blocks are endpoint-compatible:
\[
  L_1=\rho_{2a_\gamma-2},\qquad
  R_N=\rho_{2b_\gamma},\qquad
  R_i=L_{i+1}\quad(1\le i<N),
\]
and that the two-bit split of the branch word satisfies
\[
  \beta_t=\kappa(e_t)\qquad(1\le t\le m).
\]
Then the replacement window, $\beta_\gamma$, the zone endpoint pair, and
$Z_\gamma$ recover the original height halo-zone subpath.  No additional tag
word is used.
\end{proposition}

\begin{proof}
First each tagged block is a valid content-$2$ block.  Anchor blocks are
original two-step blocks of the path in \cref{lem:bc-growth-path-model}.  For
carrier tags, \cref{lem:bc-two-label-step3-diagonal-jump-two} and
\cref{cor:bc-diagonal-two-cell-has-padded-valley} give
\[
  |\nu_t|-2|\rho_t|+|\mu_t|=2-e_t,
\]
so \cref{lem:bc-bottom-padding-normalizes-carrier-valley} makes
\[
  \operatorname{Pad}_{e_t}(\nu_t,\rho_t,\mu_t)
\]
a valid content-$2$ block.

The number of tags is
\[
  |A_\gamma|+m
  =
  |Z_\gamma\setminus I|
\]
by \cref{cor:bc-height-surplus-canonical-right-anchors}.  The displayed
endpoint equalities therefore concatenate the parsed blocks to a target-size
replacement window from $\rho_{2a_\gamma-2}$ to $\rho_{2b_\gamma}$.

The deterministic parser recovers the tagged list from the replacement window
and the known zone data.  For an anchor tag
$\operatorname{anc}(r)$, this directly recovers the original block $B_r(P)$.
For a carrier tag $\operatorname{car}(t)$, the branch word recovers $e_t$
through $\kappa$, and then
\cref{lem:bc-bottom-padding-normalizes-carrier-valley} recovers
$(\nu_t,\rho_t,\mu_t)$ from the padded block and $e_t$.  Since
$E_{s_t,s_t+1}$ has one cell with entry $e_t$, this gives its cell entry and
left/bottom boundary.  Applying
\cref{cor:bc-step3-consecutive-carrier-data-recovers-boundary} with
$a=s_t$ and $b=s_t+1$ recovers the boundary segment
\[
  \rho_{2u_t-4},\rho_{2u_t-3},\rho_{2u_t-2},
  \rho_{2u_t-1},\rho_{2u_t}.
\]
By \cref{cor:bc-height-anchor-pair-step3-boundary-segment}, this is the
original two-block segment on $\{u_t-1,u_t\}$.

Finally, \cref{cor:bc-height-anchor-reduced-left-pairs} says that the blocks
of $Z_\gamma$ are partitioned into the anchor singleton blocks
$r\in A_\gamma$ and the paired intervals $\{u_t-1,u_t\}$ for
$1\le t\le m$.  Interleaving the recovered anchor blocks and recovered
two-block segments in original block order reconstructs
\[
  \rho_{2a_\gamma-2},\rho_{2a_\gamma-1},\ldots,\rho_{2b_\gamma}.
\]
\end{proof}

\begin{corollary}[Tagged splices require parser-visible tags]\label{cor:bc-tagged-splice-tags-not-free}
In the use of
\cref{prop:bc-tagged-collective-padded-splice-recovers-positive-height-zones},
the tag order and tag labels are part of the deterministic target-window
parser.  They cannot be supplied by an additional word, and they cannot be
chosen after seeing the source path unless the chosen order is itself
recoverable from the target window and the known zone data.
\end{corollary}

\begin{proof}
The only branch data used in
\cref{prop:bc-tagged-collective-padded-splice-recovers-positive-height-zones}
are the two-bit words $\beta_t=\kappa(e_t)$ for the deleted singletons.  The
proposition assumes that the replacement window is parsed as the ordered
tagged list by a deterministic parser depending only on $S$, $Z_\gamma$, and
the replacement window itself.  If the tag order were supplied by an
additional word, or chosen from the source path without being recovered by
that parser, then the map would use extra information beyond the allowed
target-size replacement window and the fixed two-bit singleton words.  Such a
map is not the hypothesis of the proposition and does not feed the
fixed-fiber count.
\end{proof}

\begin{proposition}[Grouped collective carrier splices recover positive height zones]\label{prop:bc-grouped-collective-carrier-splices-recover-positive-height-zones}
Work in the height setting of
\cref{prop:bc-height-anchor-pair-zone-recovery}.  Let
\[
  Z_\gamma\cap I=\{u_1<\cdots<u_m\},
\]
and let $A_\gamma$ be the canonical right-anchor set.  Fix an injection
\[
  \kappa:\{0,1,2\}\hookrightarrow\{0,1\}^2 .
\]
For each deleted singleton $u_t$, put $s_t=j-u_t+1$, let
\[
  \begin{matrix}
  \nu_t & \lambda_t\\
  \rho_t & \mu_t
  \end{matrix}
\]
be the corner labels of $E_{s_t,s_t+1}$, and let $e_t$ be its cell entry.

Suppose that the replacement window is parsed, by a deterministic parser
depending only on $S$, $Z_\gamma$, and that replacement window, into ordered
target subwindows
\[
  W_1,\ldots,W_L
\]
together with parser-visible tag sets
\[
  \Theta_\ell\subseteq
  \{\operatorname{anc}(r):r\in A_\gamma\}
  \sqcup
  \{\operatorname{car}(t):1\le t\le m\}.
\]
Assume the tag sets are pairwise disjoint and their union is the displayed
full tag set.  Assume also that $W_\ell$ has exactly $|\Theta_\ell|$
two-step target blocks and that the ordered subwindows concatenate from
$\rho_{2a_\gamma-2}$ to $\rho_{2b_\gamma}$.

Finally, assume that for each $\ell$ there is a fixed decoder which, from
\[
  W_\ell,\quad \Theta_\ell,\quad S,\quad Z_\gamma,\quad
  \rho_{2a_\gamma-2},\rho_{2b_\gamma},
\]
and from the words $\beta_t$ for tags $\operatorname{car}(t)\in\Theta_\ell$,
recovers:
\begin{enumerate}
\item the original block $B_r(P)$ for every
      $\operatorname{anc}(r)\in\Theta_\ell$;
\item the value $e_t$ and the carrier valley $(\nu_t,\rho_t,\mu_t)$ for every
      $\operatorname{car}(t)\in\Theta_\ell$.
\end{enumerate}
The carrier words are required to satisfy $\beta_t=\kappa(e_t)$.  Then the
replacement window, $\beta_\gamma$, the zone endpoint pair, and $Z_\gamma$
recover the original height halo-zone subpath.  No additional tag or group
word is used.
\end{proposition}

\begin{proof}
The parser-visible subwindows concatenate to a replacement window of total
length
\[
  \sum_{\ell=1}^L |\Theta_\ell|
  =
  |A_\gamma|+m
  =
  |Z_\gamma\setminus I|,
\]
where the last equality is
\cref{cor:bc-height-surplus-canonical-right-anchors}.  Thus the parsed
subwindows have the target size required by the canonical halo repair
criterion.

Apply the fixed decoder to each group.  Anchor tags directly recover the
corresponding original blocks.  For a carrier tag $\operatorname{car}(t)$,
the recovered value $e_t$ and valley $(\nu_t,\rho_t,\mu_t)$ give the unique
cell entry and left/bottom boundary of $E_{s_t,s_t+1}$.  Hence
\cref{cor:bc-step3-consecutive-carrier-data-recovers-boundary}, applied with
$a=s_t$ and $b=s_t+1$, recovers
\[
  \rho_{2u_t-4},\rho_{2u_t-3},\rho_{2u_t-2},
  \rho_{2u_t-1},\rho_{2u_t}.
\]
By \cref{cor:bc-height-anchor-pair-step3-boundary-segment}, this is the
original two-block segment on $\{u_t-1,u_t\}$.

The full tag set is recovered by the parser and is partitioned among the
groups, so all anchor blocks and all paired two-block segments are recovered.
\Cref{cor:bc-height-anchor-reduced-left-pairs} says that these pieces
partition the original blocks of $Z_\gamma$.  Interleaving them in original
block order reconstructs
\[
  \rho_{2a_\gamma-2},\rho_{2a_\gamma-1},\ldots,\rho_{2b_\gamma}.
\]
\end{proof}

\begin{proposition}[Step--3 boundary readouts supply grouped carrier decoders]\label{prop:bc-step3-boundary-readouts-supply-grouped-carrier-decoders}
Work in the setting of
\cref{prop:bc-grouped-collective-carrier-splices-recover-positive-height-zones}.
Assume that the replacement window is parsed into ordered subwindows
\[
  W_1,\ldots,W_L
\]
with parser-visible tag sets $\Theta_\ell$ satisfying the tag-partition,
length, and concatenation hypotheses of that proposition.  Fix the injection
\[
  \kappa:\{0,1,2\}\hookrightarrow\{0,1\}^2
\]
from that setting.  Suppose that, for each $\ell$, the data
\[
  W_\ell,\quad \Theta_\ell,\quad S,\quad Z_\gamma,\quad
  \rho_{2a_\gamma-2},\rho_{2b_\gamma}
\]
determine:
\begin{enumerate}
\item the original block $B_r(P)$ for every
      $\operatorname{anc}(r)\in\Theta_\ell$;
\item a Ferrers subarrangement $E_\ell$ of the Step--3 dual-RSK' half-square
      of \cref{lem:bc-krattenthaler-semistandard-path-boundary}, together
      with the partition labels on the top/right boundary of $E_\ell$;
\item for every $\operatorname{car}(t)\in\Theta_\ell$, the inclusion of the
      one-cell carrier $E_{s_t,s_t+1}$ in $E_\ell$, where
      $s_t=j-u_t+1$.
\end{enumerate}
Finally suppose that the carrier words satisfy
\[
  \beta_t=\kappa(e_t)
  \qquad(\operatorname{car}(t)\in\Theta_\ell,\ 1\le\ell\le L),
\]
where $e_t$ is the Step--3 cell entry of $E_{s_t,s_t+1}$.  Then the
replacement window, $\beta_\gamma$, the zone endpoint pair, and $Z_\gamma$
recover the original height halo-zone subpath.  No additional Step--3
boundary word is used.
\end{proposition}

\begin{proof}
Fix a group $\ell$.  By item~(2) and
\cref{cor:bc-parsed-dual-rsk-prime-boundary-recovers-subfilling}, the
displayed parser-visible data determine every cell entry and every corner
label on or inside $E_\ell$.  Hence, for each carrier tag
$\operatorname{car}(t)\in\Theta_\ell$, item~(3) identifies within the
recovered subfilling the one-cell carrier $E_{s_t,s_t+1}$ and therefore
recovers its entry $e_t$ and its left/bottom valley
\[
  (\nu_t,\rho_t,\mu_t).
\]
Item~(1) recovers the anchor blocks for the anchor tags in the same group.
Thus each group has the fixed decoder required in
\cref{prop:bc-grouped-collective-carrier-splices-recover-positive-height-zones}.
The displayed relation $\beta_t=\kappa(e_t)$ is exactly the carrier-word
hypothesis there.  Applying that proposition gives the claimed recovery of the
original zone subpath.

The top/right boundary labels of $E_\ell$ are, by hypothesis, determined by
the parsed subwindow and known zone data.  They are not supplied by an
additional branch word, so the procedure stays within the counted
target-window data.
\end{proof}

\begin{proposition}[One-cell Step--3 corner readouts supply grouped carrier decoders]\label{prop:bc-one-cell-step3-corner-readouts-supply-grouped-carrier-decoders}
Work in the setting of
\cref{prop:bc-grouped-collective-carrier-splices-recover-positive-height-zones}.
Assume that the replacement window is parsed into ordered subwindows
\[
  W_1,\ldots,W_L
\]
with parser-visible tag sets $\Theta_\ell$ satisfying the tag-partition,
length, and concatenation hypotheses of that proposition.  Let
\[
  \kappa:\{0,1,2\}\hookrightarrow\{0,1\}^2
\]
be the fixed injection from that setting.  Suppose that, for each $\ell$, the
data
\[
  W_\ell,\quad \Theta_\ell,\quad S,\quad Z_\gamma,\quad
  \rho_{2a_\gamma-2},\rho_{2b_\gamma}
\]
determine:
\begin{enumerate}
\item the original block $B_r(P)$ for every
      $\operatorname{anc}(r)\in\Theta_\ell$;
\item for every $\operatorname{car}(t)\in\Theta_\ell$, the three
      non-southwest corner labels
      \[
        \nu_t,\quad \lambda_t,\quad \mu_t
      \]
      of the one-cell Step--3 carrier
      \[
        E_{s_t,s_t+1},\qquad s_t=j-u_t+1,
      \]
      where the corner notation is
      \[
        \begin{matrix}
        \nu_t & \lambda_t\\
        \rho_t & \mu_t .
        \end{matrix}
      \]
\end{enumerate}
Finally suppose that the carrier words satisfy
\[
  \beta_t=\kappa(e_t)
  \qquad(\operatorname{car}(t)\in\Theta_\ell,\ 1\le\ell\le L),
\]
where $e_t$ is the Step--3 cell entry of $E_{s_t,s_t+1}$.
Then the replacement window, $\beta_\gamma$, the zone endpoint pair, and
$Z_\gamma$ recover the original height halo-zone subpath.  No additional
Step--3 boundary word is used.
\end{proposition}

\begin{proof}
Fix a carrier tag $\operatorname{car}(t)\in\Theta_\ell$.  By
\cref{lem:bc-dual-rsk-prime-growth-local-reversibility}, the one-cell
backward rule (B4) for the fourth arbitrary-filling dual-RSK' diagram
determines the southwest corner $\rho_t$ and the cell entry $e_t$ from the
three other corner labels
\[
  \lambda_t,\quad \mu_t,\quad \nu_t .
\]
Therefore item~(2) recovers both $e_t$ and the left/bottom carrier valley
\[
  (\nu_t,\rho_t,\mu_t)
\]
for every carrier tag in the group.  Item~(1) recovers the anchor blocks.
Thus the fixed decoders required in
\cref{prop:bc-grouped-collective-carrier-splices-recover-positive-height-zones}
are available, and the displayed relation is exactly the carrier-word
hypothesis there.  Applying that proposition gives the original height
halo-zone subpath.

The recovered corner labels are parser-visible functions of the group
subwindow and known zone data by hypothesis.  Hence this readout introduces no
new branch word beyond the already counted $\beta_\gamma$.
\end{proof}

\begin{proposition}[Padded valleys supply one-cell corner readouts]\label{prop:bc-padded-valleys-supply-one-cell-corner-readouts}
Work in the setting of
\cref{prop:bc-grouped-collective-carrier-splices-recover-positive-height-zones}.
Assume that the replacement window is parsed into ordered subwindows
\[
  W_1,\ldots,W_L
\]
with parser-visible tag sets $\Theta_\ell$ satisfying the tag-partition,
length, and concatenation hypotheses of that proposition.  Let
\[
  \kappa:\{0,1,2\}\hookrightarrow\{0,1\}^2
\]
be the fixed injection from that setting.  Suppose that, for each $\ell$, the
data
\[
  W_\ell,\quad \Theta_\ell,\quad S,\quad Z_\gamma,\quad
  \rho_{2a_\gamma-2},\rho_{2b_\gamma}
\]
determine:
\begin{enumerate}
\item the original block $B_r(P)$ for every
      $\operatorname{anc}(r)\in\Theta_\ell$;
\item for every $\operatorname{car}(t)\in\Theta_\ell$, the padded carrier
      valley
      \[
        \operatorname{Pad}_{e_t}(\nu_t,\rho_t,\mu_t)
      \]
      of the one-cell Step--3 carrier $E_{s_t,s_t+1}$, where
      $s_t=j-u_t+1$ and the corner notation is
      \[
        \begin{matrix}
        \nu_t & \lambda_t\\
        \rho_t & \mu_t .
        \end{matrix}
      \]
\end{enumerate}
Finally suppose that the carrier words satisfy
\[
  \beta_t=\kappa(e_t)
  \qquad(\operatorname{car}(t)\in\Theta_\ell,\ 1\le\ell\le L).
\]
Then the replacement window, $\beta_\gamma$, the zone endpoint pair, and
$Z_\gamma$ recover the original height halo-zone subpath.  No additional
Step--3 boundary word or corner word is used.
\end{proposition}

\begin{proof}
Fix a carrier tag $\operatorname{car}(t)\in\Theta_\ell$.  The word
$\beta_t$ determines $e_t$ because $\kappa$ is injective.  By
\cref{lem:bc-two-label-step3-diagonal-jump-two}, the carrier satisfies
\[
  |\lambda_t|-|\rho_t|=2 .
\]
Hence \cref{cor:bc-diagonal-two-cell-has-padded-valley} gives
\[
  |\nu_t|-2|\rho_t|+|\mu_t|=2-e_t .
\]
Therefore
\cref{lem:bc-bottom-padding-normalizes-carrier-valley} recovers
$(\nu_t,\rho_t,\mu_t)$ from the parser-visible padded carrier valley and the
decoded value of $e_t$.  The forward half of
\cref{lem:bc-dual-rsk-prime-growth-local-reversibility} then determines the
northeast corner $\lambda_t$ from
\[
  e_t,\quad \rho_t,\quad \mu_t,\quad \nu_t .
\]
Thus the same data determine the three non-southwest corner labels
\[
  \nu_t,\quad \lambda_t,\quad \mu_t
\]
for every carrier tag.  Together with item~(1), this supplies the hypotheses
of
\cref{prop:bc-one-cell-step3-corner-readouts-supply-grouped-carrier-decoders}.
Applying that proposition recovers the original height halo-zone subpath.
\end{proof}

\begin{corollary}[Double-successor carriers are covered by the grouped padded-valley decoder]\label{cor:bc-double-successor-covered-by-padded-valley-decoder}
In the setting of
\cref{prop:bc-padded-valleys-supply-one-cell-corner-readouts}, no additional
hypothesis $e_t\le1$ is needed.  In particular, if a carrier tag has
$e_t=2$, then the same parsed padded-valley data and the two-bit carrier word
recover that carrier valley and hence feed the grouped positive-height
decoder.
\end{corollary}

\begin{proof}
The injection
\[
  \kappa:\{0,1,2\}\hookrightarrow\{0,1\}^2
\]
is part of the grouped decoder data.  Thus the carrier word
$\beta_t=\kappa(e_t)$ recovers $e_t$ also when $e_t=2$.
\Cref{lem:bc-bottom-padding-normalizes-carrier-valley} then recovers the
unpadded valley $(\nu_t,\rho_t,\mu_t)$ from
$\operatorname{Pad}_{e_t}(\nu_t,\rho_t,\mu_t)$ and this decoded value of
$e_t$.  The proof of
\cref{prop:bc-padded-valleys-supply-one-cell-corner-readouts} uses no
additional restriction on $e_t$ beyond $e_t\in\{0,1,2\}$, so it applies to the
double-successor case exactly as stated.
\end{proof}

\begin{corollary}[Padded-valley endpoints do not determine the midpoint]\label{cor:bc-padded-valley-midpoint-not-endpoint-determined}
In the ambient one-cell dual-RSK' local category, the data
\[
  e,\qquad \operatorname{pad}_e(\nu),\qquad \mu
\]
together with the diagonal-jump condition $|\lambda|-|\rho|=2$ do not
determine the middle partition $\rho$ of the padded carrier valley
\[
  \operatorname{Pad}_e(\nu,\rho,\mu).
\]
Consequently the grouped padded-valley supplier cannot be replaced by an
endpoint-only readout of the padded left endpoint and right endpoint, unless
additional Step--3 restrictions or parser-visible data determine $\rho$ or an
equivalent midpoint datum.
\end{corollary}

\begin{proof}
Take
\[
  e=0,\qquad \nu=\mu=(2,1).
\]
There are two possible southwest corners
\[
  \rho^{(1)}=(2),\qquad \rho^{(2)}=(1,1).
\]
For each $a=1,2$, the differences $\nu/\rho^{(a)}$ and
$\mu/\rho^{(a)}$ are vertical strips.  Applying Krattenthaler's forward
algorithm (F4), in the carry notation used in the proof of
\cref{lem:bc-dual-rsk-prime-cell-size-conservation}, gives
\[
  \lambda^{(1)}=(2,1,1)
  \qquad\text{and}\qquad
  \lambda^{(2)}=(2,2).
\]
Indeed, for $\rho^{(1)}=(2)$ the carry is created in the second row and
placed in the third row, while for $\rho^{(2)}=(1,1)$ it is created in the
first row and placed in the second row.  Thus both are valid one-cell
dual-RSK' data, and in both cases
\[
  |\lambda^{(a)}|-|\rho^{(a)}|=2 .
\]
The endpoint data are identical:
\[
  \operatorname{pad}_0(\nu)=(2,1),\qquad \mu=(2,1),
\]
but the midpoint is either $(2)$ or $(1,1)$.  Hence the displayed endpoint
data and the diagonal-jump condition do not determine $\rho$.
\end{proof}

\begin{corollary}[Padded-valley midpoint defect normal form]\label{cor:bc-padded-valley-midpoint-defect-normal-form}
Let $\nu,\rho,\mu$ be partitions such that $\nu/\rho$ and $\mu/\rho$ are
vertical strips.  Put
\[
  \eta_i=\min(\nu_i,\mu_i)
\]
after extending all partitions by zero parts, and define the midpoint defect
set
\[
  D(\nu,\rho,\mu)=\{\,i:\rho_i=\eta_i-1\,\}.
\]
Then $\rho$ is determined by $\nu$, $\mu$, and $D(\nu,\rho,\mu)$ through
\[
  \rho_i=
  \begin{cases}
  \eta_i-1,& i\in D(\nu,\rho,\mu),\\
  \eta_i,& i\notin D(\nu,\rho,\mu).
  \end{cases}
\]
Moreover
\[
  |\nu|-2|\rho|+|\mu|
  =
  \sum_i |\nu_i-\mu_i|+2\,|D(\nu,\rho,\mu)|.
\]
In particular, for an actual two-label Step--3 carrier with cell entry $e$,
\[
  |D(\nu,\rho,\mu)|
  =
  \frac{2-e-\sum_i|\nu_i-\mu_i|}{2}.
\]
Thus, once $e$, $\operatorname{pad}_e(\nu)$, and $\mu$ are known, recovering
the padded carrier valley is equivalent to recovering this finite midpoint
defect set.
\end{corollary}

\begin{proof}
Because $\nu/\rho$ and $\mu/\rho$ are vertical strips, one has
\[
  \nu_i-\rho_i,\ \mu_i-\rho_i\in\{0,1\}
\]
for every row $i$.  Hence
\[
  \eta_i-\rho_i=\min(\nu_i,\mu_i)-\rho_i\in\{0,1\},
\]
which proves the displayed reconstruction formula for $\rho$ from the set
$D(\nu,\rho,\mu)$.

The size identity follows from
\[
  |\nu|-2|\rho|+|\mu|
  =
  (|\nu|+|\mu|-2|\eta|)+2(|\eta|-|\rho|).
\]
The first parenthesis is $\sum_i|\nu_i-\mu_i|$, and the second is
$2|D(\nu,\rho,\mu)|$ by the reconstruction formula.

For an actual carrier,
\cref{cor:bc-diagonal-two-cell-has-padded-valley} gives
$|\nu|-2|\rho|+|\mu|=2-e$, so the displayed cardinality formula follows.
Finally, when $e$ and $\operatorname{pad}_e(\nu)$ are known, the definition
of bottom padding recovers $\nu$ by deleting the final $e$ appended unit
rows; the reconstruction formula then shows that the only remaining datum in
the padded valley is $D(\nu,\rho,\mu)$.
\end{proof}

\begin{corollary}[Two-label midpoint defects are singleton-sized]\label{cor:bc-two-label-midpoint-defect-singleton-sized}
For an actual two-label Step--3 carrier with cell entry $e$, put
\[
  A(\nu,\mu)=\sum_i|\nu_i-\mu_i|.
\]
Then
\[
  A(\nu,\mu)+2\,|D(\nu,\rho,\mu)|=2-e.
\]
Consequently $D(\nu,\rho,\mu)$ is empty except possibly when
\[
  e=0,\qquad \nu=\mu.
\]
In that exceptional case $D(\nu,\rho,\mu)$ is a singleton.  Thus the only
midpoint information not determined by $e$, $\operatorname{pad}_e(\nu)$, and
$\mu$ is one defect row in the zero-entry equal-endpoint case.
\end{corollary}

\begin{proof}
The displayed identity is the size formula of
\cref{cor:bc-padded-valley-midpoint-defect-normal-form} combined with
\cref{cor:bc-diagonal-two-cell-has-padded-valley}.  Since
\cref{lem:bc-two-label-step3-cell-entry-three-valued} gives
$e\in\{0,1,2\}$, the right side is at most $2$.  Hence
$|D(\nu,\rho,\mu)|\le1$.

If $D(\nu,\rho,\mu)$ is nonempty, then $|D|=1$ and the displayed identity
forces $2-e-A(\nu,\mu)=2$.  Therefore $e=0$ and $A(\nu,\mu)=0$, which is
equivalent to $\nu=\mu$.  Conversely, if $e=0$ and $\nu=\mu$, then
$A(\nu,\mu)=0$, so the displayed identity gives $|D|=1$.  In all remaining
cases the defect set is empty.
\end{proof}

\begin{corollary}[The singleton defect row is the predecessor of the northeast box]\label{cor:bc-singleton-defect-row-from-northeast-box}
For an actual two-label Step--3 carrier, suppose
\[
  e=0,\qquad \nu=\mu,
\]
and write
\[
  D(\nu,\rho,\mu)=\{d\}.
\]
Then $\lambda/\nu$ consists of one box, and that box lies in row $d+1$.
Consequently the singleton defect row $d$ is determined by the row of the
unique box of $\lambda/\nu$.
\end{corollary}

\begin{proof}
Use the notation of the F4 carry computation in the proof of
\cref{lem:bc-dual-rsk-prime-cell-size-conservation}.  Since $e=0$, the
initial carry is $c_1=0$.  Since $\nu=\mu$ and $D(\nu,\rho,\mu)=\{d\}$, one
has
\[
  \min(\mu_i,\nu_i)-\rho_i=
  \begin{cases}
  1,&i=d,\\
  0,&i\ne d.
  \end{cases}
\]
For rows $i<d$ the carry is zero and the defect term is zero, so
$\lambda_i=\nu_i$.  At row $d$ the equality
$\rho_d=\nu_d-1$ gives $a_d=0$ in the F4 rule, so no box is placed in row
$d$ and the carry leaving that row is $c_{d+1}=1$.  At row $d+1$, the defect
term is zero and $\rho_{d+1}=\nu_{d+1}=\mu_{d+1}$, so $a_{d+1}=1$ and the
incoming carry is placed there:
\[
  \lambda_{d+1}=\nu_{d+1}+1,\qquad c_{d+2}=0.
\]
All later rows again have zero carry and zero defect term, hence
$\lambda_i=\nu_i$ for $i\ne d+1$.  Therefore $\lambda/\nu$ is the single box
in row $d+1$.
\end{proof}

\begin{corollary}[Exceptional northeast rows are removable-row successors]\label{cor:bc-exceptional-northeast-row-removable-normal-form}
In the exceptional case of
\cref{cor:bc-singleton-defect-row-from-northeast-box}, the defect row $d$ is a
removable-corner row of the common endpoint:
\[
  \nu_d>\nu_{d+1}.
\]
Conversely, let $\nu$ be any partition and let $d$ be any row satisfying
$\nu_d>\nu_{d+1}$.  Put $\mu=\nu$, $e=0$, and let $\rho$ be obtained from
$\nu$ by decreasing the $d$th row by one box.  Then the one-cell F4 rule
produces a valid northeast partition $\lambda$ with
\[
  \lambda/\nu=\{\text{one box in row }d+1\},
  \qquad |\lambda|-|\rho|=2 .
\]
Thus, in the zero-entry equal-endpoint case, the northeast added-box row is
exactly the successor of a removable row of the common endpoint.
\end{corollary}

\begin{proof}
Extend partitions by zero parts.  If $D(\nu,\rho,\mu)=\{d\}$ and
$\nu=\mu$, then $\rho_i=\nu_i$ for $i\ne d$ and $\rho_d=\nu_d-1$.  Since
$\rho$ is a partition,
\[
  \nu_d-1=\rho_d\ge \rho_{d+1}=\nu_{d+1},
\]
which is equivalent to $\nu_d>\nu_{d+1}$.

Conversely suppose $\nu_d>\nu_{d+1}$ and define $\rho$ by lowering row $d$.
The displayed inequality gives $\rho_d=\nu_d-1\ge\nu_{d+1}=\rho_{d+1}$, and
all other partition inequalities are unchanged, so $\rho$ is a partition.
Moreover $\nu/\rho$ and $\mu/\rho$ are vertical strips, each consisting of the
same one box in row $d$.

Run the F4 carry rule with $e=0$.  The carry is zero before row $d$.  At row
$d$, the defect term is one and $a_d=0$, so the carry leaving row $d$ is one
and no box is placed in row $d$.  At row $d+1$, the defect term is zero and
$a_{d+1}=1$, so the incoming carry is placed there and then vanishes.  Hence
\[
  \lambda_{d+1}=\nu_{d+1}+1,\qquad
  \lambda_i=\nu_i\quad(i\ne d+1).
\]
The inequality $\nu_d>\nu_{d+1}$ gives
$\lambda_d=\nu_d\ge\nu_{d+1}+1=\lambda_{d+1}$, and the remaining partition
inequalities are inherited from $\nu$.  Therefore $\lambda$ is a partition and
$\lambda/\nu$ is exactly one box in row $d+1$.  Since $|\rho|=|\nu|-1$ and
$|\lambda|=|\nu|+1$, one has $|\lambda|-|\rho|=2$.
\end{proof}

\begin{corollary}[Exceptional northeast row is not endpoint-forced]\label{cor:bc-exceptional-northeast-row-not-endpoint-forced}
The row of the exceptional northeast box is not determined by the endpoint
data
\[
  e,\qquad \operatorname{pad}_e(\nu),\qquad \mu .
\]
More precisely, with $e=0$ and $\nu=\mu=(3,2,1)$, two valid choices of the
middle partition give different northeast rows:
\[
  \rho=(2,2,1),\quad \lambda=(3,3,1),
  \qquad\text{and}\qquad
  \rho'=(3,1,1),\quad \lambda'=(3,2,2).
\]
The first has the unique box of $\lambda/\nu$ in row $2$, while the second has
the unique box of $\lambda'/\nu$ in row $3$.
\end{corollary}

\begin{proof}
For $\nu=(3,2,1)$, rows $1$ and $2$ are both removable-corner rows.  Applying
\cref{cor:bc-exceptional-northeast-row-removable-normal-form} with $d=1$
gives $\rho=(2,2,1)$ and the northeast added box in row $2$, hence
$\lambda=(3,3,1)$.  Applying the same corollary with $d=2$ gives
$\rho'=(3,1,1)$ and the northeast added box in row $3$, hence
$\lambda'=(3,2,2)$.  In both cases $e=0$ and the visible endpoint data are the
same, because $\operatorname{pad}_0(\nu)=\nu$ and $\mu=\nu$.  The two
northeast rows are different.
\end{proof}

\begin{corollary}[Exceptional removable-row ambiguity is unbounded]\label{cor:bc-exceptional-removable-row-unbounded-endpoint-ambiguity}
For every $k\ge1$ there is a common endpoint partition $\nu=\mu$ for which
the exceptional one-cell data
\[
  e=0,\qquad \operatorname{pad}_0(\nu)=\nu,\qquad \mu=\nu
\]
admit at least $k$ different removable-row choices in the sense of
\cref{cor:bc-exceptional-northeast-row-removable-normal-form}.  Consequently
no fixed local word of length $c$ bits, together with the endpoint data
$(e,\operatorname{pad}_e(\nu),\mu)$, can recover the exceptional removable
row for all ambient one-cell dual-RSK' data once $k>2^c$.
\end{corollary}

\begin{proof}
Let
\[
  \nu=\mu=(k,k-1,\ldots,2,1).
\]
For each $1\le d\le k$ one has $\nu_d>\nu_{d+1}$, with the convention
$\nu_{k+1}=0$.  Hence $d$ is a removable-corner row of the common endpoint.
By \cref{cor:bc-exceptional-northeast-row-removable-normal-form}, for each
such $d$ the partition $\rho^{(d)}$ obtained by decreasing the $d$th row of
$\nu$ by one box, together with $e=0$ and $\mu=\nu$, gives a valid one-cell
F4 datum whose northeast corner $\lambda^{(d)}$ satisfies
\[
  \lambda^{(d)}/\nu=\{\text{one box in row }d+1\},
  \qquad |\lambda^{(d)}|-|\rho^{(d)}|=2 .
\]
The $k$ choices have the same endpoint data
$(e,\operatorname{pad}_e(\nu),\mu)=(0,\nu,\nu)$ but different removable rows.
Thus a word of length $c$ bits can distinguish at most $2^c$ of them.  Choosing
$k>2^c$ proves the final assertion.
\end{proof}

\begin{corollary}[Two-bit carrier words do not recover exceptional removable rows]\label{cor:bc-two-bit-word-not-removable-row-readout}
In the ambient one-cell dual-RSK' local category, the exceptional removable
row cannot be recovered for all data from
\[
  e,\qquad \operatorname{pad}_e(\nu),\qquad \mu,
\]
and one additional two-bit word.  In particular, the counted carrier word
$\beta=\kappa(e)$ in the grouped positive-height supplier cannot by itself
serve as the removable-row readout in
\cref{prop:bc-removable-row-readouts-supply-grouped-padded-valleys}.
\end{corollary}

\begin{proof}
Apply
\cref{cor:bc-exceptional-removable-row-unbounded-endpoint-ambiguity} with
$k=5$.  The five valid exceptional one-cell data have the same values of
\[
  e=0,\qquad \operatorname{pad}_e(\nu)=\nu,\qquad \mu=\nu,
\]
but have five different removable rows.  A two-bit word has only four values,
so two of these five data have the same displayed input and the same two-bit
word.  No decoder from those data can distinguish their removable rows.

If the word is the counted carrier word $\beta=\kappa(e)$, then in the
exceptional case $e=0$ it is fixed to the single value $\kappa(0)$, so it
carries no removable-row choice at all.
\end{proof}

\begin{proposition}[Midpoint-defect readouts supply grouped padded valleys]\label{prop:bc-midpoint-defect-readouts-supply-grouped-padded-valleys}
Work in the setting of
\cref{prop:bc-grouped-collective-carrier-splices-recover-positive-height-zones}.
Assume that the replacement window is parsed into ordered subwindows
\[
  W_1,\ldots,W_L
\]
with parser-visible tag sets $\Theta_\ell$ satisfying the tag-partition,
length, and concatenation hypotheses of that proposition.  Let
\[
  \kappa:\{0,1,2\}\hookrightarrow\{0,1\}^2
\]
be the fixed injection from that setting.  Suppose that, for each $\ell$, the
data
\[
  W_\ell,\quad \Theta_\ell,\quad S,\quad Z_\gamma,\quad
  \rho_{2a_\gamma-2},\rho_{2b_\gamma}
\]
determine:
\begin{enumerate}
\item the original block $B_r(P)$ for every
      $\operatorname{anc}(r)\in\Theta_\ell$;
\item for every $\operatorname{car}(t)\in\Theta_\ell$, the padded left
      endpoint $\operatorname{pad}_{e_t}(\nu_t)$, the right endpoint $\mu_t$,
      and the midpoint defect set
      \[
        D(\nu_t,\rho_t,\mu_t)
      \]
      of the one-cell Step--3 carrier $E_{s_t,s_t+1}$, where
      $s_t=j-u_t+1$.
\end{enumerate}
Finally suppose that the carrier words satisfy
\[
  \beta_t=\kappa(e_t)
  \qquad(\operatorname{car}(t)\in\Theta_\ell,\ 1\le\ell\le L).
\]
Then the replacement window, $\beta_\gamma$, the zone endpoint pair, and
$Z_\gamma$ recover the original height halo-zone subpath.  No additional
midpoint word is used.
\end{proposition}

\begin{proof}
For a carrier tag $\operatorname{car}(t)\in\Theta_\ell$, the word $\beta_t$
determines $e_t$ through the fixed injection $\kappa$.  From
$\operatorname{pad}_{e_t}(\nu_t)$ and $e_t$ one recovers $\nu_t$ by deleting
the final $e_t$ appended unit rows.  Then
\cref{cor:bc-padded-valley-midpoint-defect-normal-form} reconstructs
$\rho_t$ from $\nu_t$, $\mu_t$, and the recovered defect set
$D(\nu_t,\rho_t,\mu_t)$.  Thus the same parser-visible data determine the
padded carrier valley
\[
  \operatorname{Pad}_{e_t}(\nu_t,\rho_t,\mu_t)
\]
for every carrier tag in the group.

Together with item~(1), these are exactly the hypotheses of
\cref{prop:bc-padded-valleys-supply-one-cell-corner-readouts}.  Applying that
proposition recovers the original height halo-zone subpath.
\end{proof}

\begin{proposition}[Singleton-defect row readouts supply grouped padded valleys]\label{prop:bc-singleton-defect-row-readouts-supply-grouped-padded-valleys}
Work in the setting of
\cref{prop:bc-grouped-collective-carrier-splices-recover-positive-height-zones}.
Assume that the replacement window is parsed into ordered subwindows
\[
  W_1,\ldots,W_L
\]
with parser-visible tag sets $\Theta_\ell$ satisfying the tag-partition,
length, and concatenation hypotheses of that proposition.  Let
\[
  \kappa:\{0,1,2\}\hookrightarrow\{0,1\}^2
\]
be the fixed injection from that setting.  Suppose that, for each $\ell$, the
data
\[
  W_\ell,\quad \Theta_\ell,\quad S,\quad Z_\gamma,\quad
  \rho_{2a_\gamma-2},\rho_{2b_\gamma}
\]
determine:
\begin{enumerate}
\item the original block $B_r(P)$ for every
      $\operatorname{anc}(r)\in\Theta_\ell$;
\item for every $\operatorname{car}(t)\in\Theta_\ell$, the padded left
      endpoint $\operatorname{pad}_{e_t}(\nu_t)$ and the right endpoint
      $\mu_t$ of the one-cell Step--3 carrier $E_{s_t,s_t+1}$, where
      $s_t=j-u_t+1$;
\item for every $\operatorname{car}(t)\in\Theta_\ell$ with
      $e_t=0$ and $\nu_t=\mu_t$, the unique row $d_t$ of the singleton defect
      set $D(\nu_t,\rho_t,\mu_t)$.
\end{enumerate}
Finally suppose that the carrier words satisfy
\[
  \beta_t=\kappa(e_t)
  \qquad(\operatorname{car}(t)\in\Theta_\ell,\ 1\le\ell\le L).
\]
Then the replacement window, $\beta_\gamma$, the zone endpoint pair, and
$Z_\gamma$ recover the original height halo-zone subpath.  No additional
defect-set word is used.
\end{proposition}

\begin{proof}
For a carrier tag $\operatorname{car}(t)\in\Theta_\ell$, the word $\beta_t$
determines $e_t$ through the fixed injection $\kappa$.  From
$\operatorname{pad}_{e_t}(\nu_t)$ and $e_t$ one recovers $\nu_t$ by deleting
the final $e_t$ appended unit rows.  The right endpoint $\mu_t$ is part of
item~(2).

By \cref{cor:bc-two-label-midpoint-defect-singleton-sized}, the defect set is
empty unless $e_t=0$ and $\nu_t=\mu_t$.  In the exceptional case item~(3)
recovers its unique row $d_t$, and hence recovers
\[
  D(\nu_t,\rho_t,\mu_t)=\{d_t\}.
\]
Thus the data determine the midpoint defect set for every carrier tag.
\Cref{prop:bc-midpoint-defect-readouts-supply-grouped-padded-valleys} then
recovers the original height halo-zone subpath.
\end{proof}

\begin{proposition}[Northeast-box row readouts supply grouped padded valleys]\label{prop:bc-northeast-box-row-readouts-supply-grouped-padded-valleys}
Work in the setting of
\cref{prop:bc-grouped-collective-carrier-splices-recover-positive-height-zones}.
Assume that the replacement window is parsed into ordered subwindows
\[
  W_1,\ldots,W_L
\]
with parser-visible tag sets $\Theta_\ell$ satisfying the tag-partition,
length, and concatenation hypotheses of that proposition.  Let
\[
  \kappa:\{0,1,2\}\hookrightarrow\{0,1\}^2
\]
be the fixed injection from that setting.  Suppose that, for each $\ell$, the
data
\[
  W_\ell,\quad \Theta_\ell,\quad S,\quad Z_\gamma,\quad
  \rho_{2a_\gamma-2},\rho_{2b_\gamma}
\]
determine:
\begin{enumerate}
\item the original block $B_r(P)$ for every
      $\operatorname{anc}(r)\in\Theta_\ell$;
\item for every $\operatorname{car}(t)\in\Theta_\ell$, the padded left
      endpoint $\operatorname{pad}_{e_t}(\nu_t)$ and the right endpoint
      $\mu_t$ of the one-cell Step--3 carrier $E_{s_t,s_t+1}$, where
      $s_t=j-u_t+1$;
\item for every $\operatorname{car}(t)\in\Theta_\ell$ with
      $e_t=0$ and $\nu_t=\mu_t$, the row $r_t$ of the unique box in
      $\lambda_t/\nu_t$.
\end{enumerate}
Finally suppose that the carrier words satisfy
\[
  \beta_t=\kappa(e_t)
  \qquad(\operatorname{car}(t)\in\Theta_\ell,\ 1\le\ell\le L).
\]
Then the replacement window, $\beta_\gamma$, the zone endpoint pair, and
$Z_\gamma$ recover the original height halo-zone subpath.  No additional
defect-row word is used.
\end{proposition}

\begin{proof}
For each carrier tag, the word $\beta_t$ determines $e_t$, and item~(2)
recovers $\nu_t$ from $\operatorname{pad}_{e_t}(\nu_t)$ and $e_t$.  If the
carrier is not in the exceptional case $e_t=0$ and $\nu_t=\mu_t$, then
\cref{cor:bc-two-label-midpoint-defect-singleton-sized} makes the defect set
empty.

In the exceptional case, item~(3) gives the row $r_t$ of the unique box of
$\lambda_t/\nu_t$.  By
\cref{cor:bc-singleton-defect-row-from-northeast-box}, the singleton defect
row is $d_t=r_t-1$.  Hence
\[
  D(\nu_t,\rho_t,\mu_t)=\{r_t-1\}.
\]
Thus the hypotheses of
\cref{prop:bc-singleton-defect-row-readouts-supply-grouped-padded-valleys}
are satisfied, and that proposition recovers the original height halo-zone
subpath.
\end{proof}

\begin{proposition}[Removable-row readouts supply grouped padded valleys]\label{prop:bc-removable-row-readouts-supply-grouped-padded-valleys}
Work in the setting of
\cref{prop:bc-grouped-collective-carrier-splices-recover-positive-height-zones}.
Assume that the replacement window is parsed into ordered subwindows
\[
  W_1,\ldots,W_L
\]
with parser-visible tag sets $\Theta_\ell$ satisfying the tag-partition,
length, and concatenation hypotheses of that proposition.  Let
\[
  \kappa:\{0,1,2\}\hookrightarrow\{0,1\}^2
\]
be the fixed injection from that setting.  Suppose that, for each $\ell$, the
data
\[
  W_\ell,\quad \Theta_\ell,\quad S,\quad Z_\gamma,\quad
  \rho_{2a_\gamma-2},\rho_{2b_\gamma}
\]
determine:
\begin{enumerate}
\item the original block $B_r(P)$ for every
      $\operatorname{anc}(r)\in\Theta_\ell$;
\item for every $\operatorname{car}(t)\in\Theta_\ell$, the padded left
      endpoint $\operatorname{pad}_{e_t}(\nu_t)$ and the right endpoint
      $\mu_t$ of the one-cell Step--3 carrier $E_{s_t,s_t+1}$, where
      $s_t=j-u_t+1$;
\item for every $\operatorname{car}(t)\in\Theta_\ell$ with
      $e_t=0$ and $\nu_t=\mu_t$, the removable-corner row $d_t$ of $\nu_t$
      whose successor row is the row of the unique box in $\lambda_t/\nu_t$.
\end{enumerate}
Finally suppose that the carrier words satisfy
\[
  \beta_t=\kappa(e_t)
  \qquad(\operatorname{car}(t)\in\Theta_\ell,\ 1\le\ell\le L).
\]
Then the replacement window, $\beta_\gamma$, the zone endpoint pair, and
$Z_\gamma$ recover the original height halo-zone subpath.  No northeast-corner
partition readout is used.
\end{proposition}

\begin{proof}
For each carrier tag, the word $\beta_t$ determines $e_t$, and item~(2)
recovers $\nu_t$ from $\operatorname{pad}_{e_t}(\nu_t)$ and $e_t$.  If the
carrier is not in the exceptional case $e_t=0$ and $\nu_t=\mu_t$, then
\cref{cor:bc-two-label-midpoint-defect-singleton-sized} makes the defect set
empty.

In the exceptional case, item~(3) gives the removable row $d_t$.  By
\cref{cor:bc-exceptional-northeast-row-removable-normal-form}, the unique box
of $\lambda_t/\nu_t$ lies in row $d_t+1$.  Thus the hypotheses of
\cref{prop:bc-northeast-box-row-readouts-supply-grouped-padded-valleys} are
satisfied, with $r_t=d_t+1$, and that proposition recovers the original height
halo-zone subpath.
\end{proof}

\begin{corollary}[The removable-row exception gate is parser-visible]\label{cor:bc-removable-row-exception-gate-parser-visible}
Work in the setting of
\cref{prop:bc-removable-row-readouts-supply-grouped-padded-valleys}.  For a
carrier tag $\operatorname{car}(t)\in\Theta_\ell$, suppose the group data
determine the padded left endpoint $\operatorname{pad}_{e_t}(\nu_t)$ and the
right endpoint $\mu_t$, and suppose the carrier word satisfies
$\beta_t=\kappa(e_t)$.  Then
\[
  e_t=0\ \text{and}\ \nu_t=\mu_t
\]
if and only if
\[
  \beta_t=\kappa(0)
  \qquad\text{and}\qquad
  \operatorname{pad}_{e_t}(\nu_t)=\mu_t .
\]
Consequently the set of carrier tags for which the removable-row readout in
item~(3) of
\cref{prop:bc-removable-row-readouts-supply-grouped-padded-valleys} is required
is determined from the parser-visible group data and the already counted carrier
words, before any removable row is read.
\end{corollary}

\begin{proof}
Since $\kappa$ is injective and $\beta_t=\kappa(e_t)$, the equality
$\beta_t=\kappa(0)$ is equivalent to $e_t=0$.  When $e_t=0$, the bottom-padding
convention in \cref{def:bc-bottom-padding} gives
\[
  \operatorname{pad}_{e_t}(\nu_t)=\operatorname{pad}_0(\nu_t)=\nu_t .
\]
Thus, under $\beta_t=\kappa(0)$, the equality
$\operatorname{pad}_{e_t}(\nu_t)=\mu_t$ is equivalent to $\nu_t=\mu_t$.
This proves both implications.  The final sentence follows because item~(2) of
\cref{prop:bc-removable-row-readouts-supply-grouped-padded-valleys} makes the
two displayed endpoint partitions part of the group readout, while $\beta_t$ is
the already counted carrier word.
\end{proof}

\begin{proposition}[Visible-gated northeast rows supply grouped padded valleys]\label{prop:bc-visible-gated-northeast-row-readouts-supply-grouped-padded-valleys}
Work in the setting of
\cref{prop:bc-grouped-collective-carrier-splices-recover-positive-height-zones}.
Assume that the replacement window is parsed into ordered subwindows
\[
  W_1,\ldots,W_L
\]
with parser-visible tag sets $\Theta_\ell$ satisfying the tag-partition,
length, and concatenation hypotheses of that proposition.  Let
\[
  \kappa:\{0,1,2\}\hookrightarrow\{0,1\}^2
\]
be the fixed injection from that setting.  Suppose that, for each $\ell$, the
data
\[
  W_\ell,\quad \Theta_\ell,\quad S,\quad Z_\gamma,\quad
  \rho_{2a_\gamma-2},\rho_{2b_\gamma}
\]
determine:
\begin{enumerate}
\item the original block $B_r(P)$ for every
      $\operatorname{anc}(r)\in\Theta_\ell$;
\item for every $\operatorname{car}(t)\in\Theta_\ell$, the padded left
      endpoint $\operatorname{pad}_{e_t}(\nu_t)$ and the right endpoint
      $\mu_t$ of the one-cell Step--3 carrier $E_{s_t,s_t+1}$, where
      $s_t=j-u_t+1$;
\item for every carrier tag satisfying the visible gate
      \[
        \beta_t=\kappa(0)
        \qquad\text{and}\qquad
        \operatorname{pad}_{e_t}(\nu_t)=\mu_t,
      \]
      the row $r_t$ of the unique box in $\lambda_t/\nu_t$.
\end{enumerate}
Finally suppose that the carrier words satisfy
\[
  \beta_t=\kappa(e_t)
  \qquad(\operatorname{car}(t)\in\Theta_\ell,\ 1\le\ell\le L).
\]
Then the replacement window, $\beta_\gamma$, the zone endpoint pair, and
$Z_\gamma$ recover the original height halo-zone subpath.  No full
visible-gated corner or boundary readout is used.
\end{proposition}

\begin{proof}
By \cref{cor:bc-removable-row-exception-gate-parser-visible}, the displayed
visible gate is equivalent, for each carrier tag, to
\[
  e_t=0,\qquad \nu_t=\mu_t .
\]
Thus item~(3) supplies the northeast-box row exactly for the carrier tags
required in item~(3) of
\cref{prop:bc-northeast-box-row-readouts-supply-grouped-padded-valleys}.
Items~(1) and (2) are the same as in that proposition, and the carrier-word
relation is also the same.  Applying
\cref{prop:bc-northeast-box-row-readouts-supply-grouped-padded-valleys}
therefore recovers the original height halo-zone subpath.
\end{proof}

\begin{corollary}[Visible-gated northeast rows are target-block midpoints]\label{cor:bc-visible-gated-northeast-row-target-block}
Let
\[
  \begin{matrix}
  \nu & \lambda\\
  \rho & \mu
  \end{matrix}
\]
be the corner labels of an actual two-label Step--3 carrier, and suppose that
its entry is $e=0$ and that $\nu=\mu$.  Let $r$ be the row of the unique box in
$\lambda/\nu$.  Define $X$ by lowering row $r-1$ of $\nu$ by one box and
leaving all other rows unchanged.  Then
\[
  X=\rho .
\]
Consequently
\[
  (\nu,X,\nu)
\]
is a valid content-$2$ target two-step block.  Conversely, this equal-endpoint
target block determines the northeast-box row $r$ as one plus the unique row in
which $\nu$ and $X$ differ.
\end{corollary}

\begin{proof}
By \cref{cor:bc-two-label-midpoint-defect-singleton-sized}, the defect set
$D(\nu,\rho,\mu)$ is a singleton.  Write it as $\{d\}$.  Since $\nu=\mu$,
\cref{cor:bc-padded-valley-midpoint-defect-normal-form} gives
\[
  \rho_i=
  \begin{cases}
  \nu_i-1,& i=d,\\
  \nu_i,& i\ne d.
  \end{cases}
\]
By \cref{cor:bc-singleton-defect-row-from-northeast-box}, the unique box of
$\lambda/\nu$ lies in row $d+1$.  Hence $d=r-1$, and the displayed formula is
exactly the definition of $X$.  Thus $X=\rho$.

The differences $\nu/\rho$ and $\mu/\rho=\nu/\rho$ are vertical strips, and
$|\rho|=|\nu|-1$.  Therefore
\[
  |\nu|-2|\rho|+|\nu|=2,
\]
so $(\nu,X,\nu)=(\nu,\rho,\mu)$ is a valid content-$2$ target two-step block.
The final assertion follows because this block has equal endpoints and its
middle partition differs from that endpoint in the single row $d=r-1$.
\end{proof}

\begin{proposition}[Visible-gated removable rows supply grouped padded valleys]\label{prop:bc-visible-gated-removable-row-readouts-supply-grouped-padded-valleys}
Work in the setting of
\cref{prop:bc-grouped-collective-carrier-splices-recover-positive-height-zones}.
Assume that the replacement window is parsed into ordered subwindows
\[
  W_1,\ldots,W_L
\]
with parser-visible tag sets $\Theta_\ell$ satisfying the tag-partition,
length, and concatenation hypotheses of that proposition.  Let
\[
  \kappa:\{0,1,2\}\hookrightarrow\{0,1\}^2
\]
be the fixed injection from that setting.  Suppose that, for each $\ell$, the
data
\[
  W_\ell,\quad \Theta_\ell,\quad S,\quad Z_\gamma,\quad
  \rho_{2a_\gamma-2},\rho_{2b_\gamma}
\]
determine:
\begin{enumerate}
\item the original block $B_r(P)$ for every
      $\operatorname{anc}(r)\in\Theta_\ell$;
\item for every $\operatorname{car}(t)\in\Theta_\ell$, the padded left
      endpoint $\operatorname{pad}_{e_t}(\nu_t)$ and the right endpoint
      $\mu_t$ of the one-cell Step--3 carrier $E_{s_t,s_t+1}$, where
      $s_t=j-u_t+1$.
\end{enumerate}
Finally suppose that the carrier words satisfy
\[
  \beta_t=\kappa(e_t)
  \qquad(\operatorname{car}(t)\in\Theta_\ell,\ 1\le\ell\le L),
\]
and that, for every carrier tag satisfying the visible gate
\[
  \beta_t=\kappa(0)
  \qquad\text{and}\qquad
  \operatorname{pad}_{e_t}(\nu_t)=\mu_t,
\]
the parser-visible group data determine the removable-corner row $d_t$ of the
common endpoint whose successor row is the row of the unique box in
$\lambda_t/\nu_t$.  Then the replacement window, $\beta_\gamma$, the zone
endpoint pair, and $Z_\gamma$ recover the original height halo-zone subpath.  No
hidden source-side exception predicate is used.
\end{proposition}

\begin{proof}
By \cref{cor:bc-removable-row-exception-gate-parser-visible}, the displayed
visible gate is equivalent, for each carrier tag, to the source condition
\[
  e_t=0,\qquad \nu_t=\mu_t .
\]
Thus the removable rows supplied in the final hypothesis are supplied exactly
for the carrier tags required in item~(3) of
\cref{prop:bc-removable-row-readouts-supply-grouped-padded-valleys}.  Items~(1)
and (2) are the same as in that proposition, and the carrier-word relation is
also the same.  Therefore
\cref{prop:bc-removable-row-readouts-supply-grouped-padded-valleys} applies and
recovers the original height halo-zone subpath.
\end{proof}

\begin{proposition}[Visible-gated midpoints supply removable rows]\label{prop:bc-visible-gated-midpoint-readouts-supply-grouped-padded-valleys}
Work in the setting of
\cref{prop:bc-grouped-collective-carrier-splices-recover-positive-height-zones}.
Assume that the replacement window is parsed into ordered subwindows
\[
  W_1,\ldots,W_L
\]
with parser-visible tag sets $\Theta_\ell$ satisfying the tag-partition,
length, and concatenation hypotheses of that proposition.  Let
\[
  \kappa:\{0,1,2\}\hookrightarrow\{0,1\}^2
\]
be the fixed injection from that setting.  Suppose that, for each $\ell$, the
data
\[
  W_\ell,\quad \Theta_\ell,\quad S,\quad Z_\gamma,\quad
  \rho_{2a_\gamma-2},\rho_{2b_\gamma}
\]
determine:
\begin{enumerate}
\item the original block $B_r(P)$ for every
      $\operatorname{anc}(r)\in\Theta_\ell$;
\item for every $\operatorname{car}(t)\in\Theta_\ell$, the padded left
      endpoint $\operatorname{pad}_{e_t}(\nu_t)$ and the right endpoint
      $\mu_t$ of the one-cell Step--3 carrier $E_{s_t,s_t+1}$, where
      $s_t=j-u_t+1$;
\item for every carrier tag satisfying the visible gate
      \[
        \beta_t=\kappa(0)
        \qquad\text{and}\qquad
        \operatorname{pad}_{e_t}(\nu_t)=\mu_t,
      \]
      the southwest corner $\rho_t$ of the same one-cell carrier.
\end{enumerate}
Finally suppose that the carrier words satisfy
\[
  \beta_t=\kappa(e_t)
  \qquad(\operatorname{car}(t)\in\Theta_\ell,\ 1\le\ell\le L).
\]
Then the replacement window, $\beta_\gamma$, the zone endpoint pair, and
$Z_\gamma$ recover the original height halo-zone subpath.  Equivalently, for
flagged carrier tags, the remaining row datum in the visible-gated local inverse
is supplied by the recovered middle partition.
\end{proposition}

\begin{proof}
Fix a flagged carrier tag.  By
\cref{cor:bc-removable-row-exception-gate-parser-visible}, the visible gate is
equivalent to
\[
  e_t=0,\qquad \nu_t=\mu_t .
\]
Since $e_t=0$, the padded left endpoint is $\nu_t$ itself.  Item~(3) gives
$\rho_t$, and \cref{cor:bc-two-label-midpoint-defect-singleton-sized} says that
the defect set $D(\nu_t,\rho_t,\mu_t)$ is a singleton.  Because
$\nu_t=\mu_t$, the definition of the defect set reduces to
\[
  D(\nu_t,\rho_t,\mu_t)
  =
  \{\,i:\rho_{t,i}=\nu_{t,i}-1\,\}.
\]
Let $d_t$ be this unique row.  Then
\cref{cor:bc-exceptional-northeast-row-removable-normal-form} says that $d_t$
is the removable-corner row of the common endpoint whose successor row is the
row of the unique box in $\lambda_t/\nu_t$.

Thus item~(3) of
\cref{prop:bc-visible-gated-removable-row-readouts-supply-grouped-padded-valleys}
is satisfied for every flagged carrier tag, while items~(1), (2), and the
carrier-word relation are exactly the remaining hypotheses there.  Applying
that proposition recovers the original height halo-zone subpath.
\end{proof}

\begin{proposition}[Visible-gated target blocks recover positive height zones]\label{prop:bc-visible-gated-target-blocks-recover-positive-height-zones}
Work in the height setting of
\cref{prop:bc-height-anchor-pair-zone-recovery}.  Let
\[
  I=I(S),
\]
fix an injection
\[
  \kappa:\{0,1,2\}\hookrightarrow\{0,1\}^2 ,
\]
and write
\[
  Z_\gamma\cap I=\{u_1<\cdots<u_m\}.
\]
Let $\beta_t$ denote the deterministic two-bit subword of $\beta_\gamma$
attached to the retained carrier slot $u_t$.
Let $A_\gamma$ be the canonical right-anchor set.  For each retained slot
$r\in Z_\gamma\setminus I$, suppose a candidate target block
\[
  C_r=(L_r,X_r,R_r)
\]
is given as follows.
\begin{enumerate}
\item If $r\in A_\gamma$, then $C_r=B_r(P)$.
\item If $r=u_t-1$ for some $t$, put $s_t=j-u_t+1$, let
      \[
      \begin{matrix}
      \nu_t & \lambda_t\\
      \rho_t & \mu_t
      \end{matrix}
      \]
      be the corner labels of $E_{s_t,s_t+1}$, and let $e_t$ be its cell
      entry.  Assume that $C_r$ is a valid content-$2$ two-step block with
      \[
        L_r=\operatorname{pad}_{e_t}(\nu_t),
        \qquad
        R_r=\mu_t .
      \]
      If the visible gate
      \[
        \beta_t=\kappa(0)
        \qquad\text{and}\qquad
        L_r=R_r
      \]
      holds, let $r_t$ be the row of the unique box in
      $\lambda_t/\nu_t$, and assume that $X_r$ is obtained from the common
      endpoint $L_r=R_r$ by lowering row $r_t-1$ by one box and leaving all
      other rows unchanged.
\end{enumerate}
Assume that the ordered blocks $C_r$ satisfy the endpoint equalities in
\cref{lem:bc-endpoint-compatible-blocks-form-replacement-window}, and that the
two-bit split of the branch word satisfies $\beta_t=\kappa(e_t)$ for every
$t$.  Then the original height halo-zone subpath is recovered from the resulting
replacement window, $\beta_\gamma$, the zone endpoint pair, and $Z_\gamma$.
\end{proposition}

\begin{proof}
By \cref{lem:bc-endpoint-compatible-blocks-form-replacement-window}, the ordered
candidate blocks concatenate to a target-size replacement window, and the
deterministic parser recovers each block $C_r=(L_r,X_r,R_r)$ in slot order.
Use singleton tag sets in that order: $\operatorname{anc}(r)$ for
$r\in A_\gamma$ and $\operatorname{car}(t)$ for $r=u_t-1$.  These tags are
determined by $S$ and $Z_\gamma$, have total size $|Z_\gamma\setminus I|$, and
partition the anchor and carrier tags required in
\cref{prop:bc-grouped-collective-carrier-splices-recover-positive-height-zones}.

For an anchor singleton, the parsed block is $B_r(P)$ by item~(1).  For a
carrier singleton, item~(2) identifies the parsed left and right endpoints with
\[
  \operatorname{pad}_{e_t}(\nu_t),\qquad \mu_t .
\]
Since $\beta_t=\kappa(e_t)$, the displayed gate $L_r=R_r$ is exactly
\[
  \operatorname{pad}_{e_t}(\nu_t)=\mu_t
\]
and hence is the visible gate of
\cref{prop:bc-visible-gated-midpoint-readouts-supply-grouped-padded-valleys}.
When that gate holds, the equality $\beta_t=\kappa(0)$ and injectivity of
$\kappa$ give $e_t=0$.  The endpoint equality then gives
\[
  L_r=R_r=\nu_t=\mu_t .
\]
By \cref{cor:bc-visible-gated-northeast-row-target-block}, the row-encoded
middle partition specified in item~(2) is the southwest corner $\rho_t$.  Thus
the same parsed singleton block supplies the flagged southwest corner required
there.

Thus the singleton parsing satisfies the hypotheses of
\cref{prop:bc-visible-gated-midpoint-readouts-supply-grouped-padded-valleys}.
Applying that proposition recovers the original height halo-zone subpath.
\end{proof}

\begin{corollary}[Singleton target blocks do not replace the grouped source]\label{cor:bc-singleton-target-blocks-do-not-replace-grouped-source}
The source theorem needed for the positive-height $C$-side supplier is not
obtained by declaring the actual one-cell carrier valleys and unchanged anchors
to be an automatically endpoint-compatible list of singleton target blocks.
The criterion of
\cref{prop:bc-visible-gated-target-blocks-recover-positive-height-zones} is a
sufficient local inverse once such blocks have been constructed; it does not
construct them from the height halo hypotheses alone.  In particular, the live
source statement remains the grouped-window readout of
\cref{prop:bc-visible-gated-midpoint-readouts-supply-grouped-padded-valleys},
or a collective target-window construction which supplies the same padded
endpoints and visible-gated middle partitions.
\end{corollary}

\begin{proof}
\Cref{prop:bc-visible-gated-target-blocks-recover-positive-height-zones}
assumes the endpoint equalities of
\cref{lem:bc-endpoint-compatible-blocks-form-replacement-window}.  Those
equalities are extra source data: by
\cref{cor:bc-pair-endpoint-padding-not-source-automatic}, the bottom-padded
one-cell carrier valley is not forced by the Step--3 one-cell boundary
convention to span the adjacent original two-block interval.

Nor can the endpoint problem be removed by a zero-padding or separate-buffer
argument.  \Cref{cor:bc-positive-height-nonzero-padding-example} gives a valid
positive-height zone with nonzero carrier padding, and with more nonzero padded
carriers than surplus anchors.  \Cref{cor:bc-positive-height-double-successor-carrier-example}
shows that the full $e=2$ carrier case occurs inside a positive-height zone.
Thus the singleton target-block criterion is a valid sufficient inverse only
after a source construction has supplied the endpoint-compatible blocks.  The
source obligation that remains attached to \cref{rem:former-post29-tail} is therefore
the collective grouped target-window readout already isolated in
\cref{prop:bc-visible-gated-midpoint-readouts-supply-grouped-padded-valleys}.
\end{proof}

\begin{corollary}[One-cell Step--3 corners supply visible-gated midpoint readouts]\label{cor:bc-one-cell-corners-supply-visible-gated-midpoint-readouts}
Work in the setting of
\cref{prop:bc-grouped-collective-carrier-splices-recover-positive-height-zones}.
Assume that the replacement window is parsed into ordered subwindows
\[
  W_1,\ldots,W_L
\]
with parser-visible tag sets $\Theta_\ell$ satisfying the tag-partition,
length, and concatenation hypotheses of that proposition.  Let
\[
  \kappa:\{0,1,2\}\hookrightarrow\{0,1\}^2
\]
be the fixed injection from that setting.  Suppose that, for each $\ell$, the
data
\[
  W_\ell,\quad \Theta_\ell,\quad S,\quad Z_\gamma,\quad
  \rho_{2a_\gamma-2},\rho_{2b_\gamma}
\]
determine:
\begin{enumerate}
\item the original block $B_r(P)$ for every
      $\operatorname{anc}(r)\in\Theta_\ell$;
\item for every $\operatorname{car}(t)\in\Theta_\ell$, the three non-southwest
      corner labels
      \[
        \nu_t,\quad \lambda_t,\quad \mu_t
      \]
      of the one-cell Step--3 carrier $E_{s_t,s_t+1}$, where
      $s_t=j-u_t+1$.
\end{enumerate}
Finally suppose that the carrier words satisfy
\[
  \beta_t=\kappa(e_t)
  \qquad(\operatorname{car}(t)\in\Theta_\ell,\ 1\le\ell\le L),
\]
where $e_t$ is the Step--3 cell entry of $E_{s_t,s_t+1}$.  Then the
hypotheses of
\cref{prop:bc-visible-gated-midpoint-readouts-supply-grouped-padded-valleys}
hold.  Consequently the replacement window, $\beta_\gamma$, the zone endpoint
pair, and $Z_\gamma$ recover the original height halo-zone subpath.
\end{corollary}

\begin{proof}
Fix a carrier tag $\operatorname{car}(t)\in\Theta_\ell$.  By
\cref{lem:bc-dual-rsk-prime-growth-local-reversibility}, the one-cell backward
rule recovers the southwest corner $\rho_t$ and the cell entry $e_t$ from the
three displayed non-southwest corners.  Hence the same parser-visible group
data determine the padded left endpoint $\operatorname{pad}_{e_t}(\nu_t)$ and
the right endpoint $\mu_t$.

If the visible gate
\[
  \beta_t=\kappa(0)
  \qquad\text{and}\qquad
  \operatorname{pad}_{e_t}(\nu_t)=\mu_t
\]
holds, then the already recovered $\rho_t$ supplies the southwest corner
required in item~(3) of
\cref{prop:bc-visible-gated-midpoint-readouts-supply-grouped-padded-valleys}.
The anchor blocks are item~(1), the padded endpoints are recovered as above,
and the carrier-word relation is the displayed hypothesis.  Thus all
hypotheses of
\cref{prop:bc-visible-gated-midpoint-readouts-supply-grouped-padded-valleys}
hold, and that proposition gives the asserted zone recovery.
\end{proof}

\begin{corollary}[Step--3 boundary readouts supply visible-gated midpoint readouts]\label{cor:bc-step3-boundary-supplies-visible-gated-midpoint-readouts}
Work in the setting of
\cref{prop:bc-grouped-collective-carrier-splices-recover-positive-height-zones}.
Assume that the replacement window is parsed into ordered subwindows
\[
  W_1,\ldots,W_L
\]
with parser-visible tag sets $\Theta_\ell$ satisfying the tag-partition,
length, and concatenation hypotheses of that proposition.  Fix the injection
\[
  \kappa:\{0,1,2\}\hookrightarrow\{0,1\}^2
\]
from that setting.  Suppose that, for each $\ell$, the data
\[
  W_\ell,\quad \Theta_\ell,\quad S,\quad Z_\gamma,\quad
  \rho_{2a_\gamma-2},\rho_{2b_\gamma}
\]
determine:
\begin{enumerate}
\item the original block $B_r(P)$ for every
      $\operatorname{anc}(r)\in\Theta_\ell$;
\item a Ferrers subarrangement $E_\ell$ of the Step--3 dual-RSK' half-square
      of \cref{lem:bc-krattenthaler-semistandard-path-boundary}, together
      with the partition labels on the top/right boundary of $E_\ell$;
\item for every $\operatorname{car}(t)\in\Theta_\ell$, the inclusion of the
      one-cell carrier $E_{s_t,s_t+1}$ in $E_\ell$, where
      $s_t=j-u_t+1$.
\end{enumerate}
Finally suppose that the carrier words satisfy
\[
  \beta_t=\kappa(e_t)
  \qquad(\operatorname{car}(t)\in\Theta_\ell,\ 1\le\ell\le L),
\]
where $e_t$ is the Step--3 cell entry of $E_{s_t,s_t+1}$.  Then the
hypotheses of
\cref{prop:bc-visible-gated-midpoint-readouts-supply-grouped-padded-valleys}
hold.  Consequently the replacement window, $\beta_\gamma$, the zone endpoint
pair, and $Z_\gamma$ recover the original height halo-zone subpath.
\end{corollary}

\begin{proof}
Fix a group $\ell$.  By item~(2) and
\cref{cor:bc-parsed-dual-rsk-prime-boundary-recovers-subfilling}, the
parser-visible data recover every cell entry and every corner label on or
inside $E_\ell$.  In particular, for each carrier tag
\(\operatorname{car}(t)\in\Theta_\ell\), item~(3) identifies the one-cell
carrier \(E_{s_t,s_t+1}\) inside this recovered subfilling, so the data recover
its non-southwest corners
\[
  \nu_t,\quad \lambda_t,\quad \mu_t .
\]
The anchor blocks are recovered by item~(1).  Thus the hypotheses of
\cref{cor:bc-one-cell-corners-supply-visible-gated-midpoint-readouts} are
satisfied for the same parser-visible grouped subwindows and the same carrier
word relation.  Applying that corollary gives the visible-gated midpoint
readouts and hence the asserted zone recovery.
\end{proof}

\begin{corollary}[Visible-gated one-cell corners are enough]\label{cor:bc-visible-gated-one-cell-corners-suffice}
Work in the setting of
\cref{prop:bc-grouped-collective-carrier-splices-recover-positive-height-zones}.
Assume that the replacement window is parsed into ordered subwindows
\[
  W_1,\ldots,W_L
\]
with parser-visible tag sets $\Theta_\ell$ satisfying the tag-partition,
length, and concatenation hypotheses of that proposition.  Fix the injection
\[
  \kappa:\{0,1,2\}\hookrightarrow\{0,1\}^2
\]
from that setting.  Suppose that, for each $\ell$, the data
\[
  W_\ell,\quad \Theta_\ell,\quad S,\quad Z_\gamma,\quad
  \rho_{2a_\gamma-2},\rho_{2b_\gamma}
\]
determine:
\begin{enumerate}
\item the original block $B_r(P)$ for every
      $\operatorname{anc}(r)\in\Theta_\ell$;
\item for every $\operatorname{car}(t)\in\Theta_\ell$, the padded left
      endpoint $\operatorname{pad}_{e_t}(\nu_t)$ and the right endpoint
      $\mu_t$ of the one-cell Step--3 carrier $E_{s_t,s_t+1}$, where
      $s_t=j-u_t+1$;
\item for every carrier tag satisfying the visible gate
      \[
        \beta_t=\kappa(0)
        \qquad\text{and}\qquad
        \operatorname{pad}_{e_t}(\nu_t)=\mu_t,
      \]
      the three non-southwest corner labels
      \[
        \nu_t,\quad \lambda_t,\quad \mu_t
      \]
      of the same one-cell carrier.
\end{enumerate}
Finally suppose that the carrier words satisfy
\[
  \beta_t=\kappa(e_t)
  \qquad(\operatorname{car}(t)\in\Theta_\ell,\ 1\le\ell\le L).
\]
Then the hypotheses of
\cref{prop:bc-visible-gated-midpoint-readouts-supply-grouped-padded-valleys}
hold.  Consequently the replacement window, $\beta_\gamma$, the zone endpoint
pair, and $Z_\gamma$ recover the original height halo-zone subpath.
\end{corollary}

\begin{proof}
Items~(1) and (2) are exactly the anchor and endpoint readouts required in
\cref{prop:bc-visible-gated-midpoint-readouts-supply-grouped-padded-valleys}.
The displayed carrier-word relation is also the same.

It remains only to supply the southwest corner for the carrier tags whose
visible gate holds.  For such a tag, item~(3) gives the three non-southwest
corners of the one-cell carrier.  By
\cref{lem:bc-dual-rsk-prime-growth-local-reversibility}, the one-cell
backward rule recovers the southwest corner $\rho_t$ from those three corners.
Thus item~(3) of
\cref{prop:bc-visible-gated-midpoint-readouts-supply-grouped-padded-valleys}
is satisfied exactly on the visible-gated carrier tags, and that proposition
gives the claimed recovery.
\end{proof}

\begin{corollary}[Visible-gated Step--3 boundary readouts are enough]\label{cor:bc-visible-gated-step3-boundary-suffices}
Work in the setting of
\cref{prop:bc-grouped-collective-carrier-splices-recover-positive-height-zones}.
Assume that the replacement window is parsed into ordered subwindows
\[
  W_1,\ldots,W_L
\]
with parser-visible tag sets $\Theta_\ell$ satisfying the tag-partition,
length, and concatenation hypotheses of that proposition.  Fix the injection
\[
  \kappa:\{0,1,2\}\hookrightarrow\{0,1\}^2
\]
from that setting.  Suppose that, for each $\ell$, the data
\[
  W_\ell,\quad \Theta_\ell,\quad S,\quad Z_\gamma,\quad
  \rho_{2a_\gamma-2},\rho_{2b_\gamma}
\]
determine:
\begin{enumerate}
\item the original block $B_r(P)$ for every
      $\operatorname{anc}(r)\in\Theta_\ell$;
\item for every $\operatorname{car}(t)\in\Theta_\ell$, the padded left
      endpoint $\operatorname{pad}_{e_t}(\nu_t)$ and the right endpoint
      $\mu_t$ of the one-cell Step--3 carrier $E_{s_t,s_t+1}$, where
      $s_t=j-u_t+1$;
\item a Ferrers subarrangement $E_\ell^{\mathrm{vis}}$ of the Step--3
      dual-RSK' half-square of
      \cref{lem:bc-krattenthaler-semistandard-path-boundary}, together with
      the partition labels on the top/right boundary of
      $E_\ell^{\mathrm{vis}}$;
\item for every carrier tag satisfying the visible gate
      \[
        \beta_t=\kappa(0)
        \qquad\text{and}\qquad
        \operatorname{pad}_{e_t}(\nu_t)=\mu_t,
      \]
      the inclusion of the one-cell carrier $E_{s_t,s_t+1}$ in
      $E_\ell^{\mathrm{vis}}$.
\end{enumerate}
Finally suppose that the carrier words satisfy
\[
  \beta_t=\kappa(e_t)
  \qquad(\operatorname{car}(t)\in\Theta_\ell,\ 1\le\ell\le L).
\]
Then the hypotheses of
\cref{prop:bc-visible-gated-midpoint-readouts-supply-grouped-padded-valleys}
hold.  Consequently the replacement window, $\beta_\gamma$, the zone endpoint
pair, and $Z_\gamma$ recover the original height halo-zone subpath.
\end{corollary}

\begin{proof}
For a fixed group $\ell$, item~(3) and
\cref{cor:bc-parsed-dual-rsk-prime-boundary-recovers-subfilling} recover every
cell entry and every corner label on or inside
$E_\ell^{\mathrm{vis}}$.  Hence, for every carrier tag whose visible gate
holds, item~(4) identifies the one-cell carrier inside this recovered
subfilling and recovers its three non-southwest corners
\[
  \nu_t,\quad \lambda_t,\quad \mu_t .
\]
Together with items~(1) and (2), these are exactly the hypotheses of
\cref{cor:bc-visible-gated-one-cell-corners-suffice}.  Applying that corollary
gives the visible-gated midpoint readouts and hence the asserted zone recovery.
\end{proof}

\begin{corollary}[Raw Step--3 corner triples are not grouped target windows]\label{cor:bc-raw-step3-corners-not-grouped-target-windows}
The source condition in
\cref{cor:bc-step3-boundary-supplies-visible-gated-midpoint-readouts} is not
discharged by taking, for each carrier tag, the three raw non-southwest corner
labels of its one-cell Step--3 carrier as a singleton target subwindow of the
grouped replacement window.  Any construction that uses the Step--3 boundary
criterion must recover those corner labels from a content-normalized or
collective grouped target subwindow, rather than identify the raw corner triple
itself with one content-$2$ target block.
\end{corollary}

\begin{proof}
Fix a carrier tag and write the four corner labels of its one-cell carrier as
\[
  \begin{matrix}
  \nu & \lambda\\
  \rho & \mu
  \end{matrix}
\]
with $e$ the cell entry.  If the raw non-southwest triple
$(\nu,\lambda,\mu)$ were a singleton target subwindow in the grouped
replacement window of
\cref{prop:bc-grouped-collective-carrier-splices-recover-positive-height-zones},
then it would be a target two-step block in the content $(2^j)$ growth-path
model, hence a content-$2$ block by \cref{lem:bc-growth-path-model}.  But
\cref{cor:bc-one-cell-top-right-corners-not-target-blocks} gives
\[
  |\nu|-2|\lambda|+|\mu|=-2-e ,
\]
which is not $2$.  This contradicts the target-block weight required in the
grouped replacement window.  Hence the raw corner-triple identification cannot
supply the source condition.
\end{proof}

\begin{corollary}[One-cell Step--3 corners supply removable-row readouts]\label{cor:bc-one-cell-corners-supply-removable-row-readouts}
Work in the setting of
\cref{prop:bc-grouped-collective-carrier-splices-recover-positive-height-zones}.
Assume that the replacement window is parsed into ordered subwindows
\[
  W_1,\ldots,W_L
\]
with parser-visible tag sets $\Theta_\ell$ satisfying the tag-partition,
length, and concatenation hypotheses of that proposition.  Let
\[
  \kappa:\{0,1,2\}\hookrightarrow\{0,1\}^2
\]
be the fixed injection from that setting.  Suppose that, for each $\ell$, the
data
\[
  W_\ell,\quad \Theta_\ell,\quad S,\quad Z_\gamma,\quad
  \rho_{2a_\gamma-2},\rho_{2b_\gamma}
\]
determine:
\begin{enumerate}
\item the original block $B_r(P)$ for every
      $\operatorname{anc}(r)\in\Theta_\ell$;
\item for every $\operatorname{car}(t)\in\Theta_\ell$, the three non-southwest
      corner labels
      \[
        \nu_t,\quad \lambda_t,\quad \mu_t
      \]
      of the one-cell Step--3 carrier $E_{s_t,s_t+1}$, where
      $s_t=j-u_t+1$.
\end{enumerate}
Finally suppose that the carrier words satisfy
\[
  \beta_t=\kappa(e_t)
  \qquad(\operatorname{car}(t)\in\Theta_\ell,\ 1\le\ell\le L),
\]
where $e_t$ is the Step--3 cell entry of $E_{s_t,s_t+1}$.  Then the
replacement window, $\beta_\gamma$, the zone endpoint pair, and $Z_\gamma$
recover the original height halo-zone subpath through the removable-row readout
criterion of \cref{prop:bc-removable-row-readouts-supply-grouped-padded-valleys}.
\end{corollary}

\begin{proof}
Fix a carrier tag $\operatorname{car}(t)\in\Theta_\ell$.  By
\cref{lem:bc-dual-rsk-prime-growth-local-reversibility}, the one-cell backward
rule recovers the southwest corner $\rho_t$ and the cell entry $e_t$ from the
three displayed corner labels.  From $e_t$ and $\nu_t$ the padded left endpoint
$\operatorname{pad}_{e_t}(\nu_t)$ is determined, while $\mu_t$ is already one of
the displayed labels.

If $e_t=0$ and $\nu_t=\mu_t$, then
\cref{cor:bc-singleton-defect-row-from-northeast-box} says that
$\lambda_t/\nu_t$ consists of one box, and its row is the successor of the
singleton defect row.  Equivalently, by
\cref{cor:bc-exceptional-northeast-row-removable-normal-form}, the predecessor
row is a removable-corner row of $\nu_t$.  Since $\lambda_t$ and $\nu_t$ are
known, the row of the unique box in $\lambda_t/\nu_t$ is known, and hence so is
the removable row required in item~(3) of
\cref{prop:bc-removable-row-readouts-supply-grouped-padded-valleys}.

Thus the hypotheses of
\cref{prop:bc-removable-row-readouts-supply-grouped-padded-valleys} are satisfied:
the anchor blocks are item~(1), the padded carrier endpoints are recovered from
item~(2) and $e_t$, and the exceptional removable rows are recovered as above.
That proposition recovers the original height halo-zone subpath.
\end{proof}

\begin{proposition}[Pair-endpoint padded carriers give endpoint-compatible placement]\label{prop:bc-pair-endpoint-padded-carriers-give-placement}
Work in the setting of
\cref{prop:bc-endpoint-compatible-padded-carriers-recover-positive-height-zones}.
For each deleted singleton $u_t$, with $s_t=j-u_t+1$, let
\[
  \begin{matrix}
  \nu_t & \lambda_t\\
  \rho_t & \mu_t
  \end{matrix}
\]
be the corner labels of $E_{s_t,s_t+1}$ and let $e_t$ be its cell entry.
Assume that the padded carrier has the pair endpoints
\[
  \operatorname{pad}_{e_t}(\nu_t)=\rho_{2u_t-4},
  \qquad
  \mu_t=\rho_{2u_t}.
\]
If the two-bit split of the branch word satisfies
$\beta_t=\kappa(e_t)$ for every $t$, then the original height halo-zone
subpath is recovered from the resulting replacement window, $\beta_\gamma$,
the zone endpoint pair, and $Z_\gamma$.
\end{proposition}

\begin{proof}
Define the candidate blocks $C_r$ as in
\cref{prop:bc-endpoint-compatible-padded-carriers-recover-positive-height-zones}.
For anchors $r\in A_\gamma$, one has $C_r=B_r(P)$, so its endpoints are
\[
  \rho_{2r-2},\qquad \rho_{2r}.
\]
For the paired slot $r=u_t-1$, the two displayed hypotheses say that
$C_r$ has endpoints
\[
  \rho_{2u_t-4},\qquad \rho_{2u_t},
\]
namely the endpoints of the original two-block interval
$\{u_t-1,u_t\}$.

By \cref{cor:bc-height-anchor-reduced-left-pairs}, the zone indices are
partitioned into the anchor singleton blocks $r\in A_\gamma$ and the adjacent
pairs $\{u_t-1,u_t\}$.  Therefore the candidate blocks, ordered by their
retained slots $r\in Z_\gamma\setminus I$, cover consecutive original index
intervals whose union is $Z_\gamma$.  Adjacent candidate blocks have matching
endpoints because the right endpoint of one covered interval is the left
endpoint of the next covered interval in the original path.  The first and
last candidate endpoints are exactly the zone endpoints
$\rho_{2a_\gamma-2}$ and $\rho_{2b_\gamma}$.

Thus the endpoint equalities in
\cref{lem:bc-endpoint-compatible-blocks-form-replacement-window} hold.
\Cref{prop:bc-endpoint-compatible-padded-carriers-recover-positive-height-zones}
then gives the claimed recovery of the original height halo-zone subpath.
\end{proof}

\begin{corollary}[Pair-endpoint padding is not source-automatic]\label{cor:bc-pair-endpoint-padding-not-source-automatic}
The pair-endpoint identities assumed in
\cref{prop:bc-pair-endpoint-padded-carriers-give-placement} do not follow
from the Step--3 one-cell boundary convention alone.  In particular, the
bottom-padded carrier valley cannot be declared to span the adjacent original
two-block interval without an additional endpoint splice or an additional
restriction beyond the ambient Step--3 construction.
\end{corollary}

\begin{proof}
Specialize the Step--3 construction to content $(2^2)$.  The Step--2 matrix
has zero diagonal and row sums equal to $2$ by
\cref{lem:bc-step2-matrix-doubled-content-degrees}, so its unique
off-diagonal entry is $2$.  By
\cref{lem:bc-step3-cells-are-step2-half-matrix}, the unique Step--3 carrier
$E_{1,2}$ therefore has cell entry $e=2$.

The full Step--3 left and bottom boundaries are labelled by empty partitions
in Krattenthaler's construction, as used in
\cref{lem:bc-step3-full-half-square-determines-rho}.  Hence the one-cell
carrier corners before applying the forward rule satisfy
\[
  \rho=\mu=\nu=\varnothing .
\]
Running the carry formula in the proof of
\cref{lem:bc-dual-rsk-prime-cell-size-conservation} with $m=2$ and these
empty boundary partitions gives $b_1=b_2=1$ and $b_i=0$ for $i\ge3$, hence
the northeast corner is
\[
  \lambda=(1,1).
\]
Thus
\[
  \operatorname{pad}_e(\nu)=\operatorname{pad}_2(\varnothing)=(1,1).
\]

On the other hand, \cref{lem:bc-growth-path-model} gives
\[
  \rho_0=\rho_4=\varnothing
\]
for the full $m=0$, content-$(2^2)$ boundary path.  Under the
label-to-block convention, the adjacent label pair $(2,1)$ is the full
two-block interval, so the left endpoint required by the identity in
\cref{prop:bc-pair-endpoint-padded-carriers-give-placement} is
$\rho_0=\varnothing$.  Therefore
\[
  \operatorname{pad}_e(\nu)=(1,1)\ne\varnothing=\rho_0 .
\]
This disproves only the source-automatic discharge of the pair-endpoint
hypothesis; the conditional placement implication of
\cref{prop:bc-pair-endpoint-padded-carriers-give-placement} remains valid.
\end{proof}

\begin{corollary}[Positive height zones do not force zero carrier padding]\label{cor:bc-positive-height-nonzero-padding-example}
The positive-height halo hypotheses do not force the two-label Step--3 carrier
entry $e_t$ to vanish.  Moreover, nonzero carrier entries in one
positive-height zone need not inject into the surplus-anchor set
$A_\gamma$.  Consequently, the replacement positive-height placement cannot
be completed by proving that all positive-height padded carrier windows have
zero padding, nor by assigning an independent surplus-anchor buffer to each
nonzero padded carrier.
\end{corollary}

\begin{proof}
Consider the semistandard tableau of content $(2^5)$
\[
T=
\begin{array}{cccc}
1&1&2&4\\
2&3&3&5\\
4\\
5
\end{array}.
\]
Its column lengths are $4,2,2,2$, so the shape is column-even.  For $q=2$ the
height witness set is
\[
  H_2(T)=\{T(1,1),T(3,1)\}=\{1,4\}.
\]
Under the block convention $s\mapsto j-s+1$, the deleted block set is
\[
  I(H_2(T))=\{5,2\}.
\]
The one-block halo is the single zone $Z=\{1,2,3,4,5\}$, with three retained
indices and two deleted indices, hence this is a positive-height zone.

The first-variant Knuth matrix associated to this tableau is
\[
M=
\begin{pmatrix}
0&1&1&0&0\\
1&0&0&0&1\\
1&0&0&1&0\\
0&0&1&0&1\\
0&1&0&1&0
\end{pmatrix}.
\]
It is symmetric, has zero diagonal, and every row sum is $2$.  The row-word
insertion check for the nonzero biletters
\[
(1,2),(1,3),(2,1),(2,5),(3,1),(3,4),
(4,3),(4,5),(5,2),(5,4)
\]
gives insertion and recording tableaux both equal to the displayed $T$.
Thus this is a valid Step--2 datum for the source construction of
\cref{lem:bc-step2-matrix-doubled-content-degrees}.

For the positive-zone deleted witness label $s=4$, the retained successor is
$s+1=5$.  By \cref{lem:bc-step3-cells-are-step2-half-matrix}, the unique
Step--3 cell entry of $E_{4,5}$ is the corresponding off-diagonal matrix
entry
\[
  e=M_{4,5}=1 .
\]
Thus a positive-height zone can have a genuinely nonzero padding parameter.

The same zone has deleted block list $u_1=2,u_2=5$.  Its surplus-anchor set
from \cref{cor:bc-height-surplus-canonical-right-anchors} is
\[
  A_\gamma=\{u_1+1\}=\{3\},
\]
because $u_2-u_1=3$ and $u_2=j$.  The other deleted witness label is $s=1$,
whose retained successor is $2$, and the same matrix gives
\[
  M_{1,2}=1 .
\]
Therefore both deleted singletons in this positive-height zone have nonzero
adjacent carrier entry, while the zone has only one surplus anchor.  No
splice rule which requires a distinct surplus-anchor buffer for each nonzero
carrier can cover this valid positive-height example.
\end{proof}

\begin{corollary}[Positive height zones may have double successor carriers]\label{cor:bc-positive-height-double-successor-carrier-example}
The positive-height carrier splice cannot be reduced to the case
$e_t\in\{0,1\}$, nor to smoothing each deleted witness through two distinct
retained graph neighbours.  There is a column-even content-$(2^7)$ tableau
\[
T=
\begin{array}{cccc}
1&1&2&2\\
3&3&4&7\\
4&5\\
5&6\\
6\\
7
\end{array}
\]
whose height witness set for $q=2$ is $H_2(T)=\{1,4\}$, whose canonical halo
zone
\[
  Z=\{3,4,5,6,7\}
\]
has positive surplus, and whose deleted witness label $s=4$ has adjacent
Step--3 carrier entry
\[
  e=M_{4,5}=2 .
\]
\end{corollary}

\begin{proof}
The rows of $T$ are weakly increasing and the columns are strictly
increasing.  Its row lengths are
\[
  (4,4,2,2,1,1),
\]
so its column lengths are $(6,4,2,2)$, all even.  Each label
$1,\ldots,7$ appears exactly twice.  The first-column entries in rows $1$ and
$3$ are $1$ and $4$, so $H_2(T)=\{1,4\}$.

The first-variant Knuth matrix associated to this tableau is
\[
M=
\begin{pmatrix}
0&0&1&0&0&1&0\\
0&0&0&0&0&0&2\\
1&0&0&0&0&1&0\\
0&0&0&0&2&0&0\\
0&0&0&2&0&0&0\\
1&0&1&0&0&0&0\\
0&2&0&0&0&0&0
\end{pmatrix}.
\]
It is symmetric, has zero diagonal, and every row sum is $2$.  The row-word
insertion check for the nonzero biletters
\[
\begin{gathered}
(1,3),(1,6),(2,7),(2,7),(3,1),(3,6),(4,5),(4,5),\\
(5,4),(5,4),(6,1),(6,3),(7,2),(7,2)
\end{gathered}
\]
gives insertion and recording tableaux both equal to the displayed $T$.
Thus this is a valid Step--2 datum for the source construction of
\cref{lem:bc-step2-matrix-doubled-content-degrees}.

Under the block convention $s\mapsto j-s+1$, the witness set
$\{1,4\}$ gives deleted block set
\[
  I=\{7,4\}.
\]
The one-block halo of $I$ is the single zone
\[
  Z=\{3,4,5,6,7\}.
\]
It contains two deleted indices and three retained indices, so the zone has
positive surplus.  For the deleted block $u=4$, the corresponding witness
label is $s=7-u+1=4$, and the retained successor label is $5$.  By
\cref{lem:bc-step3-cells-are-step2-half-matrix}, the unique Step--3 cell
entry of the adjacent two-label carrier $E_{4,5}$ is the off-diagonal matrix
entry $M_{4,5}=2$.  Hence an actual positive-height zone can require the full
two-unit padding case.
\end{proof}

\begin{corollary}[Collar-only normalization is not universal]\label{cor:bc-collar-only-normalization-obstruction}
In the ambient growth-path category, the merged positive-height window in
\cref{prop:bc-coupled-collar-carriers-recover-positive-height-zones} cannot
be supplied by leaving the retained left collar unchanged and then using only
the retained right collar, the outer endpoint pair, the known singleton
position, and a two-bit word to recover the deleted singleton block.
Equivalently, the actual merged window must carry deleted-sensitive target
data, such as a content-normalized carrier-boundary readout, or use further
fixed-height restrictions beyond the ambient two-step axioms.
\end{corollary}

\begin{proof}
Let
\[
  \mu=(5,4,3,2,1),
\]
and fix two removable corners $c_L$ and $c_R$ of $\mu$.  For every removable
corner $c$ of $\mu$, form the three-block segment
\[
  U^0=\mu,\quad U^1=\mu\setminus\{c_L\},\quad U^2=\mu,
\]
\[
  U^3=\mu\setminus\{c\},\quad U^4=\mu,
\]
\[
  U^5=\mu\setminus\{c_R\},\quad U^6=\mu .
\]
Each consecutive triple
\[
  (U^0,U^1,U^2),\qquad (U^2,U^3,U^4),\qquad (U^4,U^5,U^6)
\]
has vertical-strip adjacent differences and content weight $2$, because each
middle partition is obtained from $\mu$ by deleting one removable corner.
Thus these are valid ambient left-collar, deleted-singleton, and right-collar
blocks.

The staircase partition $\mu$ has five removable corners, so the choice of
$c$ gives five distinct deleted singleton blocks.  The retained left collar,
the retained right collar, the outer endpoints $(U^0,U^6)=(\mu,\mu)$, and the
known middle position are the same for all five choices.  A two-bit word has
only four values.  Therefore no decoder using only those data and a two-bit
word can recover the middle block for all five segments.
\end{proof}

\begin{corollary}[Right collars are parser-local after left-collar readout]\label{cor:bc-right-collars-parser-local-after-left-readout}
Work in the height setting of
\cref{prop:bc-height-anchor-pair-zone-recovery}.  Suppose that the parser
decomposes $V^\ast_\gamma$ into unchanged anchor blocks indexed by
$A_\gamma$ and merged windows $M_t$ indexed by the increasing deleted
singletons
\[
  Z_\gamma\cap I=\{u_1<\cdots<u_m\}.
\]
Assume that each $M_t$ determines the original retained left-collar block
\[
  B_{u_t-1}(P)=(\rho_{2u_t-4},\rho_{2u_t-3},\rho_{2u_t-2}).
\]
Then, for each $t$, the full parsed zone data determine the retained
right-collar block $B_{u_t+1}(P)$ whenever $u_t<j$; if $u_t=j$, the right
side is the endpoint $\rho_{2j}$.
\end{corollary}

\begin{proof}
By \cref{cor:bc-height-anchor-pair-retained-collars}, if $u_t<j$ then
$u_t+1$ is a retained index in $Z_\gamma$.  If $u_t+1\in A_\gamma$, the
parser returns $B_{u_t+1}(P)$ as an unchanged anchor block.  If
$u_t+1\notin A_\gamma$, then
\cref{cor:bc-height-anchor-reduced-left-pairs} gives
\[
  u_t+1=u_k-1
\]
for a unique $k$.  Since the deleted singleton list is increasing, this
forces $k=t+1$ and $u_{t+1}=u_t+2$.  The assumed left-collar readout from
$M_{t+1}$ therefore determines
\[
  B_{u_{t+1}-1}(P)=B_{u_t+1}(P).
\]
Thus the right collar is determined in every nonterminal case.  If $u_t=j$,
there is no right-collar block and the right endpoint is the global terminal
partition $\rho_{2j}$, which is part of the endpoint data for the terminal
zone.
\end{proof}

\begin{proposition}[Coupled collar carriers recover positive height zones]\label{prop:bc-coupled-collar-carriers-recover-positive-height-zones}
Work in the height setting of
\cref{prop:bc-height-anchor-pair-zone-recovery}.  Suppose the parser
decomposes $V^\ast_\gamma$ into unchanged anchor blocks and merged windows
$M_1,\ldots,M_m$ for the increasing deleted singleton list
\[
  Z_\gamma\cap I=\{u_1<\cdots<u_m\}.
\]
Fix an injection
\[
  \kappa:\{0,1,2\}\hookrightarrow\{0,1\}^2 .
\]
Assume the following.
\begin{enumerate}
\item Each $M_t$ determines the retained left-collar block $B_{u_t-1}(P)$.
\item After the right collars are recovered by
\cref{cor:bc-right-collars-parser-local-after-left-readout}, the parsed zone
data
\[
  S,\quad Z_\gamma,\quad V^\ast_\gamma,\quad
  \rho_{2a_\gamma-2},\rho_{2b_\gamma},\quad t
\]
determine the left/bottom boundary of the two-label carrier
$E_{s_t,s_t+1}$, where $s_t=j-u_t+1$.
\item If $e_t$ is the unique cell entry of $E_{s_t,s_t+1}$, then the
deterministic two-bit split of $\beta_\gamma$ satisfies
$\beta_t=\kappa(e_t)$ for every $t$.
\end{enumerate}
Then the original height halo-zone subpath is recovered from
$V^\ast_\gamma$, $\beta_\gamma$, the zone endpoint pair, and $Z_\gamma$.
\end{proposition}

\begin{proof}
For each $t$, item~(2) determines the left/bottom boundary of
$E_{s_t,s_t+1}$.  By
\cref{lem:bc-two-label-step3-cell-entry-three-valued}, its unique cell entry
lies in $\{0,1,2\}$, and item~(3) recovers that entry from $\beta_t$.  Hence
\cref{cor:bc-step3-consecutive-carrier-data-recovers-boundary}, applied with
$a=s_t$ and $b=s_t+1$, recovers the top/right boundary segment
\[
  \rho_{2u_t-4},\rho_{2u_t-3},\rho_{2u_t-2},
  \rho_{2u_t-1},\rho_{2u_t}.
\]
This is the original two-block segment on block indices $u_t-1,u_t$, by
\cref{cor:bc-height-anchor-pair-step3-boundary-segment}.

The anchor blocks indexed by $A_\gamma$ are recovered unchanged by the
parser.  By \cref{cor:bc-height-anchor-reduced-left-pairs}, the blocks in
$Z_\gamma$ are partitioned into those anchor singleton blocks and the pairs
$\{u_t-1,u_t\}$.  Interleaving the recovered anchor blocks and recovered
two-block segments in increasing block order reconstructs the whole subpath
\[
  \rho_{2a_\gamma-2},\rho_{2a_\gamma-1},\ldots,\rho_{2b_\gamma}.
\]
\end{proof}

\begin{corollary}[Height halo zones have singleton nonnegative packets]\label{cor:bc-height-halo-singleton-nonnegative}
Let $T$ be a semistandard tableau of content $(2^j)$ with
$\ell(\lambda(T))\ge2q$, set $S=H_q(T)$, put
$I=I(S)=\{j-s+1:s\in S\}$, and assume $j\ge2q$.  Let
$Z_\gamma=\{a_\gamma,\ldots,b_\gamma\}$ be one canonical halo zone of
\cref{lem:bc-canonical-halo-replacement-zones}, and write the maximal
consecutive components of $Z_\gamma\cap I$ as
\[
  K_t=\{u_t,\ldots,v_t\}\qquad(1\le t\le m).
\]
Then every $K_t$ is a singleton.  Moreover the local balance
\[
  \Delta_\gamma=|Z_\gamma\setminus I|-|Z_\gamma\cap I|
\]
is nonnegative.
\end{corollary}

\begin{proof}
If some $K_t$ had more than one element, then it would contain two consecutive
block indices $r$ and $r+1$.  The corresponding tableau labels
$j-r+1$ and $j-r$ would then be two consecutive elements of $S=H_q(T)$,
contradicting \cref{lem:bc-height-witness-labels-separated}.  Thus every
$K_t$ is a singleton.

Use the notation of \cref{lem:bc-canonical-halo-zone-balance}.  Since all
components are singletons, $h_t=1$ for all $t$.  If
$\epsilon_L+\epsilon_R\ge1$, then every internal gap has $g_t\ge1$, so
\[
  \Delta_\gamma
  =
  \epsilon_L+\epsilon_R+\sum_{t=1}^{m-1}g_t-m
  \ge
  \epsilon_L+\epsilon_R+(m-1)-m
  \ge0.
\]
It remains to consider $\epsilon_L=\epsilon_R=0$.  Then the first deleted
block in the zone is $1$ and the last deleted block in the zone is $j$, so
this zone contains all elements of $I$.  Hence $|Z_\gamma\cap I|=q$ and
$Z_\gamma=\{1,\ldots,j\}$.  Therefore
\[
  \Delta_\gamma=(j-q)-q=j-2q\ge0
\]
by the hypothesis $j\ge2q$.
\end{proof}

\begin{corollary}[Height zero-surplus zones are right-boundary alternating]\label{cor:bc-height-zero-surplus-zones}
In the setting of
\cref{cor:bc-height-halo-singleton-nonnegative}, write the singleton
components of $Z_\gamma\cap I$ as
\[
  K_t=\{u_t\}\qquad(1\le t\le m),
  \qquad u_1<\cdots<u_m .
\]
Then $u_1>1$.  Moreover, with
\[
  \epsilon_R=
  \begin{cases}
  1,&u_m<j,\\
  0,&u_m=j,
  \end{cases}
\]
one has the exact surplus formula
\[
  |Z_\gamma\setminus I|-|Z_\gamma\cap I|
  =
  \epsilon_R+\sum_{t=1}^{m-1}(u_{t+1}-u_t-2).
\]
Consequently the height zone has zero surplus if and only if $u_m=j$ and
\[
  u_{t+1}=u_t+2\qquad(1\le t<m).
\]
\end{corollary}

\begin{proof}
Let $t_q=T(2q-1,1)$ be the largest height witness label.  Since
$\ell(\lambda(T))\ge2q$, the first-column entry $T(2q,1)$ exists.  Column
strictness gives
\[
  t_q<T(2q,1)\le j,
\]
so $t_q<j$.  Every label in $S=H_q(T)$ is at most $t_q$, and therefore every
block index $j-s+1$ in $I$ is at least $j-t_q+1>1$.  In particular $u_1>1$.

By \cref{cor:bc-height-halo-singleton-nonnegative}, all deleted components in
the zone are singletons.  Thus, in the notation of
\cref{lem:bc-canonical-halo-zone-balance}, one has $h_t=1$ for all $t$ and
$\epsilon_L=1$ because $u_1>1$.  Also
$g_t=u_{t+1}-u_t-1$.  Substituting these values into the balance formula gives
\[
  |Z_\gamma\setminus I|-|Z_\gamma\cap I|
  =
  1+\epsilon_R+\sum_{t=1}^{m-1}(u_{t+1}-u_t-1)-m
  =
  \epsilon_R+\sum_{t=1}^{m-1}(u_{t+1}-u_t-2).
\]
The separated-label result
\cref{lem:bc-height-witness-labels-separated} gives
$u_{t+1}\ge u_t+2$, so every summand is nonnegative.  Hence the surplus is
zero exactly when $\epsilon_R=0$ and every summand vanishes, which is the
displayed right-boundary alternating condition.
\end{proof}

\begin{corollary}[Height zero-surplus zones are alternating label prefixes]\label{cor:bc-height-zero-surplus-label-prefix}
In the setting of
\cref{cor:bc-height-zero-surplus-zones}, suppose the height zone has zero
surplus and put $m=|Z_\gamma\cap I|$.  Then
\[
  Z_\gamma=\{j-2m+1,\ldots,j\},
\]
\[
  Z_\gamma\cap I=\{j-2m+2,\ j-2m+4,\ldots,j\},
  \qquad
  Z_\gamma\setminus I=\{j-2m+1,\ j-2m+3,\ldots,j-1\}.
\]
Under the label-to-block correspondence $r\mapsto j-r+1$, the deleted labels
in this zone are exactly
\[
  \{1,3,\ldots,2m-1\},
\]
and the retained labels in this zone are exactly
\[
  \{2,4,\ldots,2m\}.
\]
Thus every tight residual height zone is the terminal-block encoding of an
odd alternating label prefix.
\end{corollary}

\begin{proof}
By \cref{cor:bc-height-zero-surplus-zones}, zero surplus is equivalent to
$u_m=j$ and $u_{t+1}=u_t+2$ for $1\le t<m$.  Hence
\[
  u_t=j-2(m-t)\qquad(1\le t\le m).
\]
The deleted block set in the zone is therefore
\[
  Z_\gamma\cap I=\{u_1,\ldots,u_m\}
  =
  \{j-2m+2,\ j-2m+4,\ldots,j\}.
\]
Since $u_1>1$ and $u_m=j$, the one-block halo component has left endpoint
$u_1-1=j-2m+1$ and right endpoint $j$, proving the formula for $Z_\gamma$.
The retained blocks are the complementary alternating indices in this interval,
namely
\[
  Z_\gamma\setminus I=\{j-2m+1,\ j-2m+3,\ldots,j-1\}.
\]
Finally, a block index $r$ corresponds to label $j-r+1$.  Applying this map
to the displayed deleted and retained block sets gives the two displayed
label sets.
\end{proof}

\begin{corollary}[Positive height halo zones contain a retained anchor]\label{cor:bc-height-positive-surplus-retained-anchor}
In the setting of
\cref{cor:bc-height-zero-surplus-zones}, suppose
\[
  |Z_\gamma\setminus I|-|Z_\gamma\cap I|>0 .
\]
Then there is an index
\[
  r\in Z_\gamma\setminus I
\]
which is adjacent to exactly one deleted singleton block in $Z_\gamma\cap I$:
exactly one of $r-1$ and $r+1$ belongs to $Z_\gamma\cap I$.
More precisely, either $u_m<j$ and one may take $r=u_m+1$, or there is an
index $t<m$ with $u_{t+1}=u_t+3$ and one may take $r=u_t+1$.
\end{corollary}

\begin{proof}
By \cref{cor:bc-height-zero-surplus-zones},
\[
  |Z_\gamma\setminus I|-|Z_\gamma\cap I|
  =
  \epsilon_R+\sum_{t=1}^{m-1}(u_{t+1}-u_t-2),
\]
and every summand on the right is nonnegative.  If the left side is positive,
then either $\epsilon_R=1$ or some summand is positive.

If $\epsilon_R=1$, then $u_m<j$ and the right halo index $r=u_m+1$ lies in
$Z_\gamma\setminus I$.  It is adjacent to the deleted singleton $u_m$, and it
is not adjacent to any later deleted singleton because $u_m$ is the last
deleted index in the zone.

Otherwise, for some $t<m$ one has $u_{t+1}-u_t-2>0$.  Since
\cref{lem:bc-canonical-halo-short-gaps} gives
$u_{t+1}-u_t\in\{2,3\}$ for singleton components in one halo zone, this means
$u_{t+1}=u_t+3$.  The index $r=u_t+1$ is a retained index in the gap between
the two deleted singletons.  It is adjacent to $u_t$ and has distance two from
$u_{t+1}$; all other deleted singletons are farther away in block order.
Thus it is adjacent to exactly one deleted singleton in $Z_\gamma\cap I$.
\end{proof}

\begin{corollary}[Height surplus is carried by canonical right anchors]\label{cor:bc-height-surplus-canonical-right-anchors}
In the setting of
\cref{cor:bc-height-zero-surplus-zones}, define
\[
  A_\gamma
  =
  \{u_m+1:\ u_m<j\}
  \cup
  \{u_t+1:\ 1\le t<m,\ u_{t+1}=u_t+3\}.
\]
Then
\[
  A_\gamma\subseteq Z_\gamma\setminus I,
\]
every element of $A_\gamma$ is adjacent to exactly one deleted singleton in
$Z_\gamma\cap I$, and
\[
  |A_\gamma|
  =
  |Z_\gamma\setminus I|-|Z_\gamma\cap I|.
\]
Thus the positive surplus of a height halo zone is represented by a canonical
ordered list of one-sided retained anchors.
\end{corollary}

\begin{proof}
If $u_m<j$, then the right endpoint of the one-block halo component is
$u_m+1$, so $u_m+1\in Z_\gamma$.  It is not in $I$, since otherwise it would
be consecutive to the deleted singleton $u_m$.  It is adjacent to $u_m$ and
to no later deleted singleton.

If $u_{t+1}=u_t+3$, then the index $u_t+1$ lies in the retained gap between
the two deleted singletons $u_t$ and $u_{t+1}$.  It belongs to
$Z_\gamma\setminus I$, is adjacent to $u_t$, and has distance two from
$u_{t+1}$; the remaining deleted singletons are farther away.  Hence every
element of $A_\gamma$ is a retained index adjacent to exactly one deleted
singleton.  The displayed set is canonical because it is defined from the
increasing list of deleted singleton indices in the canonical halo zone.

It remains to count it.  The first set in the union contributes
$\epsilon_R$, and the second contributes one element exactly for each gap
with $u_{t+1}-u_t=3$.  By \cref{lem:bc-canonical-halo-short-gaps}, all gaps
between singleton components inside a canonical halo zone have size $2$ or
$3$.  Therefore
\[
  |A_\gamma|
  =
  \epsilon_R+
  \sum_{t=1}^{m-1}(u_{t+1}-u_t-2).
\]
The right side is the surplus
$|Z_\gamma\setminus I|-|Z_\gamma\cap I|$ by
\cref{cor:bc-height-zero-surplus-zones}.
\end{proof}

\begin{corollary}[Removing height surplus anchors leaves left-neighbour pairs]\label{cor:bc-height-anchor-reduced-left-pairs}
In the setting of
\cref{cor:bc-height-surplus-canonical-right-anchors}, one has
\[
  (Z_\gamma\setminus I)\setminus A_\gamma
  =
  \{u_t-1:\ 1\le t\le m\}.
\]
Consequently, after deleting the canonical surplus anchors, the retained
indices in the height zone are in canonical bijection with the deleted
singletons by
\[
  u_t\longmapsto u_t-1 .
\]
\end{corollary}

\begin{proof}
By \cref{cor:bc-height-zero-surplus-zones}, one has $u_1>1$, so $u_1-1$ is
the retained prefix index of the zone.  This is the element $u_t-1$ for
$t=1$.

Between consecutive deleted singletons $u_t<u_{t+1}$, the short-gap lemma
\cref{lem:bc-canonical-halo-short-gaps} gives
$u_{t+1}-u_t\in\{2,3\}$.  If $u_{t+1}=u_t+2$, the only retained index in the
gap is
\[
  u_t+1=u_{t+1}-1,
\]
and it is not in $A_\gamma$.  If $u_{t+1}=u_t+3$, the retained gap indices
are $u_t+1$ and $u_t+2$; the first belongs to $A_\gamma$ by definition, while
the second is $u_{t+1}-1$ and is not in $A_\gamma$.

Finally, if $u_m<j$, the retained suffix index $u_m+1$ belongs to
$A_\gamma$; if $u_m=j$, there is no suffix retained index.  Thus the retained
indices not removed by $A_\gamma$ are exactly the indices
$u_t-1$ for $1\le t\le m$.  These indices are distinct and each lies
immediately to the left of the corresponding deleted singleton $u_t$, giving
the stated bijection.
\end{proof}

\begin{corollary}[Only the terminal height halo zone can have zero surplus]\label{cor:bc-height-zero-surplus-terminal-unique}
In the setting of
\cref{cor:bc-height-halo-singleton-nonnegative}, order the canonical halo
zones increasingly by their block intervals.  At most one height halo zone has
zero surplus
\[
  |Z_\gamma\setminus I|-|Z_\gamma\cap I|=0 .
\]
If such a zone exists, it is the last zone, its right endpoint is $j$, and it
has the alternating prefix form described in
\cref{cor:bc-height-zero-surplus-label-prefix}.  Every height halo zone
strictly to its left has positive surplus.
\end{corollary}

\begin{proof}
Let $Z_\gamma$ be a zero-surplus height halo zone, and write its singleton
deleted components as in \cref{cor:bc-height-zero-surplus-zones}.  That
corollary gives $u_m=j$.  Therefore the canonical halo component containing
this last deleted block has right endpoint $j$.  Since the canonical halo
zones are disjoint intervals ordered increasingly, a zone with right endpoint
$j$ is the last zone, and no other zone can also contain the block $j$.
Hence there is at most one zero-surplus zone, and if it exists it is terminal.

\Cref{cor:bc-height-zero-surplus-label-prefix} gives the displayed
alternating prefix form for this terminal zone.  Finally,
\cref{cor:bc-height-halo-singleton-nonnegative} says that every height halo
zone has nonnegative surplus.  Any zone strictly to the left of the terminal
zero-surplus zone is not itself zero-surplus by uniqueness, and hence has
positive surplus.
\end{proof}

\begin{corollary}[Nonterminal height halo zones have retained anchors]\label{cor:bc-nonterminal-height-zones-have-anchors}
In the setting of
\cref{cor:bc-height-halo-singleton-nonnegative}, order the canonical halo
zones increasingly by their block intervals.  Every height halo zone except
possibly the unique terminal zero-surplus zone of
\cref{cor:bc-height-zero-surplus-terminal-unique} contains a retained anchor
as in \cref{cor:bc-height-positive-surplus-retained-anchor}.
\end{corollary}

\begin{proof}
By \cref{cor:bc-height-zero-surplus-terminal-unique}, every height halo zone
which is not the unique terminal zero-surplus zone has positive surplus.  The
anchor assertion is then exactly
\cref{cor:bc-height-positive-surplus-retained-anchor}.
\end{proof}

\begin{corollary}[Non-full tight height prefixes have a visible separator]\label{cor:bc-tight-height-visible-separator}
In the setting of
\cref{cor:bc-height-zero-surplus-label-prefix}, suppose $j>2m$.  Then
\[
  2m+1\notin S,
  \qquad
  2m+2\notin S\quad\text{if }2m+2\le j.
\]
Moreover the block index $j-2m$, corresponding to label $2m+1$, lies outside
the full one-block halo $H(S)$.  Hence it belongs to the unchanged complement
$R$ of \cref{lem:bc-canonical-halo-replacement-zones}, and its right endpoint
is the left endpoint of the tight zone:
\[
  \rho_{2(j-2m)}
  =
  \rho_{2(j-2m+1)-2}.
\]
\end{corollary}

\begin{proof}
By \cref{cor:bc-height-zero-surplus-label-prefix}, the tight zone is
\[
  Z_\gamma=\{j-2m+1,\ldots,j\}.
\]
If $2m+1\in S$, then the corresponding deleted block is $j-2m$.  Its
one-block halo contains $j-2m+1$, so the maximal halo component containing
$Z_\gamma$ would extend left to include $j-2m$, contradicting the displayed
left endpoint of $Z_\gamma$.

If $2m+2\in S$, then the corresponding deleted block is $j-2m-1$.  Its
one-block halo contains $j-2m$, and the tight zone contains $j-2m+1$.
These two halo pieces are consecutive, so again they would form one larger
halo component, contradicting the displayed left endpoint of $Z_\gamma$.

Now let $s\in S$.  If $s\le2m-1$, then its deleted block belongs to the tight
zone and its halo is contained in $Z_\gamma$.  If $s\ge2m+3$, then its deleted
block index is at most $j-2m-2$, so its one-block halo ends at most at
$j-2m-1$.  The two excluded cases $s=2m+1,2m+2$ have just been ruled out.
Thus no deleted block has one-block halo containing $j-2m$, and so
$j-2m\notin H(S)$.

The complement assertion follows from the definition of $R$ in
\cref{lem:bc-canonical-halo-replacement-zones}.  Finally, the right endpoint
of block $j-2m$ is
\[
  \rho_{2(j-2m)},
\]
while the left endpoint of the zone whose first block is $j-2m+1$ is
\[
  \rho_{2(j-2m+1)-2}.
\]
These are the same index.
\end{proof}

\begin{corollary}[Non-full tight height zone endpoints are parser-visible]\label{cor:bc-tight-height-zone-endpoints-visible}
Work in the setting of
\cref{cor:bc-tight-height-visible-separator}, and put
\[
  a_\gamma=j-2m+1,\qquad Z_\gamma=\{a_\gamma,\ldots,j\}.
\]
Let
\[
  P=(\rho_0,\rho_1,\ldots,\rho_{2j})
\]
be a column-even two-step growth path for the tableau.  In any
replacement-zone parser that returns unchanged blocks for indices in the
complement $R$ of the canonical halo zones, the endpoint pair required for
the right-boundary zone $Z_\gamma$,
\[
  \rho_{2a_\gamma-2},\qquad \rho_{2j},
\]
is determined by the parsed unchanged separator block $B_{a_\gamma-1}(P)$
and the known terminal endpoint $\rho_{2j}=\varnothing$.
\end{corollary}

\begin{proof}
\Cref{cor:bc-tight-height-visible-separator} says that
$a_\gamma-1=j-2m$ belongs to the unchanged complement $R$.  Hence the parser
returns
\[
  B_{a_\gamma-1}(P)
  =
  (\rho_{2a_\gamma-4},\rho_{2a_\gamma-3},\rho_{2a_\gamma-2}).
\]
Its right endpoint is exactly the left endpoint $\rho_{2a_\gamma-2}$ of the
zone.  The right endpoint of this right-boundary zone is the terminal endpoint
of the full two-step growth path, which is $\rho_{2j}=\varnothing$ by
\cref{lem:bc-growth-path-model}.  Thus both endpoints are determined from the
parsed separator block and the fixed terminal endpoint.
\end{proof}

\begin{corollary}[Column-even endpoint gate for tight height prefixes]\label{cor:bc-tight-height-prefix-column-even-gate}
In the setting of
\cref{cor:bc-height-zero-surplus-label-prefix}, let
\[
  \varnothing=\lambda^{(0)}\subset\lambda^{(1)}\subset\cdots
  \subset\lambda^{(j)}
\]
be the Pieri two-strip path of $T$.  Then
\[
  \ell(\lambda^{(t)})=t\qquad(0\le t\le 2m-1),
\]
and
\[
  \ell(\lambda^{(2r)})=2r\qquad(1\le r<m).
\]
Moreover
\[
  \ell(\lambda^{(2m)})\in\{2m-1,2m\}.
\]
If $m\ge2$ and $\lambda^{(2m)}$ is column-even, then the prefix
\[
  \lambda^{(0)}\subset\cdots\subset\lambda^{(2m)}
\]
is an odd alternating height full-zone path, with
\[
  C^h_m(\lambda^{(0)},\ldots,\lambda^{(2m)})
  =
  \{1,3,\ldots,2m-1\},
\]
and is in the domain of
\cref{prop:bc-odd-alternating-height-full-zone-decoder}.  If
$m=1$, or if $\lambda^{(2m)}$ is not column-even, then the closed full-zone
alternating decoder does not by itself supply this local zone.  In the
non-column-even case either column $1$ has odd length $2m-1$, or at least one
off-first-column column of $\lambda^{(2m)}$ has odd length.  Any target-size
decoder for this tight zone must therefore use endpoint or continuation data
beyond the closed full-zone alternating decoder.
\end{corollary}

\begin{proof}
By \cref{cor:bc-height-zero-surplus-label-prefix}, the first $m$ height
witness labels are
\[
  t_r=2r-1\qquad(1\le r\le m).
\]
For each such $r$,
\cref{lem:bc-local-crossing-states} gives
\[
  \ell(\lambda^{(2r-2)})=2r-2,
  \qquad
  \ell(\lambda^{(2r-1)})=2r-1 .
\]
This proves $\ell(\lambda^{(t)})=t$ for all odd
$t\le2m-1$ and all even $t\le2m-2$.  For $1\le r<m$, applying the same
display to $r+1$ gives
\[
  \ell(\lambda^{(2r)})=2r.
\]
Finally, the horizontal strip at time $2m$ creates at most one new row from
length $2m-1$, so
\[
  \ell(\lambda^{(2m)})\in\{2m-1,2m\}.
\]

If $m\ge2$ and $\lambda^{(2m)}$ is column-even, its first column has even
nonzero length.
Since that length is one of $2m-1$ and $2m$, it must be $2m$.  Therefore
every step in the prefix creates exactly one row.  The height statistic
\[
  \min(m,\lceil\ell(\lambda^{(t)})/2\rceil)
\]
then increases exactly at the odd times
$1,3,\ldots,2m-1$.  This proves the displayed crossing set, and the
column-even endpoint hypothesis makes the prefix a column-even Pieri path of
length $2m$.  Hence it is exactly an odd alternating height full-zone path,
so \cref{prop:bc-odd-alternating-height-full-zone-decoder} applies.

If $m=1$, the proposition just cited is not available in this local size, so
this singleton tight zone remains part of the endpoint/continuation local
decoder.  If $\lambda^{(2m)}$ is not column-even, then some column of
$\lambda^{(2m)}$ has odd length.  Column $1$ has length either $2m-1$ or
$2m$, so the odd column is either column $1$ with length $2m-1$, or it is an
off-first-column odd column.  In either case the prefix endpoint is not a
column-even endpoint of the type used in the closed full-zone alternating
decoder, so a decoder for the tight zone must incorporate the endpoint or the
later continuation parity data.
\end{proof}

\begin{lemma}[The first semistandard label is forced]\label{lem:bc-first-label-forced-domino}
Let $T$ be a semistandard tableau of content $(2^j)$ with $j\ge1$.  Then the
two boxes carrying label $1$ are exactly
\[
  (1,1)\qquad\text{and}\qquad(1,2).
\]
Equivalently, in the Pieri path of \cref{lem:bc-pieri-path-bijection}, the
first strip is the forced horizontal domino
\[
  \lambda^{(1)}=(2).
\]
\end{lemma}

\begin{proof}
No box carrying label $1$ can lie below the first row: if such a box were in
row $r>1$, then the box immediately above it would have a strictly smaller
label by column strictness, impossible.  Hence both copies of label $1$ lie
in the first row.  Since rows are weakly increasing, the boxes carrying the
smallest label are the leftmost boxes of that row.  There are exactly two
copies, so they occupy columns $1$ and $2$.

The Pieri-path statement is the same assertion under
\cref{lem:bc-pieri-path-bijection}: the skew difference
$\lambda^{(1)}/\lambda^{(0)}$ consists exactly of the two boxes carrying label
$1$.
\end{proof}

\begin{corollary}[Singleton tight height zones have forced deleted strip]\label{cor:bc-singleton-tight-height-forced-strip}
In the setting of
\cref{cor:bc-tight-height-prefix-column-even-gate}, suppose $m=1$.  Then the
tight height zone is
\[
  Z_\gamma=\{j-1,j\},\qquad
  Z_\gamma\cap I=\{j\},\qquad
  Z_\gamma\setminus I=\{j-1\},
\]
and its deleted tableau label is $1$.  The deleted Pieri strip is the forced
domino
\[
  \lambda^{(1)}/\lambda^{(0)}=\{(1,1),(1,2)\}.
\]
Thus the singleton tight C-side local problem has no remaining choice about
the deleted semistandard boxes; the remaining local data are the target-size
merge with the retained label-$2$ block and the endpoint/continuation parity
needed to recover the surrounding growth-path window.
\end{corollary}

\begin{proof}
The displayed formulas for $Z_\gamma$, $Z_\gamma\cap I$, and
$Z_\gamma\setminus I$ are the case $m=1$ of
\cref{cor:bc-height-zero-surplus-label-prefix}.  The same corollary says that
the deleted label set in the zone is $\{1\}$ and that the retained label set is
$\{2\}$.  The assertion about the deleted strip is exactly
\cref{lem:bc-first-label-forced-domino}.

Finally, this zone has one retained block and one deleted block, but
\cref{prop:bc-canonical-halo-local-repair-criterion} allows exactly
$|Z_\gamma\setminus I|=1$ replacement block for the zone.  Therefore there is
no separate positive replacement window for the deleted strip in addition to
the retained label-$2$ block; the remaining task is precisely the stated
target-size merge and endpoint/continuation recovery.
\end{proof}

\begin{lemma}[Continuation parity equals the tight-prefix endpoint defect]\label{lem:bc-tight-height-continuation-parity}
In the setting of
\cref{cor:bc-tight-height-prefix-column-even-gate}, put
\[
  O_{\mathrm{pre}}=
  \{c\ge1:\,(\lambda^{(2m)})'_c\equiv1\pmod2\}
\]
and
\[
  O_{\mathrm{cont}}=
  \{c\ge1:\,
    (\lambda^{(j)})'_c-(\lambda^{(2m)})'_c\equiv1\pmod2\}.
\]
Then
\[
  O_{\mathrm{cont}}=O_{\mathrm{pre}}.
\]
In particular, if the tight prefix endpoint is not column-even, the remaining
labels $2m+1,\ldots,j$ must add an odd number of boxes exactly in the columns
of $O_{\mathrm{pre}}$.  The set $O_{\mathrm{pre}}$ has even cardinality, and
its adjacent increasing pairing is determined by the prefix endpoint
$\lambda^{(2m)}$.
\end{lemma}

\begin{proof}
The final shape $\lambda^{(j)}$ is column-even, so
\[
  (\lambda^{(j)})'_c\equiv0\pmod2
  \qquad(c\ge1).
\]
Therefore
\[
  (\lambda^{(j)})'_c-(\lambda^{(2m)})'_c
  \equiv
  -(\lambda^{(2m)})'_c
  \equiv
  (\lambda^{(2m)})'_c
  \pmod2,
\]
which proves $O_{\mathrm{cont}}=O_{\mathrm{pre}}$.

If the prefix endpoint is not column-even, then $O_{\mathrm{pre}}$ is
nonempty, and the equality just proved says exactly that the continuation
adds odd column parity in those and only those columns.  The size
$|\lambda^{(2m)}|=4m$ is even, so the sum of the column lengths of
$\lambda^{(2m)}$ is even; hence the number of odd column lengths is even.
Thus $O_{\mathrm{pre}}$ has even cardinality.  Listing its columns
increasingly and pairing adjacent entries is therefore a canonical pairing
determined by $\lambda^{(2m)}$.
\end{proof}

\begin{lemma}[Tight-prefix endpoint defects need not fit the local slots]\label{lem:bc-tight-prefix-defects-exceed-local-slots}
There is a column-even Pieri two-strip path
\[
  \varnothing=\lambda^{(0)}\subset\lambda^{(1)}\subset\cdots
  \subset\lambda^{(7)}
\]
with $C^h_2(\lambda^\bullet)=\{1,3\}$ such that the tight height zone for
$S=\{1,3\}$ has $m=2$ deleted singleton blocks, but the endpoint defect set
of the prefix through label $4$ has three paired defects.

Consequently, the non-column-even tight-prefix continuation cannot in general
be repaired by assigning every endpoint-defect pair to one of the $m$ local
frontier slots of the tight zone.  A valid target-size decoder must either
use the continuation jointly, discard/cancel some endpoint defects before the
local marking step, or encode a smaller selected subset through the tight-zone
branch word.
\end{lemma}

\begin{proof}
Consider the path
\[
\begin{aligned}
\lambda^{(0)}&=\varnothing,\\
\lambda^{(1)}&=(2),\\
\lambda^{(2)}&=(3,1),\\
\lambda^{(3)}&=(4,1,1),\\
\lambda^{(4)}&=(6,1,1),\\
\lambda^{(5)}&=(6,3,1),\\
\lambda^{(6)}&=(6,5,1),\\
\lambda^{(7)}&=(6,6,1,1).
\end{aligned}
\]
Each successive difference is a horizontal strip of size $2$: the added
boxes are, respectively,
\[
\begin{gathered}
(1,1),(1,2);\quad
(2,1),(1,3);\quad
(3,1),(1,4);\quad
(1,5),(1,6);\\
(2,2),(2,3);\quad
(2,4),(2,5);\quad
(2,6),(4,1).
\end{gathered}
\]
The terminal conjugate partition is
\[
  (\lambda^{(7)})'=(4,2,2,2,2,2),
\]
so the terminal shape is column-even.

The lengths of the successive shapes are
\[
  0,1,2,3,3,3,3,4.
\]
Thus
\[
  \min(2,\lceil\ell(\lambda^{(t)})/2\rceil)
\]
increases exactly at $t=1$ and $t=3$, so $C^h_2(\lambda^\bullet)=\{1,3\}$.
For $j=7$ and $S=\{1,3\}$, the label-to-block map gives
\[
  I(S)=\{7,5\}.
\]
The one-block halo is the single zone $\{4,5,6,7\}$, with deleted blocks
$\{5,7\}$ and retained blocks $\{4,6\}$; hence this is the tight height zone
with $m=2$.

At the prefix endpoint one has
\[
  \lambda^{(4)}=(6,1,1),\qquad
  (\lambda^{(4)})'=(3,1,1,1,1,1).
\]
Therefore
\[
  O_{\mathrm{pre}}=\{1,2,3,4,5,6\},
\]
which has three adjacent pairs.  Since the tight zone has only $m=2$ local
frontier slots, these three endpoint-defect pairs cannot all be assigned to
distinct local slots of that tight zone.
\end{proof}

\begin{lemma}[Finite Hall injection criterion]\label{lem:bc-finite-hall-injection}
Let $A$ and $B$ be finite sets, and let $N(a)\subseteq B$ be a neighbour set
for each $a\in A$.  If
\[
  \left|\bigcup_{a\in X}N(a)\right|\ge |X|
  \qquad\text{for every }X\subseteq A,
\]
then there is an injection $\phi:A\to B$ with $\phi(a)\in N(a)$ for every
$a\in A$.
\end{lemma}

\begin{proof}
We prove the assertion by induction on $|A|$.  The cases $|A|=0,1$ are
immediate.  Suppose first that every nonempty proper subset
$X\subsetneq A$ satisfies
\[
  \left|\bigcup_{a\in X}N(a)\right|\ge |X|+1.
\]
Choose $a_0\in A$ and $b_0\in N(a_0)$.  For any
$Y\subseteq A\setminus\{a_0\}$, the neighbour set after removing $b_0$ has
size at least $|Y|$, because the original neighbour set of $Y$ has size at
least $|Y|+1$ when $Y$ is nonempty.  The induction hypothesis gives an
injection of $A\setminus\{a_0\}$ into $B\setminus\{b_0\}$ respecting
neighbours.  Sending $a_0$ to $b_0$ completes the matching.

It remains to consider the case where some nonempty proper subset
$X\subsetneq A$ has equality:
\[
  \left|\bigcup_{a\in X}N(a)\right|=|X|.
\]
Put $B_X=\bigcup_{a\in X}N(a)$.  By induction, $X$ injects into $B_X$.  For
$Y\subseteq A\setminus X$,
\[
  \left|\bigcup_{a\in Y}N(a)\setminus B_X\right|
  =
  \left|B_X\cup\bigcup_{a\in Y}N(a)\right|-|B_X|
  \ge |X\cup Y|-|X|
  = |Y|.
\]
Thus the induction hypothesis also matches $A\setminus X$ into
$B\setminus B_X$.  Combining the two injections proves the claim.
\end{proof}

\begin{lemma}[Continuation labels cover endpoint defects]\label{lem:bc-tight-continuation-label-cover}
In the setting of
\cref{lem:bc-tight-height-continuation-parity}, write
\[
  O_{\mathrm{pre}}=\{o_1<o_2<\cdots<o_{2r}\}
\]
and pair adjacent entries
\[
  P_a=\{o_{2a-1},o_{2a}\}\qquad(1\le a\le r).
\]
For each continuation label $t>2m$, let
\[
  C_t=\{c\ge1:\,(\lambda^{(t)})'_c-(\lambda^{(t-1)})'_c=1\}
\]
be the set of columns hit by the horizontal two-strip of label $t$.  Then
there is an injection
\[
  \phi:\{1,\ldots,r\}\hookrightarrow\{2m+1,\ldots,j\}
\]
such that
\[
  C_{\phi(a)}\cap P_a\ne\varnothing
  \qquad(1\le a\le r).
\]
Thus the paired endpoint defects of a tight prefix can be charged to distinct
continuation labels, although by
\cref{lem:bc-tight-prefix-defects-exceed-local-slots} they cannot in general
all be charged to the tight zone's own frontier slots.
\end{lemma}

\begin{proof}
For a subset $A_0\subseteq\{1,\ldots,r\}$, put
\[
  P(A_0)=\bigcup_{a\in A_0}P_a.
\]
Every column in $P(A_0)$ belongs to $O_{\mathrm{pre}}$.  By
\cref{lem:bc-tight-height-continuation-parity}, the continuation from
$2m+1$ through $j$ adds an odd, hence positive, number of boxes in each such
column.  Therefore the total number of incidences
\[
  (t,c)\quad\text{with}\quad t>2m,\ c\in C_t\cap P(A_0)
\]
is at least $|P(A_0)|=2|A_0|$.

Each continuation label $t$ contributes a horizontal two-strip, so
$|C_t|=2$ and it contributes at most two incidences to $P(A_0)$.  Hence the
number of continuation labels with $C_t\cap P(A_0)\ne\varnothing$ is at least
$|A_0|$.  This is exactly the Hall condition for the paired-defect set, with
neighbours those continuation labels whose strips hit the pair.  Applying
\cref{lem:bc-finite-hall-injection} gives the desired injection.
\end{proof}

\begin{lemma}[The continuation Hall charge can be chosen canonically]\label{lem:bc-canonical-continuation-hall-charge}
In the setting of \cref{lem:bc-tight-continuation-label-cover}, order the
defect pairs by $1<\cdots<r$ and the continuation labels by their usual order.
For each pair set
\[
  N(a)=\{t\in\{2m+1,\ldots,j\}: C_t\cap P_a\ne\varnothing\}.
\]
Among all injections
\[
  \phi:\{1,\ldots,r\}\hookrightarrow\{2m+1,\ldots,j\}
  \qquad\text{with}\qquad \phi(a)\in N(a),
\]
there is a lexicographically least word
\[
  (\phi(1),\ldots,\phi(r)).
\]
This canonical injection is a fixed function of the incidence family
$(N(a))_{a=1}^r$.

Consequently, if the incidence family is recovered from parser-visible
outside growth-diagram subfillings by
\cref{cor:bc-parsed-boundary-recovers-knuth-subfilling} or
\cref{cor:bc-parsed-dual-rsk-prime-boundary-recovers-subfilling}, then the
chosen Hall charge and the charged continuation labels are admissible
auxiliary data for \cref{prop:bc-right-boundary-auxiliary-global-recovery};
no additional branch word is needed to select the matching.
\end{lemma}

\begin{proof}
\Cref{lem:bc-tight-continuation-label-cover} proves that at least one such
injection exists.  Since both the defect-pair set and the continuation-label
set are finite and linearly ordered, the set of admissible image words
$(\phi(1),\ldots,\phi(r))$ is a nonempty finite set in lexicographic order.
It therefore has a unique least element.  The admissible words are determined
exactly by the neighbour sets $N(a)$ together with the injectivity condition,
so this least word is a fixed function of the incidence family.

If the incidence family is recovered from parsed outside subfillings in either
listed growth-diagram variant, then the same fixed lexicographic rule computes
the chosen matching from the already recovered outside data.  Thus the
matching is an auxiliary datum of the kind allowed in
\cref{prop:bc-right-boundary-auxiliary-global-recovery}, rather than a new
encoded choice.
\end{proof}

\begin{lemma}[Continuation incidence is semistandard-local]\label{lem:bc-continuation-incidence-semistandard-local}
Let $T$ be a semistandard tableau of content $(2^j)$, and let
\[
  \varnothing=\lambda^{(0)}\subset\lambda^{(1)}\subset\cdots\subset\lambda^{(j)}
\]
be its Pieri two-strip path under \cref{lem:bc-pieri-path-bijection}.  For a
label $t$, put
\[
  C_t=\{c\ge1:\,(\lambda^{(t)})'_c-(\lambda^{(t-1)})'_c=1\}.
\]
Then $C_t$ is exactly the set of columns occupied by the two boxes of $T$
carrying label $t$.  Consequently, in the setting of
\cref{lem:bc-tight-continuation-label-cover}, the incidence relation
\[
  C_t\cap P_a\ne\varnothing
\]
is determined by the positions of the two semistandard boxes carrying the
continuation label $t$.
\end{lemma}

\begin{proof}
By \cref{lem:bc-pieri-path-bijection}, the skew difference
$\lambda^{(t)}/\lambda^{(t-1)}$ is precisely the set of boxes of $T$ carrying
label $t$.  Passing to conjugates records column lengths.  Hence
$(\lambda^{(t)})'_c-(\lambda^{(t-1)})'_c=1$ exactly when this skew difference
contains a box in column $c$.  Since the difference is a horizontal strip,
there is at most one such box in each column, and the displayed set is exactly
the two columns occupied by label $t$.  The final assertion is the definition
of the incidence relation in
\cref{lem:bc-tight-continuation-label-cover} read with this identification.
\end{proof}

\begin{lemma}[Continuation boxes determine the Pieri continuation]\label{lem:bc-continuation-boxes-determine-pieri-segment}
Let $T$ be a semistandard tableau of content $(2^j)$, and let
\[
  \varnothing=\lambda^{(0)}\subset\lambda^{(1)}\subset\cdots\subset\lambda^{(j)}
\]
be its Pieri two-strip path.  Fix $a<j$.  If one is given $\lambda^{(a)}$ and,
for each label $t>a$, the two boxes of $T$ carrying label $t$, then the whole
Pieri continuation segment
\[
  \lambda^{(a)}\subset\lambda^{(a+1)}\subset\cdots\subset\lambda^{(j)}
\]
is determined.  Explicitly,
\[
  \lambda^{(t)}
  =
  \lambda^{(a)}
  \cup
  \bigcup_{u=a+1}^{t}\{\text{boxes of }T\text{ carrying label }u\}
  \qquad(a<t\le j).
\]
Consequently, in the tight-height setting with $a=2m$, recovery of the prefix
endpoint $\lambda^{(2m)}$ together with the labelled continuation boxes
recovers the Hall-charge auxiliary data of
\cref{lem:bc-pieri-continuation-determines-hall-charge}.
\end{lemma}

\begin{proof}
By \cref{lem:bc-pieri-path-bijection}, $\lambda^{(t)}$ is the shape formed by
all boxes whose labels are at most $t$.  The boxes whose labels are at most
$a$ form $\lambda^{(a)}$, and for $u>a$ the two boxes labelled $u$ are exactly
the skew difference $\lambda^{(u)}/\lambda^{(u-1)}$.  Thus adding, in label
order, the recovered two-box sets for $u=a+1,\ldots,t$ to
$\lambda^{(a)}$ gives the displayed formula for every $t>a$.

The final sentence is the case $a=2m$ followed by
\cref{lem:bc-pieri-continuation-determines-hall-charge}.
\end{proof}

\begin{proposition}[Recovered continuation boxes supply the tight auxiliary data]\label{prop:bc-tight-height-auxiliary-from-continuation-boxes}
Work in the non-full tight-height setting of
\cref{cor:bc-tight-height-zone-endpoints-visible}, and let
\[
  \varnothing=\lambda^{(0)}\subset\lambda^{(1)}\subset\cdots\subset\lambda^{(j)}
\]
be the Pieri path of the tableau.  Suppose the parser data outside the tight
right-boundary zone determine the prefix endpoint $\lambda^{(2m)}$ and, for
each continuation label $t=2m+1,\ldots,j$, the two boxes of the tableau
carrying label $t$.  Then those outside data determine:
\begin{enumerate}
\item the growth-path endpoint pair
\[
  \rho_{2(j-2m+1)-2},\qquad \rho_{2j}
\]
of the tight right-boundary zone;
\item the endpoint-defect pairs of $\lambda^{(2m)}$;
\item every incidence $C_t\cap P_a\ne\varnothing$ between a continuation
label and an endpoint-defect pair; and
\item the canonical continuation Hall charge of
\cref{lem:bc-canonical-continuation-hall-charge}.
\end{enumerate}
Thus, once a right-boundary tight-zone inverse is written using only these
data as its auxiliary input, the auxiliary input is of the kind allowed by
\cref{prop:bc-right-boundary-auxiliary-global-recovery}.
\end{proposition}

\begin{proof}
The endpoint pair in item (1) is determined by
\cref{cor:bc-tight-height-zone-endpoints-visible}: the parsed separator block
gives the left endpoint and the terminal endpoint is $\rho_{2j}=\varnothing$.

By \cref{lem:bc-continuation-boxes-determine-pieri-segment}, the recovered
prefix endpoint and recovered labelled continuation boxes determine the whole
Pieri continuation segment
\[
  \lambda^{(2m)}\subset\lambda^{(2m+1)}\subset\cdots\subset\lambda^{(j)} .
\]
Then \cref{lem:bc-pieri-continuation-determines-hall-charge} determines the
endpoint-defect pairs, the continuation incidence relation, and the canonical
Hall charge from that segment.  These are precisely items (2)--(4).

All data listed above are fixed functions of the outside parser data under
the stated recovery hypothesis.  Therefore they may be used as the auxiliary
datum $A(P)$ in
\cref{prop:bc-right-boundary-auxiliary-global-recovery} without adding branch
choices.
\end{proof}

\begin{lemma}[Outside data recover the strict-left growth segment]\label{lem:bc-tight-outside-recovers-left-growth-segment}
Work in the non-full tight-height setting of
\cref{cor:bc-tight-height-zone-endpoints-visible}.  Put
\[
  a_\gamma=j-2m+1,\qquad Z_\gamma=\{a_\gamma,\ldots,j\},
\]
and order the canonical halo zones increasingly, with $Z_\gamma$ the final
right-boundary zone.  In the outside datum
$\mathcal D_{\mathrm{out}}(P)$ of
\cref{prop:bc-right-boundary-auxiliary-global-recovery}, the unchanged
complement blocks and the recovered earlier zone subpaths determine the whole
initial growth segment
\[
  \rho_0,\rho_1,\ldots,\rho_{2a_\gamma-2}
  =
  \rho_0,\rho_1,\ldots,\rho_{2(j-2m)} .
\]
Under the label-to-block convention of \cref{lem:bc-growth-path-model}, this
is exactly the two-step growth-block segment corresponding to continuation
labels
\[
  2m+1,\ldots,j .
\]
\end{lemma}

\begin{proof}
Since the canonical halo zones are disjoint intervals and
$Z_\gamma=\{a_\gamma,\ldots,j\}$ is the final one, every block index
$i<a_\gamma$ lies either in the unchanged complement $R$ or in a canonical
halo zone strictly to the left of $Z_\gamma$.  The outside datum in
\cref{prop:bc-right-boundary-auxiliary-global-recovery} contains the original
two-step block $B_i(P)$ for every such $i\in R$, and it contains the whole
recovered subpath for each earlier halo zone.  Interleaving these blocks and
subpaths in increasing block order therefore reconstructs the initial segment
through the right endpoint of block $a_\gamma-1$, namely
\[
  \rho_0,\rho_1,\ldots,\rho_{2(a_\gamma-1)}
  =
  \rho_0,\rho_1,\ldots,\rho_{2a_\gamma-2}.
\]
The formula $a_\gamma=j-2m+1$ gives the displayed equality with
$\rho_{2(j-2m)}$.

Finally, \cref{lem:bc-growth-path-model} identifies tableau label $\ell$ with
block $j-\ell+1$.  As $\ell$ runs from $2m+1$ to $j$, this block index runs
from $j-2m$ down to $1$, exactly the block interval
$\{1,\ldots,j-2m\}$ reconstructed above.
\end{proof}

\begin{proposition}[Strict-left growth bridge supplies tight auxiliary data]\label{prop:bc-strict-left-growth-bridge-supplies-tight-auxiliary}
Work in the non-full tight-height setting of
\cref{lem:bc-tight-outside-recovers-left-growth-segment}.  Suppose that the
strict-left Krattenthaler growth segment
\[
  \rho_0,\rho_1,\ldots,\rho_{2(j-2m)}
\]
determines the prefix endpoint $\lambda^{(2m)}$ of the Pieri path and, for
each continuation label $t=2m+1,\ldots,j$, the two boxes of the tableau
carrying label $t$.  Then the outside datum
$\mathcal D_{\mathrm{out}}(P)$ determines all auxiliary data listed in
\cref{prop:bc-tight-height-auxiliary-from-continuation-boxes}.  Consequently,
any tight right-boundary inverse that uses only those data as its auxiliary
input satisfies the auxiliary-data hypothesis of
\cref{prop:bc-right-boundary-auxiliary-global-recovery}.
\end{proposition}

\begin{proof}
By \cref{lem:bc-tight-outside-recovers-left-growth-segment}, the outside datum
$\mathcal D_{\mathrm{out}}(P)$ determines the strict-left segment
\[
  \rho_0,\rho_1,\ldots,\rho_{2(j-2m)} .
\]
The hypothesis of the present proposition then determines
$\lambda^{(2m)}$ and the two boxes carrying every continuation label
$2m+1,\ldots,j$.  Applying
\cref{prop:bc-tight-height-auxiliary-from-continuation-boxes} to these
recovered data gives the tight-zone endpoint pair, endpoint-defect pairs,
continuation incidences, and canonical Hall charge.  These are exactly the
listed auxiliary data.  Since all of them are fixed functions of
$\mathcal D_{\mathrm{out}}(P)$, a right-boundary inverse using only these data
has auxiliary input of the form allowed in
\cref{prop:bc-right-boundary-auxiliary-global-recovery}.
\end{proof}

\begin{lemma}[Strict-left data recover continuation-column Step-3 entries]\label{lem:bc-strict-left-recovers-continuation-columns}
Work in the non-full tight-height setting of
\cref{lem:bc-tight-outside-recovers-left-growth-segment}.  Put $M=2m$.  In the
Step~3 dual-RSK' half-square of
\cref{lem:bc-krattenthaler-semistandard-path-boundary}, let $J_M$ be the
Ferrers subarrangement formed by the first $j-M$ strict-left columns, i.e. the
columns incident to the continuation labels $M+1,\ldots,j$ in the label order
$j,j-1,\ldots,1$.  Then the outside datum $\mathcal D_{\mathrm{out}}(P)$
determines every cell entry of $J_M$ and every corner label on or inside
$J_M$.
\end{lemma}

\begin{proof}
By \cref{lem:bc-tight-outside-recovers-left-growth-segment}, the outside datum
determines the strict-left growth segment
\[
  \rho_0,\rho_1,\ldots,\rho_{2(j-M)} .
\]
Under \cref{lem:bc-growth-path-model}, this is exactly the Step~3 top/right
boundary segment for the continuation labels $M+1,\ldots,j$.  These labels are
the first $j-M$ labels in Krattenthaler's Step~3 order $j,j-1,\ldots,1$, so
the displayed segment is the top/right boundary of $J_M$.  Applying
\cref{cor:bc-parsed-dual-rsk-prime-boundary-recovers-subfilling} to this
Ferrers subarrangement recovers its cell entries and internal corner labels.
\end{proof}

\begin{corollary}[Strict-left data reduce continuation recovery to the internal Q-cut]\label{cor:bc-strict-left-q-suffix-boundary-recovers-continuation-segment}
Work in the non-full tight-height setting of
\cref{lem:bc-tight-outside-recovers-left-growth-segment}.  Put $M=2m$, and let
$C_M^+$ be the Step~2 $Q$-suffix strip of
\cref{lem:bc-step2-q-suffix-strip-determines-continuation}; let $C_M$ be the
corresponding Step~2 $Q$-prefix strip.  Suppose that the
target-window data
\[
  S,\quad Z_L,\quad V^\ast_L,\quad \mathcal D_{\mathrm{out}}(P),
  \quad \rho_{2a_L-2},\rho_{2j}
\]
determine the ordered partition labels on the internal separating boundary
between the Step~2 $Q$-prefix strip $C_M$ and the suffix strip $C_M^+$.  Then
the same data determine the Pieri continuation segment
\[
  \lambda^{(2m)}\subset\lambda^{(2m+1)}\subset\cdots\subset\lambda^{(j)} .
\]
\end{corollary}

\begin{proof}
By \cref{lem:bc-strict-left-recovers-continuation-columns}, the outside datum
$\mathcal D_{\mathrm{out}}(P)$ determines every Step~3 entry in the first
$j-M$ strict-left columns.  By
\cref{lem:bc-step3-continuation-incidence-covers-q-suffix}, those entries
determine every non-diagonal entry of the Step~2 $Q$-suffix strip $C_M^+$, and
the diagonal entries of that strip are zero.  Hence the displayed target-window
data determine all cell entries of $C_M^+$.

By the hypothesis and
\cref{lem:bc-q-suffix-boundary-internal-cut}, the same data determine the full
left/bottom boundary of $C_M^+$.  Together with those boundary labels,
\cref{lem:bc-step2-q-suffix-strip-determines-continuation} determines the
$Q$-side continuation boundary
\[
  \lambda^{(M)},\lambda^{(M+1)},\ldots,\lambda^{(j)} .
\]
Since $M=2m$, this is the asserted Pieri continuation segment.
\end{proof}

\begin{corollary}[Strict-left plus low-low entries recover the internal Q-cut]\label{cor:bc-strict-left-low-low-determines-internal-q-cut}
Work in the non-full tight-height setting of
\cref{lem:bc-tight-outside-recovers-left-growth-segment}.  Put $M=2m$, and let
$L_M$ be the low-low Step~3 cell set of
\cref{lem:bc-low-label-incidence-split}.  Let $C_M$ and $C_M^+$ be the
corresponding Step~2 $Q$-prefix and $Q$-suffix strips.  Suppose that the
target-window data
\[
  S,\quad Z_L,\quad V^\ast_L,\quad \mathcal D_{\mathrm{out}}(P),
  \quad \rho_{2a_L-2},\rho_{2j}
\]
determine every Step~3 cell entry on $L_M$.  Then the same data determine the
ordered partition labels on the internal separating boundary between the
Step~2 $Q$-prefix strip $C_M$ and the suffix strip $C_M^+$.  Consequently they
determine the Pieri continuation segment
\[
  \lambda^{(2m)}\subset\lambda^{(2m+1)}\subset\cdots\subset\lambda^{(j)} .
\]
\end{corollary}

\begin{proof}
By \cref{lem:bc-strict-left-recovers-continuation-columns}, the outside datum
$\mathcal D_{\mathrm{out}}(P)$ determines every Step~3 entry in $J_M$, the
first $j-M$ strict-left columns.  By \cref{lem:bc-low-label-incidence-split},
the cross part $X_M$ of $H_M$ lies in $J_M$, while the low-low part is $L_M$.
The hypothesis supplies the entries on $L_M$, so the displayed target-window
data determine every Step~3 entry on
\[
  H_M=X_M\sqcup L_M .
\]
\Cref{cor:bc-low-label-incidence-determines-internal-q-cut} then determines
the ordered internal cut labels between $C_M$ and $C_M^+$.  Finally
\cref{cor:bc-strict-left-q-suffix-boundary-recovers-continuation-segment}
applies with this recovered internal cut and gives the displayed Pieri
continuation segment.
\end{proof}

\begin{lemma}[The strict-left dual-RSK' boundary alone does not determine the tight prefix]\label{lem:bc-strict-left-alone-prefix-obstruction}
In the $m=0$ specialization of Krattenthaler's semistandard construction,
there exist two column-even semistandard tableaux of content $(2^6)$ in the
same tight height fixed-witness fiber
\[
  H_2(T)=\{1,3\}
\]
for which the Step~3 dual-RSK' subfilling in the first two strict-left columns
is identical, and hence its top/right boundary segment
\[
  \rho_0,\rho_1,\rho_2,\rho_3,\rho_4
\]
is identical, but the Pieri prefix endpoints at label $4$ are different and
both prefix endpoints are non-column-even.
Consequently, the bridge in
\cref{prop:bc-strict-left-growth-bridge-supplies-tight-auxiliary} cannot be
proved from the recovered strict-left boundary segment together with the
tight-witness condition alone; a successful tight-zone inverse must also use
the extra constraints or data supplied by the tight right-boundary zone.
\end{lemma}

\begin{proof}
Consider the following two symmetric zero-diagonal matrices with rows and
columns indexed by $1,\ldots,6$:
\[
  A_{14}=A_{41}=A_{16}=A_{61}=A_{23}=A_{32}=A_{25}=A_{52}
  =A_{35}=A_{53}=A_{46}=A_{64}=1,
\]
all other entries being zero, and
\[
  B_{13}=B_{31}=B_{16}=B_{61}=B_{24}=B_{42}=B_{25}=B_{52}
  =B_{35}=B_{53}=B_{46}=B_{64}=1,
\]
all other entries being zero.  Every row sum of both matrices is $2$.
Applying ordinary RSK in the row-word convention used in
\cref{lem:bc-knuth-growth-local-reversibility} gives, for $A$, the tableau
\[
  T_A=
  \begin{array}{cccc}
  1&1&3&4\\
  2&2&5&6\\
  3&5&&\\
  4&6&&
  \end{array}
\]
and, for $B$, the tableau
\[
  T_B=
  \begin{array}{ccccc}
  1&1&2&3&4\\
  2&4&5&5&6\\
  3&&&&\\
  6&&&&
  \end{array}.
 \]
Indeed, the bottom words inserted under top labels $1,\ldots,6$ are
\[
  4,6,3,5,2,5,1,6,2,3,1,4
\]
for $A$, and
\[
  3,6,4,5,1,5,2,6,2,3,1,4
\]
for $B$; direct row insertion gives the displayed insertion tableau in each
case, and symmetry gives the same recording tableau.  Both final shapes are
column-even:
\[
  (4,4,2,2)'=(4,4,2,2),\qquad
  (5,5,1,1)'=(4,2,2,2,2).
\]
The heights of the Pieri prefixes for $T_A$ are
\[
  1,2,3,4,4,4,
\]
and the heights of the Pieri prefixes for $T_B$ are
\[
  1,2,3,3,3,4.
\]
Therefore, in both cases, the capped height statistic
$\min(2,\lceil\ell(\lambda^{(t)})/2\rceil)$ increases exactly at
$t=1$ and $t=3$.  Thus both tableaux lie in the tight height fixed-witness
fiber $H_2(T)=\{1,3\}$.  However the Pieri prefix endpoints after labels at
most $4$ are
\[
  \lambda_A^{(4)}=(4,2,1,1),
  \qquad
  \lambda_B^{(4)}=(5,2,1),
\]
so they are different.  Their conjugates are, respectively,
\[
  (4,2,1,1)'=(4,2,1,1),
  \qquad
  (5,2,1)'=(3,2,1,1,1),
\]
so neither prefix endpoint is column-even.

Now order the rows and columns of the Step~2 symmetric matrix as in
Krattenthaler's proof, namely $6,5,4,3,2,1$, and pass to the Step~3
strict lower-triangular dual-RSK' half-square.  The first two strict-left
columns are exactly the entries incident to labels $6$ and $5$ on the
strict-left side.  For both $A$ and $B$, these nonzero entries are the four
unit cells joining
\[
  6\text{ to }1,\qquad 6\text{ to }4,\qquad
  5\text{ to }2,\qquad 5\text{ to }3,
\]
with all other entries in those two columns equal to zero.  Therefore the
Step~3 subfilling in the first two strict-left columns is the same for $A$
and $B$.  Since the left and bottom boundaries of this subarrangement are
empty, the forward dual-RSK' local rules of
\cref{lem:bc-dual-rsk-prime-growth-local-reversibility} give the same
top/right boundary segment $\rho_0,\ldots,\rho_4$.

Thus the same strict-left boundary segment and the same tight height witness
condition are compatible with two different non-column-even Pieri prefix
endpoints at label $4$.  The final sentence follows.
\end{proof}

\begin{lemma}[Strict-left data alone do not determine the low-low triangle]\label{lem:bc-strict-left-alone-low-low-obstruction}
In the two examples of
\cref{lem:bc-strict-left-alone-prefix-obstruction}, put $M=4$.  The recovered
Step~3 strict-left continuation-column subfilling, and hence the residual
degree vector
\[
  d_a=2-\sum_{t=5}^6 u_{\{a,t\}}\qquad(1\le a\le4),
\]
is the same for $T_A$ and $T_B$.  Nevertheless the Step~3 low-low assignments
on $L_4$ are different:
\[
  T_A:\quad u_{\{1,4\}}=u_{\{2,3\}}=1,
  \qquad
  T_B:\quad u_{\{1,3\}}=u_{\{2,4\}}=1,
\]
with all other low-low entries equal to $0$ in the two cases.
Consequently, the low-low supplier in the large tight right-boundary zone
cannot be derived from the strict-left continuation subfilling or its residual
degrees alone; it must also use target-window, right-boundary, or equivalent
semistandard data.
\end{lemma}

\begin{proof}
The preceding proof shows that the first two strict-left columns of the Step~3
dual-RSK' half-square are identical for $T_A$ and $T_B$.  These are the
columns incident to the continuation labels $6$ and $5$, so they are exactly
the continuation-column subfilling for $M=4$.  Reading the displayed matrices,
the entries incident from the low labels $1,2,3,4$ to the continuation labels
$5,6$ are, in both examples,
\[
  \{1,6\},\qquad \{2,5\},\qquad \{3,5\},\qquad \{4,6\},
\]
each with multiplicity one and all other continuation incidences zero.  Hence
\[
  \sum_{t=5}^6 u_{\{a,t\}}=1
  \qquad(1\le a\le4),
\]
and the residual degree vector is $d_1=d_2=d_3=d_4=1$ for both examples.

By \cref{lem:bc-step3-cells-are-step2-half-matrix}, the Step~3 low-low
entries on $L_4$ are exactly the off-diagonal Step~2 matrix entries whose two
labels lie in $\{1,2,3,4\}$.  In the matrix $A$, the only such nonzero entries
are $A_{14}=A_{41}=1$ and $A_{23}=A_{32}=1$.  In the matrix $B$, the only such
nonzero entries are $B_{13}=B_{31}=1$ and $B_{24}=B_{42}=1$.  This gives the
displayed two different low-low assignments with the same continuation
subfilling and the same residual degrees.
\end{proof}

\begin{corollary}[Strict-left data alone do not determine the terminal boundary]\label{cor:bc-strict-left-alone-terminal-boundary-obstruction}
In the two examples of
\cref{lem:bc-strict-left-alone-prefix-obstruction}, put $M=4$.  The recovered
strict-left continuation-column subfilling and the residual degree vector are
the same, but the terminal Step~3 boundary segment
\[
  \rho_{2(j-M)},\rho_{2(j-M)+1},\ldots,\rho_{2j}
\]
cannot be the same.  Consequently, the terminal-boundary supplier for the
large tight right-boundary zone cannot be a function only of the strict-left
continuation subfilling or of its residual degrees.
\end{corollary}

\begin{proof}
For $j=6$ and $M=4$, the displayed segment is
\[
  \rho_4,\rho_5,\ldots,\rho_{12}.
\]
If the two examples had the same terminal boundary segment, then
\cref{lem:bc-low-low-terminal-corner-boundary} would recover the same Step~3
entries on $L_4$ in both examples.  This contradicts
\cref{lem:bc-strict-left-alone-low-low-obstruction}, which gives different
assignments on $L_4$ while the strict-left continuation-column subfilling and
the residual degree vector are the same.  Hence any deterministic
terminal-boundary source must use target-window, right-boundary, or equivalent
semistandard data beyond those strict-left quantities.
\end{proof}

\begin{lemma}[The strict-left obstruction is a packet endpoint obstruction]\label{lem:bc-strict-left-obstruction-packet-defects}
In the two examples of
\cref{lem:bc-strict-left-alone-prefix-obstruction}, the missing datum is
already visible in the endpoint-defect data of the singleton deleted packets.
More precisely, for the witness set $S=\{1,3\}$ the deleted packets are
$J_1=\{1\}$ and $J_2=\{3\}$.  The packet $J_1$ has the same endpoint pair and
the same local defect columns for $T_A$ and $T_B$, namely
\[
  \lambda_A^{(0)}=\lambda_B^{(0)}=\varnothing,\qquad
  \lambda_A^{(1)}=\lambda_B^{(1)}=(2),
  \qquad
  O_{J_1}(T_A)=O_{J_1}(T_B)=\{1,2\}.
\]
The packet $J_2$, however, has different endpoint pairs and different local
defect columns:
\[
  \lambda_A^{(2)}=(2,2),\quad
  \lambda_A^{(3)}=(3,2,1),\quad
  O_{J_2}(T_A)=\{1,3\},
\]
whereas
\[
  \lambda_B^{(2)}=(3,1),\quad
  \lambda_B^{(3)}=(4,1,1),\quad
  O_{J_2}(T_B)=\{1,4\}.
\]
Consequently, a tight right-boundary decoder that uses the packet inverse of
\cref{lem:bc-local-branch-recovery} on these examples must obtain the
$J_2$ endpoint-defect datum from the final-zone constraints or auxiliary
right-boundary data before the local branch word can be interpreted.
\end{lemma}

\begin{proof}
The displayed tableaux in
\cref{lem:bc-strict-left-alone-prefix-obstruction} give the Pieri prefix
shapes directly.  For both tableaux, the two boxes carrying label $1$ occupy
the first row, so
\[
  \lambda_A^{(0)}=\lambda_B^{(0)}=\varnothing,\qquad
  \lambda_A^{(1)}=\lambda_B^{(1)}=(2).
\]
Since $(2)'=(1,1)$, the local defect columns of this singleton packet are
\[
  \{c:\,(\lambda^{(1)})'_c-(\lambda^{(0)})'_c\equiv1\pmod2\}=\{1,2\}
\]
for both tableaux.

For $T_A$, after labels at most $2$ the first two rows have length $2$, and
after label $3$ the shape is $(3,2,1)$:
\[
  \lambda_A^{(2)}=(2,2),\qquad
  \lambda_A^{(3)}=(3,2,1).
\]
Thus
\[
  (\lambda_A^{(2)})'=(2,2),\qquad
  (\lambda_A^{(3)})'=(3,2,1),
\]
and the columns in which the conjugate endpoint difference is odd are
$1$ and $3$.

For $T_B$, labels at most $2$ give the shape $(3,1)$, and labels at most $3$
give the shape $(4,1,1)$:
\[
  \lambda_B^{(2)}=(3,1),\qquad
  \lambda_B^{(3)}=(4,1,1).
\]
Hence
\[
  (\lambda_B^{(2)})'=(2,1,1),\qquad
  (\lambda_B^{(3)})'=(3,1,1,1),
\]
so the odd endpoint-difference columns are $1$ and $4$.  These are exactly
the local defect columns used by item~(2) of
\cref{lem:bc-local-branch-recovery}.

By the preceding lemma, the strict-left boundary segment and the witness set
agree for the two examples, while the $J_2$ endpoint-defect datum differs.
Therefore any decoder that proceeds through the local-branch recovery lemma
must recover this datum from the right-boundary-zone side before applying the
local packet inverse.
\end{proof}

\begin{lemma}[Retained closure does not recover all tight-prefix packets]\label{lem:bc-tight-retained-closure-obstruction}
There is a non-full tight height example in which the retained boxes and their
smallest column-even Young closure do not recover the deleted singleton
packets.  More precisely, there is a column-even semistandard tableau of
content $(2^8)$ with
\[
  H_3(T)=\{1,3,5\}
\]
such that the tight right-boundary zone has deleted labels $\{1,3,5\}$ and
retained labels $\{2,4,6\}$, the prefix endpoint $\lambda^{(6)}$ is not
column-even, and after deleting labels $\{1,3,5\}$ the smallest column-even
Young diagram containing the retained boxes is strictly smaller than the
original terminal shape.

Consequently, a tight-zone decoder cannot recover all deleted packets merely
by taking the column-even closure of the retained labelled boxes, even when
the endpoint-defect parity of the tight prefix is known.  It must also recover
or encode all-deleted even-column pairs inside the tight prefix.
\end{lemma}

\begin{proof}
Consider the tableau
\[
  T=
  \begin{array}{cccc}
  1&1&4&6\\
  2&2&7&8\\
  3&3&&\\
  4&5&&\\
  5&&&\\
  6&&&\\
  7&&&\\
  8&&&
  \end{array}.
\]
It has content $(2^8)$, its rows are weakly increasing, and its columns are
strictly increasing.  Its terminal shape is
\[
  \lambda^{(8)}=(4,4,2,2,1,1,1,1),
  \qquad
  (\lambda^{(8)})'=(8,4,2,2),
\]
so the terminal shape is column-even.

The successive prefix shapes have lengths
\[
  1,2,3,4,5,6,7,8 .
\]
Therefore
\[
  \min(3,\lceil\ell(\lambda^{(t)})/2\rceil)
\]
increases exactly at $t=1,3,5$, and hence $H_3(T)=\{1,3,5\}$.  Since
$j=8$ and $m=3$, \cref{cor:bc-height-zero-surplus-label-prefix} gives the
tight right-boundary zone
\[
  Z_\gamma=\{3,4,5,6,7,8\},
\]
whose deleted labels are $\{1,3,5\}$ and whose retained in-zone labels are
$\{2,4,6\}$.  The prefix endpoint through label $6$ is
\[
  \lambda^{(6)}=(4,2,2,2,1,1),
  \qquad
  (\lambda^{(6)})'=(6,4,1,1),
\]
so it is not column-even.

Delete the boxes carrying labels $1,3,5$.  The remaining boxes include
$(8,1)$ and $(2,4)$, so any Young diagram containing them contains the shape
\[
  \mu=(4,4,1,1,1,1,1,1).
\]
Conversely, this shape contains every retained box.  Its conjugate is
\[
  \mu'=(8,2,2,2),
\]
so $\mu$ is already column-even.  Hence $\mu$ is the smallest column-even
Young diagram containing the retained boxes.

However the original terminal shape is strictly larger:
\[
  \lambda^{(8)}\setminus\mu=\{(3,2),(4,2)\}.
\]
Both of these boxes are deleted boxes, carrying labels $3$ and $5$,
respectively.  They form an all-deleted pair in column $2$, so this missing
pair is invisible to the column parity of the retained closure.  It is also
not an endpoint defect of the tight prefix, since
\[
  (\lambda^{(6)})'=(6,4,1,1)
\]
has odd columns only at $3$ and $4$.  Thus endpoint-defect parity does not
record the missing column-$2$ pair.  The final assertion follows.
\end{proof}

\begin{lemma}[Prefix endpoint and retained even labels do not recover the tight packets]\label{lem:bc-tight-even-retained-prefix-obstruction}
There are two column-even semistandard tableaux of content $(2^9)$ with the
same tight height witness set
\[
  H_4(T)=\{1,3,5,7\}
\]
such that the following data agree:
\begin{enumerate}
\item the tight-prefix endpoint $\lambda^{(8)}$;
\item the two boxes carrying each retained tight-prefix label
      $2,4,6,8$; and
\item the two boxes carrying the continuation label $9$.
\end{enumerate}
Nevertheless the deleted singleton packets for the odd labels $5$ and $7$
are different.  Consequently, a final tight-zone decoder cannot recover the
full deleted packet data from the prefix endpoint, retained even-label boxes,
and continuation boxes alone; the repaired target window or the allowed
branch word must also determine the residual odd-packet matching.
\end{lemma}

\begin{proof}
Consider the two tableaux
\[
  T_A=
  \begin{array}{cccc}
  1&1&4&7\\
  2&2&6&9\\
  3&3&&\\
  4&5&&\\
  5&8&&\\
  6&9&&\\
  7&&&\\
  8&&&
  \end{array}
  \qquad\text{and}\qquad
  T_B=
  \begin{array}{cccc}
  1&1&4&5\\
  2&2&6&9\\
  3&3&&\\
  4&7&&\\
  5&8&&\\
  6&9&&\\
  7&&&\\
  8&&&
  \end{array}.
\]
Both have weakly increasing rows, strictly increasing columns, and content
$(2^9)$.  Their common terminal shape is
\[
  \lambda^{(9)}=(4,4,2,2,2,2,1,1),
  \qquad
  (\lambda^{(9)})'=(8,6,2,2),
\]
so both terminal shapes are column-even.

For both tableaux the prefix endpoint through label $8$ is
\[
  \lambda^{(8)}=(4,3,2,2,2,1,1,1).
\]
The retained tight-prefix labels have the same positions in the two tableaux:
\[
\begin{array}{c|c}
\text{label} & \text{positions}\\ \hline
2 & (2,1),(2,2)\\
4 & (1,3),(4,1)\\
6 & (2,3),(6,1)\\
8 & (5,2),(8,1).
\end{array}
\]
The continuation label also has the same positions,
\[
  9:\ (2,4),(6,2).
\]

The deleted odd labels $1$ and $3$ agree:
\[
  1:\ (1,1),(1,2),\qquad
  3:\ (3,1),(3,2).
\]
But labels $5$ and $7$ differ.  In $T_A$ one has
\[
  5:\ (4,2),(5,1),\qquad
  7:\ (1,4),(7,1),
\]
whereas in $T_B$ one has
\[
  5:\ (1,4),(5,1),\qquad
  7:\ (4,2),(7,1).
\]
Thus the two odd-packet assignments are distinct while the displayed prefix
endpoint, retained even-label boxes, and continuation boxes are identical.

The successive prefix lengths are $1,2,\ldots,8$ in both tableaux.  Hence
\[
  \min(4,\lceil\ell(\lambda^{(t)})/2\rceil)
\]
increases exactly for $t=1,3,5,7$, proving
$H_4(T_A)=H_4(T_B)=\{1,3,5,7\}$.  Since $j=9$ and $m=4$,
\cref{cor:bc-height-zero-surplus-label-prefix} identifies the final tight
zone as the terminal block interval corresponding to labels $1,\ldots,8$.
The final sentence is exactly the displayed collision of packet data.
\end{proof}

\begin{corollary}[Hall data alone do not recover the tight odd packets]\label{cor:bc-hall-data-alone-packet-obstruction}
In the two tableaux of
\cref{lem:bc-tight-even-retained-prefix-obstruction}, the following data agree:
\begin{enumerate}
\item the tight-prefix endpoint $\lambda^{(8)}$;
\item the retained even-label boxes for labels $2,4,6,8$;
\item the labelled continuation boxes;
\item the endpoint-defect pairs of $\lambda^{(8)}$;
\item the continuation incidence relation; and
\item the canonical continuation Hall charge of
\cref{lem:bc-canonical-continuation-hall-charge}.
\end{enumerate}
Nevertheless the deleted odd-packet assignments differ.  Consequently, a
non-column-even tight-zone decoder cannot recover the residual odd-packet
matching from the prefix endpoint, retained even boxes, continuation boxes,
and Hall charge alone; any uniform large-zone carrier rule must use additional
target-window or Step~3 boundary data, or else use the already counted branch
word to select among the remaining packet assignments.
\end{corollary}

\begin{proof}
\Cref{lem:bc-tight-even-retained-prefix-obstruction} states that the two
examples have the same tight-prefix endpoint, the same retained even-label
boxes, and the same continuation boxes, while their deleted odd-packet
assignments differ.

The common prefix endpoint and common continuation boxes determine the same
Pieri continuation segment by
\cref{lem:bc-continuation-boxes-determine-pieri-segment}.  Applying
\cref{lem:bc-pieri-continuation-determines-hall-charge} to this common
segment shows that the endpoint-defect pairs, the continuation incidence
relation, and the canonical continuation Hall charge are also the same in the
two examples.  The different odd-packet assignments are exactly those
displayed in
\cref{lem:bc-tight-even-retained-prefix-obstruction}.
\end{proof}

\begin{corollary}[Direct Hall-selected blocks do not recover the tight odd packets]\label{cor:bc-direct-hall-blocks-alone-packet-obstruction}
In the two tableaux of
\cref{lem:bc-tight-even-retained-prefix-obstruction}, the direct Hall-data
package used in
\cref{cor:bc-shared-direct-hall-carrier-terminal-source} agrees up to the
Hall-selected continuation pieces: the tight-prefix endpoint, endpoint-defect
pairs, continuation-incidence family, canonical Hall charge, and Hall-selected
continuation strips are the same.  Nevertheless the deleted odd-packet
assignments differ.

Consequently the direct Hall data together with the Hall-selected continuation
pieces are not, by themselves, a terminal source for the non-column-even tight
zone.  A uniform large-zone carrier or terminal-boundary rule must use actual
target-window or Step~3 boundary data, or else use the already counted branch
word to select among the remaining packet assignments.
\end{corollary}

\begin{proof}
The agreement of the tight-prefix endpoint, retained even-label boxes, and
labelled continuation boxes is part of
\cref{lem:bc-tight-even-retained-prefix-obstruction}.  As in the proof of
\cref{cor:bc-hall-data-alone-packet-obstruction}, the common prefix endpoint and
common continuation boxes determine the same Pieri continuation segment by
\cref{lem:bc-continuation-boxes-determine-pieri-segment}; hence
\cref{lem:bc-pieri-continuation-determines-hall-charge} gives the same
endpoint-defect pairs, continuation incidences, and canonical Hall charge.

The Hall-selected continuation labels are therefore the same.  Since all
labelled continuation boxes agree, the two-box continuation strip attached to
each selected Hall label also agrees.  The deleted odd-packet assignments are
different by \cref{lem:bc-tight-even-retained-prefix-obstruction}.  Thus no
deterministic rule depending only on the displayed direct Hall package and the
selected continuation pieces can recover those packets.
\end{proof}

\begin{lemma}[Height branch rank slack]\label{lem:bc-height-branch-first-half-slack}
In the height case of \cref{lem:bc-local-branch-recovery}, the local repair
data recovered from a branch word depend only on its last $h_\alpha$ bits.
Equivalently, if
\[
  \beta,\widetilde\beta\in\{0,1\}^{2h_\alpha}
\]
have
\[
  \beta_{h_\alpha+i}=\widetilde\beta_{h_\alpha+i}
  \qquad(1\le i\le h_\alpha),
\]
then the height-case event list, local defect pairs, and selected subset
recovered by \cref{lem:bc-local-branch-recovery} are the same for
$\beta$ and $\widetilde\beta$.  Thus, in a height-zone decoder, the first
$h_\alpha$ branch bits may be read as an auxiliary rank and then replaced by
zeros before applying the local height branch recovery.
\end{lemma}

\begin{proof}
In the height case, item~(1) of \cref{lem:bc-local-branch-recovery} gives
\[
  e_i=2(a_\alpha+i-1)-1
  \qquad(1\le i\le h_\alpha),
\]
which contains no branch bit.  Item~(2) obtains the local defect columns from
the endpoint partitions $\sigma^-$ and $\sigma^+$, again with no branch bit.
Item~(3) reads the selected subset from the last $h_\alpha$ bits of the word.
Therefore two words with the same last half recover the same three pieces of
local repair data.  Replacing the first half by zeros leaves the data used by
the height local inverse unchanged, while making the word one of the
height-branch words of \cref{lem:bc-local-branch-word}.
\end{proof}

\begin{proposition}[Finite odd-packet rank fits the tight height branch budget]\label{prop:bc-tight-odd-packet-rank-criterion}
Work in the non-full tight height setting of
\cref{cor:bc-height-zero-surplus-label-prefix}.  Thus the final tight
right-boundary zone has deleted labels
\[
  1,3,\ldots,2m-1
\]
and retained tight-prefix labels
\[
  2,4,\ldots,2m .
\]
Suppose a parser-visible datum $D(P)$ determines the tight-prefix endpoint
$\lambda^{(2m)}$, the two boxes carrying each retained label
$2,4,\ldots,2m$, and the labelled continuation boxes.  Suppose also that
$D(P)$ determines a canonically ordered finite set
\[
  \Omega(D(P))
\]
of possible assignments of the deleted odd-label singleton packets, that the
true assignment $\omega(P)$ belongs to this set, and that
\[
  |\Omega(D(P))|\le 2^m .
\]
Then the first $m$ bits of the tight-zone branch word
\[
  \beta_\gamma\in\{0,1\}^{2m}
\]
can encode the lexicographic rank of $\omega(P)$ in $\Omega(D(P))$, while the
last $m$ bits remain available for the height local selected-defect data.
From $D(P)$ and $\beta_\gamma$ one recovers the full tight-prefix singleton
packet data.
\end{proposition}

\begin{proof}
Order $\Omega(D(P))$ by the canonical order in the hypothesis.  Since its
cardinality is at most $2^m$, the rank of the true assignment
$\omega(P)$ is encoded by an $m$-bit binary word.  Read this word from the
first half of $\beta_\gamma$ and choose the corresponding element of
\[
  \Omega(D(P)).
\]
This recovers the two boxes of every deleted odd label
$1,3,\ldots,2m-1$.

Together with the retained even-label boxes already contained in $D(P)$,
these recovered odd packets reconstruct the whole Pieri prefix
\[
  \lambda^{(0)}\subset\lambda^{(1)}\subset\cdots\subset\lambda^{(2m)}
\]
by adding the two boxes of labels $1,2,\ldots,2m$ in order.  Hence every
singleton deleted packet endpoint pair, its local defect columns, and the
local paired defects are determined.

The last $m$ bits of $\beta_\gamma$ have not been used in the rank encoding.
By \cref{lem:bc-height-branch-first-half-slack}, the height local recovery
depends on these last bits, together with the recovered endpoint data, exactly
as in \cref{lem:bc-local-branch-recovery}.  Thus the selected defect subset
needed by the height local inverse is still recovered, and the first half of
the branch word has carried only the auxiliary odd-packet rank.  Therefore
$D(P)$ and $\beta_\gamma$ recover the full tight-prefix singleton packet data.
\end{proof}

\begin{definition}[Tight odd-packet compatibility set]\label{def:bc-tight-odd-packet-omega}
Fix $m\ge1$.  A tight-prefix datum consists of a partition $\nu$ of size
$4m$ and, for each $1\le r\le m$, a two-box horizontal strip $E_r$ whose boxes
are intended to carry the retained even label $2r$.  Put
\[
  f_r=(2r-1,1)\qquad(1\le r\le m).
\]
Let $\Omega(\nu;E_1,\ldots,E_m)$ be the set of tuples
\[
  x=(x_1,\ldots,x_m)
\]
of boxes of $\nu$ with the following properties:
\begin{enumerate}
\item the sets
\[
  \{f_r:1\le r\le m\},\qquad
  \{x_r:1\le r\le m\},\qquad
  E_1,\ldots,E_m
\]
are pairwise disjoint and have union the Young diagram of $\nu$;
\item for every $r$, the box $x_r$ is not in column $1$, and the two boxes
$f_r,x_r$ lie in distinct columns;
\item if
\[
  \nu^{(0)}=\varnothing,\qquad
  \nu^{(2r-1)}=\nu^{(2r-2)}\cup\{f_r,x_r\},\qquad
  \nu^{(2r)}=\nu^{(2r-1)}\cup E_r ,
\]
then each $\nu^{(t)}$ is a Young diagram, each displayed difference is a
horizontal strip, and $\nu^{(2m)}=\nu$.
\end{enumerate}
The set $\Omega(\nu;E_1,\ldots,E_m)$ is ordered lexicographically by the
tuples $x_1,\ldots,x_m$ in row-column order.
\end{definition}

\begin{lemma}[The tight compatibility set contains the true odd packets]\label{lem:bc-true-tight-packets-in-omega}
Work in the non-full tight height setting of
\cref{cor:bc-height-zero-surplus-label-prefix}.  Let
\[
  \nu=\lambda^{(2m)}
\]
be the tight-prefix endpoint, and let $E_r$ be the two boxes carrying label
$2r$ for $1\le r\le m$.  For each $r$, let $x_r$ be the second box carrying
the deleted odd label $2r-1$, the first one being the forced first-column box
$f_r=(2r-1,1)$.  Then
\[
  (x_1,\ldots,x_m)\in\Omega(\nu;E_1,\ldots,E_m).
\]
Consequently the compatibility set of
\cref{def:bc-tight-odd-packet-omega} is a canonical parser-visible candidate
set for the residual odd-packet assignment once $\lambda^{(2m)}$ and the
retained even-label boxes are parser-visible.
\end{lemma}

\begin{proof}
By \cref{cor:bc-height-zero-surplus-label-prefix}, the deleted labels in the
tight prefix are exactly $1,3,\ldots,2m-1$, and the retained labels in that
prefix are exactly $2,4,\ldots,2m$.  By
\cref{cor:bc-tight-height-prefix-column-even-gate}, the odd label $2r-1$
creates the row $2r-1$.  Thus one of its two boxes is
$f_r=(2r-1,1)$; since the strip of a single label is horizontal, the other
box $x_r$ is not in column $1$.

The boxes carrying the labels $1,\ldots,2m$ are exactly the boxes of the
prefix endpoint $\lambda^{(2m)}=\nu$.  Each label occurs twice, so the forced
odd first-column boxes, the odd second boxes $x_r$, and the even strips
$E_r$ are pairwise disjoint and have union $\nu$.  The actual Pieri prefix
gives the recursive sequence
\[
  \lambda^{(0)},\lambda^{(1)},\ldots,\lambda^{(2m)} ,
\]
and its successive differences are horizontal two-strips by
\cref{lem:bc-pieri-path-bijection}.  This is exactly item~(3) in
\cref{def:bc-tight-odd-packet-omega} for the tuple $(x_1,\ldots,x_m)$.
Therefore that tuple belongs to $\Omega(\nu;E_1,\ldots,E_m)$.

The final sentence follows because the definition of $\Omega$ uses only the
displayed endpoint partition and retained even strips, together with the fixed
row-column lexicographic order.
\end{proof}

\begin{lemma}[Low-label boxes determine the tight packet datum]\label{lem:bc-low-label-boxes-determine-tight-packet-datum}
Work in the non-full tight height setting of
\cref{cor:bc-height-zero-surplus-label-prefix}, and let
\[
  \varnothing=\lambda^{(0)}\subset\lambda^{(1)}\subset\cdots\subset
  \lambda^{(j)}
\]
be the Pieri two-strip path of the underlying semistandard tableau.  Suppose
a datum determines, for each $1\le a\le2m$, the two boxes carrying
semistandard label $a$.  Then it determines the tight-prefix chain
\[
  \lambda^{(0)}\subset\lambda^{(1)}\subset\cdots\subset\lambda^{(2m)},
\]
the endpoint $\nu=\lambda^{(2m)}$, the retained even strips
\[
  E_r=\lambda^{(2r)}\setminus\lambda^{(2r-1)}
  \qquad(1\le r\le m),
\]
and the true odd-packet assignment
\[
  (x_1,\ldots,x_m)\in\Omega(\nu;E_1,\ldots,E_m).
\]
\end{lemma}

\begin{proof}
For $0\le a\le2m$, the shape $\lambda^{(a)}$ is the union of the boxes
carrying labels at most $a$ by \cref{lem:bc-pieri-path-bijection}.  Hence the
displayed low-label box data determine the whole tight-prefix chain and its
endpoint $\nu=\lambda^{(2m)}$.  The retained even strip $E_r$ is exactly the
two-box set carrying label $2r$, so all labelled retained even strips are also
determined.

It remains to identify the odd packets.  By
\cref{cor:bc-tight-height-prefix-column-even-gate}, the step from
$\lambda^{(2r-2)}$ to $\lambda^{(2r-1)}$ creates row $2r-1$ for
$1\le r\le m$.  Since this step is a horizontal two-strip, one of the two
boxes carrying label $2r-1$ is the forced first-column box
\[
  f_r=(2r-1,1),
\]
and the other box is determined by the same two-box label data; call it
$x_r$.  Applying \cref{lem:bc-true-tight-packets-in-omega} to the endpoint
$\nu$, the retained strips $E_r$, and these odd second boxes gives
\[
  (x_1,\ldots,x_m)\in\Omega(\nu;E_1,\ldots,E_m).
\]
Thus the true tight-prefix packet datum is determined.
\end{proof}

\begin{lemma}[The singleton tight odd packet is forced]\label{lem:bc-singleton-tight-omega-forced}
For $m=1$, every element of
\[
  \Omega(\nu;E_1)
\]
has $x_1=(1,2)$.  In particular
\[
  |\Omega(\nu;E_1)|\le1\le2^1 .
\]
\end{lemma}

\begin{proof}
Let $(x_1)\in\Omega(\nu;E_1)$.  In
\cref{def:bc-tight-odd-packet-omega} one has $f_1=(1,1)$, and item~(2) says
that $x_1$ is not in column $1$.  Item~(3) says that
\[
  \nu^{(1)}=\{(1,1),x_1\}
\]
is a Young diagram and that $\nu^{(1)}\setminus\nu^{(0)}$ is a horizontal
strip.  The only Young diagram with two boxes, containing $(1,1)$ and another
box outside column $1$, is $\{(1,1),(1,2)\}$.  Therefore $x_1=(1,2)$.
The cardinality bound follows immediately; if the displayed set is empty the
same bound is trivial.
\end{proof}

\begin{corollary}[The singleton tight rank leaf is closed]\label{cor:bc-singleton-tight-rank-leaf-closed}
In the setting of \cref{lem:bc-true-tight-packets-in-omega}, suppose $m=1$.
Let $D(P)$ be any parser-visible datum that determines
$\lambda^{(2)}$, the retained strip $E_1$, and the labelled continuation
boxes as in \cref{prop:bc-tight-odd-packet-rank-criterion}.  If one sets
\[
  \Omega(D(P))=\Omega(\lambda^{(2)};E_1).
\]
Then the membership, canonical-ordering, and cardinality hypotheses of
\cref{prop:bc-tight-odd-packet-rank-criterion} hold.  Thus the singleton
tight right-boundary zone needs no refined outside-data ambiguity bound for
its deleted odd packet.
\end{corollary}

\begin{proof}
\Cref{lem:bc-true-tight-packets-in-omega} puts the true deleted odd-packet
assignment in $\Omega(\lambda^{(2)};E_1)$, while
\cref{lem:bc-singleton-tight-omega-forced} gives
$|\Omega(\lambda^{(2)};E_1)|\le2^1$.  The lexicographic order from
\cref{def:bc-tight-odd-packet-omega} is canonical.  These are precisely the
membership, ordering, and cardinality hypotheses of
\cref{prop:bc-tight-odd-packet-rank-criterion}.
\end{proof}

\begin{lemma}[The first two tight odd packets have binary prefix ambiguity]\label{lem:bc-two-packet-tight-omega-binary}
Fix $m\ge2$.  Let
\[
  \Pi_2(\nu;E_1,\ldots,E_m)
\]
be the set of pairs $(x_1,x_2)$ which occur as the first two coordinates of
an element of $\Omega(\nu;E_1,\ldots,E_m)$.  Then
\[
  |\Pi_2(\nu;E_1,\ldots,E_m)|\le2 .
\]
In particular, for $m=2$,
\[
  |\Omega(\nu;E_1,E_2)|\le2\le2^2 .
\]
\end{lemma}

\begin{proof}
Let $(x_1,\ldots,x_m)\in\Omega(\nu;E_1,\ldots,E_m)$.  Applying the proof of
\cref{lem:bc-singleton-tight-omega-forced} to the first odd step gives
\[
  x_1=(1,2),\qquad \nu^{(1)}=(2).
\]

The diagram $\nu^{(3)}$ contains the forced cell $f_2=(3,1)$.  Hence, by the
Young property, it also contains $(2,1)$.  This cell is not in $\nu^{(1)}$,
is not $f_2$, and cannot be $x_2$ because $x_2$ is not in column $1$.
Therefore $(2,1)\in E_1$.

Write $E_1=\{(2,1),y\}$.  Since $E_1$ is a horizontal strip, $y$ is not in
column $1$.  Since
\[
  \nu^{(2)}=(2)\cup E_1
\]
is a Young diagram with four cells, the only possibilities are
\[
  y=(1,3)\qquad\text{or}\qquad y=(2,2).
\]
Indeed, a cell in the first row must be $(1,3)$, a cell in the second row
must be $(2,2)$, and any lower-row or farther-right cell would force an
additional missing predecessor cell.

If $y=(1,3)$, then $\nu^{(2)}=(3,1)$.  The Young diagram
\[
  \nu^{(3)}=\nu^{(2)}\cup\{(3,1),x_2\}
\]
can then have $x_2$ only at $(1,4)$ or $(2,2)$: these give the two Young
diagrams $(4,1,1)$ and $(3,2,1)$, while every other cell outside
$\nu^{(2)}$ and outside column $1$ leaves a missing predecessor.

If $y=(2,2)$, then $\nu^{(2)}=(2,2)$.  The same Young-diagram check gives
\[
  x_2\in\{(1,3),(3,2)\},
\]
corresponding to the two shapes $(3,2,1)$ and $(2,2,2)$.

Thus, for the fixed datum, $x_1$ is forced and $x_2$ has at most two possible
values among all compatible tuples.  Hence
\[
  |\Pi_2(\nu;E_1,\ldots,E_m)|\le2 .
\]

When $m=2$, an element of $\Omega(\nu;E_1,E_2)$ has no coordinates beyond
$(x_1,x_2)$, so the same bound gives
$|\Omega(\nu;E_1,E_2)|\le2$.
\end{proof}

\begin{lemma}[Binary prefix factorial bound]\label{lem:bc-binary-prefix-factorial-bound}
For every $m\ge2$,
\[
  |\Omega(\nu;E_1,\ldots,E_m)|\le 2\,(m-2)! .
\]
\end{lemma}

\begin{proof}
Fix a pair
\[
  (a_1,a_2)\in\Pi_2(\nu;E_1,\ldots,E_m).
\]
Consider an element $x=(x_1,\ldots,x_m)$ of
$\Omega(\nu;E_1,\ldots,E_m)$ with
\[
  (x_1,x_2)=(a_1,a_2).
\]
The union condition in
\cref{def:bc-tight-odd-packet-omega} forces the remaining cells
$x_3,\ldots,x_m$ to be exactly the elements of
\[
  U(a_1,a_2)=
  \nu\setminus
  \left(
    \{f_1,\ldots,f_m,a_1,a_2\}\cup E_1\cup\cdots\cup E_m
  \right).
\]
This set has cardinality $m-2$: the diagram $\nu$ has $4m$ cells, the forced
first-column cells contribute $m$ cells, the fixed even strips contribute
$2m$ cells, and the fixed prefix cells $a_1,a_2$ contribute two cells.
Therefore, after the first two coordinates are fixed, there are at most
$(m-2)!$ possible ordered tuples $(x_3,\ldots,x_m)$.

There are at most two possible first-two-coordinate pairs by
\cref{lem:bc-two-packet-tight-omega-binary}.  Multiplying gives the displayed
bound.
\end{proof}

\begin{corollary}[The two-packet tight rank leaf is closed]\label{cor:bc-two-packet-tight-rank-leaf-closed}
In the setting of \cref{lem:bc-true-tight-packets-in-omega}, suppose $m=2$.
Let $D(P)$ be any parser-visible datum that determines
$\lambda^{(4)}$, the retained strips $E_1,E_2$, and the labelled continuation
boxes as in \cref{prop:bc-tight-odd-packet-rank-criterion}.  If one sets
\[
  \Omega(D(P))=\Omega(\lambda^{(4)};E_1,E_2),
\]
then the membership, canonical-ordering, and cardinality hypotheses of
\cref{prop:bc-tight-odd-packet-rank-criterion} hold.  Thus the two-packet
tight right-boundary zone needs no refined outside-data ambiguity bound for
its deleted odd packets.
\end{corollary}

\begin{proof}
\Cref{lem:bc-true-tight-packets-in-omega} puts the true deleted odd-packet
assignment in the displayed set.  Meanwhile
\cref{lem:bc-two-packet-tight-omega-binary} gives
\[
  |\Omega(\lambda^{(4)};E_1,E_2)|\le2\le2^2 .
\]
The lexicographic order from
\cref{def:bc-tight-odd-packet-omega} is canonical.  These are precisely the
membership, ordering, and cardinality hypotheses of
\cref{prop:bc-tight-odd-packet-rank-criterion}.
\end{proof}

\begin{corollary}[Small tight rank leaves are closed]\label{cor:bc-small-tight-rank-leaves-closed}
In the setting of \cref{lem:bc-true-tight-packets-in-omega}, suppose
$2\le m\le6$.
Let $D(P)$ be any parser-visible datum that determines
$\lambda^{(2m)}$, the retained strips $E_1,\ldots,E_m$, and the labelled
continuation boxes as in \cref{prop:bc-tight-odd-packet-rank-criterion}.  If
one sets
\[
  \Omega(D(P))=\Omega(\lambda^{(2m)};E_1,\ldots,E_m),
\]
then the membership, canonical-ordering, and cardinality hypotheses of
\cref{prop:bc-tight-odd-packet-rank-criterion} hold.  Thus tight
right-boundary zones with at most six deleted odd packets need no refined
outside-data ambiguity bound for
their deleted odd packets.
\end{corollary}

\begin{proof}
The true assignment belongs to the displayed set by
\cref{lem:bc-true-tight-packets-in-omega}, and its order is the canonical
lexicographic order of \cref{def:bc-tight-odd-packet-omega}.
For the size bound, \cref{lem:bc-binary-prefix-factorial-bound} gives
\[
  |\Omega(\lambda^{(2m)};E_1,\ldots,E_m)|\le 2\,(m-2)! .
\]
For $2\le m\le6$ one checks directly that $2(m-2)!\le2^m$.  These are exactly
the membership, ordering, and cardinality hypotheses of
\cref{prop:bc-tight-odd-packet-rank-criterion}.
\end{proof}

\begin{proposition}[The seven-packet tight rank leaf is certified]\label{prop:bc-seven-packet-tight-rank-leaf-certificate}
In the setting of \cref{lem:bc-true-tight-packets-in-omega}, suppose $m=7$.
Let $D(P)$ be any parser-visible datum that determines
$\lambda^{(14)}$, the retained strips $E_1,\ldots,E_7$, and the labelled
continuation boxes as in \cref{prop:bc-tight-odd-packet-rank-criterion}.  If
one sets
\[
  \Omega(D(P))=\Omega(\lambda^{(14)};E_1,\ldots,E_7),
\]
then the membership, canonical-ordering, and cardinality hypotheses of
\cref{prop:bc-tight-odd-packet-rank-criterion} hold.  Thus tight
right-boundary zones with seven deleted odd packets need no refined
outside-data ambiguity bound for their deleted odd packets.
\end{proposition}

\begin{proof}
\Cref{lem:bc-true-tight-packets-in-omega} gives membership of the true
assignment, and the ordering is the canonical lexicographic order of
\cref{def:bc-tight-odd-packet-omega}.  It remains only to check the
cardinality bound.

The exact C++ verifier
\[
\texttt{character\_ring\_iter/post\_m29\_bc\_tight\_omega\_small.cpp}
\]
enumerates every tight-prefix datum by the recursive horizontal-two-strip
definition of \cref{def:bc-tight-odd-packet-omega}: at each odd step it adds
the forced cell $f_r=(2r-1,1)$ and one non-first-column cell, and at each even
step it adds an arbitrary horizontal two-strip $E_r$.  It groups the resulting
tuples by the pair
\[
  (\lambda^{(2m)};E_1,\ldots,E_m)
\]
and records the maximum fiber size.

The archived `optimus' run
\[
\texttt{certificates/post\_m29/bc\_tight\_omega\_small\_m7\_certificate.log}
\]
has SHA-256
\[
\hashsplit{24063143ad7f4ac92259560dd11b1d8c}{5291c544e92656e119b8b1a911b6df7d}.
\]
The source file has SHA-256
\[
\hashsplit{c6f03e177fbfcc9a3e7c055a10582445}{6e27124bc915e2c20d8e437da1aff7b7}.
\]
The log ends with \texttt{\_\_EXIT\_STATUS=0} and reports
\[
  \begin{array}{c|c|c|c|c}
  m&\texttt{states}&\texttt{total\_paths}&\texttt{max\_omega}&\texttt{pow2}\\
  \hline
  7&3160108&7565247&20&128 .
  \end{array}
\]
Thus every $m=7$ tight compatibility set has size at most $20\le2^7$.  These
are exactly the membership, ordering, and cardinality hypotheses of
\cref{prop:bc-tight-odd-packet-rank-criterion}.
\end{proof}

\begin{proposition}[The eight-packet tight rank leaf is certified]\label{prop:bc-eight-packet-tight-rank-leaf-certificate}
In the setting of \cref{lem:bc-true-tight-packets-in-omega}, suppose $m=8$.
Let $D(P)$ be any parser-visible datum that determines
$\lambda^{(16)}$, the retained strips $E_1,\ldots,E_8$, and the labelled
continuation boxes as in \cref{prop:bc-tight-odd-packet-rank-criterion}.  If
one sets
\[
  \Omega(D(P))=\Omega(\lambda^{(16)};E_1,\ldots,E_8),
\]
then the membership, canonical-ordering, and cardinality hypotheses of
\cref{prop:bc-tight-odd-packet-rank-criterion} hold.  Thus tight
right-boundary zones with eight deleted odd packets need no refined
outside-data ambiguity bound for their deleted odd packets.
\end{proposition}

\begin{proof}
\Cref{lem:bc-true-tight-packets-in-omega} gives membership of the true
assignment, and the ordering is the canonical lexicographic order of
\cref{def:bc-tight-odd-packet-omega}.  It remains only to check the
cardinality bound.

The exact C++ verifier
\[
\texttt{character\_ring\_iter/post\_m29\_bc\_tight\_omega\_packed.cpp}
\]
uses the same recursive enumeration as
\cref{prop:bc-seven-packet-tight-rank-leaf-certificate}, but stores the fixed
even strips in packed integer keys and expands each level with OpenMP on
`optimus'.  It groups the resulting tuples by
\[
  (\lambda^{(2m)};E_1,\ldots,E_m)
\]
and records the maximum fiber size.

The archived run
\[
\texttt{certificates/post\_m29/bc\_tight\_omega\_packed\_m8\_certificate.log}
\]
has SHA-256
\[
\hashsplit{1ab7eafd9308a50dc2ea4c2a946e2a9c}{69514c2661a7166356c410fb818e96e7}.
\]
The source file has SHA-256
\[
\hashsplit{2fe5f8caa1f61536aaebe97c2f0b215f}{fa7a55b7bc22c3f7b5441de42a1e9346}.
\]
The log ends with \texttt{\_\_EXIT\_STATUS=0} and reports
\[
  \begin{array}{c|c|c|c|c}
  m&\texttt{states}&\texttt{total\_paths}&\texttt{max\_omega}&\texttt{pow2}\\
  \hline
  8&44051883&152106391&45&256 .
  \end{array}
\]
Thus every $m=8$ tight compatibility set has size at most $45\le2^8$.
These are exactly the membership, ordering, and cardinality hypotheses of
\cref{prop:bc-tight-odd-packet-rank-criterion}.
\end{proof}

\begin{corollary}[The tight odd-packet ambiguity bound closes the rank leaf]\label{cor:bc-tight-omega-bound-closes-rank-leaf}
In the setting of \cref{lem:bc-true-tight-packets-in-omega}, suppose that the
canonical compatibility set satisfies
\[
  |\Omega(\lambda^{(2m)};E_1,\ldots,E_m)|\le2^m .
\]
Then the hypotheses of
\cref{prop:bc-tight-odd-packet-rank-criterion} hold for the residual
odd-packet assignment in the final tight right-boundary zone.
\end{corollary}

\begin{proof}
Take $D(P)$ to be the parser-visible datum containing
$\lambda^{(2m)}$, the retained even-label boxes $E_1,\ldots,E_m$, and the
labelled continuation boxes.  Let
\[
  \Omega(D(P))=\Omega(\lambda^{(2m)};E_1,\ldots,E_m)
\]
with the lexicographic order of \cref{def:bc-tight-odd-packet-omega}.
\Cref{lem:bc-true-tight-packets-in-omega} proves that the true residual
odd-packet assignment belongs to this set.  The displayed cardinality bound
is exactly the remaining size hypothesis of
\cref{prop:bc-tight-odd-packet-rank-criterion}.
\end{proof}

\begin{corollary}[Refined tight-packet classes are the live rank criterion]\label{cor:bc-refined-tight-packet-criterion}
Work in the setting of \cref{lem:bc-true-tight-packets-in-omega}.  Suppose a
parser-visible auxiliary datum $A(P)$, recovered without additional branch bits
from the canonical halo/right-boundary parser data already counted in the
replacement-zone package, determines a canonically ordered finite set
\[
  \Omega_A(A(P))
  \subseteq
  \Omega(\lambda^{(2m)};E_1,\ldots,E_m)
\]
such that the true residual odd-packet assignment belongs to
$\Omega_A(A(P))$ and
\[
  |\Omega_A(A(P))|\le2^m .
\]
Then the first $m$ bits of the tight-zone branch word recover the true
residual odd-packet assignment, and the last $m$ bits remain available for the
height selected-defect data.  Consequently the rank hypothesis of
\cref{prop:bc-tight-odd-packet-rank-criterion} holds for this refined
parser-visible datum.
\end{corollary}

\begin{proof}
Use as $D(P)$ the parser-visible datum consisting of
$\lambda^{(2m)}$, the retained even-label boxes $E_1,\ldots,E_m$, the
labelled continuation boxes, and $A(P)$, and put
\[
  \Omega(D(P))=\Omega_A(A(P)).
\]
The hypotheses of the present corollary are exactly the membership,
canonical-ordering, and cardinality hypotheses required in
\cref{prop:bc-tight-odd-packet-rank-criterion}.  That proposition therefore
gives the claimed recovery and leaves the last half of the branch word for
the height selected-defect data.
\end{proof}

\begin{proposition}[The refined $m\ge9$ suffix is the remaining tight rank input]\label{prop:bc-refined-tight-suffix-reduces-rank}
Work in the non-full tight height setting of
\cref{cor:bc-height-zero-surplus-label-prefix}.  Suppose that, whenever the
final tight right-boundary zone has $m\ge9$ deleted singleton packets, the
branch-free canonical halo/right-boundary parser data determine an
auxiliary datum $A(P)$ and a canonically ordered set
\[
  \Omega_A(A(P))
  \subseteq
  \Omega(\lambda^{(2m)};E_1,\ldots,E_m)
\]
such that the true residual odd-packet assignment belongs to
$\Omega_A(A(P))$ and
\[
  |\Omega_A(A(P))|\le2^m .
\]
Then the tight odd-packet rank hypothesis of
\cref{prop:bc-tight-odd-packet-rank-criterion} holds for every non-full tight
height right-boundary zone.
\end{proposition}

\begin{proof}
Let $m$ be the number of deleted singleton packets in the final tight
right-boundary zone.  If $m=1$, the assertion is
\cref{cor:bc-singleton-tight-rank-leaf-closed}.  If $2\le m\le6$, it is
\cref{cor:bc-small-tight-rank-leaves-closed}.  If $m=7$ or $m=8$, it is,
respectively,
\cref{prop:bc-seven-packet-tight-rank-leaf-certificate} and
\cref{prop:bc-eight-packet-tight-rank-leaf-certificate}.  In the remaining
case $m\ge9$, the hypothesis of the present proposition is exactly the
hypothesis of \cref{cor:bc-refined-tight-packet-criterion}, so that corollary
applies.
\end{proof}

\begin{definition}[Fixed-cell completions]\label{def:bc-tight-fixed-cell-completions}
For a tight-prefix datum $(\nu;E_1,\ldots,E_m)$ as in
\cref{def:bc-tight-odd-packet-omega}, let
\[
  U(\nu;E)=
  \nu\setminus
  \left(
    \{f_r:1\le r\le m\}\cup E_1\cup\cdots\cup E_m
  \right)
\]
be the set of unfilled cells.  Let
\[
  \mathcal C(\nu;E_1,\ldots,E_m)
\]
be the set of bijections
\[
  \tau:U(\nu;E)\longrightarrow\{1,\ldots,m\}
\]
such that the filling of $\nu$ defined by
\[
  f_r\mapsto 2r-1,\qquad
  \tau^{-1}(r)\mapsto 2r-1,\qquad
  E_r\mapsto 2r
\]
is semistandard: rows are weakly increasing from left to right and columns are
strictly increasing from top to bottom.
\end{definition}

\begin{lemma}[Omega equals completions]\label{lem:bc-omega-fixed-cell-completions}
For every tight-prefix datum, the map
\[
  (x_1,\ldots,x_m)\longmapsto
  \bigl(\tau(x_r)=r\bigr)_{r=1}^m
\]
is a bijection
\[
  \Omega(\nu;E_1,\ldots,E_m)
  \longleftrightarrow
  \mathcal C(\nu;E_1,\ldots,E_m).
\]
Consequently
\[
  |\Omega(\nu;E_1,\ldots,E_m)|
  =
  |\mathcal C(\nu;E_1,\ldots,E_m)|.
\]
\end{lemma}

\begin{proof}
Let $x=(x_1,\ldots,x_m)\in\Omega(\nu;E_1,\ldots,E_m)$.  By item~(1) of
\cref{def:bc-tight-odd-packet-omega}, the cells $x_1,\ldots,x_m$ are exactly
the cells of $U(\nu;E)$, so the displayed rule defines a bijection
$\tau_x:U(\nu;E)\to\{1,\ldots,m\}$.  The recursive chain
\[
  \nu^{(0)}\subset\nu^{(1)}\subset\cdots\subset\nu^{(2m)}=\nu
\]
in item~(3) is the sequence of shapes formed by the cells whose assigned
labels are at most $1,2,\ldots,2m$.  The successive differences are
horizontal strips.  By the Pieri-path equivalence
\cref{lem:bc-pieri-path-bijection}, the corresponding filling is
semistandard.  Hence $\tau_x\in\mathcal C(\nu;E_1,\ldots,E_m)$.

Conversely, let $\tau\in\mathcal C(\nu;E_1,\ldots,E_m)$ and write
$x_r=\tau^{-1}(r)$.  The definition of $U(\nu;E)$ gives the disjoint-union
condition in item~(1) of \cref{def:bc-tight-odd-packet-omega}.  Since
$f_r$ and $x_r$ carry the same label $2r-1$ in a semistandard filling, they
cannot lie in the same column; by definition $x_r$ is not one of the forced
first-column cells.  This gives item~(2).  Finally, applying
\cref{lem:bc-pieri-path-bijection} to the semistandard filling shows that the
sublevel shapes for labels $1,\ldots,2m$ are Young diagrams and that each
successive difference is the required horizontal two-strip.  Thus
$(x_1,\ldots,x_m)\in\Omega(\nu;E_1,\ldots,E_m)$.

The two constructions are inverse because $x_r$ is, in both directions, the
unique unfilled cell carrying the odd label $2r-1$.  The cardinality identity
follows.
\end{proof}

\begin{corollary}[Fixed-cell completion bound implies the unrefined Omega bound]\label{cor:bc-fixed-cell-bound-closes-omega}
If every tight-prefix datum satisfies
\[
  |\mathcal C(\nu;E_1,\ldots,E_m)|\le2^m ,
\]
then every tight-prefix datum satisfies
\[
  |\Omega(\nu;E_1,\ldots,E_m)|\le2^m .
\]
\end{corollary}

\begin{proof}
This is immediate from \cref{lem:bc-omega-fixed-cell-completions}.
\end{proof}

\begin{lemma}[The first column is forced in any fixed-cell completion]\label{lem:bc-fixed-cell-first-column-forced}
Let $(\nu;E_1,\ldots,E_m)$ be a tight-prefix datum with
$\mathcal C(\nu;E_1,\ldots,E_m)\ne\varnothing$.  In every completion, the
first-column cell in row $i$ carries the label $i$.  Consequently the height
of $\nu$ is either $2m-1$ or $2m$; for $1\le r<m$ one has
\[
  (2r,1)\in E_r,
\]
and, if $(2m,1)$ belongs to $\nu$, then $(2m,1)\in E_m$.
\end{lemma}

\begin{proof}
Fix a completion and let $T$ be its semistandard filling.  Since
$f_m=(2m-1,1)$ belongs to $\nu$, the first column has height at least
$2m-1$.  The entries in a column of a semistandard tableau are strictly
increasing from top to bottom, and the forced cells satisfy
\[
  T(2r-1,1)=2r-1\qquad(1\le r\le m).
\]
For $1\le r<m$, the cell $(2r,1)$ lies between the forced cells
$(2r-1,1)$ and $(2r+1,1)$ in the same column.  Its label is therefore
strictly between $2r-1$ and $2r+1$, hence is exactly $2r$.  The only cells
with label $2r$ are the two cells of $E_r$, so $(2r,1)\in E_r$.

If a row $2m$ is present, the first-column entry in that row is strictly
larger than $2m-1$.  Since all labels are at most $2m$, it is $2m$, and so
$(2m,1)\in E_m$.  A row below $2m$ would require a first-column label larger
than $2m$, impossible.  Thus the height is either $2m-1$ or $2m$, and the
first-column labels are exactly $1,2,\ldots,h$.
\end{proof}

\begin{definition}[Reduced fixed-even growth chains]\label{def:bc-reduced-fixed-even-growth-chains}
Assume that
$\mathcal C(\nu;E_1,\ldots,E_m)\ne\varnothing$.  For a non-first-column box
$c=(i,j)$ with $j>1$, write
\[
  \overline c=(i,j-1).
\]
Let $\overline\nu$ be the diagram obtained from $\nu$ by deleting the first
column, and put
\[
  \overline E_r=\{\overline c:\ c\in E_r,\ c\ne(2r,1)\}
  \qquad(1\le r\le m).
\]
Let $\varepsilon_r=1$ if $(2r,1)\in E_r$, and
$\varepsilon_r=0$ otherwise.  By
\cref{lem:bc-fixed-cell-first-column-forced}, $\varepsilon_r=1$ for
$r<m$, while $\varepsilon_m\in\{0,1\}$.

The set
\[
  \mathcal R(\overline\nu;\overline E_1,\ldots,\overline E_m)
\]
consists of tuples
\[
  y=(y_1,\ldots,y_m)
\]
of boxes of $\overline\nu$ such that the singleton sets
$\{y_r\}$ and the fixed sets $\overline E_r$ are pairwise disjoint and have
union $\overline\nu$, and such that the recursively defined sets
\[
  \alpha_0=\varnothing,\qquad
  \alpha_{r-\frac12}=\alpha_{r-1}\cup\{y_r\},\qquad
  \alpha_r=\alpha_{r-\frac12}\cup\overline E_r
\]
are Young diagrams for every $r$.  In addition,
$\alpha_{r-\frac12}$ has at most $2r-1$ rows, $\alpha_r$ has at most
$2r-1+\varepsilon_r$ rows, and each difference
$\alpha_r\setminus\alpha_{r-\frac12}$ is a horizontal strip.  The set is
ordered lexicographically by the tuple $y_1,\ldots,y_m$ in row-column order.
\end{definition}

\begin{proposition}[Fixed-cell completions equal reduced growth chains]\label{prop:bc-completions-equal-reduced-growth-chains}
For every tight-prefix datum with
$\mathcal C(\nu;E_1,\ldots,E_m)\ne\varnothing$, the map
\[
  \tau\longmapsto
  \bigl(\overline{\tau^{-1}(1)},\ldots,\overline{\tau^{-1}(m)}\bigr)
\]
is a bijection
\[
  \mathcal C(\nu;E_1,\ldots,E_m)
  \longleftrightarrow
  \mathcal R(\overline\nu;\overline E_1,\ldots,\overline E_m).
\]
Consequently
\[
  |\mathcal C(\nu;E_1,\ldots,E_m)|
  =
  |\mathcal R(\overline\nu;\overline E_1,\ldots,\overline E_m)|.
\]
\end{proposition}

\begin{proof}
Let $\tau$ be a fixed-cell completion and let $T_\tau$ be the corresponding
semistandard filling.  For each $r$, the cell $\tau^{-1}(r)$ is off the first
column: otherwise the first column would contain two cells with the same label
$2r-1$, namely $f_r$ and $\tau^{-1}(r)$, contradicting strict column
increase.  By \cref{lem:bc-pieri-path-bijection}, the sublevel
sets of labels $1,\ldots,2m$ form a Pieri growth path.  By
\cref{lem:bc-fixed-cell-first-column-forced}, the first-column cells in that
path are exactly the cells with labels $1,2,\ldots,h$, where
$h\in\{2m-1,2m\}$.  Deleting the first column from each sublevel diagram
therefore leaves a sequence of Young diagrams.  The odd step $2r-1$ loses
the forced first-column cell $f_r$ and retains exactly the single box
$\overline{\tau^{-1}(r)}$; the even step $2r$ retains exactly the off-column
fixed strip $\overline E_r$.  Hence the tuple displayed in the statement
belongs to
$\mathcal R(\overline\nu;\overline E_1,\ldots,\overline E_m)$.

Conversely, let
$y\in\mathcal R(\overline\nu;\overline E_1,\ldots,\overline E_m)$.  Lift
each $y_r=(i,j)$ to the box $x_r=(i,j+1)$ of $\nu$.  Start with the empty
diagram.  At odd step $2r-1$, adjoin $f_r=(2r-1,1)$ and $x_r$; at even step
$2r$, adjoin the two boxes of $E_r$.  The row-support conditions in
\cref{def:bc-reduced-fixed-even-growth-chains} say exactly that, after the
appropriate first column of height $2r-1$ or $2r-1+\varepsilon_r$ is restored,
the reduced diagrams $\alpha_{r-\frac12}$ and $\alpha_r$ lift to Young
diagrams.  The odd differences are horizontal because $x_r$ is off the first
column, and the even differences are horizontal by the last condition in the
definition of $\mathcal R$.  Thus the lifted sequence is a Pieri path ending
at $\nu$.

Applying \cref{lem:bc-pieri-path-bijection} to this Pieri path gives a
semistandard filling of $\nu$ in which $f_r$ and $x_r$ carry label $2r-1$
and the cells of $E_r$ carry label $2r$.  This filling defines an element of
$\mathcal C(\nu;E_1,\ldots,E_m)$ by setting $\tau(x_r)=r$.  The two
constructions are inverse, since deleting and restoring the first column
does not change any off-column box.  The cardinality identity follows.
\end{proof}

\begin{corollary}[Reduced growth-chain bound closes the fixed-cell leaf]\label{cor:bc-reduced-growth-chain-bound-closes-fixed-cell}
Suppose that every tight-prefix datum with
$\mathcal C(\nu;E_1,\ldots,E_m)\ne\varnothing$ satisfies
\[
  |\mathcal R(\overline\nu;\overline E_1,\ldots,\overline E_m)|\le2^m .
\]
Then every tight-prefix datum satisfies
\[
  |\mathcal C(\nu;E_1,\ldots,E_m)|\le2^m .
\]
\end{corollary}

\begin{proof}
If $\mathcal C(\nu;E_1,\ldots,E_m)=\varnothing$, the conclusion is
immediate.  Otherwise it follows from
\cref{prop:bc-completions-equal-reduced-growth-chains}.
\end{proof}

\begin{lemma}[Only the terminal reduced even strip may have size two]\label{lem:bc-reduced-even-strip-sizes}
For every reduced fixed-even datum arising from a nonempty fixed-cell
completion set,
\[
  |\overline E_r|=1\qquad(1\le r<m),
\]
and
\[
  |\overline E_m|=2-\varepsilon_m\in\{1,2\}.
\]
\end{lemma}

\begin{proof}
Each original even strip $E_r$ has exactly two boxes.  By
\cref{lem:bc-fixed-cell-first-column-forced}, $(2r,1)\in E_r$ for
$1\le r<m$.  Deleting the first column therefore removes exactly one box of
$E_r$, leaving $|\overline E_r|=1$.

For $r=m$, the first-column box $(2m,1)$ belongs to $E_m$ exactly when
$\varepsilon_m=1$.  Hence deleting the first column removes
$\varepsilon_m$ boxes from the two-box strip $E_m$, leaving
$|\overline E_m|=2-\varepsilon_m$.
\end{proof}

\begin{definition}[Reverse reduced corner tree]\label{def:bc-reduced-reverse-corner-tree}
Fix a reduced fixed-even datum
$(\overline\nu;\overline E_1,\ldots,\overline E_m)$ arising from a nonempty
fixed-cell completion set, with $\varepsilon_r$ as in
\cref{def:bc-reduced-fixed-even-growth-chains}.  A reverse reduced corner path
is a sequence
\[
  \overline\nu=\beta_m\supset\beta_{m-1}\supset\cdots\supset\beta_0=\varnothing
\]
together with boxes $y_m,\ldots,y_1$ such that, for each $r$,
\[
  \beta_{r-\frac12}:=\beta_r\setminus\overline E_r
\]
is a Young diagram, $\overline E_r$ has at most one box in each column of
$\beta_r\setminus\beta_{r-\frac12}$,
\[
  \ell(\beta_{r-\frac12})\le2r-1,\qquad
  \ell(\beta_r)\le2r-1+\varepsilon_r,
\]
the box $y_r$ is a removable corner of
$\beta_{r-\frac12}$, and
\[
  \beta_{r-1}=\beta_{r-\frac12}\setminus\{y_r\}.
\]
Let
\[
  \mathcal T(\overline\nu;\overline E_1,\ldots,\overline E_m)
\]
be the set of such reverse reduced corner paths.
\end{definition}

\begin{lemma}[Reduced chains are reverse corner paths]\label{lem:bc-reduced-growth-reverse-corner-tree}
For every reduced fixed-even datum arising from a nonempty fixed-cell
completion set, the assignment sending
\[
  y=(y_1,\ldots,y_m)\in
  \mathcal R(\overline\nu;\overline E_1,\ldots,\overline E_m)
\]
to the reverse sequence
\[
  \beta_r=\alpha_r(y)\qquad(0\le r\le m)
\]
is a bijection
\[
  \mathcal R(\overline\nu;\overline E_1,\ldots,\overline E_m)
  \longleftrightarrow
  \mathcal T(\overline\nu;\overline E_1,\ldots,\overline E_m).
\]
Consequently
\[
  |\mathcal R(\overline\nu;\overline E_1,\ldots,\overline E_m)|
  =
  |\mathcal T(\overline\nu;\overline E_1,\ldots,\overline E_m)|.
\]
\end{lemma}

\begin{proof}
Let
$y\in\mathcal R(\overline\nu;\overline E_1,\ldots,\overline E_m)$, and let
$\alpha_0,\alpha_{1/2},\alpha_1,\ldots,\alpha_m$ be the forward diagrams in
\cref{def:bc-reduced-fixed-even-growth-chains}.  Put $\beta_r=\alpha_r$.
Since
\[
  \alpha_r=\alpha_{r-\frac12}\cup\overline E_r
\]
and $\alpha_r\setminus\alpha_{r-\frac12}$ is a horizontal strip, deleting
$\overline E_r$ from $\beta_r$ gives the Young diagram
$\beta_{r-\frac12}=\alpha_{r-\frac12}$ and removes at most one box from each
column.  Since
\[
  \alpha_{r-\frac12}=\alpha_{r-1}\cup\{y_r\}
\]
is a Young diagram, deleting $y_r$ from it gives the Young diagram
$\alpha_{r-1}$; hence $y_r$ is a removable corner of
$\beta_{r-\frac12}$.  Thus the reversed data define an element of
$\mathcal T$.

Conversely, a reverse reduced corner path determines boxes
$y_1,\ldots,y_m$ by its removals.  Reversing the displayed recurrence gives
\[
  \alpha_{r-\frac12}=\alpha_{r-1}\cup\{y_r\},\qquad
  \alpha_r=\alpha_{r-\frac12}\cup\overline E_r,
\]
and the hypotheses in
\cref{def:bc-reduced-reverse-corner-tree} say exactly that these sets are
Young diagrams, satisfy the same row-support bounds, and that the fixed even
difference is a horizontal strip.
The disjoint-union condition follows because the reverse path removes every
box of $\overline\nu$ exactly once, either as one of the fixed sets
$\overline E_r$ or as one of the corners $y_r$.  Hence the tuple belongs to
$\mathcal R(\overline\nu;\overline E_1,\ldots,\overline E_m)$.

The two constructions are inverse term by term, so the cardinality identity
follows.
\end{proof}

\begin{corollary}[Reverse corner tree bound closes the reduced leaf]\label{cor:bc-reverse-corner-tree-bound-closes-reduced-leaf}
Suppose that every reduced fixed-even datum arising from a nonempty
fixed-cell completion set satisfies
\[
  |\mathcal T(\overline\nu;\overline E_1,\ldots,\overline E_m)|\le2^m .
\]
Then every such datum satisfies
\[
  |\mathcal R(\overline\nu;\overline E_1,\ldots,\overline E_m)|\le2^m .
\]
\end{corollary}

\begin{proof}
This is immediate from
\cref{lem:bc-reduced-growth-reverse-corner-tree}.
\end{proof}

\begin{definition}[Odd fixed-even standard tableaux]\label{def:bc-odd-fixed-even-standard-tableaux}
Let
$(\overline\nu;\overline E_1,\ldots,\overline E_m)$ be a reduced fixed-even
datum arising from a nonempty fixed-cell completion set.  Put
\[
  \gamma=\overline\nu\setminus\overline E_m .
\]
By nonemptiness and \cref{def:bc-reduced-reverse-corner-tree}, $\gamma$ is a
Young diagram.  For $1\le r<m$, let $e_r$ be the unique cell of
$\overline E_r$, whose existence is supplied by
\cref{lem:bc-reduced-even-strip-sizes}.

Let
\[
  \mathcal S(\gamma;e_1,\ldots,e_{m-1})
\]
be the set of standard Young tableaux of shape $\gamma$ with entries
$1,\ldots,2m-1$ such that the cell $e_r$ carries entry $2r$ for every
$1\le r<m$.

The same notation will also be used for any Young diagram $\delta$ of size
$2r-1$ together with distinct cells $f_1,\ldots,f_{r-1}\in\delta$: then
$\mathcal S(\delta;f_1,\ldots,f_{r-1})$ denotes the set of standard Young
tableaux of shape $\delta$ with $2s$ in $f_s$ for $1\le s<r$.  If the fixed
entries are incompatible with standardness, this set is empty.
\end{definition}

\begin{lemma}[Reverse trees are odd fixed-even standard tableaux]\label{lem:bc-reverse-tree-fixed-even-standard-tableaux}
For every reduced fixed-even datum arising from a nonempty fixed-cell
completion set, the map sending a reverse reduced corner path to the tableau
whose cell $y_r$ carries $2r-1$ and whose cell $e_r$ carries $2r$
for $1\le r<m$ is a bijection
\[
  \mathcal T(\overline\nu;\overline E_1,\ldots,\overline E_m)
  \longleftrightarrow
  \mathcal S(\gamma;e_1,\ldots,e_{m-1}).
\]
Consequently
\[
  |\mathcal T(\overline\nu;\overline E_1,\ldots,\overline E_m)|
  =
  |\mathcal S(\gamma;e_1,\ldots,e_{m-1})|.
\]
\end{lemma}

\begin{proof}
Let a reverse reduced corner path be given.  Its first step removes the fixed
terminal strip $\overline E_m$ from $\overline\nu$ and gives the Young diagram
$\gamma=\beta_{m-\frac12}$.  Thereafter the reverse removals are, in order,
\[
  y_m,\ e_{m-1},\ y_{m-1},\ e_{m-2},\ldots,e_1,\ y_1 .
\]
Equivalently, the forward construction of $\gamma$ adds
\[
  y_1,\ e_1,\ y_2,\ e_2,\ldots,y_{m-1},\ e_{m-1},\ y_m .
\]
Every intermediate set is a Young diagram by
\cref{def:bc-reduced-reverse-corner-tree}.  Therefore placing $2r-1$ in
$y_r$ for $1\le r\le m$ and $2r$ in $e_r$ for $1\le r<m$ gives a standard
Young tableau of shape $\gamma$, with the required fixed even entries.

Conversely, let
$S\in\mathcal S(\gamma;e_1,\ldots,e_{m-1})$.  Let $y_r$ be the cell carrying
$2r-1$.  The sublevel sets of a standard Young tableau are Young diagrams.
Reading these sublevel diagrams in reverse shows that the cells
\[
  y_m,\ e_{m-1},\ y_{m-1},\ldots,e_1,\ y_1
\]
are removable in that order from $\gamma$.  Restoring the fixed terminal strip
$\overline E_m$ to $\gamma$ gives the initial diagram $\overline\nu$, and the
datum already supplies that this restoration is the valid terminal fixed
horizontal strip of a reduced corner path.  Thus the same cells define an
element of
$\mathcal T(\overline\nu;\overline E_1,\ldots,\overline E_m)$.

The two constructions are inverse because the cells carrying the odd entries
and the fixed even entries are read off uniquely.  The cardinality identity
follows.
\end{proof}

\begin{corollary}[Odd fixed-even standard-tableau bound closes the reverse tree]\label{cor:bc-odd-fixed-even-standard-bound-closes-reverse-tree}
Suppose that every datum of \cref{def:bc-odd-fixed-even-standard-tableaux}
satisfies
\[
  |\mathcal S(\gamma;e_1,\ldots,e_{m-1})|\le2^m .
\]
Then every reduced fixed-even datum arising from a nonempty fixed-cell
completion set satisfies
\[
  |\mathcal T(\overline\nu;\overline E_1,\ldots,\overline E_m)|\le2^m .
\]
\end{corollary}

\begin{proof}
This is immediate from
\cref{lem:bc-reverse-tree-fixed-even-standard-tableaux}.
\end{proof}

\begin{definition}[Odd-cell interval extension poset]\label{def:bc-odd-cell-interval-extension-poset}
Let $(\gamma;e_1,\ldots,e_{m-1})$ be a datum of
\cref{def:bc-odd-fixed-even-standard-tableaux}.  Put
\[
  U=\gamma\setminus\{e_1,\ldots,e_{m-1}\}.
\]
Order boxes by the Young order:
\[
  (i,j)\preceq(i',j')\quad\Longleftrightarrow\quad i\le i'\ \text{and}\ j\le j' .
\]
For $u\in U$, define
\[
  L(u)=1+\max\bigl(\{0\}\cup\{s:\ e_s\preceq u\}\bigr),
  \qquad
  R(u)=\min\bigl(\{m\}\cup\{s:\ u\preceq e_s\}\bigr).
\]
Let
\[
  \mathcal L(\gamma;e_1,\ldots,e_{m-1})
\]
be the set of bijections $\phi:U\to\{1,\ldots,m\}$ such that
\[
  L(u)\le\phi(u)\le R(u)\qquad(u\in U),
\]
and, whenever $u\prec v$ in the Young order on $U$,
\[
  \phi(u)<\phi(v).
\]
\end{definition}

\begin{lemma}[Fixed-even tableaux are interval-constrained extensions]\label{lem:bc-fixed-even-standard-interval-extensions}
For every datum of
\cref{def:bc-odd-fixed-even-standard-tableaux}, the map
\[
  S\longmapsto\phi_S,\qquad S(u)=2\phi_S(u)-1\quad(u\in U),
\]
is a bijection
\[
  \mathcal S(\gamma;e_1,\ldots,e_{m-1})
  \longleftrightarrow
  \mathcal L(\gamma;e_1,\ldots,e_{m-1}).
\]
Consequently
\[
  |\mathcal S(\gamma;e_1,\ldots,e_{m-1})|
  =
  |\mathcal L(\gamma;e_1,\ldots,e_{m-1})|.
\]
\end{lemma}

\begin{proof}
Let $S\in\mathcal S(\gamma;e_1,\ldots,e_{m-1})$.  The non-fixed cells carry
exactly the odd entries $1,3,\ldots,2m-1$, so the displayed rule defines a
bijection $\phi_S:U\to\{1,\ldots,m\}$.  If $u\prec v$ are both in $U$, then
standardness gives $S(u)<S(v)$, hence $\phi_S(u)<\phi_S(v)$.

If $e_s\preceq u$, then standardness gives
\[
  2s=S(e_s)<S(u)=2\phi_S(u)-1,
\]
so $\phi_S(u)\ge s+1$.  Taking the largest such $s$ gives
$L(u)\le\phi_S(u)$.  If $u\preceq e_s$, then
\[
  2\phi_S(u)-1=S(u)<S(e_s)=2s,
\]
so $\phi_S(u)\le s$.  Taking the smallest such $s$ gives
$\phi_S(u)\le R(u)$.  Thus $\phi_S\in\mathcal L$.

Conversely, let $\phi\in\mathcal L(\gamma;e_1,\ldots,e_{m-1})$.  Fill
$e_s$ with $2s$ and fill $u\in U$ with $2\phi(u)-1$.  The entries are exactly
$1,\ldots,2m-1$.  If $u\prec v$ both lie in $U$, then
$\phi(u)<\phi(v)$ gives increasing entries.  If $e_s\preceq u$, the lower
bound $L(u)\le\phi(u)$ implies $\phi(u)\ge s+1$, so
$2s<2\phi(u)-1$.  If $u\preceq e_s$, the upper bound
$\phi(u)\le R(u)$ implies $\phi(u)\le s$, so $2\phi(u)-1<2s$.  Finally, if
$e_s\preceq e_t$, the datum is nonempty, so some tableau in
$\mathcal S(\gamma;e_1,\ldots,e_{m-1})$ exists; applying the previous
paragraph to that tableau gives $s<t$.  Hence the fixed even entries are also
increasing along every comparable pair.  The filling is therefore a standard
Young tableau in $\mathcal S(\gamma;e_1,\ldots,e_{m-1})$.

The two constructions are inverse because odd entries have the unique form
$2r-1$.  The cardinality identity follows.
\end{proof}

\begin{corollary}[Interval-extension bound closes the fixed-even leaf]\label{cor:bc-interval-extension-bound-closes-fixed-even-leaf}
Suppose that every datum of
\cref{def:bc-odd-cell-interval-extension-poset} satisfies
\[
  |\mathcal L(\gamma;e_1,\ldots,e_{m-1})|\le2^m .
\]
Then every datum of
\cref{def:bc-odd-fixed-even-standard-tableaux} satisfies
\[
  |\mathcal S(\gamma;e_1,\ldots,e_{m-1})|\le2^m .
\]
\end{corollary}

\begin{proof}
This is immediate from
\cref{lem:bc-fixed-even-standard-interval-extensions}.
\end{proof}

\begin{definition}[Last-pair deletion candidates]\label{def:bc-last-pair-deletion-candidates}
Let $(\gamma;e_1,\ldots,e_{m-1})$ be a datum of
\cref{def:bc-odd-fixed-even-standard-tableaux}, with $m\ge2$, and put
$e=e_{m-1}$.  Let
\[
  Y(\gamma;e)
\]
be the set of cells $y\in\gamma\setminus\{e_1,\ldots,e_{m-1}\}$ such that
$y$ is a removable corner of $\gamma$ and $e$ is a removable corner of
$\gamma\setminus\{y\}$.  For $y\in Y(\gamma;e)$, write
\[
  \gamma_y=\gamma\setminus\{y,e\}.
\]
\end{definition}

\begin{lemma}[Last-pair deletion recurrence]\label{lem:bc-last-pair-deletion-recurrence}
For every datum of \cref{def:bc-odd-fixed-even-standard-tableaux} with
$m\ge2$,
\[
  |\mathcal S(\gamma;e_1,\ldots,e_{m-1})|
  =
  \sum_{y\in Y(\gamma;e_{m-1})}
  |\mathcal S(\gamma_y;e_1,\ldots,e_{m-2})|.
\]
\end{lemma}

\begin{proof}
Let $S\in\mathcal S(\gamma;e_1,\ldots,e_{m-1})$, and let $y$ be the cell
carrying the largest entry $2m-1$.  Since sublevel sets of a standard Young
tableau are Young diagrams, $y$ is a removable corner of $\gamma$.  After
removing $y$, the largest remaining entry is the fixed entry $2m-2$ in
$e_{m-1}$, so $e_{m-1}$ is a removable corner of
$\gamma\setminus\{y\}$.  Thus $y\in Y(\gamma;e_{m-1})$.  Removing first $y$
and then $e_{m-1}$ leaves a standard tableau of shape $\gamma_y$ whose fixed
even entries are $2,4,\ldots,2m-4$ at
$e_1,\ldots,e_{m-2}$.

Conversely, fix $y\in Y(\gamma;e_{m-1})$ and a tableau in
$\mathcal S(\gamma_y;e_1,\ldots,e_{m-2})$.  Put $2m-2$ in $e_{m-1}$ and
$2m-1$ in $y$.  The definition of $Y(\gamma;e_{m-1})$ says exactly that
\[
  \gamma_y\subset \gamma_y\cup\{e_{m-1}\}
  \subset \gamma_y\cup\{e_{m-1},y\}=\gamma
\]
are Young diagrams obtained by adjoining removable corners in the displayed
order.  Therefore adjoining the two largest entries preserves standardness
and gives an element of
$\mathcal S(\gamma;e_1,\ldots,e_{m-1})$.

The two constructions are inverse, because the cell carrying $2m-1$ is
uniquely recovered as $y$.  Summing over $y$ gives the displayed recurrence.
\end{proof}

\begin{corollary}[Blocked last-pair choices are binary]\label{cor:bc-last-pair-blocked-binary}
In the setting of \cref{def:bc-last-pair-deletion-candidates}, if
$e=e_{m-1}$ is not a removable corner of $\gamma$, then
\[
  |Y(\gamma;e)|\le2 .
\]
More precisely, every $y\in Y(\gamma;e)$ is either immediately east or
immediately south of $e$.
\end{corollary}

\begin{proof}
Write $e=(i,j)$.  Since $e$ is not removable in $\gamma$, at least one of
$(i,j+1)$ and $(i+1,j)$ belongs to $\gamma$.  Since $e$ becomes removable
after deleting $y$, every cell of $\gamma$ immediately east or south of $e$
must be deleted by that same one-cell deletion.  Hence $y$ is one of those
two adjacent cells.  This gives the displayed bound.
\end{proof}

\begin{definition}[Fixed-even recursive envelope]\label{def:bc-fixed-even-recursive-envelope}
For a Young diagram $\gamma$ of odd size, define $W(\gamma)$ recursively as
follows.  If $\gamma=(1)$, put $W(\gamma)=1$.  If
$|\gamma|=2m-1$ with $m\ge2$, and if $e\in\gamma$, let
\[
  \widehat Y(\gamma;e)
\]
be the set of cells $y\in\gamma$ such that $y\ne e$, $y$ is a removable
corner of $\gamma$, and $e$ is a removable corner of $\gamma\setminus\{y\}$.
Put
\[
  W(\gamma)
  =
  \max_{e\in\gamma}
  \sum_{y\in\widehat Y(\gamma;e)}
  W(\gamma\setminus\{y,e\}).
\]
\end{definition}

\begin{lemma}[The recursive envelope bounds fixed-even fibers]\label{lem:bc-fixed-even-recursive-envelope-bounds-fibers}
For every datum of \cref{def:bc-odd-fixed-even-standard-tableaux},
\[
  |\mathcal S(\gamma;e_1,\ldots,e_{m-1})|
  \le W(\gamma).
\]
\end{lemma}

\begin{proof}
We induct on $m$.  For $m=1$, the shape has one cell and the fixed-even fiber
has cardinality at most one, which is $W((1))$.

Assume $m\ge2$ and put $e=e_{m-1}$.  By
\cref{lem:bc-last-pair-deletion-recurrence},
\[
  |\mathcal S(\gamma;e_1,\ldots,e_{m-1})|
  =
  \sum_{y\in Y(\gamma;e)}
  |\mathcal S(\gamma_y;e_1,\ldots,e_{m-2})|.
\]
For each $y\in Y(\gamma;e)$, the child shape is
\[
  \gamma_y=\gamma\setminus\{y,e\},
\]
and the induction hypothesis bounds the corresponding summand by
$W(\gamma_y)$.  Moreover
\[
  Y(\gamma;e)\subseteq\widehat Y(\gamma;e),
\]
because the latter set drops only the requirement that $y$ avoid the earlier
fixed even cells.  Hence
\[
  |\mathcal S(\gamma;e_1,\ldots,e_{m-1})|
  \le
  \sum_{y\in\widehat Y(\gamma;e)}
  W(\gamma\setminus\{y,e\})
  \le W(\gamma),
\]
by the definition of the maximum in
\cref{def:bc-fixed-even-recursive-envelope}.
\end{proof}

\begin{proposition}[The recursive fixed-even envelope is full-budget through twenty-eight]\label{prop:bc-fixed-even-envelope-m28-certificate}
For every Young diagram $\gamma$ of size $2m-1$ with $1\le m\le28$,
\[
  W(\gamma)\le2^{2m}.
\]
\end{proposition}

\begin{proof}
The exact C++ verifier
\[
\texttt{character\_ring\_iter/post\_m29\_bc\_fixed\_even\_envelope\_gmp.cpp}
\]
implements \cref{def:bc-fixed-even-recursive-envelope} by dynamic programming
over all partitions of $1,3,\ldots,55$.  For each $m$, it enumerates every
partition $\gamma$ of $2m-1$, computes $W(\gamma)$ from the previously
computed values at size $2m-3$, and records
\[
  \max_{|\gamma|=2m-1}W(\gamma).
\]

The archived `optimus' run
\[
\texttt{certificates/post\_m29/bc\_fixed\_even\_envelope\_m28\_gmp\_certificate.log}
\]
has SHA-256
\[
\hashsplit{ff511546d3ffde8c831c4076ca125ff4}{7bc6d0f9f1279f1e2d890a0d925292b0}.
\]
The source file has SHA-256
\[
\hashsplit{af7ce7efddaf8fbdc9ff0b54a54935e6}{b1bc96aa58041b335d64e220a2150b17}.
\]
The log ends with \texttt{\_\_EXIT\_STATUS=0} and
\texttt{SUMMARY failures=0}.  In the only range not already covered by the
factorial estimate, it reports
\[
\begin{array}{c|c|c}
m&\max W(\gamma)&2^{2m}\\
\hline
13&85787&67108864\\
14&370729&268435456\\
15&2037069&1073741824\\
16&10845675&4294967296\\
17&54725353&17179869184\\
18&286275138&68719476736\\
19&1788942066&274877906944\\
20&11349507850&1099511627776\\
21&67061584700&4398046511104\\
22&403479661259&17592186044416\\
23&2554342683425&70368744177664\\
24&18877135179266&281474976710656\\
25&133746685912501&1125899906842624\\
26&907074996503025&4503599627370496\\
27&6344159720142261&18014398509481984\\
28&45416126476990210&72057594037927936 .
\end{array}
\]
Every displayed maximum is at most the corresponding full budget, and the log
also checks $1\le m\le12$.
\end{proof}

\begin{definition}[Recursive fixed-even envelope frontiers]\label{def:bc-fixed-even-envelope-frontier}
For \(m\ge1\), let
\[
  \mathcal E^{\mathrm{env}}_m
  =
  \{\gamma\vdash(2m-1):W(\gamma)>2^{2m}\}.
\]
If \(m\ge2\), \(\gamma\vdash(2m-1)\), and \(e\in\gamma\), put
\[
  W^{\mathrm{top}}(\gamma;e)
  =
  \sum_{y\in\widehat Y(\gamma;e)}
  W(\gamma\setminus\{y,e\}),
\]
and let
\[
  \mathcal P^{\mathrm{env}}_m
  =
  \{(\gamma,e):\gamma\vdash(2m-1),\ e\in\gamma,\;
    W^{\mathrm{top}}(\gamma;e)>2^{2m}\}.
\]
If \(m\ge3\), \(\gamma\vdash(2m-1)\), and \(e,f\in\gamma\), put
\[
  W^{\mathrm{top},2}(\gamma;e,f)
  =
  \sum_{\substack{y\in\widehat Y(\gamma;e)\\
                  y\ne f,\ f\in\gamma\setminus\{y,e\}}}
  W^{\mathrm{top}}(\gamma\setminus\{y,e\};f),
\]
and let
\[
  \mathcal P^{\mathrm{env},2}_m
  =
  \{(\gamma,e,f):\gamma\vdash(2m-1),\ e,f\in\gamma,\;
    W^{\mathrm{top},2}(\gamma;e,f)>2^{2m}\}.
\]
If \(m\ge4\), \(\gamma\vdash(2m-1)\), and \(e,f,g\in\gamma\), put
\[
  W^{\mathrm{top},3}(\gamma;e,f,g)
  =
  \sum_{\substack{y\in\widehat Y(\gamma;e)\\
                  y\ne f,\ y\ne g,\ f,g\in\gamma\setminus\{y,e\}}}
  W^{\mathrm{top},2}(\gamma\setminus\{y,e\};f,g),
\]
and let
\[
  \mathcal P^{\mathrm{env},3}_m
  =
  \{(\gamma,e,f,g):\gamma\vdash(2m-1),\ e,f,g\in\gamma,\;
    W^{\mathrm{top},3}(\gamma;e,f,g)>2^{2m}\}.
\]
For \(d\ge4\), define recursively
\[
  W^{\mathrm{top},d}(\gamma;e_1,\ldots,e_d)
  =
  \sum_{\substack{y\in\widehat Y(\gamma;e_1)\\
                  y\notin\{e_2,\ldots,e_d\}\\
                  e_2,\ldots,e_d\in\gamma\setminus\{y,e_1\}}}
  W^{\mathrm{top},d-1}(\gamma\setminus\{y,e_1\};e_2,\ldots,e_d),
\]
and let
\[
  \mathcal P^{\mathrm{env},d}_m
  =
  \{(\gamma,e_1,\ldots,e_d):\gamma\vdash(2m-1),\
    e_1,\ldots,e_d\in\gamma,\
    W^{\mathrm{top},d}(\gamma;e_1,\ldots,e_d)>2^{2m}\}.
\]
\end{definition}

\begin{lemma}[The top-fixed recursive envelope bounds fixed-even fibers]\label{lem:bc-fixed-even-top-envelope-bounds-fibers}
For every datum of \cref{def:bc-odd-fixed-even-standard-tableaux} with
\(m\ge2\),
\[
  |\mathcal S(\gamma;e_1,\ldots,e_{m-1})|
  \le W^{\mathrm{top}}(\gamma;e_{m-1}).
\]
\end{lemma}

\begin{proof}
Put \(e=e_{m-1}\).  By
\cref{lem:bc-last-pair-deletion-recurrence},
\[
  |\mathcal S(\gamma;e_1,\ldots,e_{m-1})|
  =
  \sum_{y\in Y(\gamma;e)}
  |\mathcal S(\gamma_y;e_1,\ldots,e_{m-2})|.
\]
For each \(y\in Y(\gamma;e)\),
\[
  \gamma_y=\gamma\setminus\{y,e\},
\]
and \cref{lem:bc-fixed-even-recursive-envelope-bounds-fibers} bounds the
corresponding lower fiber by \(W(\gamma\setminus\{y,e\})\).  Since
\[
  Y(\gamma;e)\subseteq\widehat Y(\gamma;e),
\]
the displayed recurrence gives
\[
  |\mathcal S(\gamma;e_1,\ldots,e_{m-1})|
  \le
  \sum_{y\in\widehat Y(\gamma;e)}
  W(\gamma\setminus\{y,e\})
  =
  W^{\mathrm{top}}(\gamma;e).
\]
\end{proof}

\begin{lemma}[The top-two recursive envelope bounds fixed-even fibers]\label{lem:bc-fixed-even-top-two-envelope-bounds-fibers}
For every datum of \cref{def:bc-odd-fixed-even-standard-tableaux} with
\(m\ge3\),
\[
  |\mathcal S(\gamma;e_1,\ldots,e_{m-1})|
  \le W^{\mathrm{top},2}(\gamma;e_{m-1},e_{m-2}).
\]
\end{lemma}

\begin{proof}
Put \(e=e_{m-1}\) and \(f=e_{m-2}\).  By
\cref{lem:bc-last-pair-deletion-recurrence},
\[
  |\mathcal S(\gamma;e_1,\ldots,e_{m-1})|
  =
  \sum_{y\in Y(\gamma;e)}
  |\mathcal S(\gamma_y;e_1,\ldots,e_{m-2})|.
\]
Every \(y\in Y(\gamma;e)\) avoids the earlier fixed even cells, so \(y\ne f\),
and \(f\in\gamma_y=\gamma\setminus\{y,e\}\).  Applying
\cref{lem:bc-fixed-even-top-envelope-bounds-fibers} to the child datum gives
\[
  |\mathcal S(\gamma_y;e_1,\ldots,e_{m-2})|
  \le W^{\mathrm{top}}(\gamma_y;f).
\]
Also \(Y(\gamma;e)\subseteq\widehat Y(\gamma;e)\).  Summing these estimates
gives the displayed bound by the definition of
\(W^{\mathrm{top},2}(\gamma;e,f)\).
\end{proof}

\begin{lemma}[The top-three recursive envelope bounds fixed-even fibers]\label{lem:bc-fixed-even-top-three-envelope-bounds-fibers}
For every datum of \cref{def:bc-odd-fixed-even-standard-tableaux} with
\(m\ge4\),
\[
  |\mathcal S(\gamma;e_1,\ldots,e_{m-1})|
  \le W^{\mathrm{top},3}(\gamma;e_{m-1},e_{m-2},e_{m-3}).
\]
\end{lemma}

\begin{proof}
Put \(e=e_{m-1}\), \(f=e_{m-2}\), and \(g=e_{m-3}\).  The
last-pair recurrence gives
\[
  |\mathcal S(\gamma;e_1,\ldots,e_{m-1})|
  =
  \sum_{y\in Y(\gamma;e)}
  |\mathcal S(\gamma_y;e_1,\ldots,e_{m-2})|.
\]
Every \(y\in Y(\gamma;e)\) avoids \(f\) and \(g\), and both cells remain in
\(\gamma_y=\gamma\setminus\{y,e\}\).  By
\cref{lem:bc-fixed-even-top-two-envelope-bounds-fibers},
\[
  |\mathcal S(\gamma_y;e_1,\ldots,e_{m-2})|
  \le W^{\mathrm{top},2}(\gamma_y;f,g).
\]
Using again \(Y(\gamma;e)\subseteq\widehat Y(\gamma;e)\) and summing gives
the asserted \(W^{\mathrm{top},3}\) bound.
\end{proof}

\begin{lemma}[The fixed-depth recursive envelope bounds fixed-even fibers]\label{lem:bc-fixed-even-fixed-depth-envelope-bounds-fibers}
Let \(d\ge4\).  For every datum of
\cref{def:bc-odd-fixed-even-standard-tableaux} with \(m\ge d+1\),
\[
  |\mathcal S(\gamma;e_1,\ldots,e_{m-1})|
  \le
  W^{\mathrm{top},d}
  (\gamma;e_{m-1},e_{m-2},\ldots,e_{m-d}).
\]
\end{lemma}

\begin{proof}
We induct on \(d\).  For \(d=4\), put \(e=e_{m-1}\).  By
\cref{lem:bc-last-pair-deletion-recurrence},
\[
  |\mathcal S(\gamma;e_1,\ldots,e_{m-1})|
  =
  \sum_{y\in Y(\gamma;e)}
  |\mathcal S(\gamma_y;e_1,\ldots,e_{m-2})|.
\]
For \(y\in Y(\gamma;e)\), the cells
\(e_{m-2},e_{m-3},e_{m-4}\) remain in
\(\gamma_y=\gamma\setminus\{y,e\}\).  Hence
\cref{lem:bc-fixed-even-top-three-envelope-bounds-fibers} applied to the child
datum bounds the corresponding summand by
\[
  W^{\mathrm{top},3}(\gamma_y;e_{m-2},e_{m-3},e_{m-4}).
\]
Since \(Y(\gamma;e)\subseteq\widehat Y(\gamma;e)\), summing gives the
definition of
\[
  W^{\mathrm{top},4}
  (\gamma;e_{m-1},e_{m-2},e_{m-3},e_{m-4}).
\]

Assume the assertion for \(d-1\ge4\).  The same last-pair recurrence and
survival of the earlier fixed even cells give, for every \(y\in Y(\gamma;e)\),
\[
  |\mathcal S(\gamma_y;e_1,\ldots,e_{m-2})|
  \le
  W^{\mathrm{top},d-1}
  (\gamma_y;e_{m-2},\ldots,e_{m-d}).
\]
Using again \(Y(\gamma;e)\subseteq\widehat Y(\gamma;e)\), the sum of these
bounds is exactly
\[
  W^{\mathrm{top},d}
  (\gamma;e_{m-1},e_{m-2},\ldots,e_{m-d}),
\]
by \cref{def:bc-fixed-even-envelope-frontier}.
\end{proof}

\begin{lemma}[The full-depth recursive envelope is the exact fixed-even fiber]\label{lem:bc-full-depth-envelope-exact-fiber}
For every datum of \cref{def:bc-odd-fixed-even-standard-tableaux} with
\(m\ge2\), define
\[
  \mathfrak W_m(\gamma;e_1,\ldots,e_{m-1})
  =
  \begin{cases}
  W^{\mathrm{top}}(\gamma;e_1),&m=2,\\
  W^{\mathrm{top},2}(\gamma;e_2,e_1),&m=3,\\
  W^{\mathrm{top},3}(\gamma;e_3,e_2,e_1),&m=4,\\
  W^{\mathrm{top},m-1}(\gamma;e_{m-1},e_{m-2},\ldots,e_1),&m\ge5.
  \end{cases}
\]
Then
\[
  |\mathcal S(\gamma;e_1,\ldots,e_{m-1})|
  =
  \mathfrak W_m(\gamma;e_1,\ldots,e_{m-1}).
\]
\end{lemma}

\begin{proof}
We induct on \(m\).  For \(m=2\), the set \(Y(\gamma;e_1)\) of
\cref{def:bc-last-pair-deletion-candidates} is exactly
\(\widehat Y(\gamma;e_1)\).  For every \(y\in Y(\gamma;e_1)\), the child
\(\gamma\setminus\{y,e_1\}\) has one cell, so its fixed-even fiber and its
recursive envelope both have cardinality \(1\).  The last-pair recurrence
\cref{lem:bc-last-pair-deletion-recurrence} therefore gives the displayed
identity with \(W^{\mathrm{top}}(\gamma;e_1)\).

Assume \(m\ge3\) and the assertion has been proved for \(m-1\).  Put
\(e=e_{m-1}\).  In the last-pair recurrence, the summation set
\(Y(\gamma;e)\) is exactly the set of cells \(y\) appearing in the first
summation of the corresponding full-depth envelope: \(y\) is a removable
corner, \(e\) is removable after deleting \(y\), and \(y\) avoids all earlier
fixed even cells \(e_1,\ldots,e_{m-2}\), which then survive in the child
\(\gamma_y=\gamma\setminus\{y,e\}\).  For every such child, the induction
hypothesis identifies
\[
  |\mathcal S(\gamma_y;e_1,\ldots,e_{m-2})|
\]
with the lower full-depth envelope: \(W^{\mathrm{top}}(\gamma_y;e_1)\) when
\(m=3\), \(W^{\mathrm{top},2}(\gamma_y;e_2,e_1)\) when \(m=4\),
\(W^{\mathrm{top},3}(\gamma_y;e_3,e_2,e_1)\) when \(m=5\), and
\[
  W^{\mathrm{top},m-2}(\gamma_y;e_{m-2},\ldots,e_1)
\]
when \(m\ge6\).  Substituting these identities into
\cref{lem:bc-last-pair-deletion-recurrence} gives exactly the defining
recursion for \(\mathfrak W_m(\gamma;e_1,\ldots,e_{m-1})\).
\end{proof}

\begin{corollary}[The full-depth envelope is exact for broad odd-packet classes]\label{cor:bc-full-depth-envelope-exact-omega}
Suppose \(m\ge5\), and suppose target-window data determine
\[
  \lambda^{(2m)},\qquad E_1,\ldots,E_m .
\]
Let
\[
  \Omega_{\mathrm{tw}}
  =
  \Omega(\lambda^{(2m)};E_1,\ldots,E_m).
\]
If \(\Omega_{\mathrm{tw}}\ne\varnothing\), let
\((\gamma;e_1,\ldots,e_{m-1})\) be the associated fixed-even datum.  Then
\[
  |\Omega_{\mathrm{tw}}|
  =
  W^{\mathrm{top},m-1}(\gamma;e_{m-1},e_{m-2},\ldots,e_1).
\]
Consequently,
\[
  (\gamma,e_{m-1},e_{m-2},\ldots,e_1)
  \notin\mathcal P^{\mathrm{env},m-1}_m
\]
is equivalent to \(|\Omega_{\mathrm{tw}}|\le2^{2m}\).
\end{corollary}

\begin{proof}
The cardinality identifications in
\cref{lem:bc-omega-fixed-cell-completions,prop:bc-completions-equal-reduced-growth-chains,lem:bc-reduced-growth-reverse-corner-tree,lem:bc-reverse-tree-fixed-even-standard-tableaux}
give
\[
  |\Omega_{\mathrm{tw}}|
  =
  |\mathcal S(\gamma;e_1,\ldots,e_{m-1})|.
\]
Since \(m\ge5\), \cref{lem:bc-full-depth-envelope-exact-fiber} identifies the
right side with
\[
  W^{\mathrm{top},m-1}(\gamma;e_{m-1},e_{m-2},\ldots,e_1).
\]
The final equivalence is then the definition of
\(\mathcal P^{\mathrm{env},m-1}_m\).
\end{proof}

\begin{proposition}[The recursive fixed-even envelope has a certified twenty-nine frontier]\label{prop:bc-fixed-even-envelope-m29-obstruction}
Let
\[
  \gamma=(11,9,8,7,6,5,4,3,2,1,1).
\]
Then
\[
  |\mathcal E^{\mathrm{env}}_{29}|=144,
  \qquad
  |\mathcal P^{\mathrm{env}}_{29}|=346,
  \qquad
  |\mathcal P^{\mathrm{env},2}_{29}|=62,
  \qquad
  \mathcal P^{\mathrm{env},3}_{29}=\varnothing.
\]
For every partition $\eta\vdash57$ outside
$\mathcal E^{\mathrm{env}}_{29}$,
\[
  W(\eta)\le2^{58}.
\]
For every partition $\eta\vdash57$ and cell $e\in\eta$ with
\((\eta,e)\notin\mathcal P^{\mathrm{env}}_{29}\),
\[
  W^{\mathrm{top}}(\eta;e)\le2^{58}.
\]
For every partition $\eta\vdash57$ and cells $e,f\in\eta$ with
\((\eta,e,f)\notin\mathcal P^{\mathrm{env},2}_{29}\),
\[
  W^{\mathrm{top},2}(\eta;e,f)\le2^{58}.
\]
For every partition $\eta\vdash57$ and cells $e,f,g\in\eta$,
\[
  W^{\mathrm{top},3}(\eta;e,f,g)\le2^{58}.
\]
Moreover, $|\gamma|=57=2\cdot29-1$, $\gamma$ is not column-even, and
\[
  W(\gamma)=367676108829086275
  =
  \max_{\eta\vdash57}W(\eta)
  >288230376151711744=2^{58}.
\]
Consequently the uniform full-budget assertion
$W(\gamma)\le2^{2m}$ cannot be extended from
\cref{prop:bc-fixed-even-envelope-m28-certificate} to all non-column-even
shapes at $m=29$ by the same recursive envelope.  This proposition is an
obstruction to the uniform envelope route and a positive certificate for the
complement of the frontier; it gives no lower bound on the actual broad-Omega
compatibility classes or on any fixed-even tableau fiber.
\end{proposition}

\begin{proof}
The size identity is immediate by summing the displayed row lengths.  The
conjugate diagram is again
\[
  \gamma'=(11,9,8,7,6,5,4,3,2,1,1),
\]
so its first column has odd length.  Hence $\gamma$ is not column-even.

The exact C++ verifier
\[
\texttt{character\_ring\_iter/post\_m29\_bc\_fixed\_even\_envelope\_m29\_obstruction\_gmp.cpp}
\]
implements the recursion of \cref{def:bc-fixed-even-recursive-envelope} by
dynamic programming over every partition of $1,3,\ldots,57$.  It then checks
the displayed shape, its non-column-even status, the exact value of
$W(\gamma)$, the exact budget $2^{58}$, that the displayed shape is a
maximizer among partitions of $57$, and that exactly $144$ partitions of $57$
have recursive-envelope value strictly above $2^{58}$.  It also computes
\(W^{\mathrm{top}}(\eta;e)\) for every cell \(e\) of every partition
\(\eta\vdash57\), checks that exactly \(346\) pairs have top-fixed
recursive-envelope value strictly above \(2^{58}\), and prints one
\texttt{M29\_TOP\_PAIR\_OVER\_BUDGET} line for each such pair.  Since
\(W^{\mathrm{top},2}(\eta;e,f)\le W^{\mathrm{top}}(\eta;e)\), it then
computes \(W^{\mathrm{top},2}\) for every possible next fixed even cell over
the \(346\) top-fixed pairs and checks that exactly \(62\) triples remain
above budget.  Since
\(W^{\mathrm{top},3}(\eta;e,f,g)\le W^{\mathrm{top},2}(\eta;e,f)\), it then
computes \(W^{\mathrm{top},3}\) over those \(62\) triples and checks that no
quadruple remains above budget.

The archived `optimus' run
\[
\texttt{certificates/post\_m29/bc\_fixed\_even\_envelope\_m29\_obstruction\_gmp\_certificate.log}
\]
has SHA-256
\[
\hashsplit{cf30ce672a1511201db50c17ecfe15a9}{f6431fdb169cb617ef0c668940fd16a1}.
\]
The source file has SHA-256
\[
\hashsplit{2d7b99007d80e10f3a04c0c546ddb3ed}{8c54ffa8adb5382ea9898a39aadd81b8}.
\]
The log ends with \texttt{\_\_EXIT\_STATUS=0} and
\texttt{SUMMARY failures=0}.  Its frontier and target lines report
\[
\begin{gathered}
\texttt{M29\_OVER\_BUDGET count=144,\ expected\_count=144,}\\
\texttt{target\_in\_over\_budget=yes},\\
\texttt{M29\_TOP\_PAIR\_OVER\_BUDGET count=346,\ expected\_count=346,}\\
\texttt{target\_has\_top\_pair\_over\_budget=yes},\\
\texttt{M29\_TOP\_TWO\_OVER\_BUDGET count=62,\ expected\_count=62,}\\
\texttt{M29\_TOP\_THREE\_OVER\_BUDGET count=0,\ expected\_count=0,}\\
\texttt{size=57,\ column\_even=no,\ W=367676108829086275},\\
\texttt{budget=2\string^58,\ budget\_value=288230376151711744,}\\
\texttt{exceeds\_budget=yes,\ target\_is\_m29\_maximizer=yes}.
\end{gathered}
\]
These exact checks give the asserted cardinalities, the displayed maximum, and
the strict inequality.  Since the verifier enumerates every partition
$\eta\vdash57$ and records precisely those with $W(\eta)>2^{58}$, every
partition outside \(\mathcal E^{\mathrm{env}}_{29}\) satisfies
\(W(\eta)\le2^{58}\).  Since it also enumerates every pair \((\eta,e)\) with
\(\eta\vdash57\) and \(e\in\eta\), and records precisely those with
\(W^{\mathrm{top}}(\eta;e)>2^{58}\), every pair outside
\(\mathcal P^{\mathrm{env}}_{29}\) satisfies the displayed top-fixed bound.
The monotonicity inequalities in the preceding paragraph show that an
over-budget top-two triple must lie over an over-budget top pair, and an
over-budget top-three quadruple must lie over an over-budget top-two triple.
They follow directly from the definitions: \(W\) is the maximum over the
top-fixed values, and each deeper fixed-envelope sum is obtained by restricting
the preceding summation and replacing each child contribution by a value no
larger than the preceding child contribution.
The verifier records precisely these remaining over-budget triples and
quadruples, giving \(|\mathcal P^{\mathrm{env},2}_{29}|=62\) and
\(\mathcal P^{\mathrm{env},3}_{29}=\varnothing\).
\end{proof}

\begin{proposition}[The fixed-depth recursive envelope closes the thirty frontier]\label{prop:bc-fixed-even-envelope-m30-frontier-certificate}
For \(m=30\),
\[
  |\mathcal E^{\mathrm{env}}_{30}|=11090.
\]
The maximum recursive-envelope value among partitions of \(59\) is attained at
\[
  \gamma=(11,10,8,7,6,5,4,3,2,2,1),
\]
where
\[
  W(\gamma)=2963381126459933824>1152921504606846976=2^{60}.
\]
Furthermore the successive over-budget fixed-depth frontiers have
cardinalities
\[
\begin{array}{c|ccccccc}
d&1&2&3&4&5&6&7\\
\hline
|\mathcal P^{\mathrm{env},d}_{30}|
&44333&92312&60440&5438&408&18&0 .
\end{array}
\]
Consequently, for every partition \(\eta\vdash59\) and every seven distinct
cells \(e_1,\ldots,e_7\in\eta\),
\[
  W^{\mathrm{top},7}(\eta;e_1,\ldots,e_7)\le2^{60}.
\]
\end{proposition}

\begin{proof}
The exact C++ verifier
\[
\texttt{character\_ring\_iter/post\_m29\_bc\_fixed\_even\_envelope\_m30\_frontier\_gmp.cpp}
\]
implements \cref{def:bc-fixed-even-recursive-envelope,def:bc-fixed-even-envelope-frontier}
by dynamic programming over every partition of \(1,3,\ldots,59\).  It first
computes \(W(\eta)\) for every \(\eta\vdash59\), checks the displayed maximum,
and checks that exactly \(11090\) partitions have \(W(\eta)>2^{60}\).  It then
extends only the over-budget prefixes because
\[
  W^{\mathrm{top},d+1}(\eta;e_1,\ldots,e_{d+1})
  \le
  W^{\mathrm{top},d}(\eta;e_1,\ldots,e_d),
\]
which follows directly from the fixed-depth recursive definition.  Thus every
over-budget depth-\((d+1)\) prefix lies over an over-budget depth-\(d\) prefix.

The archived `optimus' run
\[
\texttt{certificates/post\_m29/bc\_fixed\_even\_envelope\_m30\_frontier\_gmp\_certificate.log}
\]
has SHA-256
\[
\hashsplit{b540fe210ff63596bf4ca3b761e32da3}{8eca42ca0c5340f1c3e25d41f09fb593}.
\]
The source file has SHA-256
\[
\hashsplit{eac9e5697b79ce291b8eabff388e41be}{7ac0858ea536223f4fe1f3b264bd5d29}.
\]
The log ends with \texttt{\_\_EXIT\_STATUS=0} and
\texttt{SUMMARY remaining\_count=0 failures=0}.  Its exact count lines are
\[
\begin{gathered}
\texttt{M30\_OVER\_BUDGET\_SHAPES count=11090,\ expected\_count=11090,}\\
\texttt{M30\_MAX W=2963381126459933824,\ expected\_W=2963381126459933824,}\\
\texttt{M30\_TOP\_DEPTH depth=1 count=44333,\ expected\_count=44333,}\\
\texttt{M30\_TOP\_DEPTH depth=2 count=92312,\ expected\_count=92312,}\\
\texttt{M30\_TOP\_DEPTH depth=3 count=60440,\ expected\_count=60440,}\\
\texttt{M30\_TOP\_DEPTH depth=4 count=5438,\ expected\_count=5438,}\\
\texttt{M30\_TOP\_DEPTH depth=5 count=408,\ expected\_count=408,}\\
\texttt{M30\_TOP\_DEPTH depth=6 count=18,\ expected\_count=18,}\\
\texttt{M30\_TOP\_DEPTH depth=7 count=0,\ expected\_count=0.}
\end{gathered}
\]
These exact checks prove the displayed cardinalities and the empty
depth-seven frontier.  The final inequality is precisely
\(\mathcal P^{\mathrm{env},7}_{30}=\varnothing\).
\end{proof}

\begin{proposition}[The fixed-depth recursive envelope closes the thirty-one frontier]\label{prop:bc-fixed-even-envelope-m31-frontier-certificate}
For \(m=31\),
\[
  |\mathcal E^{\mathrm{env}}_{31}|=67038.
\]
The maximum recursive-envelope value among partitions of \(61\) is attained at
\[
  \gamma=(12,10,9,7,6,5,4,3,2,2,1),
\]
where
\[
  W(\gamma)=22734772942354066748>4611686018427387904=2^{62}.
\]
Furthermore the successive over-budget fixed-depth frontiers have
cardinalities
\[
\begin{array}{c|cccccccccc}
d&1&2&3&4&5&6&7&8&9&10\\
\hline
|\mathcal P^{\mathrm{env},d}_{31}|
&331969&1229238&2687932&2263824&499606&49486&6142&610&24&0 .
\end{array}
\]
Consequently, for every partition \(\eta\vdash61\) and every ten distinct
cells \(e_1,\ldots,e_{10}\in\eta\),
\[
  W^{\mathrm{top},10}(\eta;e_1,\ldots,e_{10})\le2^{62}.
\]
\end{proposition}

\begin{proof}
The exact C++ verifier
\[
\texttt{character\_ring\_iter/post\_m29\_bc\_fixed\_even\_envelope\_m31\_frontier\_gmp.cpp}
\]
implements \cref{def:bc-fixed-even-recursive-envelope,def:bc-fixed-even-envelope-frontier}
by dynamic programming over every partition of \(1,3,\ldots,61\).  It first
computes \(W(\eta)\) exactly with GMP integers for every \(\eta\vdash61\),
checks the displayed maximum, and checks that exactly \(67038\) partitions
have \(W(\eta)>2^{62}\).  For the fixed-depth frontier comparisons it computes
the exact capped integer
\[
  \widetilde W^{\mathrm{top},d}
  =
  \min\{W^{\mathrm{top},d},2^{62}+1\}
\]
by the same recursive summation, stopping a summand only after an exact
partial integer sum has exceeded \(2^{62}\).  Therefore
\[
  \widetilde W^{\mathrm{top},d}>2^{62}
  \quad\Longleftrightarrow\quad
  W^{\mathrm{top},d}>2^{62},
\]
so the capped computation is exact for every over-budget membership test.
As in the \(m=30\) verifier, monotonicity gives
\[
  W^{\mathrm{top},d+1}(\eta;e_1,\ldots,e_{d+1})
  \le
  W^{\mathrm{top},d}(\eta;e_1,\ldots,e_d),
\]
so every over-budget depth-\((d+1)\) prefix lies over an over-budget
depth-\(d\) prefix.

The archived `optimus' run
\[
\texttt{certificates/post\_m29/bc\_fixed\_even\_envelope\_m31\_frontier\_gmp\_certificate.log}
\]
has SHA-256
\[
\hashsplit{7331b1cdd3ee34cd969d5456fd44ca3f}{f5f814ba2f983a0043fda5eaa4e6d69f}.
\]
The source file has SHA-256
\[
\hashsplit{654537912ac6e9333a41aa5fcc698e19}{33e59a9775c1fa1bf1f8e49a826fab0a}.
\]
The log ends with \texttt{\_\_EXIT\_STATUS=0},
\texttt{\_\_SHELL\_STATUS=0}, and
\texttt{SUMMARY remaining\_count=0 failures=0}.  Its exact count lines are
\[
\begin{gathered}
\texttt{M31\_OVER\_BUDGET\_SHAPES count=67038,\ expected\_count=67038,}\\
\texttt{M31\_MAX W=22734772942354066748,\ expected\_W=22734772942354066748,}\\
\texttt{M31\_TOP\_DEPTH depth=1 count=331969,\ expected\_count=331969,}\\
\texttt{M31\_TOP\_DEPTH depth=2 count=1229238,\ expected\_count=1229238,}\\
\texttt{M31\_TOP\_DEPTH depth=3 count=2687932,\ expected\_count=2687932,}\\
\texttt{M31\_TOP\_DEPTH depth=4 count=2263824,\ expected\_count=2263824,}\\
\texttt{M31\_TOP\_DEPTH depth=5 count=499606,\ expected\_count=499606,}\\
\texttt{M31\_TOP\_DEPTH depth=6 count=49486,\ expected\_count=49486,}\\
\texttt{M31\_TOP\_DEPTH depth=7 count=6142,\ expected\_count=6142,}\\
\texttt{M31\_TOP\_DEPTH depth=8 count=610,\ expected\_count=610,}\\
\texttt{M31\_TOP\_DEPTH depth=9 count=24,\ expected\_count=24,}\\
\texttt{M31\_TOP\_DEPTH depth=10 count=0,\ expected\_count=0.}
\end{gathered}
\]
These exact checks prove the displayed cardinalities and the empty depth-ten
frontier.  The final inequality is precisely
\(\mathcal P^{\mathrm{env},10}_{31}=\varnothing\).
\end{proof}

\begin{definition}[Commuting normal-form signature]\label{def:bc-commuting-normal-form-signature}
Let $(\gamma;e_1,\ldots,e_{m-1})$ be a datum of
\cref{def:bc-odd-fixed-even-standard-tableaux}.  For
$S\in\mathcal S(\gamma;e_1,\ldots,e_{m-1})$ define its commuting
normal-form signature $\kappa(S)$ recursively as follows.

If $m=1$, the signature is the empty word.  If $m\ge2$, let $y_m$ be the
cell carrying $2m-1$ in $S$, put $e=e_{m-1}$, and let
\[
  \gamma_{y_m}=\gamma\setminus\{y_m,e\}.
\]
Let $S^-$ be the restriction of $S$ to $\gamma_{y_m}$, with entries
$1,\ldots,2m-3$.  The first symbol of $\kappa(S)$ is $y_m$ if $e$ is a
removable corner of $\gamma$, and is the symbol $\ast$ otherwise.  The
remaining symbols are $\kappa(S^-)$, computed for the smaller datum
$(\gamma_{y_m};e_1,\ldots,e_{m-2})$.

For a word $\sigma$ in these symbols, let
\[
  \mathcal S_\sigma(\gamma;e_1,\ldots,e_{m-1})
\]
denote the set of tableaux in
$\mathcal S(\gamma;e_1,\ldots,e_{m-1})$ with commuting normal-form signature
$\sigma$.
\end{definition}

\begin{lemma}[One-step recovery of a commuting signature symbol]\label{lem:bc-one-step-commuting-signature-symbol}
Let $m\ge2$, let
$S\in\mathcal S(\gamma;e_1,\ldots,e_{m-1})$, let $y_m$ be the cell carrying
$2m-1$, put $e=e_{m-1}$, and put
\[
  \gamma^-=\gamma\setminus\{y_m,e\}.
\]
Then the first symbol of $\kappa(S)$ is $\ast$ if and only if $e$ is not a
removable corner of $\gamma$.  If $e$ is a removable corner of $\gamma$, then
the first symbol of $\kappa(S)$ is the unique cell in
\[
  \gamma\setminus(\gamma^-\cup\{e\}).
\]
Consequently, a parser that determines $\gamma$, $e$, and, in the commuting
case, the lower shape $\gamma^-$ determines the first symbol of the commuting
normal-form signature.
\end{lemma}

\begin{proof}
The first assertion is exactly the recursive definition of
\cref{def:bc-commuting-normal-form-signature}: the first symbol is $\ast$
precisely in the non-removable case, and otherwise it is the top odd cell
$y_m$.

In the removable case,
\[
  \gamma^-=\gamma\setminus\{y_m,e\}.
\]
The cells $y_m$ and $e$ are distinct, because $e$ is a fixed even cell while
$y_m$ carries the odd entry $2m-1$.  Hence
\[
  \gamma\setminus(\gamma^-\cup\{e\})=\{y_m\},
\]
which proves the displayed recovery rule.  The final sentence follows
immediately from the two cases.
\end{proof}

\begin{definition}[Odd sublevel shape chain]\label{def:bc-odd-sublevel-shape-chain}
Let $Q\in\mathcal S(\gamma;e_1,\ldots,e_{m-1})$.  For
$1\le r\le m$, define
\[
  \Gamma_r(Q)=\{c\in\gamma: Q(c)\le2r-1\}.
\]
Thus
\[
  \Gamma_1(Q)\subset\Gamma_2(Q)\subset\cdots\subset\Gamma_m(Q)=\gamma
\]
is the chain of odd sublevel shapes of the fixed-even standard tableau.
\end{definition}

\begin{proposition}[Odd sublevel shapes determine the commuting signature]\label{prop:bc-odd-sublevel-chain-determines-signature}
For every datum of \cref{def:bc-odd-fixed-even-standard-tableaux}, the
commuting normal-form signature $\kappa(Q)$ of
$Q\in\mathcal S(\gamma;e_1,\ldots,e_{m-1})$ is determined by
\[
  \gamma,\quad e_1,\ldots,e_{m-1},\quad
  \Gamma_1(Q),\ldots,\Gamma_m(Q).
\]
\end{proposition}

\begin{proof}
We argue by induction on $m$.  For $m=1$ the signature is empty by
\cref{def:bc-commuting-normal-form-signature}.

Assume $m\ge2$.  Since $\Gamma_m(Q)=\gamma$, the data determine the current
shape.  They also determine the current fixed even cell
$e=e_{m-1}$ and the lower shape
\[
  \gamma^-=\Gamma_{m-1}(Q).
\]
Indeed $\Gamma_{m-1}(Q)$ is exactly the shape obtained from $\gamma$ by
removing the cells carrying the two largest entries $2m-2$ and $2m-1$,
namely $e_{m-1}$ and the top odd cell.  Therefore
\cref{lem:bc-one-step-commuting-signature-symbol} determines the first symbol
of $\kappa(Q)$ from the displayed data.

The restriction of $Q$ to $\Gamma_{m-1}(Q)$ is an element of
\[
  \mathcal S(\Gamma_{m-1}(Q);e_1,\ldots,e_{m-2}),
\]
and its odd sublevel shape chain is
\[
  \Gamma_1(Q),\ldots,\Gamma_{m-1}(Q).
\]
By the induction hypothesis, these data determine the tail of the commuting
normal-form signature.  Together with the first symbol just determined, this
determines all of $\kappa(Q)$.
\end{proof}

\begin{corollary}[Target-window odd sublevel recovery closes the signature input]\label{cor:bc-target-window-odd-sublevel-chain-closes-signature}
Work in the setting of
\cref{prop:bc-target-window-signature-closes-refined-suffix}.  Suppose that
the data
\[
  S,\quad Z_L,\quad V^\ast_L,\quad \mathcal D_{\mathrm{out}}(P),
  \quad \rho_{2a_L-2},\rho_{2j}
\]
determine the associated fixed-even datum
\[
  (\gamma;e_1,\ldots,e_{m-1})
\]
and the odd sublevel shape chain
\[
  \Gamma_1,\ldots,\Gamma_m
\]
of the fixed-even standard tableau associated to the true odd-packet
assignment.  Then the hypothesis of
\cref{prop:bc-target-window-signature-closes-refined-suffix} holds.
\end{corollary}

\begin{proof}
By \cref{prop:bc-odd-sublevel-chain-determines-signature}, the recovered
fixed-even datum together with the recovered odd sublevel shape chain
determines the commuting normal-form signature.  Thus the
deterministic-function hypothesis of
\cref{prop:bc-target-window-signature-closes-refined-suffix} is satisfied.
\end{proof}

\begin{proposition}[Odd sublevels are reduced half-step diagrams]\label{prop:bc-odd-sublevels-are-reduced-halfsteps}
Let $(\overline\nu;\overline E_1,\ldots,\overline E_m)$ be a reduced
fixed-even datum arising from a nonempty fixed-cell completion set, and let
\[
  y=(y_1,\ldots,y_m)\in
  \mathcal R(\overline\nu;\overline E_1,\ldots,\overline E_m).
\]
Let
\[
  \alpha_0,\alpha_{\frac12},\alpha_1,\ldots,\alpha_m
\]
be the forward diagrams defined by
\cref{def:bc-reduced-fixed-even-growth-chains}.  Let
$Q\in\mathcal S(\gamma;e_1,\ldots,e_{m-1})$ be the fixed-even standard
tableau corresponding to $y$ under
\[
  \mathcal R
  \longleftrightarrow
  \mathcal T
  \longleftrightarrow
  \mathcal S
\]
from
\cref{lem:bc-reduced-growth-reverse-corner-tree,lem:bc-reverse-tree-fixed-even-standard-tableaux}.
Then, for every $1\le r\le m$,
\[
  \Gamma_r(Q)=\alpha_{r-\frac12}.
\]
\end{proposition}

\begin{proof}
By \cref{lem:bc-reverse-tree-fixed-even-standard-tableaux}, the tableau $Q$
places the odd entry $2s-1$ in $y_s$ for $1\le s\le m$ and the fixed even
entry $2s$ in $e_s$, the unique cell of $\overline E_s$, for
$1\le s<m$.  Therefore the sublevel set $\Gamma_r(Q)$ consists exactly of
the cells
\[
  y_1,\ldots,y_r,\qquad e_1,\ldots,e_{r-1}.
\]
On the other hand, the recursion in
\cref{def:bc-reduced-fixed-even-growth-chains} gives
\[
  \alpha_{r-\frac12}
  =
  \{y_1,\ldots,y_r\}\cup
  \overline E_1\cup\cdots\cup\overline E_{r-1}.
\]
For $s<m$, \cref{lem:bc-reduced-even-strip-sizes} identifies
$\overline E_s=\{e_s\}$.  Hence the two displayed sets agree for
$r<m$.

For $r=m$, the same recursion gives
\[
  \alpha_{m-\frac12}
  =
  \overline\nu\setminus\overline E_m
  =
  \gamma,
\]
because $\alpha_m=\overline\nu$ and
$\alpha_m=\alpha_{m-\frac12}\cup\overline E_m$ is a disjoint union.  Also
$\Gamma_m(Q)=\gamma$ by
\cref{def:bc-odd-sublevel-shape-chain}.  Thus the identity holds for every
$r$.
\end{proof}

\begin{lemma}[Retained even strips determine the reduced fixed-even datum]\label{lem:bc-retained-even-strips-determine-reduced-datum}
Let $(\nu;E_1,\ldots,E_m)$ be a tight-prefix datum with
\[
  \mathcal C(\nu;E_1,\ldots,E_m)\ne\varnothing .
\]
If a parser-visible datum determines $\nu$ and the labelled retained even
strips $E_1,\ldots,E_m$, then it determines the reduced fixed-even datum
\[
  (\overline\nu;\overline E_1,\ldots,\overline E_m),
\]
the indicators $\varepsilon_1,\ldots,\varepsilon_m$, and the associated
fixed-even standard-tableau datum
\[
  (\gamma;e_1,\ldots,e_{m-1}).
\]
\end{lemma}

\begin{proof}
The reduced shape $\overline\nu$ is obtained from $\nu$ by deleting the first
column and shifting every remaining cell one column left.  For each retained
even strip $E_r$, the reduced strip is the deterministic set
\[
  \overline E_r=\{\overline c:\ c\in E_r,\ c\ne(2r,1)\},
\]
and
\[
  \varepsilon_r=
  \begin{cases}
  1,&(2r,1)\in E_r,\\
  0,&(2r,1)\notin E_r.
  \end{cases}
\]
Thus $\nu$ and the labelled strips determine
$(\overline\nu;\overline E_1,\ldots,\overline E_m)$ and all
$\varepsilon_r$.

Since the completion set is nonempty,
\cref{lem:bc-reduced-even-strip-sizes} says that
$|\overline E_r|=1$ for $1\le r<m$.  Hence each fixed even cell
$e_r$ is the unique cell of $\overline E_r$.  Finally,
\[
  \gamma=\overline\nu\setminus\overline E_m
\]
by \cref{def:bc-odd-fixed-even-standard-tableaux}.  Therefore the same data
determine $(\gamma;e_1,\ldots,e_{m-1})$.
\end{proof}

\begin{lemma}[Pieri prefix shapes are reduced half-step diagrams]\label{lem:bc-pieri-prefix-shapes-are-reduced-halfsteps}
Work in the non-full tight height setting of
\cref{cor:bc-height-zero-surplus-label-prefix}.  Let
\[
  \lambda^{(0)}\subset\lambda^{(1)}\subset\cdots\subset\lambda^{(2m)}
\]
be the true tight-prefix Pieri shape chain, let
\[
  \nu=\lambda^{(2m)}
\]
be its endpoint, and let $E_r$ be the two boxes carrying the retained even
label $2r$.  For each $r$, let $x_r$ be the second box carrying the deleted
odd label $2r-1$, so that the first one is
$f_r=(2r-1,1)$, and put $y_r=\overline{x_r}$.
Let
\[
  \alpha_0,\alpha_{\frac12},\alpha_1,\ldots,\alpha_m
\]
be the reduced forward diagrams attached to
$y=(y_1,\ldots,y_m)$ by
\cref{def:bc-reduced-fixed-even-growth-chains}.  For a Young diagram $\mu$,
write
\[
  \operatorname{del}_1(\mu)
  =
  \{\overline c:\ c\in\mu,\ c\text{ is not in the first column}\}.
\]
Then, for every $1\le r\le m$,
\[
  \operatorname{del}_1(\lambda^{(2r-1)})=\alpha_{r-\frac12},
  \qquad
  \operatorname{del}_1(\lambda^{(2r)})=\alpha_r .
\]
\end{lemma}

\begin{proof}
By \cref{lem:bc-true-tight-packets-in-omega}, the tuple
$(x_1,\ldots,x_m)$ belongs to
$\Omega(\nu;E_1,\ldots,E_m)$.  Under
\cref{lem:bc-omega-fixed-cell-completions,prop:bc-completions-equal-reduced-growth-chains},
its first-column deletion is exactly the element
$y=(y_1,\ldots,y_m)$ of
$\mathcal R(\overline\nu;\overline E_1,\ldots,\overline E_m)$ used to define
the displayed diagrams $\alpha$.

For a fixed $r$, the shape $\lambda^{(2r-1)}$ consists of the forced
first-column odd boxes $f_1,\ldots,f_r$, the off-first-column odd boxes
$x_1,\ldots,x_r$, and the even strips $E_1,\ldots,E_{r-1}$.  Deleting the
first column removes the forced odd boxes and, from each earlier even strip,
removes exactly its first-column member when one is present.  Therefore
\[
  \operatorname{del}_1(\lambda^{(2r-1)})
  =
  \{y_1,\ldots,y_r\}\cup
  \overline E_1\cup\cdots\cup\overline E_{r-1}.
\]
This is $\alpha_{r-\frac12}$ by the recursion in
\cref{def:bc-reduced-fixed-even-growth-chains}.

Similarly, $\lambda^{(2r)}$ is obtained from $\lambda^{(2r-1)}$ by adjoining
the retained even strip $E_r$.  Hence
\[
  \operatorname{del}_1(\lambda^{(2r)})
  =
  \{y_1,\ldots,y_r\}\cup
  \overline E_1\cup\cdots\cup\overline E_r
  =
  \alpha_r,
\]
again by the same recursion.
\end{proof}

\begin{corollary}[Target-window tight-prefix recovery closes the signature input]\label{cor:bc-target-window-prefix-path-closes-signature}
Work in the setting of
\cref{prop:bc-target-window-signature-closes-refined-suffix}.  Suppose that
the data
\[
  S,\quad Z_L,\quad V^\ast_L,\quad \mathcal D_{\mathrm{out}}(P),
  \quad \rho_{2a_L-2},\rho_{2j}
\]
determine the true tight-prefix Pieri shape chain
\[
  \lambda^{(0)}\subset\lambda^{(1)}\subset\cdots\subset\lambda^{(2m)} .
\]
Then the hypothesis of
\cref{prop:bc-target-window-signature-closes-refined-suffix} holds.
\end{corollary}

\begin{proof}
The recovered shape chain determines
\[
  \nu=\lambda^{(2m)}
\]
and, for $1\le r\le m$, the retained even strip
\[
  E_r=\lambda^{(2r)}\setminus\lambda^{(2r-1)}.
\]
By \cref{lem:bc-pieri-prefix-shapes-are-reduced-halfsteps}, the same recovered
chain determines the reduced half-step diagrams
\[
  \alpha_{\frac12},\alpha_{\frac32},\ldots,\alpha_{m-\frac12}.
\]
Thus the hypotheses of
\cref{cor:bc-target-window-retained-even-halfstep-closes-signature} are
satisfied, and that corollary gives the deterministic commuting-signature
input.
\end{proof}

\begin{corollary}[Target-window standard even chain recovery closes the signature input]\label{cor:bc-target-window-standard-even-chain-closes-signature}
Work in the setting of
\cref{prop:bc-target-window-signature-closes-refined-suffix}, and let
$U=\operatorname{std}(T)$ be the doubled standardization of the underlying
semistandard tableau.  Let
\[
  \varnothing=\sigma^{(0)}\subset\sigma^{(1)}\subset\cdots
\]
be the standard shape chain of $U$.  Suppose that the data
\[
  S,\quad Z_L,\quad V^\ast_L,\quad \mathcal D_{\mathrm{out}}(P),
  \quad \rho_{2a_L-2},\rho_{2j}
\]
determine the even standard-time endpoint chain
\[
  \sigma^{(0)},\sigma^{(2)},\sigma^{(4)},\ldots,\sigma^{(4m)} .
\]
Then the hypothesis of
\cref{prop:bc-target-window-signature-closes-refined-suffix} holds.
\end{corollary}

\begin{proof}
By \cref{lem:bc-standard-time-pair-deletion}, the even standard-time endpoints
are exactly the semistandard Pieri shapes:
\[
  \sigma^{(2a)}=\lambda^{(a)}
  \qquad(0\le a\le2m).
\]
Thus the recovered standard even chain determines the true tight-prefix Pieri
shape chain
\[
  \lambda^{(0)}\subset\lambda^{(1)}\subset\cdots\subset\lambda^{(2m)} .
\]
Applying \cref{cor:bc-target-window-prefix-path-closes-signature} gives the
deterministic commuting-signature input.
\end{proof}

\begin{corollary}[Target-window dual-RSK' carrier recovery closes the signature input]\label{cor:bc-target-window-dual-rsk-carrier-closes-signature}
Work in the setting of
\cref{prop:bc-target-window-signature-closes-refined-suffix}.  Use the
Step~3 dual-RSK' half-square growth diagram of
\cref{lem:bc-krattenthaler-semistandard-path-boundary} for the underlying
semistandard tableau, and let
\[
  \varnothing=\sigma^{(0)}\subset\sigma^{(1)}\subset\cdots
\]
be the standard shape chain of its doubled standardization.  Suppose that the
data
\[
  S,\quad Z_L,\quad V^\ast_L,\quad \mathcal D_{\mathrm{out}}(P),
  \quad \rho_{2a_L-2},\rho_{2j}
\]
determine a Ferrers subarrangement $E$ of this Step~3 diagram and determine
the partition labels on the top/right boundary of $E$.  Suppose moreover that
the recovered cell entries and internal corner labels of $E$ determine the
even standard-time endpoints
\[
  \sigma^{(0)},\sigma^{(2)},\sigma^{(4)},\ldots,\sigma^{(4m)}
\]
by a deterministic readout fixed by the same data.  Then the hypothesis of
\cref{prop:bc-target-window-signature-closes-refined-suffix} holds.
\end{corollary}

\begin{proof}
By \cref{lem:bc-krattenthaler-semistandard-path-boundary}, the Step~3 diagram
is the fourth arbitrary-filling dual-RSK' growth diagram.  Since the displayed
data determine the shape of $E$ and its top/right boundary labels,
\cref{cor:bc-parsed-dual-rsk-prime-boundary-recovers-subfilling} recovers,
by a deterministic backward sweep, every cell entry and every corner label on
or inside $E$.  The deterministic readout in the hypothesis therefore recovers
the listed even standard-time endpoints.  Applying
\cref{cor:bc-target-window-standard-even-chain-closes-signature} gives the
deterministic commuting-signature input.
\end{proof}

\begin{corollary}[Low-label carrier readout closes the signature input]\label{cor:bc-target-window-low-label-carrier-closes-signature}
Work in the setting of
\cref{prop:bc-target-window-signature-closes-refined-suffix}.  Use the
Step~3 dual-RSK' half-square growth diagram of
\cref{lem:bc-krattenthaler-semistandard-path-boundary} for the underlying
semistandard tableau $T$.  Suppose that the data
\[
  S,\quad Z_L,\quad V^\ast_L,\quad \mathcal D_{\mathrm{out}}(P),
  \quad \rho_{2a_L-2},\rho_{2j}
\]
determine a Ferrers subarrangement $E$ of this Step~3 diagram and determine
the partition labels on the top/right boundary of $E$.  Suppose moreover that
the recovered cell entries and internal corner labels of $E$ determine, for
each semistandard label $1\le a\le2m$, the two boxes of $T$ carrying label
$a$.  Then the hypothesis of
\cref{prop:bc-target-window-signature-closes-refined-suffix} holds.
\end{corollary}

\begin{proof}
By \cref{lem:bc-krattenthaler-semistandard-path-boundary}, the Step~3 diagram
is the fourth arbitrary-filling dual-RSK' growth diagram.  The displayed data
determine $E$ and its top/right boundary labels, so
\cref{cor:bc-parsed-dual-rsk-prime-boundary-recovers-subfilling} recovers the
cell entries and internal corner labels of $E$.  The low-label readout in the
hypothesis therefore recovers the two boxes carrying each semistandard label
$1,\ldots,2m$.  Applying
\cref{lem:bc-low-label-boxes-standard-readout} with $M=2m$ gives
\[
  \sigma^{(0)},\sigma^{(2)},\sigma^{(4)},\ldots,\sigma^{(4m)}.
\]
Thus the deterministic readout hypothesis of
\cref{cor:bc-target-window-dual-rsk-carrier-closes-signature} is satisfied,
and that corollary gives the deterministic commuting-signature input.
\end{proof}

\begin{corollary}[Step-2 prefix carrier criterion for target-window signatures]\label{cor:bc-target-window-step2-prefix-carrier-closes-signature}
Work in the setting of
\cref{prop:bc-target-window-signature-closes-refined-suffix}.  Use the
Step~2 symmetric matrix and the Step~3 dual-RSK' half-square growth diagram
from \cite[Theorem 4, Steps 2--3]{Krattenthaler2016}.  Suppose that the data
\[
  S,\quad Z_L,\quad V^\ast_L,\quad \mathcal D_{\mathrm{out}}(P),
  \quad \rho_{2a_L-2},\rho_{2j}
\]
determine a Ferrers subarrangement $E$ of the Step~3 diagram and determine
the partition labels on the top/right boundary of $E$.  Suppose moreover that
the recovered cell entries and internal corner labels of $E$ determine a
Ferrers subarrangement $C$ of the Step~2 first-variant Knuth growth diagram,
the cell entries of $C$, and the left/bottom boundary labels of $C$, such
that the forward boundary sweep on $C$ contains
\[
  \lambda^{(0)},\lambda^{(1)},\ldots,\lambda^{(2m)}
\]
for the underlying semistandard tableau.  Then the hypothesis of
\cref{prop:bc-target-window-signature-closes-refined-suffix} holds.
\end{corollary}

\begin{proof}
By \cref{lem:bc-krattenthaler-semistandard-path-boundary}, the Step~3 diagram
is the fourth arbitrary-filling dual-RSK' growth diagram.  Since the displayed
data determine $E$ and its top/right boundary labels,
\cref{cor:bc-parsed-dual-rsk-prime-boundary-recovers-subfilling} recovers the
cell entries and internal corner labels of $E$.  The additional hypothesis
therefore determines the Step~2 subarrangement $C$, its cell entries, and its
left/bottom boundary labels.  Applying
\cref{lem:bc-step2-knuth-prefix-low-label-boxes} with $M=2m$ recovers the two
boxes carrying each semistandard label $1,\ldots,2m$.  Hence the hypothesis
of \cref{cor:bc-target-window-low-label-carrier-closes-signature} is
satisfied, and that corollary gives the deterministic commuting-signature
input.
\end{proof}

\begin{corollary}[Step-3 half-matrix prefix criterion for target-window signatures]\label{cor:bc-target-window-step3-half-matrix-prefix-closes-signature}
Work in the setting of
\cref{prop:bc-target-window-signature-closes-refined-suffix}.  Use the
Step~2 symmetric matrix and the Step~3 dual-RSK' half-square growth diagram
from \cite[Theorem 4, Steps 2--3]{Krattenthaler2016}.  Suppose that the data
\[
  S,\quad Z_L,\quad V^\ast_L,\quad \mathcal D_{\mathrm{out}}(P),
  \quad \rho_{2a_L-2},\rho_{2j}
\]
determine a Ferrers subarrangement $E$ of the Step~3 diagram and determine
the partition labels on the top/right boundary of $E$.  Suppose moreover that
the recovered cell entries and internal corner labels of $E$ determine a
Ferrers subarrangement $C$ of the Step~2 first-variant Knuth growth diagram and
the left/bottom boundary labels of $C$, such that every non-diagonal cell of
$C$, or its transpose in the Step~2 symmetric matrix, lies in the recovered
Step~3 subarrangement $E$, and such that the forward boundary sweep on $C$
contains
\[
  \lambda^{(0)},\lambda^{(1)},\ldots,\lambda^{(2m)}
\]
for the underlying semistandard tableau.  Then the hypothesis of
\cref{prop:bc-target-window-signature-closes-refined-suffix} holds.
\end{corollary}

\begin{proof}
By \cref{lem:bc-krattenthaler-semistandard-path-boundary}, the Step~3 diagram
is the fourth arbitrary-filling dual-RSK' growth diagram.  Since the displayed
data determine $E$ and its top/right boundary labels,
\cref{cor:bc-parsed-dual-rsk-prime-boundary-recovers-subfilling} recovers the
cell entries and internal corner labels of $E$.  The additional hypothesis
therefore determines $C$ and its left/bottom boundary labels.  By
\cref{lem:bc-step3-cells-are-step2-half-matrix}, the recovered Step~3
subfilling determines every Step~2 cell entry of $C$, using the symmetry of the
Step~2 matrix for non-diagonal cells whose transposes lie in $E$; diagonal
entries of $C$ are zero in the Step~2 matrix.  Thus the hypotheses of
\cref{cor:bc-target-window-step2-prefix-carrier-closes-signature} are
satisfied, and that corollary gives the deterministic commuting-signature
input.
\end{proof}

\begin{corollary}[Step-2 Q-prefix strip criterion for target-window signatures]\label{cor:bc-target-window-step2-q-prefix-strip-closes-signature}
Work in the setting of
\cref{prop:bc-target-window-signature-closes-refined-suffix}.  Use the
Step~2 symmetric matrix and the Step~3 dual-RSK' half-square growth diagram
from \cite[Theorem 4, Steps 2--3]{Krattenthaler2016}.  Suppose that the data
\[
  S,\quad Z_L,\quad V^\ast_L,\quad \mathcal D_{\mathrm{out}}(P),
  \quad \rho_{2a_L-2},\rho_{2j}
\]
determine a Ferrers subarrangement $E$ of the Step~3 diagram and determine
the partition labels on the top/right boundary of $E$.  Suppose moreover that,
for the Step~2 $Q$-prefix strip $C_{2m}$ of
\cref{lem:bc-step2-q-prefix-strip-determines-prefix}, every non-diagonal cell
of $C_{2m}$, or its transpose in the Step~2 symmetric matrix, lies in the
recovered Step~3 subarrangement $E$.  Then the hypothesis of
\cref{prop:bc-target-window-signature-closes-refined-suffix} holds.
\end{corollary}

\begin{proof}
The displayed data recover the Step~3 subfilling of $E$ by
\cref{cor:bc-parsed-dual-rsk-prime-boundary-recovers-subfilling}.  The
containment hypothesis and \cref{lem:bc-step3-cells-are-step2-half-matrix}
therefore determine every non-diagonal Step~2 entry in the $Q$-prefix strip
$C_{2m}$, and the diagonal entries of that strip are known to be zero.
Applying \cref{lem:bc-step2-q-prefix-strip-determines-prefix} with $M=2m$
recovers
\[
  \lambda^{(0)},\lambda^{(1)},\ldots,\lambda^{(2m)}.
\]
Then \cref{lem:bc-step2-knuth-prefix-low-label-boxes} recovers the two boxes
carrying each semistandard label $1,\ldots,2m$, so the hypothesis of
\cref{cor:bc-target-window-low-label-carrier-closes-signature} is satisfied.
That corollary gives the deterministic commuting-signature input.
\end{proof}

\begin{corollary}[Low-label incidence carrier criterion for target-window signatures]\label{cor:bc-target-window-low-label-incidence-carrier-closes-signature}
Work in the setting of
\cref{prop:bc-target-window-signature-closes-refined-suffix}.  Use the
Step~2 symmetric matrix and the Step~3 dual-RSK' half-square growth diagram
from \cite[Theorem 4, Steps 2--3]{Krattenthaler2016}.  Suppose that the data
\[
  S,\quad Z_L,\quad V^\ast_L,\quad \mathcal D_{\mathrm{out}}(P),
  \quad \rho_{2a_L-2},\rho_{2j}
\]
determine a Ferrers subarrangement $E$ of the Step~3 diagram and determine
the partition labels on the top/right boundary of $E$.  Suppose moreover that
$E$ contains the low-label incidence cell set $H_{2m}$ of
\cref{lem:bc-step3-low-label-incidence-covers-q-prefix}.  Then the hypothesis
of \cref{prop:bc-target-window-signature-closes-refined-suffix} holds.
\end{corollary}

\begin{proof}
By \cref{lem:bc-step3-low-label-incidence-covers-q-prefix}, the containment
$H_{2m}\subseteq E$ implies that every non-diagonal cell of the Step~2
$Q$-prefix strip $C_{2m}$ has its cell, or its transpose, in $E$.  Hence
\cref{cor:bc-target-window-step2-q-prefix-strip-closes-signature} applies and
gives the deterministic commuting-signature input.
\end{proof}

\begin{corollary}[Low-label incidence entries close target-window signatures]\label{cor:bc-target-window-low-label-incidence-entries-close-signature}
Work in the setting of
\cref{prop:bc-target-window-signature-closes-refined-suffix}.  Use the
Step~2 symmetric matrix and the Step~3 dual-RSK' half-square growth diagram
from \cite[Theorem 4, Steps 2--3]{Krattenthaler2016}.  Suppose that the data
\[
  S,\quad Z_L,\quad V^\ast_L,\quad \mathcal D_{\mathrm{out}}(P),
  \quad \rho_{2a_L-2},\rho_{2j}
\]
determine the Step~3 cell entries on every cell of $H_{2m}$.  Then the
hypothesis of \cref{prop:bc-target-window-signature-closes-refined-suffix}
holds.
\end{corollary}

\begin{proof}
\Cref{lem:bc-step3-low-label-incidence-covers-q-prefix} shows that the
recovered entries on $H_{2m}$ determine every non-diagonal Step~2 entry in the
$Q$-prefix strip $C_{2m}$, while the diagonal entries of that strip are known
zero by \cref{lem:bc-step3-cells-are-step2-half-matrix}.  Therefore
\cref{lem:bc-step2-q-prefix-strip-determines-prefix} with $M=2m$ recovers
\[
  \lambda^{(0)},\lambda^{(1)},\ldots,\lambda^{(2m)}.
\]
Applying \cref{lem:bc-step2-knuth-prefix-low-label-boxes} then recovers the two
boxes carrying each label $1,\ldots,2m$.  Hence
\cref{cor:bc-target-window-low-label-carrier-closes-signature} gives the
deterministic commuting-signature input.
\end{proof}

\begin{corollary}[Strict-left data reduce the incidence target to low-low entries]\label{cor:bc-strict-left-reduces-incidence-to-low-low}
Work in the non-full tight-height setting of
\cref{lem:bc-tight-outside-recovers-left-growth-segment} and in the setting of
\cref{prop:bc-target-window-signature-closes-refined-suffix}, with
$Z_L=\{j-2m+1,\ldots,j\}$.  Let $L_{2m}$ be the low-low cell set of
\cref{lem:bc-low-label-incidence-split}.  Suppose that the target-window data
\[
  S,\quad Z_L,\quad V^\ast_L,\quad \mathcal D_{\mathrm{out}}(P),
  \quad \rho_{2a_L-2},\rho_{2j}
\]
determine the Step~3 cell entries on every cell of $L_{2m}$.  Then the
hypothesis of \cref{prop:bc-target-window-signature-closes-refined-suffix}
holds.
\end{corollary}

\begin{proof}
By \cref{lem:bc-low-label-incidence-split} with $M=2m$,
\[
  H_{2m}=X_{2m}\sqcup L_{2m}.
\]
The same lemma places $X_{2m}$ inside the continuation-column subarrangement
$J_{2m}$, and \cref{lem:bc-strict-left-recovers-continuation-columns} says that
$\mathcal D_{\mathrm{out}}(P)$ determines all entries of $J_{2m}$.  Hence the
full target-window tuple determines all entries on $X_{2m}$.  By hypothesis it
also determines all entries on $L_{2m}$, and therefore all entries on
$H_{2m}$.  Applying
\cref{cor:bc-target-window-low-label-incidence-entries-close-signature} gives
the deterministic commuting-signature input.
\end{proof}

\begin{corollary}[Terminal low-label boundary closes target-window signatures]\label{cor:bc-target-window-low-low-boundary-closes-signature}
Work in the non-full tight-height setting of
\cref{lem:bc-tight-outside-recovers-left-growth-segment} and in the setting of
\cref{prop:bc-target-window-signature-closes-refined-suffix}, with
$Z_L=\{j-2m+1,\ldots,j\}$.  Suppose that the target-window data
\[
  S,\quad Z_L,\quad V^\ast_L,\quad \mathcal D_{\mathrm{out}}(P),
  \quad \rho_{2a_L-2},\rho_{2j}
\]
determine the terminal low-label Step~3 boundary segment
\[
  \rho_{2(j-2m)},\rho_{2(j-2m)+1},\ldots,\rho_{2j}.
\]
Then the hypothesis of
\cref{prop:bc-target-window-signature-closes-refined-suffix} holds.
\end{corollary}

\begin{proof}
By \cref{lem:bc-low-low-terminal-corner-boundary} with $M=2m$, the displayed
boundary segment determines every Step~3 cell entry on $L_{2m}$.  Therefore
the hypothesis of
\cref{cor:bc-strict-left-reduces-incidence-to-low-low} is satisfied, and that
corollary gives the deterministic commuting-signature input.
\end{proof}

\begin{corollary}[Branch-free terminal-zone decoding closes target-window signatures]\label{cor:bc-target-window-terminal-zone-decoder-closes-signature}
Work in the non-full tight-height setting of
\cref{lem:bc-tight-outside-recovers-left-growth-segment} and in the setting of
\cref{prop:bc-target-window-signature-closes-refined-suffix}, with
$Z_L=\{a_L,\ldots,j\}=\{j-2m+1,\ldots,j\}$.  Suppose that the target-window
data
\[
  S,\quad Z_L,\quad V^\ast_L,\quad \mathcal D_{\mathrm{out}}(P),
  \quad \rho_{2a_L-2},\rho_{2j}
\]
determine the original right-boundary zone subpath
\[
  \rho_{2a_L-2},\rho_{2a_L-1},\ldots,\rho_{2j}.
\]
Then the hypothesis of
\cref{prop:bc-target-window-signature-closes-refined-suffix} holds.
\end{corollary}

\begin{proof}
The identity $a_L=j-2m+1$ gives
\[
  2a_L-2=2(j-2m).
\]
Thus the displayed right-boundary zone subpath is exactly the terminal
low-label boundary segment
\[
  \rho_{2(j-2m)},\rho_{2(j-2m)+1},\ldots,\rho_{2j}.
\]
Applying
\cref{cor:bc-target-window-low-low-boundary-closes-signature} gives the
deterministic commuting-signature input.
\end{proof}

\begin{corollary}[Strict-left carrier criterion for target-window signatures]\label{cor:bc-target-window-strict-left-carrier-closes-signature}
Work in the non-full tight-height setting of
\cref{lem:bc-tight-outside-recovers-left-growth-segment} and in the setting of
\cref{prop:bc-target-window-signature-closes-refined-suffix}, so the final
right-boundary zone is
\[
  Z_L=\{a_L,\ldots,j\}=\{j-2m+1,\ldots,j\}.
\]
Suppose that, from
\[
  S,\quad Z_L,\quad V^\ast_L,\quad
  \rho_{2a_L-2},\rho_{2j},
  \quad
  \rho_0,\rho_1,\ldots,\rho_{2(j-2m)},
\]
a deterministic rule identifies a Ferrers subarrangement $E$ of the Step~3
dual-RSK' half-square of
\cref{lem:bc-krattenthaler-semistandard-path-boundary}, identifies the
partition labels on its top/right boundary, and identifies a deterministic
readout which, from the recovered cell entries and internal corner labels of
$E$, returns a Step~2 first-variant Knuth subarrangement $C$ and its
left/bottom boundary labels, such that every non-diagonal cell of $C$, or its
transpose in the Step~2 symmetric matrix, lies in $E$, and such that the
forward boundary sweep on $C$ contains
\[
  \lambda^{(0)},\lambda^{(1)},\ldots,\lambda^{(2m)} .
\]
Then the hypothesis of
\cref{prop:bc-target-window-signature-closes-refined-suffix} holds.
\end{corollary}

\begin{proof}
By \cref{lem:bc-tight-outside-recovers-left-growth-segment}, the outside datum
$\mathcal D_{\mathrm{out}}(P)$ determines
\[
  \rho_0,\rho_1,\ldots,\rho_{2(j-2m)}.
\]
Therefore the displayed target-window tuple determines every input to the
deterministic rule in the hypothesis once it is enlarged by
$\mathcal D_{\mathrm{out}}(P)$, as in
\cref{prop:bc-target-window-signature-closes-refined-suffix}.  It follows
that the full tuple
\[
  S,\quad Z_L,\quad V^\ast_L,\quad \mathcal D_{\mathrm{out}}(P),
  \quad \rho_{2a_L-2},\rho_{2j}
\]
determines the carrier $E$, its top/right boundary labels, and the
deterministic Step~2 half-matrix prefix readout required in
\cref{cor:bc-target-window-step3-half-matrix-prefix-closes-signature}.
Applying
\cref{cor:bc-target-window-step3-half-matrix-prefix-closes-signature} gives the
deterministic commuting-signature input.
\end{proof}

\begin{corollary}[Strict-left Q-prefix strip criterion for target-window signatures]\label{cor:bc-target-window-strict-left-q-prefix-strip-closes-signature}
Work in the non-full tight-height setting of
\cref{lem:bc-tight-outside-recovers-left-growth-segment} and in the setting of
\cref{prop:bc-target-window-signature-closes-refined-suffix}, with
$Z_L=\{j-2m+1,\ldots,j\}$.  Suppose that, from
\[
  S,\quad Z_L,\quad V^\ast_L,\quad
  \rho_{2a_L-2},\rho_{2j},
  \quad
  \rho_0,\rho_1,\ldots,\rho_{2(j-2m)},
\]
a deterministic rule identifies a Ferrers subarrangement $E$ of the Step~3
dual-RSK' half-square of
\cref{lem:bc-krattenthaler-semistandard-path-boundary}, identifies the
partition labels on its top/right boundary, and verifies that every cell of
the non-diagonal part of the Step~2 $Q$-prefix strip $C_{2m}$ has its cell, or
its transpose in the Step~2 symmetric matrix, in $E$.  Then the hypothesis of
\cref{prop:bc-target-window-signature-closes-refined-suffix} holds.
\end{corollary}

\begin{proof}
By \cref{lem:bc-tight-outside-recovers-left-growth-segment},
$\mathcal D_{\mathrm{out}}(P)$ determines
\[
  \rho_0,\rho_1,\ldots,\rho_{2(j-2m)}.
\]
Thus the full target-window tuple determines every input to the deterministic
rule.  It follows that the tuple determines $E$, its top/right boundary
labels, and the $Q$-prefix strip containment hypothesis required in
\cref{cor:bc-target-window-step2-q-prefix-strip-closes-signature}.  Applying
that corollary gives the deterministic commuting-signature input.
\end{proof}

\begin{corollary}[Strict-left low-label incidence carrier criterion]\label{cor:bc-target-window-strict-left-low-label-incidence-closes-signature}
Work in the non-full tight-height setting of
\cref{lem:bc-tight-outside-recovers-left-growth-segment} and in the setting of
\cref{prop:bc-target-window-signature-closes-refined-suffix}, with
$Z_L=\{j-2m+1,\ldots,j\}$.  Suppose that, from
\[
  S,\quad Z_L,\quad V^\ast_L,\quad
  \rho_{2a_L-2},\rho_{2j},
  \quad
  \rho_0,\rho_1,\ldots,\rho_{2(j-2m)},
\]
a deterministic rule identifies a Ferrers subarrangement $E$ of the Step~3
dual-RSK' half-square of
\cref{lem:bc-krattenthaler-semistandard-path-boundary}, identifies the
partition labels on its top/right boundary, and verifies
\[
  H_{2m}\subseteq E
\]
for the low-label incidence cell set of
\cref{lem:bc-step3-low-label-incidence-covers-q-prefix}.  Then the hypothesis
of \cref{prop:bc-target-window-signature-closes-refined-suffix} holds.
\end{corollary}

\begin{proof}
By \cref{lem:bc-tight-outside-recovers-left-growth-segment},
$\mathcal D_{\mathrm{out}}(P)$ determines
\[
  \rho_0,\rho_1,\ldots,\rho_{2(j-2m)}.
\]
Therefore the full target-window tuple determines the carrier $E$, its
top/right boundary labels, and the inclusion $H_{2m}\subseteq E$.  Applying
\cref{cor:bc-target-window-low-label-incidence-carrier-closes-signature} gives
the deterministic commuting-signature input.
\end{proof}

\begin{corollary}[Target-window half-step recovery closes the signature input]\label{cor:bc-target-window-halfstep-chain-closes-signature}
Work in the setting of
\cref{prop:bc-target-window-signature-closes-refined-suffix}.  Suppose that
the data
\[
  S,\quad Z_L,\quad V^\ast_L,\quad \mathcal D_{\mathrm{out}}(P),
  \quad \rho_{2a_L-2},\rho_{2j}
\]
determine the associated reduced fixed-even datum
\[
  (\overline\nu;\overline E_1,\ldots,\overline E_m)
\]
and the reduced half-step diagrams
\[
  \alpha_{\frac12},\alpha_{\frac32},\ldots,\alpha_{m-\frac12}
\]
of the true reduced growth chain.  Then the hypothesis of
\cref{prop:bc-target-window-signature-closes-refined-suffix} holds.
\end{corollary}

\begin{proof}
The reduced fixed-even datum determines
\[
  \gamma=\overline\nu\setminus\overline E_m
  \quad\text{and}\quad
  e_1,\ldots,e_{m-1}
\]
by \cref{def:bc-odd-fixed-even-standard-tableaux}.  By
\cref{prop:bc-odd-sublevels-are-reduced-halfsteps}, the recovered half-step
diagrams are exactly the odd sublevel shapes
\[
  \Gamma_1,\ldots,\Gamma_m
\]
of the associated fixed-even standard tableau.  Therefore
\cref{cor:bc-target-window-odd-sublevel-chain-closes-signature} applies and
gives the required deterministic signature input.
\end{proof}

\begin{corollary}[Target-window retained-even half-step recovery closes the signature input]\label{cor:bc-target-window-retained-even-halfstep-closes-signature}
Work in the setting of
\cref{prop:bc-target-window-signature-closes-refined-suffix}.  Suppose that
the data
\[
  S,\quad Z_L,\quad V^\ast_L,\quad \mathcal D_{\mathrm{out}}(P),
  \quad \rho_{2a_L-2},\rho_{2j}
\]
determine the tight-prefix endpoint
\[
  \nu=\lambda^{(2m)}
\]
and the retained even strips
\[
  E_1,\ldots,E_m
\]
of the true tight prefix, as well as the reduced half-step diagrams
\[
  \alpha_{\frac12},\alpha_{\frac32},\ldots,\alpha_{m-\frac12}
\]
of the true reduced growth chain.  Then the hypothesis of
\cref{prop:bc-target-window-signature-closes-refined-suffix} holds.
\end{corollary}

\begin{proof}
The true deleted odd-packet assignment belongs to
$\Omega(\nu;E_1,\ldots,E_m)$ by
\cref{lem:bc-true-tight-packets-in-omega}.  Hence
\cref{lem:bc-omega-fixed-cell-completions} gives
\[
  \mathcal C(\nu;E_1,\ldots,E_m)\ne\varnothing .
\]
Therefore
\cref{lem:bc-retained-even-strips-determine-reduced-datum} applies to the
recovered endpoint and retained strips, so the displayed target-window data
determine the associated reduced fixed-even datum.  Together with the
recovered half-step diagrams, this is exactly the hypothesis of
\cref{cor:bc-target-window-halfstep-chain-closes-signature}; applying that
corollary gives the deterministic commuting-signature input.
\end{proof}

\begin{proposition}[Commuting normal-form fibers are binary]\label{prop:bc-commuting-normal-form-fibers-binary}
For every datum of \cref{def:bc-odd-fixed-even-standard-tableaux} and every
signature word $\sigma$,
\[
  |\mathcal S_\sigma(\gamma;e_1,\ldots,e_{m-1})|\le2^m .
\]
\end{proposition}

\begin{proof}
We argue by induction on $m$.  For $m=1$ the shape has one cell and there is
at most one standard tableau, so the bound is immediate.

Assume $m\ge2$ and put $e=e_{m-1}$.  If $\sigma$ begins with a cell $c$, then
the fiber is empty unless $e$ is removable in $\gamma$ and
$c\in Y(\gamma;e)$.  In the nonempty case, deleting the cell $c$ carrying
$2m-1$ and then deleting $e$ gives an injection
\[
  \mathcal S_\sigma(\gamma;e_1,\ldots,e_{m-1})
  \hookrightarrow
  \mathcal S_{\sigma'}(\gamma_c;e_1,\ldots,e_{m-2}),
\]
where $\sigma'$ is the tail of $\sigma$ and
$\gamma_c=\gamma\setminus\{c,e\}$.  By induction the right-hand side has size
at most $2^{m-1}$, hence this case is bounded by $2^m$.

It remains to treat the case where $\sigma$ begins with $\ast$.  If $e$ is
removable in $\gamma$, the fiber is empty by definition of the signature.  If
$e$ is not removable, then the possible cells carrying $2m-1$ lie in
$Y(\gamma;e)$, and
\[
  |Y(\gamma;e)|\le2
\]
by \cref{cor:bc-last-pair-blocked-binary}.  Partitioning the fiber according
to this top cell $y$ and then deleting $y$ and $e$ gives injections into the
tail-signature fibers
\[
  \mathcal S_{\sigma'}(\gamma_y;e_1,\ldots,e_{m-2})
  \qquad(y\in Y(\gamma;e)).
\]
Each has size at most $2^{m-1}$ by induction, and there are at most two
choices of $y$.  Thus the $\ast$-fiber has size at most $2^m$.
\end{proof}

\begin{corollary}[Parser-visible commuting signatures close the refined rank suffix]\label{cor:bc-commuting-signature-closes-refined-suffix}
Work in the setting of \cref{prop:bc-refined-tight-suffix-reduces-rank}.
Suppose that for every final tight right-boundary zone with $m\ge9$ deleted
singleton packets, the branch-free canonical halo/right-boundary parser data
determine the commuting normal-form signature of the fixed-even standard
tableau associated to the true odd-packet assignment by
\[
  \Omega
  \longleftrightarrow
  \mathcal C
  \longleftrightarrow
  \mathcal R
  \longleftrightarrow
  \mathcal T
  \longleftrightarrow
  \mathcal S .
\]
Then the hypothesis of
\cref{prop:bc-refined-tight-suffix-reduces-rank} holds.  Consequently the
tight odd-packet rank hypothesis of
\cref{prop:bc-tight-odd-packet-rank-criterion} holds for every non-full tight
height right-boundary zone.
\end{corollary}

\begin{proof}
Fix a tight zone with $m\ge9$ and let $\kappa$ be the parser-visible
commuting normal-form signature supplied by the hypothesis.  Define
\[
  \Omega_A(A(P))
\]
to be the set of elements of
$\Omega(\lambda^{(2m)};E_1,\ldots,E_m)$ whose associated fixed-even standard
tableau has signature $\kappa$ under the displayed chain of bijections.  The
true odd-packet assignment belongs to this set by the choice of $\kappa$.
The bijections
\cref{lem:bc-omega-fixed-cell-completions,prop:bc-completions-equal-reduced-growth-chains,lem:bc-reduced-growth-reverse-corner-tree,lem:bc-reverse-tree-fixed-even-standard-tableaux}
identify this set with a commuting-normal-form fiber
\[
  \mathcal S_\kappa(\gamma;e_1,\ldots,e_{m-1}),
\]
whose cardinality is at most $2^m$ by
\cref{prop:bc-commuting-normal-form-fibers-binary}.  Thus the
$m\ge9$ hypothesis of \cref{prop:bc-refined-tight-suffix-reduces-rank} is
satisfied.  That proposition combines this case with the already closed
$m\le8$ cases and gives the claimed rank hypothesis.
\end{proof}

\begin{corollary}[Target-window refined ambiguity classes close the rank suffix]\label{cor:bc-target-window-refined-omega-closes-suffix}
Work in the right-boundary auxiliary setting of
\cref{prop:bc-right-boundary-auxiliary-global-recovery}, and suppose the last
zone $Z_L=\{a_L,\ldots,j\}$ is a non-full tight height zone with
$m\ge9$ deleted singleton packets.  Let $V^\ast_L$ be the replacement window
parsed from the repaired path, and let $\mathcal D_{\mathrm{out}}(P)$ be the
outside datum of
\cref{prop:bc-right-boundary-auxiliary-global-recovery}.  Suppose the
target-window data
\[
  S,\quad Z_L,\quad V^\ast_L,\quad \mathcal D_{\mathrm{out}}(P),
  \quad \rho_{2a_L-2},\rho_{2j}
\]
determine the tight-prefix endpoint $\lambda^{(2m)}$, the retained even-label
strips $E_1,\ldots,E_m$, the labelled continuation boxes, and a canonically
ordered finite set
\[
  \Omega_{\mathrm{tw}}
  \subseteq
  \Omega(\lambda^{(2m)};E_1,\ldots,E_m)
\]
such that the true residual odd-packet assignment belongs to
$\Omega_{\mathrm{tw}}$ and
\[
  |\Omega_{\mathrm{tw}}|\le2^m .
\]
Then the tight odd-packet rank hypothesis of
\cref{prop:bc-tight-odd-packet-rank-criterion} holds for the final
right-boundary zone.
\end{corollary}

\begin{proof}
The displayed target-window tuple is parser-visible without adding branch
words: $V^\ast_L$ is a subwindow of the already counted repaired path, the
endpoint pair is returned by the parser, and
$\mathcal D_{\mathrm{out}}(P)$ is recovered in
\cref{prop:bc-right-boundary-auxiliary-global-recovery}.  Use as $D(P)$ the
determined tight-prefix endpoint, retained even-label strips, labelled
continuation boxes, and the determined set $\Omega_{\mathrm{tw}}$.  The
membership and size hypotheses are exactly those required by
\cref{cor:bc-refined-tight-packet-criterion}.  Applying that corollary gives
the final-zone tight odd-packet rank hypothesis.
\end{proof}

\begin{corollary}[Target-window signature fibers are refined ambiguity classes]\label{cor:bc-target-window-signature-fiber-omega}
Work in the right-boundary auxiliary setting of
\cref{prop:bc-right-boundary-auxiliary-global-recovery}, and suppose the last
zone $Z_L=\{a_L,\ldots,j\}$ is a non-full tight height zone with
$m\ge9$ deleted singleton packets.  Let $V^\ast_L$ be the replacement window
parsed from the repaired path, and let $\mathcal D_{\mathrm{out}}(P)$ be the
outside datum of
\cref{prop:bc-right-boundary-auxiliary-global-recovery}.  Suppose the
target-window data
\[
  S,\quad Z_L,\quad V^\ast_L,\quad \mathcal D_{\mathrm{out}}(P),
  \quad \rho_{2a_L-2},\rho_{2j}
\]
determine $\lambda^{(2m)}$, the retained even-label strips
$E_1,\ldots,E_m$, the labelled continuation boxes, and the commuting
normal-form signature $\kappa$ of the fixed-even standard tableau associated
to the true odd-packet assignment.  Define
\[
  \Omega_{\mathrm{tw}}(\kappa)
\]
to be the set of elements of
$\Omega(\lambda^{(2m)};E_1,\ldots,E_m)$ whose image under
\[
  \Omega
  \longleftrightarrow
  \mathcal C
  \longleftrightarrow
  \mathcal R
  \longleftrightarrow
  \mathcal T
  \longleftrightarrow
  \mathcal S
\]
has commuting normal-form signature $\kappa$.  Then this
$\Omega_{\mathrm{tw}}(\kappa)$ is determined by the target-window data,
contains the true residual odd-packet assignment, and satisfies
\[
  |\Omega_{\mathrm{tw}}(\kappa)|\le2^m .
\]
Consequently the tight odd-packet rank hypothesis of
\cref{prop:bc-tight-odd-packet-rank-criterion} holds for the final
right-boundary zone.
\end{corollary}

\begin{proof}
The target-window data determine the ambient set
$\Omega(\lambda^{(2m)};E_1,\ldots,E_m)$ and the signature $\kappa$, so they
determine the displayed subset.  The true odd-packet assignment belongs to it
because $\kappa$ is, by hypothesis, the signature of the associated true
fixed-even standard tableau.

The bijections
\cref{lem:bc-omega-fixed-cell-completions,prop:bc-completions-equal-reduced-growth-chains,lem:bc-reduced-growth-reverse-corner-tree,lem:bc-reverse-tree-fixed-even-standard-tableaux}
identify $\Omega_{\mathrm{tw}}(\kappa)$ with a commuting-normal-form fiber
\[
  \mathcal S_\kappa(\gamma;e_1,\ldots,e_{m-1}).
\]
By \cref{prop:bc-commuting-normal-form-fibers-binary}, this fiber has
cardinality at most $2^m$.  Thus
\cref{cor:bc-target-window-refined-omega-closes-suffix} applies.
\end{proof}

\begin{corollary}[Bounded target-window signature lists give full-budget odd-packet classes]\label{cor:bc-bounded-signature-list-full-budget-omega}
Work in the setting of
\cref{cor:bc-target-window-signature-fiber-omega}.  Suppose the target-window
data determine \(\lambda^{(2m)}\), the retained even-label strips
\(E_1,\ldots,E_m\), the labelled continuation boxes, and a canonically ordered
finite set
\[
  \Sigma_{\mathrm{tw}}
\]
of commuting normal-form signature words such that the true fixed-even
standard tableau has signature in \(\Sigma_{\mathrm{tw}}\) and
\[
  |\Sigma_{\mathrm{tw}}|\le 2^m .
\]
Define \(\Omega_{\mathrm{tw}}(\Sigma_{\mathrm{tw}})\) to be the set of elements
of \(\Omega(\lambda^{(2m)};E_1,\ldots,E_m)\) whose image under
\[
  \Omega
  \longleftrightarrow
  \mathcal C
  \longleftrightarrow
  \mathcal R
  \longleftrightarrow
  \mathcal T
  \longleftrightarrow
  \mathcal S
\]
has commuting normal-form signature in \(\Sigma_{\mathrm{tw}}\), ordered by the
ambient lexicographic order on
\(\Omega(\lambda^{(2m)};E_1,\ldots,E_m)\).  Then
\(\Omega_{\mathrm{tw}}(\Sigma_{\mathrm{tw}})\) is determined by the
target-window data, contains the true residual odd-packet assignment, and
satisfies
\[
  |\Omega_{\mathrm{tw}}(\Sigma_{\mathrm{tw}})|\le 2^{2m}.
\]
\end{corollary}

\begin{proof}
The target-window data determine the ambient odd-packet class, the signature
list, and the displayed filtering rule, hence determine
\(\Omega_{\mathrm{tw}}(\Sigma_{\mathrm{tw}})\) with its inherited order.  The
true residual odd-packet assignment belongs to the set because the true
fixed-even standard tableau has signature in \(\Sigma_{\mathrm{tw}}\).

For each \(\sigma\in\Sigma_{\mathrm{tw}}\), the bijections
\cref{lem:bc-omega-fixed-cell-completions,prop:bc-completions-equal-reduced-growth-chains,lem:bc-reduced-growth-reverse-corner-tree,lem:bc-reverse-tree-fixed-even-standard-tableaux}
identify the \(\sigma\)-part of
\(\Omega_{\mathrm{tw}}(\Sigma_{\mathrm{tw}})\) with a subset of
\[
  \mathcal S_\sigma(\gamma;e_1,\ldots,e_{m-1}).
\]
By \cref{prop:bc-commuting-normal-form-fibers-binary}, each such fiber has
cardinality at most \(2^m\).  Summing over the signature list gives
\[
  |\Omega_{\mathrm{tw}}(\Sigma_{\mathrm{tw}})|
  \le |\Sigma_{\mathrm{tw}}|\,2^m
  \le 2^m\cdot 2^m
  =2^{2m}.
\]
\end{proof}

\begin{corollary}[All-depth overflow realizes more than \(2^m\) signatures]\label{cor:bc-all-depth-overflow-many-signatures}
Work in the setting of \cref{cor:bc-full-depth-envelope-exact-omega}, and
assume
\[
  (\gamma,e_{m-1},e_{m-2},\ldots,e_1)
  \in\mathcal P^{\mathrm{env},m-1}_m .
\]
Let \(\Sigma(\Omega_{\mathrm{tw}})\) be the set of commuting normal-form
signatures realized by the fixed-even standard tableaux corresponding to
elements of \(\Omega_{\mathrm{tw}}\).  Then
\[
  |\Omega_{\mathrm{tw}}|>2^{2m}
  \qquad\text{and}\qquad
  |\Sigma(\Omega_{\mathrm{tw}})|>2^m .
\]
\end{corollary}

\begin{proof}
The first inequality is the contrapositive of the final equivalence in
\cref{cor:bc-full-depth-envelope-exact-omega}.  If
\(|\Sigma(\Omega_{\mathrm{tw}})|\le2^m\), then the elements of
\(\Omega_{\mathrm{tw}}\) split into at most \(2^m\) commuting-signature fibers.
By \cref{prop:bc-commuting-normal-form-fibers-binary}, each fiber has
cardinality at most \(2^m\).  Hence
\[
  |\Omega_{\mathrm{tw}}|\le2^m\cdot2^m=2^{2m},
\]
contradicting the first inequality.
\end{proof}

\begin{corollary}[Target-window retained-box readout supplies the signature fiber]\label{cor:bc-target-window-retained-boxes-signature-fiber}
Work in the setting of
\cref{cor:bc-target-window-signature-fiber-omega}.  Suppose the target-window
data determine the tight-prefix endpoint $\lambda^{(2m)}$, the two boxes
carrying each retained even label
\[
  2,4,\ldots,2m,
\]
the labelled continuation boxes, and the commuting normal-form signature
$\kappa$ of the true fixed-even standard tableau.  Then the tight odd-packet
rank hypothesis of
\cref{prop:bc-tight-odd-packet-rank-criterion} holds for the final
right-boundary zone.
\end{corollary}

\begin{proof}
For each $1\le r\le m$, the two boxes carrying the retained even label $2r$
are exactly the retained strip $E_r$.  Thus the target-window data determine
\[
  \lambda^{(2m)},\quad E_1,\ldots,E_m,
\]
the labelled continuation boxes, and $\kappa$.  Applying
\cref{cor:bc-target-window-signature-fiber-omega} gives the claimed rank
hypothesis.
\end{proof}

\begin{corollary}[Target-window even-chain readout supplies the signature fiber]\label{cor:bc-target-window-even-chain-signature-fiber}
Work in the setting of
\cref{cor:bc-target-window-retained-boxes-signature-fiber}, and let
$U=\operatorname{std}(T)$ be the doubled standardization of the underlying
semistandard tableau.  Suppose the target-window data determine the even
standard-time endpoint chain
\[
  \sigma^{(0)},\sigma^{(2)},\sigma^{(4)},\ldots,\sigma^{(4m)}
\]
and the labelled continuation boxes.  Then the tight odd-packet rank
hypothesis of \cref{prop:bc-tight-odd-packet-rank-criterion} holds for the
final right-boundary zone.
\end{corollary}

\begin{proof}
By \cref{lem:bc-standard-time-pair-deletion}, the recovered even
standard-time endpoints determine the true tight-prefix Pieri chain
\[
  \lambda^{(0)}\subset\lambda^{(1)}\subset\cdots\subset\lambda^{(2m)} .
\]
In particular they determine the endpoint $\lambda^{(2m)}$ and, for
$1\le r\le m$, the retained even-label boxes
\[
  E_r=\lambda^{(2r)}\setminus\lambda^{(2r-1)} .
\]
The proof of \cref{cor:bc-target-window-prefix-path-closes-signature} then
gives a deterministic commuting normal-form signature $\kappa$ from the same
recovered prefix chain: the chain determines the retained strips and the
reduced half-step diagrams, hence the fixed-even datum and odd sublevel chain.
Thus the hypotheses of
\cref{cor:bc-target-window-retained-boxes-signature-fiber} are satisfied.
\end{proof}

\begin{corollary}[Target-window low-label boxes supply the signature fiber]\label{cor:bc-target-window-low-label-boxes-signature-fiber}
Work in the setting of
\cref{cor:bc-target-window-even-chain-signature-fiber}.  Suppose the
target-window data determine, for each semistandard label
$1\le a\le2m$, the two boxes carrying label $a$, and also determine the
labelled continuation boxes for labels $2m+1,\ldots,j$.  Then the tight
odd-packet rank hypothesis of
\cref{prop:bc-tight-odd-packet-rank-criterion} holds for the final
right-boundary zone.
\end{corollary}

\begin{proof}
By \cref{lem:bc-low-label-boxes-standard-readout} with $M=2m$, the recovered
boxes carrying labels $1,\ldots,2m$ determine
\[
  \sigma^{(0)},\sigma^{(2)},\sigma^{(4)},\ldots,\sigma^{(4m)}.
\]
Together with the recovered labelled continuation boxes, this is exactly the
hypothesis of
\cref{cor:bc-target-window-even-chain-signature-fiber}.  Applying that
corollary gives the claimed rank hypothesis.
\end{proof}

\begin{corollary}[Target-window low-label boxes close the rank suffix as a singleton]\label{cor:bc-target-window-low-label-boxes-singleton-rank}
Work in the right-boundary auxiliary setting of
\cref{prop:bc-right-boundary-auxiliary-global-recovery}, and suppose the last
zone $Z_L=\{a_L,\ldots,j\}$ is a non-full tight height zone with
$m\ge9$ deleted singleton packets.  Let $V^\ast_L$ be the replacement window
parsed from the repaired path, and let $\mathcal D_{\mathrm{out}}(P)$ be the
outside datum of
\cref{prop:bc-right-boundary-auxiliary-global-recovery}.  Suppose the
target-window data
\[
  S,\quad Z_L,\quad V^\ast_L,\quad \mathcal D_{\mathrm{out}}(P),
  \quad \rho_{2a_L-2},\rho_{2j}
\]
determine, for each semistandard label
$1\le a\le2m$, the two boxes carrying label $a$, and also determine the
labelled continuation boxes for labels $2m+1,\ldots,j$.  Then the tight
odd-packet rank hypothesis of
\cref{prop:bc-tight-odd-packet-rank-criterion} holds for the final
right-boundary zone.
\end{corollary}

\begin{proof}
\Cref{lem:bc-low-label-boxes-determine-tight-packet-datum} shows that the
recovered low-label boxes determine
\[
  \lambda^{(2m)},\qquad E_1,\ldots,E_m,
  \qquad \omega(P)=(x_1,\ldots,x_m)
  \in \Omega(\lambda^{(2m)};E_1,\ldots,E_m),
\]
where $\omega(P)$ is the true deleted odd-packet assignment.  Use the
target-window data as the parser-visible datum $D(P)$ and put
\[
  \Omega(D(P))=\{\omega(P)\}
\]
with its unique order.  This set is determined by $D(P)$, contains the true
assignment, and has size $1\le2^m$.  The same data determine the labelled
continuation boxes by hypothesis.  Therefore all hypotheses of
\cref{prop:bc-tight-odd-packet-rank-criterion} are satisfied.
\end{proof}

\begin{corollary}[Target-window continuation sweep plus low-label boxes closes the rank suffix]\label{cor:bc-target-window-continuation-sweep-singleton-rank}
Work in the setting of
\cref{cor:bc-target-window-low-label-boxes-singleton-rank}.  Suppose the
target-window data determine, for each semistandard label $1\le a\le2m$, the
two boxes carrying label $a$.  Suppose moreover that the same data determine a
Ferrers subarrangement $C$ of the Step~2 first-variant Knuth growth diagram,
the cell entries of $C$, and the left/bottom boundary labels of $C$, such that
the forward boundary sweep on $C$ contains
\[
  \lambda^{(2m)},\lambda^{(2m+1)},\ldots,\lambda^{(j)} .
\]
Then the tight odd-packet rank hypothesis of
\cref{prop:bc-tight-odd-packet-rank-criterion} holds for the final
right-boundary zone.
\end{corollary}

\begin{proof}
\Cref{lem:bc-step2-knuth-continuation-boxes} with $M=2m$ recovers the two
boxes carrying every continuation label $2m+1,\ldots,j$ from the displayed
Step~2 continuation sweep.  Together with the recovered low-label boxes, this
is exactly the hypothesis of
\cref{cor:bc-target-window-low-label-boxes-singleton-rank}.  Applying that
corollary gives the rank hypothesis.
\end{proof}

\begin{corollary}[Strict-left plus low-low entries close the singleton rank route]\label{cor:bc-strict-left-low-low-singleton-rank}
Work in the non-full tight-height setting of
\cref{lem:bc-tight-outside-recovers-left-growth-segment} and in the
right-boundary auxiliary setting of
\cref{prop:bc-right-boundary-auxiliary-global-recovery}, with
$Z_L=\{j-2m+1,\ldots,j\}$.  Let $L_{2m}$ be the low-low cell set of
\cref{lem:bc-low-label-incidence-split}.  Suppose that the target-window data
\[
  S,\quad Z_L,\quad V^\ast_L,\quad \mathcal D_{\mathrm{out}}(P),
  \quad \rho_{2a_L-2},\rho_{2j}
\]
determine the Step~3 cell entries on every cell of $L_{2m}$.  Then the tight
odd-packet rank hypothesis of
\cref{prop:bc-tight-odd-packet-rank-criterion} holds for the final
right-boundary zone.
\end{corollary}

\begin{proof}
Put $M=2m$.  By \cref{lem:bc-strict-left-recovers-continuation-columns}, the
outside datum $\mathcal D_{\mathrm{out}}(P)$ determines every Step~3 entry in
the continuation-column subarrangement $J_M$, whose cells are exactly the
off-diagonal unordered label pairs incident to some continuation label
$M+1,\ldots,j$.  The hypothesis determines every entry on $L_M$, the cells
whose two incident labels are both at most $M$.

Every off-diagonal unordered label pair in the Step~3 half-square is in
exactly one of these two classes: either both labels are at most $M$, or at
least one label is a continuation label.  Hence the target-window data
determine every Step~3 cell entry.  By
\cref{lem:bc-step3-cells-are-step2-half-matrix}, these entries determine every
off-diagonal entry of the Step~2 symmetric matrix, and the diagonal entries
are known to be zero.  Thus the entries of the whole Step~2 first-variant
Knuth diagram are determined.

Applying \cref{lem:bc-step2-q-prefix-strip-determines-prefix} with $M=j$
recovers the complete Pieri boundary path
\[
  \lambda^{(0)},\lambda^{(1)},\ldots,\lambda^{(j)}.
\]
Then \cref{lem:bc-step2-knuth-prefix-low-label-boxes} with $M=2m$ recovers the
two boxes carrying every label $1,\ldots,2m$, and
\cref{lem:bc-step2-knuth-continuation-boxes} with $M=2m$ recovers the two
boxes carrying every continuation label $2m+1,\ldots,j$.  The hypotheses of
\cref{cor:bc-target-window-low-label-boxes-singleton-rank} are therefore
satisfied, and that corollary gives the rank hypothesis.
\end{proof}

\begin{corollary}[Bounded low-low candidates fit the tight branch budget]\label{cor:bc-bounded-low-low-candidates-fit-branch-budget}
Work in the non-full tight-height setting of
\cref{lem:bc-tight-outside-recovers-left-growth-segment} and in the
right-boundary auxiliary setting of
\cref{prop:bc-right-boundary-auxiliary-global-recovery}, with
$Z_L=\{j-2m+1,\ldots,j\}$.  Let $L_{2m}$ be the low-low cell set of
\cref{lem:bc-low-label-incidence-split}.  Suppose that the target-window data
\[
  S,\quad Z_L,\quad V^\ast_L,\quad \mathcal D_{\mathrm{out}}(P),
  \quad \rho_{2a_L-2},\rho_{2j}
\]
determine a canonically ordered finite set
\[
  \mathcal L_{\mathrm{tw}}
\]
of possible Step~3 entry assignments on $L_{2m}$, that the true assignment
belongs to this set, and that
\[
  |\mathcal L_{\mathrm{tw}}|\le2^m.
\]
Then the first $m$ bits of the tight-zone branch word
\[
  \beta_\gamma\in\{0,1\}^{2m}
\]
can encode the rank of the true low-low assignment in
$\mathcal L_{\mathrm{tw}}$.  From the target-window data and
$\beta_\gamma$ one recovers the full tight-prefix singleton packet data, and
the last $m$ bits remain available for the height selected-defect data.
\end{corollary}

\begin{proof}
Use the first $m$ bits of $\beta_\gamma$ to encode the rank of the true
low-low assignment in the canonical order on $\mathcal L_{\mathrm{tw}}$.
Since $|\mathcal L_{\mathrm{tw}}|\le2^m$, this selects the true Step~3 entries
on $L_{2m}$.

Put $M=2m$.  By \cref{lem:bc-strict-left-recovers-continuation-columns}, the
outside datum determines every Step~3 entry in the continuation-column
subarrangement $J_M$.  Combining those entries with the selected entries on
$L_M$ determines the whole Step~3 half-square, by the same dichotomy used in
\cref{cor:bc-strict-left-low-low-singleton-rank}: each off-diagonal unordered
label pair either has both labels at most $M$, or is incident to a
continuation label.

\Cref{lem:bc-step3-cells-are-step2-half-matrix} then determines the whole
Step~2 symmetric matrix, with zero diagonal.  Applying
\cref{lem:bc-step2-q-prefix-strip-determines-prefix} with $M=j$ recovers
\[
  \lambda^{(0)},\lambda^{(1)},\ldots,\lambda^{(j)}.
\]
Thus \cref{lem:bc-step2-knuth-prefix-low-label-boxes} with $M=2m$ recovers
the boxes carrying labels $1,\ldots,2m$, and
\cref{lem:bc-step2-knuth-continuation-boxes} with $M=2m$ recovers the boxes
carrying continuation labels $2m+1,\ldots,j$.  Finally,
\cref{lem:bc-low-label-boxes-determine-tight-packet-datum} recovers the true
deleted odd-packet assignment and the tight-prefix endpoint and retained
even-label strips.

The last $m$ bits of $\beta_\gamma$ have not been used in the low-low rank
encoding.  By \cref{lem:bc-height-branch-first-half-slack}, these last bits
are still exactly the bits used by the height local recovery for the selected
defect subset.  Therefore the target-window data and $\beta_\gamma$ recover
the full tight-prefix singleton packet data while preserving the height
selected-defect recovery.
\end{proof}

\begin{corollary}[Bounded low-low candidates recover the tight right-boundary subpath]\label{cor:bc-bounded-low-low-candidates-recover-right-boundary}
Work in the setting of
\cref{cor:bc-bounded-low-low-candidates-fit-branch-budget}, with
\[
  Z_L=\{a_L,\ldots,j\}=\{j-2m+1,\ldots,j\}.
\]
Suppose the target-window data determine the canonically ordered set
$\mathcal L_{\mathrm{tw}}$ from
\cref{cor:bc-bounded-low-low-candidates-fit-branch-budget}.  Then from
\[
  S,\quad Z_L,\quad V^\ast_L,\quad \mathcal D_{\mathrm{out}}(P),
  \quad \rho_{2a_L-2},\rho_{2j},
  \quad \beta_L\in\{0,1\}^{2m}
\]
one recovers the original right-boundary zone subpath
\[
  \rho_{2a_L-2},\rho_{2a_L-1},\ldots,\rho_{2j}.
\]
\end{corollary}

\begin{proof}
Read the first $m$ bits of $\beta_L$ as the rank of the true low-low
assignment in the canonical order on $\mathcal L_{\mathrm{tw}}$.  Since
$|\mathcal L_{\mathrm{tw}}|\le2^m$, this selects the true Step~3 entries on
$L_{2m}$.

Put $M=2m$.  By \cref{lem:bc-strict-left-recovers-continuation-columns},
$\mathcal D_{\mathrm{out}}(P)$ determines every Step~3 entry in the
continuation-column subarrangement $J_M$.  The cell sets $L_M$ and $J_M$
cover the full Step~3 half-square: an unordered off-diagonal label pair either
has both labels at most $M$, or is incident to a continuation label
$M+1,\ldots,j$.  Hence the displayed data and the selected low-low assignment
determine every cell entry of the full Step~3 half-square.

By \cref{lem:bc-step3-full-half-square-determines-rho}, those full Step~3
entries determine the whole boundary sequence
\[
  \rho_0,\rho_1,\ldots,\rho_{2j}.
\]
Since $a_L=j-2m+1$, the desired right-boundary subpath is the terminal segment
\[
  \rho_{2(j-2m)},\rho_{2(j-2m)+1},\ldots,\rho_{2j}
  =
  \rho_{2a_L-2},\rho_{2a_L-1},\ldots,\rho_{2j}.
\]
\end{proof}

\begin{definition}[Degree-constrained target-window low-low candidates]\label{def:bc-degree-low-low-candidates}
Work in the setting of
\cref{cor:bc-bounded-low-low-candidates-fit-branch-budget}, and put
$M=2m$.  By \cref{lem:bc-strict-left-recovers-continuation-columns}, the
outside datum determines the Step~3 entries in the continuation-column
subarrangement $J_M$.  For $1\le a\le M$ define
\[
  c_a=\sum_{t=M+1}^j u_{\{a,t\}},
\]
where $u_{\{a,t\}}$ is the recovered Step~3 entry on the unordered label pair
$\{a,t\}$.  Put
\[
  d_a=2-c_a .
\]
The degree-constrained target-window low-low candidate set
\[
  \mathcal L_{\mathrm{deg}}(S,Z_L,V^\ast_L,\mathcal D_{\mathrm{out}}(P),
     \rho_{2a_L-2},\rho_{2j})
\]
is empty if some $d_a<0$.  Otherwise it is the set of all nonnegative integer
assignments
\[
  \ell=(\ell_{ab})_{1\le a<b\le M}
\]
to the low-low cell set $L_M$ satisfying
\[
  \sum_{b<a}\ell_{ba}+\sum_{a<b\le M}\ell_{ab}=d_a
  \qquad(1\le a\le M).
\]
It is ordered lexicographically by the list
\[
  (\ell_{12},\ell_{13},\ldots,\ell_{M-1,M}).
\]
\end{definition}

\begin{proposition}[Strict-left degrees prune the low-low supplier]\label{prop:bc-strict-left-degree-low-low-pruning}
In the setting of
\cref{def:bc-degree-low-low-candidates}, the displayed target-window data
determine the set $\mathcal L_{\mathrm{deg}}$ as a fixed finite combinatorial
object.  The true Step~3 assignment on $L_{2m}$ belongs to
$\mathcal L_{\mathrm{deg}}$, and $\mathcal L_{\mathrm{deg}}$ is finite.
Consequently, to apply
\cref{cor:bc-bounded-low-low-candidates-fit-branch-budget} it is enough to
find a canonical target-window subset
\[
  \mathcal L_{\mathrm{tw}}\subseteq\mathcal L_{\mathrm{deg}}
\]
which contains the true low-low assignment and has cardinality at most $2^m$.
\end{proposition}

\begin{proof}
The target-window data include $\mathcal D_{\mathrm{out}}(P)$, and
\cref{lem:bc-strict-left-recovers-continuation-columns} recovers from it every
Step~3 entry in $J_M$.  Hence the integers $c_a$ and $d_a$, and therefore the
lexicographically ordered set $\mathcal L_{\mathrm{deg}}$, are determined by
the displayed data.

Let $N=(N_{ab})$ be the true Step~2 symmetric matrix.  By
\cref{lem:bc-step3-cells-are-step2-half-matrix}, each recovered Step~3 entry
on the unordered pair $\{a,t\}$ is the corresponding off-diagonal entry
$N_{at}=N_{ta}$.  By
\cref{lem:bc-step2-matrix-doubled-content-degrees}, the $a$-th row of $N$ has
sum $2$ and $N_{aa}=0$.  Therefore, for each $1\le a\le M$,
\[
  \sum_{\substack{1\le b\le M\\ b\ne a}}N_{ab}
  =
  2-\sum_{t=M+1}^jN_{at}
  =
  d_a .
\]
The low-low Step~3 entries of the true tableau are exactly these
$N_{ab}$ with $1\le a<b\le M$, again by
\cref{lem:bc-step3-cells-are-step2-half-matrix}.  They therefore satisfy the
degree equations in \cref{def:bc-degree-low-low-candidates}, proving
membership.

Finally each $d_a$ is at most $2$ for a true target-window datum, and in any
case a nonempty candidate set has
$0\le\ell_{ab}\le\min(d_a,d_b)$.  Thus only finitely many assignments can
satisfy the displayed equations.  If a canonical subset
$\mathcal L_{\mathrm{tw}}$ of this finite ambient set contains the true
assignment and has size at most $2^m$, then it satisfies the hypotheses of
\cref{cor:bc-bounded-low-low-candidates-fit-branch-budget}, which gives the
last assertion.
\end{proof}

\begin{proposition}[Sparse low-low residual degree fits the full final-zone budget]\label{prop:bc-sparse-active-low-low-degree-full-budget}
Work in the setting of \cref{def:bc-degree-low-low-candidates}.  Let
\[
  \mathcal A_{\mathrm{tw}}=\{1\le a\le2m:d_a>0\},
  \qquad r=|\mathcal A_{\mathrm{tw}}|,\qquad
  D=\sum_{a=1}^{2m}d_a .
\]
Use the convention $(-1)!!=1$.  If
\[
  (D-1)!!\le2^{2m},
\]
then the degree-constrained set $\mathcal L_{\mathrm{deg}}$ itself is a
canonical full-budget low-low candidate set: it is determined by the
target-window data, contains the true low-low Step~3 assignment on $L_{2m}$,
and satisfies
\[
  |\mathcal L_{\mathrm{deg}}|\le2^{2m}.
\]
Consequently the final zone satisfies the right-boundary inverse hypothesis of
\cref{prop:bc-right-boundary-auxiliary-global-recovery} by
\cref{prop:bc-full-budget-low-low-candidates-supply-large-right-boundary-auxiliary}.
In particular the same conclusion holds under the simpler sufficient condition
\[
  (2r-1)!!\le2^{2m}.
\]
\end{proposition}

\begin{proof}
The target-window data determine the integers $d_a$, hence determine
$\mathcal A_{\mathrm{tw}}$, $r$, and $D$.  By
\cref{prop:bc-strict-left-degree-low-low-pruning}, they also determine the
finite set $\mathcal L_{\mathrm{deg}}$, and the true low-low assignment belongs
to it.

It remains to bound the size of $\mathcal L_{\mathrm{deg}}$.  If
$\mathcal L_{\mathrm{deg}}$ is empty there is nothing to prove, so suppose it
is nonempty.  Then $D$ is even, since every assignment in
$\mathcal L_{\mathrm{deg}}$ has total degree twice its number of low-low edges.

For each active label $a$, create $d_a$ labelled ports $(a,1),\ldots,(a,d_a)$.
Given an assignment $\ell\in\mathcal L_{\mathrm{deg}}$, order the edge-ends
incident to each vertex first by the other endpoint and then by multiplicity,
and attach the ordered ports of that vertex to those ordered edge-ends.  Pair
the two ports belonging to each low-low edge.  This gives a perfect matching of
the $D$ labelled ports.  The construction is injective: from the matched port
pairs and their vertex labels one recovers exactly the edge multiplicities
$\ell_{ab}$ for $1\le a<b\le2m$.

The number of perfect matchings of $D$ labelled ports is $(D-1)!!$.  Hence
\[
  |\mathcal L_{\mathrm{deg}}|\le (D-1)!!\le2^{2m},
\]
with the displayed convention in the case $D=0$.  Taking
$\mathcal L_{\mathrm{tw}}=\mathcal L_{\mathrm{deg}}$ in
\cref{prop:bc-full-budget-low-low-candidates-supply-large-right-boundary-auxiliary}
gives the stated right-boundary inverse.

For the final sentence, observe that $D\le2r$ because $d_a=2-c_a$ in
\cref{def:bc-degree-low-low-candidates}, with $c_a$ a sum of nonnegative
recovered Step~3 entries.  Thus $(D-1)!!\le(2r-1)!!$, and the simpler
condition implies the first displayed hypothesis.
\end{proof}

\begin{proposition}[Degree-two residual labels improve the full final-zone budget]\label{prop:bc-degree-two-symmetry-low-low-full-budget}
Work in the setting of \cref{def:bc-degree-low-low-candidates}.  Put
\[
  D=\sum_{a=1}^{2m}d_a,\qquad
  r_2=\#\{1\le a\le2m:d_a=2\},
\]
and use the convention $(-1)!!=1$.  If
\[
  (D-1)!!\le 2^{2m+\lceil r_2/2\rceil},
\]
then the degree-constrained set $\mathcal L_{\mathrm{deg}}$ itself is a
canonical full-budget low-low candidate set: it is determined by the
target-window data, contains the true low-low Step~3 assignment on $L_{2m}$,
and satisfies
\[
  |\mathcal L_{\mathrm{deg}}|\le2^{2m}.
\]
Consequently the final zone satisfies the right-boundary inverse hypothesis of
\cref{prop:bc-right-boundary-auxiliary-global-recovery} by
\cref{prop:bc-full-budget-low-low-candidates-supply-large-right-boundary-auxiliary}.
\end{proposition}

\begin{proof}
The target-window data determine the integers $d_a$, hence determine $D$ and
$r_2$.  By \cref{prop:bc-strict-left-degree-low-low-pruning}, they also
determine the finite set $\mathcal L_{\mathrm{deg}}$, and the true low-low
assignment belongs to it.

It remains to bound the size of $\mathcal L_{\mathrm{deg}}$.  If
$\mathcal L_{\mathrm{deg}}$ is empty the cardinality bound is trivial, so
suppose it is nonempty.  Then \(D\) is even.  For each label \(a\), create
\(d_a\) labelled ports.  Nonemptiness gives \(d_a\ge0\), while
\(d_a=2-c_a\) and \(c_a\) is a sum of nonnegative recovered Step~3 entries
give \(d_a\le2\).  Hence every nonempty degree vector has
\(d_a\in\{0,1,2\}\).

For an assignment \(\ell\in\mathcal L_{\mathrm{deg}}\), let
\(\mathcal P(\ell)\) be the set of perfect matchings of the \(D\) labelled
ports which realize the multiplicities \(\ell_{ab}\): a matched port pair with
vertex labels \(a\) and \(b\) contributes one unit to \(\ell_{ab}\).  These
sets \(\mathcal P(\ell)\) are pairwise disjoint subsets of the set of all
perfect matchings of the \(D\) ports.

For a fixed \(\ell\), the number of such labelled-port matchings is
\[
  |\mathcal P(\ell)|
  =
  \frac{\prod_{a=1}^{2m}d_a!}{\prod_{1\le a<b\le2m}\ell_{ab}!}.
\]
The numerator is \(2^{r_2}\).  The denominator is a power of two, with one
factor \(2\) for each doubled edge \(\ell_{ab}=2\).  A doubled edge consumes
two labels of residual degree \(2\), so there are at most
\(\lfloor r_2/2\rfloor\) doubled edges.  Therefore
\[
  |\mathcal P(\ell)|\ge
  2^{r_2-\lfloor r_2/2\rfloor}
  =
  2^{\lceil r_2/2\rceil}.
\]

The total number of perfect matchings of \(D\) labelled ports is
\((D-1)!!\).  Since the subsets \(\mathcal P(\ell)\) are disjoint,
\[
  |\mathcal L_{\mathrm{deg}}|\,2^{\lceil r_2/2\rceil}
  \le (D-1)!!.
\]
The displayed hypothesis gives
\[
  |\mathcal L_{\mathrm{deg}}|\le2^{2m}.
\]
Taking \(\mathcal L_{\mathrm{tw}}=\mathcal L_{\mathrm{deg}}\) in
\cref{prop:bc-full-budget-low-low-candidates-supply-large-right-boundary-auxiliary}
gives the stated right-boundary inverse.
\end{proof}

\begin{definition}[Exact low-low degree-profile count]\label{def:bc-exact-low-low-degree-profile-count}
For \(c\ge0\), let \(\mathsf C(c)\) be the number of loopless labelled
multigraphs on \(c\) vertices in which every vertex has degree \(2\).  Thus
\(\mathsf C(0)=1\), \(\mathsf C(1)=0\), and, for \(c\ge2\),
\[
  \mathsf C(c)=
  \sum_{s=2}^c \binom{c-1}{s-1}\operatorname{cyc}(s)\mathsf C(c-s),
\]
where \(\operatorname{cyc}(2)=1\) and
\(\operatorname{cyc}(s)=(s-1)!/2\) for \(s\ge3\).

For \(b,c\ge0\), let \(\mathsf N(b,c)\) be the number of loopless labelled
multigraphs on a fixed vertex set \(B\sqcup C\), with \(|B|=b\) and
\(|C|=c\), in which every vertex of \(B\) has degree \(1\) and every vertex of
\(C\) has degree \(2\).  Equivalently, \(\mathsf N(b,c)=0\) for odd \(b\),
\(\mathsf N(0,c)=\mathsf C(c)\), and, for even \(b\ge2\),
\[
  \mathsf N(b,c)
  =
  (b-1)\sum_{k=0}^c \binom{c}{k}k!\,\mathsf N(b-2,c-k).
\]
\end{definition}

\begin{lemma}[The profile recurrence counts exactly]\label{lem:bc-exact-low-low-degree-profile-recurrence}
The recurrences in
\cref{def:bc-exact-low-low-degree-profile-count} compute the stated labelled
multigraph counts.
\end{lemma}

\begin{proof}
We first count \(\mathsf C(c)\).  In a loopless labelled multigraph whose
vertices all have degree \(2\), every connected component is either a doubled
edge on two vertices or a simple cycle on at least three vertices.  Look at the
component containing a distinguished vertex.  If that component has size
\(s=2\), the other vertex is chosen in \(\binom{c-1}{1}\) ways and the
component is the doubled edge.  If \(s\ge3\), the other \(s-1\) vertices are
chosen in \(\binom{c-1}{s-1}\) ways and the number of undirected cycles through
the distinguished vertex on these \(s\) vertices is \((s-1)!/2\).  The
remaining \(c-s\) vertices carry an independent object counted by
\(\mathsf C(c-s)\).  Summing over \(s\) gives the displayed recurrence, with
the initial values \(\mathsf C(0)=1\) and \(\mathsf C(1)=0\).

Now suppose \(b>0\).  A graph with all vertices of \(B\) of degree \(1\) has
no component consisting only of degree-\(2\) vertices attached to a
degree-\(1\) vertex; every component containing a vertex of \(B\) is a path
whose two endpoints are vertices of \(B\) and whose internal vertices lie in
\(C\).  In particular \(b\) must be even, so \(\mathsf N(b,c)=0\) for odd
\(b\).  For even \(b\ge2\), distinguish one vertex of \(B\).  Choose the other
endpoint of its path in \(b-1\) ways.  If the path has \(k\) internal
degree-\(2\) vertices, choose those vertices in \(\binom ck\) ways and order
them along the path in \(k!\) ways.  Removing this path leaves an independent
graph with \(b-2\) degree-\(1\) vertices and \(c-k\) degree-\(2\) vertices.
This gives the displayed recurrence for \(\mathsf N(b,c)\), with base
\(\mathsf N(0,c)=\mathsf C(c)\).  The construction is reversible by taking the
component of the distinguished \(B\)-vertex, so the count is exact.
\end{proof}

\begin{proposition}[Exact degree profiles fit the full final-zone budget]\label{prop:bc-exact-degree-profile-low-low-full-budget}
Work in the setting of \cref{def:bc-degree-low-low-candidates}.  Put
\[
  r_1=\#\{1\le a\le2m:d_a=1\},\qquad
  r_2=\#\{1\le a\le2m:d_a=2\}.
\]
If
\[
  \mathsf N(r_1,r_2)\le2^{2m},
\]
then the degree-constrained set \(\mathcal L_{\mathrm{deg}}\) itself is a
canonical full-budget low-low candidate set: it is determined by the
target-window data, contains the true low-low Step~3 assignment on \(L_{2m}\),
and satisfies
\[
  |\mathcal L_{\mathrm{deg}}|\le2^{2m}.
\]
Consequently the final zone satisfies the right-boundary inverse hypothesis of
\cref{prop:bc-right-boundary-auxiliary-global-recovery} by
\cref{prop:bc-full-budget-low-low-candidates-supply-large-right-boundary-auxiliary}.
\end{proposition}

\begin{proof}
The target-window data determine the integers \(d_a\), hence determine
\(r_1\), \(r_2\), and the finite set \(\mathcal L_{\mathrm{deg}}\).  By
\cref{prop:bc-strict-left-degree-low-low-pruning}, the true low-low assignment
belongs to \(\mathcal L_{\mathrm{deg}}\).

If \(\mathcal L_{\mathrm{deg}}\) is empty there is nothing to prove.  Otherwise
the degree equations give \(d_a\ge0\), while \(d_a=2-c_a\) and the recovered
entries \(c_a\) are nonnegative give \(d_a\le2\).  Thus every \(d_a\) lies in
\(\{0,1,2\}\).  Relabel the vertices with residual degree
\(1\) as \(B\) and the vertices with residual degree \(2\) as \(C\).  An
assignment
\[
  \ell=(\ell_{ab})_{1\le a<b\le2m}\in\mathcal L_{\mathrm{deg}}
\]
is exactly a loopless labelled multigraph on \(B\sqcup C\): the edge
multiplicity between \(a\) and \(b\) is \(\ell_{ab}\), the degree equations
give degree \(1\) on \(B\) and degree \(2\) on \(C\), and labels with
\(d_a=0\) are isolated and irrelevant.  Conversely every such multigraph gives
one degree-compatible low-low assignment by reading its edge multiplicities.
Thus
\[
  |\mathcal L_{\mathrm{deg}}|=\mathsf N(r_1,r_2).
\]
The displayed hypothesis gives the full-budget bound.  Taking
\(\mathcal L_{\mathrm{tw}}=\mathcal L_{\mathrm{deg}}\) in
\cref{prop:bc-full-budget-low-low-candidates-supply-large-right-boundary-auxiliary}
gives the stated right-boundary inverse.
\end{proof}

\begin{proposition}[Partial low-low readouts with bounded residual profile fit the full budget]\label{prop:bc-partial-low-low-profile-full-budget}
Work in the setting of \cref{def:bc-degree-low-low-candidates}, and suppose
\(m\ge9\).  Let
\[
  \mathcal E_{\mathrm{tw}}
  \subseteq
  \{\{a,b\}:1\le a<b\le2m\}
\]
be a target-window-determined set of low-low cells, and suppose the same
target-window data determine nonnegative integers
\[
  p_{\{a,b\}}\qquad(\{a,b\}\in\mathcal E_{\mathrm{tw}})
\]
which agree with the true Step~3 entries on those cells.  Put
\[
  d_a^{\mathcal E}
  =
  d_a-\sum_{\substack{1\le b\le2m\\ b\ne a\\ \{a,b\}\in\mathcal E_{\mathrm{tw}}}}
  p_{\{a,b\}}
  \qquad(1\le a\le2m),
\]
and
\[
  r_1^{\mathcal E}=\#\{a:d_a^{\mathcal E}=1\},\qquad
  r_2^{\mathcal E}=\#\{a:d_a^{\mathcal E}=2\}.
\]
If
\[
  \mathsf N(r_1^{\mathcal E},r_2^{\mathcal E})\le2^{2m},
\]
then the target-window data determine a canonical full-budget low-low
candidate set containing the true low-low assignment.  Consequently the final
zone satisfies the right-boundary inverse hypothesis of
\cref{prop:bc-right-boundary-auxiliary-global-recovery} by
\cref{prop:bc-full-budget-low-low-candidates-supply-large-right-boundary-auxiliary}.
\end{proposition}

\begin{proof}
Define
\[
  \mathcal L_{\mathrm{part}}
  =
  \{\ell\in\mathcal L_{\mathrm{deg}}:
    \ell_{ab}=p_{\{a,b\}}\text{ for every }
    \{a,b\}\in\mathcal E_{\mathrm{tw}}\}.
\]
This set is a canonical function of the target-window data, with the inherited
lexicographic order from \(\mathcal L_{\mathrm{deg}}\).  The true low-low
assignment belongs to \(\mathcal L_{\mathrm{deg}}\) by
\cref{prop:bc-strict-left-degree-low-low-pruning}, and it agrees with the
displayed partial readout by hypothesis, so it belongs to
\(\mathcal L_{\mathrm{part}}\).

It remains to bound the size.  If \(\mathcal L_{\mathrm{part}}\) is empty
there is nothing to prove, so suppose it is nonempty.  The agreement equations
show that, after removing the already read entries
\(p_{\{a,b\}}\), every remaining completion has residual degrees
\(d_a^{\mathcal E}\).  Since a true completion exists, these residual degrees
are nonnegative; because \(d_a\le2\) and the removed entries are nonnegative,
they lie in \(\{0,1,2\}\).

For \(\ell\in\mathcal L_{\mathrm{part}}\), form a residual loopless multigraph
by keeping edge multiplicity \(\ell_{ab}\) on pairs not in
\(\mathcal E_{\mathrm{tw}}\) and edge multiplicity \(0\) on pairs in
\(\mathcal E_{\mathrm{tw}}\).  This residual graph has degree
\(d_a^{\mathcal E}\) at label \(a\).  The map is injective, because the
removed multiplicities \(p_{\{a,b\}}\) are fixed and adding them back on
\(\mathcal E_{\mathrm{tw}}\) recovers the original \(\ell\).  Therefore
\(\mathcal L_{\mathrm{part}}\) injects into the set of all loopless labelled
multigraphs with \(r_1^{\mathcal E}\) degree-one vertices and
\(r_2^{\mathcal E}\) degree-two vertices.  By
\cref{def:bc-exact-low-low-degree-profile-count}, that ambient set has
cardinality \(\mathsf N(r_1^{\mathcal E},r_2^{\mathcal E})\).  The displayed
hypothesis gives
\[
  |\mathcal L_{\mathrm{part}}|\le2^{2m}.
\]

Taking \(\mathcal L_{\mathrm{tw}}=\mathcal L_{\mathrm{part}}\) in
\cref{prop:bc-full-budget-low-low-candidates-supply-large-right-boundary-auxiliary}
gives the stated right-boundary inverse.
\end{proof}

\begin{corollary}[Partial low-low readouts below the certified frontier fit the full budget]\label{cor:bc-partial-low-low-frontier-full-budget}
Work in the setting of
\cref{prop:bc-partial-low-low-profile-full-budget}, with \(m\ge29\).  Let
\(A(m)\) and \(B(m)\) be the frontier functions of
\cref{prop:bc-active-twenty-low-low-profile-full-budget}, and put
\[
  D^{\mathcal E}=\sum_{a=1}^{2m}d_a^{\mathcal E}.
\]
If
\[
  r_1^{\mathcal E}+r_2^{\mathcal E}<A(m)
  \qquad\text{or}\qquad
  D^{\mathcal E}<B(m),
\]
then the same full-budget low-low candidate conclusion and right-boundary
inverse hold.
\end{corollary}

\begin{proof}
By the definitions of the frontier functions \(A(m)\) and \(B(m)\) in
\cref{prop:bc-active-twenty-low-low-profile-full-budget},
\[
  \mathsf N(r_1^{\mathcal E},r_2^{\mathcal E})\le2^{2m}
\]
under either displayed inequality.  Applying
\cref{prop:bc-partial-low-low-profile-full-budget} gives the result.
\end{proof}

\begin{corollary}[Partial low-low readouts with enough edge mass fit the full budget]\label{cor:bc-partial-low-low-edge-mass-frontier-full-budget}
Work in the setting of
\cref{prop:bc-partial-low-low-profile-full-budget}, with \(m\ge29\), and put
\[
  D=\sum_{a=1}^{2m}d_a,\qquad
  P_{\mathcal E}=
  \sum_{\{a,b\}\in\mathcal E_{\mathrm{tw}}}p_{\{a,b\}} .
\]
Let \(B(m)\) be the total-degree frontier of
\cref{prop:bc-active-twenty-low-low-profile-full-budget}.  If
\[
  2P_{\mathcal E}>D-B(m),
\]
then the same full-budget low-low candidate conclusion and right-boundary
inverse hold.
\end{corollary}

\begin{proof}
Every read low-low entry contributes once to each of its two endpoint residual
degrees.  Therefore
\[
  D^{\mathcal E}
  =
  \sum_{a=1}^{2m}d_a^{\mathcal E}
  =
  D-2P_{\mathcal E}.
\]
The displayed inequality gives \(D^{\mathcal E}<B(m)\).  Applying
\cref{cor:bc-partial-low-low-frontier-full-budget} gives the result.
\end{proof}

\begin{corollary}[Incidence-covered low labels below the certified frontier fit the full budget]\label{cor:bc-incidence-covered-low-labels-frontier-full-budget}
Work in the setting of \cref{def:bc-degree-low-low-candidates}, with
\(m\ge29\).  Let
\[
  U_{\mathrm{tw}}\subseteq\{1,\ldots,2m\}
\]
be a target-window-determined set of low labels.  Suppose the same
target-window data determine the true Step~3 entries
\[
  p_{\{a,b\}}\qquad
  (1\le a<b\le2m,\ \{a,b\}\cap U_{\mathrm{tw}}\ne\varnothing)
\]
on every low-low cell incident to \(U_{\mathrm{tw}}\).  Let \(A(m)\) and
\(B(m)\) be the frontier functions of
\cref{prop:bc-active-twenty-low-low-profile-full-budget}.  If
\[
  \#\{a\notin U_{\mathrm{tw}}:d_a>0\}<A(m)
  \qquad\text{or}\qquad
  \sum_{a\notin U_{\mathrm{tw}}}d_a<B(m),
\]
then the target-window data determine a canonical full-budget low-low
candidate set containing the true low-low assignment.  Consequently the final
zone satisfies the right-boundary inverse hypothesis of
\cref{prop:bc-right-boundary-auxiliary-global-recovery} by
\cref{prop:bc-full-budget-low-low-candidates-supply-large-right-boundary-auxiliary}.
\end{corollary}

\begin{proof}
Apply \cref{prop:bc-partial-low-low-profile-full-budget} to the incidence
readout
\[
  \mathcal E_{\mathrm{tw}}
  =
  \{\{a,b\}:1\le a<b\le2m,\ \{a,b\}\cap U_{\mathrm{tw}}\ne\varnothing\}.
\]
This is a target-window-determined low-low cell set, and the displayed
entries agree with the true Step~3 entries by hypothesis.

Let \(d_a^{\mathcal E}\) be the residual degrees after this readout.  The true
low-low assignment belongs to the corresponding partial candidate set, so
these residual degrees are nonnegative.  For \(a\in U_{\mathrm{tw}}\), all
low-low cells incident to \(a\) have been read.  Since the true low-low
assignment has low-low degree \(d_a\) at \(a\), this gives
\[
  d_a^{\mathcal E}=0 .
\]
For \(a\notin U_{\mathrm{tw}}\), the removed incident entries are
nonnegative, so
\[
  0\le d_a^{\mathcal E}\le d_a .
\]
Thus
\[
  r_1^{\mathcal E}+r_2^{\mathcal E}
  =
  \#\{a:d_a^{\mathcal E}>0\}
  \le
  \#\{a\notin U_{\mathrm{tw}}:d_a>0\}
\]
and
\[
  D^{\mathcal E}
  =
  \sum_{a=1}^{2m}d_a^{\mathcal E}
  \le
  \sum_{a\notin U_{\mathrm{tw}}}d_a .
\]
Under either displayed hypothesis, the residual active support or the
residual total degree is below the certified frontier.  Therefore
\cref{cor:bc-partial-low-low-frontier-full-budget} applies to this incidence
readout and gives the stated full-budget candidate set and right-boundary
inverse.
\end{proof}

\begin{corollary}[Large incidence-cover readouts fit the full budget]\label{cor:bc-large-incidence-cover-frontier-full-budget}
Work in the setting of
\cref{cor:bc-incidence-covered-low-labels-frontier-full-budget}, and put
\[
  D=\sum_{a=1}^{2m}d_a,\qquad
  D(U_{\mathrm{tw}})=\sum_{a\in U_{\mathrm{tw}}}d_a .
\]
If
\[
  |U_{\mathrm{tw}}|>2m-A(m)
  \qquad\text{or}\qquad
  D(U_{\mathrm{tw}})>D-B(m),
\]
then the target-window data determine a canonical full-budget low-low
candidate set containing the true low-low assignment, and the final zone
satisfies the right-boundary inverse hypothesis of
\cref{prop:bc-right-boundary-auxiliary-global-recovery}.
\end{corollary}

\begin{proof}
If \(|U_{\mathrm{tw}}|>2m-A(m)\), then
\[
  \#\{a\notin U_{\mathrm{tw}}:d_a>0\}
  \le 2m-|U_{\mathrm{tw}}|<A(m).
\]
If \(D(U_{\mathrm{tw}})>D-B(m)\), then
\[
  \sum_{a\notin U_{\mathrm{tw}}}d_a
  =
  D-D(U_{\mathrm{tw}})<B(m).
\]
Thus either displayed hypothesis implies one of the two frontier hypotheses in
\cref{cor:bc-incidence-covered-low-labels-frontier-full-budget}.  Applying
that corollary gives the asserted full-budget candidate set and right-boundary
inverse.
\end{proof}

\begin{proposition}[Exact low-low profile frontiers fit the full final-zone budget]\label{prop:bc-active-twenty-low-low-profile-full-budget}
Work in the setting of \cref{def:bc-degree-low-low-candidates}, with
\(m\ge29\), and put
\[
  r_1=\#\{1\le a\le2m:d_a=1\},\qquad
  r_2=\#\{1\le a\le2m:d_a=2\},\qquad
  D=\sum_{a=1}^{2m}d_a.
\]
If
\[
  r_1+r_2\le20,
\]
then the degree-constrained set \(\mathcal L_{\mathrm{deg}}\) itself is a
canonical full-budget low-low candidate set.  Consequently the final zone
satisfies the right-boundary inverse hypothesis of
\cref{prop:bc-right-boundary-auxiliary-global-recovery} by
\cref{prop:bc-full-budget-low-low-candidates-supply-large-right-boundary-auxiliary}.

More generally, define the exact active-support frontier
\[
  A(m)=
  \min\Bigl(
    \{\,u+v:
        0\le u,v,\ u+v\le2m,\
        \mathsf N(u,v)>2^{2m}\,\}
    \cup\{2m+1\}
  \Bigr).
\]
If \(r_1+r_2<A(m)\), then the same full-budget conclusion holds.  Hence the
remaining exact dense degree-profile branch has \(r_1+r_2\ge A(m)\).

Define also the exact total-degree frontier
\[
  B(m)=
  \min\Bigl(
    \{\,u+2v:
        0\le u,v,\ u+v\le2m,\
        \mathsf N(u,v)>2^{2m}\,\}
    \cup\{4m+1\}
  \Bigr).
\]
If \(D<B(m)\), then the same full-budget conclusion holds.  Hence the
remaining exact dense degree-profile branch also has \(D\ge B(m)\).

The active-support threshold is sharp for the degree-profile count: at
\((r_1,r_2)=(0,21)\),
\[
  \mathsf N(0,21)=4871976826254018760
  >
  288230376151711744=2^{58}.
\]
\end{proposition}

\begin{proof}
The exact C++/GMP verifier
\[
\texttt{character\_ring\_iter/post\_m29\_bc\_low\_low\_profile\_gmp.cpp}
\]
implements the recurrences of
\cref{def:bc-exact-low-low-degree-profile-count} for
\(0\le r_1,r_2\le436\), and checks all \(29\le m\le218\) with
\(r_1+r_2\le2m\).  In particular, every profile with
\(r_1+r_2\le20\) is included already in the \(m=29\) pass.  The archived
`optimus' run
\[
\texttt{certificates/post\_m29/bc\_low\_low\_profile\_gmp\_certificate.log}
\]
has SHA-256
\[
\hashsplit{cc1b41189d965a5e076d3524647caa5b}{7f5c5102ca6a1d8b03c623c15c29dc40}.
\]
The source file has SHA-256
\[
\hashsplit{e99bc1a746f569763d13ef4421797019}{a0162faea4f9b65c0ca5b51e31c3f50f}.
\]
The log ends with \texttt{\_\_EXIT\_STATUS=0} and
\texttt{SUMMARY checked\_m=190 failures=0}.  Its active-support line reports
\[
\begin{gathered}
\texttt{ACTIVE20 max\_profile=(0,20)},\\
\texttt{max\_count=237669544014377896},\\
\texttt{budget\_2\string^58=288230376151711744},\qquad
\texttt{active20\_failures=0}.
\end{gathered}
\]
Thus, whenever \(r_1+r_2\le20\),
\[
  \mathsf N(r_1,r_2)
  \le 237669544014377896
  <2^{58}
  \le2^{2m},
\]
because \(m\ge29\).  Applying
\cref{prop:bc-exact-degree-profile-low-low-full-budget} gives the asserted
full-budget low-low candidate set and right-boundary inverse.

By the definition of \(A(m)\), if \(r_1+r_2<A(m)\), then the exact profile
count \(\mathsf N(r_1,r_2)\) is at most \(2^{2m}\).  Therefore
\cref{prop:bc-exact-degree-profile-low-low-full-budget} again gives the
full-budget conclusion.  Likewise, since \(D=r_1+2r_2\), the definition of
\(B(m)\) shows that \(D<B(m)\) implies
\(\mathsf N(r_1,r_2)\le2^{2m}\), and the same proposition gives the
full-budget conclusion.

For the residual range used in the post-\(m=29\) \(B/C\) tail, the same log
reports one \texttt{ACTIVE\_FRONTIER} and one \texttt{DEGREE\_FRONTIER} line
for every interval on which these exact frontiers are constant, over
\(29\le m\le218\).  Its summary line records
\[
\texttt{global\_min\_bad\_active=21}
\quad\text{and}\quad
\texttt{global\_min\_bad\_D=34},
\]
and the interval lines end at
\[
\texttt{ACTIVE\_FRONTIER m=218 min\_dense\_active=88}
\]
and
\[
\texttt{DEGREE\_FRONTIER m=218 min\_dense\_D=152}.
\]
These printed values are not additional hypotheses; they are the exact
certificate readout of the two frontier functions just defined on the finite
residual range.

The same log reports
\[
\begin{gathered}
\texttt{FIRST\_DENSE profile=(0,21)},\\
\texttt{count=4871976826254018760},\qquad
\texttt{exceeds\_2\string^58=yes},
\end{gathered}
\]
which is the displayed sharpness statement.
\end{proof}

\begin{corollary}[Continuation-heavy profiles fit the full final-zone budget]\label{cor:bc-continuation-heavy-low-low-profile-full-budget}
Work in the setting of \cref{def:bc-degree-low-low-candidates}, with
\(m\ge29\), on a true final-zone datum.  Let \(A(m)\) and \(B(m)\) be the
active-support and total-degree frontiers defined in
\cref{prop:bc-active-twenty-low-low-profile-full-budget}.  Put
\[
  c_a=\sum_{t=2m+1}^j u_{\{a,t\}},\qquad
  C=\sum_{a=1}^{2m}c_a,\qquad
  r_0=\#\{1\le a\le2m:d_a=0\}.
\]
If
\[
  r_0>2m-A(m)
  \qquad\text{or}\qquad
  C>4m-B(m),
\]
then the degree-constrained set \(\mathcal L_{\mathrm{deg}}\) itself is a
canonical full-budget low-low candidate set.  Consequently the final zone
satisfies the right-boundary inverse hypothesis of
\cref{prop:bc-right-boundary-auxiliary-global-recovery} by
\cref{prop:bc-full-budget-low-low-candidates-supply-large-right-boundary-auxiliary}.

Equivalently, the remaining exact dense degree-profile branch has
\[
  r_0\le2m-A(m),\qquad C\le4m-B(m).
\]
Thus at least \(A(m)\) low labels have at most one continuation incidence, and
the true low-low subfilling has at least \(B(m)/2\) low-low edges.
\end{corollary}

\begin{proof}
Since we are on a true final-zone datum, the true low-low assignment belongs to
\(\mathcal L_{\mathrm{deg}}\) by
\cref{prop:bc-strict-left-degree-low-low-pruning}.  Hence the residual degrees
are \(d_a\in\{0,1,2\}\), with \(d_a=2-c_a\).  Therefore
\[
  r_1+r_2=\#\{a:d_a>0\}=2m-r_0
\]
and
\[
  D=\sum_{a=1}^{2m}d_a
   =\sum_{a=1}^{2m}(2-c_a)
   =4m-C.
\]
If \(r_0>2m-A(m)\), then \(r_1+r_2<A(m)\), so
\cref{prop:bc-active-twenty-low-low-profile-full-budget} gives the
full-budget conclusion.  If \(C>4m-B(m)\), then \(D<B(m)\), and the same
proposition again gives the full-budget conclusion.

Taking the contraposition of these two closed cases gives the displayed
conditions in the remaining exact dense branch.  For the final sentence, the
first condition says that at least \(A(m)\) labels have \(d_a>0\), equivalently
\(c_a\le1\).  The true low-low assignment has total low-low degree \(D\), so
its number of low-low edges is \(D/2\); in the remaining branch
\(D\ge B(m)\), giving at least \(B(m)/2\) such edges.
\end{proof}

\begin{corollary}[Zero-residual incidence covers are already exhausted]\label{cor:bc-zero-residual-incidence-cover-exhausted}
Work in the setting of \cref{def:bc-degree-low-low-candidates}, with
\(m\ge29\), on a true final-zone datum.  Let \(A(m)\) and \(B(m)\) be the
frontier functions of
\cref{prop:bc-active-twenty-low-low-profile-full-budget}, and put
\[
  U_0=\{1\le a\le2m:d_a=0\},\qquad
  r_i=\#\{1\le a\le2m:d_a=i\}\quad(0\le i\le2),\qquad
  D=\sum_{a=1}^{2m}d_a .
\]
Then the target-window data determine the true Step~3 entries on every low-low
cell incident to \(U_0\): all such entries are zero.

Moreover, if the incidence-cover full-budget tests of
\cref{cor:bc-large-incidence-cover-frontier-full-budget} hold for \(U_0\), or
if the two complement tests of
\cref{cor:bc-incidence-covered-low-labels-frontier-full-budget} hold for
\(U_0\), then the full-budget low-low candidate conclusion already follows
from
\cref{prop:bc-active-twenty-low-low-profile-full-budget,cor:bc-continuation-heavy-low-low-profile-full-budget}.
Equivalently, in the live exact dense branch where
\[
  \mathsf N(r_1,r_2)>2^{2m},\qquad
  r_0\le2m-A(m),\qquad
  D\ge B(m),
\]
the automatic cover \(U_0\) satisfies none of those incidence-cover
full-budget tests.
\end{corollary}

\begin{proof}
The set \(U_0\) is determined by the residual degrees \(d_a\), hence by the
target-window data.  Let \(\ell^{\mathrm{true}}\) be the true low-low Step~3
assignment.  By \cref{prop:bc-strict-left-degree-low-low-pruning}, it belongs
to \(\mathcal L_{\mathrm{deg}}\).  If \(a\in U_0\), then the low-low degree of
label \(a\) in \(\ell^{\mathrm{true}}\) is \(d_a=0\).  Since all Step~3
entries are nonnegative, every low-low entry incident to \(a\) is zero.  Thus
the target-window data determine those true incident entries.

For \(U_0\), the active complement is exactly the set of labels with
positive residual degree, so its size is \(r_1+r_2\).  If this is below
\(A(m)\), then
\cref{prop:bc-active-twenty-low-low-profile-full-budget} gives the
full-budget conclusion.  The residual degree on the complement is \(D\),
because \(d_a=0\) on \(U_0\); if \(D<B(m)\), the same proposition gives the
full-budget conclusion.  The covered-label test is
\[
  |U_0|=r_0>2m-A(m),
\]
which is one of the two closed cases in
\cref{cor:bc-continuation-heavy-low-low-profile-full-budget}.  Finally,
\[
  D(U_0)=\sum_{a\in U_0}d_a=0,
\]
so the covered-degree test \(D(U_0)>D-B(m)\) is equivalent to \(D<B(m)\),
again closed by
\cref{prop:bc-active-twenty-low-low-profile-full-budget}.

The final sentence is the contraposition of the four closed alternatives just
listed, together with the definition of the displayed live branch.
\end{proof}

\begin{corollary}[Zero-only incidence covers are not live sources]\label{cor:bc-zero-only-incidence-covers-not-live}
Work in the setting of
\cref{cor:bc-zero-residual-incidence-cover-exhausted}, and suppose the live
exact dense branch satisfies
\[
  \mathsf N(r_1,r_2)>2^{2m},\qquad
  r_0\le2m-A(m),\qquad
  D\ge B(m).
\]
Let \(U_{\mathrm{tw}}\subseteq U_0\).  Then its positive-residual projection is
empty, and hence it satisfies none of the positive-residual two-threshold
incidence-cover source tests
\[
  |U_+|>r_1+r_2-A(m)
  \qquad\text{or}\qquad
  \sum_{a\in U_+}d_a>D-B(m),
  \qquad U_+=U_{\mathrm{tw}}\cap\{a:d_a>0\}.
\]
Hence every incidence-cover source in the live branch must include at least one
positive-residual label.
\end{corollary}

\begin{proof}
Since \(U_{\mathrm{tw}}\subseteq U_0\), its positive-residual projection is
\(U_+=\varnothing\).  By the definition of \(A(m)\), the inequality
\(\mathsf N(r_1,r_2)>2^{2m}\) gives
\[
  r_1+r_2\ge A(m).
\]
Thus
\[
  |U_+|=0\le r_1+r_2-A(m),
\]
so the positive-label-count threshold fails.  Also \(D\ge B(m)\), so
\[
  \sum_{a\in U_+}d_a=0\le D-B(m),
\]
so the positive residual-degree threshold fails.  Therefore a live
incidence-cover source cannot be supported only on zero-residual labels.
\end{proof}

\begin{corollary}[Incidence covers project to positive residual labels]\label{cor:bc-incidence-cover-positive-projection}
Work in the setting of
\cref{cor:bc-large-incidence-cover-frontier-full-budget}.  Suppose the
target-window data determine a set
\[
  U_{\mathrm{tw}}\subseteq\{1,\ldots,2m\}
\]
and determine the true Step~3 entries on every low-low cell incident to
\(U_{\mathrm{tw}}\).  Put
\[
  U_+=\{a\in U_{\mathrm{tw}}:d_a>0\}.
\]
Then the same data determine \(U_+\) and determine the true Step~3 entries on
every low-low cell incident to \(U_+\).  Moreover, if any one of the four
incidence-cover tests
\[
  \#\{a\notin U_{\mathrm{tw}}:d_a>0\}<A(m),\qquad
  \sum_{a\notin U_{\mathrm{tw}}}d_a<B(m),
\]
\[
  |U_{\mathrm{tw}}|>2m-A(m),\qquad
  \sum_{a\in U_{\mathrm{tw}}}d_a>D-B(m)
\]
holds for \(U_{\mathrm{tw}}\), then \(U_+\) satisfies at least one of the three
positive-residual tests
\[
  \#\{a\notin U_+:d_a>0\}<A(m),\qquad
  \sum_{a\notin U_+}d_a<B(m),
  \qquad
  \sum_{a\in U_+}d_a>D-B(m).
\]
Thus every incidence-cover source may be stated, without loss, with
\(U_{\mathrm{tw}}\subseteq\{a:d_a>0\}\) and with only these three tests.
\end{corollary}

\begin{proof}
The residual degrees \(d_a\) are target-window determined in
\cref{def:bc-degree-low-low-candidates}, so \(U_+\) is determined.  Since
\(U_+\subseteq U_{\mathrm{tw}}\), the readout of all cells incident to
\(U_{\mathrm{tw}}\) includes all cells incident to \(U_+\).

The positive labels outside \(U_+\) are exactly the positive labels outside
\(U_{\mathrm{tw}}\).  Hence
\[
  \#\{a\notin U_+:d_a>0\}
  =
  \#\{a\notin U_{\mathrm{tw}}:d_a>0\}.
\]
Also labels in \(U_{\mathrm{tw}}\setminus U_+\) have residual degree zero, so
\[
  \sum_{a\notin U_+}d_a
  =
  \sum_{a\notin U_{\mathrm{tw}}}d_a
  \qquad\text{and}\qquad
  \sum_{a\in U_+}d_a
  =
  \sum_{a\in U_{\mathrm{tw}}}d_a .
\]
Therefore the first, second, and fourth original tests imply respectively the
first, second, and third positive-residual tests.

It remains only to handle the covered-label test.  If
\(|U_{\mathrm{tw}}|>2m-A(m)\), then
\[
  \#\{a\notin U_+:d_a>0\}
  =
  \#\{a\notin U_{\mathrm{tw}}:d_a>0\}
  \le
  2m-|U_{\mathrm{tw}}|
  <A(m),
\]
so the first positive-residual test holds.  This proves the asserted
normalization.
\end{proof}

\begin{corollary}[Positive incidence covers have two threshold tests]\label{cor:bc-positive-incidence-two-thresholds}
Work in the setting of
\cref{cor:bc-incidence-cover-positive-projection}.  Put
\[
  r_i=\#\{1\le a\le2m:d_a=i\}\quad(i=1,2),\qquad
  r_+=\#\{1\le a\le2m:d_a>0\}=r_1+r_2,\qquad
  D=\sum_{a=1}^{2m}d_a .
\]
If
\[
  U_{\mathrm{tw}}\subseteq\{1\le a\le2m:d_a>0\},
\]
then the disjunction of the three positive-residual incidence-cover tests
\[
  \#\{a\notin U_{\mathrm{tw}}:d_a>0\}<A(m),\qquad
  \sum_{a\notin U_{\mathrm{tw}}}d_a<B(m),
  \qquad
  \sum_{a\in U_{\mathrm{tw}}}d_a>D-B(m)
\]
is equivalent to the disjunction of the two threshold alternatives
\[
  |U_{\mathrm{tw}}|>r_+-A(m)
  \qquad\text{or}\qquad
  \sum_{a\in U_{\mathrm{tw}}}d_a>D-B(m).
\]
\end{corollary}

\begin{proof}
Since \(U_{\mathrm{tw}}\) contains only positive-residual labels,
\[
  \#\{a\notin U_{\mathrm{tw}}:d_a>0\}
  =
  r_+-|U_{\mathrm{tw}}|.
\]
Thus the active-complement test is equivalent to
\[
  |U_{\mathrm{tw}}|>r_+-A(m).
\]
Also
\[
  \sum_{a\notin U_{\mathrm{tw}}}d_a
  =
  D-\sum_{a\in U_{\mathrm{tw}}}d_a,
\]
so the complement residual-degree test is equivalent to the displayed
covered-degree test.  This proves the two-threshold form.
\end{proof}

\begin{lemma}[Degree-two half-square fillings have empty endpoints]\label{lem:bc-degree-two-half-square-empty-endpoints}
Use the Step~3 dual-RSK' half-square in the \(m=0\), content \((2^j)\)
specialization of \cref{lem:bc-growth-path-model}.  Let \(u\) be a
nonnegative integer filling of the Step~3 half-square, and let \(N\) be the
symmetric zero-diagonal matrix obtained by reflecting \(u\) across the
diagonal.  If every row of \(N\) has sum \(2\), then the top/right boundary
path produced from \(u\) by the forward dual-RSK' sweep satisfies
\[
  \rho_0=\rho_{2j}=\varnothing .
\]
\end{lemma}

\begin{proof}
The row-sum hypothesis says exactly that the two equal semistandard boundary
weights in Krattenthaler's Step~2 construction are \((2^j)\).  The matrix is
symmetric and has zero diagonal, so it is one of the \(m=0\) inputs in
\cite[Theorem 4, Steps 2--3]{Krattenthaler2016}: Step~3 uses one
off-diagonal half of this matrix as the dual-RSK' half-square filling, as
recorded in \cref{lem:bc-step3-cells-are-step2-half-matrix}.  The boundary
path is therefore one of the paths in
\cref{lem:bc-growth-path-model} with even-column endpoint parameter \(m=0\).
That lemma gives initial and terminal partitions equal to \(\varnothing\).
\end{proof}

\begin{definition}[Endpoint-compatible low-low candidates]\label{def:bc-endpoint-compatible-low-low-candidates}
Work in the setting of \cref{def:bc-degree-low-low-candidates}, and put
\(M=2m\).  For
\[
  \ell=(\ell_{ab})_{1\le a<b\le M}\in\mathcal L_{\mathrm{deg}},
\]
form a complete Step~3 half-square filling \(u^\ell\) as follows.  On the
low-low cell set \(L_M\), put \(u^\ell_{\{a,b\}}=\ell_{ab}\).  On each cell
of the continuation-column cell set \(J_M\), put \(u^\ell\) equal to the entry
recovered from \(\mathcal D_{\mathrm{out}}(P)\) by
\cref{lem:bc-strict-left-recovers-continuation-columns}.  These two
subarrangements cover the full Step~3 half-square.  Let
\[
  \rho^\ell_0,\rho^\ell_1,\ldots,\rho^\ell_{2j}
\]
be the top/right boundary sequence obtained from \(u^\ell\) by the deterministic
forward dual-RSK' sweep of
\cref{lem:bc-step3-full-half-square-determines-rho}.  The endpoint-compatible
low-low candidate set
\[
  \mathcal L_{\mathrm{end}}
  \subseteq\mathcal L_{\mathrm{deg}}
\]
is the set of those \(\ell\in\mathcal L_{\mathrm{deg}}\) for which
\[
  \rho^\ell_{2(j-2m)}=\rho_{2a_L-2}
  \qquad\text{and}\qquad
  \rho^\ell_{2j}=\rho_{2j}.
\]
It is ordered by the lexicographic order inherited from
\(\mathcal L_{\mathrm{deg}}\).
\end{definition}

\begin{proposition}[Endpoint compatibility gives no extra low-low pruning]\label{prop:bc-endpoint-compatible-low-low-equals-degree}
Work in the setting of
\cref{def:bc-endpoint-compatible-low-low-candidates}, on a nonempty true
final-zone target-window datum.  Then
\[
  \mathcal L_{\mathrm{end}}=\mathcal L_{\mathrm{deg}} .
\]
Consequently the endpoint-compatible filter does not give a source beyond the
exact degree-profile closures of
\cref{prop:bc-exact-degree-profile-low-low-full-budget,prop:bc-active-twenty-low-low-profile-full-budget}.
\end{proposition}

\begin{proof}
The inclusion \(\mathcal L_{\mathrm{end}}\subseteq\mathcal L_{\mathrm{deg}}\)
is part of \cref{def:bc-endpoint-compatible-low-low-candidates}.  Conversely
take \(\ell\in\mathcal L_{\mathrm{deg}}\), and form the complete Step~3
half-square filling \(u^\ell\).

First consider the endpoint at the cut between the strict-left continuation
columns and the terminal low-low corner.  The entries of \(u^\ell\) on
\(J_M\) are, by definition, the entries recovered from
\(\mathcal D_{\mathrm{out}}(P)\).  The first \(j-M\) strict-left columns of the
Step~3 half-square are exactly \(J_M\) by
\cref{lem:bc-step3-continuation-incidence-covers-q-suffix}.  The forward
dual-RSK' sweep on this Ferrers subarrangement has fixed empty incoming
boundary labels and is deterministic by
\cref{lem:bc-dual-rsk-prime-growth-local-reversibility}.  Applying the same
sweep to the true filling gives the true strict-left segment recovered in
\cref{lem:bc-strict-left-recovers-continuation-columns}; hence
\[
  \rho^\ell_{2(j-M)}=\rho_{2(j-M)}.
\]
Since \(M=2m\) and \(a_L=j-2m+1\), this is the first endpoint condition in
\cref{def:bc-endpoint-compatible-low-low-candidates}.

It remains to check the terminal endpoint.  From the entries of \(u^\ell\)
form the symmetric zero-diagonal matrix \(N^\ell\) whose off-diagonal entries
are the corresponding Step~3 half-square entries.  For \(1\le a\le M\), the
degree equations defining \(\mathcal L_{\mathrm{deg}}\) give
\[
  \sum_{b\ne a}N^\ell_{ab}
  =
  d_a+c_a
  =
  2.
\]
For a continuation label \(t>M\), every off-diagonal entry incident to \(t\)
lies in \(J_M\), so those entries are the true recovered entries.  The true
Step~2 matrix has row sum \(2\) by
\cref{lem:bc-step2-matrix-doubled-content-degrees}.  Thus every row of
\(N^\ell\) has sum \(2\).

Thus \(N^\ell\) satisfies the hypotheses of
\cref{lem:bc-degree-two-half-square-empty-endpoints}.  Hence
\[
  \rho^\ell_{2j}=\varnothing=\rho_{2j},
\]
where the last equality is the endpoint assertion of
\cref{lem:bc-growth-path-model} for the true path.

Both endpoint tests hold for the arbitrary \(\ell\in\mathcal L_{\mathrm{deg}}\),
so \(\ell\in\mathcal L_{\mathrm{end}}\).  This proves the reverse inclusion.
The final sentence follows immediately: bounding
\(\mathcal L_{\mathrm{end}}\) is the same as bounding
\(\mathcal L_{\mathrm{deg}}\), which is exactly the degree-profile budget
problem handled by the cited exact profile results.
\end{proof}

\begin{proposition}[Degree-only low-low pruning exceeds both branch budgets]\label{prop:bc-degree-low-low-ambient-too-large}
The degree equations in
\cref{def:bc-degree-low-low-candidates} do not imply the branch-budget bound
\[
  |\mathcal L_{\mathrm{deg}}|\le2^m
\]
by themselves.  More precisely, for every $m\ge3$ there is a residual degree
vector on $M=2m$ low labels for which the corresponding degree-constrained
low-low ambient set has cardinality
\[
  (2m-1)!!>2^m .
\]
For every $m\ge7$, the same examples also violate the full final-zone branch
budget:
\[
  (2m-1)!!>2^{2m}.
\]
\end{proposition}

\begin{proof}
Take
\[
  d_1=d_2=\cdots=d_{2m}=1 .
\]
An assignment
\[
  \ell=(\ell_{ab})_{1\le a<b\le2m}
\]
satisfying the degree equations of
\cref{def:bc-degree-low-low-candidates} for this vector is exactly a perfect
matching on the vertex set $\{1,\ldots,2m\}$: every vertex has degree one, so
each nonzero $\ell_{ab}$ equals $1$, and the nonzero pairs partition the
vertices.  Conversely every perfect matching gives such an assignment.

The number of perfect matchings on $2m$ labelled vertices is
\[
  (2m-1)(2m-3)\cdots3\cdot1=(2m-1)!!.
\]
For $m=3$ this is $15>8=2^3$.  Passing from $m$ to $m+1$ multiplies the
left-hand side by $2m+1$ and the right-hand side by $2$, so the strict
inequality persists for every $m\ge3$.

For the full final-zone budget, check the base case
\[
  13!!=135135>16384=2^{14}.
\]
Passing from $m$ to $m+1$ again multiplies the left-hand side by $2m+1$, while
$2^{2m}$ is multiplied by $4$.  Since $2m+1>4$ for $m\ge7$, the strict
inequality persists for every $m\ge7$.
\end{proof}

\begin{corollary}[Degree-only terminal-boundary pruning exceeds the full branch budget]\label{cor:bc-degree-terminal-boundary-ambient-too-large}
Work with the degree-constrained low-low ambient set of
\cref{def:bc-degree-low-low-candidates}.  The degree equations alone do not
give a full-budget terminal-boundary candidate list.  More precisely, for
every $m\ge7$ there is a residual degree vector on $M=2m$ low labels such that
any terminal-boundary list whose induced low-low assignments contain every
degree-compatible assignment has cardinality greater than $2^{2m}$.
\end{corollary}

\begin{proof}
Use the vector
\[
  d_1=d_2=\cdots=d_{2m}=1
\]
from \cref{prop:bc-degree-low-low-ambient-too-large}.  The corresponding
degree-compatible low-low assignments are exactly the perfect matchings on
\(\{1,\ldots,2m\}\), so there are
\[
  (2m-1)!!
\]
of them, and this number is greater than \(2^{2m}\) for \(m\ge7\).

By \cref{lem:bc-low-low-terminal-corner-boundary}, a terminal Step~3 boundary
segment determines every Step~3 entry on \(L_{2m}\).  Hence one terminal
boundary segment can induce at most one low-low assignment.  Therefore any
terminal-boundary list whose induced assignments contain all assignments in
the displayed degree ambient set has at least \((2m-1)!!\) elements, which is
larger than the full final-zone budget \(2^{2m}\) for \(m\ge7\).
\end{proof}

\begin{corollary}[Refined odd-packet classes induce bounded low-low candidates]\label{cor:bc-refined-omega-induces-low-low-candidates}
Work in the setting of
\cref{def:bc-degree-low-low-candidates}.  Suppose the target-window data
determine
\[
  \lambda^{(2m)},\qquad E_1,\ldots,E_m,
\]
the labelled continuation boxes, and a canonically ordered set
\[
  \Omega_{\mathrm{tw}}\subseteq
  \Omega(\lambda^{(2m)};E_1,\ldots,E_m)
\]
which contains the true deleted odd-packet assignment and satisfies
\[
  |\Omega_{\mathrm{tw}}|\le2^m .
\]
Let $\mathcal L_{\mathrm{tw}}$ be the set of candidates
$\ell\in\mathcal L_{\mathrm{deg}}$ with the following property: after
combining $\ell$ with the recovered continuation-column subfilling $J_{2m}$,
the resulting full Step~3 half-square and Step~2 symmetric matrix reconstruct
a semistandard tableau whose boxes for labels $2m+1,\ldots,j$ are the
displayed continuation boxes, whose tight-prefix endpoint and retained even
strips are the displayed $\lambda^{(2m)}$ and $E_1,\ldots,E_m$, and whose
deleted odd-packet assignment belongs to $\Omega_{\mathrm{tw}}$.  Order
$\mathcal L_{\mathrm{tw}}$ by first ordering its image in
$\Omega_{\mathrm{tw}}$ and breaking ties lexicographically as in
\cref{def:bc-degree-low-low-candidates}.  Then
$\mathcal L_{\mathrm{tw}}$ is determined by the target-window data, contains
the true low-low assignment, and satisfies
\[
  |\mathcal L_{\mathrm{tw}}|\le2^m .
\]
Consequently the hypotheses of
\cref{cor:bc-bounded-low-low-candidates-fit-branch-budget} hold for the final
right-boundary zone.
\end{corollary}

\begin{proof}
Every object used in the definition of $\mathcal L_{\mathrm{tw}}$ is
determined by the target-window data: the ambient set
$\mathcal L_{\mathrm{deg}}$ by
\cref{prop:bc-strict-left-degree-low-low-pruning}, the continuation-column
subfilling by \cref{lem:bc-strict-left-recovers-continuation-columns}, and the
displayed endpoint, retained strips, continuation boxes, and
$\Omega_{\mathrm{tw}}$ by hypothesis.  Thus the set and its order are
canonical functions of the target-window data.

The true low-low assignment lies in $\mathcal L_{\mathrm{deg}}$ by
\cref{prop:bc-strict-left-degree-low-low-pruning}.  Combining it with the true
continuation-column subfilling gives the true Step~3 half-square, hence the
true Step~2 matrix by \cref{lem:bc-step3-cells-are-step2-half-matrix}.  The
reconstructed tableau is therefore the original one.  Its displayed base data
are the true base data, and its deleted odd-packet assignment belongs to
$\Omega_{\mathrm{tw}}$ by hypothesis.  Hence the true low-low assignment
belongs to $\mathcal L_{\mathrm{tw}}$.

It remains to prove the cardinality bound.  Send
$\ell\in\mathcal L_{\mathrm{tw}}$ to the deleted odd-packet assignment of the
semistandard tableau reconstructed from the completed matrix.  This map lands
in $\Omega_{\mathrm{tw}}$ by definition.  It is injective.  Indeed, if two
candidates have the same image in $\Omega_{\mathrm{tw}}$, then they have the
same deleted odd-label boxes; by the definition of
$\mathcal L_{\mathrm{tw}}$ they also have the same retained even-label boxes
$E_1,\ldots,E_m$ and the same labelled continuation boxes.  Therefore they
determine the same labelled semistandard tableau.  By
\cref{lem:bc-step2-matrix-unique-from-tableau}, that tableau determines a
unique Step~2 symmetric matrix.  The low-low entries are entries of this
matrix, so the two candidates are equal.

Thus
\[
  |\mathcal L_{\mathrm{tw}}|\le|\Omega_{\mathrm{tw}}|\le2^m .
\]
The membership, canonical-ordering, and cardinality hypotheses of
\cref{cor:bc-bounded-low-low-candidates-fit-branch-budget} are now satisfied.
\end{proof}

\begin{corollary}[Full-budget odd-packet classes induce full-budget low-low candidates]\label{cor:bc-full-budget-omega-induces-low-low-candidates}
Work in the setting of
\cref{def:bc-degree-low-low-candidates}.  Suppose the target-window data
determine
\[
  \lambda^{(2m)},\qquad E_1,\ldots,E_m,
\]
the labelled continuation boxes, and a canonically ordered set
\[
  \Omega_{\mathrm{src}}\subseteq
  \Omega(\lambda^{(2m)};E_1,\ldots,E_m)
\]
which contains the true deleted odd-packet assignment and satisfies
\[
  |\Omega_{\mathrm{src}}|\le2^{2m}.
\]
Then the same target-window data determine a canonically ordered set
\[
  \mathcal L_{\mathrm{tw}}\subseteq\mathcal L_{\mathrm{deg}}
\]
of possible Step~3 entry assignments on \(L_{2m}\), the true low-low
assignment belongs to this set, and
\[
  |\mathcal L_{\mathrm{tw}}|\le2^{2m}.
\]
Consequently, whenever the corresponding right-boundary or graph terminal
setting is in force, this is a full-budget low-low candidate source in the
sense of
\cref{prop:bc-full-budget-low-low-candidates-supply-large-right-boundary-auxiliary}
and
\cref{cor:bc-full-budget-low-low-candidate-list-recovers-large-deleted-only-graph}.
\end{corollary}

\begin{proof}
Define \(\mathcal L_{\mathrm{tw}}\) exactly as in
\cref{cor:bc-refined-omega-induces-low-low-candidates}, with
\(\Omega_{\mathrm{src}}\) in place of \(\Omega_{\mathrm{tw}}\).  The
determinacy and true-membership parts of that proof use only that the
target-window data determine the displayed base data and the displayed
odd-packet class, and that the true odd-packet assignment belongs to the
class.  Hence the same argument shows that
\(\mathcal L_{\mathrm{tw}}\) is determined and contains the true low-low
assignment.

The same injection used there sends each candidate low-low assignment to the
deleted odd-packet assignment of the reconstructed semistandard tableau.  Its
image lies in \(\Omega_{\mathrm{src}}\), and injectivity follows again from
\cref{lem:bc-step2-matrix-unique-from-tableau}.  Therefore
\[
  |\mathcal L_{\mathrm{tw}}|
  \le|\Omega_{\mathrm{src}}|
  \le2^{2m}.
\]
This is exactly the full-budget candidate-list hypothesis in the cited
right-boundary and graph statements.
\end{proof}

\begin{corollary}[Bounded signature lists induce full-budget low-low candidates]\label{cor:bc-bounded-signature-list-low-low-candidates}
Work simultaneously in the settings of
\cref{cor:bc-bounded-signature-list-full-budget-omega} and
\cref{def:bc-degree-low-low-candidates}.  Under the hypotheses of
\cref{cor:bc-bounded-signature-list-full-budget-omega}, the target-window data
determine a canonically ordered set
\[
  \mathcal L_{\mathrm{tw}}\subseteq\mathcal L_{\mathrm{deg}}
\]
of possible Step~3 entry assignments on \(L_{2m}\), the true low-low
assignment belongs to this set, and
\[
  |\mathcal L_{\mathrm{tw}}|\le2^{2m}.
\]
\end{corollary}

\begin{proof}
\Cref{cor:bc-bounded-signature-list-full-budget-omega} constructs a
target-window-determined odd-packet class
\[
  \Omega_{\mathrm{tw}}(\Sigma_{\mathrm{tw}})
  \subseteq \Omega(\lambda^{(2m)};E_1,\ldots,E_m)
\]
which contains the true deleted odd-packet assignment and has cardinality at
most \(2^{2m}\).  Applying
\cref{cor:bc-full-budget-omega-induces-low-low-candidates} to this
odd-packet class gives the asserted low-low candidate set.
\end{proof}

\begin{corollary}[Target-window signature fibers induce bounded low-low candidates]\label{cor:bc-target-window-signature-fiber-low-low-candidates}
Work in the right-boundary auxiliary setting of
\cref{prop:bc-right-boundary-auxiliary-global-recovery}, and suppose the last
zone $Z_L=\{a_L,\ldots,j\}$ is a non-full tight height zone with
$m\ge9$ deleted singleton packets.  Let $V^\ast_L$ be the replacement window
parsed from the repaired path, and let $\mathcal D_{\mathrm{out}}(P)$ be the
outside datum of
\cref{prop:bc-right-boundary-auxiliary-global-recovery}.  Suppose the
target-window data
\[
  S,\quad Z_L,\quad V^\ast_L,\quad \mathcal D_{\mathrm{out}}(P),
  \quad \rho_{2a_L-2},\rho_{2j}
\]
determine $\lambda^{(2m)}$, the retained even-label strips
$E_1,\ldots,E_m$, the labelled continuation boxes, and the commuting
normal-form signature $\kappa$ of the fixed-even standard tableau associated
to the true odd-packet assignment.  Then the same target-window data determine
a canonical set
\[
  \mathcal L_{\mathrm{tw}}\subseteq\mathcal L_{\mathrm{deg}}
\]
of possible low-low Step~3 assignments which contains the true low-low
assignment and satisfies
\[
  |\mathcal L_{\mathrm{tw}}|\le2^m .
\]
Consequently the hypotheses of
\cref{cor:bc-bounded-low-low-candidates-fit-branch-budget} hold for the final
right-boundary zone.
\end{corollary}

\begin{proof}
\Cref{cor:bc-target-window-signature-fiber-omega} defines the target-window
signature fiber
\[
  \Omega_{\mathrm{tw}}(\kappa)
  \subseteq
  \Omega(\lambda^{(2m)};E_1,\ldots,E_m)
\]
from the displayed data and proves that it is canonical, contains the true
odd-packet assignment, and has cardinality at most $2^m$.  Apply
\cref{cor:bc-refined-omega-induces-low-low-candidates} with
$\Omega_{\mathrm{tw}}=\Omega_{\mathrm{tw}}(\kappa)$.  The resulting set
$\mathcal L_{\mathrm{tw}}$ is the required bounded low-low candidate class,
and \cref{cor:bc-bounded-low-low-candidates-fit-branch-budget} then applies.
\end{proof}

\begin{corollary}[Target-window signature fibers give the tight right-boundary inverse]\label{cor:bc-target-window-signature-fiber-right-boundary-inverse}
Work in the setting of
\cref{cor:bc-target-window-signature-fiber-low-low-candidates}, with
\[
  Z_L=\{a_L,\ldots,j\}=\{j-2m+1,\ldots,j\}.
\]
Under the hypotheses of
\cref{cor:bc-target-window-signature-fiber-low-low-candidates}, from
\[
  S,\quad Z_L,\quad V^\ast_L,\quad \mathcal D_{\mathrm{out}}(P),
  \quad \rho_{2a_L-2},\rho_{2j},
  \quad \beta_L\in\{0,1\}^{2m}
\]
one recovers the original right-boundary zone subpath
\[
  \rho_{2a_L-2},\rho_{2a_L-1},\ldots,\rho_{2j}.
\]
\end{corollary}

\begin{proof}
\Cref{cor:bc-target-window-signature-fiber-low-low-candidates} gives a
canonical low-low candidate set $\mathcal L_{\mathrm{tw}}$ containing the true
low-low assignment and satisfying
\[
  |\mathcal L_{\mathrm{tw}}|\le2^m .
\]
Applying
\cref{cor:bc-bounded-low-low-candidates-recover-right-boundary} to this set
recovers the displayed right-boundary subpath.
\end{proof}

\begin{proposition}[Target-window signatures supply the right-boundary auxiliary inverse]\label{prop:bc-target-window-signature-supplies-right-boundary-auxiliary}
Work in the non-full tight-height setting of
\cref{cor:bc-height-zero-surplus-label-prefix} and in the right-boundary
auxiliary setting of
\cref{prop:bc-right-boundary-auxiliary-global-recovery}.  Suppose the final
zone is
\[
  Z_L=\{a_L,\ldots,j\}=\{j-2m+1,\ldots,j\}
\]
with $m\ge9$ deleted singleton packets, and suppose the target-window data
\[
  S,\quad Z_L,\quad V^\ast_L,\quad \mathcal D_{\mathrm{out}}(P),
  \quad \rho_{2a_L-2},\rho_{2j}
\]
determine $\lambda^{(2m)}$, the retained even-label strips
$E_1,\ldots,E_m$, the labelled continuation boxes, and the commuting
normal-form signature $\kappa$ of the true fixed-even standard tableau.  Then
the final zone satisfies the right-boundary inverse hypothesis in
\cref{prop:bc-right-boundary-auxiliary-global-recovery}: taking
\[
  A(P)=\mathcal D_{\mathrm{out}}(P),
\]
there is a fixed inverse which recovers the original subpath on $Z_L$ from
\[
  V^\ast_L,\quad \beta_L,\quad \rho_{2a_L-2},\rho_{2j},\quad Z_L,\quad A(P),
  \qquad \beta_L\in\{0,1\}^{2|Z_L\cap I(S)|}.
\]
\end{proposition}

\begin{proof}
By \cref{cor:bc-height-zero-surplus-label-prefix}, the final tight zone has
\[
  |Z_L\cap I(S)|=m .
\]
The datum $A(P)=\mathcal D_{\mathrm{out}}(P)$ is a fixed function of
$S$ and $\mathcal D_{\mathrm{out}}(P)$, as required in
\cref{prop:bc-right-boundary-auxiliary-global-recovery}.  The inverse is
allowed to use $V^\ast_L$, the endpoint pair, the known interval $Z_L$, and
this auxiliary datum.  Together these are exactly the target-window data in
the hypothesis, and hence determine the base endpoint, retained strips,
continuation boxes, and $\kappa$.

Applying
\cref{cor:bc-target-window-signature-fiber-right-boundary-inverse} recovers
\[
  \rho_{2a_L-2},\rho_{2a_L-1},\ldots,\rho_{2j}
\]
from these data and the word $\beta_L\in\{0,1\}^{2m}$.  Since
$m=|Z_L\cap I(S)|$, this is precisely a word of the length allowed in
\cref{prop:bc-right-boundary-auxiliary-global-recovery}.  Thus the displayed
procedure is the required fixed right-boundary inverse.
\end{proof}

\begin{proposition}[Even standard-time chains supply the large right-boundary auxiliary inverse]\label{prop:bc-even-chain-supplies-large-right-boundary-auxiliary}
Work in the non-full tight-height setting of
\cref{cor:bc-height-zero-surplus-label-prefix} and in the right-boundary
auxiliary setting of
\cref{prop:bc-right-boundary-auxiliary-global-recovery}.  Suppose the final
zone is
\[
  Z_L=\{a_L,\ldots,j\}=\{j-2m+1,\ldots,j\}
\]
with $m\ge9$ deleted singleton packets.  Let $U=\operatorname{std}(T)$ be the
doubled standardization of the underlying semistandard tableau, and let
\[
  \varnothing=\sigma^{(0)}\subset\sigma^{(1)}\subset\cdots
\]
be the standard shape chain of $U$.  Suppose the target-window data
\[
  S,\quad Z_L,\quad V^\ast_L,\quad \mathcal D_{\mathrm{out}}(P),
  \quad \rho_{2a_L-2},\rho_{2j}
\]
determine the even standard-time endpoint chain
\[
  \sigma^{(0)},\sigma^{(2)},\sigma^{(4)},\ldots,\sigma^{(4m)}
\]
and the labelled continuation boxes for labels $2m+1,\ldots,j$.  Then the
final zone satisfies the right-boundary inverse hypothesis in
\cref{prop:bc-right-boundary-auxiliary-global-recovery}: taking
\[
  A(P)=\mathcal D_{\mathrm{out}}(P),
\]
there is a fixed inverse which recovers the original subpath on $Z_L$ from
\[
  V^\ast_L,\quad \beta_L,\quad \rho_{2a_L-2},\rho_{2j},\quad Z_L,\quad A(P),
  \qquad \beta_L\in\{0,1\}^{2|Z_L\cap I(S)|}.
\]
\end{proposition}

\begin{proof}
By \cref{lem:bc-standard-time-pair-deletion}, the recovered even
standard-time endpoints determine the true Pieri prefix
\[
  \lambda^{(0)},\lambda^{(1)},\ldots,\lambda^{(2m)} .
\]
Therefore the same data determine
\[
  \lambda^{(2m)}
  \quad\text{and}\quad
  E_r=\lambda^{(2r)}\setminus\lambda^{(2r-1)}
  \qquad(1\le r\le m).
\]
The proof of \cref{cor:bc-target-window-prefix-path-closes-signature} applied
to this recovered prefix chain gives the commuting normal-form signature
$\kappa$ of the true fixed-even standard tableau.

Thus the target-window data determine $\lambda^{(2m)}$, the retained strips
$E_1,\ldots,E_m$, the labelled continuation boxes, and $\kappa$.  Applying
\cref{prop:bc-target-window-signature-supplies-right-boundary-auxiliary} gives
the required fixed right-boundary auxiliary inverse.
\end{proof}

\begin{proposition}[Low-label box readouts supply the large right-boundary auxiliary inverse]\label{prop:bc-low-label-boxes-supply-large-right-boundary-auxiliary}
Work in the setting of
\cref{prop:bc-even-chain-supplies-large-right-boundary-auxiliary}.  Suppose
the target-window data determine, for each semistandard label
$1\le a\le2m$, the two boxes carrying label $a$, and also determine the
labelled continuation boxes for labels $2m+1,\ldots,j$.  Then the final zone
satisfies the right-boundary inverse hypothesis in
\cref{prop:bc-right-boundary-auxiliary-global-recovery} with
$A(P)=\mathcal D_{\mathrm{out}}(P)$.
\end{proposition}

\begin{proof}
By \cref{lem:bc-low-label-boxes-standard-readout} with $M=2m$, the recovered
boxes carrying labels $1,\ldots,2m$ determine the even standard-time endpoint
chain
\[
  \sigma^{(0)},\sigma^{(2)},\sigma^{(4)},\ldots,\sigma^{(4m)} .
\]
Together with the recovered labelled continuation boxes, this is exactly the
hypothesis of
\cref{prop:bc-even-chain-supplies-large-right-boundary-auxiliary}.  Applying
that proposition gives the claimed right-boundary auxiliary inverse.
\end{proof}

\begin{corollary}[Signature-route readouts recover the large deleted-only graph]\label{cor:bc-signature-route-recovers-large-deleted-only-graph}
Work in the setting of
\cref{cor:bc-full-budget-low-low-candidate-list-recovers-large-deleted-only-graph},
and suppose \(m\ge9\).  Suppose the target-window data
\[
  S,\quad Z_L,\quad V^\ast_L,\quad \mathcal D_{\mathrm{out}}(P),
  \quad \rho_{2a_L-2},\rho_{2j}
\]
determine \(\lambda^{(2m)}\), the retained even-label strips
\(E_1,\ldots,E_m\), and the labelled continuation boxes.  Suppose in addition
that the same data determine one of the following signature-route readouts:
the commuting normal-form signature \(\kappa\) of the true fixed-even standard
tableau; the even standard-time endpoint chain
\[
  \sigma^{(0)},\sigma^{(2)},\sigma^{(4)},\ldots,\sigma^{(4m)}
\]
of the doubled standardization; or, for every semistandard label
\(1\le a\le2m\), the two boxes carrying label \(a\).
Then \(\Gamma\) is recovered from
\[
  S,\quad \Gamma^\ast,\quad \rho(v)\in\{0,1\}^q,\quad
  \iota(A)\in\{0,1\}^q,
\]
the recovered outside datum, these target-window data, and the already
budgeted final-zone word \(\beta_L\in\{0,1\}^{2m}\).
\end{corollary}

\begin{proof}
First suppose the target-window data determine \(\kappa\).  Then
\cref{cor:bc-target-window-signature-fiber-low-low-candidates} constructs a
canonical low-low candidate set \(\mathcal L_{\mathrm{tw}}\) containing the
true low-low assignment and satisfying
\[
  |\mathcal L_{\mathrm{tw}}|\le2^m\le2^{2m}.
\]
Applying
\cref{cor:bc-full-budget-low-low-candidate-list-recovers-large-deleted-only-graph}
to this candidate set recovers \(\Gamma\) from the displayed graph data and
the same final-zone word.

If the target-window data determine the even standard-time endpoint chain,
\cref{cor:bc-target-window-standard-even-chain-closes-signature} shows that
the hypothesis of \cref{prop:bc-target-window-signature-closes-refined-suffix}
holds.  That hypothesis is precisely deterministic recovery of the true
commuting normal-form signature from the displayed target-window tuple.  This
reduces the even-chain case to the preceding paragraph.

Finally suppose the target-window data determine the two boxes carrying each
label \(1\le a\le2m\).  By
\cref{lem:bc-low-label-boxes-standard-readout}, with \(M=2m\), these boxes
determine the even standard-time endpoint chain.  The previous paragraph then
reduces this case to the direct-signature case.
\end{proof}

\begin{proposition}[Low-low Step-3 entries supply the large right-boundary auxiliary inverse]\label{prop:bc-low-low-entries-supply-large-right-boundary-auxiliary}
Work in the setting of
\cref{prop:bc-low-label-boxes-supply-large-right-boundary-auxiliary}.  Let
$L_{2m}$ be the low-low cell set of
\cref{lem:bc-low-label-incidence-split}.  Suppose the target-window data
\[
  S,\quad Z_L,\quad V^\ast_L,\quad \mathcal D_{\mathrm{out}}(P),
  \quad \rho_{2a_L-2},\rho_{2j}
\]
determine the Step~3 cell entries on every cell of $L_{2m}$.  Then the final
zone satisfies the right-boundary inverse hypothesis in
\cref{prop:bc-right-boundary-auxiliary-global-recovery} with
$A(P)=\mathcal D_{\mathrm{out}}(P)$.
\end{proposition}

\begin{proof}
Put $M=2m$.  By \cref{lem:bc-strict-left-recovers-continuation-columns}, the
outside datum $\mathcal D_{\mathrm{out}}(P)$ determines every Step~3 entry in
the continuation-column subarrangement $J_M$.  The hypothesis determines every
entry on $L_M$.  These two sets cover the full Step~3 half-square: an
off-diagonal unordered label pair either has both labels at most $M$, or is
incident to a continuation label $M+1,\ldots,j$.

Thus the displayed target-window data determine the whole Step~3 half-square.
By \cref{lem:bc-step3-cells-are-step2-half-matrix}, this determines every
off-diagonal entry of the Step~2 symmetric matrix, and the diagonal entries
are zero.  Applying \cref{lem:bc-step2-q-prefix-strip-determines-prefix} with
$M=j$ recovers the complete Pieri boundary path
\[
  \lambda^{(0)},\lambda^{(1)},\ldots,\lambda^{(j)}.
\]
Then \cref{lem:bc-step2-knuth-prefix-low-label-boxes} with $M=2m$ recovers
the two boxes carrying every label $1,\ldots,2m$, and
\cref{lem:bc-step2-knuth-continuation-boxes} with $M=2m$ recovers the two
boxes carrying every continuation label $2m+1,\ldots,j$.  Applying
\cref{prop:bc-low-label-boxes-supply-large-right-boundary-auxiliary} gives the
claimed right-boundary auxiliary inverse.
\end{proof}

\begin{proposition}[Terminal boundary supplies the large right-boundary auxiliary inverse]\label{prop:bc-terminal-boundary-supplies-large-right-boundary-auxiliary}
Work in the setting of
\cref{prop:bc-low-low-entries-supply-large-right-boundary-auxiliary}.  Suppose
the target-window data
\[
  S,\quad Z_L,\quad V^\ast_L,\quad \mathcal D_{\mathrm{out}}(P),
  \quad \rho_{2a_L-2},\rho_{2j}
\]
determine the terminal low-label Step~3 boundary segment
\[
  \rho_{2(j-2m)},\rho_{2(j-2m)+1},\ldots,\rho_{2j}.
\]
Then the final zone satisfies the right-boundary inverse hypothesis in
\cref{prop:bc-right-boundary-auxiliary-global-recovery} with
$A(P)=\mathcal D_{\mathrm{out}}(P)$.
\end{proposition}

\begin{proof}
By \cref{lem:bc-low-low-terminal-corner-boundary} with $M=2m$, the displayed
terminal boundary segment determines every Step~3 cell entry on $L_{2m}$.
Thus the hypothesis of
\cref{prop:bc-low-low-entries-supply-large-right-boundary-auxiliary} is
satisfied, and that proposition gives the required fixed right-boundary
auxiliary inverse.
\end{proof}

\begin{proposition}[Full branch-budget terminal boundaries supply the large right-boundary auxiliary inverse]\label{prop:bc-full-budget-terminal-boundaries-supply-large-right-boundary-auxiliary}
Work in the non-full tight-height setting of
\cref{cor:bc-height-zero-surplus-label-prefix} and in the right-boundary
auxiliary setting of
\cref{prop:bc-right-boundary-auxiliary-global-recovery}.  Suppose the final
zone is
\[
  Z_L=\{a_L,\ldots,j\}=\{j-2m+1,\ldots,j\}
\]
with $m\ge9$ deleted singleton packets.  Suppose the target-window data
\[
  S,\quad Z_L,\quad V^\ast_L,\quad \mathcal D_{\mathrm{out}}(P),
  \quad \rho_{2a_L-2},\rho_{2j}
\]
determine a canonically ordered finite set
\[
  \mathcal B_{\mathrm{tw}}
\]
of possible terminal low-label Step~3 boundary segments
\[
  \tau=(\tau_{2(j-2m)},\tau_{2(j-2m)+1},\ldots,\tau_{2j}),
\]
that the true segment
\[
  (\rho_{2(j-2m)},\rho_{2(j-2m)+1},\ldots,\rho_{2j})
\]
belongs to this set, and that
\[
  |\mathcal B_{\mathrm{tw}}|\le2^{2m}.
\]
Then the final zone satisfies the right-boundary inverse hypothesis in
\cref{prop:bc-right-boundary-auxiliary-global-recovery}: taking
\[
  A(P)=\mathcal D_{\mathrm{out}}(P),
\]
there is a fixed inverse which recovers the original subpath on $Z_L$ from
\[
  V^\ast_L,\quad \beta_L,\quad \rho_{2a_L-2},\rho_{2j},\quad Z_L,\quad A(P),
  \qquad \beta_L\in\{0,1\}^{2|Z_L\cap I(S)|}.
\]
\end{proposition}

\begin{proof}
By \cref{cor:bc-height-zero-surplus-label-prefix},
\[
  |Z_L\cap I(S)|=m .
\]
Use the $2m$ bits of $\beta_L$ to encode the rank of the true terminal
boundary segment in the canonical order on $\mathcal B_{\mathrm{tw}}$.  Since
$|\mathcal B_{\mathrm{tw}}|\le2^{2m}$, this selects the true segment.

The identity $a_L=j-2m+1$ gives
\[
  2a_L-2=2(j-2m).
\]
Thus the selected segment is exactly the original right-boundary zone subpath
\[
  \rho_{2a_L-2},\rho_{2a_L-1},\ldots,\rho_{2j}.
\]
With $A(P)=\mathcal D_{\mathrm{out}}(P)$, this is the required fixed
right-boundary inverse.
\end{proof}

\begin{corollary}[Full branch-budget terminal boundaries induce low-low candidates]\label{cor:bc-full-budget-terminal-boundaries-induce-low-low-candidates}
Work in the non-full tight-height setting of
\cref{cor:bc-height-zero-surplus-label-prefix}.  Suppose the final zone is
\[
  Z_L=\{a_L,\ldots,j\}=\{j-2m+1,\ldots,j\}
\]
with $m\ge9$ deleted singleton packets, and let $L_{2m}$ be the low-low cell
set of \cref{lem:bc-low-label-incidence-split}.  Suppose the target-window data
\[
  S,\quad Z_L,\quad V^\ast_L,\quad \mathcal D_{\mathrm{out}}(P),
  \quad \rho_{2a_L-2},\rho_{2j}
\]
determine a canonically ordered finite set
\[
  \mathcal B_{\mathrm{tw}}
\]
of possible terminal low-label Step~3 boundary segments
\[
  \tau=(\tau_{2(j-2m)},\tau_{2(j-2m)+1},\ldots,\tau_{2j}),
\]
that the true segment
\[
  (\rho_{2(j-2m)},\rho_{2(j-2m)+1},\ldots,\rho_{2j})
\]
belongs to this set, and that
\[
  |\mathcal B_{\mathrm{tw}}|\le2^{2m}.
\]
Then the same target-window data determine a canonically ordered finite set
\[
  \mathcal L_{\mathrm{tw}}
\]
of possible Step~3 entry assignments on $L_{2m}$, the true low-low assignment
belongs to this set, and
\[
  |\mathcal L_{\mathrm{tw}}|\le2^{2m}.
\]
\end{corollary}

\begin{proof}
For each boundary segment in $\mathcal B_{\mathrm{tw}}$, apply the deterministic
backward sweep of
\cref{cor:bc-parsed-dual-rsk-prime-boundary-recovers-subfilling} to the
terminal Ferrers subarrangement $L_{2m}$ identified by
\cref{lem:bc-low-low-terminal-corner-boundary}.  This gives a canonically
ordered image set
\[
  \mathcal L_{\mathrm{tw}},
\]
ordered by the first boundary segment which produces the image, with duplicate
images removed.  Thus $\mathcal L_{\mathrm{tw}}$ is determined by the same
target-window data and has cardinality at most
$|\mathcal B_{\mathrm{tw}}|\le2^{2m}$.

Since the true terminal boundary segment belongs to $\mathcal B_{\mathrm{tw}}$,
\cref{lem:bc-low-low-terminal-corner-boundary}, applied with $M=2m$, shows that
its image under the same backward sweep is exactly the true Step~3 assignment
on $L_{2m}$.
Therefore the true low-low assignment belongs to $\mathcal L_{\mathrm{tw}}$.
\end{proof}

\begin{proposition}[Bounded terminal boundaries supply the large right-boundary auxiliary inverse]\label{prop:bc-bounded-terminal-boundaries-supply-large-right-boundary-auxiliary}
Work in the non-full tight-height setting of
\cref{cor:bc-height-zero-surplus-label-prefix} and in the right-boundary
auxiliary setting of
\cref{prop:bc-right-boundary-auxiliary-global-recovery}.  Suppose the final
zone is
\[
  Z_L=\{a_L,\ldots,j\}=\{j-2m+1,\ldots,j\}
\]
with $m\ge9$ deleted singleton packets.  Suppose the target-window data
\[
  S,\quad Z_L,\quad V^\ast_L,\quad \mathcal D_{\mathrm{out}}(P),
  \quad \rho_{2a_L-2},\rho_{2j}
\]
determine a canonically ordered finite set
\[
  \mathcal B_{\mathrm{tw}}
\]
of possible terminal low-label Step~3 boundary segments
\[
  \tau=(\tau_{2(j-2m)},\tau_{2(j-2m)+1},\ldots,\tau_{2j}),
\]
that the true segment
\[
  (\rho_{2(j-2m)},\rho_{2(j-2m)+1},\ldots,\rho_{2j})
\]
belongs to this set, and that
\[
  |\mathcal B_{\mathrm{tw}}|\le2^m .
\]
Then the final zone satisfies the right-boundary inverse hypothesis in
\cref{prop:bc-right-boundary-auxiliary-global-recovery}: taking
\[
  A(P)=\mathcal D_{\mathrm{out}}(P),
\]
there is a fixed inverse which recovers the original subpath on $Z_L$ from
\[
  V^\ast_L,\quad \beta_L,\quad \rho_{2a_L-2},\rho_{2j},\quad Z_L,\quad A(P),
  \qquad \beta_L\in\{0,1\}^{2|Z_L\cap I(S)|}.
\]
\end{proposition}

\begin{proof}
For each boundary segment in $\mathcal B_{\mathrm{tw}}$, apply the deterministic
backward sweep of
\cref{cor:bc-parsed-dual-rsk-prime-boundary-recovers-subfilling} to the
terminal Ferrers subarrangement $L_{2m}$ identified in
\cref{lem:bc-low-low-terminal-corner-boundary}.  This gives a canonically
ordered image set
\[
  \mathcal L_{\mathrm{tw}}
\]
of Step~3 entry assignments on $L_{2m}$, ordered by first occurrence from
$\mathcal B_{\mathrm{tw}}$ and with duplicate images removed.

The set $\mathcal L_{\mathrm{tw}}$ is determined by the target-window data.
It has cardinality at most $|\mathcal B_{\mathrm{tw}}|\le2^m$.  Since the true
terminal boundary segment belongs to $\mathcal B_{\mathrm{tw}}$,
\cref{lem:bc-low-low-terminal-corner-boundary} shows that the true low-low
Step~3 assignment belongs to $\mathcal L_{\mathrm{tw}}$.  Therefore the
hypotheses of
\cref{prop:bc-bounded-low-low-candidates-supply-large-right-boundary-auxiliary}
are satisfied, and that proposition gives the required fixed right-boundary
auxiliary inverse with $A(P)=\mathcal D_{\mathrm{out}}(P)$.
\end{proof}

\begin{proposition}[Bounded low-low candidates supply the large right-boundary auxiliary inverse]\label{prop:bc-bounded-low-low-candidates-supply-large-right-boundary-auxiliary}
Work in the non-full tight-height setting of
\cref{cor:bc-height-zero-surplus-label-prefix} and in the right-boundary
auxiliary setting of
\cref{prop:bc-right-boundary-auxiliary-global-recovery}.  Suppose the final
zone is
\[
  Z_L=\{a_L,\ldots,j\}=\{j-2m+1,\ldots,j\}
\]
with $m\ge9$ deleted singleton packets.  Let $L_{2m}$ be the low-low cell set
of \cref{lem:bc-low-label-incidence-split}.  Suppose the target-window data
\[
  S,\quad Z_L,\quad V^\ast_L,\quad \mathcal D_{\mathrm{out}}(P),
  \quad \rho_{2a_L-2},\rho_{2j}
\]
determine a canonically ordered finite set
\[
  \mathcal L_{\mathrm{tw}}
\]
of possible Step~3 entry assignments on $L_{2m}$, that the true assignment
belongs to this set, and that
\[
  |\mathcal L_{\mathrm{tw}}|\le2^m .
\]
Then the final zone satisfies the right-boundary inverse hypothesis in
\cref{prop:bc-right-boundary-auxiliary-global-recovery}: taking
\[
  A(P)=\mathcal D_{\mathrm{out}}(P),
\]
there is a fixed inverse which recovers the original subpath on $Z_L$ from
\[
  V^\ast_L,\quad \beta_L,\quad \rho_{2a_L-2},\rho_{2j},\quad Z_L,\quad A(P),
  \qquad \beta_L\in\{0,1\}^{2|Z_L\cap I(S)|}.
\]
\end{proposition}

\begin{proof}
The displayed hypotheses are exactly the bounded low-low candidate hypotheses
of \cref{cor:bc-bounded-low-low-candidates-fit-branch-budget}.  Therefore
\cref{cor:bc-bounded-low-low-candidates-recover-right-boundary} applies to the
canonically ordered set $\mathcal L_{\mathrm{tw}}$ and recovers
\[
  \rho_{2a_L-2},\rho_{2a_L-1},\ldots,\rho_{2j}
\]
from the target-window data and the word $\beta_L\in\{0,1\}^{2m}$.
By \cref{cor:bc-height-zero-surplus-label-prefix},
\[
  |Z_L\cap I(S)|=m ,
\]
so this word has exactly the length allowed by
\cref{prop:bc-right-boundary-auxiliary-global-recovery}.  With
$A(P)=\mathcal D_{\mathrm{out}}(P)$, the displayed recovery procedure is the
required fixed right-boundary inverse.
\end{proof}

\begin{proposition}[Full branch-budget low-low candidates supply the large right-boundary auxiliary inverse]\label{prop:bc-full-budget-low-low-candidates-supply-large-right-boundary-auxiliary}
Work in the non-full tight-height setting of
\cref{cor:bc-height-zero-surplus-label-prefix} and in the right-boundary
auxiliary setting of
\cref{prop:bc-right-boundary-auxiliary-global-recovery}.  Suppose the final
zone is
\[
  Z_L=\{a_L,\ldots,j\}=\{j-2m+1,\ldots,j\}
\]
with $m\ge9$ deleted singleton packets.  Let $L_{2m}$ be the low-low cell set
of \cref{lem:bc-low-label-incidence-split}.  Suppose the target-window data
\[
  S,\quad Z_L,\quad V^\ast_L,\quad \mathcal D_{\mathrm{out}}(P),
  \quad \rho_{2a_L-2},\rho_{2j}
\]
determine a canonically ordered finite set
\[
  \mathcal L_{\mathrm{tw}}
\]
of possible Step~3 entry assignments on $L_{2m}$, that the true assignment
belongs to this set, and that
\[
  |\mathcal L_{\mathrm{tw}}|\le2^{2m}.
\]
Then the final zone satisfies the right-boundary inverse hypothesis in
\cref{prop:bc-right-boundary-auxiliary-global-recovery}: taking
\[
  A(P)=\mathcal D_{\mathrm{out}}(P),
\]
there is a fixed inverse which recovers the original subpath on $Z_L$ from
\[
  V^\ast_L,\quad \beta_L,\quad \rho_{2a_L-2},\rho_{2j},\quad Z_L,\quad A(P),
  \qquad \beta_L\in\{0,1\}^{2|Z_L\cap I(S)|}.
\]
\end{proposition}

\begin{proof}
By \cref{cor:bc-height-zero-surplus-label-prefix},
\[
  |Z_L\cap I(S)|=m .
\]
Use the $2m$ bits of $\beta_L$ to encode the rank of the true low-low
assignment in the canonical order on $\mathcal L_{\mathrm{tw}}$.  Since
$|\mathcal L_{\mathrm{tw}}|\le2^{2m}$, this selects the true Step~3 entries on
$L_{2m}$.

Put $M=2m$.  By \cref{lem:bc-strict-left-recovers-continuation-columns},
$\mathcal D_{\mathrm{out}}(P)$ determines every Step~3 entry in the
continuation-column subarrangement $J_M$.  The selected entries determine the
entries on $L_M$.  The cell sets $L_M$ and $J_M$ cover the full Step~3
half-square: an unordered off-diagonal label pair either has both labels at
most $M$, or is incident to a continuation label $M+1,\ldots,j$.

Therefore the displayed data and $\beta_L$ determine every cell entry of the
full Step~3 half-square.  By
\cref{lem:bc-step3-full-half-square-determines-rho}, these entries determine
the whole boundary sequence
\[
  \rho_0,\rho_1,\ldots,\rho_{2j}.
\]
Since $a_L=j-2m+1$, the desired right-boundary zone subpath is the terminal
segment
\[
  \rho_{2a_L-2},\rho_{2a_L-1},\ldots,\rho_{2j}.
\]
With $A(P)=\mathcal D_{\mathrm{out}}(P)$, the displayed recovery procedure is
the required fixed right-boundary inverse.
\end{proof}

\begin{proposition}[Refined odd-packet classes supply the large right-boundary auxiliary inverse]\label{prop:bc-refined-omega-supplies-large-right-boundary-auxiliary}
Work in the non-full tight-height setting of
\cref{cor:bc-height-zero-surplus-label-prefix} and in the right-boundary
auxiliary setting of
\cref{prop:bc-right-boundary-auxiliary-global-recovery}.  Suppose the final
zone is
\[
  Z_L=\{a_L,\ldots,j\}=\{j-2m+1,\ldots,j\}
\]
with $m\ge9$ deleted singleton packets.  Suppose the target-window data
\[
  S,\quad Z_L,\quad V^\ast_L,\quad \mathcal D_{\mathrm{out}}(P),
  \quad \rho_{2a_L-2},\rho_{2j}
\]
determine $\lambda^{(2m)}$, the retained even-label strips
$E_1,\ldots,E_m$, the labelled continuation boxes, and a canonically ordered
set
\[
  \Omega_{\mathrm{tw}}\subseteq
  \Omega(\lambda^{(2m)};E_1,\ldots,E_m)
\]
which contains the true deleted odd-packet assignment and satisfies
\[
  |\Omega_{\mathrm{tw}}|\le2^m .
\]
Then the final zone satisfies the right-boundary inverse hypothesis in
\cref{prop:bc-right-boundary-auxiliary-global-recovery}: taking
\[
  A(P)=\mathcal D_{\mathrm{out}}(P),
\]
there is a fixed inverse which recovers the original subpath on $Z_L$ from
\[
  V^\ast_L,\quad \beta_L,\quad \rho_{2a_L-2},\rho_{2j},\quad Z_L,\quad A(P),
  \qquad \beta_L\in\{0,1\}^{2|Z_L\cap I(S)|}.
\]
\end{proposition}

\begin{proof}
The hypotheses are exactly the inputs needed for
\cref{cor:bc-refined-omega-induces-low-low-candidates}.  That corollary
constructs from the displayed target-window data a canonically ordered
low-low candidate set
\[
  \mathcal L_{\mathrm{tw}}
\]
containing the true low-low assignment and satisfying
$|\mathcal L_{\mathrm{tw}}|\le2^m$.  Applying
\cref{prop:bc-bounded-low-low-candidates-supply-large-right-boundary-auxiliary}
to this set gives the required fixed right-boundary auxiliary inverse.
\end{proof}

\begin{proposition}[Full-budget odd-packet classes supply the large right-boundary auxiliary inverse]\label{prop:bc-full-budget-omega-supplies-large-right-boundary-auxiliary}
Work in the non-full tight-height setting of
\cref{cor:bc-height-zero-surplus-label-prefix} and in the right-boundary
auxiliary setting of
\cref{prop:bc-right-boundary-auxiliary-global-recovery}.  Suppose the final
zone is
\[
  Z_L=\{a_L,\ldots,j\}=\{j-2m+1,\ldots,j\}
\]
with $m\ge9$ deleted singleton packets.  Suppose the target-window data
\[
  S,\quad Z_L,\quad V^\ast_L,\quad \mathcal D_{\mathrm{out}}(P),
  \quad \rho_{2a_L-2},\rho_{2j}
\]
determine $\lambda^{(2m)}$, the retained even-label strips
$E_1,\ldots,E_m$, the labelled continuation boxes, and a canonically ordered
set
\[
  \Omega_{\mathrm{tw}}\subseteq
  \Omega(\lambda^{(2m)};E_1,\ldots,E_m)
\]
which contains the true deleted odd-packet assignment and satisfies
\[
  |\Omega_{\mathrm{tw}}|\le2^{2m}.
\]
Then the final zone satisfies the right-boundary inverse hypothesis in
\cref{prop:bc-right-boundary-auxiliary-global-recovery}: taking
\[
  A(P)=\mathcal D_{\mathrm{out}}(P),
\]
there is a fixed inverse which recovers the original subpath on $Z_L$ from
\[
  V^\ast_L,\quad \beta_L,\quad \rho_{2a_L-2},\rho_{2j},\quad Z_L,\quad A(P),
  \qquad \beta_L\in\{0,1\}^{2|Z_L\cap I(S)|}.
\]
\end{proposition}

\begin{proof}
Define $\mathcal L_{\mathrm{tw}}\subseteq\mathcal L_{\mathrm{deg}}$ exactly
as in \cref{cor:bc-refined-omega-induces-low-low-candidates}, using the
displayed set $\Omega_{\mathrm{tw}}$.  The proof of that corollary shows that
$\mathcal L_{\mathrm{tw}}$ is determined by the target-window data, contains
the true low-low assignment, and injects into $\Omega_{\mathrm{tw}}$.  Hence
\[
  |\mathcal L_{\mathrm{tw}}|\le |\Omega_{\mathrm{tw}}|\le2^{2m}.
\]
Applying
\cref{prop:bc-full-budget-low-low-candidates-supply-large-right-boundary-auxiliary}
to this low-low candidate set gives the required fixed right-boundary
auxiliary inverse.
\end{proof}

\begin{proposition}[Bounded signature lists supply the large right-boundary auxiliary inverse]\label{prop:bc-bounded-signature-list-supplies-large-right-boundary-auxiliary}
Work in the setting of
\cref{prop:bc-full-budget-omega-supplies-large-right-boundary-auxiliary}.
Suppose the target-window data determine \(\lambda^{(2m)}\), the retained
even-label strips \(E_1,\ldots,E_m\), the labelled continuation boxes, and a
canonically ordered finite set
\[
  \Sigma_{\mathrm{tw}}
\]
of commuting normal-form signature words such that the true fixed-even standard
tableau has signature in \(\Sigma_{\mathrm{tw}}\) and
\[
  |\Sigma_{\mathrm{tw}}|\le2^m .
\]
Then the final zone satisfies the right-boundary inverse hypothesis in
\cref{prop:bc-right-boundary-auxiliary-global-recovery} with
\(A(P)=\mathcal D_{\mathrm{out}}(P)\).
\end{proposition}

\begin{proof}
\Cref{cor:bc-bounded-signature-list-full-budget-omega} constructs from the
displayed target-window data a canonically ordered set
\[
  \Omega_{\mathrm{tw}}(\Sigma_{\mathrm{tw}})
  \subseteq \Omega(\lambda^{(2m)};E_1,\ldots,E_m)
\]
which contains the true deleted odd-packet assignment and has cardinality at
most \(2^{2m}\).  Therefore the hypotheses of
\cref{prop:bc-full-budget-omega-supplies-large-right-boundary-auxiliary} hold
with \(\Omega_{\mathrm{tw}}=\Omega_{\mathrm{tw}}(\Sigma_{\mathrm{tw}})\), and
that proposition gives the required fixed right-boundary auxiliary inverse.
\end{proof}

\begin{proposition}[Nine-to-twelve broad odd-packet classes supply the large right-boundary auxiliary inverse]\label{prop:bc-nine-to-twelve-broad-omega-supplies-large-right-boundary-auxiliary}
Work in the setting of
\cref{prop:bc-full-budget-omega-supplies-large-right-boundary-auxiliary}, and
suppose
\[
  9\le m\le12 .
\]
Suppose the target-window data determine $\lambda^{(2m)}$, the retained
even-label strips $E_1,\ldots,E_m$, and the labelled continuation boxes.
Then the final zone satisfies the right-boundary inverse hypothesis in
\cref{prop:bc-right-boundary-auxiliary-global-recovery} with
$A(P)=\mathcal D_{\mathrm{out}}(P)$.
\end{proposition}

\begin{proof}
Let
\[
  \Omega_{\mathrm{tw}}=\Omega(\lambda^{(2m)};E_1,\ldots,E_m)
\]
with the lexicographic order of \cref{def:bc-tight-odd-packet-omega}.  The
target-window data determine this set, and
\cref{lem:bc-true-tight-packets-in-omega} puts the true deleted odd-packet
assignment in it.

By \cref{lem:bc-binary-prefix-factorial-bound},
\[
  |\Omega_{\mathrm{tw}}|\le 2(m-2)! .
\]
For $9\le m\le12$ one has
\[
  2(m-2)!\le 2^{2m};
\]
explicitly the left side is, for $m=9,10,11,12$,
\[
  10080,\quad80640,\quad725760,\quad7257600,
\]
while $2^{2m}$ is respectively
\[
  262144,\quad1048576,\quad4194304,\quad16777216.
\]
Thus the hypotheses of
\cref{prop:bc-full-budget-omega-supplies-large-right-boundary-auxiliary} are
satisfied by the unrefined compatibility set $\Omega_{\mathrm{tw}}$, and that
proposition gives the stated inverse.
\end{proof}

\begin{proposition}[Thirteen-to-twenty-eight broad odd-packet classes supply the large right-boundary auxiliary inverse]\label{prop:bc-thirteen-to-twenty-eight-envelope-omega-supplies-large-right-boundary-auxiliary}
Work in the setting of
\cref{prop:bc-full-budget-omega-supplies-large-right-boundary-auxiliary}, and
suppose
\[
  13\le m\le28 .
\]
Suppose the target-window data determine $\lambda^{(2m)}$, the retained
even-label strips $E_1,\ldots,E_m$, and the labelled continuation boxes.
Then the final zone satisfies the right-boundary inverse hypothesis in
\cref{prop:bc-right-boundary-auxiliary-global-recovery} with
$A(P)=\mathcal D_{\mathrm{out}}(P)$.
\end{proposition}

\begin{proof}
Let
\[
  \Omega_{\mathrm{tw}}=\Omega(\lambda^{(2m)};E_1,\ldots,E_m)
\]
with the lexicographic order of \cref{def:bc-tight-odd-packet-omega}.  The
target-window data determine this set, and
\cref{lem:bc-true-tight-packets-in-omega} puts the true deleted odd-packet
assignment in it.

If $\Omega_{\mathrm{tw}}$ is empty there is nothing to prove.  Otherwise
\cref{lem:bc-omega-fixed-cell-completions,prop:bc-completions-equal-reduced-growth-chains,lem:bc-reduced-growth-reverse-corner-tree,lem:bc-reverse-tree-fixed-even-standard-tableaux}
identify its cardinality with the cardinality of the associated fixed-even
standard-tableau fiber
\[
  \mathcal S(\gamma;e_1,\ldots,e_{m-1}).
\]
By \cref{lem:bc-fixed-even-recursive-envelope-bounds-fibers},
\[
  |\Omega_{\mathrm{tw}}|
  =
  |\mathcal S(\gamma;e_1,\ldots,e_{m-1})|
  \le W(\gamma).
\]
Since $13\le m\le28$,
\cref{prop:bc-fixed-even-envelope-m28-certificate} gives
\[
  W(\gamma)\le2^{2m}.
\]
Thus the hypotheses of
\cref{prop:bc-full-budget-omega-supplies-large-right-boundary-auxiliary} are
satisfied by the unrefined compatibility set $\Omega_{\mathrm{tw}}$, and that
proposition gives the stated inverse.
\end{proof}

\begin{proposition}[Top-fixed envelope non-frontier broad odd-packet classes supply the large right-boundary auxiliary inverse]\label{prop:bc-envelope-nonfrontier-omega-supplies-large-right-boundary-auxiliary}
Work in the setting of
\cref{prop:bc-full-budget-omega-supplies-large-right-boundary-auxiliary}, and
suppose \(m\ge13\).  Suppose the target-window data determine
\(\lambda^{(2m)}\), the retained even-label strips
\(E_1,\ldots,E_m\), and the labelled continuation boxes.  Let
\[
  \Omega_{\mathrm{tw}}
  =
  \Omega(\lambda^{(2m)};E_1,\ldots,E_m)
\]
with the lexicographic order of \cref{def:bc-tight-odd-packet-omega}.  If
\(\Omega_{\mathrm{tw}}\ne\varnothing\), let
\((\gamma;e_1,\ldots,e_{m-1})\) be the associated fixed-even datum of
\cref{def:bc-odd-fixed-even-standard-tableaux}, and assume
\[
  (\gamma,e_{m-1})\notin\mathcal P^{\mathrm{env}}_m.
\]
Then the final zone satisfies the right-boundary inverse hypothesis in
\cref{prop:bc-right-boundary-auxiliary-global-recovery} with
\(A(P)=\mathcal D_{\mathrm{out}}(P)\).
\end{proposition}

\begin{proof}
The target-window data determine \(\Omega_{\mathrm{tw}}\), and
\cref{lem:bc-true-tight-packets-in-omega} puts the true deleted odd-packet
assignment in it.  If \(\Omega_{\mathrm{tw}}=\varnothing\), the assertion is
vacuous.  Otherwise,
\cref{lem:bc-omega-fixed-cell-completions,prop:bc-completions-equal-reduced-growth-chains,lem:bc-reduced-growth-reverse-corner-tree,lem:bc-reverse-tree-fixed-even-standard-tableaux}
identify its cardinality with
\[
  |\mathcal S(\gamma;e_1,\ldots,e_{m-1})|.
\]
By \cref{lem:bc-fixed-even-top-envelope-bounds-fibers},
\[
  |\Omega_{\mathrm{tw}}|\le W^{\mathrm{top}}(\gamma;e_{m-1}).
\]
Since \((\gamma,e_{m-1})\notin\mathcal P^{\mathrm{env}}_m\),
\cref{def:bc-fixed-even-envelope-frontier} gives
\[
  W^{\mathrm{top}}(\gamma;e_{m-1})\le2^{2m}.
\]
Thus the hypotheses of
\cref{prop:bc-full-budget-omega-supplies-large-right-boundary-auxiliary} are
satisfied by the unrefined compatibility set \(\Omega_{\mathrm{tw}}\), and
that proposition gives the stated inverse.
\end{proof}

\begin{proposition}[Top-two envelope non-frontier broad odd-packet classes supply the large right-boundary auxiliary inverse]\label{prop:bc-envelope-top-two-nonfrontier-omega-supplies-large-right-boundary-auxiliary}
Work in the setting of
\cref{prop:bc-full-budget-omega-supplies-large-right-boundary-auxiliary}, and
suppose \(m\ge13\).  Suppose the target-window data determine
\(\lambda^{(2m)}\), the retained even-label strips
\(E_1,\ldots,E_m\), and the labelled continuation boxes.  Let
\[
  \Omega_{\mathrm{tw}}
  =
  \Omega(\lambda^{(2m)};E_1,\ldots,E_m)
\]
with the lexicographic order of \cref{def:bc-tight-odd-packet-omega}.  If
\(\Omega_{\mathrm{tw}}\ne\varnothing\), let
\((\gamma;e_1,\ldots,e_{m-1})\) be the associated fixed-even datum of
\cref{def:bc-odd-fixed-even-standard-tableaux}, and assume
\[
  (\gamma,e_{m-1},e_{m-2})\notin\mathcal P^{\mathrm{env},2}_m.
\]
Then the final zone satisfies the right-boundary inverse hypothesis in
\cref{prop:bc-right-boundary-auxiliary-global-recovery} with
\(A(P)=\mathcal D_{\mathrm{out}}(P)\).
\end{proposition}

\begin{proof}
The target-window data determine \(\Omega_{\mathrm{tw}}\), and
\cref{lem:bc-true-tight-packets-in-omega} puts the true deleted odd-packet
assignment in it.  If \(\Omega_{\mathrm{tw}}=\varnothing\), the assertion is
vacuous.  Otherwise,
\cref{lem:bc-omega-fixed-cell-completions,prop:bc-completions-equal-reduced-growth-chains,lem:bc-reduced-growth-reverse-corner-tree,lem:bc-reverse-tree-fixed-even-standard-tableaux}
identify its cardinality with
\[
  |\mathcal S(\gamma;e_1,\ldots,e_{m-1})|.
\]
By \cref{lem:bc-fixed-even-top-two-envelope-bounds-fibers},
\[
  |\Omega_{\mathrm{tw}}|
  \le W^{\mathrm{top},2}(\gamma;e_{m-1},e_{m-2}).
\]
Since \((\gamma,e_{m-1},e_{m-2})\notin\mathcal P^{\mathrm{env},2}_m\),
\cref{def:bc-fixed-even-envelope-frontier} gives
\[
  W^{\mathrm{top},2}(\gamma;e_{m-1},e_{m-2})\le2^{2m}.
\]
Thus the hypotheses of
\cref{prop:bc-full-budget-omega-supplies-large-right-boundary-auxiliary} are
satisfied by the unrefined compatibility set \(\Omega_{\mathrm{tw}}\), and
that proposition gives the stated inverse.
\end{proof}

\begin{proposition}[Fixed-depth envelope non-frontier broad odd-packet classes supply the large right-boundary auxiliary inverse]\label{prop:bc-fixed-depth-nonfrontier-omega-supplies-large-right-boundary-auxiliary}
Work in the setting of
\cref{prop:bc-full-budget-omega-supplies-large-right-boundary-auxiliary}.  Let
\(d\ge3\), and suppose \(m\ge d+1\).  Suppose the target-window data determine
\(\lambda^{(2m)}\), the retained even-label strips
\(E_1,\ldots,E_m\), and the labelled continuation boxes.  Let
\[
  \Omega_{\mathrm{tw}}
  =
  \Omega(\lambda^{(2m)};E_1,\ldots,E_m)
\]
with the lexicographic order of \cref{def:bc-tight-odd-packet-omega}.  If
\(\Omega_{\mathrm{tw}}\ne\varnothing\), let
\((\gamma;e_1,\ldots,e_{m-1})\) be the associated fixed-even datum of
\cref{def:bc-odd-fixed-even-standard-tableaux}, and assume
\[
  (\gamma,e_{m-1},e_{m-2},\ldots,e_{m-d})
  \notin\mathcal P^{\mathrm{env},d}_m.
\]
Then the final zone satisfies the right-boundary inverse hypothesis in
\cref{prop:bc-right-boundary-auxiliary-global-recovery} with
\(A(P)=\mathcal D_{\mathrm{out}}(P)\).
\end{proposition}

\begin{proof}
The target-window data determine \(\Omega_{\mathrm{tw}}\), and
\cref{lem:bc-true-tight-packets-in-omega} puts the true deleted odd-packet
assignment in it.  If \(\Omega_{\mathrm{tw}}=\varnothing\), the assertion is
vacuous.  Otherwise,
\cref{lem:bc-omega-fixed-cell-completions,prop:bc-completions-equal-reduced-growth-chains,lem:bc-reduced-growth-reverse-corner-tree,lem:bc-reverse-tree-fixed-even-standard-tableaux}
identify its cardinality with
\[
  |\mathcal S(\gamma;e_1,\ldots,e_{m-1})|.
\]
If \(d=3\), then \cref{lem:bc-fixed-even-top-three-envelope-bounds-fibers}
gives
\[
  |\Omega_{\mathrm{tw}}|
  \le W^{\mathrm{top},3}(\gamma;e_{m-1},e_{m-2},e_{m-3}).
\]
If \(d\ge4\), then
\cref{lem:bc-fixed-even-fixed-depth-envelope-bounds-fibers} gives
\[
  |\Omega_{\mathrm{tw}}|
  \le
  W^{\mathrm{top},d}(\gamma;e_{m-1},e_{m-2},\ldots,e_{m-d}).
\]
In both cases, the displayed non-frontier hypothesis and
\cref{def:bc-fixed-even-envelope-frontier} give
\[
  |\Omega_{\mathrm{tw}}|\le2^{2m}.
\]
Thus the hypotheses of
\cref{prop:bc-full-budget-omega-supplies-large-right-boundary-auxiliary} are
satisfied by the unrefined compatibility set \(\Omega_{\mathrm{tw}}\), and
that proposition gives the stated inverse.
\end{proof}

\begin{proposition}[Twenty-nine top-three envelope broad odd-packet classes supply the large right-boundary auxiliary inverse]\label{prop:bc-twenty-nine-top-three-envelope-omega-supplies-large-right-boundary-auxiliary}
Work in the setting of
\cref{prop:bc-full-budget-omega-supplies-large-right-boundary-auxiliary}.
Suppose the target-window data determine \(\lambda^{(58)}\), the retained
even-label strips \(E_1,\ldots,E_{29}\), and the labelled continuation boxes.
Then the final zone satisfies the right-boundary inverse hypothesis in
\cref{prop:bc-right-boundary-auxiliary-global-recovery} with
\(A(P)=\mathcal D_{\mathrm{out}}(P)\).
\end{proposition}

\begin{proof}
Let
\[
  \Omega_{\mathrm{tw}}=\Omega(\lambda^{(58)};E_1,\ldots,E_{29})
\]
with the lexicographic order of \cref{def:bc-tight-odd-packet-omega}.  The
target-window data determine this set, and
\cref{lem:bc-true-tight-packets-in-omega} puts the true deleted odd-packet
assignment in it.  If \(\Omega_{\mathrm{tw}}=\varnothing\), the assertion is
vacuous.

Assume \(\Omega_{\mathrm{tw}}\ne\varnothing\), and let
\((\gamma;e_1,\ldots,e_{28})\) be the associated fixed-even datum.  As above,
\cref{lem:bc-omega-fixed-cell-completions,prop:bc-completions-equal-reduced-growth-chains,lem:bc-reduced-growth-reverse-corner-tree,lem:bc-reverse-tree-fixed-even-standard-tableaux}
identify its cardinality with
\[
  |\mathcal S(\gamma;e_1,\ldots,e_{28})|.
\]
By \cref{lem:bc-fixed-even-top-three-envelope-bounds-fibers} and
\cref{prop:bc-fixed-even-envelope-m29-obstruction},
\[
  |\Omega_{\mathrm{tw}}|
  \le W^{\mathrm{top},3}(\gamma;e_{28},e_{27},e_{26})
  \le 2^{58}.
\]
Thus the hypotheses of
\cref{prop:bc-full-budget-omega-supplies-large-right-boundary-auxiliary} are
satisfied by the unrefined compatibility set \(\Omega_{\mathrm{tw}}\), and
that proposition gives the stated inverse.
\end{proof}

\begin{proposition}[Thirty fixed-depth envelope broad odd-packet classes supply the large right-boundary auxiliary inverse]\label{prop:bc-thirty-fixed-depth-envelope-omega-supplies-large-right-boundary-auxiliary}
Work in the setting of
\cref{prop:bc-full-budget-omega-supplies-large-right-boundary-auxiliary}.
Suppose the target-window data determine \(\lambda^{(60)}\), the retained
even-label strips \(E_1,\ldots,E_{30}\), and the labelled continuation boxes.
Then the final zone satisfies the right-boundary inverse hypothesis in
\cref{prop:bc-right-boundary-auxiliary-global-recovery} with
\(A(P)=\mathcal D_{\mathrm{out}}(P)\).
\end{proposition}

\begin{proof}
Let
\[
  \Omega_{\mathrm{tw}}=\Omega(\lambda^{(60)};E_1,\ldots,E_{30})
\]
with the lexicographic order of \cref{def:bc-tight-odd-packet-omega}.  The
target-window data determine this set, and
\cref{lem:bc-true-tight-packets-in-omega} puts the true deleted odd-packet
assignment in it.  If \(\Omega_{\mathrm{tw}}=\varnothing\), the assertion is
vacuous.

Assume \(\Omega_{\mathrm{tw}}\ne\varnothing\), and let
\((\gamma;e_1,\ldots,e_{29})\) be the associated fixed-even datum.  As above,
\cref{lem:bc-omega-fixed-cell-completions,prop:bc-completions-equal-reduced-growth-chains,lem:bc-reduced-growth-reverse-corner-tree,lem:bc-reverse-tree-fixed-even-standard-tableaux}
identify its cardinality with
\[
  |\mathcal S(\gamma;e_1,\ldots,e_{29})|.
\]
By \cref{lem:bc-fixed-even-fixed-depth-envelope-bounds-fibers} and
\cref{prop:bc-fixed-even-envelope-m30-frontier-certificate},
\[
  |\Omega_{\mathrm{tw}}|
  \le W^{\mathrm{top},7}
  (\gamma;e_{29},e_{28},e_{27},e_{26},e_{25},e_{24},e_{23})
  \le 2^{60}.
\]
Thus the hypotheses of
\cref{prop:bc-full-budget-omega-supplies-large-right-boundary-auxiliary} are
satisfied by the unrefined compatibility set \(\Omega_{\mathrm{tw}}\), and
that proposition gives the stated inverse.
\end{proof}

\begin{proposition}[Thirty-one fixed-depth envelope broad odd-packet classes supply the large right-boundary auxiliary inverse]\label{prop:bc-thirty-one-fixed-depth-envelope-omega-supplies-large-right-boundary-auxiliary}
Work in the setting of
\cref{prop:bc-full-budget-omega-supplies-large-right-boundary-auxiliary}.
Suppose the target-window data determine \(\lambda^{(62)}\), the retained
even-label strips \(E_1,\ldots,E_{31}\), and the labelled continuation boxes.
Then the final zone satisfies the right-boundary inverse hypothesis in
\cref{prop:bc-right-boundary-auxiliary-global-recovery} with
\(A(P)=\mathcal D_{\mathrm{out}}(P)\).
\end{proposition}

\begin{proof}
Let
\[
  \Omega_{\mathrm{tw}}=\Omega(\lambda^{(62)};E_1,\ldots,E_{31})
\]
with the lexicographic order of \cref{def:bc-tight-odd-packet-omega}.  The
target-window data determine this set, and
\cref{lem:bc-true-tight-packets-in-omega} puts the true deleted odd-packet
assignment in it.  If \(\Omega_{\mathrm{tw}}=\varnothing\), the assertion is
vacuous.

Assume \(\Omega_{\mathrm{tw}}\ne\varnothing\), and let
\((\gamma;e_1,\ldots,e_{30})\) be the associated fixed-even datum.  As above,
\cref{lem:bc-omega-fixed-cell-completions,prop:bc-completions-equal-reduced-growth-chains,lem:bc-reduced-growth-reverse-corner-tree,lem:bc-reverse-tree-fixed-even-standard-tableaux}
identify its cardinality with
\[
  |\mathcal S(\gamma;e_1,\ldots,e_{30})|.
\]
By \cref{lem:bc-fixed-even-fixed-depth-envelope-bounds-fibers} and
\cref{prop:bc-fixed-even-envelope-m31-frontier-certificate},
\[
  |\Omega_{\mathrm{tw}}|
  \le W^{\mathrm{top},10}
  (\gamma;e_{30},e_{29},e_{28},e_{27},e_{26},
  e_{25},e_{24},e_{23},e_{22},e_{21})
  \le 2^{62}.
\]
Thus the hypotheses of
\cref{prop:bc-full-budget-omega-supplies-large-right-boundary-auxiliary} are
satisfied by the unrefined compatibility set \(\Omega_{\mathrm{tw}}\), and
that proposition gives the stated inverse.
\end{proof}

\begin{lemma}[The binary-prefix full-budget certificate stops at twelve]\label{lem:bc-binary-prefix-full-budget-stops-at-twelve}
For every integer $m\ge13$,
\[
  2(m-2)!>2^{2m}.
\]
Consequently, the binary-prefix factorial estimate of
\cref{lem:bc-binary-prefix-factorial-bound} by itself does not certify the
full-branch budget inequality
$|\Omega(\lambda^{(2m)};E_1,\ldots,E_m)|\le2^{2m}$ in the range $m\ge13$.
\end{lemma}

\begin{proof}
At $m=13$,
\[
  2(m-2)!=2\cdot 11! = 79833600
  > 67108864 = 2^{26}=2^{2m}.
\]
If
\[
  2(m-2)!>2^{2m}
\]
holds for some $m\ge13$, then
\[
  \frac{2(m-1)!}{2^{2m+2}}
  =
  \frac{m-1}{4}\cdot\frac{2(m-2)!}{2^{2m}}
  >
  \frac{m-1}{4}
  \ge 3
  >1 .
\]
Thus the displayed strict inequality propagates from $m$ to $m+1$, and
induction proves it for every $m\ge13$.

\Cref{lem:bc-binary-prefix-factorial-bound} gives only the upper estimate
$|\Omega|\le2(m-2)!$.  Since this upper estimate is already larger than the
full budget $2^{2m}$ for $m\ge13$, that estimate alone cannot supply the
displayed budget inequality in this range.
\end{proof}

\begin{proposition}[Strict-left low-label incidence carriers supply the large right-boundary auxiliary inverse]\label{prop:bc-strict-left-incidence-carrier-supplies-large-right-boundary-auxiliary}
Work in the non-full tight-height setting of
\cref{lem:bc-tight-outside-recovers-left-growth-segment} and in the
right-boundary auxiliary setting of
\cref{prop:bc-right-boundary-auxiliary-global-recovery}.  Suppose the final
zone is
\[
  Z_L=\{a_L,\ldots,j\}=\{j-2m+1,\ldots,j\}
\]
with $m\ge9$ deleted singleton packets.  Suppose that, from
\[
  S,\quad Z_L,\quad V^\ast_L,\quad
  \rho_{2a_L-2},\rho_{2j},
  \quad
  \rho_0,\rho_1,\ldots,\rho_{2(j-2m)},
\]
a deterministic rule identifies a Ferrers subarrangement $E$ of the Step~3
dual-RSK' half-square of
\cref{lem:bc-krattenthaler-semistandard-path-boundary}, identifies the
partition labels on the top/right boundary of $E$, and verifies
\[
  H_{2m}\subseteq E
\]
for the low-label incidence cell set of
\cref{lem:bc-step3-low-label-incidence-covers-q-prefix}.  Then the final zone
satisfies the right-boundary inverse hypothesis in
\cref{prop:bc-right-boundary-auxiliary-global-recovery}: taking
\[
  A(P)=\mathcal D_{\mathrm{out}}(P),
\]
there is a fixed inverse which recovers the original subpath on $Z_L$ from
\[
  V^\ast_L,\quad \beta_L,\quad \rho_{2a_L-2},\rho_{2j},\quad Z_L,\quad A(P),
  \qquad \beta_L\in\{0,1\}^{2|Z_L\cap I(S)|}.
\]
\end{proposition}

\begin{proof}
By \cref{lem:bc-tight-outside-recovers-left-growth-segment}, the outside datum
$\mathcal D_{\mathrm{out}}(P)$ determines
\[
  \rho_0,\rho_1,\ldots,\rho_{2(j-2m)}.
\]
Therefore the full target-window tuple determines every input to the
deterministic rule in the hypothesis.  It follows that the tuple determines
the carrier $E$, its top/right boundary labels, and the inclusion
$H_{2m}\subseteq E$.

Since the top/right boundary labels of $E$ are determined,
\cref{cor:bc-parsed-dual-rsk-prime-boundary-recovers-subfilling} recovers
every Step~3 cell entry on $E$.  The inclusion $H_{2m}\subseteq E$ therefore
gives every entry on $H_{2m}$.  By \cref{lem:bc-low-label-incidence-split},
\[
  H_{2m}=X_{2m}\sqcup L_{2m},
\]
so the entries on the low-low cell set $L_{2m}$ are determined.  Applying
\cref{prop:bc-low-low-entries-supply-large-right-boundary-auxiliary} gives the
required fixed right-boundary auxiliary inverse with
$A(P)=\mathcal D_{\mathrm{out}}(P)$.
\end{proof}

\begin{proposition}[Strict-left low-low carriers supply the large right-boundary auxiliary inverse]\label{prop:bc-strict-left-low-low-carrier-supplies-large-right-boundary-auxiliary}
Work in the non-full tight-height setting of
\cref{lem:bc-tight-outside-recovers-left-growth-segment} and in the
right-boundary auxiliary setting of
\cref{prop:bc-right-boundary-auxiliary-global-recovery}.  Suppose the final
zone is
\[
  Z_L=\{a_L,\ldots,j\}=\{j-2m+1,\ldots,j\}
\]
with $m\ge9$ deleted singleton packets.  Suppose that, from
\[
  S,\quad Z_L,\quad V^\ast_L,\quad
  \rho_{2a_L-2},\rho_{2j},
  \quad
  \rho_0,\rho_1,\ldots,\rho_{2(j-2m)},
\]
a deterministic rule identifies a Ferrers subarrangement $E$ of the Step~3
dual-RSK' half-square of
\cref{lem:bc-krattenthaler-semistandard-path-boundary}, identifies the
partition labels on the top/right boundary of $E$, and verifies
\[
  L_{2m}\subseteq E
\]
for the low-low cell set of \cref{lem:bc-low-label-incidence-split}.  Then the
final zone satisfies the right-boundary inverse hypothesis in
\cref{prop:bc-right-boundary-auxiliary-global-recovery}: taking
\[
  A(P)=\mathcal D_{\mathrm{out}}(P),
\]
there is a fixed inverse which recovers the original subpath on $Z_L$ from
\[
  V^\ast_L,\quad \beta_L,\quad \rho_{2a_L-2},\rho_{2j},\quad Z_L,\quad A(P),
  \qquad \beta_L\in\{0,1\}^{2|Z_L\cap I(S)|}.
\]
\end{proposition}

\begin{proof}
By \cref{lem:bc-tight-outside-recovers-left-growth-segment}, the outside datum
$\mathcal D_{\mathrm{out}}(P)$ determines
\[
  \rho_0,\rho_1,\ldots,\rho_{2(j-2m)}.
\]
Therefore the full target-window tuple determines every input to the
deterministic rule in the hypothesis.  It follows that the tuple determines
the carrier $E$, its top/right boundary labels, and the inclusion
$L_{2m}\subseteq E$.

Since the top/right boundary labels of $E$ are determined,
\cref{cor:bc-parsed-dual-rsk-prime-boundary-recovers-subfilling} recovers
every Step~3 cell entry on $E$.  The inclusion $L_{2m}\subseteq E$ therefore
gives every entry on $L_{2m}$.  Applying
\cref{prop:bc-low-low-entries-supply-large-right-boundary-auxiliary} gives the
required fixed right-boundary auxiliary inverse with
$A(P)=\mathcal D_{\mathrm{out}}(P)$.
\end{proof}

\begin{proposition}[Strict-left incidence-cover carriers supply the large right-boundary auxiliary inverse]\label{prop:bc-strict-left-incidence-cover-carrier-supplies-large-right-boundary-auxiliary}
Work in the non-full tight-height setting of
\cref{lem:bc-tight-outside-recovers-left-growth-segment} and in the
right-boundary auxiliary setting of
\cref{prop:bc-right-boundary-auxiliary-global-recovery}.  Suppose the final
zone is
\[
  Z_L=\{a_L,\ldots,j\}=\{j-2m+1,\ldots,j\}
\]
with \(m\ge29\) deleted singleton packets.  Let \(A(m)\) and \(B(m)\) be the
frontier functions of
\cref{prop:bc-active-twenty-low-low-profile-full-budget}.  Let \(d_a\) be the
residual low-low degrees of \cref{def:bc-degree-low-low-candidates}.  Suppose
that, from
\[
  S,\quad Z_L,\quad V^\ast_L,\quad
  \rho_{2a_L-2},\rho_{2j},
  \quad
  \rho_0,\rho_1,\ldots,\rho_{2(j-2m)},
\]
a deterministic rule identifies a target-window-determined set
\[
  U_{\mathrm{tw}}\subseteq\{1\le a\le2m:d_a>0\},
\]
identifies a Ferrers subarrangement \(E\) of the Step~3 dual-RSK' half-square
of \cref{lem:bc-krattenthaler-semistandard-path-boundary}, identifies the
partition labels on the top/right boundary of \(E\), and verifies that \(E\)
contains every low-low cell incident to \(U_{\mathrm{tw}}\):
\[
  \{\{a,b\}:1\le a<b\le2m,\ \{a,b\}\cap U_{\mathrm{tw}}\ne\varnothing\}
  \subseteq E .
\]
Assume moreover that one positive-residual incidence-cover threshold holds.
Namely, put
\[
  r_i=\#\{1\le a\le2m:d_a=i\}\quad(i=1,2),\qquad
  D=\sum_{a=1}^{2m}d_a,\qquad
  D(U_{\mathrm{tw}})=\sum_{a\in U_{\mathrm{tw}}}d_a.
\]
Assume one of
\[
  |U_{\mathrm{tw}}|>r_1+r_2-A(m)
  \qquad\text{or}\qquad
  D(U_{\mathrm{tw}})>D-B(m).
\]
Then the final zone satisfies the right-boundary inverse hypothesis in
\cref{prop:bc-right-boundary-auxiliary-global-recovery}: taking
\[
  A(P)=\mathcal D_{\mathrm{out}}(P),
\]
there is a fixed inverse which recovers the original subpath on \(Z_L\) from
\[
  V^\ast_L,\quad \beta_L,\quad \rho_{2a_L-2},\rho_{2j},\quad Z_L,\quad A(P),
  \qquad \beta_L\in\{0,1\}^{2|Z_L\cap I(S)|}.
\]
\end{proposition}

\begin{proof}
By \cref{lem:bc-tight-outside-recovers-left-growth-segment}, the outside datum
\(\mathcal D_{\mathrm{out}}(P)\) determines
\[
  \rho_0,\rho_1,\ldots,\rho_{2(j-2m)}.
\]
Therefore the full target-window tuple determines every input to the
deterministic rule in the hypothesis.  It follows that the tuple determines
\(U_{\mathrm{tw}}\), the carrier \(E\), its top/right boundary labels, and the
displayed incidence inclusion.

Since the top/right boundary labels of \(E\) are determined,
\cref{cor:bc-parsed-dual-rsk-prime-boundary-recovers-subfilling} recovers every
Step~3 cell entry on \(E\).  The displayed inclusion therefore gives a
target-window partial low-low readout on every low-low cell incident to
\(U_{\mathrm{tw}}\).  If
\[
  |U_{\mathrm{tw}}|>r_1+r_2-A(m),
\]
then, because \(U_{\mathrm{tw}}\) contains only positive-residual labels,
\[
  \#\{a\notin U_{\mathrm{tw}}:d_a>0\}
  =
  r_1+r_2-|U_{\mathrm{tw}}|
  <A(m),
\]
and \cref{cor:bc-incidence-covered-low-labels-frontier-full-budget} supplies a
canonical full-budget low-low candidate set containing the true low-low
assignment.  If instead \(D(U_{\mathrm{tw}})>D-B(m)\), then
\cref{cor:bc-large-incidence-cover-frontier-full-budget} supplies the same
candidate set.  In all cases
\cref{prop:bc-full-budget-low-low-candidates-supply-large-right-boundary-auxiliary}
then gives the required fixed right-boundary inverse with
\(A(P)=\mathcal D_{\mathrm{out}}(P)\).
\end{proof}

\begin{proposition}[Column-even tight prefixes supply the large right-boundary auxiliary inverse]\label{prop:bc-column-even-prefix-supplies-large-right-boundary-auxiliary}
Work in the non-full tight-height setting of
\cref{cor:bc-height-zero-surplus-label-prefix} and in the right-boundary
auxiliary setting of
\cref{prop:bc-right-boundary-auxiliary-global-recovery}.  Suppose the final
zone is
\[
  Z_L=\{a_L,\ldots,j\}=\{j-2m+1,\ldots,j\}
\]
with $m\ge2$ deleted singleton packets, and let
\[
  \lambda^{(0)}\subset\lambda^{(1)}\subset\cdots\subset\lambda^{(2m)}
\]
be the true tight-prefix Pieri path.  Assume that $\lambda^{(2m)}$ is
column-even.  Choose the final-zone replacement window $V^\ast_L$ and branch
word
\[
  \beta_L\in\{0,1\}^{2m}
\]
by the fixed odd-alternating height full-zone decoder of
\cref{prop:bc-odd-alternating-height-full-zone-decoder} applied to this tight
prefix.  Then the final zone satisfies the right-boundary inverse hypothesis in
\cref{prop:bc-right-boundary-auxiliary-global-recovery}: taking
\[
  A(P)=\mathcal D_{\mathrm{out}}(P),
\]
there is a fixed inverse which recovers the original subpath on $Z_L$ from
\[
  V^\ast_L,\quad \beta_L,\quad \rho_{2a_L-2},\rho_{2j},\quad Z_L,\quad A(P).
\]
\end{proposition}

\begin{proof}
By \cref{cor:bc-tight-height-prefix-column-even-gate}, the displayed prefix is
an odd alternating height full-zone path with crossing set
\[
  \{1,3,\ldots,2m-1\}.
\]
Since $m\ge2$ and $\lambda^{(2m)}$ is column-even,
\cref{prop:bc-odd-alternating-height-full-zone-decoder} supplies a target path
$V^\ast_L$ with $m$ two-step blocks, a word $\beta_L\in\{0,1\}^{2m}$, and a
fixed decoder recovering the full prefix Pieri path from
$(V^\ast_L,\beta_L)$.

The recovered prefix path determines the boxes carrying the semistandard labels
$1,\ldots,2m$ by \cref{lem:bc-pieri-path-bijection}.  Restricting to these
labels gives the same labelled prefix tableau as in the original tableau.
Therefore \cref{lem:bc-step2-matrix-unique-from-tableau} determines the Step~2
symmetric matrix entries whose two labels are at most $2m$, and
\cref{lem:bc-step3-cells-are-step2-half-matrix} identifies these entries with
the Step~3 cell entries on the low-low cell set $L_{2m}$.

The outside datum $A(P)=\mathcal D_{\mathrm{out}}(P)$ determines every Step~3
entry in the continuation-column subarrangement by
\cref{lem:bc-strict-left-recovers-continuation-columns}.  Combining those
entries with the recovered entries on $L_{2m}$ determines every entry of the
full Step~3 half-square: an unordered off-diagonal label pair either has both
labels at most $2m$, or is incident to a continuation label.  Applying
\cref{lem:bc-step3-full-half-square-determines-rho} recovers
\[
  \rho_0,\rho_1,\ldots,\rho_{2j}.
\]
The required right-boundary subpath is its terminal segment
\[
  \rho_{2a_L-2},\rho_{2a_L-1},\ldots,\rho_{2j},
\]
because $a_L=j-2m+1$.
\end{proof}

\begin{proposition}[Target-window Hall-augmented non-column-even carriers supply the large right-boundary auxiliary inverse]\label{prop:bc-target-window-hall-augmented-low-low-carrier-supplies-large-right-boundary-auxiliary}
Work in the non-full tight-height setting of
\cref{cor:bc-tight-height-zone-endpoints-visible,lem:bc-tight-outside-recovers-left-growth-segment}
and in the right-boundary auxiliary setting of
\cref{prop:bc-right-boundary-auxiliary-global-recovery}.  Suppose the final
zone is
\[
  Z_L=\{a_L,\ldots,j\}=\{j-2m+1,\ldots,j\}
\]
with $m\ge9$ deleted singleton packets, and let
\[
  \lambda^{(0)}\subset\lambda^{(1)}\subset\cdots\subset\lambda^{(2m)}
\]
be the true tight-prefix Pieri path.  Assume that $\lambda^{(2m)}$ is not
column-even.  Suppose that the target-window data
\[
  S,\quad Z_L,\quad V^\ast_L,\quad \mathcal D_{\mathrm{out}}(P),
  \quad \rho_{2a_L-2},\rho_{2j}
\]
determine $\lambda^{(2m)}$, the paired endpoint-defect set of
\cref{lem:bc-tight-height-continuation-parity}, the continuation-incidence
sets
\[
  N(a)=\{t\in\{2m+1,\ldots,j\}: C_t\cap P_a\ne\varnothing\},
\]
where $P_a$ are the adjacent endpoint-defect pairs and $C_t$ is the
two-column set hit by continuation label $t$, as in
\cref{lem:bc-tight-continuation-label-cover,lem:bc-continuation-incidence-semistandard-local},
and the lexicographically least canonical Hall charge of
\cref{lem:bc-canonical-continuation-hall-charge}.  Suppose that, from
\[
  S,\quad Z_L,\quad V^\ast_L,\quad
  \rho_{2a_L-2},\rho_{2j},\quad
  \rho_0,\rho_1,\ldots,\rho_{2(j-2m)},
\]
together with $\lambda^{(2m)}$, the endpoint-defect pairs, the continuation
incidences, the canonical continuation Hall charge, and the Hall-charged
continuation blocks recovered by
\cref{cor:bc-direct-hall-charge-selected-blocks-recovered}, a deterministic rule
identifies a Ferrers subarrangement $E$ of the Step~3 dual-RSK' half-square of
\cref{lem:bc-krattenthaler-semistandard-path-boundary}, identifies the
partition labels on the top/right boundary of $E$, and verifies
\[
  L_{2m}\subseteq E
\]
for the low-low cell set of \cref{lem:bc-low-label-incidence-split}.  Then the
final zone satisfies the right-boundary inverse hypothesis in
\cref{prop:bc-right-boundary-auxiliary-global-recovery}: taking
\[
  A(P)=\mathcal D_{\mathrm{out}}(P),
\]
there is a fixed inverse which recovers the original subpath on $Z_L$ from
\[
  V^\ast_L,\quad \beta_L,\quad \rho_{2a_L-2},\rho_{2j},\quad Z_L,\quad A(P),
  \qquad \beta_L\in\{0,1\}^{2|Z_L\cap I(S)|}.
\]
\end{proposition}

\begin{proof}
The first recovery hypothesis makes $\lambda^{(2m)}$, the endpoint-defect
pairs, the continuation-incidence sets, and the canonical Hall charge fixed
functions of the displayed target-window data.  The outside component
$\mathcal D_{\mathrm{out}}(P)$ also determines
\[
  \rho_0,\rho_1,\ldots,\rho_{2(j-2m)}
\]
by \cref{lem:bc-tight-outside-recovers-left-growth-segment}.  Therefore the
displayed target-window data determine the Hall-charged continuation blocks by
\cref{cor:bc-direct-hall-charge-selected-blocks-recovered}, and hence
determine every input to the deterministic Hall-augmented carrier rule.  They
determine the carrier $E$, its top/right boundary labels, and the inclusion
$L_{2m}\subseteq E$.

Since the top/right boundary labels of $E$ are determined,
\cref{cor:bc-parsed-dual-rsk-prime-boundary-recovers-subfilling} recovers every
Step~3 cell entry on $E$.  The inclusion $L_{2m}\subseteq E$ therefore gives
every entry on the low-low cell set $L_{2m}$.  Applying
\cref{prop:bc-low-low-entries-supply-large-right-boundary-auxiliary} gives the
required fixed right-boundary auxiliary inverse with
$A(P)=\mathcal D_{\mathrm{out}}(P)$.
\end{proof}

\begin{proposition}[Hall-augmented incidence-cover carriers supply the large right-boundary auxiliary inverse]\label{prop:bc-target-window-hall-augmented-incidence-cover-carrier-supplies-large-right-boundary-auxiliary}
Work in the non-full tight-height setting of
\cref{cor:bc-tight-height-zone-endpoints-visible,lem:bc-tight-outside-recovers-left-growth-segment}
and in the right-boundary auxiliary setting of
\cref{prop:bc-right-boundary-auxiliary-global-recovery}.  Suppose the final
zone is
\[
  Z_L=\{a_L,\ldots,j\}=\{j-2m+1,\ldots,j\}
\]
with \(m\ge29\) deleted singleton packets, and let
\[
  \lambda^{(0)}\subset\lambda^{(1)}\subset\cdots\subset\lambda^{(2m)}
\]
be the true tight-prefix Pieri path.  Assume that \(\lambda^{(2m)}\) is not
column-even.  Suppose that the target-window data
\[
  S,\quad Z_L,\quad V^\ast_L,\quad \mathcal D_{\mathrm{out}}(P),
  \quad \rho_{2a_L-2},\rho_{2j}
\]
determine \(\lambda^{(2m)}\), the paired endpoint-defect set of
\cref{lem:bc-tight-height-continuation-parity}, the continuation-incidence
sets
\[
  N(a)=\{t\in\{2m+1,\ldots,j\}: C_t\cap P_a\ne\varnothing\},
\]
where \(P_a\) are the adjacent endpoint-defect pairs and \(C_t\) is the
two-column set hit by continuation label \(t\), as in
\cref{lem:bc-tight-continuation-label-cover,lem:bc-continuation-incidence-semistandard-local},
and the lexicographically least canonical Hall charge of
\cref{lem:bc-canonical-continuation-hall-charge}.  Let \(A(m)\) and \(B(m)\)
be the frontier functions of
\cref{prop:bc-active-twenty-low-low-profile-full-budget}.  Let \(d_a\) be the
residual low-low degrees of \cref{def:bc-degree-low-low-candidates}.  Suppose
that, from
\[
  S,\quad Z_L,\quad V^\ast_L,\quad
  \rho_{2a_L-2},\rho_{2j},\quad
  \rho_0,\rho_1,\ldots,\rho_{2(j-2m)},
\]
together with \(\lambda^{(2m)}\), the endpoint-defect pairs, the continuation
incidences, the canonical continuation Hall charge, and the Hall-charged
continuation blocks recovered by
\cref{cor:bc-direct-hall-charge-selected-blocks-recovered}, a deterministic
rule identifies a target-window-determined set
\[
  U_{\mathrm{tw}}\subseteq\{1\le a\le2m:d_a>0\},
\]
identifies a Ferrers subarrangement \(E\) of the Step~3 dual-RSK' half-square
of \cref{lem:bc-krattenthaler-semistandard-path-boundary}, identifies the
partition labels on the top/right boundary of \(E\), and verifies
\[
  \{\{a,b\}:1\le a<b\le2m,\ \{a,b\}\cap U_{\mathrm{tw}}\ne\varnothing\}
  \subseteq E .
\]
Assume moreover that one positive-residual incidence-cover threshold holds.
Namely, put
\[
  r_i=\#\{1\le a\le2m:d_a=i\}\quad(i=1,2),\qquad
  D=\sum_{a=1}^{2m}d_a,\qquad
  D(U_{\mathrm{tw}})=\sum_{a\in U_{\mathrm{tw}}}d_a.
\]
Assume one of
\[
  |U_{\mathrm{tw}}|>r_1+r_2-A(m)
  \qquad\text{or}\qquad
  D(U_{\mathrm{tw}})>D-B(m).
\]
Then the final zone satisfies the right-boundary inverse hypothesis in
\cref{prop:bc-right-boundary-auxiliary-global-recovery}: taking
\[
  A(P)=\mathcal D_{\mathrm{out}}(P),
\]
there is a fixed inverse which recovers the original subpath on \(Z_L\) from
\[
  V^\ast_L,\quad \beta_L,\quad \rho_{2a_L-2},\rho_{2j},\quad Z_L,\quad A(P),
  \qquad \beta_L\in\{0,1\}^{2|Z_L\cap I(S)|}.
\]
\end{proposition}

\begin{proof}
The first recovery hypothesis makes \(\lambda^{(2m)}\), the endpoint-defect
pairs, the continuation-incidence sets, and the canonical Hall charge fixed
functions of the displayed target-window data.  The outside component
\(\mathcal D_{\mathrm{out}}(P)\) also determines
\[
  \rho_0,\rho_1,\ldots,\rho_{2(j-2m)}
\]
by \cref{lem:bc-tight-outside-recovers-left-growth-segment}.  Therefore the
displayed target-window data determine the Hall-charged continuation blocks by
\cref{cor:bc-direct-hall-charge-selected-blocks-recovered}, and hence every
input to the deterministic Hall-augmented carrier rule.  They determine
\(U_{\mathrm{tw}}\), the carrier \(E\), its top/right boundary labels, and the
displayed incidence inclusion.

Since the top/right boundary labels of \(E\) are determined,
\cref{cor:bc-parsed-dual-rsk-prime-boundary-recovers-subfilling} recovers every
Step~3 cell entry on \(E\).  Hence the target-window data determine the true
Step~3 entries on every low-low cell incident to \(U_{\mathrm{tw}}\).  In the
positive-label-count threshold subcase, the positivity of \(U_{\mathrm{tw}}\)
gives
\[
  \#\{a\notin U_{\mathrm{tw}}:d_a>0\}
  =
  r_1+r_2-|U_{\mathrm{tw}}|
  <A(m),
\]
so \cref{cor:bc-incidence-covered-low-labels-frontier-full-budget} gives a
canonical full-budget low-low candidate set containing the true low-low
assignment.  In the positive residual-degree threshold subcase,
\cref{cor:bc-large-incidence-cover-frontier-full-budget} gives the same
candidate set.  Finally
\cref{prop:bc-full-budget-low-low-candidates-supply-large-right-boundary-auxiliary}
gives the required fixed right-boundary inverse with
\(A(P)=\mathcal D_{\mathrm{out}}(P)\).
\end{proof}

\begin{proposition}[Target-window continuation segments supply the Hall-augmented carrier inverse]\label{prop:bc-target-window-continuation-segment-hall-carrier-supplies-large-right-boundary-auxiliary}
Work in the setting of
\cref{prop:bc-target-window-hall-augmented-low-low-carrier-supplies-large-right-boundary-auxiliary}.
Suppose that the target-window data
\[
  S,\quad Z_L,\quad V^\ast_L,\quad \mathcal D_{\mathrm{out}}(P),
  \quad \rho_{2a_L-2},\rho_{2j}
\]
determine the Pieri continuation segment
\[
  \lambda^{(2m)}\subset\lambda^{(2m+1)}\subset\cdots\subset\lambda^{(j)} .
\]
Using this recovered segment, form the endpoint-defect pairs, the
continuation incidences, and the canonical continuation Hall charge of
\cref{lem:bc-pieri-continuation-determines-hall-charge}, and recover the
Hall-charged continuation blocks of
\cref{cor:bc-hall-charge-selected-blocks-recovered}.  Suppose that, from
\[
  S,\quad Z_L,\quad V^\ast_L,\quad
  \rho_{2a_L-2},\rho_{2j},\quad
  \rho_0,\rho_1,\ldots,\rho_{2(j-2m)},
\]
together with the recovered continuation segment, its canonical Hall data, and
the Hall-charged continuation blocks,
a deterministic rule identifies a Ferrers subarrangement $E$ of the Step~3
dual-RSK' half-square of
\cref{lem:bc-krattenthaler-semistandard-path-boundary}, identifies the
partition labels on the top/right boundary of $E$, and verifies
\[
  L_{2m}\subseteq E .
\]
Then the final zone satisfies the right-boundary inverse hypothesis in
\cref{prop:bc-right-boundary-auxiliary-global-recovery} with
$A(P)=\mathcal D_{\mathrm{out}}(P)$.
\end{proposition}

\begin{proof}
The first partition of the recovered continuation segment is
$\lambda^{(2m)}$.  For each continuation label $t=2m+1,\ldots,j$, the
difference
\[
  \lambda^{(t)}\setminus\lambda^{(t-1)}
\]
is exactly the pair of boxes carrying label $t$ by
\cref{lem:bc-pieri-path-bijection}.  Hence the target-window data determine
the prefix endpoint and all labelled continuation boxes.

\Cref{lem:bc-pieri-continuation-determines-hall-charge} shows that the same
segment determines the endpoint-defect set, its canonical adjacent pairing,
the continuation incidences, and the canonical Hall charge.  Moreover,
\cref{cor:bc-hall-charge-selected-blocks-recovered} shows that the
Hall-charged continuation blocks are fixed functions of this segment and the
recovered outside datum.
By \cref{lem:bc-tight-outside-recovers-left-growth-segment},
$\mathcal D_{\mathrm{out}}(P)$ determines
\[
  \rho_0,\rho_1,\ldots,\rho_{2(j-2m)} .
\]
Therefore the displayed target-window data determine every input to the
deterministic carrier rule in the present proposition.  They determine the
carrier $E$, its top/right boundary labels, and the inclusion
$L_{2m}\subseteq E$.

The top/right boundary labels of $E$ determine every Step~3 entry on $E$ by
\cref{cor:bc-parsed-dual-rsk-prime-boundary-recovers-subfilling}.  Hence every
entry on the low-low cell set $L_{2m}$ is determined.  Applying
\cref{prop:bc-low-low-entries-supply-large-right-boundary-auxiliary} gives the
claimed right-boundary inverse with $A(P)=\mathcal D_{\mathrm{out}}(P)$.
\end{proof}

\begin{corollary}[Internal Q-cuts supply the Hall-augmented carrier inverse]\label{cor:bc-target-window-q-suffix-hall-carrier-supplies-large-right-boundary-auxiliary}
Work in the setting of
\cref{prop:bc-target-window-continuation-segment-hall-carrier-supplies-large-right-boundary-auxiliary}.
Let $C_{2m}^+$ be the Step~2 $Q$-suffix strip of
\cref{lem:bc-step2-q-suffix-strip-determines-continuation}, and let $C_{2m}$
be the corresponding Step~2 $Q$-prefix strip.  Suppose that the
target-window data
\[
  S,\quad Z_L,\quad V^\ast_L,\quad \mathcal D_{\mathrm{out}}(P),
  \quad \rho_{2a_L-2},\rho_{2j}
\]
determine the ordered partition labels on the internal separating boundary
between the Step~2 $Q$-prefix strip $C_{2m}$ and the suffix strip
$C_{2m}^+$.  Use
\cref{cor:bc-strict-left-q-suffix-boundary-recovers-continuation-segment} to
recover the Pieri continuation segment and then form its canonical Hall data
and Hall-charged continuation blocks.  Suppose that, from
\[
  S,\quad Z_L,\quad V^\ast_L,\quad
  \rho_{2a_L-2},\rho_{2j},\quad
  \rho_0,\rho_1,\ldots,\rho_{2(j-2m)},
\]
together with the recovered continuation segment, its canonical Hall data, and
the Hall-charged continuation blocks, a deterministic rule identifies a Ferrers
subarrangement $E$ of the Step~3 dual-RSK' half-square of
\cref{lem:bc-krattenthaler-semistandard-path-boundary}, identifies the partition
labels on the top/right boundary of $E$, and verifies
\[
  L_{2m}\subseteq E .
\]
Then the final zone satisfies the right-boundary inverse hypothesis in
\cref{prop:bc-right-boundary-auxiliary-global-recovery} with
$A(P)=\mathcal D_{\mathrm{out}}(P)$.
\end{corollary}

\begin{proof}
\Cref{cor:bc-strict-left-q-suffix-boundary-recovers-continuation-segment}
recovers the Pieri continuation segment from the displayed target-window data.
The remaining hypotheses are then exactly the Hall-augmented carrier hypotheses
of
\cref{prop:bc-target-window-continuation-segment-hall-carrier-supplies-large-right-boundary-auxiliary}.
Applying that proposition gives the asserted right-boundary inverse.
\end{proof}

\begin{corollary}[Internal Q-cut Hall carriers are the shared terminal source]\label{cor:bc-shared-q-cut-hall-carrier-terminal-source}
Work simultaneously in the non-full tight-height right-boundary setting of
\cref{cor:bc-target-window-q-suffix-hall-carrier-supplies-large-right-boundary-auxiliary}
and in the large deleted-only graph setting of
\cref{cor:bc-q-suffix-hall-carrier-recovers-large-deleted-only-graph}.  Assume
the target-window data determine the ordered partition labels on the internal
separating boundary between the Step~2 $Q$-prefix strip $C_{2m}$ and the
suffix strip $C_{2m}^+$.  Use
\cref{cor:bc-strict-left-q-suffix-boundary-recovers-continuation-segment} to
recover the Pieri continuation segment, form its canonical Hall data, and
recover the Hall-charged continuation blocks.  Suppose that the same recovered
data and strict-left segment determine a Ferrers subarrangement $E$ of the
Step~3 dual-RSK' half-square, determine the partition labels on the top/right
boundary of $E$, and verify
\[
  L_{2m}\subseteq E .
\]
Then the same source artifact has both consequences:
\begin{enumerate}
\item the final zone satisfies the right-boundary inverse hypothesis in
\cref{prop:bc-right-boundary-auxiliary-global-recovery} with
$A(P)=\mathcal D_{\mathrm{out}}(P)$;
\item in the large-complement graph branch, the deleted-only graph is recovered
from
\[
  S,\quad \Gamma^\ast,\quad \rho(v)\in\{0,1\}^q,\quad
  \iota(A)\in\{0,1\}^q,
\]
the recovered outside datum, the recovered internal $Q$-cut, and this
Hall-augmented carrier rule.
\end{enumerate}
\end{corollary}

\begin{proof}
The first assertion is exactly
\cref{cor:bc-target-window-q-suffix-hall-carrier-supplies-large-right-boundary-auxiliary}
with the displayed internal cut and Hall-augmented carrier rule.

The second assertion is exactly
\cref{cor:bc-q-suffix-hall-carrier-recovers-large-deleted-only-graph} with the
same internal cut and the same Hall-augmented carrier rule.  Thus no separate
continuation-segment or deleted-only graph source is used.
\end{proof}

\begin{corollary}[Direct Hall-data carriers are a shared terminal source]\label{cor:bc-shared-direct-hall-carrier-terminal-source}
Work simultaneously in the non-full tight-height right-boundary setting of
\cref{prop:bc-target-window-hall-augmented-low-low-carrier-supplies-large-right-boundary-auxiliary}
and in the large deleted-only graph setting of
\cref{cor:bc-hall-augmented-low-low-carrier-recovers-large-deleted-only-graph}.
Assume that the target-window data
\[
  S,\quad Z_L,\quad V^\ast_L,\quad \mathcal D_{\mathrm{out}}(P),
  \quad \rho_{2a_L-2},\rho_{2j}
\]
determine \(\lambda^{(2m)}\), the paired endpoint-defect set, the
continuation-incidence sets \(N(a)\), and the canonical Hall charge.  Let the
Hall-charged continuation blocks be the blocks determined from this Hall charge
and the outside datum by
\cref{cor:bc-direct-hall-charge-selected-blocks-recovered}.  Suppose that, from
\[
  S,\quad Z_L,\quad V^\ast_L,\quad
  \rho_{2a_L-2},\rho_{2j},\quad
  \rho_0,\rho_1,\ldots,\rho_{2(j-2m)},
\]
together with those direct Hall data and the recovered Hall-charged continuation
blocks, the same deterministic rule determines a Ferrers subarrangement \(E\) of
the Step~3 dual-RSK' half-square, determines the partition labels on the
top/right boundary of \(E\), and verifies
\[
  L_{2m}\subseteq E .
\]
Then the same source artifact has both consequences:
\begin{enumerate}
\item the final zone satisfies the right-boundary inverse hypothesis in
\cref{prop:bc-right-boundary-auxiliary-global-recovery} with
\(A(P)=\mathcal D_{\mathrm{out}}(P)\);
\item in the large-complement graph branch, the deleted-only graph is recovered
from
\[
  S,\quad \Gamma^\ast,\quad \rho(v)\in\{0,1\}^q,\quad
  \iota(A)\in\{0,1\}^q,
\]
the recovered outside datum and this direct Hall-data carrier rule.
\end{enumerate}
\end{corollary}

\begin{proof}
The first assertion is exactly
\cref{prop:bc-target-window-hall-augmented-low-low-carrier-supplies-large-right-boundary-auxiliary}
with the displayed direct Hall-data carrier rule.  The second assertion is
exactly
\cref{cor:bc-hall-augmented-low-low-carrier-recovers-large-deleted-only-graph}
with the same direct Hall data, the same recovered Hall-charged continuation
blocks, and the same carrier rule.
\end{proof}

\begin{corollary}[Direct Hall-data incidence-cover carriers are a shared terminal source]\label{cor:bc-shared-direct-hall-incidence-cover-carrier-terminal-source}
Work simultaneously in the non-full tight-height right-boundary setting of
\cref{prop:bc-target-window-hall-augmented-incidence-cover-carrier-supplies-large-right-boundary-auxiliary}
and in the large deleted-only graph setting of
\cref{cor:bc-hall-augmented-incidence-cover-carrier-recovers-large-deleted-only-graph}.
Assume that the target-window data
\[
  S,\quad Z_L,\quad V^\ast_L,\quad \mathcal D_{\mathrm{out}}(P),
  \quad \rho_{2a_L-2},\rho_{2j}
\]
determine \(\lambda^{(2m)}\), the paired endpoint-defect set, the
continuation-incidence sets \(N(a)\), and the canonical Hall charge.  Let the
Hall-charged continuation blocks be the blocks determined from this Hall charge
and the outside datum by
\cref{cor:bc-direct-hall-charge-selected-blocks-recovered}.  Let \(d_a\) be
the residual low-low degrees from the incidence-cover carrier setting.  Suppose
that, from
\[
  S,\quad Z_L,\quad V^\ast_L,\quad
  \rho_{2a_L-2},\rho_{2j},\quad
  \rho_0,\rho_1,\ldots,\rho_{2(j-2m)},
\]
together with those direct Hall data and the recovered Hall-charged
continuation blocks, the same deterministic rule determines a
target-window-determined set
\[
  U_{\mathrm{tw}}\subseteq\{1\le a\le2m:d_a>0\},
\]
determines a Ferrers subarrangement \(E\) of the Step~3 dual-RSK' half-square,
determines the partition labels on the top/right boundary of \(E\), verifies
\[
  \{\{a,b\}:1\le a<b\le2m,\ \{a,b\}\cap U_{\mathrm{tw}}\ne\varnothing\}
  \subseteq E ,
\]
and verifies one of the positive-residual incidence-cover thresholds in
\cref{prop:bc-target-window-hall-augmented-incidence-cover-carrier-supplies-large-right-boundary-auxiliary}
for this \(U_{\mathrm{tw}}\).
Then the same source artifact has both consequences:
\begin{enumerate}
\item the final zone satisfies the right-boundary inverse hypothesis in
\cref{prop:bc-right-boundary-auxiliary-global-recovery} with
\(A(P)=\mathcal D_{\mathrm{out}}(P)\);
\item in the large-complement graph branch, the deleted-only graph is recovered
from
\[
  S,\quad \Gamma^\ast,\quad \rho(v)\in\{0,1\}^q,\quad
  \iota(A)\in\{0,1\}^q,
\]
the recovered outside datum and this direct Hall-data incidence-cover carrier
rule.
\end{enumerate}
\end{corollary}

\begin{proof}
The first assertion is exactly
\cref{prop:bc-target-window-hall-augmented-incidence-cover-carrier-supplies-large-right-boundary-auxiliary}
with the displayed direct Hall-data incidence-cover carrier rule.  The second
assertion is exactly
\cref{cor:bc-hall-augmented-incidence-cover-carrier-recovers-large-deleted-only-graph}
with the same direct Hall data, recovered Hall-charged continuation blocks,
carrier, incidence cover, and incidence-cover frontier verification.
\end{proof}

\begin{corollary}[Strict-left low-low carriers are a shared terminal source]\label{cor:bc-shared-low-low-carrier-terminal-source}
Work simultaneously in the non-full tight-height right-boundary setting of
\cref{prop:bc-strict-left-low-low-carrier-supplies-large-right-boundary-auxiliary}
and in the large deleted-only graph setting of
\cref{cor:bc-strict-left-low-low-carrier-recovers-large-deleted-only-graph}, and
in the target-window signature setting of
\cref{prop:bc-target-window-signature-closes-refined-suffix}.
Suppose that, from
\[
  S,\quad Z_L,\quad V^\ast_L,\quad
  \rho_{2a_L-2},\rho_{2j},
  \quad
  \rho_0,\rho_1,\ldots,\rho_{2(j-2m)},
\]
a deterministic rule identifies a Ferrers subarrangement $E$ of the Step~3
dual-RSK' half-square, identifies the partition labels on the top/right
boundary of $E$, and verifies
\[
  L_{2m}\subseteq E .
\]
Then the same low-low carrier source has all five consequences:
\begin{enumerate}
\item the final zone satisfies the right-boundary inverse hypothesis in
\cref{prop:bc-right-boundary-auxiliary-global-recovery} with
$A(P)=\mathcal D_{\mathrm{out}}(P)$;
\item in the large-complement graph branch, the deleted-only graph is recovered
from
\[
  S,\quad \Gamma^\ast,\quad \rho(v)\in\{0,1\}^q,\quad
  \iota(A)\in\{0,1\}^q,
\]
the recovered outside datum, the carrier $E$, its recovered top/right boundary
labels, and the verified inclusion $L_{2m}\subseteq E$;
\item the target-window data determine the terminal low-low boundary segment
\[
  \rho_{2(j-2m)},\rho_{2(j-2m)+1},\ldots,\rho_{2j};
\]
\item the target-window data determine the internal separating boundary between
the Step~2 $Q$-prefix strip $C_{2m}$ and the suffix strip $C_{2m}^+$, and hence
determine the Pieri continuation segment.
\item the target-window data satisfy the hypothesis of
\cref{prop:bc-target-window-signature-closes-refined-suffix}.
\end{enumerate}
\end{corollary}

\begin{proof}
The first assertion is
\cref{prop:bc-strict-left-low-low-carrier-supplies-large-right-boundary-auxiliary},
and the second is
\cref{cor:bc-strict-left-low-low-carrier-recovers-large-deleted-only-graph}.

The third assertion is \cref{cor:bc-low-low-carrier-forces-terminal-boundary}
with \(M=2m\).

For the fourth assertion, the recovered top/right boundary labels of $E$
determine every Step~3 cell entry on $E$ by
\cref{cor:bc-parsed-dual-rsk-prime-boundary-recovers-subfilling}; the inclusion
$L_{2m}\subseteq E$ therefore determines every entry on $L_{2m}$.  Applying
\cref{cor:bc-strict-left-low-low-determines-internal-q-cut} gives the internal
$Q$-cut and the Pieri continuation segment.

By \cref{lem:bc-tight-outside-recovers-left-growth-segment}, the outside datum
$\mathcal D_{\mathrm{out}}(P)$ determines the strict-left segment
$\rho_0,\rho_1,\ldots,\rho_{2(j-2m)}$.  Hence the full target-window tuple
determines every input to the same carrier rule, and therefore determines the
same entries on $L_{2m}$.  Thus
\cref{cor:bc-strict-left-reduces-incidence-to-low-low} gives the fifth
assertion.
\end{proof}

\begin{corollary}[Small tight compatibility sets recover the tight right-boundary subpath]\label{cor:bc-small-omega-right-boundary-inverse}
Work in the right-boundary auxiliary setting of
\cref{prop:bc-right-boundary-auxiliary-global-recovery} and the non-full
tight-height setting of \cref{cor:bc-height-zero-surplus-label-prefix}.
Suppose the final zone is
\[
  Z_L=\{a_L,\ldots,j\}=\{j-2m+1,\ldots,j\}
\]
with $1\le m\le8$ deleted singleton packets.  Suppose the target-window data
\[
  S,\quad Z_L,\quad V^\ast_L,\quad \mathcal D_{\mathrm{out}}(P),
  \quad \rho_{2a_L-2},\rho_{2j}
\]
determine $\lambda^{(2m)}$, the retained even-label strips
$E_1,\ldots,E_m$, and the labelled continuation boxes.  Then from these data
and
\[
  \beta_L\in\{0,1\}^{2m}
\]
one recovers the original right-boundary zone subpath
\[
  \rho_{2a_L-2},\rho_{2a_L-1},\ldots,\rho_{2j}.
\]
\end{corollary}

\begin{proof}
Set
\[
  \Omega_{\mathrm{tw}}
  =
  \Omega(\lambda^{(2m)};E_1,\ldots,E_m)
\]
with the lexicographic order of \cref{def:bc-tight-odd-packet-omega}.  This
set is a canonical function of the displayed target-window data.  The true
deleted odd-packet assignment belongs to it by
\cref{lem:bc-true-tight-packets-in-omega}.

It remains to record the small-zone cardinality bound.  If $m=1$, then
\cref{lem:bc-singleton-tight-omega-forced} gives
$|\Omega_{\mathrm{tw}}|\le1\le2^m$.  If $2\le m\le6$, then
\cref{lem:bc-binary-prefix-factorial-bound} gives
\[
  |\Omega_{\mathrm{tw}}|\le 2(m-2)!\le2^m .
\]
For $m=7$ and $m=8$, apply respectively
\cref{prop:bc-seven-packet-tight-rank-leaf-certificate,prop:bc-eight-packet-tight-rank-leaf-certificate}
to the displayed target-window datum as $D(P)$; the hypotheses on $D(P)$ are
exactly the displayed recovery of $\lambda^{(2m)}$, $E_1,\ldots,E_m$, and the
labelled continuation boxes.  These propositions give
$|\Omega_{\mathrm{tw}}|\le2^m$ in the two certified cases.

Thus $\Omega_{\mathrm{tw}}$ satisfies the membership, canonical-ordering, and
cardinality hypotheses needed in
\cref{cor:bc-refined-omega-induces-low-low-candidates}.  That corollary gives
a canonically ordered low-low candidate set
$\mathcal L_{\mathrm{tw}}$ containing the true low-low assignment and
satisfying
\[
  |\mathcal L_{\mathrm{tw}}|\le2^m .
\]
Applying \cref{cor:bc-bounded-low-low-candidates-recover-right-boundary} to
this candidate set recovers the displayed right-boundary subpath from the
target-window data and $\beta_L$.
\end{proof}

\begin{proposition}[Small tight compatibility sets supply the right-boundary auxiliary inverse]\label{prop:bc-small-omega-supplies-right-boundary-auxiliary}
Work in the non-full tight-height setting of
\cref{cor:bc-height-zero-surplus-label-prefix} and in the right-boundary
auxiliary setting of
\cref{prop:bc-right-boundary-auxiliary-global-recovery}.  Suppose the final
zone is
\[
  Z_L=\{a_L,\ldots,j\}=\{j-2m+1,\ldots,j\}
\]
with $1\le m\le8$ deleted singleton packets, and suppose the target-window
data
\[
  S,\quad Z_L,\quad V^\ast_L,\quad \mathcal D_{\mathrm{out}}(P),
  \quad \rho_{2a_L-2},\rho_{2j}
\]
determine $\lambda^{(2m)}$, the retained even-label strips
$E_1,\ldots,E_m$, and the labelled continuation boxes.  Then the final zone
satisfies the right-boundary inverse hypothesis in
\cref{prop:bc-right-boundary-auxiliary-global-recovery}: taking
\[
  A(P)=\mathcal D_{\mathrm{out}}(P),
\]
there is a fixed inverse which recovers the original subpath on $Z_L$ from
\[
  V^\ast_L,\quad \beta_L,\quad \rho_{2a_L-2},\rho_{2j},\quad Z_L,\quad A(P),
  \qquad \beta_L\in\{0,1\}^{2|Z_L\cap I(S)|}.
\]
\end{proposition}

\begin{proof}
By \cref{cor:bc-height-zero-surplus-label-prefix}, the final tight zone has
\[
  |Z_L\cap I(S)|=m .
\]
The datum $A(P)=\mathcal D_{\mathrm{out}}(P)$ is an allowed auxiliary input in
\cref{prop:bc-right-boundary-auxiliary-global-recovery}.  Applying
\cref{cor:bc-small-omega-right-boundary-inverse} recovers
\[
  \rho_{2a_L-2},\rho_{2a_L-1},\ldots,\rho_{2j}
\]
from the allowed data and the word $\beta_L\in\{0,1\}^{2m}$.  Since
$m=|Z_L\cap I(S)|$, this is the word length allowed by the global
right-boundary auxiliary criterion.
\end{proof}

\begin{proposition}[Split target-window readouts supply the right-boundary auxiliary inverse]\label{prop:bc-split-target-window-supplies-right-boundary-auxiliary}
Work in the non-full tight-height setting of
\cref{cor:bc-height-zero-surplus-label-prefix} and in the right-boundary
auxiliary setting of
\cref{prop:bc-right-boundary-auxiliary-global-recovery}.  Suppose the final
zone is
\[
  Z_L=\{a_L,\ldots,j\}=\{j-2m+1,\ldots,j\}
\]
with $m\ge1$ deleted singleton packets.  Suppose the target-window data
\[
  S,\quad Z_L,\quad V^\ast_L,\quad \mathcal D_{\mathrm{out}}(P),
  \quad \rho_{2a_L-2},\rho_{2j}
\]
determine $\lambda^{(2m)}$, the retained even-label strips
$E_1,\ldots,E_m$, and the labelled continuation boxes.  If \(m\ge29\), let
\[
  \Omega_{\mathrm{tw}}
  =
  \Omega(\lambda^{(2m)};E_1,\ldots,E_m)
\]
with the lexicographic order of \cref{def:bc-tight-odd-packet-omega}; when
\(\Omega_{\mathrm{tw}}\ne\varnothing\), let
\((\gamma;e_1,\ldots,e_{m-1})\) be the associated fixed-even datum of
\cref{def:bc-odd-fixed-even-standard-tableaux}.  In the \(m\ge29\) cases not
already covered by the condition
\[
  \Omega_{\mathrm{tw}}=\varnothing
  \quad\text{or}\quad
  \bigl(\Omega_{\mathrm{tw}}\ne\varnothing
  \quad\text{and}\quad
  \bigl(m=29
  \quad\text{or}\quad
  m=30
  \quad\text{or}\quad
  m=31
  \quad\text{or}\quad
  (\gamma,e_{m-1})\notin\mathcal P^{\mathrm{env}}_m
  \quad\text{or}\quad
  (\gamma,e_{m-1},e_{m-2})\notin
  \mathcal P^{\mathrm{env},2}_m
  \quad\text{or}\quad
  \exists d,\ 3\le d\le m-1,\
  (\gamma,e_{m-1},e_{m-2},\ldots,e_{m-d})
  \notin\mathcal P^{\mathrm{env},d}_m\bigr)\bigr),
\]
suppose moreover that the same target-window data determine either the
commuting normal-form signature \(\kappa\) of the true fixed-even standard
tableau, or a canonically ordered finite set
\[
  \Sigma_{\mathrm{tw}}
\]
of commuting normal-form signature words such that the true fixed-even standard
tableau has signature in \(\Sigma_{\mathrm{tw}}\) and
\[
  |\Sigma_{\mathrm{tw}}|\le2^m .
\]
Then
the final zone satisfies the right-boundary inverse hypothesis in
\cref{prop:bc-right-boundary-auxiliary-global-recovery}: taking
\[
  A(P)=\mathcal D_{\mathrm{out}}(P),
\]
there is a fixed inverse which recovers the original subpath on $Z_L$ from
\[
  V^\ast_L,\quad \beta_L,\quad \rho_{2a_L-2},\rho_{2j},\quad Z_L,\quad A(P),
  \qquad \beta_L\in\{0,1\}^{2|Z_L\cap I(S)|}.
\]
\end{proposition}

\begin{proof}
If $1\le m\le8$, the hypotheses are exactly those of
\cref{prop:bc-small-omega-supplies-right-boundary-auxiliary}; applying that
proposition gives the required fixed inverse.

If $9\le m\le12$, the hypotheses are exactly those of
\cref{prop:bc-nine-to-twelve-broad-omega-supplies-large-right-boundary-auxiliary};
applying that proposition gives the required fixed inverse.

If $13\le m\le28$, the hypotheses are exactly those of
\cref{prop:bc-thirteen-to-twenty-eight-envelope-omega-supplies-large-right-boundary-auxiliary};
applying that proposition gives the required fixed inverse.

Now assume \(m\ge29\).  The target-window data determine the unrefined
compatibility set.  If this set is empty, there is no current fiber.  If
\(m=29\), then
\cref{prop:bc-twenty-nine-top-three-envelope-omega-supplies-large-right-boundary-auxiliary}
supplies the required fixed inverse.  If \(m=30\), then
\cref{prop:bc-thirty-fixed-depth-envelope-omega-supplies-large-right-boundary-auxiliary}
supplies it.  If \(m=31\), then
\cref{prop:bc-thirty-one-fixed-depth-envelope-omega-supplies-large-right-boundary-auxiliary}
supplies it.  If \(m\ge32\) and
\((\gamma,e_{m-1})\notin\mathcal P^{\mathrm{env}}_m\), then
\cref{prop:bc-envelope-nonfrontier-omega-supplies-large-right-boundary-auxiliary}
supplies it.  If \(m\ge32\) and
\((\gamma,e_{m-1},e_{m-2})\notin\mathcal P^{\mathrm{env},2}_m\), then
\cref{prop:bc-envelope-top-two-nonfrontier-omega-supplies-large-right-boundary-auxiliary}
supplies it.  If \(m\ge32\) and some \(3\le d\le m-1\) satisfies
\[
  (\gamma,e_{m-1},e_{m-2},\ldots,e_{m-d})
  \notin\mathcal P^{\mathrm{env},d}_m,
\]
then
\cref{prop:bc-fixed-depth-nonfrontier-omega-supplies-large-right-boundary-auxiliary}
supplies it.  In the remaining case, taking \(d=m-1\) and applying
\cref{cor:bc-full-depth-envelope-exact-omega} shows that the unrefined
compatibility set itself has cardinality \(>2^{2m}\).  Thus the additional
signature or signature-list hypothesis supplies the required refinement input.
If the data determine \(\kappa\), then
\cref{prop:bc-target-window-signature-supplies-right-boundary-auxiliary}
applies.  If they determine \(\Sigma_{\mathrm{tw}}\), then
\cref{prop:bc-bounded-signature-list-supplies-large-right-boundary-auxiliary}
applies.  In both cases we get the same fixed inverse conclusion with
\(A(P)=\mathcal D_{\mathrm{out}}(P)\).
\end{proof}

\begin{proposition}[Right-boundary target-window signatures close the refined suffix]\label{prop:bc-target-window-signature-closes-refined-suffix}
Work in the right-boundary auxiliary setting of
\cref{prop:bc-right-boundary-auxiliary-global-recovery}, and suppose the last
zone $Z_L=\{a_L,\ldots,j\}$ is a non-full tight height zone with
$m\ge9$ deleted singleton packets.  Let
\[
  V^\ast_L
\]
be the replacement window parsed from the repaired path, and let
\[
  \mathcal D_{\mathrm{out}}(P)
\]
be the outside datum of
\cref{prop:bc-right-boundary-auxiliary-global-recovery}.  Suppose the
commuting normal-form signature of the fixed-even standard tableau associated
to the true odd-packet assignment is a deterministic function of
\[
  S,\quad Z_L,\quad V^\ast_L,\quad \mathcal D_{\mathrm{out}}(P),
  \quad \rho_{2a_L-2},\rho_{2j}.
\]
Then the hypothesis of
\cref{prop:bc-refined-tight-suffix-reduces-rank} holds for the final
right-boundary zone.
\end{proposition}

\begin{proof}
The replacement window is a subwindow of the already counted repaired path
$P^\ast$, and the endpoint pair is returned by the parser.  The outside datum
$\mathcal D_{\mathrm{out}}(P)$ is recovered in
\cref{prop:bc-right-boundary-auxiliary-global-recovery} from the unchanged
blocks and earlier recovered zones, using only branch words already present in
that recovery.  Thus the deterministic function in the hypothesis supplies the
commuting normal-form signature without adding a new branch word.  Applying
\cref{cor:bc-commuting-signature-closes-refined-suffix} gives the required
refined suffix hypothesis.
\end{proof}

\begin{lemma}[The broad fixed-even bound is false]\label{lem:bc-broad-fixed-even-bound-false}
The estimate
\[
  |\mathcal S(\gamma;e_1,\ldots,e_{m-1})|\le2^m
\]
does not hold for all compatible fixed-even data.  More precisely, there is a
reduced fixed-even datum with $m=26$ for which
\[
  |\mathcal S(\gamma;e_1,\ldots,e_{25})|=254531860>2^{26}.
\]
\end{lemma}

\begin{proof}
Take
\[
  \gamma=(12,10,8,6,4,3,3,2,2,1)
\]
and fix the even cells
\[
\begin{array}{c|cccccccccc}
r&1&2&3&4&5&6&7&8&9&10\\ \hline
e_r&(1,2)&(2,1)&(2,2)&(3,2)&(4,1)&(2,4)&(1,5)&(6,1)&(1,7)&(1,8)
\end{array}
\]
\[
\begin{array}{c|cccccccccc}
r&11&12&13&14&15&16&17&18&19&20\\ \hline
e_r&(5,2)&(7,1)&(1,9)&(5,3)&(2,6)&(3,5)&(4,5)&(2,7)&(2,8)&(4,6)
\end{array}
\]
\[
\begin{array}{c|ccccc}
r&21&22&23&24&25\\ \hline
e_r&(5,4)&(1,11)&(9,1)&(3,8)&(10,1).
\end{array}
\]
This datum is nonempty.  One witness fills the odd labels
$1,3,\ldots,51$ in the cells
\[
\begin{array}{c|cccccccccc}
r&1&2&3&4&5&6&7&8&9&10\\ \hline
y_r&(1,1)&(1,3)&(1,4)&(3,1)&(2,3)&(4,2)&(5,1)&(1,6)&(3,3)&(4,3)
\end{array}
\]
\[
\begin{array}{c|cccccccccc}
r&11&12&13&14&15&16&17&18&19&20\\ \hline
y_r&(3,4)&(4,4)&(8,1)&(6,2)&(2,5)&(7,2)&(1,10)&(8,2)&(3,6)&(6,3)
\end{array}
\]
\[
\begin{array}{c|cccccc}
r&21&22&23&24&25&26\\ \hline
y_r&(2,9)&(3,7)&(7,3)&(9,2)&(1,12)&(2,10).
\end{array}
\]
The sublevel diagrams are Young diagrams, so this is a standard tableau with
the displayed fixed even entries.  Adjoining the terminal singleton fixed
strip $\overline E_{26}=\{(1,13)\}$ lifts it to a reduced fixed-even datum of
\cref{def:bc-odd-fixed-even-standard-tableaux}.

The exact verifier
\[
\texttt{character\_ring\_iter/post\_m29\_bc\_fixed\_even\_counterexample.cpp}
\]
evaluates the recurrence of
\cref{lem:bc-last-pair-deletion-recurrence} for this datum.  The archived
`optimus' log
\[
\texttt{certificates/post\_m29/bc\_fixed\_even\_counterexample\_certificate.log}
\]
has SHA-256
\[
\hashsplit{69921ed52cfc8ad47d922dad9baf4fda}{d4bbc2cc95393e3f9c21c6a586166566},
\]
ends with \texttt{\_\_EXIT\_STATUS=0}, verifies the displayed witness path,
and reports
\[
  \texttt{fixed\_even\_count=254531860},\qquad
  \texttt{pow2\_budget=67108864}.
\]
Thus the displayed compatible datum has more than $2^{26}$ fixed-even
standard tableaux.
\end{proof}

\begin{definition}[Viable active odd cells]\label{def:bc-viable-active-odd-cells}
In the setting of
\cref{def:bc-odd-cell-interval-extension-poset}, let $I\subseteq U$ and
$0\le r<m$.  Call $I$ viable at depth $r$ if $|I|=r$ and there exists
$\phi\in\mathcal L(\gamma;e_1,\ldots,e_{m-1})$ such that
\[
  I=\{u\in U:\phi(u)\le r\}.
\]
For a viable $I$ of depth $r$, define the viable active set
\[
  A_{\mathrm{viab}}(I,r+1)
\]
to be the set of cells $u\in U\setminus I$ for which there exists
$\phi\in\mathcal L(\gamma;e_1,\ldots,e_{m-1})$ satisfying
\[
  I=\{v\in U:\phi(v)\le r\},
  \qquad
  \phi(u)=r+1.
\]
\end{definition}

\begin{proposition}[Binary viable branching closes the interval leaf]\label{prop:bc-binary-viable-branching-closes-interval-leaf}
Suppose that, for every datum of
\cref{def:bc-odd-cell-interval-extension-poset}, every viable set $I$ of depth
$0\le r<m$ satisfies
\[
  |A_{\mathrm{viab}}(I,r+1)|\le2 .
\]
Then every such datum satisfies
\[
  |\mathcal L(\gamma;e_1,\ldots,e_{m-1})|\le2^m .
\]
\end{proposition}

\begin{proof}
Build a rooted tree whose vertices at depth $r$ are the viable sets $I$ of
depth $r$, with an edge
\[
  I\longrightarrow I\cup\{u\}
\]
for each $u\in A_{\mathrm{viab}}(I,r+1)$.  By definition, every
$\phi\in\mathcal L(\gamma;e_1,\ldots,e_{m-1})$ determines exactly one path
from the root $\varnothing$ to a depth-$m$ leaf, namely its sublevel sets
\[
  \{u:\phi(u)\le r\}\qquad(0\le r\le m).
\]
Conversely, a depth-$m$ path records, at each depth, the unique cell assigned
the next label; since every edge is required to be viable, the resulting final
bijection is an element of $\mathcal L(\gamma;e_1,\ldots,e_{m-1})$.
Thus the leaves of this tree are in bijection with
$\mathcal L(\gamma;e_1,\ldots,e_{m-1})$.

The hypothesis says that every vertex has at most two children.  A rooted tree
of depth $m$ with branching at most two has at most $2^m$ leaves.  Therefore
$|\mathcal L(\gamma;e_1,\ldots,e_{m-1})|\le2^m$.
\end{proof}

\begin{lemma}[Binary viable branching is too strong]\label{lem:bc-binary-viable-branching-obstruction}
The hypothesis of
\cref{prop:bc-binary-viable-branching-closes-interval-leaf} is not satisfied
by all reduced fixed-even data arising from nonempty fixed-cell completion
sets.  More precisely, for $m=6$ there is such a datum and a viable depth-$3$
set $I$ for which
\[
  |A_{\mathrm{viab}}(I,4)|=3 .
\]
\end{lemma}

\begin{proof}
Consider the reduced fixed-even endpoint
\[
  \overline\nu=(5,3,3,1,1)
\]
and the reduced fixed even strips
\[
  \overline E_1=\{(1,2)\},\quad
  \overline E_2=\{(2,1)\},\quad
  \overline E_3=\{(3,1)\},
\]
\[
  \overline E_4=\{(2,3)\},\quad
  \overline E_5=\{(1,5)\},\quad
  \overline E_6=\{(3,3),(5,1)\}.
\]
This reduced datum is obtained by deleting the first column from the
tight-prefix datum with endpoint
\[
  \nu=(6,4,4,2,2,1,1,1,1,1,1)
\]
and even strips
\[
\begin{array}{c|c}
r&E_r\\ \hline
1&\{(1,3),(2,1)\}\\
2&\{(2,2),(4,1)\}\\
3&\{(3,2),(6,1)\}\\
4&\{(2,4),(8,1)\}\\
5&\{(1,6),(10,1)\}\\
6&\{(3,4),(5,2)\}.
\end{array}
\]
It is nonempty, for instance by taking odd second boxes
\[
  (1,2),(1,4),(2,3),(1,5),(4,2),(3,3)
\]
for the labels $1,3,5,7,9,11$.

After removing the terminal fixed strip one has
\[
  \gamma=\overline\nu\setminus\overline E_6=(5,3,2,1),
\]
with fixed even cells
\[
  e_1=(1,2),\quad e_2=(2,1),\quad e_3=(3,1),\quad
  e_4=(2,3),\quad e_5=(1,5).
\]
The odd-cell set is
\[
  U=\{(1,1),(1,3),(1,4),(2,2),(3,2),(4,1)\}.
\]
The interval bounds of
\cref{def:bc-odd-cell-interval-extension-poset} are
\[
\begin{array}{c|c}
u&[L(u),R(u)]\\ \hline
(1,1)&[1,1]\\
(1,3)&[2,4]\\
(1,4)&[2,5]\\
(2,2)&[3,4]\\
(3,2)&[4,6]\\
(4,1)&[4,6].
\end{array}
\]
Let
\[
  I=\{(1,1),(1,3),(2,2)\}.
\]
It is viable at depth $3$: the assignment
\[
  (1,1)\mapsto1,\quad (1,3)\mapsto2,\quad (2,2)\mapsto3,\quad
  (1,4)\mapsto4,\quad (3,2)\mapsto5,\quad (4,1)\mapsto6
\]
belongs to $\mathcal L(\gamma;e_1,\ldots,e_5)$ and has depth-$3$ sublevel set
$I$.

Moreover each of $(1,4)$, $(3,2)$, and $(4,1)$ is viable as the next cell.
Indeed, the first displayed assignment has $(1,4)$ carrying $4$; the
assignment
\[
  (1,1)\mapsto1,\quad (1,3)\mapsto2,\quad (2,2)\mapsto3,\quad
  (3,2)\mapsto4,\quad (1,4)\mapsto5,\quad (4,1)\mapsto6
\]
has $(3,2)$ carrying $4$; and the assignment
\[
  (1,1)\mapsto1,\quad (1,3)\mapsto2,\quad (2,2)\mapsto3,\quad
  (4,1)\mapsto4,\quad (1,4)\mapsto5,\quad (3,2)\mapsto6
\]
has $(4,1)$ carrying $4$.  Each assignment respects the displayed intervals
and the Young-order inequalities on $U$.  Thus
\[
  \{(1,4),(3,2),(4,1)\}\subseteq A_{\mathrm{viab}}(I,4).
\]
No other cell remains outside $I$, so this viable active set has cardinality
$3$.
\end{proof}

\begin{lemma}[Reverse corner steps are commuting or fixed-blocked]\label{lem:bc-reverse-corner-commuting-or-blocked}
Let
\[
  \beta_r\supset\beta_{r-\frac12}\supset\beta_{r-1}
\]
be the $r$th step of a reverse reduced corner path, and write
\[
  \overline E_r=\beta_r\setminus\beta_{r-\frac12},\qquad
  \{y_r\}=\beta_{r-\frac12}\setminus\beta_{r-1}.
\]
If $y_r$ is a removable corner of $\beta_r$, then the two removals commute:
\[
  \beta_r\setminus\{y_r\},\qquad
  (\beta_r\setminus\{y_r\})\setminus\overline E_r=\beta_{r-1}
\]
are Young diagrams, and $\overline E_r$ is still a horizontal strip in the
intermediate diagram.

If $y_r=(i,j)$ is not a removable corner of $\beta_r$, then at least one of
\[
  (i,j+1),\qquad (i+1,j)
\]
belongs to $\overline E_r$.  In particular, when $\overline E_r$ is a
singleton, there are at most two non-commuting possible positions for
$y_r$.
\end{lemma}

\begin{proof}
Assume first that $y_r$ is a removable corner of $\beta_r$.  Removing a
removable corner from a Young diagram leaves a Young diagram, so
$\beta_r\setminus\{y_r\}$ is a Young diagram.  Since
$y_r\in\beta_{r-\frac12}$, it is disjoint from $\overline E_r$, and hence
\[
  (\beta_r\setminus\{y_r\})\setminus\overline E_r
  =
  \beta_{r-\frac12}\setminus\{y_r\}
  =
  \beta_{r-1}.
\]
The last diagram is a Young diagram by the definition of the reverse corner
path.  The set $\overline E_r$ has at most one box in each column by
\cref{def:bc-reduced-reverse-corner-tree}; this property is unchanged by
removing the disjoint box $y_r$ first.  Thus $\overline E_r$ remains a
horizontal strip for the swapped order.

Now suppose that $y_r=(i,j)$ is not a removable corner of $\beta_r$.  Because
$y_r$ is a removable corner of $\beta_{r-\frac12}$, neither $(i,j+1)$ nor
$(i+1,j)$ belongs to $\beta_{r-\frac12}$.  Since $y_r$ is not removable in
$\beta_r$, at least one of those two adjacent boxes belongs to $\beta_r$.
The difference between $\beta_r$ and $\beta_{r-\frac12}$ is exactly
$\overline E_r$, so that adjacent box belongs to $\overline E_r$.  If
$\overline E_r$ is a singleton, the only possible non-commuting positions are
therefore the box immediately west of it and the box immediately north of it.
\end{proof}

\begin{corollary}[Branch overflow is commuting]\label{cor:bc-branch-overflow-commuting}
In the setting of \cref{lem:bc-reverse-corner-commuting-or-blocked}, assume
that $\overline E_r$ is a singleton.  Put
\[
  \beta_{r-\frac12}=\beta_r\setminus\overline E_r
\]
and let $B_r(\beta_r)$ be the set of removable corners $y$ of
$\beta_{r-\frac12}$.  Then the subset
\[
  N_r(\beta_r)=\{y\in B_r(\beta_r):y\text{ is not a removable corner of }
  \beta_r\}
\]
has cardinality at most two.  Equivalently, if a reverse state has more than
two possible odd removals after the fixed singleton removal, all but at most
two of those removals commute with the fixed removal.
\end{corollary}

\begin{proof}
Write the singleton as $\overline E_r=\{e\}$, with
$e=(i,j)$.  By \cref{lem:bc-reverse-corner-commuting-or-blocked}, every
$y\in N_r(\beta_r)$ must be one of the two adjacent positions
\[
  (i,j-1),\qquad (i-1,j),
\]
with the nonexistent row or column omitted.  Hence $|N_r(\beta_r)|\le2$.
The final sentence is the same statement together with the commuting
alternative in \cref{lem:bc-reverse-corner-commuting-or-blocked}.
\end{proof}

\begin{lemma}[The Pieri continuation segment determines the tight Hall charge]\label{lem:bc-pieri-continuation-determines-hall-charge}
In the setting of \cref{lem:bc-tight-height-continuation-parity}, suppose one
is given the Pieri continuation segment
\[
  \lambda^{(2m)}\subset\lambda^{(2m+1)}\subset\cdots\subset\lambda^{(j)} .
\]
Then this segment determines the endpoint-defect set $O_{\mathrm{pre}}$, its
canonical adjacent pairing $P_1,\ldots,P_r$, every continuation column set
$C_t$ for $2m<t\le j$, the incidence relation
\[
  C_t\cap P_a\ne\varnothing,
\]
and the canonical Hall charge of
\cref{lem:bc-canonical-continuation-hall-charge}.
\end{lemma}

\begin{proof}
The first partition in the displayed segment is $\lambda^{(2m)}$, so it
determines
\[
  O_{\mathrm{pre}}
  =
  \{c\ge1:\,(\lambda^{(2m)})'_c\equiv1\pmod2\}.
\]
Listing this finite set increasingly determines its adjacent pairing, as in
\cref{lem:bc-tight-height-continuation-parity}.  For each continuation label
$t>2m$, the two consecutive partitions
$\lambda^{(t-1)}\subset\lambda^{(t)}$ determine
\[
  C_t=\{c\ge1:\,(\lambda^{(t)})'_c-(\lambda^{(t-1)})'_c=1\}.
\]
Therefore the incidence relation $C_t\cap P_a\ne\varnothing$ is determined.
Finally, \cref{lem:bc-canonical-continuation-hall-charge} defines the Hall
charge as the lexicographically least admissible injection for this incidence
family, so that charge is determined by the same segment.
\end{proof}

\begin{lemma}[Tight continuation labels are external to the tight zone]\label{lem:bc-tight-continuation-labels-external}
In the setting of
\cref{cor:bc-height-zero-surplus-label-prefix}, write the tight height zone as
\[
  Z_\gamma=\{a_\gamma,\ldots,j\},
  \qquad a_\gamma=j-2m+1 .
\]
Let $H(S)$ and $R$ be the one-block halo and unchanged complement of
\cref{lem:bc-canonical-halo-replacement-zones}.  For a continuation label
$t\in\{2m+1,\ldots,j\}$, put
\[
  r(t)=j-t+1 .
\]
Then $r(t)\notin Z_\gamma$.  Moreover, either $r(t)\in R$, or $r(t)$ belongs
to a canonical halo zone $Z_\eta$ with $b_\eta<a_\gamma$.  Consequently, every
label selected by the injection of
\cref{lem:bc-tight-continuation-label-cover} corresponds to a block outside
the tight zone: an unchanged complement block or a block carried by a strictly
earlier halo-zone window.
\end{lemma}

\begin{proof}
For $t\ge2m+1$ one has
\[
  r(t)=j-t+1\le j-2m=a_\gamma-1,
\]
so $r(t)\notin Z_\gamma$.

If $r(t)\notin H(S)$, then $r(t)\in R$ by the definition of the complement in
\cref{lem:bc-canonical-halo-replacement-zones}.  If $r(t)\in H(S)$, it lies in
a unique maximal consecutive component $Z_\eta=\{a_\eta,\ldots,b_\eta\}$ of
$H(S)$.  This component cannot be $Z_\gamma$, because $r(t)<a_\gamma$.
The canonical halo zones are disjoint intervals, and $Z_\gamma$ begins at
$a_\gamma$, so any different component containing $r(t)$ must satisfy
$b_\eta<a_\gamma$.  Applying this to the labels in the image of the injection
from \cref{lem:bc-tight-continuation-label-cover} proves the final sentence.
\end{proof}

\begin{corollary}[Hall-charged continuation blocks are outside data]\label{cor:bc-hall-charge-selected-blocks-recovered}
Work in the non-full tight-height right-boundary setting of
\cref{lem:bc-tight-outside-recovers-left-growth-segment}.  Suppose that the
Pieri continuation segment
\[
  \lambda^{(2m)}\subset\lambda^{(2m+1)}\subset\cdots\subset\lambda^{(j)}
\]
has been recovered, and let
\[
  \phi:\{1,\ldots,r\}\hookrightarrow\{2m+1,\ldots,j\}
\]
be the canonical continuation Hall charge determined by
\cref{lem:bc-pieri-continuation-determines-hall-charge}.  For each
endpoint-defect pair $P_a$ put
\[
  b_a=j-\phi(a)+1 .
\]
Then the outside datum $\mathcal D_{\mathrm{out}}(P)$ determines the original
two-step growth block
\[
  B_{b_a}(P)
  =
  (\rho_{2b_a-2},\rho_{2b_a-1},\rho_{2b_a}).
\]
Consequently the canonically ordered list of Hall-charged continuation blocks is
a fixed function of the recovered continuation segment and the outside datum.
\end{corollary}

\begin{proof}
The recovered continuation segment determines the canonical Hall charge by
\cref{lem:bc-pieri-continuation-determines-hall-charge}.  Since
$\phi(a)\ge2m+1$, the corresponding block index satisfies
\[
  b_a=j-\phi(a)+1\le j-2m .
\]
Also $b_a\ge1$ because $\phi(a)\le j$.  Thus the three boundary labels
\[
  \rho_{2b_a-2},\qquad \rho_{2b_a-1},\qquad \rho_{2b_a}
\]
all lie in the strict-left segment
\[
  \rho_0,\rho_1,\ldots,\rho_{2(j-2m)} .
\]
By \cref{lem:bc-tight-outside-recovers-left-growth-segment}, this whole segment
is determined by $\mathcal D_{\mathrm{out}}(P)$.  Therefore every block
$B_{b_a}(P)$ selected by the Hall charge is recovered, and the final assertion
follows by listing the defect pairs in their canonical order.
\end{proof}

\begin{corollary}[Direct Hall charges select outside blocks]\label{cor:bc-direct-hall-charge-selected-blocks-recovered}
Work in the non-full tight-height right-boundary setting of
\cref{lem:bc-tight-outside-recovers-left-growth-segment}.  Suppose that the
paired endpoint-defect set and the canonical continuation Hall charge
\[
  \phi:\{1,\ldots,r\}\hookrightarrow\{2m+1,\ldots,j\}
\]
are recovered.  For each endpoint-defect pair $P_a$ put
\[
  b_a=j-\phi(a)+1 .
\]
Then the outside datum $\mathcal D_{\mathrm{out}}(P)$ determines the original
two-step growth block
\[
  B_{b_a}(P)
  =
  (\rho_{2b_a-2},\rho_{2b_a-1},\rho_{2b_a}).
\]
Consequently the canonically ordered list of Hall-charged continuation blocks
is a fixed function of the recovered Hall charge and the outside datum, and
does not require the full Pieri continuation segment.
\end{corollary}

\begin{proof}
The proof is the same block-location argument as in
\cref{cor:bc-hall-charge-selected-blocks-recovered}, but the Hall charge is
now an input rather than being computed from the continuation segment.  Since
\[
  2m+1\le\phi(a)\le j,
\]
the corresponding block index satisfies
\[
  1\le b_a=j-\phi(a)+1\le j-2m .
\]
Thus the whole block $B_{b_a}(P)$ lies in the strict-left segment
\[
  \rho_0,\rho_1,\ldots,\rho_{2(j-2m)} .
\]
By \cref{lem:bc-tight-outside-recovers-left-growth-segment}, this segment is
determined by $\mathcal D_{\mathrm{out}}(P)$.  Listing the defect pairs in
their canonical order gives the stated ordered list of blocks.
\end{proof}

\begin{lemma}[Canonical halo zones carry frontier-event packets]\label{lem:bc-halo-zones-frontier-packets}
Work in either the width or height setting of
\cref{lem:bc-frontier-events-per-joint}.  Let
\[
  J_\beta=\{a_\beta,a_\beta+1,\ldots,b_\beta\}
\]
be the maximal consecutive witness-label intervals in $S$, and let
$E_\beta$ denote the corresponding ordered frontier-event packet
$E^w_\beta$ or $E^h_\beta$.  Put
\[
  K_\beta=\{j-b_\beta+1,\ldots,j-a_\beta+1\}\subset I(S)
\]
for the corresponding deleted block interval, and let
\[
  \widehat Z_\beta
  =
  \{\max(1,j-b_\beta),\ldots,\min(j,j-a_\beta+2)\}
\]
be its one-block halo.  Let $Z_\gamma$ be the canonical halo replacement
zones of \cref{lem:bc-canonical-halo-replacement-zones}.

Then each interval $K_\beta$ is contained in a unique $Z_\gamma$, and for
each $\gamma$ the set $Z_\gamma\cap I(S)$ is the disjoint union of the
intervals $K_\beta$ contained in $Z_\gamma$.  Moreover
\[
  \sum_{K_\beta\subseteq Z_\gamma}|E_\beta|
  =
  |Z_\gamma\cap I(S)|.
\]
The assignment of witness-label packets and deleted block intervals to a halo
zone is determined by $S$; within each packet the frontier events occur in
the same order as the labels.
\end{lemma}

\begin{proof}
By \cref{lem:bc-frontier-events-per-joint}, the maximal witness-label
interval $J_\beta=\{a_\beta,\ldots,b_\beta\}$ maps under the label-to-block
conversion $s\mapsto j-s+1$ to the deleted block interval
$K_\beta=\{j-b_\beta+1,\ldots,j-a_\beta+1\}$.  The sets $K_\beta$ are exactly
the maximal consecutive components of $I(S)$, in the reversed order.

The one-block halo of $K_\beta=\{u,\ldots,v\}$ inside $\{1,\ldots,j\}$ is
$\{\max(1,u-1),\ldots,\min(j,v+1)\}$, which is precisely the displayed
$\widehat Z_\beta$.  The full halo $H(S)$ of
\cref{lem:bc-canonical-halo-replacement-zones} is the union of these
$\widehat Z_\beta$ over all maximal deleted block intervals.  Therefore each
$K_\beta\subset H(S)$ lies in a unique maximal consecutive component
$Z_\gamma$ of $H(S)$, and $Z_\gamma\cap I(S)$ is the disjoint union of the
$K_\beta$ contained in that component.

Again by \cref{lem:bc-frontier-events-per-joint},
$|E_\beta|=|J_\beta|$, and the label-to-block map is a bijection from
$J_\beta$ to $K_\beta$, so $|K_\beta|=|J_\beta|=|E_\beta|$.  Summing over
the intervals contained in a fixed $Z_\gamma$ gives the displayed equality.
The final order statement is also part of
\cref{lem:bc-frontier-events-per-joint}: the elements of each
$E_\beta$ occur in the same order as the labels in $J_\beta$.
\end{proof}

\begin{lemma}[Zone branch words split into packets]\label{lem:bc-zone-branch-word-splitting}
Work in the setting of \cref{lem:bc-halo-zones-frontier-packets}.  Fix one
canonical halo zone $Z_\gamma$, and let
\[
  J_{\beta_1},\ldots,J_{\beta_m}
\]
be the witness-label packets whose deleted block intervals
$K_{\beta_r}$ are contained in $Z_\gamma$, listed in increasing witness-label
order.  Put
\[
  h_r=|E_{\beta_r}|=|J_{\beta_r}|
  \qquad\text{and}\qquad
  e_\gamma=\sum_{r=1}^m h_r.
\]
Then the packet lengths $h_1,\ldots,h_m$ and their order are determined by
$S$.  Moreover concatenation gives a bijection
\[
  \prod_{r=1}^m\{0,1\}^{2h_r}
  \longleftrightarrow
  \{0,1\}^{2e_\gamma}.
\]
Consequently a single zone branch word in $\{0,1\}^{2e_\gamma}$ carries
exactly the same branch data as the packet-local branch words used in
\cref{lem:bc-local-branch-word,lem:bc-local-branch-recovery} for the packets
inside $Z_\gamma$.
\end{lemma}

\begin{proof}
The set $S$ determines its maximal consecutive witness-label intervals
$J_\beta$ and their increasing order.  By
\cref{lem:bc-halo-zones-frontier-packets}, $S$ also determines which of these
packets are assigned to the fixed halo zone $Z_\gamma$, and
$|E_{\beta_r}|=|J_{\beta_r}|$ for every $r$.  Hence the ordered list
$h_1,\ldots,h_m$ is determined by $S$.

Concatenation sends
\[
  (\beta_{\beta_1},\ldots,\beta_{\beta_m})
  \in \prod_{r=1}^m\{0,1\}^{2h_r}
\]
to the word obtained by writing the packet words in the displayed order.  The
length of the concatenated word is
\[
  \sum_{r=1}^m 2h_r=2e_\gamma.
\]
Conversely, because the cut lengths $2h_1,\ldots,2h_m$ are determined by
$S$, any word in $\{0,1\}^{2e_\gamma}$ is split uniquely into consecutive
subwords of those lengths.  These two operations are inverse to one another.
The final sentence is exactly this bijection applied to the packet-local
branch words of
\cref{lem:bc-local-branch-word,lem:bc-local-branch-recovery}.
\end{proof}

\begin{lemma}[Retained halo blocks expose packet endpoints]\label{lem:bc-zone-retained-blocks-expose-packet-endpoints}
Let
\[
  P=(\rho_0,\rho_1,\ldots,\rho_{2j})
\]
be a two-step growth path, let $S\subset\{1,\ldots,j\}$, put
$I=I(S)$, and let
\[
  Z_\gamma=\{a_\gamma,\ldots,b_\gamma\}
\]
be one canonical halo zone.  Write the maximal consecutive components of
$Z_\gamma\cap I$ as
\[
  K_t=\{u_t,\ldots,v_t\}\qquad(1\le t\le m).
\]
Suppose one is given $S$, the zone endpoint pair
\[
  \rho_{2a_\gamma-2},\qquad \rho_{2b_\gamma},
\]
and, for every retained halo index $r\in Z_\gamma\setminus I$, the original
two-step block
\[
  B_r(P)=(\rho_{2r-2},\rho_{2r-1},\rho_{2r}).
\]
Then, for every packet $K_t$, the packet endpoint pair
\[
  \rho_{2u_t-2},\qquad \rho_{2v_t}
\]
is determined by these data.
\end{lemma}

\begin{proof}
The set $S$ determines $I$, the canonical halo zones, and the components
$K_t=\{u_t,\ldots,v_t\}$ of $Z_\gamma\cap I$.

Fix one component $K_t$.  For its left endpoint, if $u_t=a_\gamma$, then
$\rho_{2u_t-2}=\rho_{2a_\gamma-2}$ is the given left zone endpoint.  If
$u_t>a_\gamma$, then the preceding index $u_t-1$ lies in $Z_\gamma$.  By
maximality of $K_t$ as a component of $Z_\gamma\cap I$, this index is not in
$I$, so the retained block $B_{u_t-1}(P)$ is among the given data.  Its right
endpoint is
\[
  \rho_{2(u_t-1)}=\rho_{2u_t-2},
\]
which determines the desired left packet endpoint.

For the right endpoint, if $v_t=b_\gamma$, then
$\rho_{2v_t}=\rho_{2b_\gamma}$ is the given right zone endpoint.  If
$v_t<b_\gamma$, then the following index $v_t+1$ lies in $Z_\gamma\setminus I$
by the same maximality argument.  The given retained block $B_{v_t+1}(P)$ has
left endpoint
\[
  \rho_{2(v_t+1)-2}=\rho_{2v_t},
\]
which determines the desired right packet endpoint.
\end{proof}

\begin{lemma}[Packet endpoints need only adjacent halo blocks]\label{lem:bc-packet-endpoints-adjacent-halo}
In the setting of \cref{lem:bc-zone-retained-blocks-expose-packet-endpoints},
fix one packet $K_t=\{u_t,\ldots,v_t\}$ in a canonical halo zone
$Z_\gamma=\{a_\gamma,\ldots,b_\gamma\}$.  Its endpoint pair
\[
  \rho_{2u_t-2},\qquad \rho_{2v_t}
\]
is determined by $S$, the zone endpoint pair
\[
  \rho_{2a_\gamma-2},\qquad \rho_{2b_\gamma},
\]
and only the following adjacent retained halo blocks, when they exist:
\[
  B_{u_t-1}(P)\quad\text{if }u_t>a_\gamma,
  \qquad
  B_{v_t+1}(P)\quad\text{if }v_t<b_\gamma.
\]
Consequently, a zone-local inverse that calls packet inverses does not need
all retained halo blocks merely to determine packet endpoints; it needs, for
each packet, only the retained block immediately to its left or right, unless
the corresponding side is the zone boundary.
\end{lemma}

\begin{proof}
The proof is the packetwise part of
\cref{lem:bc-zone-retained-blocks-expose-packet-endpoints}.  If
$u_t=a_\gamma$, the left packet endpoint is the left zone endpoint.  If
$u_t>a_\gamma$, then $u_t-1\in Z_\gamma\setminus I$ by maximality of
$K_t$ as a component of $Z_\gamma\cap I$, and the right endpoint of
$B_{u_t-1}(P)$ is
\[
  \rho_{2(u_t-1)}=\rho_{2u_t-2}.
\]
This determines the left endpoint.  The right endpoint is identical: if
$v_t=b_\gamma$, use the right zone endpoint, and if $v_t<b_\gamma$, then
$v_t+1\in Z_\gamma\setminus I$ and the left endpoint of $B_{v_t+1}(P)$ is
\[
  \rho_{2(v_t+1)-2}=\rho_{2v_t}.
\]
No other retained halo block enters these formulas.
\end{proof}

\begin{lemma}[All retained halo blocks are packet-adjacent]\label{lem:bc-all-retained-halo-packet-adjacent}
Let $S\subset\{1,\ldots,j\}$, put $I=I(S)$, and let
$Z_\gamma=\{a_\gamma,\ldots,b_\gamma\}$ be one canonical halo zone.  Write the
maximal consecutive components of $Z_\gamma\cap I$ as
\[
  K_t=\{u_t,\ldots,v_t\}\qquad(1\le t\le m).
\]
Then
\[
  Z_\gamma\setminus I
  =
  \bigl\{u_t-1:\ 1\le t\le m,\ a_\gamma\le u_t-1\bigr\}
  \cup
  \bigl\{v_t+1:\ 1\le t\le m,\ v_t+1\le b_\gamma\bigr\}.
\]
Consequently, if a zone-local inverse uses
\cref{lem:bc-packet-endpoints-adjacent-halo} to determine the endpoints of
every packet $K_t$, then the adjacent retained block indices it may need are
exactly all retained halo indices in $Z_\gamma\setminus I$.  Any saving in
visible retained blocks must therefore come from a merged rule that recovers
some endpoints without exposing the corresponding retained block, or from a
packet mechanism that is not called for every deleted component as a separate
packet inverse.
\end{lemma}

\begin{proof}
The inclusion from right to left is immediate: if $a_\gamma\le u_t-1$, then
$u_t-1\in Z_\gamma$ and it is not in $I$ because $K_t$ is a component of
$Z_\gamma\cap I$; similarly $v_t+1\in Z_\gamma\setminus I$ whenever
$v_t+1\le b_\gamma$.

For the reverse inclusion, let $r\in Z_\gamma\setminus I$.  Since
$Z_\gamma$ is a component of the one-block halo of $I$, there is some
$i\in I$ with $|r-i|=1$.  The index $i$ lies in a unique component
$K_t=\{u_t,\ldots,v_t\}$ of $Z_\gamma\cap I$.  If $r=i-1$, then $r\notin I$
forces $i=u_t$, so $r=u_t-1$.  If $r=i+1$, then $r\notin I$ forces
$i=v_t$, so $r=v_t+1$.  This proves the displayed equality.

The final statement is just the displayed equality read together with
\cref{lem:bc-packet-endpoints-adjacent-halo}: the left and right endpoint
lookups over all packets use precisely the retained indices appearing in the
two displayed sets, whose union is $Z_\gamma\setminus I$.
\end{proof}

\begin{proposition}[Packet inverses assemble inside a halo zone]\label{prop:bc-zone-packet-inverses-assemble}
Let
\[
  P=(\rho_0,\rho_1,\ldots,\rho_{2j})
\]
be a two-step growth path, let $S\subset\{1,\ldots,j\}$, put
$I=I(S)$, and let
\[
  Z_\gamma=\{a_\gamma,\ldots,b_\gamma\}
\]
be one canonical halo zone.  Write the maximal consecutive components of
$Z_\gamma\cap I$ in increasing block order as
\[
  K_t=\{u_t,\ldots,v_t\}\qquad(1\le t\le m).
\]
Suppose a zone-local datum consists of a replacement output $V^\ast_\gamma$,
a zone branch word $\beta_\gamma\in\{0,1\}^{2e_\gamma}$, the endpoint pair
$\rho_{2a_\gamma-2},\rho_{2b_\gamma}$, and the interval $Z_\gamma$, with the
following properties.
\begin{enumerate}
\item It recovers the original retained halo blocks
\[
  B_r(P)=(\rho_{2r-2},\rho_{2r-1},\rho_{2r})
  \qquad(r\in Z_\gamma\setminus I).
\]
\item Using the deterministic splitting of
\cref{lem:bc-zone-branch-word-splitting}, it recovers a packet branch word
$\beta_{K_t}$ for each component $K_t$.
\item For each $K_t$, it recovers a packet output $U^\ast_t$ such that a fixed
packet inverse recovers the deleted packet subpath
\[
  (\rho_{2u_t-2},\rho_{2u_t-1},\ldots,\rho_{2v_t})
\]
from $U^\ast_t$, $\beta_{K_t}$, the packet endpoint pair
$\rho_{2u_t-2},\rho_{2v_t}$, and $K_t$.
\end{enumerate}
Then the whole original zone subpath
\[
  (\rho_{2a_\gamma-2},\rho_{2a_\gamma-1},\ldots,\rho_{2b_\gamma})
\]
is recoverable from the zone-local datum.
\end{proposition}

\begin{proof}
The set $S$ determines $I$, the canonical halo zones, and the ordered list of
components $K_t$ of $Z_\gamma\cap I$.  By the first hypothesis, every retained
halo block with index in $Z_\gamma\setminus I$ is recovered.  Together with
the given zone endpoint pair, these retained blocks determine the endpoint pair
of every deleted packet by
\cref{lem:bc-zone-retained-blocks-expose-packet-endpoints}.

By the second hypothesis, the zone branch word gives the branch word attached
to each packet.  Therefore the third hypothesis applies to every $K_t$ and
recovers each deleted packet subpath
\[
  (\rho_{2u_t-2},\rho_{2u_t-1},\ldots,\rho_{2v_t}).
\]
The packet intervals $K_t$ and the retained singleton indices
$Z_\gamma\setminus I$ partition $Z_\gamma$.  Interleaving the recovered packet
subpaths and the recovered retained blocks in increasing block order
reconstructs
\[
  \rho_{2a_\gamma-2},\rho_{2a_\gamma-1},\ldots,\rho_{2b_\gamma}.
\]
The adjacent endpoints agree because each recovered piece carries the endpoint
pair it had in the original path.  Hence the zone subpath is recovered.
\end{proof}

\begin{lemma}[Halo-zone outputs have no spare packet windows]\label{lem:bc-zone-no-spare-packet-windows}
Let $S\subset\{1,\ldots,j\}$, put $I=I(S)$, and let
\[
  Z_\gamma=\{a_\gamma,\ldots,b_\gamma\}
\]
be one canonical halo zone.  Write the maximal consecutive components of
$Z_\gamma\cap I$ as $K_1,\ldots,K_m$.  Suppose a target replacement window
$V^\ast_\gamma$ for this zone has exactly
\[
  |Z_\gamma\setminus I|
\]
two-step blocks.  If a parser of $V^\ast_\gamma$ returns, as pairwise disjoint
subwindows,
\begin{enumerate}
\item one unchanged retained halo block for every $r\in Z_\gamma\setminus I$;
and
\item packet output windows $U^\ast_t$ attached to the deleted packets $K_t$,
\end{enumerate}
then every packet output window $U^\ast_t$ has half-length zero.

Consequently, a canonical halo-zone repair cannot expose all retained halo
blocks as unchanged full windows and also expose any positive-length packet
HKKO output as a separate subwindow.  Any positive packet information must be
encoded by modifying or merging retained halo windows, or by a zero-window
inverse whose branch and endpoint data are sufficient for that restricted
packet family.
\end{lemma}

\begin{proof}
There are exactly $|Z_\gamma\setminus I|$ retained halo indices.  The parser in
item (1) therefore already returns $|Z_\gamma\setminus I|$ pairwise disjoint
two-step windows.  Since $V^\ast_\gamma$ itself has exactly this many two-step
blocks, no additional disjoint packet output window can have positive
half-length.  Hence every $U^\ast_t$ has half-length zero.

The consequence is the contrapositive applied to a proposed positive packet
output: a positive separate packet subwindow would add at least one two-step
block beyond the zone target count unless some retained halo window is not
kept as an unchanged full window, or unless the packet is represented by a
zero-window inverse rather than by a positive output window.
\end{proof}

\begin{corollary}[Unchanged-retained zero-window packet parsers are not a uniform halo supplier]\label{cor:bc-unchanged-retained-zero-window-packet-obstruction}
In the ambient two-step growth-path category, the fixed-witness canonical-halo
supplier cannot be obtained by the following universal parser rule:
\begin{enumerate}
\item expose every retained halo block in $Z_\gamma\setminus I(S)$ unchanged
as a disjoint subwindow of $V^\ast_\gamma$;
\item attach no positive-length packet output to any deleted packet; and
\item recover each deleted packet only from its endpoint pair, its known
packet interval, and its packet branch word.
\end{enumerate}
Consequently, any canonical-halo supplier using the below-wall packet-splice
interface must either merge or modify retained halo windows, use additional
fixed-witness restrictions in the packet inverse, or recover packet-sensitive
data from the replacement window.
\end{corollary}

\begin{proof}
If every retained halo block is exposed unchanged as a disjoint subwindow,
\cref{lem:bc-zone-no-spare-packet-windows} forces every disjoint packet output
to have half-length zero.  Thus the proposed rule would need, already for a
singleton packet, a fixed inverse from only the common endpoint pair, the
known singleton interval, and a two-bit packet branch word.

\Cref{cor:bc-singleton-zero-window-packet-obstruction} gives an explicit
ambient family with common endpoint pair
\[
  (5,4,3,2,1),\qquad (5,4,3,2,1),
\]
known singleton interval, and five possible deleted midpoints.  A two-bit
branch word has only four values, so such an endpoint-only zero-window inverse
cannot be injective on that family.  This contradicts the proposed universal
parser rule.  The final sentence lists the remaining alternatives after this
endpoint-only zero-window route is removed.
\end{proof}

\begin{lemma}[Halo-zone parser budget]\label{lem:bc-halo-zone-parser-budget}
Let $S\subset\{1,\ldots,j\}$, put $I=I(S)$, and let
\[
  Z_\gamma=\{a_\gamma,\ldots,b_\gamma\}
\]
be one canonical halo zone.  Write the maximal consecutive components of
$Z_\gamma\cap I$ as $K_1,\ldots,K_m$.  Suppose a target replacement window
$V^\ast_\gamma$ for this zone has exactly
\[
  R_\gamma=|Z_\gamma\setminus I|
\]
two-step blocks.  If a parser of $V^\ast_\gamma$ returns, as pairwise
disjoint subwindows,
\begin{enumerate}
\item unchanged retained halo blocks for the indices in a set
$A\subseteq Z_\gamma\setminus I$; and
\item packet output windows $U^\ast_t$ of half-length $N_t\ge0$ for
$1\le t\le m$,
\end{enumerate}
then
\[
  |A|+\sum_{t=1}^m N_t\le R_\gamma .
\]
Consequently, if the packet output attached to $K_t$ has the raw HKKO
half-length $|K_t|$ for every $t$, then such disjoint outputs can fit only if
\[
  |A|+|Z_\gamma\cap I|\le |Z_\gamma\setminus I|.
\]
In particular, a zone with $|Z_\gamma\cap I|>|Z_\gamma\setminus I|$ cannot
expose raw packet outputs for every deleted packet as disjoint target
subwindows, even if no retained halo block is exposed unchanged.
\end{lemma}

\begin{proof}
The window $V^\ast_\gamma$ contains exactly $R_\gamma$ two-step blocks.  The
subwindows returned by the parser are pairwise disjoint.  The unchanged
retained blocks in item (1) therefore occupy $|A|$ two-step blocks, and the
packet output windows in item (2) occupy $\sum_t N_t$ two-step blocks.  Since
all these subwindows lie inside $V^\ast_\gamma$, their total number of
two-step blocks is at most $R_\gamma$, proving the first inequality.

If $N_t=|K_t|$ for every packet, then the packet intervals $K_t$ are exactly
the components of $Z_\gamma\cap I$, so
\[
  \sum_{t=1}^m N_t
  =
  \sum_{t=1}^m |K_t|
  =
  |Z_\gamma\cap I|.
\]
Substitution gives the second displayed inequality.  If
$|Z_\gamma\cap I|>|Z_\gamma\setminus I|$, this inequality fails even for
$A=\varnothing$, giving the final assertion.
\end{proof}

\begin{corollary}[Packet compression budget]\label{cor:bc-packet-compression-budget}
In the setting of \cref{lem:bc-halo-zone-parser-budget}, suppose that every
packet output half-length satisfies
\[
  0\le N_t\le |K_t|.
\]
Define the total packet compression in the zone by
\[
  C_\gamma=\sum_{t=1}^m (|K_t|-N_t),
\]
and put
\[
  \Delta_\gamma=|Z_\gamma\setminus I|-|Z_\gamma\cap I|.
\]
Then any parser exposing unchanged retained halo blocks for
$A\subseteq Z_\gamma\setminus I$ satisfies
\[
  |A|\le \Delta_\gamma+C_\gamma .
\]
Equivalently, exposing $a$ retained halo blocks requires packet compression at
least $a-\Delta_\gamma$ whenever $a>\Delta_\gamma$.

In particular, a negative-balance zone needs at least $-\Delta_\gamma$ packet
compression even if it exposes no retained halo block, and exposing every
retained halo block requires
\[
  C_\gamma\ge |Z_\gamma\cap I|.
\]
\end{corollary}

\begin{proof}
Since the packet intervals $K_t$ are the components of $Z_\gamma\cap I$,
\[
  \sum_{t=1}^m |K_t|=|Z_\gamma\cap I|.
\]
The definition of $C_\gamma$ gives
\[
  \sum_{t=1}^m N_t=|Z_\gamma\cap I|-C_\gamma .
\]
Substituting this into the parser-budget inequality of
\cref{lem:bc-halo-zone-parser-budget} yields
\[
  |A|+|Z_\gamma\cap I|-C_\gamma
  \le |Z_\gamma\setminus I|.
\]
Rearranging gives $|A|\le\Delta_\gamma+C_\gamma$.  The equivalent lower bound
on $C_\gamma$ is the same inequality rearranged with $|A|=a$.

For $A=\varnothing$ in a zone with $\Delta_\gamma<0$, the inequality gives
$0\le\Delta_\gamma+C_\gamma$, hence $C_\gamma\ge-\Delta_\gamma$.  If every
retained halo block is exposed, then $|A|=|Z_\gamma\setminus I|$, so
\[
  C_\gamma\ge |Z_\gamma\setminus I|-\Delta_\gamma
  =|Z_\gamma\cap I|.
\]
\end{proof}

\begin{corollary}[Raw packet outputs expose only the zone surplus]\label{cor:bc-raw-packet-zone-surplus}
Work in the setting of \cref{lem:bc-halo-zone-parser-budget}.  Suppose every
component $K_t$ is assigned its raw HKKO packet output, and that this output
has half-length $|K_t|$.
Put
\[
  \Delta_\gamma=|Z_\gamma\setminus I|-|Z_\gamma\cap I|.
\]
Then any parser returning these raw packet outputs as pairwise disjoint
subwindows can expose at most $\Delta_\gamma$ unchanged retained halo blocks:
\[
  |A|\le \Delta_\gamma .
\]
In particular, if $\Delta_\gamma<0$, raw packet outputs for all deleted
components cannot fit even with no visible retained block; if
$\Delta_\gamma=0$, no unchanged retained halo block can be exposed.
\end{corollary}

\begin{proof}
The second displayed inequality in \cref{lem:bc-halo-zone-parser-budget} gives
\[
  |A|+|Z_\gamma\cap I|\le |Z_\gamma\setminus I|.
\]
Rearranging gives $|A|\le\Delta_\gamma$.  If $\Delta_\gamma<0$, this is
impossible because $|A|\ge0$.  If $\Delta_\gamma=0$, it forces $|A|=0$.
\end{proof}

\begin{corollary}[Tight raw packet zones hide retained halo blocks]\label{cor:bc-tight-zone-raw-packets-hide-retained-blocks}
In the setting of \cref{lem:bc-halo-zone-parser-budget}, suppose
\[
  |Z_\gamma\cap I|=|Z_\gamma\setminus I|
\]
and suppose the packet output attached to every component $K_t$ has raw HKKO
half-length $|K_t|$.  Then any parser that returns these packet outputs as
pairwise disjoint subwindows returns no unchanged retained halo block as a
disjoint subwindow.

Consequently, in such a tight zone the endpoint-recovery mechanism of
\cref{lem:bc-zone-retained-blocks-expose-packet-endpoints} is unavailable
unless the same target-size replacement blocks also decode the retained halo
blocks by a merged rule.
\end{corollary}

\begin{proof}
The raw packet hypothesis gives
\[
  \sum_t N_t=|Z_\gamma\cap I|.
\]
Since $|Z_\gamma\cap I|=|Z_\gamma\setminus I|=R_\gamma$, the budget inequality
of \cref{lem:bc-halo-zone-parser-budget} gives
\[
  |A|+R_\gamma\le R_\gamma,
\]
and hence $|A|=0$.  Thus no unchanged retained halo block is returned as a
disjoint subwindow.

The consequence follows because
\cref{lem:bc-zone-retained-blocks-expose-packet-endpoints} requires the
retained halo blocks inside $Z_\gamma\setminus I$ as input.  If none of those
blocks is exposed unchanged as a disjoint subwindow, that lemma can be used
only after an additional merged decoder recovers the retained blocks from the
same replacement window.
\end{proof}

\begin{corollary}[Separate Step--3 carrier windows cannot be appended to unchanged retained blocks]\label{cor:bc-separate-step3-carrier-window-obstruction}
Work in the setting of \cref{lem:bc-zone-no-spare-packet-windows}.  Suppose a
replacement-window parser tries to satisfy the retained-block input of
\cref{prop:bc-step3-carrier-packet-readouts-give-halo-zone-inverses} by
returning every retained halo block $B_r(P)$, $r\in Z_\gamma\setminus I$, as an
unchanged pairwise disjoint subwindow of $V^\ast_\gamma$.  Suppose also that,
for some deleted packet $K_t$, the same parser tries to satisfy the Step--3
carrier/readout input of that proposition by returning a positive-half-length
pairwise disjoint packet subwindow from which the carrier $E_t$, its
top/right boundary, and the deterministic Step--2 readout are decoded.  Then
no such parser exists.

Consequently the carrier/readout hypotheses of
\cref{prop:bc-step3-carrier-halo-local-repair-criterion} cannot be supplied by
appending separate positive Step--3 carrier windows to unchanged retained halo
blocks.  Any target-size source theorem for that criterion must recover at
least one of the retained blocks and the carrier/readout data by a merged or
modified replacement-window decoder, or else use a zero-window packet inverse
whose endpoint and branch data are already sufficient for the restricted packet
family.
\end{corollary}

\begin{proof}
The proposed positive Step--3 carrier subwindow is a packet output window in
the sense of \cref{lem:bc-zone-no-spare-packet-windows}.  Since all retained
halo blocks are also returned as unchanged pairwise disjoint subwindows, that
lemma forces every packet output window to have half-length zero.  This
contradicts the assumed positive half-length of the carrier subwindow.

The consequence is the same budget obstruction applied to the packet-carrier
input in
\cref{prop:bc-step3-carrier-packet-readouts-give-halo-zone-inverses}.  The only
remaining target-size possibilities are to decode retained blocks and carrier
data from overlapping or modified replacement-window structure, or to replace
the positive carrier output by a zero-window inverse whose data already
separate the relevant packet family.
\end{proof}

\begin{proposition}[Alternating full-zone decoder criterion]\label{prop:bc-alternating-full-zone-decoder-criterion}
Fix $q\ge1$, put $j=2q$, and let
\[
  S\in
  \bigl\{\{1,3,\ldots,2q-1\},\{2,4,\ldots,2q\}\bigr\}.
\]
Consider either fixed-witness fiber $\mathcal F^w(S)$ or
$\mathcal F^h(S)$ of \cref{prop:bc-jointwise-repair-criterion}.  Suppose
that, for every tableau $T$ in the chosen fiber, its growth path $P(T)$ is
assigned data
\[
  P^\ast(T)\in\mathcal P_q^{\mathrm{even}},
  \qquad
  \beta(T)\in\{0,1\}^{2q},
\]
and that a fixed decoder depending only on $S$ recovers $P(T)$ from
$P^\ast(T)$ and $\beta(T)$.  Then the chosen fixed-witness fiber has
cardinality at most
\[
  2^{2q}s_q .
\]

Moreover, for the canonical halo zones of $S$, this decoder is a single
full-zone decoder: the only zone is $Z=\{1,\ldots,2q\}$, its endpoint pair is
$\varnothing,\varnothing$, its replacement output has $q$ two-step blocks, and
the $q$ deleted components are singleton packets.
\end{proposition}

\begin{proof}
The recovery hypothesis makes the assignment
\[
  T\longmapsto\bigl(P^\ast(T),\beta(T)\bigr)
\]
injective on the chosen fixed-witness fiber.  By
\cref{lem:bc-littlewood-graph-model,lem:bc-growth-path-model},
$|\mathcal P_q^{\mathrm{even}}|=s_q$, and there are $2^{2q}$ possible branch
words.  Thus the chosen fiber has size at most $2^{2q}s_q$.

By \cref{lem:bc-alternating-witness-full-halo-zone}, the canonical halo
construction for this $S$ has the single zone $Z=\{1,\ldots,2q\}$, with
$|Z\setminus I(S)|=q$, and the deleted components are $q$ singletons.  Since a
growth path starts and ends at $\varnothing$ by \cref{lem:bc-growth-path-model},
the endpoint pair of this full zone is $\varnothing,\varnothing$.  The target
path $P^\ast(T)\in\mathcal P_q^{\mathrm{even}}$ has exactly $q$ two-step
blocks, so it is precisely the target-size replacement output for this single
zone.
\end{proof}

\begin{proposition}[Count-compatible halo-zone packages give local inverses]\label{prop:bc-count-compatible-halo-zone-package}
Let
\[
  P=(\rho_0,\rho_1,\ldots,\rho_{2j})
\]
be a two-step growth path, let $S\subset\{1,\ldots,j\}$, put
$I=I(S)$, and let
\[
  Z_\gamma=\{a_\gamma,\ldots,b_\gamma\}
\]
be one canonical halo zone.  Write the maximal consecutive components of
$Z_\gamma\cap I$ in increasing block order as
\[
  K_t=\{u_t,\ldots,v_t\}\qquad(1\le t\le m).
\]
Put $e_\gamma=|Z_\gamma\cap I|$.  Suppose a replacement output
$V^\ast_\gamma$ has exactly $|Z_\gamma\setminus I|$ two-step blocks, and
suppose that from
\[
  S,\ Z_\gamma,\ V^\ast_\gamma,\ \beta_\gamma,\quad
  \rho_{2a_\gamma-2},\rho_{2b_\gamma},
  \qquad
  \beta_\gamma\in\{0,1\}^{2e_\gamma},
\]
a deterministic decoder recovers the following data.
\begin{enumerate}
\item The original retained halo blocks
\[
  B_r(P)=(\rho_{2r-2},\rho_{2r-1},\rho_{2r})
  \qquad(r\in Z_\gamma\setminus I).
\]
\item For every packet $K_t$, a packet output $U^\ast_t$, possibly of
half-length zero.
\item For every packet $K_t$, after splitting $\beta_\gamma$ by the
deterministic cuts of \cref{lem:bc-zone-branch-word-splitting}, a fixed
packet inverse recovers
\[
  (\rho_{2u_t-2},\rho_{2u_t-1},\ldots,\rho_{2v_t})
\]
from $U^\ast_t$, the packet branch word $\beta_{K_t}$, the packet endpoint
pair $\rho_{2u_t-2},\rho_{2v_t}$ determined from the retained blocks and the
zone endpoint pair, and the known packet interval $K_t$.
\end{enumerate}
Then a fixed zone-local inverse recovers the whole original subpath
\[
  (\rho_{2a_\gamma-2},\rho_{2a_\gamma-1},\ldots,\rho_{2b_\gamma})
\]
from $V^\ast_\gamma$, $\beta_\gamma$, the zone endpoint pair, and
$Z_\gamma$.
\end{proposition}

\begin{proof}
The decoder in item (1) recovers every retained halo block.  Hence
\cref{lem:bc-zone-retained-blocks-expose-packet-endpoints} determines the
endpoint pair for every packet $K_t$ from the recovered retained blocks and
the zone endpoint pair.  By
\cref{lem:bc-zone-branch-word-splitting}, the single word $\beta_\gamma$
splits into the packet branch words $\beta_{K_t}$ with no extra branch data.
Item (3) then recovers every deleted packet subpath.

The packet components $K_t$ and the retained singleton indices
$Z_\gamma\setminus I$ partition $Z_\gamma$.  Interleaving the recovered
packet subpaths and retained blocks in increasing block order reconstructs the
displayed zone subpath.  This procedure depends only on the fixed decoder, the
deterministic packet ordering, and the packet inverses, so it is the required
zone-local inverse.
\end{proof}

\begin{proposition}[Anchor-pair packages recover positive height halo zones]\label{prop:bc-height-anchor-pair-zone-recovery}
Work in the height setting of
\cref{cor:bc-height-surplus-canonical-right-anchors}, and let
\[
  P=(\rho_0,\rho_1,\ldots,\rho_{2j})
\]
be the corresponding growth path.  Write
\[
  Z_\gamma\cap I=\{u_1,\ldots,u_m\},\qquad u_1<\cdots<u_m,
\]
and let $A_\gamma$ be the canonical right-anchor set of
\cref{cor:bc-height-surplus-canonical-right-anchors}.  Suppose a replacement
window $V^\ast_\gamma$ has exactly $|Z_\gamma\setminus I|$ two-step blocks
and that, from
\[
  S,\ Z_\gamma,\ V^\ast_\gamma,\ \beta_\gamma,\quad
  \rho_{2a_\gamma-2},\rho_{2b_\gamma},
  \qquad
  \beta_\gamma\in\{0,1\}^{2m},
\]
a deterministic parser decomposes $V^\ast_\gamma$ into
\begin{enumerate}
\item unchanged anchor blocks
\[
  B_r(P)=(\rho_{2r-2},\rho_{2r-1},\rho_{2r})
  \qquad(r\in A_\gamma),
\]
in their block order; and
\item one merged two-step window $M_t$ for each deleted singleton $u_t$,
listed in increasing $t$.
\end{enumerate}
Assume that the deterministic split of $\beta_\gamma$ into consecutive
two-bit words $\beta_t$ has the following local inverse property: for each
$t$, a fixed pair inverse recovers the consecutive original two-block segment
\[
  \rho_{2u_t-4},\rho_{2u_t-3},\rho_{2u_t-2},
  \rho_{2u_t-1},\rho_{2u_t}
\]
from $M_t$, $\beta_t$, the zone endpoint pair, $S$, $Z_\gamma$, and $t$.
Then the whole original zone subpath
\[
  \rho_{2a_\gamma-2},\rho_{2a_\gamma-1},\ldots,\rho_{2b_\gamma}
\]
is recoverable from $V^\ast_\gamma$, $\beta_\gamma$, the zone endpoint pair,
and $Z_\gamma$.
\end{proposition}

\begin{proof}
The branch word split into two-bit subwords is determined by the ordered
singleton list $u_1<\cdots<u_m$, hence by $S$ and $Z_\gamma$.
Also
\[
  |A_\gamma|+m=|Z_\gamma\setminus I|
\]
by \cref{cor:bc-height-surplus-canonical-right-anchors}, so the displayed
parser uses exactly the target-size number of two-step windows.
By \cref{cor:bc-height-anchor-reduced-left-pairs},
\[
  (Z_\gamma\setminus I)\setminus A_\gamma
  =
  \{u_t-1:\ 1\le t\le m\}.
\]
Together with $Z_\gamma\cap I=\{u_1,\ldots,u_m\}$, this says that the block
indices in $Z_\gamma$ are partitioned into the anchor singleton blocks
$A_\gamma$ and the consecutive pairs
\[
  \{u_t-1,u_t\}\qquad(1\le t\le m).
\]

The parser recovers the original anchor blocks.  For each $t$, the assumed
pair inverse recovers the original two consecutive blocks with indices
$u_t-1$ and $u_t$.  Interleaving the recovered anchor blocks and recovered
two-block pair segments in increasing block order therefore reconstructs every
block of the original zone.  The endpoint overlaps agree because each
recovered piece is recovered as a subpath of the original growth path with its
original endpoints.  Thus the displayed zone subpath is uniquely recovered.
\end{proof}

\begin{corollary}[Retained-only anchor-pair windows are not universal]\label{cor:bc-retained-only-anchor-pair-obstruction}
The pair inverse in \cref{prop:bc-height-anchor-pair-zone-recovery} cannot,
in the ambient growth-path category, be replaced by the rule which takes
$M_t$ to be only the unchanged left-neighbour retained block
\[
  B_{u_t-1}(P)
\]
and then tries to recover the deleted singleton block from the retained block,
the outer endpoint pair, the known pair position, and a two-bit word.
Equivalently, a valid merged pair-window construction must either modify the
left-neighbour block, carry deleted-sensitive target data, or use additional
fixed-witness restrictions beyond the ambient growth-path axioms.
\end{corollary}

\begin{proof}
Let $\mu=(5,4,3,2,1)$, and fix one removable corner $c_0$ of $\mu$.  Let
\[
  U^0=\mu,\qquad U^1=\mu\setminus\{c_0\},\qquad U^2=\mu
\]
be the retained left-neighbour block.  This block has vertical-strip adjacent
moves and weight
\[
  |U^0|-2|U^1|+|U^2|=2 .
\]
Now append a deleted singleton block from $\mu$ back to $\mu$ by choosing any
one of the five removable corners $c$ of $\mu$:
\[
  U^2=\mu,\qquad U^3=\mu\setminus\{c\},\qquad U^4=\mu .
\]
Each choice of $c$ gives a valid two-step deleted singleton block of weight
$2$, and all five choices have the same retained block $B_{u_t-1}$ and the
same outer endpoint pair $(U^0,U^4)=(\mu,\mu)$.  A two-bit word has only four
values, so those data cannot encode all five possibilities.  Therefore an
inverse using only the unchanged retained block, the endpoint data, the pair
position, and two bits cannot recover all such two-block segments.
\end{proof}

\begin{corollary}[One-block ambient pair windows are not universal]\label{cor:bc-one-block-anchor-pair-obstruction}
Even if the merged pair window in
\cref{prop:bc-height-anchor-pair-zone-recovery} is allowed to be an arbitrary
single two-step growth block with the same outer endpoint pair as the original
two-block segment, no ambient inverse can recover all such two-block segments
from that one-block window and a two-bit word.
Equivalently, a target-size one-block anchor-pair construction must use
additional fixed-witness restrictions, neighbouring parser data, or nonlocal
right-boundary data; a generic same-endpoint one-block replacement is not
enough in the ambient growth-path category.
\end{corollary}

\begin{proof}
Again let
\[
  \mu=(5,4,3,2,1).
\]
For every ordered pair $(c,d)$ of removable corners of $\mu$, define a
two-block segment
\[
  U^0=\mu,\qquad U^1=\mu\setminus\{c\},\qquad U^2=\mu,\qquad
  U^3=\mu\setminus\{d\},\qquad U^4=\mu .
\]
Both blocks have vertical-strip adjacent moves and weight $2$, so this gives
$25$ distinct ambient two-block segments with common outer endpoint pair
$(\mu,\mu)$.

Now consider any single two-step growth block
\[
  M=(\mu,X,\mu)
\]
with the same outer endpoint pair and weight $2$.  The weight condition gives
\[
  |\mu|-2|X|+|\mu|=2,
\]
so $|X|=|\mu|-1$.  Since both adjacent moves are vertical-strip moves and
$X\subseteq\mu$, the partition $X$ is obtained from $\mu$ by deleting one
removable corner.  Thus there are only five possible same-endpoint one-block
windows $M$.

Together with a two-bit word, such data have at most $5\cdot4=20$ possible
values, while the ambient two-block family above has $25$ members.  Therefore
no inverse using only a same-endpoint one-block replacement window and two
branch bits can recover all ambient two-block segments.  The final statement
is the same counting obstruction applied to a proposed target-size
anchor-pair merger before any fixed-witness or extra parser data have been
used.
\end{proof}

\begin{proposition}[Below-wall packet splices give count-compatible halo-zone inverses]\label{prop:bc-below-wall-packet-splices-give-halo-zone-inverses}
Work in the setting of
\cref{prop:bc-count-compatible-halo-zone-package}.  For each deleted packet
$K_t=\{u_t,\ldots,v_t\}$ in $Z_\gamma\cap I$, let $J_t$ be the corresponding
maximal witness-label interval, and put $h_t=|J_t|$.  Interpret each packet
in the local branch-recovery coordinates of
\cref{lem:bc-local-branch-recovery}; the packet branch word obtained from
\cref{lem:bc-zone-branch-word-splitting} is the local branch word for this
datum.  Suppose that the
replacement-window decoder, from
\[
  S,\ Z_\gamma,\ V^\ast_\gamma,\ \beta_\gamma,\quad
  \rho_{2a_\gamma-2},\rho_{2b_\gamma},
\]
recovers the retained halo blocks
\[
  B_r(P)\qquad(r\in Z_\gamma\setminus I),
\]
so that the packet endpoint pair
\[
  \mu_t=\rho_{2u_t-2},\qquad \nu_t=\rho_{2v_t}
\]
is determined for every $K_t$ by
\cref{lem:bc-zone-retained-blocks-expose-packet-endpoints}.  Suppose further
that, for every packet $K_t$, the canonical marked packet preimage
\[
  W_t=W(\mu_t,\nu_t,\beta_{K_t})
\]
of \cref{lem:bc-canonical-below-wall-packet-output} admits a fixed packet
readout which recovers the deleted packet subpath
\[
  (\rho_{2u_t-2},\rho_{2u_t-1},\ldots,\rho_{2v_t})
\]
from $W_t$, the packet endpoint pair $\rho_{2u_t-2},\rho_{2v_t}$, and the
known packet interval $K_t$.  No packet-output datum is required from the
replacement-window decoder.
Then a fixed zone-local inverse recovers the whole original subpath
\[
  (\rho_{2a_\gamma-2},\rho_{2a_\gamma-1},\ldots,\rho_{2b_\gamma})
\]
from $V^\ast_\gamma$, $\beta_\gamma$, the zone endpoint pair, and
$Z_\gamma$.
\end{proposition}

\begin{proof}
By \cref{lem:bc-zone-branch-word-splitting}, the zone branch word
$\beta_\gamma$ splits, with cuts determined by $S$, into packet branch words
\[
  \beta_{K_t}\in\{0,1\}^{2h_t}.
\]
The recovered retained halo blocks and the zone endpoint pair determine
\[
  \mu_t=\rho_{2u_t-2},\qquad \nu_t=\rho_{2v_t}
\]
for every packet by
\cref{lem:bc-zone-retained-blocks-expose-packet-endpoints}.
For a fixed packet $K_t$, let $w_t=w(\mu_t,\nu_t)$ be the canonical wall of
\cref{lem:bc-local-endpoints-canonical-wall}.  Applying
\cref{lem:bc-canonical-below-wall-packet-output} to the packet endpoint pair
$\mu_t,\nu_t$ and the branch word $\beta_{K_t}$ gives the canonical output
$W^\ast_t$ and its marked preimage $W_t$.  The same lemma recovers $W_t$ from
\[
  W^\ast_t,\quad \beta_{K_t},\quad w_t,\quad h_t,\quad \mu_t .
\]
The fixed packet readout in the hypothesis then recovers the deleted packet
subpath from $W_t$, the packet endpoint pair, and $K_t$.

Thus every packet has the fixed packet inverse required in item (3) of
\cref{prop:bc-count-compatible-halo-zone-package}.  The retained halo blocks
are recovered by hypothesis, and the packet outputs for item (2) are the
canonically constructed $W^\ast_t$.  Applying
\cref{prop:bc-count-compatible-halo-zone-package} gives the desired
zone-local inverse.
\end{proof}

\begin{proposition}[Label-box packet readouts give halo-zone inverses]\label{prop:bc-label-box-packet-readouts-give-halo-zone-inverses}
Work in the setting of
\cref{prop:bc-count-compatible-halo-zone-package}.  For each deleted packet
$K_t=\{u_t,\ldots,v_t\}$ in $Z_\gamma\cap I$, let $J_t$ be the corresponding
maximal witness-label interval, and put $h_t=|J_t|$.  Interpret each packet in
the local branch-recovery coordinates of \cref{lem:bc-local-branch-recovery};
the packet branch word obtained from \cref{lem:bc-zone-branch-word-splitting}
is the local branch word for this datum.

Suppose that the replacement-window decoder, from
\[
  S,\ Z_\gamma,\ V^\ast_\gamma,\ \beta_\gamma,\quad
  \rho_{2a_\gamma-2},\rho_{2b_\gamma},
\]
recovers the retained halo blocks
\[
  B_r(P)\qquad(r\in Z_\gamma\setminus I),
\]
and, for every packet $K_t$, recovers a datum $A_t$ which determines, in the
semistandard tableau encoded by the local packet, the two boxes carrying each
deleted label in $J_t$.  Then a fixed zone-local inverse recovers the whole
original subpath
\[
  (\rho_{2a_\gamma-2},\rho_{2a_\gamma-1},\ldots,\rho_{2b_\gamma})
\]
from $V^\ast_\gamma$, $\beta_\gamma$, the zone endpoint pair, and $Z_\gamma$.
\end{proposition}

\begin{proof}
By \cref{lem:bc-zone-branch-word-splitting}, the zone branch word
$\beta_\gamma$ splits, with cuts determined by $S$, into packet branch words
\[
  \beta_{K_t}\in\{0,1\}^{2h_t}.
\]
The recovered retained halo blocks and the zone endpoint pair determine
\[
  \rho_{2u_t-2},\qquad \rho_{2v_t}
\]
for every packet by \cref{lem:bc-zone-retained-blocks-expose-packet-endpoints}.
For a fixed packet $K_t$, the known interval $J_t$, the left endpoint
$\rho_{2u_t-2}$ in local coordinates, the local branch word $\beta_{K_t}$, and
the recovered label-box datum $A_t$ satisfy the hypotheses of
\cref{cor:bc-local-label-boxes-supply-second-copy-readout}.  That corollary
recovers the whole standardized local packet subpath.  Under the local
branch-recovery interpretation of the packet, this is the deleted subpath
\[
  (\rho_{2u_t-2},\rho_{2u_t-1},\ldots,\rho_{2v_t}).
\]

Thus every deleted packet subpath is recovered from the replacement-window
decoder output and the already budgeted branch word.  The packet components
$K_t$ and the retained singleton indices $Z_\gamma\setminus I$ partition
$Z_\gamma$.  Interleaving the recovered packet subpaths and retained blocks in
increasing block order reconstructs the displayed zone subpath.
\end{proof}

\begin{proposition}[Label-box halo repairs give the fixed-witness bound]\label{prop:bc-label-box-halo-local-repair-criterion}
Fix $q\ge1$, $j\ge q$, and a $q$-element set
$S\subset\{1,\ldots,j\}$.  Let $Z_\gamma$ be the canonical halo replacement
zones of \cref{lem:bc-canonical-halo-replacement-zones}, with each zone written
as $Z_\gamma=\{a_\gamma,\ldots,b_\gamma\}$, and consider either fixed-witness
fiber $\mathcal F^w(S)$ or $\mathcal F^h(S)$ of
\cref{prop:bc-jointwise-repair-criterion}.  For each zone put
\[
  e_\gamma=|Z_\gamma\cap I(S)|.
\]
Suppose that, for every $T$ in the chosen fiber, there is a repaired path
$P^\ast(T)\in\mathcal P_{j-q}^{\mathrm{even}}$ and a parser depending only on
$S$ with the following properties.
\begin{enumerate}
\item The parser decomposes $P^\ast(T)$ into the original retained blocks
with indices outside the halo zones, unchanged and in their original order,
and local replacement windows $V^\ast_\gamma(T)$ for the zones $Z_\gamma$.
\item Each $V^\ast_\gamma(T)$ has exactly
$|Z_\gamma\setminus I(S)|$ two-step blocks.
\item For each zone there is a branch word
\[
  \beta_\gamma(T)\in\{0,1\}^{2e_\gamma}
\]
and, for every zone, the endpoint pair
$\rho_{2a_\gamma-2},\rho_{2b_\gamma}$ is the one determined from the parsed
complement blocks and the fixed endpoints of the growth-path model by
\cref{lem:bc-canonical-halo-complement-exposes-zone-endpoints}.
\item For each zone, the replacement-window decoder satisfies the hypotheses of
\cref{prop:bc-label-box-packet-readouts-give-halo-zone-inverses} for
$V^\ast_\gamma(T)$, $\beta_\gamma(T)$, the zone endpoint pair, and the known
interval $Z_\gamma$.
\end{enumerate}
Then the chosen fixed-witness fiber has cardinality at most
\[
  2^{2q}s_{j-q}.
\]
\end{proposition}

\begin{proof}
By \cref{lem:bc-halo-zones-frontier-packets}, for each canonical halo zone
the number $e_\gamma=|Z_\gamma\cap I(S)|$ is also the sum of the frontier-event
packet lengths used in
\cref{prop:bc-canonical-halo-local-repair-criterion}.  The endpoint pair of
each zone is determined from the parsed complement blocks by
\cref{lem:bc-canonical-halo-complement-exposes-zone-endpoints}.  For each zone,
item~(4) and
\cref{prop:bc-label-box-packet-readouts-give-halo-zone-inverses} give a fixed
local inverse recovering the original subpath on $Z_\gamma$ from
$V^\ast_\gamma(T)$, $\beta_\gamma(T)$, the determined endpoint pair, and the
known zone interval.  Therefore the hypotheses of
\cref{prop:bc-canonical-halo-local-repair-criterion} hold, and that proposition
gives the displayed fixed-witness bound.
\end{proof}

\begin{proposition}[Below-wall packet-splice halo repairs give the fixed-witness bound]\label{prop:bc-below-wall-packet-splice-halo-local-repair-criterion}
Fix $q\ge1$, $j\ge q$, and a $q$-element set
$S\subset\{1,\ldots,j\}$.  Let $Z_\gamma$ be the canonical halo replacement
zones of \cref{lem:bc-canonical-halo-replacement-zones}, with each zone written
as $Z_\gamma=\{a_\gamma,\ldots,b_\gamma\}$, and consider either fixed-witness
fiber $\mathcal F^w(S)$ or $\mathcal F^h(S)$ of
\cref{prop:bc-jointwise-repair-criterion}.  For each zone put
\[
  e_\gamma=|Z_\gamma\cap I(S)|.
\]
Suppose that, for every $T$ in the chosen fiber, there is a repaired path
$P^\ast(T)\in\mathcal P_{j-q}^{\mathrm{even}}$ and a parser depending only on
$S$ with the following properties.
\begin{enumerate}
\item The parser decomposes $P^\ast(T)$ into the original retained blocks
with indices outside the halo zones, unchanged and in their original order,
and local replacement windows $V^\ast_\gamma(T)$ for the zones $Z_\gamma$.
\item Each $V^\ast_\gamma(T)$ has exactly
$|Z_\gamma\setminus I(S)|$ two-step blocks.
\item For each zone there is a branch word
\[
  \beta_\gamma(T)\in\{0,1\}^{2e_\gamma}
\]
and, for every zone, the endpoint pair
$\rho_{2a_\gamma-2},\rho_{2b_\gamma}$ is the one determined from the parsed
complement blocks and the fixed endpoints of the growth-path model by
\cref{lem:bc-canonical-halo-complement-exposes-zone-endpoints}.
\item For each zone, the replacement-window decoder satisfies the hypotheses of
\cref{prop:bc-below-wall-packet-splices-give-halo-zone-inverses} for
$V^\ast_\gamma(T)$, $\beta_\gamma(T)$, the zone endpoint pair, and the known
interval $Z_\gamma$.
\end{enumerate}
Then the chosen fixed-witness fiber has cardinality at most
\[
  2^{2q}s_{j-q}.
\]
\end{proposition}

\begin{proof}
By \cref{lem:bc-halo-zones-frontier-packets}, for each canonical halo zone
the number $e_\gamma=|Z_\gamma\cap I(S)|$ is also the sum of the frontier-event
packet lengths used in
\cref{prop:bc-canonical-halo-local-repair-criterion}.
\Cref{lem:bc-canonical-halo-complement-exposes-zone-endpoints} applied to the
parsed complement blocks and the fixed endpoints of
\cref{lem:bc-growth-path-model} determines the endpoint pair of every
canonical halo zone.
For each zone, item (4) and
\cref{prop:bc-below-wall-packet-splices-give-halo-zone-inverses} give a fixed
local inverse recovering the original subpath on $Z_\gamma$ from
$V^\ast_\gamma(T)$, $\beta_\gamma(T)$, the zone endpoint pair, and the known
zone interval.  Augmenting the parser by the endpoint pairs determined above,
items (1)--(3) are exactly the parser, target-window, branch-word, and local
inverse hypotheses of
\cref{prop:bc-canonical-halo-local-repair-criterion}.  That proposition
applies and gives the displayed fixed-witness bound.
\end{proof}

\begin{corollary}[Below-wall packet splices supply canonical halo repair data]\label{cor:bc-below-wall-packet-splice-supplies-canonical-halo-repair}
Work in the setting of
\cref{prop:bc-below-wall-packet-splice-halo-local-repair-criterion}.  If the
parser, replacement windows, branch words, and packet readouts satisfy the
hypotheses of that proposition, then the same data satisfy the hypotheses of
\cref{prop:bc-canonical-halo-local-repair-criterion}.
\end{corollary}

\begin{proof}
By \cref{lem:bc-halo-zones-frontier-packets}, for each canonical halo zone
one has
\[
  e_\gamma=|Z_\gamma\cap I(S)|.
\]
Thus the branch words have the lengths required in
\cref{prop:bc-canonical-halo-local-repair-criterion}.  The parser and
replacement-window length hypotheses are identical in the two criteria.
The zone endpoint pair is determined from the parsed complement blocks by
\cref{lem:bc-canonical-halo-complement-exposes-zone-endpoints}.  Finally,
item~(4) of
\cref{prop:bc-below-wall-packet-splice-halo-local-repair-criterion} and
\cref{prop:bc-below-wall-packet-splices-give-halo-zone-inverses} give, for
each zone, a fixed local inverse recovering the original zone subpath from the
replacement window, the branch word, the determined endpoint pair, and the
known zone interval.  These are exactly the remaining local-inverse hypotheses
of \cref{prop:bc-canonical-halo-local-repair-criterion}.
\end{proof}

\begin{proposition}[Canonical halo local repair criterion]\label{prop:bc-canonical-halo-local-repair-criterion}
Fix $q\ge1$, $j\ge q$, and a $q$-element set
$S\subset\{1,\ldots,j\}$.  Let $Z_\gamma$ be the canonical halo replacement
zones of \cref{lem:bc-canonical-halo-replacement-zones}, with each zone written
as $Z_\gamma=\{a_\gamma,\ldots,b_\gamma\}$, and consider either fixed-witness
fiber $\mathcal F^w(S)$ or $\mathcal F^h(S)$ of
\cref{prop:bc-jointwise-repair-criterion}.  For a zone $Z_\gamma$, let
\[
  e_\gamma=\sum_{K_\beta\subseteq Z_\gamma}|E_\beta|
\]
using the frontier-event packets assigned by
\cref{lem:bc-halo-zones-frontier-packets}.  Suppose that, for every
$T$ in the chosen fiber, there is a repaired path
$P^\ast(T)\in\mathcal P_{j-q}^{\mathrm{even}}$ and a parser depending only on
$S$ with the following properties.
\begin{enumerate}
\item The parser decomposes $P^\ast(T)$ into the original retained blocks
with indices outside the halo zones, unchanged and in their original order, and
local replacement windows $V^\ast_\gamma(T)$ for the zones $Z_\gamma$.
\item Each $V^\ast_\gamma(T)$ has exactly
$|Z_\gamma\setminus I(S)|$ two-step blocks.
\item For each $\gamma$ there is a branch word
\[
  \beta_\gamma(T)\in\{0,1\}^{2e_\gamma}
\]
and a fixed local inverse recovers the original subpath on $Z_\gamma$ from
$V^\ast_\gamma(T)$, $\beta_\gamma(T)$, the endpoint pair
$\rho_{2a_\gamma-2},\rho_{2b_\gamma}$ determined from the parsed complement
blocks by \cref{lem:bc-canonical-halo-complement-exposes-zone-endpoints}, and
the known interval $Z_\gamma$.
\end{enumerate}
Then the chosen fixed-witness fiber has cardinality at most
\[
  2^{2q}s_{j-q}.
\]
\end{proposition}

\begin{proof}
By \cref{lem:bc-halo-zones-frontier-packets},
\[
  e_\gamma=|Z_\gamma\cap I(S)|
\]
for each canonical halo zone.  Hence the branch words in the statement have
the form required in \cref{prop:bc-replacement-zone-fiber-criterion}.  The
canonical halo zones satisfy the deterministic-zone and block-count hypotheses
by \cref{lem:bc-canonical-halo-replacement-zones}.  The three displayed
properties, together with
\cref{lem:bc-canonical-halo-complement-exposes-zone-endpoints}, canonically
augment the parser by the endpoint pair of every zone.  After this deterministic
augmentation they are exactly the remaining parser, target, and local-inverse
hypotheses of \cref{prop:bc-replacement-zone-fiber-criterion}.  Applying that
proposition gives the claimed bound.
\end{proof}

\begin{proposition}[Replacement-zone criterion for fixed-witness fibers]\label{prop:bc-replacement-zone-fiber-criterion}
Fix $q\ge1$, $j\ge q$, and a $q$-element set
$S\subset\{1,\ldots,j\}$.  Consider either fixed-witness fiber
$\mathcal F^w(S)$ or $\mathcal F^h(S)$ of
\cref{prop:bc-jointwise-repair-criterion}.  Suppose that, for every tableau
$T$ in the chosen fiber, its growth path $P(T)$ is assigned deterministic
replacement zones $Z_\alpha$ satisfying the hypotheses of
\cref{prop:bc-replacement-zones-global-recovery}, with
\[
  I(S)=\{j-s+1:s\in S\}
\]
and with branch words
\[
  \beta_\alpha(T)\in
  \{0,1\}^{2|Z_\alpha\cap I(S)|}.
\]
Assume also that the repaired path $P^\ast(T)$ lies in
$\mathcal P_{j-q}^{\mathrm{even}}$.  Then the chosen fixed-witness fiber has
cardinality at most
\[
  2^{2q}s_{j-q}.
\]
\end{proposition}

\begin{proof}
By \cref{prop:bc-replacement-zones-global-recovery}, the data
\[
  \Bigl(P^\ast(T),(\beta_\alpha(T))_\alpha\Bigr)
\]
recover the original growth path $P(T)$ from $S$ and hence recover $T$ by the
growth-path encoding.  Therefore the assignment from the chosen fixed-witness
fiber to these data is injective.

The zones are disjoint and cover every element of $I(S)$ exactly once, so
\[
  \sum_\alpha |Z_\alpha\cap I(S)|=|I(S)|=q.
\]
Thus the total number of possible branch-word tuples is
\[
  \prod_\alpha 2^{2|Z_\alpha\cap I(S)|}=2^{2q}.
\]
The possible repaired paths are contained in
$\mathcal P_{j-q}^{\mathrm{even}}$, whose cardinality is $s_{j-q}$ by
\cref{lem:bc-littlewood-graph-model,lem:bc-growth-path-model}.  Multiplying
these two counts gives the claimed bound.
\end{proof}

\begin{proposition}[Right-boundary auxiliary fixed-witness criterion]\label{prop:bc-right-boundary-auxiliary-fiber-criterion}
Fix $q\ge1$, $j\ge q$, and a $q$-element set
$S\subset\{1,\ldots,j\}$.  Consider either fixed-witness fiber
$\mathcal F^w(S)$ or $\mathcal F^h(S)$ of
\cref{prop:bc-jointwise-repair-criterion}.  Suppose that, for every tableau
$T$ in the chosen fiber, its growth path $P(T)$ is assigned deterministic
replacement zones
\[
  Z_1,\ldots,Z_L
\]
which satisfy the hypotheses of
\cref{prop:bc-right-boundary-auxiliary-global-recovery}, with
\[
  I(S)=\{j-s+1:s\in S\},
  \qquad Z_L=\{a_L,\ldots,j\},
\]
and with branch words
\[
  \beta_\alpha(T)\in
  \{0,1\}^{2|Z_\alpha\cap I(S)|}.
\]
Assume also that the repaired path $P^\ast(T)$ lies in
$\mathcal P_{j-q}^{\mathrm{even}}$.  Then the chosen fixed-witness fiber has
cardinality at most
\[
  2^{2q}s_{j-q}.
\]
\end{proposition}

\begin{proof}
By \cref{prop:bc-right-boundary-auxiliary-global-recovery}, the data
\[
  \Bigl(P^\ast(T),(\beta_\alpha(T))_{\alpha=1}^L\Bigr)
\]
recover the original growth path $P(T)$ from $S$ and hence recover $T$.
Thus the assignment from the chosen fixed-witness fiber to these data is
injective.

The replacement zones are disjoint and cover every element of $I(S)$ exactly
once, so
\[
  \sum_{\alpha=1}^L |Z_\alpha\cap I(S)|=|I(S)|=q.
\]
The total number of branch-word tuples is therefore
\[
  \prod_{\alpha=1}^L 2^{2|Z_\alpha\cap I(S)|}=2^{2q}.
\]
The repaired paths lie in
$\mathcal P_{j-q}^{\mathrm{even}}$, whose cardinality is $s_{j-q}$ by
\cref{lem:bc-littlewood-graph-model,lem:bc-growth-path-model}.  Multiplying
gives the displayed bound.
\end{proof}

\begin{proposition}[Jointwise repair criterion for fixed witnesses]\label{prop:bc-jointwise-repair-criterion}
Fix $q\ge1$, $j\ge q$, and a $q$-element set
$S\subset\{1,\ldots,j\}$.  Define the two fixed-witness fibers
\[
 \mathcal F^w(S)=
 \{T:\lambda(T)' \text{ even},\ \lambda_1(T)\ge2q,\ W_q(T)=S\}
\]
and
\[
 \mathcal F^h(S)=
 \{T:\lambda(T)' \text{ even},\ \ell(\lambda(T))\ge2q,\ H_q(T)=S\}.
\]
Let $\mathcal P_{j-q}^{\mathrm{even}}$ be the set of column-even two-step
growth paths corresponding to semistandard tableaux of content $(2^{j-q})$.

Consider either $\mathcal F^w(S)$ or $\mathcal F^h(S)$.  Suppose that for
every $T$ in the chosen fiber, after decomposing $S$ into maximal consecutive
label intervals $J_\alpha$ and taking the corresponding per-joint event sets
$E_\alpha(T)$ of \cref{lem:bc-frontier-events-per-joint}, there is an
assignment
\[
  T\longmapsto
  \Bigl(P^\ast(T),(\beta_\alpha(T))_{\alpha=1}^L\Bigr)
\]
such that $P^\ast(T)\in\mathcal P_{j-q}^{\mathrm{even}}$,
\[
  \beta_\alpha(T)\in\{0,1\}^{2|J_\alpha|},
\]
and $T$ is uniquely recoverable from $S$, $P^\ast(T)$, and the words
$\beta_\alpha(T)$.  Then the chosen fixed-witness fiber has cardinality at
most
\[
  2^{2q}s_{j-q}.
\]
\end{proposition}

\begin{proof}
By \cref{lem:bc-frontier-events-per-joint}, the sets $E_\alpha(T)$ contain
one frontier event for each label in the corresponding interval $J_\alpha$.
Since the intervals $J_\alpha$ partition $S$, one has
\[
  \sum_{\alpha=1}^L |E_\alpha(T)|
  =
  \sum_{\alpha=1}^L |J_\alpha|
  =
  |S|=q .
\]
Therefore the total number of possible branch-word tuples is
\[
  \prod_{\alpha=1}^L 2^{2|J_\alpha|}
  =
  2^{2q}.
\]
The recovery hypothesis makes the displayed assignment injective from the
chosen fiber into
\[
  \mathcal P_{j-q}^{\mathrm{even}}
  \times
  \prod_{\alpha=1}^L\{0,1\}^{2|J_\alpha|}.
\]
Finally,
\cref{lem:bc-littlewood-graph-model,lem:bc-growth-path-model} identify
$|\mathcal P_{j-q}^{\mathrm{even}}|$ with the stable moment $s_{j-q}$.
Multiplying the two counts gives the claimed bound.
\end{proof}

\begin{corollary}[Defect bound for fixed residual witnesses]\label{cor:bc-witness-defect}
In either fixed-witness fiber appearing in
\cref{prop:bc-bad-shape-fiber-reduction}, deleting the two copies of every
label in the fixed witness set $S$ leaves a labelled retained subdiagram of
content $(2^{j-q})$ whose remaining row and column inequalities are inherited
from the original tableau and whose column occupancies are odd in at most
$2q$ columns.
\end{corollary}

\begin{proof}
Apply \cref{lem:bc-deletion-defect} with $|S|=q$.
\end{proof}

\begin{lemma}[Alternating frontiers after witness deletion]\label{lem:bc-witness-frontier}
Let $T$ be as in \cref{lem:bad-shape-witnesses}.
If $\lambda_1(T)\ge2q$ and $S=W_q(T)$, then in the retained subdiagram
$T\setminus S$ no box remains at any of the top-row positions
\[
  (1,1),(1,3),\ldots,(1,2q-1).
\]
Consequently, among the first $2q$ positions of the top row, at most $q$
boxes remain.

If $\ell(\lambda(T))\ge2q$ and $S=H_q(T)$, then in the retained subdiagram
$T\setminus S$ no box remains at any of the first-column positions
\[
  (1,1),(3,1),\ldots,(2q-1,1).
\]
Consequently, among the first $2q$ positions of the first column, at most
$q$ boxes remain.
\end{lemma}

\begin{proof}
In the width case, the labels occupying
$(1,1),(1,3),\ldots,(1,2q-1)$ are exactly the labels in $W_q(T)=S$ by
definition, so those boxes are deleted when forming $T\setminus S$.  The only
top-row positions among the first $2q$ that can remain are therefore the even
positions $2,4,\ldots,2q$, of which there are $q$.

The height case is identical with the first column in place of the first row:
the boxes at $(1,1),(3,1),\ldots,(2q-1,1)$ carry exactly the labels in
$H_q(T)=S$, and only the even row positions among $1,\ldots,2q$ can remain.
\end{proof}

\begin{lemma}[Alternating frontier retention is local]\label{lem:bc-alternating-frontier-retention}
Let $T$ be as in \cref{lem:bad-shape-witnesses}.

In the width case, assume $\lambda_1(T)\ge2q$, put
$s_r=T(1,2r-1)$ for $1\le r\le q$, and set $S=W_q(T)$.  For
$1\le r\le q$, put $a_r=T(1,2r)$.  After deleting the labels in $S$, the box
$(1,2r-1)$ is absent.  The box $(1,2r)$ is absent if and only if either
$a_r=s_r$ or $r<q$ and $a_r=s_{r+1}$; otherwise $(1,2r)$ remains.

In the height case, assume $\ell(\lambda(T))\ge2q$, put
$t_r=T(2r-1,1)$ for $1\le r\le q$, and set $S=H_q(T)$.  After deleting the
labels in $S$, the boxes
\[
  (1,1),(3,1),\ldots,(2q-1,1)
\]
are absent, while the boxes
\[
  (2,1),(4,1),\ldots,(2q,1)
\]
all remain.  Thus exactly $q$ boxes remain among the first $2q$ first-column
positions.
\end{lemma}

\begin{proof}
The odd width boxes carry precisely the labels in $S$ by definition of
$W_q(T)$, so they are absent after deleting $S$.  Since the first row is weakly
increasing,
\[
  s_r\le a_r\le s_{r+1}\qquad(1\le r<q),
\]
and $s_q\le a_q$.  The labels $s_1,\ldots,s_q$ are strictly increasing by
\cref{lem:bad-shape-witnesses}.  Hence, for $r<q$, the only labels from $S$
which can equal $a_r$ are $s_r$ and $s_{r+1}$; for $r=q$, the only possible
label from $S$ equal to $a_q$ is $s_q$.  A box is deleted exactly when its
label lies in $S$, proving the width statement.

For height, the odd first-column boxes carry exactly the labels in $S$ and are
therefore deleted.  Let $b_r=T(2r,1)$ for $1\le r\le q$.  The first column is
strictly increasing.  Thus, for $r<q$,
\[
  t_r<b_r<t_{r+1},
\]
and for $r=q$ one has $t_q<b_q$.  Therefore no $b_r$ belongs to
$\{t_1,\ldots,t_q\}=S$.  The even first-column boxes are not deleted, and the
last assertion follows.
\end{proof}

\begin{corollary}[Width orientation forces retention]\label{cor:bc-width-orientation-retention}
In the width setting of \cref{lem:bc-alternating-frontier-retention}, let
$\epsilon_r^w(T)$ be the orientation bits of
\cref{lem:bc-standardized-witness-pairs}.  For $1\le r<q$, if the box
$(1,2r)$ is deleted when passing from $T$ to $T\setminus W_q(T)$, then
\[
  \epsilon_r^w(T)=0\quad\text{or}\quad \epsilon_{r+1}^w(T)=1.
\]
Equivalently, if $\epsilon_r^w(T)=1$ and $\epsilon_{r+1}^w(T)=0$, then
$(1,2r)$ remains in $T\setminus W_q(T)$.  For the final slot, if
$(1,2q)$ is deleted then $\epsilon_q^w(T)=0$; equivalently
$\epsilon_q^w(T)=1$ forces $(1,2q)$ to remain.
\end{corollary}

\begin{proof}
Write $s_r=T(1,2r-1)$ and $a_r=T(1,2r)$.  By
\cref{lem:bc-alternating-frontier-retention}, deletion of $(1,2r)$ with
$r<q$ means that $a_r=s_r$ or $a_r=s_{r+1}$.  If $a_r=s_r$, then the two
first-row copies of the label $s_r$ occupy columns $2r-1$ and $2r$, so the
frontier copy at $(1,2r-1)$ is the leftmost copy.  Hence
\cref{lem:bc-standardized-witness-pairs} gives $\epsilon_r^w(T)=0$.
If $a_r=s_{r+1}$, then the two first-row copies of the label $s_{r+1}$ occupy
columns $2r$ and $2r+1$, so the frontier copy at $(1,2r+1)$ is the rightmost
copy.  Thus $\epsilon_{r+1}^w(T)=1$ by the same lemma.

For $r=q$, \cref{lem:bc-alternating-frontier-retention} says that deletion of
$(1,2q)$ can only occur when $a_q=s_q$.  Then the frontier copy at
$(1,2q-1)$ is the leftmost copy, so $\epsilon_q^w(T)=0$.
The equivalent retention statements are the contrapositives.
\end{proof}

\begin{corollary}[Orientation bits charge width boundary deletions]\label{cor:bc-width-boundary-charge}
In the width setting of \cref{lem:bc-alternating-frontier-retention}, decompose
$S=W_q(T)$ into maximal consecutive label intervals $J_\alpha$.  Let
$s_r=T(1,2r-1)$ and let $\epsilon_r^w(T)$ be as in
\cref{lem:bc-standardized-witness-pairs}.  If the second copy of $s_r$ lies
in the first row, then the additional deleted first-row box carrying $s_r$ is
\[
  (1,2r)\quad\text{when }\epsilon_r^w(T)=0,
\]
and is
\[
  (1,2r-2)\quad\text{when }\epsilon_r^w(T)=1,
\]
with the second alternative omitted for $r=1$.  Hence every deleted even
top-row box whose label lies in $J_\alpha$ is charged to the unique
frontier-event slot of that label in $J_\alpha$, and the possible adjacent
position is determined by the first half of the local branch word.

In the height setting, no even first-column box among
$(2,1),(4,1),\ldots,(2q,1)$ is deleted by $H_q(T)$.
\end{corollary}

\begin{proof}
The first assertion is exactly the local second-copy statement in
\cref{lem:bc-standardized-witness-pairs}: the frontier copy of $s_r$ is the
leftmost copy when $\epsilon_r^w(T)=0$ and the rightmost copy when
$\epsilon_r^w(T)=1$, and the only possible first-row positions for the other
copy are the adjacent columns $2r$ and $2r-2$.  Since the witness labels in
$J_\alpha$ are distinct, each deleted even top-row box carrying a label from
$J_\alpha$ belongs to exactly one such label, and therefore to its unique
frontier-event slot.  The first half of the local branch word records these
orientation bits by \cref{lem:bc-local-branch-word} and
\cref{lem:bc-local-branch-recovery}.  The height assertion is the height half
of \cref{lem:bc-alternating-frontier-retention}.
\end{proof}

\begin{lemma}[Local position of width-witness second copies]\label{lem:bc-width-witness-second-copy}
Let $T$ be a semistandard tableau of content $(2^j)$ with
$\lambda_1(T)\ge2q$.  Put $s_r=T(1,2r-1)$ for $1\le r\le q$.
If the second copy of $s_r$ occurs in the first row, then it occurs in column
$2r-2$ or column $2r$, omitting column $0$ when $r=1$.
\end{lemma}

\begin{proof}
The first row is weakly increasing and the label $s_r$ occurs exactly twice in
the whole tableau.  Suppose first that another copy of $s_r$ occurs in the
first row at a column $c<2r-1$.  If $c\le2r-3$, then weak increase forces every
entry from column $c$ through column $2r-1$ to equal $s_r$, giving at least
three copies.  Hence $c=2r-2$.  Similarly, if another copy occurs at
$c>2r-1$ and $c\ge2r+1$, then every entry from column $2r-1$ through column
$c$ equals $s_r$, again giving at least three copies.  Hence $c=2r$.
\end{proof}

\begin{lemma}[Height-witness second copies avoid the first column]\label{lem:bc-height-witness-second-copy}
Let $T$ be a semistandard tableau of content $(2^j)$ with
$\ell(\lambda(T))\ge2q$.  Put $s_r=T(2r-1,1)$ for $1\le r\le q$.
The second copy of $s_r$ is not in the first column.
\end{lemma}

\begin{proof}
Columns in a semistandard tableau are strictly increasing, so the first-column
entries are pairwise distinct.  Since $s_r$ already occurs at $(2r-1,1)$, its
second copy cannot occur elsewhere in that column.
\end{proof}

\begin{lemma}[Local branch words determine boundary witness readouts]\label{lem:bc-local-branch-boundary-readout}
Work in the setting of \cref{lem:bc-frontier-events-per-joint} and fix a
maximal witness-label interval
\[
  J_\alpha=\{a_\alpha,\ldots,b_\alpha\}\subset S .
\]
Let $h_\alpha=|J_\alpha|$ and let
\[
  \beta_\alpha=(\beta_1,\ldots,\beta_{2h_\alpha})
\]
be the local branch word of \cref{lem:bc-local-branch-word}.  In the width
case, write $S=W_q(T)=\{s_1<\cdots<s_q\}$ and let $p_\alpha$ be determined by
$s_{p_\alpha}=a_\alpha$.  For $1\le i\le h_\alpha$, put
$r=p_\alpha+i-1$.  The deleted witness label $s_r=a_\alpha+i-1$ always
deletes the bad-boundary box
\[
  (1,2r-1).
\]
If its second copy is also in the first row, then the only additional deleted
bad-boundary box for this label is
\[
  (1,2r)\quad\text{when }\beta_i=0,
  \qquad
  (1,2r-2)\quad\text{when }\beta_i=1,
\]
with the second alternative omitted for $r=1$.

In the height case, write $S=H_q(T)=\{t_1<\cdots<t_q\}$ and let
$p_\alpha$ be determined by $t_{p_\alpha}=a_\alpha$.  For
$1\le i\le h_\alpha$, put $r=p_\alpha+i-1$.  The deleted witness label
$t_r=a_\alpha+i-1$ deletes the bad-boundary box
\[
  (2r-1,1),
\]
and its second copy is not in the first column.  Thus the bad-boundary part of
each local deleted packet uses no branch data beyond the first half of
$\beta_\alpha$ in the width case, and no branch data in the height case.
\end{lemma}

\begin{proof}
In the width case,
\cref{lem:bc-consecutive-witness-ranks} gives
$s_{p_\alpha+i-1}=a_\alpha+i-1$ for $1\le i\le h_\alpha$.  The box
$(1,2r-1)$ carries $s_r$ by the definition of $W_q(T)$, so it is deleted when
the witness labels in $S$ are deleted.  By
\cref{lem:bc-local-branch-word}, the first half of $\beta_\alpha$ is exactly
the orientation word $\epsilon^w_r(T)$ along the ordered labels of
$J_\alpha$.  Hence $\beta_i=\epsilon^w_r(T)$.  If the second copy of $s_r$ is
also in the first row, \cref{lem:bc-standardized-witness-pairs} places it at
$(1,2r)$ when $\epsilon^w_r(T)=0$ and at $(1,2r-2)$ when
$\epsilon^w_r(T)=1$, with the unavailable column $0$ omitted.  This proves the
width assertions.

The height proof is identical up to the boundary location.  Now
\cref{lem:bc-consecutive-witness-ranks} gives
$t_{p_\alpha+i-1}=a_\alpha+i-1$, and the definition of $H_q(T)$ says that
$(2r-1,1)$ carries $t_r$, so it is deleted.  The second copy of $t_r$ is not
in the first column by \cref{lem:bc-height-witness-second-copy}.  Therefore no
additional first-column boundary position has to be selected.
\end{proof}

\begin{corollary}[Bad-boundary packets need no interior branch source]\label{cor:bc-bad-boundary-packets-closed}
In the setting of \cref{lem:bc-local-branch-boundary-readout}, the data
\[
  J_\alpha,\qquad \beta_\alpha,\qquad S
\]
determine the entire bad-boundary part of the deleted local packet:
\begin{enumerate}
\item in the width case, for each label $s_r\in J_\alpha$, the forced odd
frontier box $(1,2r-1)$ is deleted, and the first half of
$\beta_\alpha$ determines whether a possible additional first-row copy of the
same label is at $(1,2r)$ or at $(1,2r-2)$;
\item in the height case, for each label $t_r\in J_\alpha$, the forced odd
frontier box $(2r-1,1)$ is deleted, and there is no additional first-column
copy to select.
\end{enumerate}
Consequently any remaining local packet theorem for the fixed-witness repair
only has to reconstruct the off-boundary/interior part of the packet and the
joint parity repair; the bad row or column frontier consumes no branch data
beyond the already allocated local word.
\end{corollary}

\begin{proof}
In the width case, \cref{lem:bc-local-branch-boundary-readout} identifies the
rank $r$ of every label in $J_\alpha$ from $S$ and the interval
$J_\alpha$, proves that $(1,2r-1)$ is always deleted, and gives the two
possible adjacent positions for a second first-row copy as the two values of
the corresponding first-half bit of $\beta_\alpha$.  In the height case, the
same lemma identifies $r$, proves that $(2r-1,1)$ is deleted, and proves that
the second copy is not in column $1$.  These are exactly the listed
bad-boundary readouts, and the final sentence is the same statement expressed
as a restriction on the remaining local packet theorem.
\end{proof}

\begin{corollary}[Boundary-supported width packets have a fixed local inverse]\label{cor:bc-boundary-supported-width-packet-inverse}
Work in the width case of \cref{lem:bc-local-branch-boundary-readout}.  Write
\[
  J_\alpha=\{a_\alpha,\ldots,b_\alpha\},\qquad
  S=W_q(T)=\{s_1<\cdots<s_q\},
\]
and let $p_\alpha$ be determined by $s_{p_\alpha}=a_\alpha$.  Suppose that,
for every $1\le i\le h_\alpha=|J_\alpha|$, the second copy of
$s_{p_\alpha+i-1}$ also lies in the first row.  Then the data
\[
  S,\quad J_\alpha,\quad \sigma^{(2a_\alpha-2)},\quad
  \beta_\alpha\in\{0,1\}^{2h_\alpha}
\]
determine the whole standardized local packet subpath
\[
  \sigma^{(2a_\alpha-2)},\sigma^{(2a_\alpha-1)},\ldots,\sigma^{(2b_\alpha)}.
\]
Consequently a width local packet with no off-boundary deleted witness copy
has no remaining packet-inverse ambiguity after the local branch word is
known.
\end{corollary}

\begin{proof}
For $1\le i\le h_\alpha$, put $r=p_\alpha+i-1$ and
$s=s_r=a_\alpha+i-1$.  By
\cref{lem:bc-consecutive-witness-ranks}, these are exactly the witness labels
in the interval $J_\alpha$, in increasing order.  By
\cref{lem:bc-local-branch-boundary-readout}, the first half of
$\beta_\alpha$ determines the orientation bit
$\epsilon_r^w(T)$.  The frontier box $(1,2r-1)$ is deleted.  Under the
boundary-supported hypothesis the other deleted box carrying $s$ is also in
the first row, and
\cref{lem:bc-standardized-witness-pairs} gives the complete standardized
placement:
\[
\begin{array}{c|c|c}
\epsilon_r^w(T) & \text{box carrying }2s-1 & \text{box carrying }2s\\ \hline
0 & (1,2r-1) & (1,2r)\\
1 & (1,2r-2) & (1,2r-1).
\end{array}
\]
The second row of the table is unavailable when $r=1$ and the displayed
column $2r-2$ is not a box; this is simply an inconsistent branch value for an
actual boundary-supported packet.

Thus the displayed data determine, for each standard time
$2s-1,2s$ with $s\in J_\alpha$, the unique box added at that time.  Starting
from the known shape $\sigma^{(2a_\alpha-2)}$ and adjoining these determined
boxes in the standard-time order
\[
  2a_\alpha-1,2a_\alpha,\ldots,2b_\alpha
\]
recovers every intermediate shape through $\sigma^{(2b_\alpha)}$.  This is
the claimed local inverse.  The final sentence restates that, in this
subcase, \cref{cor:bc-bad-boundary-packets-closed} has already recovered all
deleted witness boxes, so there is no off-boundary packet datum left to
decode.
\end{proof}

\begin{corollary}[Second-copy readout]\label{cor:bc-local-packets-reduce-to-second-copy-readout}
Work in the setting of \cref{lem:bc-local-branch-boundary-readout}.  In the
width case, write
\[
  J_\alpha=\{a_\alpha,\ldots,b_\alpha\},\qquad
  S=W_q(T)=\{s_1<\cdots<s_q\},
\]
and let $p_\alpha$ be determined by $s_{p_\alpha}=a_\alpha$.  Suppose that an
auxiliary datum $A_\alpha$ determines the subset
\[
  R_\alpha^w(T)=
  \{\,s_r\in J_\alpha:\hbox{ the second copy of }s_r
      \hbox{ is not in the first row}\,\}
\]
and, for every $s_r\in R_\alpha^w(T)$, the box occupied by this second copy.
Then
\[
  S,\quad J_\alpha,\quad \sigma^{(2a_\alpha-2)},\quad
  \beta_\alpha,\quad A_\alpha
\]
determine the whole standardized local packet subpath
\[
  \sigma^{(2a_\alpha-2)},\sigma^{(2a_\alpha-1)},\ldots,
  \sigma^{(2b_\alpha)}.
\]

In the height case, write
\[
  J_\alpha=\{a_\alpha,\ldots,b_\alpha\},\qquad
  S=H_q(T)=\{t_1<\cdots<t_q\},
\]
and let $p_\alpha$ be determined by $t_{p_\alpha}=a_\alpha$.  Suppose that an
auxiliary datum $A_\alpha$ determines, for every
$t_r\in J_\alpha$, the box occupied by the second copy of $t_r$.  Then the
same data determine the whole standardized local packet subpath.

Consequently the packet-sensitive part not supplied by the bad-boundary
readout is exactly the second-copy readout: in width, only the off-first-row
second copies have to be exposed, and in height all second copies have to be
exposed.
\end{corollary}

\begin{proof}
Consider first the width case.  For $1\le i\le h_\alpha$, put
$r=p_\alpha+i-1$ and $s=s_r=a_\alpha+i-1$.  By
\cref{lem:bc-consecutive-witness-ranks}, these are exactly the labels of
$J_\alpha$ in increasing order.  The first half of $\beta_\alpha$ is the
orientation word, so
\[
  \epsilon_r^w(T)=\beta_i
\]
by \cref{lem:bc-local-branch-boundary-readout}.  The frontier box
$(1,2r-1)$ is always deleted.  If $s_r\notin R_\alpha^w(T)$, then the second
copy is also in the first row, and
\cref{lem:bc-standardized-witness-pairs} places the adjacent standard pair as
\[
\begin{array}{c|c|c}
\epsilon_r^w(T) & \text{box carrying }2s-1 & \text{box carrying }2s\\ \hline
0 & (1,2r-1) & (1,2r)\\
1 & (1,2r-2) & (1,2r-1).
\end{array}
\]
If $s_r\in R_\alpha^w(T)$, then $A_\alpha$ supplies the off-first-row second
box, call it $c_r$.  The definition of $\epsilon_r^w(T)$ gives the remaining
standardized placement:
\[
\begin{array}{c|c|c}
\epsilon_r^w(T) & \text{box carrying }2s-1 & \text{box carrying }2s\\ \hline
0 & (1,2r-1) & c_r\\
1 & c_r & (1,2r-1).
\end{array}
\]
Thus the displayed data determine, for every $s\in J_\alpha$, the unique box
added at the two standard times $2s-1$ and $2s$.

In the height case, the same rank argument gives
$t_{p_\alpha+i-1}=a_\alpha+i-1$.  By
\cref{lem:bc-standardized-witness-pairs}, the frontier box
$(2r-1,1)$ carries $2t_r-1$, and
\cref{lem:bc-height-witness-second-copy} says the second copy is not in the
first column.  The datum $A_\alpha$ supplies its box; that box carries
$2t_r$.  Hence the data determine the unique box added at the two standard
times $2t_r-1$ and $2t_r$ for every label in $J_\alpha$.

Starting from the known left endpoint $\sigma^{(2a_\alpha-2)}$ and adjoining
the determined boxes in the order
\[
  2a_\alpha-1,2a_\alpha,\ldots,2b_\alpha
\]
recovers every intermediate shape through $\sigma^{(2b_\alpha)}$ in both
cases.  The final sentence is just the same reconstruction statement with the
already recovered bad-boundary boxes removed from the list of missing data.
\end{proof}

\begin{corollary}[Label-box readout supplies second-copy readout]\label{cor:bc-local-label-boxes-supply-second-copy-readout}
Work in the setting of \cref{lem:bc-local-branch-boundary-readout}.  In the
width case, suppose that an auxiliary datum $A_\alpha$ determines, for every
label $s\in J_\alpha$, the two boxes of $T$ carrying the semistandard label
$s$.  Then
\[
  S,\quad J_\alpha,\quad \sigma^{(2a_\alpha-2)},\quad
  \beta_\alpha,\quad A_\alpha
\]
determine the whole standardized local packet subpath.

In the height case, suppose instead that $A_\alpha$ determines, for every label
$t\in J_\alpha$, the two boxes of $T$ carrying the semistandard label $t$.
Then the same data determine the whole standardized local packet subpath.

Consequently a target-size parser which exposes the two deleted semistandard
boxes for every label in the local joint supplies the second-copy readout of
\cref{cor:bc-local-packets-reduce-to-second-copy-readout}; no further
packet-local branch source is needed for the original standardized subpath.
\end{corollary}

\begin{proof}
Consider first the width case.  Write
\[
  S=W_q(T)=\{s_1<\cdots<s_q\},
\]
and let $p_\alpha$ be determined by $s_{p_\alpha}=a_\alpha$.  For
$s_r\in J_\alpha$, the rank $r$ is determined by $S$ and $J_\alpha$, and the
frontier box is
\[
  f_r=(1,2r-1)
\]
by \cref{lem:bc-local-branch-boundary-readout}.  The datum $A_\alpha$ supplies
the two boxes carrying $s_r$; removing the known frontier box $f_r$ from this
two-element set determines the other box $c_r$.  The label $s_r$ belongs to
$R_\alpha^w(T)$ exactly when $c_r$ is not in the first row, and in that case
$A_\alpha$ determines the required off-first-row box $c_r$.  Thus
$A_\alpha$ supplies the width second-copy readout required by
\cref{cor:bc-local-packets-reduce-to-second-copy-readout}, and that corollary
recovers the local packet subpath.

The height case is the same.  Write
\[
  S=H_q(T)=\{t_1<\cdots<t_q\},
\]
and let $p_\alpha$ be determined by $t_{p_\alpha}=a_\alpha$.  For
$t_r\in J_\alpha$, the rank $r$ is determined by $S$ and $J_\alpha$, and the
frontier box is
\[
  f_r=(2r-1,1).
\]
By \cref{lem:bc-standardized-witness-pairs}, this frontier box carries the
standard entry $2t_r-1$.  Hence the other box among the two boxes supplied by
$A_\alpha$ is the box carrying $2t_r$ in the doubled standardization.  This is
exactly the height second-copy readout required by
\cref{cor:bc-local-packets-reduce-to-second-copy-readout}.  Applying that
corollary recovers the local packet subpath.  The final sentence is the same
implication stated in parser-facing language.
\end{proof}

\begin{corollary}[Step-2 interval readout supplies local packet inverses]\label{cor:bc-local-step2-interval-supplies-label-box-readout}
Work in the setting of \cref{lem:bc-local-branch-boundary-readout}, and let
\[
  J_\alpha=\{a_\alpha,\ldots,b_\alpha\}.
\]
Suppose that an auxiliary datum $A_\alpha$ determines the Step~2 Knuth boundary
interval
\[
  \lambda^{(a_\alpha-1)},\lambda^{(a_\alpha)},\ldots,\lambda^{(b_\alpha)}
\]
of the underlying semistandard tableau.  Then
\[
  S,\quad J_\alpha,\quad \sigma^{(2a_\alpha-2)},\quad
  \beta_\alpha,\quad A_\alpha
\]
determine the whole standardized local packet subpath in both the width and
height cases.

The same conclusion holds if $A_\alpha$ determines a Step~2 first-variant
Ferrers subarrangement whose forward boundary sweep contains the displayed
interval.
\end{corollary}

\begin{proof}
By \cref{lem:bc-step2-knuth-interval-label-boxes}, the displayed Step~2
boundary interval determines the two semistandard boxes carrying every label in
$J_\alpha$.  The same is true in the subarrangement form by the final assertion
of that lemma.  Applying
\cref{cor:bc-local-label-boxes-supply-second-copy-readout} gives the local
packet inverse in either family.
\end{proof}

\begin{proposition}[Step-2 interval packet readouts give halo-zone inverses]\label{prop:bc-step2-interval-packet-readouts-give-halo-zone-inverses}
Work in the setting of
\cref{prop:bc-count-compatible-halo-zone-package}.  For each deleted packet
$K_t=\{u_t,\ldots,v_t\}$ in $Z_\gamma\cap I$, let
\[
  J_t=\{c_t,\ldots,d_t\}
\]
be the corresponding maximal witness-label interval, and put $h_t=|J_t|$.
Interpret each packet in the local branch-recovery coordinates of
\cref{lem:bc-local-branch-recovery}; the packet branch word obtained from
\cref{lem:bc-zone-branch-word-splitting} is the local branch word for this
datum.

Suppose that the replacement-window decoder, from
\[
  S,\ Z_\gamma,\ V^\ast_\gamma,\ \beta_\gamma,\quad
  \rho_{2a_\gamma-2},\rho_{2b_\gamma},
\]
recovers the retained halo blocks
\[
  B_r(P)\qquad(r\in Z_\gamma\setminus I),
\]
and, for every packet $K_t$, recovers a datum $A_t$ which determines either the
Step~2 Knuth boundary interval
\[
  \lambda^{(c_t-1)},\lambda^{(c_t)},\ldots,\lambda^{(d_t)}
\]
of the semistandard tableau encoded by the local packet, or a Step~2
first-variant Ferrers subarrangement whose forward boundary sweep contains this
interval.  Then a fixed zone-local inverse recovers the whole original subpath
\[
  (\rho_{2a_\gamma-2},\rho_{2a_\gamma-1},\ldots,\rho_{2b_\gamma})
\]
from $V^\ast_\gamma$, $\beta_\gamma$, the zone endpoint pair, and
$Z_\gamma$.
\end{proposition}

\begin{proof}
By \cref{lem:bc-zone-branch-word-splitting}, the zone branch word
$\beta_\gamma$ splits, with cuts determined by $S$, into packet branch words
\[
  \beta_{K_t}\in\{0,1\}^{2h_t}.
\]
The recovered retained halo blocks and the zone endpoint pair determine
\[
  \rho_{2u_t-2},\qquad \rho_{2v_t}
\]
for every packet by \cref{lem:bc-zone-retained-blocks-expose-packet-endpoints}.
For a fixed packet $K_t$, the known interval $J_t$, the left endpoint
$\rho_{2u_t-2}$ in local coordinates, the local branch word $\beta_{K_t}$, and
the recovered Step~2 interval datum $A_t$ satisfy the hypotheses of
\cref{cor:bc-local-step2-interval-supplies-label-box-readout}.  That corollary
recovers the whole standardized local packet subpath.  Under the local
branch-recovery interpretation of the packet, this is the deleted subpath
\[
  (\rho_{2u_t-2},\rho_{2u_t-1},\ldots,\rho_{2v_t}).
\]

Thus every deleted packet subpath is recovered from the replacement-window
decoder output and the already budgeted branch word.  The packet components
$K_t$ and the retained singleton indices $Z_\gamma\setminus I$ partition
$Z_\gamma$.  Interleaving the recovered packet subpaths and retained blocks in
increasing block order reconstructs the displayed zone subpath.
\end{proof}

\begin{proposition}[Step-2 interval halo repairs give the fixed-witness bound]\label{prop:bc-step2-interval-halo-local-repair-criterion}
Fix $q\ge1$, $j\ge q$, and a $q$-element set
$S\subset\{1,\ldots,j\}$.  Let $Z_\gamma$ be the canonical halo replacement
zones of \cref{lem:bc-canonical-halo-replacement-zones}, with each zone written
as $Z_\gamma=\{a_\gamma,\ldots,b_\gamma\}$, and consider either fixed-witness
fiber $\mathcal F^w(S)$ or $\mathcal F^h(S)$ of
\cref{prop:bc-jointwise-repair-criterion}.  For each zone put
\[
  e_\gamma=|Z_\gamma\cap I(S)|.
\]
Suppose that, for every $T$ in the chosen fiber, there is a repaired path
$P^\ast(T)\in\mathcal P_{j-q}^{\mathrm{even}}$ and a parser depending only on
$S$ with the following properties.
\begin{enumerate}
\item The parser decomposes $P^\ast(T)$ into the original retained blocks
with indices outside the halo zones, unchanged and in their original order,
and local replacement windows $V^\ast_\gamma(T)$ for the zones $Z_\gamma$.
\item Each $V^\ast_\gamma(T)$ has exactly
$|Z_\gamma\setminus I(S)|$ two-step blocks.
\item For each zone there is a branch word
\[
  \beta_\gamma(T)\in\{0,1\}^{2e_\gamma}
\]
and, for every zone, the endpoint pair
$\rho_{2a_\gamma-2},\rho_{2b_\gamma}$ is the one determined from the parsed
complement blocks and the fixed endpoints of the growth-path model by
\cref{lem:bc-canonical-halo-complement-exposes-zone-endpoints}.
\item For each zone, the replacement-window decoder satisfies the hypotheses of
\cref{prop:bc-step2-interval-packet-readouts-give-halo-zone-inverses} for
$V^\ast_\gamma(T)$, $\beta_\gamma(T)$, the zone endpoint pair, and the known
interval $Z_\gamma$.
\end{enumerate}
Then the chosen fixed-witness fiber has cardinality at most
\[
  2^{2q}s_{j-q}.
\]
\end{proposition}

\begin{proof}
By \cref{lem:bc-halo-zones-frontier-packets}, for each canonical halo zone
the number $e_\gamma=|Z_\gamma\cap I(S)|$ is also the sum of the frontier-event
packet lengths used in
\cref{prop:bc-canonical-halo-local-repair-criterion}.  The endpoint pair of
each zone is determined from the parsed complement blocks by
\cref{lem:bc-canonical-halo-complement-exposes-zone-endpoints}.  For each zone,
item~(4) and
\cref{prop:bc-step2-interval-packet-readouts-give-halo-zone-inverses} give a
fixed local inverse recovering the original subpath on $Z_\gamma$ from
$V^\ast_\gamma(T)$, $\beta_\gamma(T)$, the determined endpoint pair, and the
known zone interval.  Therefore the hypotheses of
\cref{prop:bc-canonical-halo-local-repair-criterion} hold, and that proposition
gives the displayed fixed-witness bound.
\end{proof}

\begin{proposition}[Step-3 carrier packet readouts give halo-zone inverses]\label{prop:bc-step3-carrier-packet-readouts-give-halo-zone-inverses}
Work in the setting of
\cref{prop:bc-count-compatible-halo-zone-package}.  For each deleted packet
$K_t=\{u_t,\ldots,v_t\}$ in $Z_\gamma\cap I$, let
\[
  J_t=\{c_t,\ldots,d_t\}
\]
be the corresponding maximal witness-label interval, and put $h_t=|J_t|$.
Interpret each packet in the local branch-recovery coordinates of
\cref{lem:bc-local-branch-recovery}; the packet branch word obtained from
\cref{lem:bc-zone-branch-word-splitting} is the local branch word for this
datum.

Suppose that the replacement-window decoder, from
\[
  S,\ Z_\gamma,\ V^\ast_\gamma,\ \beta_\gamma,\quad
  \rho_{2a_\gamma-2},\rho_{2b_\gamma},
\]
recovers the retained halo blocks
\[
  B_r(P)\qquad(r\in Z_\gamma\setminus I).
\]
Suppose further that, for every packet $K_t$, the same decoder recovers a
Ferrers subarrangement $E_t$ of the Step~3 dual-RSK' half-square and the
partition labels on the top/right boundary of $E_t$, and identifies a
deterministic readout which, from the recovered cell entries and internal corner
labels of $E_t$, returns a Step~2 first-variant Knuth subarrangement $C_t$ and
its left/bottom boundary labels, such that every non-diagonal cell of $C_t$, or
its transpose in the Step~2 symmetric matrix, lies in $E_t$, and such that the
forward boundary sweep on $C_t$ contains
\[
  \lambda^{(c_t-1)},\lambda^{(c_t)},\ldots,\lambda^{(d_t)} .
\]
Then a fixed zone-local inverse recovers the whole original subpath
\[
  (\rho_{2a_\gamma-2},\rho_{2a_\gamma-1},\ldots,\rho_{2b_\gamma})
\]
from $V^\ast_\gamma$, $\beta_\gamma$, the zone endpoint pair, and
$Z_\gamma$.
\end{proposition}

\begin{proof}
The recovered top/right boundary of $E_t$ determines every cell entry and
internal corner label of $E_t$ by
\cref{cor:bc-parsed-dual-rsk-prime-boundary-recovers-subfilling}.  The
deterministic readout in the hypothesis therefore determines $C_t$ and its
left/bottom boundary labels.  By
\cref{lem:bc-step3-cells-are-step2-half-matrix}, the recovered Step~3
subfilling determines every non-diagonal Step~2 entry of $C_t$, using symmetry
when the transpose cell lies in $E_t$; diagonal entries are zero by
\cref{lem:bc-step2-matrix-doubled-content-degrees}.  Hence the decoder
determines a Step~2 first-variant subarrangement with known entries and known
left/bottom boundary labels whose forward boundary sweep contains
\[
  \lambda^{(c_t-1)},\lambda^{(c_t)},\ldots,\lambda^{(d_t)}
\]
for every packet $K_t$.

These are exactly the packet interval data required in
\cref{prop:bc-step2-interval-packet-readouts-give-halo-zone-inverses}.  The
retained halo blocks are recovered by hypothesis, so applying that proposition
gives the asserted fixed zone-local inverse.
\end{proof}

\begin{proposition}[Step-3 carrier halo repairs give the fixed-witness bound]\label{prop:bc-step3-carrier-halo-local-repair-criterion}
Fix $q\ge1$, $j\ge q$, and a $q$-element set
$S\subset\{1,\ldots,j\}$.  Let $Z_\gamma$ be the canonical halo replacement
zones of \cref{lem:bc-canonical-halo-replacement-zones}, with each zone written
as $Z_\gamma=\{a_\gamma,\ldots,b_\gamma\}$, and consider either fixed-witness
fiber $\mathcal F^w(S)$ or $\mathcal F^h(S)$ of
\cref{prop:bc-jointwise-repair-criterion}.  For each zone put
\[
  e_\gamma=|Z_\gamma\cap I(S)|.
\]
Suppose that, for every $T$ in the chosen fiber, there is a repaired path
$P^\ast(T)\in\mathcal P_{j-q}^{\mathrm{even}}$ and a parser depending only on
$S$ with the following properties.
\begin{enumerate}
\item The parser decomposes $P^\ast(T)$ into the original retained blocks
with indices outside the halo zones, unchanged and in their original order,
and local replacement windows $V^\ast_\gamma(T)$ for the zones $Z_\gamma$.
\item Each $V^\ast_\gamma(T)$ has exactly
$|Z_\gamma\setminus I(S)|$ two-step blocks.
\item For each zone there is a branch word
\[
  \beta_\gamma(T)\in\{0,1\}^{2e_\gamma}
\]
and, for every zone, the endpoint pair
$\rho_{2a_\gamma-2},\rho_{2b_\gamma}$ is the one determined from the parsed
complement blocks and the fixed endpoints of the growth-path model by
\cref{lem:bc-canonical-halo-complement-exposes-zone-endpoints}.
\item For each zone, the replacement-window decoder satisfies the hypotheses of
\cref{prop:bc-step3-carrier-packet-readouts-give-halo-zone-inverses} for
$V^\ast_\gamma(T)$, $\beta_\gamma(T)$, the zone endpoint pair, and the known
interval $Z_\gamma$.
\end{enumerate}
Then the chosen fixed-witness fiber has cardinality at most
\[
  2^{2q}s_{j-q}.
\]
\end{proposition}

\begin{proof}
For each zone, item~(4) and
\cref{prop:bc-step3-carrier-packet-readouts-give-halo-zone-inverses} give a
fixed local inverse recovering the original subpath on $Z_\gamma$ from
$V^\ast_\gamma(T)$, $\beta_\gamma(T)$, the endpoint pair, and the known zone
interval.  The other items are exactly the parser, target-size, and endpoint
hypotheses of \cref{prop:bc-canonical-halo-local-repair-criterion}.  Applying
that proposition gives the displayed fixed-witness bound.
\end{proof}

\begin{corollary}[Local branch words leave only segment realization]\label{cor:bc-local-branch-leaves-only-segment-realization}
Work in the setting of \cref{lem:bc-local-branch-boundary-readout}, and put
\[
  \sigma^-=\sigma^{(2a_\alpha-2)},\qquad
  \sigma^+=\sigma^{(2b_\alpha)} .
\]
The data
\[
  S,\quad J_\alpha,\quad \sigma^-,\quad \sigma^+,\quad
  \beta_\alpha\in\{0,1\}^{2|J_\alpha|}
\]
determine all discrete local packet data used by the HKKO repair:
\begin{enumerate}
\item the bad-boundary packet readout of
\cref{cor:bc-bad-boundary-packets-closed};
\item the ordered frontier-event list;
\item the paired local parity-defect slots; and
\item the selected subset of those paired defect slots.
\end{enumerate}
Consequently the remaining local packet theorem has no further event,
defect, selection, or bad-boundary data to store.  It must construct the
zero-free local segment, place the recovered selected defect slots on eligible
boundary-up steps as in \cref{lem:bc-frontier-events-supply-marks}, and prove
the off-boundary/interior packet inverse.
\end{corollary}

\begin{proof}
\Cref{cor:bc-bad-boundary-packets-closed} gives item (1) from
$S$, $J_\alpha$, and $\beta_\alpha$.  The three conclusions of
\cref{lem:bc-local-branch-recovery}, applied with the displayed endpoints
$\sigma^-$ and $\sigma^+$, give items (2), (3), and (4).  Finally,
\cref{lem:bc-frontier-events-supply-marks} states exactly what remains once
those recovered event and selection data are fixed: a zero-free segment with
the selected slots realized at eligible boundary-up steps.  The off-boundary
packet inverse is the part not covered by the bad-boundary readout.
\end{proof}

\begin{corollary}[Below-wall packets discharge the HKKO segment leaf]\label{cor:bc-below-wall-discharges-segment-realization}
Work in the setting of
\cref{cor:bc-local-branch-leaves-only-segment-realization}.  Let
$\mu,\nu$ be an endpoint pair for the local deleted packet, put
$h=|J_\alpha|$, and set
\[
  w=w(\mu,\nu)
\]
as in \cref{lem:bc-local-endpoints-canonical-wall}.  Then
\[
  \mu,\quad \nu,\quad J_\alpha,\quad \beta_\alpha
\]
canonically determine a below-wall $w$-up-down packet segment
$W(\mu,\nu,\beta_\alpha)$, the selected HKKO event indices corresponding to
the selected paired-defect subset recovered from $\beta_\alpha$, and the
unmarked HKKO boundary-event output
$W^\ast(\mu,\nu,\beta_\alpha)$.  The marked preimage
$W(\mu,\nu,\beta_\alpha)$ is recoverable from
\[
  W^\ast(\mu,\nu,\beta_\alpha),\quad
  \beta_\alpha,\quad w,\quad h,\quad \mu .
\]
After the coefficient extraction of \cref{lem:bc-hkko-updown-boundary-events},
the selected event indices are eligible boundary-up marks in the
vacillating model of \cref{lem:bc-hkko-marked-boundary}.

Consequently, once a target-size parser has recovered the packet endpoint
pair, the local theorem no longer has to construct a separate HKKO segment,
marking injection, or wall parameter.  The remaining packet-specific input is
a target-size parser/merger and fixed off-boundary/interior readout that
recovers the original deleted packet subpath from the recovered target-window
data and the canonical marked packet preimage.
\end{corollary}

\begin{proof}
\Cref{cor:bc-local-branch-leaves-only-segment-realization} and
\cref{lem:bc-local-branch-recovery} recover the selected subset
$M_\alpha$ of paired defect slots from the second half of
$\beta_\alpha$.  By \cref{lem:bc-local-endpoints-canonical-wall}, the displayed
integer $w=w(\mu,\nu)$ satisfies
\[
  \ell(\mu)<w,\qquad \ell(\nu)<w,
\]
and is determined by the endpoint pair.  Therefore
\cref{prop:bc-below-wall-local-hkko-splice-package} applies with this
endpoint pair, $h=h_\alpha$, and the branch word $\beta_\alpha$.  It supplies
the canonical assigned event-index list, the below-wall segment
$W(\mu,\nu,\beta_\alpha)$ realizing the subset $M_\alpha$, the unmarked
boundary-event splice output $W^\ast(\mu,\nu,\beta_\alpha)$, and the displayed
recovery of the marked preimage.

\Cref{lem:bc-hkko-updown-boundary-events} identifies the corresponding
coefficient-extracted marked-vacillating events with eligible boundary-up
steps, and \cref{lem:bc-hkko-marked-boundary} records the same admissibility
condition in the marked vacillating class.  Thus the segment and marking part
of the local packet theorem is supplied by the cited below-wall package after
the packet endpoints are known.  What is not supplied by that package is a
target-size replacement-window decoder or a readout of the original
off-boundary packet subpath; those are exactly the remaining inputs stated in
the consequence.
\end{proof}

\begin{proposition}[Bad-shape fiber reduction]\label{prop:bc-bad-shape-fiber-reduction}
Fix residual $q=b+1$.  Suppose that, for every $j$ in the residual window,
every $q$-element subset $S\subset\{1,\ldots,j\}$ satisfies the two fiber
estimates
\[
 \#\{T:\lambda(T)' \text{ even},\ \lambda_1(T)\ge2q,\ W_q(T)=S\}
 \le 2^{2q} s_{j-q},
\]
and
\[
 \#\{T:\lambda(T)' \text{ even},\ \ell(\lambda(T))\ge2q,\ H_q(T)=S\}
 \le 2^{2q} s_{j-q}.
\]
Then
\[
  \sum_{\substack{\lambda':\,\mathrm{even}\\ \lambda_1\ge2q}}a_j(\lambda)
  \le 2^{2q}\binom jq s_{j-q},
  \qquad
  \sum_{\substack{\lambda':\,\mathrm{even}\\ \ell(\lambda)\ge2q}}a_j(\lambda)
  \le 2^{2q}\binom jq s_{j-q}.
\]
\end{proposition}

\begin{proof}
By \cref{lem:bad-shape-witnesses}, each width-bad tableau has exactly one
canonical witness $W_q(T)$, and each height-bad tableau has exactly one
canonical witness $H_q(T)$.  Partitioning the two bad sets by that witness
and applying the displayed fiber estimates gives
\[
  \sum_{\substack{\lambda':\,\mathrm{even}\\ \lambda_1\ge2q}}a_j(\lambda)
  \le
  \sum_{\substack{S\subset\{1,\ldots,j\}\\ |S|=q}}2^{2q} s_{j-q}
  =
  2^{2q}\binom jq s_{j-q},
\]
and the height estimate is identical with $H_q$ in place of $W_q$.
\end{proof}

\begin{proposition}[The $2^q$ fixed-height-fiber target is false]\label{prop:bc-fixed-witness-qbit-obstruction}
Let $q=15$, $j=30$, and
\[
  S=\{1,3,5,\ldots,29\}.
\]
Then the height fixed-witness fiber satisfies
\[
\#\{T:\lambda(T)' \text{ even},\ \ell(\lambda(T))\ge30,\ H_{15}(T)=S\}
  =
  6190283353629375,
\]
whereas
\[
  2^{15}s_{15}=4825650970886144.
\]
In particular, the estimate
\[
  \#\{T:\lambda(T)' \text{ even},\ \ell(\lambda(T))\ge2q,\ H_q(T)=S\}
  \le 2^q s_{j-q}
\]
is false even at the first residual boundary $q=15$, $j=30$.
\end{proposition}

\begin{proof}
Write the displayed height fiber as $\mathcal F^h(S)$.
By \cref{cor:bc-fixed-fibers-pieri-paths}, the displayed fiber is
equinumerous with the column-even Pieri two-strip paths of length $30$ with
crossing set $C^h_{15}=S$.  The exact GMP/OpenMP obstruction replay
\[
\texttt{certificates/post\_m29/bc\_fixed\_witness\_2q\_obstruction\_gmp\_certificate.log}
\]
enumerates this path fiber and independently recomputes $s_{15}$.  It runs on
\texttt{optimus} with $64$ OpenMP threads, uses exact \texttt{mpz\_class}
arithmetic, reports
\[
\begin{aligned}
s_{15}&=147267180508,\\
2^{15}s_{15}&=4825650970886144,\\
\#\mathcal F^h(S)&=6190283353629375,
\end{aligned}
\]
and verifies \texttt{height\_exceeds\_2powq\_rhs=yes},
\texttt{SUMMARY failures=0}, and \texttt{\_\_EXIT\_STATUS=0}.  Its SHA-256 is
\[
\hashsplit{cc303f3bffd77ba4d8954054c31bc99c}{cd5c406f08e3a397ca0f37b759398436}.
\]
The final displayed inequality is therefore contradicted by this exact
integer count.
\end{proof}

\begin{lemma}[$B$-boundary width-bad paths have a three-step encoding]\label{lem:bc-b-boundary-width-bad-three-step}
Let $q\ge1$, and let $\lambda^\bullet$ be a column-even Pieri two-strip path
of length $2q$ such that
\[
  \lambda^{(2q)}_1\ge2q .
\]
Then
\[
  \lambda^{(2q)}=(2q,2q),
\]
and the number of such paths is at most \(3^{2q}\).
\end{lemma}

\begin{proof}
The terminal partition has size \(4q\).  Since
\((\lambda^{(2q)})'\) has only even parts, every nonzero column has height at
least \(2\).  Hence
\[
  4q=|\lambda^{(2q)}|\ge2\lambda^{(2q)}_1\ge4q .
\]
Thus \(\lambda^{(2q)}_1=2q\), every nonzero column has height exactly \(2\),
and the terminal partition is \((2q,2q)\).

Every intermediate partition is contained in this two-row rectangle.  At each
Pieri step, the horizontal two-strip can therefore add its two boxes in only
one of three row patterns: two boxes in the first row, one box in each row, or
two boxes in the second row.  Sending a path to this word in a three-letter
alphabet is injective, because the current partition and the row pattern
determine the next partition.  Hence there are at most \(3^{2q}\) such paths.
\end{proof}

\begin{lemma}[Stable moments dominate the boundary width count]\label{lem:bc-stable-dominates-boundary-width-count}
For every \(q\ge15\),
\[
  2\cdot3^{2q}\le s_{2q}.
\]
\end{lemma}

\begin{proof}
By \cref{lem:bc-littlewood-graph-model}, \(s_{2q}\) counts loopless labelled
degree-two multigraphs on \(2q\) vertices.  It therefore contains all labelled
simple cycles, whose number is
\[
  \frac{(2q-1)!}{2}.
\]
It is enough to prove
\[
  4\cdot3^{2q}\le(2q-1)! .
\]
At \(q=15\),
\[
  4\cdot3^{30}=823564528378596
  <
  8841761993739701954543616000000
  =
  29! .
\]
If the displayed factorial inequality holds at \(q\), then at \(q+1\) the
left side is multiplied by \(9\), while the right side is multiplied by
\[
  (2q)(2q+1)>9 .
\]
Thus induction proves the inequality for every \(q\ge15\).
\end{proof}

\begin{proposition}[Residual $B$-boundary lower correction]\label{prop:bc-b-boundary-lower-correction}
For every residual row \(B_b\) with \(14\le b\le123\), put \(q=b+1\) and
write \(s_j=m_j^{\mathrm{st}}\) and \(m_j^{B_b}=s_j+\delta_j^{B_b}\).  Then
\[
  -\frac12 s_{2q}\le \delta_{2q}^{B_b}.
\]
\end{proposition}

\begin{proof}
By \cref{lem:bc-correction-bad-shapes} and the coefficient interpretation in
\cref{lem:bad-shape-witnesses}, \(|\delta_{2q}^{B_b}|\) is the number of
semistandard tableaux \(T\) of content \((2^{2q})\) whose shape
\(\lambda(T)\) has \(\lambda(T)'\) even and \(\lambda_1(T)\ge2q\).  Under
\cref{lem:bc-pieri-path-bijection}, these tableaux are exactly the
column-even Pieri two-strip paths of length \(2q\) with terminal first row at
least \(2q\).  By \cref{lem:bc-b-boundary-width-bad-three-step}, their number
is at most \(3^{2q}\).  Since \(q\ge15\),
\cref{lem:bc-stable-dominates-boundary-width-count} gives
\[
  |\delta_{2q}^{B_b}|\le3^{2q}\le\frac12s_{2q}.
\]
This proves the stated lower bound.
\end{proof}

\begin{lemma}[Stable moments dominate the first $B$ nonboundary width count]\label{lem:bc-stable-dominates-first-b-nonboundary-width-count}
For every \(q\ge15\),
\[
  2\cdot10^{2q+1}\le s_{2q+1}.
\]
\end{lemma}

\begin{proof}
By \cref{lem:bc-littlewood-graph-model}, \(s_{2q+1}\) counts loopless
labelled degree-two multigraphs on \(2q+1\) vertices.  It therefore contains
all labelled simple cycles, whose number is
\[
  \frac{(2q)!}{2}.
\]
It is enough to prove
\[
  4\cdot10^{2q+1}\le(2q)! .
\]
At \(q=15\),
\[
  4\cdot10^{31}
  <
  265252859812191058636308480000000
  =
  30! .
\]
If the displayed factorial inequality holds at \(q\), then at \(q+1\) the
left side is multiplied by \(100\), while the right side is multiplied by
\[
  (2q+1)(2q+2)\ge31\cdot32>100 .
\]
Thus induction proves the inequality for every \(q\ge15\).
\end{proof}

\begin{lemma}[First $B$ nonboundary width-bad paths have a ten-step encoding]\label{lem:bc-b-first-nonboundary-width-bad-ten-step}
Let \(q\ge1\), and let \(\lambda^\bullet\) be a column-even Pieri two-strip
path of length \(2q+1\) such that
\[
  \lambda^{(2q+1)}_1\ge2q .
\]
Then the terminal partition is either
\[
  (2q+1,2q+1)
  \qquad\text{or}\qquad
  (2q,2q,1,1),
\]
and the number of such paths is at most \(10^{2q+1}\).
\end{lemma}

\begin{proof}
The terminal partition has size \(4q+2\).  Since
\((\lambda^{(2q+1)})'\) has only even parts, every nonzero column has height
at least \(2\).  Hence the terminal first row is either \(2q\) or \(2q+1\).
If it is \(2q+1\), all columns have height exactly \(2\), so the terminal
partition is \((2q+1,2q+1)\).  If it is \(2q\), the columns of height at least
\(2\) account for \(4q\) boxes.  The remaining two boxes must form one more
even column-height contribution, hence the terminal partition is
\((2q,2q,1,1)\).

Thus every intermediate partition has at most four rows.  At each Pieri step,
record the unordered pair of rows, with repetition allowed, receiving the two
new boxes.  There are \(\binom{4+1}{2}=10\) possible row patterns, and the
current partition together with this row pattern determines the next partition:
the boxes are added at the right ends of the selected rows.  Therefore the map
from paths to row-pattern words is injective, giving at most \(10^{2q+1}\)
paths.
\end{proof}

\begin{proposition}[Residual $B$ first nonboundary lower correction]\label{prop:bc-b-first-nonboundary-lower-correction}
For every residual row \(B_b\) with \(14\le b\le123\), put \(q=b+1\) and
write \(s_j=m_j^{\mathrm{st}}\) and \(m_j^{B_b}=s_j+\delta_j^{B_b}\).  Then
\[
  -\frac12 s_{2q+1}\le \delta_{2q+1}^{B_b}.
\]
\end{proposition}

\begin{proof}
By \cref{lem:bc-correction-bad-shapes} and the coefficient interpretation in
\cref{lem:bad-shape-witnesses}, \(|\delta_{2q+1}^{B_b}|\) is the number of
semistandard tableaux \(T\) of content \((2^{2q+1})\) whose shape
\(\lambda(T)\) has \(\lambda(T)'\) even and \(\lambda_1(T)\ge2q\).  Under
\cref{lem:bc-pieri-path-bijection}, these tableaux are exactly the
column-even Pieri two-strip paths of length \(2q+1\) with terminal first row
at least \(2q\).  By \cref{lem:bc-b-first-nonboundary-width-bad-ten-step},
their number is at most \(10^{2q+1}\).  Since \(q\ge15\),
\cref{lem:bc-stable-dominates-first-b-nonboundary-width-count} gives
\[
  |\delta_{2q+1}^{B_b}|\le10^{2q+1}\le\frac12s_{2q+1}.
\]
This proves the stated lower bound.
\end{proof}

\begin{lemma}[Second $B$ nonboundary width-bad paths have sparse exceptional rows]\label{lem:bc-b-second-nonboundary-width-bad-sparse-row-code}
Let \(q\ge1\), and let \(\lambda^\bullet\) be a column-even Pieri two-strip
path of length \(2q+2\) such that
\[
  \lambda^{(2q+2)}_1\ge2q .
\]
Then the number of such paths is at most
\[
  5(2q+2)^4\,21^4\,3^{2q+2}.
\]
\end{lemma}

\begin{proof}
The terminal partition has size \(4q+4\).  Since
\((\lambda^{(2q+2)})'\) has only even parts, every nonzero column has height
at least \(2\).  Hence the terminal first row \(W=\lambda^{(2q+2)}_1\)
satisfies
\[
  2q\le W\le2q+2 .
\]
The number of terminal boxes below the second row is
\[
  |\lambda^{(2q+2)}|-2W\le4 .
\]
Thus the terminal partition has at most six nonzero rows, and at most four
Pieri steps add a box below the second row.

Encode a path as follows.  First record the exceptional set \(E\) of steps
which add at least one box below the second row.  It has size at most \(4\).
For a nonexceptional step, record which unordered pair of rows among
\(\{1,2\}\), with repetition allowed, receives the two boxes; there are
three choices.  For an exceptional step, record the unordered pair of rows
among the at most six terminal rows; there are
\(\binom{6+1}{2}=21\) choices.  The current partition together with this row
pair determines the next partition, because the new boxes are added at the
right ends of the selected rows.  Therefore the encoding is injective.

For \(n=2q+2\), the number of codes is at most
\[
  \sum_{e=0}^4 \binom ne 3^{n-e}21^e
  \le
  5n^4\,21^4\,3^n,
\]
which is the displayed bound.
\end{proof}

\begin{lemma}[Stable moments dominate the second $B$ nonboundary width count]\label{lem:bc-stable-dominates-second-b-nonboundary-width-count}
For every \(q\ge15\),
\[
  10(2q+2)^4\,21^4\,3^{2q+2}\le s_{2q+2}.
\]
\end{lemma}

\begin{proof}
By \cref{lem:bc-littlewood-graph-model}, \(s_{2q+2}\) counts loopless labelled
degree-two multigraphs on \(2q+2\) vertices.  It therefore contains all
labelled simple cycles, whose number is
\[
  \frac{(2q+1)!}{2}.
\]
It is enough to prove
\[
  20(2q+2)^4\,21^4\,3^{2q+2}\le(2q+1)! .
\]
At \(q=15\),
\[
\begin{aligned}
  20\cdot32^4\cdot21^4\cdot3^{32}
\ &=\
  7557658063102958937487441920,\\
  31!
\ &=\
  8222838654177922817725562880000000 .
\end{aligned}
\]
The first integer is smaller than the second.
If the displayed factorial inequality holds at \(q\), then at \(q+1\) the
left side is multiplied by
\[
  9\biggl(\frac{2q+4}{2q+2}\biggr)^4<12,
\]
while the right side is multiplied by
\[
  (2q+2)(2q+3)\ge32\cdot33>12 .
\]
Thus induction proves the inequality for every \(q\ge15\).
\end{proof}

\begin{proposition}[Residual $B$ second nonboundary lower correction]\label{prop:bc-b-second-nonboundary-lower-correction}
For every residual row \(B_b\) with \(14\le b\le123\), put \(q=b+1\) and
write \(s_j=m_j^{\mathrm{st}}\) and \(m_j^{B_b}=s_j+\delta_j^{B_b}\).  Then
\[
  -\frac12 s_{2q+2}\le \delta_{2q+2}^{B_b}.
\]
\end{proposition}

\begin{proof}
By \cref{lem:bc-correction-bad-shapes} and the coefficient interpretation in
\cref{lem:bad-shape-witnesses}, \(|\delta_{2q+2}^{B_b}|\) is the number of
semistandard tableaux \(T\) of content \((2^{2q+2})\) whose shape
\(\lambda(T)\) has \(\lambda(T)'\) even and \(\lambda_1(T)\ge2q\).  Under
\cref{lem:bc-pieri-path-bijection}, these tableaux are exactly the
column-even Pieri two-strip paths of length \(2q+2\) with terminal first row
at least \(2q\).  By
\cref{lem:bc-b-second-nonboundary-width-bad-sparse-row-code}, their number is
at most
\[
  5(2q+2)^4\,21^4\,3^{2q+2}.
\]
Since \(q\ge15\),
\cref{lem:bc-stable-dominates-second-b-nonboundary-width-count} gives
\[
  |\delta_{2q+2}^{B_b}|
  \le
  5(2q+2)^4\,21^4\,3^{2q+2}
  \le
  \frac12s_{2q+2}.
\]
This proves the stated lower bound.
\end{proof}

\begin{lemma}[Third $B$ nonboundary width-bad paths have sparse exceptional rows]\label{lem:bc-b-third-nonboundary-width-bad-sparse-row-code}
Let \(q\ge1\), and let \(\lambda^\bullet\) be a column-even Pieri two-strip
path of length \(2q+3\) such that
\[
  \lambda^{(2q+3)}_1\ge2q .
\]
Then the number of such paths is at most
\[
  \frac7{720}(2q+3)^6\,12^6\,3^{2q+3}.
\]
\end{lemma}

\begin{proof}
Put \(n=2q+3\).  The terminal partition has size \(2n=4q+6\).  Since its
columns have even lengths and its first row has length at least \(2q\), the
number of boxes below the second row is at most
\[
  4q+6-2(2q)=6 .
\]
Hence the terminal partition has at most eight nonzero rows, and at most six
Pieri steps add a box below the second row.

As in \cref{lem:bc-b-second-nonboundary-width-bad-sparse-row-code}, record
the exceptional set \(E\) of below-second-row steps.  At nonexceptional steps
there are three unordered row-pair choices in the first two rows.  At
exceptional steps there are at most \(\binom{8+1}{2}=36\) unordered row-pair
choices.  This row-pair word reconstructs the path recursively, so the
encoding is injective.  Therefore the number of paths is at most
\[
  \sum_{e=0}^6\binom ne3^{n-e}36^e
  =
  3^n\sum_{e=0}^6\binom ne12^e .
\]
Since \(\binom ne\le n^e/e!\), it remains to bound the elementary sum with
these larger summands.  For \(0\le e\le6\), the terms
\[
  \frac{n^e12^e}{e!}
\]
increase with \(e\), since \(12n/(e+1)>1\).  Thus the last display is at most
\[
  3^n\cdot 7\,\frac{n^6 12^6}{6!}
  =
  \frac7{720}n^6\,12^6\,3^n,
\]
as claimed.
\end{proof}

\begin{lemma}[Stable moments dominate the third $B$ nonboundary width count]\label{lem:bc-stable-dominates-third-b-nonboundary-width-count}
For every \(q\ge15\),
\[
  \frac7{360}(2q+3)^6\,12^6\,3^{2q+3}\le s_{2q+3}.
\]
\end{lemma}

\begin{proof}
By \cref{lem:bc-littlewood-graph-model}, \(s_{2q+3}\) contains all labelled
simple cycles on \(2q+3\) vertices, whose number is \((2q+2)!/2\).  It is
enough to prove
\[
  7(2q+3)^6\,12^6\,3^{2q+3}\le180(2q+2)! .
\]
At \(q=15\),
\[
\begin{aligned}
  7\cdot33^6\cdot12^6\cdot3^{33}
\ &=\
  150061941592511735750704505081856,\\
  180\cdot32!
\ &=\
  47363550648064835430099242188800000000 .
\end{aligned}
\]
The first integer is smaller than the second.
If the displayed factorial inequality holds at \(q\), then at \(q+1\) the
left side is multiplied by
\[
  9\biggl(\frac{2q+5}{2q+3}\biggr)^6<13,
\]
while the right side is multiplied by
\[
  (2q+3)(2q+4)\ge33\cdot34>13 .
\]
Induction proves the result for every \(q\ge15\).
\end{proof}

\begin{proposition}[Residual $B$ third nonboundary lower correction]\label{prop:bc-b-third-nonboundary-lower-correction}
For every residual row \(B_b\) with \(14\le b\le123\), put \(q=b+1\) and
write \(s_j=m_j^{\mathrm{st}}\) and \(m_j^{B_b}=s_j+\delta_j^{B_b}\).  Then
\[
  -\frac12 s_{2q+3}\le \delta_{2q+3}^{B_b}.
\]
\end{proposition}

\begin{proof}
By \cref{lem:bc-correction-bad-shapes}, the coefficient interpretation in
\cref{lem:bad-shape-witnesses}, and the Pieri-path bijection
\cref{lem:bc-pieri-path-bijection}, \(|\delta_{2q+3}^{B_b}|\) is the number
of column-even Pieri two-strip paths of length \(2q+3\) with terminal first
row at least \(2q\).  By
\cref{lem:bc-b-third-nonboundary-width-bad-sparse-row-code}, this number is at
most
\[
  \frac7{720}(2q+3)^6\,12^6\,3^{2q+3}.
\]
Since \(q\ge15\),
\cref{lem:bc-stable-dominates-third-b-nonboundary-width-count} gives
\[
  |\delta_{2q+3}^{B_b}|\le\frac12s_{2q+3}.
\]
This proves the stated lower bound.
\end{proof}

\begin{lemma}[Fourth $B$ nonboundary width-bad paths have sparse exceptional rows]\label{lem:bc-b-fourth-nonboundary-width-bad-sparse-row-code}
Let \(q\ge1\), and let \(\lambda^\bullet\) be a column-even Pieri two-strip
path of length \(2q+4\) such that
\[
  \lambda^{(2q+4)}_1\ge2q .
\]
Then the number of such paths is at most
\[
  \frac9{8!}(2q+4)^8\,55^8\,3^{2q-4}.
\]
\end{lemma}

\begin{proof}
Put \(n=2q+4\).  The terminal partition has size \(2n=4q+8\).  Since its
columns have even lengths and its first row has length at least \(2q\), the
number of boxes below the second row is at most
\[
  4q+8-2(2q)=8 .
\]
Hence the terminal partition has at most ten nonzero rows, and at most eight
Pieri steps add a box below the second row.

Record the exceptional set \(E\) of below-second-row steps.  At a
nonexceptional step there are three unordered row-pair choices in the first
two rows.  At an exceptional step there are at most
\(\binom{10+1}{2}=55\) unordered row-pair choices among terminal rows.  The
current partition and the recorded row pair determine the next partition, so
the encoding is injective.  Therefore the number of paths is at most
\[
  \sum_{e=0}^8\binom ne3^{n-e}55^e
  =
  3^n\sum_{e=0}^8\binom ne\biggl(\frac{55}{3}\biggr)^e .
\]
Using \(\binom ne\le n^e/e!\), the \(e\)-th summand is at most
\[
  \frac{n^e55^e3^{n-e}}{e!}.
\]
These upper summands increase for \(0\le e<8\), because
\[
  \frac{55n}{3(e+1)}>1.
\]
Thus the sum is bounded by nine times its final upper summand, giving
\[
  \frac9{8!}n^8\,55^8\,3^{n-8}
  =
  \frac9{8!}(2q+4)^8\,55^8\,3^{2q-4}.
\]
\end{proof}

\begin{lemma}[Stable moments dominate the fourth $B$ nonboundary width count]\label{lem:bc-stable-dominates-fourth-b-nonboundary-width-count}
For every \(q\ge15\),
\[
  \frac{18}{8!}(2q+4)^8\,55^8\,3^{2q-4}\le s_{2q+4}.
\]
\end{lemma}

\begin{proof}
By \cref{lem:bc-littlewood-graph-model}, \(s_{2q+4}\) contains all labelled
simple cycles on \(2q+4\) vertices, whose number is \((2q+3)!/2\).  It is
enough to prove
\[
  18(2q+4)^8\,55^8\,3^{2q-4}\le \frac{8!}{2}(2q+3)! .
\]
At \(q=15\),
\[
\begin{aligned}
  18\cdot34^8\cdot55^8\cdot3^{26}
  &=
  6841604740407802210384021379536200000000,\\
  \frac{8!}{2}\,33!
  &=
  175055683195247631749646799129804800000000 .
\end{aligned}
\]
The first integer is smaller than the second.
If the displayed factorial inequality holds at \(q\), then at \(q+1\) the
left side is multiplied by
\[
  9\biggl(\frac{2q+6}{2q+4}\biggr)^8<16,
\]
while the right side is multiplied by
\[
  (2q+4)(2q+5)\ge34\cdot35>16 .
\]
Induction proves the result for every \(q\ge15\).
\end{proof}

\begin{proposition}[Residual $B$ fourth nonboundary lower correction]\label{prop:bc-b-fourth-nonboundary-lower-correction}
For every residual row \(B_b\) with \(14\le b\le123\), put \(q=b+1\) and
write \(s_j=m_j^{\mathrm{st}}\) and \(m_j^{B_b}=s_j+\delta_j^{B_b}\).  Then
\[
  -\frac12 s_{2q+4}\le \delta_{2q+4}^{B_b}.
\]
\end{proposition}

\begin{proof}
By \cref{lem:bc-correction-bad-shapes}, the coefficient interpretation in
\cref{lem:bad-shape-witnesses}, and the Pieri-path bijection
\cref{lem:bc-pieri-path-bijection}, \(|\delta_{2q+4}^{B_b}|\) is the number
of column-even Pieri two-strip paths of length \(2q+4\) with terminal first
row at least \(2q\).  By
\cref{lem:bc-b-fourth-nonboundary-width-bad-sparse-row-code}, this number is
at most
\[
  \frac9{8!}(2q+4)^8\,55^8\,3^{2q-4}.
\]
Since \(q\ge15\),
\cref{lem:bc-stable-dominates-fourth-b-nonboundary-width-count} gives
\[
  |\delta_{2q+4}^{B_b}|\le\frac12s_{2q+4}.
\]
This proves the stated lower bound.
\end{proof}

\begin{lemma}[Fifth $B$ nonboundary width-bad paths have a lower-box code]\label{lem:bc-b-fifth-nonboundary-width-bad-lower-box-code}
For every \(q\ge15\), the number of column-even Pieri two-strip paths
\[
  \varnothing=\lambda^{(0)}\subset\lambda^{(1)}\subset\cdots
  \subset\lambda^{(2q+5)}
\]
with
\[
  \lambda^{(2q+5)}_1\ge2q
\]
is at most
\[
  36864\binom{2q+5}{10}20^{10}3^{2q-5}.
\]
\end{lemma}

\begin{proof}
Put \(n=2q+5\).  The terminal partition has size \(2n=4q+10\).  Since its
columns have even lengths and its first row has length at least \(2q\), the
number of boxes below the second row is at most
\[
  4q+10-2(2q)=10 .
\]
Hence the terminal partition has at most twelve nonzero rows.

Encode a path by its row-pair word, but keep track of the number of boxes put
below row \(2\).  A step with no such box has one of the three top-row pair
types.  A step with exactly one such box has at most \(2\cdot10=20\) choices:
one top row and one of the ten possible lower terminal rows.  A step with two
such boxes has at most \(\binom{10+1}{2}=55\) choices.  The current partition
and this row-pair type reconstruct the next partition, so the encoding is
injective.  Thus the number of paths is bounded by the sum of the coefficients
of \(x^0,\ldots,x^{10}\) in
\[
  (3+20x+55x^2)^n .
\]

For a monomial with \(a\) one-lower-box steps and \(c\) two-lower-box steps,
\(a+2c\le10\), its contribution is
\[
  \binom{n}{a+c}\binom{a+c}{c}3^{n-a-c}20^a55^c .
\]
Since \(n\ge35\) and \(a+c\le10\),
\[
  \binom{n}{a+c}\le\binom n{10},
  \qquad
  \binom{a+c}{c}\le2^{10}.
\]
Also \(55\le20^2/3\), so
\[
  3^{n-a-c}20^a55^c
  \le
  20^{a+2c}3^{n-a-2c}
  \le
  20^{10}3^{n-10}.
\]
There are \(36\) pairs \((a,c)\) with \(a,c\ge0\) and \(a+2c\le10\).
Multiplying these bounds gives
\[
  36\cdot2^{10}\binom n{10}20^{10}3^{n-10}
  =
  36864\binom{2q+5}{10}20^{10}3^{2q-5}.
\]
\end{proof}

\begin{lemma}[Stable moments dominate the fifth $B$ nonboundary width count]\label{lem:bc-stable-dominates-fifth-b-nonboundary-width-count}
For every \(q\ge15\),
\[
  73728\binom{2q+5}{10}20^{10}3^{2q-5}\le s_{2q+5}.
\]
\end{lemma}

\begin{proof}
By \cref{lem:bc-littlewood-graph-model}, \(s_{2q+5}\) contains all labelled
simple cycles on \(2q+5\) vertices, whose number is \((2q+4)!/2\).  It is
enough to prove
\[
  147456\binom{2q+5}{10}20^{10}3^{2q-5}\le (2q+4)! .
\]
At \(q=15\),
\[
\begin{aligned}
  147456\binom{35}{10}20^{10}3^{25}
  &=
  234864679708823602547990200320000000000,\\
  34!
  &=
  295232799039604140847618609643520000000 .
\end{aligned}
\]
The first integer is smaller than the second.
If the displayed factorial inequality holds at \(q\), then at \(q+1\) the
left side is multiplied by
\[
  9\frac{(2q+7)(2q+6)}{(2q-3)(2q-4)}<18,
\]
while the right side is multiplied by
\[
  (2q+5)(2q+6)\ge35\cdot36>18 .
\]
Induction proves the result for every \(q\ge15\).
\end{proof}

\begin{proposition}[Residual $B$ fifth nonboundary lower correction]\label{prop:bc-b-fifth-nonboundary-lower-correction}
For every residual row \(B_b\) with \(14\le b\le123\), put \(q=b+1\) and
write \(s_j=m_j^{\mathrm{st}}\) and \(m_j^{B_b}=s_j+\delta_j^{B_b}\).  Then
\[
  -\frac12 s_{2q+5}\le \delta_{2q+5}^{B_b}.
\]
\end{proposition}

\begin{proof}
By \cref{lem:bc-correction-bad-shapes}, the coefficient interpretation in
\cref{lem:bad-shape-witnesses}, and the Pieri-path bijection
\cref{lem:bc-pieri-path-bijection}, \(|\delta_{2q+5}^{B_b}|\) is the number
of column-even Pieri two-strip paths of length \(2q+5\) with terminal first
row at least \(2q\).  By
\cref{lem:bc-b-fifth-nonboundary-width-bad-lower-box-code}, this number is at
most
\[
  36864\binom{2q+5}{10}20^{10}3^{2q-5}.
\]
Since \(q\ge15\),
\cref{lem:bc-stable-dominates-fifth-b-nonboundary-width-count} gives
\[
  |\delta_{2q+5}^{B_b}|\le\frac12s_{2q+5}.
\]
This proves the stated lower bound.
\end{proof}

\begin{lemma}[Sixth $B$ nonboundary width-bad paths have a lower-box code]\label{lem:bc-b-sixth-nonboundary-width-bad-lower-box-code}
For every \(q\ge15\), the number of column-even Pieri two-strip paths
\[
  \varnothing=\lambda^{(0)}\subset\lambda^{(1)}\subset\cdots
  \subset\lambda^{(2q+6)}
\]
with
\[
  \lambda^{(2q+6)}_1\ge2q
\]
is at most
\[
  6\binom{2q+6}{12}24^{12}3^{2q-6}.
\]
\end{lemma}

\begin{proof}
Put \(n=2q+6\).  The terminal partition has size \(2n=4q+12\).  Since its
columns have even lengths and its first row has length at least \(2q\), the
number of boxes below the second row is at most
\[
  4q+12-2(2q)=12 .
\]
Hence the terminal partition has at most fourteen nonzero rows.

Encode a path by its row-pair word, keeping the number of boxes placed below
row \(2\).  A step with no lower box has one of three top-row choices.  A step
with exactly one lower box has at most \(2\cdot12=24\) choices.  A step with
two lower boxes has at most \(\binom{12+1}{2}=78\) choices.  The current
partition and the row-pair choice recover the next partition, so the encoding
is injective.  Thus the number of paths is bounded by the sum of the
coefficients of \(x^0,\ldots,x^{12}\) in
\[
  (3+24x+78x^2)^n .
\]

The monomial with \(a\) one-lower-box steps and \(c\) two-lower-box steps,
where \(a+2c\le12\), contributes
\[
  \binom{n}{a+c}\binom{a+c}{c}3^{n-a-c}24^a78^c .
\]
Since \(n\ge36\) and \(a+c\le12\),
\[
  \frac{\binom{n}{a+c}\binom{a+c}{c}}{\binom n{12}}
  \le
  \frac{12!}{a!\,c!\,25^{12-a-c}}.
\]
Also
\[
  3^{n-a-c}24^a78^c
  =
  24^{12}3^{n-12}
  \biggl(\frac{13}{32}\biggr)^c 8^{a+2c-12}.
\]
Therefore the coefficient sum is at most
\[
  \binom n{12}24^{12}3^{n-12}
  \sum_{\substack{a,c\ge0\\ a+2c\le12}}
  \frac{12!}{a!\,c!\,25^{12-a-c}}
  \biggl(\frac{13}{32}\biggr)^c 8^{a+2c-12}.
\]
The finite rational sum is
\[
  \frac{880987349156210452990561}
       {160000000000000000000000}
  <6 .
\]
This gives the displayed bound.
\end{proof}

\begin{lemma}[Stable moments dominate the sixth $B$ nonboundary width count]\label{lem:bc-stable-dominates-sixth-b-nonboundary-width-count}
For every \(q\ge15\),
\[
  12\binom{2q+6}{12}24^{12}3^{2q-6}\le s_{2q+6}.
\]
\end{lemma}

\begin{proof}
By \cref{lem:bc-littlewood-graph-model}, \(s_{2q+6}\) contains all labelled
simple cycles on \(2q+6\) vertices, whose number is \((2q+5)!/2\).  It is
enough to prove
\[
  24\binom{2q+6}{12}24^{12}3^{2q-6}\le (2q+5)! .
\]
At \(q=15\),
\[
\begin{aligned}
  24\binom{36}{12}24^{12}3^{24}
  &=
  309848052206342328100817975974350028800,\\
  35!
  &=
  10333147966386144929666651337523200000000 .
\end{aligned}
\]
The first integer is smaller than the second.  If the displayed factorial
inequality holds at \(q\), then at \(q+1\) the left side is multiplied by
\[
  9\frac{(2q+8)(2q+7)}{(2q-4)(2q-5)}<20,
\]
because \(44q^2-630q-104>0\) for \(q\ge15\).  The right side is multiplied by
\[
  (2q+6)(2q+7)\ge36\cdot37>20 .
\]
Induction proves the result.
\end{proof}

\begin{proposition}[Residual $B$ sixth nonboundary lower correction]\label{prop:bc-b-sixth-nonboundary-lower-correction}
For every residual row \(B_b\) with \(14\le b\le123\), put \(q=b+1\) and
write \(s_j=m_j^{\mathrm{st}}\) and \(m_j^{B_b}=s_j+\delta_j^{B_b}\).  Then
\[
  -\frac12 s_{2q+6}\le \delta_{2q+6}^{B_b}.
\]
\end{proposition}

\begin{proof}
By \cref{lem:bc-correction-bad-shapes}, the coefficient interpretation in
\cref{lem:bad-shape-witnesses}, and the Pieri-path bijection
\cref{lem:bc-pieri-path-bijection}, \(|\delta_{2q+6}^{B_b}|\) is the number
of column-even Pieri two-strip paths of length \(2q+6\) with terminal first
row at least \(2q\).  By
\cref{lem:bc-b-sixth-nonboundary-width-bad-lower-box-code}, this number is at
most
\[
  6\binom{2q+6}{12}24^{12}3^{2q-6}.
\]
Since \(q\ge15\),
\cref{lem:bc-stable-dominates-sixth-b-nonboundary-width-count} gives
\[
  |\delta_{2q+6}^{B_b}|\le\frac12s_{2q+6}.
\]
This proves the stated lower bound.
\end{proof}

\begin{lemma}[Seventh $B$ nonboundary width-bad paths have a lower-box code]\label{lem:bc-b-seventh-nonboundary-width-bad-lower-box-code}
For every \(q\ge15\), the number of column-even Pieri two-strip paths
\[
  \varnothing=\lambda^{(0)}\subset\lambda^{(1)}\subset\cdots
  \subset\lambda^{(2q+7)}
\]
with
\[
  \lambda^{(2q+7)}_1\ge2q
\]
is at most
\[
  11\binom{2q+7}{14}28^{14}3^{2q-7}.
\]
\end{lemma}

\begin{proof}
Put \(n=2q+7\).  The terminal partition has size \(2n=4q+14\).  Since its
columns have even lengths and its first row has length at least \(2q\), the
number of boxes below the second row is at most
\[
  4q+14-2(2q)=14 .
\]
Hence the terminal partition has at most sixteen nonzero rows.

As in \cref{lem:bc-b-sixth-nonboundary-width-bad-lower-box-code}, encode a
path by the row-pair word and track lower boxes.  A step with no lower box has
three choices, a step with one lower box has at most \(2\cdot14=28\) choices,
and a step with two lower boxes has at most \(\binom{14+1}{2}=105\) choices.
Thus the number of paths is bounded by the sum of the coefficients of
\(x^0,\ldots,x^{14}\) in
\[
  (3+28x+105x^2)^n .
\]

The monomial with \(a\) one-lower-box steps and \(c\) two-lower-box steps,
where \(a+2c\le14\), contributes
\[
  \binom{n}{a+c}\binom{a+c}{c}3^{n-a-c}28^a105^c .
\]
Since \(n\ge37\), \(a+c\le14\), and the fourteen factors in the quotient
\(\binom n{14}/\binom n{a+c}\) which are not cancelled are all at least
\(24\),
\[
  \frac{\binom{n}{a+c}\binom{a+c}{c}}{\binom n{14}}
  \le
  \frac{14!}{a!\,c!\,24^{14-a-c}}.
\]
Also
\[
  3^{n-a-c}28^a105^c
  =
  28^{14}3^{n-14}
  \biggl(\frac{45}{112}\biggr)^c
  \biggl(\frac{28}{3}\biggr)^{a+2c-14}.
\]
Therefore the coefficient sum is at most
\[
  \binom n{14}28^{14}3^{n-14}
  \sum_{\substack{a,c\ge0\\ a+2c\le14}}
  \frac{14!}{a!\,c!\,24^{14-a-c}}
  \biggl(\frac{45}{112}\biggr)^c
  \biggl(\frac{28}{3}\biggr)^{a+2c-14}.
\]
The finite rational sum is
\[
  \frac{79869422716689772849072285277}
       {7978958832736206792937177088}
  <11 .
\]
This gives the displayed bound.
\end{proof}

\begin{lemma}[Stable moments dominate the seventh $B$ nonboundary width count above the first row]\label{lem:bc-stable-dominates-seventh-b-nonboundary-width-count}
For every \(q\ge16\),
\[
  22\binom{2q+7}{14}28^{14}3^{2q-7}\le s_{2q+7}.
\]
\end{lemma}

\begin{proof}
By \cref{lem:bc-littlewood-graph-model}, \(s_{2q+7}\) contains all labelled
simple cycles on \(2q+7\) vertices, whose number is \((2q+6)!/2\).  It is
enough to prove
\[
  44\binom{2q+7}{14}28^{14}3^{2q-7}\le (2q+6)! .
\]
At \(q=16\),
\[
\begin{aligned}
  44\binom{39}{14}28^{14}3^{25}
  &=
  102382924139334956961091944497434635490295808,\\
  38!
  &=
  523022617466601111760007224100074291200000000 .
\end{aligned}
\]
The first integer is smaller than the second.  If the displayed factorial
inequality holds at \(q\), then at \(q+1\) the left side is multiplied by
\[
  9\frac{(2q+9)(2q+8)}{(2q-5)(2q-6)}<23,
\]
because \(56q^2-812q+42>0\) for \(q\ge16\).  The right side is multiplied by
\[
  (2q+7)(2q+8)\ge39\cdot40>23 .
\]
Induction proves the result.
\end{proof}

\begin{lemma}[Exact $B_{14}$ seventh nonboundary bad-shape count]\label{lem:bc-b14-seventh-nonboundary-badshape-certificate}
For \(q=15\), the number of column-even Pieri two-strip paths of length
\(37\) with terminal first row at least \(30\) is
\[
  6170018716859356895648705996 .
\]
Moreover
\[
  2\cdot6170018716859356895648705996
  \le
  s_{37}.
\]
\end{lemma}

\begin{proof}
The exact C++/GMP replay
\[
\texttt{certificates/post\_m29/bc\_b14\_j37\_badshape\_gmp\_certificate.log}
\]
implements the following finite reverse Pieri recursion.  It initializes the
terminal set to the \(45\) partitions
\[
  (\mu_1,\mu_1,\mu_2,\mu_2,\ldots),
  \qquad |\mu|=37,\qquad \mu_1\ge30 .
\]
At each reverse step it removes either two boxes from one row, or one box from
two rows whose removed boxes have different columns, and keeps exactly those
predecessors which remain partitions.  This is precisely the inverse of
adding a horizontal two-strip, so after \(37\) reverse steps the coefficient of
the empty partition is the displayed number of Pieri paths.

The replay exits with status \(0\), reports \texttt{SUMMARY failures=0}, and
has SHA-256
\[
\hashsplit{bf212c9d15e9b1953f99b01aff69af74}{b7e4a0a78ae24b769ba2fcd502fa046b}.
\]
It reports
\[
  s_{37}=991021752128621887393916677761923440691808
\]
from the recurrence of \cref{lem:bc-stable-moment-recurrence}, and verifies
the positive margin
\[
  s_{37}
  -2\cdot6170018716859356895648705996
  =
  991021752128609547356482959048132143279816 .
\]
This proves both assertions.
\end{proof}

\begin{proposition}[Residual $B$ seventh nonboundary lower correction]\label{prop:bc-b-seventh-nonboundary-lower-correction}
For every residual row \(B_b\) with \(14\le b\le123\), put \(q=b+1\) and
write \(s_j=m_j^{\mathrm{st}}\) and \(m_j^{B_b}=s_j+\delta_j^{B_b}\).  Then
\[
  -\frac12 s_{2q+7}\le \delta_{2q+7}^{B_b}.
\]
\end{proposition}

\begin{proof}
By \cref{lem:bc-correction-bad-shapes}, the coefficient interpretation in
\cref{lem:bad-shape-witnesses}, and the Pieri-path bijection
\cref{lem:bc-pieri-path-bijection}, \(|\delta_{2q+7}^{B_b}|\) is the number
of column-even Pieri two-strip paths of length \(2q+7\) with terminal first
row at least \(2q\).  If \(q=15\), this is bounded by \(s_{37}/2\) by
\cref{lem:bc-b14-seventh-nonboundary-badshape-certificate}.  If \(q\ge16\),
\cref{lem:bc-b-seventh-nonboundary-width-bad-lower-box-code} bounds this
number by
\[
  11\binom{2q+7}{14}28^{14}3^{2q-7},
\]
and \cref{lem:bc-stable-dominates-seventh-b-nonboundary-width-count} absorbs
that count into \(s_{2q+7}/2\).  In all cases
\(|\delta_{2q+7}^{B_b}|\le s_{2q+7}/2\), which proves the stated lower bound.
\end{proof}

\begin{lemma}[Eighth $B$ nonboundary width-bad paths have a lower-box code]\label{lem:bc-b-eighth-nonboundary-width-bad-lower-box-code}
For every \(q\ge15\), the number of column-even Pieri two-strip paths
\[
  \varnothing=\lambda^{(0)}\subset\lambda^{(1)}\subset\cdots
  \subset\lambda^{(2q+8)}
\]
with
\[
  \lambda^{(2q+8)}_1\ge2q
\]
is at most
\[
  21\binom{2q+8}{16}32^{16}3^{2q-8}.
\]
\end{lemma}

\begin{proof}
Put \(n=2q+8\).  The terminal partition has size \(2n=4q+16\).  Since its
columns have even lengths and its first row has length at least \(2q\), the
number of boxes below the second row is at most
\[
  4q+16-2(2q)=16 .
\]
Hence the terminal partition has at most eighteen nonzero rows.

As in \cref{lem:bc-b-sixth-nonboundary-width-bad-lower-box-code}, encode a
path by the row-pair word and track lower boxes.  A step with no lower box has
three choices, a step with one lower box has at most \(2\cdot16=32\) choices,
and a step with two lower boxes has at most \(\binom{16+1}{2}=136\) choices.
Thus the number of paths is bounded by the sum of the coefficients of
\(x^0,\ldots,x^{16}\) in
\[
  (3+32x+136x^2)^n .
\]

The monomial with \(a\) one-lower-box steps and \(c\) two-lower-box steps,
where \(a+2c\le16\), contributes
\[
  \binom{n}{a+c}\binom{a+c}{c}3^{n-a-c}32^a136^c .
\]
Since \(n\ge38\), \(a+c\le16\), and the sixteen factors in the quotient
\(\binom n{16}/\binom n{a+c}\) which are not cancelled are all at least
\(23\),
\[
  \frac{\binom{n}{a+c}\binom{a+c}{c}}{\binom n{16}}
  \le
  \frac{16!}{a!\,c!\,23^{16-a-c}}.
\]
Therefore the coefficient sum is at most
\[
  \binom n{16}32^{16}3^{n-16}
  \sum_{\substack{a,c\ge0\\ a+2c\le16}}
  \frac{16!}{a!\,c!\,23^{16-a-c}}
  \frac{3^{16-a-c}136^c}{32^{16-a}}.
\]
The finite rational sum is
\[
  \frac{4526978891280486763545625577579279602498403}
       {226253389683665863434374437380034188541952}
  <21 .
\]
This gives the displayed bound.
\end{proof}

\begin{lemma}[Stable moments dominate the eighth $B$ nonboundary width count above the first row]\label{lem:bc-stable-dominates-eighth-b-nonboundary-width-count}
For every \(q\ge18\),
\[
  42\binom{2q+8}{16}32^{16}3^{2q-8}\le s_{2q+8}.
\]
\end{lemma}

\begin{proof}
By \cref{lem:bc-littlewood-graph-model}, \(s_{2q+8}\) contains all labelled
simple cycles on \(2q+8\) vertices, whose number is \((2q+7)!/2\).  It is
enough to prove
\[
  84\binom{2q+8}{16}32^{16}3^{2q-8}\le (2q+7)! .
\]
At \(q=18\),
\[
\begin{aligned}
  84\binom{44}{16}32^{16}3^{28}
  &=
  968083911280252724764260401389533398116759593025536,\\
  43!
  &=
  60415263063373835637355132068513997507264512000000000 .
\end{aligned}
\]
The first integer is smaller than the second.  If the displayed factorial
inequality holds at \(q\), then at \(q+1\) the left side is multiplied by
\[
  9\frac{(2q+10)(2q+9)}{(2q-6)(2q-7)}<23,
\]
because \(56q^2-940q+156>0\) for \(q\ge18\).  The right side is multiplied by
\[
  (2q+8)(2q+9)\ge44\cdot45>23 .
\]
Induction proves the result.
\end{proof}

\begin{lemma}[Bounded Littlewood determinant for finite type-$B$ adjoint moments]
\label{lem:bc-b-bounded-littlewood-determinant}
Work in the completed symmetric-function ring over \(\mathbb Q\), put
\(e_r=0\) for \(r<0\), and define
\[
 f_r^e=\sum_{k\in\mathbb Z}e_ke_{k+r},\qquad
 e=\sum_{k\ge0}e_k,\qquad
 \bar e=\sum_{k\ge0}(-1)^ke_k.
\]
For \(b\ge1\), set
\[
 \mathcal B_b=\frac12\left\{
 e\det_{1\le i,j\le b}(f^e_{i-j}-f^e_{i+j-1})
 +\bar e\det_{1\le i,j\le b}(f^e_{i-j}+f^e_{i+j-1})
 \right\}.
\]
Then, for every \(r\ge0\),
\[
 m_r^{B_b}=[m_{(2^r)}]\mathcal B_b.
\]
Moreover, if \(F\) is homogeneous of symmetric-function degree \(2r\),
and \(\pi_2(p_1)=p\), \(\pi_2(p_2)=z p^2\),
\(\pi_2(p_k)=0\) for \(k\ge3\), then
\[
 [m_{(2^r)}]F
 =r!\sum_{a=0}^r(2(r-a)-1)!!
   [p^{2r}z^a]\pi_2(F).
\]
\end{lemma}

\begin{proof}
By \cref{lem:bc-correction-bad-shapes,lem:bc-littlewood-graph-model}, the
finite type-$B$ moment is the stable column-even sum with precisely the
width-excluded terms removed.  Transpose the resulting admitted partitions.
The involution \(\omega\) interchanges
\(e_2\) and \(h_2\), so Hall duality gives
\[
 m_r^{B_b}
 =\left\langle h_2^r,
   \sum_{\substack{\lambda_1\le2b+1\\c(\lambda)=0}}s_\lambda
  \right\rangle,
\]
where \(c(\lambda)\) is the number of odd columns.  The two identities in
\cite[Theorem~1.3, equations~(1.9)--(1.10)]{HKKO2025}, specialized at
\(u=0\), give the sum and difference of the sectors
\(c(\lambda)=0\) and \(c(\lambda)=2b+1\).  Their half-sum is
\(\mathcal B_b\), and
\(\langle h_2^r,F\rangle=[m_{(2^r)}]F\), proving the first assertion.

For the second, expand \(F\) in power sums.  A power-sum monomial involving
some \(p_k\), \(k\ge3\), cannot contribute to \(m_{(2^r)}\).  In
\(p_1^{2(r-a)}p_2^a\), the \(a\) ordered \(p_2\)-factors choose distinct
variables and the remaining ordered \(p_1\)-factors hit every other
variable twice.  The coefficient of \(x_1^2\cdots x_r^2\) is therefore
\[
 \frac{r!}{(r-a)!}\frac{(2(r-a))!}{2^{r-a}}
 =r!(2(r-a)-1)!!,
\]
which is exactly the displayed functional after substituting \(q=zp^2\).
\end{proof}

\begin{lemma}[Unified exact type-$B$ frontier certificate, H8--H27]
\label{lem:bc-b-h8-h27-bounded-littlewood-certificate}
Let \(\mathcal F_B\) consist of the integer pairs \((h,q)\) in
\[
 \begin{split}
 &h=8,\quad 15\le q\le17;\\
 &9\le h\le23,\quad 15\le q\le\lceil3h/2\rceil+4;\\
 &24\le h\le27,\quad 15\le q\le2h-7.
 \end{split}
\]
For every \((h,q)\in\mathcal F_B\), the number of column-even Pieri
two-strip paths of length \(2q+h\) with terminal first row at least \(2q\)
is strictly less than \(s_{2q+h}/2\).
\end{lemma}

\begin{proof}
By \cref{lem:bc-correction-bad-shapes,lem:bc-b-bounded-littlewood-determinant},
the required bad-path count is exactly
\[
 \#\operatorname{bad}(q,j)=s_j-m_j^{B_{q-1}},\qquad j=2q+h.
\]
The current exact verifier
\[
\texttt{character\_ring\_iter/verify\_active\_bc\_b\_frontiers\_bounded\_littlewood\_gmp.cpp}
\]
generates \(\mathcal F_B\) from the three displayed ranges and rejects any
scope other than its \(337\) distinct cases.  It specializes the determinant
of \cref{lem:bc-b-bounded-littlewood-determinant} at the integer nodes
\(z=0,1,\ldots,121\), using
\[
 \sum_{r\ge0}\pi_2(e_r)u^r=\exp(pu-zp^2u^2/2)
\]
to construct every entry as a truncated exact power series.  Lagrange
interpolation applies the coefficient functional in that lemma.

The independent prime-field evaluations are merged by exact GMP CRT.  Since
the finite moment counts a subset of the unrestricted stable objects,
\(0\le m_j^{B_b}\le s_j\).  The verifier chooses primes until their product
is strictly larger than \(s_{121}\), and rejects a reconstructed value
outside its corresponding stable bound.  Thus every CRT representative is
the unique integer moment, rather than a probable reconstruction.

The accepted transcript
\[
\texttt{certificates/post\_m29/bc\_b\_h8\_h27\_bounded\_littlewood\_gmp\_certificate.log}
\]
has SHA-256
\[
\hashsplit{79a6d2721185cc71a8884987774da1dd}{0e90f38482c4f8af1b26bf3c7e695327}.
\]
It records the per-\(h\) scope and minimum margin, reports
\texttt{rows=33 cases=337 maximum\_moment=121 failures=0}, and ends with
\texttt{ACTIVE\_B\_FRONTIER\_BOUNDED\_LITTLEWOOD VERIFICATION: ALL PASS}.
In particular, every exact value
\(s_j-2\#\operatorname{bad}(q,j)=2m_j^{B_{q-1}}-s_j\) is positive.
The older reverse-Pieri transcripts cited below remain independent controls,
but are not premises of this unified calculation.
\end{proof}

\begin{lemma}[Exact first three eighth $B$ nonboundary bad-shape counts]\label{lem:bc-b-eighth-frontier-badshape-certificate}
For \(q=15,16,17\), the number of column-even Pieri two-strip paths of length
\(2q+8\) with terminal first row at least \(2q\) is at most \(s_{2q+8}/2\).
\end{lemma}

\begin{proof}
The exact C++/GMP/OpenMP replay
\[
\texttt{certificates/post\_m29/bc\_b\_eighth\_frontier\_gmp\_certificate.log}
\]
implements the same reverse Pieri recursion as
\cref{lem:bc-b14-seventh-nonboundary-badshape-certificate}, now for the three
cases \(q=15,16,17\), \(j=2q+8\).  It initializes the terminal set to the
column-even width-bad partitions
\[
  (\mu_1,\mu_1,\mu_2,\mu_2,\ldots),
  \qquad |\mu|=j,\qquad \mu_1\ge2q ,
\]
and reverses horizontal two-strip additions for \(j\) steps.

The replay exits with status \(0\), reports \texttt{SUMMARY failures=0}, and
has SHA-256
\[
\hashsplit{5d05cb22e6586256710a9d062f898f67}{f3537c88425ad5b3d3d3ebcecd8972ed}.
\]
It verifies the following exact margins:
\[
\begin{array}{c|c|c|c}
q & j & \#\mathrm{bad}(q,j) & s_j-2\#\mathrm{bad}(q,j)\\ \hline
15&38&
727372931790442036273061869531&
37163028392636523845676450412420959371591786\\
16&40&
15500236532483430857257682568105&
56514900379404414074195085438168663587012160846\\
17&42&
314516359551758315471982172229150&
94986302333375731438360035751682822280371771872356 .
\end{array}
\]
Each margin is positive, proving the claimed half-stable bound.
\end{proof}

\begin{proposition}[Residual $B$ eighth nonboundary lower correction]\label{prop:bc-b-eighth-nonboundary-lower-correction}
For every residual row \(B_b\) with \(14\le b\le123\), put \(q=b+1\) and
write \(s_j=m_j^{\mathrm{st}}\) and \(m_j^{B_b}=s_j+\delta_j^{B_b}\).  Then
\[
  -\frac12 s_{2q+8}\le \delta_{2q+8}^{B_b}.
\]
\end{proposition}

\begin{proof}
By \cref{lem:bc-correction-bad-shapes}, the coefficient interpretation in
\cref{lem:bad-shape-witnesses}, and the Pieri-path bijection
\cref{lem:bc-pieri-path-bijection}, \(|\delta_{2q+8}^{B_b}|\) is the number
of column-even Pieri two-strip paths of length \(2q+8\) with terminal first
row at least \(2q\).  If \(q=15,16,17\), this is bounded by \(s_{2q+8}/2\) by
\cref{lem:bc-b-h8-h27-bounded-littlewood-certificate}.  If \(q\ge18\),
\cref{lem:bc-b-eighth-nonboundary-width-bad-lower-box-code} bounds this
number by
\[
  21\binom{2q+8}{16}32^{16}3^{2q-8},
\]
and \cref{lem:bc-stable-dominates-eighth-b-nonboundary-width-count} absorbs
that count into \(s_{2q+8}/2\).  In all cases
\(|\delta_{2q+8}^{B_b}|\le s_{2q+8}/2\), which proves the stated lower bound.
\end{proof}

\begin{lemma}[Ninth $B$ nonboundary width-bad paths have a lower-box code]\label{lem:bc-b-ninth-nonboundary-width-bad-lower-box-code}
For every \(q\ge15\), the number of column-even Pieri two-strip paths
\[
  \varnothing=\lambda^{(0)}\subset\lambda^{(1)}\subset\cdots
  \subset\lambda^{(2q+9)}
\]
with
\[
  \lambda^{(2q+9)}_1\ge2q
\]
is at most
\[
  45\binom{2q+9}{18}36^{18}3^{2q-9}.
\]
\end{lemma}

\begin{proof}
Put \(n=2q+9\).  The terminal partition has size \(2n=4q+18\).  Since its
columns have even lengths and its first row has length at least \(2q\), the
number of boxes below the second row is at most
\[
  4q+18-2(2q)=18 .
\]
Hence the terminal partition has at most twenty nonzero rows.

As in \cref{lem:bc-b-sixth-nonboundary-width-bad-lower-box-code}, encode a
path by the row-pair word and track lower boxes.  A step with no lower box has
three choices, a step with one lower box has at most \(2\cdot18=36\) choices,
and a step with two lower boxes has at most \(\binom{18+1}{2}=171\) choices.
Thus the number of paths is bounded by the sum of the coefficients of
\(x^0,\ldots,x^{18}\) in
\[
  (3+36x+171x^2)^n .
\]

The monomial with \(a\) one-lower-box steps and \(c\) two-lower-box steps,
where \(a+2c\le18\), contributes
\[
  \binom{n}{a+c}\binom{a+c}{c}3^{n-a-c}36^a171^c .
\]
Since \(n\ge39\), \(a+c\le18\), and the eighteen factors in the quotient
\(\binom n{18}/\binom n{a+c}\) which are not cancelled are all at least
\(22\),
\[
  \frac{\binom{n}{a+c}\binom{a+c}{c}}{\binom n{18}}
  \le
  \frac{18!}{a!\,c!\,22^{18-a-c}}.
\]
Therefore the coefficient sum is at most
\[
  \binom n{18}36^{18}3^{n-18}
  \sum_{\substack{a,c\ge0\\ a+2c\le18}}
  \frac{18!}{a!\,c!\,22^{18-a-c}}
  \frac{3^{18-a-c}171^c}{36^{18-a}}.
\]
The finite rational sum is
\[
  \frac{361917809283112963557622284723886093}
       {8204044929474043719393332089061376}
  <45 .
\]
This gives the displayed bound.
\end{proof}

\begin{lemma}[Stable moments dominate the ninth $B$ nonboundary width count above the first row]\label{lem:bc-stable-dominates-ninth-b-nonboundary-width-count}
For every \(q\ge19\),
\[
  90\binom{2q+9}{18}36^{18}3^{2q-9}\le s_{2q+9}.
\]
\end{lemma}

\begin{proof}
By \cref{lem:bc-littlewood-graph-model}, \(s_{2q+9}\) contains all labelled
simple cycles on \(2q+9\) vertices, whose number is \((2q+8)!/2\).  It is
enough to prove
\[
  180\binom{2q+9}{18}36^{18}3^{2q-9}\le (2q+8)! .
\]
At \(q=19\),
\[
\begin{aligned}
  180\binom{47}{18}36^{18}3^{29}
  &=
  582132191395269676126538459972138264700438046490335641600,\\
  46!
  &=
  5502622159812088949850305428800254892961651752960000000000 .
\end{aligned}
\]
The first integer is smaller than the second.  If the displayed factorial
inequality holds at \(q\), then at \(q+1\) the left side is multiplied by
\[
  9\frac{(2q+11)(2q+10)}{(2q-7)(2q-8)}<23,
\]
because \(56q^2-1068q+298>0\) for \(q\ge19\).  The right side is multiplied
by
\[
  (2q+9)(2q+10)\ge47\cdot48>23 .
\]
Induction proves the result.
\end{proof}

\begin{lemma}[Exact first four ninth $B$ nonboundary bad-shape counts]\label{lem:bc-b-ninth-frontier-badshape-certificate}
For \(q=15,16,17,18\), the number of column-even Pieri two-strip paths of
length \(2q+9\) with terminal first row at least \(2q\) is at most
\(s_{2q+9}/2\).
\end{lemma}

\begin{proof}
The exact C++/GMP/OpenMP replay
\[
\texttt{certificates/post\_m29/bc\_ninth\_frontier\_gmp\_certificate.log}
\]
implements the reverse Pieri recursion of
\cref{lem:bc-b-eighth-frontier-badshape-certificate} for
\(q=15,16,17,18\), \(j=2q+9\).  It exits with status \(0\), reports
\texttt{SUMMARY failures=0}, and has SHA-256
\[
\hashsplit{3459b8f1c779c878c9a1666f18b27a03}{548202729f4ddaa23399b6f2c11c242c}.
\]
It verifies the following exact margins:
\[
\begin{array}{c|c|c|c}
q & j & \#\mathrm{bad}(q,j) & s_j-2\#\mathrm{bad}(q,j)\\ \hline
15&39&
81596655286601105266416444390197&
1430766394625961748437041706576357337819220214\\
16&41&
1926342911785452502271099095125800&
2288839508814961324669405817236421254438863593520\\
17&43&
43077718528836304303270590063965610&
4036896628145341970914689003048977156793546166686156\\
18&45&
917849214459918430646886583313786946&
7814348742732519466761141196680536729520495482426158780 .
\end{array}
\]
Each margin is positive, proving the claimed half-stable bound.
\end{proof}

\begin{proposition}[Residual $B$ ninth nonboundary lower correction]\label{prop:bc-b-ninth-nonboundary-lower-correction}
For every residual row \(B_b\) with \(14\le b\le123\), put \(q=b+1\) and
write \(s_j=m_j^{\mathrm{st}}\) and \(m_j^{B_b}=s_j+\delta_j^{B_b}\).  Then
\[
  -\frac12 s_{2q+9}\le \delta_{2q+9}^{B_b}.
\]
\end{proposition}

\begin{proof}
By \cref{lem:bc-correction-bad-shapes}, the coefficient interpretation in
\cref{lem:bad-shape-witnesses}, and the Pieri-path bijection
\cref{lem:bc-pieri-path-bijection}, \(|\delta_{2q+9}^{B_b}|\) is the number
of column-even Pieri two-strip paths of length \(2q+9\) with terminal first
row at least \(2q\).  If \(q=15,16,17,18\), this is bounded by
\(s_{2q+9}/2\) by
\cref{lem:bc-b-h8-h27-bounded-littlewood-certificate}.  If \(q\ge19\),
\cref{lem:bc-b-ninth-nonboundary-width-bad-lower-box-code} bounds this
number by
\[
  45\binom{2q+9}{18}36^{18}3^{2q-9},
\]
and \cref{lem:bc-stable-dominates-ninth-b-nonboundary-width-count} absorbs
that count into \(s_{2q+9}/2\).  In all cases
\(|\delta_{2q+9}^{B_b}|\le s_{2q+9}/2\), which proves the stated lower bound.
\end{proof}

\begin{lemma}[Tenth $B$ nonboundary width-bad paths have a lower-box code]\label{lem:bc-b-tenth-nonboundary-width-bad-lower-box-code}
For every \(q\ge15\), the number of column-even Pieri two-strip paths
\[
  \varnothing=\lambda^{(0)}\subset\lambda^{(1)}\subset\cdots
  \subset\lambda^{(2q+10)}
\]
with
\[
  \lambda^{(2q+10)}_1\ge2q
\]
is at most
\[
  108\binom{2q+10}{20}40^{20}3^{2q-10}.
\]
\end{lemma}

\begin{proof}
Put \(n=2q+10\).  The terminal partition has size \(2n=4q+20\).  Since its
columns have even lengths and its first row has length at least \(2q\), the
number of boxes below the second row is at most
\[
  4q+20-2(2q)=20 .
\]
Hence the terminal partition has at most twenty-two nonzero rows.

As in \cref{lem:bc-b-sixth-nonboundary-width-bad-lower-box-code}, encode a
path by the row-pair word and track lower boxes.  A step with no lower box has
three choices, a step with one lower box has at most \(2\cdot20=40\) choices,
and a step with two lower boxes has at most \(\binom{20+1}{2}=210\) choices.
Thus the number of paths is bounded by the sum of the coefficients of
\(x^0,\ldots,x^{20}\) in
\[
  (3+40x+210x^2)^n .
\]

The monomial with \(a\) one-lower-box steps and \(c\) two-lower-box steps,
where \(a+2c\le20\), contributes
\[
  \binom{n}{a+c}\binom{a+c}{c}3^{n-a-c}40^a210^c .
\]
Since \(n\ge40\), \(a+c\le20\), and the twenty factors in the quotient
\(\binom n{20}/\binom n{a+c}\) which are not cancelled are all at least
\(21\),
\[
  \frac{\binom{n}{a+c}\binom{a+c}{c}}{\binom n{20}}
  \le
  \frac{20!}{a!\,c!\,21^{20-a-c}}.
\]
Therefore the coefficient sum is at most
\[
  \binom n{20}40^{20}3^{n-20}
  \sum_{\substack{a,c\ge0\\ a+2c\le20}}
  \frac{20!}{a!\,c!\,21^{20-a-c}}
  \frac{3^{20-a-c}210^c}{40^{20-a}}.
\]
The finite rational sum is
\[
  \frac{117722605237864079789280365964598334253029}
       {1092809866779230061199360000000000000000}
  <108 .
\]
This gives the displayed bound.
\end{proof}

\begin{lemma}[Stable moments dominate the tenth $B$ nonboundary width count above the first row]\label{lem:bc-stable-dominates-tenth-b-nonboundary-width-count}
For every \(q\ge20\),
\[
  216\binom{2q+10}{20}40^{20}3^{2q-10}\le s_{2q+10}.
\]
\end{lemma}

\begin{proof}
By \cref{lem:bc-littlewood-graph-model}, \(s_{2q+10}\) contains all labelled
simple cycles on \(2q+10\) vertices, whose number is \((2q+9)!/2\).  It is
enough to prove
\[
  432\binom{2q+10}{20}40^{20}3^{2q-10}\le (2q+9)! .
\]
At \(q=20\),
\[
\begin{aligned}
  432\binom{50}{20}40^{20}3^{30}
  &=
  460904974698595369328663350383438282620928000000000000000000000,\\
  49!
  &=
  608281864034267560872252163321295376887552831379210240000000000 .
\end{aligned}
\]
The first integer is smaller than the second.  If the displayed factorial
inequality holds at \(q\), then at \(q+1\) the left side is multiplied by
\[
  9\frac{(2q+12)(2q+11)}{(2q-8)(2q-9)}<25,
\]
because \(64q^2-1264q+612>0\) for \(q\ge20\).  The right side is multiplied
by
\[
  (2q+10)(2q+11)\ge50\cdot51>25 .
\]
Induction proves the result.
\end{proof}

\begin{lemma}[Exact first five tenth $B$ nonboundary bad-shape counts]\label{lem:bc-b-tenth-frontier-badshape-certificate}
For \(q=15,16,17,18,19\), the number of column-even Pieri two-strip paths of
length \(2q+10\) with terminal first row at least \(2q\) is at most
\(s_{2q+10}/2\).
\end{lemma}

\begin{proof}
The exact C++/GMP/OpenMP replay
\[
\texttt{certificates/post\_m29/bc\_tenth\_frontier\_gmp\_certificate.log}
\]
implements the reverse Pieri recursion of
\cref{lem:bc-b-eighth-frontier-badshape-certificate} for
\(q=15,16,17,18,19\), \(j=2q+10\).  It exits with status \(0\), reports
\texttt{SUMMARY failures=0}, and has SHA-256
\[
\hashsplit{abf7ac5edc1f162794cea23d37343d9c}{7dda9415a30f8eab1a622884ebe7b49c}.
\]
It verifies the following exact margins:
\[
\begin{array}{c|c|c|c}
q & j & \#\mathrm{bad}(q,j) & s_j-2\#\mathrm{bad}(q,j)\\ \hline
15&40&
8794106710689188840998694148435610&
56514900379386856861246772027348380714080425836\\
16&42&
229375444845227744316056074438167665&
94986302333375273316503064399710821112187239995326\\
17&44&
5638833915013514621203873625181397292&
175604143934124940100093706905280814409929173685827432\\
18&46&
131485817886217717873410741865804465230&
355551352238135109999385991741420958490914665526325828260\\
19&48&
2923804498135795385020064717727104507466&
785317984503258803849482363853038965238893379037105540047660 .
\end{array}
\]
Each margin is positive, proving the claimed half-stable bound.
\end{proof}

\begin{proposition}[Residual $B$ tenth nonboundary lower correction]\label{prop:bc-b-tenth-nonboundary-lower-correction}
For every residual row \(B_b\) with \(14\le b\le123\), put \(q=b+1\) and
write \(s_j=m_j^{\mathrm{st}}\) and \(m_j^{B_b}=s_j+\delta_j^{B_b}\).  Then
\[
  -\frac12 s_{2q+10}\le \delta_{2q+10}^{B_b}.
\]
\end{proposition}

\begin{proof}
By \cref{lem:bc-correction-bad-shapes}, the coefficient interpretation in
\cref{lem:bad-shape-witnesses}, and the Pieri-path bijection
\cref{lem:bc-pieri-path-bijection}, \(|\delta_{2q+10}^{B_b}|\) is the number
of column-even Pieri two-strip paths of length \(2q+10\) with terminal first
row at least \(2q\).  If \(q=15,16,17,18,19\), this is bounded by
\(s_{2q+10}/2\) by
\cref{lem:bc-b-h8-h27-bounded-littlewood-certificate}.  If \(q\ge20\),
\cref{lem:bc-b-tenth-nonboundary-width-bad-lower-box-code} bounds this
number by
\[
  108\binom{2q+10}{20}40^{20}3^{2q-10},
\]
and \cref{lem:bc-stable-dominates-tenth-b-nonboundary-width-count} absorbs
that count into \(s_{2q+10}/2\).  In all cases
\(|\delta_{2q+10}^{B_b}|\le s_{2q+10}/2\), which proves the stated lower
bound.
\end{proof}

\begin{lemma}[Eleventh $B$ nonboundary width-bad paths have a lower-box code]\label{lem:bc-b-eleventh-nonboundary-width-bad-lower-box-code}
For every \(q\ge15\), the number of column-even Pieri two-strip paths
\[
  \varnothing=\lambda^{(0)}\subset\lambda^{(1)}\subset\cdots
  \subset\lambda^{(2q+11)}
\]
with
\[
  \lambda^{(2q+11)}_1\ge2q
\]
is at most
\[
  293\binom{2q+11}{22}44^{22}3^{2q-11}.
\]
\end{lemma}

\begin{proof}
Put \(n=2q+11\).  The terminal partition has size \(2n=4q+22\).  Since its
columns have even lengths and its first row has length at least \(2q\), the
number of boxes below the second row is at most
\[
  4q+22-2(2q)=22 .
\]
Hence the terminal partition has at most twenty-four nonzero rows.

As in \cref{lem:bc-b-sixth-nonboundary-width-bad-lower-box-code}, encode a
path by the row-pair word and track lower boxes.  A step with no lower box has
three choices, a step with one lower box has at most \(2\cdot22=44\) choices,
and a step with two lower boxes has at most \(\binom{22+1}{2}=253\) choices.
Thus the number of paths is bounded by the sum of the coefficients of
\(x^0,\ldots,x^{22}\) in
\[
  (3+44x+253x^2)^n .
\]

The monomial with \(a\) one-lower-box steps and \(c\) two-lower-box steps,
where \(a+2c\le22\), contributes
\[
  \binom{n}{a+c}\binom{a+c}{c}3^{n-a-c}44^a253^c .
\]
Since \(n\ge41\), \(a+c\le22\), and the twenty-two factors in the quotient
\(\binom n{22}/\binom n{a+c}\) which are not cancelled are all at least
\(20\),
\[
  \frac{\binom{n}{a+c}\binom{a+c}{c}}{\binom n{22}}
  \le
  \frac{22!}{a!\,c!\,20^{22-a-c}}.
\]
Therefore the coefficient sum is at most
\[
  \binom n{22}44^{22}3^{n-22}
  \sum_{\substack{a,c\ge0\\ a+2c\le22}}
  \frac{22!}{a!\,c!\,20^{22-a-c}}
  \frac{3^{22-a-c}253^c}{44^{22-a}}.
\]
The finite rational sum is
\[
  \frac{443599768741143824547304978674911318281273800637175688599}
       {1514898313245890164857636960597966848000000000000000000}
  <293 .
\]
This gives the displayed bound.
\end{proof}

\begin{lemma}[Stable moments dominate the eleventh $B$ nonboundary width count above the first row]\label{lem:bc-stable-dominates-eleventh-b-nonboundary-width-count}
For every \(q\ge22\),
\[
  586\binom{2q+11}{22}44^{22}3^{2q-11}\le s_{2q+11}.
\]
\end{lemma}

\begin{proof}
By \cref{lem:bc-littlewood-graph-model}, \(s_{2q+11}\) contains all labelled
simple cycles on \(2q+11\) vertices, whose number is \((2q+10)!/2\).  It is
enough to prove
\[
  1172\binom{2q+11}{22}44^{22}3^{2q-11}\le (2q+10)! .
\]
At \(q=22\),
\[
\begin{aligned}
  1172\binom{55}{22}44^{22}3^{33}
  &=
  12137139171939900106689615815324886810578010253530380203172753414553600,\\
  54!
  &=
  230843697339241380472092742683027581083278564571807941132288000000000000 .
\end{aligned}
\]
The first integer is smaller than the second.  If the displayed factorial
inequality holds at \(q\), then at \(q+1\) the left side is multiplied by
\[
  9\frac{(2q+13)(2q+12)}{(2q-9)(2q-10)}<27,
\]
because \(72q^2-1476q+1026>0\) for \(q\ge22\).  The right side is multiplied
by
\[
  (2q+11)(2q+12)\ge55\cdot56>27 .
\]
Induction proves the result.
\end{proof}

\begin{lemma}[Exact first seven eleventh $B$ nonboundary bad-shape counts]\label{lem:bc-b-eleventh-frontier-badshape-certificate}
For \(15\le q\le21\), the number of column-even Pieri two-strip paths
of length \(2q+11\) with terminal first row at least \(2q\) is at most
\(s_{2q+11}/2\).
\end{lemma}

\begin{proof}
The exact C++/GMP/OpenMP replay
\[
\texttt{certificates/post\_m29/bc\_eleventh\_frontier\_gmp\_certificate.log}
\]
implements the reverse Pieri recursion of
\cref{lem:bc-b-eighth-frontier-badshape-certificate} for
\(15\le q\le21\), \(j=2q+11\).  It exits with status \(0\), reports
\texttt{SUMMARY failures=0}, and has SHA-256
\[
\hashsplit{de69669b8df7784d87b33232131fee8f}{a6eb22c830fdcc8eeb066d851a409c0c}.
\]
It verifies the following exact margins:
\[
\begin{array}{c|c|c|c}
q & j & \#\mathrm{bad}(q,j) & s_j-2\#\mathrm{bad}(q,j)\\ \hline
15&41&
917406516916434909318039084336626480&
2288839508813130364321396518322789718468380592160\\
16&43&
26369959387638196716572586608961318825&
4036896628145289317151350784328152618161508371979726\\
17&45&
710983441612846257292453314350422789848&
7814348742732518046629956399907859005907639948208152976\\
18&47&
18104254770555721644360125274088694749197&
16533071956148091606739412756578627122586773576653757319078\\
19&49&
437907984655019292950173937746451968524840&
38087788784726914578665583593910361469571031312818908657357232\\
20&51&
10111325553336459132478426731966206572544775&
95209341861812270694950717642580556511339581646654413332606512050\\
21&53&
223820420460688222608586493645429584684736700&
257420799961522256421436385667843545103560322026256827976578728322824 .
\end{array}
\]
Each margin is positive, proving the claimed half-stable bound.
\end{proof}

\begin{proposition}[Residual $B$ eleventh nonboundary lower correction]\label{prop:bc-b-eleventh-nonboundary-lower-correction}
For every residual row \(B_b\) with \(14\le b\le123\), put \(q=b+1\) and
write \(s_j=m_j^{\mathrm{st}}\) and \(m_j^{B_b}=s_j+\delta_j^{B_b}\).  Then
\[
  -\frac12 s_{2q+11}\le \delta_{2q+11}^{B_b}.
\]
\end{proposition}

\begin{proof}
By \cref{lem:bc-correction-bad-shapes}, the coefficient interpretation in
\cref{lem:bad-shape-witnesses}, and the Pieri-path bijection
\cref{lem:bc-pieri-path-bijection}, \(|\delta_{2q+11}^{B_b}|\) is the number
of column-even Pieri two-strip paths of length \(2q+11\) with terminal first
row at least \(2q\).  If \(15\le q\le21\), this is bounded by
\(s_{2q+11}/2\) by
\cref{lem:bc-b-h8-h27-bounded-littlewood-certificate}.  If \(q\ge22\),
\cref{lem:bc-b-eleventh-nonboundary-width-bad-lower-box-code} bounds this
number by
\[
  293\binom{2q+11}{22}44^{22}3^{2q-11},
\]
and \cref{lem:bc-stable-dominates-eleventh-b-nonboundary-width-count} absorbs
that count into \(s_{2q+11}/2\).  In all cases
\(|\delta_{2q+11}^{B_b}|\le s_{2q+11}/2\), which proves the stated lower
bound.
\end{proof}

\begin{lemma}[Twelfth $B$ nonboundary width-bad paths have a lower-box code]\label{lem:bc-b-twelfth-nonboundary-width-bad-lower-box-code}
For every \(q\ge15\), the number of column-even Pieri two-strip paths
\[
  \varnothing=\lambda^{(0)}\subset\lambda^{(1)}\subset\cdots
  \subset\lambda^{(2q+12)}
\]
with
\[
  \lambda^{(2q+12)}_1\ge2q
\]
is at most
\[
  892\binom{2q+12}{24}48^{24}3^{2q-12}.
\]
\end{lemma}

\begin{proof}
Put \(n=2q+12\).  The terminal partition has size \(2n=4q+24\).  Since its
columns have even lengths and its first row has length at least \(2q\), the
number of boxes below the second row is at most
\[
  4q+24-2(2q)=24 .
\]
Hence the terminal partition has at most twenty-six nonzero rows.

As in \cref{lem:bc-b-sixth-nonboundary-width-bad-lower-box-code}, encode a
path by the row-pair word and track lower boxes.  A step with no lower box has
three choices, a step with one lower box has at most \(2\cdot24=48\) choices,
and a step with two lower boxes has at most \(\binom{24+1}{2}=300\) choices.
Thus the number of paths is bounded by the sum of the coefficients of
\(x^0,\ldots,x^{24}\) in
\[
  (3+48x+300x^2)^n .
\]

The monomial with \(a\) one-lower-box steps and \(c\) two-lower-box steps,
where \(a+2c\le24\), contributes
\[
  \binom{n}{a+c}\binom{a+c}{c}3^{n-a-c}48^a300^c .
\]
Since \(n\ge42\), \(a+c\le24\), and the twenty-four factors in the quotient
\(\binom n{24}/\binom n{a+c}\) which are not cancelled are all at least
\(19\),
\[
  \frac{\binom{n}{a+c}\binom{a+c}{c}}{\binom n{24}}
  \le
  \frac{24!}{a!\,c!\,19^{24-a-c}}.
\]
Therefore the coefficient sum is at most
\[
  \binom n{24}48^{24}3^{n-24}
  \sum_{\substack{a,c\ge0\\ a+2c\le24}}
  \frac{24!}{a!\,c!\,19^{24-a-c}}
  \frac{3^{24-a-c}300^c}{48^{24-a}}.
\]
The finite rational sum is
\[
  \frac{4341923919103137819106764321052880725980324389105539729}
       {4870263973198688168897112743383371412545482895917056}
  <892 .
\]
This gives the displayed bound.
\end{proof}

\begin{lemma}[Stable moments dominate the twelfth $B$ nonboundary width count above the first row]\label{lem:bc-stable-dominates-twelfth-b-nonboundary-width-count}
For every \(q\ge23\),
\[
  1784\binom{2q+12}{24}48^{24}3^{2q-12}\le s_{2q+12}.
\]
\end{lemma}

\begin{proof}
By \cref{lem:bc-littlewood-graph-model}, \(s_{2q+12}\) contains all labelled
simple cycles on \(2q+12\) vertices, whose number is \((2q+11)!/2\).  It is
enough to prove
\[
  3568\binom{2q+12}{24}48^{24}3^{2q-12}\le (2q+11)! .
\]
At \(q=23\),
\[
\begin{aligned}
  3568\binom{58}{24}48^{24}3^{34}
  &=
  17085925815413297276694044456077086105282622055811413510717226451649062502400,\\
  57!
  &=
  40526919504877216755680601905432322134980384796226602145184481280000000000000 .
\end{aligned}
\]
The first integer is smaller than the second.  If the displayed factorial
inequality holds at \(q\), then at \(q+1\) the left side is multiplied by
\[
  9\frac{(2q+14)(2q+13)}{(2q-10)(2q-11)}<27,
\]
because \(8q^2-180q+148>0\) for \(q\ge23\).  The right side is multiplied by
\[
  (2q+12)(2q+13)\ge58\cdot59>27 .
\]
Induction proves the result.
\end{proof}

\begin{lemma}[Exact first eight twelfth $B$ nonboundary bad-shape counts]\label{lem:bc-b-twelfth-frontier-badshape-certificate}
For \(15\le q\le22\), the number of column-even Pieri two-strip
paths of length \(2q+12\) with terminal first row at least \(2q\) is at most
\(s_{2q+12}/2\).
\end{lemma}

\begin{proof}
The exact C++/GMP/OpenMP replay
\[
\texttt{certificates/post\_m29/bc\_twelfth\_frontier\_gmp\_certificate.log}
\]
implements the reverse Pieri recursion of
\cref{lem:bc-b-eighth-frontier-badshape-certificate} for
\(15\le q\le22\), \(j=2q+12\).  It exits with status \(0\),
reports \texttt{SUMMARY failures=0}, and has SHA-256
\[
\hashsplit{027740d3f1fc0d56fcb27fa4e8d1f32a}{35bf7f4d0938ae970eeaadef8c4ba875}.
\]
It verifies the following exact margins:
\[
\begin{array}{c|c|c|c}
q & j & \#\mathrm{bad}(q,j) & s_j-2\#\mathrm{bad}(q,j)\\ \hline
15&42&
93189699623293897062273042455009001300&
94986302333189352668146167061074907139426098328056\\
16&44&
2944970066850265258433435574323853045000&
175604143934119061437627836401793189946527776342532016\\
17&46&
86893906123913310142005123703778431425750&
355551352238134936474545379687236110227488741701071907220\\
18&48&
2411360700038598865341416512503696751025757&
785317984503258799032608572772112825326100483465166247011078\\
19&50&
63325607009226634814812586877483266881504446&
1885339339071848339887681776626879027529213088016666282409649732\\
20&52&
1582080760407312019816362520545354254202411288&
4903266810544563254692169491537040013954573481350200637091863636048\\
21&54&
37772835182389352001906800037722047354320693120&
13771977065833964133817259347577277238395987638845438215079109063115904\\
22&56&
865247986062810751415300150963091108347321056935&
41656588897766692204622473783877622959545445539284227823934198657668262450 .
\end{array}
\]
Each margin is positive, proving the claimed half-stable bound.
\end{proof}

\begin{proposition}[Residual $B$ twelfth nonboundary lower correction]\label{prop:bc-b-twelfth-nonboundary-lower-correction}
For every residual row \(B_b\) with \(14\le b\le123\), put \(q=b+1\) and
write \(s_j=m_j^{\mathrm{st}}\) and \(m_j^{B_b}=s_j+\delta_j^{B_b}\).  Then
\[
  -\frac12 s_{2q+12}\le \delta_{2q+12}^{B_b}.
\]
\end{proposition}

\begin{proof}
By \cref{lem:bc-correction-bad-shapes}, the coefficient interpretation in
\cref{lem:bad-shape-witnesses}, and the Pieri-path bijection
\cref{lem:bc-pieri-path-bijection}, \(|\delta_{2q+12}^{B_b}|\) is the number
of column-even Pieri two-strip paths of length \(2q+12\) with terminal first
row at least \(2q\).  If \(15\le q\le22\), this is bounded by
\(s_{2q+12}/2\) by
\cref{lem:bc-b-h8-h27-bounded-littlewood-certificate}.  If \(q\ge23\),
\cref{lem:bc-b-twelfth-nonboundary-width-bad-lower-box-code} bounds this
number by
\[
  892\binom{2q+12}{24}48^{24}3^{2q-12},
\]
and \cref{lem:bc-stable-dominates-twelfth-b-nonboundary-width-count} absorbs
that count into \(s_{2q+12}/2\).  In all cases
\(|\delta_{2q+12}^{B_b}|\le s_{2q+12}/2\), which proves the stated lower
bound.
\end{proof}

\begin{lemma}[Thirteenth $B$ nonboundary width-bad paths have a lower-box code]\label{lem:bc-b-thirteenth-nonboundary-width-bad-lower-box-code}
For every \(q\ge15\), the number of column-even Pieri two-strip paths
\[
  \varnothing=\lambda^{(0)}\subset\lambda^{(1)}\subset\cdots
  \subset\lambda^{(2q+13)}
\]
with
\[
  \lambda^{(2q+13)}_1\ge2q
\]
is at most
\[
  3063\binom{2q+13}{26}52^{26}3^{2q-13}.
\]
\end{lemma}

\begin{proof}
Put \(n=2q+13\).  The terminal partition has size \(2n=4q+26\).  Since its
columns have even lengths and its first row has length at least \(2q\), the
number of boxes below the second row is at most
\[
  4q+26-2(2q)=26 .
\]
Hence the terminal partition has at most twenty-eight nonzero rows.

As in \cref{lem:bc-b-sixth-nonboundary-width-bad-lower-box-code}, encode a
path by the row-pair word and track lower boxes.  A step with no lower box has
three choices, a step with one lower box has at most \(2\cdot26=52\) choices,
and a step with two lower boxes has at most \(\binom{26+1}{2}=351\) choices.
Thus the number of paths is bounded by the sum of the coefficients of
\(x^0,\ldots,x^{26}\) in
\[
  (3+52x+351x^2)^n .
\]

The monomial with \(a\) one-lower-box steps and \(c\) two-lower-box steps,
where \(a+2c\le26\), contributes
\[
  \binom{n}{a+c}\binom{a+c}{c}3^{n-a-c}52^a351^c .
\]
Since \(n\ge43\), \(a+c\le26\), and the twenty-six factors in the quotient
\(\binom n{26}/\binom n{a+c}\) which are not cancelled are all at least
\(18\),
\[
  \frac{\binom{n}{a+c}\binom{a+c}{c}}{\binom n{26}}
  \le
  \frac{26!}{a!\,c!\,18^{26-a-c}}.
\]
Therefore the coefficient sum is at most
\[
  \binom n{26}52^{26}3^{n-26}
  \sum_{\substack{a,c\ge0\\ a+2c\le26}}
  \frac{26!}{a!\,c!\,18^{26-a-c}}
  \frac{3^{26-a-c}351^c}{52^{26-a}}.
\]
The exact replay
\[
\texttt{certificates/post\_m29/bc\_thirteenth\_frontier\_gmp\_certificate.log}
\]
prints this finite rational sum as
\[
  \frac{2578305522094518768889219360281912999537046155826506667}
       {841841424745431191739613263078552619810374977323008}
  <3063 .
\]
This gives the displayed bound.
\end{proof}

\begin{lemma}[Stable moments dominate the thirteenth $B$ nonboundary width count above the first row]\label{lem:bc-stable-dominates-thirteenth-b-nonboundary-width-count}
For every \(q\ge25\),
\[
  6126\binom{2q+13}{26}52^{26}3^{2q-13}\le s_{2q+13}.
\]
\end{lemma}

\begin{proof}
By \cref{lem:bc-littlewood-graph-model}, \(s_{2q+13}\) contains all labelled
simple cycles on \(2q+13\) vertices, whose number is \((2q+12)!/2\).  It is
enough to prove
\[
  12252\binom{2q+13}{26}52^{26}3^{2q-13}\le (2q+12)! .
\]
At \(q=25\),
\[
\begin{aligned}
  12252\binom{63}{26}52^{26}3^{37}
  &=
  814069287992760839673942139741184436931225733248105750664747413404836079102741970944,\\
  62!
  &=
  31469973260387937525653122354950764088012280797258232192163168247821107200000000000000 .
\end{aligned}
\]
The first integer is smaller than the second.  If the displayed factorial
inequality holds at \(q\), then at \(q+1\) the left side is multiplied by
\[
  9\frac{(2q+15)(2q+14)}{(2q-11)(2q-12)}<27,
\]
because \(8q^2-196q+186>0\) for \(q\ge25\).  The right side is multiplied by
\[
  (2q+13)(2q+14)\ge63\cdot64>27 .
\]
Induction proves the result.
\end{proof}

\begin{lemma}[Exact first ten thirteenth $B$ nonboundary bad-shape counts]\label{lem:bc-b-thirteenth-frontier-badshape-certificate}
For \(15\le q\le24\), the number of column-even Pieri
two-strip paths of length \(2q+13\) with terminal first row at least \(2q\) is
at most \(s_{2q+13}/2\).
\end{lemma}

\begin{proof}
The exact C++/GMP/OpenMP replay
\[
\texttt{certificates/post\_m29/bc\_thirteenth\_frontier\_gmp\_certificate.log}
\]
implements the reverse Pieri recursion of
\cref{lem:bc-b-eighth-frontier-badshape-certificate} for
\(15\le q\le24\), \(j=2q+13\).  It exits with status \(0\),
reports \texttt{SUMMARY failures=0}, and has SHA-256
\[
\hashsplit{6f6ffcfb724ca85fa51310916cd2451e}{4b38d08bcf0dd224719ce3d8b809e3e9}.
\]
It verifies the following exact margins:
\[
\begin{array}{c|c|c|c}
q & j & \#\mathrm{bad}(q,j) & s_j-2\#\mathrm{bad}(q,j)\\ \hline
15&43&
9261959617731322839939927285357323722740&
4036896628126818137834663415041705908764011647171896\\
16&45&
321085018576437107884314070871667147756360&
7814348742731877298559686751384604962672525314758219952\\
17&47&
10346577483470970761767098836420684066750148&
16533071956148070949792955355748546877109351283463013317176\\
18&49&
312316199642079406316926021804798245675155352&
38087788784726913954909000279061587421619335578715321244096208\\
19&51&
8889079614417722669504600448536456057226101910&
95209341861812270677192781064851784090595337603045433631299397780\\
20&53&
239889630231558154001173136014867176902705974938&
257420799961522256420957054048221350172003192927214384481942685846348\\
21&55&
6168070230594284508518829816148238282986612855740&
750570909947066995055412314636054651557333090807838685620396595934109448\\
22&57&
151734536869315297557940713054485939491583508893664&
2353592104550586210415883410043731104479495613189519840996942501170379956800\\
23&59&
3584278928823864849821504396330622763576445058140532&
7916863333164200256454983208040696668961580499915095656840060749633498643359384\\
24&61&
81563273739758344575205491568003548506654429986144278&
28498624571845034398726355851767445640622747933084436782816994579034437113257216724 .
\end{array}
\]
Each margin is positive, proving the claimed half-stable bound.
\end{proof}

\begin{proposition}[Residual $B$ thirteenth nonboundary lower correction]\label{prop:bc-b-thirteenth-nonboundary-lower-correction}
For every residual row \(B_b\) with \(14\le b\le123\), put \(q=b+1\) and
write \(s_j=m_j^{\mathrm{st}}\) and \(m_j^{B_b}=s_j+\delta_j^{B_b}\).  Then
\[
  -\frac12 s_{2q+13}\le \delta_{2q+13}^{B_b}.
\]
\end{proposition}

\begin{proof}
By \cref{lem:bc-correction-bad-shapes}, the coefficient interpretation in
\cref{lem:bad-shape-witnesses}, and the Pieri-path bijection
\cref{lem:bc-pieri-path-bijection}, \(|\delta_{2q+13}^{B_b}|\) is the number
of column-even Pieri two-strip paths of length \(2q+13\) with terminal first
row at least \(2q\).  If \(15\le q\le24\), this is bounded
by \(s_{2q+13}/2\) by
\cref{lem:bc-b-h8-h27-bounded-littlewood-certificate}.  If \(q\ge25\),
\cref{lem:bc-b-thirteenth-nonboundary-width-bad-lower-box-code} bounds this
number by
\[
  3063\binom{2q+13}{26}52^{26}3^{2q-13},
\]
and \cref{lem:bc-stable-dominates-thirteenth-b-nonboundary-width-count}
absorbs that count into \(s_{2q+13}/2\).  In all cases
\(|\delta_{2q+13}^{B_b}|\le s_{2q+13}/2\), which proves the stated lower
bound.
\end{proof}

\begin{lemma}[Fourteenth $B$ nonboundary width-bad paths have a lower-box code]\label{lem:bc-b-fourteenth-nonboundary-width-bad-lower-box-code}
For every \(q\ge15\), the number of column-even Pieri two-strip paths
\[
  \varnothing=\lambda^{(0)}\subset\lambda^{(1)}\subset\cdots
  \subset\lambda^{(2q+14)}
\]
with
\[
  \lambda^{(2q+14)}_1\ge2q
\]
is at most
\[
  11979\binom{2q+14}{28}56^{28}3^{2q-14}.
\]
\end{lemma}

\begin{proof}
Put \(n=2q+14\).  The terminal partition has size \(2n=4q+28\).  Since its
columns have even lengths and its first row has length at least \(2q\), the
number of boxes below the second row is at most
\[
  4q+28-2(2q)=28 .
\]
Hence the terminal partition has at most thirty nonzero rows.

As in \cref{lem:bc-b-sixth-nonboundary-width-bad-lower-box-code}, encode a
path by the row-pair word and track lower boxes.  A step with no lower box has
three choices, a step with one lower box has at most \(2\cdot28=56\) choices,
and a step with two lower boxes has at most \(\binom{28+1}{2}=406\) choices.
Thus the number of paths is bounded by the sum of the coefficients of
\(x^0,\ldots,x^{28}\) in
\[
  (3+56x+406x^2)^n .
\]

The monomial with \(a\) one-lower-box steps and \(c\) two-lower-box steps,
where \(a+2c\le28\), contributes
\[
  \binom{n}{a+c}\binom{a+c}{c}3^{n-a-c}56^a406^c .
\]
Since \(n\ge44\), \(a+c\le28\), and the twenty-eight factors in the quotient
\(\binom n{28}/\binom n{a+c}\) which are not cancelled are all at least
\(17\),
\[
  \frac{\binom{n}{a+c}\binom{a+c}{c}}{\binom n{28}}
  \le
  \frac{28!}{a!\,c!\,17^{28-a-c}}.
\]
Therefore the coefficient sum is at most
\[
  \binom n{28}56^{28}3^{n-28}
  \sum_{\substack{a,c\ge0\\ a+2c\le28}}
  \frac{28!}{a!\,c!\,17^{28-a-c}}
  \frac{3^{28-a-c}406^c}{56^{28-a}}.
\]
The exact replay
\[
\texttt{certificates/post\_m29/bc\_fourteenth\_frontier\_gmp\_certificate.log}
\]
prints this finite rational sum as
\[
  \frac{2206251862522460035843841746796935620083943371000444640299878209563705578207}
       {184180472053899552005790127718892583988774500112337828541287665105895424}
  <11979 .
\]
This gives the displayed bound.
\end{proof}

\begin{lemma}[Stable moments dominate the fourteenth $B$ nonboundary width count above the first row]\label{lem:bc-stable-dominates-fourteenth-b-nonboundary-width-count}
For every \(q\ge26\),
\[
  23958\binom{2q+14}{28}56^{28}3^{2q-14}\le s_{2q+14}.
\]
\end{lemma}

\begin{proof}
By \cref{lem:bc-littlewood-graph-model}, \(s_{2q+14}\) contains all labelled
simple cycles on \(2q+14\) vertices, whose number is \((2q+13)!/2\).  It is
enough to prove
\[
  47916\binom{2q+14}{28}56^{28}3^{2q-14}\le (2q+13)! .
\]
At \(q=26\),
\[
\begin{aligned}
  47916\binom{66}{28}56^{28}3^{38}
  &=
  1965920218394549173531718720316381901912412245557438940932955044811503832120275636615905280,\\
  65!
  &=
  8247650592082470666723170306785496252186258551345437492922123134388955774976000000000000000 .
\end{aligned}
\]
The first integer is smaller than the second.  If the displayed factorial
inequality holds at \(q\), then at \(q+1\) the left side is multiplied by
\[
  9\frac{(2q+16)(2q+15)}{(2q-12)(2q-13)}<27,
\]
because \(8q^2-212q+228>0\) for \(q\ge26\).  The right side is multiplied by
\[
  (2q+14)(2q+15)\ge66\cdot67>27 .
\]
Induction proves the result.
\end{proof}

\begin{lemma}[Exact first eleven fourteenth $B$ nonboundary bad-shape counts]\label{lem:bc-b-fourteenth-frontier-badshape-certificate}
For \(15\le q\le25\), the number of column-even Pieri
two-strip paths of length \(2q+14\) with terminal first row at least \(2q\) is
at most \(s_{2q+14}/2\).
\end{lemma}

\begin{proof}
The exact C++/GMP/OpenMP replay
\[
\texttt{certificates/post\_m29/bc\_fourteenth\_frontier\_gmp\_certificate.log}
\]
implements the reverse Pieri recursion of
\cref{lem:bc-b-eighth-frontier-badshape-certificate} for
\(15\le q\le25\), \(j=2q+14\).  It exits with status
\(0\), reports \texttt{SUMMARY failures=0}, and has SHA-256
\[
\hashsplit{315dcc3f51bd178dba5b186d4cb4f3e0}{0c636284f3b43d8428915d46a90af82b}.
\]
It verifies the following exact margins:
\[
\begin{array}{c|c|c|c}
q & j & \#\mathrm{bad}(q,j) & s_j-2\#\mathrm{bad}(q,j)\\ \hline
15&44&
904241499510894118795869663481520405775720&
175604143932316468378739748694718317490713383237070576\\
16&46&
34316545579074576964762425743302944534892140&
355551352238066477171199478359926869386249543368864974440\\
17&48&
1205344532427555321349524816513886234415737008&
785317984503256393166265117738667856959300480700090917588576\\
18&50&
39505255996167644461782446721379011147962343549&
1885339339071848261003820998310043373593944819013610520247971526\\
19&52&
1216565694103918841752398350801977668931714967796&
4903266810544563252262202264850016954489409504787336007736838523032\\
20&54&
35409205390444139344776401846463410341856783154108&
13771977065833964133746516482466753738410438648752586838490104138193928\\
21&56&
979048556886031667100003765498860377735567368606600&
41656588897766692204620517417259823021832748362353532029360944217573163120\\
22&58&
25829478471615497730368142419286411973382961058385799&
135331264354280596525260343372952908465292572072808415083728761256012171892850\\
23&60&
652716666090191470295272747661134251183413102876318760&
471052485162704317319812794664974957053085989881678780643648524177163566557148336\\
24&62&
15852716879610122946050266806220734180535400568392792238&
1752662440577781575209503668206421439574564463157839529802025182010505987910969154980\\
25&64&
371154395327141239391075310169894302656286034605646431864&
6955857144727139970751971374392861008233250096453707124089936965869531818002462773466128 .
\end{array}
\]
Each margin is positive, proving the claimed half-stable bound.
\end{proof}

\begin{proposition}[Residual $B$ fourteenth nonboundary lower correction]\label{prop:bc-b-fourteenth-nonboundary-lower-correction}
For every residual row \(B_b\) with \(14\le b\le123\), put \(q=b+1\) and
write \(s_j=m_j^{\mathrm{st}}\) and \(m_j^{B_b}=s_j+\delta_j^{B_b}\).  Then
\[
  -\frac12 s_{2q+14}\le \delta_{2q+14}^{B_b}.
\]
\end{proposition}

\begin{proof}
By \cref{lem:bc-correction-bad-shapes}, the coefficient interpretation in
\cref{lem:bad-shape-witnesses}, and the Pieri-path bijection
\cref{lem:bc-pieri-path-bijection}, \(|\delta_{2q+14}^{B_b}|\) is the number
of column-even Pieri two-strip paths of length \(2q+14\) with terminal first
row at least \(2q\).  If \(15\le q\le25\), this is
bounded by \(s_{2q+14}/2\) by
\cref{lem:bc-b-h8-h27-bounded-littlewood-certificate}.  If \(q\ge26\),
\cref{lem:bc-b-fourteenth-nonboundary-width-bad-lower-box-code} bounds this
number by
\[
  11979\binom{2q+14}{28}56^{28}3^{2q-14},
\]
and \cref{lem:bc-stable-dominates-fourteenth-b-nonboundary-width-count}
absorbs that count into \(s_{2q+14}/2\).  In all cases
\(|\delta_{2q+14}^{B_b}|\le s_{2q+14}/2\), which proves the stated lower
bound.
\end{proof}

\begin{lemma}[Fifteenth $B$ nonboundary width-bad paths have a lower-box code]\label{lem:bc-b-fifteenth-nonboundary-width-bad-lower-box-code}
For every \(q\ge15\), the number of column-even Pieri two-strip paths
\[
  \varnothing=\lambda^{(0)}\subset\lambda^{(1)}\subset\cdots
  \subset\lambda^{(2q+15)}
\]
with
\[
  \lambda^{(2q+15)}_1\ge2q
\]
is at most
\[
  53918\binom{2q+15}{30}60^{30}3^{2q-15}.
\]
\end{lemma}

\begin{proof}
Put \(n=2q+15\).  The terminal partition has size \(2n=4q+30\).  Since its
columns have even lengths and its first row has length at least \(2q\), the
number of boxes below the second row is at most
\[
  4q+30-2(2q)=30 .
\]
Hence the terminal partition has at most thirty-two nonzero rows.

As in \cref{lem:bc-b-sixth-nonboundary-width-bad-lower-box-code}, encode a
path by the row-pair word and track lower boxes.  A step with no lower box has
three choices, a step with one lower box has at most \(2\cdot30=60\) choices,
and a step with two lower boxes has at most \(\binom{30+1}{2}=465\) choices.
Thus the number of paths is bounded by the sum of the coefficients of
\(x^0,\ldots,x^{30}\) in
\[
  (3+60x+465x^2)^n .
\]

The monomial with \(a\) one-lower-box steps and \(c\) two-lower-box steps,
where \(a+2c\le30\), contributes
\[
  \binom{n}{a+c}\binom{a+c}{c}3^{n-a-c}60^a465^c .
\]
Since \(n\ge45\), \(a+c\le30\), and the thirty factors in the quotient
\(\binom n{30}/\binom n{a+c}\) which are not cancelled are all at least
\(16\),
\[
  \frac{\binom{n}{a+c}\binom{a+c}{c}}{\binom n{30}}
  \le
  \frac{30!}{a!\,c!\,16^{30-a-c}}.
\]
Therefore the coefficient sum is at most
\[
  \binom n{30}60^{30}3^{n-30}
  \sum_{\substack{a,c\ge0\\ a+2c\le30}}
  \frac{30!}{a!\,c!\,16^{30-a-c}}
  \frac{3^{30-a-c}465^c}{60^{30-a}}.
\]
The exact replay
\[
\texttt{certificates/post\_m29/bc\_fifteenth\_frontier\_gmp\_certificate.log}
\]
prints this finite rational sum as
\[
  \frac{14677725404098265068620724855748750166905094501498612189320160692921}
       {272225893536750770770699685945414569164800000000000000000000000}
  <53918 .
\]
This gives the displayed bound.
\end{proof}

\begin{lemma}[Stable moments dominate the fifteenth $B$ nonboundary width count above the first row]\label{lem:bc-stable-dominates-fifteenth-b-nonboundary-width-count}
For every \(q\ge28\),
\[
  107836\binom{2q+15}{30}60^{30}3^{2q-15}\le s_{2q+15}.
\]
\end{lemma}

\begin{proof}
By \cref{lem:bc-littlewood-graph-model}, \(s_{2q+15}\) contains all labelled
simple cycles on \(2q+15\) vertices, whose number is \((2q+14)!/2\).  It is
enough to prove
\[
  215672\binom{2q+15}{30}60^{30}3^{2q-15}\le (2q+14)! .
\]
At \(q=28\),
\[
\begin{aligned}
  215672\binom{71}{30}60^{30}3^{41}
  &=
  166677548926834328215778899144982282353333778858586866364152952127488000000000000000000000000000000,\\
  70!
  &=
  11978571669969891796072783721689098736458938142546425857555362864628009582789845319680000000000000000 .
\end{aligned}
\]
The first integer is smaller than the second.  If the displayed factorial
inequality holds at \(q\), then at \(q+1\) the left side is multiplied by
\[
  9\frac{(2q+17)(2q+16)}{(2q-13)(2q-14)}<27,
\]
because \(8q^2-228q+274>0\) for \(q\ge28\).  The right side is multiplied by
\[
  (2q+15)(2q+16)\ge71\cdot72>27 .
\]
Induction proves the result.
\end{proof}

\begin{lemma}[Exact first thirteen fifteenth $B$ nonboundary bad-shape counts]\label{lem:bc-b-fifteenth-frontier-badshape-certificate}
For \(15\le q\le27\), the number of
column-even Pieri two-strip paths of length \(2q+15\) with terminal first row
at least \(2q\) is at most \(s_{2q+15}/2\).
\end{lemma}

\begin{proof}
The exact C++/GMP/OpenMP replay
\[
\texttt{certificates/post\_m29/bc\_fifteenth\_frontier\_gmp\_certificate.log}
\]
implements the reverse Pieri recursion of
\cref{lem:bc-b-eighth-frontier-badshape-certificate} for
\(15\le q\le27\), \(j=2q+15\).  It exits with
status \(0\), reports \texttt{SUMMARY failures=0}, and has SHA-256
\[
\hashsplit{5f829fb541bf8d84f097ff14aa6be272}{508554407503bd74b37532e62c1484b9}.
\]
It verifies the following exact margins:
\[
\begin{array}{c|c|c|c}
q & j & \#\mathrm{bad}(q,j) & s_j-2\#\mathrm{bad}(q,j)\\ \hline
15&45&
87003920996552906597257682496265484471080770&
7814348742558511626603733812405858225821737680111571132\\
16&47&
3607536125230226666736550792452423000450904113&
16533071956140876570697461844356597309722119278830245009246\\
17&49&
137866801862434269213014351322071153961572141304&
38087788784726638845937674694681974026768735046003889450124304\\
18&51&
4897938044505836438911625788422682007179764173625&
95209341861812260899094851282014351606352961654754331386223254350\\
19&53&
162938890183785316560360745867975853262168456585696&
257420799961522256095559052941113833359284047463292412311411184624832\\
20&55&
5107253627350791334756002476104885485190100308831416&
750570909947066995045210143521814257456838123515261211126582368542158096\\
21&57&
151642695284223916189224144734975201369045203750659540&
2353592104550586210415580428122236395277712280781475999566083393929896425048\\
22&59&
4284903864622671457145546904202741371329180364475710597&
7916863333164200256454974645401525281266395908464296044018563618425659808219254\\
23&61&
115692900449087199221605049442686667091810004951400197760&
28498624571845034398726355620544771289927866179024748880579667492427736070429109760\\
24&63&
2995482956282684852310456475723829592541709520358042539396&
109541225789834175931658495490412553634851132915530955435886810144674617471147133041144\\
25&65&
74609001886847997188501399507450228099031611764174645505536&
448652128232881537732129952277643155760149930451547351191388653100147393271097763359032320\\
26&67&
1792659116038401166877320912707051090756263525315000428518861&
1954211164681443257926522972458387998748966177409228472342412732878018653927895874991504704870\\
27&69&
41655688006298221943407882203261078530441457103239477150150316&
9035761317088612444238991255468019372218605875608789660379001392822365211344779616325486454823080 .
\end{array}
\]
Each margin is positive, proving the claimed half-stable bound.
\end{proof}

\begin{proposition}[Residual $B$ fifteenth nonboundary lower correction]\label{prop:bc-b-fifteenth-nonboundary-lower-correction}
For every residual row \(B_b\) with \(14\le b\le123\), put \(q=b+1\) and
write \(s_j=m_j^{\mathrm{st}}\) and \(m_j^{B_b}=s_j+\delta_j^{B_b}\).  Then
\[
  -\frac12 s_{2q+15}\le \delta_{2q+15}^{B_b}.
\]
\end{proposition}

\begin{proof}
By \cref{lem:bc-correction-bad-shapes}, the coefficient interpretation in
\cref{lem:bad-shape-witnesses}, and the Pieri-path bijection
\cref{lem:bc-pieri-path-bijection}, \(|\delta_{2q+15}^{B_b}|\) is the number
of column-even Pieri two-strip paths of length \(2q+15\) with terminal first
row at least \(2q\).  If \(15\le q\le27\), this
is bounded by \(s_{2q+15}/2\) by
\cref{lem:bc-b-h8-h27-bounded-littlewood-certificate}.  If \(q\ge28\),
\cref{lem:bc-b-fifteenth-nonboundary-width-bad-lower-box-code} bounds this
number by
\[
  53918\binom{2q+15}{30}60^{30}3^{2q-15},
\]
and \cref{lem:bc-stable-dominates-fifteenth-b-nonboundary-width-count}
absorbs that count into \(s_{2q+15}/2\).  In all cases
\(|\delta_{2q+15}^{B_b}|\le s_{2q+15}/2\), which proves the stated lower
bound.
\end{proof}

\begin{lemma}[Sixteenth $B$ nonboundary width-bad paths have a lower-box code]\label{lem:bc-b-sixteenth-nonboundary-width-bad-lower-box-code}
For every \(q\ge15\), the number of column-even Pieri two-strip paths
\[
  \varnothing=\lambda^{(0)}\subset\lambda^{(1)}\subset\cdots
  \subset\lambda^{(2q+16)}
\]
with
\[
  \lambda^{(2q+16)}_1\ge2q
\]
is at most
\[
  282957\binom{2q+16}{32}64^{32}3^{2q-16}.
\]
\end{lemma}

\begin{proof}
Put \(n=2q+16\).  The terminal partition has size \(2n=4q+32\).  Since its
columns have even lengths and its first row has length at least \(2q\), the
number of boxes below the second row is at most
\[
  4q+32-2(2q)=32 .
\]
Hence the terminal partition has at most thirty-four nonzero rows.

As in \cref{lem:bc-b-sixth-nonboundary-width-bad-lower-box-code}, encode a
path by the row-pair word and track lower boxes.  A step with no lower box has
three choices, a step with one lower box has at most \(2\cdot32=64\) choices,
and a step with two lower boxes has at most \(\binom{32+1}{2}=528\) choices.
Thus the number of paths is bounded by the sum of the coefficients of
\(x^0,\ldots,x^{32}\) in
\[
  (3+64x+528x^2)^n .
\]

The monomial with \(a\) one-lower-box steps and \(c\) two-lower-box steps,
where \(a+2c\le32\), contributes
\[
  \binom{n}{a+c}\binom{a+c}{c}3^{n-a-c}64^a528^c .
\]
Since \(n\ge46\), \(a+c\le32\), and the thirty-two factors in the quotient
\(\binom n{32}/\binom n{a+c}\) which are not cancelled are all at least
\(15\),
\[
  \frac{\binom{n}{a+c}\binom{a+c}{c}}{\binom n{32}}
  \le
  \frac{32!}{a!\,c!\,15^{32-a-c}}.
\]
Therefore the coefficient sum is at most
\[
  \binom n{32}64^{32}3^{n-32}
  \sum_{\substack{a,c\ge0\\ a+2c\le32}}
  \frac{32!}{a!\,c!\,15^{32-a-c}}
  \frac{3^{32-a-c}528^c}{64^{32-a}}.
\]
The exact replay
\[
\texttt{certificates/post\_m29/bc\_sixteenth\_frontier\_gmp\_certificate.log}
\]
prints this finite rational sum as
\[
  \frac{246489541168394136468483634206198590453034992597677288160669567006157511}
       {871122859317602466466238995025326621327360000000000000000000000000}
  <282957 .
\]
This gives the displayed bound.
\end{proof}

\begin{lemma}[Stable moments dominate the sixteenth $B$ nonboundary width count above the first row]\label{lem:bc-stable-dominates-sixteenth-b-nonboundary-width-count}
For every \(q\ge29\),
\[
  565914\binom{2q+16}{32}64^{32}3^{2q-16}\le s_{2q+16}.
\]
\end{lemma}

\begin{proof}
By \cref{lem:bc-littlewood-graph-model}, \(s_{2q+16}\) contains all labelled
simple cycles on \(2q+16\) vertices, whose number is \((2q+15)!/2\).  It is
enough to prove
\[
  1131828\binom{2q+16}{32}64^{32}3^{2q-16}\le (2q+15)! .
\]
At \(q=29\),
\[
\begin{aligned}
  1131828\binom{74}{32}64^{32}3^{42}
  &=
  695557115329569709224150063682778061838037342308385804626953595512549529256644715070119059644382734974976,\\
  73!
  &=
  4470115461512684340891257138125051110076800700282905015819080092370422104067183317016903680000000000000000 .
\end{aligned}
\]
The first integer is smaller than the second.  If the displayed factorial
inequality holds at \(q\), then at \(q+1\) the left side is multiplied by
\[
  9\frac{(2q+18)(2q+17)}{(2q-14)(2q-15)}<30,
\]
because \(84q^2-2370q+3546>0\) for \(q\ge29\).  The right side is multiplied by
\[
  (2q+16)(2q+17)\ge74\cdot75>30 .
\]
Induction proves the result.
\end{proof}

\begin{lemma}[Exact first fourteen sixteenth $B$ nonboundary bad-shape counts]\label{lem:bc-b-sixteenth-frontier-badshape-certificate}
For \(15\le q\le28\), the number of
column-even Pieri two-strip paths of length \(2q+16\) with terminal first row
at least \(2q\) is at most \(s_{2q+16}/2\).
\end{lemma}

\begin{proof}
The exact C++/GMP/OpenMP replay
\[
\texttt{certificates/post\_m29/bc\_sixteenth\_frontier\_gmp\_certificate.log}
\]
implements the reverse Pieri recursion of
\cref{lem:bc-b-eighth-frontier-badshape-certificate} for
\(15\le q\le28\), \(j=2q+16\).  It exits with
status \(0\), reports \texttt{SUMMARY failures=0}, and has SHA-256
\[
\hashsplit{9276304d5e5aff2105d4267604664e9f}{14dbcc42df60e0a643c335b028250009}.
\]
It verifies the following exact margins:
\[
\begin{array}{c|c|c|c}
q & j & \#\mathrm{bad}(q,j) & s_j-2\#\mathrm{bad}(q,j)\\ \hline
15&46&
8272983354280527593468052465966387975898819215&
355551352221589143553796572326920289305803373306137120290\\
16&48&
374098023898022085704780517499191373641504735747&
785317984502510607807533928677900994973935125725276739591098\\
17&50&
15528496562219269103802590570633462837432040098832&
1885339339071817283021208552107124691977696994845957952092460960\\
18&52&
597029860046538795531886347547448881528802399385354&
4903266810544562060635613559980263574221511111493528287995469687916\\
19&54&
21423293633682332357729677671046428807332880396861126&
13771977065833964090970747625882977301640636110352656044508056910779892\\
20&56&
722136871516149082858615658006670079745042801148811140&
41656588897766692203178201771341296919449717053871189590626329750012754040\\
21&58&
22994558956786999982315266809317628071410159794070017603&
135331264354280596525214405913996277696323402275474618400409887702346148629242\\
22&60&
695043507649755527397768535849816623690533710220113245909&
471052485162704317319811405883392989722414134935152576332669645476569332083294038\\
23&62&
20027377511752658955477085619459584786840264485171632719887&
1752662440577781575209503628183371849828466798095768824495297076691047818704489299682\\
24&64&
552183555309122185320016983329451126500219486553090833065783&
6955857144727139970751971373289236206405660008292455308051374501474405416965492400198290\\
25&66&
14615891637018864121415359376808915816497044412622226036971340&
29386673299901708964749475366730337406328925449922948207233125442364428752735436190800851560\\
26&68&
372499764620854488855056835541977547662825473236779430400948538&
131909085289227211701789514679798183330284404718062127535641201435022833027609889543630800130956\\
27&70&
9164865545096087327418774390016633212464178987982649947882888584&
627984678371728270376818250246812063500373154593071883732187585120608215213669982361566356229572336\\
28&72&
218197243528751072881974686848428246042581074443937961025749863594&
3165506839235730069039027353629446009170081519012503274768331305040780096639467941388157166991803039916 .
\end{array}
\]
Each margin is positive, proving the claimed half-stable bound.
\end{proof}

\begin{proposition}[Residual $B$ sixteenth nonboundary lower correction]\label{prop:bc-b-sixteenth-nonboundary-lower-correction}
For every residual row \(B_b\) with \(14\le b\le123\), put \(q=b+1\) and
write \(s_j=m_j^{\mathrm{st}}\) and \(m_j^{B_b}=s_j+\delta_j^{B_b}\).  Then
\[
  -\frac12 s_{2q+16}\le \delta_{2q+16}^{B_b}.
\]
\end{proposition}

\begin{proof}
By \cref{lem:bc-correction-bad-shapes}, the coefficient interpretation in
\cref{lem:bad-shape-witnesses}, and the Pieri-path bijection
\cref{lem:bc-pieri-path-bijection}, \(|\delta_{2q+16}^{B_b}|\) is the number
of column-even Pieri two-strip paths of length \(2q+16\) with terminal first
row at least \(2q\).  If \(15\le q\le28\), this
is bounded by \(s_{2q+16}/2\) by
\cref{lem:bc-b-h8-h27-bounded-littlewood-certificate}.  If \(q\ge29\),
\cref{lem:bc-b-sixteenth-nonboundary-width-bad-lower-box-code} bounds this
number by
\[
  282957\binom{2q+16}{32}64^{32}3^{2q-16},
\]
and \cref{lem:bc-stable-dominates-sixteenth-b-nonboundary-width-count}
absorbs that count into \(s_{2q+16}/2\).  In all cases
\(|\delta_{2q+16}^{B_b}|\le s_{2q+16}/2\), which proves the stated lower
bound.
\end{proof}

\begin{lemma}[Seventeenth $B$ nonboundary width-bad paths have a lower-box code]\label{lem:bc-b-seventeenth-nonboundary-width-bad-lower-box-code}
For every \(q\ge15\), the number of column-even Pieri two-strip paths
\[
  \varnothing=\lambda^{(0)}\subset\lambda^{(1)}\subset\cdots
  \subset\lambda^{(2q+17)}
\]
with
\[
  \lambda^{(2q+17)}_1\ge2q
\]
is at most
\[
  1759061\binom{2q+17}{34}68^{34}3^{2q-17}.
\]
\end{lemma}

\begin{proof}
Put \(n=2q+17\).  The terminal partition has size \(2n=4q+34\).  Since its
columns have even lengths and its first row has length at least \(2q\), the
number of boxes below the second row is at most
\[
  4q+34-2(2q)=34 .
\]
Hence the terminal partition has at most thirty-six nonzero rows.

As in \cref{lem:bc-b-sixth-nonboundary-width-bad-lower-box-code}, encode a
path by the row-pair word and track lower boxes.  A step with no lower box has
three choices, a step with one lower box has at most \(2\cdot34=68\) choices,
and a step with two lower boxes has at most \(\binom{34+1}{2}=595\) choices.
Thus the number of paths is bounded by the sum of the coefficients of
\(x^0,\ldots,x^{34}\) in
\[
  (3+68x+595x^2)^n .
\]

The monomial with \(a\) one-lower-box steps and \(c\) two-lower-box steps,
where \(a+2c\le34\), contributes
\[
  \binom{n}{a+c}\binom{a+c}{c}3^{n-a-c}68^a595^c .
\]
Since \(n\ge47\), \(a+c\le34\), and the thirty-four factors in the quotient
\(\binom n{34}/\binom n{a+c}\) which are not cancelled are all at least
\(14\),
\[
  \frac{\binom{n}{a+c}\binom{a+c}{c}}{\binom n{34}}
  \le
  \frac{34!}{a!\,c!\,14^{34-a-c}}.
\]
Therefore the coefficient sum is at most
\[
  \binom n{34}68^{34}3^{n-34}
  \sum_{\substack{a,c\ge0\\ a+2c\le34}}
  \frac{34!}{a!\,c!\,14^{34-a-c}}
  \frac{3^{34-a-c}595^c}{68^{34-a}}.
\]
The exact replay
\[
\texttt{certificates/post\_m29/bc\_seventeenth\_frontier\_gmp\_certificate.log}
\]
prints this finite rational sum as
\[
  \frac{110837640989552376651072260993779108203216051363825886336192281558360154094119521897428034229}
       {63009550437898855669184790546695665525474533007607651696497530576675910920262370983936}
  <1759061 .
\]
This gives the displayed bound.
\end{proof}

\begin{lemma}[Stable moments dominate the seventeenth $B$ nonboundary width count above the first row]\label{lem:bc-stable-dominates-seventeenth-b-nonboundary-width-count}
For every \(q\ge31\),
\[
  3518122\binom{2q+17}{34}68^{34}3^{2q-17}\le s_{2q+17}.
\]
\end{lemma}

\begin{proof}
By \cref{lem:bc-littlewood-graph-model}, \(s_{2q+17}\) contains all labelled
simple cycles on \(2q+17\) vertices, whose number is \((2q+16)!/2\).  It is
enough to prove
\[
  7036244\binom{2q+17}{34}68^{34}3^{2q-17}\le (2q+16)! .
\]
At \(q=31\),
\[
\begin{aligned}
  7036244\binom{79}{34}68^{34}3^{45}
  &=
  106355916925065248010489952553960701708665860237767524896895162233289562428870769777751176444776643393303389143040,\\
  78!
  &=
  11324281178206297831457521158732046228731749579488251990048962825668835325234200766245086213177344000000000000000000 .
\end{aligned}
\]
The first integer is smaller than the second.  If the displayed factorial
inequality holds at \(q\), then at \(q+1\) the left side is multiplied by
\[
  9\frac{(2q+19)(2q+18)}{(2q-15)(2q-16)}<30,
\]
because \(84q^2-2526q+4122>0\) for \(q\ge31\).  The right side is multiplied
by
\[
  (2q+17)(2q+18)\ge79\cdot80>30 .
\]
Induction proves the result.
\end{proof}

\begin{lemma}[Exact first sixteen seventeenth $B$ nonboundary bad-shape counts]\label{lem:bc-b-seventeenth-frontier-badshape-certificate}
For \(15\le q\le30\), the number of
column-even Pieri two-strip paths of length \(2q+17\) with terminal first row
at least \(2q\) is at most \(s_{2q+17}/2\).
\end{lemma}

\begin{proof}
The exact C++/GMP/OpenMP replay
\[
\texttt{certificates/post\_m29/bc\_seventeenth\_frontier\_gmp\_certificate.log}
\]
implements the reverse Pieri recursion of
\cref{lem:bc-b-eighth-frontier-badshape-certificate} for
\(15\le q\le30\), \(j=2q+17\).  It exits
with status \(0\), reports \texttt{SUMMARY failures=0}, and has SHA-256
\[
\hashsplit{a1241f73bdcbce52caf77e4ec67c6ef4}{44661548ce303dfae09cf932a1a100a1}.
\]
It verifies the following exact margins:
\[
\begin{array}{c|c|c|c}
q & j & \#\mathrm{bad}(q,j) & s_j-2\#\mathrm{bad}(q,j)\\ \hline
15&47&
779230830328282566194944718451483846088992348189&
16533071954589629982291357165300180974404056432653162121094\\
16&49&
38360387184944232391876469096785250109518667246360&
38087788784650193805171511098436647117277808688092775259914192\\
17&51&
1726692115725507500481070313173742205607315188935400&
95209341861808817310739489278686267288978191015707131115373730800\\
18&53&
71735386427900995474187612996166772277800826683386414&
257420799961522112950663977506693518104779546865699563234094731023396\\
19&55&
2772571726861316706277840896757520378655743571097478742&
750570909947066989510281197053882427570668334952430224785475426964863444\\
20&57&
100370231018083149740520465136960860287328877535131642104&
2353592104550586210215143251476638543629049798797024229394163729267134459920\\
21&59&
3423218908643343883538262352580350165137385093654041899679&
7916863333164200256448136777391967838842233674852943749171031506599080675841090\\
22&61&
110553640659569321617408193271531514298313876556448290129830&
28498624571845034398726134744649253049459029805848304702924405048294633076649245620\\
23&63&
3395819242242053528800626128774330586968419239274333313082655&
109541225789834175931658488704765035063309445018899610838673295392919557963196591954626\\
24&65&
99596157540943834973370831784468736492783160020751304567686720&
448652128232881537732129952078600058682035956499182690421466080570779136453123503514669952\\
25&67&
2798759702591853244197525420893125257124671570527511663162590028&
1954211164681443257926522972452794064662014547723167176142451896465950823313891481666036562536\\
26&69&
75586918786910147053521816336294231373415979198115500561207963742&
9035761317088612444238991255467868281692408067911126503562093210881775440269297591803318339196228\\
27&71&
1967329436741387607092007799580668990476115949062045477478124752560&
44272870337243698081214185227207888743419022428245108656246490836170225063557283909493476317385922720\\
28&73&
49468245069462752490490464045768006352591715268463435477641465370248&
229499010251335973773574243350048229184393644354037258389257905108061108935923733607882120367436505963248\\
29&75&
1204367164349240206919084144787606267144533904932810904500278351782883&
1256676780468867042847317084437511663920163009273552752219055923258693566158256435112508748529070610537183802\\
30&77&
28447654093505078593061960117489559777962687009967279148071918716382106&
7258238000562610837751823684753927597930978364363379670268537714181494663142776845562488510899955352263791025356 .
\end{array}
\]
Each margin is positive, proving the claimed half-stable bound.
\end{proof}

\begin{proposition}[Residual $B$ seventeenth nonboundary lower correction]\label{prop:bc-b-seventeenth-nonboundary-lower-correction}
For every residual row \(B_b\) with \(14\le b\le123\), put \(q=b+1\) and
write \(s_j=m_j^{\mathrm{st}}\) and \(m_j^{B_b}=s_j+\delta_j^{B_b}\).  Then
\[
  -\frac12 s_{2q+17}\le \delta_{2q+17}^{B_b}.
\]
\end{proposition}

\begin{proof}
By \cref{lem:bc-correction-bad-shapes}, the coefficient interpretation in
\cref{lem:bad-shape-witnesses}, and the Pieri-path bijection
\cref{lem:bc-pieri-path-bijection}, \(|\delta_{2q+17}^{B_b}|\) is the number
of column-even Pieri two-strip paths of length \(2q+17\) with terminal first
row at least \(2q\).  If
\(15\le q\le30\), this is bounded by
\(s_{2q+17}/2\) by
\cref{lem:bc-b-h8-h27-bounded-littlewood-certificate}.  If \(q\ge31\),
\cref{lem:bc-b-seventeenth-nonboundary-width-bad-lower-box-code} bounds this
number by
\[
  1759061\binom{2q+17}{34}68^{34}3^{2q-17},
\]
and \cref{lem:bc-stable-dominates-seventeenth-b-nonboundary-width-count}
absorbs that count into \(s_{2q+17}/2\).  In all cases
\(|\delta_{2q+17}^{B_b}|\le s_{2q+17}/2\), which proves the stated lower
bound.
\end{proof}

\begin{lemma}[Eighteenth $B$ nonboundary width-bad paths have a lower-box code]\label{lem:bc-b-eighteenth-nonboundary-width-bad-lower-box-code}
For every \(q\ge15\), the number of column-even Pieri two-strip paths
\[
  \varnothing=\lambda^{(0)}\subset\lambda^{(1)}\subset\cdots
  \subset\lambda^{(2q+18)}
\]
with
\[
  \lambda^{(2q+18)}_1\ge2q
\]
is at most
\[
  13210035\binom{2q+18}{36}72^{36}3^{2q-18}.
\]
\end{lemma}

\begin{proof}
Put \(n=2q+18\).  The terminal partition has size \(2n=4q+36\).  Since its
columns have even lengths and its first row has length at least \(2q\), the
number of boxes below the second row is at most
\[
  4q+36-2(2q)=36 .
\]
Hence the terminal partition has at most thirty-eight nonzero rows.

As in \cref{lem:bc-b-sixth-nonboundary-width-bad-lower-box-code}, encode a
path by the row-pair word and track lower boxes.  A step with no lower box has
three choices, a step with one lower box has at most \(2\cdot36=72\) choices,
and a step with two lower boxes has at most \(\binom{36+1}{2}=666\) choices.
Thus the number of paths is bounded by the sum of the coefficients of
\(x^0,\ldots,x^{36}\) in
\[
  (3+72x+666x^2)^n .
\]

The monomial with \(a\) one-lower-box steps and \(c\) two-lower-box steps,
where \(a+2c\le36\), contributes
\[
  \binom{n}{a+c}\binom{a+c}{c}3^{n-a-c}72^a666^c .
\]
Since \(n\ge48\), \(a+c\le36\), and the thirty-six factors in the quotient
\(\binom n{36}/\binom n{a+c}\) which are not cancelled are all at least
\(13\),
\[
  \frac{\binom{n}{a+c}\binom{a+c}{c}}{\binom n{36}}
  \le
  \frac{36!}{a!\,c!\,13^{36-a-c}}.
\]
Therefore the coefficient sum is at most
\[
  \binom n{36}72^{36}3^{n-36}
  \sum_{\substack{a,c\ge0\\ a+2c\le36}}
  \frac{36!}{a!\,c!\,13^{36-a-c}}
  \frac{3^{36-a-c}666^c}{72^{36-a}}.
\]
The exact replay
\[
\texttt{certificates/post\_m29/bc\_eighteenth\_frontier\_gmp\_certificate.log}
\]
prints this finite rational sum as
\[
  \frac{21702084974242567341489161043883418505255308117330855762905856283433984937273}
       {1642848430676142439453686466493927838235149844032662468884099699310592}
  <13210035 .
\]
This gives the displayed bound.
\end{proof}

\begin{lemma}[Stable moments dominate the eighteenth $B$ nonboundary width count above the first row]\label{lem:bc-stable-dominates-eighteenth-b-nonboundary-width-count}
For every \(q\ge32\),
\[
  26420070\binom{2q+18}{36}72^{36}3^{2q-18}\le s_{2q+18}.
\]
\end{lemma}

\begin{proof}
By \cref{lem:bc-littlewood-graph-model}, \(s_{2q+18}\) contains all labelled
simple cycles on \(2q+18\) vertices, whose number is \((2q+17)!/2\).  It is
enough to prove
\[
  52840140\binom{2q+18}{36}72^{36}3^{2q-18}\le (2q+17)! .
\]
At \(q=32\),
\[
\begin{aligned}
  52840140\binom{82}{36}72^{36}3^{46}
  &=
  795121139679439680748590041906633953130656619006560341265646903093524728495031336078488722393875949486369900038245580800,\\
  81!
  &=
  5797126020747367985879734231578109105412357244731625958745865049716390179693892056256184534249745940480000000000000000000 .
\end{aligned}
\]
The first integer is smaller than the second.  If the displayed factorial
inequality holds at \(q\), then at \(q+1\) the left side is multiplied by
\[
  9\frac{(2q+20)(2q+19)}{(2q-16)(2q-17)}<30,
\]
because \(84q^2-2682q+4740>0\) for \(q\ge32\).  The right side is multiplied
by
\[
  (2q+18)(2q+19)\ge82\cdot83>30 .
\]
Induction proves the result.
\end{proof}

\begin{lemma}[Exact first seventeen eighteenth $B$ nonboundary bad-shape counts]\label{lem:bc-b-eighteenth-frontier-badshape-certificate}
For \(15\le q\le31\), the number of
column-even Pieri two-strip paths of length \(2q+18\) with terminal first row
at least \(2q\) is at most \(s_{2q+18}/2\).
\end{lemma}

\begin{proof}
The exact C++/GMP/OpenMP replay
\[
\texttt{certificates/post\_m29/bc\_eighteenth\_frontier\_gmp\_certificate.log}
\]
implements the reverse Pieri recursion of
\cref{lem:bc-b-eighth-frontier-badshape-certificate} for
\(15\le q\le31\), \(j=2q+18\).  It exits
with status \(0\), reports \texttt{SUMMARY failures=0}, and has SHA-256
\[
\hashsplit{2c1768d73f6e84f175648fe543b44e67}{f791c9221e749ff42253c8a69a6ed7a0}.
\]
It verifies the following exact margins:
\[
\begin{array}{c|c|c|c}
q & j & \#\mathrm{bad}(q,j) & s_j-2\#\mathrm{bad}(q,j)\\ \hline
15&48&
72847209805019332586002508503098120976428207913334&
785317984357564384245291307677305539002737266519703333235924\\
16&50&
3897671958418006847914295217951096754733373051015315&
1885339339064052996097496976949503706722936068262166070070627994\\
17&52&
189956451904271257973367176544772686298983735393295632&
4903266810544183341791525110541907903641116661018693378129481867360\\
18&54&
8515222682240817384288986784949072823253006669037635654&
13771977065833947103371970411612873439126421554299867153160479629230836\\
19&56&
354009388401361339571139085869943006414905915046298093568&
41656588897766691496603698711650915942888776629998516920304585259714189184\\
20&58&
13745840672301965992126714036868470336812979278255456156804&
135331264354280596497768713687305919712034604735356312982926749465423376350840\\
21&60&
501537005858674333330084220305394292725502806422776823786938&
471052485162704317318809721958690940566808762031613487380466020931144218662211980\\
22&62&
17286220248116620360011696433092366670786327505855441698376089&
1752662440577781575209469095797630640093064685657073878681125077716565078164357987278\\
23&64&
565412260020227317349094931173175026051291206417829480409450806&
6955857144727139970751970243569083276569269680742626928360224651892431554412713247428244\\
24&66&
17622233793457998802606163139381904487039919418495500636414785527&
29386673299901708964749475331515101602686965572953452647222934299917584004569679370045223186\\
25&68&
525227120153705052445522264916256422275333828462705468427530311688&
131909085289227211701789514678748474089506236322148793119479772545567491021630952165636541404656\\
26&70&
15017983315275271514943772179439551173264045729072445515932158299334&
627984678371728270376818250246782045863473694242216651025377486051528112050569813435834387678750836\\
27&72&
413137929816968194764470941137967440190873604534173610774411341339804&
3165506839235730069039027353629445183330616372133615891590398402802756206977421021928811540220620087496\\
28&74&
10962547901933364527319535364053039700674396423197629467092052579347946&
16868160644365457941953337882591367932501288052108501658663011812498906384019434162381687962487460918312492\\
29&76&
281232323362466735177620928843846456166533731604285601669187439420430570&
94879010830044980231973613148711343553951153789424079612257058436848957410956389901125932638346783683138095660\\
30&78&
6989805140556529916322470879509894906924056264106263867936760916148947281&
562512973617582593017381142149439039365735839082232808233625389602883250292824675325085827160750052651887880339806\\
31&80&
168630546042345400047058712680532671697137285643783300438293211852492794558&
3510497287734778986037296352864815062600027653536998711627383460335903484235467656770090691961806799561483596096398980 .
\end{array}
\]
Each margin is positive, proving the claimed half-stable bound.
\end{proof}

\begin{proposition}[Residual $B$ eighteenth nonboundary lower correction]\label{prop:bc-b-eighteenth-nonboundary-lower-correction}
For every residual row \(B_b\) with \(14\le b\le123\), put \(q=b+1\) and
write \(s_j=m_j^{\mathrm{st}}\) and \(m_j^{B_b}=s_j+\delta_j^{B_b}\).  Then
\[
  -\frac12 s_{2q+18}\le \delta_{2q+18}^{B_b}.
\]
\end{proposition}

\begin{proof}
By \cref{lem:bc-correction-bad-shapes}, the coefficient interpretation in
\cref{lem:bad-shape-witnesses}, and the Pieri-path bijection
\cref{lem:bc-pieri-path-bijection}, \(|\delta_{2q+18}^{B_b}|\) is the number
of column-even Pieri two-strip paths of length \(2q+18\) with terminal first
row at least \(2q\).  If
\(15\le q\le31\), this is bounded by
\(s_{2q+18}/2\) by
\cref{lem:bc-b-h8-h27-bounded-littlewood-certificate}.  If \(q\ge32\),
\cref{lem:bc-b-eighteenth-nonboundary-width-bad-lower-box-code} bounds this
number by
\[
  13210035\binom{2q+18}{36}72^{36}3^{2q-18},
\]
and \cref{lem:bc-stable-dominates-eighteenth-b-nonboundary-width-count}
absorbs that count into \(s_{2q+18}/2\).  In all cases
\(|\delta_{2q+18}^{B_b}|\le s_{2q+18}/2\), which proves the stated lower
bound.
\end{proof}

\begin{lemma}[Nineteenth $B$ nonboundary width-bad paths have a lower-box code]\label{lem:bc-b-nineteenth-nonboundary-width-bad-lower-box-code}
For every \(q\ge15\), the number of column-even Pieri two-strip paths
\[
  \varnothing=\lambda^{(0)}\subset\lambda^{(1)}\subset\cdots
  \subset\lambda^{(2q+19)}
\]
with
\[
  \lambda^{(2q+19)}_1\ge2q
\]
is at most
\[
  122788879\binom{2q+19}{38}76^{38}3^{2q-19}.
\]
\end{lemma}

\begin{proof}
Put \(n=2q+19\).  The terminal partition has size \(2n=4q+38\).  Since its
columns have even lengths and its first row has length at least \(2q\), the
number of boxes below the second row is at most
\[
  4q+38-2(2q)=38 .
\]
Hence the terminal partition has at most forty nonzero rows.

As in \cref{lem:bc-b-sixth-nonboundary-width-bad-lower-box-code}, encode a
path by the row-pair word and track lower boxes.  A step with no lower box has
three choices, a step with one lower box has at most \(2\cdot38=76\) choices,
and a step with two lower boxes has at most \(\binom{38+1}{2}=741\) choices.
Thus the number of paths is bounded by the sum of the coefficients of
\(x^0,\ldots,x^{38}\) in
\[
  (3+76x+741x^2)^n .
\]

The monomial with \(a\) one-lower-box steps and \(c\) two-lower-box steps,
where \(a+2c\le38\), contributes
\[
  \binom{n}{a+c}\binom{a+c}{c}3^{n-a-c}76^a741^c .
\]
Since \(n\ge49\), \(a+c\le38\), and the thirty-eight factors in the quotient
\(\binom n{38}/\binom n{a+c}\) which are not cancelled are all at least
\(12\),
\[
  \frac{\binom{n}{a+c}\binom{a+c}{c}}{\binom n{38}}
  \le
  \frac{38!}{a!\,c!\,12^{38-a-c}}.
\]
Therefore the coefficient sum is at most
\[
  \binom n{38}76^{38}3^{n-38}
  \sum_{\substack{a,c\ge0\\ a+2c\le38}}
  \frac{38!}{a!\,c!\,12^{38-a-c}}
  \frac{3^{38-a-c}741^c}{76^{38-a}}.
\]
The exact replay
\[
\texttt{certificates/post\_m29/bc\_nineteenth\_frontier\_gmp\_certificate.log}
\]
prints this finite rational sum as
\[
  \frac{221206640788254239145767560481888827762009440417078851871455636769429079125950524556363491}
       {1801520162107641037790961357887041840549124446276873034587394818667487083450335232}
  <122788879 .
\]
This gives the displayed bound.
\end{proof}

\begin{lemma}[Stable moments dominate the nineteenth $B$ nonboundary width count above the first row]\label{lem:bc-stable-dominates-nineteenth-b-nonboundary-width-count}
For every \(q\ge34\),
\[
  245577758\binom{2q+19}{38}76^{38}3^{2q-19}\le s_{2q+19}.
\]
\end{lemma}

\begin{proof}
By \cref{lem:bc-littlewood-graph-model}, \(s_{2q+19}\) contains all labelled
simple cycles on \(2q+19\) vertices, whose number is \((2q+18)!/2\).  It is
enough to prove
\[
  491155516\binom{2q+19}{38}76^{38}3^{2q-19}\le (2q+18)! .
\]
At \(q=34\),
\[
\begin{aligned}
  491155516\binom{87}{38}76^{38}3^{49}
  &=
  230288514110446621380410114809443281593262146551241782588645544977111171092717096776659883403545939985052017964291967108408934400,\\
  86!
  &=
  24227095383672732381765523203441259715284870552429381750838764496720162249742450276789464634901319465571660595200000000000000000000 .
\end{aligned}
\]
The first integer is smaller than the second by
\[
  23996806869562285760385113088631816433691608405878139968250118951743051078649733180012804751497773525586608577235708032891591065600 .
\]
If the displayed factorial inequality holds at \(q\), then at \(q+1\) the
left side is multiplied by
\[
  9\frac{(2q+21)(2q+20)}{(2q-17)(2q-18)}<30,
\]
because \(84q^2-2838q+5400>0\) for \(q\ge34\).  The right side is multiplied
by
\[
  (2q+19)(2q+20)\ge87\cdot88>30 .
\]
Induction proves the result.
\end{proof}

\begin{lemma}[Exact first nineteen nineteenth $B$ nonboundary bad-shape counts]\label{lem:bc-b-nineteenth-frontier-badshape-certificate}
For \(15\le q\le33\), the number of
column-even Pieri two-strip paths of length \(2q+19\) with terminal first row
at least \(2q\) is at most \(s_{2q+19}/2\).
\end{lemma}

\begin{proof}
The exact C++/GMP/OpenMP replay
\[
\texttt{certificates/post\_m29/bc\_nineteenth\_frontier\_gmp\_certificate.log}
\]
implements the reverse Pieri recursion of
\cref{lem:bc-b-eighth-frontier-badshape-certificate} for
\(15\le q\le33\), \(j=2q+19\).
It exits with status \(0\), reports \texttt{SUMMARY failures=0}, and has
SHA-256
\[
\hashsplit{519d2680d84e75ce37567e5b89e5b1d7}{3642f0ddf74f02f115aeb728df4168dd}.
\]
It verifies the following exact margins:
\[
\begin{array}{c|c|c|c}
q & j & \#\mathrm{bad}(q,j) & s_j-2\#\mathrm{bad}(q,j)\\ \hline
15&49&
6770832467644019079564113437221206469804889256556352&
38087788771185249644253361404092173181028966248702034081294208\\
16&51&
393122822418047174334190323048832027972488753247225875&
95209341861026025050134845945018848783506874444173368239257149850\\
17&53&
20713609885788249141195428390914016422967880019324583232&
257420799961480829201665256810402075623223711166398183075709448629760\\
18&55&
1000513078990890186889315546888119042022674618893259218256&
750570909947064994029266668996142061495256352229386936747724782641384416\\
19&57&
44683654959956359626290956640694641622888432782519135205912&
2353592104550586121048573793600085590528177447681662704191955919299127332304\\
20&59&
1858704376151413331518423923037635366002147196947105461916934&
7916863333164200252737574462906427863572462353482833717497011882892177835806580\\
21&61&
72467756181056332993253942006277920932854289443127449327382990&
28498624571845034398581420339568459522115756738222291924087293097161491074574739300\\
22&63&
2662742091452354373351765052132387893149373255693957805007608594&
109541225789834175931653170012220614838669799090047603611548170583246648596253202902748\\
23&65&
92649248834862294871171056177379692469942680647355334065212980416&
448652128232881537732129766979294704039333884103811999231018613670984161783957982224082560\\
24&67&
3065580383564886079043046133623641868554505565519696386236636779095&
1954211164681443257926522966327230816937426096125469959736954409871189035415553732519088184402\\
25&69&
96819085800680404234392380763172336679155449656053883968023472031752&
9035761317088612444238991255274381283928621079736448785668421125986211372915586054868393811060208\\
26&71&
2928382315052539397305817789315759122543700932218718904280851982211510&
44272870337243698081214185227202035913447790832225711204683458479263118613924744595775869569671004820\\
27&73&
85074927202479445497105355919110590965014947669202147642568533027500384&
229499010251335973773574243350048059133475729534071769159526994978415191611211825740513706185653381702976\\
28&75&
2380343018946518105665481853782505830247110834040589189899097933395525850&
1256676780468867042847317084437511659161885705709215021301930383983257118198323834841195990539875300449697868\\
29&77&
64296161734029197024019023043412600623223069241190599527445629315215210058&
7258238000562610837751823684753927597802442936203508286377685792014904441015886632454127246403360237470793369452\\
30&79&
1680283079261180654459758454356345266839516313276107630481313587080286524901&
44157232836692700825063754574781471258469541645948190214366928899160825727087749511175294024150102699910285977546806\\
31&81&
42568609323106085115368512864181140363369184496892098629888240297497728680920&
282594820649086173955053884185975495344284216876528810508051194103365072761732733894645251911812624014809789051928811600\\
32&83&
1047339729772837773239753504133535275882447037159770705867392734527969102456219&
1900094241916970117931209240138935340233911256943486387578560442324438888619246289896495761665556156934238477666017200951498\\
33&85&
25066345850201404826353089982459865488949766200157456123373154639121945830921160&
13406572350450220333476929689717057966687138989769754346719747084213925159133552597982651773589939258393987499432673350091612272 .
\end{array}
\]
Each margin is positive, proving the claimed half-stable bound.
\end{proof}

\begin{proposition}[Residual $B$ nineteenth nonboundary lower correction]\label{prop:bc-b-nineteenth-nonboundary-lower-correction}
For every residual row \(B_b\) with \(14\le b\le123\), put \(q=b+1\) and
write \(s_j=m_j^{\mathrm{st}}\) and \(m_j^{B_b}=s_j+\delta_j^{B_b}\).  Then
\[
  -\frac12 s_{2q+19}\le \delta_{2q+19}^{B_b}.
\]
\end{proposition}

\begin{proof}
By \cref{lem:bc-correction-bad-shapes}, the coefficient interpretation in
\cref{lem:bad-shape-witnesses}, and the Pieri-path bijection
\cref{lem:bc-pieri-path-bijection}, \(|\delta_{2q+19}^{B_b}|\) is the number
of column-even Pieri two-strip paths of length \(2q+19\) with terminal first
row at least \(2q\).  If
\(15\le q\le33\), this is
bounded by \(s_{2q+19}/2\) by
\cref{lem:bc-b-h8-h27-bounded-littlewood-certificate}.  If \(q\ge34\),
\cref{lem:bc-b-nineteenth-nonboundary-width-bad-lower-box-code} bounds this
number by
\[
  122788879\binom{2q+19}{38}76^{38}3^{2q-19},
\]
and \cref{lem:bc-stable-dominates-nineteenth-b-nonboundary-width-count}
absorbs that count into \(s_{2q+19}/2\).  In all cases
\(|\delta_{2q+19}^{B_b}|\le s_{2q+19}/2\), which proves the stated lower
bound.
\end{proof}

\begin{lemma}[Twentieth $B$ nonboundary width-bad paths have a lower-box code]\label{lem:bc-b-twentieth-nonboundary-width-bad-lower-box-code}
For every \(q\ge15\), the number of column-even Pieri two-strip paths
\[
  \varnothing=\lambda^{(0)}\subset\lambda^{(1)}\subset\cdots
  \subset\lambda^{(2q+20)}
\]
with
\[
  \lambda^{(2q+20)}_1\ge2q
\]
is at most
\[
  1456821633\binom{2q+20}{40}80^{40}3^{2q-20}.
\]
\end{lemma}

\begin{proof}
Put \(n=2q+20\).  The terminal partition has size \(2n=4q+40\).  Since its
columns have even lengths and its first row has length at least \(2q\), the
number of boxes below the second row is at most
\[
  4q+40-2(2q)=40 .
\]
Hence the terminal partition has at most forty-two nonzero rows.

As in \cref{lem:bc-b-sixth-nonboundary-width-bad-lower-box-code}, encode a
path by the row-pair word and track lower boxes.  A step with no lower box has
three choices, a step with one lower box has at most \(2\cdot40=80\) choices,
and a step with two lower boxes has at most \(\binom{40+1}{2}=820\) choices.
Thus the number of paths is bounded by the sum of the coefficients of
\(x^0,\ldots,x^{40}\) in
\[
  (3+80x+820x^2)^n .
\]

The monomial with \(a\) one-lower-box steps and \(c\) two-lower-box steps,
where \(a+2c\le40\), contributes
\[
  \binom{n}{a+c}\binom{a+c}{c}3^{n-a-c}80^a820^c .
\]
Since \(n\ge50\), \(a+c\le40\), and the forty factors in the quotient
\(\binom n{40}/\binom n{a+c}\) which are not cancelled are all at least
\(11\),
\[
  \frac{\binom{n}{a+c}\binom{a+c}{c}}{\binom n{40}}
  \le
  \frac{40!}{a!\,c!\,11^{40-a-c}}.
\]
Therefore the coefficient sum is at most
\[
  \binom n{40}80^{40}3^{n-40}
  \sum_{\substack{a,c\ge0\\ a+2c\le40}}
  \frac{40!}{a!\,c!\,11^{40-a-c}}
  \frac{3^{40-a-c}820^c}{80^{40-a}}.
\]
The exact replay
\[
\texttt{certificates/post\_m29/bc\_twentieth\_frontier\_gmp\_certificate.log}
\]
prints this finite rational sum as
\[
  \frac{12264937461209716714954521537636349092481663378799225339977493272890654593450478195835053181380571504481551}
       {8418969889721232717238251192735471087049751188699919746359969710080000000000000000000000000000000}
  <1456821633 .
\]
This gives the displayed bound.
\end{proof}

\begin{lemma}[Stable moments dominate the twentieth $B$ nonboundary width count above the first row]\label{lem:bc-stable-dominates-twentieth-b-nonboundary-width-count}
For every \(q\ge35\),
\[
  2913643266\binom{2q+20}{40}80^{40}3^{2q-20}\le s_{2q+20}.
\]
\end{lemma}

\begin{proof}
By \cref{lem:bc-littlewood-graph-model}, \(s_{2q+20}\) contains all labelled
simple cycles on \(2q+20\) vertices, whose number is \((2q+19)!/2\).  It is
enough to prove
\[
  5827286532\binom{2q+20}{40}80^{40}3^{2q-20}\le (2q+19)! .
\]
At \(q=35\),
\[
\begin{aligned}
  5827286532\binom{90}{40}80^{40}3^{50}
  &=
  3329234896895568379204224134529516379707320732634239565033354088031683478673430657166483065405440000000000000000000000000000000000000000,\\
  89!
  &=
  16507955160908461081216919262453619309839666236496541854913520707833171034378509739399912570787600662729080382999756800000000000000000000 .
\end{aligned}
\]
The first integer is smaller than the second by
\[
  13178720264012892702012695127924102930132345503862302289880166619801487555705079082233429505382160662729080382999756800000000000000000000 .
\]
If the displayed factorial inequality holds at \(q\), then at \(q+1\) the
left side is multiplied by
\[
  9\frac{(2q+22)(2q+21)}{(2q-18)(2q-19)}<30,
\]
because \(84q^2-2994q+6102>0\) for \(q\ge35\).  The right side is multiplied
by
\[
  (2q+20)(2q+21)\ge90\cdot91>30 .
\]
Induction proves the result.
\end{proof}

\begin{lemma}[Exact first twenty twentieth $B$ nonboundary bad-shape counts]\label{lem:bc-b-twentieth-frontier-badshape-certificate}
For \(15\le q\le34\), the
number of column-even Pieri two-strip paths of length \(2q+20\) with terminal
first row at least \(2q\) is at most \(s_{2q+20}/2\).
\end{lemma}

\begin{proof}
The exact C++/GMP/OpenMP replay
\[
\texttt{certificates/post\_m29/bc\_twentieth\_frontier\_gmp\_certificate.log}
\]
implements the reverse Pieri recursion of
\cref{lem:bc-b-eighth-frontier-badshape-certificate} for
\(15\le q\le34\),
\(j=2q+20\).  It exits with status \(0\), reports
\texttt{SUMMARY failures=0}, and has SHA-256
\[
\hashsplit{431f9748488a0ab0453777cc498a2260}{996f1ddcf8b393545d30eaed4c291232}.
\]
It verifies the following exact margins:
\[
\begin{array}{c|c|c|c}
q & j & \#\mathrm{bad}(q,j) & s_j-2\#\mathrm{bad}(q,j)\\ \hline
15&50&
626601681181840133819216429897487931989295147764905498&
1885339337818644977650652723006899437363862397793042520642847628\\
16&52&
39420612603703867392627355783715290645400771435211691290&
4903266810465722029487925918272599926426775625100489802729845076044\\
17&54&
2242446146464775064275562338286456935754106993349472655980&
13771977065829479241524405343119090892423418538574005445187118759190184\\
18&56&
116558337698592557180247131336960516811482291550872041365600&
41656588897766459087947078329259234590904274448850906785533313608227645120\\
19&58&
5585210170121879526561125922429194385885745060805002368176460&
135331264354280585354840054788150798574036187950704481885062586411929552311528\\
20&60&
248599509132885821679577291704863685290633807256444573042201325&
471052485162704316822613777704636645874314347062496905384649412031100626225383206\\
21&62&
10345700647858013251402199244799072050890211369889928665516259466&
1752662440577781575188812266942410846830980310560341919312686227631796931716722220524\\
22&64&
404836600016733576711103055510191898670645733628688399470753645445&
6955857144727139970751161701193569849870482172821468890912935463007587013272732559038966\\
23&66&
14969678123642353715447551104423952372380690664212299390659112081830&
29386673299901708964749445427403321904895532283063570078081998512616094416961899324650630580\\
24&68&
525354764503926380135728268608709883616923564834003848150244395360974&
131909085289227211701789513629089399321960886155583300431892518156271029010548666802002811306084\\
25&70&
17566197089705483436441868884632067546570010121532303339011809074961848&
627984678371728270376818250211679687650693277912362800800472230060734619898963351648842633845425808\\
26&72&
561539692065488545749386866580206855352144207401671492062135109927629398&
3165506839235730069039027353628322930222345028978506646799120265026932300309826747291908818823447508308\\
27&74&
17214551458069968350433713444653484281812625806975750382245432515629949524&
16868160644365457941953337882591333525323467716038529845875193233636422160116613057276182405806534817109336\\
28&76&
507484204295461464937454166627315706838195136741042158389386084365845856475&
94879010830044980231973613148711342539545209845226083207703967039905236646899183882250187062912989830287243850\\
29&78&
14422556961003995248296886797113645544851487860035467470064996179205139251148&
562512973617582593017381142149439039336904704770505930796864260950414978992935547717543104748355933815309899732072\\
30&80&
396035212304441030625036839495664726356475522204673041377484731929117410978937&
3510497287734778986037296352864815062599235920373481914256933480774337515848098100000252633445652706684049066260030222\\
31&82&
10528981632129397811017258557478103262553686570594685961906226176015171593146138&
23031461318491612837555179058823773176080157469237699341298896496597194266834837470553712350530985944750395962058801782220\\
32&84&
271525326374832491253485074328530657045777479745138428600205017352812738426359874&
158657763193874483403838398425529131802372689295122009945235536093107998553790399053774178262194660874425019778752931695953788\\
33&86&
6803720890658077113591301726343365291891097789940838630961777707393423990642895000&
1146261223216104476113850870504823904206973429562561380739229608352457616031721661870486480423783207958751382467015901495460849360\\
34&88&
165910482327388070198645559912579287410075267024122936747501078593601271337767995276&
8675754396856524876368104879696901763819780995644414859957047940170808748356386666889663394705181155539134163048350281405869496609000 .
\end{array}
\]
Each margin is positive, proving the claimed half-stable bound.
\end{proof}

\begin{proposition}[Residual $B$ twentieth nonboundary lower correction]\label{prop:bc-b-twentieth-nonboundary-lower-correction}
For every residual row \(B_b\) with \(14\le b\le123\), put \(q=b+1\) and
write \(s_j=m_j^{\mathrm{st}}\) and \(m_j^{B_b}=s_j+\delta_j^{B_b}\).  Then
\[
  -\frac12 s_{2q+20}\le \delta_{2q+20}^{B_b}.
\]
\end{proposition}

\begin{proof}
By \cref{lem:bc-correction-bad-shapes}, the coefficient interpretation in
\cref{lem:bad-shape-witnesses}, and the Pieri-path bijection
\cref{lem:bc-pieri-path-bijection}, \(|\delta_{2q+20}^{B_b}|\) is the number
of column-even Pieri two-strip paths of length \(2q+20\) with terminal first
row at least \(2q\).  If
\(15\le q\le34\), this is
bounded by \(s_{2q+20}/2\) by
\cref{lem:bc-b-h8-h27-bounded-littlewood-certificate}.  If \(q\ge35\),
\cref{lem:bc-b-twentieth-nonboundary-width-bad-lower-box-code} bounds this
number by
\[
  1456821633\binom{2q+20}{40}80^{40}3^{2q-20},
\]
and \cref{lem:bc-stable-dominates-twentieth-b-nonboundary-width-count}
absorbs that count into \(s_{2q+20}/2\).  In all cases
\(|\delta_{2q+20}^{B_b}|\le s_{2q+20}/2\), which proves the stated lower
bound.
\end{proof}

\begin{lemma}[Twenty-first $B$ nonboundary width-bad paths have a lower-box code]\label{lem:bc-b-twentyfirst-nonboundary-width-bad-lower-box-code}
For every \(q\ge15\), the number of column-even Pieri two-strip paths
\[
  \varnothing=\lambda^{(0)}\subset\lambda^{(1)}\subset\cdots
  \subset\lambda^{(2q+21)}
\]
with
\[
  \lambda^{(2q+21)}_1\ge2q
\]
is at most
\[
  22953167364\binom{2q+21}{42}84^{42}3^{2q-21}.
\]
\end{lemma}

\begin{proof}
Put \(n=2q+21\).  The terminal partition has size \(2n=4q+42\).  Since its
columns have even lengths and its first row has length at least \(2q\), the
number of boxes below the second row is at most
\[
  4q+42-2(2q)=42 .
\]
Hence the terminal partition has at most forty-four nonzero rows.

As in \cref{lem:bc-b-sixth-nonboundary-width-bad-lower-box-code}, encode a
path by the row-pair word and track lower boxes.  A step with no lower box has
three choices, a step with one lower box has at most \(2\cdot42=84\) choices,
and a step with two lower boxes has at most \(\binom{42+1}{2}=903\) choices.
Thus the number of paths is bounded by the sum of the coefficients of
\(x^0,\ldots,x^{42}\) in
\[
  (3+84x+903x^2)^n .
\]

The monomial with \(a\) one-lower-box steps and \(c\) two-lower-box steps,
where \(a+2c\le42\), contributes
\[
  \binom{n}{a+c}\binom{a+c}{c}3^{n-a-c}84^a903^c .
\]
Since \(n\ge51\), \(a+c\le42\), and the forty-two factors in the quotient
\(\binom n{42}/\binom n{a+c}\) which are not cancelled are all at least
\(10\),
\[
  \frac{\binom{n}{a+c}\binom{a+c}{c}}{\binom n{42}}
  \le
  \frac{42!}{a!\,c!\,10^{42-a-c}}.
\]
Therefore the coefficient sum is at most
\[
  \binom n{42}84^{42}3^{n-42}
  \sum_{\substack{a,c\ge0\\ a+2c\le42}}
  \frac{42!}{a!\,c!\,10^{42-a-c}}
  \frac{3^{42-a-c}903^c}{84^{42-a}}.
\]
The exact replay
\[
\texttt{certificates/post\_m29/bc\_twentyfirst\_frontier\_gmp\_certificate.log}
\]
prints this finite rational sum as
\[
  \frac{1096457568590064650301360788749759942430223535847866944918020801891074989549823633723676719}
       {47769336197201542165283531601931133238466576384000000000000000000000000000000000}
  <22953167364 .
\]
This gives the displayed bound.
\end{proof}

\begin{lemma}[Stable moments dominate the twenty-first $B$ nonboundary width count above the first row]\label{lem:bc-stable-dominates-twentyfirst-b-nonboundary-width-count}
For every \(q\ge37\),
\[
  45906334728\binom{2q+21}{42}84^{42}3^{2q-21}\le s_{2q+21}.
\]
\end{lemma}

\begin{proof}
By \cref{lem:bc-littlewood-graph-model}, \(s_{2q+21}\) contains all labelled
simple cycles on \(2q+21\) vertices, whose number is \((2q+20)!/2\).  It is
enough to prove
\[
  91812669456\binom{2q+21}{42}84^{42}3^{2q-21}\le (2q+20)! .
\]
At \(q=37\),
\[
\begin{aligned}
  91812669456\binom{95}{42}84^{42}3^{53}
  &=
  2020952329702038580056883671765353532240490512901670259202664882974226156991067846672183897991598578931407451396781826335819778114297216850984960,\\
  94!
  &=
  108736615665674308027365285256786601004186803580182872307497374434045199869417927630229109214583415458560865651202385340530688000000000000000000 .
\end{aligned}
\]
The first integer is smaller than the second by
\[
  106715663335972269447308401585021247471946313067281202048294709551070973712426859783556925316591816879629458199805603514194868221885702783149015040 .
\]
If the displayed factorial inequality holds at \(q\), then at \(q+1\) the
left side is multiplied by
\[
  9\frac{(2q+23)(2q+22)}{(2q-19)(2q-20)}<30,
\]
because \(84q^2-3150q+6846>0\) for \(q\ge37\).  The right side is multiplied
by
\[
  (2q+21)(2q+22)\ge95\cdot96>30 .
\]
Induction proves the result.
\end{proof}

\begin{lemma}[Exact first twenty-two twenty-first $B$ nonboundary bad-shape counts]\label{lem:bc-b-twentyfirst-frontier-badshape-certificate}
For \(15\le q\le36\), the
number of column-even Pieri two-strip paths of length \(2q+21\) with terminal
first row at least \(2q\) is at most \(s_{2q+21}/2\).
\end{lemma}

\begin{proof}
The exact C++/GMP/OpenMP replay
\[
\texttt{certificates/post\_m29/bc\_twentyfirst\_frontier\_gmp\_certificate.log}
\]
implements the reverse Pieri recursion of
\cref{lem:bc-b-eighth-frontier-badshape-certificate} for
\(15\le q\le36\),
\(j=2q+21\).  It exits with status \(0\), reports
\texttt{SUMMARY failures=0}, and has SHA-256
\[
\hashsplit{a5593e2a97e9cadcf10f028940d40366}{500494e1c24910ccb6d39bf66334659c}.
\]
It verifies the following exact margins:
\[
\begin{array}{c|c|c|c}
q & j & \#\mathrm{bad}(q,j) & s_j-2\#\mathrm{bad}(q,j)\\ \hline
15&51&
57811543977979238208912981519511377483207388012807997250&
95209341746189182739012463875861266390581783533703569720135607100\\
16&53&
3935271444298477943782795499450299610411496135409790836638&
257420799953651713532839877421118875481104939978421126564928516122948\\
17&55&
241359790403996572016624122823037347552325345755026539476832&
750570909946584275474616657632482591880704053772366331405452516080867264\\
18&57&
13483281043990464341359276425635930103373985735409424547427888&
2353592104550559243853795732584122124557239457210739201997350665488302888352\\
19&59&
692386586161110504840591158478703586672029913275563393651758692&
7916863333164198871681810892988244845426993242150931105441479725659601456123064\\
20&61&
32940040994191698899954268857577120885406919305555183359844313880&
28498624571845034332846273863547174388193726907080605995139163064937379253540877520\\
21&63&
1461676232296968493641854994839517358525314295674909346240332096712&
109541225789834175928735143031809582560132792630472833670283840738408217819382553926512\\
22&65&
60851792748979218669089820745968269933777857355269431707076117488512&
448652128232881537732008248692294415326585446804432417450535997841634917631211960415066368\\
23&67&
2388995341463768143872970292900698372394645088033353666095569744142086&
1954211164681443257926518194467708656530910508270976425587493357690024099747614313852873458420\\
24&69&
88845432633654924642683731090191307380513016866522699304144400735686930&
9035761317088612444238991077777154188220132602837771366812151038318488539182295214515639283749852\\
25&71&
3142436098248435450700340835367906594534784980700914457176590699368430280&
44272870337243698081214185220923020481581025009619641169526276808438636054387353119231249874898567280\\
26&73&
106084616069733522991706312324435603742287840792452964287794263878255227336&
229499010251335973773574243349836060051190667446979350745589961992112545959522258216233402794962926249072\\
27&75&
3429043748809728020773262400360296509803773781663534433422854159153451438292&
1256676780468867042847317084437504805835074124146209686108093370955249171144982175853507524629752860337872984\\
28&77&
106430978497767131004211612045042701677382712228295965458538173449368693544050&
7258238000562610837751823684753927385069078264137304672002499748016326287496908314344577528381904597363836701468\\
29&79&
3180209606591683066722938312879631206989185811412007356779413941353748937520548&
44157232836692700825063754574781471252112482998923346442229971790310276003643056920977831525852237444376948675555512\\
30&81&
91695611909720121588087263800545921240894401059068291433938283249932971054409680&
282594820649086173955053884185975495344100910789928016477045250312789709281531671830896109113142007224790518105277354080\\
31&83&
2556620978845164037472442365575197853546037465279694388543124435075230258839889258&
1900094241916970117931209240138935340233906145796208156796161043919214746490609749586459276596320482420153796261437726085420\\
32&85&
69063341728972379136179074439017761543446116942228451918639682453152822122876000818&
13406572350450220333476929689717057966687138851693203580475391621508482461062949242068317421533350333361368902405272996001452956\\
33&87&
1810751404550169304474456329896888502972674997950825632252450415579385587625360391760&
99151536294250092657629161470485213770197136711695005130820868468289390882972150353866185554243407324584724810085001079294885522272\\
34&89&
46153312083789859378846255075109804869964874276638518172691114855747721346017012231552&
767803834154796949339700774490437588122613781359175679743046676360270546915333507312483096869169646263181127702593714178182956063737088\\
35&91&
1145317127352153227002215852642636270280448741968539720685823104969892401118932253522220&
6219012485561720169125175711747072746616101573033679459338926359090199247402202707719759233715461730950242190951110480674107051219062384040\\
36&93&
27709125359939135035710749040141509792026747440364037302700734332288341785668113871241340&
52636114508684281023057462983247168448627959118498202940361626572032044893632019966982062829183893176533714744456478219369391755716957274222344 .
\end{array}
\]
Each margin is positive, proving the claimed half-stable bound.
\end{proof}

\begin{proposition}[Residual $B$ twenty-first nonboundary lower correction]\label{prop:bc-b-twentyfirst-nonboundary-lower-correction}
For every residual row \(B_b\) with \(14\le b\le123\), put \(q=b+1\) and
write \(s_j=m_j^{\mathrm{st}}\) and \(m_j^{B_b}=s_j+\delta_j^{B_b}\).  Then
\[
  -\frac12 s_{2q+21}\le \delta_{2q+21}^{B_b}.
\]
\end{proposition}

\begin{proof}
By \cref{lem:bc-correction-bad-shapes}, the coefficient interpretation in
\cref{lem:bad-shape-witnesses}, and the Pieri-path bijection
\cref{lem:bc-pieri-path-bijection}, \(|\delta_{2q+21}^{B_b}|\) is the number
of column-even Pieri two-strip paths of length \(2q+21\) with terminal first
row at least \(2q\).  If
\(15\le q\le36\), this is
bounded by \(s_{2q+21}/2\) by
\cref{lem:bc-b-h8-h27-bounded-littlewood-certificate}.  If \(q\ge37\),
\cref{lem:bc-b-twentyfirst-nonboundary-width-bad-lower-box-code} bounds this
number by
\[
  22953167364\binom{2q+21}{42}84^{42}3^{2q-21},
\]
and \cref{lem:bc-stable-dominates-twentyfirst-b-nonboundary-width-count}
absorbs that count into \(s_{2q+21}/2\).  In all cases
\(|\delta_{2q+21}^{B_b}|\le s_{2q+21}/2\), which proves the stated lower
bound.
\end{proof}

\begin{lemma}[Twenty-second $B$ nonboundary width-bad paths have a lower-box code]\label{lem:bc-b-twentysecond-nonboundary-width-bad-lower-box-code}
For every \(q\ge15\), the number of column-even Pieri two-strip paths
\[
  \varnothing=\lambda^{(0)}\subset\lambda^{(1)}\subset\cdots
  \subset\lambda^{(2q+22)}
\]
with
\[
  \lambda^{(2q+22)}_1\ge2q
\]
is at most
\[
  505942192849\binom{2q+22}{44}88^{44}3^{2q-22}.
\]
\end{lemma}

\begin{proof}
Put \(n=2q+22\).  The terminal partition has size \(2n=4q+44\).  Since its
columns have even lengths and its first row has length at least \(2q\), the
number of boxes below the second row is at most
\[
  4q+44-2(2q)=44 .
\]
Hence the terminal partition has at most forty-six nonzero rows.

As in \cref{lem:bc-b-sixth-nonboundary-width-bad-lower-box-code}, encode a
path by the row-pair word and track lower boxes.  A step with no lower box has
three choices, a step with one lower box has at most \(2\cdot44=88\) choices,
and a step with two lower boxes has at most \(\binom{44+1}{2}=990\) choices.
Thus the number of paths is bounded by the sum of the coefficients of
\(x^0,\ldots,x^{44}\) in
\[
  (3+88x+990x^2)^n .
\]

The monomial with \(a\) one-lower-box steps and \(c\) two-lower-box steps,
where \(a+2c\le44\), contributes
\[
  \binom{n}{a+c}\binom{a+c}{c}3^{n-a-c}88^a990^c .
\]
Since \(n\ge52\), \(a+c\le44\), and the forty-four factors in the quotient
\(\binom n{44}/\binom n{a+c}\) which are not cancelled are all at least
\(9\),
\[
  \frac{\binom{n}{a+c}\binom{a+c}{c}}{\binom n{44}}
  \le
  \frac{44!}{a!\,c!\,9^{44-a-c}}.
\]
Therefore the coefficient sum is at most
\[
  \binom n{44}88^{44}3^{n-44}
  \sum_{\substack{a,c\ge0\\ a+2c\le44}}
  \frac{44!}{a!\,c!\,9^{44-a-c}}
  \frac{3^{44-a-c}990^c}{88^{44-a}}.
\]
The exact replay
\[
\texttt{certificates/post\_m29/bc\_twentysecond\_frontier\_gmp\_certificate.log}
\]
prints this finite rational sum as
\[
  \frac{480362701389146293910605320070060795822066632477782402229407083999759642061754427537952186041}
       {949441869406064352643569097126964670279561525103351211617261691143903975323992064}
  <505942192849 .
\]
This gives the displayed bound.
\end{proof}

\begin{lemma}[Stable moments dominate the twenty-second $B$ nonboundary width count above the first row]\label{lem:bc-stable-dominates-twentysecond-b-nonboundary-width-count}
For every \(q\ge38\),
\[
  1011884385698\binom{2q+22}{44}88^{44}3^{2q-22}\le s_{2q+22}.
\]
\end{lemma}

\begin{proof}
By \cref{lem:bc-littlewood-graph-model}, \(s_{2q+22}\) contains all labelled
simple cycles on \(2q+22\) vertices, whose number is \((2q+21)!/2\).  It is
enough to prove
\[
  2023768771396\binom{2q+22}{44}88^{44}3^{2q-22}\le (2q+21)! .
\]
At \(q=38\),
\[
\begin{aligned}
  2023768771396\binom{98}{44}88^{44}3^{54}
  &=
  65222480362836809952971880455666389153083387061070787188578773411432174644663894576003463013085469685050442643017163501545982315458846527106223293071360,\\
  97!
  &=
  96192759682482119853328425949563698712343813919172976158104477319333745612481875498805879175589072651261284189679678167647067832320000000000000000000000 .
\end{aligned}
\]
The first integer is smaller than the second by
\[
  30970279319645309900356545493897309559260426858102188969525703907901570967817980922802416162503602966210841546662514666101085516861153472893776706928640 .
\]
If the displayed factorial inequality holds at \(q\), then at \(q+1\) the
left side is multiplied by
\[
  9\frac{(2q+24)(2q+23)}{(2q-20)(2q-21)}<30,
\]
because \(84q^2-3306q+7632>0\) for \(q\ge38\).  The right side is multiplied
by
\[
  (2q+22)(2q+23)\ge98\cdot99>30 .
\]
Induction proves the result.
\end{proof}

\begin{lemma}[Exact first twenty-three twenty-second $B$ nonboundary bad-shape counts]\label{lem:bc-b-twentysecond-frontier-badshape-certificate}
For \(15\le q\le37\), the
number of column-even Pieri two-strip paths of length \(2q+22\) with terminal
first row at least \(2q\) is at most \(s_{2q+22}/2\).
\end{lemma}

\begin{proof}
The exact C++/GMP/OpenMP replay
\[
\texttt{certificates/post\_m29/bc\_twentysecond\_frontier\_gmp\_certificate.log}
\]
implements the reverse Pieri recursion of
\cref{lem:bc-b-eighth-frontier-badshape-certificate} for
\(15\le q\le37\),
\(j=2q+22\).  It exits with status \(0\), reports
\texttt{SUMMARY failures=0}, and has SHA-256
\[
\hashsplit{c2db04c00a2fcf719ada1acdb7dadc65}{e468bbaa1c87e3dc9a5b9acd34d653a4}.
\]
It verifies the following exact margins:
\[
\begin{array}{c|c|c|c}
q & j & \#\mathrm{bad}(q,j) & s_j-2\#\mathrm{bad}(q,j)\\ \hline
15&52&
5323446653597160265034771100537264465709989137912479705596&
4903266799897669947501013122988312436919677274971313069775309047432\\
16&54&
391554862972481386596424143010396044394295401319384620857534&
13771977065050854407872372120054793731079199363656922856535048462787076\\
17&56&
25859547932757013356832635997528874181577973416133615996397600&
41656588897714973108756961487659929813171890621521373803284148120317581120\\
18&58&
1550742762553659348667394009433698976116366730276697043739919936&
135331264354277495039735287713212516908269165411141020923092154627846808824576\\
19&60&
85242895433240203918806916229151814537253279743436583143929952152&
471052485162704146834021929490000451619636472168595201459357539670823484449881552\\
20&62&
4329962869966090601536843587787934103162815823735156489170666541261&
1752662440577781566549577928305945670260097533474071857088835002901263810706421656934\\
21&64&
204663057519903431651033174009112472574474979975894933913926071072724&
6955857144727139970342645259353796453721838030914264329561326794523054522243821924184408\\
22&66&
9056239950446667193743574468809630675156045150849828616746594261440965&
29386673299901708964731362886858675855215476029228159664636431183695723184327187454351912310\\
23&68&
377141522624646355140328840518497494512116732940003787624669230532222948&
131909085289227211701788760396753679037103365769358800654323261156652278670981113764030537582136\\
24&70&
14850078186731138423985248674917656424793646087051943648719983451992745202&
627984678371728270376818220546655708367827396814749188734423515566582466038140660886899348009859100\\
25&72&
555149081642615737288848298400293225447491129358822235748371195596699510096&
3165506839235730069039027352519147846321244531492308823731694227842654330006985618779290697849903746912\\
26&74&
19776026674693781692086274041223555144906782028524551930055135052232191346902&
16868160644365457941953337882551815901076996292591058165219635430315172221311177904916836626567101694314580\\
27&76&
673505255032373499948674138000900279194668169367082847349408761371139557634200&
94879010830044980231973613148709996544003553689149115734336298493978260986950723198639805024162416282863688400\\
28&78&
21993526020951550616792268240249754700707318613095492392253147020571791206369995&
562512973617582593017381142149438995378697776789412687708921554045132868668001297246629255182191885030137765494378\\
29&80&
690485903339611044427013997950184870025933019447977651797016217715895573752135316&
3510497287734778986037296352864815061219056184118868707464155558552959105248945012148706675934375240716116153577717464\\
30&82&
20891324507642358784525690594672942219409198560531818203698435293279516475388186937&
23031461318491612837555179058823773176038395878185678882525467149924963338602544180805789903495512886616188959451211700622\\
31&84&
610489260918378676975516456004524172821749473928682746517197893039353161784257326908&
158657763193874483403838398425529131802371468859650825937546567567166138162506069646381280386978483679048975778054840034019720\\
32&86&
17264609322988471549439305743530068322150090948951837637129588102542524243091881345590&
1146261223216104476113850870504823904206973395046950176543602663701028732424271748154086778099989610961498732796754263292983948180\\
33&88&
473362465274461427761811856572603603536037011412988634669670804610877073925734374493952&
8675754396856524876368104879696901763819780994698021750372779860787582326331004618392409522413252132073287555983783336097076283611648\\
34&90&
12604196146271047822033601895687909401492047354315172216664616565100506719981511047858780&
68718405964968875842382854467623549502768495062184845403898826790226678245916587358200182336580633196515379172537144866740531009361978184\\
35&92&
326428585335249770483542621184451051418483064466405760480428587999234699059605687441469656&
569039354428064809584501754853999021023455072762059298881668478459525096027119327247557229194896870590495014745923098434061567967681405578832\\
36&94&
8234325934363867839569348497508591288886234966633995120675510643422203724946419246833102165&
4921474373856215123552185218948554597467557837257766347124203491683905873869384362726367414290228374205912667636064652615876092492060437704512854\\
37&96&
202583023176814053897477663683490322926799724426356872929411083825682875168306496912004395540&
44415035728604135032970506742036555025782948717026354322736716471060309077044850251831931031314322554041263773407435284320835328219725104370585021400 .
\end{array}
\]
Each margin is positive, proving the claimed half-stable bound.
\end{proof}

\begin{proposition}[Residual $B$ twenty-second nonboundary lower correction]\label{prop:bc-b-twentysecond-nonboundary-lower-correction}
For every residual row \(B_b\) with \(14\le b\le123\), put \(q=b+1\) and
write \(s_j=m_j^{\mathrm{st}}\) and \(m_j^{B_b}=s_j+\delta_j^{B_b}\).  Then
\[
  -\frac12 s_{2q+22}\le \delta_{2q+22}^{B_b}.
\]
\end{proposition}

\begin{proof}
By \cref{lem:bc-correction-bad-shapes}, the coefficient interpretation in
\cref{lem:bad-shape-witnesses}, and the Pieri-path bijection
\cref{lem:bc-pieri-path-bijection}, \(|\delta_{2q+22}^{B_b}|\) is the number
of column-even Pieri two-strip paths of length \(2q+22\) with terminal first
row at least \(2q\).  If
\(15\le q\le37\), this is
bounded by \(s_{2q+22}/2\) by
\cref{lem:bc-b-h8-h27-bounded-littlewood-certificate}.  If \(q\ge38\),
\cref{lem:bc-b-twentysecond-nonboundary-width-bad-lower-box-code} bounds this
number by
\[
  505942192849\binom{2q+22}{44}88^{44}3^{2q-22},
\]
and \cref{lem:bc-stable-dominates-twentysecond-b-nonboundary-width-count}
absorbs that count into \(s_{2q+22}/2\).  In all cases
\(|\delta_{2q+22}^{B_b}|\le s_{2q+22}/2\), which proves the stated lower
bound.
\end{proof}

\begin{lemma}[Twenty-third $B$ nonboundary width-bad paths have a lower-box code]\label{lem:bc-b-twentythird-nonboundary-width-bad-lower-box-code}
For every \(q\ge15\), the number of column-even Pieri two-strip paths
\[
  \varnothing=\lambda^{(0)}\subset\lambda^{(1)}\subset\cdots
  \subset\lambda^{(2q+23)}
\]
with
\[
  \lambda^{(2q+23)}_1\ge2q
\]
is at most
\[
  16741606354773\binom{2q+23}{46}92^{46}3^{2q-23}.
\]
\end{lemma}

\begin{proof}
Put \(n=2q+23\).  The terminal partition has size \(2n=4q+46\).  Since its
columns have even lengths and its first row has length at least \(2q\), the
number of boxes below the second row is at most
\[
  4q+46-2(2q)=46 .
\]
Hence the terminal partition has at most forty-eight nonzero rows.

As in \cref{lem:bc-b-sixth-nonboundary-width-bad-lower-box-code}, encode a
path by the row-pair word and track lower boxes.  A step with no lower box has
three choices, a step with one lower box has at most \(2\cdot46=92\) choices,
and a step with two lower boxes has at most \(\binom{46+1}{2}=1081\) choices.
Thus the number of paths is bounded by the sum of the coefficients of
\(x^0,\ldots,x^{46}\) in
\[
  (3+92x+1081x^2)^n .
\]

The monomial with \(a\) one-lower-box steps and \(c\) two-lower-box steps,
where \(a+2c\le46\), contributes
\[
  \binom{n}{a+c}\binom{a+c}{c}3^{n-a-c}92^a1081^c .
\]
Since \(n\ge53\), \(a+c\le46\), and the forty-six factors in the quotient
\(\binom n{46}/\binom n{a+c}\) which are not cancelled are all at least
\(8\),
\[
  \frac{\binom{n}{a+c}\binom{a+c}{c}}{\binom n{46}}
  \le
  \frac{46!}{a!\,c!\,8^{46-a-c}}.
\]
Therefore the coefficient sum is at most
\[
  \binom n{46}92^{46}3^{n-46}
  \sum_{\substack{a,c\ge0\\ a+2c\le46}}
  \frac{46!}{a!\,c!\,8^{46-a-c}}
  \frac{3^{46-a-c}1081^c}{92^{46-a}}.
\]
The exact replay
\[
\texttt{certificates/post\_m29/bc\_twentythird\_frontier\_gmp\_certificate.log}
\]
prints this finite rational sum as
\[
  \frac{5413287532978551318097678891722127088655784535993669973030234234250361370376877380234841755619418027188786488040818356483827177779}
       {323343376869893514353888006024488507396488924654735353434601397268712399548508175475572934010793064806941220715626496}
  <16741606354773 .
\]
This gives the displayed bound.
\end{proof}

\begin{lemma}[Stable moments dominate the twenty-third $B$ nonboundary width count above the first row]\label{lem:bc-stable-dominates-twentythird-b-nonboundary-width-count}
For every \(q\ge40\),
\[
  33483212709546\binom{2q+23}{46}92^{46}3^{2q-23}\le s_{2q+23}.
\]
\end{lemma}

\begin{proof}
By \cref{lem:bc-littlewood-graph-model}, \(s_{2q+23}\) contains all labelled
simple cycles on \(2q+23\) vertices, whose number is \((2q+22)!/2\).  It is
enough to prove
\[
  66966425419092\binom{2q+23}{46}92^{46}3^{2q-23}\le (2q+22)! .
\]
At \(q=40\),
\[
\begin{aligned}
  66966425419092\binom{103}{46}92^{46}3^{57}
  &=
  100799373482071927756100447624964335717029457547995464555007687269102068139512017998152522370498302975581376576947997835914189117769465601974511263680494378680320,\\
  102!
  &=
  961446671503512660926865558697259548455355905059659464369444714048531715130254590603314961882364451384985595980362059157503710042865532928000000000000000000000000 .
\end{aligned}
\]
The first integer is smaller than the second by
\[
  860647298021440733170765111072295212738326447511663999814437026779429646990742572605162439511866148409404219403414061321589520925096067326025488736319505621319680 .
\]
If the displayed factorial inequality holds at \(q\), then at \(q+1\) the
left side is multiplied by
\[
  9\frac{(2q+25)(2q+24)}{(2q-21)(2q-22)}<30,
\]
because \(84q^2-3498q+8100>0\) for \(q\ge40\).  The right side is multiplied
by
\[
  (2q+23)(2q+24)\ge103\cdot104>30 .
\]
Induction proves the result.
\end{proof}

\begin{lemma}[Exact first twenty-five twenty-third $B$ nonboundary bad-shape counts]\label{lem:bc-b-twentythird-frontier-badshape-certificate}
For \(15\le q\le39\), the
number of column-even Pieri two-strip paths of length \(2q+23\) with terminal
first row at least \(2q\) is at most \(s_{2q+23}/2\).
\end{lemma}

\begin{proof}
The exact C++/GMP/OpenMP replay
\[
\texttt{certificates/post\_m29/bc\_twentythird\_frontier\_gmp\_certificate.log}
\]
implements the reverse Pieri recursion of
\cref{lem:bc-b-eighth-frontier-badshape-certificate} for
\(15\le q\le39\),
\(j=2q+23\).  It exits with status \(0\), reports
\texttt{SUMMARY failures=0} in each partition, and has SHA-256
\[
\hashsplit{e095f8a87dbf4718ab5ac066d438ca1b}{f6b32452d617172d90acd926094f717d}.
\]
It verifies the following exact margins:
\[
\begin{array}{c|c|c|c}
q & j & \#\mathrm{bad}(q,j) & s_j-2\#\mathrm{bad}(q,j)\\ \hline
15&53&
489720029079333857613134125714875124856232927111053146129896&
257420798982082198262769118082416215050255289486778264613641805536432\\
16&55&
38870934384506456534228679227956458714632949649039646028420588&
750570909869325126286411737708058481670437211038205082798883277102979752\\
17&57&
2760972531623856623586360234296180553041385386441903402894257432&
2353592104545064265352636000265632122641498367964863179195937677531609229264\\
18&59&
177529317063313594727013993839354605142600809019023903692310588212&
7916863333163845197820856588019800498621631490347819247883268228979004138464024\\
19&61&
10434766599703968803394804221878935599357976145166175185279618980446&
28498624571845013529193156443992965398493820864363649050000711343697375413991544388\\
20&63&
565345216504047943127435592612320105186245465113366840852839369823961&
109541225789834174800968062488307633292545317395511658014843539103024354806184478472014\\
21&65&
28436452122494546784503973932702955134312420993458055393779840223891968&
448652128232881537675257048032803280194915678580518443721778711569429345707066432202259456\\
22&67&
1336195933930798511146548004105498833028709271000610258638238376308898757&
1954211164681443257923850580590529987044905158203351229318180728438198945937670028239743945078\\
23&69&
58974092244345858142200133092969801790833911541634513030662754818190165252&
9035761317088612444238873307283530763813697487939047607591184131521439003201632497294804374793208\\
24&71&
2456569840020508957073775728150479642048185471841740547517048174713983263458&
44272870337243698081214180314068212637059981763468866539302805901137262332308086999488081845668900924\\
25&73&
96987495191309572433445415443111879196267468508045127977093262561694838645616&
229499010251335973773574243156073238900710989626071932484015074807062184625016908190622466199329759412512\\
26&75&
3643014674651632199039247569020925167825601643487772678803453906243107112389334&
1256676780468867042847317084430225634573268479204172737494852241212617575405569957564767462525584953015970900\\
27&77&
130626558654627780336785493356319144285233277019286563030236447231223641928726792&
7258238000562610837751823684753666344813726004110893109439011199813159175708294197810447972563789048817366335984\\
28&79&
4484824408626513520837247625031122901074027004243550353290153127575853286119227960&
44157232836692700825063754574781462288824084959079670901180598353823736269567420056701139659104810175377874312140688\\
29&81&
147839993704151891097673114935981535868645820113678338225334706966064166592768329640&
282594820649086173955053884185975495048604314605443674524875194968427729386721820405657569245340469859162050861849514160\\
30&83&
4690810571265634452525230604643372066452349462118325152714379364500745954517147932817&
1900094241916970117931209240138935340224529637895634578219185527594659152753412142737153185679668809940022454812921109998302\\
31&85&
143579375783345675129915376399568094232061561510997014964985636510055593581476759259076&
13406572350450220333476929689717057966686851831068320347069890063113832202910008205837529311960324199367713697523754288234936440\\
32&87&
4248316050658717351663544273621748138012165488984367502227485861088874554402804651204618&
99151536294250092657629161470485213770197128218684406622486503750149756299269651335480557581410054134117903463494663448936303896556\\
33&89&
121741397701642803189252770380994743350843307316009993882328516506636481279153495650343536&
767803834154796949339700774490437588122613781115785190963928649700522698663494240220536411985702935534869476919032247062567998787513120\\
34&91&
3384581136537593661899995441061156486986342192260755133345315750613170930637552174185364260&
6219012485561720169125175711747072746616101573026912587700105876072853260951785680018327110228424157762992931096094078597634184735198699960\\
35&93&
91433429026135885153051601448881562353933632458095125897796643809931783337877280108972533346&
52636114508684281023057462983247168448627959118498020128921825020140008861850620284140374545972471867009993756570327020379399572492967071638332\\
36&95&
2403659572979022768206316866500445519920780823246805116765721775040200192632674950996587156935&
465079114888095472747378954931122536275227942736325900402962100465988261414976812349575413009986465099649036282868029083392176506913941910130566258\\
37&97&
61573662580028538121138960264906234970216539335277488869959450210655901086720201850042978311904&
4286049101830469787585519062310491636993132895599321887180676975724911052612077176435617242577630148038138756746834614683804694344830243614261847507520\\
38&99&
1538912441716013944498604310705441375163609522812461895354428981430433150377490598407851632893185&
41162110219623827986714077844387249387887643477052640993124155319052740791524607978035634866636003958801759691664111820942413257561151865783948334858360318\\
39&101&
37569326911962077279113652639608882499403096094587221496093838457684817231354420086304461819573884&
411610493361212659403148807521076157229850906012077840841555796551366675987937672908423250486030980537421626324382116463164316234745031427737901636673153889032 .
\end{array}
\]
Each margin is positive, proving the claimed half-stable bound.
\end{proof}

\begin{proposition}[Residual $B$ twenty-third nonboundary lower correction]\label{prop:bc-b-twentythird-nonboundary-lower-correction}
For every residual row \(B_b\) with \(14\le b\le123\), put \(q=b+1\) and
write \(s_j=m_j^{\mathrm{st}}\) and \(m_j^{B_b}=s_j+\delta_j^{B_b}\).  Then
\[
  -\frac12 s_{2q+23}\le \delta_{2q+23}^{B_b}.
\]
\end{proposition}

\begin{proof}
By \cref{lem:bc-correction-bad-shapes}, the coefficient interpretation in
\cref{lem:bad-shape-witnesses}, and the Pieri-path bijection
\cref{lem:bc-pieri-path-bijection}, \(|\delta_{2q+23}^{B_b}|\) is the number
of column-even Pieri two-strip paths of length \(2q+23\) with terminal first
row at least \(2q\).  If
\(15\le q\le39\), this is
bounded by \(s_{2q+23}/2\) by
\cref{lem:bc-b-h8-h27-bounded-littlewood-certificate}.  If \(q\ge40\),
\cref{lem:bc-b-twentythird-nonboundary-width-bad-lower-box-code} bounds this
number by
\[
  16741606354773\binom{2q+23}{46}92^{46}3^{2q-23},
\]
and \cref{lem:bc-stable-dominates-twentythird-b-nonboundary-width-count}
absorbs that count into \(s_{2q+23}/2\).  In all cases
\(|\delta_{2q+23}^{B_b}|\le s_{2q+23}/2\), which proves the stated lower
bound.
\end{proof}

\begin{lemma}[Twenty-fourth $B$ nonboundary width-bad paths have a lower-box code]\label{lem:bc-b-twentyfourth-nonboundary-width-bad-lower-box-code}
For every \(q\ge15\), the number of column-even Pieri two-strip paths
\[
  \varnothing=\lambda^{(0)}\subset\lambda^{(1)}\subset\cdots
  \subset\lambda^{(2q+24)}
\]
with
\[
  \lambda^{(2q+24)}_1\ge2q
\]
is at most
\[
  917878401790345\binom{2q+24}{48}96^{48}3^{2q-24}.
\]
\end{lemma}

\begin{proof}
Put \(n=2q+24\).  The terminal partition has size \(2n=4q+48\).  Since its
columns have even lengths and its first row has length at least \(2q\), the
number of boxes below the second row is at most
\[
  4q+48-2(2q)=48 .
\]
Hence the terminal partition has at most fifty nonzero rows.

As in \cref{lem:bc-b-sixth-nonboundary-width-bad-lower-box-code}, encode a
path by the row-pair word and track lower boxes.  A step with no lower box has
three choices, a step with one lower box has at most \(2\cdot48=96\) choices,
and a step with two lower boxes has at most \(\binom{48+1}{2}=1176\) choices.
Thus the number of paths is bounded by the sum of the coefficients of
\(x^0,\ldots,x^{48}\) in
\[
  (3+96x+1176x^2)^n .
\]

The monomial with \(a\) one-lower-box steps and \(c\) two-lower-box steps,
where \(a+2c\le48\), contributes
\[
  \binom{n}{a+c}\binom{a+c}{c}3^{n-a-c}96^a1176^c .
\]
Since \(n\ge54\), \(a+c\le48\), and the forty-eight factors in the quotient
\(\binom n{48}/\binom n{a+c}\) which are not cancelled are all at least
\(7\),
\[
  \frac{\binom{n}{a+c}\binom{a+c}{c}}{\binom n{48}}
  \le
  \frac{48!}{a!\,c!\,7^{48-a-c}}.
\]
Therefore the coefficient sum is at most
\[
  \binom n{48}96^{48}3^{n-48}
  \sum_{\substack{a,c\ge0\\ a+2c\le48}}
  \frac{48!}{a!\,c!\,7^{48-a-c}}
  \frac{3^{48-a-c}1176^c}{96^{48-a}}.
\]
The exact replay
\[
\texttt{certificates/post\_m29/bc\_twentyfourth\_frontier\_gmp\_certificate.log}
\]
prints this finite rational sum as
\[
  \frac{7189885764136456983655939976530114679706233205135043948665828743454022081720455490347644395574905142522888623}
       {7833157148171707308011570939888586941509981933966218053272652442206028867904865324503574511616}
  <917878401790345 .
\]
This gives the displayed bound.
\end{proof}

\begin{lemma}[Stable moments dominate the twenty-fourth $B$ nonboundary width count above the first row]\label{lem:bc-stable-dominates-twentyfourth-b-nonboundary-width-count}
For every \(q\ge42\),
\[
  1835756803580690\binom{2q+24}{48}96^{48}3^{2q-24}\le s_{2q+24}.
\]
\end{lemma}

\begin{proof}
By \cref{lem:bc-littlewood-graph-model}, \(s_{2q+24}\) contains all labelled
simple cycles on \(2q+24\) vertices, whose number is \((2q+23)!/2\).  It is
enough to prove
\[
  3671513607161380\binom{2q+24}{48}96^{48}3^{2q-24}\le (2q+23)! .
\]
At \(q=42\), the left side is
\[
  281290092234071577973503421044138432694077953662272593063201152617396786044386533989846255433558737317238612037947479045428174071804158800497377420219097550803872147046400,
\]
whereas the right side is
\[
  12265202031961379393517517010387338887131568154382945052653251412013535324922144249034658613287059061933743916719318560380966506520420000368175349760000000000000000000000000.
\]
The second integer exceeds the first by
\[
  11983911939727307815544013589343200454437490200720672459590050259396138538877757715044812357853500324616505304681371081335538332448615841567677972339780902449196127852953600.
\]
If the displayed factorial inequality holds at \(q\), then at \(q+1\) the
left side is multiplied by
\[
  9\frac{(2q+26)(2q+25)}{(2q-22)(2q-23)}<30,
\]
because \(84q^2-3618q+9330>0\) for \(q\ge42\).  The right side is multiplied
by
\[
  (2q+24)(2q+25)\ge108\cdot109>30 .
\]
Induction proves the result.
\end{proof}

\begin{lemma}[Exact first twenty-seven twenty-fourth $B$ nonboundary bad-shape counts]\label{lem:bc-b-twentyfourth-frontier-badshape-certificate}
For \(15\le q\le41\), the
number of column-even Pieri two-strip paths of length \(2q+24\) with terminal
first row at least \(2q\) is at most \(s_{2q+24}/2\).
\end{lemma}

\begin{proof}
The exact C++/GMP/OpenMP replay
\[
\texttt{certificates/post\_m29/bc\_twentyfourth\_frontier\_gmp\_certificate.log}
\]
implements the reverse Pieri recursion of
\cref{lem:bc-b-eighth-frontier-badshape-certificate} for
\(15\le q\le41\),
\(j=2q+24\).  It exits with status \(0\), reports
\texttt{SUMMARY failures=0} in each partition, and has SHA-256
\[
\hashsplit{ccd56620d91d6a3c2c9e23673211e25d}{cb5e1ae5d1a4ddc71e699a7bf80fab4a}.
\]
It verifies the following exact margins:
\[
\begin{array}{c|c|c|c}
q & j & \#\mathrm{bad}(q,j) & s_j-2\#\mathrm{bad}(q,j)\\ \hline
15&54&
45045263939682490798625888586980598343275540820709885621748726&
13771976975743436254452353295995864843138794765894432017754046461004692\\
16&56&
3853595618586362180667774183531853844560042311864307437609532120&
41656588890059500967449751153038046718103240680764445126387800477091312080\\
17&58&
294039258448642901148097883499242498063650162962918317768929232232&
135331264353692518008363109229613655929289547812965953330626871386396430199984\\
18&60&
20249959191843444305833437598392302381278777523904378211166683503100&
471052485162663817401429109081796622358272145867461718410869217787567438942779656\\
19&62&
1271390358886551502663758823963550766753045144774576745906496942946037&
1752662440577779032428785895135121545816136781948406557324177101218084976053868847382\\
20&64&
73398462046113262012901298312547242425376312199951172643980925945125739&
6955857144727139823955047282167079291221307753838004423957652354572499102109822176078378\\
21&66&
3925088466731500848265635962988754101732680429696664513864789639730385500&
29386673299901708956899298433296567493071691252188270722521382414604093813831101363414023240\\
22&68&
195679580685942783913621964210302707886038922705442508986856265743410077222&
131909085289227211701398155518427042761986402498619232233540207544707273660582650571004781873588\\
23&70&
9145458684858075239962878536000760042303659534014128431080163108797587710904&
627984678371728270376799959329442365679624319028174537049651757834806611885165798000648656819927696\\
24&72&
402688018352700462807220414025748876981815909448209153041190692851610246248562&
3165506839235730069039026548253409304205550391629177369034526715105817692306324007894647385822810269980\\
25&74&
16777584876016509681520581246827057147288538396438702152281311848682584186066200&
16868160644365457941953337849036198202393364492934068219647968246027908992490822704214323199306397704875984\\
26&76&
664009147497707641509103242761462672991328960454795898549530340200634149939763300&
94879010830044980231973613147383325259518203153130806597089374948553992402379865567235443161283890262099430200\\
27&78&
25050198929822899527976135341552742144793230642789794608233058749368692925605228167&
562512973617582593017381142149388938967890172893457969022774929060352683621352943848397573570987188787868967778034\\
28&80&
903642111799618870549970764934666658338114161601112856370226200733039852660845378845&
3510497287734778986037296352864813255315804391560349696376653685120012169072487848818949239076006210068201979391230406\\
29&82&
31257421937463111361362659840762127231565148079150260679881109803052912678346680439266&
23031461318491612837555179058823773113565334652274740877369198849590054759911066419626332180140690151096922635708627195964\\
30&84&
1039415550823820707263415883628582850385866940261125283184326131683062536884205967833172&
158657763193874483403838398425529131800293858736525021280373686832820981510079979264806887186102865811468929410609996613007192\\
31&86&
33305171647009003918559435158958772551890489064287365979814292201828764290165127165152930&
1146261223216104476113850870504823904206906819232874804512708643709322301939304611476140101271705256636271280353222419222416333500\\
32&88&
1030478358849312857260381531499271084474857808305628541123443156025862460289958735586143358&
8675754396856524876368104879696901763819778934688028982295988195548143040934042875748866933982147154526317113480616904031073860312836\\
33&90&
30846845539926195880506199003540429333265201502561122989757264051124621076396995822469779794&
68718405964968875842382854467623549502768495000516362716338977124858347442627104510472763426167017561434180303418103727386502386518136156\\
34&92&
894947929676122083872901537759812156760573140029404763312322181492051760951134902103369589150&
569039354428064809584501754853999021023455072760270055879486904791320260036842071836138919880966993875391331240114993381557417374849549339844\\
35&94&
25206262614608521048539077211064939054294161813570017964872057998825172838285337199978546453799&
4921474373856215123552185218948554597467557837257715951067626143369544474853659227865441403739071167437973164871088289114606971710498974277809586\\
36&96&
690233574462789692594850732767703353446501514145841953617365179647040640257618499139020032683275&
44415035728604135032970506742036555025782948717026352942674733591834551995138340043792204784164893711210070284535898157890920563319339820154528445930\\
37&98&
18401969786379726509397550928547608668849022060215237021158376690029263548754575191596046896961224&
417889612984311197972881691783935928096937375597305445986352710733921306024321448990581987404762569297015318862011729845392458401771406744610731284750960\\
38&100&
478269003844562569604369062258328562466998476454805786786548381394337263460214408898997275248369306&
4095628359235226170344431641109679151351419816735792386999621199938760398734934237313071615798391265669544725547987428986378191175226993278828223708741378764\\
39&102&
12132175928801974779624112816786884758210596625799885849133358631452145248544991708615133324294344258&
41778449637156136111589911062266613865092602309151061119068884600495926616345005830809526260844732371588482361350868475689895706862246077865474658190283617108860\\
40&104&
300708818660876522548302103782613253132638102651326194486615013540965687730620281067692837157016372928&
443216814209338571157347357905009055945908874869036427451842429019139480523150218138988446315765810468777082647789394591092589885307010867932145203289945850205737600\\
41&106&
7290277856693512788026993904674025738843409377869210602852462788859436376147123184482197231103638539560&
4886351313102937424311841336468296380298521371436515901130417647337153440439863468297240623754779583629242523324475917211904198485224788159926080471121817552593787389360 .
\end{array}
\]
Each margin is positive, proving the claimed half-stable bound.
\end{proof}

\begin{proposition}[Residual $B$ twenty-fourth nonboundary lower correction]\label{prop:bc-b-twentyfourth-nonboundary-lower-correction}
For every residual row \(B_b\) with \(14\le b\le123\), put \(q=b+1\) and
write \(s_j=m_j^{\mathrm{st}}\) and \(m_j^{B_b}=s_j+\delta_j^{B_b}\).  Then
\[
  -\frac12 s_{2q+24}\le \delta_{2q+24}^{B_b}.
\]
\end{proposition}

\begin{proof}
By \cref{lem:bc-correction-bad-shapes}, the coefficient interpretation in
\cref{lem:bad-shape-witnesses}, and the Pieri-path bijection
\cref{lem:bc-pieri-path-bijection}, \(|\delta_{2q+24}^{B_b}|\) is the number
of column-even Pieri two-strip paths of length \(2q+24\) with terminal first
row at least \(2q\).  If
\(15\le q\le41\), this is
bounded by \(s_{2q+24}/2\) by
\cref{lem:bc-b-h8-h27-bounded-littlewood-certificate}.  If \(q\ge42\),
\cref{lem:bc-b-twentyfourth-nonboundary-width-bad-lower-box-code} bounds this
number by
\[
  917878401790345\binom{2q+24}{48}96^{48}3^{2q-24},
\]
and \cref{lem:bc-stable-dominates-twentyfourth-b-nonboundary-width-count}
absorbs that count into \(s_{2q+24}/2\).  In all cases
\(|\delta_{2q+24}^{B_b}|\le s_{2q+24}/2\), which proves the stated lower
bound.
\end{proof}


\begin{lemma}[Twenty-fifth $B$ nonboundary width-bad paths have a lower-box code]\label{lem:bc-b-twentyfifth-nonboundary-width-bad-lower-box-code}
For every \(q\ge15\), the number of column-even Pieri two-strip paths
\[
  \varnothing=\lambda^{(0)}\subset\lambda^{(1)}\subset\cdots
  \subset\lambda^{(2q+25)}
\]
with
\[
  \lambda^{(2q+25)}_1\ge2q
\]
is at most
\[
  96319792639339148\binom{2q+25}{50}100^{50}3^{2q-25}.
\]
\end{lemma}

\begin{proof}
Put \(n=2q+25\).  The terminal partition has size \(2n=4q+50\).  Since its
columns have even lengths and its first row has length at least \(2q\), the
number of boxes below the second row is at most
\[
  4q+50-2(2q)=50 .
\]
Hence the terminal partition has at most fifty-two nonzero rows.

As in \cref{lem:bc-b-sixth-nonboundary-width-bad-lower-box-code}, encode a
path by the row-pair word and track lower boxes.  A step with no lower box has
three choices, a step with one lower box has at most \(2\cdot50=100\) choices,
and a step with two lower boxes has at most \(\binom{50+1}{2}=1275\) choices.
Thus the number of paths is bounded by the sum of the coefficients of
\(x^0,\ldots,x^{50}\) in
\[
  (3+100x+1275x^2)^n .
\]

The monomial with \(a\) one-lower-box steps and \(c\) two-lower-box steps,
where \(a+2c\le50\), contributes
\[
  \binom{n}{a+c}\binom{a+c}{c}3^{n-a-c}100^a1275^c .
\]
Since \(n\ge55\), \(a+c\le50\), and the fifty factors in the quotient
\(\binom n{50}/\binom n{a+c}\) which are not cancelled are all at least
\(6\),
\[
  \frac{\binom{n}{a+c}\binom{a+c}{c}}{\binom n{50}}
  \le
  \frac{50!}{a!\,c!\,6^{50-a-c}}.
\]
Therefore the coefficient sum is at most
\[
  \binom n{50}100^{50}3^{n-50}
  \sum_{\substack{a,c\ge0\\ a+2c\le50}}
  \frac{50!}{a!\,c!\,6^{50-a-c}}
  \frac{3^{50-a-c}1275^c}{100^{50-a}}.
\]
The exact replay
\[
\texttt{certificates/post\_m29/bc\_twentyfifth\_frontier\_gmp\_certificate.log}
\]
prints this finite rational sum as
\[
  \frac{31562069652058651724050564331116427866898930990810385836335683810816344863875220220375406369645430784395169909}
       {327680000000000000000000000000000000000000000000000000000000000000000000000000000000000000000}
  <96319792639339148 .
\]
This gives the displayed bound.
\end{proof}

\begin{lemma}[Stable moments dominate the twenty-fifth $B$ nonboundary width count above the first row]\label{lem:bc-stable-dominates-twentyfifth-b-nonboundary-width-count}
For every \(q\ge44\),
\[
  192639585278678296\binom{2q+25}{50}100^{50}3^{2q-25}\le s_{2q+25}.
\]
\end{lemma}

\begin{proof}
By \cref{lem:bc-littlewood-graph-model}, \(s_{2q+25}\) contains all labelled
simple cycles on \(2q+25\) vertices, whose number is \((2q+24)!/2\).  It is
enough to prove
\[
  385279170557356592\binom{2q+25}{50}100^{50}3^{2q-25}\le (2q+24)! .
\]
At \(q=44\), the left side is
\[
  1631698339603928414165194585139136415150069360512059164383537511536124994126576640000000000000000000000000000000000000000000000000000000000000000000000000000000000000000000000000000,
\]
whereas the right side is
\[
  197450685722107402353682037275992488341277868034975337796656295094902858969771811440894224355027779366597957338237853638272334919686385621811850780464277094400000000000000000000000000.
\]
The second integer exceeds the first by
\[
  195818987382503473939516842690853351926127798674463278632272757583366733975645234800894224355027779366597957338237853638272334919686385621811850780464277094400000000000000000000000000 .
\]
If the displayed factorial inequality holds at \(q\), then at \(q+1\) the
left side is multiplied by
\[
  9\frac{(2q+27)(2q+26)}{(2q-23)(2q-24)}<30,
\]
because \(84q^2-3774q+10242>0\) for \(q\ge44\).  The right side is multiplied
by
\[
  (2q+25)(2q+26)\ge113\cdot114>30 .
\]
Induction proves the result.
\end{proof}

\begin{lemma}[Exact first twenty-nine twenty-fifth $B$ nonboundary bad-shape counts]\label{lem:bc-b-twentyfifth-frontier-badshape-certificate}
For \(15\le q\le43\), the
number of column-even Pieri two-strip paths of length \(2q+25\) with terminal
first row at least \(2q\) is at most \(s_{2q+25}/2\).
\end{lemma}

\begin{proof}
The exact C++/GMP/OpenMP replay
\[
\texttt{certificates/post\_m29/bc\_twentyfifth\_frontier\_gmp\_certificate.log}
\]
implements the reverse Pieri recursion of
\cref{lem:bc-b-eighth-frontier-badshape-certificate} for
\(15\le q\le43\),
\(j=2q+25\).  It exits with status \(0\), reports
\texttt{SUMMARY failures=0} in each B partition, and has SHA-256
\[
\hashsplit{1bd9d79c07c867c4d2d40c446cc351f3}{a2fbd100861b423ff9e32c0bfc41ed91}.
\]
It verifies the following exact margins:
\[
\begin{array}{c|c|c|c}
q & j & \#\mathrm{bad}(q,j) & s_j-2\#\mathrm{bad}(q,j)\\ \hline
15&55&
4145931517461905557858975498460147368518189723816827644791522010&
750570901655203960131613535058564843206055391431091534463307279576776908\\
16&57&
381825147256219851645212554161959212798369909875907096220856947944&
2353592103786935915903444010222379734786172303474206130217007291895683848240\\
17&59&
31262241721162037041241938421324025149111199451620716290903589710968&
7916863333101675773012659141126770649766662149259882050598064844204581580218512\\
18&61&
2303513684154537786512596537819715085137130432054160558264571499792840&
28498624571840427371358046776357546995026625192064573505088893354931216830229919600\\
19&63&
154331061421702794859302866377271197814492106835326405573603958567232701&
109541225789833867269535652090813800941683748077756239403120798676946889303946083654534\\
20&65&
9484773290242459222057202928554516041650527645222023360762614503785816096&
448652128232881518762583371792873929649517769336892270689348263112298734969397105078411200\\
21&67&
538765609706049811735637740933424335650982597262166834434347342128085157340&
1954211164681443256848991753046291960595922772344713555682272952455866497586251820736191427912\\
22&69&
28472554463555590153447734908385862845491061556375669678159717963940642743534&
9035761317088612444182046146540908275223086418388461821503783676231770932871374386876559469636644\\
23&71&
1407999823983756896957975237043370441410192679057679962625726147821044728467240&
44272870337243698081211369227559925164283979960546236099379269612150090655863930581288789184178493360\\
24&73&
65482513596139903435397956148551264351869178182153461043978953484141568164518232&
229499010251335973773574112385021037003523263697050466267710129461240757334185076186902023039583107667280\\
25&75&
2877000945746070073513592998559095898913681456515908712069980686077444195516288104&
1256676780468867042847317078683509772430431596575065235514702293720905865660728078782412998183182776208173360\\
26&77&
119888756626776632089808575502343737859968173294893746848818950024528282238973333612&
7258238000562610837751823684514150084677481995491949529421036362381793295672545277238870545409194931623277122344\\
27&79&
4755433453302025540806807094223694719079936824229588911800371943793040789295397920941&
44157232836692700825063754574771960391566298161025098961487401026631378543972969365978245495523479245505855754754726\\
28&81&
180122261004411799850782300882024397519889210366410200764974162892136015474289706437920&
282594820649086173955053884185975135099762293190147755155619660791595761345592727812612715747684099519259435467973297600\\
29&83&
6533814889183911337086297668211913451588726206575274598861572338785122326920449835991954&
1900094241916970117931209240138935327166281480670343172951640652380118993708864428510840638261952891098779292881055733880028\\
30&85&
227576088845539798983484324827616614867063917040530292163791999671686785188166588936377358&
13406572350450220333476929689717057966231986812128807440453181244211398109364344494879490721662670171297360238334584063880699876\\
31&87&
7629243383042814630196006231888299575515378431998099393346712401121861451922225008034154330&
99151536294250092657629161470485213770181878228550422310660815064774526943614896602947539351358365164287381918339927804529537997132\\
32&89&
246703306024066097803373166592171570021148679547491900777820203439034237196358256869385924964&
767803834154796949339700774490437588122613287952655938235018649332695055082343684624863931633921367659119631863830816904361251316350264\\
33&91&
7710283187893013835784344358602922029642325400191151143795062012174421286757753724395062592270&
6219012485561720169125175711747072746616101557613115374187153391827964534628063933707648994230643380439559538248477846943401840293444243940\\
34&93&
233324004198302830434854312274352246658842166040188021897341170251820615304989830564682131613676&
52636114508684281023057462983247168448627959118031554987383271631040606340504813553947398081157012013467107009354305653336095665923820753477672\\
35&95&
6848108068247151965795627395858831467260786613484327219811180368085085014244250337430578107544895&
465079114888095472747378954931122536275227942736312208994144752120102206572818827687531931278321142938819647453575408993764073271588982747089790338\\
36&97&
195243007857585300168028972108265141626487986830143959392983018240227775048035735368323967882264544&
4286049101830469787585519062310491636993132895599321496817808585714367792796410880435146459542089566420774948520717034649566400446799910666412039602240\\
37&99&
5414954288680224882813114357431128639023531061721428660661337850502110298527677652548871833485814200&
41162110219623827986714077844387249387887643477052640982297324566575723053787376471794788471339269054404527294132144977901053526806551541883020371152518288\\
38&101&
146286237117273352894547699227239310976202885033796329383086113103977967509724924541573727667937924550&
411610493361212659403148807521076157229850906012077840841263299215785953206303135737274051081843573573546222840058936424633275669360044287494926790260917187700\\
39&103&
3854208145235496907201420896374848214274728530488552179052146178754920439872389117617201368228523392174&
4282289551634348924637847939334988424968279064470634107325419718782403448969788932329158846270183968871956233820510706247424515595115005373062980142668165291697828\\
40&105&
99148506255084971187483653672387066123475937428263609916771914087900986624591806679913530709013275727664&
46316141414884575423512785844162280630709788031116921318799723858789467181092314057828795308755402478761745913159139545969528577914828290659084552753349830883202224544\\
41&107&
2492980620950423184030911529930656255902260845725546617956542306845227660381862220931430772388347754440216&
520396248607814162823456898918379242662159520014922891460847538243510608836810618504477719528362625247830625823194172205046442989452872079233110065580070694962664054412240\\
42&109&
61328864800734725452304195789836622444459461915696851008253808474360134995239026619488849442774457977304792&
6069768027148561641469987059305686919203463935482378676309774609902944012654714526503114749223079939184754909539962195692260581280715316990819031184935102761658602563986149968\\
43&111&
1477500817551458801892338405894396207044666417511868155552612189932909325352000260541969253705692496658370785&
73442633136598658548124378676446602820000468916554313883241994137346871713887965773926070832575361847114241104749542170651189054330469960695541408770995495301695217983881475890750 .
\end{array}
\]
Each margin is positive, proving the claimed half-stable bound.
\end{proof}

\begin{proposition}[Residual $B$ twenty-fifth nonboundary lower correction]\label{prop:bc-b-twentyfifth-nonboundary-lower-correction}
For every residual row \(B_b\) with \(14\le b\le123\), put \(q=b+1\) and
write \(s_j=m_j^{\mathrm{st}}\) and \(m_j^{B_b}=s_j+\delta_j^{B_b}\).  Then
\[
  -\frac12 s_{2q+25}\le \delta_{2q+25}^{B_b}.
\]
\end{proposition}

\begin{proof}
By \cref{lem:bc-correction-bad-shapes}, the coefficient interpretation in
\cref{lem:bad-shape-witnesses}, and the Pieri-path bijection
\cref{lem:bc-pieri-path-bijection}, \(|\delta_{2q+25}^{B_b}|\) is the number
of column-even Pieri two-strip paths of length \(2q+25\) with terminal first
row at least \(2q\).  If
\(15\le q\le43\), this is
bounded by \(s_{2q+25}/2\) by
\cref{lem:bc-b-h8-h27-bounded-littlewood-certificate}.  If \(q\ge44\),
\cref{lem:bc-b-twentyfifth-nonboundary-width-bad-lower-box-code} bounds this
number by
\[
  96319792639339148\binom{2q+25}{50}100^{50}3^{2q-25},
\]
and \cref{lem:bc-stable-dominates-twentyfifth-b-nonboundary-width-count}
absorbs that count into \(s_{2q+25}/2\).  In all cases
\(|\delta_{2q+25}^{B_b}|\le s_{2q+25}/2\), which proves the stated lower
bound.
\end{proof}

\begin{lemma}[Twenty-sixth $B$ nonboundary width-bad paths have a lower-box code]\label{lem:bc-b-twentysixth-nonboundary-width-bad-lower-box-code}
For every \(q\ge15\), the number of column-even Pieri two-strip paths
\[
  \varnothing=\lambda^{(0)}\subset\lambda^{(1)}\subset\cdots
  \subset\lambda^{(2q+26)}
\]
with
\[
  \lambda^{(2q+26)}_1\ge2q
\]
is at most
\[
  24185591292956608915\binom{2q+26}{52}104^{52}3^{2q-26}.
\]
\end{lemma}

\begin{proof}
Put \(n=2q+26\).  The terminal partition has size \(2n=4q+52\).  Since its
columns have even lengths and its first row has length at least \(2q\), the
number of boxes below the second row is at most
\[
  4q+52-2(2q)=52 .
\]
Hence the terminal partition has at most fifty-four nonzero rows.

As in \cref{lem:bc-b-sixth-nonboundary-width-bad-lower-box-code}, encode a
path by the row-pair word and track lower boxes.  A step with no lower box has
three choices, a step with one lower box has at most \(2\cdot52=104\) choices,
and a step with two lower boxes has at most \(\binom{52+1}{2}=1378\) choices.
Thus the number of paths is bounded by the sum of the coefficients of
\(x^0,\ldots,x^{52}\) in
\[
  (3+104x+1378x^2)^n .
\]

The monomial with \(a\) one-lower-box steps and \(c\) two-lower-box steps,
where \(a+2c\le52\), contributes
\[
  \binom{n}{a+c}\binom{a+c}{c}3^{n-a-c}104^a1378^c .
\]
Since \(n\ge56\), \(a+c\le52\), and the factors in the quotient
\(\binom n{52}/\binom n{a+c}\) which are not cancelled are all at least
\(5\),
\[
  \frac{\binom{n}{a+c}\binom{a+c}{c}}{\binom n{52}}
  \le
  \frac{52!}{a!\,c!\,5^{52-a-c}}.
\]
Therefore the coefficient sum is at most
\[
  \binom n{52}104^{52}3^{n-52}
  \sum_{\substack{a,c\ge0\\ a+2c\le52}}
  \frac{52!}{a!\,c!\,5^{52-a-c}}
  \frac{3^{52-a-c}1378^c}{104^{52-a}}.
\]
The exact replay
\[
\texttt{certificates/post\_m29/bc\_twentysixth\_frontier\_gmp\_certificate.log}
\]
prints this finite rational sum as
\[
  \frac{10515921361294745138317260438626617216441934753953658489934041385438143468232284823955926238983677843094936348125098252699625877660467}
       {434801086064710736754792110167239250958093807148894244421109245230847098880000000000000000000000000000000000000000}
  <24185591292956608915 .
\]
This gives the displayed bound.
\end{proof}

\begin{lemma}[Stable moments dominate the twenty-sixth $B$ nonboundary width count above the first row]\label{lem:bc-stable-dominates-twentysixth-b-nonboundary-width-count}
For every \(q\ge46\),
\[
  48371182585913217830\binom{2q+26}{52}104^{52}3^{2q-26}\le s_{2q+26}.
\]
\end{lemma}

\begin{proof}
By \cref{lem:bc-littlewood-graph-model}, \(s_{2q+26}\) contains all labelled
simple cycles on \(2q+26\) vertices, whose number is \((2q+25)!/2\).  It is
enough to prove
\[
  96742365171826435660\binom{2q+26}{52}104^{52}3^{2q-26}\le (2q+25)! .
\]
At \(q=46\), the left side is
\[
  24518630791967223663879224636187908518451788630001417644052433922417408816879166518893832643245474863896645192993141244974868343737632188923129425593309956864566879309733993729763149631979520,
\]
whereas the right side is
\[
  3969937160808720895401959629498630647790406360168322301129748464310422041758630649341780708631240196854767624444057168110272995649603642560353748940315749184568295424000000000000000000000000000.
\]
The second integer exceeds the first by
\[
  3945418530016753671738080404862442739271954571538320883485696030388004632941751482822886875987994721990870979251064026865298127305866010371430619514722439227703728544690266006270236850368020480 .
\]
If the displayed factorial inequality holds at \(q\), then at \(q+1\) the
left side is multiplied by
\[
  9\frac{(2q+28)(2q+27)}{(2q-24)(2q-25)}<30,
\]
because \(84q^2-3930q+11196>0\) for \(q\ge46\).  The right side is multiplied
by
\[
  (2q+26)(2q+27)\ge118\cdot119>30 .
\]
Induction proves the result.
\end{proof}

\begin{lemma}[Exact first thirty-one twenty-sixth $B$ nonboundary bad-shape counts]\label{lem:bc-b-twentysixth-frontier-badshape-certificate}
For \(15\le q\le45\), the
number of column-even Pieri two-strip paths of length \(2q+26\) with terminal
first row at least \(2q\) is at most \(s_{2q+26}/2\).
\end{lemma}

\begin{proof}
The exact C++/GMP/OpenMP replay
\[
\texttt{certificates/post\_m29/bc\_twentysixth\_frontier\_gmp\_certificate.log}
\]
implements the reverse Pieri recursion of
\cref{lem:bc-b-eighth-frontier-badshape-certificate} for
\(15\le q\le45\),
\(j=2q+26\).  It exits with status \(0\) in each B partition, reports
\texttt{SUMMARY failures=0} in each B partition, and has SHA-256
\[
\hashsplit{6a107c7c20c107e6a71db9aa65cf2884}{df0e1e94db82c20d7d3f1ae3acc4c07d}.
\]
The transcript contains exactly thirty-one \texttt{B\_MARGIN} records, one for
each displayed \(q\), and each record is the exact positive integer
\[
  s_{2q+26}-2\#\mathrm{bad}(q,2q+26).
\]
Thus every displayed bad-shape count is at most \(s_{2q+26}/2\).
\end{proof}

\begin{proposition}[Residual $B$ twenty-sixth nonboundary lower correction]\label{prop:bc-b-twentysixth-nonboundary-lower-correction}
For every residual row \(B_b\) with \(14\le b\le123\), put \(q=b+1\) and
write \(s_j=m_j^{\mathrm{st}}\) and \(m_j^{B_b}=s_j+\delta_j^{B_b}\).  Then
\[
  -\frac12 s_{2q+26}\le \delta_{2q+26}^{B_b}.
\]
\end{proposition}

\begin{proof}
By \cref{lem:bc-correction-bad-shapes}, the coefficient interpretation in
\cref{lem:bad-shape-witnesses}, and the Pieri-path bijection
\cref{lem:bc-pieri-path-bijection}, \(|\delta_{2q+26}^{B_b}|\) is the number
of column-even Pieri two-strip paths of length \(2q+26\) with terminal first
row at least \(2q\).  If
\(15\le q\le45\), this is
bounded by \(s_{2q+26}/2\) by
\cref{lem:bc-b-h8-h27-bounded-littlewood-certificate}.  If \(q\ge46\),
\cref{lem:bc-b-twentysixth-nonboundary-width-bad-lower-box-code} bounds this
number by
\[
  24185591292956608915\binom{2q+26}{52}104^{52}3^{2q-26},
\]
and \cref{lem:bc-stable-dominates-twentysixth-b-nonboundary-width-count}
absorbs that count into \(s_{2q+26}/2\).  In all cases
\(|\delta_{2q+26}^{B_b}|\le s_{2q+26}/2\), which proves the stated lower
bound.
\end{proof}

\begin{lemma}[Twenty-seventh $B$ nonboundary width-bad paths have a lower-box code]\label{lem:bc-b-twentyseventh-nonboundary-width-bad-lower-box-code}
For every \(q\ge15\), the number of column-even Pieri two-strip paths
\[
  \varnothing=\lambda^{(0)}\subset\lambda^{(1)}\subset\cdots
  \subset\lambda^{(2q+27)}
\]
with
\[
  \lambda^{(2q+27)}_1\ge2q
\]
is at most
\[
  21113589594277165270589\binom{2q+27}{54}108^{54}3^{2q-27}.
\]
\end{lemma}

\begin{proof}
Put \(n=2q+27\).  The terminal partition has size \(2n=4q+54\).  Since its
columns have even lengths and its first row has length at least \(2q\), the
number of boxes below the second row is at most
\[
  4q+54-2(2q)=54 .
\]
Hence the terminal partition has at most fifty-six nonzero rows.

As in \cref{lem:bc-b-sixth-nonboundary-width-bad-lower-box-code}, encode a
path by the row-pair word and track lower boxes.  A step with no lower box has
three choices, a step with one lower box has at most \(2\cdot54=108\) choices,
and a step with two lower boxes has at most \(\binom{54+1}{2}=1485\) choices.
Thus the number of paths is bounded by the sum of the coefficients of
\(x^0,\ldots,x^{54}\) in
\[
  (3+108x+1485x^2)^n .
\]

The monomial with \(a\) one-lower-box steps and \(c\) two-lower-box steps,
where \(a+2c\le54\), contributes
\[
  \binom{n}{a+c}\binom{a+c}{c}3^{n-a-c}108^a1485^c .
\]
Since \(n\ge57\), \(a+c\le54\), and the factors in the quotient
\(\binom n{54}/\binom n{a+c}\) which are not cancelled are all at least
\(4\),
\[
  \frac{\binom{n}{a+c}\binom{a+c}{c}}{\binom n{54}}
  \le
  \frac{54!}{a!\,c!\,4^{54-a-c}}.
\]
Therefore the coefficient sum is at most
\[
  \binom n{54}108^{54}3^{n-54}
  \sum_{\substack{a,c\ge0\\ a+2c\le54}}
  \frac{54!}{a!\,c!\,4^{54-a-c}}
  \frac{3^{54-a-c}1485^c}{108^{54-a}}.
\]
The exact replay
\[
\texttt{certificates/post\_m29/bc\_twentyseventh\_frontier\_gmp\_certificate.log}
\]
prints this finite rational sum as
\[
  \frac{2627146205381501605105934114888073677477007311327886945300155746611238154833874468430700303859194480319027224499}
       {124429159411793678131320841300679691068154374989700419393399303806437927962865201395531776}
  <21113589594277165270589 .
\]
This gives the displayed bound.
\end{proof}

\begin{lemma}[Stable moments dominate the twenty-seventh $B$ nonboundary width count above the first row]\label{lem:bc-stable-dominates-twentyseventh-b-nonboundary-width-count}
For every \(q\ge48\),
\[
  42227179188554330541178\binom{2q+27}{54}108^{54}3^{2q-27}\le s_{2q+27}.
\]
\end{lemma}

\begin{proof}
By \cref{lem:bc-littlewood-graph-model}, \(s_{2q+27}\) contains all labelled
simple cycles on \(2q+27\) vertices, whose number is \((2q+26)!/2\).  It is
enough to prove
\[
  84454358377108661082356\binom{2q+27}{54}108^{54}3^{2q-27}\le (2q+26)! .
\]
At \(q=48\), the left side is
\[
  1382580989788954876046671229620184847513314093131940228783286920916351208602873186022008689840612544163943709058433179375497089498785389185389770907252631190072242900215336209818152644872892879061647360,
\]
whereas the right side is
\[
  98750442008336013624115798714482080125644041369783596059584700502676714572050143649033796427745042294071023050579626404736512939596842694895821378210620013388054747214795243520000000000000000000000000000.
\]
The second integer exceeds the first by
\[
  97367861018547058748069127484861895278130727276651655830801413581760363363447270463011787737904429749907079341521193225361015850098057305710431607303367382197982504314579907310181847355127107120938352640.
\]
If the displayed factorial inequality holds at \(q\), then at \(q+1\) the
left side is multiplied by
\[
  9\frac{(2q+29)(2q+28)}{(2q-25)(2q-26)}<30,
\]
because \(84q^2-4086q+12192>0\) for \(q\ge48\).  The right side is multiplied
by
\[
  (2q+27)(2q+28)\ge123\cdot124>30 .
\]
Induction proves the result.
\end{proof}

\begin{lemma}[Exact first thirty-three twenty-seventh $B$ nonboundary bad-shape counts]\label{lem:bc-b-twentyseventh-frontier-badshape-certificate}
For \(15\le q\le47\), the
number of column-even Pieri two-strip paths of length \(2q+27\) with terminal
first row at least \(2q\) is at most \(s_{2q+27}/2\).
\end{lemma}

\begin{proof}
The exact C++/GMP/OpenMP replay
\[
\texttt{certificates/post\_m29/bc\_twentyseventh\_frontier\_gmp\_certificate.log}
\]
implements the reverse Pieri recursion of
\cref{lem:bc-b-eighth-frontier-badshape-certificate} for
\(15\le q\le47\),
\(j=2q+27\).  It exits with status \(0\) in each B partition, reports
\texttt{SUMMARY failures=0} in each B partition, and has SHA-256
\[
\hashsplit{67df773882805ff4321cd3659e16b1b9}{985db56a76a94d7136448011ef30f80d}.
\]
The transcript contains exactly thirty-three \texttt{B\_MARGIN} records, one
for each displayed \(q\), and each record is the exact positive integer
\[
  s_{2q+27}-2\#\mathrm{bad}(q,2q+27).
\]
Thus every displayed bad-shape count is at most \(s_{2q+27}/2\).
\end{proof}

\begin{proposition}[Residual $B$ twenty-seventh nonboundary lower correction]\label{prop:bc-b-twentyseventh-nonboundary-lower-correction}
For every residual row \(B_b\) with \(14\le b\le123\), put \(q=b+1\) and
write \(s_j=m_j^{\mathrm{st}}\) and \(m_j^{B_b}=s_j+\delta_j^{B_b}\).  Then
\[
  -\frac12 s_{2q+27}\le \delta_{2q+27}^{B_b}.
\]
\end{proposition}

\begin{proof}
By \cref{lem:bc-correction-bad-shapes}, the coefficient interpretation in
\cref{lem:bad-shape-witnesses}, and the Pieri-path bijection
\cref{lem:bc-pieri-path-bijection}, \(|\delta_{2q+27}^{B_b}|\) is the number
of column-even Pieri two-strip paths of length \(2q+27\) with terminal first
row at least \(2q\).  If
\(15\le q\le47\), this is
bounded by \(s_{2q+27}/2\) by
\cref{lem:bc-b-h8-h27-bounded-littlewood-certificate}.  If \(q\ge48\),
\cref{lem:bc-b-twentyseventh-nonboundary-width-bad-lower-box-code} bounds this
number by
\[
  21113589594277165270589\binom{2q+27}{54}108^{54}3^{2q-27},
\]
and \cref{lem:bc-stable-dominates-twentyseventh-b-nonboundary-width-count}
absorbs that count into \(s_{2q+27}/2\).  In all cases
\(|\delta_{2q+27}^{B_b}|\le s_{2q+27}/2\), which proves the stated lower
bound.
\end{proof}

\begin{lemma}[Twenty-eighth $B$ nonboundary width-bad paths have a lower-box code]\label{lem:bc-b-twentyeighth-nonboundary-width-bad-lower-box-code}
For every \(q\ge15\), the number of column-even Pieri two-strip paths
\[
  \varnothing=\lambda^{(0)}\subset\lambda^{(1)}\subset\cdots
  \subset\lambda^{(2q+28)}
\]
with
\[
  \lambda^{(2q+28)}_1\ge2q
\]
is at most
\[
  129062907305803474509175447\binom{2q+28}{56}112^{56}3^{2q-28}.
\]
\end{lemma}

\begin{proof}
Put \(n=2q+28\).  The terminal partition has size \(2n=4q+56\).  Since its
columns have even lengths and its first row has length at least \(2q\), the
number of boxes below the second row is at most
\[
  4q+56-2(2q)=56 .
\]
Hence the terminal partition has at most fifty-eight nonzero rows.

As in \cref{lem:bc-b-sixth-nonboundary-width-bad-lower-box-code}, encode a
path by the row-pair word and track lower boxes.  A step with no lower box has
three choices, a step with one lower box has at most \(2\cdot56=112\) choices,
and a step with two lower boxes has at most \(\binom{56+1}{2}=1596\) choices.
Thus the number of paths is bounded by the sum of the coefficients of
\(x^0,\ldots,x^{56}\) in
\[
  (3+112x+1596x^2)^n .
\]

The monomial with \(a\) one-lower-box steps and \(c\) two-lower-box steps,
where \(a+2c\le56\), contributes
\[
  \binom{n}{a+c}\binom{a+c}{c}3^{n-a-c}112^a1596^c .
\]
Since \(n\ge58\), \(a+c\le56\), and the factors in the quotient
\(\binom n{56}/\binom n{a+c}\) which are not cancelled are all at least
\(3\),
\[
  \frac{\binom{n}{a+c}\binom{a+c}{c}}{\binom n{56}}
  \le
  \frac{56!}{a!\,c!\,3^{56-a-c}}.
\]
Therefore the coefficient sum is at most
\[
  \binom n{56}112^{56}3^{n-56}
  \sum_{\substack{a,c\ge0\\ a+2c\le56}}
  \frac{56!}{a!\,c!\,3^{56-a-c}}
  \frac{3^{56-a-c}1596^c}{112^{56-a}}.
\]
The exact replay
\[
\texttt{certificates/post\_m29/bc\_twentyeighth\_b\_absorption\_gmp\_certificate.log}
\]
prints this finite rational sum as
\[
  \frac{2025529548764936535663219997915257739427053097450973064895015969532428326966504020337406927050447841928760131584542621}
       {15694126151719317999571618173922000017876418395539549210232909869689551256045946062675902464}
  <129062907305803474509175447 .
\]
This gives the displayed bound.
\end{proof}

\begin{lemma}[Exact twenty-eighth $B$ variable lower-box absorption]\label{lem:bc-b-twentyeighth-variable-lower-box-absorption-certificate}
For \(42\le q\le49\), put
\[
  R_q=
  \sum_{\substack{a,c\ge0\\ a+2c\le56}}
  \frac{56!}{a!\,c!\,(2q-27)^{56-a-c}}
  \frac{3^{56-a-c}1596^c}{112^{56-a}} .
\]
Then
\[
  2R_q\binom{2q+28}{56}112^{56}3^{2q-28}
  \le s_{2q+28}.
\]
\end{lemma}

\begin{proof}
For a fixed displayed \(q\), the proof of
\cref{lem:bc-b-twentyeighth-nonboundary-width-bad-lower-box-code} may use the
sharper lower bound \(2q-27\) for every uncancelled quotient factor.  Hence
the corresponding bad-shape count is at most
\[
  R_q\binom{2q+28}{56}112^{56}3^{2q-28}.
\]
The exact C++/GMP replay
\[
\texttt{certificates/post\_m29/bc\_twentyeighth\_b\_absorption\_gmp\_certificate.log}
\]
cross-multiplies the eight rational inequalities exactly, exits with status
\(0\), reports \texttt{SUMMARY failures=0}, and has SHA-256
\[
\hashsplit{c1127264a1dacd0dab0341a2c0afbfb1}{1f76063c9f8531c5981953d7a104aa4a}.
\]
Its transcript contains exactly eight \texttt{B\_VARIABLE\_ABSORPTION}
records, one for each displayed \(q\), and each record prints a positive exact
margin numerator.  This proves the displayed inequalities.
\end{proof}

\begin{lemma}[Stable moments dominate the twenty-eighth $B$ nonboundary width count above the first row]\label{lem:bc-stable-dominates-twentyeighth-b-nonboundary-width-count}
For every \(q\ge50\),
\[
  258125814611606949018350894\binom{2q+28}{56}112^{56}3^{2q-28}\le s_{2q+28}.
\]
\end{lemma}

\begin{proof}
By \cref{lem:bc-littlewood-graph-model}, \(s_{2q+28}\) contains all labelled
simple cycles on \(2q+28\) vertices, whose number is \((2q+27)!/2\).  It is
enough to prove
\[
  516251629223213898036701788\binom{2q+28}{56}112^{56}3^{2q-28}\le (2q+27)! .
\]
At \(q=50\), the left side is
\[
  587619932792485841297698709732493321980746498038495779458663513421569462665494436057433168011562205529350986467465004414564735037011434723317921965808368421534260528034746952978553319112791180346822600232114585600,
\]
whereas the right side is
\[
  3012660018457659544809977077527059692324164918673621799053346900596667207618480809067860692097713761984609779945772783965563851033300772326297773087851869982500270661791244122597621760000000000000000000000000000000.
\]
The second integer exceeds the first by
\[
  2425040085665173703512278367794566370343418420635126019594683387175097744952986373010427524086151556455258793478307779550999115996289337602979851122043501560966010133756497169619068440887208819653177399767885414400.
\]
If the displayed factorial inequality holds at \(q\), then at \(q+1\) the
left side is multiplied by
\[
  9\frac{(2q+30)(2q+29)}{(2q-26)(2q-27)}<30,
\]
because \(84q^2-4242q+13230>0\) for \(q\ge50\).  The right side is multiplied
by
\[
  (2q+28)(2q+29)\ge128\cdot129>30 .
\]
Induction proves the result.
\end{proof}

\begin{lemma}[Exact first twenty-seven twenty-eighth $B$ nonboundary bad-shape counts]\label{lem:bc-b-twentyeighth-frontier-badshape-certificate}
For \(15\le q\le41\), the
number of column-even Pieri two-strip paths of length \(2q+28\) with terminal
first row at least \(2q\) is at most \(s_{2q+28}/2\).
\end{lemma}

\begin{proof}
The exact C++/GMP/OpenMP replay
\[
\texttt{certificates/post\_m29/bc\_twentyeighth\_frontier\_gmp\_certificate.log}
\]
implements the reverse Pieri recursion of
\cref{lem:bc-b-eighth-frontier-badshape-certificate} for
\(15\le q\le41\),
\(j=2q+28\).  It exits with status \(0\) in each B partition, reports
\texttt{SUMMARY failures=0} in each B partition, and has SHA-256
\[
\hashsplit{400e6b4c50b78e1c490ce93daca73864}{a2ebfc3380e93c50ac7d502205299b0e}.
\]
The transcript contains exactly twenty-seven \texttt{B\_MARGIN} records, one
for each displayed \(q\), and each record is the exact positive integer
\[
  s_{2q+28}-2\#\mathrm{bad}(q,2q+28).
\]
Thus every displayed bad-shape count is at most \(s_{2q+28}/2\).
\end{proof}

\begin{proposition}[Residual $B$ twenty-eighth nonboundary lower correction]\label{prop:bc-b-twentyeighth-nonboundary-lower-correction}
For every residual row \(B_b\) with \(14\le b\le123\), put \(q=b+1\) and
write \(s_j=m_j^{\mathrm{st}}\) and \(m_j^{B_b}=s_j+\delta_j^{B_b}\).  Then
\[
  -\frac12 s_{2q+28}\le \delta_{2q+28}^{B_b}.
\]
\end{proposition}

\begin{proof}
By \cref{lem:bc-correction-bad-shapes}, the coefficient interpretation in
\cref{lem:bad-shape-witnesses}, and the Pieri-path bijection
\cref{lem:bc-pieri-path-bijection}, \(|\delta_{2q+28}^{B_b}|\) is the number
of column-even Pieri two-strip paths of length \(2q+28\) with terminal first
row at least \(2q\).  If
\(15\le q\le41\), this is
bounded by \(s_{2q+28}/2\) by
\cref{lem:bc-b-twentyeighth-frontier-badshape-certificate}.  If
\(42\le q\le49\), the same bound follows from
\cref{lem:bc-b-twentyeighth-variable-lower-box-absorption-certificate}.  If
\(q\ge50\),
\cref{lem:bc-b-twentyeighth-nonboundary-width-bad-lower-box-code} bounds this
number by
\[
  129062907305803474509175447\binom{2q+28}{56}112^{56}3^{2q-28},
\]
and \cref{lem:bc-stable-dominates-twentyeighth-b-nonboundary-width-count}
absorbs that count into \(s_{2q+28}/2\).  In all cases
\(|\delta_{2q+28}^{B_b}|\le s_{2q+28}/2\), which proves the stated lower
bound.
\end{proof}

\begin{lemma}[Twenty-ninth $B$ nonboundary width-bad paths have a lower-box code]\label{lem:bc-b-twentyninth-nonboundary-width-bad-lower-box-code}
For every \(q\ge15\), the number of column-even Pieri two-strip paths
\[
  \varnothing=\lambda^{(0)}\subset\lambda^{(1)}\subset\cdots
  \subset\lambda^{(2q+29)}
\]
with
\[
  \lambda^{(2q+29)}_1\ge2q
\]
is at most
\[
  26657057892939326657049273688667\binom{2q+29}{58}116^{58}3^{2q-29}.
\]
\end{lemma}

\begin{proof}
Put \(n=2q+29\).  The terminal partition has size \(2n=4q+58\).  Since its
columns have even lengths and its first row has length at least \(2q\), the
number of boxes below the second row is at most
\[
  4q+58-2(2q)=58 .
\]
Hence the terminal partition has at most sixty nonzero rows.

As in \cref{lem:bc-b-sixth-nonboundary-width-bad-lower-box-code}, encode a
path by the row-pair word and track lower boxes.  A step with no lower box has
three choices, a step with one lower box has at most \(2\cdot58=116\) choices,
and a step with two lower boxes has at most \(\binom{58+1}{2}=1711\) choices.
Thus the number of paths is bounded by the sum of the coefficients of
\(x^0,\ldots,x^{58}\) in
\[
  (3+116x+1711x^2)^n .
\]

The monomial with \(a\) one-lower-box steps and \(c\) two-lower-box steps,
where \(a+2c\le58\), contributes
\[
  \binom{n}{a+c}\binom{a+c}{c}3^{n-a-c}116^a1711^c .
\]
Since \(n\ge59\), \(a+c\le58\), and the factors in the quotient
\(\binom n{58}/\binom n{a+c}\) which are not cancelled are all at least
\(2\),
\[
  \frac{\binom{n}{a+c}\binom{a+c}{c}}{\binom n{58}}
  \le
  \frac{58!}{a!\,c!\,2^{58-a-c}}.
\]
Therefore the coefficient sum is at most
\[
  \binom n{58}116^{58}3^{n-58}
  \sum_{\substack{a,c\ge0\\ a+2c\le58}}
  \frac{58!}{a!\,c!\,2^{58-a-c}}
  \frac{3^{58-a-c}1711^c}{116^{58-a}}.
\]
The exact replay
\[
\texttt{certificates/post\_m29/bc\_twentyninth\_b\_absorption\_gmp\_certificate.log}
\]
prints this finite rational sum as
\[
  \frac{277779106855808268919410844062670919179911198214031840717487498638936112277059652420979367222105109134880580581288858691889356285997892023429345943573}
       {10420471305251725205780667230686101341858409721635912006716989825031102329848667800732237308856347868892719711199952896}
  <26657057892939326657049273688667 .
\]
This gives the displayed bound.
\end{proof}

\begin{lemma}[Exact twenty-ninth $B$ variable lower-box absorption]\label{lem:bc-b-twentyninth-variable-lower-box-absorption-certificate}
For \(43\le q\le52\), put
\[
  R_q=
  \sum_{\substack{a,c\ge0\\ a+2c\le58}}
  \frac{58!}{a!\,c!\,(2q-28)^{58-a-c}}
  \frac{3^{58-a-c}1711^c}{116^{58-a}} .
\]
Then
\[
  2R_q\binom{2q+29}{58}116^{58}3^{2q-29}
  \le s_{2q+29}.
\]
\end{lemma}

\begin{proof}
For a fixed displayed \(q\), the proof of
\cref{lem:bc-b-twentyninth-nonboundary-width-bad-lower-box-code} may use the
sharper lower bound \(2q-28\) for every uncancelled quotient factor.  Hence
the corresponding bad-shape count is at most
\[
  R_q\binom{2q+29}{58}116^{58}3^{2q-29}.
\]
The exact C++/GMP replay
\[
\texttt{certificates/post\_m29/bc\_twentyninth\_b\_absorption\_gmp\_certificate.log}
\]
cross-multiplies the ten rational inequalities exactly, exits with status
\(0\), reports \texttt{SUMMARY failures=0}, and has SHA-256
\[
\hashsplit{8fba7d987f90866756d4e1edaca8395d}{780c116c775f14c9d8718e848732f355}.
\]
Its transcript contains exactly ten \texttt{B\_VARIABLE\_ABSORPTION} records,
one for each displayed \(q\), and each record prints a positive exact margin
numerator.  This proves the displayed inequalities.
\end{proof}

\begin{lemma}[Stable moments dominate the twenty-ninth $B$ nonboundary width count above the first row]\label{lem:bc-stable-dominates-twentyninth-b-nonboundary-width-count}
For every \(q\ge53\),
\[
  53314115785878653314098547377334\binom{2q+29}{58}116^{58}3^{2q-29}\le s_{2q+29}.
\]
\end{lemma}

\begin{proof}
By \cref{lem:bc-littlewood-graph-model}, \(s_{2q+29}\) contains all labelled
simple cycles on \(2q+29\) vertices, whose number is \((2q+28)!/2\).  It is
enough to prove
\[
  106628231571757306628197094754668\binom{2q+29}{58}116^{58}3^{2q-29}\le (2q+28)! .
\]
At \(q=53\), the left side is
\[
  252065354206657676926114802925973737564636826587243684665665452735913512758742494339799108660499914732866265504420398625899669935265304050254641154393177320811027368879086081678397678468712605716783232149947189375611469352140800,
\]
whereas the right side is
\[
  1992942746161518876737324194182948445222558439641379420268181650403332765909000151088003004705990626788174304488982644980314383013435909872030129915047718433336536425741860482713199069412263880294400000000000000000000000000000000.
\]
The second integer exceeds the first by
\[
  1740877391954861199811209391256974707657921613054135735602516197667419253150257656748203896045490712055308038984562246354414713078170605821775488760654541112525509056862774401034801390943551274577616767850052810624388530647859200.
\]
If the displayed factorial inequality holds at \(q\), then at \(q+1\) the
left side is multiplied by
\[
  9\frac{(2q+31)(2q+30)}{(2q-27)(2q-28)}<30,
\]
because \(84q^2-4398q+14310>0\) for \(q\ge53\).  The right side is multiplied
by
\[
  (2q+29)(2q+30)\ge135\cdot136>30 .
\]
Induction proves the result.
\end{proof}

\begin{proposition}[Residual $B$ twenty-ninth nonboundary lower correction above the exact frontier]\label{prop:bc-b-twentyninth-high-nonboundary-lower-correction}
For every residual row \(B_b\) with \(42\le b\le123\), put \(q=b+1\) and
write \(s_j=m_j^{\mathrm{st}}\) and \(m_j^{B_b}=s_j+\delta_j^{B_b}\).  Then
\[
  -\frac12 s_{2q+29}\le \delta_{2q+29}^{B_b}.
\]
\end{proposition}

\begin{proof}
By \cref{lem:bc-correction-bad-shapes}, the coefficient interpretation in
\cref{lem:bad-shape-witnesses}, and the Pieri-path bijection
\cref{lem:bc-pieri-path-bijection}, \(|\delta_{2q+29}^{B_b}|\) is the number
of column-even Pieri two-strip paths of length \(2q+29\) with terminal first
row at least \(2q\).  If \(43\le q\le52\), this number is
bounded by \(s_{2q+29}/2\) by
\cref{lem:bc-b-twentyninth-variable-lower-box-absorption-certificate}.  If
\(q\ge53\),
\cref{lem:bc-b-twentyninth-nonboundary-width-bad-lower-box-code} bounds this
number by
\[
  26657057892939326657049273688667\binom{2q+29}{58}116^{58}3^{2q-29},
\]
and \cref{lem:bc-stable-dominates-twentyninth-b-nonboundary-width-count}
absorbs that count into \(s_{2q+29}/2\).  In all cases
\(|\delta_{2q+29}^{B_b}|\le s_{2q+29}/2\), which proves the stated lower
bound.
\end{proof}

\begin{lemma}[$C$-boundary height-bad paths are alternating]\label{lem:bc-c-boundary-height-bad-alternating}
Let $q\ge1$, and let $\lambda^\bullet$ be a Pieri two-strip path of length
$2q$ such that
\[
  \ell(\lambda^{(2q)})\ge2q .
\]
Then
\[
  \ell(\lambda^{(t)})=t\qquad(0\le t\le2q)
\]
and
\[
  C^h_q(\lambda^\bullet)=\{1,3,\ldots,2q-1\}.
\]
\end{lemma}

\begin{proof}
A horizontal two-strip can create at most one new row.  Thus the height
sequence starts at $0$, increases by at most $1$ at each of the $2q$ steps,
and ends at height at least $2q$.  Hence every step creates exactly one new
row, giving $\ell(\lambda^{(t)})=t$ for all $t$.

With this height formula,
\[
  \min(q,\lceil\ell(\lambda^{(t)})/2\rceil)=\min(q,\lceil t/2\rceil).
\]
The displayed statistic increases exactly at the odd times
$t=1,3,\ldots,2q-1$ and is constant at the even times.  This is precisely the
claimed height-crossing set.
\end{proof}

\begin{proposition}[Residual $C$-boundary lower correction]\label{prop:bc-c-boundary-lower-correction}
For every residual row $C_b$ with $20\le b\le217$, put $q=b+1$ and write
$s_j=m_j^{\mathrm{st}}$ and $m_j^{C_b}=s_j+\delta_j^{C_b}$.  Then
\[
  -\frac12 s_{2q}\le \delta_{2q}^{C_b}.
\]
\end{proposition}

\begin{proof}
By \cref{lem:bc-correction-bad-shapes} and the coefficient interpretation in
\cref{lem:bad-shape-witnesses}, $|\delta_{2q}^{C_b}|$ is the number of
semistandard tableaux $T$ of content $(2^{2q})$ whose shape
$\lambda(T)$ has $\lambda(T)'$ even and $\ell(\lambda(T))\ge2q$.  Under the
Pieri-path bijection of \cref{lem:bc-pieri-path-bijection}, every such tableau
gives a path satisfying the hypothesis of
\cref{lem:bc-c-boundary-height-bad-alternating}.  By
\cref{lem:bc-witness-frontier-crossings}, all of these tableaux therefore lie
in the single height fixed-witness fiber with
\[
  S=\{1,3,\ldots,2q-1\}.
\]
Conversely, the defining height fixed-witness fiber already has
$\ell(\lambda(T))\ge2q$.  Hence the bad-shape set is exactly that fiber.

Since $q\ge21$, \cref{cor:bc-height-alternating-full-zone-decoder} gives
\[
  |\delta_{2q}^{C_b}|\le 2^{2q}s_q .
\]
The exact GMP/OpenMP replay
\[
\texttt{certificates/post\_m29/bc\_pieri\_powerloss\_2q\_gmp\_certificate.log}
\]
checks the stronger boundary inequality
\[
  2^{2q+1}\binom{2q}{q}s_q\le s_{2q}
\]
for every residual $C_b$ row, reports \texttt{failures=0}, and has SHA-256
\[
\hashsplit{124068c5daa1123088a3cd79c3f6b839}{fee45a1aaef637ed82161c465d86bf45}.
\]
Since $\binom{2q}{q}\ge1$, the preceding two inequalities give
$|\delta_{2q}^{C_b}|\le s_{2q}/2$.  This implies the stated lower bound.
\end{proof}

\begin{lemma}[Even-column standard tableaux have a stable graph code]\label{lem:bc-even-column-standard-tableaux-2n-bound}
Let \(s_n\) be the stable Littlewood moment sequence of
\cref{lem:bc-littlewood-graph-model}.  For every \(n\ge2\), the number of
standard Young tableaux of size \(2n\) whose columns all have even length is
at most
\[
  2^{2n}s_n .
\]
\end{lemma}

\begin{proof}
Let \(U\) be such a tableau.  The Robinson--Schensted inverse applied to
\((U,U)\) gives an involution \(\pi\) on \(\{1,\ldots,2n\}\).  By
Rains' fixed-point/odd-column form of Schensted's theorem
\cite[Lemma 3.3]{Rains1997}, the fixed points of \(\pi\) are counted by the
odd columns of \(U\).  Hence \(\pi\) is fixed-point-free.

Contract the adjacent pairs
\[
  \{1,2\},\{3,4\},\ldots,\{2n-1,2n\}
\]
to vertices \(1,\ldots,n\).  The matching \(\pi\) becomes a labelled
degree-two multigraph \(\Gamma^\circ\) on \(n\) vertices, with loops allowed.
Let \(L\) be its loop set.  By
\cref{lem:bc-looped-degree-two-graph-repair}, \(\Gamma^\circ\) is recoverable
from the loopless repaired graph \(\mathcal R(\Gamma^\circ)\) and \(L\).
There are \(s_n\) choices for the repaired graph by
\cref{lem:bc-littlewood-graph-model}, and \(2^n\) choices for \(L\).

After \(\Gamma^\circ\) is known, one further bit at each vertex records which
of the two canonically ordered incident edge-slots receives the odd port
\(2v-1\); the port \(2v\) receives the other slot.  These \(n\) bits recover
the original matching \(\pi\), hence recover \(U\) by Schensted insertion.
Thus \(U\) injects into a loopless stable graph, a loop-set word, and a slot
word, giving at most \(s_n2^n2^n=2^{2n}s_n\) possibilities.
\end{proof}

\begin{lemma}[Stable moments are monotone after two]\label{lem:bc-stable-moments-monotone-after-two}
For the stable Littlewood moment sequence \(s_n\) of
\cref{lem:bc-littlewood-graph-model},
\[
  s_n\le s_{n+1}\qquad(n\ge2).
\]
\end{lemma}

\begin{proof}
By \cref{lem:bc-littlewood-graph-model}, \(s_n\) counts loopless labelled
degree-two multigraphs on \(\{1,\ldots,n\}\), with parallel edges allowed.
For such a graph, choose the lexicographically least edge \(ab\), with
parallel edges ordered by multiplicity.  Delete this edge and insert the new
vertex \(n+1\) by adding the two edges \(a(n+1)\) and \((n+1)b\).  The result
is again loopless and degree two at every vertex.

The construction is injective: from its output, delete the vertex \(n+1\) and
its two incident edges, then restore one edge between its two neighbours.  This
recovers the original graph.  Hence the number of graphs on \(n\) vertices is
at most the number on \(n+1\) vertices.
\end{proof}

\begin{lemma}[$C$ first-nonboundary height-bad paths have a one-pause code]\label{lem:bc-c-first-nonboundary-one-pause-code}
For \(q\ge1\), the number of column-even Pieri two-strip paths
\[
  \varnothing=\lambda^{(0)}\subset\lambda^{(1)}\subset\cdots
  \subset\lambda^{(2q+1)}
\]
with
\[
  \ell(\lambda^{(2q+1)})\ge2q
\]
is at most
\[
  (2q+1)2^{2q+2}s_{q+1}.
\]
\end{lemma}

\begin{proof}
The terminal height is even, because \((\lambda^{(2q+1)})'\) has only even
parts.  A Pieri two-strip can create at most one new row, so the terminal
height is at most \(2q+1\).  The displayed lower bound therefore forces
\[
  \ell(\lambda^{(2q+1)})=2q .
\]
Consequently exactly one step, call it \(\tau\), does not create a new row,
and every other step creates one new row.

Delete the first column from every \(\lambda^{(t)}\), shift the remaining
boxes one column to the left, and call the resulting partition
\(\nu^{(t)}\).  If \(t\ne\tau\), the step \(t\) creates a new first-column
box and one off-column box, so \(\nu^{(t)}/\nu^{(t-1)}\) has size \(1\).  At
\(t=\tau\), no first-column box is created and the two boxes of the Pieri
strip are both off the first column, so
\(\nu^{(\tau)}/\nu^{(\tau-1)}\) is a horizontal strip of size \(2\).  Filling
the boxes of each skew difference with the label \(t\) gives a semistandard
tableau \(U\) of size \(2q+2\), with the label \(\tau\) appearing twice and
every other label appearing once.  Its shape has only even columns, since it
is obtained from the column-even terminal shape by deleting the first column,
whose height is \(2q\).

The pair \((\tau,U)\) recovers the original path: the subtableau of labels at
most \(t\) recovers \(\nu^{(t)}\), and the first column has height \(t\) for
\(t<\tau\) and height \(t-1\) for \(t\ge\tau\).  Adding that first column back
and shifting \(\nu^{(t)}\) one column to the right reconstructs
\(\lambda^{(t)}\).

Standardize \(U\) by replacing the two \(\tau\)-entries by \(\tau\) and
\(\tau+1\), in left-to-right order, and increasing every later label by one.
This gives a standard Young tableau of the same even-column shape.  Since
\(\tau\) is recorded, merging the entries \(\tau,\tau+1\) and decreasing later
entries recovers \(U\).  Thus the original paths inject into the choice of
\(\tau\in\{1,\ldots,2q+1\}\) and an even-column standard tableau of size
\(2q+2\).  Applying
\cref{lem:bc-even-column-standard-tableaux-2n-bound} with \(n=q+1\) gives the
claimed bound.
\end{proof}

\begin{lemma}[The central binomial factor absorbs the one-pause loss]\label{lem:bc-first-c-pause-binomial-absorption}
For every \(q\ge3\),
\[
  4(2q+1)\le \binom{2q+1}{q}.
\]
\end{lemma}

\begin{proof}
At \(q=3\), this is \(28\le35\).  If the inequality holds at \(q\), then
\[
 \frac{\binom{2q+3}{q+1}}{\binom{2q+1}{q}}
 =
 \frac{(2q+3)(2q+2)}{(q+1)(q+2)}
 \ge
 \frac{2q+3}{2q+1},
\]
because \(2(2q+1)\ge q+2\).  Hence
\[
  \binom{2q+3}{q+1}
  \ge 4(2q+1)\frac{2q+3}{2q+1}
  =4(2q+3).
\]
Induction proves the claim.
\end{proof}

\begin{proposition}[Residual $C$ first nonboundary lower correction]\label{prop:bc-c-first-nonboundary-lower-correction}
For every residual row \(C_b\) with \(20\le b\le217\), put \(q=b+1\) and
write \(s_j=m_j^{\mathrm{st}}\) and \(m_j^{C_b}=s_j+\delta_j^{C_b}\).  Then
\[
  -\frac12 s_{2q+1}\le \delta_{2q+1}^{C_b}.
\]
\end{proposition}

\begin{proof}
By \cref{lem:bc-correction-bad-shapes} and the coefficient interpretation in
\cref{lem:bad-shape-witnesses}, \(|\delta_{2q+1}^{C_b}|\) is the number of
semistandard tableaux \(T\) of content \((2^{2q+1})\) whose shape
\(\lambda(T)\) has \(\lambda(T)'\) even and
\(\ell(\lambda(T))\ge2q\).  Under \cref{lem:bc-pieri-path-bijection}, these
tableaux are exactly the column-even Pieri two-strip paths of length
\(2q+1\) with terminal height at least \(2q\).  Therefore
\cref{lem:bc-c-first-nonboundary-one-pause-code} gives
\[
  |\delta_{2q+1}^{C_b}|
  \le
  (2q+1)2^{2q+2}s_{q+1}.
\]
Since \(q\ge21\), \cref{lem:bc-first-c-pause-binomial-absorption} gives
\[
  (2q+1)2^{2q+2}s_{q+1}
  \le
  2^{2q}\binom{2q+1}{q}s_{q+1}.
\]
The exact GMP/OpenMP replay
\[
\texttt{certificates/post\_m29/bc\_pieri\_powerloss\_2q\_gmp\_certificate.log}
\]
checks the inequality
\[
  2^{2q+1}\binom{2q+1}{q}s_{q+1}\le s_{2q+1}
\]
on this residual row and moment as part of its all-row residual-window audit;
it reports \texttt{failures=0} and has SHA-256
\[
\hashsplit{124068c5daa1123088a3cd79c3f6b839}{fee45a1aaef637ed82161c465d86bf45}.
\]
Combining the displayed inequalities gives
\(|\delta_{2q+1}^{C_b}|\le s_{2q+1}/2\), and hence the stated lower bound.
\end{proof}

\begin{lemma}[$C$ second-nonboundary height-bad paths have a two-pause code]\label{lem:bc-c-second-nonboundary-two-pause-code}
For \(q\ge1\), the number of column-even Pieri two-strip paths
\[
  \varnothing=\lambda^{(0)}\subset\lambda^{(1)}\subset\cdots
  \subset\lambda^{(2q+2)}
\]
with
\[
  \ell(\lambda^{(2q+2)})\ge2q
\]
is at most
\[
  \biggl(4+16\binom{2q+2}{2}\biggr)2^{2q}s_{q+2}.
\]
\end{lemma}

\begin{proof}
The terminal height is even and at most \(2q+2\).  Thus the terminal height is
either \(2q+2\) or \(2q\).

If the terminal height is \(2q+2\), every step creates a new row.  Deleting the
first column and shifting the remaining boxes one column left gives an
even-column standard tableau of size \(2q+2\), and this deletion map is
injective by the same reconstruction used in
\cref{lem:bc-c-first-nonboundary-one-pause-code}.  Therefore this case has at
most
\[
  2^{2q+2}s_{q+1}\le 2^{2q+2}s_{q+2}
\]
paths, by
\cref{lem:bc-even-column-standard-tableaux-2n-bound,lem:bc-stable-moments-monotone-after-two}.

If the terminal height is \(2q\), exactly two steps fail to create a new row.
Let \(P\subset\{1,\ldots,2q+2\}\) be this two-element pause set.  Deleting the
first column and shifting left gives a semistandard tableau of size \(2q+4\)
whose shape has only even columns, whose labels in \(P\) occur twice, and
whose other labels occur once.  The pair \((P,U)\) recovers the original path:
the labels at most \(t\) recover the shifted off-column shape at time \(t\),
and the first-column height is \(t-\#(P\cap\{1,\ldots,t\})\).

For fixed \(P\), standardize \(U\) by replacing the two equal entries for each
label in \(P\) by consecutive labels, in left-to-right order, and shifting
later labels accordingly.  Since \(P\) is recorded, merging those consecutive
pairs recovers \(U\).  Hence, for each \(P\), the paths inject into
even-column standard tableaux of size \(2q+4\).  By
\cref{lem:bc-even-column-standard-tableaux-2n-bound}, this gives at most
\[
  \binom{2q+2}{2}2^{2q+4}s_{q+2}
\]
paths in the two-pause case.

Adding the no-pause and two-pause bounds gives the displayed estimate.
\end{proof}

\begin{lemma}[The central binomial factor absorbs the two-pause loss]\label{lem:bc-second-c-pause-binomial-absorption}
For every \(q\ge6\),
\[
  4+16\binom{2q+2}{2}\le \binom{2q+2}{q}.
\]
\end{lemma}

\begin{proof}
At \(q=6\), the inequality is
\[
  1460\le3003=\binom{14}{6}.
\]
Put \(R_q=4+16\binom{2q+2}{2}=32q^2+48q+20\).  For \(q\ge1\),
\[
  \frac{\binom{2q+4}{q+1}}{\binom{2q+2}{q}}
  =
  \frac{(2q+4)(2q+3)}{(q+2)^2}
  =
  4-\frac{2}{q+2}
  >3,
\]
whereas
\[
  R_{q+1}<3R_q
\]
is equivalent to \(64q^2+32q-40>0\).  Thus the binomial side grows by a
larger factor than \(R_q\) from \(q\) to \(q+1\), and induction from \(q=6\)
proves the claim.
\end{proof}

\begin{proposition}[Residual $C$ second nonboundary lower correction]\label{prop:bc-c-second-nonboundary-lower-correction}
For every residual row \(C_b\) with \(20\le b\le217\), put \(q=b+1\) and
write \(s_j=m_j^{\mathrm{st}}\) and \(m_j^{C_b}=s_j+\delta_j^{C_b}\).  Then
\[
  -\frac12 s_{2q+2}\le \delta_{2q+2}^{C_b}.
\]
\end{proposition}

\begin{proof}
As in the proof of
\cref{prop:bc-c-first-nonboundary-lower-correction},
\(|\delta_{2q+2}^{C_b}|\) is the number of column-even Pieri two-strip paths
of length \(2q+2\) with terminal height at least \(2q\).  By
\cref{lem:bc-c-second-nonboundary-two-pause-code},
\[
  |\delta_{2q+2}^{C_b}|
  \le
  \biggl(4+16\binom{2q+2}{2}\biggr)2^{2q}s_{q+2}.
\]
Since \(q\ge21\), \cref{lem:bc-second-c-pause-binomial-absorption} gives
\[
  |\delta_{2q+2}^{C_b}|
  \le
  2^{2q}\binom{2q+2}{q}s_{q+2}.
\]
The exact GMP/OpenMP replay
\[
\texttt{certificates/post\_m29/bc\_pieri\_powerloss\_2q\_gmp\_certificate.log}
\]
checks the inequality
\[
  2^{2q+1}\binom{2q+2}{q}s_{q+2}\le s_{2q+2}
\]
on this residual row and moment as part of its all-row residual-window audit;
it reports \texttt{failures=0} and has SHA-256
\[
\hashsplit{124068c5daa1123088a3cd79c3f6b839}{fee45a1aaef637ed82161c465d86bf45}.
\]
Thus \(|\delta_{2q+2}^{C_b}|\le s_{2q+2}/2\), which implies the claimed lower
bound.
\end{proof}

\begin{lemma}[$C$ third-nonboundary height-bad paths have a three-pause code]\label{lem:bc-c-third-nonboundary-three-pause-code}
For \(q\ge1\), the number of column-even Pieri two-strip paths
\[
  \varnothing=\lambda^{(0)}\subset\lambda^{(1)}\subset\cdots
  \subset\lambda^{(2q+3)}
\]
with
\[
  \ell(\lambda^{(2q+3)})\ge2q
\]
is at most
\[
  \biggl(16(2q+3)+64\binom{2q+3}{3}\biggr)2^{2q}s_{q+3}.
\]
\end{lemma}

\begin{proof}
The terminal height is even and at most \(2q+3\).  Thus it is either
\(2q+2\) or \(2q\).

If the terminal height is \(2q+2\), exactly one step fails to create a new
row.  Recording this pause, deleting the first column, shifting left, and
standardizing the repeated pause label gives an injection into the choice of
a pause and an even-column standard tableau of size \(2q+4\), just as in
\cref{lem:bc-c-first-nonboundary-one-pause-code}.  This contributes at most
\[
  (2q+3)2^{2q+4}s_{q+2}
  \le
  16(2q+3)2^{2q}s_{q+3}
\]
paths, using
\cref{lem:bc-even-column-standard-tableaux-2n-bound,lem:bc-stable-moments-monotone-after-two}.

If the terminal height is \(2q\), exactly three steps fail to create a new
row.  Recording the three-element pause set, deleting the first column,
shifting left, and standardizing the three repeated labels gives an injection
into the choice of that pause set and an even-column standard tableau of size
\(2q+6\).  This contributes at most
\[
  \binom{2q+3}{3}2^{2q+6}s_{q+3}
  =
  64\binom{2q+3}{3}2^{2q}s_{q+3}
\]
paths.  Adding the two cases gives the claim.
\end{proof}

\begin{lemma}[The central binomial factor absorbs the three-pause loss]\label{lem:bc-third-c-pause-binomial-absorption}
For every \(q\ge21\),
\[
  16(2q+3)+64\binom{2q+3}{3}\le \binom{2q+3}{q}.
\]
\end{lemma}

\begin{proof}
At \(q=21\),
\[
  16\cdot45+64\binom{45}{3}=908880
  <
  3773655750150
  =
  \binom{45}{21}.
\]
Let
\[
  P(q)=16(2q+3)+64\binom{2q+3}{3}.
\]
For \(q\ge21\),
\[
  P(q+1)<2P(q),
\]
because
\[
  2P(q)-P(q+1)=\frac{16}{3}(16q^3-94q-93)>0.
\]
On the other hand,
\[
  \frac{\binom{2q+5}{q+1}}{\binom{2q+3}{q}}
  =
  \frac{(2q+5)(2q+4)}{(q+1)(q+4)}
  >3 .
\]
Thus the binomial side grows by a larger factor than \(P(q)\), and induction
from \(q=21\) proves the inequality.
\end{proof}

\begin{proposition}[Residual $C$ third nonboundary lower correction]\label{prop:bc-c-third-nonboundary-lower-correction}
For every residual row \(C_b\) with \(20\le b\le217\), put \(q=b+1\) and
write \(s_j=m_j^{\mathrm{st}}\) and \(m_j^{C_b}=s_j+\delta_j^{C_b}\).  Then
\[
  -\frac12 s_{2q+3}\le \delta_{2q+3}^{C_b}.
\]
\end{proposition}

\begin{proof}
As in the proof of
\cref{prop:bc-c-first-nonboundary-lower-correction},
\(|\delta_{2q+3}^{C_b}|\) is the number of column-even Pieri two-strip paths
of length \(2q+3\) with terminal height at least \(2q\).  By
\cref{lem:bc-c-third-nonboundary-three-pause-code},
\[
  |\delta_{2q+3}^{C_b}|
  \le
  \biggl(16(2q+3)+64\binom{2q+3}{3}\biggr)2^{2q}s_{q+3}.
\]
Since \(q\ge21\), \cref{lem:bc-third-c-pause-binomial-absorption} gives
\[
  |\delta_{2q+3}^{C_b}|
  \le
  2^{2q}\binom{2q+3}{q}s_{q+3}.
\]
The exact GMP/OpenMP replay
\[
\texttt{certificates/post\_m29/bc\_pieri\_powerloss\_2q\_gmp\_certificate.log}
\]
checks the inequality
\[
  2^{2q+1}\binom{2q+3}{q}s_{q+3}\le s_{2q+3}
\]
on this residual row and moment as part of its all-row residual-window audit;
it reports \texttt{failures=0} and has SHA-256
\[
\hashsplit{124068c5daa1123088a3cd79c3f6b839}{fee45a1aaef637ed82161c465d86bf45}.
\]
Thus \(|\delta_{2q+3}^{C_b}|\le s_{2q+3}/2\), which implies the claimed lower
bound.
\end{proof}

\begin{lemma}[$C$ fourth-nonboundary height-bad paths have a four-pause code]\label{lem:bc-c-fourth-nonboundary-four-pause-code}
For \(q\ge1\), the number of column-even Pieri two-strip paths
\[
  \varnothing=\lambda^{(0)}\subset\lambda^{(1)}\subset\cdots
  \subset\lambda^{(2q+4)}
\]
with
\[
  \ell(\lambda^{(2q+4)})\ge2q
\]
is at most
\[
  \biggl(16+64\binom{2q+4}{2}+256\binom{2q+4}{4}\biggr)2^{2q}s_{q+4}.
\]
\end{lemma}

\begin{proof}
The terminal height is even and at most \(2q+4\).  Thus it is
\[
  2q+4,\qquad 2q+2,\qquad\text{or}\qquad 2q .
\]

If the terminal height is \(2q+4\), every step creates a new row.  Deleting
the first column and shifting left gives an injection into even-column
standard tableaux of size \(2q+4\), as in
\cref{lem:bc-c-second-nonboundary-two-pause-code}.  This gives at most
\[
  2^{2q+4}s_{q+2}\le16\,2^{2q}s_{q+4}
\]
paths by
\cref{lem:bc-even-column-standard-tableaux-2n-bound,lem:bc-stable-moments-monotone-after-two}.

If the terminal height is \(2q+2\), exactly two steps fail to create a new
row.  Recording the two-element pause set, deleting the first column, shifting
left, and standardizing the two repeated pause labels injects the paths into
the choice of the pause set and an even-column standard tableau of size
\(2q+6\).  This contributes at most
\[
  \binom{2q+4}{2}2^{2q+6}s_{q+3}
  \le
  64\binom{2q+4}{2}2^{2q}s_{q+4}
\]
paths.

If the terminal height is \(2q\), exactly four steps fail to create a new row.
The same deletion and standardization construction with the four-element
pause set injects the paths into the choice of that pause set and an
even-column standard tableau of size \(2q+8\).  This contributes at most
\[
  \binom{2q+4}{4}2^{2q+8}s_{q+4}
  =
  256\binom{2q+4}{4}2^{2q}s_{q+4}
\]
paths.  Adding the three cases gives the claim.
\end{proof}

\begin{lemma}[The central binomial factor absorbs the four-pause loss]\label{lem:bc-fourth-c-pause-binomial-absorption}
For every \(q\ge21\),
\[
  16+64\binom{2q+4}{2}+256\binom{2q+4}{4}\le \binom{2q+4}{q}.
\]
\end{lemma}

\begin{proof}
At \(q=21\),
\[
  16+64\binom{46}{2}+256\binom{46}{4}
  =
  41841616
  <
  6943526580276
  =
  \binom{46}{21}.
\]
Let
\[
  P(q)=16+64\binom{2q+4}{2}+256\binom{2q+4}{4}.
\]
For \(q\ge21\),
\[
  P(q+1)<2P(q),
\]
because
\[
  2P(q)-P(q+1)
  =
  \frac{16}{3}\bigl(32q^4+32q^3-368q^2-932q-657\bigr)>0 .
\]
Indeed the quartic in parentheses is positive at \(q=21\), where it equals
\(6337227\), and its derivative is increasing and positive for \(q\ge21\).
On the other hand,
\[
  \frac{\binom{2q+6}{q+1}}{\binom{2q+4}{q}}
  =
  \frac{(2q+6)(2q+5)}{(q+1)(q+5)}
  >3,
\]
since the numerator minus three times the denominator is \(q^2+4q+15\).
Thus the binomial side grows by a larger factor than \(P(q)\), and induction
from \(q=21\) proves the inequality.
\end{proof}

\begin{proposition}[Residual $C$ fourth nonboundary lower correction]\label{prop:bc-c-fourth-nonboundary-lower-correction}
For every residual row \(C_b\) with \(20\le b\le217\), put \(q=b+1\) and
write \(s_j=m_j^{\mathrm{st}}\) and \(m_j^{C_b}=s_j+\delta_j^{C_b}\).  Then
\[
  -\frac12 s_{2q+4}\le \delta_{2q+4}^{C_b}.
\]
\end{proposition}

\begin{proof}
As in the proof of
\cref{prop:bc-c-first-nonboundary-lower-correction},
\(|\delta_{2q+4}^{C_b}|\) is the number of column-even Pieri two-strip paths
of length \(2q+4\) with terminal height at least \(2q\).  By
\cref{lem:bc-c-fourth-nonboundary-four-pause-code},
\[
  |\delta_{2q+4}^{C_b}|
  \le
  \biggl(16+64\binom{2q+4}{2}+256\binom{2q+4}{4}\biggr)2^{2q}s_{q+4}.
\]
Since \(q\ge21\), \cref{lem:bc-fourth-c-pause-binomial-absorption} gives
\[
  |\delta_{2q+4}^{C_b}|
  \le
  2^{2q}\binom{2q+4}{q}s_{q+4}.
\]
The exact GMP/OpenMP replay
\[
\texttt{certificates/post\_m29/bc\_pieri\_powerloss\_2q\_gmp\_certificate.log}
\]
checks the inequality
\[
  2^{2q+1}\binom{2q+4}{q}s_{q+4}\le s_{2q+4}
\]
on this residual row and moment as part of its all-row residual-window audit;
it reports \texttt{failures=0} and has SHA-256
\[
\hashsplit{124068c5daa1123088a3cd79c3f6b839}{fee45a1aaef637ed82161c465d86bf45}.
\]
Thus \(|\delta_{2q+4}^{C_b}|\le s_{2q+4}/2\), which implies the claimed lower
bound.
\end{proof}

\begin{lemma}[$C$ fifth-nonboundary height-bad paths have a five-pause code]\label{lem:bc-c-fifth-nonboundary-five-pause-code}
For \(q\ge1\), the number of column-even Pieri two-strip paths
\[
  \varnothing=\lambda^{(0)}\subset\lambda^{(1)}\subset\cdots
  \subset\lambda^{(2q+5)}
\]
with
\[
  \ell(\lambda^{(2q+5)})\ge2q
\]
is at most
\[
  \biggl(64(2q+5)+256\binom{2q+5}{3}
    +1024\binom{2q+5}{5}\biggr)2^{2q}s_{q+5}.
\]
\end{lemma}

\begin{proof}
The terminal height is even and at most \(2q+5\).  Thus it is
\[
  2q+4,\qquad 2q+2,\qquad\text{or}\qquad 2q .
\]

If the terminal height is \(2q+4\), exactly one step fails to create a new
row.  Recording the pause, deleting the first column, shifting left, and
standardizing the repeated pause label injects the paths into the choice of
that pause and an even-column standard tableau of size \(2q+6\).  This
contributes at most
\[
  (2q+5)2^{2q+6}s_{q+3}
  \le
  64(2q+5)2^{2q}s_{q+5}
\]
paths.

If the terminal height is \(2q+2\), exactly three steps fail to create a new
row.  The same deletion and standardization construction gives at most
\[
  \binom{2q+5}{3}2^{2q+8}s_{q+4}
  \le
  256\binom{2q+5}{3}2^{2q}s_{q+5}
\]
paths.

If the terminal height is \(2q\), exactly five steps fail to create a new row.
Recording the five-element pause set and standardizing the five repeated
labels injects the paths into the choice of that set and an even-column
standard tableau of size \(2q+10\).  This contributes at most
\[
  \binom{2q+5}{5}2^{2q+10}s_{q+5}
  =
  1024\binom{2q+5}{5}2^{2q}s_{q+5}
\]
paths.  Adding the three cases gives the claim.
\end{proof}

\begin{lemma}[The central binomial factor absorbs the five-pause loss]\label{lem:bc-fifth-c-pause-binomial-absorption}
For every \(q\ge21\),
\[
  64(2q+5)+256\binom{2q+5}{3}+1024\binom{2q+5}{5}
  \le \binom{2q+5}{q}.
\]
\end{lemma}

\begin{proof}
At \(q=21\),
\[
  64\cdot47+256\binom{47}{3}+1024\binom{47}{5}
  =
  1574907584
  <
  12551759587422
  =
  \binom{47}{21}.
\]
Let
\[
  P(q)=64(2q+5)+256\binom{2q+5}{3}+1024\binom{2q+5}{5}.
\]
For \(q\ge21\),
\[
  P(q+1)<2P(q),
\]
because
\[
  2P(q)-P(q+1)
  =
  \frac{64}{15}
  (64q^5+160q^4-1120q^3-5560q^2-9054q-5415)>0 .
\]
Indeed the polynomial in parentheses is positive at \(q=21\), where it equals
\(279479595\).  Its derivative equals \(66436626\) at \(q=21\), and its
second derivative is positive for \(q\ge21\), so the polynomial remains
positive thereafter.
On the other hand,
\[
  \frac{\binom{2q+7}{q+1}}{\binom{2q+5}{q}}
  =
  \frac{(2q+7)(2q+6)}{(q+1)(q+6)}
  >3,
\]
since the numerator minus three times the denominator is \(q^2+5q+24\).
Thus the binomial side grows by a larger factor than \(P(q)\), and induction
from \(q=21\) proves the inequality.
\end{proof}

\begin{proposition}[Residual $C$ fifth nonboundary lower correction]\label{prop:bc-c-fifth-nonboundary-lower-correction}
For every residual row \(C_b\) with \(20\le b\le217\), put \(q=b+1\) and
write \(s_j=m_j^{\mathrm{st}}\) and \(m_j^{C_b}=s_j+\delta_j^{C_b}\).  Then
\[
  -\frac12 s_{2q+5}\le \delta_{2q+5}^{C_b}.
\]
\end{proposition}

\begin{proof}
As in the proof of
\cref{prop:bc-c-first-nonboundary-lower-correction},
\(|\delta_{2q+5}^{C_b}|\) is the number of column-even Pieri two-strip paths
of length \(2q+5\) with terminal height at least \(2q\).  By
\cref{lem:bc-c-fifth-nonboundary-five-pause-code},
\[
  |\delta_{2q+5}^{C_b}|
  \le
  \biggl(64(2q+5)+256\binom{2q+5}{3}
    +1024\binom{2q+5}{5}\biggr)2^{2q}s_{q+5}.
\]
Since \(q\ge21\), \cref{lem:bc-fifth-c-pause-binomial-absorption} gives
\[
  |\delta_{2q+5}^{C_b}|
  \le
  2^{2q}\binom{2q+5}{q}s_{q+5}.
\]
The exact GMP/OpenMP replay
\[
\texttt{certificates/post\_m29/bc\_pieri\_powerloss\_2q\_gmp\_certificate.log}
\]
checks the inequality
\[
  2^{2q+1}\binom{2q+5}{q}s_{q+5}\le s_{2q+5}
\]
on this residual row and moment as part of its all-row residual-window audit;
it reports \texttt{failures=0} and has SHA-256
\[
\hashsplit{124068c5daa1123088a3cd79c3f6b839}{fee45a1aaef637ed82161c465d86bf45}.
\]
Thus \(|\delta_{2q+5}^{C_b}|\le s_{2q+5}/2\), which implies the claimed lower
bound.
\end{proof}

\begin{lemma}[$C$ sixth-nonboundary height-bad paths have a six-pause code]\label{lem:bc-c-sixth-nonboundary-six-pause-code}
For \(q\ge1\), the number of column-even Pieri two-strip paths
\[
  \varnothing=\lambda^{(0)}\subset\lambda^{(1)}\subset\cdots
  \subset\lambda^{(2q+6)}
\]
with
\[
  \ell(\lambda^{(2q+6)})\ge2q
\]
is at most
\[
  \biggl(64+256\binom{2q+6}{2}+1024\binom{2q+6}{4}
    +4096\binom{2q+6}{6}\biggr)2^{2q}s_{q+6}.
\]
\end{lemma}

\begin{proof}
The terminal height is even and at most \(2q+6\).  Thus it is
\[
  2q+6,\qquad 2q+4,\qquad 2q+2,\qquad\text{or}\qquad 2q .
\]

If the terminal height is \(2q+6\), every step creates a new row.  Deleting
the first column and shifting left gives an injection into even-column
standard tableaux of size \(2q+6\).  This gives at most
\[
  2^{2q+6}s_{q+3}\le64\,2^{2q}s_{q+6}
\]
paths, by
\cref{lem:bc-even-column-standard-tableaux-2n-bound,lem:bc-stable-moments-monotone-after-two}.

If the terminal height is \(2q+4\), exactly two steps fail to create a new
row.  Recording the two-element pause set, deleting the first column, shifting
left, and standardizing the two repeated pause labels injects the paths into
the choice of the pause set and an even-column standard tableau of size
\(2q+8\).  This contributes at most
\[
  \binom{2q+6}{2}2^{2q+8}s_{q+4}
  \le
  256\binom{2q+6}{2}2^{2q}s_{q+6}
\]
paths, by the same two lemmas.

If the terminal height is \(2q+2\), exactly four steps fail to create a new
row.  The same deletion and standardization construction gives at most
\[
  \binom{2q+6}{4}2^{2q+10}s_{q+5}
  \le
  1024\binom{2q+6}{4}2^{2q}s_{q+6}
\]
paths, again by the same two lemmas.

If the terminal height is \(2q\), exactly six steps fail to create a new row.
Recording the six-element pause set and standardizing the six repeated labels
injects the paths into the choice of that set and an even-column standard
tableau of size \(2q+12\).  This contributes at most
\[
  \binom{2q+6}{6}2^{2q+12}s_{q+6}
  =
  4096\binom{2q+6}{6}2^{2q}s_{q+6}
\]
paths, by \cref{lem:bc-even-column-standard-tableaux-2n-bound}.  Adding the
four cases gives the claim.
\end{proof}

\begin{lemma}[The central binomial factor absorbs the six-pause loss]\label{lem:bc-sixth-c-pause-binomial-absorption}
For every \(q\ge21\),
\[
  64+256\binom{2q+6}{2}+1024\binom{2q+6}{4}
    +4096\binom{2q+6}{6}
  \le \binom{2q+6}{q}.
\]
\end{lemma}

\begin{proof}
At \(q=21\),
\[
  64+256\binom{48}{2}+1024\binom{48}{4}+4096\binom{48}{6}
  =
  50463651904
  <
  22314239266528
  =
  \binom{48}{21}.
\]
Let
\[
  P(q)=64+256\binom{2q+6}{2}+1024\binom{2q+6}{4}
    +4096\binom{2q+6}{6}.
\]
For \(q\ge21\),
\[
  P(q+1)<2P(q),
\]
because
\[
  2P(q)-P(q+1)
  =
  \frac{64}{45}
  (256q^6+1152q^5-5600q^4-50880q^3
  -143696q^2-189852q-103275)>0 .
\]
Indeed the polynomial in parentheses equals \(25033257945\) at \(q=21\),
its derivative equals \(7112403972\) at \(q=21\), and its second derivative is
\[
  32(240q^4+720q^3-2100q^2-9540q-8981),
\]
which is positive for \(q\ge21\).  On the other hand,
\[
  \frac{\binom{2q+8}{q+1}}{\binom{2q+6}{q}}
  =
  \frac{(2q+8)(2q+7)}{(q+1)(q+7)}
  >3,
\]
since the numerator minus three times the denominator is \(q^2+6q+35\).
Thus the binomial side grows by a larger factor than \(P(q)\), and induction
from \(q=21\) proves the inequality.
\end{proof}

\begin{proposition}[Residual $C$ sixth nonboundary lower correction]\label{prop:bc-c-sixth-nonboundary-lower-correction}
For every residual row \(C_b\) with \(20\le b\le217\), put \(q=b+1\) and
write \(s_j=m_j^{\mathrm{st}}\) and \(m_j^{C_b}=s_j+\delta_j^{C_b}\).  Then
\[
  -\frac12 s_{2q+6}\le \delta_{2q+6}^{C_b}.
\]
\end{proposition}

\begin{proof}
As in the proof of
\cref{prop:bc-c-first-nonboundary-lower-correction},
\(|\delta_{2q+6}^{C_b}|\) is the number of column-even Pieri two-strip paths
of length \(2q+6\) with terminal height at least \(2q\).  By
\cref{lem:bc-c-sixth-nonboundary-six-pause-code},
\[
  |\delta_{2q+6}^{C_b}|
  \le
  \biggl(64+256\binom{2q+6}{2}+1024\binom{2q+6}{4}
    +4096\binom{2q+6}{6}\biggr)2^{2q}s_{q+6}.
\]
Since \(q\ge21\), \cref{lem:bc-sixth-c-pause-binomial-absorption} gives
\[
  |\delta_{2q+6}^{C_b}|
  \le
  2^{2q}\binom{2q+6}{q}s_{q+6}.
\]
The exact GMP/OpenMP replay
\[
\texttt{certificates/post\_m29/bc\_pieri\_powerloss\_2q\_gmp\_certificate.log}
\]
checks the inequality
\[
  2^{2q+1}\binom{2q+6}{q}s_{q+6}\le s_{2q+6}
\]
on this residual row and moment as part of its all-row residual-window audit;
it reports \texttt{failures=0} and has SHA-256
\[
\hashsplit{124068c5daa1123088a3cd79c3f6b839}{fee45a1aaef637ed82161c465d86bf45}.
\]
Thus \(|\delta_{2q+6}^{C_b}|\le s_{2q+6}/2\), which implies the claimed lower
bound.
\end{proof}

\begin{lemma}[$C$ seventh-nonboundary height-bad paths have a seven-pause code]\label{lem:bc-c-seventh-nonboundary-seven-pause-code}
For \(q\ge1\), the number of column-even Pieri two-strip paths
\[
  \varnothing=\lambda^{(0)}\subset\lambda^{(1)}\subset\cdots
  \subset\lambda^{(2q+7)}
\]
with
\[
  \ell(\lambda^{(2q+7)})\ge2q
\]
is at most
\[
  \biggl(256(2q+7)+1024\binom{2q+7}{3}
    +4096\binom{2q+7}{5}+16384\binom{2q+7}{7}\biggr)
  2^{2q}s_{q+7}.
\]
\end{lemma}

\begin{proof}
The terminal height is even and at most \(2q+7\).  Thus it is
\[
  2q+6,\qquad 2q+4,\qquad 2q+2,\qquad\text{or}\qquad 2q .
\]

If the terminal height is \(2q+6\), exactly one step fails to create a new
row.  Deleting the first column, shifting left, and standardizing the repeated
pause label injects the paths into the choice of that pause and an
even-column standard tableau of size \(2q+8\).  This contributes at most
\[
  (2q+7)2^{2q+8}s_{q+4}
  \le
  256(2q+7)2^{2q}s_{q+7}
\]
paths by
\cref{lem:bc-even-column-standard-tableaux-2n-bound,lem:bc-stable-moments-monotone-after-two}.

If the terminal height is \(2q+4\), exactly three steps fail to create a new
row.  The same deletion and standardization construction gives at most
\[
  \binom{2q+7}{3}2^{2q+10}s_{q+5}
  \le
  1024\binom{2q+7}{3}2^{2q}s_{q+7}
\]
paths, by the same two lemmas.

If the terminal height is \(2q+2\), exactly five steps fail to create a new
row.  The same construction gives at most
\[
  \binom{2q+7}{5}2^{2q+12}s_{q+6}
  \le
  4096\binom{2q+7}{5}2^{2q}s_{q+7}
\]
paths, by the same two lemmas.

If the terminal height is \(2q\), exactly seven steps fail to create a new
row.  Recording the seven-element pause set and standardizing the seven
repeated labels injects the paths into the choice of that set and an
even-column standard tableau of size \(2q+14\).  This contributes at most
\[
  \binom{2q+7}{7}2^{2q+14}s_{q+7}
  =
  16384\binom{2q+7}{7}2^{2q}s_{q+7}
\]
paths by \cref{lem:bc-even-column-standard-tableaux-2n-bound}.  Adding the
four cases gives the claim.
\end{proof}

\begin{lemma}[The central binomial factor absorbs the seven-pause loss]\label{lem:bc-seventh-c-pause-binomial-absorption}
For every \(q\ge21\),
\[
  256(2q+7)+1024\binom{2q+7}{3}
    +4096\binom{2q+7}{5}+16384\binom{2q+7}{7}
  \le \binom{2q+7}{q}.
\]
\end{lemma}

\begin{proof}
At \(q=21\),
\[
  256\cdot49+1024\binom{49}{3}
    +4096\binom{49}{5}+16384\binom{49}{7}
  =
  1415224643840
  <
  39049918716424
  =
  \binom{49}{21}.
\]
Let
\[
  P(q)=256(2q+7)+1024\binom{2q+7}{3}
    +4096\binom{2q+7}{5}+16384\binom{2q+7}{7}.
\]
For \(q\ge21\),
\[
  P(q+1)<2P(q),
\]
because
\[
  2P(q)-P(q+1)
  =
  \frac{256}{315}
  (512q^7+3584q^6-11200q^5-196000q^4-855232q^3
  -1876504q^2-2193990q-1124865)>0 .
\]
Indeed the polynomial in parentheses equals \(1136887862985\) at \(q=21\),
its derivative equals \(375846246474\) at \(q=21\), and its second derivative
is
\[
  112(192q^5+960q^4-2000q^3-21000q^2-45816q-33509),
\]
which is positive for \(q\ge21\): the inner polynomial equals \(942070507\)
at \(q=21\), its derivative equals \(218690184\) at \(q=21\), and its second
derivative is \(240(2q+5)(8q^2+4q-35)>0\) for \(q\ge21\).  On the other hand,
\[
  \frac{\binom{2q+9}{q+1}}{\binom{2q+7}{q}}
  =
  \frac{(2q+9)(2q+8)}{(q+1)(q+8)}
  >3,
\]
since the numerator minus three times the denominator is \(q^2+7q+48\).
Thus the binomial side grows by a larger factor than \(P(q)\), and induction
from \(q=21\) proves the inequality.
\end{proof}

\begin{proposition}[Residual $C$ seventh nonboundary lower correction]\label{prop:bc-c-seventh-nonboundary-lower-correction}
For every residual row \(C_b\) with \(20\le b\le217\), put \(q=b+1\) and
write \(s_j=m_j^{\mathrm{st}}\) and \(m_j^{C_b}=s_j+\delta_j^{C_b}\).  Then
\[
  -\frac12 s_{2q+7}\le \delta_{2q+7}^{C_b}.
\]
\end{proposition}

\begin{proof}
As in the proof of
\cref{prop:bc-c-first-nonboundary-lower-correction},
\(|\delta_{2q+7}^{C_b}|\) is the number of column-even Pieri two-strip paths
of length \(2q+7\) with terminal height at least \(2q\).  By
\cref{lem:bc-c-seventh-nonboundary-seven-pause-code},
\[
  |\delta_{2q+7}^{C_b}|
  \le
  \biggl(256(2q+7)+1024\binom{2q+7}{3}
    +4096\binom{2q+7}{5}+16384\binom{2q+7}{7}\biggr)
  2^{2q}s_{q+7}.
\]
Since \(q\ge21\), \cref{lem:bc-seventh-c-pause-binomial-absorption} gives
\[
  |\delta_{2q+7}^{C_b}|
  \le
  2^{2q}\binom{2q+7}{q}s_{q+7}.
\]
The exact GMP/OpenMP replay
\[
\texttt{certificates/post\_m29/bc\_pieri\_powerloss\_2q\_gmp\_certificate.log}
\]
checks the inequality
\[
  2^{2q+1}\binom{2q+7}{q}s_{q+7}\le s_{2q+7}
\]
on this residual row and moment as part of its all-row residual-window audit;
it reports \texttt{failures=0} and has SHA-256
\[
\hashsplit{124068c5daa1123088a3cd79c3f6b839}{fee45a1aaef637ed82161c465d86bf45}.
\]
Thus \(|\delta_{2q+7}^{C_b}|\le s_{2q+7}/2\), which implies the claimed lower
bound.
\end{proof}

\begin{lemma}[$C$ eighth-nonboundary height-bad paths have an eight-pause code]\label{lem:bc-c-eighth-nonboundary-eight-pause-code}
For \(q\ge1\), the number of column-even Pieri two-strip paths
\[
  \varnothing=\lambda^{(0)}\subset\lambda^{(1)}\subset\cdots
  \subset\lambda^{(2q+8)}
\]
with
\[
  \ell(\lambda^{(2q+8)})\ge2q
\]
is at most
\[
  \biggl(256+1024\binom{2q+8}{2}
    +4096\binom{2q+8}{4}+16384\binom{2q+8}{6}
    +65536\binom{2q+8}{8}\biggr)2^{2q}s_{q+8}.
\]
\end{lemma}

\begin{proof}
The terminal height is even and at most \(2q+8\).  Thus it is
\[
  2q+8,\qquad 2q+6,\qquad 2q+4,\qquad 2q+2,\qquad\text{or}\qquad 2q .
\]
These cases have respectively \(0,2,4,6,8\) steps which fail to create a new
row.  Deleting the first column, shifting left, and standardizing the repeated
pause labels injects the paths with \(p\) pauses into the choice of the
\(p\)-element pause set and an even-column standard tableau of size
\(2q+8+p\).  Therefore the five cases contribute at most
\[
  2^{2q+8}s_{q+4},
  \quad
  \binom{2q+8}{2}2^{2q+10}s_{q+5},
  \quad
  \binom{2q+8}{4}2^{2q+12}s_{q+6},
\]
\[
  \binom{2q+8}{6}2^{2q+14}s_{q+7},
  \quad
  \binom{2q+8}{8}2^{2q+16}s_{q+8},
\]
by
\cref{lem:bc-even-column-standard-tableaux-2n-bound,lem:bc-stable-moments-monotone-after-two}.
After replacing \(s_{q+4},\ldots,s_{q+7}\) by \(s_{q+8}\), the sum is the
displayed bound.
\end{proof}

\begin{lemma}[The central binomial factor absorbs the eight-pause loss]\label{lem:bc-eighth-c-pause-binomial-absorption}
For every \(q\ge21\),
\[
  256+1024\binom{2q+8}{2}
    +4096\binom{2q+8}{4}+16384\binom{2q+8}{6}
    +65536\binom{2q+8}{8}
  \le \binom{2q+8}{q}.
\]
\end{lemma}

\begin{proof}
At \(q=21\),
\[
  256+1024\binom{50}{2}
    +4096\binom{50}{4}+16384\binom{50}{6}+65536\binom{50}{8}
  =
  35446176998656
  <
  67327446062800
  =
  \binom{50}{21}.
\]
Let
\[
  P(q)=256+1024\binom{2q+8}{2}
    +4096\binom{2q+8}{4}+16384\binom{2q+8}{6}
    +65536\binom{2q+8}{8}.
\]
For \(q\ge21\), \(P(q+1)<2P(q)\).  Indeed, with \(x=q-21\),
\[
\begin{aligned}
315(2P(q)-P(q+1))
={}&131072x^8+23330816x^7+1809317888x^6\\
&+79798304768x^5+2187606622208x^4\\
&+38136956125184x^3+412400070438912x^2\\
&+2525280659192832x+6690395111258880,
\end{aligned}
\]
which is positive for \(x\ge0\).  On the other hand,
\[
  \frac{\binom{2q+10}{q+1}}{\binom{2q+8}{q}}
  =
  \frac{(2q+10)(2q+9)}{(q+1)(q+9)}
  >3,
\]
since the numerator minus three times the denominator is \(q^2+8q+63\).
Thus the binomial side grows by a larger factor than \(P(q)\), and induction
from \(q=21\) proves the inequality.
\end{proof}

\begin{proposition}[Residual $C$ eighth nonboundary lower correction]\label{prop:bc-c-eighth-nonboundary-lower-correction}
For every residual row \(C_b\) with \(20\le b\le217\), put \(q=b+1\) and
write \(s_j=m_j^{\mathrm{st}}\) and \(m_j^{C_b}=s_j+\delta_j^{C_b}\).  Then
\[
  -\frac12 s_{2q+8}\le \delta_{2q+8}^{C_b}.
\]
\end{proposition}

\begin{proof}
As in the proof of
\cref{prop:bc-c-first-nonboundary-lower-correction},
\(|\delta_{2q+8}^{C_b}|\) is the number of column-even Pieri two-strip paths
of length \(2q+8\) with terminal height at least \(2q\).  By
\cref{lem:bc-c-eighth-nonboundary-eight-pause-code},
\[
  |\delta_{2q+8}^{C_b}|
  \le
  \biggl(256+1024\binom{2q+8}{2}
    +4096\binom{2q+8}{4}+16384\binom{2q+8}{6}
    +65536\binom{2q+8}{8}\biggr)2^{2q}s_{q+8}.
\]
Since \(q\ge21\), \cref{lem:bc-eighth-c-pause-binomial-absorption} gives
\[
  |\delta_{2q+8}^{C_b}|
  \le
  2^{2q}\binom{2q+8}{q}s_{q+8}.
\]
The exact GMP/OpenMP replay
\[
\texttt{certificates/post\_m29/bc\_pieri\_powerloss\_2q\_gmp\_certificate.log}
\]
checks the inequality
\[
  2^{2q+1}\binom{2q+8}{q}s_{q+8}\le s_{2q+8}
\]
on this residual row and moment as part of its all-row residual-window audit;
it reports \texttt{failures=0} and has SHA-256
\[
\hashsplit{124068c5daa1123088a3cd79c3f6b839}{fee45a1aaef637ed82161c465d86bf45}.
\]
Thus \(|\delta_{2q+8}^{C_b}|\le s_{2q+8}/2\), which implies the claimed lower
bound.
\end{proof}

\begin{lemma}[$C$ ninth-nonboundary height-bad paths have a nine-pause code]\label{lem:bc-c-ninth-nonboundary-nine-pause-code}
For \(q\ge1\), the number of column-even Pieri two-strip paths
\[
  \varnothing=\lambda^{(0)}\subset\lambda^{(1)}\subset\cdots
  \subset\lambda^{(2q+9)}
\]
with
\[
  \ell(\lambda^{(2q+9)})\ge2q
\]
is at most
\[
  \biggl(1024(2q+9)+4096\binom{2q+9}{3}
    +16384\binom{2q+9}{5}+65536\binom{2q+9}{7}
    +262144\binom{2q+9}{9}\biggr)2^{2q}s_{q+9}.
\]
\end{lemma}

\begin{proof}
The terminal height is even and at most \(2q+9\).  Thus it is
\[
  2q+8,\qquad 2q+6,\qquad 2q+4,\qquad 2q+2,\qquad\text{or}\qquad 2q .
\]
These cases have respectively \(1,3,5,7,9\) steps which fail to create a new
row.  Deleting the first column, shifting left, and standardizing the
repeated pause labels injects the paths with \(p\) pauses into the choice of
the \(p\)-element pause set and an even-column standard tableau of size
\(2q+9+p\).  Therefore the five cases contribute at most
\[
  (2q+9)2^{2q+10}s_{q+5},
  \quad
  \binom{2q+9}{3}2^{2q+12}s_{q+6},
  \quad
  \binom{2q+9}{5}2^{2q+14}s_{q+7},
\]
\[
  \binom{2q+9}{7}2^{2q+16}s_{q+8},
  \quad
  \binom{2q+9}{9}2^{2q+18}s_{q+9},
\]
by
\cref{lem:bc-even-column-standard-tableaux-2n-bound,lem:bc-stable-moments-monotone-after-two}.
After replacing \(s_{q+5},\ldots,s_{q+8}\) by \(s_{q+9}\), the sum is the
displayed bound.
\end{proof}

\begin{lemma}[The central binomial factor absorbs the nine-pause loss]\label{lem:bc-ninth-c-pause-binomial-absorption}
For every \(q\ge23\),
\[
  1024(2q+9)+4096\binom{2q+9}{3}
    +16384\binom{2q+9}{5}+65536\binom{2q+9}{7}
    +262144\binom{2q+9}{9}
  \le \binom{2q+9}{q}.
\]
\end{lemma}

\begin{proof}
At \(q=23\),
\[
  1024\cdot55+4096\binom{55}{3}
    +16384\binom{55}{5}+65536\binom{55}{7}
    +262144\binom{55}{9}
  =
  1680173121915904
  <
  1866442158555975
  =
  \binom{55}{23}.
\]
Let
\[
  P(q)=1024(2q+9)+4096\binom{2q+9}{3}
    +16384\binom{2q+9}{5}+65536\binom{2q+9}{7}
    +262144\binom{2q+9}{9}.
\]
For \(q\ge23\), \(P(q+1)<2P(q)\).  Indeed, with \(x=q-23\),
\[
\begin{aligned}
2835(2P(q)-P(q+1))
={}&1048576x^9+231211008x^8+22580035584x^7\\
&+1281321861120x^6+46534796378112x^5\\
&+1120976919724032x^4+17896390057066496x^3\\
&+182409816397455360x^2+1075667566906386432x\\
&+2791290177480268800,
\end{aligned}
\]
which is positive for \(x\ge0\).  On the other hand,
\[
  \frac{\binom{2q+11}{q+1}}{\binom{2q+9}{q}}
  =
  \frac{(2q+11)(2q+10)}{(q+1)(q+10)}
  >3,
\]
since the numerator minus three times the denominator is \(q^2+9q+80\).
Thus the binomial side grows by a larger factor than \(P(q)\), and induction
from \(q=23\) proves the inequality.
\end{proof}

\begin{lemma}[Exact first two ninth $C$ pause-arithmetic checks]\label{lem:bc-c-ninth-frontier-pause-arithmetic-certificate}
For \(q=21,22\), put
\[
  P(q)=1024(2q+9)+4096\binom{2q+9}{3}
    +16384\binom{2q+9}{5}+65536\binom{2q+9}{7}
    +262144\binom{2q+9}{9}.
\]
Then
\[
  2^{2q+1}P(q)s_{q+9}\le s_{2q+9}.
\]
\end{lemma}

\begin{proof}
The exact C++/GMP/OpenMP replay
\[
\texttt{certificates/post\_m29/bc\_ninth\_frontier\_gmp\_certificate.log}
\]
computes the displayed integer inequality for \(q=21,22\).  It exits with
status \(0\), reports \texttt{SUMMARY failures=0}, and has SHA-256
\[
\hashsplit{3459b8f1c779c878c9a1666f18b27a03}{548202729f4ddaa23399b6f2c11c242c}.
\]
It verifies the exact positive margins
\[
\begin{array}{c|c|c}
q & P(q) & s_{2q+9}-2^{2q+1}P(q)s_{q+9}\\ \hline
21&
805149937982464&
95209191753417477281958974480308259512574357282473262219794837952\\
22&
1171869936268288&
257420773307602153259394842419604777034268451060763250279932333172864 .
\end{array}
\]
This proves the claim.
\end{proof}

\begin{proposition}[Residual $C$ ninth nonboundary lower correction]\label{prop:bc-c-ninth-nonboundary-lower-correction}
For every residual row \(C_b\) with \(20\le b\le217\), put \(q=b+1\) and
write \(s_j=m_j^{\mathrm{st}}\) and \(m_j^{C_b}=s_j+\delta_j^{C_b}\).  Then
\[
  -\frac12 s_{2q+9}\le \delta_{2q+9}^{C_b}.
\]
\end{proposition}

\begin{proof}
As in the proof of
\cref{prop:bc-c-first-nonboundary-lower-correction},
\(|\delta_{2q+9}^{C_b}|\) is the number of column-even Pieri two-strip paths
of length \(2q+9\) with terminal height at least \(2q\).  By
\cref{lem:bc-c-ninth-nonboundary-nine-pause-code}, with \(P(q)\) denoting
the polynomial from \cref{lem:bc-c-ninth-frontier-pause-arithmetic-certificate},
\[
  |\delta_{2q+9}^{C_b}|
  \le
  P(q)\,2^{2q}s_{q+9},
\]
If \(q=21,22\),
\cref{lem:bc-c-ninth-frontier-pause-arithmetic-certificate} gives
\[
  |\delta_{2q+9}^{C_b}|\le s_{2q+9}/2 .
\]
If \(q\ge23\), \cref{lem:bc-ninth-c-pause-binomial-absorption} gives
\[
  |\delta_{2q+9}^{C_b}|
  \le
  2^{2q}\binom{2q+9}{q}s_{q+9}.
\]
The exact GMP/OpenMP replay
\[
\texttt{certificates/post\_m29/bc\_pieri\_powerloss\_2q\_gmp\_certificate.log}
\]
checks the inequality
\[
  2^{2q+1}\binom{2q+9}{q}s_{q+9}\le s_{2q+9}
\]
on this residual row and moment as part of its all-row residual-window audit;
it reports \texttt{failures=0} and has SHA-256
\[
\hashsplit{124068c5daa1123088a3cd79c3f6b839}{fee45a1aaef637ed82161c465d86bf45}.
\]
Thus \(|\delta_{2q+9}^{C_b}|\le s_{2q+9}/2\) also for \(q\ge23\), which
implies the claimed lower bound.
\end{proof}

\begin{lemma}[$C$ tenth-nonboundary height-bad paths have a ten-pause code]\label{lem:bc-c-tenth-nonboundary-ten-pause-code}
For \(q\ge1\), the number of column-even Pieri two-strip paths
\[
  \varnothing=\lambda^{(0)}\subset\lambda^{(1)}\subset\cdots
  \subset\lambda^{(2q+10)}
\]
with
\[
  \ell(\lambda^{(2q+10)})\ge2q
\]
is at most
\[
  \biggl(1024+4096\binom{2q+10}{2}
    +16384\binom{2q+10}{4}+65536\binom{2q+10}{6}
    +262144\binom{2q+10}{8}+1048576\binom{2q+10}{10}
  \biggr)2^{2q}s_{q+10}.
\]
\end{lemma}

\begin{proof}
The terminal height is even and at most \(2q+10\).  Thus it is
\[
  2q+10,\quad 2q+8,\quad 2q+6,\quad 2q+4,\quad 2q+2,\quad\text{or}\quad 2q .
\]
These cases have respectively \(0,2,4,6,8,10\) steps which fail to create a
new row.  Deleting the first column, shifting left, and standardizing the
repeated pause labels injects the paths with \(p\) pauses into the choice of
the \(p\)-element pause set and an even-column standard tableau of size
\(2q+10+p\).  Therefore the six cases contribute at most
\[
  2^{2q+10}s_{q+5},
  \quad
  \binom{2q+10}{2}2^{2q+12}s_{q+6},
  \quad
  \binom{2q+10}{4}2^{2q+14}s_{q+7},
\]
\[
  \binom{2q+10}{6}2^{2q+16}s_{q+8},
  \quad
  \binom{2q+10}{8}2^{2q+18}s_{q+9},
  \quad
  \binom{2q+10}{10}2^{2q+20}s_{q+10},
\]
by
\cref{lem:bc-even-column-standard-tableaux-2n-bound,lem:bc-stable-moments-monotone-after-two}.
After replacing \(s_{q+5},\ldots,s_{q+9}\) by \(s_{q+10}\), the sum is the
displayed bound.
\end{proof}

\begin{lemma}[The central binomial factor absorbs the ten-pause loss]\label{lem:bc-tenth-c-pause-binomial-absorption}
For every \(q\ge26\),
\[
  1024+4096\binom{2q+10}{2}
    +16384\binom{2q+10}{4}+65536\binom{2q+10}{6}
    +262144\binom{2q+10}{8}+1048576\binom{2q+10}{10}
  \le \binom{2q+10}{q}.
\]
\end{lemma}

\begin{proof}
At \(q=26\),
\[
  1024+4096\binom{62}{2}
    +16384\binom{62}{4}+65536\binom{62}{6}
    +262144\binom{62}{8}+1048576\binom{62}{10}
  =
  113632146094453760
  <
  209769429934732479
  =
  \binom{62}{26}.
\]
Let
\[
  P(q)=1024+4096\binom{2q+10}{2}
    +16384\binom{2q+10}{4}+65536\binom{2q+10}{6}
    +262144\binom{2q+10}{8}+1048576\binom{2q+10}{10}.
\]
For \(q\ge26\), \(P(q+1)<2P(q)\).  Indeed, with \(x=q-26\),
\[
\begin{aligned}
14175(2P(q)-P(q+1))
={}&4194304x^{10}+1163919360x^9+144947281920x^8\\
&+10664395407360x^7+513170922012672x^6\\
&+16868613697044480x^5+383409676620922880x^4\\
&+5946352734974115840x^3+60177808034062516224x^2\\
&+358500237510556200960x+953522452127742336000,
\end{aligned}
\]
which is positive for \(x\ge0\).  On the other hand,
\[
  \frac{\binom{2q+12}{q+1}}{\binom{2q+10}{q}}
  =
  \frac{(2q+12)(2q+11)}{(q+1)(q+11)}
  >3,
\]
since the numerator minus three times the denominator is \(q^2+10q+99\).
Thus the binomial side grows by a larger factor than \(P(q)\), and induction
from \(q=26\) proves the inequality.
\end{proof}

\begin{lemma}[Exact first five tenth $C$ pause-arithmetic checks]\label{lem:bc-c-tenth-frontier-pause-arithmetic-certificate}
For \(q=21,22,23,24,25\), put
\[
  P(q)=1024+4096\binom{2q+10}{2}
    +16384\binom{2q+10}{4}+65536\binom{2q+10}{6}
    +262144\binom{2q+10}{8}+1048576\binom{2q+10}{10}.
\]
Then
\[
  2^{2q+1}P(q)s_{q+10}\le s_{2q+10}.
\]
\end{lemma}

\begin{proof}
The exact C++/GMP/OpenMP replay
\[
\texttt{certificates/post\_m29/bc\_tenth\_frontier\_gmp\_certificate.log}
\]
computes the displayed integer inequality for \(q=21,22,23,24,25\).  It exits
with status \(0\), reports \texttt{SUMMARY failures=0}, and has SHA-256
\[
\hashsplit{abf7ac5edc1f162794cea23d37343d9c}{7dda9415a30f8eab1a622884ebe7b49c}.
\]
It verifies the exact positive margins
\[
\begin{array}{c|c|c}
q & P(q) & s_{2q+10}-2^{2q+1}P(q)s_{q+10}\\ \hline
21&
16787109734261760&
4903171355778686064964164515745269611213904591723677230177156061824\\
22&
25367621153403904&
13771958891177334015216861406002714565114571278870369402339849495106432\\
23&
37711207486899200&
41656585385438952634870696811804338216451333106527331667025870754751489920\\
24&
55219289355244544&
135331263665127919819047410128118678155782794529945205279339252342307465590656\\
25&
79730383933629440&
471052485025387657456032265569888807150614594170760972150823054365241098991856896 .
\end{array}
\]
This proves the claim.
\end{proof}

\begin{proposition}[Residual $C$ tenth nonboundary lower correction]\label{prop:bc-c-tenth-nonboundary-lower-correction}
For every residual row \(C_b\) with \(20\le b\le217\), put \(q=b+1\) and
write \(s_j=m_j^{\mathrm{st}}\) and \(m_j^{C_b}=s_j+\delta_j^{C_b}\).  Then
\[
  -\frac12 s_{2q+10}\le \delta_{2q+10}^{C_b}.
\]
\end{proposition}

\begin{proof}
As in the proof of
\cref{prop:bc-c-first-nonboundary-lower-correction},
\(|\delta_{2q+10}^{C_b}|\) is the number of column-even Pieri two-strip paths
of length \(2q+10\) with terminal height at least \(2q\).  By
\cref{lem:bc-c-tenth-nonboundary-ten-pause-code}, with \(P(q)\) denoting the
polynomial from \cref{lem:bc-c-tenth-frontier-pause-arithmetic-certificate},
\[
  |\delta_{2q+10}^{C_b}|
  \le
  P(q)\,2^{2q}s_{q+10}.
\]
If \(q=21,22,23,24,25\),
\cref{lem:bc-c-tenth-frontier-pause-arithmetic-certificate} gives
\[
  |\delta_{2q+10}^{C_b}|\le s_{2q+10}/2 .
\]
If \(q\ge26\), \cref{lem:bc-tenth-c-pause-binomial-absorption} gives
\[
  |\delta_{2q+10}^{C_b}|
  \le
  2^{2q}\binom{2q+10}{q}s_{q+10}.
\]
The exact GMP/OpenMP replay
\[
\texttt{certificates/post\_m29/bc\_pieri\_powerloss\_2q\_gmp\_certificate.log}
\]
checks the inequality
\[
  2^{2q+1}\binom{2q+10}{q}s_{q+10}\le s_{2q+10}
\]
on this residual row and moment as part of its all-row residual-window audit;
it reports \texttt{failures=0} and has SHA-256
\[
\hashsplit{124068c5daa1123088a3cd79c3f6b839}{fee45a1aaef637ed82161c465d86bf45}.
\]
Thus \(|\delta_{2q+10}^{C_b}|\le s_{2q+10}/2\) also for \(q\ge26\), which
implies the claimed lower bound.
\end{proof}

\begin{lemma}[$C$ eleventh-nonboundary height-bad paths have an eleven-pause code]\label{lem:bc-c-eleventh-nonboundary-eleven-pause-code}
For \(q\ge1\), the number of column-even Pieri two-strip paths
\[
  \varnothing=\lambda^{(0)}\subset\lambda^{(1)}\subset\cdots
  \subset\lambda^{(2q+11)}
\]
with
\[
  \ell(\lambda^{(2q+11)})\ge2q
\]
is at most
\[
  \biggl(4096\binom{2q+11}{1}+16384\binom{2q+11}{3}
    +65536\binom{2q+11}{5}+262144\binom{2q+11}{7}
    +1048576\binom{2q+11}{9}+4194304\binom{2q+11}{11}
  \biggr)2^{2q}s_{q+11}.
\]
\end{lemma}

\begin{proof}
The terminal height is even and at most \(2q+11\).  Thus it is
\[
  2q+10,\quad 2q+8,\quad 2q+6,\quad 2q+4,\quad 2q+2,\quad\text{or}\quad 2q .
\]
These cases have respectively \(1,3,5,7,9,11\) steps which fail to create a
new row.  Deleting the first column, shifting left, and standardizing the
repeated pause labels injects the paths with \(p\) pauses into the choice of
the \(p\)-element pause set and an even-column standard tableau of size
\(2q+11+p\).  Therefore the six cases contribute at most
\[
  \binom{2q+11}{1}2^{2q+12}s_{q+6},
  \quad
  \binom{2q+11}{3}2^{2q+14}s_{q+7},
  \quad
  \binom{2q+11}{5}2^{2q+16}s_{q+8},
\]
\[
  \binom{2q+11}{7}2^{2q+18}s_{q+9},
  \quad
  \binom{2q+11}{9}2^{2q+20}s_{q+10},
  \quad
  \binom{2q+11}{11}2^{2q+22}s_{q+11},
\]
by
\cref{lem:bc-even-column-standard-tableaux-2n-bound,lem:bc-stable-moments-monotone-after-two}.
After replacing \(s_{q+6},\ldots,s_{q+10}\) by \(s_{q+11}\), the sum is the
displayed bound.
\end{proof}

\begin{lemma}[The central binomial factor absorbs the eleven-pause loss]\label{lem:bc-eleventh-c-pause-binomial-absorption}
For every \(q\ge28\),
\[
  4096\binom{2q+11}{1}+16384\binom{2q+11}{3}
    +65536\binom{2q+11}{5}+262144\binom{2q+11}{7}
    +1048576\binom{2q+11}{9}+4194304\binom{2q+11}{11}
  \le \binom{2q+11}{q}.
\]
\end{lemma}

\begin{proof}
At \(q=28\),
\[
  4096\binom{67}{1}+16384\binom{67}{3}
    +65536\binom{67}{5}+262144\binom{67}{7}
    +1048576\binom{67}{9}+4194304\binom{67}{11}
  =
  5435009637461815296
  <
  5864393356544251760
  =
  \binom{67}{28}.
\]
Let
\[
  P(q)=4096\binom{2q+11}{1}+16384\binom{2q+11}{3}
    +65536\binom{2q+11}{5}+262144\binom{2q+11}{7}
    +1048576\binom{2q+11}{9}+4194304\binom{2q+11}{11}.
\]
For \(q\ge28\), \(P(q+1)<2P(q)\).  Indeed, with \(x=q-28\),
\[
\begin{aligned}
155925(2P(q)-P(q+1))
={}&33554432x^{11}+11072962560x^{10}+1657022709760x^9\\
&+148393961717760x^8+8834179154313216x^7\\
&+366968144391045120x^6+10849539239009320960x^5\\
&+228200591272315453440x^4+3344494768984406884352x^3\\
&+32506342800225239531520x^2+188410100820385114644480x\\
&+492831435774247034572800,
\end{aligned}
\]
which is positive for \(x\ge0\).  On the other hand,
\[
  \frac{\binom{2q+13}{q+1}}{\binom{2q+11}{q}}
  =
  \frac{(2q+13)(2q+12)}{(q+1)(q+12)}
  >3,
\]
since the numerator minus three times the denominator is \(q^2+11q+120\).
Thus the binomial side grows by a larger factor than \(P(q)\), and induction
from \(q=28\) proves the inequality.
\end{proof}

\begin{lemma}[Exact first seven eleventh $C$ pause-arithmetic checks]\label{lem:bc-c-eleventh-frontier-pause-arithmetic-certificate}
For \(21\le q\le27\), put
\[
  P(q)=4096\binom{2q+11}{1}+16384\binom{2q+11}{3}
    +65536\binom{2q+11}{5}+262144\binom{2q+11}{7}
    +1048576\binom{2q+11}{9}+4194304\binom{2q+11}{11}.
\]
Then
\[
  2^{2q+1}P(q)s_{q+11}\le s_{2q+11}.
\]
\end{lemma}

\begin{proof}
The exact C++/GMP/OpenMP replay
\[
\texttt{certificates/post\_m29/bc\_eleventh\_frontier\_gmp\_certificate.log}
\]
computes the displayed integer inequality for \(21\le q\le27\).  It
exits with status \(0\), reports \texttt{SUMMARY failures=0}, and has SHA-256
\[
\hashsplit{de69669b8df7784d87b33232131fee8f}{a6eb22c830fdcc8eeb066d851a409c0c}.
\]
It verifies the exact positive margins
\[
\begin{array}{c|c|c}
q & P(q) & s_{2q+11}-2^{2q+1}P(q)s_{q+11}\\ \hline
21&
324393072099643392&
257362697029215877963949523297153447302766944466517748161416488855680\\
22&
508584122346582016&
750559067883129880349273491327444158124656900055892284235763406948255360\\
23&
783391012291121152&
2353589660314073337128089738870306304406389791628228318393371452705364754944\\
24&
1187133100402733056&
7916862822025459545211278821831496164228888938237813980332978702989743176972672\\
25&
1771912890992054272&
28498624463510700078311199299089653133614164147508330272882543415215643039488592640\\
26&
2607774406749368320&
109541225766556320001716025709278150174781456225226335841962103028640068955141950940416\\
27&
3787894083819409408&
448652128227809777832975850852768721615133669475553568921414615289569671480179369963642880 .
\end{array}
\]
This proves the claim.
\end{proof}

\begin{proposition}[Residual $C$ eleventh nonboundary lower correction]\label{prop:bc-c-eleventh-nonboundary-lower-correction}
For every residual row \(C_b\) with \(20\le b\le217\), put \(q=b+1\) and
write \(s_j=m_j^{\mathrm{st}}\) and \(m_j^{C_b}=s_j+\delta_j^{C_b}\).  Then
\[
  -\frac12 s_{2q+11}\le \delta_{2q+11}^{C_b}.
\]
\end{proposition}

\begin{proof}
As in the proof of
\cref{prop:bc-c-first-nonboundary-lower-correction},
\(|\delta_{2q+11}^{C_b}|\) is the number of column-even Pieri two-strip paths
of length \(2q+11\) with terminal height at least \(2q\).  By
\cref{lem:bc-c-eleventh-nonboundary-eleven-pause-code}, with \(P(q)\) denoting
the polynomial from
\cref{lem:bc-c-eleventh-frontier-pause-arithmetic-certificate},
\[
  |\delta_{2q+11}^{C_b}|
  \le
  P(q)\,2^{2q}s_{q+11}.
\]
If \(21\le q\le27\),
\cref{lem:bc-c-eleventh-frontier-pause-arithmetic-certificate} gives
\[
  |\delta_{2q+11}^{C_b}|\le s_{2q+11}/2 .
\]
If \(q\ge28\), \cref{lem:bc-eleventh-c-pause-binomial-absorption} gives
\[
  |\delta_{2q+11}^{C_b}|
  \le
  2^{2q}\binom{2q+11}{q}s_{q+11}.
\]
The exact GMP/OpenMP replay
\[
\texttt{certificates/post\_m29/bc\_pieri\_powerloss\_2q\_gmp\_certificate.log}
\]
checks the inequality
\[
  2^{2q+1}\binom{2q+11}{q}s_{q+11}\le s_{2q+11}
\]
on this residual row and moment as part of its all-row residual-window audit;
it reports \texttt{failures=0} and has SHA-256
\[
\hashsplit{124068c5daa1123088a3cd79c3f6b839}{fee45a1aaef637ed82161c465d86bf45}.
\]
Thus \(|\delta_{2q+11}^{C_b}|\le s_{2q+11}/2\) also for \(q\ge28\), which
implies the claimed lower bound.
\end{proof}

\begin{lemma}[$C$ twelfth-nonboundary height-bad paths have a twelve-pause code]\label{lem:bc-c-twelfth-nonboundary-twelve-pause-code}
For \(q\ge1\), the number of column-even Pieri two-strip paths
\[
  \varnothing=\lambda^{(0)}\subset\lambda^{(1)}\subset\cdots
  \subset\lambda^{(2q+12)}
\]
with
\[
  \ell(\lambda^{(2q+12)})\ge2q
\]
is at most
\[
  \biggl(4096+16384\binom{2q+12}{2}
    +65536\binom{2q+12}{4}+262144\binom{2q+12}{6}
    +1048576\binom{2q+12}{8}+4194304\binom{2q+12}{10}
    +16777216\binom{2q+12}{12}
  \biggr)2^{2q}s_{q+12}.
\]
\end{lemma}

\begin{proof}
The terminal height is even and at most \(2q+12\).  Thus it is
\[
  2q+12,\quad 2q+10,\quad 2q+8,\quad 2q+6,\quad
  2q+4,\quad 2q+2,\quad\text{or}\quad 2q .
\]
These cases have respectively \(0,2,4,6,8,10,12\) steps which fail to create a
new row.  Deleting the first column, shifting left, and standardizing the
repeated pause labels injects the paths with \(p\) pauses into the choice of
the \(p\)-element pause set and an even-column standard tableau of size
\(2q+12+p\).  Therefore the seven cases contribute at most
\[
  2^{2q+12}s_{q+6},
  \quad
  \binom{2q+12}{2}2^{2q+14}s_{q+7},
  \quad
  \binom{2q+12}{4}2^{2q+16}s_{q+8},
\]
\[
  \binom{2q+12}{6}2^{2q+18}s_{q+9},
  \quad
  \binom{2q+12}{8}2^{2q+20}s_{q+10},
  \quad
  \binom{2q+12}{10}2^{2q+22}s_{q+11},
\]
\[
  \binom{2q+12}{12}2^{2q+24}s_{q+12},
\]
by
\cref{lem:bc-even-column-standard-tableaux-2n-bound,lem:bc-stable-moments-monotone-after-two}.
After replacing \(s_{q+6},\ldots,s_{q+11}\) by \(s_{q+12}\), the sum is the
displayed bound.
\end{proof}

\begin{lemma}[The central binomial factor absorbs the twelve-pause loss]\label{lem:bc-twelfth-c-pause-binomial-absorption}
For every \(q\ge31\),
\[
  4096+16384\binom{2q+12}{2}
    +65536\binom{2q+12}{4}+262144\binom{2q+12}{6}
    +1048576\binom{2q+12}{8}+4194304\binom{2q+12}{10}
    +16777216\binom{2q+12}{12}
  \le \binom{2q+12}{q}.
\]
\end{lemma}

\begin{proof}
At \(q=31\),
\[
  4096+16384\binom{74}{2}
    +65536\binom{74}{4}+262144\binom{74}{6}
    +1048576\binom{74}{8}+4194304\binom{74}{10}
    +16777216\binom{74}{12}
  =
  371189423168378720256
  <
  665858611887899825472
  =
  \binom{74}{31}.
\]
Let
\[
  P(q)=4096+16384\binom{2q+12}{2}
    +65536\binom{2q+12}{4}+262144\binom{2q+12}{6}
    +1048576\binom{2q+12}{8}+4194304\binom{2q+12}{10}
    +16777216\binom{2q+12}{12}.
\]
For \(q\ge31\), \(P(q+1)<2P(q)\).  Indeed, with \(x=q-31\),
\[
\begin{aligned}
467775(2P(q)-P(q+1))
={}&67108864x^{12}+26776436736x^{11}+4887236575232x^{10}\\
&+539476120043520x^9+40103480490196992x^8\\
&+2114600225942274048x^7+81074772196129243136x^6\\
&+2276668272958085529600x^5+46454921071694256799744x^4\\
&+671432360820605377314816x^3+6521508020414066359468032x^2\\
&+38194452157657256601108480x+101923801617413045181542400,
\end{aligned}
\]
which is positive for \(x\ge0\).  On the other hand,
\[
  \frac{\binom{2q+14}{q+1}}{\binom{2q+12}{q}}
  =
  \frac{(2q+14)(2q+13)}{(q+1)(q+13)}
  >3,
\]
since the numerator minus three times the denominator is \(q^2+12q+143\).
Thus the binomial side grows by a larger factor than \(P(q)\), and induction
from \(q=31\) proves the inequality.
\end{proof}

\begin{lemma}[Exact first ten twelfth $C$ pause-arithmetic checks]\label{lem:bc-c-twelfth-frontier-pause-arithmetic-certificate}
For \(21\le q\le30\), put
\[
  P(q)=4096+16384\binom{2q+12}{2}
    +65536\binom{2q+12}{4}+262144\binom{2q+12}{6}
    +1048576\binom{2q+12}{8}+4194304\binom{2q+12}{10}
    +16777216\binom{2q+12}{12}.
\]
Then
\[
  2^{2q+1}P(q)s_{q+12}\le s_{2q+12}.
\]
\end{lemma}

\begin{proof}
The exact C++/GMP/OpenMP replay
\[
\texttt{certificates/post\_m29/bc\_twelfth\_frontier\_gmp\_certificate.log}
\]
computes the displayed integer inequality for
\(21\le q\le30\).  It exits with status \(0\), reports
\texttt{SUMMARY failures=0}, and has SHA-256
\[
\hashsplit{027740d3f1fc0d56fcb27fa4e8d1f32a}{35bf7f4d0938ae970eeaadef8c4ba875}.
\]
It verifies the exact positive margins
\[
\begin{array}{c|c|c}
q & P(q) & s_{2q+12}-2^{2q+1}P(q)s_{q+12}\\ \hline
21&
5856171146995879936&
13737887743093491029874945629543498680225743342441692732099355335971712\\
22&
9518962187314073600&
41649163935462476049958505836798783090818697479057364159395837163257243520\\
23&
15182710965852655616&
135329630065808428579424292893519281225757899062377328264131211583751558920064\\
24&
23796267139390607360&
471052121438170656227510842657968048509244598133358450149784571617082575423565056\\
25&
36695884124388970496&
1752662358687885188987478205394018051759346782906849304190277599933097187266142260992\\
26&
55739965040872656896&
6955857126069033829258591712286204229109865339121174753159805133951078138602493998424320\\
27&
83483175402122792960&
29386673295598252119315606146972896461970012686091855147040625434158518678060782984083104000\\
28&
123399400507042336768&
131909085288222166096090487444534010029062356675207583449019800324247261789364850799927165421056\\
29&
180164856350448209920&
627984678371490556154911685166126439591413733396445863377606242170773209583438050710786209392034304\\
30&
260014766495085432832&
3165506839235673119694504216537225971338421328159148455404020879282993502672554389276865198890883013120 .
\end{array}
\]
This proves the claim.
\end{proof}

\begin{proposition}[Residual $C$ twelfth nonboundary lower correction]\label{prop:bc-c-twelfth-nonboundary-lower-correction}
For every residual row \(C_b\) with \(20\le b\le217\), put \(q=b+1\) and
write \(s_j=m_j^{\mathrm{st}}\) and \(m_j^{C_b}=s_j+\delta_j^{C_b}\).  Then
\[
  -\frac12 s_{2q+12}\le \delta_{2q+12}^{C_b}.
\]
\end{proposition}

\begin{proof}
As in the proof of
\cref{prop:bc-c-first-nonboundary-lower-correction},
\(|\delta_{2q+12}^{C_b}|\) is the number of column-even Pieri two-strip paths
of length \(2q+12\) with terminal height at least \(2q\).  By
\cref{lem:bc-c-twelfth-nonboundary-twelve-pause-code}, with \(P(q)\) denoting
the polynomial from
\cref{lem:bc-c-twelfth-frontier-pause-arithmetic-certificate},
\[
  |\delta_{2q+12}^{C_b}|
  \le
  P(q)\,2^{2q}s_{q+12}.
\]
If \(21\le q\le30\),
\cref{lem:bc-c-twelfth-frontier-pause-arithmetic-certificate} gives
\[
  |\delta_{2q+12}^{C_b}|\le s_{2q+12}/2 .
\]
If \(q\ge31\), \cref{lem:bc-twelfth-c-pause-binomial-absorption} gives
\[
  |\delta_{2q+12}^{C_b}|
  \le
  2^{2q}\binom{2q+12}{q}s_{q+12}.
\]
The exact GMP/OpenMP replay
\[
\texttt{certificates/post\_m29/bc\_pieri\_powerloss\_2q\_gmp\_certificate.log}
\]
checks the inequality
\[
  2^{2q+1}\binom{2q+12}{q}s_{q+12}\le s_{2q+12}
\]
on this residual row and moment as part of its all-row residual-window audit;
it reports \texttt{failures=0} and has SHA-256
\[
\hashsplit{124068c5daa1123088a3cd79c3f6b839}{fee45a1aaef637ed82161c465d86bf45}.
\]
Thus \(|\delta_{2q+12}^{C_b}|\le s_{2q+12}/2\) also for \(q\ge31\), which
implies the claimed lower bound.
\end{proof}

\begin{lemma}[$C$ thirteenth-nonboundary height-bad paths have a thirteen-pause code]\label{lem:bc-c-thirteenth-nonboundary-thirteen-pause-code}
For \(q\ge1\), the number of column-even Pieri two-strip paths
\[
  \varnothing=\lambda^{(0)}\subset\lambda^{(1)}\subset\cdots
  \subset\lambda^{(2q+13)}
\]
with
\[
  \ell(\lambda^{(2q+13)})\ge2q
\]
is at most
\[
  \biggl(16384\binom{2q+13}{1}+65536\binom{2q+13}{3}
    +262144\binom{2q+13}{5}+1048576\binom{2q+13}{7}
    +4194304\binom{2q+13}{9}+16777216\binom{2q+13}{11}
    +67108864\binom{2q+13}{13}
  \biggr)2^{2q}s_{q+13}.
\]
\end{lemma}

\begin{proof}
The terminal height is even and at most \(2q+13\).  Thus it is
\[
  2q+12,\quad 2q+10,\quad 2q+8,\quad 2q+6,\quad
  2q+4,\quad 2q+2,\quad\text{or}\quad 2q .
\]
These cases have respectively \(1,3,5,7,9,11,13\) steps which fail to create a
new row.  Deleting the first column, shifting left, and standardizing the
repeated pause labels injects the paths with \(p\) pauses into the choice of
the \(p\)-element pause set and an even-column standard tableau of size
\(2q+13+p\).  Therefore the seven cases contribute at most
\[
  \binom{2q+13}{1}2^{2q+14}s_{q+7},
  \quad
  \binom{2q+13}{3}2^{2q+16}s_{q+8},
  \quad
  \binom{2q+13}{5}2^{2q+18}s_{q+9},
\]
\[
  \binom{2q+13}{7}2^{2q+20}s_{q+10},
  \quad
  \binom{2q+13}{9}2^{2q+22}s_{q+11},
  \quad
  \binom{2q+13}{11}2^{2q+24}s_{q+12},
\]
\[
  \binom{2q+13}{13}2^{2q+26}s_{q+13},
\]
by
\cref{lem:bc-even-column-standard-tableaux-2n-bound,lem:bc-stable-moments-monotone-after-two}.
After replacing \(s_{q+7},\ldots,s_{q+12}\) by \(s_{q+13}\), the sum is the
displayed bound.
\end{proof}

\begin{lemma}[The central binomial factor absorbs the thirteen-pause loss]\label{lem:bc-thirteenth-c-pause-binomial-absorption}
For every \(q\ge33\),
\[
  16384\binom{2q+13}{1}+65536\binom{2q+13}{3}
    +262144\binom{2q+13}{5}+1048576\binom{2q+13}{7}
    +4194304\binom{2q+13}{9}+16777216\binom{2q+13}{11}
    +67108864\binom{2q+13}{13}
  \le \binom{2q+13}{q}.
\]
\end{lemma}

\begin{proof}
At \(q=33\),
\[
  16384\binom{79}{1}+65536\binom{79}{3}
    +262144\binom{79}{5}+1048576\binom{79}{7}
    +4194304\binom{79}{9}+16777216\binom{79}{11}
    +67108864\binom{79}{13}
  =
  17864298942435109027840
  <
  18723298265888530356770
  =
  \binom{79}{33}.
\]
Let
\[
  P(q)=16384\binom{2q+13}{1}+65536\binom{2q+13}{3}
    +262144\binom{2q+13}{5}+1048576\binom{2q+13}{7}
    +4194304\binom{2q+13}{9}+16777216\binom{2q+13}{11}
    +67108864\binom{2q+13}{13}.
\]
For \(q\ge33\), \(P(q+1)<2P(q)\).  Indeed, with \(x=q-33\),
\[
\begin{aligned}
6081075(2P(q)-P(q+1))
={}&536870912x^{13}+247765925888x^{12}+52683411030016x^{11}\\
&+6833111170023424x^{10}+603052064423018496x^9\\
&+38235592030949474304x^8+1791431987680749027328x^7\\
&+62782980715340357435392x^6+1645392427080320194445312x^5\\
&+31837793611414160915038208x^4+441932813329921918349869056x^3\\
&+4165225135993797091904323584x^2\\
&+23877555284696150243189882880x+62833122000287032433950310400,
\end{aligned}
\]
which is positive for \(x\ge0\).  On the other hand,
\[
  \frac{\binom{2q+15}{q+1}}{\binom{2q+13}{q}}
  =
  \frac{(2q+15)(2q+14)}{(q+1)(q+14)}
  >3,
\]
since the numerator minus three times the denominator is \(q^2+13q+168\).
Thus the binomial side grows by a larger factor than \(P(q)\), and induction
from \(q=33\) proves the inequality.
\end{proof}

\begin{lemma}[Exact first twelve thirteenth $C$ pause-arithmetic checks]\label{lem:bc-c-thirteenth-frontier-pause-arithmetic-certificate}
For \(21\le q\le32\), put
\[
  P(q)=16384\binom{2q+13}{1}+65536\binom{2q+13}{3}
    +262144\binom{2q+13}{5}+1048576\binom{2q+13}{7}
    +4194304\binom{2q+13}{9}+16777216\binom{2q+13}{11}
    +67108864\binom{2q+13}{13}.
\]
Then
\[
  2^{2q+1}P(q)s_{q+13}\le s_{2q+13}.
\]
\end{lemma}

\begin{proof}
The exact C++/GMP/OpenMP replay
\[
\texttt{certificates/post\_m29/bc\_thirteenth\_frontier\_gmp\_certificate.log}
\]
computes the displayed integer inequality for
\(21\le q\le32\).  It exits with status \(0\),
reports \texttt{SUMMARY failures=0}, and has SHA-256
\[
\hashsplit{6f6ffcfb724ca85fa51310916cd2451e}{4b38d08bcf0dd224719ce3d8b809e3e9}.
\]
It verifies the exact positive margins
\[
\begin{array}{c|c|c}
q & P(q) & s_{2q+13}-2^{2q+1}P(q)s_{q+13}\\ \hline
21&
99421578468163108864&
731183253555922246274145306849776724283935343412240521338981033999058560\\
22&
167435930009130385408&
2349086344497982960380567450134420546006293822947033888650532606552947540480\\
23&
276363361539293691904&
7915807283761492701741559988361827399990308007320009827870256865677681228151168\\
24&
447739369047974821888&
28498374779999066839524516000230618733300521031838837186021836373553965642135948032\\
25&
712955835787213717504&
109541166127051607981301979498068684490946462361460615532644080461796654207629385278720\\
26&
1117151165606205079552&
448652113835933266852742777764745511772674732526011424622707823204933417946550301817528320\\
27&
1724410263482076872704&
1954211161170261374923664714989792376243604267813600314045649818913786405050060276086835708160\\
28&
2624645313339026522112&
9035761316222856032195612149457571806209086085064215683959163378490951478285117127173212392044032\\
29&
3942614982711209148416&
44272870337027816746858558646361549656504231610775669646028860443921659475011087072106629992943662592\\
30&
5849642632546285928448&
229499010251281522668378046504096718321829854693072240153475734580282679151017260965635067349105902389248\\
31&
8578715686145537064960&
1256676780468853148224502722041231129903054470193866885549402798554126364708883036010445696292421349143986688\\
32&
12443791119168545046528&
7258238000562607250224480367208885061023620028393530822223165594504067877769292050894174201678791838370956768256 .
\end{array}
\]
This proves the claim.
\end{proof}

\begin{proposition}[Residual $C$ thirteenth nonboundary lower correction]\label{prop:bc-c-thirteenth-nonboundary-lower-correction}
For every residual row \(C_b\) with \(20\le b\le217\), put \(q=b+1\) and
write \(s_j=m_j^{\mathrm{st}}\) and \(m_j^{C_b}=s_j+\delta_j^{C_b}\).  Then
\[
  -\frac12 s_{2q+13}\le \delta_{2q+13}^{C_b}.
\]
\end{proposition}

\begin{proof}
As in the proof of
\cref{prop:bc-c-first-nonboundary-lower-correction},
\(|\delta_{2q+13}^{C_b}|\) is the number of column-even Pieri two-strip paths
of length \(2q+13\) with terminal height at least \(2q\).  By
\cref{lem:bc-c-thirteenth-nonboundary-thirteen-pause-code}, with \(P(q)\)
denoting the polynomial from
\cref{lem:bc-c-thirteenth-frontier-pause-arithmetic-certificate},
\[
  |\delta_{2q+13}^{C_b}|
  \le
  P(q)\,2^{2q}s_{q+13}.
\]
If \(21\le q\le32\),
\cref{lem:bc-c-thirteenth-frontier-pause-arithmetic-certificate} gives
\[
  |\delta_{2q+13}^{C_b}|\le s_{2q+13}/2 .
\]
If \(q\ge33\), \cref{lem:bc-thirteenth-c-pause-binomial-absorption} gives
\[
  |\delta_{2q+13}^{C_b}|
  \le
  2^{2q}\binom{2q+13}{q}s_{q+13}.
\]
The exact GMP/OpenMP replay
\[
\texttt{certificates/post\_m29/bc\_pieri\_powerloss\_2q\_gmp\_certificate.log}
\]
checks the inequality
\[
  2^{2q+1}\binom{2q+13}{q}s_{q+13}\le s_{2q+13}
\]
on this residual row and moment as part of its all-row residual-window audit;
it reports \texttt{failures=0} and has SHA-256
\[
\hashsplit{124068c5daa1123088a3cd79c3f6b839}{fee45a1aaef637ed82161c465d86bf45}.
\]
Thus \(|\delta_{2q+13}^{C_b}|\le s_{2q+13}/2\) also for \(q\ge33\), which
implies the claimed lower bound.
\end{proof}

\begin{lemma}[$C$ fourteenth-nonboundary height-bad paths have a fourteen-pause code]\label{lem:bc-c-fourteenth-nonboundary-fourteen-pause-code}
For \(q\ge1\), the number of column-even Pieri two-strip paths
\[
  \varnothing=\lambda^{(0)}\subset\lambda^{(1)}\subset\cdots
  \subset\lambda^{(2q+14)}
\]
with
\[
  \ell(\lambda^{(2q+14)})\ge2q
\]
is at most
\[
\begin{aligned}
  \biggl(&16384\binom{2q+14}{0}+65536\binom{2q+14}{2}
    +262144\binom{2q+14}{4}+1048576\binom{2q+14}{6}\\
  &\quad+4194304\binom{2q+14}{8}+16777216\binom{2q+14}{10}
    +67108864\binom{2q+14}{12}
    +268435456\binom{2q+14}{14}\biggr)2^{2q}s_{q+14}.
\end{aligned}
\]
\end{lemma}

\begin{proof}
The terminal height is even and at most \(2q+14\).  Thus it is
\[
  2q+14,\quad 2q+12,\quad 2q+10,\quad 2q+8,\quad
  2q+6,\quad 2q+4,\quad 2q+2,\quad\text{or}\quad 2q .
\]
These cases have respectively \(0,2,4,6,8,10,12,14\) steps which fail to
create a new row.  Deleting the first column, shifting left, and standardizing
the repeated pause labels injects the paths with \(p\) pauses into the choice
of the \(p\)-element pause set and an even-column standard tableau of size
\(2q+14+p\).  Therefore the eight cases contribute at most
\[
  \binom{2q+14}{0}2^{2q+14}s_{q+7},
  \quad
  \binom{2q+14}{2}2^{2q+16}s_{q+8},
  \quad
  \binom{2q+14}{4}2^{2q+18}s_{q+9},
\]
\[
  \binom{2q+14}{6}2^{2q+20}s_{q+10},
  \quad
  \binom{2q+14}{8}2^{2q+22}s_{q+11},
  \quad
  \binom{2q+14}{10}2^{2q+24}s_{q+12},
\]
\[
  \binom{2q+14}{12}2^{2q+26}s_{q+13},
  \quad
  \binom{2q+14}{14}2^{2q+28}s_{q+14},
\]
by
\cref{lem:bc-even-column-standard-tableaux-2n-bound,lem:bc-stable-moments-monotone-after-two}.
After replacing \(s_{q+7},\ldots,s_{q+13}\) by \(s_{q+14}\), the sum is the
displayed bound.
\end{proof}

\begin{lemma}[The central binomial factor absorbs the fourteen-pause loss]\label{lem:bc-fourteenth-c-pause-binomial-absorption}
For every \(q\ge36\),
\[
\begin{aligned}
  &16384\binom{2q+14}{0}+65536\binom{2q+14}{2}
    +262144\binom{2q+14}{4}+1048576\binom{2q+14}{6}\\
  &\quad+4194304\binom{2q+14}{8}+16777216\binom{2q+14}{10}
    +67108864\binom{2q+14}{12}
    +268435456\binom{2q+14}{14}
  \le \binom{2q+14}{q}.
\end{aligned}
\]
\end{lemma}

\begin{proof}
At \(q=36\), the left side is
\[
  1228572339217567202885632
  <
  2141368094445751179086598
  =
  \binom{86}{36}.
\]
Let \(P(q)\) denote the displayed left side.  For \(0\le p\le14\),
\[
  \frac{\binom{2q+16}{p}}{\binom{2q+14}{p}}
  \le
  \frac{(2q+16)(2q+15)}{(2q+2)(2q+1)}
  <2
  \qquad(q\ge36),
\]
because \(4q^2-50q-236>0\).  Hence \(P(q+1)<2P(q)\).  On the other hand,
\[
  \frac{\binom{2q+16}{q+1}}{\binom{2q+14}{q}}
  =
  \frac{(2q+16)(2q+15)}{(q+1)(q+15)}
  >3,
\]
since the numerator minus three times the denominator is \(q^2+14q+195\).
Thus the binomial side grows by a larger factor than \(P(q)\), and induction
from \(q=36\) proves the inequality.
\end{proof}

\begin{lemma}[Exact first fifteen fourteenth $C$ pause-arithmetic checks]\label{lem:bc-c-fourteenth-frontier-pause-arithmetic-certificate}
For \(21\le q\le35\), put
\[
\begin{aligned}
  P(q)={}&16384\binom{2q+14}{0}+65536\binom{2q+14}{2}
    +262144\binom{2q+14}{4}+1048576\binom{2q+14}{6}\\
  &+4194304\binom{2q+14}{8}+16777216\binom{2q+14}{10}
    +67108864\binom{2q+14}{12}
    +268435456\binom{2q+14}{14}.
\end{aligned}
\]
Then
\[
  2^{2q+1}P(q)s_{q+14}\le s_{2q+14}.
\]
\end{lemma}

\begin{proof}
The exact C++/GMP/OpenMP replay
\[
\texttt{certificates/post\_m29/bc\_fourteenth\_frontier\_gmp\_certificate.log}
\]
computes the displayed integer inequality for
\(21\le q\le35\).  It exits with status
\(0\), reports \texttt{SUMMARY failures=0}, and has SHA-256
\[
\hashsplit{315dcc3f51bd178dba5b186d4cb4f3e0}{0c636284f3b43d8428915d46a90af82b}.
\]
It verifies the exact positive margins
\[
\begin{array}{c|c|c}
q & P(q) & s_{2q+14}-2^{2q+1}P(q)s_{q+14}\\ \hline
21&
1596271720409684787200&
30917529858094002525675633417861850565638275718143447666636062145035675520\\
22&
2783455765485702692864&
132672204573136145950500640822645551216920842036009669542684930183208960915328\\
23&
4751439282346409738240&
470389783336939185412775373842016791528832025449402724385708691553062943501348096\\
24&
7952613960499958595584&
1752496064536743822300264310789194720241001029061494820038965372279215802546612195072\\
25&
13069126500807444807680&
6955815038621903173192248300411331900227468094245033781979107973384315273703513778708736\\
26&
21114496385003122933760&
29386662551758018082154847284853106262575512858314384278735038682345181630868645711657918720\\
27&
33574047686425773228032&
131909082520572441945981645075716186740477689687535040935905402509018503590021353816558633190912\\
28&
52596819785025308016640&
627984677651732044024185506034241629251735008508738949711974505481755667181203759271981181971137024\\
29&
81255101945107809845248&
3165506839046640308276374664830570485603569956580991759627801929422266577900904929727336361298654190080\\
30&
123892006741556730347520&
16868160644315292137234793151061476357142219647945787848677390455912307942796862950762802857240886238943744\\
31&
186582701460026967343104&
94879010830031532349875873741874771299670970548527039813930028602579942534927063585645793842574676669159095296\\
32&
277741219368227086221312&
562512973617578949738715051413602370438296814822790768369280613204437717984504076997473538041680135009142325285888\\
33&
408912361098124892585984&
3510497287734777988352154395214204024810663544838658577411248004504946986095211218898758792608835147398481712965860352\\
34&
595797281870381039828992&
23031461318491612561361483118389880070206975243636714799442705876023076415533799740908938692502777310236898914266080918528\\
35&
859572181508965409112064&
158657763193874483326535139324059261929095073029658493666643098115540993728569364080099703614686993911702096986997787803625472 .
\end{array}
\]
This proves the claim.
\end{proof}

\begin{proposition}[Residual $C$ fourteenth nonboundary lower correction]\label{prop:bc-c-fourteenth-nonboundary-lower-correction}
For every residual row \(C_b\) with \(20\le b\le217\), put \(q=b+1\) and
write \(s_j=m_j^{\mathrm{st}}\) and \(m_j^{C_b}=s_j+\delta_j^{C_b}\).  Then
\[
  -\frac12 s_{2q+14}\le \delta_{2q+14}^{C_b}.
\]
\end{proposition}

\begin{proof}
As in the proof of
\cref{prop:bc-c-first-nonboundary-lower-correction},
\(|\delta_{2q+14}^{C_b}|\) is the number of column-even Pieri two-strip paths
of length \(2q+14\) with terminal height at least \(2q\).  By
\cref{lem:bc-c-fourteenth-nonboundary-fourteen-pause-code}, with \(P(q)\)
denoting the polynomial from
\cref{lem:bc-c-fourteenth-frontier-pause-arithmetic-certificate},
\[
  |\delta_{2q+14}^{C_b}|
  \le
  P(q)\,2^{2q}s_{q+14}.
\]
If \(21\le q\le35\),
\cref{lem:bc-c-fourteenth-frontier-pause-arithmetic-certificate} gives
\[
  |\delta_{2q+14}^{C_b}|\le s_{2q+14}/2 .
\]
If \(q\ge36\), \cref{lem:bc-fourteenth-c-pause-binomial-absorption} gives
\[
  |\delta_{2q+14}^{C_b}|
  \le
  2^{2q}\binom{2q+14}{q}s_{q+14}.
\]
The exact GMP/OpenMP replay
\[
\texttt{certificates/post\_m29/bc\_pieri\_powerloss\_2q\_gmp\_certificate.log}
\]
checks the inequality
\[
  2^{2q+1}\binom{2q+14}{q}s_{q+14}\le s_{2q+14}
\]
on this residual row and moment as part of its all-row residual-window audit;
it reports \texttt{failures=0} and has SHA-256
\[
\hashsplit{124068c5daa1123088a3cd79c3f6b839}{fee45a1aaef637ed82161c465d86bf45}.
\]
Thus \(|\delta_{2q+14}^{C_b}|\le s_{2q+14}/2\) also for \(q\ge36\), which
implies the claimed lower bound.
\end{proof}

\begin{lemma}[Even-column standard tableaux have a fixed-point-free involution count]\label{lem:bc-even-column-standard-tableaux-fpf-count}
For every \(n\ge1\), the number of standard Young tableaux of size \(2n\)
whose columns all have even length is
\[
  (2n-1)!!=\frac{(2n)!}{2^n n!}.
\]
\end{lemma}

\begin{proof}
Under the Robinson--Schensted correspondence, involutions on
\(\{1,\ldots,2n\}\) are in bijection with standard Young tableaux of size
\(2n\).  Rains \cite[Lemma 3.3]{Rains1997} records the Schensted fact that
the number of fixed points of the involution is the number of odd columns of
the tableau.  Thus the tableaux with only even columns are exactly the
fixed-point-free involutions.  These are perfect matchings of \(2n\) labelled
points, counted by
\[
  (2n-1)(2n-3)\cdots3\cdot1=\frac{(2n)!}{2^n n!}.
\]
\end{proof}

\begin{lemma}[$C$ fifteenth-nonboundary height-bad paths have a fifteen-pause code]\label{lem:bc-c-fifteenth-nonboundary-fifteen-pause-code}
For \(q\ge1\), the number of column-even Pieri two-strip paths
\[
  \varnothing=\lambda^{(0)}\subset\lambda^{(1)}\subset\cdots
  \subset\lambda^{(2q+15)}
\]
with
\[
  \ell(\lambda^{(2q+15)})\ge2q
\]
is at most
\[
\begin{aligned}
  \biggl(&65536\binom{2q+15}{1}+262144\binom{2q+15}{3}
    +1048576\binom{2q+15}{5}+4194304\binom{2q+15}{7}\\
  &\quad+16777216\binom{2q+15}{9}+67108864\binom{2q+15}{11}
    +268435456\binom{2q+15}{13}
    +1073741824\binom{2q+15}{15}\biggr)2^{2q}s_{q+15}.
\end{aligned}
\]
\end{lemma}

\begin{proof}
The terminal height is even and at most \(2q+15\).  Thus it is
\[
  2q+14,\quad 2q+12,\quad 2q+10,\quad 2q+8,
  \quad 2q+6,\quad 2q+4,\quad 2q+2,\quad\text{or}\quad 2q .
\]
These cases have respectively \(1,3,5,7,9,11,13,15\) steps which fail to
create a new row.  Deleting the first column, shifting left, and
standardizing the repeated pause labels injects the paths with \(p\) pauses
into the choice of the \(p\)-element pause set and an even-column standard
tableau of size \(2q+15+p\).  Therefore the eight cases contribute at most
\[
  \binom{2q+15}{1}2^{2q+16}s_{q+8},
  \quad
  \binom{2q+15}{3}2^{2q+18}s_{q+9},
  \quad
  \binom{2q+15}{5}2^{2q+20}s_{q+10},
\]
\[
  \binom{2q+15}{7}2^{2q+22}s_{q+11},
  \quad
  \binom{2q+15}{9}2^{2q+24}s_{q+12},
  \quad
  \binom{2q+15}{11}2^{2q+26}s_{q+13},
\]
\[
  \binom{2q+15}{13}2^{2q+28}s_{q+14},
  \quad
  \binom{2q+15}{15}2^{2q+30}s_{q+15},
\]
by
\cref{lem:bc-even-column-standard-tableaux-2n-bound,lem:bc-stable-moments-monotone-after-two}.
After replacing \(s_{q+8},\ldots,s_{q+14}\) by \(s_{q+15}\), the sum is the
displayed bound.
\end{proof}

\begin{lemma}[The first fifteenth $C$ pause case is absorbed by the fixed-point-free count]\label{lem:bc-fifteenth-c-q21-fpf-pause-certificate}
For \(q=21\), the number of column-even Pieri two-strip paths of length
\(2q+15\) with terminal height at least \(2q\) is at most
\[
  52920726139837306981793811162459494689272852275711217409665468750
  \le \frac12 s_{57}.
\]
\end{lemma}

\begin{proof}
Use the pause-set injection from
\cref{lem:bc-c-fifteenth-nonboundary-fifteen-pause-code}, but count the
target even-column standard tableaux exactly by
\cref{lem:bc-even-column-standard-tableaux-fpf-count}.  For \(q=21\), this
gives the upper bound
\[
  \sum_{\substack{1\le p\le15\\ p\ {\rm odd}}}
  \binom{57}{p}(56+p)!!
  =
  52920726139837306981793811162459494689272852275711217409665468750 .
\]
The exact C++/GMP/OpenMP replay
\[
\texttt{certificates/post\_m29/bc\_fifteenth\_frontier\_gmp\_certificate.log}
\]
checks the displayed integer sum and verifies
\[
\begin{aligned}
  &s_{57}
  -2\cdot
  52920726139837306981793811162459494689272852275711217409665468750\\
  &\qquad =
  2353592104444744758136209099549217220785171739692400245417399049518066806628
  >0 .
\end{aligned}
\]
The replay exits with status \(0\), reports \texttt{SUMMARY failures=0}, and
has SHA-256
\[
\hashsplit{5f829fb541bf8d84f097ff14aa6be272}{508554407503bd74b37532e62c1484b9}.
\]
\end{proof}

\begin{lemma}[The central binomial factor absorbs the fifteen-pause loss]\label{lem:bc-fifteenth-c-pause-binomial-absorption}
For every \(q\ge38\),
\[
\begin{aligned}
  &65536\binom{2q+15}{1}+262144\binom{2q+15}{3}
    +1048576\binom{2q+15}{5}+4194304\binom{2q+15}{7}\\
  &\quad+16777216\binom{2q+15}{9}+67108864\binom{2q+15}{11}
    +268435456\binom{2q+15}{13}
    +1073741824\binom{2q+15}{15}
  \le \binom{2q+15}{q}.
\end{aligned}
\]
\end{lemma}

\begin{proof}
At \(q=38\), the left side is
\[
  59395730540209942865838080
  <
  60468953239450827809517780
  =
  \binom{91}{38}.
\]
Let \(P(q)\) denote the displayed left side.  For \(1\le p\le15\),
\[
  \frac{\binom{2q+17}{p}}{\binom{2q+15}{p}}
  \le
  \frac{(2q+17)(2q+16)}{(2q+2)(2q+1)}
  <2
  \qquad(q\ge38),
\]
because \(4q^2-54q-268>0\).  Hence \(P(q+1)<2P(q)\).  On the other hand,
\[
  \frac{\binom{2q+17}{q+1}}{\binom{2q+15}{q}}
  =
  \frac{(2q+17)(2q+16)}{(q+1)(q+16)}
  >3,
\]
since the numerator minus three times the denominator is \(q^2+15q+224\).
Thus the binomial side grows by a larger factor than \(P(q)\), and induction
from \(q=38\) proves the inequality.
\end{proof}

\begin{lemma}[Exact fifteenth $C$ pause-arithmetic checks after the first row]\label{lem:bc-c-fifteenth-frontier-pause-arithmetic-certificate}
For \(22\le q\le37\), put
\[
\begin{aligned}
  P(q)={}&65536\binom{2q+15}{1}+262144\binom{2q+15}{3}
    +1048576\binom{2q+15}{5}+4194304\binom{2q+15}{7}\\
  &+16777216\binom{2q+15}{9}+67108864\binom{2q+15}{11}
    +268435456\binom{2q+15}{13}
    +1073741824\binom{2q+15}{15}.
\end{aligned}
\]
Then
\[
  2^{2q+1}P(q)s_{q+15}\le s_{2q+15}.
\]
\end{lemma}

\begin{proof}
The exact C++/GMP/OpenMP replay
\[
\texttt{certificates/post\_m29/bc\_fifteenth\_frontier\_gmp\_certificate.log}
\]
computes the displayed integer inequality for
\(22\le q\le37\).  It exits with status
\(0\), reports \texttt{SUMMARY failures=0}, and has SHA-256
\[
\hashsplit{5f829fb541bf8d84f097ff14aa6be272}{508554407503bd74b37532e62c1484b9}.
\]
It verifies the exact positive margins
\[
\begin{array}{c|c|c}
q & P(q) & s_{2q+15}-2^{2q+1}P(q)s_{q+15}\\ \hline
22&
43942992213014570008576&
6384638072633941694505837874695804644380836826145144184602242817435341577156992\\
23&
77532692687157148844032&
28093110657666714909057277466061899267359556118402886615164731950107952805050108672\\
24&
133989774466389220458496&
109433303742710598263276391321287752551411989136851033518696417662108760558860099134720\\
25&
227135493100253359439872&
448623222915029924737440591476055129159320823405207288109005870309482339727342404079149056\\
26&
378177045530579061637120&
1954203368170796853773594264062320589909216147679262173635093966623002616542398755564634177792\\
27&
619178862806272017498112&
9035759198108574868114204314427314125944410433316935740711085947218950678406318049415253887202816\\
28&
997959096073050057277440&
44272869756652934516466706008181558607533958678344603437118091479604466403261918234480788648856860160\\
29&
1584918071910533189664768&
229499010090896465645712757794233230267780943277870916185560831482062042822298289774311766288716060723200\\
30&
2482459843722246457589760&
1256676780424136439434570598877690300477266610257380710030351114922366397990844287348484853177097779871728128\\
31&
3837862004424133603295232&
7258238000550025020641273629755535532744340586317480271914940646430973318818221399020564801617450002227297973248\\
32&
5860691101463253572321280&
44157232836689126022129312225382056021822413875402709397622545669276564133063527278884565304776200276824534556351488\\
33&
8846161207169290968563712&
282594820649085148751501837552075377224088586036858394452218557433550485230437025354535128980184244547955074857087049728\\
34&
13206202986736261774376960&
1900094241916969821015326783899098957068686376895632009991267385626884091921948782430230615100558163036290402493373876509696\\
35&
19510463367286790841761792&
13406572350450220246623533421950480016032212094880317356032197960348111543033308580701452562064817221218366308625135077652318208\\
36&
28540007123376295459880960&
99151536294250092631965221915300023892901248805730607655964371666677718184985273915999787809330644818264446243230358332146284193792\\
37&
41357159172102519624630272&
767803834154796949332039766767729572150243961827925415210051698803158660857002361916356753846249297137223116592514977515242497715314688 .
\end{array}
\]
This proves the claim.
\end{proof}

\begin{proposition}[Residual $C$ fifteenth nonboundary lower correction]\label{prop:bc-c-fifteenth-nonboundary-lower-correction}
For every residual row \(C_b\) with \(20\le b\le217\), put \(q=b+1\) and
write \(s_j=m_j^{\mathrm{st}}\) and \(m_j^{C_b}=s_j+\delta_j^{C_b}\).  Then
\[
  -\frac12 s_{2q+15}\le \delta_{2q+15}^{C_b}.
\]
\end{proposition}

\begin{proof}
As in the proof of
\cref{prop:bc-c-first-nonboundary-lower-correction},
\(|\delta_{2q+15}^{C_b}|\) is the number of column-even Pieri two-strip paths
of length \(2q+15\) with terminal height at least \(2q\).  If \(q=21\),
\cref{lem:bc-fifteenth-c-q21-fpf-pause-certificate} gives
\[
  |\delta_{2q+15}^{C_b}|\le s_{2q+15}/2 .
\]
For \(q\ge22\), \cref{lem:bc-c-fifteenth-nonboundary-fifteen-pause-code},
with \(P(q)\) denoting the polynomial from
\cref{lem:bc-c-fifteenth-frontier-pause-arithmetic-certificate}, gives
\[
  |\delta_{2q+15}^{C_b}|
  \le
  P(q)\,2^{2q}s_{q+15}.
\]
If \(22\le q\le37\),
\cref{lem:bc-c-fifteenth-frontier-pause-arithmetic-certificate} gives
\[
  |\delta_{2q+15}^{C_b}|\le s_{2q+15}/2 .
\]
If \(q\ge38\), \cref{lem:bc-fifteenth-c-pause-binomial-absorption} gives
\[
  |\delta_{2q+15}^{C_b}|
  \le
  2^{2q}\binom{2q+15}{q}s_{q+15}.
\]
The exact GMP/OpenMP replay
\[
\texttt{certificates/post\_m29/bc\_pieri\_powerloss\_2q\_gmp\_certificate.log}
\]
checks the inequality
\[
  2^{2q+1}\binom{2q+15}{q}s_{q+15}\le s_{2q+15}
\]
on this residual row and moment as part of its all-row residual-window audit;
it reports \texttt{failures=0} and has SHA-256
\[
\hashsplit{124068c5daa1123088a3cd79c3f6b839}{fee45a1aaef637ed82161c465d86bf45}.
\]
Thus \(|\delta_{2q+15}^{C_b}|\le s_{2q+15}/2\) also for \(q\ge38\), which
implies the claimed lower bound.
\end{proof}

\begin{lemma}[$C$ sixteenth-nonboundary height-bad paths have a sixteen-pause code]\label{lem:bc-c-sixteenth-nonboundary-sixteen-pause-code}
For \(q\ge1\), the number of column-even Pieri two-strip paths
\[
  \varnothing=\lambda^{(0)}\subset\lambda^{(1)}\subset\cdots
  \subset\lambda^{(2q+16)}
\]
with
\[
  \ell(\lambda^{(2q+16)})\ge2q
\]
is at most
\[
\begin{aligned}
  \biggl(&65536\binom{2q+16}{0}+262144\binom{2q+16}{2}
    +1048576\binom{2q+16}{4}+4194304\binom{2q+16}{6}\\
  &\quad+16777216\binom{2q+16}{8}+67108864\binom{2q+16}{10}
    +268435456\binom{2q+16}{12}\\
  &\quad+1073741824\binom{2q+16}{14}
    +4294967296\binom{2q+16}{16}\biggr)2^{2q}s_{q+16}.
\end{aligned}
\]
\end{lemma}

\begin{proof}
The terminal height is even and at most \(2q+16\).  Thus it is
\[
  2q+16,\quad 2q+14,\quad 2q+12,\quad 2q+10,\quad 2q+8,
  \quad 2q+6,\quad 2q+4,\quad 2q+2,\quad\text{or}\quad 2q .
\]
These cases have respectively \(0,2,4,6,8,10,12,14,16\) steps which fail to
create a new row.  Deleting the first column, shifting left, and
standardizing the repeated pause labels injects the paths with \(p\) pauses
into the choice of the \(p\)-element pause set and an even-column standard
tableau of size \(2q+16+p\).  Therefore the nine cases contribute at most
\[
  \binom{2q+16}{0}2^{2q+16}s_{q+8},
  \quad
  \binom{2q+16}{2}2^{2q+18}s_{q+9},
  \quad
  \binom{2q+16}{4}2^{2q+20}s_{q+10},
\]
\[
  \binom{2q+16}{6}2^{2q+22}s_{q+11},
  \quad
  \binom{2q+16}{8}2^{2q+24}s_{q+12},
  \quad
  \binom{2q+16}{10}2^{2q+26}s_{q+13},
\]
\[
  \binom{2q+16}{12}2^{2q+28}s_{q+14},
  \quad
  \binom{2q+16}{14}2^{2q+30}s_{q+15},
  \quad
  \binom{2q+16}{16}2^{2q+32}s_{q+16},
\]
by
\cref{lem:bc-even-column-standard-tableaux-2n-bound,lem:bc-stable-moments-monotone-after-two}.
After replacing \(s_{q+8},\ldots,s_{q+15}\) by \(s_{q+16}\), the sum is the
displayed bound.
\end{proof}

\begin{lemma}[The first two sixteenth $C$ pause cases are absorbed by the fixed-point-free count]\label{lem:bc-sixteenth-c-q21-q22-fpf-pause-certificate}
For \(q=21,22\), the number of column-even Pieri two-strip paths of length
\(2q+16\) with terminal height at least \(2q\) is at most \(s_{2q+16}/2\).
\end{lemma}

\begin{proof}
Use the pause-set injection from
\cref{lem:bc-c-sixteenth-nonboundary-sixteen-pause-code}, but count the target
even-column standard tableaux exactly by
\cref{lem:bc-even-column-standard-tableaux-fpf-count}.  For \(q=21,22\), this
gives the upper bound
\[
  \sum_{\substack{0\le p\le16\\ p\ {\rm even}}}
  \binom{2q+16}{p}(2q+15+p)!! .
\]
The exact C++/GMP/OpenMP replay
\[
\texttt{certificates/post\_m29/bc\_sixteenth\_frontier\_gmp\_certificate.log}
\]
checks the displayed integer sums and verifies
\[
\begin{array}{c|c|c|c}
q & j & \mathrm{FPF}(q,j) & s_j-2\mathrm{FPF}(q,j)\\ \hline
21&58&
14006591587695753581293514913909039719080789693690052832937036640625&
135331264326267413349868887869322821868469953370931674269172602356060215383198\\
22&60&
1965135565412826288161569628583208390623978721644886751970133067109375&
471052485158774046188987143394085149976302513690976318522627253040049506175567106 .
\end{array}
\]
The replay exits with status \(0\), reports \texttt{SUMMARY failures=0}, and
has SHA-256
\[
\hashsplit{9276304d5e5aff2105d4267604664e9f}{14dbcc42df60e0a643c335b028250009}.
\]
\end{proof}

\begin{lemma}[The central binomial factor absorbs the sixteen-pause loss]\label{lem:bc-sixteenth-c-pause-binomial-absorption}
For every \(q\ge41\),
\[
\begin{aligned}
  &65536\binom{2q+16}{0}+262144\binom{2q+16}{2}
    +1048576\binom{2q+16}{4}+4194304\binom{2q+16}{6}\\
  &\quad+16777216\binom{2q+16}{8}+67108864\binom{2q+16}{10}
    +268435456\binom{2q+16}{12}\\
  &\quad+1073741824\binom{2q+16}{14}
    +4294967296\binom{2q+16}{16}
  \le \binom{2q+16}{q}.
\end{aligned}
\]
\end{lemma}

\begin{proof}
At \(q=41\), the left side is
\[
  4106074088782728938205282304
  <
  6953379637389094006997330256
  =
  \binom{98}{41}.
\]
Let \(P(q)\) denote the displayed left side.  For \(0\le p\le16\),
\[
  \frac{\binom{2q+18}{p}}{\binom{2q+16}{p}}
  \le
  \frac{(2q+18)(2q+17)}{(2q+2)(2q+1)}
  <2
  \qquad(q\ge41),
\]
because \(4q^2-58q-302>0\).  Hence \(P(q+1)<2P(q)\).  On the other hand,
\[
  \frac{\binom{2q+18}{q+1}}{\binom{2q+16}{q}}
  =
  \frac{(2q+18)(2q+17)}{(q+1)(q+17)}
  >3,
\]
since the numerator minus three times the denominator is \(q^2+16q+255\).
Thus the binomial side grows by a larger factor than \(P(q)\), and induction
from \(q=41\) proves the inequality.
\end{proof}

\begin{lemma}[Exact sixteenth $C$ pause-arithmetic checks after the first two rows]\label{lem:bc-c-sixteenth-frontier-pause-arithmetic-certificate}
For \(23\le q\le40\), put
\[
\begin{aligned}
  P(q)={}&65536\binom{2q+16}{0}+262144\binom{2q+16}{2}
    +1048576\binom{2q+16}{4}+4194304\binom{2q+16}{6}\\
  &+16777216\binom{2q+16}{8}+67108864\binom{2q+16}{10}
    +268435456\binom{2q+16}{12}\\
  &+1073741824\binom{2q+16}{14}
    +4294967296\binom{2q+16}{16}.
\end{aligned}
\]
Then
\[
  2^{2q+1}P(q)s_{q+16}\le s_{2q+16}.
\]
\end{lemma}

\begin{proof}
The exact C++/GMP/OpenMP replay
\[
\texttt{certificates/post\_m29/bc\_sixteenth\_frontier\_gmp\_certificate.log}
\]
computes the displayed integer inequality for
\(23\le q\le40\).  It exits with
status \(0\), reports \texttt{SUMMARY failures=0}, and has SHA-256
\[
\hashsplit{9276304d5e5aff2105d4267604664e9f}{14dbcc42df60e0a643c335b028250009}.
\]
It verifies the exact positive margins
\[
\begin{array}{c|c|c}
q & P(q) & s_{2q+16}-2^{2q+1}P(q)s_{q+16}\\ \hline
23&
1205807351611961202507776&
1509858095356523969157575750967910253082932258610380162720969783022376907880344178432\\
24&
2150483723975075628646400&
6887439374866769116850907359416051398468346401926423260145098621747260172624887220429056\\
25&
3758461644764183387242496&
29367302157090546319937862855313938422330258031617816279038033437555273144503561727521153280\\
26&
6446046066848765913399296&
131903570306498698303191973572111336293213641721840023392678650626740716482555395230962178557440\\
27&
10862292206316129681735680&
627983098510752268297381811601656593400968420701540151719427015916832474624724607008670529749451264\\
28&
18004415858340125211688960&
3165506383593940812468228722041173064141891020166206247078653265267460916650382989359944786840604957184\\
29&
29383678802502665771483136&
16868160512001818037962791490485900780051661543048128204573475148033250265173410976979575043152094991129088\\
30&
47261085444415729746640896&
94879010791298187280337279767813612901594464804195239499568072174104637185659235242519350402043059657655106560\\
31&
74978658256448367095513088&
562512973606149046916957038460615249681855162365167296181745038752772515224273173045181013557820531175807000417280\\
32&
117420327558262762294542336&
3510497287731376952825487778351155018445004309094586720427210166473009943659609766816499788004630733426037417146432512\\
33&
181647019438047091761676288&
23031461318490591842248538642177456855384557959341538185195457548228614660390539828570591712060418628420937657600481737728\\
34&
277763886822011089319428096&
158657763193874174278210820082600662167588545781561877476243192916863780523554044405574573232077161704573954546870411113637888\\
35&
420094438856161969507729408&
1146261223216104381673805361166568220687909854353528630247868414286977064356805899615664456043856543189889053552940393928314664960\\
36&
628757338415691257940606976&
8675754396856524847250323812657370220257430919244929339951771850070937018943755082721469025674495507686966397278040733024039425030144\\
37&
931767749965878824904491008&
68718405964968875833321335891201862224993158588072239675646218556428209629218997623557804536218766867319872635571876223450378612674213888\\
38&
1367817381467253149340860416&
569039354428064809581655104841603985107343828899994063684659330770025354114924528949837165102483062905369652375645316754042484144682497888256\\
39&
1989927006546544129956708352&
4921474373856215123551282404435396424010032555189249540385631351130693692793295114168251678989705870335566866580160241189491462836436345101971456\\
40&
2870213714624525227454955520&
44415035728604135032970217655684809370220791946857006628208622398005414380915524485038453845820483617257238753156086416783366815115957157569424199680 .
\end{array}
\]
This proves the claim.
\end{proof}

\begin{proposition}[Residual $C$ sixteenth nonboundary lower correction]\label{prop:bc-c-sixteenth-nonboundary-lower-correction}
For every residual row \(C_b\) with \(20\le b\le217\), put \(q=b+1\) and
write \(s_j=m_j^{\mathrm{st}}\) and \(m_j^{C_b}=s_j+\delta_j^{C_b}\).  Then
\[
  -\frac12 s_{2q+16}\le \delta_{2q+16}^{C_b}.
\]
\end{proposition}

\begin{proof}
As in the proof of
\cref{prop:bc-c-first-nonboundary-lower-correction},
\(|\delta_{2q+16}^{C_b}|\) is the number of column-even Pieri two-strip paths
of length \(2q+16\) with terminal height at least \(2q\).  If \(q=21,22\),
\cref{lem:bc-sixteenth-c-q21-q22-fpf-pause-certificate} gives
\[
  |\delta_{2q+16}^{C_b}|\le s_{2q+16}/2 .
\]
For \(q\ge23\), \cref{lem:bc-c-sixteenth-nonboundary-sixteen-pause-code},
with \(P(q)\) denoting the polynomial from
\cref{lem:bc-c-sixteenth-frontier-pause-arithmetic-certificate}, gives
\[
  |\delta_{2q+16}^{C_b}|
  \le
  P(q)\,2^{2q}s_{q+16}.
\]
If \(23\le q\le40\),
\cref{lem:bc-c-sixteenth-frontier-pause-arithmetic-certificate} gives
\[
  |\delta_{2q+16}^{C_b}|\le s_{2q+16}/2 .
\]
If \(q\ge41\), \cref{lem:bc-sixteenth-c-pause-binomial-absorption} gives
\[
  |\delta_{2q+16}^{C_b}|
  \le
  2^{2q}\binom{2q+16}{q}s_{q+16}.
\]
The exact GMP/OpenMP replay
\[
\texttt{certificates/post\_m29/bc\_pieri\_powerloss\_2q\_gmp\_certificate.log}
\]
checks the inequality
\[
  2^{2q+1}\binom{2q+16}{q}s_{q+16}\le s_{2q+16}
\]
on this residual row and moment as part of its all-row residual-window audit;
it reports \texttt{failures=0} and has SHA-256
\[
\hashsplit{124068c5daa1123088a3cd79c3f6b839}{fee45a1aaef637ed82161c465d86bf45}.
\]
Thus \(|\delta_{2q+16}^{C_b}|\le s_{2q+16}/2\) also for \(q\ge41\), which
implies the claimed lower bound.
\end{proof}

\begin{lemma}[$C$ seventeenth-nonboundary height-bad paths have a seventeen-pause code]\label{lem:bc-c-seventeenth-nonboundary-seventeen-pause-code}
For \(q\ge1\), the number of column-even Pieri two-strip paths
\[
  \varnothing=\lambda^{(0)}\subset\lambda^{(1)}\subset\cdots
  \subset\lambda^{(2q+17)}
\]
with
\[
  \ell(\lambda^{(2q+17)})\ge2q
\]
is at most
\[
\begin{aligned}
  \biggl(&262144\binom{2q+17}{1}+1048576\binom{2q+17}{3}
    +4194304\binom{2q+17}{5}+16777216\binom{2q+17}{7}\\
  &\quad+67108864\binom{2q+17}{9}+268435456\binom{2q+17}{11}
    +1073741824\binom{2q+17}{13}\\
  &\quad+4294967296\binom{2q+17}{15}
    +17179869184\binom{2q+17}{17}\biggr)2^{2q}s_{q+17}.
\end{aligned}
\]
\end{lemma}

\begin{proof}
The terminal height is even and at most \(2q+17\).  Thus it is
\[
  2q+16,\quad 2q+14,\quad 2q+12,\quad 2q+10,\quad 2q+8,
  \quad 2q+6,\quad 2q+4,\quad 2q+2,\quad\text{or}\quad 2q .
\]
These cases have respectively \(1,3,5,7,9,11,13,15,17\) steps which fail to
create a new row.  Deleting the first column, shifting left, and
standardizing the repeated pause labels injects the paths with \(p\) pauses
into the choice of the \(p\)-element pause set and an even-column standard
tableau of size \(2q+17+p\).  Therefore the nine cases contribute at most
\[
  \binom{2q+17}{1}2^{2q+18}s_{q+9},
  \quad
  \binom{2q+17}{3}2^{2q+20}s_{q+10},
  \quad
  \binom{2q+17}{5}2^{2q+22}s_{q+11},
\]
\[
  \binom{2q+17}{7}2^{2q+24}s_{q+12},
  \quad
  \binom{2q+17}{9}2^{2q+26}s_{q+13},
  \quad
  \binom{2q+17}{11}2^{2q+28}s_{q+14},
\]
\[
  \binom{2q+17}{13}2^{2q+30}s_{q+15},
  \quad
  \binom{2q+17}{15}2^{2q+32}s_{q+16},
  \quad
  \binom{2q+17}{17}2^{2q+34}s_{q+17},
\]
by
\cref{lem:bc-even-column-standard-tableaux-2n-bound,lem:bc-stable-moments-monotone-after-two}.
After replacing \(s_{q+9},\ldots,s_{q+16}\) by \(s_{q+17}\), the sum is the
displayed bound.
\end{proof}

\begin{lemma}[The first three seventeenth $C$ pause cases are absorbed by the fixed-point-free count]\label{lem:bc-seventeenth-c-q21-q22-q23-fpf-pause-certificate}
For \(q=21,22,23\), the number of column-even Pieri two-strip paths of length
\(2q+17\) with terminal height at least \(2q\) is at most \(s_{2q+17}/2\).
\end{lemma}

\begin{proof}
Use the pause-set injection from
\cref{lem:bc-c-seventeenth-nonboundary-seventeen-pause-code}, but count the
target even-column standard tableaux exactly by
\cref{lem:bc-even-column-standard-tableaux-fpf-count}.  For \(q=21,22,23\),
this gives the upper bound
\[
  \sum_{\substack{1\le p\le17\\ p\ {\rm odd}}}
  \binom{2q+17}{p}(2q+16+p)!! .
\]
The exact C++/GMP/OpenMP replay
\[
\texttt{certificates/post\_m29/bc\_seventeenth\_frontier\_gmp\_certificate.log}
\]
checks the displayed integer sums and verifies
\[
\begin{array}{c|c|c|c}
q & j & \mathrm{FPF}(q,j) & s_j-2\mathrm{FPF}(q,j)\\ \hline
21&59&
3646487365485267435348823185481307926680796826335393924646909705390625&
7916863325871225525484448344511608155646694346196510796830518427492569348859198\\
22&61&
543042690038776462819349775683735234740329154328714344834795232013984375&
28498624570758949018648802926291872636734794152754189457295572986378155509201536530\\
23&63&
80936530223570392233014003681608918802344862242611517611589595216614921875&
109541225789672302871211354711937491540053284782547178662849246294534857321429988276186 .
\end{array}
\]
The replay exits with status \(0\), reports \texttt{SUMMARY failures=0}, and
has SHA-256
\[
\hashsplit{a1241f73bdcbce52caf77e4ec67c6ef4}{44661548ce303dfae09cf932a1a100a1}.
\]
\end{proof}

\begin{lemma}[The central binomial factor absorbs the seventeen-pause loss]\label{lem:bc-seventeenth-c-pause-binomial-absorption}
For every \(q\ge44\),
\[
\begin{aligned}
  &262144\binom{2q+17}{1}+1048576\binom{2q+17}{3}
    +4194304\binom{2q+17}{5}+16777216\binom{2q+17}{7}\\
  &\quad+67108864\binom{2q+17}{9}+268435456\binom{2q+17}{11}
    +1073741824\binom{2q+17}{13}\\
  &\quad+4294967296\binom{2q+17}{15}
    +17179869184\binom{2q+17}{17}
  \le \binom{2q+17}{q}.
\end{aligned}
\]
\end{lemma}

\begin{proof}
At \(q=44\), the left side is
\[
  284005981622630843449868550144
  <
  801458485881381465643043319600
  =
  \binom{105}{44}.
\]
Let \(P(q)\) denote the displayed left side.  For \(1\le p\le17\),
\[
  \frac{\binom{2q+19}{p}}{\binom{2q+17}{p}}
  \le
  \frac{(2q+19)(2q+18)}{(2q+2)(2q+1)}
  <2
  \qquad(q\ge44),
\]
because \(4q^2-62q-338>0\).  Hence \(P(q+1)<2P(q)\).  On the other hand,
\[
  \frac{\binom{2q+19}{q+1}}{\binom{2q+17}{q}}
  =
  \frac{(2q+19)(2q+18)}{(q+1)(q+18)}
  >3,
\]
since the numerator minus three times the denominator is \(q^2+17q+288\).
Thus the binomial side grows by a larger factor than \(P(q)\), and induction
from \(q=44\) proves the inequality.
\end{proof}

\begin{lemma}[Exact seventeenth $C$ pause-arithmetic checks after the first three rows]\label{lem:bc-c-seventeenth-frontier-pause-arithmetic-certificate}
For \(24\le q\le43\), put
\[
\begin{aligned}
  P(q)={}&262144\binom{2q+17}{1}+1048576\binom{2q+17}{3}
    +4194304\binom{2q+17}{5}+16777216\binom{2q+17}{7}\\
  &+67108864\binom{2q+17}{9}+268435456\binom{2q+17}{11}
    +1073741824\binom{2q+17}{13}\\
  &+4294967296\binom{2q+17}{15}
    +17179869184\binom{2q+17}{17}.
\end{aligned}
\]
Then
\[
  2^{2q+1}P(q)s_{q+17}\le s_{2q+17}.
\]
\end{lemma}

\begin{proof}
The exact C++/GMP/OpenMP replay
\[
\texttt{certificates/post\_m29/bc\_seventeenth\_frontier\_gmp\_certificate.log}
\]
computes the displayed integer inequality for
\(24\le q\le43\).  It exits
with status \(0\), reports \texttt{SUMMARY failures=0}, and has SHA-256
\[
\hashsplit{a1241f73bdcbce52caf77e4ec67c6ef4}{44661548ce303dfae09cf932a1a100a1}.
\]
It verifies the exact positive margins
\[
\begin{array}{c|c|c}
q & P(q) & s_{2q+17}-2^{2q+1}P(q)s_{q+17}\\ \hline
24&
32998770011970420546469888&
406133143944118295207614240726718448502652325289719017791476084366503925692809970595356672\\
25&
59432295280479833893371904&
1941499182848369406719839490518623819407961920803338186507903767856284985075833970054084624640\\
26&
104949831086780696214372352&
9031945222395551182913192167976209468851868232302089682305445821776409870033578901180701443798528\\
27&
181941306932409381388288000&
44271719229880890189547872064109949525701403361559606499910776314872736408892370485018870394678883840\\
28&
310009288261961582165360640&
229498661129310833547231936224305438372558721366368822440852078362829082699391190525256472750860834473984\\
29&
519720029531414377302589440&
1256676673946322185324179381725643201647556271171114628526512576061321133074616189635261271593410393514951168\\
30&
858089355587184271954280448&
7258237967849954884931042357231583835023780403411800254704455350097305906761294820608346704844363133786011923456\\
31&
1396512829567985071065661440&
44157232826577367463066310769097404935543041447060462616239168375501541075669014522999300413448063226800971749792768\\
32&
2242104392410675112264859648&
282594820645935588180606434340648656948538070201354651900555743018447857962611806536203980023318890837733402268871467008\\
33&
3553741968600343190861250560&
1900094241915981371041842174158949860514688227184176473508776365234209320123587584120550457574615370757135590943323242863616\\
34&
5564553815388651717026643968&
13406572350449907596228493749260500982499361952017590125882132323276823171683987748495924295842451134934399752657927118730708992\\
35&
8613141912361358819620290560&
99151536294249992938880297472579366309800773550293850709355141394036929696778156090420365629994503489283327867410263638212840220672\\
36&
13186560505712868155787051008&
767803834154796917279602379262274067098099810859656697933211098066247212563468397527059996962310916734459925523186258372175378763849728\\
37&
19978987958686258441028894720&
6219012485561720158730296415209176480235169093635443707357945941299397925258623737956824977251408429175755006309554416848552204030717450240\\
38&
29971195321780924549902368768&
52636114508684281019658038026429125714487520644936565969511426459184031133799338961330958328524645955609481224725571333523390937278352961089536\\
39&
44537382011599040538580090880&
465079114888095472746257509756462428220107222726481931184292362018946686954802126127939516408279551958947220911963307970936286875122253708863459328\\
40&
65587785241707120173479362560&
4286049101830469787585145826064000317264620522713040807151224853474572863140823397090827911938405137037628615146236916602459850403706257452358927187968\\
41&
95757755839909279244180258816&
41162110219623827986713952512453953258690562246881118283934277469268456583091093250575747832004523306959225351920567914455615250061826130952059064530898944\\
42&
138656824216398390088785199104&
411610493361212659403148765054814801562910207570797881088266812485092456976604116455115892512067925145910330139553773678408593345734093956515364107748362182656\\
43&
199194769184163601401791643648&
4282289551634348924637847924815314882271058168493559037093783432851176072866241963481233049813442652660078840873862478967639526137222901228961502751701895318822912 .
\end{array}
\]
This proves the claim.
\end{proof}

\begin{proposition}[Residual $C$ seventeenth nonboundary lower correction]\label{prop:bc-c-seventeenth-nonboundary-lower-correction}
For every residual row \(C_b\) with \(20\le b\le217\), put \(q=b+1\) and
write \(s_j=m_j^{\mathrm{st}}\) and \(m_j^{C_b}=s_j+\delta_j^{C_b}\).  Then
\[
  -\frac12 s_{2q+17}\le \delta_{2q+17}^{C_b}.
\]
\end{proposition}

\begin{proof}
As in the proof of
\cref{prop:bc-c-first-nonboundary-lower-correction},
\(|\delta_{2q+17}^{C_b}|\) is the number of column-even Pieri two-strip paths
of length \(2q+17\) with terminal height at least \(2q\).  If \(q=21,22,23\),
\cref{lem:bc-seventeenth-c-q21-q22-q23-fpf-pause-certificate} gives
\[
  |\delta_{2q+17}^{C_b}|\le s_{2q+17}/2 .
\]
For \(q\ge24\), \cref{lem:bc-c-seventeenth-nonboundary-seventeen-pause-code},
with \(P(q)\) denoting the polynomial from
\cref{lem:bc-c-seventeenth-frontier-pause-arithmetic-certificate}, gives
\[
  |\delta_{2q+17}^{C_b}|
  \le
  P(q)\,2^{2q}s_{q+17}.
\]
If \(24\le q\le43\),
\cref{lem:bc-c-seventeenth-frontier-pause-arithmetic-certificate} gives
\[
  |\delta_{2q+17}^{C_b}|\le s_{2q+17}/2 .
\]
If \(q\ge44\), \cref{lem:bc-seventeenth-c-pause-binomial-absorption} gives
\[
  |\delta_{2q+17}^{C_b}|
  \le
  2^{2q}\binom{2q+17}{q}s_{q+17}.
\]
The exact GMP/OpenMP replay
\[
\texttt{certificates/post\_m29/bc\_pieri\_powerloss\_2q\_gmp\_certificate.log}
\]
checks the inequality
\[
  2^{2q+1}\binom{2q+17}{q}s_{q+17}\le s_{2q+17}
\]
on this residual row and moment as part of its all-row residual-window audit;
it reports \texttt{failures=0} and has SHA-256
\[
\hashsplit{124068c5daa1123088a3cd79c3f6b839}{fee45a1aaef637ed82161c465d86bf45}.
\]
Thus \(|\delta_{2q+17}^{C_b}|\le s_{2q+17}/2\) also for \(q\ge44\), which
implies the claimed lower bound.
\end{proof}

\begin{lemma}[$C$ eighteenth-nonboundary height-bad paths have an eighteen-pause code]\label{lem:bc-c-eighteenth-nonboundary-eighteen-pause-code}
For \(q\ge1\), the number of column-even Pieri two-strip paths
\[
  \varnothing=\lambda^{(0)}\subset\lambda^{(1)}\subset\cdots
  \subset\lambda^{(2q+18)}
\]
with
\[
  \ell(\lambda^{(2q+18)})\ge2q
\]
is at most
\[
\begin{aligned}
  \biggl(&262144\binom{2q+18}{0}+1048576\binom{2q+18}{2}
    +4194304\binom{2q+18}{4}+16777216\binom{2q+18}{6}\\
  &\quad+67108864\binom{2q+18}{8}+268435456\binom{2q+18}{10}
    +1073741824\binom{2q+18}{12}\\
  &\quad+4294967296\binom{2q+18}{14}
    +17179869184\binom{2q+18}{16}
    +68719476736\binom{2q+18}{18}\biggr)2^{2q}s_{q+18}.
\end{aligned}
\]
\end{lemma}

\begin{proof}
The terminal height is even and at most \(2q+18\).  Thus it is
\[
  2q+18,\quad 2q+16,\quad 2q+14,\quad 2q+12,\quad 2q+10,
  \quad 2q+8,\quad 2q+6,\quad 2q+4,\quad 2q+2,\quad\text{or}\quad 2q .
\]
These cases have respectively \(0,2,4,6,8,10,12,14,16,18\) steps which fail
to create a new row.  Deleting the first column, shifting left, and
standardizing the repeated pause labels injects the paths with \(p\) pauses
into the choice of the \(p\)-element pause set and an even-column standard
tableau of size \(2q+18+p\).  Therefore the ten cases contribute at most
\[
  \binom{2q+18}{0}2^{2q+18}s_{q+9},
  \quad
  \binom{2q+18}{2}2^{2q+20}s_{q+10},
  \quad
  \binom{2q+18}{4}2^{2q+22}s_{q+11},
\]
\[
  \binom{2q+18}{6}2^{2q+24}s_{q+12},
  \quad
  \binom{2q+18}{8}2^{2q+26}s_{q+13},
  \quad
  \binom{2q+18}{10}2^{2q+28}s_{q+14},
\]
\[
  \binom{2q+18}{12}2^{2q+30}s_{q+15},
  \quad
  \binom{2q+18}{14}2^{2q+32}s_{q+16},
  \quad
  \binom{2q+18}{16}2^{2q+34}s_{q+17},
\]
\[
  \binom{2q+18}{18}2^{2q+36}s_{q+18},
\]
by
\cref{lem:bc-even-column-standard-tableaux-2n-bound,lem:bc-stable-moments-monotone-after-two}.
After replacing \(s_{q+9},\ldots,s_{q+17}\) by \(s_{q+18}\), the sum is the
displayed bound.
\end{proof}

\begin{lemma}[The first three eighteenth $C$ pause cases are absorbed by the fixed-point-free count]\label{lem:bc-eighteenth-c-q21-q22-q23-fpf-pause-certificate}
For \(q=21,22,23\), the number of column-even Pieri two-strip paths of length
\(2q+18\) with terminal height at least \(2q\) is at most \(s_{2q+18}/2\).
\end{lemma}

\begin{proof}
Use the pause-set injection from
\cref{lem:bc-c-eighteenth-nonboundary-eighteen-pause-code}, but count the
target even-column standard tableaux exactly by
\cref{lem:bc-even-column-standard-tableaux-fpf-count}.  For \(q=21,22,23\),
this gives the upper bound
\[
  \sum_{\substack{0\le p\le18\\ p\ {\rm even}}}
  \binom{2q+18}{p}(2q+17+p)!! .
\]
The exact C++/GMP/OpenMP replay
\[
\texttt{certificates/post\_m29/bc\_eighteenth\_frontier\_gmp\_certificate.log}
\]
checks the displayed integer sums and verifies
\[
\begin{array}{c|c|c|c}
q & j & \mathrm{FPF}(q,j) & s_j-2\mathrm{FPF}(q,j)\\ \hline
21&60&
936103804833627766619558658829368839169497669297772127543300376719453125&
471052483290496707652557262731290971574731252133228937370372771457389018870879606\\
22&62&
147792316180051306675315989515099722426677110364389754267875360964914531250&
1752662440282196942849401054887494894303585264196585842685687141835826067117925676956\\
23&64&
23313182866657367477677417211957455715304495181498386549564779747042798203125&
6955857144680513605018656639438248482599989467501880284347279981375714830577588469923606 .
\end{array}
\]
The replay exits with status \(0\), reports \texttt{SUMMARY failures=0}, and
has SHA-256
\[
\hashsplit{2c1768d73f6e84f175648fe543b44e67}{f791c9221e749ff42253c8a69a6ed7a0}.
\]
\end{proof}

\begin{lemma}[The central binomial factor absorbs the eighteen-pause loss]\label{lem:bc-eighteenth-c-pause-binomial-absorption}
For every \(q\ge46\),
\[
\begin{aligned}
  &262144\binom{2q+18}{0}+1048576\binom{2q+18}{2}
    +4194304\binom{2q+18}{4}+16777216\binom{2q+18}{6}\\
  &\quad+67108864\binom{2q+18}{8}+268435456\binom{2q+18}{10}
    +1073741824\binom{2q+18}{12}\\
  &\quad+4294967296\binom{2q+18}{14}
    +17179869184\binom{2q+18}{16}
    +68719476736\binom{2q+18}{18}
  \le \binom{2q+18}{q}.
\end{aligned}
\]
\end{lemma}

\begin{proof}
At \(q=46\), the left side is
\[
  13826118675637230186548178452480
  <
  22747362824110665179416185383175
  =
  \binom{110}{46}.
\]
Let \(P(q)\) denote the displayed left side.  For \(0\le p\le18\),
\[
  \frac{\binom{2q+20}{p}}{\binom{2q+18}{p}}
  \le
  \frac{(2q+20)(2q+19)}{(2q+2)(2q+1)}
  <2
  \qquad(q\ge46),
\]
because \(4q^2-66q-376>0\).  Hence \(P(q+1)<2P(q)\).  On the other hand,
\[
  \frac{\binom{2q+20}{q+1}}{\binom{2q+18}{q}}
  =
  \frac{(2q+20)(2q+19)}{(q+1)(q+19)}
  >3,
\]
since the numerator minus three times the denominator is \(q^2+18q+323\).
Thus the binomial side grows by a larger factor than \(P(q)\), and induction
from \(q=46\) proves the inequality.
\end{proof}

\begin{lemma}[Exact eighteenth $C$ pause-arithmetic checks after the first three rows]\label{lem:bc-c-eighteenth-frontier-pause-arithmetic-certificate}
For \(24\le q\le45\),
put
\[
\begin{aligned}
  P(q)={}&262144\binom{2q+18}{0}+1048576\binom{2q+18}{2}
    +4194304\binom{2q+18}{4}+16777216\binom{2q+18}{6}\\
  &+67108864\binom{2q+18}{8}+268435456\binom{2q+18}{10}
    +1073741824\binom{2q+18}{12}\\
  &+4294967296\binom{2q+18}{14}
    +17179869184\binom{2q+18}{16}
    +68719476736\binom{2q+18}{18}.
\end{aligned}
\]
Then
\[
  2^{2q+1}P(q)s_{q+18}\le s_{2q+18}.
\]
\end{lemma}

\begin{proof}
The exact C++/GMP/OpenMP replay
\[
\texttt{certificates/post\_m29/bc\_eighteenth\_frontier\_gmp\_certificate.log}
\]
computes the displayed integer inequality for
\(24\le q\le45\).  It exits
with status \(0\), reports \texttt{SUMMARY failures=0}, and has SHA-256
\[
\hashsplit{2c1768d73f6e84f175648fe543b44e67}{f791c9221e749ff42253c8a69a6ed7a0}.
\]
It verifies the exact positive margins
\[
\begin{array}{c|c|c}
q & P(q) & s_{2q+18}-2^{2q+1}P(q)s_{q+18}\\ \hline
24&
485690524492741371143389184&
3415569985153575367274533197969017388038334138152550172056441941023544842155192750098702592\\
25&
901013500420353108094484480&
123718617515897586083287220196259693056833787006661964394946443203878730757386403376819652756992\\
26&
1637475412045018423195074560&
625394681032748308818491724462077828049815837548051194740929889237728311134709854075006194072093184\\
27&
2919209192203235399393607680&
3164684960453723419033988171777179335593364112200150921402526606711948345746203707148891262097133197824\\
28&
5111212844565486053329469440&
16867898744030385967552373953954359844277505317807397078701856006530416894479655755013361259277853426358784\\
29&
8798834219976758419406454784&
94878926971285181538362280432357436386352886938210148090373554441858993703952358346277164182311821723573861376\\
30&
14907371697563580919094444032&
562512946623016005557246402126450655392812786058689979945775339139007064297727803540268437148796723348746175943680\\
31&
24879815802004287614091526144&
3510497278994553285352124342714737035974598674202364154782135377114916374403314125803840200958448825657422499966467072\\
32&
40937925700357484315497725952&
23031461315644103996802830544002057424805173049182115216330189013059475845295474945514034743505845226678652754150313999360\\
33&
66461509138151476374200909824&
158657763192940671359587645926155393584322449555853777732149024169128998067201988327914212410814814936932576676042471785998336\\
34&
106533717472108512681303212032&
1146261223215796127281153047067350611449763178692448923941319717385463465830104573061717280272938022298371243734483970747397126144\\
35&
168717341749300791726141865984&
8675754396856422326985058396204698935972320285039604537325007431446181745610465588672873452556769296343249163706982015629959205238784\\
36&
264149708380352676845388562432&
68718405964968841483731398906052289974250895191315055265268798505297825076375599885681674624931303886893466708197283025059823822024030208\\
37&
409073328050388683175147864064&
569039354428064797984913670253094170681299410754373761692490833218806022921157367406202814323741402252155125718796725562049124835847063359488\\
38&
626957812521160029184913047552&
4921474373856215119605507272183515287971251329422263637050499515179821579622753202390863736087838752303061601368975952103575187622073710273327104\\
39&
951418036509207647414246440960&
44415035728604135031616963184031838155041986870000616121704484654814288749672989373715732481361386194193420854811113890872159621802739070075973079040\\
40&
1430196899008872205162802053120&
417889612984311197972413715741748323772344422906441677583908278778108905721753151975371538494906635256729786947618421020890410173026737895363840681381888\\
41&
2130561760307622776104959606784&
4095628359235226170344268509952848370247001712060145269184776415125987077056785411510818545864925263991570640882173727642204791341532558420075830159180062720\\
42&
3146565859795394230612091207680&
41778449637156136111589853722521634686017301903581847378885849827856525293907642812020576950191088240405075026665981979218121837877137192564493707361175613218816\\
43&
4608754783361488019049362489344&
443216814209338571157347337580637569052962560356503201840060461059507360951427126728892198626465957510509130212037216979092424360355571471425307335054956552396906496\\
44&
6697059394790832303211325685760&
4886351313102937424311841329203021597323852241640700645228150343195504189732562797678949425782069485468123700774333533211215477761493812931761484720621722531497970933760\\
45&
9657817815894010546150315393024&
55942579353574319299284514196127136548296562054101092404013697242428863607039004976195082562661116297743699395165386419324954141421672360147943644470193623896920549659361280 .
\end{array}
\]
This proves the claim.
\end{proof}

\begin{proposition}[Residual $C$ eighteenth nonboundary lower correction]\label{prop:bc-c-eighteenth-nonboundary-lower-correction}
For every residual row \(C_b\) with \(20\le b\le217\), put \(q=b+1\) and
write \(s_j=m_j^{\mathrm{st}}\) and \(m_j^{C_b}=s_j+\delta_j^{C_b}\).  Then
\[
  -\frac12 s_{2q+18}\le \delta_{2q+18}^{C_b}.
\]
\end{proposition}

\begin{proof}
As in the proof of
\cref{prop:bc-c-first-nonboundary-lower-correction},
\(|\delta_{2q+18}^{C_b}|\) is the number of column-even Pieri two-strip paths
of length \(2q+18\) with terminal height at least \(2q\).  If \(q=21,22,23\),
\cref{lem:bc-eighteenth-c-q21-q22-q23-fpf-pause-certificate} gives
\[
  |\delta_{2q+18}^{C_b}|\le s_{2q+18}/2 .
\]
For \(q\ge24\), \cref{lem:bc-c-eighteenth-nonboundary-eighteen-pause-code},
with \(P(q)\) denoting the polynomial from
\cref{lem:bc-c-eighteenth-frontier-pause-arithmetic-certificate}, gives
\[
  |\delta_{2q+18}^{C_b}|
  \le
  P(q)\,2^{2q}s_{q+18}.
\]
If \(24\le q\le45\),
\cref{lem:bc-c-eighteenth-frontier-pause-arithmetic-certificate} gives
\[
  |\delta_{2q+18}^{C_b}|\le s_{2q+18}/2 .
\]
If \(q\ge46\), \cref{lem:bc-eighteenth-c-pause-binomial-absorption} gives
\[
  |\delta_{2q+18}^{C_b}|
  \le
  2^{2q}\binom{2q+18}{q}s_{q+18}.
\]
The exact GMP/OpenMP replay
\[
\texttt{certificates/post\_m29/bc\_pieri\_powerloss\_2q\_gmp\_certificate.log}
\]
checks the inequality
\[
  2^{2q+1}\binom{2q+18}{q}s_{q+18}\le s_{2q+18}
\]
on this residual row and moment as part of its all-row residual-window audit;
it reports \texttt{failures=0} and has SHA-256
\[
\hashsplit{124068c5daa1123088a3cd79c3f6b839}{fee45a1aaef637ed82161c465d86bf45}.
\]
Thus \(|\delta_{2q+18}^{C_b}|\le s_{2q+18}/2\) also for \(q\ge46\), which
implies the claimed lower bound.
\end{proof}

\begin{lemma}[$C$ nineteenth-nonboundary height-bad paths have a nineteen-pause code]\label{lem:bc-c-nineteenth-nonboundary-nineteen-pause-code}
For \(q\ge1\), the number of column-even Pieri two-strip paths
\[
  \varnothing=\lambda^{(0)}\subset\lambda^{(1)}\subset\cdots
  \subset\lambda^{(2q+19)}
\]
with
\[
  \ell(\lambda^{(2q+19)})\ge2q
\]
is at most
\[
\begin{aligned}
  \biggl(&1048576\binom{2q+19}{1}+4194304\binom{2q+19}{3}
    +16777216\binom{2q+19}{5}\\
  &\quad+67108864\binom{2q+19}{7}+268435456\binom{2q+19}{9}
    +1073741824\binom{2q+19}{11}\\
  &\quad+4294967296\binom{2q+19}{13}
    +17179869184\binom{2q+19}{15}\\
  &\quad+68719476736\binom{2q+19}{17}
    +274877906944\binom{2q+19}{19}\biggr)2^{2q}s_{q+19}.
\end{aligned}
\]
\end{lemma}

\begin{proof}
The terminal height is even and at most \(2q+19\).  Thus it is
\[
  2q+18,\quad 2q+16,\quad 2q+14,\quad 2q+12,\quad 2q+10,
  \quad 2q+8,\quad 2q+6,\quad 2q+4,\quad 2q+2,\quad\text{or}\quad 2q .
\]
These cases have respectively \(1,3,5,7,9,11,13,15,17,19\) steps which fail
to create a new row.  Deleting the first column, shifting left, and
standardizing the repeated pause labels injects the paths with \(p\) pauses
into the choice of the \(p\)-element pause set and an even-column standard
tableau of size \(2q+19+p\).  Therefore the ten cases contribute at most
\[
  \binom{2q+19}{1}2^{2q+20}s_{q+10},
  \quad
  \binom{2q+19}{3}2^{2q+22}s_{q+11},
  \quad
  \binom{2q+19}{5}2^{2q+24}s_{q+12},
\]
\[
  \binom{2q+19}{7}2^{2q+26}s_{q+13},
  \quad
  \binom{2q+19}{9}2^{2q+28}s_{q+14},
  \quad
  \binom{2q+19}{11}2^{2q+30}s_{q+15},
\]
\[
  \binom{2q+19}{13}2^{2q+32}s_{q+16},
  \quad
  \binom{2q+19}{15}2^{2q+34}s_{q+17},
  \quad
  \binom{2q+19}{17}2^{2q+36}s_{q+18},
\]
\[
  \binom{2q+19}{19}2^{2q+38}s_{q+19},
\]
by
\cref{lem:bc-even-column-standard-tableaux-2n-bound,lem:bc-stable-moments-monotone-after-two}.
After replacing \(s_{q+10},\ldots,s_{q+18}\) by \(s_{q+19}\), the sum is the
displayed bound.
\end{proof}

\begin{lemma}[The first four nineteenth $C$ pause cases are absorbed by the fixed-point-free count]\label{lem:bc-nineteenth-c-q21-q22-q23-q24-fpf-pause-certificate}
For \(q=21,22,23,24\), the number of column-even Pieri two-strip paths of
length \(2q+19\) with terminal height at least \(2q\) is at most
\(s_{2q+19}/2\).
\end{lemma}

\begin{proof}
Use the pause-set injection from
\cref{lem:bc-c-nineteenth-nonboundary-nineteen-pause-code}, but count the
target even-column standard tableaux exactly by
\cref{lem:bc-even-column-standard-tableaux-fpf-count}.  For \(q=21,22,23,24\),
this gives the upper bound
\[
  \sum_{\substack{1\le p\le19\\ p\ {\rm odd}}}
  \binom{2q+19}{p}(2q+18+p)!! .
\]
The exact C++/GMP/OpenMP replay
\[
\texttt{certificates/post\_m29/bc\_nineteenth\_frontier\_gmp\_certificate.log}
\]
checks the displayed integer sums and verifies
\[
\begin{array}{c|c|c|c}
q & j & \mathrm{FPF}(q,j) & s_j-2\mathrm{FPF}(q,j)\\ \hline
21&61&
237470108332329713171082643084030364430708345727294881364488502819416328125&
28498624096904817734066929509765286020041535760818156311363238947070740334396849030\\
22&63&
39700601747085018754484911634048710728947977468719606083510193151637269218750&
109541225710432972437488457987433696279319081162255913449895257162737650208588679682436\\
23&65&
6620723654392057477395380854479066940610645630038027813080726841442956808203125&
448652128219640090423345837323001612054965492564703062730327497930184002811740199033637142\\
24&67&
1102869093764874549303216415018746428293136265004334004265436419435902759984375000&
1954211164679237519738993223359785151237160761354975954454229478971669293616074699472392992592 .
\end{array}
\]
The replay exits with status \(0\), reports \texttt{SUMMARY failures=0}, and
has SHA-256
\[
\hashsplit{519d2680d84e75ce37567e5b89e5b1d7}{3642f0ddf74f02f115aeb728df4168dd}.
\]
\end{proof}

\begin{lemma}[The central binomial factor absorbs the nineteen-pause loss]\label{lem:bc-nineteenth-c-pause-binomial-absorption}
For every \(q\ge49\),
\[
\begin{aligned}
  &1048576\binom{2q+19}{1}+4194304\binom{2q+19}{3}
    +16777216\binom{2q+19}{5}\\
  &\quad+67108864\binom{2q+19}{7}+268435456\binom{2q+19}{9}
    +1073741824\binom{2q+19}{11}\\
  &\quad+4294967296\binom{2q+19}{13}
    +17179869184\binom{2q+19}{15}\\
  &\quad+68719476736\binom{2q+19}{17}
    +274877906944\binom{2q+19}{19}
  \le \binom{2q+19}{q}.
\end{aligned}
\]
\end{lemma}

\begin{proof}
At \(q=49\), the left side is
\[
  959886034421822120920097112457216
  <
  2631605948503014979355959106417475
  =
  \binom{117}{49}.
\]
Let \(P(q)\) denote the displayed left side.  For \(1\le p\le19\),
\[
  \frac{\binom{2q+21}{p}}{\binom{2q+19}{p}}
  \le
  \frac{(2q+21)(2q+20)}{(2q+2)(2q+1)}
  <2
  \qquad(q\ge49),
\]
because \(4q^2-70q-416>0\).  Hence \(P(q+1)<2P(q)\).  On the other hand,
\[
  \frac{\binom{2q+21}{q+1}}{\binom{2q+19}{q}}
  =
  \frac{(2q+21)(2q+20)}{(q+1)(q+20)}
  >3,
\]
since the numerator minus three times the denominator is \(q^2+19q+360\).
Thus the binomial side grows by a larger factor than \(P(q)\), and induction
from \(q=49\) proves the inequality.
\end{proof}

\begin{lemma}[Exact nineteenth $C$ pause-arithmetic checks after the first four rows]\label{lem:bc-c-nineteenth-frontier-pause-arithmetic-certificate}
For \(25\le q\le48\),
put
\[
\begin{aligned}
  P(q)={}&1048576\binom{2q+19}{1}+4194304\binom{2q+19}{3}
    +16777216\binom{2q+19}{5}\\
  &+67108864\binom{2q+19}{7}+268435456\binom{2q+19}{9}
    +1073741824\binom{2q+19}{11}\\
  &+4294967296\binom{2q+19}{13}
    +17179869184\binom{2q+19}{15}\\
  &+68719476736\binom{2q+19}{17}
    +274877906944\binom{2q+19}{19}.
\end{aligned}
\]
Then
\[
  2^{2q+1}P(q)s_{q+19}\le s_{2q+19}.
\]
\end{lemma}

\begin{proof}
The exact C++/GMP/OpenMP replay
\[
\texttt{certificates/post\_m29/bc\_nineteenth\_frontier\_gmp\_certificate.log}
\]
computes the displayed integer inequality for
\(25\le q\le48\).
It exits with status \(0\), reports \texttt{SUMMARY failures=0}, and has
SHA-256
\[
\hashsplit{519d2680d84e75ce37567e5b89e5b1d7}{3642f0ddf74f02f115aeb728df4168dd}.
\]
It verifies the exact positive margins
\[
\begin{array}{c|c|c}
q & P(q) & s_{2q+19}-2^{2q+1}P(q)s_{q+19}\\ \hline
25&
13133656758990092272525115392&
3842380120851551944142473147494552077681333347048494952440197066852371374503365116937097548576256\\
26&
24554262757961748518655754240&
42544608825317433825017194072171563221014792045929523136375250871080764354516346378152381902243182080\\
27&
44996875384669081611592007680&
228922596340449498597709272052905627396620123708835017028698459928221956938741047037001684394505925330944\\
28&
80926429529001584723581992960&
1256483959754054882893658262044304324534475077506191506224914323532222922862475479568718579330232663657660928\\
29&
143001582816313022041992200192&
7258173263031613370440356372340613277392551551121377908802219761516476232283048126136018279949264780368380097536\\
30&
248531126707426721043479265280&
44157211009580330347799852365610389867646375399980869260837407956655471123510405335703346449862803360041499264293888\\
31&
425225447850373260819160891392&
282594813254761709829596426296813127388409677795124834347777328446962215948880445410153649680606686689313802356801970176\\
32&
716855183174662333546904944640&
1900094239398941214011041594489397140198908647419788401645661251705045828808782696643946980715692190935359115928224915495936\\
33&
1191687384782268653486732214272&
13406572349587921939227856292443348922331560285939150037786070042479321733339567833058205414911318838003001756814446538540075008\\
34&
1954924263512918327464731607040&
99151536293953033275577402527213970301292818983692703180432292122298796313450985069923488170781024804207575620602585446945727670272\\
35&
3166852726909542920159941885952&
767803834154693969111441063073248733253484113702637951830009933315836735283213525532997041978901467597701890997890436243798336028909568\\
36&
5069065425860958444430457569280&
6219012485561684234815308755293489245290611952063363904592571478250765036838029158742019756523969181319495597983229484671359842652228147200\\
37&
8021988242925087774775178166272&
52636114508684268398522185144146687794616579709169747400742525957491221275634629893719559369982111069191703146698594889510674349422317297250304\\
38&
12558111617945873215456044646400&
465079114888095468280900911393862494883254664227088421555258592053187518620789596123454586082158443011484736251024310918198276707854260417171017728\\
39&
19456857797833949746152734720000&
4286049101830469785993894980348880248457611926002771098965364284788320044849264995689920858827118083458305537023652897763538610343833892233575482851328\\
40&
29849028426037613464679481344000&
41162110219623827986142713682581476779460775136967730397659528658517488946880722663220898289704131309362223206341041876417959103036743621919134021636722688\\
41&
45361398415060047205612274057216&
411610493361212659402942152779612426035671413503991194745996214371682855617877400942854625714083379577612360765418708044785145376427765560092741084947008479232\\
42&
68315416105960826809088440532992&
4282289551634348924637772622467391571052467599720276217270259836464965846273656436276081095407377104121802416263129745110015125003409433024359161898911655254777856\\
43&
101998338065906113339054674149376&
46316141414884575423512758181052631642077578418249613878781970668380438640450664053709209657765008846809126867873962396578394707199903601592712393830542561227664687104\\
44&
151030717518270678269547382833152&
520396248607814162823456888678083574419340659794434686095011665287540075473832731821460177752393256807330106483766578985757206447421175601487977177258999265323742719582208\\
45&
221861281805804846522903362011136&
6069768027148561641469987055484821179900569544002511222699863695645693775147353430401562009777811585923752404130006219706714588038484040173569056915379457859318318461889757184\\
46&
323429246142932688920017619648512&
73442633136598658548124378675009530522796332345920669968955863270601680097171534430050096731843267205526088096581132111831528144658814838049197342622924998895013179392542669258752\\
47&
468045466676533938884971133403136&
921245524013205551953790796064948432422207505327530462415741404473788266278303523601102448802363164678225300392139893987795028748170393657917251102898282026062915325631567166665785344\\
48&
672558077360155695305232772759552&
11972270293412322391994921932842237116039497959401574743677868601647816015977507065502280690490490564115364142207230527762613568271988347220584619482063464562880063648323635618825263759360 .
\end{array}
\]
This proves the claim.
\end{proof}

\begin{proposition}[Residual $C$ nineteenth nonboundary lower correction]\label{prop:bc-c-nineteenth-nonboundary-lower-correction}
For every residual row \(C_b\) with \(20\le b\le217\), put \(q=b+1\) and
write \(s_j=m_j^{\mathrm{st}}\) and \(m_j^{C_b}=s_j+\delta_j^{C_b}\).  Then
\[
  -\frac12 s_{2q+19}\le \delta_{2q+19}^{C_b}.
\]
\end{proposition}

\begin{proof}
As in the proof of
\cref{prop:bc-c-first-nonboundary-lower-correction},
\(|\delta_{2q+19}^{C_b}|\) is the number of column-even Pieri two-strip paths
of length \(2q+19\) with terminal height at least \(2q\).  If
\(q=21,22,23,24\),
\cref{lem:bc-nineteenth-c-q21-q22-q23-q24-fpf-pause-certificate} gives
\[
  |\delta_{2q+19}^{C_b}|\le s_{2q+19}/2 .
\]
For \(q\ge25\), \cref{lem:bc-c-nineteenth-nonboundary-nineteen-pause-code},
with \(P(q)\) denoting the polynomial from
\cref{lem:bc-c-nineteenth-frontier-pause-arithmetic-certificate}, gives
\[
  |\delta_{2q+19}^{C_b}|
  \le
  P(q)\,2^{2q}s_{q+19}.
\]
If \(25\le q\le48\),
\cref{lem:bc-c-nineteenth-frontier-pause-arithmetic-certificate} gives
\[
  |\delta_{2q+19}^{C_b}|\le s_{2q+19}/2 .
\]
If \(q\ge49\), \cref{lem:bc-nineteenth-c-pause-binomial-absorption} gives
\[
  |\delta_{2q+19}^{C_b}|
  \le
  2^{2q}\binom{2q+19}{q}s_{q+19}.
\]
The exact GMP/OpenMP replay
\[
\texttt{certificates/post\_m29/bc\_pieri\_powerloss\_2q\_gmp\_certificate.log}
\]
checks the inequality
\[
  2^{2q+1}\binom{2q+19}{q}s_{q+19}\le s_{2q+19}
\]
on this residual row and moment as part of its all-row residual-window audit;
it reports \texttt{failures=0} and has SHA-256
\[
\hashsplit{124068c5daa1123088a3cd79c3f6b839}{fee45a1aaef637ed82161c465d86bf45}.
\]
Thus \(|\delta_{2q+19}^{C_b}|\le s_{2q+19}/2\) also for \(q\ge49\), which
implies the claimed lower bound.
\end{proof}

\begin{lemma}[$C$ twentieth-nonboundary height-bad paths have a twenty-pause code]\label{lem:bc-c-twentieth-nonboundary-twenty-pause-code}
For \(q\ge1\), the number of column-even Pieri two-strip paths
\[
  \varnothing=\lambda^{(0)}\subset\lambda^{(1)}\subset\cdots
  \subset\lambda^{(2q+20)}
\]
with
\[
  \ell(\lambda^{(2q+20)})\ge2q
\]
is at most
\[
\begin{aligned}
  \biggl(&1048576\binom{2q+20}{0}
    +4194304\binom{2q+20}{2}
    +16777216\binom{2q+20}{4}\\
  &\quad+67108864\binom{2q+20}{6}
    +268435456\binom{2q+20}{8}
    +1073741824\binom{2q+20}{10}\\
  &\quad+4294967296\binom{2q+20}{12}
    +17179869184\binom{2q+20}{14}\\
  &\quad+68719476736\binom{2q+20}{16}
    +274877906944\binom{2q+20}{18}
    +1099511627776\binom{2q+20}{20}\biggr)2^{2q}s_{q+20}.
\end{aligned}
\]
\end{lemma}

\begin{proof}
The terminal height is even and at most \(2q+20\).  Thus it is
\[
  2q+20,\quad 2q+18,\quad 2q+16,\quad 2q+14,\quad 2q+12,
  \quad 2q+10,\quad 2q+8,\quad 2q+6,\quad 2q+4,\quad 2q+2,\quad\text{or}\quad 2q .
\]
These cases have respectively \(0,2,4,6,8,10,12,14,16,18,20\) steps which
fail to create a new row.  Deleting the first column, shifting left, and
standardizing the repeated pause labels injects the paths with \(p\) pauses
into the choice of the \(p\)-element pause set and an even-column standard
tableau of size \(2q+20+p\).  Therefore the eleven cases contribute at most
\[
  \binom{2q+20}{0}2^{2q+20}s_{q+10},
  \quad
  \binom{2q+20}{2}2^{2q+22}s_{q+11},
  \quad
  \binom{2q+20}{4}2^{2q+24}s_{q+12},
\]
\[
  \binom{2q+20}{6}2^{2q+26}s_{q+13},
  \quad
  \binom{2q+20}{8}2^{2q+28}s_{q+14},
  \quad
  \binom{2q+20}{10}2^{2q+30}s_{q+15},
\]
\[
  \binom{2q+20}{12}2^{2q+32}s_{q+16},
  \quad
  \binom{2q+20}{14}2^{2q+34}s_{q+17},
\]
\[
  \binom{2q+20}{16}2^{2q+36}s_{q+18},
  \quad
  \binom{2q+20}{18}2^{2q+38}s_{q+19},
  \quad
  \binom{2q+20}{20}2^{2q+40}s_{q+20},
\]
by
\cref{lem:bc-even-column-standard-tableaux-2n-bound,lem:bc-stable-moments-monotone-after-two}.
After replacing \(s_{q+10},\ldots,s_{q+19}\) by \(s_{q+20}\), the sum is the
displayed bound.
\end{proof}

\begin{lemma}[The first five twentieth $C$ pause cases are absorbed by the fixed-point-free count]\label{lem:bc-twentieth-c-q21-q22-q23-q24-q25-fpf-pause-certificate}
For \(q=21,22,23,24,25\), the number of column-even Pieri two-strip paths of
length \(2q+20\) with terminal height at least \(2q\) is at most
\(s_{2q+20}/2\).
\end{lemma}

\begin{proof}
Use the pause-set injection from
\cref{lem:bc-c-twentieth-nonboundary-twenty-pause-code}, but count the target
even-column standard tableaux exactly by
\cref{lem:bc-even-column-standard-tableaux-fpf-count}.  For
\(q=21,22,23,24,25\), this gives the upper bound
\[
  \sum_{\substack{0\le p\le20\\ p\ {\rm even}}}
  \binom{2q+20}{p}(2q+19+p)!! .
\]
The exact C++/GMP/OpenMP replay
\[
\texttt{certificates/post\_m29/bc\_twentieth\_frontier\_gmp\_certificate.log}
\]
checks the displayed integer sums and verifies
\[
\begin{array}{c|c|c|c}
q & j & \mathrm{FPF}(q,j) & s_j-2\mathrm{FPF}(q,j)\\ \hline
21&62&
59640178498959305514120228939240938313534246071836194292902921336161643046875&
1752662321297424577290892639997668994851908081981447919742078064565734116724468645706\\
22&64&
10546296264465058052491781427837459696489984515461788143252233884332368579453125&
6955857123634547441821855269410040461348984985952520243786700467970376621406936907423606\\
23&66&
1857401551044838044392852676786302182793147532481339984781861148229496918789062500&
29386673296186905862659799277973863836030358812579483861024063304433700545101686805296669240\\
24&68&
326274573471185200824831553398486097764741139539400805519386410474790347556680468750&
131909085288574662554847144278149265223016674231510355370233483779079385857606782407378241090532\\
25&70&
57237304288628801488251204476412623107980093291037502641817121101346066389375667187500&
627984678371613795768240992643835579421151419539030578383154290148591005899825266194087500660974504 .
\end{array}
\]
The replay exits with status \(0\), reports \texttt{SUMMARY failures=0}, and
has SHA-256
\[
\hashsplit{431f9748488a0ab0453777cc498a2260}{996f1ddcf8b393545d30eaed4c291232}.
\]
\end{proof}

\begin{lemma}[The central binomial factor absorbs the twenty-pause loss]\label{lem:bc-twentieth-c-pause-binomial-absorption}
For every \(q\ge51\),
\[
\begin{aligned}
  &1048576\binom{2q+20}{0}
    +4194304\binom{2q+20}{2}
    +16777216\binom{2q+20}{4}\\
  &\quad+67108864\binom{2q+20}{6}
    +268435456\binom{2q+20}{8}
    +1073741824\binom{2q+20}{10}\\
  &\quad+4294967296\binom{2q+20}{12}
    +17179869184\binom{2q+20}{14}\\
  &\quad+68719476736\binom{2q+20}{16}
    +274877906944\binom{2q+20}{18}
    +1099511627776\binom{2q+20}{20}
  \le \binom{2q+20}{q}.
\end{aligned}
\]
\end{lemma}

\begin{proof}
At \(q=51\), the left side is
\[
  46832788955329562701706417826430976
  <
  74856689240948403197057601352346232
  =
  \binom{122}{51}.
\]
Let \(P(q)\) denote the displayed left side.  For \(1\le p\le20\),
\[
  \frac{\binom{2q+22}{p}}{\binom{2q+20}{p}}
  \le
  \frac{(2q+22)(2q+21)}{(2q+2)(2q+1)}
  <2
  \qquad(q\ge51),
\]
because \(4q^2-74q-458>0\).  The \(p=0\) summand has ratio \(1\), so
\(P(q+1)<2P(q)\).  On the other hand,
\[
  \frac{\binom{2q+22}{q+1}}{\binom{2q+20}{q}}
  =
  \frac{(2q+22)(2q+21)}{(q+1)(q+21)}
  >3,
\]
since the numerator minus three times the denominator is \(q^2+20q+399\).
Thus the binomial side grows by a larger factor than \(P(q)\), and induction
from \(q=51\) proves the inequality.
\end{proof}

\begin{lemma}[Exact twentieth $C$ pause-arithmetic checks after the first five rows]\label{lem:bc-c-twentieth-frontier-pause-arithmetic-certificate}
For \(26\le q\le50\),
put
\[
\begin{aligned}
  P(q)={}&1048576\binom{2q+20}{0}
    +4194304\binom{2q+20}{2}
    +16777216\binom{2q+20}{4}\\
  &+67108864\binom{2q+20}{6}
    +268435456\binom{2q+20}{8}
    +1073741824\binom{2q+20}{10}\\
  &+4294967296\binom{2q+20}{12}
    +17179869184\binom{2q+20}{14}\\
  &+68719476736\binom{2q+20}{16}
    +274877906944\binom{2q+20}{18}
    +1099511627776\binom{2q+20}{20}.
\end{aligned}
\]
Then
\[
  2^{2q+1}P(q)s_{q+20}\le s_{2q+20}.
\]
\end{lemma}

\begin{proof}
The exact C++/GMP/OpenMP replay
\[
\texttt{certificates/post\_m29/bc\_twentieth\_frontier\_gmp\_certificate.log}
\]
computes the displayed integer inequality for
\(26\le q\le50\).
It exits with status \(0\), reports \texttt{SUMMARY failures=0}, and has
SHA-256
\[
\hashsplit{431f9748488a0ab0453777cc498a2260}{996f1ddcf8b393545d30eaed4c291232}.
\]
It verifies the exact positive margins
\[
\begin{array}{c|c|c}
q & P(q) & s_{2q+20}-2^{2q+1}P(q)s_{q+20}\\ \hline
26&
354778401204155339898251902976&
2029321248636848132156229336256516289542273932570394327907502732535074464896495180780860172670218852864\\
27&
668046057206256226400690241536&
16470227858291556398868251953084445249011371801453753871162632928765387742332661304825544355924844558064128\\
28&
1233678087925221127336074674176&
94739387771535650729044632822516217732154802554740302342213832584467959027972591565288745136376182913113701376\\
29&
2236909402936097168961508999168&
562463859776645134721777408681391508043653821256254339964732794455721571638853258638459822769907471937211924648960\\
30&
3986640469434493642603832016896&
3510479956626479945064065093422685541292954668941382938283717075048742433937775444835212124123584113313494890413145088\\
31&
6990366999405411127331654729728&
23031455179892574060782309578898181395903222640563653328774210532276095875378614529615852789872015683095394932605117316096\\
32&
12070201653596990046658930671616&
158657761010391709928635933825879720014712852712073244608307670442668130031117226778978817473473585747790460038708177250920448\\
33&
20540316585644715652565573828608&
1146261222435806504222778451299208601166705583012514909657009562955583359670235855046916129914704307012836690340197544826148829184\\
34&
34475171214194326069382575292416&
8675754396576257580484489516247180010264966731523629612851197647094670914521681322218341646884687895051098503694139034331675589883904\\
35&
57110515561481856763043840524288&
68718405964867662561904265761192598509676134512816715127234748064104426081861403054347005524062565333813465621844609512483362290630735872\\
36&
93436643634283881488961284079616&
569039354428028048305125528256416769925482817745425010937525872867883145721910344150248856572253383726847704527847292505155439949686697775104\\
37&
151067376328878364756360343584768&
4921474373856201691215779968823141681781070217450565831646060308813907493185792261003037411457430418525374388607043173375760716448729015439552512\\
38&
241500944272106790013050091995136&
44415035728604130094111381673734028437943733224721232844576881782370385048998090673231228820108791056631562677240757154011985894251611947057038168064\\
39&
381933118070631065623644026699776&
417889612984311196145159757609244052940786047117277660521046813847855539611165828196138705286819557264874768639678794385676266980195285401264201818640384\\
40&
597842195094790018058475044601856&
4095628359235226169663528226641720275352265015692941967199710360533629510006236451097708164910355021586777313423004997897464440572239901608627322159762530304\\
41&
926644394250348431285693883351040&
41778449637156136111334507680682619050342599001186727105782802836361858860912943936296248029821397585233723366523096658006234188789438324382011452093718432587776\\
42&
1422822669341308738068607948292096&
443216814209338571157250886481406979474905746941980123757723379764527347403853309376043718891693511733571433672248196123874243820659577068008611816259655729215873024\\
43&
2165069297334773836476255541657600&
4886351313102937424311804637063514269142801222360408920750974420230475310776237230430790833270890535720334871326791916987350075999262451330911831063433982028819318087680\\
44&
3266162087164285945781253393350656&
55942579353574319299284500136393250019253802996105767741717827327750023959806505303037207821440900962673973606705673647141080068618186498072870306737127232347646415078277120\\
45&
4886527256556420549226075243151360&
664639403789427225561471728776951853213131700423439820015699549221259574648825594145736675540347744091130636990556912333149491891418212169359254452149986608480755678679176003584\\
46&
7252743326889686377601445345099776&
8188851318177335552057955122474484220669788144328372380064385876691784910041920651064776549594771008032396196789683388361323422691106803592343436618548320690259347370517086087921664\\
47&
10682627600811694024363297152696320&
104561339429957869976868377815387334966954930490003196027519571902672409696488366605908183254951799481771548678036940038653355370276412471889609121530153787850275475197490060488129101824\\
48&
15619041865727584281493799321993216&
1382796873366955667733698787057210090554400524921852847079418204016369853723477421040920207577688871959694107167306467123879148341547199311427242878603156125689114580950476133409875551748096\\
49&
22675183846767937430683325589094400&
18928751602053607333207431922902761406058558552419382632691632458421281957173923901607181230033775141234959749975745287719800608602511100883284529879312119407326828567600736159251147259345338368\\
50&
32694928538565825284208708373446656&
268045196852917110920858864292769236947110205860353984103918140593860870508962594303438435038934331300021389007635651385858812594502889092383930755683908658212795237567567873604653202227239233486848 .
\end{array}
\]
This proves the claim.
\end{proof}

\begin{proposition}[Residual $C$ twentieth nonboundary lower correction]\label{prop:bc-c-twentieth-nonboundary-lower-correction}
For every residual row \(C_b\) with \(20\le b\le217\), put \(q=b+1\) and
write \(s_j=m_j^{\mathrm{st}}\) and \(m_j^{C_b}=s_j+\delta_j^{C_b}\).  Then
\[
  -\frac12 s_{2q+20}\le \delta_{2q+20}^{C_b}.
\]
\end{proposition}

\begin{proof}
As in the proof of
\cref{prop:bc-c-first-nonboundary-lower-correction},
\(|\delta_{2q+20}^{C_b}|\) is the number of column-even Pieri two-strip paths
of length \(2q+20\) with terminal height at least \(2q\).  If
\(q=21,22,23,24,25\),
\cref{lem:bc-twentieth-c-q21-q22-q23-q24-q25-fpf-pause-certificate} gives
\[
  |\delta_{2q+20}^{C_b}|\le s_{2q+20}/2 .
\]
For \(q\ge26\), \cref{lem:bc-c-twentieth-nonboundary-twenty-pause-code},
with \(P(q)\) denoting the polynomial from
\cref{lem:bc-c-twentieth-frontier-pause-arithmetic-certificate}, gives
\[
  |\delta_{2q+20}^{C_b}|
  \le
  P(q)\,2^{2q}s_{q+20}.
\]
If \(26\le q\le50\),
\cref{lem:bc-c-twentieth-frontier-pause-arithmetic-certificate} gives
\[
  |\delta_{2q+20}^{C_b}|\le s_{2q+20}/2 .
\]
If \(q\ge51\), \cref{lem:bc-twentieth-c-pause-binomial-absorption} gives
\[
  |\delta_{2q+20}^{C_b}|
  \le
  2^{2q}\binom{2q+20}{q}s_{q+20}.
\]
The exact GMP/OpenMP replay
\[
\texttt{certificates/post\_m29/bc\_pieri\_powerloss\_2q\_gmp\_certificate.log}
\]
checks the inequality
\[
  2^{2q+1}\binom{2q+20}{q}s_{q+20}\le s_{2q+20}
\]
on this residual row and moment as part of its all-row residual-window audit;
it reports \texttt{failures=0} and has SHA-256
\[
\hashsplit{124068c5daa1123088a3cd79c3f6b839}{fee45a1aaef637ed82161c465d86bf45}.
\]
Thus \(|\delta_{2q+20}^{C_b}|\le s_{2q+20}/2\) also for \(q\ge51\), which
implies the claimed lower bound.
\end{proof}

\begin{lemma}[$C$ twenty-first-nonboundary height-bad paths have a twenty-one-pause code]\label{lem:bc-c-twentyfirst-nonboundary-twentyone-pause-code}
For \(q\ge1\), the number of column-even Pieri two-strip paths
\[
  \varnothing=\lambda^{(0)}\subset\lambda^{(1)}\subset\cdots
  \subset\lambda^{(2q+21)}
\]
with
\[
  \ell(\lambda^{(2q+21)})\ge2q
\]
is at most
\[
\begin{aligned}
  \biggl(&4194304\binom{2q+21}{1}
    +16777216\binom{2q+21}{3}
    +67108864\binom{2q+21}{5}\\
  &\quad+268435456\binom{2q+21}{7}
    +1073741824\binom{2q+21}{9}
    +4294967296\binom{2q+21}{11}\\
  &\quad+17179869184\binom{2q+21}{13}
    +68719476736\binom{2q+21}{15}\\
  &\quad+274877906944\binom{2q+21}{17}
    +1099511627776\binom{2q+21}{19}\\
  &\quad+4398046511104\binom{2q+21}{21}\biggr)2^{2q}s_{q+21}.
\end{aligned}
\]
\end{lemma}

\begin{proof}
The terminal height is even and at most \(2q+20\).  Thus it is
\[
  2q+20,\quad 2q+18,\quad 2q+16,\quad 2q+14,\quad 2q+12,
  \quad 2q+10,\quad 2q+8,\quad 2q+6,\quad 2q+4,\quad 2q+2,\quad\text{or}\quad 2q .
\]
These cases have respectively \(1,3,5,7,9,11,13,15,17,19,21\) steps which
fail to create a new row.  Deleting the first column, shifting left, and
standardizing the repeated pause labels injects the paths with \(p\) pauses
into the choice of the \(p\)-element pause set and an even-column standard
tableau of size \(2q+21+p\).  Therefore the eleven cases contribute at most
\[
  \binom{2q+21}{1}2^{2q+22}s_{q+11},
  \quad
  \binom{2q+21}{3}2^{2q+24}s_{q+12},
  \quad
  \binom{2q+21}{5}2^{2q+26}s_{q+13},
\]
\[
  \binom{2q+21}{7}2^{2q+28}s_{q+14},
  \quad
  \binom{2q+21}{9}2^{2q+30}s_{q+15},
  \quad
  \binom{2q+21}{11}2^{2q+32}s_{q+16},
\]
\[
  \binom{2q+21}{13}2^{2q+34}s_{q+17},
  \quad
  \binom{2q+21}{15}2^{2q+36}s_{q+18},
\]
\[
  \binom{2q+21}{17}2^{2q+38}s_{q+19},
  \quad
  \binom{2q+21}{19}2^{2q+40}s_{q+20},
  \quad
  \binom{2q+21}{21}2^{2q+42}s_{q+21},
\]
by
\cref{lem:bc-even-column-standard-tableaux-2n-bound,lem:bc-stable-moments-monotone-after-two}.
After replacing \(s_{q+11},\ldots,s_{q+20}\) by \(s_{q+21}\), the sum is the
displayed bound.
\end{proof}

\begin{lemma}[The first six twenty-first $C$ pause cases are absorbed by the fixed-point-free count]\label{lem:bc-twentyfirst-c-q21-q22-q23-q24-q25-q26-fpf-pause-certificate}
For \(21\le q\le26\), the number of column-even Pieri two-strip paths
of length \(2q+21\) with terminal height at least \(2q\) is at most
\(s_{2q+21}/2\).
\end{lemma}

\begin{proof}
Use the pause-set injection from
\cref{lem:bc-c-twentyfirst-nonboundary-twentyone-pause-code}, but count the
target even-column standard tableaux exactly by
\cref{lem:bc-even-column-standard-tableaux-fpf-count}.  For
\(21\le q\le26\), this gives the upper bound
\[
  \sum_{\substack{1\le p\le21\\ p\ {\rm odd}}}
  \binom{2q+21}{p}(2q+20+p)!! .
\]
The exact C++/GMP/OpenMP replay
\[
\texttt{certificates/post\_m29/bc\_twentyfirst\_frontier\_gmp\_certificate.log}
\]
checks the displayed integer sums and verifies
\[
\begin{array}{c|c|c|c}
q & j & \mathrm{FPF}(q,j) & s_j-2\mathrm{FPF}(q,j)\\ \hline
21&63&
14853304761259119316408268236115782519914732648875198936203636640975622669609375&
109541196083224653413419862679867047315851462790686570636936596922484754560617878901186\\
22&65&
2775167462979685571034677507459521918382520508108542832850330574595439382589453125&
448652122682546611772758810208437358844879789680883337773317887855684307303747347471137142\\
23&67&
515643299177238772508084006779623797418449131836770274856971010298098212396703125000&
1954211163650156659572045427442223570507950659374663963310697597266258145858750080198955492592\\
24&69&
95425704873602417500267836363618545745961310075692241308086946707711680690633486718750&
9035761316897761034491786420467483782802744796713310927043381316728348378812270461423173781686212\\
25&71&
17612308918214414042573132661483881464445143576594538854748006859427149621063875057812500&
44272870337208473463377756399122746412754928148091432553108900919798710002070327734342303523519802840\\
26&73&
3245715862664763040869103601604861498089285847566683990032167861739402848257414371356250000&
229499010251329482341848913823966491076126924769966584798543699831617057305383685339112476493976724203744 .
\end{array}
\]
The replay exits with status \(0\), reports \texttt{SUMMARY failures=0}, and
has SHA-256
\[
\hashsplit{a5593e2a97e9cadcf10f028940d40366}{500494e1c24910ccb6d39bf66334659c}.
\]
\end{proof}

\begin{lemma}[The central binomial factor absorbs the twenty-one-pause loss]\label{lem:bc-twentyfirst-c-pause-binomial-absorption}
For every \(q\ge54\),
\[
\begin{aligned}
  &4194304\binom{2q+21}{1}
    +16777216\binom{2q+21}{3}
    +67108864\binom{2q+21}{5}\\
  &\quad+268435456\binom{2q+21}{7}
    +1073741824\binom{2q+21}{9}
    +4294967296\binom{2q+21}{11}\\
  &\quad+17179869184\binom{2q+21}{13}
    +68719476736\binom{2q+21}{15}\\
  &\quad+274877906944\binom{2q+21}{17}
    +1099511627776\binom{2q+21}{19}
    +4398046511104\binom{2q+21}{21}
  \le \binom{2q+21}{q}.
\end{aligned}
\]
\end{lemma}

\begin{proof}
At \(q=54\), the left side is
\[
  3261219301280669347028169833976954880
  <
  8686004662321561726483889296581293760
  =
  \binom{129}{54}.
\]
Let \(P(q)\) denote the displayed left side.  For \(1\le p\le21\),
\[
  \frac{\binom{2q+23}{p}}{\binom{2q+21}{p}}
  \le
  \frac{(2q+23)(2q+22)}{(2q+2)(2q+1)}
  <2
  \qquad(q\ge54),
\]
because \(4q^2-82q-502>0\).  Hence \(P(q+1)<2P(q)\).  On the other hand,
\[
  \frac{\binom{2q+23}{q+1}}{\binom{2q+21}{q}}
  =
  \frac{(2q+23)(2q+22)}{(q+1)(q+22)}
  >3,
\]
since the numerator minus three times the denominator is \(q^2+21q+440\).
Thus the binomial side grows by a larger factor than \(P(q)\), and induction
from \(q=54\) proves the inequality.
\end{proof}

\begin{lemma}[Exact twenty-first $C$ pause-arithmetic checks after the first six rows]\label{lem:bc-c-twentyfirst-frontier-pause-arithmetic-certificate}
For \(27\le q\le53\),
put
\[
\begin{aligned}
  P(q)={}&4194304\binom{2q+21}{1}
    +16777216\binom{2q+21}{3}
    +67108864\binom{2q+21}{5}\\
  &+268435456\binom{2q+21}{7}
    +1073741824\binom{2q+21}{9}
    +4294967296\binom{2q+21}{11}\\
  &+17179869184\binom{2q+21}{13}
    +68719476736\binom{2q+21}{15}\\
  &+274877906944\binom{2q+21}{17}
    +1099511627776\binom{2q+21}{19}
    +4398046511104\binom{2q+21}{21}.
\end{aligned}
\]
Then
\[
  2^{2q+1}P(q)s_{q+21}\le s_{2q+21}.
\]
\end{lemma}

\begin{proof}
The exact C++/GMP/OpenMP replay
\[
\texttt{certificates/post\_m29/bc\_twentyfirst\_frontier\_gmp\_certificate.log}
\]
computes the displayed integer inequality for
\(27\le q\le53\).
It exits with status \(0\), reports \texttt{SUMMARY failures=0}, and has
SHA-256
\[
\hashsplit{a5593e2a97e9cadcf10f028940d40366}{500494e1c24910ccb6d39bf66334659c}.
\]
It verifies the exact positive margins
\[
\begin{array}{c|c|c}
q & P(q) & s_{2q+21}-2^{2q+1}P(q)s_{q+21}\\ \hline
27&
9575151515658781798713933168640&
985756847345365397668549487566928101202076121316509614773458892525721909334173324823800152168696519104905728\\
28&
18149753346596525461825218674688&
7158613480929762268448578947015233616259200760118484622452008356339934295745671832573210225646317695966635815936\\
29&
33757003245024294460845415464960&
44120544910882978721202463250197610593423144328107849325051457714354837298459106119318176500002942147786488526937088\\
30&
61673628973181416064980403355648&
282581280956550544245902388664810830472849239087675220981796874701933208175009681903504692235627848615039916457727139840\\
31&
110792957802818732910527066931200&
1900089231343788999457655862281293726410243695825185147085100958512450452135588410019209568532250630002737425851101313473536\\
32&
195884888100800055947048033189888&
13406570490101813922018918189286451995570875976146293397721077830650280881295367612048770932348077151989216336599429795824683008\\
33&
341140422629441482711784908390400&
99151535600921383911565449474296514088484211000218384693673187906296230264541841980058648203681457949706172485338995008571746772992\\
34&
585663004126789398027409805017088&
767803833895313764326235266814275193563047625851401244038831046389056682051747169658886346826304346253313066959885699612377886736359424\\
35&
991877484988091565409201459036160&
6219012485464160250675640067306720643382853351710774079025192427642045817909714769622082087040159484335381141565790416914558656285666170880\\
36&
1658265587248587282210004410564608&
52636114508647419356166723022649218158505616690899857107851748355726480996916376230376601818458115970203040338570196147029654984636923862102016\\
37&
2738452785992812795864484308582400&
465079114888081471940289531648985472257412313538510479133327670261543342836070296766671018659891071562531351537598771608274494691753346478791446528\\
38&
4469530554282533144521588330201088&
4286049101830464440406495420109395031352184405316405765232512254118723890648838508922916455277402916881167657458652528105893626761155031379100672000000\\
39&
7213685774082238473572296788279296&
41162110219623825932734233843601208328244031510942473643645999673220062414159622962599623642949315785137013910206388827775437555030430641183869871915929600\\
40&
11518838880023956802605225443590144&
411610493361212658609438075514938007634401958271237047730983098090314390859983980123617643307938811174731755930873962223083674536089370308191059126797667164160\\
41&
18206211659385783497600980125483008&
4282289551634348924329239706730221927448619999726851143512536260039107673812641791340906021596074618433248591945285593607784491053429855502204991712173465248161792\\
42&
28495746356420880173204428663291904&
46316141414884575423392030248557324373347325062095570344235381917801863515870883225873510751105223475181076803224145416339184981913216938954396074847474257008629284864\\
43&
44184327326044573051589598210162688&
520396248607814162823409340377204289423884505013387801432379105605759316422422772263344832852251178362441094223806886645553310780138514876615527628983685898646679493746688\\
44&
67897132543066529765526867231637504&
6069768027148561641469968204121477225392587834421204200319935034455051444628527600774726608600034576509190781323760816421657347083690230906723463358250976247744896370331303936\\
45&
103439570086625016614969015384670208&
73442633136598658548124371150402878780418925317719173754823781157768800231555088625820275366047569203753463140761080174217107996317324264158042584075646414540008862401898571612160\\
46&
156286648763455584613060296231092224&
921245524013205551953790793040788704469227867486519229414132271679590770819733900881763224757941771074895208163079441738001066902956981888865641346434818510776139017905718779864547328\\
47&
234258942103124593785807135318212608&
11972270293412322391994921931618335297324123090001871438912085148481002653952144852197350460078416907811449979858628414106791753929314838209815288040538475845166144128257893060512461209600\\
48&
348450347791182776912873805530005504&
161095796530604867861834758162350912768004867949482092683233794758296912819818344505858193204479649301413475234827580628029280764588690022696949547732290599270861392001864705314215920490545152\\
49&
514493642014577122423618856443445248&
2243056546216084944751807745638494487795511108402461697692038010895203957275122305700540804864135054948937866971399581912136423639155355415999523093805870238736018467778793229609840817849089032192\\
50&
754276652865419534823025158519783424&
32299439121457320747514179580679347285477123267840979907021820508147995918657351519455895498157028869409313165113431597572787105302564960829477825501250590157674270183153800297502833245176389900042240\\
51&
1098256307260626870214220222983831552&
480736573437627632444442336999710451279209290573111695742968139549979963787136574785586233584619561334575411074277790549878241774693597957150767990068780604458283993443317895087741144498332202068961492992\\
52&
1588561791218806248281669312079462400&
7391682389000089802339052812312253432891207688078572676092652130829134490225275535427940381802945635126455768704119046862514125047041133491380667739344626287931386744851356466887954376470515971521003007967232\\
53&
2283134001974655790967430957034897408&
117348456612109982890333740601748666125536761977221951143034700801029267699557908273501065501788081486464104857575988336392926173121182550684645021410313364752493141763532200784298388022736314368175040141366460416 .
\end{array}
\]
This proves the claim.
\end{proof}

\begin{proposition}[Residual $C$ twenty-first nonboundary lower correction]\label{prop:bc-c-twentyfirst-nonboundary-lower-correction}
For every residual row \(C_b\) with \(20\le b\le217\), put \(q=b+1\) and
write \(s_j=m_j^{\mathrm{st}}\) and \(m_j^{C_b}=s_j+\delta_j^{C_b}\).  Then
\[
  -\frac12 s_{2q+21}\le \delta_{2q+21}^{C_b}.
\]
\end{proposition}

\begin{proof}
As in the proof of
\cref{prop:bc-c-first-nonboundary-lower-correction},
\(|\delta_{2q+21}^{C_b}|\) is the number of column-even Pieri two-strip paths
of length \(2q+21\) with terminal height at least \(2q\).  If
\(21\le q\le26\),
\cref{lem:bc-twentyfirst-c-q21-q22-q23-q24-q25-q26-fpf-pause-certificate}
gives
\[
  |\delta_{2q+21}^{C_b}|\le s_{2q+21}/2 .
\]
For \(q\ge27\), \cref{lem:bc-c-twentyfirst-nonboundary-twentyone-pause-code},
with \(P(q)\) denoting the polynomial from
\cref{lem:bc-c-twentyfirst-frontier-pause-arithmetic-certificate}, gives
\[
  |\delta_{2q+21}^{C_b}|
  \le
  P(q)\,2^{2q}s_{q+21}.
\]
If \(27\le q\le53\),
\cref{lem:bc-c-twentyfirst-frontier-pause-arithmetic-certificate} gives
\[
  |\delta_{2q+21}^{C_b}|\le s_{2q+21}/2 .
\]
If \(q\ge54\), \cref{lem:bc-twentyfirst-c-pause-binomial-absorption} gives
\[
  |\delta_{2q+21}^{C_b}|
  \le
  2^{2q}\binom{2q+21}{q}s_{q+21}.
\]
The exact GMP/OpenMP replay
\[
\texttt{certificates/post\_m29/bc\_pieri\_powerloss\_2q\_gmp\_certificate.log}
\]
checks the inequality
\[
  2^{2q+1}\binom{2q+21}{q}s_{q+21}\le s_{2q+21}
\]
on this residual row and moment as part of its all-row residual-window audit;
it reports \texttt{failures=0} and has SHA-256
\[
\hashsplit{124068c5daa1123088a3cd79c3f6b839}{fee45a1aaef637ed82161c465d86bf45}.
\]
Thus \(|\delta_{2q+21}^{C_b}|\le s_{2q+21}/2\) also for \(q\ge54\), which
implies the claimed lower bound.
\end{proof}

\begin{lemma}[$C$ twenty-second-nonboundary height-bad paths have a twenty-two-pause code]\label{lem:bc-c-twentysecond-nonboundary-twentytwo-pause-code}
For \(q\ge1\), the number of column-even Pieri two-strip paths
\[
  \varnothing=\lambda^{(0)}\subset\lambda^{(1)}\subset\cdots
  \subset\lambda^{(2q+22)}
\]
with
\[
  \ell(\lambda^{(2q+22)})\ge2q
\]
is at most
\[
\begin{aligned}
  \biggl(&4194304
    +16777216\binom{2q+22}{2}
    +67108864\binom{2q+22}{4}\\
  &\quad+268435456\binom{2q+22}{6}
    +1073741824\binom{2q+22}{8}\\
  &\quad+4294967296\binom{2q+22}{10}
    +17179869184\binom{2q+22}{12}\\
  &\quad+68719476736\binom{2q+22}{14}
    +274877906944\binom{2q+22}{16}\\
  &\quad+1099511627776\binom{2q+22}{18}
    +4398046511104\binom{2q+22}{20}\\
  &\quad+17592186044416\binom{2q+22}{22}\biggr)2^{2q}s_{q+22}.
\end{aligned}
\]
\end{lemma}

\begin{proof}
The terminal height is even and at most \(2q+22\).  Thus it is
\[
  2q+22,\quad 2q+20,\quad 2q+18,\quad 2q+16,\quad 2q+14,
  \quad 2q+12,\quad 2q+10,\quad 2q+8,\quad 2q+6,
  \quad 2q+4,\quad 2q+2,\quad\text{or}\quad 2q .
\]
These cases have respectively \(0,2,4,6,8,10,12,14,16,18,20,22\) steps
which fail to create a new row.  Deleting the first column, shifting left, and
standardizing the repeated pause labels injects the paths with \(p\) pauses
into the choice of the \(p\)-element pause set and an even-column standard
tableau of size \(2q+22+p\).  Therefore the twelve cases contribute at most
\[
  \binom{2q+22}{0}2^{2q+22}s_{q+11},
  \quad
  \binom{2q+22}{2}2^{2q+24}s_{q+12},
  \quad
  \binom{2q+22}{4}2^{2q+26}s_{q+13},
\]
\[
  \binom{2q+22}{6}2^{2q+28}s_{q+14},
  \quad
  \binom{2q+22}{8}2^{2q+30}s_{q+15},
  \quad
  \binom{2q+22}{10}2^{2q+32}s_{q+16},
\]
\[
  \binom{2q+22}{12}2^{2q+34}s_{q+17},
  \quad
  \binom{2q+22}{14}2^{2q+36}s_{q+18},
\]
\[
  \binom{2q+22}{16}2^{2q+38}s_{q+19},
  \quad
  \binom{2q+22}{18}2^{2q+40}s_{q+20},
\]
\[
  \binom{2q+22}{20}2^{2q+42}s_{q+21},
  \quad
  \binom{2q+22}{22}2^{2q+44}s_{q+22},
\]
by
\cref{lem:bc-even-column-standard-tableaux-2n-bound,lem:bc-stable-moments-monotone-after-two}.
After replacing \(s_{q+11},\ldots,s_{q+21}\) by \(s_{q+22}\), the sum is the
displayed bound.
\end{proof}

\begin{lemma}[The first seven twenty-second $C$ pause cases are absorbed by the fixed-point-free count]\label{lem:bc-twentysecond-c-q21-q22-q23-q24-q25-q26-q27-fpf-pause-certificate}
For \(21\le q\le27\), the number of column-even Pieri two-strip
paths of length \(2q+22\) with terminal height at least \(2q\) is at most
\(s_{2q+22}/2\).
\end{lemma}

\begin{proof}
Use the pause-set injection from
\cref{lem:bc-c-twentysecond-nonboundary-twentytwo-pause-code}, but count the
target even-column standard tableaux exactly by
\cref{lem:bc-even-column-standard-tableaux-fpf-count}.  For
\(21\le q\le27\), this gives the upper bound
\[
  \sum_{\substack{0\le p\le22\\ p\ {\rm even}}}
  \binom{2q+22}{p}(2q+21+p)!! .
\]
The exact C++/GMP/OpenMP replay
\[
\texttt{certificates/post\_m29/bc\_twentysecond\_frontier\_gmp\_certificate.log}
\]
checks the displayed integer sums and verifies
\[
\begin{array}{c|c|c|c}
q & j & \mathrm{FPF}(q,j) & s_j-2\mathrm{FPF}(q,j)\\ \hline
21&64&
3673546597525622629835972748936048001604247648717791222500228955528081485766953125&
6955849797633944919506711702448105444927901170437191839127831754016933333908702532423606\\
22&66&
724448100514946445082843487720702486409131621639371094099504071412022564047695312500&
29386671851005507934856585201072593748197990359902535647244555074989280017515552547484169240\\
23&68&
141871669249161368402940347157334384670147654103190649625545483370186803889439492968750&
131909085005483873203466777873918234015144877086744529442930986139027191938182755323612616090532\\
24&70&
27634119982394172339974879118216636692043911713376867907724295133823961896204868831250000&
627984678316460030412029905566862323593670971401158715142982629338426049874380034534456514332849504\\
25&72&
5361259500752364359897462463162369710276578512029857877855952385609760942353222041075000000&
3165506839225007550037522624909651084680151266649452263508220962981924683953410568391326644961152767104\\
26&74&
1037279652297170971042023343593013599978416878321007651590281460717021843855347766405181250000&
16868160644363383382648743540649283907739197828775273879545440525301805222446792965089236201138755714508384\\
27&76&
200357862346821966866479160795239368809935057938118084156631161999245725604543961093357734375000&
94879010830044579516248919504777610596192027957411393212496424058368337060965458872883294753762971846510206800 .
\end{array}
\]
The replay exits with status \(0\), reports \texttt{SUMMARY failures=0}, and
has SHA-256
\[
\hashsplit{c2db04c00a2fcf719ada1acdb7dadc65}{e468bbaa1c87e3dc9a5b9acd34d653a4}.
\]
\end{proof}

\begin{lemma}[The central binomial factor absorbs the twenty-two-pause loss]\label{lem:bc-twentysecond-c-pause-binomial-absorption}
For every \(q\ge56\),
\[
\begin{aligned}
  &4194304
    +16777216\binom{2q+22}{2}
    +67108864\binom{2q+22}{4}\\
  &\quad+268435456\binom{2q+22}{6}
    +1073741824\binom{2q+22}{8}
    +4294967296\binom{2q+22}{10}\\
  &\quad+17179869184\binom{2q+22}{12}
    +68719476736\binom{2q+22}{14}\\
  &\quad+274877906944\binom{2q+22}{16}
    +1099511627776\binom{2q+22}{18}\\
  &\quad+4398046511104\binom{2q+22}{20}
    +17592186044416\binom{2q+22}{22}
  \le \binom{2q+22}{q}.
\end{aligned}
\]
\end{lemma}

\begin{proof}
At \(q=56\), the left side is
\[
  159401924255242975364410651900970532864
  <
  247522931562325802835549014143162387440
  =
  \binom{134}{56}.
\]
The margin in this base comparison is
\[
  88121007307082827471138362242191854576 .
\]
Let \(P(q)\) denote the displayed left side.  For \(1\le p\le22\),
\[
  \frac{\binom{2q+24}{p}}{\binom{2q+22}{p}}
  \le
  \frac{(2q+24)(2q+23)}{(2q+2)(2q+1)}
  <2
  \qquad(q\ge56),
\]
because \(4q^2-82q-548>0\).  The \(p=0\) summand is constant, so
\(P(q+1)<2P(q)\).  On the other hand,
\[
  \frac{\binom{2q+24}{q+1}}{\binom{2q+22}{q}}
  =
  \frac{(2q+24)(2q+23)}{(q+1)(q+23)}
  >3,
\]
since the numerator minus three times the denominator is \(q^2+22q+483\).
Thus the binomial side grows by a larger factor than \(P(q)\), and induction
from \(q=56\) proves the inequality.
\end{proof}

\begin{lemma}[Exact twenty-second $C$ pause-arithmetic checks after the first seven rows]\label{lem:bc-c-twentysecond-frontier-pause-arithmetic-certificate}
For \(28\le q\le55\),
put
\[
\begin{aligned}
  P(q)={}&4194304
    +16777216\binom{2q+22}{2}
    +67108864\binom{2q+22}{4}\\
  &+268435456\binom{2q+22}{6}
    +1073741824\binom{2q+22}{8}
    +4294967296\binom{2q+22}{10}\\
  &+17179869184\binom{2q+22}{12}
    +68719476736\binom{2q+22}{14}\\
  &+274877906944\binom{2q+22}{16}
    +1099511627776\binom{2q+22}{18}\\
  &+4398046511104\binom{2q+22}{20}
    +17592186044416\binom{2q+22}{22}.
\end{aligned}
\]
Then
\[
  2^{2q+1}P(q)s_{q+22}\le s_{2q+22}.
\]
\end{lemma}

\begin{proof}
The exact C++/GMP/OpenMP replay
\[
\texttt{certificates/post\_m29/bc\_twentysecond\_frontier\_gmp\_certificate.log}
\]
computes the displayed integer inequality for
\(28\le q\le55\).
It exits with status \(0\), reports \texttt{SUMMARY failures=0}, and has
SHA-256
\[
\hashsplit{c2db04c00a2fcf719ada1acdb7dadc65}{e468bbaa1c87e3dc9a5b9acd34d653a4}.
\]
It verifies the exact positive margins
\[
\begin{array}{c|c|c}
q & P(q) & s_{2q+22}-2^{2q+1}P(q)s_{q+22}\\ \hline
28&
258232049829169720156803397844992&
492349767652154619284693481786288236237095661454957232098550035795270121549384242662073768384508577547526308795392\\
29&
492497525342386521140182577577984&
3483466832503280899440696727311462501859329870562046674826263147602027156248074999507174034923285076608985660313680896\\
30&
922100309616885951378364492152832&
23021035901630077745090780235767623748628262348446900102731034684383355912289547211532393139652769256132053981105193794560\\
31&
1696598100048949236625274551926784&
158653734981690342641968115177593341474947050157045800007738266761342885201012596839444267358036127005508076604328738254071808\\
32&
3070553978397797843455115926700032&
1146259663079073097306649009879055070333578319532153953047054468587448893082794094141322410050692192209058271041018414794310498304\\
33&
5471032294063014254708346012565504&
8675753790859187255712452925703518902648658778434713496189514486069017491610845330889282049022569281121595780823032473306550480418816\\
34&
9604733587650220221817503248023552&
68718405728791263253484986922521028203919796920417966346542649695134802135107818487443516702903452389799510920930377584476819313273944064\\
35&
16625985218571243036882413471924224&
569039354335669849819008724579438469567142552046944983975232083374057359886892591671233129940426017607192272678504658484351539179176662028288\\
36&
28397123618411399190739068211167232&
4921474373819918838683817296764244584777652081219805826664610970133397418874362450249976850500456182072442333371476474691242392451199903922614272\\
37&
47887667885091847721553177182470144&
44415035728589812279713961893426735092172720379532192786694968595671509511430731568027943561845742821062085339556719266251982264347818846018985201664\\
38&
79779897210998404237517082352156672&
417889612984305518965438070040311560961682595139748428620620496601259413473948915178980011868502142614972383370048260382306062008566079544888620030562304\\
39&
131378453514526214976117152231194624&
4095628359235223907171693595308447469685570290256003419980443287620150570215116701754057883771055467620817324098051565097422202099071276405893028422027862016\\
40&
213963689433201229108023355066286080&
41778449637156135204880743980620049571726476214323010731457798253478559865629334122986823597490076283409222181556211534045305735822695321102779051188247054524416\\
41&
344787072400281922188260093936533504&
443216814209338570792073820551216354419973112758402700338092272454352348596724497056836886578504665717212600329057875964593464077596315238810447825914438752603643904\\
42&
549987880542188023764797796070195200&
4886351313102937424163844168436993005369033802655085360800939373298610015956158588367731328155910975754084907719311852731338450735790070239123350159963463854245030830080\\
43&
868821390393188226793894493084975104&
55942579353574319299224195792255582674764724022074526287177488596497496454085998525810084385577411281816539885431999477349922235188568687166823504138484803217270484688027648\\
44&
1359739857343091889785609376644464640&
664639403789427225561447001257947088689596883194726062726764916014453948769926032949660602233088806325443799349840701811653413945303865214746461812035408280020595571965659398144\\
45&
2109071981931189293917744295831404544&
8188851318177335552057944920067218736961260629180059375073073831568531584640283434479098811517956149189419564211439790004871185452306550851967206677127407892155854900611086429962240\\
46&
3243321278808598988780603540158545920&
104561339429957869976868373579255172011486290623368195306618721369753278347677539018842384040878915322483438369998714641325618451290557286305762968418028605245083618700982482883560587264\\
47&
4946470774020608634677247868261105664&
1382796873366955667733698785286980074157267178875875009827842356900495714534498979670691906462326499095594377512415918035149174228171124814684607860324735857888746044981734609236485822119936\\
48&
7484168886498429501476337523349258240&
18928751602053607333207431922158155511380168384351702050049851409072414021915098555592542513334707916076451323223062626090117125759494679195067731890424030322737468780917517426791668431713435648\\
49&
11237315081066759481974355165471309824&
268045196852917110920858864292453955721628047290841296537909620843269638828018884291466095874231486883221822347063069993588841213358044098104087940256347085193274303754733181433547765556279492182016\\
50&
16748409447373062556011217099654430720&
3924381011990010004171741527608164978530497770328942545279534321141635845996580399288371378142292912158949545428989308238780575834481797389619120737978038693740120789100928933597441050984972807914356736\\
51&
24785135282172804514497676198302187520&
59370954705574675743610932769466315206466503907905141886507701961765487474125475516394128490557956737635659706359250187603604941972386279628758590982808875081565464849590043216124284499591297379482818379776\\
52&
36427080369369564745217083061571682304&
927655959519066445608888540023467196375398960443697367332179503419659131760171301641400034247109897251509193238322104155202394251583784209240902686384555013003831791878095058868402456183887015322265348665704448\\
53&
53183361546847072514389722488322916352&
14961925445802499547156199139453750585162613576092731610569356185714633640126882589729405206798510690213251303787180327029784781597086269309420092302838918018746668809333889373083060424970909729714224614870077210624\\
54&
77151311285441864814811229339500150784&
248977571795961719585950638301180814987846650628626904010669237221170333286101822758486894115008628284986898588281356442069723383176481067975883126730727507661999142058502771176147292000203655218840591177770868374831104\\
55&
111229454697260018616582403493170511872&
4272640406230412894616231425874304234341228457638137625507562967790043849931104883915472404656047215926768659250609166161501882550473575298039440765850867027979580999584752754517841396553135819466820130913129644173222608896 .
\end{array}
\]
This proves the claim.
\end{proof}

\begin{proposition}[Residual $C$ twenty-second nonboundary lower correction]\label{prop:bc-c-twentysecond-nonboundary-lower-correction}
For every residual row \(C_b\) with \(20\le b\le217\), put \(q=b+1\) and
write \(s_j=m_j^{\mathrm{st}}\) and \(m_j^{C_b}=s_j+\delta_j^{C_b}\).  Then
\[
  -\frac12 s_{2q+22}\le \delta_{2q+22}^{C_b}.
\]
\end{proposition}

\begin{proof}
As in the proof of
\cref{prop:bc-c-first-nonboundary-lower-correction},
\(|\delta_{2q+22}^{C_b}|\) is the number of column-even Pieri two-strip paths
of length \(2q+22\) with terminal height at least \(2q\).  If
\(21\le q\le27\),
\cref{lem:bc-twentysecond-c-q21-q22-q23-q24-q25-q26-q27-fpf-pause-certificate}
gives
\[
  |\delta_{2q+22}^{C_b}|\le s_{2q+22}/2 .
\]
For \(q\ge28\), \cref{lem:bc-c-twentysecond-nonboundary-twentytwo-pause-code},
with \(P(q)\) denoting the polynomial from
\cref{lem:bc-c-twentysecond-frontier-pause-arithmetic-certificate}, gives
\[
  |\delta_{2q+22}^{C_b}|
  \le
  P(q)\,2^{2q}s_{q+22}.
\]
If \(28\le q\le55\),
\cref{lem:bc-c-twentysecond-frontier-pause-arithmetic-certificate} gives
\[
  |\delta_{2q+22}^{C_b}|\le s_{2q+22}/2 .
\]
If \(q\ge56\), \cref{lem:bc-twentysecond-c-pause-binomial-absorption} gives
\[
  |\delta_{2q+22}^{C_b}|
  \le
  2^{2q}\binom{2q+22}{q}s_{q+22}.
\]
The exact GMP/OpenMP replay
\[
\texttt{certificates/post\_m29/bc\_pieri\_powerloss\_2q\_gmp\_certificate.log}
\]
checks the inequality
\[
  2^{2q+1}\binom{2q+22}{q}s_{q+22}\le s_{2q+22}
\]
on this residual row and moment as part of its all-row residual-window audit;
it reports \texttt{failures=0} and has SHA-256
\[
\hashsplit{124068c5daa1123088a3cd79c3f6b839}{fee45a1aaef637ed82161c465d86bf45}.
\]
Thus \(|\delta_{2q+22}^{C_b}|\le s_{2q+22}/2\) also for \(q\ge56\), which
implies the claimed lower bound.
\end{proof}

\begin{lemma}[$C$ twenty-third-nonboundary height-bad paths have a twenty-three-pause code]\label{lem:bc-c-twentythird-nonboundary-twentythree-pause-code}
For \(q\ge1\), the number of column-even Pieri two-strip paths
\[
  \varnothing=\lambda^{(0)}\subset\lambda^{(1)}\subset\cdots
  \subset\lambda^{(2q+23)}
\]
with
\[
  \ell(\lambda^{(2q+23)})\ge2q
\]
is at most
\[
\begin{aligned}
  \biggl(&16777216\binom{2q+23}{1}
    +67108864\binom{2q+23}{3}
    +268435456\binom{2q+23}{5}\\
  &\quad+1073741824\binom{2q+23}{7}
    +4294967296\binom{2q+23}{9}\\
  &\quad+17179869184\binom{2q+23}{11}
    +68719476736\binom{2q+23}{13}\\
  &\quad+274877906944\binom{2q+23}{15}
    +1099511627776\binom{2q+23}{17}\\
  &\quad+4398046511104\binom{2q+23}{19}
    +17592186044416\binom{2q+23}{21}\\
  &\quad+70368744177664\binom{2q+23}{23}\biggr)2^{2q}s_{q+23}.
\end{aligned}
\]
\end{lemma}

\begin{proof}
The terminal height is even and at most \(2q+22\).  Thus it is
\[
  2q+22,\quad 2q+20,\quad 2q+18,\quad 2q+16,\quad 2q+14,
  \quad 2q+12,\quad 2q+10,\quad 2q+8,\quad 2q+6,
  \quad 2q+4,\quad 2q+2,\quad\text{or}\quad 2q .
\]
These cases have respectively \(1,3,5,7,9,11,13,15,17,19,21,23\) steps
which fail to create a new row.  Deleting the first column, shifting left, and
standardizing the repeated pause labels injects the paths with \(p\) pauses
into the choice of the \(p\)-element pause set and an even-column standard
tableau of size \(2q+23+p\).  Therefore the twelve cases contribute at most
\[
  \binom{2q+23}{1}2^{2q+24}s_{q+12},
  \quad
  \binom{2q+23}{3}2^{2q+26}s_{q+13},
  \quad
  \binom{2q+23}{5}2^{2q+28}s_{q+14},
\]
\[
  \binom{2q+23}{7}2^{2q+30}s_{q+15},
  \quad
  \binom{2q+23}{9}2^{2q+32}s_{q+16},
  \quad
  \binom{2q+23}{11}2^{2q+34}s_{q+17},
\]
\[
  \binom{2q+23}{13}2^{2q+36}s_{q+18},
  \quad
  \binom{2q+23}{15}2^{2q+38}s_{q+19},
\]
\[
  \binom{2q+23}{17}2^{2q+40}s_{q+20},
  \quad
  \binom{2q+23}{19}2^{2q+42}s_{q+21},
\]
\[
  \binom{2q+23}{21}2^{2q+44}s_{q+22},
  \quad
  \binom{2q+23}{23}2^{2q+46}s_{q+23},
\]
by
\cref{lem:bc-even-column-standard-tableaux-2n-bound,lem:bc-stable-moments-monotone-after-two}.
After replacing \(s_{q+12},\ldots,s_{q+22}\) by \(s_{q+23}\), the sum is the
displayed bound.
\end{proof}

\begin{lemma}[The first eight twenty-third $C$ pause cases are absorbed by the fixed-point-free count]\label{lem:bc-twentythird-c-q21-q22-q23-q24-q25-q26-q27-q28-fpf-pause-certificate}
For \(21\le q\le28\), the number of column-even Pieri two-strip
paths of length \(2q+23\) with terminal height at least \(2q\) is at most
\(s_{2q+23}/2\).
\end{lemma}

\begin{proof}
Use the pause-set injection from
\cref{lem:bc-c-twentythird-nonboundary-twentythree-pause-code}, but count the
target even-column standard tableaux exactly by
\cref{lem:bc-even-column-standard-tableaux-fpf-count}.  For
\(21\le q\le28\), this gives the upper bound
\[
  \sum_{\substack{1\le p\le23\\ p\ {\rm odd}}}
  \binom{2q+23}{p}(2q+22+p)!! .
\]
The exact C++/GMP/OpenMP replay
\[
\texttt{certificates/post\_m29/bc\_twentythird\_frontier\_gmp\_certificate.log}
\]
checks the displayed integer sums and verifies
\[
\begin{array}{c|c|c|c}
q & j & \mathrm{FPF}(q,j) & s_j-2\mathrm{FPF}(q,j)\\ \hline
21&65&
903395458925233862459538367562195635591789894032637801584533085396134624499776953125&
448650321441963687264405033201057249372652442866136288714800384490174664225377113096137142\\
22&67&
187855136252296598892955089953165371907925245701375187830773864045211057311947484375000&
1954210788971170753333325186548211677736454438361070824233862485432472076032831881097392992592\\
23&69&
38737318451337022431677803596238118856006496549152483169175211416709091377624641299218750&
9035761239613975541564946392112412263053744176192240448889799460994099438809511067555158156686212\\
24&71&
7934575627332601543808973311549601252211573686508191850366035064164098120241572629940625000&
44272870321374546826548982139589946054978693406597175467245706296776135887460985793101286013754177840\\
25&73&
1616755154291510388197995200467896602543581622197201447085467210087604018787127794570575000000&
229499010248102463464991222573652238882394341287875600125844430305426187220931956107234735733578286703744\\
26&75&
328136620018632690136818267699230899854470272882286110463320468592371808038005330891447656250000&
1256676780468210769607279819057238027387173281802542292087129519327342137755360759306298359676288271928249568\\
27&77&
66413080832039621078408360898635664063082364466853853287713048631560118881554333641485725859375000&
7258238000562478011590159605511770781209237988343440515696790130395154216165069650698745336790968524649505039568\\
28&79&
13417863213593947455210578777994636773263119312159516026950171056539369741242161635850168708281250000&
44157232836692673989337327386886560837315346222833166049437224097037484517279315465062363755087793626747029988096608 .
\end{array}
\]
The replay exits with status \(0\), reports \texttt{SUMMARY failures=0}, and
has SHA-256
\[
\hashsplit{e095f8a87dbf4718ab5ac066d438ca1b}{f6b32452d617172d90acd926094f717d}.
\]
\end{proof}

\begin{lemma}[The central binomial factor absorbs the twenty-three-pause loss]\label{lem:bc-twentythird-c-pause-binomial-absorption}
For every \(q\ge59\),
\[
\begin{aligned}
  &16777216\binom{2q+23}{1}
    +67108864\binom{2q+23}{3}
    +268435456\binom{2q+23}{5}\\
  &\quad+1073741824\binom{2q+23}{7}
    +4294967296\binom{2q+23}{9}
    +17179869184\binom{2q+23}{11}\\
  &\quad+68719476736\binom{2q+23}{13}
    +274877906944\binom{2q+23}{15}\\
  &\quad+1099511627776\binom{2q+23}{17}
    +4398046511104\binom{2q+23}{19}\\
  &\quad+17592186044416\binom{2q+23}{21}
    +70368744177664\binom{2q+23}{23}
  \le \binom{2q+23}{q}.
\end{aligned}
\]
\end{lemma}

\begin{proof}
At \(q=59\), the left side is
\[
  11127725043353600952651764558062923284480
  <
  28792469044176159473161828154011690184400
  =
  \binom{141}{59}.
\]
The margin in this base comparison is
\[
  17664744000822558520510063595948766899920 .
\]
Let \(P(q)\) denote the displayed left side.  For \(1\le p\le23\),
\[
  \frac{\binom{2q+25}{p}}{\binom{2q+23}{p}}
  \le
  \frac{(2q+25)(2q+24)}{(2q+2)(2q+1)}
  <2
  \qquad(q\ge59),
\]
because \(4q^2-86q-596>0\).  Hence \(P(q+1)<2P(q)\).  On the other hand,
\[
  \frac{\binom{2q+25}{q+1}}{\binom{2q+23}{q}}
  =
  \frac{(2q+25)(2q+24)}{(q+1)(q+24)}
  >3,
\]
since the numerator minus three times the denominator is \(q^2+23q+528\).
Thus the binomial side grows by a larger factor than \(P(q)\), and induction
from \(q=59\) proves the inequality.
\end{proof}

\begin{lemma}[Exact twenty-third $C$ pause-arithmetic checks after the first eight rows]\label{lem:bc-c-twentythird-frontier-pause-arithmetic-certificate}
For \(29\le q\le58\),
put
\[
\begin{aligned}
  P(q)={}&16777216\binom{2q+23}{1}
    +67108864\binom{2q+23}{3}
    +268435456\binom{2q+23}{5}\\
  &+1073741824\binom{2q+23}{7}
    +4294967296\binom{2q+23}{9}
    +17179869184\binom{2q+23}{11}\\
  &+68719476736\binom{2q+23}{13}
    +274877906944\binom{2q+23}{15}\\
  &+1099511627776\binom{2q+23}{17}
    +4398046511104\binom{2q+23}{19}\\
  &+17592186044416\binom{2q+23}{21}
    +70368744177664\binom{2q+23}{23}.
\end{aligned}
\]
Then
\[
  2^{2q+1}P(q)s_{q+23}\le s_{2q+23}.
\]
\end{lemma}

\begin{proof}
The exact C++/GMP/OpenMP replay
\[
\texttt{certificates/post\_m29/bc\_twentythird\_frontier\_gmp\_certificate.log}
\]
computes the displayed integer inequality for
\(29\le q\le58\).
It exits with status \(0\), reports \texttt{SUMMARY failures=0}, and has
SHA-256
\[
\hashsplit{e095f8a87dbf4718ab5ac066d438ca1b}{f6b32452d617172d90acd926094f717d}.
\]
It verifies the exact positive margins
\[
\begin{array}{c|c|c}
q & P(q) & s_{2q+23}-2^{2q+1}P(q)s_{q+23}\\ \hline
29&
6959847218612122581155773302177792&
262922528006441399142918828714267913476240941026880200453666150140515503922968654849511691381689387860747668371565223936\\
30&
13349871631976307535143002785710080&
1892170132553195255385066987251681003257253233938376946794730614650933439139381536060337297887916915077038793292293500804096\\
31&
25149969107522801701576766130225152&
13403377699060053823326297425495888171169641624155236894853080179073898100276887289882919721260803741044277077572913129163737088\\
32&
46580242725092383554064925012787200&
99150246432412329359110239867329923011459192488748269341318123080059828273442363353162494250882430439994970226462670807207920494592\\
33&
84890254315525433868343134579064832&
767803312298115206380124590668612460791274271862995854727843082867091063826624240691954764532928491414937207520650052984921343060713472\\
34&
152357514950698565274444603682979840&
6219012273889054177034092486988637274704739875148842753030412471980785838598511677571292654154841980142843945402545248304669930662407690240\\
35&
269495356939459044522956688113795072&
52636114422569219004636129831169799132744683673390939573364040719629053243787870968559597926987305851214436239395378243551911491927636317351936\\
36&
470140096490772841533509163234099200&
465079114852941847109337451532761244551657005922639408120056162898473524839335483180893314956745356991369976299307356592016696596429156918280572928\\
37&
809426194995622618572102831756541952&
4286049101816065352548816412291165840736929919932053033562139352115348522378338693864707924050357567766085932781402410495866424166034454015541524365312\\
38&
1376152883815089322159894077290577920&
41162110219617901459938343912727016995032964405718622103152182328431309860201648692830782045089087510021594915880343847359751470524721781452447226109759488\\
39&
2311761539545983226391674546527666176&
411610493361210210278036077216449246138717985604567682817755037715590356945504900132723968061308007186773211343651116355590493055514696891731003676472012726272\\
40&
3839169632931209750414177779960512512&
4282289551634347907816754155702309537420341739733932749455332650336185199996541562922057933792687213828764334222267524154107507646881537891794193660483552093888512\\
41&
6306166825583292591958642849445052416&
46316141414884574999278514354751190723683581587119411306516820118331046817057041749061927161393316662469528250007192849882506291454443207541753731276688588647892680704\\
42&
10250137928218062699305754359393943552&
520396248607814162645551522457218886000865448647499950977610131637324894849314023735577938113626570881665247771377589166129637712336954223196431663672309527955139077545984\\
43&
16493759814146855344831476732711141376&
6069768027148561641394983873006235848758633131329060018188780965648660991304548118709174862174951542140599406612359545027161788828471247303697820683932754176760780541456760832\\
44&
26285324552202525986867921870160658432&
73442633136598658548092583997206106244129174619756220499063477959640329719278812421882687208268535534747162966382810056474726357803539373668862114529391698886709427292551794900992\\
45&
41502867286266040072370371990657171456&
921245524013205551953777241599390509538033643838635537858724440216766205151030737106199475416558020952643195672431491630645201243760244807930963838362248003732892803572924638528667648\\
46&
64948851136519542628806235572228063232&
11972270293412322391994916120839570364588466494434545579204000751779310045645561126319238554558813958457046659734240577640501538048460257507585892239217486398860293974401087078667404820480\\
47&
100772474255034418205651304353225506816&
161095796530604867861834755655929337116832809482232977447217696241572720858675336379385221376004734206837014945705637570677799870025510126092123324762454928914784168526611307414727296326598656\\
48&
155070619782362929786856023651539484672&
2243056546216084944751807744550829821437083443578837475650888116534031151963895999913280925727581331323280352358389887252881610986210833932846662174401863431525322250557906298994635856291922608128\\
49&
236737242158663771683872963272214839296&
32299439121457320747514179580204447478536294093536632919278349022770251264859372565817296877455521722053168509855264311989953107779487328574724529339707177671473609272332813331451421920588725507719168\\
50&
358656091291801559574794297533133750272&
480736573437627632444442336999501802805755996801583315801384527688198597503928749544131647440241477723182700554271834130032342822099517110352662739530519037412643531988839667104360900947678056046412267520\\
51&
539365072633013965767348183551809945600&
7391682389000089802339052812312161182018208476891540798328492101603557949501945206358782661892505978036863459971499597026992005569884320157167074154459232558304870357220376082129956691372000203058185383903232\\
52&
805364726928469583756578139051112005632&
117348456612109982890333740601748625076815645934981681377185477259115506924646139921757826444013957158394337167279004355581966235040023322792971525788818976101815076207701796976685813626706130492474124166910771200\\
53&
1194301475271651671216217023645136977920&
1922607071870099047776047162593793145592773628626072961895351827009711892130090998431432884637000474239943340163933837619543574148599943214272620603215396738506279148253608052534802540106562918916934564589422187118592\\
54&
1759332456734695880182256141017506906112&
32491567507955082463681564380506379353319468845032910194418789643177304843769025445998468879835810297202419594889959121319796925928702813330102984652722155794792947323374944189599248658270279058688799122755945889615904768\\
55&
2575078096059693015444602118006783344640&
566124760447716687947008385494882991678864446076761149394571930976456960410609033369704293260812043563934665210799271341102690487997456224691017721311838704364558034715785658390916552648533894078965897569040137338076385378304\\
56&
3745697405479725315540839646260669448192&
10165191425994752876422337440716385533946880872422015491545245926867342861225075660298553306122004331610801138552774695576281167402950706831210383230666808778113818700821527377373669236113759299023385134796437454437422956969066496\\
57&
5415787503093167348264225623949610844160&
188012781990683033386658460821740358242778748176665131804483924314196340966432997952525302602383176767503528742780793537648809060014067002204485854529494320580445704666997291627966385065197699421166550718325734760740888836933251760128\\
58&
7785022976611373546066487491729788239872&
3580467372446220814152304752927059603961612402007422412364145090431399938181342628210340495226086617922725443098177687766978775250286104160884924615340097997872838051027943034756480292897594729869171895988990327969980597424148683188994048 .
\end{array}
\]
This proves the claim.
\end{proof}

\begin{proposition}[Residual $C$ twenty-third nonboundary lower correction]\label{prop:bc-c-twentythird-nonboundary-lower-correction}
For every residual row \(C_b\) with \(20\le b\le217\), put \(q=b+1\) and
write \(s_j=m_j^{\mathrm{st}}\) and \(m_j^{C_b}=s_j+\delta_j^{C_b}\).  Then
\[
  -\frac12 s_{2q+23}\le \delta_{2q+23}^{C_b}.
\]
\end{proposition}

\begin{proof}
As in the proof of
\cref{prop:bc-c-first-nonboundary-lower-correction},
\(|\delta_{2q+23}^{C_b}|\) is the number of column-even Pieri two-strip paths
of length \(2q+23\) with terminal height at least \(2q\).  If
\(21\le q\le28\),
\cref{lem:bc-twentythird-c-q21-q22-q23-q24-q25-q26-q27-q28-fpf-pause-certificate}
gives
\[
  |\delta_{2q+23}^{C_b}|\le s_{2q+23}/2 .
\]
For \(q\ge29\), \cref{lem:bc-c-twentythird-nonboundary-twentythree-pause-code},
with \(P(q)\) denoting the polynomial from
\cref{lem:bc-c-twentythird-frontier-pause-arithmetic-certificate}, gives
\[
  |\delta_{2q+23}^{C_b}|
  \le
  P(q)\,2^{2q}s_{q+23}.
\]
If \(29\le q\le58\),
\cref{lem:bc-c-twentythird-frontier-pause-arithmetic-certificate} gives
\[
  |\delta_{2q+23}^{C_b}|\le s_{2q+23}/2 .
\]
If \(q\ge59\), \cref{lem:bc-twentythird-c-pause-binomial-absorption} gives
\[
  |\delta_{2q+23}^{C_b}|
  \le
  2^{2q}\binom{2q+23}{q}s_{q+23}.
\]
The exact GMP/OpenMP replay
\[
\texttt{certificates/post\_m29/bc\_pieri\_powerloss\_2q\_gmp\_certificate.log}
\]
checks the inequality
\[
  2^{2q+1}\binom{2q+23}{q}s_{q+23}\le s_{2q+23}
\]
on this residual row and moment as part of its all-row residual-window audit;
it reports \texttt{failures=0} and has SHA-256
\[
\hashsplit{124068c5daa1123088a3cd79c3f6b839}{fee45a1aaef637ed82161c465d86bf45}.
\]
Thus \(|\delta_{2q+23}^{C_b}|\le s_{2q+23}/2\) also for \(q\ge59\), which
implies the claimed lower bound.
\end{proof}

\begin{lemma}[$C$ twenty-fourth-nonboundary height-bad paths have a twenty-four-pause code]\label{lem:bc-c-twentyfourth-nonboundary-twentyfour-pause-code}
For \(q\ge1\), the number of column-even Pieri two-strip paths
\[
  \varnothing=\lambda^{(0)}\subset\lambda^{(1)}\subset\cdots
  \subset\lambda^{(2q+24)}
\]
with
\[
  \ell(\lambda^{(2q+24)})\ge2q
\]
is at most
\[
\begin{aligned}
  \biggl(&16777216\binom{2q+24}{0}
    +67108864\binom{2q+24}{2}
    +268435456\binom{2q+24}{4}\\
  &\quad+1073741824\binom{2q+24}{6}
    +4294967296\binom{2q+24}{8}\\
  &\quad+17179869184\binom{2q+24}{10}
    +68719476736\binom{2q+24}{12}\\
  &\quad+274877906944\binom{2q+24}{14}
    +1099511627776\binom{2q+24}{16}\\
  &\quad+4398046511104\binom{2q+24}{18}
    +17592186044416\binom{2q+24}{20}\\
  &\quad+70368744177664\binom{2q+24}{22}
    +281474976710656\binom{2q+24}{24}\biggr)2^{2q}s_{q+24}.
\end{aligned}
\]
\end{lemma}

\begin{proof}
The terminal height is even and at most \(2q+24\).  Thus it is
\[
  2q+24,\quad 2q+22,\quad 2q+20,\quad 2q+18,\quad 2q+16,
  \quad 2q+14,\quad 2q+12,\quad 2q+10,\quad 2q+8,
  \quad 2q+6,\quad 2q+4,\quad 2q+2,\quad\text{or}\quad 2q .
\]
These cases have respectively \(0,2,4,6,8,10,12,14,16,18,20,22,24\) steps
which fail to create a new row.  Deleting the first column, shifting left, and
standardizing the repeated pause labels injects the paths with \(p\) pauses
into the choice of the \(p\)-element pause set and an even-column standard
tableau of size \(2q+24+p\).  Therefore the thirteen cases contribute at most
\[
  \binom{2q+24}{0}2^{2q+24}s_{q+12},
  \quad
  \binom{2q+24}{2}2^{2q+26}s_{q+13},
  \quad
  \binom{2q+24}{4}2^{2q+28}s_{q+14},
\]
\[
  \binom{2q+24}{6}2^{2q+30}s_{q+15},
  \quad
  \binom{2q+24}{8}2^{2q+32}s_{q+16},
  \quad
  \binom{2q+24}{10}2^{2q+34}s_{q+17},
\]
\[
  \binom{2q+24}{12}2^{2q+36}s_{q+18},
  \quad
  \binom{2q+24}{14}2^{2q+38}s_{q+19},
\]
\[
  \binom{2q+24}{16}2^{2q+40}s_{q+20},
  \quad
  \binom{2q+24}{18}2^{2q+42}s_{q+21},
\]
\[
  \binom{2q+24}{20}2^{2q+44}s_{q+22},
  \quad
  \binom{2q+24}{22}2^{2q+46}s_{q+23},
  \quad
  \binom{2q+24}{24}2^{2q+48}s_{q+24},
\]
by
\cref{lem:bc-even-column-standard-tableaux-2n-bound,lem:bc-stable-moments-monotone-after-two}.
After replacing \(s_{q+12},\ldots,s_{q+23}\) by \(s_{q+24}\), the sum is the
displayed bound.
\end{proof}

\begin{lemma}[The first nine twenty-fourth $C$ pause cases are absorbed by the fixed-point-free count]\label{lem:bc-twentyfourth-c-q21-q22-q23-q24-q25-q26-q27-q28-q29-fpf-pause-certificate}
For \(21\le q\le29\), the number of column-even Pieri two-strip
paths of length \(2q+24\) with terminal height at least \(2q\) is at most
\(s_{2q+24}/2\).
\end{lemma}

\begin{proof}
Use the pause-set injection from
\cref{lem:bc-c-twentyfourth-nonboundary-twentyfour-pause-code}, but count the
target even-column standard tableaux exactly by
\cref{lem:bc-even-column-standard-tableaux-fpf-count}.  For
\(21\le q\le29\), this gives the upper bound
\[
  \sum_{\substack{0\le p\le24\\ p\ {\rm even}}}
  \binom{2q+24}{p}(2q+23+p)!! .
\]
The exact C++/GMP/OpenMP replay
\[
\texttt{certificates/post\_m29/bc\_twentyfourth\_frontier\_gmp\_certificate.log}
\]
checks the displayed integer sums and verifies
\[
\begin{array}{c|c|c|c}
q & j & \mathrm{FPF}(q,j) & s_j-2\mathrm{FPF}(q,j)\\ \hline
21&66&
221151264435901656238544121628617366307975586291767907173573858313983738121479335937500&
29386230997372837161436998278516311954870347226993195390172396126280794874084437684202919240\\
22&68&
48444237602328436386124860047000459772495816702949179309459451070871911425056949258593750&
131908988400752007044916742430078834328894101435406431750953666471216016934733512988593084840532\\
23&70&
10509247527460596851185631453984174845850214095312956407091681807621698204689554487971875000&
627984657353233215455624547875549173861754553084818347943823550970511024898631548947757276051599504\\
24&72&
2261695805248852696343321513446275820373208213247535091658572763717440987569198217182676562500&
3165506834712338458541321960942802982713924365323588994037210495420491061289748115137636654677949642104\\
25&74&
483611060862391083361785820148019768484334416474890938047383312006436030727572497310988189062500&
16868160643398235820228555715867796314130344319006561880352300664510219161355354947321801902049589698883384\\
26&76&
102883760855128117126028139996669303388886138139515950558646075783774976009571960057627366523437500&
94879010829839212710263356914459287274520279829371240806333268393419358171721907412315359921569903828932081800\\
27&78&
21802604109980250245943707808188627322588633384955492241589185041273962279804658011997949891132812500&
562512973617538987809161181648947151950133441437868744026688123560159091712555975513054430372461930273937912609368\\
28&80&
4607287018232315624867407379250784363152357786636938322685439776958548701986208857622629213478867187500&
3510497287734769771463259888233565327785269489229364491602610280576683474421257133947258007110692430889480343347613096\\
29&82&
971787874854834355775051826541159058553705103847742341196288915488789687821798939394191225757695546875000&
23031461318491610893979429349155061625976525444882846492684310835629626830463531600347477896305030968820296477010894324496 .
\end{array}
\]
The replay exits with status \(0\), reports \texttt{SUMMARY failures=0}, and
has SHA-256
\[
\hashsplit{ccd56620d91d6a3c2c9e23673211e25d}{cb5e1ae5d1a4ddc71e699a7bf80fab4a}.
\]
\end{proof}

\begin{lemma}[The central binomial factor absorbs the twenty-four-pause loss]\label{lem:bc-twentyfourth-c-pause-binomial-absorption}
For every \(q\ge61\),
\[
\begin{aligned}
  &16777216\binom{2q+24}{0}
    +67108864\binom{2q+24}{2}
    +268435456\binom{2q+24}{4}\\
  &\quad+1073741824\binom{2q+24}{6}
    +4294967296\binom{2q+24}{8}
    +17179869184\binom{2q+24}{10}\\
  &\quad+68719476736\binom{2q+24}{12}
    +274877906944\binom{2q+24}{14}\\
  &\quad+1099511627776\binom{2q+24}{16}
    +4398046511104\binom{2q+24}{18}\\
  &\quad+17592186044416\binom{2q+24}{20}
    +70368744177664\binom{2q+24}{22}
    +281474976710656\binom{2q+24}{24}
  \le \binom{2q+24}{q}.
\end{aligned}
\]
\end{lemma}

\begin{proof}
At \(q=61\), the left side is
\[
  544718913148137508039167628581708985532416
  <
  821730205038148901357321269744201244708160
  =
  \binom{146}{61}.
\]
The margin in this base comparison is
\[
  277011291890011393318153641162492259175744 .
\]
Let \(P(q)\) denote the displayed left side.  For \(0\le p\le24\),
\[
  \frac{\binom{2q+26}{p}}{\binom{2q+24}{p}}
  \le
  \frac{(2q+26)(2q+25)}{(2q+2)(2q+1)}
  <2
  \qquad(q\ge61),
\]
because \(4q^2-90q-646>0\).  Hence \(P(q+1)<2P(q)\).  On the other hand,
\[
  \frac{\binom{2q+26}{q+1}}{\binom{2q+24}{q}}
  =
  \frac{(2q+26)(2q+25)}{(q+1)(q+25)}
  >3,
\]
since the numerator minus three times the denominator is \(q^2+24q+575\).
Thus the binomial side grows by a larger factor than \(P(q)\), and induction
from \(q=61\) proves the inequality.
\end{proof}

\begin{lemma}[Exact twenty-fourth $C$ pause-arithmetic checks after the first nine rows]\label{lem:bc-c-twentyfourth-frontier-pause-arithmetic-certificate}
For \(30\le q\le60\),
put
\[
\begin{aligned}
  P(q)={}&16777216\binom{2q+24}{0}
    +67108864\binom{2q+24}{2}
    +268435456\binom{2q+24}{4}\\
  &+1073741824\binom{2q+24}{6}
    +4294967296\binom{2q+24}{8}
    +17179869184\binom{2q+24}{10}\\
  &+68719476736\binom{2q+24}{12}
    +274877906944\binom{2q+24}{14}\\
  &+1099511627776\binom{2q+24}{16}
    +4398046511104\binom{2q+24}{18}\\
  &+17592186044416\binom{2q+24}{20}
    +70368744177664\binom{2q+24}{22}
    +281474976710656\binom{2q+24}{24}.
\end{aligned}
\]
Then
\[
  2^{2q+1}P(q)s_{q+24}\le s_{2q+24}.
\]
\end{lemma}

\begin{proof}
The exact C++/GMP/OpenMP replay
\[
\texttt{certificates/post\_m29/bc\_twentyfourth\_frontier\_gmp\_certificate.log}
\]
computes the displayed integer inequality for
\(30\le q\le60\).
It exits with status \(0\), reports \texttt{SUMMARY failures=0}, and has
SHA-256
\[
\hashsplit{ccd56620d91d6a3c2c9e23673211e25d}{cb5e1ae5d1a4ddc71e699a7bf80fab4a}.
\]
It verifies the exact positive margins
\[
\begin{array}{c|c|c}
q & P(q) & s_{2q+24}-2^{2q+1}P(q)s_{q+24}\\ \hline
30&
187480189610906253395513621459501056&
152704139097871013142492293033505400204450787722883583472906050548331311475764735005352067097574849785076655327088181524932608\\
31&
361534372870305479222123357670473728&
1143758394960371245824775272138777286197289961127712644006871745955823899716077881608915871881023064108931577623916442620700735488\\
32&
685047451016679782159360902945570816&
8674701576976002221340231670277008106211846902476029819141182616454829383827423296141980454215545348175084055401841452803912290547712\\
33&
1276632968835875025030772443377041408&
68717962553457316450163404628290001830418790062539983147014073651171664844139435380666010990910886657581546429970698206330588303294543872\\
34&
2341817351039060023677249220872503296&
569039167350865492708009571033614789086355217943909095763027482581285546713478937410490743215664810430298854029513202517521731680933802037248\\
35&
4231764443057592491270784110232076288&
4921474294751144034097965894068037494532894778122830609219036972770117815348456633246982381248669855194609966701734773780274142113579384696696832\\
36&
7538531237089192836487032324726194176&
44415035695065473896178353901038559626357468942020052875814250827407796446093372793764425908144752760286941584525536875302579274571065508312100835328\\
37&
13247738110892704221519004143859007488&
417889612970048052584838800345578251023938118249645145308732633218462357627139197413789731047380341548699659967893576284569528541661734550608515016101888\\
38&
22980482826037444330997883444549124096&
4095628359229139677790038545688124387171460094377123147086100355618200738154337207887266653473368634147100419996259154525445667751484644340129805782617186304\\
39&
39372519886172890702513405381071339520&
41778449637153529114988857676920442216434294820583551155303705617283849502220728999700911080583796811511858074713729772019113902426963035230344049853311574614016\\
40&
66662457918691349040098453939694862336&
443216814209337450012757743685336657936846792416452111469453292540771902870321088648864489914993361670611702690660893193079319969217922510918757316999783678645018624\\
41&
111595199364953179020755435232270745600&
4886351313102936940089431518329040071265989632165263817250627352079929803482069725282933337662572091192563888518216893101958444005951532932676711798760644351906084986880\\
42&
184796496702022672234964705370299170816&
55942579353574319089200121285804183052364702460701355050031430265656095048183347996058690570092644767757844738449620019786837271458656742323612799747115143354520763894710272\\
43&
302845432968906503364914207342628700160&
664639403789427225469891417386216893082563736850625297308327739656035774053930804706232181854421856374729538204123701396077875271442978374414386388001637320109565305493531869184\\
44&
491372191669037241023287832277602533376&
8188851318177335552017835681331553521625823830781908333652171621113540647458560250228121805877747120829475283379658046339080133859498482993158249283469910028404252050618882778972160\\
45&
789649978248266532899566899469557432320&
104561339429957869976850712192175081857147421476484780389061643703911519085056807261854251118938188013171199693841009638836550765298603332228991619457007223121688865548132983818834919424\\
46&
1257347646879485290109527127404845203456&
1382796873366955667733690967286843289987756054446816766576306142225977851802283676018058207181478995158353945163645617169612775882059945288277316355606057458636272131807416829792162204516352\\
47&
1984383908535430897599901596338471567360&
18928751602053607333207428442632680590630278194344844191327581739462380415346384804482035375812513526928445438587250165859941025626381579194999341590383000375782799509489626024548721085820141568\\
48&
3105202164610359463871128385177199312896&
268045196852917110920858862735216711572063370817254142509937383503318988311628113834360690049627929249456804003219236246555964396658406311356616801508849616168364547390102893840621428747416912232448\\
49&
4819303060582601485630522408707377070080&
3924381011990010004171741526907272586495603870423361575392707824762601471585680924553350870861239559700866686574877569056358603034039575504451563512173433948640478037121021242996500490315935122277236736\\
50&
7420577239759044961332697315377976180736&
59370954705574675743610932769149025465094289810529648106620169776007660610556910037344232180976623135577124841747554271684563774016932445259164116138560932914791515937422523018030042940298179964487048429568\\
51&
11338935656308391815239519964619620745216&
927655959519066445608888540023322715448747914219128683282635294869591395692061194841427738089327973005580665179266955557076837192895161539607023283101987835212100901642493531445922514576581493767038266907754496\\
52&
17199020185826235449368671748043369873408&
14961925445802499547156199139453684401279518482264563401424189166222589424479010156630396068268404871891208739432982833721131117716272311608051343551231908351519506336160508468874054838159713449019081760279254597632\\
53&
25902498203249895658868172377391299035136&
248977571795961719585950638301180784486799423456352646420822065655939772346125968183993343870566658365093321610647117879421404109833500428096376864853831449495281315669040256649352151167018085250655233369027843639541760\\
54&
38742736876560392863289241836099351871488&
4272640406230412894616231425874304220198669125066632079275818659721857857874177360526533028844979016924453049306754765019679374880365495383738540194548780751821493553638049054888081589560303686059241204283801258202895876096\\
55&
57563690369881776688980740324262890438656&
75577643333990134698069877278328635352321569315538900459216908366767409544237231339934088845470704543840641139952621207584932354338773240398207336956751357443158125682140101348367251432277661046719959581130241609583500028149760\\
56&
84978838825635691607476887063932717498368&
1377383225901130709058774297082485196193647821316210892873165874812255280770769610403573420585038148052677808687218856863625869339304535321929707807768019097471330223083123948450179889247172169638062092234766147823974979332032036864\\
57&
124671275785201319272217273168479362482176&
25851753711347463363316510715681015703504022257923548656702277033750767963944257680535643073473161572427713087490666596424649025118435583782164772815108315930482094312588863771235523325040744136755959383234056418447936408023796331184128\\
58&
181802910737646130354949964751042351464448&
499475127943829349601459914348599634357544242939569285100818023472633461738484986402944050685573262176246738556897543676399772756881237235189485775302590445531590212646520909031594569174238005396547000445248699515668215518024066853056020480\\
59&
263569690794744095879283258861947239530496&
9929937435148106573450272984529735153079341429560554937441580616711155591949846482760110409023522938245758537492265826004725889124081619563145629155185055606983026447540559828830088225706340140717996067747634995377954360326970534301965755875328\\
60&
379951322839051563495119197660121481084928&
203054754927418813474247295975011105304294871765659314381761097132555522291610788891316070363046108534532858882992348624485109727223433365009887997161843926580058493908104167693058821991885990763093354213323789313758483330780174811396397212309913600 .
\end{array}
\]
This proves the claim.
\end{proof}

\begin{proposition}[Residual $C$ twenty-fourth nonboundary lower correction]\label{prop:bc-c-twentyfourth-nonboundary-lower-correction}
For every residual row \(C_b\) with \(20\le b\le217\), put \(q=b+1\) and
write \(s_j=m_j^{\mathrm{st}}\) and \(m_j^{C_b}=s_j+\delta_j^{C_b}\).  Then
\[
  -\frac12 s_{2q+24}\le \delta_{2q+24}^{C_b}.
\]
\end{proposition}

\begin{proof}
As in the proof of
\cref{prop:bc-c-first-nonboundary-lower-correction},
\(|\delta_{2q+24}^{C_b}|\) is the number of column-even Pieri two-strip paths
of length \(2q+24\) with terminal height at least \(2q\).  If
\(21\le q\le29\),
\cref{lem:bc-twentyfourth-c-q21-q22-q23-q24-q25-q26-q27-q28-q29-fpf-pause-certificate}
gives
\[
  |\delta_{2q+24}^{C_b}|\le s_{2q+24}/2 .
\]
For \(q\ge30\), \cref{lem:bc-c-twentyfourth-nonboundary-twentyfour-pause-code},
with \(P(q)\) denoting the polynomial from
\cref{lem:bc-c-twentyfourth-frontier-pause-arithmetic-certificate}, gives
\[
  |\delta_{2q+24}^{C_b}|
  \le
  P(q)\,2^{2q}s_{q+24}.
\]
If \(30\le q\le60\),
\cref{lem:bc-c-twentyfourth-frontier-pause-arithmetic-certificate} gives
\[
  |\delta_{2q+24}^{C_b}|\le s_{2q+24}/2 .
\]
If \(q\ge61\), \cref{lem:bc-twentyfourth-c-pause-binomial-absorption} gives
\[
  |\delta_{2q+24}^{C_b}|
  \le
  2^{2q}\binom{2q+24}{q}s_{q+24}.
\]
The exact GMP/OpenMP replay
\[
\texttt{certificates/post\_m29/bc\_pieri\_powerloss\_2q\_gmp\_certificate.log}
\]
checks the inequality
\[
  2^{2q+1}\binom{2q+24}{q}s_{q+24}\le s_{2q+24}
\]
on this residual row and moment as part of its all-row residual-window audit;
it reports \texttt{failures=0} and has SHA-256
\[
\hashsplit{124068c5daa1123088a3cd79c3f6b839}{fee45a1aaef637ed82161c465d86bf45}.
\]
Thus \(|\delta_{2q+24}^{C_b}|\le s_{2q+24}/2\) also for \(q\ge61\), which
implies the claimed lower bound.
\end{proof}


\begin{lemma}[$C$ twenty-fifth-nonboundary height-bad paths have a twenty-five-pause code]\label{lem:bc-c-twentyfifth-nonboundary-twentyfive-pause-code}
For \(q\ge1\), the number of column-even Pieri two-strip paths
\[
  \varnothing=\lambda^{(0)}\subset\lambda^{(1)}\subset\cdots
  \subset\lambda^{(2q+25)}
\]
with
\[
  \ell(\lambda^{(2q+25)})\ge2q
\]
is at most
\[
\begin{aligned}
  \biggl(&67108864\binom{2q+25}{1}
    +268435456\binom{2q+25}{3}
    +1073741824\binom{2q+25}{5}\\
  &\quad+4294967296\binom{2q+25}{7}
    +17179869184\binom{2q+25}{9}\\
  &\quad+68719476736\binom{2q+25}{11}
    +274877906944\binom{2q+25}{13}\\
  &\quad+1099511627776\binom{2q+25}{15}
    +4398046511104\binom{2q+25}{17}\\
  &\quad+17592186044416\binom{2q+25}{19}
    +70368744177664\binom{2q+25}{21}\\
  &\quad+281474976710656\binom{2q+25}{23}
    +1125899906842624\binom{2q+25}{25}\biggr)2^{2q}s_{q+25}.
\end{aligned}
\]
\end{lemma}

\begin{proof}
The terminal height is even and at most \(2q+24\).  Thus it is
\[
  2q+24,\quad 2q+22,\quad 2q+20,\quad 2q+18,\quad 2q+16,
  \quad 2q+14,\quad 2q+12,\quad 2q+10,\quad 2q+8,
  \quad 2q+6,\quad 2q+4,\quad 2q+2,\quad\text{or}\quad 2q .
\]
These cases have respectively \(1,3,5,7,9,11,13,15,17,19,21,23,25\) steps
which fail to create a new row.  Deleting the first column, shifting left, and
standardizing the repeated pause labels injects the paths with \(p\) pauses
into the choice of the \(p\)-element pause set and an even-column standard
tableau of size \(2q+25+p\).  Therefore the thirteen cases contribute at most
\[
  \binom{2q+25}{1}2^{2q+26}s_{q+13},
  \quad
  \binom{2q+25}{3}2^{2q+28}s_{q+14},
  \quad
  \binom{2q+25}{5}2^{2q+30}s_{q+15},
\]
\[
  \binom{2q+25}{7}2^{2q+32}s_{q+16},
  \quad
  \binom{2q+25}{9}2^{2q+34}s_{q+17},
  \quad
  \binom{2q+25}{11}2^{2q+36}s_{q+18},
\]
\[
  \binom{2q+25}{13}2^{2q+38}s_{q+19},
  \quad
  \binom{2q+25}{15}2^{2q+40}s_{q+20},
\]
\[
  \binom{2q+25}{17}2^{2q+42}s_{q+21},
  \quad
  \binom{2q+25}{19}2^{2q+44}s_{q+22},
\]
\[
  \binom{2q+25}{21}2^{2q+46}s_{q+23},
  \quad
  \binom{2q+25}{23}2^{2q+48}s_{q+24},
  \quad
  \binom{2q+25}{25}2^{2q+50}s_{q+25},
\]
by
\cref{lem:bc-even-column-standard-tableaux-2n-bound,lem:bc-stable-moments-monotone-after-two}.
After replacing \(s_{q+13},\ldots,s_{q+24}\) by \(s_{q+25}\), the sum is the
displayed bound.
\end{proof}

\begin{lemma}[The first ten twenty-fifth $C$ pause cases are absorbed by the fixed-point-free count]\label{lem:bc-twentyfifth-c-q21-q22-q23-q24-q25-q26-q27-q28-q29-q30-fpf-pause-certificate}
For \(21\le q\le30\), the number of column-even Pieri two-strip
paths of length \(2q+25\) with terminal height at least \(2q\) is at most
\(s_{2q+25}/2\).
\end{lemma}

\begin{proof}
Use the pause-set injection from
\cref{lem:bc-c-twentyfifth-nonboundary-twentyfive-pause-code}, but count the
target even-column standard tableaux exactly by
\cref{lem:bc-even-column-standard-tableaux-fpf-count}.  For
\(21\le q\le30\), this gives the upper bound
\[
  \sum_{\substack{1\le p\le25\\ p\ {\rm odd}}}
  \binom{2q+25}{p}(2q+24+p)!! .
\]
The exact C++/GMP/OpenMP replay
\[
\texttt{certificates/post\_m29/bc\_twentyfifth\_frontier\_gmp\_certificate.log}
\]
checks the displayed integer sums and verifies
\[
\begin{array}{c|c|c|c}
q & j & \mathrm{FPF}(q,j) & s_j-2\mathrm{FPF}(q,j)\\ \hline
21&67&
53945547104126407383322291602955043744711552210652322566892979376723601047458376000000000&
1954103273587235005111756327875185673979708831107140922339104361021446719252851588240361742592\\
22&69&
12436988606154125333980540152151971543403099228006940852375342097616651767138901047939843750&
9035736443111400135988323294387715151586895082006777533313061048660327038924159545002344875436212\\
23&71&
2835873872964951684296367117210385126083085889221674751522197771933392568825887299874081250000&
44272864665495952151310816634473658257307643744848544396912586953112662149004044381809831525472927840\\
24&73&
640700951891983448419183023400924950368312337844681703863135696851715480648574704218477801562500&
229499008969934069989607346511682182481480233756338087680875425473325727937676202847659582885763833578744\\
25&75&
143652549639416524670187725918613977100508202153063555438013850540670917446501998745632600976562500&
1256676780181561943568484035097136212085343789401234828326766980672241077611203668029370372846805965287624568\\
26&77&
32009166083826444919190417813722339499428387619084085979451889481573554295618138132951285179960937500&
7258238000498592505584170794915546762303590580672709905187555666142825863299185662345272169192348925741301914568\\
27&79&
7097141878747663521542393958084376639726674765305168806462532450413011248167599752568323745342070312500&
44157232836678506541306259247738386470556733458827259226145238078456613352720602521305511039905928679593762409971608\\
28&81&
1567567585023481192490907927926316227559466912440613907406572763297522025365414992850474541903447617187500&
282594820649083038819883837223590513528428449381292337786396543612575987978486865588963514866024182842206577152151798440\\
29&83&
345252126968540163371614107551115155897950374319690614908709560770110350594373567232369892938377255703125000&
1900094241916969427426955303058608597005696156807854051223487349422716079467195301440688647237963103929238069967443999613936\\
30&85&
75893035884435531801065507087379609797949033754764854559411596003010676376030494156952934584007376399218750000&
13406572350450220181690857920845994364556124815060667442522082384336922220275280322939198553929265856734864640696164443315954592 .
\end{array}
\]
The replay exits with status \(0\), reports \texttt{SUMMARY failures=0}, and
has SHA-256
\[
\hashsplit{1bd9d79c07c867c4d2d40c446cc351f3}{a2fbd100861b423ff9e32c0bfc41ed91}.
\]
\end{proof}

\begin{lemma}[The central binomial factor absorbs the twenty-five-pause loss]\label{lem:bc-twentyfifth-c-pause-binomial-absorption}
For every \(q\ge64\),
\[
\begin{aligned}
  &67108864\binom{2q+25}{1}
    +268435456\binom{2q+25}{3}
    +1073741824\binom{2q+25}{5}\\
  &\quad+4294967296\binom{2q+25}{7}
    +17179869184\binom{2q+25}{9}
    +68719476736\binom{2q+25}{11}\\
  &\quad+274877906944\binom{2q+25}{13}
    +1099511627776\binom{2q+25}{15}\\
  &\quad+4398046511104\binom{2q+25}{17}
    +17592186044416\binom{2q+25}{19}\\
  &\quad+70368744177664\binom{2q+25}{21}
    +281474976710656\binom{2q+25}{23}
    +1125899906842624\binom{2q+25}{25}
  \le \binom{2q+25}{q}.
\end{aligned}
\]
\end{lemma}

\begin{proof}
At \(q=64\), the left side is
\[
  38106144974042111076149000047527760958062592
  <
  95784491736452493294473263363363867621614750
  =
  \binom{153}{64}.
\]
The margin in this base comparison is
\[
  57678346762410382218324263315836106663552158 .
\]
Let \(P(q)\) denote the displayed left side.  For \(1\le p\le25\),
\[
  \frac{\binom{2q+27}{p}}{\binom{2q+25}{p}}
  \le
  \frac{(2q+27)(2q+26)}{(2q+2)(2q+1)}
  <2
  \qquad(q\ge64),
\]
because \(4q^2-94q-698>0\).  Hence \(P(q+1)<2P(q)\).  On the other hand,
\[
  \frac{\binom{2q+27}{q+1}}{\binom{2q+25}{q}}
  =
  \frac{(2q+27)(2q+26)}{(q+1)(q+26)}
  >3,
\]
since the numerator minus three times the denominator is \(q^2+25q+624\).
Thus the binomial side grows by a larger factor than \(P(q)\), and induction
from \(q=64\) proves the inequality.
\end{proof}

\begin{lemma}[Exact twenty-fifth $C$ pause-arithmetic checks after the first ten rows]\label{lem:bc-c-twentyfifth-frontier-pause-arithmetic-certificate}
For \(31\le q\le63\),
put
\[
\begin{aligned}
  P(q)={}&67108864\binom{2q+25}{1}
    +268435456\binom{2q+25}{3}
    +1073741824\binom{2q+25}{5}\\
  &+4294967296\binom{2q+25}{7}
    +17179869184\binom{2q+25}{9}
    +68719476736\binom{2q+25}{11}\\
  &+274877906944\binom{2q+25}{13}
    +1099511627776\binom{2q+25}{15}\\
  &+4398046511104\binom{2q+25}{17}
    +17592186044416\binom{2q+25}{19}\\
  &+70368744177664\binom{2q+25}{21}
    +281474976710656\binom{2q+25}{23}
    +1125899906842624\binom{2q+25}{25}.
\end{aligned}
\]
Then
\[
  2^{2q+1}P(q)s_{q+25}\le s_{2q+25}.
\]
\end{lemma}

\begin{proof}
The exact C++/GMP/OpenMP replay
\[
\texttt{certificates/post\_m29/bc\_twentyfifth\_frontier\_gmp\_certificate.log}
\]
computes the displayed integer inequality for
\(31\le q\le63\).
It exits with status \(0\), reports \texttt{SUMMARY failures=0}, and has
SHA-256
\[
\hashsplit{1bd9d79c07c867c4d2d40c446cc351f3}{a2fbd100861b423ff9e32c0bfc41ed91}.
\]
It verifies the exact positive margins
\[
\begin{array}{c|c|c}
q & P(q) & s_{2q+25}-2^{2q+1}P(q)s_{q+25}\\ \hline
31&
5047910995737832421064635984025485312&
97212057122575785915759525641356783823015024808426402433408095662647193929202815446110778772212386010016690567990609326507159939072\\
32&
9783006720269792541474484407385456640&
766954353939393667569022992733765659627320447644051088035886803401787878320182866061951295169165970091790231820563514577699107988938752\\
33&
18637826401908385799827173801136750592&
6218640262642298623471717586378681599886035423752070670239226479836892062882930723067076544314900919400470699573351001037903803081716652032\\
34&
34934644971009688944854053909189099520&
52635951248913380363096554670940018998687904181126439744860647599390334254131045206609871659195754850200902899681575443089135964970022965653504\\
35&
64476834804439084162653019781224988672&
465079043174301178184341342065280398037363358103801609164006057017319303807352772327207631445866712638835105197486933485620701727521928112464187392\\
36&
117262843157300261390013131767288954880&
4286049070267780244518691917077741591716599511447976648288701620850991830045484746836775910594577952463729420787301739274381220880340765557565978902528\\
37&
210293928924511933265011519610056343552&
41162110205699485727215339801851569360003072364737654866917945016614751335571302428572672472798228228711291271476228604848252895925404917830532364590518272\\
38&
372119345391770058758436281929800089600&
411610493355052818279014621341250612631165594694676111232575854973800296983548758460568261542246356358629772292125569504547950917834663202207777683562435993600\\
39&
650111319848417452756276732233928146944&
4282289551631615494114688941497730836944521750231747800052976634031574713602713948849735091626026989751657091606628419354322738950073704573212257614722887140204544\\
40&
1121984269071048889070437519272118845440&
46316141414883358323830991089558890165390225909696775048318027260804124004205204517583514728370573190721470757390521348826419219367353011835469893188271706081902886912\\
41&
1913847194781341639765142715539163971584&
520396248607813618888056039970573737576416291340395281490594120292579772502192842603237098490209560027663514896065142965844760721482200834085989460537478851693417009979392\\
42&
3228216961546034014455968288751254241280&
6069768027148561397417163478880042011925471821600363021004117538289020366099964328475453167457782340779278129204724950767844356450563316479059783686633218731847336200460255232\\
43&
5387080268270936355857978319768905580544&
73442633136598658438163607192844136475707651745106740958765256070499076775127961249335956052117258511331063033472811270436778910496225407164003459074427953452926495052654312144896\\
44&
8897492608787993868254394228067155312640&
921245524013205551904028317863614222246461577870892453852060285820248778078754660373931401545439176523407938139700609748110410402480627055246962113183739202592674854077332786853576704\\
45&
14550646558193869490293003471973100027904&
11972270293412322391972298372881474972085005273394744721007237567823874685518333821840780991436392594167382261878408470787919634600269140872721519305117362453639004378093553889982574346240\\
46&
23570246115045439129108320722056212643840&
161095796530604867861824423617067152648031497712093719386881574559529007608349199329781763181124243339030954893722313330215338715067674829545443756428457008831494197652298776309589290820599808\\
47&
37832957795248335618617112623112377073664&
2243056546216084944751803001437049750979070591224284350737505927179198261499391693157746798390083114365256346365526907849146018844682726368164164309735661165889189340361632518946794316479695585280\\
48&
60193445195613404452996443326105339822080&
32299439121457320747514177391710501466402092013167029819537121997675012086905136846153307770441919376514822370157341648520337832282033997695320641249503788354839009498550571069608190232062149340037120\\
49&
94960073805735595991098655183408918953984&
480736573437627632444442335984447992758354598717036525772465499694966097182063008643230425156751454414918321002551296174079102396480233844441394778693326607787941698384630667798698320081033742881625571328\\
50&
148586194263139130728563524945902588395520&
7391682389000089802339052811838849786799161367624831879592272371152075216279428155128810413595237593926056365691267517347572277772116461285241814936498613248305165130191867065090020008937044521827778026995712\\
51&
230667838351478041798202113726030595227648&
117348456612109982890333740601526720425634039231943514917811519681247869748671876191084801320015732361202214980943820531875682865413198777596986053270162399161544278752412159590327995033056968977269327412359135232\\
52&
355374164207232412053046237485559762124800&
1922607071870099047776047162593688531479358066518748369246194954730854888092945685178853701385041334706893164259079373685120146923773528088807053189989346528763513830541244285549829051141276279812942327357524297121792\\
53&
543485326228885150584532144126367962759168&
32491567507955082463681564380506329756099361850786123605078230688004342363121789624083883531168901810822933010391881362621541071865472893154598898264224694491620874781379870349527441819467624069729627135545497384388722688\\
54&
825277901933613159407929980558882667233280&
566124760447716687947008385494882968030422223724738218851281121371358940877288842482231544590966245332113470500812557495446042670764535058543187273532245544870224535589925949609113656707493662860365386231723749130006894608384\\
55&
1244586175355712284672662269264877257228288&
10165191425994752876422337440716385522605723819592550729513028500215155175192921427722956972046854885385937400190163185442477018643508316492391326576879283861269740390446388176969310416712669107041559568470212003681397813032648704\\
56&
1864485727473122212483823931436168286044160&
188012781990683033386658460821740358237307957685411817590725441491092560213796608265641197240689771422190045382237578537049242206596189477928506009016149968348824263444879476718798924817215488957217930139312987883480563982282019831808\\
57&
2775203333705242089089334506800809412395008&
3580467372446220814152304752927059603958957759728714717147664963961223152885589177654770551741669607637811670657243819612277939584357826341648955599856718236096642572078066287964077616899528630241087526389498003046604030327927160032067584\\
58&
4105066251980629155191870424941471612272640&
70176245780667169039660641024950160871814512823772855044674054024142339441971273403262985812659257151710402437723248175993895202994953619792695850924242539208302196882729896265159600504407155498256600381672385080263056788469466295044219863040\\
59&
6035580177644924306983708262908042755964928&
1415015897186040195814710359626221483641510356304576359437403210437788218086343920080774700404849706913793543228190221865190813197916310392245894008023846536238144237555869876458939223054380406225148700756855569502945946175081253496396124536963072\\
60&
8822088353065283435060801970710521598443520&
29341408362910814023208684901539056904326733923073805664012069255904450899146645447130641572453654341904998696268843792579400652808960655321886463893138720432761065009541439744228660820847402449458619482872522491780259203009878747259150551038025531392\\
61&
12821939899535888663003786654326644694581248&
625433971413586789951272620870545357688048891354032157322278449762852737797980863532905728627516721165489285111497510465167355125870679556099090045558938890224274076497116789371673717536597241881722814919070900043582052109922419638514253462917872432709632\\
62&
18532715539394626734029159870149216226508800&
13699346116150342953286001332419072291155123064090619006771840983382965352684994833634945501050462429269621670438555819957805858762949285838961068087306789612807502765254114172662733184099825479528595127580714149719364531177577768150764748951569985734363840512\\
63&
26643864127563813718862612702497080520736768&
308231792878096742189917625001089406902867307567147701783229431502508469269082006751052137576273510382282813560168170544798567917743462138864167361895574640290273721012328235634333610674470073931773144502044554253748726884868826347977693548123023697889900683067392 .
\end{array}
\]
This proves the claim.
\end{proof}

\begin{proposition}[Residual $C$ twenty-fifth nonboundary lower correction]\label{prop:bc-c-twentyfifth-nonboundary-lower-correction}
For every residual row \(C_b\) with \(20\le b\le217\), put \(q=b+1\) and
write \(s_j=m_j^{\mathrm{st}}\) and \(m_j^{C_b}=s_j+\delta_j^{C_b}\).  Then
\[
  -\frac12 s_{2q+25}\le \delta_{2q+25}^{C_b}.
\]
\end{proposition}

\begin{proof}
As in the proof of
\cref{prop:bc-c-first-nonboundary-lower-correction},
\(|\delta_{2q+25}^{C_b}|\) is the number of column-even Pieri two-strip paths
of length \(2q+25\) with terminal height at least \(2q\).  If
\(21\le q\le30\),
\cref{lem:bc-twentyfifth-c-q21-q22-q23-q24-q25-q26-q27-q28-q29-q30-fpf-pause-certificate}
gives
\[
  |\delta_{2q+25}^{C_b}|\le s_{2q+25}/2 .
\]
For \(q\ge31\), \cref{lem:bc-c-twentyfifth-nonboundary-twentyfive-pause-code},
with \(P(q)\) denoting the polynomial from
\cref{lem:bc-c-twentyfifth-frontier-pause-arithmetic-certificate}, gives
\[
  |\delta_{2q+25}^{C_b}|
  \le
  P(q)\,2^{2q}s_{q+25}.
\]
If \(31\le q\le63\),
\cref{lem:bc-c-twentyfifth-frontier-pause-arithmetic-certificate} gives
\[
  |\delta_{2q+25}^{C_b}|\le s_{2q+25}/2 .
\]
If \(q\ge64\), \cref{lem:bc-twentyfifth-c-pause-binomial-absorption} gives
\[
  |\delta_{2q+25}^{C_b}|
  \le
  2^{2q}\binom{2q+25}{q}s_{q+25}.
\]
The exact GMP/OpenMP replay
\[
\texttt{certificates/post\_m29/bc\_pieri\_powerloss\_2q\_gmp\_certificate.log}
\]
checks the inequality
\[
  2^{2q+1}\binom{2q+25}{q}s_{q+25}\le s_{2q+25}
\]
on this residual row and moment as part of its all-row residual-window audit;
it reports \texttt{failures=0} and has SHA-256
\[
\hashsplit{124068c5daa1123088a3cd79c3f6b839}{fee45a1aaef637ed82161c465d86bf45}.
\]
Thus \(|\delta_{2q+25}^{C_b}|\le s_{2q+25}/2\) also for \(q\ge64\), which
implies the claimed lower bound.
\end{proof}

\begin{lemma}[$C$ twenty-sixth-nonboundary height-bad paths have a twenty-six-pause code]\label{lem:bc-c-twentysixth-nonboundary-twentysix-pause-code}
For \(q\ge1\), the number of column-even Pieri two-strip paths
\[
  \varnothing=\lambda^{(0)}\subset\lambda^{(1)}\subset\cdots
  \subset\lambda^{(2q+26)}
\]
with
\[
  \ell(\lambda^{(2q+26)})\ge2q
\]
is at most
\[
\begin{aligned}
  \biggl(&67108864\binom{2q+26}{0}
    +268435456\binom{2q+26}{2}
    +1073741824\binom{2q+26}{4}\\
  &\quad+4294967296\binom{2q+26}{6}
    +17179869184\binom{2q+26}{8}\\
  &\quad+68719476736\binom{2q+26}{10}
    +274877906944\binom{2q+26}{12}\\
  &\quad+1099511627776\binom{2q+26}{14}
    +4398046511104\binom{2q+26}{16}\\
  &\quad+17592186044416\binom{2q+26}{18}
    +70368744177664\binom{2q+26}{20}\\
  &\quad+281474976710656\binom{2q+26}{22}
    +1125899906842624\binom{2q+26}{24}\\
  &\quad+4503599627370496\binom{2q+26}{26}\biggr)2^{2q}s_{q+26}.
\end{aligned}
\]
\end{lemma}

\begin{proof}
The terminal height is even and at most \(2q+26\).  Thus it is
\[
  2q+26,\quad 2q+24,\quad 2q+22,\quad 2q+20,\quad 2q+18,
  \quad 2q+16,\quad 2q+14,\quad 2q+12,\quad 2q+10,
  \quad 2q+8,\quad 2q+6,\quad 2q+4,\quad 2q+2,
  \quad\text{or}\quad 2q .
\]
These cases have respectively \(0,2,4,\ldots,26\) steps which fail to create a
new row.  Deleting the first column, shifting left, and standardizing the
repeated pause labels injects the paths with \(p\) pauses into the choice of
the \(p\)-element pause set and an even-column standard tableau of size
\(2q+26+p\).  Therefore the fourteen cases contribute at most
\[
  \binom{2q+26}{0}2^{2q+26}s_{q+13},
  \quad
  \binom{2q+26}{2}2^{2q+28}s_{q+14},
  \quad
  \binom{2q+26}{4}2^{2q+30}s_{q+15},
\]
\[
  \binom{2q+26}{6}2^{2q+32}s_{q+16},
  \quad
  \binom{2q+26}{8}2^{2q+34}s_{q+17},
  \quad
  \binom{2q+26}{10}2^{2q+36}s_{q+18},
\]
\[
  \binom{2q+26}{12}2^{2q+38}s_{q+19},
  \quad
  \binom{2q+26}{14}2^{2q+40}s_{q+20},
\]
\[
  \binom{2q+26}{16}2^{2q+42}s_{q+21},
  \quad
  \binom{2q+26}{18}2^{2q+44}s_{q+22},
\]
\[
  \binom{2q+26}{20}2^{2q+46}s_{q+23},
  \quad
  \binom{2q+26}{22}2^{2q+48}s_{q+24},
\]
\[
  \binom{2q+26}{24}2^{2q+50}s_{q+25},
  \quad
  \binom{2q+26}{26}2^{2q+52}s_{q+26},
\]
by
\cref{lem:bc-even-column-standard-tableaux-2n-bound,lem:bc-stable-moments-monotone-after-two}.
After replacing \(s_{q+13},\ldots,s_{q+25}\) by \(s_{q+26}\), the sum is the
displayed bound.
\end{proof}

\begin{lemma}[The first eleven twenty-sixth $C$ pause cases are absorbed by the fixed-point-free count]\label{lem:bc-twentysixth-c-q21-q22-q23-q24-q25-q26-q27-q28-q29-q30-q31-fpf-pause-certificate}
For \(21\le q\le31\), the number of column-even Pieri
two-strip paths of length \(2q+26\) with terminal height at least \(2q\) is at
most \(s_{2q+26}/2\).
\end{lemma}

\begin{proof}
Use the pause-set injection from
\cref{lem:bc-c-twentysixth-nonboundary-twentysix-pause-code}, but count the
target even-column standard tableaux exactly by
\cref{lem:bc-even-column-standard-tableaux-fpf-count}.  For
\(21\le q\le31\), this gives the upper bound
\[
  \sum_{\substack{0\le p\le26\\ p\ {\rm even}}}
  \binom{2q+26}{p}(2q+25+p)!! .
\]
The exact C++/GMP/OpenMP replay
\[
\texttt{certificates/post\_m29/bc\_twentysixth\_frontier\_gmp\_certificate.log}
\]
checks the displayed integer sums, exits with status \(0\), reports
\texttt{SUMMARY failures=0}, and has SHA-256
\[
\hashsplit{6a107c7c20c107e6a71db9aa65cf2884}{df0e1e94db82c20d7d3f1ae3acc4c07d}.
\]
Its transcript contains exactly eleven \texttt{C\_FPF\_MARGIN} records, one
for each displayed \(q\), and each record is the exact positive integer
\[
  s_{2q+26}-2\,\mathrm{FPF}(q,2q+26).
\]
This proves the half-stable bound for these first eleven cases.
\end{proof}

\begin{lemma}[The central binomial factor absorbs the twenty-six-pause loss]\label{lem:bc-twentysixth-c-pause-binomial-absorption}
For every \(q\ge66\),
\[
\begin{aligned}
  &67108864\binom{2q+26}{0}
    +268435456\binom{2q+26}{2}
    +1073741824\binom{2q+26}{4}\\
  &\quad+4294967296\binom{2q+26}{6}
    +17179869184\binom{2q+26}{8}
    +68719476736\binom{2q+26}{10}\\
  &\quad+274877906944\binom{2q+26}{12}
    +1099511627776\binom{2q+26}{14}\\
  &\quad+4398046511104\binom{2q+26}{16}
    +17592186044416\binom{2q+26}{18}\\
  &\quad+70368744177664\binom{2q+26}{20}
    +281474976710656\binom{2q+26}{22}\\
  &\quad+1125899906842624\binom{2q+26}{24}
    +4503599627370496\binom{2q+26}{26}
  \le \binom{2q+26}{q}.
\end{aligned}
\]
\end{lemma}

\begin{proof}
At \(q=66\), the left side is
\[
  1867725731873348752568016234429315107522609152
  <
  2737158422982075697084497097760673619727797850
  =
  \binom{158}{66}.
\]
The margin in this base comparison is
\[
  869432691108726944516480863331358512205188698 .
\]
Let \(P(q)\) denote the displayed left side.  For \(0\le p\le26\),
\[
  \frac{\binom{2q+28}{p}}{\binom{2q+26}{p}}
  \le
  \frac{(2q+28)(2q+27)}{(2q+2)(2q+1)}
  <2
  \qquad(q\ge66),
\]
where the \(p=0\) term has ratio \(1\), and the displayed strict inequality
follows from \(4q^2-98q-752>0\).  Hence \(P(q+1)<2P(q)\).  On the other hand,
\[
  \frac{\binom{2q+28}{q+1}}{\binom{2q+26}{q}}
  =
  \frac{(2q+28)(2q+27)}{(q+1)(q+27)}
  >3,
\]
since the numerator minus three times the denominator is \(q^2+26q+675\).
Thus the binomial side grows by a larger factor than \(P(q)\), and induction
from \(q=66\) proves the inequality.
\end{proof}

\begin{lemma}[Exact twenty-sixth $C$ pause-arithmetic checks after the first eleven rows]\label{lem:bc-c-twentysixth-frontier-pause-arithmetic-certificate}
For \(32\le q\le65\),
put
\[
\begin{aligned}
  P(q)={}&67108864\binom{2q+26}{0}
    +268435456\binom{2q+26}{2}
    +1073741824\binom{2q+26}{4}\\
  &+4294967296\binom{2q+26}{6}
    +17179869184\binom{2q+26}{8}
    +68719476736\binom{2q+26}{10}\\
  &+274877906944\binom{2q+26}{12}
    +1099511627776\binom{2q+26}{14}\\
  &+4398046511104\binom{2q+26}{16}
    +17592186044416\binom{2q+26}{18}\\
  &+70368744177664\binom{2q+26}{20}
    +281474976710656\binom{2q+26}{22}\\
  &+1125899906842624\binom{2q+26}{24}
    +4503599627370496\binom{2q+26}{26}.
\end{aligned}
\]
Then
\[
  2^{2q+1}P(q)s_{q+26}\le s_{2q+26}.
\]
\end{lemma}

\begin{proof}
The exact C++/GMP/OpenMP replay
\[
\texttt{certificates/post\_m29/bc\_twentysixth\_frontier\_gmp\_certificate.log}
\]
computes the displayed integer inequality for
\(32\le q\le65\).
It exits with status \(0\), reports \texttt{SUMMARY failures=0}, and has
SHA-256
\[
\hashsplit{6a107c7c20c107e6a71db9aa65cf2884}{df0e1e94db82c20d7d3f1ae3acc4c07d}.
\]
Its transcript contains exactly thirty-four \texttt{C\_MARGIN} records, one
for each displayed \(q\), and each record is the exact positive integer
\[
  s_{2q+26}-2^{2q+1}P(q)s_{q+26}.
\]
This proves the claim.
\end{proof}

\begin{proposition}[Residual $C$ twenty-sixth nonboundary lower correction]\label{prop:bc-c-twentysixth-nonboundary-lower-correction}
For every residual row \(C_b\) with \(20\le b\le217\), put \(q=b+1\) and
write \(s_j=m_j^{\mathrm{st}}\) and \(m_j^{C_b}=s_j+\delta_j^{C_b}\).  Then
\[
  -\frac12 s_{2q+26}\le \delta_{2q+26}^{C_b}.
\]
\end{proposition}

\begin{proof}
As in the proof of
\cref{prop:bc-c-first-nonboundary-lower-correction},
\(|\delta_{2q+26}^{C_b}|\) is the number of column-even Pieri two-strip paths
of length \(2q+26\) with terminal height at least \(2q\).  If
\(21\le q\le31\),
\cref{lem:bc-twentysixth-c-q21-q22-q23-q24-q25-q26-q27-q28-q29-q30-q31-fpf-pause-certificate}
gives
\[
  |\delta_{2q+26}^{C_b}|\le s_{2q+26}/2 .
\]
For \(q\ge32\), \cref{lem:bc-c-twentysixth-nonboundary-twentysix-pause-code},
with \(P(q)\) denoting the polynomial from
\cref{lem:bc-c-twentysixth-frontier-pause-arithmetic-certificate}, gives
\[
  |\delta_{2q+26}^{C_b}|
  \le
  P(q)\,2^{2q}s_{q+26}.
\]
If \(32\le q\le65\),
\cref{lem:bc-c-twentysixth-frontier-pause-arithmetic-certificate} gives
\[
  |\delta_{2q+26}^{C_b}|\le s_{2q+26}/2 .
\]
If \(q\ge66\), \cref{lem:bc-twentysixth-c-pause-binomial-absorption} gives
\[
  |\delta_{2q+26}^{C_b}|
  \le
  2^{2q}\binom{2q+26}{q}s_{q+26}.
\]
The exact GMP/OpenMP replay
\[
\texttt{certificates/post\_m29/bc\_pieri\_powerloss\_2q\_gmp\_certificate.log}
\]
checks the inequality
\[
  2^{2q+1}\binom{2q+26}{q}s_{q+26}\le s_{2q+26}
\]
on this residual row and moment as part of its all-row residual-window audit;
it reports \texttt{failures=0} and has SHA-256
\[
\hashsplit{124068c5daa1123088a3cd79c3f6b839}{fee45a1aaef637ed82161c465d86bf45}.
\]
Thus \(|\delta_{2q+26}^{C_b}|\le s_{2q+26}/2\) also for \(q\ge66\), which
implies the claimed lower bound.
\end{proof}

\begin{lemma}[$C$ twenty-seventh-nonboundary height-bad paths have a twenty-seven-pause code]\label{lem:bc-c-twentyseventh-nonboundary-twentyseven-pause-code}
For \(q\ge1\), the number of column-even Pieri two-strip paths
\[
  \varnothing=\lambda^{(0)}\subset\lambda^{(1)}\subset\cdots
  \subset\lambda^{(2q+27)}
\]
with
\[
  \ell(\lambda^{(2q+27)})\ge2q
\]
is at most
\[
\begin{aligned}
  \biggl(&268435456\binom{2q+27}{1}
    +1073741824\binom{2q+27}{3}
    +4294967296\binom{2q+27}{5}\\
  &\quad+17179869184\binom{2q+27}{7}
    +68719476736\binom{2q+27}{9}\\
  &\quad+274877906944\binom{2q+27}{11}
    +1099511627776\binom{2q+27}{13}\\
  &\quad+4398046511104\binom{2q+27}{15}
    +17592186044416\binom{2q+27}{17}\\
  &\quad+70368744177664\binom{2q+27}{19}
    +281474976710656\binom{2q+27}{21}\\
  &\quad+1125899906842624\binom{2q+27}{23}
    +4503599627370496\binom{2q+27}{25}\\
  &\quad+18014398509481984\binom{2q+27}{27}\biggr)2^{2q}s_{q+27}.
\end{aligned}
\]
\end{lemma}

\begin{proof}
The terminal height is even and at most \(2q+27\).  Thus it is
\[
  2q+26,\quad 2q+24,\quad 2q+22,\quad 2q+20,\quad 2q+18,
  \quad 2q+16,\quad 2q+14,\quad 2q+12,\quad 2q+10,
  \quad 2q+8,\quad 2q+6,\quad 2q+4,\quad 2q+2,
  \quad\text{or}\quad 2q .
\]
These cases have respectively \(1,3,5,\ldots,27\) steps which fail to create
a new row.  Deleting the first column, shifting left, and standardizing the
repeated pause labels injects the paths with \(p\) pauses into the choice of
the \(p\)-element pause set and an even-column standard tableau of size
\(2q+27+p\).  Therefore the fourteen cases contribute at most
\[
  \binom{2q+27}{1}2^{2q+28}s_{q+14},
  \quad
  \binom{2q+27}{3}2^{2q+30}s_{q+15},
  \quad
  \binom{2q+27}{5}2^{2q+32}s_{q+16},
\]
\[
  \binom{2q+27}{7}2^{2q+34}s_{q+17},
  \quad
  \binom{2q+27}{9}2^{2q+36}s_{q+18},
  \quad
  \binom{2q+27}{11}2^{2q+38}s_{q+19},
\]
\[
  \binom{2q+27}{13}2^{2q+40}s_{q+20},
  \quad
  \binom{2q+27}{15}2^{2q+42}s_{q+21},
\]
\[
  \binom{2q+27}{17}2^{2q+44}s_{q+22},
  \quad
  \binom{2q+27}{19}2^{2q+46}s_{q+23},
\]
\[
  \binom{2q+27}{21}2^{2q+48}s_{q+24},
  \quad
  \binom{2q+27}{23}2^{2q+50}s_{q+25},
\]
\[
  \binom{2q+27}{25}2^{2q+52}s_{q+26},
  \quad
  \binom{2q+27}{27}2^{2q+54}s_{q+27},
\]
by
\cref{lem:bc-even-column-standard-tableaux-2n-bound,lem:bc-stable-moments-monotone-after-two}.
After replacing \(s_{q+14},\ldots,s_{q+26}\) by \(s_{q+27}\), the sum is the
displayed bound.
\end{proof}

\begin{lemma}[The first twelve twenty-seventh $C$ pause cases are absorbed by the fixed-point-free count]\label{lem:bc-twentyseventh-c-q21-q22-q23-q24-q25-q26-q27-q28-q29-q30-q31-q32-fpf-pause-certificate}
For \(21\le q\le32\), the number of column-even
Pieri two-strip paths of length \(2q+27\) with terminal height at least
\(2q\) is at most \(s_{2q+27}/2\).
\end{lemma}

\begin{proof}
Use the pause-set injection from
\cref{lem:bc-c-twentyseventh-nonboundary-twentyseven-pause-code}, but count
the target even-column standard tableaux exactly by
\cref{lem:bc-even-column-standard-tableaux-fpf-count}.  For
\(21\le q\le32\), this gives the upper bound
\[
  \sum_{\substack{1\le p\le27\\ p\ {\rm odd}}}
  \binom{2q+27}{p}(2q+26+p)!! .
\]
The exact C++/GMP/OpenMP replay
\[
\texttt{certificates/post\_m29/bc\_twentyseventh\_frontier\_gmp\_certificate.log}
\]
checks the displayed integer sums, exits with status \(0\), reports
\texttt{SUMMARY failures=0}, and has SHA-256
\[
\hashsplit{67df773882805ff4321cd3659e16b1b9}{985db56a76a94d7136448011ef30f80d}.
\]
Its transcript contains exactly twelve \texttt{C\_FPF\_MARGIN} records, one
for each displayed \(q\), and each record is the exact positive integer
\[
  s_{2q+27}-2\,\mathrm{FPF}(q,2q+27).
\]
This proves the half-stable bound for these first twelve cases.
\end{proof}

\begin{lemma}[The central binomial factor absorbs the twenty-seven-pause loss]\label{lem:bc-twentyseventh-c-pause-binomial-absorption}
For every \(q\ge69\),
\[
\begin{aligned}
  &268435456\binom{2q+27}{1}
    +1073741824\binom{2q+27}{3}
    +4294967296\binom{2q+27}{5}\\
  &\quad+17179869184\binom{2q+27}{7}
    +68719476736\binom{2q+27}{9}
    +274877906944\binom{2q+27}{11}\\
  &\quad+1099511627776\binom{2q+27}{13}
    +4398046511104\binom{2q+27}{15}\\
  &\quad+17592186044416\binom{2q+27}{17}
    +70368744177664\binom{2q+27}{19}\\
  &\quad+281474976710656\binom{2q+27}{21}
    +1125899906842624\binom{2q+27}{23}\\
  &\quad+4503599627370496\binom{2q+27}{25}
    +18014398509481984\binom{2q+27}{27}
  \le \binom{2q+27}{q}.
\end{aligned}
\]
\end{lemma}

\begin{proof}
At \(q=69\), the left side is
\[
  130891286985753943793524421574917285603272818688
  <
  319620688063970456986413452647207180184213345250
  =
  \binom{165}{69}.
\]
The margin in this base comparison is
\[
  188729401078216513192889031072289894580940526562 .
\]
Let \(P(q)\) denote the displayed left side.  For \(1\le p\le27\),
\[
  \frac{\binom{2q+29}{p}}{\binom{2q+27}{p}}
  \le
  \frac{(2q+29)(2q+28)}{(2q+2)(2q+1)}
  <2
  \qquad(q\ge69),
\]
and the displayed strict inequality follows from \(4q^2-102q-808>0\).
Hence \(P(q+1)<2P(q)\).  On the other hand,
\[
  \frac{\binom{2q+29}{q+1}}{\binom{2q+27}{q}}
  =
  \frac{(2q+29)(2q+28)}{(q+1)(q+28)}
  >3,
\]
since the numerator minus three times the denominator is \(q^2+27q+728\).
Thus the binomial side grows by a larger factor than \(P(q)\), and induction
from \(q=69\) proves the inequality.
\end{proof}

\begin{lemma}[Exact twenty-seventh $C$ pause-arithmetic checks after the first twelve rows]\label{lem:bc-c-twentyseventh-frontier-pause-arithmetic-certificate}
For \(33\le q\le68\),
put
\[
\begin{aligned}
  P(q)={}&268435456\binom{2q+27}{1}
    +1073741824\binom{2q+27}{3}
    +4294967296\binom{2q+27}{5}\\
  &+17179869184\binom{2q+27}{7}
    +68719476736\binom{2q+27}{9}
    +274877906944\binom{2q+27}{11}\\
  &+1099511627776\binom{2q+27}{13}
    +4398046511104\binom{2q+27}{15}\\
  &+17592186044416\binom{2q+27}{17}
    +70368744177664\binom{2q+27}{19}\\
  &+281474976710656\binom{2q+27}{21}
    +1125899906842624\binom{2q+27}{23}\\
  &+4503599627370496\binom{2q+27}{25}
    +18014398509481984\binom{2q+27}{27}.
\end{aligned}
\]
Then
\[
  2^{2q+1}P(q)s_{q+27}\le s_{2q+27}.
\]
\end{lemma}

\begin{proof}
The exact C++/GMP/OpenMP replay
\[
\texttt{certificates/post\_m29/bc\_twentyseventh\_frontier\_gmp\_certificate.log}
\]
computes the displayed integer inequality for
\(33\le q\le68\).
It exits with status \(0\), reports \texttt{SUMMARY failures=0}, and has
SHA-256
\[
\hashsplit{67df773882805ff4321cd3659e16b1b9}{985db56a76a94d7136448011ef30f80d}.
\]
Its transcript contains exactly thirty-six \texttt{C\_MARGIN} records, one
for each displayed \(q\), and each record is the exact positive integer
\[
  s_{2q+27}-2^{2q+1}P(q)s_{q+27}.
\]
This proves the claim.
\end{proof}

\begin{proposition}[Residual $C$ twenty-seventh nonboundary lower correction]\label{prop:bc-c-twentyseventh-nonboundary-lower-correction}
For every residual row \(C_b\) with \(20\le b\le217\), put \(q=b+1\) and
write \(s_j=m_j^{\mathrm{st}}\) and \(m_j^{C_b}=s_j+\delta_j^{C_b}\).  Then
\[
  -\frac12 s_{2q+27}\le \delta_{2q+27}^{C_b}.
\]
\end{proposition}

\begin{proof}
As in the proof of
\cref{prop:bc-c-first-nonboundary-lower-correction},
\(|\delta_{2q+27}^{C_b}|\) is the number of column-even Pieri two-strip paths
of length \(2q+27\) with terminal height at least \(2q\).  If
\(21\le q\le32\),
\cref{lem:bc-twentyseventh-c-q21-q22-q23-q24-q25-q26-q27-q28-q29-q30-q31-q32-fpf-pause-certificate}
gives
\[
  |\delta_{2q+27}^{C_b}|\le s_{2q+27}/2 .
\]
For \(q\ge33\), \cref{lem:bc-c-twentyseventh-nonboundary-twentyseven-pause-code},
with \(P(q)\) denoting the polynomial from
\cref{lem:bc-c-twentyseventh-frontier-pause-arithmetic-certificate}, gives
\[
  |\delta_{2q+27}^{C_b}|
  \le
  P(q)\,2^{2q}s_{q+27}.
\]
If \(33\le q\le68\),
\cref{lem:bc-c-twentyseventh-frontier-pause-arithmetic-certificate} gives
\[
  |\delta_{2q+27}^{C_b}|\le s_{2q+27}/2 .
\]
If \(q\ge69\), \cref{lem:bc-twentyseventh-c-pause-binomial-absorption} gives
\[
  |\delta_{2q+27}^{C_b}|
  \le
  2^{2q}\binom{2q+27}{q}s_{q+27}.
\]
The exact GMP/OpenMP replay
\[
\texttt{certificates/post\_m29/bc\_pieri\_powerloss\_2q\_gmp\_certificate.log}
\]
checks the inequality
\[
  2^{2q+1}\binom{2q+27}{q}s_{q+27}\le s_{2q+27}
\]
on this residual row and moment as part of its all-row residual-window audit;
it reports \texttt{failures=0} and has SHA-256
\[
\hashsplit{124068c5daa1123088a3cd79c3f6b839}{fee45a1aaef637ed82161c465d86bf45}.
\]
Thus \(|\delta_{2q+27}^{C_b}|\le s_{2q+27}/2\) also for \(q\ge69\), which
implies the claimed lower bound.
\end{proof}

\begin{lemma}[$C$ twenty-eighth-nonboundary height-bad paths have a twenty-eight-pause code]\label{lem:bc-c-twentyeighth-nonboundary-twentyeight-pause-code}
For \(q\ge1\), the number of column-even Pieri two-strip paths
\[
  \varnothing=\lambda^{(0)}\subset\lambda^{(1)}\subset\cdots
  \subset\lambda^{(2q+28)}
\]
with
\[
  \ell(\lambda^{(2q+28)})\ge2q
\]
is at most
\[
\begin{aligned}
  \biggl(&268435456\binom{2q+28}{0}
    +1073741824\binom{2q+28}{2}
    +4294967296\binom{2q+28}{4}\\
  &\quad+17179869184\binom{2q+28}{6}
    +68719476736\binom{2q+28}{8}\\
  &\quad+274877906944\binom{2q+28}{10}
    +1099511627776\binom{2q+28}{12}\\
  &\quad+4398046511104\binom{2q+28}{14}
    +17592186044416\binom{2q+28}{16}\\
  &\quad+70368744177664\binom{2q+28}{18}
    +281474976710656\binom{2q+28}{20}\\
  &\quad+1125899906842624\binom{2q+28}{22}
    +4503599627370496\binom{2q+28}{24}\\
  &\quad+18014398509481984\binom{2q+28}{26}
    +72057594037927936\binom{2q+28}{28}\biggr)2^{2q}s_{q+28}.
\end{aligned}
\]
\end{lemma}

\begin{proof}
The terminal height is even and at most \(2q+28\).  Thus it is one of
\[
  2q+28,\ 2q+26,\ 2q+24,\ 2q+22,\ 2q+20,\ 2q+18,\ 2q+16,\
  2q+14,\ 2q+12,\ 2q+10,\ 2q+8,\ 2q+6,\ 2q+4,\ 2q+2,\ 2q .
\]
These cases have respectively \(0,2,4,\ldots,28\) steps which fail to create
a new row.  Deleting the first column, shifting left, and standardizing the
repeated pause labels injects the paths with \(p\) pauses into the choice of
the \(p\)-element pause set and an even-column standard tableau of size
\(2q+28+p\).  Hence the \(p\)-pause contribution is at most
\[
  \binom{2q+28}{p}2^{2q+28+p}s_{q+14+p/2}
  \qquad(p=0,2,4,\ldots,28)
\]
by
\cref{lem:bc-even-column-standard-tableaux-2n-bound,lem:bc-stable-moments-monotone-after-two}.
After replacing \(s_{q+14},s_{q+15},\ldots,s_{q+27}\) by \(s_{q+28}\), the
sum is the displayed bound.
\end{proof}

\begin{lemma}[The first thirteen twenty-eighth $C$ pause cases are absorbed by the fixed-point-free count]\label{lem:bc-twentyeighth-c-q21-q22-q23-q24-q25-q26-q27-q28-q29-q30-q31-q32-q33-fpf-pause-certificate}
For \(21\le q\le33\), the number of column-even
Pieri two-strip paths of length \(2q+28\) with terminal height at least
\(2q\) is at most \(s_{2q+28}/2\).
\end{lemma}

\begin{proof}
Use the pause-set injection from
\cref{lem:bc-c-twentyeighth-nonboundary-twentyeight-pause-code}, but count
the target even-column standard tableaux exactly by
\cref{lem:bc-even-column-standard-tableaux-fpf-count}.  For
\(21\le q\le33\), this gives the upper bound
\[
  \sum_{\substack{0\le p\le28\\ p\ {\rm even}}}
  \binom{2q+28}{p}(2q+27+p)!! .
\]
The exact C++/GMP/OpenMP replay
\[
\texttt{certificates/post\_m29/bc\_twentyeighth\_frontier\_gmp\_certificate.log}
\]
checks the displayed integer sums, exits with status \(0\), reports
\texttt{SUMMARY failures=0}, and has SHA-256
\[
\hashsplit{400e6b4c50b78e1c490ce93daca73864}{a2ebfc3380e93c50ac7d502205299b0e}.
\]
Its transcript contains exactly thirteen \texttt{C\_FPF\_MARGIN} records,
one for each displayed \(q\), and each record is the exact positive integer
\[
  s_{2q+28}-2\,\mathrm{FPF}(q,2q+28).
\]
This proves the half-stable bound for these first thirteen cases.
\end{proof}

\begin{lemma}[The central binomial factor absorbs the twenty-eight-pause loss]\label{lem:bc-twentyeighth-c-pause-binomial-absorption}
For every \(q\ge71\),
\[
\begin{aligned}
  &268435456\binom{2q+28}{0}
    +1073741824\binom{2q+28}{2}
    +4294967296\binom{2q+28}{4}\\
  &\quad+17179869184\binom{2q+28}{6}
    +68719476736\binom{2q+28}{8}
    +274877906944\binom{2q+28}{10}\\
  &\quad+1099511627776\binom{2q+28}{12}
    +4398046511104\binom{2q+28}{14}\\
  &\quad+17592186044416\binom{2q+28}{16}
    +70368744177664\binom{2q+28}{18}\\
  &\quad+281474976710656\binom{2q+28}{20}
    +1125899906842624\binom{2q+28}{22}\\
  &\quad+4503599627370496\binom{2q+28}{24}
    +18014398509481984\binom{2q+28}{26}
    +72057594037927936\binom{2q+28}{28}
  \le \binom{2q+28}{q}.
\end{aligned}
\]
\end{lemma}

\begin{proof}
At \(q=71\), the left side is
\[
  6422467389983358344404895865219565878768457220096
  <
  9143552349859580429564574407161791083419805934000
  =
  \binom{170}{71}.
\]
The margin in this base comparison is
\[
  2721084959876222085159678541942225204651348713904 .
\]
Let \(P(q)\) denote the displayed left side.  For \(0\le p\le28\),
\[
  \frac{\binom{2q+30}{p}}{\binom{2q+28}{p}}
  \le
  \frac{(2q+30)(2q+29)}{(2q+2)(2q+1)}
  <2
  \qquad(q\ge71),
\]
and the displayed strict inequality follows from \(4q^2-106q-866>0\).
Hence \(P(q+1)<2P(q)\).  On the other hand,
\[
  \frac{\binom{2q+30}{q+1}}{\binom{2q+28}{q}}
  =
  \frac{(2q+30)(2q+29)}{(q+1)(q+29)}
  >3,
\]
since the numerator minus three times the denominator is \(q^2+28q+783\).
Thus the binomial side grows by a larger factor than \(P(q)\), and induction
from \(q=71\) proves the inequality.
\end{proof}

\begin{lemma}[Exact twenty-eighth $C$ pause-arithmetic checks after the first thirteen rows]\label{lem:bc-c-twentyeighth-frontier-pause-arithmetic-certificate}
For \(34\le q\le70\),
put
\[
\begin{aligned}
  P(q)={}&268435456\binom{2q+28}{0}
    +1073741824\binom{2q+28}{2}
    +4294967296\binom{2q+28}{4}\\
  &+17179869184\binom{2q+28}{6}
    +68719476736\binom{2q+28}{8}
    +274877906944\binom{2q+28}{10}\\
  &+1099511627776\binom{2q+28}{12}
    +4398046511104\binom{2q+28}{14}\\
  &+17592186044416\binom{2q+28}{16}
    +70368744177664\binom{2q+28}{18}\\
  &+281474976710656\binom{2q+28}{20}
    +1125899906842624\binom{2q+28}{22}\\
  &+4503599627370496\binom{2q+28}{24}
    +18014398509481984\binom{2q+28}{26}
    +72057594037927936\binom{2q+28}{28}.
\end{aligned}
\]
Then
\[
  2^{2q+1}P(q)s_{q+28}\le s_{2q+28}.
\]
\end{lemma}

\begin{proof}
The exact C++/GMP/OpenMP replay
\[
\texttt{certificates/post\_m29/bc\_twentyeighth\_frontier\_gmp\_certificate.log}
\]
computes the displayed integer inequality for
\(34\le q\le70\).
It exits with status \(0\), reports \texttt{SUMMARY failures=0}, and has
SHA-256
\[
\hashsplit{400e6b4c50b78e1c490ce93daca73864}{a2ebfc3380e93c50ac7d502205299b0e}.
\]
Its transcript contains exactly thirty-seven \texttt{C\_MARGIN} records, one
for each displayed \(q\), and each record is the exact positive integer
\[
  s_{2q+28}-2^{2q+1}P(q)s_{q+28}.
\]
This proves the claim.
\end{proof}

\begin{proposition}[Residual $C$ twenty-eighth nonboundary lower correction]\label{prop:bc-c-twentyeighth-nonboundary-lower-correction}
For every residual row \(C_b\) with \(20\le b\le217\), put \(q=b+1\) and
write \(s_j=m_j^{\mathrm{st}}\) and \(m_j^{C_b}=s_j+\delta_j^{C_b}\).  Then
\[
  -\frac12 s_{2q+28}\le \delta_{2q+28}^{C_b}.
\]
\end{proposition}

\begin{proof}
As in the proof of
\cref{prop:bc-c-first-nonboundary-lower-correction},
\(|\delta_{2q+28}^{C_b}|\) is the number of column-even Pieri two-strip paths
of length \(2q+28\) with terminal height at least \(2q\).  If
\(21\le q\le33\),
\cref{lem:bc-twentyeighth-c-q21-q22-q23-q24-q25-q26-q27-q28-q29-q30-q31-q32-q33-fpf-pause-certificate}
gives
\[
  |\delta_{2q+28}^{C_b}|\le s_{2q+28}/2 .
\]
For \(q\ge34\), \cref{lem:bc-c-twentyeighth-nonboundary-twentyeight-pause-code},
with \(P(q)\) denoting the polynomial from
\cref{lem:bc-c-twentyeighth-frontier-pause-arithmetic-certificate}, gives
\[
  |\delta_{2q+28}^{C_b}|
  \le
  P(q)\,2^{2q}s_{q+28}.
\]
If \(34\le q\le70\),
\cref{lem:bc-c-twentyeighth-frontier-pause-arithmetic-certificate} gives
\[
  |\delta_{2q+28}^{C_b}|\le s_{2q+28}/2 .
\]
If \(q\ge71\), \cref{lem:bc-twentyeighth-c-pause-binomial-absorption} gives
\[
  |\delta_{2q+28}^{C_b}|
  \le
  2^{2q}\binom{2q+28}{q}s_{q+28}.
\]
The exact GMP/OpenMP replay
\[
\texttt{certificates/post\_m29/bc\_pieri\_powerloss\_2q\_gmp\_certificate.log}
\]
checks the inequality
\[
  2^{2q+1}\binom{2q+28}{q}s_{q+28}\le s_{2q+28}
\]
on this residual row and moment as part of its all-row residual-window audit;
it reports \texttt{failures=0} and has SHA-256
\[
\hashsplit{124068c5daa1123088a3cd79c3f6b839}{fee45a1aaef637ed82161c465d86bf45}.
\]
Thus \(|\delta_{2q+28}^{C_b}|\le s_{2q+28}/2\) also for \(q\ge71\), which
implies the claimed lower bound.
\end{proof}

\begin{lemma}[$C$ twenty-ninth-nonboundary height-bad paths have a twenty-nine-pause code]\label{lem:bc-c-twentyninth-nonboundary-twentynine-pause-code}
For \(q\ge1\), the number of column-even Pieri two-strip paths
\[
  \varnothing=\lambda^{(0)}\subset\lambda^{(1)}\subset\cdots
  \subset\lambda^{(2q+29)}
\]
with
\[
  \ell(\lambda^{(2q+29)})\ge2q
\]
is at most
\[
\begin{aligned}
  \biggl(&1073741824\binom{2q+29}{1}
    +4294967296\binom{2q+29}{3}
    +17179869184\binom{2q+29}{5}\\
  &\quad+68719476736\binom{2q+29}{7}
    +274877906944\binom{2q+29}{9}\\
  &\quad+1099511627776\binom{2q+29}{11}
    +4398046511104\binom{2q+29}{13}\\
  &\quad+17592186044416\binom{2q+29}{15}
    +70368744177664\binom{2q+29}{17}\\
  &\quad+281474976710656\binom{2q+29}{19}
    +1125899906842624\binom{2q+29}{21}\\
  &\quad+4503599627370496\binom{2q+29}{23}
    +18014398509481984\binom{2q+29}{25}\\
  &\quad+72057594037927936\binom{2q+29}{27}
    +288230376151711744\binom{2q+29}{29}\biggr)2^{2q}s_{q+29}.
\end{aligned}
\]
\end{lemma}

\begin{proof}
The terminal height is even and at most \(2q+29\).  Thus it is one of
\[
  2q+28,\ 2q+26,\ 2q+24,\ 2q+22,\ 2q+20,\ 2q+18,\ 2q+16,\
  2q+14,\ 2q+12,\ 2q+10,\ 2q+8,\ 2q+6,\ 2q+4,\ 2q+2,\ 2q .
\]
These cases have respectively \(1,3,5,\ldots,29\) steps which fail to create
a new row.  Deleting the first column, shifting left, and standardizing the
repeated pause labels injects the paths with \(p\) pauses into the choice of
the \(p\)-element pause set and an even-column standard tableau of size
\(2q+29+p\).  Hence the \(p\)-pause contribution is at most
\[
  \binom{2q+29}{p}2^{2q+29+p}s_{q+14+(p+1)/2}
  \qquad(p=1,3,5,\ldots,29)
\]
by
\cref{lem:bc-even-column-standard-tableaux-2n-bound,lem:bc-stable-moments-monotone-after-two}.
After replacing \(s_{q+15},s_{q+16},\ldots,s_{q+28}\) by \(s_{q+29}\), the
sum is the displayed bound.
\end{proof}

\begin{lemma}[The first fourteen twenty-ninth $C$ pause cases are absorbed by the fixed-point-free count]\label{lem:bc-twentyninth-c-q21-q34-fpf-pause-certificate}
For \(21\le q\le34\), the number of column-even
Pieri two-strip paths of length \(2q+29\) with terminal height at least
\(2q\) is at most \(s_{2q+29}/2\).
\end{lemma}

\begin{proof}
Use the pause-set injection from
\cref{lem:bc-c-twentyninth-nonboundary-twentynine-pause-code}, but count the
target even-column standard tableaux exactly by
\cref{lem:bc-even-column-standard-tableaux-fpf-count}.  For
\(21\le q\le34\), this gives the upper bound
\[
  \sum_{\substack{1\le p\le29\\ p\ {\rm odd}}}
  \binom{2q+29}{p}(2q+28+p)!! .
\]
The exact C++/GMP replay
\[
\texttt{certificates/post\_m29/bc\_twentyninth\_c\_frontier\_gmp\_certificate.log}
\]
checks the displayed integer sums, exits with status \(0\), reports
\texttt{SUMMARY failures=0}, and has SHA-256
\[
\hashsplit{deaa91314bb4c09ba9e41a2c5b3586dc}{b434b76f15943a8870c4acb5b8cdf72f}.
\]
Its transcript contains exactly fourteen \texttt{C\_FPF\_MARGIN} records,
one for each displayed \(q\), and each record is the exact positive integer
\[
  s_{2q+29}-2\,\mathrm{FPF}(q,2q+29).
\]
This proves the half-stable bound for these first fourteen cases.
\end{proof}

\begin{lemma}[The central binomial factor absorbs the twenty-nine-pause loss]\label{lem:bc-twentyninth-c-pause-binomial-absorption}
For every \(q\ge74\),
\[
\begin{aligned}
  &1073741824\binom{2q+29}{1}
    +4294967296\binom{2q+29}{3}
    +17179869184\binom{2q+29}{5}\\
  &\quad+68719476736\binom{2q+29}{7}
    +274877906944\binom{2q+29}{9}
    +1099511627776\binom{2q+29}{11}\\
  &\quad+4398046511104\binom{2q+29}{13}
    +17592186044416\binom{2q+29}{15}\\
  &\quad+70368744177664\binom{2q+29}{17}
    +281474976710656\binom{2q+29}{19}\\
  &\quad+1125899906842624\binom{2q+29}{21}
    +4503599627370496\binom{2q+29}{23}\\
  &\quad+18014398509481984\binom{2q+29}{25}
    +72057594037927936\binom{2q+29}{27}
    +288230376151711744\binom{2q+29}{29}
  \le \binom{2q+29}{q}.
\end{aligned}
\]
\end{lemma}

\begin{proof}
At \(q=74\), the left side is
\[
  450783285331425602968260277664824699027879842807808
  <
  1069333076921382630709408439483450169458590305354000
  =
  \binom{177}{74}.
\]
The margin in this base comparison is
\[
  618549791589957027741148161818625470430710462546192 .
\]
Let \(P(q)\) denote the displayed left side.  For \(1\le p\le29\),
\[
  \frac{\binom{2q+31}{p}}{\binom{2q+29}{p}}
  \le
  \frac{(2q+31)(2q+30)}{(2q+2)(2q+1)}
  <2
  \qquad(q\ge74),
\]
and the displayed strict inequality follows from \(4q^2-110q-926>0\).
Hence \(P(q+1)<2P(q)\).  On the other hand,
\[
  \frac{\binom{2q+31}{q+1}}{\binom{2q+29}{q}}
  =
  \frac{(2q+31)(2q+30)}{(q+1)(q+30)}
  >3,
\]
since the numerator minus three times the denominator is \(q^2+29q+840\).
Thus the binomial side grows by a larger factor than \(P(q)\), and induction
from \(q=74\) proves the inequality.
\end{proof}

\begin{lemma}[Exact twenty-ninth $C$ pause-arithmetic checks after the first fourteen rows]\label{lem:bc-c-twentyninth-frontier-pause-arithmetic-certificate}
For \(35\le q\le73\),
put
\[
\begin{aligned}
  P(q)={}&1073741824\binom{2q+29}{1}
    +4294967296\binom{2q+29}{3}
    +17179869184\binom{2q+29}{5}\\
  &+68719476736\binom{2q+29}{7}
    +274877906944\binom{2q+29}{9}
    +1099511627776\binom{2q+29}{11}\\
  &+4398046511104\binom{2q+29}{13}
    +17592186044416\binom{2q+29}{15}\\
  &+70368744177664\binom{2q+29}{17}
    +281474976710656\binom{2q+29}{19}\\
  &+1125899906842624\binom{2q+29}{21}
    +4503599627370496\binom{2q+29}{23}\\
  &+18014398509481984\binom{2q+29}{25}
    +72057594037927936\binom{2q+29}{27}
    +288230376151711744\binom{2q+29}{29}.
\end{aligned}
\]
Then
\[
  2^{2q+1}P(q)s_{q+29}\le s_{2q+29}.
\]
\end{lemma}

\begin{proof}
The exact C++/GMP replay
\[
\texttt{certificates/post\_m29/bc\_twentyninth\_c\_frontier\_gmp\_certificate.log}
\]
computes the displayed integer inequality for
\(35\le q\le73\).
It exits with status \(0\), reports \texttt{SUMMARY failures=0}, and has
SHA-256
\[
\hashsplit{deaa91314bb4c09ba9e41a2c5b3586dc}{b434b76f15943a8870c4acb5b8cdf72f}.
\]
Its transcript contains exactly thirty-nine \texttt{C\_MARGIN} records, one
for each displayed \(q\), and each record is the exact positive integer
\[
  s_{2q+29}-2^{2q+1}P(q)s_{q+29}.
\]
This proves the claim.
\end{proof}

\begin{proposition}[Residual $C$ twenty-ninth nonboundary lower correction]\label{prop:bc-c-twentyninth-nonboundary-lower-correction}
For every residual row \(C_b\) with \(20\le b\le217\), put \(q=b+1\) and
write \(s_j=m_j^{\mathrm{st}}\) and \(m_j^{C_b}=s_j+\delta_j^{C_b}\).  Then
\[
  -\frac12 s_{2q+29}\le \delta_{2q+29}^{C_b}.
\]
\end{proposition}

\begin{proof}
As in the proof of
\cref{prop:bc-c-first-nonboundary-lower-correction},
\(|\delta_{2q+29}^{C_b}|\) is the number of column-even Pieri two-strip paths
of length \(2q+29\) with terminal height at least \(2q\).  If
\(21\le q\le34\),
\cref{lem:bc-twentyninth-c-q21-q34-fpf-pause-certificate} gives
\[
  |\delta_{2q+29}^{C_b}|\le s_{2q+29}/2 .
\]
For \(q\ge35\), \cref{lem:bc-c-twentyninth-nonboundary-twentynine-pause-code},
with \(P(q)\) denoting the polynomial from
\cref{lem:bc-c-twentyninth-frontier-pause-arithmetic-certificate}, gives
\[
  |\delta_{2q+29}^{C_b}|
  \le
  P(q)\,2^{2q}s_{q+29}.
\]
If \(35\le q\le73\),
\cref{lem:bc-c-twentyninth-frontier-pause-arithmetic-certificate} gives
\[
  |\delta_{2q+29}^{C_b}|\le s_{2q+29}/2 .
\]
If \(q\ge74\), \cref{lem:bc-twentyninth-c-pause-binomial-absorption} gives
\[
  |\delta_{2q+29}^{C_b}|
  \le
  2^{2q}\binom{2q+29}{q}s_{q+29}.
\]
The exact GMP/OpenMP replay
\[
\texttt{certificates/post\_m29/bc\_pieri\_powerloss\_2q\_gmp\_certificate.log}
\]
checks the inequality
\[
  2^{2q+1}\binom{2q+29}{q}s_{q+29}\le s_{2q+29}
\]
on this residual row and moment as part of its all-row residual-window audit;
it reports \texttt{failures=0} and has SHA-256
\[
\hashsplit{124068c5daa1123088a3cd79c3f6b839}{fee45a1aaef637ed82161c465d86bf45}.
\]
Thus \(|\delta_{2q+29}^{C_b}|\le s_{2q+29}/2\) also for \(q\ge74\), which
implies the claimed lower bound.
\end{proof}

\ifdefined\GinibreDevelopmentArchive
\subsection*{Development archive: superseded residual-tail routes}

\begin{remark}[Status of the following development archive]
The results from this point through the return marker below record conditional
compression suppliers, obstructions, and superseded decoder constructions.
They are retained to document why those routes were abandoned.  No result in
this block is an incoming node of
\cref{prop:bc-active-correction-prefix-contract,prop:post29-bc-residual-closure,thm:classical}.
In particular, a proposition in this block whose statement
begins with ``suppose'' remains only a conditional implication and is not an
accepted proof premise.  The active proof resumes at the explicitly marked
half-stable bridge.
\end{remark}

\begin{lemma}[Rank-free Pieri compression fails in low rank]\label{lem:bc-rank-free-pieri-compression-obstruction}
Let \(s_n\) be the stable Littlewood moment sequence of
\cref{lem:bc-littlewood-graph-model}.  In content \((2^7)\), the two bad-shape
sums
\[
  A^w=\sum_{\substack{\lambda':\,\mathrm{even}\\ \lambda_1\ge6}}a_7(\lambda),
  \qquad
  A^h=\sum_{\substack{\lambda':\,\mathrm{even}\\ \ell(\lambda)\ge6}}a_7(\lambda)
\]
both equal \(211\), whereas
\[
  \binom{7}{3}s_4=210 .
\]
Equivalently, by \cref{lem:bc-correction-bad-shapes},
\[
  |\delta_7^{B_2}|=|\delta_7^{C_2}|=211 .
\]
Consequently the rank-plus-one Pieri-compression estimate
\[
  \sum_{\substack{\lambda':\,\mathrm{even}\\ \lambda_1\ge2q}}a_j(\lambda)
  \le \binom jq\,s_{j-q},
  \qquad
  \sum_{\substack{\lambda':\,\mathrm{even}\\ \ell(\lambda)\ge2q}}a_j(\lambda)
  \le \binom jq\,s_{j-q},
\]
cannot hold as a rank-free theorem for all \(q,j\).
\end{lemma}

\begin{proof}
By \cref{lem:bc-pieri-path-bijection}, \(a_7(\lambda)\) is the number of
seven-step horizontal-two-strip Pieri paths from \(\varnothing\) to
\(\lambda\).  Applying this finite recursion gives the following nonzero
column-even terminal contributions with \(\lambda_1\ge6\):
\[
\begin{array}{c|c}
\lambda & a_7(\lambda)\\ \hline
(7,7)&36\\
(6,6,1,1)&175 .
\end{array}
\]
Thus \(A^w=36+175=211\).  The same recursion gives the following nonzero
column-even terminal contributions with \(\ell(\lambda)\ge6\):
\[
\begin{array}{c|c}
\lambda & a_7(\lambda)\\ \hline
(5,5,1,1,1,1)&49\\
(4,4,2,2,1,1)&126\\
(3,3,2,2,2,2)&15\\
(3,3,3,3,1,1)&21 .
\end{array}
\]
Hence \(A^h=49+126+15+21=211\).

Finally, \cref{lem:bc-stable-moment-recurrence} gives
\[
  s_0=1,\quad s_1=0,\quad s_2=1,\quad s_3=1,\quad s_4=6,
\]
so \(\binom{7}{3}s_4=35\cdot6=210\).  Taking \(q=3\) and \(j=7\) contradicts
both displayed rank-free estimates.
\end{proof}

\begin{proposition}[Rank-plus-one Pieri compression would imply the residual $B/C$ tail]\label{prop:bc-pieri-compression-arithmetic-tail}
For the rows
\[
  B_{14},\ldots,B_{61},\qquad C_{20},\ldots,C_{61},
\]
put $q=b+1$ for $G=B_b$ or $G=C_b$.  Suppose that, for every such row and
every $j$ in the residual window of \cref{rem:former-post29-tail}, the correction
satisfies
\[
  |\delta_j^G|\le \binom jq\,s_{j-q}.
\]
Then \cref{rem:former-post29-tail} holds.
\end{proposition}

\begin{proof}
The exact GMP/OpenMP replay
\[
\texttt{certificates/post\_m29/bc\_pieri\_compression\_gmp\_certificate.log}
\]
checks, over all $402$ rows of the original $B/C$ range and all corresponding
moments, the
integer inequalities
\[
  2\binom jq\,s_{j-q}\le s_j.
\]
It checks $3{,}904{,}626$ inequalities, reports \texttt{failures=0}, and has
SHA-256
\[
\hashsplit{af7d338f58854db62c6f17ca23975e17}{5aee5e2206ef7c2386ead79875007558}.
\]
The replay reports its largest ratio diagnostic at $B_{14}$, moment $30$:
$\log_{10}(2\binom jq\,s_{j-q}/s_j)\approx-11.6664415829297177$.  The exact
checked integer inequality, rather than this rounded diagnostic, shows that the
displayed compression hypothesis gives $|\delta_j^G|\le s_j/2$ throughout
the residual windows.  Since \cref{prop:bc-pieri-correction-sign} gives
$\delta_j^G\le0$, this is the lower bound $\delta_j^G\ge -s_j/2$ required in
\cref{rem:former-post29-tail}.
\end{proof}

\begin{proposition}[$q$-bit-loss Pieri compression would imply the residual $B/C$ tail]\label{prop:bc-pieri-compression-qbit-arithmetic-tail}
For the rows
\[
  B_{14},\ldots,B_{61},\qquad C_{20},\ldots,C_{61},
\]
put $q=b+1$ for $G=B_b$ or $G=C_b$.  Suppose that, for every such row and
every $j$ in the residual window of \cref{rem:former-post29-tail}, the correction
satisfies
\[
  |\delta_j^G|\le 2^q\binom jq\,s_{j-q}.
\]
Then \cref{rem:former-post29-tail} holds.
\end{proposition}

\begin{proof}
The exact GMP/OpenMP replay
\[
\texttt{certificates/post\_m29/bc\_pieri\_compression\_pow2\_gmp\_certificate.log}
\]
checks, over all $402$ rows of the original $B/C$ range and all corresponding
moments, the
integer inequalities
\[
  2^{q+1}\binom jq\,s_{j-q}\le s_j.
\]
It checks $3{,}904{,}626$ inequalities, reports \texttt{failures=0}, and has
SHA-256
\[
\hashsplit{a98540c3b38d65a7c414898302214411}{2663801ab59fe2cb8deec1b34f6c5f7e}.
\]
The replay again reports its largest ratio diagnostic at $B_{14}$, moment $30$:
$\log_{10}(2^{q+1}\binom jq\,s_{j-q}/s_j)\approx-7.15099164796999887$.  The
exact checked integer inequality, rather than this rounded diagnostic, shows that the
displayed controlled-loss compression hypothesis gives
$|\delta_j^G|\le s_j/2$ throughout the residual windows.  Combining this with
\cref{prop:bc-pieri-correction-sign} gives the lower bound
$\delta_j^G\ge -s_j/2$, which is exactly \cref{rem:former-post29-tail}.
\end{proof}

\begin{proposition}[Residual $B/C$ bad-shape reduction]\label{prop:bc-bad-shape-reduction}
For the rows
\[
  B_{14},\ldots,B_{61},\qquad C_{20},\ldots,C_{61},
\]
put $q=b+1$ for $G=B_b$ or $G=C_b$.  Suppose that, for every such row and
every $j$ in the residual window of \cref{rem:former-post29-tail}, the corresponding
bad-shape estimate
\[
  \sum_{\substack{\lambda':\,\mathrm{even}\\ \lambda_1\ge 2q}}
    a_j(\lambda)
  \le
  2^{2q}\binom jq\,s_{j-q}
\]
holds on the $B_b$ side, and
\[
  \sum_{\substack{\lambda':\,\mathrm{even}\\ \ell(\lambda)\ge 2q}}
    a_j(\lambda)
  \le
  2^{2q}\binom jq\,s_{j-q}
\]
holds on the $C_b$ side.  Then \cref{rem:former-post29-tail} holds.
\end{proposition}

\begin{proof}
By \cref{lem:bc-correction-bad-shapes}, the displayed bad-shape estimates
give
\[
  |\delta_j^G|\le 2^{2q}\binom jq\,s_{j-q}
\]
for the relevant residual row and moment.  The controlled-loss exact
GMP/OpenMP replay
\[
\texttt{certificates/post\_m29/bc\_pieri\_powerloss\_2q\_gmp\_certificate.log}
\]
checks, over all residual rows and windows, the integer inequalities
\[
  2^{2q+1}\binom jq\,s_{j-q}\le s_j.
\]
It reports \texttt{failures=0} and has SHA-256
\[
\hashsplit{124068c5daa1123088a3cd79c3f6b839}{fee45a1aaef637ed82161c465d86bf45}.
\]
The replay reports its largest ratio diagnostic at $B_{14}$, moment $30$, as
the approximation
$\log_{10}(2^{2q+1}\binom jq\,s_{j-q}/s_j)
\approx-2.63554171301028006$.  The conclusion uses the exact checked integer
inequality, not this rounded display, and therefore gives
$|\delta_j^G|\le s_j/2$.  Together with the sign
$\delta_j^G\le0$ from \cref{prop:bc-pieri-correction-sign}, this is exactly
$-s_j/2\le\delta_j^G$ on the residual window.
\end{proof}

\begin{corollary}[Pieri crossing estimates imply the residual $B/C$ tail]\label{cor:bc-pieri-crossing-estimates-tail}
For the rows
\[
  B_{14},\ldots,B_{61},\qquad C_{20},\ldots,C_{61},
\]
put $q=b+1$ for $G=B_b$ or $G=C_b$.  Suppose that, for every such row, every
$j$ in the residual window of \cref{rem:former-post29-tail}, and every
$q$-element subset $S\subset\{1,\ldots,j\}$, the following direct
Pieri-crossing estimates hold:
\[
\#\{\lambda^\bullet:\lambda^\bullet\text{ is a column-even Pieri two-strip
path of length }j,\ C^w_q(\lambda^\bullet)=S,\ \lambda^{(j)}_1\ge2q\}
\le 2^{2q}s_{j-q}
\]
for the $B_b$ rows, and
\[
\#\{\lambda^\bullet:\lambda^\bullet\text{ is a column-even Pieri two-strip
path of length }j,\ C^h_q(\lambda^\bullet)=S\}
\le 2^{2q}s_{j-q}
\]
for the $C_b$ rows.  Then \cref{rem:former-post29-tail} holds.
\end{corollary}

\begin{proof}
Fix a residual $B_b$ row and an admissible $j$.  By
\cref{cor:bc-fixed-fibers-pieri-paths}, the width fixed-witness fiber with
witness $S$ is equinumerous with the first displayed Pieri-crossing set.
Summing the assumed bound over all $q$-element subsets
$S\subset\{1,\ldots,j\}$ gives
\[
  \sum_{\substack{\lambda':\,\mathrm{even}\\ \lambda_1\ge 2q}}
    a_j(\lambda)
  \le
  2^{2q}\binom jq\,s_{j-q}.
\]
Thus the $B_b$ bad-shape estimate required in
\cref{prop:bc-bad-shape-reduction} holds.

The $C_b$ case is the same, using the height version of
\cref{cor:bc-fixed-fibers-pieri-paths}.  For each witness set $S$, the height
fixed-witness fiber is equinumerous with the second displayed Pieri-crossing
set, so summing over $S$ gives
\[
  \sum_{\substack{\lambda':\,\mathrm{even}\\ \ell(\lambda)\ge 2q}}
    a_j(\lambda)
  \le
  2^{2q}\binom jq\,s_{j-q}.
\]
Hence the $C_b$ bad-shape estimate required in
\cref{prop:bc-bad-shape-reduction} also holds.  Since the residual row and
moment were arbitrary, that proposition gives \cref{rem:former-post29-tail}.
\end{proof}

\begin{proposition}[Replacement-zone supplier implies the residual $B/C$ tail]\label{prop:bc-replacement-zone-tail-supplier}
For each row
\[
  B_{14},\ldots,B_{61},\qquad C_{20},\ldots,C_{61},
\]
put $q=b+1$ for $G=B_b$ or $G=C_b$.  Suppose that for every such row, every
$j$ in the residual window of \cref{rem:former-post29-tail}, and every
$q$-element subset $S\subset\{1,\ldots,j\}$, the following replacement-zone
repair data exist for the width fixed-witness fiber when $G=B_b$, and for the
height fixed-witness fiber when $G=C_b$:
\begin{enumerate}
\item deterministic replacement zones $Z_\alpha$ satisfying the hypotheses of
\cref{prop:bc-replacement-zones-global-recovery};
\item a repaired path $P^\ast(T)\in\mathcal P_{j-q}^{\mathrm{even}}$ for each
tableau $T$ in the fiber; and
\item branch words
\[
  \beta_\alpha(T)\in
  \{0,1\}^{2|Z_\alpha\cap I(S)|},
  \qquad I(S)=\{j-s+1:s\in S\},
\]
for which the local inverses required in
\cref{prop:bc-replacement-zones-global-recovery} recover the original
replacement-zone subpaths.
\end{enumerate}
Then \cref{rem:former-post29-tail} holds.
\end{proposition}

\begin{proof}
Fix a $B_b$ residual row, an admissible $j$, and one
$q$-element witness set $S$.  Applying
\cref{prop:bc-replacement-zone-fiber-criterion} to the width fiber gives
\[
 \#\{T:\lambda(T)' \text{ even},\ \lambda_1(T)\ge2q,\ W_q(T)=S\}
 \le 2^{2q}s_{j-q},
\]
and summing over $S$ as in
\cref{prop:bc-bad-shape-fiber-reduction} gives the $B_b$ bad-shape estimate.

The same argument for a $C_b$ residual row applies
\cref{prop:bc-replacement-zone-fiber-criterion} to the height fiber:
\[
 \#\{T:\lambda(T)' \text{ even},\ \ell(\lambda(T))\ge2q,\ H_q(T)=S\}
 \le 2^{2q}s_{j-q}.
\]
Summing over $S$ gives the $C_b$ bad-shape estimate.  Since the row and moment
were arbitrary in the residual lists and residual windows,
\cref{prop:bc-bad-shape-reduction} applies and yields
\cref{rem:former-post29-tail}.
\end{proof}

\begin{proposition}[Canonical halo repairs imply the residual $B/C$ tail]\label{prop:bc-canonical-halo-tail-supplier}
For each row
\[
  B_{14},\ldots,B_{61},\qquad C_{20},\ldots,C_{61},
\]
put $q=b+1$ for $G=B_b$ or $G=C_b$.  Suppose that for every such row, every
$j$ in the residual window of \cref{rem:former-post29-tail}, every
$q$-element subset $S\subset\{1,\ldots,j\}$, the canonical halo zones of
\cref{lem:bc-canonical-halo-replacement-zones} satisfy the following repair
requirements.
\begin{enumerate}
\item For the width fixed-witness fiber when $G=B_b$, they admit the parser,
target path, replacement windows, branch words, and local inverses listed in
\cref{prop:bc-canonical-halo-local-repair-criterion}.
\item For the height fixed-witness fiber when $G=C_b$, either the same local
repair data exist, or the canonical halo zones can be ordered as
$Z_1,\ldots,Z_L$ with $Z_L=\{a_L,\ldots,j\}$ and satisfy the right-boundary
auxiliary hypotheses of
\cref{prop:bc-right-boundary-auxiliary-fiber-criterion}.
\end{enumerate}
Then \cref{rem:former-post29-tail} holds.
\end{proposition}

\begin{proof}
Fix one residual row, one admissible $j$, and one witness set $S$.  In the
$B_b$ case, apply \cref{prop:bc-canonical-halo-local-repair-criterion} to the
width fixed-witness fiber.  It gives
\[
 \#\{T:\lambda(T)' \text{ even},\ \lambda_1(T)\ge2q,\ W_q(T)=S\}
 \le 2^{2q}s_{j-q}.
\]
In the $C_b$ case, apply
\cref{prop:bc-canonical-halo-local-repair-criterion} if the local repair data
of item (2) are available, and otherwise apply
\cref{prop:bc-right-boundary-auxiliary-fiber-criterion} to the ordered
canonical halo zones supplied by item (2).  In both alternatives one obtains
\[
 \#\{T:\lambda(T)' \text{ even},\ \ell(\lambda(T))\ge2q,\ H_q(T)=S\}
 \le 2^{2q}s_{j-q}.
\]
Since these estimates hold for every residual row, every admissible $j$, and
every $S$, \cref{prop:bc-bad-shape-fiber-reduction} gives the corresponding
bad-shape estimates.  Then \cref{prop:bc-bad-shape-reduction} yields
\cref{rem:former-post29-tail}.
\end{proof}

\begin{proposition}[Split target-window canonical halo repairs imply the residual $B/C$ tail]\label{prop:bc-split-target-window-halo-tail-supplier}
For each row
\[
  B_{14},\ldots,B_{61},\qquad C_{20},\ldots,C_{61},
\]
put $q=b+1$ for $G=B_b$ or $G=C_b$.  Suppose that for every such row, every
$j$ in the residual window of \cref{rem:former-post29-tail}, every
$q$-element subset $S\subset\{1,\ldots,j\}$, the canonical halo zones of
\cref{lem:bc-canonical-halo-replacement-zones} satisfy the following repair
requirements.
\begin{enumerate}
\item For the width fixed-witness fiber when $G=B_b$, they admit the parser,
target path, replacement windows, branch words, and local inverses listed in
\cref{prop:bc-canonical-halo-local-repair-criterion}.
\item For the height fixed-witness fiber when $G=C_b$, either the same local
repair data exist, or the canonical halo zones can be ordered as
$Z_1,\ldots,Z_L$ with $Z_L=\{a_L,\ldots,j\}$ so that the hypotheses of
\cref{prop:bc-right-boundary-auxiliary-global-recovery} hold for the parser,
target path, endpoint pairs, zones strictly left of $Z_L$, and their local
inverses.  For the final zone, write
\[
  Z_L=\{j-2m+1,\ldots,j\}.
\]
Assume it is in the non-full tight-height setting of
\cref{cor:bc-height-zero-surplus-label-prefix}, and that the target-window data
\[
  S,\quad Z_L,\quad V^\ast_L,\quad \mathcal D_{\mathrm{out}}(P),
  \quad \rho_{2a_L-2},\rho_{2j}
\]
determine $\lambda^{(2m)}$, the retained even-label strips
$E_1,\ldots,E_m$, and the labelled continuation boxes.  If \(m\ge29\), let
\[
  \Omega_{\mathrm{tw}}
  =
  \Omega(\lambda^{(2m)};E_1,\ldots,E_m)
\]
with the lexicographic order of \cref{def:bc-tight-odd-packet-omega}; when
\(\Omega_{\mathrm{tw}}\ne\varnothing\), let
\((\gamma;e_1,\ldots,e_{m-1})\) be the associated fixed-even datum of
\cref{def:bc-odd-fixed-even-standard-tableaux}.  In the \(m\ge29\) cases not
already covered by the condition
\[
  \Omega_{\mathrm{tw}}=\varnothing
  \quad\text{or}\quad
  \bigl(\Omega_{\mathrm{tw}}\ne\varnothing
  \quad\text{and}\quad
  \bigl(m=29
  \quad\text{or}\quad
  m=30
  \quad\text{or}\quad
  m=31
  \quad\text{or}\quad
  (\gamma,e_{m-1})\notin\mathcal P^{\mathrm{env}}_m
  \quad\text{or}\quad
  (\gamma,e_{m-1},e_{m-2})\notin
  \mathcal P^{\mathrm{env},2}_m
  \quad\text{or}\quad
  \exists d,\ 3\le d\le m-1,\
  (\gamma,e_{m-1},e_{m-2},\ldots,e_{m-d})
  \notin\mathcal P^{\mathrm{env},d}_m\bigr)\bigr),
\]
assume in addition that the same data determine either the commuting
normal-form signature \(\kappa\) of the true fixed-even standard tableau, or a
canonically ordered finite set
\[
  \Sigma_{\mathrm{tw}}
\]
of commuting normal-form signature words such that the true fixed-even standard
tableau has signature in \(\Sigma_{\mathrm{tw}}\) and
\[
  |\Sigma_{\mathrm{tw}}|\le2^m .
\]
\end{enumerate}
Then \cref{rem:former-post29-tail} holds.
\end{proposition}

\begin{proof}
We check the two alternatives in
\cref{prop:bc-canonical-halo-tail-supplier}.  The $B_b$ width hypothesis is
exactly item~(1) of that proposition.

For a $C_b$ height fiber, the first alternative in item~(2) is again exactly
the local-repair alternative of
\cref{prop:bc-canonical-halo-tail-supplier}.  In the right-boundary alternative,
\cref{prop:bc-split-target-window-supplies-right-boundary-auxiliary} applies to
the displayed final-zone target-window data.  It supplies the fixed inverse for
the final right-boundary zone with
\[
  A(P)=\mathcal D_{\mathrm{out}}(P).
\]
Together with the assumed parser, endpoint pairs, and local inverses for
$Z_1,\ldots,Z_{L-1}$, this gives the right-boundary auxiliary hypotheses
required in item~(2) of \cref{prop:bc-canonical-halo-tail-supplier}.  Applying
that proposition yields \cref{rem:former-post29-tail}.
\end{proof}

\begin{proposition}[Terminal-source target-window canonical halo repairs imply the residual $B/C$ tail]\label{prop:bc-carrier-target-window-halo-tail-supplier}
For each row
\[
  B_{14},\ldots,B_{61},\qquad C_{20},\ldots,C_{61},
\]
put $q=b+1$ for $G=B_b$ or $G=C_b$.  Suppose that for every such row, every
$j$ in the residual window of \cref{rem:former-post29-tail}, every
$q$-element subset $S\subset\{1,\ldots,j\}$, the canonical halo zones of
\cref{lem:bc-canonical-halo-replacement-zones} satisfy the following repair
requirements.
\begin{enumerate}
\item For the width fixed-witness fiber when $G=B_b$, they admit the parser,
target path, replacement windows, branch words, and local inverses listed in
\cref{prop:bc-canonical-halo-local-repair-criterion}.
\item For the height fixed-witness fiber when $G=C_b$, either the same local
repair data exist, or the canonical halo zones can be ordered as
$Z_1,\ldots,Z_L$ with $Z_L=\{a_L,\ldots,j\}$ so that the hypotheses of
\cref{prop:bc-right-boundary-auxiliary-global-recovery} hold for the parser,
target path, endpoint pairs, zones strictly left of $Z_L$, and their local
inverses.  For the final zone, write
\[
  Z_L=\{j-2m+1,\ldots,j\}.
\]
Assume it is in the non-full tight-height setting of
\cref{cor:bc-height-zero-surplus-label-prefix}.  If $m\le28$, assume the
target-window data
\[
  S,\quad Z_L,\quad V^\ast_L,\quad \mathcal D_{\mathrm{out}}(P),
  \quad \rho_{2a_L-2},\rho_{2j}
\]
determine $\lambda^{(2m)}$, the retained even-label strips
$E_1,\ldots,E_m$, and the labelled continuation boxes.  If $m\ge29$ and
$\lambda^{(2m)}$ is column-even, assume the final-zone replacement window and
branch word are chosen by the decoder of
\cref{prop:bc-column-even-prefix-supplies-large-right-boundary-auxiliary}.  If
$m\ge29$ and $\lambda^{(2m)}$ is not column-even, assume that the
target-window data
\[
  S,\quad Z_L,\quad V^\ast_L,\quad \mathcal D_{\mathrm{out}}(P),
  \quad \rho_{2a_L-2},\rho_{2j}
\]
determine \(\lambda^{(2m)}\), the retained even-label strips
\(E_1,\ldots,E_m\), and the labelled continuation boxes.  Let
\[
  \Omega_{\mathrm{tw}}
  =
  \Omega(\lambda^{(2m)};E_1,\ldots,E_m)
\]
with the lexicographic order of \cref{def:bc-tight-odd-packet-omega}; when
\(\Omega_{\mathrm{tw}}\ne\varnothing\), let
\((\gamma;e_1,\ldots,e_{m-1})\) be the associated fixed-even datum of
\cref{def:bc-odd-fixed-even-standard-tableaux}.  In the \(m\ge29\) cases not
already covered by the condition
\[
  \Omega_{\mathrm{tw}}=\varnothing
  \quad\text{or}\quad
  \bigl(\Omega_{\mathrm{tw}}\ne\varnothing
  \quad\text{and}\quad
  \bigl(m=29
  \quad\text{or}\quad
  m=30
  \quad\text{or}\quad
  m=31
  \quad\text{or}\quad
  (\gamma,e_{m-1})\notin\mathcal P^{\mathrm{env}}_m
  \quad\text{or}\quad
  (\gamma,e_{m-1},e_{m-2})\notin
  \mathcal P^{\mathrm{env},2}_m
  \quad\text{or}\quad
  \exists d,\ 3\le d\le m-1,\
  (\gamma,e_{m-1},e_{m-2},\ldots,e_{m-d})
  \notin\mathcal P^{\mathrm{env},d}_m\bigr)\bigr),
\]
form from the target-window data the degree-constrained ambient set and
residual profile
\[
  r_1=\#\{1\le a\le2m:d_a=1\},\qquad
  r_2=\#\{1\le a\le2m:d_a=2\}
\]
of \cref{def:bc-degree-low-low-candidates}.  Also put
\[
  r_0=\#\{1\le a\le2m:d_a=0\},\qquad
  D=\sum_{a=1}^{2m}d_a,\qquad
  C=\sum_{a=1}^{2m}\sum_{t=2m+1}^j u_{\{a,t\}},
\]
and let \(A(m)\) and \(B(m)\) be the frontier functions of
\cref{prop:bc-active-twenty-low-low-profile-full-budget}.  If
\[
  \mathsf N(r_1,r_2)>2^{2m},\qquad
  r_0\le2m-A(m),\qquad
  C\le4m-B(m),
\]
assume in addition that one of the following six source hypotheses holds.
First, the same target-window data determine a set
\[
  U_{\mathrm{tw}}\subseteq\{1\le a\le2m:d_a>0\}
\]
and determine the true Step~3 entries on every low-low cell incident to
\(U_{\mathrm{tw}}\), with
\[
  D(U_{\mathrm{tw}})=\sum_{a\in U_{\mathrm{tw}}}d_a,
\]
and with one of the following positive-residual threshold tests:
\[
  |U_{\mathrm{tw}}|>r_1+r_2-A(m)
  \qquad\text{or}\qquad
  D(U_{\mathrm{tw}})>D-B(m).
\]
Second, the same target-window data determine a partial low-low readout
\[
  \mathcal E_{\mathrm{tw}},\qquad
  p_{\{a,b\}}\quad(\{a,b\}\in\mathcal E_{\mathrm{tw}})
\]
as in \cref{prop:bc-partial-low-low-profile-full-budget}, agreeing with the
true entries on \(\mathcal E_{\mathrm{tw}}\), with residual profile satisfying
\[
  \mathsf N(r_1^{\mathcal E},r_2^{\mathcal E})\le2^{2m},
\]
or, with
\[
  P_{\mathcal E}=
  \sum_{\{a,b\}\in\mathcal E_{\mathrm{tw}}}p_{\{a,b\}},
\]
satisfying
\[
  2P_{\mathcal E}>D-B(m).
\]
Third, the same target-window data determine a canonically ordered set
\[
  \Omega_{\mathrm{src}}\subseteq
  \Omega(\lambda^{(2m)};E_1,\ldots,E_m)
\]
containing the true deleted odd-packet assignment and satisfying
\[
  |\Omega_{\mathrm{src}}|\le2^{2m}.
\]
Fourth, the same target-window data determine a canonically ordered finite set
\[
  \Sigma_{\mathrm{tw}}
\]
of commuting normal-form signature words such that the true fixed-even standard
tableau has signature in \(\Sigma_{\mathrm{tw}}\) and
\[
  |\Sigma_{\mathrm{tw}}|\le2^m .
\]
Fifth, the same target-window data determine a canonically ordered finite set
\[
  \mathcal L_{\mathrm{tw}}
\]
of possible Step~3 entry assignments on $L_{2m}$, the true assignment belongs
to this set, and
\[
  |\mathcal L_{\mathrm{tw}}|\le2^{2m}.
\]
Sixth, the same target-window data determine one of the following
signature-route readouts:
the commuting normal-form signature \(\kappa\) of the true fixed-even standard
tableau; or the even standard-time endpoint chain
\[
  \sigma^{(0)},\sigma^{(2)},\sigma^{(4)},\ldots,\sigma^{(4m)}
\]
of the doubled standardization; or, for each semistandard label
\(1\le a\le2m\), the two boxes carrying label \(a\).
\end{enumerate}
Then \cref{rem:former-post29-tail} holds.
\end{proposition}

\begin{proof}
As in the proof of
\cref{prop:bc-split-target-window-halo-tail-supplier}, it is enough to verify
the right-boundary alternative in item~(2) of
\cref{prop:bc-canonical-halo-tail-supplier}; all other cases are exactly the
local-repair hypotheses of that proposition.

In the right-boundary alternative with $m\le28$,
\cref{prop:bc-split-target-window-supplies-right-boundary-auxiliary} applies to
the displayed base target-window data.  No commuting-signature hypothesis is
needed in this range, by the small-compatibility and full-branch broad-$\Omega$
cases built into that proposition.

In the right-boundary alternative with $m\ge29$ and column-even prefix endpoint,
\cref{prop:bc-column-even-prefix-supplies-large-right-boundary-auxiliary}
applies to the displayed final-zone replacement data.  In the remaining
right-boundary alternative, where $m\ge29$ and the prefix endpoint is not
column-even, the displayed target-window data determine the unrefined
compatibility set.  If this set is empty, there is no current fiber.  If
\(m=29\), then
\cref{prop:bc-twenty-nine-top-three-envelope-omega-supplies-large-right-boundary-auxiliary}
supplies the fixed right-boundary inverse.  If \(m=30\), then
\cref{prop:bc-thirty-fixed-depth-envelope-omega-supplies-large-right-boundary-auxiliary}
supplies the fixed right-boundary inverse.  If \(m=31\), then
\cref{prop:bc-thirty-one-fixed-depth-envelope-omega-supplies-large-right-boundary-auxiliary}
supplies the fixed right-boundary inverse.  If \(m\ge32\) and the
associated fixed-even top pair
\((\gamma,e_{m-1})\notin\mathcal P^{\mathrm{env}}_m\), then
\cref{prop:bc-envelope-nonfrontier-omega-supplies-large-right-boundary-auxiliary}
supplies the fixed right-boundary inverse.  Hence it remains only to handle
the top-fixed envelope-frontier cases.  If \(m\ge32\) and
\((\gamma,e_{m-1},e_{m-2})\notin\mathcal P^{\mathrm{env},2}_m\), then
\cref{prop:bc-envelope-top-two-nonfrontier-omega-supplies-large-right-boundary-auxiliary}
supplies the fixed right-boundary inverse.  Hence it remains only to handle
the top-two fixed envelope-frontier cases.  If \(m\ge32\) and there is a
\(d\) with \(3\le d\le m-1\) such that
\[
  (\gamma,e_{m-1},e_{m-2},\ldots,e_{m-d})
  \notin\mathcal P^{\mathrm{env},d}_m,
\]
then
\cref{prop:bc-fixed-depth-nonfrontier-omega-supplies-large-right-boundary-auxiliary}
supplies the fixed right-boundary inverse.  Hence it remains only to handle
the cases that stay in every fixed-depth envelope frontier:
\[
  m\ge32,\qquad
  (\gamma,e_{m-1},e_{m-2})\in\mathcal P^{\mathrm{env},2}_m,
  \qquad
  (\gamma,e_{m-1},e_{m-2},\ldots,e_{m-d})
  \in\mathcal P^{\mathrm{env},d}_m
  \quad(3\le d\le m-1).
\]
Taking \(d=m-1\) and applying
\cref{cor:bc-full-depth-envelope-exact-omega}, this remaining branch has
\[
  |\Omega_{\mathrm{tw}}|>2^{2m}.
\]
By \cref{cor:bc-all-depth-overflow-many-signatures}, it also realizes more
than \(2^m\) commuting normal-form signatures.  Thus the unrefined broad
odd-packet class itself is genuinely over budget; the following source
alternatives must refine it, find a bounded signature list, or replace it by an
equivalent low-low candidate source.
In those cases form the degree-constrained ambient set and residual degree-one
and degree-two counts \(r_1,r_2\), together with \(r_0\) and \(C\), from
\cref{def:bc-degree-low-low-candidates}.  If
\[
  \mathsf N(r_1,r_2)\le2^{2m},
\]
\cref{prop:bc-exact-degree-profile-low-low-full-budget} supplies the fixed
right-boundary inverse.  If instead \(\mathsf N(r_1,r_2)>2^{2m}\) but either
\[
  r_0>2m-A(m)
  \qquad\text{or}\qquad
  C>4m-B(m),
\]
then
\cref{cor:bc-continuation-heavy-low-low-profile-full-budget} supplies the same
fixed right-boundary inverse.  If the exact profile and continuation-frontier
closures both fail, then we are in the displayed live branch.  By
\cref{cor:bc-incidence-cover-positive-projection,cor:bc-positive-incidence-two-thresholds},
this two-threshold positive-residual form is equivalent to the earlier
four-test incidence-cover form.  Under the incidence-cover source hypothesis,
\cref{cor:bc-incidence-covered-low-labels-frontier-full-budget} supplies the
fixed right-boundary inverse in the positive-label-count threshold subcase, while
\cref{cor:bc-large-incidence-cover-frontier-full-budget} supplies it in the
positive residual-degree threshold subcase.  Under the partial residual-profile
source hypothesis, \cref{prop:bc-partial-low-low-profile-full-budget} supplies
the fixed right-boundary inverse in the exact residual-profile subcase, while
\cref{cor:bc-partial-low-low-edge-mass-frontier-full-budget} supplies it in
the read-edge-mass subcase.  Under the full odd-packet-class source
hypothesis,
\cref{prop:bc-full-budget-omega-supplies-large-right-boundary-auxiliary}
applies to \(\Omega_{\mathrm{src}}\); equivalently,
\cref{cor:bc-full-budget-omega-induces-low-low-candidates} converts the same
source into a full-budget low-low candidate list.  Under the bounded
signature-list source
hypothesis,
\cref{prop:bc-bounded-signature-list-supplies-large-right-boundary-auxiliary}
applies to the displayed signature list, and
\cref{cor:bc-bounded-signature-list-low-low-candidates} gives the same
candidate-list form.  Under the full candidate-list source
hypothesis,
\cref{prop:bc-full-budget-low-low-candidates-supply-large-right-boundary-auxiliary}
applies to the displayed candidate-list hypothesis.  Under the signature
source hypothesis, there are three subcases.  If the source is \(\kappa\),
\cref{prop:bc-target-window-signature-supplies-right-boundary-auxiliary}
applies to the already recovered base target-window datum.  If the source is
the even standard-time endpoint chain, then
\cref{prop:bc-even-chain-supplies-large-right-boundary-auxiliary} applies,
using the labelled continuation boxes already recovered in the base
target-window datum.  If the source is the low-label box readout, then
\cref{prop:bc-low-label-boxes-supply-large-right-boundary-auxiliary} applies,
again using the recovered labelled continuation boxes.  In all three subcases
we obtain the fixed right-boundary inverse with
\(A(P)=\mathcal D_{\mathrm{out}}(P)\).  The earlier carrier and
Hall formulations are not additional alternatives in the present statement.
The direct low-low carrier route is reduced by
\cref{cor:bc-shared-low-low-carrier-terminal-source}; in particular,
\cref{cor:bc-low-low-carrier-forces-terminal-boundary} turns any parsed carrier
containing $L_{2m}$ into a terminal-boundary source after the dual-RSK' backward
sweep.  The internal-$Q$-cut and direct Hall-data routes are reduced by
\cref{cor:bc-shared-q-cut-hall-carrier-terminal-source,cor:bc-shared-direct-hall-carrier-terminal-source},
and the continuation-segment Hall route feeds the same hypotheses through
\cref{prop:bc-target-window-continuation-segment-hall-carrier-supplies-large-right-boundary-auxiliary,cor:bc-continuation-segment-hall-carrier-recovers-large-deleted-only-graph}.
The incidence-cover carrier relaxation is the partial-carrier version of the
same reduction: \cref{prop:bc-strict-left-incidence-cover-carrier-supplies-large-right-boundary-auxiliary}
and
\cref{cor:bc-strict-left-incidence-cover-carrier-recovers-large-deleted-only-graph}
feed the positive-residual two-threshold incidence source into the full-budget
low-low candidate-list source, after the general four-test form has been
projected by
\cref{cor:bc-incidence-cover-positive-projection,cor:bc-positive-incidence-two-thresholds},
and
\cref{cor:bc-shared-direct-hall-incidence-cover-carrier-terminal-source}
records the direct Hall-data shared form.
A full-budget terminal-boundary candidate list is not a separate terminal
artifact: by
\cref{cor:bc-full-budget-terminal-boundaries-induce-low-low-candidates} it
induces a full-budget low-low candidate list.  The graph-side counterparts are
\cref{cor:bc-full-budget-terminal-boundary-list-recovers-large-deleted-only-graph,cor:bc-full-budget-low-low-candidate-list-recovers-large-deleted-only-graph}.
Full-budget odd-packet classes feed both sides through
\cref{prop:bc-full-budget-omega-supplies-large-right-boundary-auxiliary,cor:bc-full-budget-omega-recovers-large-deleted-only-graph},
and they feed the low-low candidate-list source through
\cref{cor:bc-full-budget-omega-induces-low-low-candidates}.
Bounded signature lists feed both sides through
\cref{prop:bc-bounded-signature-list-supplies-large-right-boundary-auxiliary,cor:bc-bounded-signature-list-recovers-large-deleted-only-graph}
and induce the same full-budget low-low candidate-list source by
\cref{cor:bc-bounded-signature-list-low-low-candidates}.
Finally, \cref{cor:bc-target-window-signature-fiber-low-low-candidates} turns a
recovered signature fiber into a bounded low-low candidate class, while
\cref{prop:bc-even-chain-supplies-large-right-boundary-auxiliary,prop:bc-low-label-boxes-supply-large-right-boundary-auxiliary}
record concrete parser-facing ways to recover the same signature-route
inverse.  The matching graph-side statement is
\cref{cor:bc-signature-route-recovers-large-deleted-only-graph}.  Thus the
live source obligation is an incidence-cover readout below
the certified positive-label-count or residual-degree frontier, a
partial low-low readout below the
residual profile frontier, a full-budget odd-packet class, a bounded signature
list, bounded low-low candidate recovery, or one of these signature-route
readouts, in the exact
degree-profile dense branch where \(\mathsf N(r_1,r_2)>2^{2m}\), inside every
fixed-depth recursive-envelope frontier with \(m\ge32\); by
\cref{cor:bc-partial-low-low-edge-mass-frontier-full-budget,cor:bc-large-incidence-cover-frontier-full-budget},
the first two alternatives include the covered-degree and read-edge-mass tests
stated there, with covered-label tests normalized by
\cref{cor:bc-incidence-cover-positive-projection,cor:bc-positive-incidence-two-thresholds}.  By
\cref{prop:bc-active-twenty-low-low-profile-full-budget}, this branch has at
least \(A(m)\) positive residual labels and total residual degree at least
\(B(m)\).  By
\cref{cor:bc-continuation-heavy-low-low-profile-full-budget}, the same branch
has at most \(2m-A(m)\) fully continued low labels and total
low-continuation incidence at most \(4m-B(m)\).
Exact terminal-boundary recovery is a singleton terminal-boundary list, hence
a special case of this low-low candidate-list input.  We obtain the fixed
inverse for the final right-boundary zone with
\[
  A(P)=\mathcal D_{\mathrm{out}}(P).
\]
Together with the assumed parser, endpoint pairs, and local inverses for
$Z_1,\ldots,Z_{L-1}$, this gives the right-boundary auxiliary hypotheses
required in item~(2) of \cref{prop:bc-canonical-halo-tail-supplier}.  Applying
that proposition yields \cref{rem:former-post29-tail}.
\end{proof}

\begin{proposition}[Below-wall packet-splice carrier repairs imply the residual $B/C$ tail]\label{prop:bc-below-wall-carrier-target-window-halo-tail-supplier}
For each row
\[
  B_{14},\ldots,B_{61},\qquad C_{20},\ldots,C_{61},
\]
put $q=b+1$ for $G=B_b$ or $G=C_b$.  Suppose that for every such row, every
$j$ in the residual window of \cref{rem:former-post29-tail}, every
$q$-element subset $S\subset\{1,\ldots,j\}$, the canonical halo zones of
\cref{lem:bc-canonical-halo-replacement-zones} satisfy the following repair
requirements.
\begin{enumerate}
\item For the width fixed-witness fiber when $G=B_b$, they satisfy the
hypotheses of
\cref{prop:bc-below-wall-packet-splice-halo-local-repair-criterion}.
\item For the height fixed-witness fiber when $G=C_b$, either they satisfy the
hypotheses of
\cref{prop:bc-below-wall-packet-splice-halo-local-repair-criterion}, or they
satisfy the right-boundary alternative listed in item~(2) of
\cref{prop:bc-carrier-target-window-halo-tail-supplier}.
\end{enumerate}
Then \cref{rem:former-post29-tail} holds.
\end{proposition}

\begin{proof}
By
\cref{cor:bc-below-wall-packet-splice-supplies-canonical-halo-repair}, the
below-wall packet-splice hypotheses in item~(1) supply the local-repair
hypotheses required in item~(1) of
\cref{prop:bc-carrier-target-window-halo-tail-supplier}.  The same corollary
supplies the first height alternative in item~(2) of that proposition.  The
second height alternative is exactly assumed in item~(2) above.  Therefore the
hypotheses of \cref{prop:bc-carrier-target-window-halo-tail-supplier} hold,
and that proposition gives \cref{rem:former-post29-tail}.
\end{proof}

\begin{proposition}[Step--3 carrier packet readouts plus the terminal target imply the residual $B/C$ tail]\label{prop:bc-step3-carrier-terminal-halo-tail-supplier}
For each row
\[
  B_{14},\ldots,B_{61},\qquad C_{20},\ldots,C_{61},
\]
put $q=b+1$ for $G=B_b$ or $G=C_b$.  Suppose that for every such row, every
$j$ in the residual window of \cref{rem:former-post29-tail}, and every
$q$-element subset $S\subset\{1,\ldots,j\}$, the following repair data are
available.
\begin{enumerate}
\item For the width fixed-witness fiber when $G=B_b$, the canonical halo zones
satisfy the hypotheses of
\cref{prop:bc-step3-carrier-halo-local-repair-criterion}.
\item For the height fixed-witness fiber when $G=C_b$, order the canonical halo
zones increasingly.  If no height halo zone has zero surplus, then the canonical
halo zones satisfy the hypotheses of
\cref{prop:bc-step3-carrier-halo-local-repair-criterion}.
\item If a height zero-surplus zone exists and is the full alternating zone
$j=2q$ with $S=\{1,3,\ldots,2q-1\}$, use the fixed full-zone decoder of
\cref{cor:bc-height-alternating-full-zone-decoder}.
\item If a height zero-surplus zone exists and is not full, then, with this
terminal zone written $Z_L=\{a_L,\ldots,j\}$, the global parser, target path,
replacement windows, branch words, endpoint pairs, and unchanged complement
blocks satisfy the right-boundary auxiliary hypotheses of
\cref{prop:bc-right-boundary-auxiliary-global-recovery} except for the local
inverse assertions on the zones strictly left of $Z_L$.  Each zone strictly left
of $Z_L$ satisfies the Step--3 carrier packet-readout hypotheses of
\cref{prop:bc-step3-carrier-packet-readouts-give-halo-zone-inverses}, using the
parser, window, branch word, and endpoint data supplied by this right-boundary
package.  The final zone satisfies the final-zone target-window/right-boundary
alternative listed in item~(2) of
\cref{prop:bc-carrier-target-window-halo-tail-supplier}.
\end{enumerate}
Then \cref{rem:former-post29-tail} holds.
\end{proposition}

\begin{proof}
For a $B_b$ row, item~(1) and
\cref{prop:bc-step3-carrier-halo-local-repair-criterion} give
\[
 \#\{T:\lambda(T)' \text{ even},\ \lambda_1(T)\ge2q,\ W_q(T)=S\}
 \le 2^{2q}s_{j-q}.
\]

For a $C_b$ row with no zero-surplus height halo zone, item~(2) gives the
corresponding height fixed-fiber estimate by the same carrier local-repair
criterion.  If the zero-surplus zone is full, item~(3) and
\cref{cor:bc-height-alternating-full-zone-decoder} give the height estimate, as
in \cref{prop:bc-terminal-height-split-halo-tail-supplier}.

It remains to treat the non-full zero-surplus case.  By
\cref{cor:bc-height-zero-surplus-terminal-unique}, the zero-surplus zone is
terminal and every zone strictly to its left has positive surplus.  For each
left zone, item~(4) and
\cref{prop:bc-step3-carrier-packet-readouts-give-halo-zone-inverses} supply the
missing local inverse in the right-boundary auxiliary package.  The final-zone
target-window/right-boundary alternative supplies the fixed inverse for $Z_L$.
Therefore \cref{prop:bc-right-boundary-auxiliary-fiber-criterion} gives
\[
 \#\{T:\lambda(T)' \text{ even},\ \ell(\lambda(T))\ge2q,\ H_q(T)=S\}
 \le 2^{2q}s_{j-q}
\]
also in this case.

Thus the required width and height fixed-fiber estimates hold for every
residual row, admissible $j$, and witness set $S$.  Applying
\cref{prop:bc-bad-shape-fiber-reduction} and then
\cref{prop:bc-bad-shape-reduction} gives \cref{rem:former-post29-tail}.
\end{proof}

\begin{proposition}[Terminal height-zone split implies the residual $B/C$ tail]\label{prop:bc-terminal-height-split-halo-tail-supplier}
For each row
\[
  B_{14},\ldots,B_{61},\qquad C_{20},\ldots,C_{61},
\]
put $q=b+1$ for $G=B_b$ or $G=C_b$.  Suppose that for every such row, every
$j$ in the residual window of \cref{rem:former-post29-tail}, and every
$q$-element subset $S\subset\{1,\ldots,j\}$, the following repair data are
available.
\begin{enumerate}
\item For the width fixed-witness fiber when $G=B_b$, the canonical halo
zones satisfy the hypotheses of
\cref{prop:bc-below-wall-packet-splice-halo-local-repair-criterion}.
\item For the height fixed-witness fiber when $G=C_b$, order the canonical
halo zones increasingly.  If no height halo zone has zero surplus, then the
canonical halo zones satisfy the hypotheses of
\cref{prop:bc-below-wall-packet-splice-halo-local-repair-criterion}.
\item If a height zero-surplus zone exists and is the full alternating zone
$j=2q$ with $S=\{1,3,\ldots,2q-1\}$, use the fixed full-zone decoder of
\cref{cor:bc-height-alternating-full-zone-decoder}.
\item If a height zero-surplus zone exists and is not full, then, with this
terminal zone written $Z_L=\{a_L,\ldots,j\}$, the zones strictly left of
$Z_L$ are supplied by below-wall packet-splice local inverses, and the final
zone satisfies the final-zone target-window/right-boundary alternative listed
in item~(2) of
\cref{prop:bc-carrier-target-window-halo-tail-supplier}.
\end{enumerate}
Then \cref{rem:former-post29-tail} holds.
\end{proposition}

\begin{proof}
Fix a residual row, an admissible $j$, and a witness set $S$.

For a $B_b$ row, item~(1) and
\cref{prop:bc-below-wall-packet-splice-halo-local-repair-criterion} give the
width fixed-fiber estimate
\[
 \#\{T:\lambda(T)' \text{ even},\ \lambda_1(T)\ge2q,\ W_q(T)=S\}
 \le 2^{2q}s_{j-q}.
\]

Now consider a $C_b$ row.  By
\cref{cor:bc-height-zero-surplus-terminal-unique}, there is at most one
zero-surplus height halo zone.  If no such zone exists, item~(2) and
\cref{prop:bc-below-wall-packet-splice-halo-local-repair-criterion} give the
height fixed-fiber estimate.

If the zero-surplus zone is full, then
\cref{cor:bc-height-zero-surplus-label-prefix} writes it as
$Z_\gamma=\{j-2m+1,\ldots,j\}$.  Since the full-zone hypothesis says
$Z_\gamma=\{1,\ldots,j\}$, one has $j=2m$.  The full zone contains every
deleted block, so $m=|I(S)|=q$.  The same corollary then gives
\[
  j=2q,\qquad S=\{1,3,\ldots,2q-1\}.
\]
Since residual rows have $q\ge15$, \cref{cor:bc-height-alternating-full-zone-decoder}
applies and gives the same height fixed-fiber estimate.

It remains to treat the non-full zero-surplus case.  By
\cref{cor:bc-height-zero-surplus-terminal-unique}, the zero-surplus zone is
the terminal right-boundary zone and every earlier zone has positive surplus.
Item~(4), together with
\cref{cor:bc-below-wall-packet-splice-supplies-canonical-halo-repair}, gives
the right-boundary auxiliary hypotheses used in the proof of
\cref{prop:bc-carrier-target-window-halo-tail-supplier}; equivalently,
\cref{prop:bc-right-boundary-auxiliary-fiber-criterion} applies to the height
fixed-witness fiber.  Hence
\[
 \#\{T:\lambda(T)' \text{ even},\ \ell(\lambda(T))\ge2q,\ H_q(T)=S\}
 \le 2^{2q}s_{j-q}
\]
also holds in this case.

Thus the required width and height fixed-fiber estimates hold for every
residual row, admissible $j$, and witness set $S$.  Applying
\cref{prop:bc-bad-shape-fiber-reduction} and then
\cref{prop:bc-bad-shape-reduction} gives \cref{rem:former-post29-tail}.
\end{proof}

\begin{proposition}[Anchor-pair positive height zones imply the residual $B/C$ tail]\label{prop:bc-anchor-pair-positive-height-tail-supplier}
For each row
\[
  B_{14},\ldots,B_{61},\qquad C_{20},\ldots,C_{61},
\]
put $q=b+1$ for $G=B_b$ or $G=C_b$.  Suppose that for every such row, every
$j$ in the residual window of \cref{rem:former-post29-tail}, and every
$q$-element subset $S\subset\{1,\ldots,j\}$, the following repair data are
available.
\begin{enumerate}
\item For the width fixed-witness fiber when $G=B_b$, the canonical halo
zones satisfy the hypotheses of
\cref{prop:bc-below-wall-packet-splice-halo-local-repair-criterion}.
\item For the height fixed-witness fiber when $G=C_b$, order the canonical
halo zones increasingly.  If no height halo zone has zero surplus, then the
global parser, target path, replacement windows, branch words, and endpoint
pairs satisfy the hypotheses of
\cref{prop:bc-canonical-halo-local-repair-criterion} except for the local
inverse assertion, and every zone satisfies the anchor-pair hypotheses of
\cref{prop:bc-height-anchor-pair-zone-recovery}.
\item If a height zero-surplus zone exists and is the full alternating zone
$j=2q$ with $S=\{1,3,\ldots,2q-1\}$, use the fixed full-zone decoder of
\cref{cor:bc-height-alternating-full-zone-decoder}.
\item If a height zero-surplus zone exists and is not full, then, with this
terminal zone written $Z_L=\{a_L,\ldots,j\}$, the global parser, target path,
replacement windows, branch words, endpoint pairs, and unchanged complement
blocks satisfy the right-boundary auxiliary hypotheses of
\cref{prop:bc-right-boundary-auxiliary-global-recovery} except for the local
inverse assertions on the zones strictly left of $Z_L$.  Each zone strictly
left of $Z_L$ satisfies the anchor-pair hypotheses of
\cref{prop:bc-height-anchor-pair-zone-recovery}, and the final zone satisfies
the final-zone target-window/right-boundary alternative listed in item~(2) of
\cref{prop:bc-carrier-target-window-halo-tail-supplier}.
\end{enumerate}
Then \cref{rem:former-post29-tail} holds.
\end{proposition}

\begin{proof}
Fix a residual row, an admissible $j$, and a witness set $S$.

The $B_b$ width estimate is exactly as in
\cref{prop:bc-terminal-height-split-halo-tail-supplier}: item~(1) and
\cref{prop:bc-below-wall-packet-splice-halo-local-repair-criterion} give
\[
 \#\{T:\lambda(T)' \text{ even},\ \lambda_1(T)\ge2q,\ W_q(T)=S\}
 \le 2^{2q}s_{j-q}.
\]

Now let $G=C_b$.  By
\cref{cor:bc-height-zero-surplus-terminal-unique}, there is at most one
zero-surplus height halo zone.

If no zero-surplus height halo zone exists, then every height halo zone has
positive surplus.  For each zone,
\cref{prop:bc-height-anchor-pair-zone-recovery} supplies the missing local
inverse in
\cref{prop:bc-canonical-halo-local-repair-criterion}.  Therefore that
canonical local-repair criterion applies to the height fixed-witness fiber and
gives
\[
 \#\{T:\lambda(T)' \text{ even},\ \ell(\lambda(T))\ge2q,\ H_q(T)=S\}
 \le 2^{2q}s_{j-q}.
\]

If the zero-surplus zone is full, then the same argument as in
\cref{prop:bc-terminal-height-split-halo-tail-supplier} applies
\cref{cor:bc-height-alternating-full-zone-decoder} and gives the same height
fixed-fiber estimate.

It remains to treat the non-full zero-surplus case.  By
\cref{cor:bc-height-zero-surplus-terminal-unique}, the zero-surplus zone is
terminal and every zone strictly to its left has positive surplus.  Applying
\cref{prop:bc-height-anchor-pair-zone-recovery} to each such left zone
supplies the local inverses required in
\cref{prop:bc-right-boundary-auxiliary-global-recovery}.  The final-zone
target-window/right-boundary alternative supplies the fixed inverse for
$Z_L$ with the auxiliary datum specified there.  Hence the hypotheses of
\cref{prop:bc-right-boundary-auxiliary-fiber-criterion} hold for the height
fixed-witness fiber, and this again gives
\[
 \#\{T:\lambda(T)' \text{ even},\ \ell(\lambda(T))\ge2q,\ H_q(T)=S\}
 \le 2^{2q}s_{j-q}.
\]

Thus the required width and height fixed-fiber estimates hold for every
residual row, admissible $j$, and witness set $S$.  Applying
\cref{prop:bc-bad-shape-fiber-reduction} and then
\cref{prop:bc-bad-shape-reduction} gives \cref{rem:former-post29-tail}.
\end{proof}

\begin{proposition}[Two-label carrier positive height zones imply the residual $B/C$ tail]\label{prop:bc-two-label-carrier-positive-height-tail-supplier}
For each row
\[
  B_{14},\ldots,B_{61},\qquad C_{20},\ldots,C_{61},
\]
put $q=b+1$ for $G=B_b$ or $G=C_b$.  Fix once and for all an injection
\[
  \kappa:\{0,1,2\}\hookrightarrow\{0,1\}^2 .
\]
Suppose that for every such row, every $j$ in the residual window of
\cref{rem:former-post29-tail}, and every
$q$-element subset $S\subset\{1,\ldots,j\}$, the following repair data are
available.
\begin{enumerate}
\item For the width fixed-witness fiber when $G=B_b$, the canonical halo
zones satisfy the hypotheses of
\cref{prop:bc-below-wall-packet-splice-halo-local-repair-criterion}.
\item For the height fixed-witness fiber when $G=C_b$, order the canonical
halo zones increasingly.  If no height halo zone has zero surplus, then the
global parser, target path, replacement windows, branch words, and endpoint
pairs satisfy the hypotheses of
\cref{prop:bc-canonical-halo-local-repair-criterion} except for the local
inverse assertion.  For every such height zone $Z_\gamma$, the parser
decomposes $V^\ast_\gamma$ into the unchanged surplus-anchor blocks and one
merged window $M_t$ for each deleted singleton $u_t$, as in
\cref{prop:bc-height-anchor-pair-zone-recovery}.  Moreover, with
$s_t=j-u_t+1$, with $E_{s_t,s_t+1}$ the two-label Step--3 carrier, and with
$e_t$ its unique cell entry, the data
\[
  M_t,\quad S,\quad Z_\gamma,\quad t,\quad
  \rho_{2a_\gamma-2},\rho_{2b_\gamma}
\]
determine the left/bottom boundary of $E_{s_t,s_t+1}$ and the corresponding
two-bit word satisfies $\beta_t=\kappa(e_t)$.
\item If a height zero-surplus zone exists and is the full alternating zone
$j=2q$ with $S=\{1,3,\ldots,2q-1\}$, use the fixed full-zone decoder of
\cref{cor:bc-height-alternating-full-zone-decoder}.
\item If a height zero-surplus zone exists and is not full, then, with this
terminal zone written $Z_L=\{a_L,\ldots,j\}$, the global parser, target path,
replacement windows, branch words, endpoint pairs, and unchanged complement
blocks satisfy the right-boundary auxiliary hypotheses of
\cref{prop:bc-right-boundary-auxiliary-global-recovery} except for the local
inverse assertions on the zones strictly left of $Z_L$.  Each zone strictly
left of $Z_L$ satisfies the two-label carrier hypotheses stated in item~(2),
and the final zone satisfies the final-zone target-window/right-boundary
alternative listed in item~(2) of
\cref{prop:bc-carrier-target-window-halo-tail-supplier}.
\end{enumerate}
Then \cref{rem:former-post29-tail} holds.
\end{proposition}

\begin{proof}
The $B_b$ width estimate is supplied by item~(1), exactly as in
\cref{prop:bc-anchor-pair-positive-height-tail-supplier}.

Now consider a $C_b$ row.  If no zero-surplus height halo zone exists, then
every height halo zone has positive surplus by
\cref{cor:bc-height-halo-singleton-nonnegative}.  For each such zone,
item~(2) gives the parser decomposition, left/bottom two-label carrier
boundary, and cell-entry branch encoding required by
\cref{prop:bc-two-label-left-bottom-carrier-zone-recovery}.  Hence that
proposition supplies the missing local inverse for every zone in
\cref{prop:bc-canonical-halo-local-repair-criterion}, and the height
fixed-fiber estimate follows from that criterion.

If the zero-surplus zone is full, item~(3) and
\cref{cor:bc-height-alternating-full-zone-decoder} give the height
fixed-fiber estimate as in
\cref{prop:bc-anchor-pair-positive-height-tail-supplier}.

It remains to treat the non-full zero-surplus case.  By
\cref{cor:bc-height-zero-surplus-terminal-unique}, the zero-surplus zone is
terminal and every zone strictly to its left has positive surplus.  Item~(4)
and \cref{prop:bc-two-label-left-bottom-carrier-zone-recovery} supply the
local inverses for those left zones, while the final-zone
target-window/right-boundary alternative supplies the fixed inverse for
$Z_L$.  Thus the hypotheses of
\cref{prop:bc-right-boundary-auxiliary-fiber-criterion} hold for the height
fixed-witness fiber, giving the same height fixed-fiber estimate.

Therefore the width and height fixed-fiber estimates hold for every residual
row, admissible $j$, and witness set $S$.  Applying
\cref{prop:bc-bad-shape-fiber-reduction} and then
\cref{prop:bc-bad-shape-reduction} gives \cref{rem:former-post29-tail}.
\end{proof}

\begin{proposition}[Coupled collar carriers imply the residual $B/C$ tail]\label{prop:bc-coupled-collar-carrier-positive-height-tail-supplier}
For each row
\[
  B_{14},\ldots,B_{61},\qquad C_{20},\ldots,C_{61},
\]
put $q=b+1$ for $G=B_b$ or $G=C_b$.  Fix once and for all an injection
\[
  \kappa:\{0,1,2\}\hookrightarrow\{0,1\}^2 .
\]
Suppose that for every such row, every $j$ in the residual window of
\cref{rem:former-post29-tail}, and every
$q$-element subset $S\subset\{1,\ldots,j\}$, the following repair data are
available.
\begin{enumerate}
\item For the width fixed-witness fiber when $G=B_b$, the canonical halo
zones satisfy the hypotheses of
\cref{prop:bc-below-wall-packet-splice-halo-local-repair-criterion}.
\item For the height fixed-witness fiber when $G=C_b$, order the canonical
halo zones increasingly.  If no height halo zone has zero surplus, then the
global parser, target path, replacement windows, branch words, and endpoint
pairs satisfy the hypotheses of
\cref{prop:bc-canonical-halo-local-repair-criterion} except for the local
inverse assertion.  For every such height zone $Z_\gamma$, the parser
decomposes $V^\ast_\gamma$ into unchanged surplus-anchor blocks and merged
windows $M_1,\ldots,M_m$ for the increasing deleted singleton list
\[
  Z_\gamma\cap I=\{u_1<\cdots<u_m\}.
\]
These merged windows satisfy the coupled collar hypotheses of
\cref{prop:bc-coupled-collar-carriers-recover-positive-height-zones}: each
$M_t$ determines the left-collar block $B_{u_t-1}(P)$, the parsed zone data
determine the left/bottom boundary of $E_{s_t,s_t+1}$ after the
parser-local right collars are recovered, and the two-bit word satisfies
$\beta_t=\kappa(e_t)$ for the unique carrier entry $e_t$.
\item If a height zero-surplus zone exists and is the full alternating zone
$j=2q$ with $S=\{1,3,\ldots,2q-1\}$, use the fixed full-zone decoder of
\cref{cor:bc-height-alternating-full-zone-decoder}.
\item If a height zero-surplus zone exists and is not full, then, with this
terminal zone written $Z_L=\{a_L,\ldots,j\}$, the global parser, target path,
replacement windows, branch words, endpoint pairs, and unchanged complement
blocks satisfy the right-boundary auxiliary hypotheses of
\cref{prop:bc-right-boundary-auxiliary-global-recovery} except for the local
inverse assertions on the zones strictly left of $Z_L$.  Each zone strictly
left of $Z_L$ satisfies the coupled collar hypotheses stated in item~(2), and
the final zone satisfies the final-zone target-window/right-boundary
alternative listed in item~(2) of
\cref{prop:bc-carrier-target-window-halo-tail-supplier}.
\end{enumerate}
Then \cref{rem:former-post29-tail} holds.
\end{proposition}

\begin{proof}
The $B_b$ width estimate is supplied by item~(1), exactly as in
\cref{prop:bc-two-label-carrier-positive-height-tail-supplier}.

Now consider a $C_b$ row.  If no zero-surplus height halo zone exists, then
every height halo zone has positive surplus by
\cref{cor:bc-height-halo-singleton-nonnegative}.  For each such zone,
item~(2) supplies the hypotheses of
\cref{prop:bc-coupled-collar-carriers-recover-positive-height-zones}.  That
proposition gives the missing local inverse for every zone in
\cref{prop:bc-canonical-halo-local-repair-criterion}, and the height
fixed-fiber estimate follows.

If the zero-surplus zone is full, item~(3) and
\cref{cor:bc-height-alternating-full-zone-decoder} give the height
fixed-fiber estimate as in
\cref{prop:bc-two-label-carrier-positive-height-tail-supplier}.

It remains to treat the non-full zero-surplus case.  By
\cref{cor:bc-height-zero-surplus-terminal-unique}, the zero-surplus zone is
terminal and every zone strictly to its left has positive surplus.  Item~(4)
and \cref{prop:bc-coupled-collar-carriers-recover-positive-height-zones}
supply the local inverses for those left zones, while the final-zone
target-window/right-boundary alternative supplies the fixed inverse for
$Z_L$.  Therefore the hypotheses of
\cref{prop:bc-right-boundary-auxiliary-fiber-criterion} hold for the height
fixed-witness fiber, giving the same height fixed-fiber estimate.

The width and height fixed-fiber estimates now hold for every residual row,
admissible $j$, and witness set $S$.  Applying
\cref{prop:bc-bad-shape-fiber-reduction} and then
\cref{prop:bc-bad-shape-reduction} gives \cref{rem:former-post29-tail}.
\end{proof}

\begin{proposition}[Padded carrier valleys imply the residual $B/C$ tail]\label{prop:bc-padded-carrier-valley-positive-height-tail-supplier}
For each row
\[
  B_{14},\ldots,B_{61},\qquad C_{20},\ldots,C_{61},
\]
put $q=b+1$ for $G=B_b$ or $G=C_b$.  Fix once and for all an injection
\[
  \kappa:\{0,1,2\}\hookrightarrow\{0,1\}^2 .
\]
Suppose that for every such row, every $j$ in the residual window of
\cref{rem:former-post29-tail}, and every
$q$-element subset $S\subset\{1,\ldots,j\}$, the following repair data are
available.
\begin{enumerate}
\item For the width fixed-witness fiber when $G=B_b$, the canonical halo
zones satisfy the hypotheses of
\cref{prop:bc-below-wall-packet-splice-halo-local-repair-criterion}.
\item For the height fixed-witness fiber when $G=C_b$, order the canonical
halo zones increasingly.  If no height halo zone has zero surplus, then the
global parser, target path, replacement windows, branch words, and endpoint
pairs satisfy the hypotheses of
\cref{prop:bc-canonical-halo-local-repair-criterion} except for the local
inverse assertion.  For every such height zone $Z_\gamma$, the parser
decomposes $V^\ast_\gamma$ into unchanged surplus-anchor blocks and merged
windows $M_t$ for the increasing deleted singleton list.  Each merged window
is the padded carrier valley of
\cref{prop:bc-padded-carrier-valleys-recover-positive-height-zones}: for
$s_t=j-u_t+1$, if
\[
  \begin{matrix}
  \nu_t & \lambda_t\\
  \rho_t & \mu_t
  \end{matrix}
\]
are the corner labels of $E_{s_t,s_t+1}$ and $e_t$ is its cell entry, then
\[
  M_t=\operatorname{Pad}_{e_t}(\nu_t,\rho_t,\mu_t),
\]
and the two-bit word satisfies $\beta_t=\kappa(e_t)$.
\item If a height zero-surplus zone exists and is the full alternating zone
$j=2q$ with $S=\{1,3,\ldots,2q-1\}$, use the fixed full-zone decoder of
\cref{cor:bc-height-alternating-full-zone-decoder}.
\item If a height zero-surplus zone exists and is not full, then, with this
terminal zone written $Z_L=\{a_L,\ldots,j\}$, the global parser, target path,
replacement windows, branch words, endpoint pairs, and unchanged complement
blocks satisfy the right-boundary auxiliary hypotheses of
\cref{prop:bc-right-boundary-auxiliary-global-recovery} except for the local
inverse assertions on the zones strictly left of $Z_L$.  Each zone strictly
left of $Z_L$ satisfies the padded-carrier hypotheses stated in item~(2), and
the final zone satisfies the final-zone target-window/right-boundary
alternative listed in item~(2) of
\cref{prop:bc-carrier-target-window-halo-tail-supplier}.
\end{enumerate}
Then \cref{rem:former-post29-tail} holds.
\end{proposition}

\begin{proof}
The proof is the same reduction as
\cref{prop:bc-coupled-collar-carrier-positive-height-tail-supplier}, with
\cref{prop:bc-padded-carrier-valleys-recover-positive-height-zones} replacing
the coupled-collar local inverse.  Item~(1) supplies the $B_b$ width estimate.

For a $C_b$ row with no zero-surplus height halo zone, every height halo zone
has positive surplus by \cref{cor:bc-height-halo-singleton-nonnegative}.
Item~(2) and
\cref{prop:bc-padded-carrier-valleys-recover-positive-height-zones} supply the
missing local inverse for every zone in
\cref{prop:bc-canonical-halo-local-repair-criterion}, giving the height
fixed-fiber estimate.

If the zero-surplus zone is full, item~(3) and
\cref{cor:bc-height-alternating-full-zone-decoder} give the height
fixed-fiber estimate.  In the remaining non-full zero-surplus case,
\cref{cor:bc-height-zero-surplus-terminal-unique} makes that zero-surplus zone
terminal and every zone strictly to its left positive.  Item~(4) supplies the
padded-carrier inverses for those left zones and the final-zone
target-window/right-boundary inverse for $Z_L$, so
\cref{prop:bc-right-boundary-auxiliary-fiber-criterion} gives the height
fixed-fiber estimate.

Thus the width and height fixed-fiber estimates hold for every residual row,
admissible $j$, and witness set $S$.  Applying
\cref{prop:bc-bad-shape-fiber-reduction} and then
\cref{prop:bc-bad-shape-reduction} gives \cref{rem:former-post29-tail}.
\end{proof}

\begin{proposition}[Tagged collective padded splices imply the residual $B/C$ tail]\label{prop:bc-tagged-collective-padded-splice-positive-height-tail-supplier}
For each row
\[
  B_{14},\ldots,B_{61},\qquad C_{20},\ldots,C_{61},
\]
put $q=b+1$ for $G=B_b$ or $G=C_b$.  Fix once and for all an injection
\[
  \kappa:\{0,1,2\}\hookrightarrow\{0,1\}^2 .
\]
Suppose that for every such row, every $j$ in the residual window of
\cref{rem:former-post29-tail}, and every
$q$-element subset $S\subset\{1,\ldots,j\}$, the following repair data are
available.
\begin{enumerate}
\item For the width fixed-witness fiber when $G=B_b$, the canonical halo
zones satisfy the hypotheses of
\cref{prop:bc-below-wall-packet-splice-halo-local-repair-criterion}.
\item For the height fixed-witness fiber when $G=C_b$, order the canonical
halo zones increasingly.  If no height halo zone has zero surplus, then the
global parser, target path, replacement windows, branch words, and endpoint
pairs satisfy the hypotheses of
\cref{prop:bc-canonical-halo-local-repair-criterion} except for the local
inverse assertion.  For every such height zone $Z_\gamma$, the parser
decomposes $V^\ast_\gamma$ as a tagged collective padded splice satisfying
\cref{prop:bc-tagged-collective-padded-splice-recovers-positive-height-zones},
with the same fixed injection $\kappa$.
\item If a height zero-surplus zone exists and is the full alternating zone
$j=2q$ with $S=\{1,3,\ldots,2q-1\}$, use the fixed full-zone decoder of
\cref{cor:bc-height-alternating-full-zone-decoder}.
\item If a height zero-surplus zone exists and is not full, then, with this
terminal zone written $Z_L=\{a_L,\ldots,j\}$, the global parser, target path,
replacement windows, branch words, endpoint pairs, and unchanged complement
blocks satisfy the right-boundary auxiliary hypotheses of
\cref{prop:bc-right-boundary-auxiliary-global-recovery} except for the local
inverse assertions on the zones strictly left of $Z_L$.  Each zone strictly
left of $Z_L$ satisfies the tagged collective padded-splice hypothesis stated
in item~(2), and the final zone satisfies the final-zone
target-window/right-boundary alternative listed in item~(2) of
\cref{prop:bc-carrier-target-window-halo-tail-supplier}.
\end{enumerate}
Then \cref{rem:former-post29-tail} holds.
\end{proposition}

\begin{proof}
The reduction is the same as
\cref{prop:bc-padded-carrier-valley-positive-height-tail-supplier}, with
\cref{prop:bc-tagged-collective-padded-splice-recovers-positive-height-zones}
as the positive-height local inverse.  Item~(1) supplies the $B_b$ width
estimate.

For a $C_b$ row with no zero-surplus height halo zone, every height halo zone
has positive surplus by \cref{cor:bc-height-halo-singleton-nonnegative}.
Item~(2) and
\cref{prop:bc-tagged-collective-padded-splice-recovers-positive-height-zones}
supply the missing local inverse for every zone in
\cref{prop:bc-canonical-halo-local-repair-criterion}, giving the height
fixed-fiber estimate.

If the zero-surplus zone is full, item~(3) and
\cref{cor:bc-height-alternating-full-zone-decoder} give the height
fixed-fiber estimate.  In the remaining non-full zero-surplus case,
\cref{cor:bc-height-zero-surplus-terminal-unique} makes that zero-surplus zone
terminal and every zone strictly to its left positive.  Item~(4) supplies the
tagged collective padded-splice inverses for those left zones and the
final-zone target-window/right-boundary inverse for $Z_L$, so
\cref{prop:bc-right-boundary-auxiliary-fiber-criterion} gives the height
fixed-fiber estimate.

Thus the width and height fixed-fiber estimates hold for every residual row,
admissible $j$, and witness set $S$.  Applying
\cref{prop:bc-bad-shape-fiber-reduction} and then
\cref{prop:bc-bad-shape-reduction} gives \cref{rem:former-post29-tail}.
\end{proof}

\begin{proposition}[Grouped collective carrier splices imply the residual $B/C$ tail]\label{prop:bc-grouped-collective-carrier-splice-positive-height-tail-supplier}
For each row
\[
  B_{14},\ldots,B_{61},\qquad C_{20},\ldots,C_{61},
\]
put $q=b+1$ for $G=B_b$ or $G=C_b$.  Fix once and for all an injection
\[
  \kappa:\{0,1,2\}\hookrightarrow\{0,1\}^2 .
\]
Suppose that for every such row, every $j$ in the residual window of
\cref{rem:former-post29-tail}, and every
$q$-element subset $S\subset\{1,\ldots,j\}$, the following repair data are
available.
\begin{enumerate}
\item For the width fixed-witness fiber when $G=B_b$, the canonical halo
zones satisfy the hypotheses of
\cref{prop:bc-below-wall-packet-splice-halo-local-repair-criterion}.
\item For the height fixed-witness fiber when $G=C_b$, order the canonical
halo zones increasingly.  If no height halo zone has zero surplus, then the
global parser, target path, replacement windows, branch words, and endpoint
pairs satisfy the hypotheses of
\cref{prop:bc-canonical-halo-local-repair-criterion} except for the local
inverse assertion.  For every such height zone $Z_\gamma$, the parser
decomposes $V^\ast_\gamma$ as a grouped collective carrier splice satisfying
\cref{prop:bc-grouped-collective-carrier-splices-recover-positive-height-zones},
with the same fixed injection $\kappa$.
\item If a height zero-surplus zone exists and is the full alternating zone
$j=2q$ with $S=\{1,3,\ldots,2q-1\}$, use the fixed full-zone decoder of
\cref{cor:bc-height-alternating-full-zone-decoder}.
\item If a height zero-surplus zone exists and is not full, then, with this
terminal zone written $Z_L=\{a_L,\ldots,j\}$, the global parser, target path,
replacement windows, branch words, endpoint pairs, and unchanged complement
blocks satisfy the right-boundary auxiliary hypotheses of
\cref{prop:bc-right-boundary-auxiliary-global-recovery} except for the local
inverse assertions on the zones strictly left of $Z_L$.  Each zone strictly
left of $Z_L$ satisfies the grouped collective carrier-splice hypothesis
stated in item~(2), and the final zone satisfies the final-zone
target-window/right-boundary alternative listed in item~(2) of
\cref{prop:bc-carrier-target-window-halo-tail-supplier}.
\end{enumerate}
Then \cref{rem:former-post29-tail} holds.
\end{proposition}

\begin{proof}
The reduction is the same as
\cref{prop:bc-tagged-collective-padded-splice-positive-height-tail-supplier},
with
\cref{prop:bc-grouped-collective-carrier-splices-recover-positive-height-zones}
as the positive-height local inverse.  The $B_b$ width estimate is supplied by
item~(1).

For a $C_b$ row with no zero-surplus height halo zone, every height halo zone
has positive surplus by \cref{cor:bc-height-halo-singleton-nonnegative}.
Item~(2) and
\cref{prop:bc-grouped-collective-carrier-splices-recover-positive-height-zones}
supply the missing local inverse for every zone in
\cref{prop:bc-canonical-halo-local-repair-criterion}, giving the height
fixed-fiber estimate.

If the zero-surplus zone is full, item~(3) and
\cref{cor:bc-height-alternating-full-zone-decoder} give the height
fixed-fiber estimate.  In the remaining non-full zero-surplus case,
\cref{cor:bc-height-zero-surplus-terminal-unique} makes that zero-surplus zone
terminal and every zone strictly to its left positive.  Item~(4) supplies the
grouped collective carrier-splice inverses for those left zones and the
final-zone target-window/right-boundary inverse for $Z_L$, so
\cref{prop:bc-right-boundary-auxiliary-fiber-criterion} gives the height
fixed-fiber estimate.

Thus the width and height fixed-fiber estimates hold for every residual row,
admissible $j$, and witness set $S$.  Applying
\cref{prop:bc-bad-shape-fiber-reduction} and then
\cref{prop:bc-bad-shape-reduction} gives \cref{rem:former-post29-tail}.
\end{proof}

\begin{proposition}[Step--3 boundary grouped splices imply the residual $B/C$ tail]\label{prop:bc-step3-boundary-grouped-splice-positive-height-tail-supplier}
For each row
\[
  B_{14},\ldots,B_{61},\qquad C_{20},\ldots,C_{61},
\]
put $q=b+1$ for $G=B_b$ or $G=C_b$.  Fix once and for all an injection
\[
  \kappa:\{0,1,2\}\hookrightarrow\{0,1\}^2 .
\]
Suppose that for every such row, every $j$ in the residual window of
\cref{rem:former-post29-tail}, and every
$q$-element subset $S\subset\{1,\ldots,j\}$, the following repair data are
available.
\begin{enumerate}
\item For the width fixed-witness fiber when $G=B_b$, the canonical halo
zones satisfy the hypotheses of
\cref{prop:bc-below-wall-packet-splice-halo-local-repair-criterion}.
\item For the height fixed-witness fiber when $G=C_b$, order the canonical
halo zones increasingly.  If no height halo zone has zero surplus, then the
global parser, target path, replacement windows, branch words, and endpoint
pairs satisfy the hypotheses of
\cref{prop:bc-canonical-halo-local-repair-criterion} except for the local
inverse assertion.  For every such height zone $Z_\gamma$, the parser
decomposes $V^\ast_\gamma$ into grouped subwindows and parser-visible tag
sets satisfying the hypotheses of
\cref{prop:bc-step3-boundary-readouts-supply-grouped-carrier-decoders}, with
the same fixed injection $\kappa$.
\item If a height zero-surplus zone exists and is the full alternating zone
$j=2q$ with $S=\{1,3,\ldots,2q-1\}$, use the fixed full-zone decoder of
\cref{cor:bc-height-alternating-full-zone-decoder}.
\item If a height zero-surplus zone exists and is not full, then, with this
terminal zone written $Z_L=\{a_L,\ldots,j\}$, the global parser, target path,
replacement windows, branch words, endpoint pairs, and unchanged complement
blocks satisfy the right-boundary auxiliary hypotheses of
\cref{prop:bc-right-boundary-auxiliary-global-recovery} except for the local
inverse assertions on the zones strictly left of $Z_L$.  Each zone strictly
left of $Z_L$ satisfies the Step--3 boundary grouped-splice hypothesis stated
in item~(2), and the final zone satisfies the final-zone
target-window/right-boundary alternative listed in item~(2) of
\cref{prop:bc-carrier-target-window-halo-tail-supplier}.
\end{enumerate}
Then \cref{rem:former-post29-tail} holds.
\end{proposition}

\begin{proof}
The reduction is the same as
\cref{prop:bc-grouped-collective-carrier-splice-positive-height-tail-supplier},
with \cref{prop:bc-step3-boundary-readouts-supply-grouped-carrier-decoders}
providing the positive-height local inverse.  Item~(1) supplies the $B_b$
width estimate.

For a $C_b$ row with no zero-surplus height halo zone, every height halo zone
has positive surplus by \cref{cor:bc-height-halo-singleton-nonnegative}.
Item~(2) and
\cref{prop:bc-step3-boundary-readouts-supply-grouped-carrier-decoders} supply
the missing local inverse for every zone in
\cref{prop:bc-canonical-halo-local-repair-criterion}, giving the height
fixed-fiber estimate.

If the zero-surplus zone is full, item~(3) and
\cref{cor:bc-height-alternating-full-zone-decoder} give the height
fixed-fiber estimate.  In the remaining non-full zero-surplus case,
\cref{cor:bc-height-zero-surplus-terminal-unique} makes that zero-surplus zone
terminal and every zone strictly to its left positive.  Item~(4) supplies the
Step--3 boundary grouped-splice inverses for those left zones and the
final-zone target-window/right-boundary inverse for $Z_L$, so
\cref{prop:bc-right-boundary-auxiliary-fiber-criterion} gives the height
fixed-fiber estimate.

Thus the width and height fixed-fiber estimates hold for every residual row,
admissible $j$, and witness set $S$.  Applying
\cref{prop:bc-bad-shape-fiber-reduction} and then
\cref{prop:bc-bad-shape-reduction} gives \cref{rem:former-post29-tail}.
\end{proof}

\begin{proposition}[One-cell corner grouped splices imply the residual $B/C$ tail]\label{prop:bc-one-cell-corner-grouped-splice-positive-height-tail-supplier}
For each row
\[
  B_{14},\ldots,B_{61},\qquad C_{20},\ldots,C_{61},
\]
put $q=b+1$ for $G=B_b$ or $G=C_b$.  Fix once and for all an injection
\[
  \kappa:\{0,1,2\}\hookrightarrow\{0,1\}^2 .
\]
Suppose that for every such row, every $j$ in the residual window of
\cref{rem:former-post29-tail}, and every
$q$-element subset $S\subset\{1,\ldots,j\}$, the following repair data are
available.
\begin{enumerate}
\item For the width fixed-witness fiber when $G=B_b$, the canonical halo
zones satisfy the hypotheses of
\cref{prop:bc-below-wall-packet-splice-halo-local-repair-criterion}.
\item For the height fixed-witness fiber when $G=C_b$, order the canonical
halo zones increasingly.  If no height halo zone has zero surplus, then the
global parser, target path, replacement windows, branch words, and endpoint
pairs satisfy the hypotheses of
\cref{prop:bc-canonical-halo-local-repair-criterion} except for the local
inverse assertion.  For every such height zone $Z_\gamma$, the parser
decomposes $V^\ast_\gamma$ into grouped subwindows and parser-visible tag
sets satisfying the hypotheses of
\cref{prop:bc-one-cell-step3-corner-readouts-supply-grouped-carrier-decoders},
with the same fixed injection $\kappa$.
\item If a height zero-surplus zone exists and is the full alternating zone
$j=2q$ with $S=\{1,3,\ldots,2q-1\}$, use the fixed full-zone decoder of
\cref{cor:bc-height-alternating-full-zone-decoder}.
\item If a height zero-surplus zone exists and is not full, then, with this
terminal zone written $Z_L=\{a_L,\ldots,j\}$, the global parser, target path,
replacement windows, branch words, endpoint pairs, and unchanged complement
blocks satisfy the right-boundary auxiliary hypotheses of
\cref{prop:bc-right-boundary-auxiliary-global-recovery} except for the local
inverse assertions on the zones strictly left of $Z_L$.  Each zone strictly
left of $Z_L$ satisfies the one-cell corner grouped-splice hypothesis stated
in item~(2), and the final zone satisfies the final-zone
target-window/right-boundary alternative listed in item~(2) of
\cref{prop:bc-carrier-target-window-halo-tail-supplier}.
\end{enumerate}
Then \cref{rem:former-post29-tail} holds.
\end{proposition}

\begin{proof}
The proof is identical to
\cref{prop:bc-step3-boundary-grouped-splice-positive-height-tail-supplier},
with
\cref{prop:bc-one-cell-step3-corner-readouts-supply-grouped-carrier-decoders}
as the positive-height local inverse.  Item~(1) gives the $B_b$ width
estimate.  In the no-zero-surplus $C_b$ case, item~(2) supplies the local
inverse for every positive height zone in
\cref{prop:bc-canonical-halo-local-repair-criterion}.  In the full
zero-surplus case, item~(3) supplies the fixed full-zone decoder.  In the
remaining terminal zero-surplus case, item~(4) supplies the left-zone
one-cell corner inverses and the final-zone target-window/right-boundary
inverse, so \cref{prop:bc-right-boundary-auxiliary-fiber-criterion} applies.
Thus the width and height fixed-fiber estimates hold for every residual row,
admissible $j$, and witness set $S$, and
\cref{prop:bc-bad-shape-fiber-reduction,prop:bc-bad-shape-reduction} give
\cref{rem:former-post29-tail}.
\end{proof}

\begin{proposition}[Grouped padded-valley readouts imply the residual $B/C$ tail]\label{prop:bc-padded-valley-grouped-splice-positive-height-tail-supplier}
For each row
\[
  B_{14},\ldots,B_{61},\qquad C_{20},\ldots,C_{61},
\]
put $q=b+1$ for $G=B_b$ or $G=C_b$.  Fix once and for all an injection
\[
  \kappa:\{0,1,2\}\hookrightarrow\{0,1\}^2 .
\]
Suppose that for every such row, every $j$ in the residual window of
\cref{rem:former-post29-tail}, and every
$q$-element subset $S\subset\{1,\ldots,j\}$, the following repair data are
available.
\begin{enumerate}
\item For the width fixed-witness fiber when $G=B_b$, the canonical halo
zones satisfy the hypotheses of
\cref{prop:bc-below-wall-packet-splice-halo-local-repair-criterion}.
\item For the height fixed-witness fiber when $G=C_b$, order the canonical
halo zones increasingly.  If no height halo zone has zero surplus, then the
global parser, target path, replacement windows, branch words, and endpoint
pairs satisfy the hypotheses of
\cref{prop:bc-canonical-halo-local-repair-criterion} except for the local
inverse assertion.  For every such height zone $Z_\gamma$, the parser
decomposes $V^\ast_\gamma$ into grouped subwindows and parser-visible tag
sets satisfying the hypotheses of
\cref{prop:bc-padded-valleys-supply-one-cell-corner-readouts}, with the same
fixed injection $\kappa$.
\item If a height zero-surplus zone exists and is the full alternating zone
$j=2q$ with $S=\{1,3,\ldots,2q-1\}$, use the fixed full-zone decoder of
\cref{cor:bc-height-alternating-full-zone-decoder}.
\item If a height zero-surplus zone exists and is not full, then, with this
terminal zone written $Z_L=\{a_L,\ldots,j\}$, the global parser, target path,
replacement windows, branch words, endpoint pairs, and unchanged complement
blocks satisfy the right-boundary auxiliary hypotheses of
\cref{prop:bc-right-boundary-auxiliary-global-recovery} except for the local
inverse assertions on the zones strictly left of $Z_L$.  Each zone strictly
left of $Z_L$ satisfies the grouped padded-valley hypothesis stated in
item~(2), and the final zone satisfies the final-zone
target-window/right-boundary alternative listed in item~(2) of
\cref{prop:bc-carrier-target-window-halo-tail-supplier}.
\end{enumerate}
Then \cref{rem:former-post29-tail} holds.
\end{proposition}

\begin{proof}
The proof is identical to
\cref{prop:bc-one-cell-corner-grouped-splice-positive-height-tail-supplier},
with
\cref{prop:bc-padded-valleys-supply-one-cell-corner-readouts}
as the positive-height local inverse.  Item~(1) gives the $B_b$ width
estimate.  In the no-zero-surplus $C_b$ case, item~(2) supplies the local
inverse for every positive height zone in
\cref{prop:bc-canonical-halo-local-repair-criterion}.  In the full
zero-surplus case, item~(3) supplies the fixed full-zone decoder.  In the
remaining terminal zero-surplus case, item~(4) supplies the left-zone
grouped padded-valley inverses and the final-zone target-window/right-boundary
inverse, so \cref{prop:bc-right-boundary-auxiliary-fiber-criterion} applies.
Thus the width and height fixed-fiber estimates hold for every residual row,
admissible $j$, and witness set $S$, and
\cref{prop:bc-bad-shape-fiber-reduction,prop:bc-bad-shape-reduction} give
\cref{rem:former-post29-tail}.
\end{proof}

\begin{proposition}[Grouped midpoint-defect readouts imply the residual $B/C$ tail]\label{prop:bc-midpoint-defect-grouped-splice-positive-height-tail-supplier}
For each row
\[
  B_{14},\ldots,B_{61},\qquad C_{20},\ldots,C_{61},
\]
put $q=b+1$ for $G=B_b$ or $G=C_b$.  Fix once and for all an injection
\[
  \kappa:\{0,1,2\}\hookrightarrow\{0,1\}^2 .
\]
Suppose that for every such row, every $j$ in the residual window of
\cref{rem:former-post29-tail}, and every
$q$-element subset $S\subset\{1,\ldots,j\}$, the following repair data are
available.
\begin{enumerate}
\item For the width fixed-witness fiber when $G=B_b$, the canonical halo
zones satisfy the hypotheses of
\cref{prop:bc-below-wall-packet-splice-halo-local-repair-criterion}.
\item For the height fixed-witness fiber when $G=C_b$, order the canonical
halo zones increasingly.  If no height halo zone has zero surplus, then the
global parser, target path, replacement windows, branch words, and endpoint
pairs satisfy the hypotheses of
\cref{prop:bc-canonical-halo-local-repair-criterion} except for the local
inverse assertion.  For every such height zone $Z_\gamma$, the parser
decomposes $V^\ast_\gamma$ into grouped subwindows and parser-visible tag
sets satisfying the hypotheses of
\cref{prop:bc-midpoint-defect-readouts-supply-grouped-padded-valleys}, with
the same fixed injection $\kappa$.
\item If a height zero-surplus zone exists and is the full alternating zone
$j=2q$ with $S=\{1,3,\ldots,2q-1\}$, use the fixed full-zone decoder of
\cref{cor:bc-height-alternating-full-zone-decoder}.
\item If a height zero-surplus zone exists and is not full, then, with this
terminal zone written $Z_L=\{a_L,\ldots,j\}$, the global parser, target path,
replacement windows, branch words, endpoint pairs, and unchanged complement
blocks satisfy the right-boundary auxiliary hypotheses of
\cref{prop:bc-right-boundary-auxiliary-global-recovery} except for the local
inverse assertions on the zones strictly left of $Z_L$.  Each zone strictly
left of $Z_L$ satisfies the grouped midpoint-defect hypothesis stated in
item~(2), and the final zone satisfies the final-zone
target-window/right-boundary alternative listed in item~(2) of
\cref{prop:bc-carrier-target-window-halo-tail-supplier}.
\end{enumerate}
Then \cref{rem:former-post29-tail} holds.
\end{proposition}

\begin{proof}
The proof is identical to
\cref{prop:bc-padded-valley-grouped-splice-positive-height-tail-supplier},
with
\cref{prop:bc-midpoint-defect-readouts-supply-grouped-padded-valleys}
as the positive-height local inverse.  Item~(1) gives the $B_b$ width
estimate.  In the no-zero-surplus $C_b$ case, item~(2) supplies the local
inverse for every positive height zone in
\cref{prop:bc-canonical-halo-local-repair-criterion}.  In the full
zero-surplus case, item~(3) supplies the fixed full-zone decoder.  In the
remaining terminal zero-surplus case, item~(4) supplies the left-zone
grouped midpoint-defect inverses and the final-zone target-window/right-boundary
inverse, so \cref{prop:bc-right-boundary-auxiliary-fiber-criterion} applies.
Thus the width and height fixed-fiber estimates hold for every residual row,
admissible $j$, and witness set $S$, and
\cref{prop:bc-bad-shape-fiber-reduction,prop:bc-bad-shape-reduction} give
\cref{rem:former-post29-tail}.
\end{proof}

\begin{proposition}[Grouped singleton-defect row readouts imply the residual $B/C$ tail]\label{prop:bc-singleton-defect-row-grouped-splice-positive-height-tail-supplier}
For each row
\[
  B_{14},\ldots,B_{61},\qquad C_{20},\ldots,C_{61},
\]
put $q=b+1$ for $G=B_b$ or $G=C_b$.  Fix once and for all an injection
\[
  \kappa:\{0,1,2\}\hookrightarrow\{0,1\}^2 .
\]
Suppose that for every such row, every $j$ in the residual window of
\cref{rem:former-post29-tail}, and every
$q$-element subset $S\subset\{1,\ldots,j\}$, the following repair data are
available.
\begin{enumerate}
\item For the width fixed-witness fiber when $G=B_b$, the canonical halo
zones satisfy the hypotheses of
\cref{prop:bc-below-wall-packet-splice-halo-local-repair-criterion}.
\item For the height fixed-witness fiber when $G=C_b$, order the canonical
halo zones increasingly.  If no height halo zone has zero surplus, then the
global parser, target path, replacement windows, branch words, and endpoint
pairs satisfy the hypotheses of
\cref{prop:bc-canonical-halo-local-repair-criterion} except for the local
inverse assertion.  For every such height zone $Z_\gamma$, the parser
decomposes $V^\ast_\gamma$ into grouped subwindows and parser-visible tag
sets satisfying the hypotheses of
\cref{prop:bc-singleton-defect-row-readouts-supply-grouped-padded-valleys},
with the same fixed injection $\kappa$.
\item If a height zero-surplus zone exists and is the full alternating zone
$j=2q$ with $S=\{1,3,\ldots,2q-1\}$, use the fixed full-zone decoder of
\cref{cor:bc-height-alternating-full-zone-decoder}.
\item If a height zero-surplus zone exists and is not full, then, with this
terminal zone written $Z_L=\{a_L,\ldots,j\}$, the global parser, target path,
replacement windows, branch words, endpoint pairs, and unchanged complement
blocks satisfy the right-boundary auxiliary hypotheses of
\cref{prop:bc-right-boundary-auxiliary-global-recovery} except for the local
inverse assertions on the zones strictly left of $Z_L$.  Each zone strictly
left of $Z_L$ satisfies the grouped singleton-defect-row hypothesis stated in
item~(2), and the final zone satisfies the final-zone
target-window/right-boundary alternative listed in item~(2) of
\cref{prop:bc-carrier-target-window-halo-tail-supplier}.
\end{enumerate}
Then \cref{rem:former-post29-tail} holds.
\end{proposition}

\begin{proof}
The proof is identical to
\cref{prop:bc-midpoint-defect-grouped-splice-positive-height-tail-supplier},
with
\cref{prop:bc-singleton-defect-row-readouts-supply-grouped-padded-valleys}
as the positive-height local inverse.  Item~(1) gives the $B_b$ width
estimate.  In the no-zero-surplus $C_b$ case, item~(2) supplies the local
inverse for every positive height zone in
\cref{prop:bc-canonical-halo-local-repair-criterion}.  In the full
zero-surplus case, item~(3) supplies the fixed full-zone decoder.  In the
remaining terminal zero-surplus case, item~(4) supplies the left-zone grouped
singleton-defect-row inverses and the final-zone target-window/right-boundary
inverse, so \cref{prop:bc-right-boundary-auxiliary-fiber-criterion} applies.
Thus the width and height fixed-fiber estimates hold for every residual row,
admissible $j$, and witness set $S$, and
\cref{prop:bc-bad-shape-fiber-reduction,prop:bc-bad-shape-reduction} give
\cref{rem:former-post29-tail}.
\end{proof}

\begin{proposition}[Grouped northeast-box row readouts imply the residual $B/C$ tail]\label{prop:bc-northeast-box-row-grouped-splice-positive-height-tail-supplier}
For each row
\[
  B_{14},\ldots,B_{61},\qquad C_{20},\ldots,C_{61},
\]
put $q=b+1$ for $G=B_b$ or $G=C_b$.  Fix once and for all an injection
\[
  \kappa:\{0,1,2\}\hookrightarrow\{0,1\}^2 .
\]
Suppose that for every such row, every $j$ in the residual window of
\cref{rem:former-post29-tail}, and every
$q$-element subset $S\subset\{1,\ldots,j\}$, the following repair data are
available.
\begin{enumerate}
\item For the width fixed-witness fiber when $G=B_b$, the canonical halo
zones satisfy the hypotheses of
\cref{prop:bc-below-wall-packet-splice-halo-local-repair-criterion}.
\item For the height fixed-witness fiber when $G=C_b$, order the canonical
halo zones increasingly.  If no height halo zone has zero surplus, then the
global parser, target path, replacement windows, branch words, and endpoint
pairs satisfy the hypotheses of
\cref{prop:bc-canonical-halo-local-repair-criterion} except for the local
inverse assertion.  For every such height zone $Z_\gamma$, the parser
decomposes $V^\ast_\gamma$ into grouped subwindows and parser-visible tag
sets satisfying the hypotheses of
\cref{prop:bc-northeast-box-row-readouts-supply-grouped-padded-valleys}, with
the same fixed injection $\kappa$.
\item If a height zero-surplus zone exists and is the full alternating zone
$j=2q$ with $S=\{1,3,\ldots,2q-1\}$, use the fixed full-zone decoder of
\cref{cor:bc-height-alternating-full-zone-decoder}.
\item If a height zero-surplus zone exists and is not full, then, with this
terminal zone written $Z_L=\{a_L,\ldots,j\}$, the global parser, target path,
replacement windows, branch words, endpoint pairs, and unchanged complement
blocks satisfy the right-boundary auxiliary hypotheses of
\cref{prop:bc-right-boundary-auxiliary-global-recovery} except for the local
inverse assertions on the zones strictly left of $Z_L$.  Each zone strictly
left of $Z_L$ satisfies the grouped northeast-box-row hypothesis stated in
item~(2), and the final zone satisfies the final-zone
target-window/right-boundary alternative listed in item~(2) of
\cref{prop:bc-carrier-target-window-halo-tail-supplier}.
\end{enumerate}
Then \cref{rem:former-post29-tail} holds.
\end{proposition}

\begin{proof}
The proof is identical to
\cref{prop:bc-singleton-defect-row-grouped-splice-positive-height-tail-supplier},
with
\cref{prop:bc-northeast-box-row-readouts-supply-grouped-padded-valleys}
as the positive-height local inverse.  Item~(1) gives the $B_b$ width
estimate.  In the no-zero-surplus $C_b$ case, item~(2) supplies the local
inverse for every positive height zone in
\cref{prop:bc-canonical-halo-local-repair-criterion}.  In the full
zero-surplus case, item~(3) supplies the fixed full-zone decoder.  In the
remaining terminal zero-surplus case, item~(4) supplies the left-zone grouped
northeast-box-row inverses and the final-zone target-window/right-boundary
inverse, so \cref{prop:bc-right-boundary-auxiliary-fiber-criterion} applies.
Thus the width and height fixed-fiber estimates hold for every residual row,
admissible $j$, and witness set $S$, and
\cref{prop:bc-bad-shape-fiber-reduction,prop:bc-bad-shape-reduction} give
\cref{rem:former-post29-tail}.
\end{proof}

\begin{proposition}[Visible-gated northeast-box row readouts imply the residual $B/C$ tail]\label{prop:bc-visible-gated-northeast-box-row-positive-height-tail-supplier}
For each row
\[
  B_{14},\ldots,B_{61},\qquad C_{20},\ldots,C_{61},
\]
put $q=b+1$ for $G=B_b$ or $G=C_b$.  Fix once and for all an injection
\[
  \kappa:\{0,1,2\}\hookrightarrow\{0,1\}^2 .
\]
Suppose that for every such row, every $j$ in the residual window of
\cref{rem:former-post29-tail}, and every
$q$-element subset $S\subset\{1,\ldots,j\}$, the following repair data are
available.
\begin{enumerate}
\item For the width fixed-witness fiber when $G=B_b$, the canonical halo
zones satisfy the hypotheses of
\cref{prop:bc-below-wall-packet-splice-halo-local-repair-criterion}.
\item For the height fixed-witness fiber when $G=C_b$, order the canonical
halo zones increasingly.  If no height halo zone has zero surplus, then the
global parser, target path, replacement windows, branch words, and endpoint
pairs satisfy the hypotheses of
\cref{prop:bc-canonical-halo-local-repair-criterion} except for the local
inverse assertion.  For every such height zone $Z_\gamma$, the parser
decomposes $V^\ast_\gamma$ into grouped subwindows and parser-visible tag
sets satisfying the hypotheses of
\cref{prop:bc-visible-gated-northeast-row-readouts-supply-grouped-padded-valleys},
with the same fixed injection $\kappa$.
\item If a height zero-surplus zone exists and is the full alternating zone
$j=2q$ with $S=\{1,3,\ldots,2q-1\}$, use the fixed full-zone decoder of
\cref{cor:bc-height-alternating-full-zone-decoder}.
\item If a height zero-surplus zone exists and is not full, then, with this
terminal zone written $Z_L=\{a_L,\ldots,j\}$, the global parser, target path,
replacement windows, branch words, endpoint pairs, and unchanged complement
blocks satisfy the right-boundary auxiliary hypotheses of
\cref{prop:bc-right-boundary-auxiliary-global-recovery} except for the local
inverse assertions on the zones strictly left of $Z_L$.  Each zone strictly
left of $Z_L$ satisfies the visible-gated northeast-box row hypothesis stated
in item~(2), and the final zone satisfies the final-zone
target-window/right-boundary alternative listed in item~(2) of
\cref{prop:bc-carrier-target-window-halo-tail-supplier}.
\end{enumerate}
Then \cref{rem:former-post29-tail} holds.
\end{proposition}

\begin{proof}
The proof is identical to
\cref{prop:bc-northeast-box-row-grouped-splice-positive-height-tail-supplier},
with
\cref{prop:bc-visible-gated-northeast-row-readouts-supply-grouped-padded-valleys}
as the positive-height local inverse.  Item~(1) gives the $B_b$ width
estimate.  In the no-zero-surplus $C_b$ case, item~(2) supplies the local
inverse for every positive height zone in
\cref{prop:bc-canonical-halo-local-repair-criterion}.  In the full
zero-surplus case, item~(3) supplies the fixed full-zone decoder.  In the
remaining terminal zero-surplus case, item~(4) supplies the left-zone
visible-gated northeast-box row inverses and the final-zone
target-window/right-boundary inverse, so
\cref{prop:bc-right-boundary-auxiliary-fiber-criterion} applies.  Thus the
width and height fixed-fiber estimates hold for every residual row, admissible
$j$, and witness set $S$, and
\cref{prop:bc-bad-shape-fiber-reduction,prop:bc-bad-shape-reduction} give
\cref{rem:former-post29-tail}.
\end{proof}

\begin{proposition}[Grouped removable-row readouts imply the residual $B/C$ tail]\label{prop:bc-removable-row-grouped-splice-positive-height-tail-supplier}
For each row
\[
  B_{14},\ldots,B_{61},\qquad C_{20},\ldots,C_{61},
\]
put $q=b+1$ for $G=B_b$ or $G=C_b$.  Fix once and for all an injection
\[
  \kappa:\{0,1,2\}\hookrightarrow\{0,1\}^2 .
\]
Suppose that for every such row, every $j$ in the residual window of
\cref{rem:former-post29-tail}, and every
$q$-element subset $S\subset\{1,\ldots,j\}$, the following repair data are
available.
\begin{enumerate}
\item For the width fixed-witness fiber when $G=B_b$, the canonical halo
zones satisfy the hypotheses of
\cref{prop:bc-below-wall-packet-splice-halo-local-repair-criterion}.
\item For the height fixed-witness fiber when $G=C_b$, order the canonical
halo zones increasingly.  If no height halo zone has zero surplus, then the
global parser, target path, replacement windows, branch words, and endpoint
pairs satisfy the hypotheses of
\cref{prop:bc-canonical-halo-local-repair-criterion} except for the local
inverse assertion.  For every such height zone $Z_\gamma$, the parser
decomposes $V^\ast_\gamma$ into grouped subwindows and parser-visible tag
sets satisfying the hypotheses of
\cref{prop:bc-removable-row-readouts-supply-grouped-padded-valleys}, with the
same fixed injection $\kappa$.
\item If a height zero-surplus zone exists and is the full alternating zone
$j=2q$ with $S=\{1,3,\ldots,2q-1\}$, use the fixed full-zone decoder of
\cref{cor:bc-height-alternating-full-zone-decoder}.
\item If a height zero-surplus zone exists and is not full, then, with this
terminal zone written $Z_L=\{a_L,\ldots,j\}$, the global parser, target path,
replacement windows, branch words, endpoint pairs, and unchanged complement
blocks satisfy the right-boundary auxiliary hypotheses of
\cref{prop:bc-right-boundary-auxiliary-global-recovery} except for the local
inverse assertions on the zones strictly left of $Z_L$.  Each zone strictly
left of $Z_L$ satisfies the grouped removable-row hypothesis stated in
item~(2), and the final zone satisfies the final-zone
target-window/right-boundary alternative listed in item~(2) of
\cref{prop:bc-carrier-target-window-halo-tail-supplier}.
\end{enumerate}
Then \cref{rem:former-post29-tail} holds.
\end{proposition}

\begin{proof}
The proof is identical to
\cref{prop:bc-northeast-box-row-grouped-splice-positive-height-tail-supplier},
with
\cref{prop:bc-removable-row-readouts-supply-grouped-padded-valleys}
as the positive-height local inverse.  Item~(1) gives the $B_b$ width
estimate.  In the no-zero-surplus $C_b$ case, item~(2) supplies the local
inverse for every positive height zone in
\cref{prop:bc-canonical-halo-local-repair-criterion}.  In the full
zero-surplus case, item~(3) supplies the fixed full-zone decoder.  In the
remaining terminal zero-surplus case, item~(4) supplies the left-zone grouped
removable-row inverses and the final-zone target-window/right-boundary inverse,
so \cref{prop:bc-right-boundary-auxiliary-fiber-criterion} applies.  Thus the
width and height fixed-fiber estimates hold for every residual row, admissible
$j$, and witness set $S$, and
\cref{prop:bc-bad-shape-fiber-reduction,prop:bc-bad-shape-reduction} give
\cref{rem:former-post29-tail}.
\end{proof}

\begin{proposition}[Visible-gated grouped removable-row readouts imply the residual $B/C$ tail]\label{prop:bc-visible-gated-removable-row-grouped-splice-positive-height-tail-supplier}
For each row
\[
  B_{14},\ldots,B_{61},\qquad C_{20},\ldots,C_{61},
\]
put $q=b+1$ for $G=B_b$ or $G=C_b$.  Fix once and for all an injection
\[
  \kappa:\{0,1,2\}\hookrightarrow\{0,1\}^2 .
\]
Suppose that for every such row, every $j$ in the residual window of
\cref{rem:former-post29-tail}, and every
$q$-element subset $S\subset\{1,\ldots,j\}$, the following repair data are
available.
\begin{enumerate}
\item For the width fixed-witness fiber when $G=B_b$, the canonical halo
zones satisfy the hypotheses of
\cref{prop:bc-below-wall-packet-splice-halo-local-repair-criterion}.
\item For the height fixed-witness fiber when $G=C_b$, order the canonical
halo zones increasingly.  If no height halo zone has zero surplus, then the
global parser, target path, replacement windows, branch words, and endpoint
pairs satisfy the hypotheses of
\cref{prop:bc-canonical-halo-local-repair-criterion} except for the local
inverse assertion.  For every such height zone $Z_\gamma$, the parser
decomposes $V^\ast_\gamma$ into grouped subwindows and parser-visible tag
sets satisfying the hypotheses of
\cref{prop:bc-visible-gated-removable-row-readouts-supply-grouped-padded-valleys},
with the same fixed injection $\kappa$.
\item If a height zero-surplus zone exists and is the full alternating zone
$j=2q$ with $S=\{1,3,\ldots,2q-1\}$, use the fixed full-zone decoder of
\cref{cor:bc-height-alternating-full-zone-decoder}.
\item If a height zero-surplus zone exists and is not full, then, with this
terminal zone written $Z_L=\{a_L,\ldots,j\}$, the global parser, target path,
replacement windows, branch words, endpoint pairs, and unchanged complement
blocks satisfy the right-boundary auxiliary hypotheses of
\cref{prop:bc-right-boundary-auxiliary-global-recovery} except for the local
inverse assertions on the zones strictly left of $Z_L$.  Each zone strictly
left of $Z_L$ satisfies the visible-gated grouped removable-row hypothesis
stated in item~(2), and the final zone satisfies the final-zone
target-window/right-boundary alternative listed in item~(2) of
\cref{prop:bc-carrier-target-window-halo-tail-supplier}.
\end{enumerate}
Then \cref{rem:former-post29-tail} holds.
\end{proposition}

\begin{proof}
The proof is identical to
\cref{prop:bc-removable-row-grouped-splice-positive-height-tail-supplier}, with
\cref{prop:bc-visible-gated-removable-row-readouts-supply-grouped-padded-valleys}
as the positive-height local inverse.  Item~(1) gives the $B_b$ width
estimate.  In the no-zero-surplus $C_b$ case, item~(2) supplies the local
inverse for every positive height zone in
\cref{prop:bc-canonical-halo-local-repair-criterion}.  In the full
zero-surplus case, item~(3) supplies the fixed full-zone decoder.  In the
remaining terminal zero-surplus case, item~(4) supplies the left-zone
visible-gated removable-row inverses and the final-zone
target-window/right-boundary inverse, so
\cref{prop:bc-right-boundary-auxiliary-fiber-criterion} applies.  Thus the
width and height fixed-fiber estimates hold for every residual row, admissible
$j$, and witness set $S$, and
\cref{prop:bc-bad-shape-fiber-reduction,prop:bc-bad-shape-reduction} give
\cref{rem:former-post29-tail}.
\end{proof}

\begin{proposition}[Visible-gated midpoint readouts imply the residual $B/C$ tail]\label{prop:bc-visible-gated-midpoint-grouped-splice-positive-height-tail-supplier}
For each row
\[
  B_{14},\ldots,B_{61},\qquad C_{20},\ldots,C_{61},
\]
put $q=b+1$ for $G=B_b$ or $G=C_b$.  Fix once and for all an injection
\[
  \kappa:\{0,1,2\}\hookrightarrow\{0,1\}^2 .
\]
Suppose that for every such row, every $j$ in the residual window of
\cref{rem:former-post29-tail}, and every
$q$-element subset $S\subset\{1,\ldots,j\}$, the following repair data are
available.
\begin{enumerate}
\item For the width fixed-witness fiber when $G=B_b$, the canonical halo
zones satisfy the hypotheses of
\cref{prop:bc-below-wall-packet-splice-halo-local-repair-criterion}.
\item For the height fixed-witness fiber when $G=C_b$, order the canonical
halo zones increasingly.  If no height halo zone has zero surplus, then the
global parser, target path, replacement windows, branch words, and endpoint
pairs satisfy the hypotheses of
\cref{prop:bc-canonical-halo-local-repair-criterion} except for the local
inverse assertion.  For every such height zone $Z_\gamma$, the parser
decomposes $V^\ast_\gamma$ into grouped subwindows and parser-visible tag
sets satisfying the hypotheses of
\cref{prop:bc-visible-gated-midpoint-readouts-supply-grouped-padded-valleys},
with the same fixed injection $\kappa$.
\item If a height zero-surplus zone exists and is the full alternating zone
$j=2q$ with $S=\{1,3,\ldots,2q-1\}$, use the fixed full-zone decoder of
\cref{cor:bc-height-alternating-full-zone-decoder}.
\item If a height zero-surplus zone exists and is not full, then, with this
terminal zone written $Z_L=\{a_L,\ldots,j\}$, the global parser, target path,
replacement windows, branch words, endpoint pairs, and unchanged complement
blocks satisfy the right-boundary auxiliary hypotheses of
\cref{prop:bc-right-boundary-auxiliary-global-recovery} except for the local
inverse assertions on the zones strictly left of $Z_L$.  Each zone strictly
left of $Z_L$ satisfies the visible-gated midpoint grouped-splice hypothesis
stated in item~(2), and the final zone satisfies the final-zone
target-window/right-boundary alternative listed in item~(2) of
\cref{prop:bc-carrier-target-window-halo-tail-supplier}.
\end{enumerate}
Then \cref{rem:former-post29-tail} holds.
\end{proposition}

\begin{proof}
The proof is identical to
\cref{prop:bc-visible-gated-removable-row-grouped-splice-positive-height-tail-supplier},
with
\cref{prop:bc-visible-gated-midpoint-readouts-supply-grouped-padded-valleys}
as the positive-height local inverse.  Item~(1) gives the $B_b$ width
estimate.  In the no-zero-surplus $C_b$ case, item~(2) supplies the local
inverse for every positive height zone in
\cref{prop:bc-canonical-halo-local-repair-criterion}.  In the full
zero-surplus case, item~(3) supplies the fixed full-zone decoder.  In the
remaining terminal zero-surplus case, item~(4) supplies the left-zone
visible-gated midpoint inverses and the final-zone target-window/right-boundary
inverse, so \cref{prop:bc-right-boundary-auxiliary-fiber-criterion} applies.
Thus the width and height fixed-fiber estimates hold for every residual row,
admissible $j$, and witness set $S$, and
\cref{prop:bc-bad-shape-fiber-reduction,prop:bc-bad-shape-reduction} give
\cref{rem:former-post29-tail}.
\end{proof}

\begin{proposition}[Visible-gated Step--3 boundary readouts imply the residual $B/C$ tail]\label{prop:bc-visible-gated-step3-boundary-positive-height-tail-supplier}
For each row
\[
  B_{14},\ldots,B_{61},\qquad C_{20},\ldots,C_{61},
\]
put $q=b+1$ for $G=B_b$ or $G=C_b$.  Fix once and for all an injection
\[
  \kappa:\{0,1,2\}\hookrightarrow\{0,1\}^2 .
\]
Suppose that for every such row, every $j$ in the residual window of
\cref{rem:former-post29-tail}, and every
$q$-element subset $S\subset\{1,\ldots,j\}$, the following repair data are
available.
\begin{enumerate}
\item For the width fixed-witness fiber when $G=B_b$, the canonical halo
zones satisfy the hypotheses of
\cref{prop:bc-below-wall-packet-splice-halo-local-repair-criterion}.
\item For the height fixed-witness fiber when $G=C_b$, order the canonical
halo zones increasingly.  If no height halo zone has zero surplus, then the
global parser, target path, replacement windows, branch words, and endpoint
pairs satisfy the hypotheses of
\cref{prop:bc-canonical-halo-local-repair-criterion} except for the local
inverse assertion.  For every such height zone $Z_\gamma$, the parser
decomposes $V^\ast_\gamma$ into grouped subwindows and parser-visible tag
sets satisfying the hypotheses of
\cref{cor:bc-visible-gated-step3-boundary-suffices}, with the same fixed
injection $\kappa$.
\item If a height zero-surplus zone exists and is the full alternating zone
$j=2q$ with $S=\{1,3,\ldots,2q-1\}$, use the fixed full-zone decoder of
\cref{cor:bc-height-alternating-full-zone-decoder}.
\item If a height zero-surplus zone exists and is not full, then, with this
terminal zone written $Z_L=\{a_L,\ldots,j\}$, the global parser, target path,
replacement windows, branch words, endpoint pairs, and unchanged complement
blocks satisfy the right-boundary auxiliary hypotheses of
\cref{prop:bc-right-boundary-auxiliary-global-recovery} except for the local
inverse assertions on the zones strictly left of $Z_L$.  Each zone strictly
left of $Z_L$ satisfies the visible-gated Step--3 boundary hypothesis stated
in item~(2), and the final zone satisfies the final-zone
target-window/right-boundary alternative listed in item~(2) of
\cref{prop:bc-carrier-target-window-halo-tail-supplier}.
\end{enumerate}
Then \cref{rem:former-post29-tail} holds.
\end{proposition}

\begin{proof}
For every positive height zone appearing in items~(2) and (4),
\cref{cor:bc-visible-gated-step3-boundary-suffices} converts the stated
visible-gated Step--3 boundary readout into the hypotheses of
\cref{prop:bc-visible-gated-midpoint-readouts-supply-grouped-padded-valleys}.
Thus the hypotheses of
\cref{prop:bc-visible-gated-midpoint-grouped-splice-positive-height-tail-supplier}
are satisfied with the same width data, the same full-zone decoder, and the
same terminal right-boundary alternative.  Applying that proposition gives
\cref{rem:former-post29-tail}.
\end{proof}

\begin{proposition}[Visible-gated target blocks imply the residual $B/C$ tail]\label{prop:bc-visible-gated-target-block-positive-height-tail-supplier}
For each row
\[
  B_{14},\ldots,B_{61},\qquad C_{20},\ldots,C_{61},
\]
put $q=b+1$ for $G=B_b$ or $G=C_b$.  Fix once and for all an injection
\[
  \kappa:\{0,1,2\}\hookrightarrow\{0,1\}^2 .
\]
Suppose that for every such row, every $j$ in the residual window of
\cref{rem:former-post29-tail}, and every
$q$-element subset $S\subset\{1,\ldots,j\}$, the following repair data are
available.
\begin{enumerate}
\item For the width fixed-witness fiber when $G=B_b$, the canonical halo
zones satisfy the hypotheses of
\cref{prop:bc-below-wall-packet-splice-halo-local-repair-criterion}.
\item For the height fixed-witness fiber when $G=C_b$, order the canonical
halo zones increasingly.  If no height halo zone has zero surplus, then the
global parser, target path, replacement windows, branch words, and endpoint
pairs satisfy the hypotheses of
\cref{prop:bc-canonical-halo-local-repair-criterion} except for the local
inverse assertion.  Every such height zone $Z_\gamma$ satisfies the
visible-gated target-block hypotheses of
\cref{prop:bc-visible-gated-target-blocks-recover-positive-height-zones}, with
the same fixed injection $\kappa$.
\item If a height zero-surplus zone exists and is the full alternating zone
$j=2q$ with $S=\{1,3,\ldots,2q-1\}$, use the fixed full-zone decoder of
\cref{cor:bc-height-alternating-full-zone-decoder}.
\item If a height zero-surplus zone exists and is not full, then, with this
terminal zone written $Z_L=\{a_L,\ldots,j\}$, the global parser, target path,
replacement windows, branch words, endpoint pairs, and unchanged complement
blocks satisfy the right-boundary auxiliary hypotheses of
\cref{prop:bc-right-boundary-auxiliary-global-recovery} except for the local
inverse assertions on the zones strictly left of $Z_L$.  Each zone strictly
left of $Z_L$ satisfies the visible-gated target-block hypothesis stated in
item~(2), and the final zone satisfies the final-zone
target-window/right-boundary alternative listed in item~(2) of
\cref{prop:bc-carrier-target-window-halo-tail-supplier}.
\end{enumerate}
Then \cref{rem:former-post29-tail} holds.
\end{proposition}

\begin{proof}
The proof is identical to
\cref{prop:bc-visible-gated-midpoint-grouped-splice-positive-height-tail-supplier},
with \cref{prop:bc-visible-gated-target-blocks-recover-positive-height-zones}
as the positive-height local inverse.  Item~(1) gives the $B_b$ width
estimate.  In the no-zero-surplus $C_b$ case, item~(2) supplies the local
inverse for every positive height zone in
\cref{prop:bc-canonical-halo-local-repair-criterion}.  In the full
zero-surplus case, item~(3) supplies the fixed full-zone decoder.  In the
remaining terminal zero-surplus case, item~(4) supplies the left-zone
visible-gated target-block inverses and the final-zone
target-window/right-boundary inverse, so
\cref{prop:bc-right-boundary-auxiliary-fiber-criterion} applies.  Thus the
width and height fixed-fiber estimates hold for every residual row, admissible
$j$, and witness set $S$, and
\cref{prop:bc-bad-shape-fiber-reduction,prop:bc-bad-shape-reduction} give
\cref{rem:former-post29-tail}.
\end{proof}

\fi

\subsection*{Return to the active proof: the half-stable bridge}

\begin{proposition}[Post-$m=29$ $B/C$ half-stable bridge reduction]\label{prop:post29-bc-half-bridge}
For the rows
\[
  B_{14},\ldots,B_{123},\qquad
  C_{20},\ldots,C_{217},
\]
let $r_G$ be the rank of $G$ and let $N_{\mathrm{dir}}^{BC}(G)$ be the
direct-tail onset of \cref{prop:post29-bc-interval-direct}.  Put
$s_j=m_j^{\mathrm{st}}$ and $m_j^G=s_j+\delta_j^G$.  Suppose that
\[
  -\frac12 s_j\le\delta_j^{B_b}
\]
for every \(14\le b\le41\) and every integer \(j\) with
\[
  2b+31\le j\le N_{\mathrm{dir}}^{BC}(B_b)+2,
\]
and that
\[
  -\frac12 s_j\le\delta_j^{B_b}
\]
for every \(42\le b\le123\) and every integer \(j\) with
\[
  2b+32\le j\le N_{\mathrm{dir}}^{BC}(B_b)+2,
\]
and that
\[
  -\frac12s_j\le\delta_j^{C_b}
\]
for every \(20\le b\le217\) and every integer \(j\) with
\[
  2b+32\le j\le N_{\mathrm{dir}}^{BC}(C_b)+2.
\]
Then, for every integer $m$ with
\[
  N_{\mathrm{hit}}^{(29)}(G)\le 2m+1<N_{\mathrm{dir}}^{BC}(G),
\]
one has
\[
  D_G(2m+1)\ge0.
\]
More locally, for a fixed displayed row and a fixed such $m$, the same
conclusion follows if
\[
 -\frac12s_j\le\delta_j^G\le0
 \qquad(2r_G+2\le j\le2m+5)
\]
without any hypothesis on later correction variables.
\end{proposition}

\begin{proof}
First combine the hypotheses with
\cref{prop:bc-b-boundary-lower-correction,prop:bc-b-first-nonboundary-lower-correction,prop:bc-b-second-nonboundary-lower-correction,prop:bc-b-third-nonboundary-lower-correction,prop:bc-b-fourth-nonboundary-lower-correction,prop:bc-b-fifth-nonboundary-lower-correction,prop:bc-b-sixth-nonboundary-lower-correction,prop:bc-b-seventh-nonboundary-lower-correction,prop:bc-b-eighth-nonboundary-lower-correction,prop:bc-b-ninth-nonboundary-lower-correction,prop:bc-b-tenth-nonboundary-lower-correction,prop:bc-b-eleventh-nonboundary-lower-correction,prop:bc-b-twelfth-nonboundary-lower-correction,prop:bc-b-thirteenth-nonboundary-lower-correction,prop:bc-b-fourteenth-nonboundary-lower-correction,prop:bc-b-fifteenth-nonboundary-lower-correction,prop:bc-b-sixteenth-nonboundary-lower-correction,prop:bc-b-seventeenth-nonboundary-lower-correction,prop:bc-b-eighteenth-nonboundary-lower-correction,prop:bc-b-nineteenth-nonboundary-lower-correction,prop:bc-c-boundary-lower-correction,prop:bc-c-first-nonboundary-lower-correction,prop:bc-c-second-nonboundary-lower-correction,prop:bc-c-third-nonboundary-lower-correction,prop:bc-c-fourth-nonboundary-lower-correction,prop:bc-c-fifth-nonboundary-lower-correction,prop:bc-c-sixth-nonboundary-lower-correction,prop:bc-c-seventh-nonboundary-lower-correction,prop:bc-c-eighth-nonboundary-lower-correction,prop:bc-c-ninth-nonboundary-lower-correction,prop:bc-c-tenth-nonboundary-lower-correction,prop:bc-c-eleventh-nonboundary-lower-correction,prop:bc-c-twelfth-nonboundary-lower-correction,prop:bc-c-thirteenth-nonboundary-lower-correction,prop:bc-c-fourteenth-nonboundary-lower-correction,prop:bc-c-fifteenth-nonboundary-lower-correction,prop:bc-c-sixteenth-nonboundary-lower-correction,prop:bc-c-seventeenth-nonboundary-lower-correction,prop:bc-c-eighteenth-nonboundary-lower-correction,prop:bc-c-nineteenth-nonboundary-lower-correction}.
Also use the twentieth, twenty-first, twenty-second, twenty-third, twenty-fourth, twenty-fifth, twenty-sixth, twenty-seventh, twenty-eighth,
B-side high-row twenty-ninth, and C-side twenty-ninth
nonboundary theorems
\cref{prop:bc-b-twentieth-nonboundary-lower-correction,prop:bc-c-twentieth-nonboundary-lower-correction,prop:bc-b-twentyfirst-nonboundary-lower-correction,prop:bc-c-twentyfirst-nonboundary-lower-correction,prop:bc-b-twentysecond-nonboundary-lower-correction,prop:bc-c-twentysecond-nonboundary-lower-correction,prop:bc-b-twentythird-nonboundary-lower-correction,prop:bc-c-twentythird-nonboundary-lower-correction,prop:bc-b-twentyfourth-nonboundary-lower-correction,prop:bc-c-twentyfourth-nonboundary-lower-correction,prop:bc-b-twentyfifth-nonboundary-lower-correction,prop:bc-c-twentyfifth-nonboundary-lower-correction,prop:bc-b-twentysixth-nonboundary-lower-correction,prop:bc-c-twentysixth-nonboundary-lower-correction,prop:bc-b-twentyseventh-nonboundary-lower-correction,prop:bc-c-twentyseventh-nonboundary-lower-correction,prop:bc-b-twentyeighth-nonboundary-lower-correction,prop:bc-c-twentyeighth-nonboundary-lower-correction,prop:bc-b-twentyninth-high-nonboundary-lower-correction,prop:bc-c-twentyninth-nonboundary-lower-correction}.
For \(B_b\) rows, these propositions supply \(j=2b+2\), \(j=2b+3\),
\(j=2b+4\), \(j=2b+5\), \(j=2b+6\), \(j=2b+7\), \(j=2b+8\), and
\(j=2b+9\), \(j=2b+10\), \(j=2b+11\), \(j=2b+12\), \(j=2b+13\),
\(j=2b+14\), \(j=2b+15\), \(j=2b+16\), \(j=2b+17\), \(j=2b+18\), and
\(j=2b+19\), \(j=2b+20\), \(j=2b+21\), \(j=2b+22\), \(j=2b+23\),
\(j=2b+24\), \(j=2b+25\), \(j=2b+26\), \(j=2b+27\), \(j=2b+28\),
\(j=2b+29\), and \(j=2b+30\).  For \(42\le b\le123\),
\cref{prop:bc-b-twentyninth-high-nonboundary-lower-correction} supplies
\(j=2b+31\), while the hypothesis supplies every value \(j\ge2b+32\).  For
\(14\le b\le41\), the hypothesis supplies every value \(j\ge2b+31\).  For
\(C_b\) rows,
\cref{prop:bc-c-boundary-lower-correction,prop:bc-c-first-nonboundary-lower-correction,prop:bc-c-second-nonboundary-lower-correction,prop:bc-c-third-nonboundary-lower-correction,prop:bc-c-fourth-nonboundary-lower-correction,prop:bc-c-fifth-nonboundary-lower-correction,prop:bc-c-sixth-nonboundary-lower-correction,prop:bc-c-seventh-nonboundary-lower-correction,prop:bc-c-eighth-nonboundary-lower-correction,prop:bc-c-ninth-nonboundary-lower-correction,prop:bc-c-tenth-nonboundary-lower-correction,prop:bc-c-eleventh-nonboundary-lower-correction,prop:bc-c-twelfth-nonboundary-lower-correction,prop:bc-c-thirteenth-nonboundary-lower-correction,prop:bc-c-fourteenth-nonboundary-lower-correction,prop:bc-c-fifteenth-nonboundary-lower-correction,prop:bc-c-sixteenth-nonboundary-lower-correction,prop:bc-c-seventeenth-nonboundary-lower-correction,prop:bc-c-eighteenth-nonboundary-lower-correction,prop:bc-c-nineteenth-nonboundary-lower-correction},
together with
\cref{prop:bc-c-twentieth-nonboundary-lower-correction,prop:bc-c-twentyfirst-nonboundary-lower-correction,prop:bc-c-twentysecond-nonboundary-lower-correction,prop:bc-c-twentythird-nonboundary-lower-correction,prop:bc-c-twentyfourth-nonboundary-lower-correction,prop:bc-c-twentyfifth-nonboundary-lower-correction,prop:bc-c-twentysixth-nonboundary-lower-correction,prop:bc-c-twentyseventh-nonboundary-lower-correction,prop:bc-c-twentyeighth-nonboundary-lower-correction,prop:bc-c-twentyninth-nonboundary-lower-correction},
supplies \(j=2b+2\), \(j=2b+3\), \(j=2b+4\), \(j=2b+5\), \(j=2b+6\), and
\(j=2b+7\), \(j=2b+8\), \(j=2b+9\), \(j=2b+10\), \(j=2b+11\),
\(j=2b+12\), \(j=2b+13\), \(j=2b+14\), \(j=2b+15\), \(j=2b+16\),
\(j=2b+17\), \(j=2b+18\), \(j=2b+19\), \(j=2b+20\), \(j=2b+21\),
\(j=2b+22\), \(j=2b+23\), \(j=2b+24\), \(j=2b+25\), \(j=2b+26\),
\(j=2b+27\), \(j=2b+28\), \(j=2b+29\), \(j=2b+30\), and \(j=2b+31\), while
the hypothesis supplies every value \(j\ge2b+32\).  Thus
every displayed row has
\[
  -\frac12s_j\le\delta_j^G
\]
for every correction variable in the full interval
$2r_G+2\le j\le N_{\mathrm{dir}}^{BC}(G)+2$.

Expand the Chain difference in the stable moments and correction variables
as in the proof of \cref{prop:post29}.  By
\cref{prop:bc-pieri-correction-sign}, the $B/C$ corrections are
nonpositive.  Terms with negative linear coefficient are therefore
nonnegative and may be discarded, while a positive linear coefficient
$L_{m,j}$ contributes at worst $-\frac12 L_{m,j}s_j$.  Likewise positive
quadratic terms may be discarded and a negative quadratic coefficient
$Q_{m,i,j}$ contributes at worst $\frac14 Q_{m,i,j}s_is_j$.

For each Chain index $m$, the verifier takes the smallest active boundary
$2\rank(G)+2$ among rows whose finite bridge still includes $D_G(2m+1)$.
Using this frontier only adds possible negative terms, so it gives a common
lower bound for every active row.  The exact GMP/OpenMP replay
\[
\texttt{certificates/post\_m29/bc\_lambda\_half\_bridge\_gmp\_certificate.log}
\]
computes a still smaller rigorous bound, using the separable linear
majorant and generating-function reduction recorded in
\cref{prop:post29-bc-local-half-bridge}, for all $m=31,\ldots,15447$ with
$\lambda=1/2$.  It exits with status $0$, reports \texttt{failures=0}, and
has SHA-256
\[
\hashsplit{fa267f4eed4eda48405e7eb468c8b022}{f997ad0bee5ce6324bc8618affcb0e6a}.
\]
Its worst positive lower bound is attained at $m=15447$, odd exponent
$30897$, with frontier rank $217$ and log$_{10}$ lower bound
$125316.825861430509>0$.  Hence every active bridge Chain difference is
nonnegative under the displayed assumptions and the cited boundary and initial
nonboundary theorems.  For a fixed $m$, the two moment-formula diagonals in
$D_G(2m+1)$ involve moment indices only through $2m+5$.  The same coefficient
calculation may therefore be truncated at that index, proving the local
clause.
\end{proof}

\begin{remark}[Superseded half-stable route]\label{rem:former-post29-tail}
Before \cref{prop:post29-bc-layered-mgf}, the only unresolved route used the
following two-row condition.  It is recorded solely for provenance and is
not an assumption of any theorem in this paper.  For $b=14,15$, let
$N_{\mathrm{dir}}^{BC}(B_b)$ be the direct-tail onset of
\cref{prop:post29-bc-interval-direct}.  With $s_j=m_j^{\mathrm{st}}$ and
$m_j^{B_b}=s_j+\delta_j^{B_b}$, the former condition was
\[
 -\frac12s_j\le\delta_j^{B_b}
 \qquad
 \left(2b+31\le j\le N_{\mathrm{dir}}^{BC}(B_b)+2\right).
\]
The earlier values are theorems: \(j=2b+2\) for \(B_b\) rows by
\cref{prop:bc-b-boundary-lower-correction}, \(j=2b+3\) for \(B_b\) rows by
\cref{prop:bc-b-first-nonboundary-lower-correction}, \(j=2b+4\) for \(B_b\)
rows by \cref{prop:bc-b-second-nonboundary-lower-correction}, \(j=2b+5\) for
\(B_b\) rows by \cref{prop:bc-b-third-nonboundary-lower-correction},
\(j=2b+6\) for \(B_b\) rows by
\cref{prop:bc-b-fourth-nonboundary-lower-correction}, \(j=2b+7\) for
\(B_b\) rows by \cref{prop:bc-b-fifth-nonboundary-lower-correction}, and
\(j=2b+8\) for \(B_b\) rows by
\cref{prop:bc-b-sixth-nonboundary-lower-correction}, and \(j=2b+9\) for
\(B_b\) rows by \cref{prop:bc-b-seventh-nonboundary-lower-correction}, and
\(j=2b+10\) for \(B_b\) rows by
\cref{prop:bc-b-eighth-nonboundary-lower-correction}, and \(j=2b+11\) for
\(B_b\) rows by
\cref{prop:bc-b-ninth-nonboundary-lower-correction}, and \(j=2b+12\) for
\(B_b\) rows by
\cref{prop:bc-b-tenth-nonboundary-lower-correction}, and \(j=2b+13\) for
\(B_b\) rows by
\cref{prop:bc-b-eleventh-nonboundary-lower-correction}, and \(j=2b+14\) for
\(B_b\) rows by
\cref{prop:bc-b-twelfth-nonboundary-lower-correction}, and \(j=2b+15\) for
\(B_b\) rows by
\cref{prop:bc-b-thirteenth-nonboundary-lower-correction}, and \(j=2b+16\)
for \(B_b\) rows by
\cref{prop:bc-b-fourteenth-nonboundary-lower-correction}, and \(j=2b+17\)
for \(B_b\) rows by
\cref{prop:bc-b-fifteenth-nonboundary-lower-correction}, and \(j=2b+18\) for
\(B_b\) rows by
\cref{prop:bc-b-sixteenth-nonboundary-lower-correction}, and \(j=2b+19\) for
\(B_b\) rows by
\cref{prop:bc-b-seventeenth-nonboundary-lower-correction}, and \(j=2b+20\)
for \(B_b\) rows by
\cref{prop:bc-b-eighteenth-nonboundary-lower-correction}, and \(j=2b+21\)
for \(B_b\) rows by
\cref{prop:bc-b-nineteenth-nonboundary-lower-correction}, and \(j=2b+22\)
for \(B_b\) rows by
\cref{prop:bc-b-twentieth-nonboundary-lower-correction},
and \(j=2b+23\) for \(B_b\) rows by
\cref{prop:bc-b-twentyfirst-nonboundary-lower-correction},
and \(j=2b+24\) for \(B_b\) rows by
\cref{prop:bc-b-twentysecond-nonboundary-lower-correction},
and \(j=2b+25\) for \(B_b\) rows by
\cref{prop:bc-b-twentythird-nonboundary-lower-correction},
and \(j=2b+26\) for \(B_b\) rows by
\cref{prop:bc-b-twentyfourth-nonboundary-lower-correction},
and \(j=2b+27\) for \(B_b\) rows by
\cref{prop:bc-b-twentyfifth-nonboundary-lower-correction},
and \(j=2b+28\) for \(B_b\) rows by
\cref{prop:bc-b-twentysixth-nonboundary-lower-correction},
and \(j=2b+29\) for \(B_b\) rows by
\cref{prop:bc-b-twentyseventh-nonboundary-lower-correction},
and \(j=2b+30\) for \(B_b\) rows by
\cref{prop:bc-b-twentyeighth-nonboundary-lower-correction}.
\end{remark}

\fi

\begin{proposition}[Local $B/C$ half-stable bridge]
\label{prop:post29-bc-local-half-bridge}
Fix one of the rows
\[
  B_{14},\ldots,B_{123},\qquad C_{20},\ldots,C_{217},
\]
and an integer $m$ for which
\[
 N_{\mathrm{hit}}^{(29)}(G)\le 2m+1<N_{\mathrm{dir}}^{BC}(G).
\]
Put $s_j=m_j^{\mathrm{st}}$ and $m_j^G=s_j+\delta_j^G$.  If
\[
 -\frac12s_j\le\delta_j^G\le0
 \qquad(2\rank(G)+2\le j\le2m+5),
\]
then $D_G(2m+1)\ge0$.
\end{proposition}

\begin{proof}
Substitution in \cref{lem:moment-formula,def:chain-diff} gives
\[
 D_G(2m+1)=D_{\rm st}(2m+1)
 +\sum_jL_{m,j}\delta_j^G
 +\sum_{i,j}Q_{m,i,j}\delta_i^G\delta_j^G .
\]
The two moment-formula diagonals use indices only through $2m+5$.
Under the displayed box constraints, discarding sign-favourable terms gives
the exact sign-truncated lower bound
\[
 \mathcal L_m=D_{\rm st}(2m+1)
 -\frac12\sum_{L_{m,j}>0}L_{m,j}s_j
 +\frac14\sum_{Q_{m,i,j}<0}Q_{m,i,j}s_is_j .
\]
The replay uses the still smaller bound \(\widetilde{\mathcal L}_m\) obtained
from \([x+y]_+\le[x]_++[y]_+\) on the two linear diagonals; thus
\(D_G(2m+1)\ge\mathcal L_m\ge\widetilde{\mathcal L}_m\).

Here is the analytic reduction used by the source.  Put \(S=n+2\) and
\[
 A_n(k)=\binom n{k-2}+\binom nk-2\binom n{k-1}
 =\frac{n!}{k!(S-k)!}\bigl((2k-S)^2-S\bigr),
 \qquad
 q_n=\sum_{k=0}^{S}A_n(k)s_ks_{S-k}.
\]
If \(M(x)=\sum_{j\ge0}s_jx^j/j!\), the stable recurrence gives
\[
 M(x)=\frac{\exp(-x/2+x^2/4)}{\sqrt{1-x}},\qquad
 \sum_{n\ge0}q_n\frac{x^n}{n!}
 =2\bigl(M''M-(M')^2\bigr)
 =e^{-x+x^2/2}\left((1-x)^{-3}+(1-x)^{-1}\right).
\]
Thus all stable diagonal values are generated by a fixed finite recurrence;
equivalently, for \(Q(x)=\sum q_nx^n/n!\),
\[
 (-x^3+3x^2-4x+2)Q'
 =(-x^4+4x^3-6x^2+4x+2)Q,
 \qquad q_0=q_1=2.
\]
Moreover, \(A_n(k)<0\) only when \(|2k-S|<\sqrt S\).  Consequently the
negative parts require only \(O(\sqrt n)\) central terms.  The positive
active parts are recovered from \(q_n\) by subtracting that central sum and
the excluded boundary sum \(k<f_m\), where \(f_m\le436\) is the active
frontier.  Finally, since
\[
 L_{m,k}s_k=2A_{2m+3}(k)s_ks_{2m+5-k}
             -8A_{2m+1}(k)s_ks_{2m+3-k},
\]
the inequality \([x+y]_+\le[x]_++[y]_+\) supplies the asserted separable
linear majorant, while the negative quadratic part is retained exactly.
This replaces the former full diagonal traversal by
\(O(\sqrt m+f_m)\) exact terms per index.
The GMP/OpenMP source
\path{character_ring_iter/post_m29_bc_interval_bridge_frontier_gmp.cpp}
checks $\widetilde{\mathcal L}_m>0$ for every $m=31,\ldots,15447$.  Its accepted replay
reports no failures and has SHA-256
\[
\hashsplit{fa267f4eed4eda48405e7eb468c8b022}{f997ad0bee5ce6324bc8618affcb0e6a}.
\]
This range contains every bridge index in the displayed rows, and proves the
claim.
\end{proof}

\begin{proposition}[Active residual $B/C$ correction-prefix contract]
\label{prop:bc-active-correction-prefix-contract}
Put $q=b+1$, write $s_j=m_j^{\mathrm{st}}$, and write
$m_j^G=s_j+\delta_j^G$.  The following bounds hold:
\[
 -\frac12s_{2q+r}\le\delta_{2q+r}^{B_b}\le0
 \quad(14\le b\le123, 0\le r\le27),
 \tag{\ref{prop:bc-active-correction-prefix-contract}.1}
\]
and
\[
 -\frac12s_{2q+r}\le\delta_{2q+r}^{C_b}\le0
 \quad(20\le b\le217, 0\le r\le27).
 \tag{\ref{prop:bc-active-correction-prefix-contract}.2}
\]
The lower bound at each offset is supplied by the exact result in the same
row of the following table; the common upper bound is
\cref{prop:bc-pieri-correction-sign}.

\begingroup\scriptsize
\setlength{\LTleft}{0pt}
\setlength{\LTright}{0pt}
\begin{longtable}{cL{0.43\linewidth}L{0.43\linewidth}}
\toprule
$r$ & type $B$ lower-bound result & type $C$ lower-bound result\\
\midrule
\endfirsthead
\toprule
$r$ & type $B$ lower-bound result & type $C$ lower-bound result\\
\midrule
\endhead
0&\contractref{prop:bc-b-boundary-lower-correction}&\contractref{prop:bc-c-boundary-lower-correction}\\
1&\contractref{prop:bc-b-first-nonboundary-lower-correction}&\contractref{prop:bc-c-first-nonboundary-lower-correction}\\
2&\contractref{prop:bc-b-second-nonboundary-lower-correction}&\contractref{prop:bc-c-second-nonboundary-lower-correction}\\
3&\contractref{prop:bc-b-third-nonboundary-lower-correction}&\contractref{prop:bc-c-third-nonboundary-lower-correction}\\
4&\contractref{prop:bc-b-fourth-nonboundary-lower-correction}&\contractref{prop:bc-c-fourth-nonboundary-lower-correction}\\
5&\contractref{prop:bc-b-fifth-nonboundary-lower-correction}&\contractref{prop:bc-c-fifth-nonboundary-lower-correction}\\
6&\contractref{prop:bc-b-sixth-nonboundary-lower-correction}&\contractref{prop:bc-c-sixth-nonboundary-lower-correction}\\
7&\contractref{prop:bc-b-seventh-nonboundary-lower-correction}&\contractref{prop:bc-c-seventh-nonboundary-lower-correction}\\
8&\contractref{prop:bc-b-eighth-nonboundary-lower-correction}&\contractref{prop:bc-c-eighth-nonboundary-lower-correction}\\
9&\contractref{prop:bc-b-ninth-nonboundary-lower-correction}&\contractref{prop:bc-c-ninth-nonboundary-lower-correction}\\
10&\contractref{prop:bc-b-tenth-nonboundary-lower-correction}&\contractref{prop:bc-c-tenth-nonboundary-lower-correction}\\
11&\contractref{prop:bc-b-eleventh-nonboundary-lower-correction}&\contractref{prop:bc-c-eleventh-nonboundary-lower-correction}\\
12&\contractref{prop:bc-b-twelfth-nonboundary-lower-correction}&\contractref{prop:bc-c-twelfth-nonboundary-lower-correction}\\
13&\contractref{prop:bc-b-thirteenth-nonboundary-lower-correction}&\contractref{prop:bc-c-thirteenth-nonboundary-lower-correction}\\
14&\contractref{prop:bc-b-fourteenth-nonboundary-lower-correction}&\contractref{prop:bc-c-fourteenth-nonboundary-lower-correction}\\
15&\contractref{prop:bc-b-fifteenth-nonboundary-lower-correction}&\contractref{prop:bc-c-fifteenth-nonboundary-lower-correction}\\
16&\contractref{prop:bc-b-sixteenth-nonboundary-lower-correction}&\contractref{prop:bc-c-sixteenth-nonboundary-lower-correction}\\
17&\contractref{prop:bc-b-seventeenth-nonboundary-lower-correction}&\contractref{prop:bc-c-seventeenth-nonboundary-lower-correction}\\
18&\contractref{prop:bc-b-eighteenth-nonboundary-lower-correction}&\contractref{prop:bc-c-eighteenth-nonboundary-lower-correction}\\
19&\contractref{prop:bc-b-nineteenth-nonboundary-lower-correction}&\contractref{prop:bc-c-nineteenth-nonboundary-lower-correction}\\
20&\contractref{prop:bc-b-twentieth-nonboundary-lower-correction}&\contractref{prop:bc-c-twentieth-nonboundary-lower-correction}\\
21&\contractref{prop:bc-b-twentyfirst-nonboundary-lower-correction}&\contractref{prop:bc-c-twentyfirst-nonboundary-lower-correction}\\
22&\contractref{prop:bc-b-twentysecond-nonboundary-lower-correction}&\contractref{prop:bc-c-twentysecond-nonboundary-lower-correction}\\
23&\contractref{prop:bc-b-twentythird-nonboundary-lower-correction}&\contractref{prop:bc-c-twentythird-nonboundary-lower-correction}\\
24&\contractref{prop:bc-b-twentyfourth-nonboundary-lower-correction}&\contractref{prop:bc-c-twentyfourth-nonboundary-lower-correction}\\
25&\contractref{prop:bc-b-twentyfifth-nonboundary-lower-correction}&\contractref{prop:bc-c-twentyfifth-nonboundary-lower-correction}\\
26&\contractref{prop:bc-b-twentysixth-nonboundary-lower-correction}&\contractref{prop:bc-c-twentysixth-nonboundary-lower-correction}\\
27&\contractref{prop:bc-b-twentyseventh-nonboundary-lower-correction}&\contractref{prop:bc-c-twentyseventh-nonboundary-lower-correction}\\
\bottomrule
\end{longtable}
\endgroup
The additional offset-$28$ and offset-$29$ results in the source archive are
not needed by the main-theorem propagation below.
\end{proposition}

\begin{proof}
For each table row, the stated lower-bound result has exactly the family,
rank interval, and moment index printed in
\eqref{prop:bc-active-correction-prefix-contract}.1 or
\eqref{prop:bc-active-correction-prefix-contract}.2.
\ifginibrefullproof
Those results and their proofs occur earlier in this detailed build.
\else
Their exact supplement labels are printed in the table, so the import
boundary does not depend on an unindexed collection of certificates.
\fi
The Pieri omitted-summand formula in
\cref{prop:bc-pieri-correction-sign} gives $\delta_j^G\le0$ for both
families.  Combining the two sides proves the contract.
\end{proof}

\begin{proposition}[Residual $B/C$ post-$m=29$ closure]\label{prop:post29-bc-residual-closure}
For
\[
  G=B_b\quad(14\le b\le123),
  \qquad\text{or}\qquad
  G=C_b\quad(20\le b\le217),
\]
one has
\[
  \Qthree^G(n)\ge0
\]
for every odd integer
$n\ge N_{\mathrm{hit}}^{(29)}(G)$.
\end{proposition}

\begin{proof}
Write $m_j^G=s_j+\delta_j^G$, where $s_j$ is the stable moment.  The complete
active correction import is \cref{prop:bc-active-correction-prefix-contract};
in particular, its table replaces the former prose reference to a suite of
unnamed contracts.

We record why this prefix is sufficient.  The analytic suppliers used below
start no later than the odd exponent
\[
 N_{\rm tail}(b)=2b+29.
\]
To propagate to an odd target $t<N_{\rm tail}(b)$, the last Chain difference
used is $D_G(t-2)$.  By \cref{lem:moment-formula,def:chain-diff}, that
difference involves moments only through $m_{t+2}^G$.  Since
$t\le2b+27$, its largest possible correction index is
\[
 t+2\le2b+29=2q+27.
\]
Thus the offsets $0\le r\le27$ in
\cref{prop:bc-active-correction-prefix-contract} are exactly the correction
prefix consumed by the main theorem.  The offset-$28$ and offset-$29$
results and the longer conditional development routes retained in the source
archive are not premises of this proposition.

Apply the local clause of the supplement's half-stable bridge reduction
\contractref{prop:post29-bc-local-half-bridge}.  To make the imported finite
predicate explicit, write the exact moment-formula expansion as
\[
 D_G(2m+1)=D_{\rm st}(2m+1)
 +\sum_j L_{m,j}\delta_j^G
 +\sum_{i,j}Q_{m,i,j}\delta_i^G\delta_j^G,
\]
where the integer coefficients are obtained by substituting
$m_j^G=s_j+\delta_j^G$ into
\cref{lem:moment-formula,def:chain-diff}.  Under
$-s_j/2\le\delta_j^G\le0$, sign truncation gives the exact lower bound
\(\mathcal L_m\) displayed in
\contractref{prop:post29-bc-local-half-bridge}.  The separable linear
majorant proved there then gives the certified bound
\[
 \widetilde{\mathcal L}_m\le\mathcal L_m\le D_G(2m+1).
\]
Its generating-function recurrence and central-band sign formula reduce the
exact evaluation to \(O(\sqrt m+f_m)\) terms.

\noindent The exact GMP/OpenMP source
\begin{center}
\small\path{character_ring_iter/post_m29_bc_interval_bridge_frontier_gmp.cpp}
\end{center}
checks $\widetilde{\mathcal L}_m>0$ for $m=31,\ldots,15447$, reports no failures, and
has replay SHA-256
\[
\hashsplit{fa267f4eed4eda48405e7eb468c8b022}{f997ad0bee5ce6324bc8618affcb0e6a}.
\]
Together with the first-hit proposition \cref{prop:post29} and Chain
propagation, this reaches the applicable analytic supplier.  The layered-MGF
proposition \cref{prop:post29-bc-layered-mgf} supplies
$B_{14},\ldots,B_{61}$ and $C_{20},\ldots,C_{61}$; the tilted-MGF
propositions \cref{prop:post29-b-tilted-mgf,prop:post29-c-tilted-mgf} and the
direct interval proposition \cref{prop:post29-bc-interval-direct} supply
$B_{62},\ldots,B_{123}$ and $C_{62},\ldots,C_{217}$.  For
$B_{14},\ldots,B_{17}$ the layered-MGF onset
already precedes $N_{\mathrm{hit}}^{(29)}(G)$.  These ranges are disjoint and
exhaust the displayed rows, proving every stated odd exponent.  No inactive
archival route is used.
\end{proof}

\begin{theorem}[Classical $B,C,D$ closure]
\label{thm:classical}
For every compact connected simple group whose Lie algebra has type $B_b$,
$C_b$, or $D_b$, and for every $n\ge0$,
\[
  \Qthree^G(n)\ge0.
\]
\end{theorem}

\begin{proof}
\ifginibrefullproof
Even $n$ is \cref{lem:even}.  For odd $n$, the moment formula and
\cref{lem:adjoint-moment-positive} give the strict base value
$\Qthree^G(1)=2m_3^G>0$ in each adjoint classical row.  The stable edge
\cref{lem:classical-stable-edge} and the row-gated bridge
\cref{prop:classical-bridge} supply the needed nonnegative Chain
differences before the post-bridge hit, so \cref{lem:chain-propagation}
gives the pre-tail odd exponents.  The high-rank rows are covered by
\cref{prop:classical-high}.  The post-bridge finite rows listed in
\cref{prop:first-hit} are covered by the trace-tail criterion.  The low-rank
post-$m=29$ rows in \cref{prop:post29-direct} are covered by their direct
pushforward-tail certificates and \cref{prop:d-low-exact-bridge}, with
$D_{19}$ also supplied by \cref{prop:post29-d19-exact}.  The rows
$D_{25},\ldots,D_{52}$ are covered unconditionally by
\cref{prop:post29-d25-52-exact}.  Every row
$D_{53},\ldots,D_{295}$ is covered unconditionally by
\cref{prop:post29-d-a-free}; $D_{297}$ is covered by
\cref{prop:first-hit}, and all other higher type-$D$ rows are covered by
\cref{prop:classical-high}.  The high-edge
rows $B_b,C_b$ for $218\le b\le275$ or $b=277$ are covered by
\cref{prop:post29-chernoff}.  The rows $B_{124},\ldots,B_{217}$ are covered by
\cref{prop:post29-b-unitary-square}.  The remaining rows
$B_{14},\ldots,B_{123}$ and $C_{20},\ldots,C_{217}$ are exactly
\cref{prop:post29-bc-residual-closure}.  This proposition imports the minimal
offset-$0$--$27$ correction contract, the local half-stable bridge, and the
layered/tilted MGF or direct interval supplier.  No surplus correction offset
or development-archive route is used.  These cases exhaust every non-$A$
classical row and every odd exponent.
\else
Even $n$ is \cref{lem:even}.  For odd $n$, the moment formula and
\cref{lem:adjoint-moment-positive} give the strict base value
$\Qthree^G(1)=2m_3^G>0$ in each adjoint classical row.  The stable edge
\cref{lem:classical-stable-edge} and row-gated bridge
\cref{prop:classical-bridge}, followed by \cref{lem:chain-propagation},
supply all pre-tail odd exponents.  The high-rank rows are
\cref{prop:classical-high}, and the post-bridge finite rows listed in
\cref{prop:first-hit} follow from the trace-tail criterion.

It remains only to partition the post-$m=29$ rows.  The low-rank rows are
\cref{prop:post29-direct}; the type-$D$ rows $D_{25},\ldots,D_{52}$ are
\cref{prop:post29-d25-52-exact}, the rows $D_{53},\ldots,D_{295}$ are
\cref{prop:post29-d-a-free}, while $D_{297}$ is \cref{prop:first-hit} and
the remaining higher type-$D$ rows are \cref{prop:classical-high}.  The
high-edge $B/C$ rows are \cref{prop:post29-chernoff}.  The residual ranges
$B_{14},\ldots,B_{123}$ and $C_{20},\ldots,C_{217}$ are
\cref{prop:post29-bc-residual-closure}.  Finally,
$B_{124},\ldots,B_{217}$ are \cref{prop:post29-b-unitary-square}.
These alternatives exhaust every non-$A$ classical row and every odd
exponent.
\fi
\end{proof}

%======================================================================
\section{Exceptional Lie Groups}
%======================================================================

The exceptional proof uses the exact Chain difference of
\cref{def:chain-diff}.  For every exceptional adjoint family, the moment
formula and \cref{lem:adjoint-moment-positive} give
$\Qthree^G(1)=2m_3^G>0$.  Nonnegative Chain differences therefore propagate
positivity by \cref{lem:chain-propagation}.

\begin{table}[ht]
\centering
\small
\begin{tabular}{ccl}
\toprule
Group & exact prefix & tail/bridge closure \\
\midrule
$\Gtwo$ & Chain $m=0,\ldots,97$ & direct rectangular tail from odd $n\ge21$ \\
$\Ffour$ & Chain $m=0,\ldots,107$ & direct rectangular tail from odd $n\ge65$ \\
$\Esix$ & Chain $m=0,\ldots,37$ & direct Chain tail from odd $n\ge75$ \\
$\Eseven$ & Chain $m=0,\ldots,32$ & direct Chain tail from odd $n\ge65$ \\
$\Eeight$ & Chain $m=0,\ldots,32$ & tail from odd $n\ge133$ plus bridges to $n=77$ \\
\bottomrule
\end{tabular}
\caption{Exceptional exact prefixes and tail certificates.}
\label{tab:exceptional}
\end{table}

\begin{lemma}[Exceptional adjoint endpoint and Chain negative region]
\label{lem:exceptional-negative-region}
Let $G$ be a compact exceptional simple group of rank $r$, and put
$C_G=r$.  Then
\[
  \chi_{\rm ad}(g)\ge-C_G\qquad(g\in G).
\]
For odd $n$, with the pushforward measure $\psi_G$ of
\cref{lem:pushforward-tail},
\[
 D_G(n)=\int_{\mathbb R}u^n(u^2-4)\,d\psi_G(u).
\]
Consequently, if $L(n)$ is any lower bound for a positive part of this
integral, then
\[
 D_G(n)\ge L(n)-U_G(n),\qquad
 U_G(n)=2\bigl((2C_G)^2+4\bigr)(2C_G)^n.
 \tag{\ref{lem:exceptional-negative-region}.1}
\]
\end{lemma}

\begin{proof}
Garibaldi--Guralnick--Rains
\cite[Theorem~2.1]{GaribaldiGuralnickRains2025} prove
$\operatorname{Tr}\operatorname{Ad}(g)\ge-\operatorname{rank}G$ for every
element of every compact Lie group.  The displayed identity follows from
\cref{def:chain-diff,lem:pushforward-tail}.  For odd $n$, its integrand is
negative only on
\[
 [-2C_G,-2]\cup[0,2].
\]
On this union its absolute value is at most
$((2C_G)^2+4)(2C_G)^n$.  Since
$\psi_G(\mathbb R)=2$ by \cref{lem:pushforward-tail}, integration gives
\eqref{lem:exceptional-negative-region}.1.
\end{proof}

\begin{lemma}[Normalized exceptional Weyl--rectangle coefficient]
\label{lem:exceptional-weyl-normalization}
Normalize long roots to have squared length two.  Let $Q^\vee\subset
P^\vee$ be the coroot and coweight lattices, put
\[
 z_R=\prod_{\alpha\in R_+}\frac{|\alpha|^2}{2},\qquad
 \nu_R=\frac{|P^\vee/Q^\vee|^2}
 {\operatorname{covol}(Q^\vee)^2},
\]
and let $d_i$ be the invariant degrees.  For the exceptional root systems,
\[
\begin{array}{c|ccccc}
R&G_2&F_4&E_6&E_7&E_8\\
\hline
z_R&1/27&1/4096&1&1&1\\
\operatorname{covol}(Q^\vee)^2&3&4&3&2&1\\
|P^\vee/Q^\vee|&1&1&3&2&1\\
\nu_R&1/3&1/4&3&2&1\\
\min_{0\ne\lambda\in P^\vee}|\lambda|^2&2&2&4/3&3/2&2
\end{array}
\tag{\ref{lem:exceptional-weyl-normalization}.1}
\]
Let $a=d/2$, where $d=\dim G$, and let
\[
 A_R=\nu_Rz_R^2
 \frac{8ad^2(d/\kappa)^d\prod_i((d_i-1)!)^2}{(2\pi)^r}.
 \tag{\ref{lem:exceptional-weyl-normalization}.2}
\]
For a retained radial cell
$[s_0,s_1]\times[t_0,t_1]$ and its transpose, suppose the local character
bounds are $g_R,h_R$ and the product of the certified Weyl-sine and optional
radial exponential losses is $\eta_R(n)>0$.  Then their combined positive
contribution to $D_G(n)$ is at least
\[
 \frac{A_R(2d)^n}{n^{d+2}}
 \frac{g_R(n)^2\eta_R(n)}{a\Gamma(a)^2}
 \left(\int_{s_0}^{s_1}s^{a-1}\,ds\right)
 \left(\int_{t_0}^{t_1}t^{a-1}\,dt\right)
 h_R(n)^n\bigl((2dh_R(n))^2-4\bigr).
 \tag{\ref{lem:exceptional-weyl-normalization}.3}
\]
All exceptional caps used below lie in an injective exponential chart of the
adjoint maximal torus.
\end{lemma}

\begin{proof}
For the adjoint global form the exponential kernel is $2\pi P^\vee$.
Thus normalized torus Haar measure in Euclidean exponential coordinates is
\[
 \frac{dv}{(2\pi)^r\operatorname{covol}(P^\vee)},\qquad
 \operatorname{covol}(P^\vee)
 =\frac{\operatorname{covol}(Q^\vee)}{|P^\vee/Q^\vee|}.
 \tag{\ref{lem:exceptional-weyl-normalization}.4}
\]
Apply the Weyl integration formula
\cite[Chapter~IV, (1.11)]{BrockertomDieck} in both variables.  If
$\Delta_R(v)=\prod_{\alpha>0}\alpha(v)$ and $\Delta_{\rm std}$ uses the
squared-length-two normal to every reflecting hyperplane, then
$\Delta_R^2=z_R\Delta_{\rm std}^2$.  The Macdonald--Mehta identity
\cite[\S2.1 and Theorem~3.1(i)]{Etingof2009}, scaled by
$(d/\kappa)^{1/2}$, gives
\[
 \int e^{-\kappa|v|^2/(2d)}\Delta_R(v)^2\,dv
 =z_R(2\pi)^{r/2}(d/\kappa)^a\prod_i d_i!.
\]
The square of this identity, the two factors in
\eqref{lem:exceptional-weyl-normalization}.4, and
$|W|=\prod_i d_i$ give
\[
 \nu_Rz_R^2\frac{(d/\kappa)^d
       \prod_i((d_i-1)!)^2}{(2\pi)^r}.
\]
In the scaled radial variable $s=n\kappa|v|^2/(2d)$, homogeneity of
$\Delta_R^2$ gives the radial density
$s^{a-1}/\Gamma(a)$.  The character-gap square contributes
$(2d/n)^2g_R^2$; adjoining the transposed cell contributes the factor two.
Writing $8d^2$ as $(8ad^2)/a$ yields exactly
\eqref{lem:exceptional-weyl-normalization}.3.  The retained cells have
$s_0\ge t_1$, so the cells and their transposes are disjoint up to null
boundaries.

The Cartan and simple-coroot Gram determinants give the middle three rows of
\eqref{lem:exceptional-weyl-normalization}.1; the short-root products give
the first row.  For $E_6,E_7$, inversion of the Cartan matrix and exact
enumeration of the simple-root pairings in $\{-1,0,1\}$ give the displayed
coweight minima.  This enumeration is global: if one pairing has absolute
value at least two, Cauchy's inequality and $|\alpha_i|^2=2$ give
$|\lambda|^2\ge2$.  The same calculation gives the $G_2,F_4$ minima; for
$E_8$, determinant one and the even integral Cartan form give minimum two.
These determinant and minimum calculations are repeated in the exact C++
checkers.

Finally, a cap point satisfies $|v|^2<\delta^2/2$.  For $G_2,F_4$ one has
$\delta\le\pi$; for $E_6,E_7,E_8$ one has $\delta\le14/3$.  The last row of
\eqref{lem:exceptional-weyl-normalization}.1 and the rational inequality
$\pi>333/106$ show that the diameter of every cap is strictly smaller than
$2\pi\min_{0\ne\lambda\in P^\vee}|\lambda|$.  Hence no two cap points
differ by a nonzero exponential-kernel vector.
\end{proof}

\begin{lemma}[Rectangular Chain certificate template]\label{lem:rect-template}
Let $G$ be exceptional, let $d=\dim G$, and write
$D_G(n)=\Qthree^G(n+2)-4\Qthree^G(n)$.  Suppose a finite union of rational
rectangles $R=[s_0,s_1]\times[t_0,t_1]$ in radial variables satisfies, for
all odd $n\ge n_0$,
\[
  \chi(y)-\chi(x)\ge \frac{2d}{n}g_R,\qquad
  \frac{\chi(x)+\chi(y)}{2d}\ge h_R(n),
\]
with $g_R>0$, $h_R(n)>0$, and with every rectangle contained in the
specified root-angle cap $|\alpha\cdot v|\le\delta$ on which the stated
scalar trigonometric minorants have been proved.  Let $L(n)>0$ be the
resulting sum of integrated rectangular lower bounds, and suppose an explicit
negative-region estimate $U(n)>0$ satisfies, for every odd $n\ge n_0$,
\[
  D_G(n)\ge L(n)-U(n),\qquad L(n_0)\ge U(n_0),\qquad
  L(n+2)U(n)\ge L(n)U(n+2)
  \tag{\ref{lem:rect-template}.1}
\]
Then $D_G(n)\ge0$ for every such $n$.
\end{lemma}

\begin{proof}
The two local bounds supply the squared gap and positive-power factors; the
cap makes the stated scalar bounds uniform.  The middle
comparison gives $L(n_0)/U(n_0)\ge1$, and the last makes this ratio
nondecreasing along odd $n$.  Hence $L(n)\ge U(n)$ and $D_G(n)\ge0$.
\end{proof}

\begin{proposition}[Direct rectangular tails for $E_6$ and $E_7$]\label{prop:e6e7-rect}
The following exact rectangular certificates hold:
\[
\begin{array}{c|c|c|c|c|c}
G & n_0 & \text{grid}&\text{retained cells}&
 \log_{10}(L(n_0)/U(n_0))&\text{odd-step lower bound}\\
\hline
E_6&75&1/10&23088&8.904&>10^{0.573}\\
E_7&65&1/20&66925&>0.42&>101/100
\end{array}
\]
Consequently $D_{E_6}(n)\ge0$ for every odd $n\ge75$, and
$D_{E_7}(n)\ge0$ for every odd $n\ge65$.
\end{proposition}

\begin{proof}
\ifginibrefullproof
Put $\delta=14/3$.  The root constants are
\[
\begin{array}{c|ccccc}
G&d&r&|R_+|&\kappa:\sum_{\alpha>0}(\alpha\cdot u)^2
 =\kappa|u|^2&q:\sum_{\alpha>0}(\alpha\cdot u)^4=Q(u)^2/q\\
\hline
E_6&78&6&36&12&16\\
E_7&133&7&63&18&27.
\end{array}
\]
The checker generates the $36$ and $63$ positive roots from the two Cartan
matrices by exact simple reflections, expands both degree-four polynomials
coefficientwise in the simple-root coordinates, and verifies
$16\sum(\alpha\cdot u)^4=Q(u)^2$ and
$27\sum(\alpha\cdot u)^4=Q(u)^2$, respectively.

On $|z|\le\delta$, exact Bernstein coefficients applied to alternating
Taylor minorants prove
\[
\begin{aligned}
2(1-\cos z)&\ge a_1z^2+a_2z^4+a_3z^6,\\
2\cos z&\ge2-z^2+cz^4,\\
\frac{\sin(z/2)}{z/2}&\ge
 \exp\{-z^2/24-z^4/2880-z^6/103800\},
\end{aligned}
\tag{\ref{prop:e6e7-rect}.1}
\]
where
\[
(a_1,a_2,a_3)=
\left(\frac{99119225}{10^8},-\frac{7802858}{10^8},
\frac{169078}{10^8}\right),\qquad c=\frac{83}{2000}.
\]
Since $\sum z^6\ge(\sum z^4)^2/\sum z^2$ and
$\sum z^6\le\delta^2\sum z^4$, these scalar inequalities give, on a cell
$R=[s_0,s_1]\times[t_0,t_1]$,
\[
\begin{aligned}
g_R(n)={}&a_1s_0-t_1
 +a_2\frac{2d}{q}\frac{s_1^2}{n}
 +a_3\frac{(2d)^2}{q^2}\frac{s_0^3}{n^2}
 +c\frac{2d}{q}\frac{t_0^2}{n},\\
h_R(n)={}&1-\frac{s_1+t_1}{n}
 +c\frac{2d}{q}\frac{s_0^2+t_0^2}{n^2}.
\end{aligned}
\tag{\ref{prop:e6e7-rect}.2}
\]
The sine-density loss is bounded using
\[
2\frac{2d}{24n}(s_1+t_1)
+2\left(\frac1{2880}+\frac{\delta^2}{103800}\right)
 \frac{(2d)^2}{q n^2}(s_1^2+t_1^2).
\tag{\ref{prop:e6e7-rect}.3}
\]

For $E_6$ the checker partitions
$[10,30]\times[5,20]$ into $1/10$-cells.  The cap value at $n=75$ is
$2450/39>30$.  It retains exactly $23088$ cells.  It uses
\cref{lem:exceptional-negative-region} with $C_{E_6}=6$, hence
$U_{E_6}(n)=296\cdot12^n$.

For $E_7$, let
\[
 P(S)=\frac{T_{26}((2S-273)/259)}{T_{26}(-301/259)},\qquad
 I=\int S^{14}(S^2+4)P(S)^2\,d\psi_{E_7}(S).
\]
Exact Bernstein coefficients on $[-14,-2]$ and $[0,2]$ prove at $n=65$
\[
 |S^{65}(S^2-4)|\le
 \frac{24}{25}14^{51}S^{14}(S^2+4)P(S)^2.
\tag{\ref{prop:e6e7-rect}.4}
\]
The integral $I$ is an exact linear combination of
$\Qthree^{E_7}(0),\ldots,\Qthree^{E_7}(68)$ and therefore uses only the
moments $m_0,\ldots,m_{70}$ regenerated by
\cref{lem:racah-speiser-step,lem:character-pairing}.  Thus
\[
 U_{E_7}(n)=\frac{24}{25}14^{n-14}I
\]
is a valid negative-region bound for every later odd $n$: multiplication by
$14^2$ preserves \eqref{prop:e6e7-rect}.4 because $|S|\le14$ on both
negative-sign intervals.  The checker partitions
$[15,35]\times[8,28]$ into $1/20$-cells; the cap is $910/19>35$, and
exactly $66925$ cells survive.

For every retained cell the checker verifies the two inequalities in
\cref{lem:rect-template}, positivity of the direct factor, and the cleared
monotonicity inequalities for $g_R(n)$ and $h_R(n)$.  It also verifies
cellwise
\[
 \left(\frac d{C_G}h_R(n_0)\right)^2
 \left(\frac{n_0}{n_0+2}\right)^\alpha>1
\]
(and $>101/100$ for $E_7$), while the cap and sine factors improve.
The positive sums use the normalized Weyl factors
$\nu_{E_6}=3$ and $\nu_{E_7}=2$ from
\cref{lem:exceptional-weyl-normalization}; the checker verifies the Cartan
determinants, coweight minima, and disjointness from every transposed cell
before summing.
OpenMP parallelizes the cell sum, but every retained-cell value and both
final comparisons use GMP rationals.  The resulting base-ratio logarithms
are those in the table, so \cref{lem:rect-template} proves both tails.
\else
Put $\delta=14/3$.  Exact reflection generation gives
\[
(d,r,|R_+|,\kappa,q)=(78,6,36,12,16)\quad(E_6),
\qquad
(133,7,63,18,27)\quad(E_7),
\]
where $Q=\kappa|u|^2$ and
$\sum_{\alpha>0}(\alpha\cdot u)^4=Q(u)^2/q$; the checker verifies both
polynomial identities coefficientwise.  It also proves the rational
wide-cap cosine and Weyl-sine minorants by exact Bernstein coefficients;
with Cauchy's radial sixth-power bounds these give exact $g_R(n),h_R(n)$
satisfying \cref{lem:rect-template} on every retained cell.

For $E_6$ use $U(n)=296\cdot12^n$.  For $E_7$ put
\[
 P(S)=\frac{T_{26}((2S-273)/259)}{T_{26}(-301/259)},\quad
 I=\int S^{14}(S^2+4)P(S)^2\,d\psi_{E_7}(S),\quad
 U(n)=\frac{24}{25}14^{n-14}I.
\]
Exact Bernstein coefficients prove the associated pointwise majorant on
$[-14,-2]\cup[0,2]$ at $n=65$; it propagates since $|S|\le14$, and $I$
uses only $m_0,\ldots,m_{70}$.  The OpenMP/GMP checker verifies the cell
counts, the normalized Weyl factors $3,2$, the coweight minima,
transposed-cell disjointness, all cleared monotonicity inequalities, and the
stated base and odd-step bounds.  Hence \eqref{lem:rect-template}.1 holds for
both tails.
\fi
\end{proof}

\begin{proposition}[$E_8$ rectangular bridge and tail ledger]\label{prop:e8-rect-ledger}
For $E_8$, the following exact rectangular certificates prove
$D_{E_8}(n)\ge0$ at the listed odd exponents:
\ifginibrefullproof
\[
\begin{array}{c|c|c}
n & \text{retained cells} & \log_{10}(\text{base ratio})\\
\hline
77&57817&0.19\\
79&68095&0.79\\
81&79060&0.27\\
83&91138&0.28\\
85&103495&0.29\\
87&116526&0.44\\
89&130740&0.47\\
91&69502&0.13\\
93&53822&0.04\\
95&29784&0.70\\
97&35869&0.27\\
99&42662&1.70\\
101&49478&0.68\\
103&56387&0.68\\
105&63660&0.19\\
107&70562&0.08\\
109&77454&1.15\\
111&84702&0.39\\
113&91581&1.89\\
115&98455&0.52\\
117&111427&0.05\\
119&15000&0.66\\
121&74053&0.10\\
123&28121&1.10\\
125&8497&1.11\\
127&10528&1.04\\
129&12719&0.35\\
131&17797&1.59
\end{array}
\]
\else
\[
\begin{array}{c|r|c@{\qquad}c|r|c}
n&\text{cells}&\log_{10}\text{ ratio}&n&\text{cells}&\log_{10}\text{ ratio}\\
\hline
77&57817&.19&105&63660&.19\\
79&68095&.79&107&70562&.08\\
81&79060&.27&109&77454&1.15\\
83&91138&.28&111&84702&.39\\
85&103495&.29&113&91581&1.89\\
87&116526&.44&115&98455&.52\\
89&130740&.47&117&111427&.05\\
91&69502&.13&119&15000&.66\\
93&53822&.04&121&74053&.10\\
95&29784&.70&123&28121&1.10\\
97&35869&.27&125&8497&1.11\\
99&42662&1.70&127&10528&1.04\\
101&49478&.68&129&12719&.35\\
103&56387&.68&131&17797&1.59
\end{array}
\]
\fi
The analytic tail from $n\ge133$ uses $13558$ retained cells, base-ratio
logarithm $0.02$, and odd-step logarithm $0.85$.
\end{proposition}

\begin{proof}
The $E_8$ constants are $d=248$, $|R_+|=120$, rank $8$, and $C_G=8$.
For the analytic tail the negative-region estimate is
$U(n)=2((2C_G)^2+4)(2C_G)^n=520\cdot16^n$.
The certificates use the quartic Weyl identity
\[
  \sum_{\alpha>0}(\alpha\cdot u)^4=\frac{1}{50}
  \left(\sum_{\alpha>0}(\alpha\cdot u)^2\right)^2.
\]
On each retained rational cell the checker proves the two local character
bounds of \cref{lem:rect-template}; all retained cells are inside the
root-angle cap.  It also verifies the factor-one normalized Weyl coefficient
and transposed-cell disjointness required by
\cref{lem:exceptional-weyl-normalization}.  For the tail certificate at
$n=133$, for instance, the
cell formula is
\[
  g_R=s_0-t_1-\frac{248}{300n}s_1^2,\qquad
  h_R(n)=1-\frac{s_1+t_1}{n}
       +\frac{1}{18}\frac{496}{50}\frac{s_0^2+t_0^2}{n^2},
\]
and the retained cells have
\[
  \min h_R(133)=\frac{176378749}{318402000},\qquad
  \min(496^2h_R(133)^2-4)>0.
\]
\ifginibrefullproof
Before the table is consumed,
\texttt{verify\_e8\_negative\_bounds\_gmp.cpp} checks its exact $29$-row
schema.  For $n=77,79,\ldots,113$, it reads the locally verified and ancillary
moment chain through $m_{100}$, reconstructs the rational square-majorant from
the pinned factor witness, proves both pointwise inequalities by exact
Bernstein coefficients, and integrates the witness with GMP.  Only $n=81$
and $n=85$ require subdivision, of depths three and five respectively.  For
$n=115,117,\ldots,133$, the audit recomputes $520\cdot16^n$ from
\cref{lem:exceptional-negative-region}.  Every resulting rational must equal
the corresponding table entry; no historical Python row module is active.

Only then does the OpenMP/GMP replay
\texttt{direct\_chain\_rect\_e8\_parallel\_certificate.cpp} consume that
table.  It checks every retained-cell count, the $n=133$ tail row, and the two
exact rational comparisons for each case.  The base comparisons and tail
odd-step comparison are exact rational checks.  For completeness, on the
\else
An exact source audit first rederives all $29$ negative-bound rows from their
moment sources and pointwise majorants and requires exact equality with the
manifest-pinned table.  The OpenMP/GMP replay then consumes it and checks all
cell counts, the factor-one normalized $E_8$ Weyl coefficient of
\cref{lem:exceptional-weyl-normalization}, transposed-cell disjointness, and
rational base comparisons.  For the tail propagation, on the
\fi
fixed $n=133$ tail cells one has $s_1<40$, $t_1\le31$, and
$s_1+t_1>63$.  With $A_R=s_1+t_1$ and
$B_R=(124/225)(s_0^2+t_0^2)$, one has
$133A_R>8379>2824>2B_R$, so $h_R(n)=1-A_R/n+B_R/n^2$ increases thereafter.
The gap, cap, sine-density, and direct factors also improve.  Hence every
cell has odd-step quotient at least
$((248/8)\min_R h_R(133))^2(133/135)^{250}>1$, the logged exact value.
Thus \eqref{lem:rect-template}.1 holds throughout the tail, proving the
listed differences and every odd $n\ge133$.
\end{proof}

\begin{proposition}[$G_2$ and $F_4$]\label{prop:g2-f4}
For $G=\Gtwo$ or $G=\Ffour$ and every $n\ge0$,
\[
  \Qthree^G(n)\ge0.
\]
\end{proposition}

\begin{proof}
The base value $\Qthree^G(1)>0$ is
\cref{lem:adjoint-moment-positive,lem:moment-formula}.
Put $Q(u)=\sum_{\alpha>0}(\alpha\cdot u)^2$.  Exact expansion of the positive
roots gives
\[
\begin{array}{c|cccccc}
G&d&r&|R_+|&\kappa:Q=\kappa|u|^2&q:
 \sum_{\alpha>0}(\alpha\cdot u)^4=Q(u)^2/q&(d_i)\\
\hline
G_2&14&2&6&4&16/5&(2,6)\\
F_4&52&4&24&9&54/5&(2,6,8,12).
\end{array}
\]
The Macdonald--Mehta identity
\cite[Theorem~3.1(i)]{Etingof2009}, whose root normals have squared length
$2$, gives the positive-saddle normalization
\[
 A_0=\nu_Gz_G^2\frac{8ad^2(d/\kappa)^d
       \prod_i((d_i-1)!)^2}{(2\pi)^r},\qquad a=d/2.
\]
Here the actual short roots change the squared discriminant by
$z_{G_2}=1/27$ and $z_{F_4}=1/4096$; these are respectively the products
of $|\alpha|^2/2$ over the three and twelve positive short roots.
The normalized-torus factors from
\cref{lem:exceptional-weyl-normalization} are
$\nu_{G_2}=1/3$ and $\nu_{F_4}=1/4$.

Apply \cref{lem:rect-template} with $C=2,4$ and with
\[
\begin{array}{c|c|c|c|c}
G&n_0&[s_0,s_1]\times[t_0,t_1]&g&h(n_0)\\
\hline
G_2&21&[19/5,21/5]\times[7/5,9/5]&111/80&1661/2268\\
F_4&65&[51/4,53/4]\times[31/4,33/4]&3023/1296&44063/63180.
\end{array}
\]
Writing $q$ for the preceding quartic denominator, the exact local bounds are
\[
 g=s_0-t_1-\frac{d s_1^2}{6qn_0},\qquad
 h(n)=1-\frac{s_1+t_1}{n}
 +\frac{2d(s_0^2+t_0^2)}{18q n^2}.
\]
The normalized rectangular radial mass is bounded below by
\[
 g^2\,3^{-\lceil s_1+t_1\rceil}
 \frac{(s_1^a-s_0^a)(t_1^a-t_0^a)}{a^3((a-1)!)^2}.
\]
Using $\pi<355/113$, the Weyl-sine lower factors are $3^{-1}$ and $3^{-4}$.
The resulting exact base comparisons with
$U(n)=2((2C)^2+4)(2C)^n$ exceed $1$ and $10^{13}$, respectively.  Moreover
the odd-step quotients are bounded below by
\[
 \left(\frac dC h(n_0)\right)^2
 \left(\frac{n_0}{n_0+2}\right)^{d+2}>6,10,
\]
respectively.  The exact checks also give
$n_0(s_1+t_1)>2d(s_0^2+t_0^2)/(9q)$, so $h(n)$ increases; the gap and cap
improve with $n$.  They also verify the coroot Gram and Cartan determinants,
the coweight minima, and transposed-cell disjointness required by
\cref{lem:exceptional-weyl-normalization}.  Thus the rectangular Chain tails
begin at odd $n=21,65$.
The exact prefixes reach odd $n=197,217$ and overlap both tails.
\Cref{lem:chain-propagation} closes all odd exponents; even exponents are
\cref{lem:even}.
\end{proof}

\begin{proposition}[$E_6$ and $E_7$]\label{prop:e6-e7}
For $G=\Esix$ or $G=\Eseven$ and every $n\ge0$,
\[
  \Qthree^G(n)\ge0.
\]
\end{proposition}

\begin{proof}
The base value $\Qthree^G(1)>0$ is
\cref{lem:adjoint-moment-positive,lem:moment-formula}.
The exact prefixes verify the Chain differences through $m=37$ for $E_6$
and $m=32$ for $E_7$, reaching the odd exponents $75$ and $65$,
respectively.  The direct rectangular certificates of
\cref{prop:e6e7-rect} prove the corresponding Chain tails from those same
odd exponents onward.  The prefix endpoint and the tail start overlap
exactly, so \cref{lem:chain-propagation} gives $\Qthree^G(n)\ge0$ for
every odd $n$.  Even $n$ is \cref{lem:even}.
\end{proof}

\begin{proposition}[$E_8$ finite prefix and bridge]\label{prop:e8-bridge}
For $E_8$, the finite prefix Chain differences
\[
  D_{E_8}(n)\quad(n=1,3,\ldots,65)
\]
are positive.  The finite bridge values
\[
  \Qthree^{E_8}(n)\quad (n=69,71,73,75,77)
\]
are positive, and the Chain differences
\[
  D_{E_8}(n)\quad (n=67,69,71,73,75,77,79,81,83,85,87,89,91,93,95)
\]
are nonnegative.
\end{proposition}

\begin{proof}
The exact source-pairing replay applies
\cref{lem:racah-speiser-step,lem:character-pairing} through
$\ad^{\otimes50}$ and independently regenerates
$m_0,\ldots,m_{100}$ from the $E_8$ Cartan datum.  It checks every value
against the archived local prefix and the Bourjaily--Plesser--Vergu
comparison table, with zero mismatches.  For
odd $n\le98$, \cref{lem:moment-formula} computes $\Qthree^{E_8}(n)$
using only these moments.  Exact integer evaluation gives
\[
\begin{array}{ll}
\min D_{E_8}(1,3,\ldots,65)                   & =20,\\
\log_{10}\min \Qthree^{E_8}(69,71,73,75,77) & \approx100.94,\\
\log_{10}\min D_{E_8}(67,69,\ldots,77)      & \approx100.94,\\
\log_{10}\min D_{E_8}(67,69,\ldots,95)      & \approx100.94.
\end{array}
\]
The replay flags are
\[
\begin{array}{ll}
\texttt{moment\_log\_0\_100\_arxiv\_extension} & \texttt{OK},\\
\texttt{prefix\_chain\_diff\_1\_65\_positive} & \texttt{OK},\\
\texttt{bridge\_q3\_69\_77\_positive} & \texttt{OK},\\
\texttt{chain\_diff\_67\_77\_positive} & \texttt{OK},\\
\texttt{extended\_chain\_diff\_67\_95\_positive} & \texttt{OK}.
\end{array}
\]
\end{proof}

\begin{proposition}[$E_8$ tail]\label{prop:e8-tail}
For $E_8$, the Chain difference is nonnegative for every odd
$n\ge77$.
\end{proposition}

\begin{proof}
The direct rectangular certificates in \cref{prop:e8-rect-ledger} prove
the Chain tail for odd $n\ge133$ and the successive bridge steps
\[
  n=131,129,127,125,123,121,119,117,115,113,111,109,107,105,103,101,
\]
\[
  n=99,97,95,93,91,89,87,85,83,81,79,77.
\]
Together with \cref{prop:e8-bridge}, these bridge steps connect the exact
prefix to the $n=133$ tail.  Therefore $D_{E_8}(n)\ge0$ for every odd
$n\ge77$.
\end{proof}

\begin{theorem}[Exceptional closure]\label{thm:exceptional}
For every $G\in\{\Gtwo,\Ffour,\Esix,\Eseven,\Eeight\}$ and every
$n\ge0$,
\[
  \Qthree^G(n)\ge0.
\]
\end{theorem}

\begin{proof}
Use \cref{prop:g2-f4} and \cref{prop:e6-e7} for
$G_2,F_4,E_6,E_7$.  For $E_8$, \cref{prop:e8-bridge} supplies the prefix
Chain differences through $n=65$, the bridge values at
$69,71,73,75,77$, and the bridge Chain differences through $95$; then
\cref{prop:e8-tail} supplies every later Chain difference.  Applying
\cref{lem:chain-propagation} from the strict base value
$\Qthree^{E_8}(1)>0$ of
\cref{lem:adjoint-moment-positive,lem:moment-formula} gives every odd
exponent.  Even exponents are
covered by \cref{lem:even}.
\end{proof}

%======================================================================
\section{Assembly of the Classification}
%======================================================================

\begin{proof}[Proof of \cref{thm:main}]
By \cref{lem:global-form}, $\Qthree^G(n)$ depends only on the compact
simple Lie algebra.  In the irreducible-root-system classification,
$B_1=C_1=A_1$ and $D_3=A_3$ are already in the type-$A$ branch,
$D_2$ is not simple, and the remaining classical types have
$B_b,C_b$ with $b\ge2$ and $D_b$ with $b\ge4$.  The classification
therefore leaves type $A$, the non-$A$ classical types $B,C,D$, and the exceptional types
$\Gtwo,\Ffour,\Esix,\Eseven,\Eeight$.  Type $A$ is
\cref{thm:type-a}; types $B,C,D$ are \cref{thm:classical}; the exceptional
types are \cref{thm:exceptional}.  These cases exhaust all compact
connected simple Lie groups.
\end{proof}

\section{Conclusion and limitations}

The proof establishes one uniform sign inequality for the adjoint tensor
category of every compact simple type.  Its common structure is the passage
from the moment identity in \cref{lem:moment-formula} to a stable prefix, a
finite-rank correction interval, and an analytic tail.  The classification is
needed because the mechanism controlling the middle interval differs between
type~$A$, the three other classical families, and the exceptional groups.

The theorem does not compare unequal characters and does not prove Ginibre's
condition on the cone of all central positive-definite functions.  It also
does not assert that the exactly-two-minus argument extends directly to four
or more minus factors.  Part~III addresses the latter hierarchy only on the
closed multiplicative convex cone generated by $\chiad$, using the present
theorem as an input.  These scope boundaries separate the mathematical
conclusion from possible applications to larger character cones or to
model-specific correlation inequalities.

\begin{remark}[Final-facing dependency map]\label{rem:final-dependency-map}
The proof graph is the following; it deliberately omits superseded diagnostic
routes retained in the expanded archive.  All incoming nodes are proved in
the manuscript or discharged by the authenticated exact replays.
\begin{center}
\begin{tabular}{p{0.16\linewidth}|p{0.73\linewidth}}
closure & load-bearing incoming results\\
\hline
\cref{thm:type-a} &
\cref{thm:su-stable-rank,thm:su2-row,thm:su-low-rank-ratios,lem:su-transition-strip,thm:su-ratio,cor:su-finite}\\
\cref{thm:classical} &
\cref{lem:classical-stable-edge,prop:classical-high,prop:classical-bridge,prop:first-hit,prop:post29-direct,prop:post29-chernoff,prop:post29-d-a-free,prop:post29-b-unitary-square,prop:bc-active-correction-prefix-contract},
\contractref{prop:post29-bc-local-half-bridge},
\cref{prop:post29-bc-residual-closure}\\
\cref{thm:exceptional} &
\cref{prop:g2-f4,prop:e6-e7,prop:e8-bridge,prop:e8-tail}\\
\cref{thm:main} &
\cref{lem:global-form,thm:type-a,thm:classical,thm:exceptional}
\end{tabular}
\end{center}
Inside \cref{prop:post29-bc-residual-closure}, the only live residual route is
the proved correction-prefix contract and local half-stable bridge, followed
by the layered/tilted MGF or direct interval supplier.  Local fixed-witness
lemmas used to prove the
boundary-through-twenty-seventh correction contract are represented in the
accepted correction-prefix suite.  The expanded archive supplies a
line-by-line audit of those contracts.  Global all-row
fixed-witness/compression interfaces explicitly marked superseded have no
incoming edge from this map.
\end{remark}

%======================================================================
\appendix
\ifginibrefullproof
\section{Exact Certificate Protocol}
%======================================================================

All finite certificates used above are exact.  Moment computations use either
the hook-length partition formula in type $A$ or the exact
Racah--Speiser iteration of \cref{lem:racah-speiser-step} together with the
pairing identities of \cref{lem:character-pairing}.  Signs are decided only after
reducing the target inequality to an integer or rational comparison.  The
replay outputs quoted in the paper are therefore certificate outputs, not
floating-point tests.

The common checker contract is the following.
\[
\begin{array}{ll}
\text{Input} &
\text{root datum, highest weights, exact moment or delta logs, and row list},\\
\text{Moment step} &
\text{compute }m_k=\mult(\mathbf 1,\ad^{\otimes k})
  \text{ by integer arithmetic},\\
\text{Target step} &
\text{form }\Qthree(n)\text{ or }D_G(n)
  \text{ from \cref{lem:moment-formula} and \cref{def:chain-diff}},\\
\text{Reduction step} &
\text{rewrite every tail or bridge claim as an integer/rational inequality},\\
\text{Decision step} &
\text{accept only if every required integer/rational margin is } >0.
\end{array}
\]
For every classical or exceptional source table passed to an active
root-datum source-replay stage, the checker
reconstructs the relevant Cartan matrix and its full root set, verifies the
adjoint root multiplicities and rank-dimensional zero-weight space, rejects
negative output multiplicities and coordinate overflow, regenerates the
moments through half degree, and compares every claimed moment or full
correction.  In
particular, the $E_8$ bridge regenerates $m_0,\ldots,m_{100}$ from
$\ad^{\otimes0},\ldots,\ad^{\otimes50}$.  For each classical source check,
the checker independently regenerates the stable sequence from
\cref{lem:bc-stable-moment-recurrence}, verifies every archived OEIS value
through the largest required degree, and uses the regenerated value in the
finite correction.  Thus the OEIS and arXiv tables are independent comparisons
rather than accepted source axioms.

The finite row-gated bridge is now source-generated.  The standard command
runs the Pieri/CRT verifier of
\cref{prop:bc-source-row-gated-bridge} for all $870$ $B/C$ steps and the
exact-GMP interval mode of \cref{prop:d-a-free-row-gated-bridge} for the
remaining $36$ high-rank type-$D$ steps.  Together with the active
$D_4$--$D_{52}$ replays, these computations replace the historical
row-gated ledger completely.  Its old raw logs remain authenticated
provenance, but no theorem consumes their classifications, statuses, or
printed margins.

For type $D$, a bare historical line
\texttt{D\_b Delta\_j = B\_j} with $j>2b$ records only the nonnegative
determinant component $B_j$, not the full correction $A_j+B_j$.  Both the
source checker and the downstream interval checker reject such a line as an
exact moment.  Bare lines with $j\le2b$, where $A_j=0$, and explicit full or
exact-Weyl moment lines are regenerated and compared exactly.

\begin{longtable}{@{}L{0.18\textwidth}L{0.48\textwidth}@{}}
\toprule
Certificate & Output used in the proof \\
\midrule
\endhead
Type $A$ finite-overlap replay &
The endpoint constants, rows $N=6,\ldots,18$, and the final finite-overlap
band are checked; the final line is
\texttt{All SU(N) repair certificates verified.} \\
\addlinespace
Classical Rains-square cutoff replay &
The cutoff inequalities are checked for $C_G=296,\ldots,10000$, the
pre-cutoff bad range is $248,\ldots,295$, and the replay exits with
status $0$. \\
\addlinespace
Classical first-hit trace replay &
The replay tests 881 finite rows, closes exactly 57 first-hit rows, and
records the minimum margin at $D_{281}$. \\
\addlinespace
Source-generated classical row-gated bridge &
The common Pieri traversal regenerates $832$ exact $B/C$ corrections and
checks all $870$ row-gated inequalities: $416$ are direct exact evaluations
and $454$ use proved one-sided correction intervals.  It reports zero
failures and minimum lower margin $1046$.  The separate type-$D$ mode checks
all $36$ nonvacuous $D_{53}$--$D_{63}$ steps with zero failures.  Both
programs are rebuilt and executed in the isolated replay; the historical raw
bridge archive is provenance only. \\
\addlinespace
Classical post-$m=29$ first-hit replay &
The replay checks the first post-$m=29$ Chain step in 796 rows, generates
stable moments through $m_{557}$, and exits with status $0$.  The five bucket
row counts are $14,26,490,25,241$.  The replay is used only for the first-hit
claim of \cref{prop:post29}; the all-later $B/C$ closure is the separate
correction-prefix and analytic-supplier argument of
\cref{prop:post29-bc-residual-closure}.  Its input delta-log manifest and
artifact classification are
recorded in \texttt{certificates/post\_m29}.  The input, accepted, diagnostic,
and classification manifest hashes are, respectively,
\hashsplit{f8ca9e149f896c28b22ad5c53ba4bb8d}{e877aca153cc56d6f04f21c04c25c95e},
\hashsplit{ce34ba6dd902f8058923219f5e3b9b29}{faf66a541e875af4782b854cc9fb0540},
\hashsplit{01f053f2d50173900e3b7d2d773629b6}{ec6b3a46bad5fd1eed9cdb4fb50f640f},
and
\hashsplit{31e37c889291d5ab2504472163f89d0e}{2ada74f621eea2315291350f09733455}. \\
\addlinespace
Classical $B_{11}$ post-$m=29$ bridge &
The replay checks 28 copied exact $B_{11}$ boundary-delta logs
$\Delta_{24},\ldots,\Delta_{51}$, verifies that each records
\texttt{\_\_EXIT\_STATUS=0}, recomputes Chain margins for
$m=31,\ldots,36$, and exits with status $0$.  The minimum margin is attained
at $m=31$ and is positive.  The delta-log manifest hash is
\hashsplit{c689442193378fd04a2e3db6cbfa9b315}{3839d88b37f06055ca13e71ac9aa129}. \\
\addlinespace
Classical $B_{12}$ post-$m=29$ finite tail &
The replay checks one existing exact $B_{12}$ log for $\Delta_{26},\ldots,
\Delta_{39}$ and 14 copied exact $B_{12}$ boundary-delta logs
$\Delta_{40},\ldots,\Delta_{53}$, verifies that each records
\texttt{\_\_EXIT\_STATUS=0}, recomputes Chain margins for
$m=31,\ldots,38$, and checks the direct Chebyshev tail from odd exponent
$81$ onward.  The minimum bridge margin is attained at $m=31$ and is
positive.  The delta-log manifest hash is
\hashsplit{480b77eb57b215cded54e4c195a15e}{61c119894aae01f3dc04c58eb2d22f2e09}. \\
\addlinespace
Classical $B_{13}$ post-$m=29$ finite tail &
The replay checks one existing exact $B_{13}$ log for $\Delta_{28},\ldots,
\Delta_{39}$ and 18 copied exact $B_{13}$ boundary-delta logs
$\Delta_{40},\ldots,\Delta_{57}$, verifies that each records
\texttt{\_\_EXIT\_STATUS=0}, recomputes Chain margins for
$m=31,\ldots,43$, checks stable-envelope bounds for the missing negative
quadratic pairs at $m=42,43$, and checks the direct Chebyshev tail from odd
exponent $91$ onward.  The delta-log manifest and replay log hashes are
\hashsplit{0fe4e21a1533ba290a7bfad3c07609bf}{ff8d70a0d7ad6d46976faf60d7ed3b37}
and
\hashsplit{08b9666bc771897e9702c0132620f4ca}{e9215b30c92c160ed52795fe6db096ee}. \\
\addlinespace
Classical $C_{13}$ post-$m=29$ finite tail &
The replay checks one existing exact $C_{13}$ log for $\Delta_{28},\ldots,
\Delta_{39}$ and four copied exact $C_{13}$ boundary-delta logs
$\Delta_{40},\ldots,\Delta_{43}$, verifies that each records
\texttt{\_\_EXIT\_STATUS=0}, recomputes Chain margins for
$m=31,32,33$, and checks the direct Chebyshev tail from odd exponent
$71$ onward.  The minimum bridge margin is attained at $m=31$ and is
positive.  The delta-log manifest hash is
\hashsplit{e2c407483c5f0894d4ceee5b887c2e22}{aa6e5604b85bde037ace25923c2f0320}. \\
\addlinespace
Classical $C_{14}$ post-$m=29$ finite tail &
The replay checks one existing exact $C_{14}$ log for $\Delta_{30},\ldots,
\Delta_{39}$ and 12 copied exact $C_{14}$ boundary-delta logs
$\Delta_{40},\ldots,\Delta_{51}$, verifies that each records
\texttt{\_\_EXIT\_STATUS=0}, recomputes Chain margins for
$m=31,\ldots,39$, and checks the direct Chebyshev tail from odd exponent
$83$ onward.  The minimum bridge margin is attained at $m=31$ and is
positive.  The delta-log manifest hash is
\hashsplit{aa13f907722ad5e1bd384c243ac042ce}{1fddd429782c641c246c146acabf4756};
the replay log has SHA-256
\hashsplit{6dc9a3f519c64568b378b8689dd78867}{e5311abdf1257268f37ba79408feb1c3}. \\
\addlinespace
Classical $C_{15}$ post-$m=29$ direct tail &
The replay checks the existing exact $C_{15}$ log for
$\Delta_{32},\ldots,\Delta_{39}$, verifies that it records
\texttt{\_\_EXIT\_STATUS=0}, and checks the direct Chebyshev tail from odd
exponent $63$ onward.  The replay log has SHA-256
\hashsplit{68fd43b912006354f77fb2c051475767}{f8c6d85b3355b2a3ba3c438cd7a2ad68}. \\
\addlinespace
Classical $C_{16}$ post-$m=29$ direct tail &
The replay checks the existing exact $C_{16}$ log for
$\Delta_{34},\ldots,\Delta_{39}$, verifies that it records
\texttt{\_\_EXIT\_STATUS=0}, and checks the direct Chebyshev tail from odd
exponent $63$ onward.  The replay log has SHA-256
\hashsplit{9e15444c6e789d3aba9f83feaea03a22}{f5dd423e9c2a138b68f959bdb22c719f}. \\
\addlinespace
Classical $C_{17}$ post-$m=29$ finite tail &
The replay checks exact $C_{17}$ boundary-delta logs
$\Delta_{36},\ldots,\Delta_{45}$, verifies that each records
\texttt{\_\_EXIT\_STATUS=0}, recomputes Chain margins for
$m=31,\ldots,34$, and checks the direct Chebyshev tail from odd exponent
$73$ onward.  The delta-log manifest and replay log hashes are
\hashsplit{47b4939674358a15b1c66ecb8d3c33a8}{f684a33a53a0c3ef25c1e7004749ccff}
and
\hashsplit{dcd5eb8e782be93d926a6486f18d8a5e}{c445d33f8784ba251f60dd0203d8eb32}. \\
\addlinespace
Classical $C_{18}$ post-$m=29$ finite tail &
The replay checks exact $C_{18}$ boundary-delta logs
$\Delta_{38},\ldots,\Delta_{50}$, verifies that each records
\texttt{\_\_EXIT\_STATUS=0}, recomputes Chain margins for
$m=31,\ldots,40$, and checks the direct Chebyshev tail from odd exponent
$85$ onward.  The delta-log manifest and replay log hashes are
\hashsplit{5e61b65a8c745b131b16839fab178eb31}{73b29c90b3dcd922942ea64fa26a001}
and
\hashsplit{e25bfaef2b3b6289631431fd26717704}{dc82c78009cc02cf9c8e8dc8861b5541}. \\
\addlinespace
Classical $C_{19}$ post-$m=29$ finite tail &
The replay checks exact $C_{19}$ boundary-delta logs
$\Delta_{40},\ldots,\Delta_{58}$, verifies that each records
\texttt{\_\_EXIT\_STATUS=0}, recomputes Chain margins for
$m=31,\ldots,47$, and checks the direct Chebyshev tail from odd exponent
$99$ onward.  The delta-log manifest and replay log hashes are
\hashsplit{5567ba76b394f5579832d7fce8991e71}{70d98cef1f3cf4cbb43a57bd06450448}
and
\hashsplit{034d55810933481b342966686cf249a33}{649fac5333ad335df9b94e8f95b8273}. \\
\addlinespace
Superseded determinant-only $D_{12}$--$D_{97}$ tail routes &
These historical replays treated the determinant-isotypic contribution as the
full finite-rank correction after the omitted-even-row term became nonzero.
They are diagnostic artifacts and no numbered result consumes them.  The
accepted replacements are the full-signed-moment $D_{12}$--$D_{24}$ onset-$63$
certificate, the two-sided $D_{25}$--$D_{52}$ bridge and exact/mixed MGF
certificate, and the $A$-free $D_{53}$--$D_{295}$ frontier. \\
\addlinespace
Classical high-edge no-TV Chernoff trace &
The directed-rounding C++/MPFR replay checks the one-trace source range, the Chernoff admissibility
window, and the trace-pushforward inequality for $B_b,C_b$ with
$218\le b\le275$ or $b=277$, using the three checked B/C parameter sets.  It
exits with status $0$ and has SHA-256
\hashsplit{ca3a938c9d45d65eb892f6ba637ed8c3}{a7b1d32a571a5be52e2f5efb80477176}. \\
\addlinespace
Classical middle type-$B$ unitary square trace &
The $384$-bit directed-rounding MPFR replay checks the four strict logarithmic
comparisons in \cref{prop:post29-b-unitary-square} for all $94$ rows
$B_{124},\ldots,B_{217}$.  It reports \texttt{rows\_checked=94},
\texttt{failures=0}, and minimum log margin greater than $126$.  Its SHA-256
is
\hashsplit{7780bfe3253bf9e030a9a544dcd7327f}{cd905aed8a2890ded191dc38efae8230}. \\
\addlinespace
Classical middle type-$B/C$ exponentially tilted MGFs &
The $384$-bit directed-rounding MPFR replay checks the CJL source domains,
positive tilted interval masses, and strict comparisons in
\cref{prop:post29-b-tilted-mgf,prop:post29-c-tilted-mgf} for all $62$ rows
$B_{62},\ldots,B_{123}$ and all $156$ rows
$C_{62},\ldots,C_{217}$.  The type-$C$ part additionally checks both
orthogonal-component source domains from the Rains square decomposition.
It reports \texttt{rows\_checked=218}, \texttt{failures=0}, and has SHA-256
\hashsplit{ca6860204466e4a6b38d9c5597511cd1}{1cfebb4b1b65cab88db6ee45f9477111}. \\
\addlinespace
Classical residual layered $B/C/D$ MGFs &
The $384$-bit directed-rounding MPFR replay checks exact low-rank
Toeplitz--Hankel determinants, Bessel-series remainder bounds, positive
interval-Cholesky pivots, the extended CJL trace-norm domains, eight nested
tilted-mass bounds, exact stable push moments, one-sided polynomial
expectations, and the final strict comparison in
\cref{prop:post29-bc-layered-mgf,prop:post29-d-a-free}.  It covers $48$
type-$B$, $42$ type-$C$, and $243$ type-$D$ rows, reports
\texttt{failures=0}, and has minimum natural-log margin
$0.0078804615513483267\ldots$ at $D_{220}$.  Its SHA-256 is
\hashsplit{3e4900e5e2de2ec1a7fb04830ac59dea}{7369f8eebb1932ec08ac9ef73ba11654}. \\
\addlinespace
Classical $D_{25}$--$D_{52}$ exact/mixed finite-rank MGFs &
The $384$-bit directed-rounding verifier checks $28$ exact/mixed MGF onsets,
including exact normalized SO-even defining determinants at ranks $25$--$28$
and exact Rains square-component determinants at every rank.  It reports
\texttt{failures=0}, minimum final log margin
$0.013751090874926124\ldots$, and has SHA-256
\hashsplit{f87ab4bce0b189878ef4c3528ce34687}{1721bff05cb00ccb553a6c4e84e1f049}. \\
\addlinespace
Classical $D_{25}$--$D_{52}$ correction-aware bridge &
The exact C++/GMP/OpenMP verifier checks all $707$ finite Chain steps using
the proved two-sided correction boxes.  It reports \texttt{failures=0},
minimum exact lower bound $4695986280539928492451072$, and has SHA-256
\hashsplit{e47196c8097fb2a087b4673d47a5e34c}{aae1cd1a316058759e17b13708e2d888}. \\
\addlinespace
Classical $D_{53}$--$D_{295}$ $A$-free frontier &
The exact GMP/OpenMP verifier checks every Chain step $m=30,\ldots,163$,
requires every direct onset to satisfy $N_D(b)+2\le2b$, reports minimum
$A$-free slack $1$ and \texttt{failures=0}, and has SHA-256
\hashsplit{876d3f752d03b8f35183f0b28af7c15}{e9e28d104396fc65c215e5d854d8ae31f}. \\
\addlinespace
Type-$D$ legacy-envelope counterexample &
The exact machine-C C++/GMP replay computes
$B_{18}(4)-A_{18}(4)=-80781084484224$ and has SHA-256
\hashsplit{463e3c9632bb92936620fa65c03d3ff5}{abdb86940624f1621ffa976156d7d02f}. \\
\addlinespace
Classical $B/C$ interval direct tail &
The exact GMP negative-tail verifier checks the degree-$8$, scale-$12000$,
cutoff-$80$ Chebyshev input for all $402$ residual $B/C$ rows and reports
\texttt{failures=0}.  The directed-rounding MPFR verifier then checks the
direct pushforward-tail inequality at the logged row-dependent onset for
each residual row and reports \texttt{failures=0}.  The two replay SHA-256
hashes are
\hashsplit{80fce36378ec2192dca2c0b616c083b46168cf}{5134b66e73775156b995ea303b}
and
\hashsplit{53530cfe3d4b1947a77bf45c5c5aaa56}{a052ec1cde93c0704b68e004f22f2cb5}. \\
\addlinespace
\ifginibrefullproof
Classical $B/C$ correction-prefix bridge and superseded tail routes &
The exact GMP/OpenMP verifier checks the residual $B/C$ finite bridge under
the Pieri sign theorem and the lower bound $\delta_j\ge-s_j/2$.
The expanded correction propositions prove the prefix hypotheses used in
\cref{prop:post29-bc-residual-closure}.
It covers
$m=31,\ldots,15447$, reports \texttt{failures=0}, and its worst positive
lower bound occurs at $B_{217}$, odd exponent $30897$.  The replay SHA-256
hash is
\hashsplit{fa267f4eed4eda48405e7eb468c8b022}{f997ad0bee5ce6324bc8618affcb0e6a}.
\iffalse
The separate Pieri-compression arithmetic verifier checks that the candidate
bound
$|\delta_j|\le\binom{j}{\rank(G)+1}s_{j-\rank(G)-1}$ would imply
$\delta_j\ge-s_j/2$ throughout the residual $B/C$ windows.  It checks
$3{,}904{,}626$ exact GMP inequalities, reports \texttt{failures=0}, and has
SHA-256
\hashsplit{af7d338f58854db62c6f17ca23975e1}{75aee5e2206ef7c2386ead79875007558}.
This is only an arithmetic implication.  The source audit of the possible
Pieri-compression supplier reads Rains' \S3 tableaux/involution criteria,
Okada's Theorems 5.3--5.4 and Corollary 5.5, Huh--Kim--Krattenthaler--Okada's
Theorems 7.15 and 8.7 with Definitions 8.5--8.6, and Wenzl/Ram--Wenzl's
permissible Brauer quotient and matrix-unit statements.  These primary sources
identify the relevant bounded tableaux, vacillating tableaux, and permissible
path models, but they do not state the residual-rank compression inequality;
the displayed compression bound would have to be new proof content.
The controlled-loss arithmetic verifier checks the corresponding implication
for the weaker candidate
$|\delta_j|\le 2^{\rank(G)+1}\binom{j}{\rank(G)+1}s_{j-\rank(G)-1}$ on the
same residual windows.  It checks the same $3{,}904{,}626$ inequalities,
reports \texttt{failures=0}, and has SHA-256
\hashsplit{a98540c3b38d65a7c4148983022144112}{663801ab59fe2cb8deec1b34f6c5f7e}.
The repaired fixed-witness arithmetic verifier checks the larger-loss
implication used in \cref{prop:bc-bad-shape-reduction}:
\[
  2^{2q+1}\binom jq\,s_{j-q}\le s_j
\]
for all residual $B/C$ rows and residual moments, with
$q=\rank(G)+1$.  It checks the same $3{,}904{,}626$ inequalities, reports
\texttt{failures=0}, and reports its largest logarithmic-ratio diagnostic as
$-2.63554171301028006$ at $B_{14}$, moment $30$, and has SHA-256
\hashsplit{124068c5daa1123088a3cd79c3f6b839}{fee45a1aaef637ed82161c465d86bf45}.
The exact fixed-witness obstruction verifier checks the first residual height
fiber $q=15$, $j=30$, $S=\{1,3,\ldots,29\}$ and reports
\[
  \#\mathcal F^h(S)=6190283353629375
  >
  4825650970886144=2^{15}s_{15}.
\]
It reports \texttt{failures=0}, has SHA-256
\hashsplit{cc303f3bffd77ba4d8954054c31bc99}{ccd5c406f08e3a397ca0f37b759398436},
and is used only for the obstruction in
\cref{prop:bc-fixed-witness-qbit-obstruction}: it rules out the smaller
$2^q$ fixed-height-fiber target but does not discharge the remaining
combinatorial $2^{2q}$ fixed-witness supplier.
The exact complement-frontier verifier computes the maximal deleted-only
complement size absorbed by the existing $2q$ branch word for every residual
rank parameter $15\le q\le218$, using the residual window bounds from
\cref{cor:post29-bc-window-size-strata}.  It reports \texttt{failures=0},
\texttt{q\_checked=204}, \texttt{min\_frontier=9},
\texttt{max\_frontier=83}, \texttt{max\_frontier\_minus\_q=-6}, and has
SHA-256
\hashsplit{09ae87166af7f93ae6ef4df0322fcacde}{c86fc6ae5701ce0117851a655643b70}.
The low-low degree-profile verifier computes the exact profile count
\(\mathsf N(r_1,r_2)\) used in
\cref{prop:bc-active-twenty-low-low-profile-full-budget}.  It reports
\texttt{SUMMARY checked\_m=190 failures=0}, verifies
\texttt{active20\_failures=0}, and has SHA-256
\hashsplit{cc1b41189d965a5e076d3524647caa5b}{7f5c5102ca6a1d8b03c623c15c29dc40}.
\fi
The Pieri and fixed-witness arithmetic verifiers above are conditional
arithmetic reductions: they do not discharge their combinatorial
fixed-witness or compression premises.  Their detailed ledger is retained in
the source archive but excluded from the publication proof ledger because
those routes are not used after \cref{prop:post29-bc-layered-mgf}. \\
\else
Classical $B/C$ correction-prefix bridge &
The compact residual-closure proof states the Pieri sign and the sixty
quantified boundary-through-twenty-ninth correction contracts.  The
clean-room domain audit matches every contract to its callers, and the exact
frontier logs record strictly positive margins.  The GMP/OpenMP bridge verifier
checks $m=31,\ldots,15447$, reports \texttt{failures=0}, and has SHA-256
\hashsplit{fa267f4eed4eda48405e7eb468c8b022}{f997ad0bee5ce6324bc8618affcb0e6a}.
The longer fixed-witness and compression routes are proof-development history
and are omitted from the main paper. \\
\fi
\addlinespace
$E_8$ finite bridge replay &
The replay checks the $m_0,\ldots,m_{100}$ source table, verifies
the prefix Chain differences from $1$ through $65$, verifies
$\Qthree^{E_8}(69),\ldots,\Qthree^{E_8}(77)>0$, and verifies the extended
Chain differences from $67$ through $95$. \\
\addlinespace
$E_6,E_7$ rectangular tail replays &
The combined C++/OpenMP/GMP replay verifies the rank-corrected negative
endpoints, exact root identities, scalar and Chebyshev Bernstein predicates,
retained-cell counts, base-ratio lower bounds, and uniform odd-step bounds for
both tails, exiting with \texttt{RESULT: ALL PASS}. \\
\addlinespace
$E_8$ rectangular OpenMP replay &
The exact source audit first reconstructs all $29$ negative-bound rows from
their pointwise majorants, moment sources, and moment integrals and reports
\texttt{E8\_NEGATIVE\_BOUND\_AUDIT\_GMP: ALL PASS}.  The C++/OpenMP replay then
consumes the authenticated table, verifies every expected retained-cell count,
prints all base-ratio lower bounds and the $n=133$ odd-step ratio, and exits
with \texttt{RESULT: ALL PASS} and status $0$. \\
\bottomrule
\end{longtable}

The following table is an identity ledger, not the active-build manifest.
For historical source certificates, the first hash is the recorded generator
identity and need not identify a current source snapshot; the accepted output
is the immutable output boundary stated in the replay contract.  Current
sources rebuilt by the standard command are instead pinned by
\texttt{replay\_sources.sha256}.  The repaired
type-$A$ and first-hit trace rows below give both the current source hash and a
newly archived accepted output hash.  This table authenticates identities; it
does not override the accepted/diagnostic/superseded classification, and a
superseded recorded output is not proof evidence.

\begin{longtable}{L{0.20\linewidth}L{0.34\linewidth}L{0.34\linewidth}}
\toprule
Replay & Generation/checker identity & Recorded output identity \\
\midrule
\endhead
Type $A$ finite-overlap &
\hashsplit{3a0cacb5c1f402dfc954e07513da4335}{219d96cd9335cc15ade6356d54d784a0} &
\hashsplit{7339140167ea9e12101a7e0a316abc9d}{9be1e52c021a44c23e77f46ab776ff05} \\
\addlinespace
Classical Rains-square cutoff &
\hashsplit{ea02f81ac518eaaf849039112fd06781f}{908c4094619d54f0685987a7b2424c3} &
\hashsplit{ca3a938c9d45d65eb892f6ba637ed8c3}{a7b1d32a571a5be52e2f5efb80477176} \\
\addlinespace
Classical first-hit trace &
\hashsplit{ea02f81ac518eaaf849039112fd06781f}{908c4094619d54f0685987a7b2424c3} &
\hashsplit{ca3a938c9d45d65eb892f6ba637ed8c3}{a7b1d32a571a5be52e2f5efb80477176} \\
\addlinespace
Classical $B/C$ source-generated row-gated bridge &
\hashsplit{fbd5565215a7b976b67a4529ad21daa2}{f235f2988da463c502370e3dab3ad797} &
\hashsplit{d019c183f7dedff1049199b31a50d681}{3e0c17a0c565df1e845119e2a9fba054} \\
\addlinespace
Classical $D_{53}$--$D_{63}$ source-generated row-gated bridge &
\hashsplit{1e574e0c53e624df30e9c69bc24280b}{8e6c56a099923428f0a27036b74a7454b} &
\hashsplit{a84df851316bf376f00933ff958b7ad0d}{ac444acfa8f45d05d4221f17ddadd12} \\
\addlinespace
Classical post-$m=29$ first-hit &
\hashsplit{4499d62b37fc54d700fbe80973d07dca}{4f3e91a6ee17ed7b397595aadf84863f} &
\hashsplit{21642c87b90affa8e618bb9f97bf6ad3}{46289380832690546488c6dcc7a7300e} \\
\addlinespace
Classical $B/C$ interval direct tail &
\hashsplit{cba8addb7581d99cbd7d92d33c11f01a}{2c78267b0dd8d17d39f7c938bf565876} &
\begin{tabular}{@{}l@{}}
\hashsplit{80fce36378ec2192dca2c0b616c083b461}{68cf5134b66e73775156b995ea303b}\\[-1pt]
\hashsplit{53530cfe3d4b1947a77bf45c5c5aaa56}{a052ec1cde93c0704b68e004f22f2cb5}
\end{tabular} \\
\addlinespace
Classical middle type-$B$ unitary square trace &
\hashsplit{d0d503ffc4bd6cb0485a97a4cef2178a}{930f3f00c38ffd4486cb66aa01652aa7} &
\hashsplit{7780bfe3253bf9e030a9a544dcd7327f}{cd905aed8a2890ded191dc38efae8230} \\
\addlinespace
Classical middle type-$B/C$ exponentially tilted MGFs &
\hashsplit{0f5ad1037bcf738ab4ddd9c5b12f151a}{adc1bff2b5c4726c60622f608ef9e0e7} &
\hashsplit{ca6860204466e4a6b38d9c5597511cd1}{1cfebb4b1b65cab88db6ee45f9477111} \\
\addlinespace
Classical residual layered $B/C/D$ MGFs &
\hashsplit{71cc2293c4ad5ba50d216c8f939cd1ec}{604ce7b86d6a19e9be1fcb27adc421fe} &
\hashsplit{3e4900e5e2de2ec1a7fb04830ac59dea}{7369f8eebb1932ec08ac9ef73ba11654} \\
\addlinespace
Classical $D_{25}$--$D_{52}$ exact/mixed finite-rank MGFs &
\hashsplit{71cc2293c4ad5ba50d216c8f939cd1ec}{604ce7b86d6a19e9be1fcb27adc421fe} &
\hashsplit{f87ab4bce0b189878ef4c3528ce34687}{1721bff05cb00ccb553a6c4e84e1f049} \\
\addlinespace
Classical $D_{25}$--$D_{52}$ correction-aware bridge &
\hashsplit{9a6cf4da5b866d3b8fc2ce6d1d1a442f}{d3b419702fda2296da83d2fb788379e4} &
\hashsplit{e47196c8097fb2a087b4673d47a5e34c}{aae1cd1a316058759e17b13708e2d888} \\
\addlinespace
Classical $D_{53}$--$D_{295}$ $A$-free frontier &
\hashsplit{444cf4dbe59d5000a3315f31054ef07e}{c1cf710edf91a7c685f973644d65ad74} &
\hashsplit{876d3f752d03b8f35183f0b28af7c15}{e9e28d104396fc65c215e5d854d8ae31f} \\
\addlinespace
Type-$D$ legacy-envelope counterexample &
\hashsplit{bc22efcafafe2d40b5f5430cb8f9cb046}{bc6b9b188da684e61dc6fdfaa27077e} &
\hashsplit{463e3c9632bb92936620fa65c03d3ff5}{abdb86940624f1621ffa976156d7d02f} \\
\addlinespace
Classical $B_{11}$ post-$m=29$ bridge &
\hashsplit{6ae0d3bf506737b32480c336a6213ef6}{820e656149789997731199daa9177f56} &
\hashsplit{14f8c9813d3fda536e02a5c5809cae13}{6e6f531669d9f20806352415702db9f3} \\
\addlinespace
Classical $B_{12}$ post-$m=29$ finite tail &
\hashsplit{1afb4bb14311864103dd5ebd381842bd}{238f09f90c2064257f0b05e34e705443} &
\hashsplit{a8e6a7c58eccf3565f578c8ed9d9187a}{507000b497efdca6593c40013c4f6854} \\
\addlinespace
Classical $B_{13}$ post-$m=29$ finite tail &
\hashsplit{d6e6f57279872ce94c81d31ed45e5e7b}{c0948993ef335a45d9717bc64f7a9031} &
\hashsplit{08b9666bc771897e9702c0132620f4ca}{e9215b30c92c160ed52795fe6db096ee} \\
\addlinespace
Classical $C_{13}$ post-$m=29$ finite tail &
\hashsplit{1511406f2e1015184027f9d90ae380f8}{9789c5d2ac5b8f49a599608d62ca1f1f} &
\hashsplit{1f66ffebfe9e722f443073031ec74401}{06b87116a1d1db279d154967a7e45049} \\
\addlinespace
Classical $C_{14}$ post-$m=29$ finite tail &
\hashsplit{000e286d0dd04595a66c9efeeafaae9b}{a38318cab396f08f282540b6713d696f} &
\hashsplit{6dc9a3f519c64568b378b8689dd78867}{e5311abdf1257268f37ba79408feb1c3} \\
\addlinespace
Classical $C_{15}$ post-$m=29$ direct tail &
\hashsplit{d7ec8b228ee53f304d4de0060c371121}{b9e21c3c456f8263346ad862b20fa4c1} &
\hashsplit{68fd43b912006354f77fb2c051475767}{f8c6d85b3355b2a3ba3c438cd7a2ad68} \\
\addlinespace
Classical $C_{16}$ post-$m=29$ direct tail &
\hashsplit{23af44fbdc597cf0332c422d716cb2c}{483a47bafd1756f15b48486062b77a1b9} &
\hashsplit{9e15444c6e789d3aba9f83feaea03a22}{f5dd423e9c2a138b68f959bdb22c719f} \\
\addlinespace
Classical $C_{17}$ post-$m=29$ finite tail &
\hashsplit{710f53d2736eec18956752a744f2239bc}{1461627736ddf8fb1c17842c9bb7ff5} &
\hashsplit{dcd5eb8e782be93d926a6486f18d8a5e}{c445d33f8784ba251f60dd0203d8eb32} \\
\addlinespace
Classical $C_{18}$ post-$m=29$ finite tail &
\hashsplit{43f9efa815f1cf4ec325932889e383a26}{cfd8e829b2848043feb2efec41fea68} &
\hashsplit{e25bfaef2b3b6289631431fd26717704}{dc82c78009cc02cf9c8e8dc8861b5541} \\
\addlinespace
Classical $C_{19}$ post-$m=29$ finite tail &
\hashsplit{282f6e3ab2c952155473cfa46c4f82e4}{cf331462f4de406d76baab4302c1f307} &
\hashsplit{034d55810933481b342966686cf249a33}{649fac5333ad335df9b94e8f95b8273} \\
\addlinespace
Classical $D_{12}$ post-$m=29$ finite tail &
\hashsplit{6455ea1539e414087cba2335aa75bcfe}{176d936edffc15536a7c84321147aa0f} &
\hashsplit{f9b0a67a24912120f30b23f7781c88d5}{a2126144a88e341d5f9b7a0df1988c2a} \\
\addlinespace
Classical $D_{13}$ post-$m=29$ finite tail &
\hashsplit{6455ea1539e414087cba2335aa75bcfe}{176d936edffc15536a7c84321147aa0f} &
\hashsplit{6f6ee3ffe7f92d0127c89ac1ce429e30}{bb59a3a83cd113b54a60faa8c37b088a} \\
\addlinespace
Classical $D_{14}$ post-$m=29$ finite tail &
\hashsplit{6455ea1539e414087cba2335aa75bcfe}{176d936edffc15536a7c84321147aa0f} &
\hashsplit{aa4ad190afbe3a76383ad6999501f547}{4098036724a8bd48144620bc4fa045dc} \\
\addlinespace
Classical $D_{15}$ post-$m=29$ finite tail &
\hashsplit{6455ea1539e414087cba2335aa75bcfe}{176d936edffc15536a7c84321147aa0f} &
\hashsplit{d7977837cc415e9ab1550123a5829832}{e07f9cc9c3163b54ceb980fab46d8108} \\
\addlinespace
Classical $D_{16}$ post-$m=29$ finite tail &
\hashsplit{6455ea1539e414087cba2335aa75bcfe}{176d936edffc15536a7c84321147aa0f} &
\hashsplit{2a037f4f94b10af7f9da01425c6c22be}{eee1f492e4dda55e26d6962ee37a0ca8} \\
\addlinespace
Classical $D_{17}$ post-$m=29$ finite tail &
\hashsplit{6455ea1539e414087cba2335aa75bcfe}{176d936edffc15536a7c84321147aa0f} &
\hashsplit{dc513cfa308133f6e0229107cd44ee}{ebdaa8093ec31491febbde2b223641a3a4} \\
\addlinespace
Classical $D_{18}$ post-$m=29$ finite tail &
\hashsplit{6455ea1539e414087cba2335aa75bcfe}{176d936edffc15536a7c84321147aa0f} &
\hashsplit{7241b564ab367ced09f6faed98e3876}{a5cf65d5b4de6e06a09d445b4e5bee455} \\
\addlinespace
Classical $D_{19}$ exact tail and bridge &
\begin{tabular}{@{}l@{}}
\hashsplit{1f296ce3a32756fbee27bdc44eb9749d}{03dc88e74107eb7f15d91f8fda3a15e5}\\[-1pt]
\hashsplit{edec7e651f112a342909230068651c9a}{59dc44a58bffcb5fa9e0420bf9f7ad8b}
\end{tabular} &
\begin{tabular}{@{}l@{}}
\hashsplit{42e987f476bacc0e7422adb3825562227}{aa0511a58a6984f61adcf4c08b3e4b5}\\[-1pt]
\hashsplit{a0a46e0274a2fa4ded548edfbec2f2d8}{3cc7d5cecc13be155aa37d77fe9e11a0}
\end{tabular} \\
\addlinespace
Classical $D_{20}$ post-$m=29$ finite tail &
\hashsplit{6455ea1539e414087cba2335aa75bcfe}{176d936edffc15536a7c84321147aa0f} &
\hashsplit{d114e6a029422bde33cf360dab227cce}{8703fcf746c8fe928fcb24529e09736b} \\
\addlinespace
Classical $D_{21}$ post-$m=29$ finite tail &
\hashsplit{6455ea1539e414087cba2335aa75bcfe}{176d936edffc15536a7c84321147aa0f} &
\hashsplit{87cb3ddbec928e5b22ae74f2d4a4a0}{e70cc5a6001774558db105a0fdc1c7c235} \\
\addlinespace
Classical $D_{22}$ post-$m=29$ finite tail &
\hashsplit{6455ea1539e414087cba2335aa75bcfe}{176d936edffc15536a7c84321147aa0f} &
\hashsplit{8b321bacfdb915016fea96fd04b3672d}{33a1758b0163ebe06d68daab36034c64} \\
\addlinespace
Classical $D_{23}$ post-$m=29$ finite tail &
\hashsplit{6455ea1539e414087cba2335aa75bcfe}{176d936edffc15536a7c84321147aa0f} &
\hashsplit{366460e47742ef94d5ac2371169d70742}{d96e17426b89a1216bdcc6a8346e593} \\
\addlinespace
Classical $D_{24}$ post-$m=29$ finite tail &
\hashsplit{6455ea1539e414087cba2335aa75bcfe}{176d936edffc15536a7c84321147aa0f} &
\hashsplit{32e4c416211b54709ebb88d9458a0bc6}{3a5e99cd144680e0d9d8fac05c9ba761} \\
\addlinespace
Classical $D_{25}$ post-$m=29$ finite tail &
\hashsplit{6455ea1539e414087cba2335aa75bcfe}{176d936edffc15536a7c84321147aa0f} &
\hashsplit{49cfdb89f74c85d5cd7b0fd90dc9d3bf}{e81cc219872d39bd4751e6327efa0d77} \\
\addlinespace
Classical $D_{26}$ post-$m=29$ finite tail &
\hashsplit{6455ea1539e414087cba2335aa75bcfe}{176d936edffc15536a7c84321147aa0f} &
\hashsplit{3becbf73915fac1362361aa8181e60af5}{7a3456497b1751d8b22f42e266c5b1b} \\
\addlinespace
Classical $D_{27}$ post-$m=29$ finite tail &
\hashsplit{6455ea1539e414087cba2335aa75bcfe}{176d936edffc15536a7c84321147aa0f} &
\hashsplit{6c2405c8fa6304bd3dd215ed05d00eda}{a9be0cef7bf6fdcae4998ed0de6fd290} \\
\addlinespace
Classical $D_{28}$ post-$m=29$ finite tail &
\hashsplit{6455ea1539e414087cba2335aa75bcfe}{176d936edffc15536a7c84321147aa0f} &
\hashsplit{20d8b9bebc9b09cecd90c6e860c30079}{ac13cec71ae13d716625d351da18ef7d} \\
\addlinespace
Classical $D_{29}$ post-$m=29$ finite tail &
\hashsplit{6455ea1539e414087cba2335aa75bcfe}{176d936edffc15536a7c84321147aa0f} &
\hashsplit{3340f19beacc7fcd2efa4bfdc4776c4f}{6a21cbecf97d68de7a83aa95f2d8402e} \\
\addlinespace
Classical $D_{30}$ post-$m=29$ finite tail &
\hashsplit{6455ea1539e414087cba2335aa75bcfe}{176d936edffc15536a7c84321147aa0f} &
\hashsplit{c6afbd1a77b10d874d82e5ff142128b}{837d7c09602f0dda5c3c7da92893a8e66} \\
\addlinespace
Classical $D_{31}$ post-$m=29$ finite tail &
\hashsplit{6455ea1539e414087cba2335aa75bcfe}{176d936edffc15536a7c84321147aa0f} &
\hashsplit{beb8d9008c08783096f4ad91753676bd}{2ffff514189cc42603cbfe5b1faf9f1f} \\
\addlinespace
Classical $D_{32}$ post-$m=29$ finite tail &
\hashsplit{6455ea1539e414087cba2335aa75bcfe}{176d936edffc15536a7c84321147aa0f} &
\hashsplit{e245d7ee1124af187c02e53365a1bd3d}{dc4f72893c1b8fa280ccb527a3357df2} \\
\addlinespace
Classical $D_{33}$ post-$m=29$ finite tail &
\hashsplit{930c0b9a705e079599af852d13bd05a5}{228337f36ec2fd981ea899f5c39e34ec} &
\hashsplit{b7d813cde42faaf8ca92a803946a20ce}{1b93322f9102f81a095e6771d14d2f9e} \\
\addlinespace
Classical $D_{34}$ post-$m=29$ finite tail &
\hashsplit{3cdcda5f9bd87351035730d42c459d40}{63545cc85a734c87ed29970f6c62f72e} &
\hashsplit{6700bcf5bd99f8a4b512a79e98708ae8}{b76ca06a6d1cbbd795f5acb87c96e800} \\
\addlinespace
Classical $D_{35}$ post-$m=29$ finite tail &
\hashsplit{e4b5f031f4adbde095ac9fd2ee463d62}{eadb03849825bef593d3caff6b9053fb} &
\hashsplit{1100c50808f2b21a9977602e41b24166}{9da29afcdb2e36da9db445550bcd6bbc} \\
\addlinespace
Classical $D_{36}$ post-$m=29$ finite tail &
\hashsplit{8b9bc37a35ee3d03433876a1cb24f9a1}{0e10977c095d1f08a5fde35416f58e27} &
\hashsplit{a9a3b4aa024930a41cdaf0eb2e97d89c}{16dac9f72f58e3016283185d1a8a03fa} \\
\addlinespace
Classical $D_{37}$ post-$m=29$ finite tail &
\hashsplit{66642943003e5bce9d3ae2e182fcd310}{42a3131d237f6c10bd1e932dc16aa225} &
\hashsplit{8b7f4c16de0354d97d3a6c6233883ad}{4e0f9640c9ac74c47bbdf1b760009e15b} \\
\addlinespace
Classical $D_{38}$ post-$m=29$ finite tail &
\hashsplit{d143978e45e6ba1819d0dcdc9658cfc5}{dc02ca5df5478c0726023220969edbb6} &
\hashsplit{44641f3b21f4d23764e653f2bda16d90}{6bd9d664e13f6928bb5ba961f3930ca3} \\
\addlinespace
Classical $D_{39}$ post-$m=29$ finite tail &
\hashsplit{9985e56e5313cb20f7e1b368bacde6a0}{b3112b31dc2dcc149fd636f5dc5cbc0f} &
\hashsplit{6f052631b33f7be04fb0e9e32dc34d6}{3b1b2bf401a8ee3b479870d53ce9240a7} \\
\addlinespace
Classical $D_{40}$ post-$m=29$ finite tail &
\hashsplit{3fd9763dc1afc89af81a27758e0098a1}{0af7d28ddc1b3101e36ed101d01b6220} &
\hashsplit{d3e94418ad9f6fd1e0911a89e19c2f1e}{5168d7727d5dc57be9ddd2b491a5e098} \\
\addlinespace
Classical $D_{41}$ post-$m=29$ finite tail &
\hashsplit{8352225e575eb2a427244fb81df383e9}{030498b60b85463707fe2edd5bcf9577} &
\hashsplit{475339350b7802bb40f28f76de4f10aa}{9587942ac1e0df86873c0119463e4e51} \\
\addlinespace
Classical $D_{42}$ post-$m=29$ finite tail &
\hashsplit{b903e1adf2c11e389d32bb0bc5f0001f}{ba4c5d089c937c34fadae21f6c749ca9} &
\hashsplit{22ff552c35e6800d669e62379cf7e690}{235cf90c96cfe21c0d1205ba161569fa} \\
\addlinespace
Classical $D_{43}$ post-$m=29$ finite tail &
\hashsplit{6c448a0b987d603c81037325fb96bba2}{0bf2eeff52ed66f9bd713fd45f836f5a} &
\hashsplit{d42e5f1fa39b939d92d21b4b8632b6a5}{1c8a646d517a6900bbde580b3ccfdb09} \\
\addlinespace
Classical $D_{44}$ post-$m=29$ finite tail &
\hashsplit{b6ac9839cedf8c1a9457f0d3f9fe5a37}{5da947589749585716dfb3df07d9a3d6} &
\hashsplit{4a6e05c3837db8d3aa30cb7f46c78735}{9337b394808c563544e6ace72a4740b8} \\
\addlinespace
Classical $D_{45}$ post-$m=29$ finite tail &
\hashsplit{964825593c8c7a3314d4f12587c3ca55}{b5a5d77fcd928c6f7fa5d98b1f946efa} &
\hashsplit{1264b3c773d580ff2d08a8355b118316}{64b1bea28f3595f8ccbef028c41fd93b} \\
\addlinespace
Classical $D_{46}$ post-$m=29$ finite tail &
\hashsplit{29fc071c7d01d990c17532c1594f20a7}{f9a8a3b0f07900a9082818bbb8508c18} &
\hashsplit{ef7655165b98bcd29261a1a3dbd39336}{a48e038df122b738bb63f39dc5668f41} \\
\addlinespace
Classical $D_{47}$ post-$m=29$ finite tail &
\hashsplit{1b663aef0410dfba5ca577df7bacc51e}{8dd1cfac397d6730774d02089dd3717c} &
\hashsplit{897adf7d7720993a66ff9ecd2e9ee58b5}{a8cfc4211005b3f0b652a0a75eec147} \\
\addlinespace
Classical $D_{48}$ post-$m=29$ finite tail &
\hashsplit{6e38bb995342e904702ca544397c0ecc}{9d31fcda8b043ab5b16525bf3becdbda} &
\hashsplit{d66d077c5538c716da9e6551903db589}{67d4a41df3779aaad8fcd3ba544c6cbd} \\
\addlinespace
Classical $D_{49}$ post-$m=29$ finite tail &
\hashsplit{cf9f373a83a0dce0b398ee8421d7c449}{887cb2998f3e5d9c756a884973ce36b8} &
\hashsplit{1c4a2651719271c321e11075425f3904}{e080fb06339cc31c2630300ed9fdec3d} \\
\addlinespace
Classical $D_{50}$ post-$m=29$ finite tail &
\hashsplit{54d814fe2824c0a7a7fd09a2fa8973ea}{584e694099e4b6f2ac426a113f5ca14d} &
\hashsplit{1d88e66f53c82abd5d2162a3c746d4e6}{dd6c25657ec1ce66b89d8a4cbf256c12} \\
\addlinespace
Classical $D_{51}$ post-$m=29$ finite tail &
\hashsplit{95f5e39b8ad38e5f197a09f22aa52445}{0471a61e53f3ef49011fbbd033327158} &
\hashsplit{828271646ebe60b7b8ad37af84860882}{c96526486cad8be2d70a600eaaa1c49f} \\
\addlinespace
Classical $D_{52}$ post-$m=29$ finite tail &
\hashsplit{d05e95b2414c036e5c6f88077ced3287}{4e9af1e153492e4597a4f933c9e47f35} &
\hashsplit{e8e7dc7aae4c08d7c571e41dd2939213}{bf17b4ed0cc322adf86734505976803b} \\
\addlinespace
Classical $D_{53}$ post-$m=29$ finite tail &
\hashsplit{bc9d5a804fcc9e72bf1b49b394de72c1}{a8269ad0beac4d1d7365845af2701bab} &
\hashsplit{a3c2559355281edd74889039564bbae1}{3dfba5ce10ec92d957fa3a33a6227e48} \\
\addlinespace
Classical $D_{54}$ post-$m=29$ finite tail &
\hashsplit{43df015ed783f51dfb134d19c778968a}{c0eee616769fd97d5bf625af4ca087b6} &
\hashsplit{0f54bb8b1e66209fcc79179573d49dc2}{93e1e79fdb36a396e081b261a1b5b303} \\
\addlinespace
Classical $D_{55}$ post-$m=29$ finite tail &
\hashsplit{54e76fa1b690aee3e2f2c02214850c38}{d2a69e2eac388a90217ea43df2a3eb93} &
\hashsplit{6d5a489a0f135317a3305e66cc0cd6fe}{da03a6cfcf37a04337da5872bdcbc127} \\
\addlinespace
Classical $D_{56}$ post-$m=29$ finite tail &
\hashsplit{066731d6ee5d0e150645591e84c2a0cb}{dc1da3bc2a4c86855ffba5f1d9a57817} &
\hashsplit{26806ffc35f50a94493c9ee7f0c196c2}{92616ca1cf00a22acdbf1e2a0415336d} \\
\addlinespace
Classical $D_{57}$ post-$m=29$ finite tail &
\hashsplit{bc130849ef80bf0902757e96f3c8cc6}{c84c4ef72db2fa965eca1d4467c968dde} &
\hashsplit{c6e69b001e3632bd578195126f266b6e}{01f8bbc80ae9e9451243ee3afbaf5dac} \\
\addlinespace
Classical $D_{58}$ post-$m=29$ finite tail &
\hashsplit{063c93acf978ef1816b1a0c8c927b721}{7add401a80704d6111f3c1fb26de4554} &
\hashsplit{6347fb99063973f4f19b3b15354aa73f}{68d9e471f749adafcd8a60d3fac027ab} \\
\addlinespace
Classical $D_{59}$ post-$m=29$ finite tail &
\hashsplit{5cbce0bacd52cfee3d85bec546f79a37}{f3c1e119f15901c073387978e37ea273} &
\hashsplit{aeb95cd95e84ef81a5bd97dae7b05839}{2e0acb24a23ad2f3f7e2c4ecee0fd401} \\
\addlinespace
Classical $D_{60}$ post-$m=29$ finite tail &
\hashsplit{9b196b4efba5c85abafe14a3b51c3038}{a3440c6fd590485a922469e1cdb4327d} &
\hashsplit{38370d31bad7b33df73ecc854ad4c3}{a97a2bcb8e7e85dda7758aaad252e56a07} \\
\addlinespace
Classical $D_{61}$ post-$m=29$ finite tail &
\hashsplit{9982522dc6db42b69563e313dd51c576}{e230db1d2635f72d1dec277200db968b} &
\hashsplit{3920dbf5295facfac7ca14dba3f03792}{9eb9b80a07f650da1300f6e35d81c419} \\
\addlinespace
Classical $D_{62}$ post-$m=29$ finite tail &
\hashsplit{3aaef754cca0cf37fbc475cf88685351}{395a3568c1dd8c7f82e7d3cf4caa2398} &
\hashsplit{1656ddcb8569128be98151eed8f8fb6d}{798d249e837b6879aed8afd8d7a5dfa9} \\
\addlinespace
Classical $D_{63}$ post-$m=29$ finite tail &
\hashsplit{247b133c4bcd160998fe093f6bb429ac}{6989e50590642ccd328ddd9c56639f66} &
\hashsplit{9df3f216e7718fcf3c33278b98523017}{41f6f98f681e17ed008584085220c9c1} \\
\addlinespace
Classical $D_{64}$ post-$m=29$ finite tail &
\hashsplit{9c4afa3e54ca488f848cb3b9e2043cda}{d9afbab05fb018b3510db7e162e9719b} &
\hashsplit{d4dbd28548a01801d0f4fb42fef3aeae}{9eb112e3f2cac518437abbf10ddfc532} \\
\addlinespace
Classical $D_{65}$ post-$m=29$ finite tail &
\hashsplit{350c95515bbe01694ac24071f0483159}{50470f5b31d9ab7a7510a7b55bf896c1} &
\hashsplit{022308c3011e70f9743db140e618fa92}{8d8ebaaf808ba7ace61a06133b968e4e} \\
\addlinespace
Classical $D_{66}$ post-$m=29$ finite tail &
\hashsplit{d5ff976201f0e668893e78254dcd9822}{c86af1d4d5505c50049ca0464e8476d9} &
\hashsplit{8d8b8df1ee57bd6567f1a3d08e6f39f4}{6ac24447c12248beaadd5ecb1d505f1c} \\
\addlinespace
Classical $D_{67}$ post-$m=29$ finite tail &
\hashsplit{807ed743b04aee27029719abdcd350b5}{ae03eb8941f2f6daa4346703f4e747da} &
\hashsplit{fb5499210f42feb3db19d0bbccf3589c}{b8622f6b0aa8edae43f917709e68876a} \\
\addlinespace
Classical $D_{68}$ post-$m=29$ finite tail &
\hashsplit{e52df93f12a68fd3dfd92ea1a489aee4}{be887fc1be7cf545f2fd6e21bd89e15d} &
\hashsplit{e6ca2fce3bde6cf2fbfb79308053aa}{fe39ad3a891a3cc0d7fe37084a24177055} \\
\addlinespace
Classical $D_{69}$ post-$m=29$ finite tail &
\hashsplit{e6b77aad80fe398dfe2b5c9770548904}{ccdf2b6f4c8e5810ba5aa95694d59910} &
\hashsplit{86bd7a3e527be1c8e91fd24859ccfe66}{3918950c79483e0fb7d24a5d1f84691e} \\
\addlinespace
Classical $D_{70}$ post-$m=29$ finite tail &
\hashsplit{4a42ae7ba0e238fed18f221b44d8f347}{5b539524f0ffe32ada8c8bcf8f351572} &
\hashsplit{510fa3957eaf02a26a9e4e69945bb4a1}{b34d33e36738a9efda17ab36b10dccf6} \\
\addlinespace
Classical $D_{71}$ post-$m=29$ finite tail &
\hashsplit{b07c0663d7d4b70129d2cf4d56373458}{2afbce31b82b78c98a5d2253aa854150} &
\hashsplit{602b99dd2bd22b16e4a336b26c15b491}{1f044367e7c7ff421c4bea53b358fb53} \\
\addlinespace
Classical $D_{72}$ post-$m=29$ finite tail &
\hashsplit{650d9d3cb069d9f220100577d4c2250}{a8ff8a85041b743b45a5a3f227a982c9e} &
\hashsplit{6c1b1070bf528fcbf5e7cfcdbbbe670}{443c122ed2837ab8e2c944abcbbb97ddb} \\
\addlinespace
Classical $D_{73}$ post-$m=29$ finite tail &
\hashsplit{322b1d38392ccdba434f172e51d2f657}{c8efabf0a0652296051889e31cf50fb8} &
\hashsplit{b2ee7080c19dae3f253a08a57de0351}{2a4d0f4da5c27429fd034c37498f6a12d} \\
\addlinespace
Classical $D_{74}$ post-$m=29$ finite tail &
\hashsplit{2a9b94175879a4b2456c1bfc6ff0fe96}{12e14aee280aff9acdf61d1ff5ac5fe4} &
\hashsplit{eee9ed7239c9671ffdcf480e2b3d0890}{7332adef7e930af5930181c826eb46f8} \\
\addlinespace
Classical $D_{75}$ post-$m=29$ finite tail &
\hashsplit{9901fdae225b08bead8b02665fefc17d}{5c476f5172d7403ba3171954ba916c30} &
\hashsplit{c804e87fec00b21fa05056729c74be7a}{ece3f28f2f69ff4d8c64691c8af212db} \\
\addlinespace
Classical $D_{76}$ post-$m=29$ finite tail &
\hashsplit{210cb83e64c1ec7b78e3873d8d991d2}{681a183be71452c87c261ebdb586ebc86} &
\hashsplit{0f29482007a812ffbc79483f77fa0aa}{f69c033d2c2f288813568e90ad0955230} \\
\addlinespace
Classical $D_{77}$ post-$m=29$ finite tail &
\hashsplit{4dcbbcdcfeef7b6b4fb73c48b1920184}{a4cba920baac9b89cf832d4515c9e98e} &
\hashsplit{46d8bbe29c15c4dc3315ba9ebf7be9a}{7db533ecae0ec87c36d29a184dd481f34} \\
\addlinespace
Classical $D_{78}$ post-$m=29$ finite tail &
\hashsplit{f74b68a2c3a5746d82acf30e5b10b641}{b0ab66416c566b769e5dacded19e21c9} &
\hashsplit{796de51fc5a01c17bbe0f6eafa4c127e}{29389b521e38bcc4dda32610f55f6b79} \\
\addlinespace
Classical $D_{79}$ post-$m=29$ finite tail &
\hashsplit{96061823694f2984e928f7bfa37b0f3}{0106b9d212356c39d5add8493617baede} &
\hashsplit{a2475a6f521776dac099fc21ccc46974cc}{2abb2cc86dba75fc3aef8af6568fd9} \\
\addlinespace
Classical $D_{80}$ post-$m=29$ finite tail &
\hashsplit{df1d12680cea80faeda9b017deee3bc4}{07cc423e34e447a7fc87f4857f192a38} &
\hashsplit{8c1264fa4292d8f7f67ef978d9109d99}{8fa39f759a4db622a58030cefa854851} \\
\addlinespace
Classical $D_{81}$ post-$m=29$ finite tail &
\hashsplit{457989dcabfe1c5aa97e8dde03da4b5f}{2fff9eb48aba4d8b0204d44b36685b96} &
\hashsplit{a295f235c7e29ad7719bd8f1552f0831}{c5220c6c1f08a6c4b2445fad822be839} \\
\addlinespace
Classical $D_{82}$ post-$m=29$ finite tail &
\hashsplit{d114c05a9e22475cb14e643c7e4db545}{7aa911ff10656c657d1e9d3216df2db6} &
\hashsplit{1e7af36a394f21b0057c5c845571681}{06c7068239b15c22b2ce9480871c14ad8} \\
\addlinespace
Classical $D_{83}$ post-$m=29$ finite tail &
\hashsplit{f8ae51c0443ca3a14b3529367cc9f641}{303180eae84e20b9148212076a922582} &
\hashsplit{b14ea35f764564dd5d7685f89adc9669}{48c39e414b9fe46752ef7f21a3b1d43f} \\
\addlinespace
Classical $D_{84}$ post-$m=29$ finite tail &
\hashsplit{2acea25cd9b982e5010ec51e8d049590}{bbe0230c6e099c76265a1f83b05a18a5} &
\hashsplit{92693ae7cb8f56d258f9824442f3804e}{c8ca7a2e8f15b4da3ceec0b725920e66} \\
\addlinespace
Classical $D_{85}$ post-$m=29$ finite tail &
\hashsplit{6b109272136ec349ff69c00b1f270660}{1e95c1efbbe01dd159c2f5dbd222237a} &
\hashsplit{0d8ec01ede79008e9a2dcffe97fa8ce8}{c1c81307bef53038f7a3036ddd4baa5f} \\
\addlinespace
Classical $D_{86}$ post-$m=29$ finite tail &
\hashsplit{c4288a10881cb4c628091fb91354852f}{e73d8499158fdf52bb69230e546b7ec3} &
\hashsplit{5ab292178239b92c9e0e4e4445e9e0f9}{42b75fb3ebe00e3e5706cdce3e2d71d4} \\
\addlinespace
Classical $D_{87}$ post-$m=29$ finite tail &
\hashsplit{5f1f823bf1133549db47e3e9220c71299}{353f1fdc2a0c41c4b9b39bf8d595cbe} &
\hashsplit{9de45735f34728dc14228c0f95a42f88f}{84be059311c8cf9dbb2577ce6d706f3} \\
\addlinespace
Classical $D_{88}$ post-$m=29$ finite tail &
\hashsplit{adcb516f50700d151d05cf926b484531}{a6a9224a32bd2bea197b006f81a49b19} &
\hashsplit{4bb7b5143445fd7b5cc613c14bdf3800}{f0e5cd8e302adfdf178ef4db540422f1} \\
\addlinespace
Classical $D_{89}$ post-$m=29$ finite tail &
\hashsplit{920265217898e51e8825c43a16a47ce0}{bf4590ae92a00b676bd3d89018a71561} &
\hashsplit{7cf2e4118d1c3815b907bcb3849a847}{0598de3ec0d96f0b066c7385f713b5f5a} \\
\addlinespace
Classical $D_{90}$ post-$m=29$ finite tail &
\hashsplit{5b579ff50cdac23aa775f319400fd54095}{defefe2704410cf69f8255450b7550} &
\hashsplit{ba3e023245c270566c72eba6723612d39}{9c979a46e25b1d56ea17b8ea08696b8} \\
\addlinespace
Classical $D_{91}$ post-$m=29$ finite tail &
\hashsplit{518e1358dbf4de57ba79592cc90d4117}{d5c99ffe6db660e23f8c75b312f82a8e} &
\hashsplit{71c6d6b1b8bc17b21422f55cc2c1dbb4}{001a5d5e7bdf7868ec91ca62c2a22abb} \\
\addlinespace
Classical $D_{92}$ post-$m=29$ finite tail &
\hashsplit{62ae14c767b1b01b3081b6093ada47ba}{0e0f36e3886f9d0e5e29212fc67d9b39} &
\hashsplit{d5c20243f621cfb443c7ca6f1619e3c}{ed149dbffbb92ab179320be96699deae3} \\
\addlinespace
Classical $D_{93}$ post-$m=29$ finite tail &
\hashsplit{282efa573810e3ea137038b47a447e56}{f1bb0ba67023d853d9255c146b25f8fe} &
\hashsplit{27d61a5fff05d9bf688cfbddfbffbc21}{d5b90ade28e3be86d83e039268cbe7a3} \\
\addlinespace
Classical $D_{94}$ post-$m=29$ finite tail &
\hashsplit{9f63c2508b241a3534130fa35681eebf}{322401c6ddfe3ae08e0ece045655341d} &
\hashsplit{690fea30cfb9ad79d887d57653d24291}{1b5dfed68157669ab38b8a351331d545} \\
\addlinespace
Classical $D_{95}$ post-$m=29$ finite tail &
\hashsplit{d1c2aa46419f77c0503dae8ca16cf7c}{0c2e349564dca90a1a7fb0229607a6f2c} &
\hashsplit{b37110f3c85f43bd736599b0715f512}{bf8cf48e2ef964939dd64657845c96ceb} \\
\addlinespace
Classical $D_{96}$ post-$m=29$ finite tail &
\hashsplit{23569a4501d98a11e41af838660ddc4c}{c39a0ed580fe53eb632b2fcc295c699a} &
\hashsplit{1ac17e2524c6501556f5fd55a61ef95a}{fe744a47ccd469c109492f20467e02b7} \\
\addlinespace
Classical $D_{97}$ post-$m=29$ finite tail &
\hashsplit{171c1586935f6f787cef2bdea17e4771}{3bb13c48f1a0b67dbb374aba58a10a16} &
\hashsplit{26be4b2020bf5eb25a1abd7565520a55}{0e67abf90a452893ca7c22132ec90ca9} \\
\addlinespace
Classical high-edge no-TV Chernoff trace &
\hashsplit{ea02f81ac518eaaf849039112fd06781f}{908c4094619d54f0685987a7b2424c3} &
\hashsplit{ca3a938c9d45d65eb892f6ba637ed8c3}{a7b1d32a571a5be52e2f5efb80477176} \\
\addlinespace
$E_8$ finite bridge &
\hashsplit{46900d26a211d572bd41cbe92dd6e59}{3936e4255af0e1560824a0f29ba299129} &
\hashsplit{17cd75bbedc88d6e9e6ac85ca8ae8c4}{f4825018e5e8de5a5cfdf1cf60ed2414c} \\
\addlinespace
$E_6/E_7$ rank-corrected rectangular tails &
\hashsplit{20c89e8befe8a7172455f62e4b5467c6}{19f9f7f01d93e717f78ca60c832b4cb3} &
\hashsplit{8ded845d7746e5d46d2cb834b09f0faa}{ac5e2e8cefc0b983cf734627eba5713a} \\
\addlinespace
$E_8$ rectangular OpenMP source &
\hashsplit{2024a73d5ad98c33d80dbc9c12820927}{37d0172a3088e7db527aa7cd7406d764} &
\hashsplit{6bc3c5bf1774e5c50017360f9c7dda09}{932454f2f4fe73f7c8cfbba0650ee025} \\
\addlinespace
$E_8$ negative-bound exact audit &
\hashsplit{7758379e442ca9b65844db4de7fbcb9d}{49cd0105b142d958c0ca5300c8678046} &
\hashsplit{ffd490645e0fe1858980f7ab9bbda23a}{2e1cc56764d1c1180460271e57ff3789} \\
\bottomrule
\end{longtable}

The classical no-TV Chernoff high-edge replay is a load-bearing proof replay
for \cref{prop:post29-chernoff}.  Its archived log hash is
\hashsplit{ca3a938c9d45d65eb892f6ba637ed8c3}{a7b1d32a571a5be52e2f5efb80477176}.

\section{Post-\texorpdfstring{$m=29$}{m=29} Classical Replay Transcript}

The post-$m=29$ first-hit proof is the union of the active $530$-row $B/C$
replay and the two exact type-$D$ bridge replays.  Its composite ledger is:
\[
\begin{array}{ll}
\text{finite cutoff} & 296,\\
\text{rows checked} & 783,\\
\text{$B/C$ stable moments generated} & m_0,\ldots,m_{557},\\
\text{scope} & \text{one Chain step at }N_{\mathrm{hit}}^{(29)}(G),\\
\text{status} & \text{all three accepted replays exit with status }0.
\end{array}
\]
The bucket decomposition is:
\[
\begin{array}{c|c|c}
\text{bucket} & \text{rows} & \text{minimum margin row}\\
\hline
\text{BC exact }m=30\text{ window} & 14 & B_{11}\\
\text{BC stable-tail }m=30 & 26 & B_{19}\\
\text{BC stable-tail sloped} & 490 & B_{32}\\
\text{D exact two-sided interval} & 28 & D_{25}\\
\text{D }A\text{-free frontier} & 225 & D_{53}.
\end{array}
\]
The row count identity is
\[
  14+26+490+28+225=783.
\]
The $B/C$ replay is not a pushforward-tail replay and is not used by itself
to propagate any row beyond its first post-$m=29$ hit.  The type-$D$ bridge
replays check every intervening Chain step through their independently
certified monotone-tail onsets.

\section{E8 Finite Bridge Transcript}

The final finite bridge replay for $E_8$ reports:
\[
\begin{array}{ll}
\text{prefix Chain odd exponents} & 1,3,\ldots,65,\\
\text{bridge odd exponents} & 69,71,73,75,77,\\
\text{Chain odd exponents} & 67,69,71,73,75,77,\\
\text{extended Chain odd exponents} & 67,69,\ldots,95,\\
\text{minimum prefix Chain difference} & 20,\\
\log_{10}\text{ minimum bridge value} & 100.94,\\
\log_{10}\text{ minimum Chain difference} & 100.94,\\
\log_{10}\text{ minimum extended Chain difference} & 100.94.
\end{array}
\]
The exact replay flags are all \texttt{OK}.

\section{Type A Certificate Details}

The type $A$ replay is organized around a finite set of endpoint and
finite-overlap inequalities.  The stable-rank theorem covers
$N\ge n+2$.  The finite-rank theorem covers $N<n+2$ by separating the
following cases:
\[
\begin{array}{c|c|c}
\text{case} & \text{rank range} & \text{closure mechanism}\\
\hline
\text{small ranks} & N=2,3,4,5 & \text{finite transition-ratio identities}\\
\text{finite overlap} & 6\le N\le18 & \text{integer moment replay}\\
\text{middle endpoint} & 24\le N\le27 & \text{row endpoint certificate}\\
\text{large endpoint} & N\ge30 & \text{monotone endpoint comparison}\\
\text{interpolation} & \text{all remaining }N\ge6 & \text{ratio monotonicity}
\end{array}
\]
For the small ranks the proof is not a numerical exception.  For $N=2$,
Weyl integration and the change of variables
$\theta=\alpha+\beta$, $\phi=\alpha-\beta$ identify
\[
  \frac12\Qthree^{\SU(2)}(n)
  =
  \int_{\mathrm{USp}(4)}
  \operatorname{tr}\bigl((\Lambda^2 V)^{\otimes n}\bigr)\,dA,
\]
which is the multiplicity of the trivial representation in
$(\Lambda^2V)^{\otimes n}$ and hence is a nonnegative integer.  For
$N=3,4,5$, Gessel's determinant gives the following adjoint-moment
recurrences:
\[
\begin{array}{c|l}
N & \text{recurrence}\\
\hline
3 &
(k+3)^2m_{k+1}=k(7k+11)m_k+8k(k+1)m_{k-1}\\
4 &
\begin{aligned}
&(k^3+22k^2+161k+392)m_{k+4}\\
&-(16k^3+230k^2+1054k+1524)m_{k+3}\\
&+(10k^3-340k-750)m_{k+2}\\
&+(72k^3+522k^2+1242k+972)m_{k+1}\\
&+(45k^3+270k^2+495k+270)m_k=0
\end{aligned}\\
5 &
\begin{aligned}
&192(k+1)(k+2)(k+3)(k+6)m_k\\
&+16(k+2)(k+3)(22k^2+184k+309)m_{k+1}\\
&+(k+3)(129k^3+1245k^2+2570k-1112)m_{k+2}\\
&-(k+3)(30k^3+642k^2+4487k+10196)m_{k+3}\\
&+(k+8)^2(k+10)^2m_{k+4}=0 .
\end{aligned}
\end{array}
\]
Substitution into the boundary reduction proves the ratio monotonicity
needed for every odd $n\ge5$ in these rows; the finite initial cases are
checked by exact integer moments.

The replay computes adjoint moments by iterating tensor multiplication with
the adjoint representation and reducing each tensor product by the Weyl
dot-action.  For each row it evaluates the exact integer target
\[
  R_{N,k}:=m_{k+2}^{(N)}m_k^{(N)}-(m_{k+1}^{(N)})^2.
\]
The sign check $R_{N,k}\ge0$ is exactly
$\rho_k^{(N)}\le\rho_{k+1}^{(N)}$.  The replay output contains the lines
\[
\begin{array}{rl}
\texttt{Propositions 64-65 endpoint constants:} & \texttt{ok},\\
\texttt{Proposition 66 row N=24:} & \texttt{ok},\\
\texttt{Proposition 66 row N=25:} & \texttt{ok},\\
\texttt{Proposition 66 row N=26:} & \texttt{ok},\\
\texttt{Proposition 66 row N=27:} & \texttt{ok},\\
\texttt{Proposition 67 band certificate:} & \texttt{ok},\\
\texttt{Proposition 67 tail and endpoint certificates:} & \texttt{ok},\\
\texttt{Proposition 68 row N=6:} & \texttt{ok},\\
\texttt{Proposition 68 row N=7:} & \texttt{ok},\\
\texttt{Proposition 68 row N=8:} & \texttt{ok},\\
\texttt{Proposition 68 row N=9:} & \texttt{ok},\\
\texttt{Proposition 68 row N=10:} & \texttt{ok},\\
\texttt{Proposition 68 row N=11:} & \texttt{ok},\\
\texttt{Proposition 68 row N=12:} & \texttt{ok},\\
\texttt{Proposition 68 row N=13:} & \texttt{ok},\\
\texttt{Proposition 68 row N=14:} & \texttt{ok},\\
\texttt{Proposition 68 row N=15:} & \texttt{ok},\\
\texttt{Proposition 68 row N=16:} & \texttt{ok},\\
\texttt{Proposition 68 row N=17:} & \texttt{ok},\\
\texttt{Proposition 68 row N=18:} & \texttt{ok}.
\end{array}
\]
The terminal line is
\[
  \texttt{All SU(N) repair certificates verified.}
\]

\section{Classical Row Ledger}

This appendix records the final row coverage used in
\cref{thm:classical}.  The high-rank theorem covers all rows with
$C_G\ge296$.  The finite rows with $C_G<296$ are divided by the bridge
and post-bridge certificates.

\subsection*{High-rank and trace rows}

\[
\begin{array}{c|c|c}
\text{family} & \text{rank condition} & \text{certificate}\\
\hline
B,C,D & C_G\ge296 & \text{Rains-square trace cutoff}\\
B,C & b=276,278,\ldots,295 & \text{first-hit trace}\\
D & b=276,278,280,\ldots,295,297 & \text{first-hit trace}
\end{array}
\]

\subsection*{Direct post-\texorpdfstring{$m=29$}{m=29} low-rank tails}

\[
\begin{array}{c|c|c}
\text{family} & \text{rank condition} & \text{certificate}\\
\hline
B & 2\le b\le13 & \cref{prop:post29-direct}\\
C & 2\le b\le19 & \cref{prop:post29-direct}\\
D & 4\le b\le24 & \cref{prop:post29-direct}
\end{array}
\]

\subsection*{Post-\texorpdfstring{$m=29$}{m=29} first-hit rows}

\[
\begin{array}{c|c|c}
\text{family} & \text{rank interval} & \text{number of rows}\\
\hline
B & 11\le b\le275,\ b=277 & 266\\
C & 13\le b\le275,\ b=277 & 264\\
D & 25\le b\le275,\ b=277,279 & 253\\
\hline
\text{total} & & 783
\end{array}
\]
The bucket decomposition is:
\[
\begin{array}{c|c|c|c}
\text{bucket} & B\text{ ranks} & C\text{ ranks} & D\text{ ranks}\\
\hline
\text{BC exact }m=30\text{ window} & 11,\ldots,18 & 13,\ldots,18 & -\\
\text{BC stable-tail }m=30 & 19,\ldots,31 & 19,\ldots,31 & -\\
\text{BC stable-tail sloped} & 32,\ldots,275,277 & 32,\ldots,275,277 & -\\
\text{D exact two-sided interval} & - & - & 25,\ldots,52\\
\text{D }A\text{-free frontier} & - & - & 53,\ldots,275,277,279
\end{array}
\]
The replay checks that every row in this table has a strictly positive
integer lower margin at its first post-$m=29$ hit.  The subsequent odd
exponents for $B_{11},B_{12},B_{13}$ and for $C_{13},\ldots,C_{19}$ are
covered by \cref{prop:post29-direct}; $D_{25},\ldots,D_{52}$ are covered by
\cref{prop:post29-d25-52-exact}, and $D_{53},\ldots,D_{295}$ by
\cref{prop:post29-d-a-free}; the $B/C$ high-edge rows with $218\le b\le275$
or $b=277$ are covered by
\cref{prop:post29-chernoff}; the $D$ interval rows
$D_{98},\ldots,D_{275},D_{277},D_{279}$ are covered by
\cref{prop:post29-d-interval}.  The remaining $B/C$ rows are exhausted by
the explicit split in the proof of \cref{thm:classical}.  Namely,
\cref{prop:post29-bc-layered-mgf} supplies every odd exponent
$n\ge2b+29$ for $B_{14},\ldots,B_{61}$ and
$C_{20},\ldots,C_{61}$.  It starts directly at odd exponents
$57,59,61,63$ for $B_{14},\ldots,B_{17}$; for the other rows, the proved
boundary and nonboundary correction theorems and the exact half-stable bridge
are assembled in \cref{prop:post29-bc-residual-closure}, which supplies the
required earlier prefixes.
The rows $B_{62},\ldots,B_{123}$ and $C_{62},\ldots,C_{217}$ are closed by
the same finite-prefix mechanism followed by
\cref{prop:post29-bc-interval-direct,prop:post29-b-tilted-mgf,prop:post29-c-tilted-mgf};
$B_{124},\ldots,B_{217}$ are covered by
\cref{prop:post29-b-unitary-square}.  The exact bridge
\cref{prop:d-low-exact-bridge} and the terminal-tail proposition
\cref{prop:post29-d10-24-exact-tail} close the remaining
$D_4,\ldots,D_{24}$ rows.

\section{Exceptional Certificate Ledger}

The exceptional proof uses exact prefixes and direct Chain tails.  The
following table records the exact moment source and the terminal Chain
coverage.
\[
\begin{array}{c|c|c|c}
G & \text{moments checked} & \text{Chain prefix} & \text{tail start}\\
\hline
\Gtwo & m_0,\ldots,m_{200} & m=0,\ldots,97 & \text{inside prefix}\\
\Ffour & m_0,\ldots,m_{220} & m=0,\ldots,107 & \text{inside prefix}\\
\Esix & m_0,\ldots,m_{80} & m=0,\ldots,37 & n\ge75\\
\Eseven & m_0,\ldots,m_{70} & m=0,\ldots,32 & n\ge65\\
\Eeight & m_0,\ldots,m_{100} & m=0,\ldots,32 & n\ge133
\end{array}
\]
For $E_8$, the bridge chain consists of the odd exponents
\[
131,129,127,125,123,121,119,117,115,113,111,109,107,105,103,101,
\]
\[
99,97,95,93,91,89,87,85,83,81,79,77.
\]
The finite bridge supplies $Q_3(n)>0$ at $n=69,71,73,75,77$ and
nonnegative Chain differences from $67$ through $95$.  Therefore the
finite bridge, the explicit bridge chain down to $77$, and the tail from
$133$ form one continuous Chain-propagation interval.

\section{External Inputs and Replay Boundary}

The proof uses two kinds of external data.  First, it cites published
theorems with the specific loci used above: Weyl integration, character
orthogonality, classification through irreducible root systems, root-system
conventions, and classical compact groups in
\cite[Chapters II, IV--VI]{BrockertomDieck}, the finite-Coxeter
Macdonald--Mehta integral and squared-length-two root-normal convention in
\cite[\S2.1 and Theorem~3.1(i)]{Etingof2009}, Ginibre's condition (Q3) in
\cite[\S1]{Ginibre1970}, the universal compact-group adjoint trace bound in
\cite[Theorem~2.1]{GaribaldiGuralnickRains2025},
the unitary Schur-expansion moment calculation in
\cite[\S1]{Rains1997}, Gessel's determinant in
\cite[\S7, Theorem 16 and the formula after (26)]{Gessel1990}, the exact
orthogonal and symplectic joint trace moments in
\cite[Theorems~4 and 6]{DiaconisShahshahani1994}, Rains'
orthogonal and symplectic Schur-integral stabilization in
\cite[\S3, Lemmas~3.1--3.2 and Theorem~3.4]{Rains1997}, Albion's
column-even Littlewood identity in
\cite[Introduction, (1.1c)]{Albion2024}, Krattenthaler's semistandard
growth-diagram conversion in \cite[Theorem 4]{Krattenthaler2016},
Krattenthaler's arbitrary-filling growth-diagram local rules in
\cite[\S4.1, Theorem 7]{Krattenthaler2006},
Huh--Kim--Krattenthaler--Okada's marked even-bound vacillating-tableau model
in \cite[Definitions 8.1, 8.5, 8.6 and Theorem 8.7]{HKKO2025}, the
orthogonal and symplectic
power-map decomposition
of Rains \cite[Theorem 1.5, equations (22)--(23), and Theorem 1.7,
equation (37)]{RainsPowerMaps}, and the
Courteaut--Johansson--Lambert trace estimates
\cite[Theorems 1.1, 1.4, 4.1 and \S4.2]{CJL2022}.  The ancillary table of
Bourjaily--Plesser--Vergu
\cite[Appendix B and ancillary file
\texttt{adjoint\_tensor\_ranks.tex}]{BPV2024} is used only as an independent
comparison for the $E_8$ values $m_{71},\ldots,m_{100}$.  The proof input is
regenerated from the $E_8$ Cartan datum by the exact Racah--Speiser iteration
and character-pairing identities of
\cref{lem:racah-speiser-step,lem:character-pairing}; the replay requires
termwise equality with the local transcription of the published table.

The source notes in this repository are not invoked as additional analytic
lemmas.  They are part of the computer-assisted certificate boundary: the
manuscript states the theorem, reductions, certificate criteria, row ledgers,
replay transcripts, and artifact hashes, while the replay implementations and
hashed ledgers check the exact integer/rational certificate statements printed
above.

The cited primary-source files used in the source audit are vendored in
\path{ginibre_q3/references}.  In particular, the local folder contains
Ginibre's paper, Rains' 1998 and 2003 papers, Gessel's author PDF,
Albion's arXiv PDF, Krattenthaler's arXiv PDF,
Krattenthaler's 2006 arXiv PDF on arbitrary-filling growth diagrams,
Huh--Kim--Krattenthaler--Okada's author PDF,
Courteaut--Johansson--Lambert's arXiv PDF,
Bröcker--tom Dieck's GTM text, the Bourjaily--Plesser--Vergu arXiv paper, and
the BPV ancillary
\path{adjoint_tensor_ranks.tex} file used for the independent $E_8$
comparison.
The folder also contains the primary Etingof arXiv PDF and the primary
Garibaldi--Guralnick--Rains journal PDF; these replace the former
Markdown-only mappings for those two citations.
The fourteen bibliography keys and fifteen corresponding local source files
are mapped to their verified theorem/equation loci in
\path{references/active_paper_sources.tsv}.  The complete directory is pinned
by \path{references/references_manifest.sha256}; the clean-room replay rejects
a missing, changed, unlisted, or newly added reference file.

\else

\section{Certificate and Replay Contract}

\begin{definition}[Accepted certificate]\label{def:accepted-certificate}
An accepted certificate has manifested source and output, recorded status
zero, and either exact integer/rational arithmetic or an outward-rounded
interval with positive lower endpoint.  Diagnostic, interrupted, nonzero, and
superseded artifacts are excluded.  A certificate may verify a conditional
arithmetic implication only when its hypotheses are stated in a numbered
result and discharged at every main-theorem-reachable invocation; an
undischarged implication is not accepted proof evidence.  Every
character-moment source passed to
an active root-datum source-replay stage is regenerated during the clean-room
replay and compared termwise with the consumed ledger.  The finite row-gated
bridge is likewise regenerated by the exact Pieri/CRT and type-$D$ interval
verifiers of
\cref{prop:bc-source-row-gated-bridge,prop:d-a-free-row-gated-bridge}.
The superseded historical bridge archive is provenance only.  A content hash
establishes integrity and provenance only; it never by itself decides a sign.
\end{definition}

\begin{proposition}[Certificate coverage and replay contract]\label{prop:certificate-sign-sound}
Every computer-assisted sign or threshold claim used in
\cref{thm:type-a,thm:classical,thm:exceptional} is decided by an accepted
certificate in the sense of \cref{def:accepted-certificate}; no unenclosed
floating-point value and no hash-only source decides a sign.  Moreover,
\texttt{make -C ginibre\_q3 clean-room-replay} returns success only after the
manifest, source-regeneration, status, source-structure, accepted-margin,
fresh-build, main-theorem certificate-reachability coverage, and
three-pass document checks stated in \texttt{REPLAY.md} all pass.
\end{proposition}

\begin{proof}
The discrete checkers form \cref{lem:moment-formula,def:chain-diff} with GMP
integers or rationals; analytic checkers use directed MPFR and report positive
lower margins.  The character-source checker first reconstructs every root
system from the stated Cartan matrix, verifies the adjoint weight multiset,
performs the exact Racah--Speiser iterations, and applies
\cref{lem:character-pairing} before any stored moment is consumed.  The
isolated driver deletes copied build products, rebuilds the active checkers and
both documents, extracts all certificate logs cited on the main-theorem
dependency walk, and requires exact equality with the set of recomputed
certificate stages.  It also requires every such log to be hash-authenticated
and every reachable post-$m=29$ log to have accepted classification.  It
rejects any failed prerequisite listed in the statement.
\end{proof}

\fi

%======================================================================
\begingroup\scriptsize
\begin{thebibliography}{99}
\scriptsize
\setlength{\itemsep}{0pt}
\setlength{\parsep}{0pt}

\bibitem{BrockertomDieck}
T.~Brocker and T.~tom~Dieck,
\emph{Representations of Compact Lie Groups},
Graduate Texts in Math. \textbf{98}, Springer, 1985.

\bibitem{Ginibre1970}
J.~Ginibre,
\emph{General formulation of Griffiths' inequalities},
Comm. Math. Phys. \textbf{16} (1970), 310--328.

\bibitem{Gessel1990}
I.~M. Gessel,
\emph{Symmetric functions and P-recursiveness},
J. Combin. Theory Ser. A \textbf{53} (1990), 257--285.

\bibitem{DiaconisShahshahani1994}
P.~Diaconis and M.~Shahshahani,
\emph{On the eigenvalues of random matrices},
J. Appl. Probab. \textbf{31} (1994), 49--62.

\bibitem{Rains1997}
E.~M.~Rains,
\emph{Increasing subsequences and the classical groups},
Electron. J. Combin. \textbf{5} (1998), Research Paper 12.

\bibitem{RainsPowerMaps}
E.~M.~Rains,
\emph{Images of eigenvalue distributions under power maps},
Probab. Theory Related Fields \textbf{125} (2003), 522--538.

\bibitem{Albion2024}
S.~P.~Albion,
\emph{Proof of some Littlewood identities conjectured by Lee, Rains and
Warnaar},
Proc. Amer. Math. Soc. Ser. B \textbf{11} (2024), 133--146.

\bibitem{Krattenthaler2016}
C.~Krattenthaler,
\emph{Bijections between oscillating tableaux and (semi)standard tableaux via
growth diagrams},
J. Combin. Theory Ser. A \textbf{144} (2016), 277--291.

\bibitem{Krattenthaler2006}
C.~Krattenthaler,
\emph{Growth diagrams, and increasing and decreasing chains in fillings of
Ferrers shapes},
Adv. Appl. Math. \textbf{37} (2006), 404--431.

\bibitem{HKKO2025}
J.~Huh, J.~S.~Kim, C.~Krattenthaler, and S.~Okada,
\emph{Bounded Littlewood identities with fixed number of odd rows or odd
columns},
Electron. J. Combin. \textbf{33} (2026), no.~3, Paper~P3.5.

\bibitem{CJL2022}
C.~Courteaut, K.~Johansson, and G.~Lambert,
\emph{From Berry--Esseen to super-exponential},
arXiv:2204.03282.

\bibitem{Etingof2009}
P.~Etingof,
\emph{A uniform proof of the Macdonald--Mehta--Opdam identity for finite
Coxeter groups},
arXiv:0903.5084, 2009.

\bibitem{GaribaldiGuralnickRains2025}
S.~Garibaldi, R.~M.~Guralnick, and E.~M.~Rains,
\emph{Invariant derivations and trace bounds},
J. \'{E}c. polytech. Math. \textbf{12} (2025), 1229--1287.

\bibitem{BPV2024}
J.~L. Bourjaily, M.~R. Plesser, and C.~Vergu,
\emph{The many colours of amplitudes},
arXiv:2412.21189.

\end{thebibliography}
\endgroup

\end{document}
